% Status and semantic-level disputes are settled by the concordance (concordance.tex).
% This file is the master chapter for the twisted holography / quantum gravity
% programme.  It assumes the theorem environments and typography of main.tex.

\chapter{Twisted Holography and Quantum Gravity}
% label removed: chap:twisted-holography-quantum-gravity
\index{twisted holography|textbf}
\index{quantum gravity!modular Koszul approach|textbf}
\index{holographic modular Koszul datum!quantum gravity programme|textbf}

% ----------------------------------------------------------------------
% Local macros: provided only if absent from the ambient manuscript.
% ----------------------------------------------------------------------
\providecommand{\Bbbk}{\mathbb{k}}
\providecommand{\HS}{\operatorname{HS}}
\providecommand{\wt}{\operatorname{wt}}
\providecommand{\CE}{\operatorname{CE}}
\providecommand{\End}{\operatorname{End}}
\providecommand{\Hom}{\operatorname{Hom}}
\providecommand{\Sym}{\operatorname{Sym}}
\providecommand{\id}{\mathrm{id}}
\providecommand{\Tr}{\operatorname{Tr}}
\providecommand{\Det}{\operatorname{Det}}
\providecommand{\Aut}{\operatorname{Aut}}
\providecommand{\Fred}{\operatorname{Fred}}
\providecommand{\Spec}{\operatorname{Spec}}
\providecommand{\Res}{\operatorname{Res}}
\providecommand{\rank}{\operatorname{rank}}
\providecommand{\ad}{\operatorname{ad}}
\providecommand{\Ainf}{\mathsf{A}_{\infty}}
\providecommand{\Linf}{L_\infty}
\providecommand{\PV}{\operatorname{PV}}
\providecommand{\RHom}{\operatorname{RHom}}
\providecommand{\BBar}{\operatorname{Bar}}
\providecommand{\Cobar}{\operatorname{\Omega}}
\providecommand{\barB}{\overline{B}}
\providecommand{\Diff}{\operatorname{Diff}}
\providecommand{\tr}{\operatorname{tr}}
\providecommand{\Zder}{\operatorname{Z}_{\mathrm{der}}}
\providecommand{\Abulk}{A_{\mathrm{bulk}}}
\providecommand{\Bbound}{B_{\partial}}
\providecommand{\sn}{\mathrm{SN}}
\providecommand{\hSym}{\widehat{\Sym}}
\providecommand{\lines}{\mathsf{Lines}}
\providecommand{\ModEnv}{\operatorname{ModEnv}}
\providecommand{\ohat}{\widehat{\otimes}}
\providecommand{\Mbar}{\overline{\mathcal M}}
\providecommand{\Perf}{\operatorname{Perf}}
\providecommand{\dCrit}{\operatorname{dCrit}}
\providecommand{\gr}{\operatorname{gr}}
\providecommand{\coker}{\operatorname{coker}}
\providecommand{\Spf}{\operatorname{Spf}}

% ======================================================================
%
%  SECTION 1: THE HOLOGRAPHIC MODULAR KOSZUL DATUM
%
% ======================================================================

\section{The holographic modular Koszul datum}
% label removed: sec:thqg-holographic-datum

\subsection{Standing hypotheses}
% label removed: subsec:thqg-standing

\begin{definition}[Standing hypotheses]
% label removed: def:thqg-standing-hypotheses
A \emph{gravitational input} is a quadruple $(\cA, X, \langle\cdot,\cdot\rangle, k)$ where:
\begin{enumerate}[label=\textup{(\roman*)}]
\item $\cA$ is a chiral algebra on a smooth projective curve $X$, chirally Koszul
  in the sense of Theorem~\ref*{V1-thm:koszul-equivalences-meta} (any of the twelve
  equivalent characterizations);
\item $\cA$ is cyclically admissible
  (Definition~\ref{def:cyclically-admissible}): conformal grading, exhaustive
  filtration, bounded poles, invariant pairing $\langle\cdot,\cdot\rangle$;
\item the HS-sewing criterion holds: polynomial OPE growth and
  subexponential sector growth (Theorem~\ref{thm:general-hs-sewing});
\item the level $k$ avoids the critical value (for $\widehat{\mathfrak{g}}$:
  $k \ne -h^\vee$; for universal $\mathcal{W}$-algebras: generic level).
\end{enumerate}
All results below assume these hypotheses unless explicitly relaxed.
% label removed: fwk:thqg-standing-hypotheses%
\end{definition}

\begin{remark}[Scope of the hypotheses]
% label removed: rem:thqg-scope
The four conditions are satisfied by:
\begin{itemize}
\item all affine Kac--Moody algebras $V_k(\mathfrak{g})$ at non-critical level;
\item all principal $\mathcal{W}$-algebras $\mathcal{W}_k(\mathfrak{g}, f_{\mathrm{prin}})$ at generic $k$;
\item all free-field algebras (Heisenberg, $\beta\gamma$, $bc$, lattice);
\item the Virasoro algebra at all central charges $c$;
\item hook-type $\mathcal{W}$-algebras in type~$A$ (Fehily, Creutzig--Linshaw--Nakatsuka--Sato).
\end{itemize}
The critical exclusion in~(iv) removes the locus where the Sugawara construction is undefined.
The HS-sewing criterion in~(iii) is automatic for the standard landscape
(Volume~I, Theorem~\ref{thm:general-hs-sewing}).
\end{remark}


\subsection{The Swiss-cheese convolution algebra}
% label removed: subsec:thqg-swiss-cheese-conv

Fix a gravitational input $(\cA, X, \langle\cdot,\cdot\rangle, k)$.
The holomorphic-topological geometry on $\C \times \R_{\ge 0}$ produces
a Swiss-cheese operad $\mathsf{SC}^{\mathrm{ch,top}}$
(Chapter~\ref{chap:ht-bulk-boundary-line}).
Its Koszul dual cooperad $(\mathsf{SC}^{\mathrm{ch,top}})^{!}$
determines a convolution dg Lie algebra.

\begin{definition}[Swiss-cheese convolution algebra]
% label removed: def:thqg-swiss-cheese-conv
\index{Swiss-cheese convolution algebra}
The \emph{Swiss-cheese convolution algebra} of $\cA$ is
\begin{equation}% label removed: eq:thqg-gSC
\mathfrak{g}^{\mathrm{SC}}_\cA
\;:=\;
\Convstr\bigl(
(\mathsf{SC}^{\mathrm{ch,top}})^{!},\;
\End_\cA
\bigr),
\end{equation}
the strict convolution dg Lie algebra of the Koszul dual cooperad
evaluated in the endomorphism operad of $\cA$.  Its differential is
the convolution differential $D_\cA$; its bracket is the
operadic composition bracket.
\end{definition}

\begin{proposition}[Factorization of the Swiss-cheese algebra;
\ClaimStatusProvedHere]
% label removed: prop:thqg-gSC-factorization
The Swiss-cheese convolution algebra decomposes as a direct product
\begin{equation}% label removed: eq:thqg-gSC-factorization
\mathfrak{g}^{\mathrm{SC}}_\cA
\;\cong\;
\gAmod \;\times\; \mathfrak{g}^{\R}_\cA,
\end{equation}
where:
\begin{enumerate}[label=\textup{(\alph*)}]
\item $\gAmod = \Convstr(\mathcal{C}^!_{\mathrm{ch}},
  \mathcal{P}^{\mathrm{ch}})$ is the modular convolution dg Lie
  algebra (Definition~\ref*{V1-def:modular-convolution-dg-lie}),
  controlling the chiral ($\C$) direction;
\item $\mathfrak{g}^{\R}_\cA$ is the topological ($\R$) convolution
  algebra, controlling the $E_1$-coalgebra structure along the
  topological direction.
\end{enumerate}
The two factors commute: $[\gAmod, \mathfrak{g}^{\R}_\cA] = 0$.
\end{proposition}

\begin{proof}
The Swiss-cheese operad factors as a semidirect product of the
chiral operad (governing disk operations in $\C$) and the
associative operad (governing interval operations along $\R$).
The Koszul dual cooperad inherits the same factorization, and the
convolution functor preserves products.
\end{proof}


\subsection{The six components}
% label removed: subsec:thqg-six-components

\begin{definition}[Holographic modular Koszul datum]
% label removed: def:thqg-holographic-datum
\index{holographic modular Koszul datum|textbf}
Given a gravitational input $(\cA, X, \langle\cdot,\cdot\rangle, k)$, the
\emph{holographic modular Koszul datum} is the sextuple
\begin{equation}% label removed: eq:thqg-holographic-datum
\mathcal{H}(\cA) \;=\; \bigl(\cA,\;\cA^!,\;\mathcal{C},\;r(z),\;
\Theta_\cA,\;\nabla^{\mathrm{hol}}\bigr),
\end{equation}
whose six components are defined as follows.
\end{definition}

\begin{enumerate}[label=\textbf{Component (\alph*):}, leftmargin=4.5em, itemsep=1.5ex]

\item \textbf{The boundary chiral algebra $\cA$.}
% label removed: item:thqg-component-a
This is the input datum: a chirally Koszul, cyclically admissible
chiral algebra on $X$.  In the holographic context, $\cA$ is the
algebra of boundary local operators for a 3d holomorphic-topological
theory on $\C \times \R_{\ge 0}$.  Its OPE data are the initial
conditions for the entire holographic system.

\item \textbf{The Koszul dual algebra $\cA^!$.}
% label removed: item:thqg-component-b
\index{Koszul dual!holographic}
The Koszul dual is
\begin{equation}% label removed: eq:thqg-koszul-dual
\cA^! \;:=\; \bigl(H^\bullet(B(\cA))\bigr)^\vee,
\end{equation}
the linear dual of the bar cohomology.  The bar complex $B(\cA)$ is
the cofree coalgebra on the desuspended generators of $\cA$, with
differential encoding the OPE.  The Verdier duality converts the
coalgebra structure on $H^\bullet(B(\cA))$ into an algebra structure
on $\cA^!$.

The Koszul dual is \emph{not} the cobar construction: $\Omega(B(\cA))
\simeq \cA$ is bar-cobar inversion (Theorem~B), while $\cA^!$ is
obtained by the composite $B \to H^\bullet \to (-)^\vee$.

\begin{center}
\renewcommand{\arraystretch}{1.2}
\small
\begin{tabular}{lll}
\textbf{Family} & $\cA$ & $\cA^!$ \\
\hline
Heisenberg & $\mathcal{H}$ & Free boson (not self-dual) \\
Affine $\widehat{\mathfrak{g}}_k$ & $V_k(\mathfrak{g})$
  & $V_{-k-2h^\vee}(\mathfrak{g})$ \\
Virasoro & $\mathrm{Vir}_c$ & $\mathrm{Vir}_{26-c}$ \\
$\beta\gamma$ & $\beta\gamma$ & $bc$ \\
Lattice $V_\Lambda$ & $V_\Lambda$ & $V_{\Lambda^*}$
\end{tabular}
\end{center}

\item \textbf{The line-operator category $\mathcal{C}$.}
% label removed: item:thqg-component-c
\index{line-operator category!holographic}
The category of line operators along $\R$ is the module category of
the Koszul dual:
\begin{equation}% label removed: eq:thqg-line-category
\mathcal{C} \;\simeq\; \cA^!\text{-}\mathsf{mod}.
\end{equation}
This is a dg category with compact generator given by $\cA^!$ itself.
The identification follows from the line-operator theorem
(Theorem~\ref{thm:lines_as_modules}): the $E_1$-coalgebra structure
along $\R$ produces a monoidal category whose modules are precisely
line operators, and the perturbative construction of Dimofte--Niu--Py
identifies this monoidal category with $\cA^!$-modules.

\item \textbf{The collision residue $r(z)$.}
% label removed: item:thqg-component-d
\index{collision residue|textbf}
\index{classical Yang--Baxter equation!from collision residue}
The collision residue is the genus-$0$, arity-$2$ projection of the
universal MC element:
\begin{equation}% label removed: eq:thqg-collision-residue
r(z) \;:=\; \Res^{\mathrm{coll}}_{0,2}(\Theta_\cA)
\;\in\;
\cA^! \ohat \cA^!\bigl[\!\bigl[z^{-1}, z\bigr]\!\bigr].
\end{equation}
The notation $\Res^{\mathrm{coll}}_{0,2}$ denotes the residue of the
MC element at the codimension-$1$ boundary stratum of $\overline{C}_2(X)$
where the two points collide.  The spectral parameter $z$ is the
relative coordinate of the two insertions.

The collision residue satisfies the classical Yang--Baxter equation:
\begin{equation}% label removed: eq:thqg-cybe
[r_{12}(z_1-z_2),\, r_{13}(z_1-z_3)]
+ [r_{12}(z_1-z_2),\, r_{23}(z_2-z_3)]
+ [r_{13}(z_1-z_3),\, r_{23}(z_2-z_3)]
= 0.
\end{equation}
This is the genus-$0$, arity-$3$ MC equation, evaluated on the
codimension-$2$ boundary of $\overline{C}_3(X)$.  The three terms
correspond to the three boundary strata where pairs of points
collide, and the vanishing is the Arnold relation.

\item \textbf{The universal MC element $\Theta_\cA$.}
% label removed: item:thqg-component-e
\index{Maurer--Cartan element!universal|textbf}
The bar-intrinsic MC element
(Volume~I, Theorem~\ref*{V1-thm:mc2-bar-intrinsic}) is
\begin{equation}% label removed: eq:thqg-theta
\Theta_\cA \;:=\; D_\cA - d_0
\;\in\;
\MC\bigl(\gAmod\bigr),
\end{equation}
where $D_\cA$ is the full bar differential and $d_0$ is the linear
part.  The MC equation
\begin{equation}% label removed: eq:thqg-mc-equation
D_\cA \Theta_\cA + \tfrac{1}{2}[\Theta_\cA, \Theta_\cA] = 0
\end{equation}
holds because $D_\cA^2 = 0$
(Volume~I, Theorem~\ref{thm:convolution-d-squared-zero}).

The element $\Theta_\cA$ decomposes by genus and arity:
\begin{equation}% label removed: eq:thqg-theta-decomposition
\Theta_\cA
= \sum_{g \ge 0} \sum_{r \ge 2} \hbar^g\, \Theta_\cA^{(g,r)},
\end{equation}
where $\Theta_\cA^{(g,r)}$ is the component with loop genus $g$
and arity $r$ (number of inputs in the graph sum).  The finite-order
truncations $\Theta_\cA^{\le r} := \sum_{s=2}^{r}\sum_g
\hbar^g\,\Theta_\cA^{(g,s)}$ are the shadow obstruction tower.

\item \textbf{The shadow connection $\nabla^{\mathrm{hol}}$.}
% label removed: item:thqg-component-f
\index{shadow connection|textbf}
\index{modular shadow connection}
The modular shadow connection is
\begin{equation}% label removed: eq:thqg-shadow-connection
\nabla^{\mathrm{hol}}_{g,n}
\;:=\; d - \operatorname{Sh}_{g,n}(\Theta_\cA),
\end{equation}
where $\operatorname{Sh}_{g,n}$ is the shadow projection
(Volume~I, Corollary~\ref*{V1-cor:shadow-extraction}) that evaluates
$\Theta_\cA$ on the genus-$g$, $n$-pointed stratum of the
modular operad.  The shadow connection is a connection on the
moduli space $\Mbar_{g,n}$ valued in the cyclic deformation
complex $\Defcyc^{\mathrm{mod}}(\cA)$.

\end{enumerate}


\subsection{Projection theorem}
% label removed: subsec:thqg-projection

\begin{theorem}[All components are projections of $\Theta_\cA$;
\ClaimStatusProvedHere]
% label removed: thm:thqg-projection
\index{holographic modular Koszul datum!projection theorem}
Every component of the holographic modular Koszul datum
$\mathcal{H}(\cA) = (\cA, \cA^!, \mathcal{C}, r(z),
\Theta_\cA, \nabla^{\mathrm{hol}})$ is determined by $\Theta_\cA$:
\begin{center}
\renewcommand{\arraystretch}{1.4}
\small
\begin{tabular}{lll}
\textbf{Component} & \textbf{Projection} & \textbf{Formula} \\
\hline
$\cA$ & Input & $\cA$ is the input \\[2pt]
$\cA^!$ & $H^\bullet(D_\cA) = H^\bullet(d_0 + \Theta_\cA)$
  & Bar cohomology + Verdier dual \\[2pt]
$\mathcal{C}$ & $\cA^!\text{-}\mathsf{mod}$
  & Module category of $\cA^!$ \\[2pt]
$r(z)$ & $\Res^{\mathrm{coll}}_{0,2}(\Theta_\cA)$
  & Genus-$0$, arity-$2$ boundary residue \\[2pt]
$\Theta_\cA$ & Identity & The MC element itself \\[2pt]
$\nabla^{\mathrm{hol}}_{g,n}$
  & $d - \operatorname{Sh}_{g,n}(\Theta_\cA)$
  & Shadow projection to $(g,n)$
\end{tabular}
\end{center}
The components $\cA^!$, $\mathcal{C}$, $r(z)$, and
$\nabla^{\mathrm{hol}}$ are all functorial in $\cA$ through
$\Theta_\cA$.
\end{theorem}

\begin{proof}
The bar differential is $D_\cA = d_0 + \Theta_\cA$, so:
\begin{enumerate}[label=(\roman*)]
\item $\cA^! = (H^\bullet(B(\cA)))^\vee = (H^\bullet(D_\cA))^\vee$
  is determined by the differential, hence by $\Theta_\cA$.
\item $\mathcal{C} \simeq \cA^!\text{-}\mathsf{mod}$ depends only
  on $\cA^!$.
\item $r(z) = \Res^{\mathrm{coll}}_{0,2}(\Theta_\cA)$ is
  the boundary residue of $\Theta_\cA$ at the collision
  stratum of $\overline{C}_2(X)$, a direct projection.
\item $\nabla^{\mathrm{hol}}_{g,n} = d - \operatorname{Sh}_{g,n}(\Theta_\cA)$
  is the shadow projection, a direct projection.
\end{enumerate}
Functoriality: a quasi-isomorphism $\cA \xrightarrow{\sim} \cA'$ of
chirally Koszul algebras induces a gauge equivalence
$\Theta_\cA \sim \Theta_{\cA'}$ in $\MC(\gAmod)$
(Volume~I, Theorem~\ref{thm:shadow-homotopy-invariance}),
and all projections are gauge-equivariant.
\end{proof}


\subsection{Flatness of the shadow connection}
% label removed: subsec:thqg-flatness

\begin{theorem}[Flatness; \ClaimStatusProvedHere]
% label removed: thm:thqg-flatness
\index{shadow connection!flatness}
The shadow connection $\nabla^{\mathrm{hol}}_{g,n}$ is flat:
\begin{equation}% label removed: eq:thqg-flatness
(\nabla^{\mathrm{hol}}_{g,n})^2 = 0.
\end{equation}
\end{theorem}

\begin{proof}
Expanding:
\begin{align}
(\nabla^{\mathrm{hol}}_{g,n})^2
&= (d - \operatorname{Sh}_{g,n}(\Theta_\cA))^2 \notag \\
&= - d\,\operatorname{Sh}_{g,n}(\Theta_\cA)
   - \operatorname{Sh}_{g,n}(\Theta_\cA)\,d
   + \operatorname{Sh}_{g,n}(\Theta_\cA)^2.
   % label removed: eq:thqg-flatness-expand
\end{align}
The shadow projection $\operatorname{Sh}_{g,n}$ is a Lie algebra
homomorphism on the $(g,n)$-stratum, so:
\[
\operatorname{Sh}_{g,n}(\Theta_\cA)^2
= \tfrac{1}{2}\operatorname{Sh}_{g,n}([\Theta_\cA, \Theta_\cA]).
\]
The MC equation~\eqref{eq:thqg-mc-equation} gives
$[\Theta_\cA, \Theta_\cA] = -2\,D_\cA\Theta_\cA$, and the shadow
projection intertwines $D_\cA$ with $d$:
$\operatorname{Sh}_{g,n}(D_\cA\Theta_\cA)
= d\,\operatorname{Sh}_{g,n}(\Theta_\cA)
+ \operatorname{Sh}_{g,n}(\Theta_\cA)\,d$.
Substituting:
\[
(\nabla^{\mathrm{hol}}_{g,n})^2
= - \operatorname{Sh}_{g,n}(D_\cA\Theta_\cA)
  + \tfrac{1}{2}\operatorname{Sh}_{g,n}([\Theta_\cA, \Theta_\cA])
= - \operatorname{Sh}_{g,n}\!\bigl(D_\cA\Theta_\cA
  + \tfrac{1}{2}[\Theta_\cA, \Theta_\cA]\bigr)
= 0.\qedhere
\]
\end{proof}

\begin{remark}[Genus-$0$ shadow flatness and affine KZ comparison]
% label removed: rem:thqg-kz
\index{KZ connection!affine comparison surface}
\index{shadow connection!genus-$0$ flatness}
At genus~$0$, the shadow connection with $n$ insertions is
\begin{equation}% label removed: eq:thqg-kz
\nabla^{\mathrm{hol}}_{0,n}
= d - \sum_{i<j} r^{(ij)}(z_i - z_j)\, d\log(z_i - z_j),
\end{equation}
where $r^{(ij)}$ denotes $r(z)$ acting on the $i$-th and $j$-th
tensor factors.  Flatness of this connection is the statement that
$r(z)$ satisfies the CYBE~\eqref{eq:thqg-cybe}, which in turn
follows from the Arnold relations on $\overline{C}_3(X)$
(Volume~I, Theorem~\ref{thm:collision-residue-twisting}).
On the affine Kac--Moody comparison surface, this genus-$0$
shadow connection identifies with the KZ connection.
On a single \emph{primary line} (a one-dimensional
subspace of the cyclic deformation complex spanned by a
single primary field; Volume~I,
Definition~\ref{def:shadow-metric}), the shadow connection
is the logarithmic
connection of the shadow metric~$Q_L$
(Theorem~10.22.26 of Vol~I); its monodromy
is the Koszul sign~$(-1)$.
\end{remark}


\subsection{The shadow obstruction tower}
% label removed: subsec:thqg-shadow-tower

\begin{definition}[Shadow depth and archetype classification]
% label removed: def:thqg-shadow-archetype
\index{shadow depth|textbf}
\index{shadow archetype|textbf}
The \emph{shadow depth} $r_{\max}(\cA) \in \{2,3,4,\ldots\} \cup \{\infty\}$ is
the smallest $r$ such that the obstruction class
$o^{(r+1)}(\cA) \in H^2(F^{r+1}\gAmod / F^{r+2}\gAmod, d_2)$ vanishes
(or $\infty$ if no such $r$ exists).  The four archetype classes
(Volume~I, Definition~\ref{def:shadow-depth-classification}) are:
\begin{center}
\renewcommand{\arraystretch}{1.2}
\small
\begin{tabular}{clll}
\textbf{Class} & \textbf{Name} & $r_{\max}$ & \textbf{Standard examples} \\
\hline
G & Gaussian & $2$ & Heisenberg, free fermion, lattice, KS on CY$_3$ \\
L & Lie/tree & $3$ & Affine $\widehat{\mathfrak{g}}$, symplectic fermion \\
C & Contact & $4$ & $\beta\gamma$, $bc$ \\
M & Mixed & $\infty$ & Virasoro, $\mathcal{W}_N$ ($N \ge 3$)
\end{tabular}
\end{center}
\end{definition}

\begin{proposition}[Named shadows; \ClaimStatusProvedHere]
% label removed: prop:thqg-named-shadows
\index{modular characteristic!as shadow}
\index{cubic shadow}
\index{quartic resonance class}
The first three layers of the shadow obstruction tower are:
\begin{enumerate}[label=\textup{(\roman*)}]
\item \textbf{Modular characteristic} ($r=2$):
  $\kappa(\cA) := \Theta_\cA^{\le 2}|_{\text{scalar}}$, the
  supertrace of the identity on weight-$1$ generators.  This is
  Theorem~D of Volume~I.
\item \textbf{Cubic shadow} ($r=3$):
  $\mathfrak{C}(\cA) := \Theta_\cA^{\le 3} - \Theta_\cA^{\le 2}$,
  the tree-level $3$-point coupling.  Vanishes for class~G algebras;
  nontrivial for class~L and above.
\item \textbf{Quartic resonance class} ($r=4$):
  $\mathfrak{Q}(\cA) := [\Theta'^{\le 4}_\cA]$, the equivalence class
  of the quartic correction modulo gauge equivalence.  For class~G
  and~L, $\mathfrak{Q}(\cA) = 0$.  The first nontrivial quartic
  class is $\mathfrak{Q}(\beta\gamma)$.
  The Virasoro contact invariant is
  $Q^{\mathrm{contact}}_{\mathrm{Vir}} = 10/[c(5c+22)]$
  (Volume~I, Theorem~\ref{thm:virasoro-quartic-class}).
\end{enumerate}
All three are determined by $\Theta_\cA$.
\end{proposition}

\begin{proof}
By the shadow extraction corollary
(Volume~I, Corollary~\ref*{V1-cor:shadow-extraction}), each finite-order
truncation $\Theta_\cA^{\le r}$ is a well-defined MC element in the
truncated quotient $\gAmod / F^{r+1}\gAmod$.  The named shadows are
the differences between consecutive truncations, hence determined
by $\Theta_\cA$.  The specific identifications with the modular
characteristic, cubic OPE coupling, and quartic resonance class
follow from the shadow-formality identification
(Volume~I, Proposition~\ref{prop:shadow-formality-low-arity}).
\end{proof}


\subsection{The gravitational phase space}
% label removed: subsec:thqg-phase-space

\begin{definition}[Gravitational phase space]
% label removed: def:thqg-gravitational-phase-space
\index{gravitational phase space|textbf}
The \emph{gravitational phase space} at genus $g$ is the ambient complex
$\mathbf{C}_g(\cA)$ equipped with the shifted-symplectic Lagrangian splitting
\begin{equation}% label removed: eq:thqg-lagrangian-splitting
\mathbf{C}_g(\cA) \;\simeq\; \mathbf{Q}_g(\cA) \oplus \mathbf{Q}_g(\cA^!)
\end{equation}
from the complementarity theorem (Volume~I, Theorem~\ref{thm:quantum-complementarity-main})
and its shifted-symplectic upgrade
(Volume~I, Theorem~\ref{thm:shifted-symplectic-complementarity}).
The summands $\mathbf{Q}_g(\cA)$ and $\mathbf{Q}_g(\cA^!)$ are the bulk and defect
sectors.  The Verdier involution $\sigma$ interchanges them.
\end{definition}

\begin{proposition}[Phase space at genus $0$; \ClaimStatusProvedHere]
% label removed: prop:thqg-phase-space-genus0
At genus $0$, the gravitational phase space is
\begin{equation}% label removed: eq:thqg-phase-genus0
\mathbf{C}_0(\cA)
\;\simeq\;
\operatorname{Tot}\bigl(
T^*[-1]\,\MC(\mathfrak{g}^{\mathrm{tree}}_\cA)
\bigr),
\end{equation}
the $(-1)$-shifted cotangent bundle of the tree-level MC moduli.
The Lagrangian splitting identifies $\mathbf{Q}_0(\cA)$ with the
MC locus and $\mathbf{Q}_0(\cA^!)$ with the shifted cotangent
fiber.
\end{proposition}

\begin{proof}
At genus~$0$, only tree graphs contribute.  The convolution algebra
$\mathfrak{g}^{\mathrm{tree}}_\cA$ is the genus-$0$ truncation of
$\gAmod$.  The complementarity theorem at genus~$0$ reduces to the
classical statement that the MC locus is Lagrangian in the
$(-1)$-shifted cotangent bundle.
\end{proof}


\subsection{Summary table}
% label removed: subsec:thqg-summary-table

\begin{center}
\renewcommand{\arraystretch}{1.35}
\small
\begin{tabular}{p{2.8cm}p{4.2cm}p{4.8cm}}
\textbf{Component}
  & \textbf{Geometric origin}
  & \textbf{Physical meaning} \\
\hline
$\cA$
  & Boundary vertex algebra
  & Boundary local operators \\
$\cA^!$
  & Bar cohomology + Verdier
  & Koszul dual / defect algebra \\
$\mathcal{C}$
  & $\cA^!\text{-}\mathsf{mod}$
  & Line operators along $\R$ \\
$r(z)$
  & Collision residue on $\overline{C}_2(X)$
  & Gravitational scattering kernel \\
$\Theta_\cA$
  & $D_\cA - d_0 \in \MC(\gAmod)$
  & All-genus, all-arity bulk coupling \\
$\nabla^{\mathrm{hol}}_{g,n}$
  & $d - \operatorname{Sh}_{g,n}(\Theta_\cA)$
  & Modular shadow connection \\[4pt]
\hline
$\kappa(\cA)$
  & Arity-$2$ shadow
  & Modular characteristic / gravitational anomaly \\
$\mathfrak{C}(\cA)$
  & Arity-$3$ shadow
  & Cubic gravitational vertex \\
$\mathfrak{Q}(\cA)$
  & Arity-$4$ shadow
  & Quartic resonance / contact term \\
$\mathbf{C}_g(\cA)$
  & Ambient complex at genus $g$
  & Gravitational phase space \\
\end{tabular}
\end{center}


\subsection{Functoriality}
% label removed: subsec:thqg-functoriality

\begin{theorem}[Functoriality of $\mathcal{H}$; \ClaimStatusProvedHere]
% label removed: thm:thqg-functoriality
\index{holographic modular Koszul datum!functoriality}
The assignment $\cA \mapsto \mathcal{H}(\cA)$ is functorial with
respect to:
\begin{enumerate}[label=\textup{(\roman*)}]
\item \textbf{Quasi-isomorphisms}: If $f: \cA \xrightarrow{\sim} \cA'$
  is a quasi-isomorphism of chirally Koszul algebras, then
  $\mathcal{H}(\cA) \cong \mathcal{H}(\cA')$ via a canonical
  equivalence of MC moduli.
\item \textbf{Level deformation}: For a family $\cA_k$ parameterized
  by level $k$, the datum $\mathcal{H}(\cA_k)$ varies algebraically
  in $k$ away from the critical locus.
\item \textbf{Drinfeld--Sokolov reduction}: For principal DS reduction
  of the total current-plus-ghost system
  $V_k(\mathfrak{g})\otimes F_{\mathrm{gh}}$,
  the holographic datum transforms as
  $\mathcal{H}(\mathcal{W}_k)
  = \mathrm{DS}(\mathcal{H}(V_k\otimes F_{\mathrm{gh}}))$.
\end{enumerate}
\end{theorem}

\begin{proof}
(i)~follows from the homotopy invariance of the shadow algebra
(Volume~I, Theorem~\ref{thm:shadow-homotopy-invariance}).
(ii)~follows from the algebraic family rigidity theorem
(Volume~I, Theorem~\ref{thm:algebraic-family-rigidity}).
(iii)~follows from the principal-current-plus-ghost DS descent
statement for the modular Koszul datum and universal MC class
(Volume~I, Theorem~\ref{thm:ds-platonic-functor} and
Corollary~\ref{cor:ds-theta-descent}).
\end{proof}


% ======================================================================
%
%  SECTION 2: THE TEN GRAVITATIONAL THEOREMS
%
% ======================================================================

\section{The ten gravitational theorems}
% label removed: sec:thqg-ten-theorems

The ten theorems below encode the gravitational content of a holographic
modular Koszul datum $\mathcal{H}(\cA)$.  Each theorem is a projection
of the single MC equation~\eqref{eq:thqg-mc-equation} to a specific
stratum of $\gAmod$.  The dependency structure is recorded in the
dependency theorem (Theorem~\ref{thm:thqg-dependency}).

\subsection{G1: Perturbative finiteness}
% label removed: subsec:thqg-g1-statement

\begin{theorem}[\textbf{G1}: Perturbative finiteness; \ClaimStatusProvedHere]
% label removed: thm:thqg-g1-finiteness% label removed: item:g1
\index{perturbative finiteness|textbf}
\index{UV finiteness!twisted holography}
Let $(\cA, X, \langle\cdot,\cdot\rangle, k)$ be a gravitational input
(Definition~\ref{def:thqg-standing-hypotheses}; in particular,
$\cA$ has polynomial OPE growth and subexponential sector growth
by hypothesis~\textup{(iii)}).
For every $g \ge 0$ and $n \ge 1$, the genus-$g$ amplitude
$\langle \Theta_\cA^{\otimes n} \rangle_g$ converges absolutely.
No UV renormalization is needed: the Fulton--MacPherson compactification
$\overline{C}_n(X)$ resolves all collision singularities, and
the arity-cutoff lemma (Volume~I, Lemma~\ref{lem:arity-cutoff}) bounds the
contributing graphs at each genus.
\end{theorem}

\begin{proof}[Proof sketch]
The amplitude $\langle \Theta_\cA^{\otimes n} \rangle_g$ is a sum over
stable graphs $\Gamma$ of genus $g$ with $n$ external legs.  The arity-cutoff
lemma shows that for fixed $g$ and $n$, only finitely many graphs $\Gamma$
contribute (the strong filtration axiom bounds the vertex valency).
Each graph amplitude is computed by integration over the interior
$\overline{C}_{|\Gamma|}(X)$, which is a compact smooth manifold with
corners.  The integrand is a smooth form (the collision singularities are
resolved by the blowup construction of Fulton--MacPherson), so each
integral converges.  The sum over graphs is finite.

The HS-sewing criterion (Volume~I, Theorem~\ref{thm:general-hs-sewing})
extends this to the all-genera sewing amplitudes: polynomial OPE growth
and subexponential sector growth imply that the sewing kernel
$K_g(\cA)$ is trace-class for every $g \ge 1$.
\end{proof}


\subsection{G2: Gravitational complexity}
% label removed: subsec:thqg-g2-statement

\begin{theorem}[\textbf{G2}: Gravitational complexity; \ClaimStatusProvedHere]
% label removed: thm:thqg-g2-complexity
\index{gravitational complexity|textbf}
\index{shadow depth!gravitational interpretation}
The shadow depth $r_{\max}(\cA)$ equals the minimal
homotopy-algebraic complexity of the bulk gravitational theory:
the bulk $\Linf$-algebra has nontrivial brackets $\ell_j$
for $2 \le j \le r_{\max}(\cA)$ and $\ell_j = 0$ for $j > r_{\max}$.
In particular:
\begin{enumerate}[label=\textup{(\roman*)}]
\item $r_{\max} = 2$: free bulk theory (Gaussian);
\item $r_{\max} = 3$: semistrict higher-spin gravity (Lie/tree);
\item $r_{\max} = 4$: quartic gravitational vertices (contact);
\item $r_{\max} = \infty$: full gravitational $\Linf$-algebra (mixed).
\end{enumerate}
\end{theorem}

\begin{proof}[Proof sketch]
The bulk $\Linf$-algebra is the modular tangent complex
$T_{\Theta_\cA}\gAmod$ (Volume~I,
Construction~\ref{const:vol1-modular-tangent-complex}).  Its
$k$-ary bracket $\ell_k$ is the linearization of the MC equation
at arity $k$.  The shadow-formality identification
(Volume~I, Proposition~\ref{prop:shadow-formality-low-arity})
identifies the shadow obstruction tower with the $\Linf$-formality obstruction
tower at arities $2$, $3$, $4$.  At arity $r > r_{\max}$, the
obstruction class $o^{(r+1)}(\cA) = 0$ implies that $\ell_r$ is
homotopically trivial; at arity $r \le r_{\max}$, $o^{(r)}(\cA) \ne 0$
implies $\ell_r$ is nontrivial.

The four archetype classes (Definition~\ref{def:thqg-shadow-archetype})
correspond to the four qualitative types of gravitational theory: free,
cubic, quartic, and fully interacting.
\end{proof}


\subsection{G3: Symplectic polarization}
% label removed: subsec:thqg-g3-statement

\begin{theorem}[\textbf{G3}: Symplectic polarization; \ClaimStatusProvedHere]
% label removed: thm:thqg-g3-polarization
\index{symplectic polarization!gravitational}
\index{Lagrangian decomposition!complementarity}
Let $(\cA, X, \langle\cdot,\cdot\rangle, k)$ be a gravitational input.
\textbf{Additional hypotheses}: (a)~the cyclic pairing
$\langle\cdot,\cdot\rangle$ is nondegenerate at the chain level
(not merely on cohomology); (b)~the ambient complex $\mathbf{C}_g(\cA)$
is perfect as a chain complex.
Under these hypotheses, the complementarity splitting
\eqref{eq:thqg-lagrangian-splitting} is a shifted-symplectic Lagrangian
decomposition: both summands are Lagrangian with respect to the
$(-1)$-shifted symplectic form on $\mathbf{C}_g(\cA)$.  The Verdier
involution acts as gravitational CPT conjugation.

\smallskip\noindent
\textit{Conditionality note.}
The perfectness and chain-level nondegeneracy hypotheses are satisfied
for all standard families in the landscape
(Definition~\ref{def:thqg-standing-hypotheses}), but may fail for
exotic or non-standard inputs.  The unsymplectified complementarity
splitting $\mathbf{C}_g = \mathbf{Q}_g \oplus \mathbf{Q}_g^!$ holds
unconditionally; only the Lagrangian upgrade requires the additional
hypotheses.
\end{theorem}

\begin{proof}[Proof sketch]
The complementarity theorem (Volume~I, Theorem~\ref{thm:quantum-complementarity-main})
provides the splitting~\eqref{eq:thqg-lagrangian-splitting} at the level
of chain complexes.  The shifted-symplectic upgrade
(Volume~I, Theorem~\ref{thm:shifted-symplectic-complementarity}) equips
$\mathbf{C}_g(\cA)$ with a $(-1)$-shifted symplectic form $\omega_{-1}$
arising from the invariant pairing $\langle\cdot,\cdot\rangle$ and the
Serre duality on $\Mbar_{g,n}$.  The two summands are Lagrangian because:
(a)~$\omega_{-1}|_{\mathbf{Q}_g(\cA)} = 0$ by cyclicity of $\cA$;
(b)~$\omega_{-1}|_{\mathbf{Q}_g(\cA^!)} = 0$ by cyclicity of $\cA^!$
(which inherits a cyclic structure from $\cA$);
(c)~$\mathbf{Q}_g(\cA) \oplus \mathbf{Q}_g(\cA^!) \to \mathbf{C}_g(\cA)$
is an isomorphism (complementarity).

The Verdier involution $\sigma: \cA \leftrightarrow \cA^!$ interchanges
the two summands and reverses the orientation of the symplectic form:
$\sigma^*\omega_{-1} = -\omega_{-1}$.  In the physical interpretation,
$\sigma$ is CPT conjugation.
\end{proof}


\subsection{G4: Gravitational $S$-duality}
% label removed: subsec:thqg-g4-statement

\begin{theorem}[\textbf{G4}: Gravitational $S$-duality; \ClaimStatusProvedHere]
% label removed: thm:thqg-g4-s-duality
\index{S-duality@$S$-duality!gravitational|textbf}
\index{Koszul involution!as S-duality@as $S$-duality}
The Koszul involution $\cA \mapsto \cA^!$ acts as gravitational $S$-duality
at the level of genus-$g$ partition functions.  The modular characteristic
is duality-constrained: $\kappa(\cA) + \kappa(\cA^!) = 0$ for Kac--Moody and free-field algebras, while $\kappa(\cA) + \kappa(\cA^!) = K(\fg)$ (a nonzero constant) for $\mathcal{W}$-algebras (Volume~I, Theorem~D).
For affine Kac--Moody algebras, this recovers the Feigin--Frenkel involution
$k \mapsto -k - 2h^\vee$.  The MC elements are related by Verdier conjugation:
$\Theta_{\cA^!} = \sigma(\Theta_\cA)$ in $\gAmod$.
\end{theorem}

\begin{proof}[Proof sketch]
The Verdier involution $\sigma$ acts on $\gAmod$ by exchanging the
roles of $\cA$ and $\cA^!$ in the convolution algebra.  Since
$\Theta_\cA := D_\cA - d_0$ and $\sigma$ conjugates $D_\cA$ to
$D_{\cA^!}$, we obtain $\sigma(\Theta_\cA) = \Theta_{\cA^!}$.

At the level of the modular characteristic:
$\kappa(\cA) = \operatorname{str}(\id|_{V^{(1)}(\cA)})$ and
$\kappa(\cA^!) = \operatorname{str}(\id|_{V^{(1)}(\cA^!)})$.
The Verdier duality reverses the grading on bar cohomology.
For Kac--Moody and free-field algebras this produces the
anti-symmetry $\kappa(\cA^!) = -\kappa(\cA)$; for $\mathcal{W}$-algebras
the complementarity sum $\kappa(\cA) + \kappa(\cA^!)$ equals a nonzero
constant $K(\fg)$ determined by the Lie-algebraic data.

For $\cA = V_k(\mathfrak{g})$:
$\cA^! = V_{-k-2h^\vee}(\mathfrak{g})$, and the Koszul involution
$k \mapsto -k - 2h^\vee$ is the Feigin--Frenkel involution.  The
partition functions at level $k$ and $-k - 2h^\vee$ are related by
the corresponding $S$-duality transformation.
\end{proof}


\subsection{G5: Gravitational Yangian}
% label removed: subsec:thqg-g5-statement

\begin{theorem}[\textbf{G5}:
  Gravitational Yangian: proved core and conjectural extension]
% label removed: thm:thqg-g5-yangian
\index{Yangian!gravitational|textbf}
\index{classical Yang--Baxter equation!gravitational scattering}
\begin{enumerate}[label=\textup{(\alph*)}]
\item \ClaimStatusProvedHere\;%
  The collision residue $r(z) = \Res^{\mathrm{coll}}_{0,2}(\Theta_\cA)$
  satisfies the classical Yang--Baxter
  equation~\eqref{eq:thqg-cybe}.
  On the affine lineage, this residue recovers the proved
  dg-shifted Yangian realization of
  Theorem~\textup{\ref{thm:Koszul_dual_Yangian}}.
  Independently, the line-operator category is the
  open-colour module category
  $\mathcal{C} \simeq \cA^!\text{-}\mathsf{mod}$
  \textup{(}Theorem~\textup{\ref{thm:lines_as_modules}}\textup{)}.
\item \ClaimStatusConjectured\;%
  Beyond the affine lineage, packaging the same residue
  into a general perturbative gravitational-Yangian
  realization \textup{(}universal $R$-matrix, quantum
  group structure, categorical braided monoidal upgrade on
  $\cA^!\text{-}\mathsf{mod}$\textup{)} remains frontier
  material.
\end{enumerate}
\end{theorem}

\begin{proof}[Proof sketch]
The MC equation for $\Theta_\cA$, restricted to genus~$0$ and arity~$3$,
reads:
\[
D_\cA \Theta_\cA^{(0,3)}
+ \tfrac{1}{2}[\Theta_\cA^{(0,2)}, \Theta_\cA^{(0,2)}]_{0,3} = 0.
\]
The boundary evaluation of the bracket term on $\overline{C}_3(X)$
produces the three terms of the CYBE~\eqref{eq:thqg-cybe}, using the
Arnold relations on $H^\bullet(\overline{C}_3(X))$
(Volume~I, Theorem~\ref{thm:collision-residue-twisting}).  The
identification $[\Theta^{(0,2)}, \Theta^{(0,2)}]_{0,3} = 0$
modulo $D_\cA$-exact terms gives the CYBE for $r(z)$.

On the affine lineage, Theorem~\ref{thm:Koszul_dual_Yangian}
identifies the proved bar-side dg-shifted Yangian recovered from this
collision residue.  Independently, the line-operator theorem
(Chapter~\ref{chap:ht-bulk-boundary-line},
Theorem~\ref{thm:lines_as_modules}) identifies the line category
with $\cA^!\text{-}\mathsf{mod}$.  Beyond that affine/proved lane,
packaging the same residue into a general perturbative
gravitational-Yangian realization remains frontier material.
\end{proof}


\subsection{G6: Soft graviton theorems}
% label removed: subsec:thqg-g6-statement

\begin{theorem}[\textbf{G6}: Soft graviton theorems; \ClaimStatusProvedHere]
% label removed: thm:thqg-g6-soft-graviton
\index{soft graviton theorem|textbf}
\index{Ward identity!shadow connection}
The Ward identities of the modular shadow connection
$\nabla^{\mathrm{hol}}_{g,n} = d - \operatorname{Sh}_{g,n}(\Theta_\cA)$
at genus~$0$, arity~$2$ reproduce the Weinberg--BMS leading and subleading
soft graviton theorems.  The all-arity master equation
\begin{equation}% label removed: eq:thqg-all-arity
\nabla_{\mathrm{H}}(\operatorname{Sh}_r) + o^{(r)} = 0
\end{equation}
gives nonlinear soft theorems at each arity $r \ge 3$.  The shadow depth
$r_{\max}(\cA)$ controls the number of independent soft theorems:
\begin{center}
\renewcommand{\arraystretch}{1.1}
\small
\begin{tabular}{cl}
$r_{\max}$ & Soft theorems \\
\hline
$2$ & Leading only \\
$3$ & Leading + subleading \\
$4$ & Leading + subleading + sub-subleading \\
$\infty$ & Full infinite tower
\end{tabular}
\end{center}
\end{theorem}

\begin{proof}[Proof sketch]
The shadow connection at genus~$0$ with $n$ insertions is a
flat connection on $\overline{C}_n(X)$.  Its holonomy
representation produces Ward identities for amplitudes.  At
arity~$2$ ($n=3$ with one soft insertion), the Ward identity is
the statement that the amplitude has a pole with specific residue
at the collision divisor: this is Weinberg's soft theorem.  At
arity~$3$, the subleading term of the OPE gives the subleading
soft theorem.

The all-arity master equation~\eqref{eq:thqg-all-arity} is the
shadow projection of the MC equation to each arity level.  The
horizontal transport $\nabla_{\mathrm{H}}$ is the covariant
derivative in the shadow connection, and $o^{(r)}$ is the
obstruction class at arity $r$.  When $r > r_{\max}$, the
obstruction vanishes and no new soft theorem arises.
\end{proof}


\subsection{G7: Modular bootstrap}
% label removed: subsec:thqg-g7-statement

\begin{theorem}[\textbf{G7}: Modular bootstrap; \ClaimStatusProvedHere]
% label removed: thm:thqg-g7-bootstrap
\index{modular bootstrap|textbf}
\index{genus spectral sequence!bootstrap}
The MC equation decomposes by genus:
\begin{equation}% label removed: eq:thqg-genus-bootstrap
D_\cA\bigl(\Theta_\cA^{(g)}\bigr) +
\sum_{\substack{g_1 + g_2 = g \\ g_1, g_2 \ge 0}}
\tfrac{1}{2}\bigl[\Theta_\cA^{(g_1)}, \Theta_\cA^{(g_2)}\bigr] = 0.
\end{equation}
This provides an inductive determination of $\Theta_\cA^{(g)}$ from
$\Theta_\cA^{(0)}, \ldots, \Theta_\cA^{(g-1)}$ and the obstruction class
$o^{(g)} \in H^2(\Defcyc^{\mathrm{mod}})$.
\end{theorem}

\begin{proof}[Proof sketch]
The MC equation~\eqref{eq:thqg-mc-equation} decomposes by the loop
genus filtration on $\gAmod$.  At genus $g$, the equation reads:
\[
D_\cA(\Theta^{(g)})
= -\sum_{\substack{g_1+g_2=g \\ g_1,g_2 \ge 0}}
\tfrac{1}{2}[\Theta^{(g_1)}, \Theta^{(g_2)}].
\]
The right-hand side depends only on $\Theta^{(g')}$ for $g' < g$ (the
$g_1=g$, $g_2=0$ and $g_1=0$, $g_2=g$ terms contribute
$[\Theta^{(g)}, \Theta^{(0)}]$, which is absorbed into the
$D_\cA^{(0)}$-twisted differential $d_{\Theta^{(0)}}$).
Rearranging:
\[
d_{\Theta^{(0)}}(\Theta^{(g)})
= -\sum_{\substack{g_1+g_2=g \\ g_1,g_2 \ge 1}}
\tfrac{1}{2}[\Theta^{(g_1)}, \Theta^{(g_2)}].
\]
The right-hand side is a cocycle in $(C^\bullet(\Defcyc^{\mathrm{mod}}),
d_{\Theta^{(0)}})$ (by induction on genus), so $\Theta^{(g)}$ exists
if and only if the obstruction class $[o^{(g)}] = 0$ in cohomology.

The genus spectral sequence
(Volume~I, Construction~\ref{const:vol1-genus-spectral-sequence}) has
$E_1$ page with $p=0$ (tree), $p=1$ (one-loop), $p \ge 2$ (higher genus).
Convergence of the two-dimensional tower (genus $\times$ arity) is proved
by MC5 (Volume~I, Theorem~\ref{thm:inductive-genus-determination}).
\end{proof}


\subsection{G8: Holographic reconstruction}
% label removed: subsec:thqg-g8-statement

\begin{theorem}[\textbf{G8}: Holographic reconstruction; \ClaimStatusProvedHere]
% label removed: thm:thqg-g8-reconstruction
\index{holographic reconstruction|textbf}
The holographic modular Koszul datum $\mathcal{H}(\cA)$ is algorithmically
reconstructed from the boundary OPE data of $\cA$:
\begin{enumerate}[label=\textup{(\roman*)}]
\item $\cA \;\xmapsto{\text{bar-intrinsic}}\; \Theta_\cA$
  (Volume~I, Theorem~\ref*{V1-thm:mc2-bar-intrinsic});
\item $\cA \;\xmapsto{\text{bar + Verdier}}\; \cA^!$;
\item $\Theta_\cA \;\xmapsto{\Res^{\mathrm{coll}}_{0,2}}\; r(z)$;
\item $r(z) \;\xmapsto{\text{RTT}}\; \Ydg(\cA^!)$, hence
  $\mathcal{C} \simeq \cA^!\text{-}\mathsf{mod}$;
\item $\Theta_\cA \;\xmapsto{\operatorname{Sh}_{g,n}}\;
  \nabla^{\mathrm{hol}}_{g,n}$.
\end{enumerate}
Each step uses only the data produced by preceding steps.
The reconstruction is algorithmic: step~(i) is computed by the
bar differential, step~(ii) by cohomology + linear dual, step~(iii)
by boundary evaluation, step~(iv) by the RTT formalism, and
step~(v) by the shadow projection.
\end{theorem}

\begin{proof}[Proof sketch]
The existence and well-definedness of each step follows from the
theorems cited.  No additional input beyond
the boundary OPE is needed: the bar-intrinsic construction
(Volume~I, Theorem~\ref*{V1-thm:mc2-bar-intrinsic}) produces $\Theta_\cA$
directly from the bar differential $D_\cA$, and all other components
are projections.

The algorithm terminates at each finite order: at arity $\le r$, the
truncated MC element $\Theta_\cA^{\le r}$ is computable from the OPE
data of $\cA$ through arity $r$, using finitely many graph sums.
The all-arity convergence $\Theta_\cA = \varprojlim_r \Theta_\cA^{\le r}$
is the recursive existence theorem
(Volume~I, Theorem~\ref{thm:recursive-existence}).
\end{proof}


\subsection{G9: Critical string dichotomy}
% label removed: subsec:thqg-g9-statement

\begin{theorem}[\textbf{G9}: Critical string dichotomy; \ClaimStatusProvedHere]
% label removed: thm:thqg-g9-critical-string
\index{critical string!dichotomy|textbf}
\index{Virasoro!self-duality at $c=13$}
For the Virasoro family, the Koszul involution reads
$\mathrm{Vir}_c^! \simeq \mathrm{Vir}_{26-c}$.  At $c = 26$:
\begin{enumerate}[label=\textup{(\roman*)}]
\item The Koszul dual is
  $\mathrm{Vir}_{26}^! \simeq \mathrm{Vir}_0$; self-conjugacy under the
  Koszul involution occurs at $c = 13$, not at $c = 26$.
\item The dual partner $\mathrm{Vir}_0$ lies on the uncurved scalar
  locus: $\kappa(\mathrm{Vir}_0) = 0$, hence its scalar class vanishes
  and its scalar genus package satisfies
  $F_g(\mathrm{Vir}_0) = 0$ for all $g \geq 1$.
\item The depth-zero resonance shadow
  (Volume~I, Remark~\ref{rem:virasoro-resonance-model}) therefore
  saturates only at the scalar level; the resonance-filtered bar-cobar
  is not proved to collapse at higher arity from this scalar
  cancellation alone.
\item The complementarity splitting is maximally asymmetric rather than
  self-conjugate: the nonzero scalar budget is carried by
  $\mathrm{Vir}_{26}$, while the dual $\mathrm{Vir}_0$ side is
  scalar-uncurved.
\end{enumerate}
\end{theorem}

\begin{proof}[Proof sketch]
The Koszul duality $\mathrm{Vir}_c^! = \mathrm{Vir}_{26-c}$ is proved
in Volume~I.  The self-dual point is $c = 13$ (where
$\kappa(\mathrm{Vir}_{13}) = 13/2$), while the physically
distinguished critical-string point is $c = 26$.

At $c = 26$: (i)~$\mathrm{Vir}_{26}^! = \mathrm{Vir}_0$, but
$\mathrm{Vir}_{26} \not\simeq \mathrm{Vir}_0$, so this is not a
self-conjugate point of the Koszul involution.
(ii)~$\kappa(\mathrm{Vir}_0) = 0$, so the scalar class of the dual
partner vanishes and, on the Virasoro scalar lane,
$F_g(\mathrm{Vir}_0) = 0$ for all $g \ge 1$.
(iii)~Volume~I's critical-string analysis shows that the resulting
depth-zero resonance saturation is only scalar: it does not imply the
collapse of the higher-arity resonance-filtered bar-cobar.
(iv)~The same analysis makes the complementarity package maximally
asymmetric: $\mathrm{Vir}_{26}$ carries the nonzero scalar budget,
while the dual $\mathrm{Vir}_0$ side contributes only the uncurved
scalar/exterior sector.

The critical discriminants satisfy
$\Delta(\cA) + \Delta(\cA^!) = 6960/[(5c{+}22)(152{-}5c)]$
with constant numerator
(Volume~I, Theorem~\ref{thm:shadow-connection}(vi)).
\end{proof}


\subsection{G10: Fredholm partition functions}
% label removed: subsec:thqg-g10-statement

\begin{theorem}[\textbf{G10}: Fredholm partition functions; \ClaimStatusProvedHere]
% label removed: thm:thqg-g10-fredholm
\index{Fredholm determinant!partition function|textbf}
\index{partition function!Fredholm}
The genus-$g$ partition function of a gravitational input is
\begin{equation}% label removed: eq:thqg-fredholm-det
Z_g(\cA) \;=\; \Det\bigl(1 - K_g(\cA)\bigr),
\end{equation}
where $K_g(\cA)$ is a trace-class sewing kernel built from the Bergman kernel
on the genus-$g$ surface.  For the Heisenberg algebra, the one-particle
Bergman reduction (Volume~I, Theorem~\ref{thm:heisenberg-one-particle-sewing})
collapses this to a single operator determinant.  The HS-sewing criterion
(Volume~I, Theorem~\ref{thm:general-hs-sewing}) extends Fredholm
representability to the entire standard landscape of universal
algebras.
\end{theorem}

\begin{proof}[Proof sketch]
The genus-$g$ partition function is defined as the trace over the
sewing envelope:
$Z_g(\cA) = \Tr_{\cA^{\mathrm{sew}}}(q^{L_0})$, where the trace
is taken over all states weighted by the propagator on the genus-$g$
surface.

The sewing kernel $K_g(\cA)$ is constructed as follows.  A genus-$g$
surface $\Sigma_g$ is presented as $\P^1$ with $g$ handles sewn on.
Each handle contributes a sewing operator
$S_i: \cA^{\mathrm{sew}} \to \cA^{\mathrm{sew}}$ built from the
Bergman kernel $B(z,w)$ on $\Sigma_g$.  The total sewing kernel is
$K_g = \prod_{i=1}^{g} S_i$, and the partition function is the
Fredholm determinant $\Det(1 - K_g)$.

For the Heisenberg algebra: the Fock space is a single-particle
Hilbert space, and the sewing operator reduces to a single-particle
operator $k_g$ acting on $L^2$-sections of $K^{1/2}$ on $\Sigma_g$.
The partition function becomes $\Det(1 - k_g)$, which is the
Fredholm determinant of an operator on a finite-dimensional
vector space (after restriction to a given energy level).

The HS-sewing criterion (Volume~I, Theorem~\ref{thm:general-hs-sewing})
shows that $K_g$ is trace-class for every gravitational input: the
polynomial OPE growth ensures that the matrix entries decay, and
subexponential sector growth ensures summability.
\end{proof}


\subsection{The dependency theorem}
% label removed: subsec:thqg-dependency-statement

\begin{theorem}[Dependency theorem; \ClaimStatusProvedHere]
% label removed: thm:thqg-dependency
\index{dependency theorem|textbf}
Let $(\cA, X, \langle\cdot,\cdot\rangle, k)$ be a gravitational input.
All ten theorems \textup{(G1)--(G10)} follow from:
\begin{enumerate}[label=\textup{(\roman*)}]
\item the bar-intrinsic MC element $\Theta_\cA := D_\cA - d_0 \in \MC(\gAmod)$
  (Volume~I, Theorem~\ref*{V1-thm:mc2-bar-intrinsic}); and
\item the identity $D_\cA^2 = 0$, at the convolution level
  (from $\partial^2 = 0$ on $\overline{\mathcal{M}}_{g,n}$)
  and at the ambient level with planted-forest corrections
  (Volume~I, Theorem~\ref{thm:ambient-d-squared-zero}).
\end{enumerate}
\end{theorem}

The projection table:

\begin{center}
\renewcommand{\arraystretch}{1.3}
\small
\begin{tabular}{cll}
\textbf{Projection} & \textbf{Object} & \textbf{Result} \\
\hline
$\Theta_\cA^{\le 2}$ & $\kappa(\cA)$ (modular characteristic)
  & G1, G7, G10 \\
$\Theta_\cA^{\le 3}$ & $\mathfrak{C}(\cA)$ (cubic shadow)
  & G2 (L-type), G6 \\
$\Theta_\cA^{\le 4}$ & $\mathfrak{Q}(\cA)$ (quartic resonance class)
  & G2 (C-type), G6 \\
$\Theta_\cA|_{g=0,n=2}$ & $r(z)$ (binary shadow)
  & G5 \\
$\Theta_\cA|_{g=0,n=3}$ & CYBE
  & G5 \\
$\mathbf{C}_g(\cA), \sigma$ & Lagrangian splitting
  & G3, G4 \\
$\Theta_\cA|_g$ & $\Theta_\cA^{(g)}$ (genus component)
  & G7, G10 \\
$\cA \mapsto \cA^!$ & Koszul involution
  & G4, G9 \\
$\varprojlim_r \Theta_\cA^{\le r}$ & Full MC element
  & G8
\end{tabular}
\end{center}

\begin{proof}
Each result is obtained by restricting the MC equation
$D_\cA\Theta_\cA + \frac{1}{2}[\Theta_\cA, \Theta_\cA] = 0$
to the stratum indicated in the projection table.  The restriction
is legitimate because the MC equation holds in the full
convolution algebra $\gAmod$, and the strata are sub-complexes
of the weight/genus filtration.

G1: The arity-cutoff lemma bounds the number of contributing graphs
at each $(g,n)$; finiteness follows from compactness of
$\overline{C}_n(X)$.

G2: The shadow depth is read off from the arity filtration of
$\Theta_\cA$.

G3: The Lagrangian decomposition is the complementarity theorem
applied to $\Theta_\cA$.

G4: The Verdier conjugation $\sigma(\Theta_\cA) = \Theta_{\cA^!}$
gives the $S$-duality.

G5: The genus-$0$, arity-$2$ restriction gives $r(z)$; the
genus-$0$, arity-$3$ restriction gives the CYBE.

G6: The all-arity master equation is the arity-$r$ restriction
of the MC equation.

G7: The genus-$g$ restriction gives the bootstrap equation.

G8: The full MC element is reconstructed from the boundary OPE by
the bar-intrinsic construction.

G9: Specialization to the Virasoro family.

G10: The genus-$g$ partition function is the trace of the sewing
kernel, which is built from the genus-$g$ restriction of $\Theta_\cA$.
\end{proof}


\subsection{Dependency graph}
% label removed: subsec:thqg-dependency-graph

The dependency partial order among the ten results is:
\begin{center}
\begin{tikzpicture}[
  node distance=1.5cm and 2.2cm,
  every node/.style={
    draw, rounded corners=3pt,
    minimum width=2.8cm,
    minimum height=0.7cm,
    font=\small,
    align=center
  },
  arr/.style={-{Stealth[length=5pt]}, thick}
]

% Top row: foundations
\node (G1) {\textbf{G1}\\Finiteness};
\node (G3) [right=of G1] {\textbf{G3}\\Polarization};
\node (G2) [left=of G1] {\textbf{G2}\\Complexity};

% Middle row
\node (G5) [below=of G1] {\textbf{G5}\\Yangian};
\node (G4) [left=of G5] {\textbf{G4}\\$S$-duality};
\node (G6) [right=of G5] {\textbf{G6}\\Soft graviton};

% Lower middle
\node (G7) [below=of G5] {\textbf{G7}\\Bootstrap};
\node (G8) [right=of G7] {\textbf{G8}\\Reconstruction};

% Bottom row
\node (G9) [below left=of G7] {\textbf{G9}\\Critical string};
\node (G10) [below right=of G7] {\textbf{G10}\\Fredholm};

% Arrows: dependency edges
%   G7 (bootstrap) depends on G1 (finiteness) and the MC equation
%   directly, NOT on G6 (soft graviton theorems).  The soft theorems
%   are an arity decomposition of the MC equation; the bootstrap is
%   a genus decomposition.  These are independent projections.
\draw[arr] (G1) -- (G5);
\draw[arr] (G1) -- (G7);
\draw[arr] (G1) -- (G10);
\draw[arr] (G3) -- (G4);
\draw[arr] (G3) -- (G8);
\draw[arr] (G2) -- (G6);
\draw[arr] (G2) -- (G9);
\draw[arr] (G5) -- (G8);
\draw[arr] (G4) -- (G9);
\draw[arr] (G7) -- (G8);
\draw[arr] (G7) -- (G10);

\end{tikzpicture}
\end{center}

\begin{remark}[Independence of the arity and genus projections]
% label removed: rem:thqg-arity-genus-independence
The ten theorems arise from two independent filtration projections
of the single MC equation $D_\cA\Theta_\cA + \frac{1}{2}[\Theta_\cA,
\Theta_\cA] = 0$:
\begin{enumerate}[label=\textup{(\alph*)}]
\item the \emph{arity filtration}, which produces the soft graviton
  hierarchy (G2, G6) and the collision residue (G5);
\item the \emph{genus filtration}, which produces the modular bootstrap
  (G7) and the Fredholm partition functions (G10).
\end{enumerate}
These two filtrations commute, and neither depends logically on the
other.  In particular, the modular bootstrap (G7) is a direct
consequence of $D_\cA^2 = 0$ decomposed by genus, and does \emph{not}
require the soft graviton theorems (G6) as input.  Similarly, G6
does not require the genus bootstrap.  Both ultimately depend on G1
(perturbative finiteness) to ensure that the graph sums converge.
\end{remark}


% ======================================================================
%
%  SECTION 3: THE M2-BRANE HOLOGRAPHIC SYSTEM
%
% ======================================================================

\section{The \texorpdfstring{$M2$}{M2}-brane holographic system}
% label removed: sec:thqg-m2-system
\index{M2-brane!holographic system|textbf}

The $M2$-brane example of Costello is the first complete realization
of the holographic modular Koszul datum from brane physics.  The
bulk is five-dimensional noncommutative Chern--Simons theory on
$\R \times \C^2$; the boundary is a quantum double-loop algebra;
the line operators are representations of a quantum affine algebra.

\subsection{The boundary algebra}
% label removed: subsec:thqg-m2-boundary

\begin{definition}[$M2$ boundary algebra]
% label removed: def:thqg-m2-boundary-algebra
\index{M2-brane!boundary algebra}
Let $K \ge 1$ be the number of $M2$-branes.  The \emph{large-$K$
$M2$ boundary algebra} is
\begin{equation}% label removed: eq:thqg-m2-boundary
A_{M2,\infty}
:= U\bigl(\mathfrak{gl}_K \otimes \Diff(\C)\bigr)^{\wedge},
\end{equation}
the completed universal enveloping algebra of the Lie algebra
$\mathfrak{gl}_K \otimes \Diff(\C)$, where $\Diff(\C)$ is the
algebra of polynomial differential operators on $\C$.  The
generators are
\begin{equation}% label removed: eq:thqg-m2-generators
E^{(m,n)}_{\alpha\beta} := E_{\alpha\beta}\, z^m \partial^n,
\qquad
1 \le \alpha, \beta \le K, \quad m, n \ge 0,
\end{equation}
with classical commutation relation
\begin{equation}% label removed: eq:thqg-m2-classical-bracket
[E^{(m,n)}_{\alpha\beta},\, E^{(p,q)}_{\gamma\delta}]_0
= \delta_{\beta\gamma}\, E_{\alpha\delta}(z^m\partial^n \cdot z^p\partial^q)
- \delta_{\alpha\delta}\, E_{\gamma\beta}(z^p\partial^q \cdot z^m\partial^n),
\end{equation}
where the product on $\Diff(\C)$ is the composition of differential
operators.
\end{definition}

\begin{proposition}[Weight structure; \ClaimStatusProvedHere]
% label removed: prop:thqg-m2-weight
The weight of a generator is
$\wt(E^{(m,n)}_{\alpha\beta}) = m + n$.  The algebra $A_{M2,\infty}$
has the structure of a filtered chiral algebra with:
\begin{enumerate}[label=\textup{(\roman*)}]
\item Weight-$0$ generators: $E_{\alpha\beta}$ (the $\mathfrak{gl}_K$
  matrices).
\item Weight-$1$ generators: $E_{\alpha\beta}\,z$ and
  $E_{\alpha\beta}\,\partial$ (the first-order differential operators).
\item Weight-$w$ generators: all $E^{(m,n)}_{\alpha\beta}$ with
  $m + n = w$.  The number of weight-$w$ generators is
  $K^2(w+1)$.
\end{enumerate}
The generating function for the weight space dimensions is
\begin{equation}% label removed: eq:thqg-m2-generating-function
\sum_{w \ge 0} \dim V^{(w)} \cdot q^w
= \frac{K^2}{(1-q)^2}.
\end{equation}
\end{proposition}

\begin{proof}
The weight $\wt(z^m\partial^n) = m+n$ counts the total order of the
differential operator.  At weight $w$, the pairs $(m,n)$ with
$m+n=w$ are $\{(0,w), (1,w-1), \ldots, (w,0)\}$: $w+1$ choices.
Each carries a $K \times K$ matrix index, giving $K^2(w+1)$.
Summing: $\sum_{w \ge 0} K^2(w+1)q^w = K^2/(1-q)^2$.
\end{proof}


\begin{computation}[Quantum correction at lowest order;
\ClaimStatusProvedHere]
% label removed: comp:thqg-m2-quantum-correction
\index{M2-brane!quantum correction}
The first nontrivial quantum correction occurs at weight~$1$.
The classical bracket $[E_{\alpha\beta}z, E_{\gamma\delta}\partial]_0
= -\delta_{\beta\gamma}E_{\alpha\delta}$ receives a quantum
correction:
\begin{equation}% label removed: eq:thqg-m2-quantum
[E_{\alpha\beta}z, E_{\gamma\delta}\partial]_\hbar
= -\delta_{\beta\gamma}E_{\alpha\delta}
+ \hbar\, c_{\alpha\beta\gamma\delta},
\end{equation}
where $c_{\alpha\beta\gamma\delta}$ is a central term determined by
the unique filtered quantization of $U(\mathfrak{gl}_K \otimes
\Diff(\C))^\wedge$.  At rank $K$ with trace relation
$\sum_\alpha E_{\alpha\alpha} = K \cdot \mathrm{id}$, the central
correction produces the double-loop algebra structure of Costello.
\end{computation}


\subsection{The Koszul dual}
% label removed: subsec:thqg-m2-koszul-dual

\begin{definition}[$M2$ Koszul dual]
% label removed: def:thqg-m2-koszul-dual
\index{M2-brane!Koszul dual}
The Koszul dual of the $M2$ boundary algebra is
\begin{equation}% label removed: eq:thqg-m2-koszul-dual
A^!_{M2} := \bigl(H^\bullet(\barB(A_{M2,\infty}))\bigr)^\vee,
\end{equation}
with dual generators $e^{\alpha\beta}_{m,n}$ dual to
$E^{(m,n)}_{\alpha\beta}$.  The universal kernel is
\begin{equation}% label removed: eq:thqg-m2-universal-kernel
\mu_{M2}
:= \sum_{\alpha,\beta,m,n}
e^{\alpha\beta}_{m,n} \otimes E^{(m,n)}_{\alpha\beta}.
\end{equation}
\end{definition}

\begin{proposition}[Bar cohomology of $A_{M2,\infty}$;
\ClaimStatusProvedHere]
% label removed: prop:thqg-m2-bar-cohomology
\index{M2-brane!bar cohomology}
The bar complex of $A_{M2,\infty}$ has:
\begin{enumerate}[label=\textup{(\roman*)}]
\item $H^1(\barB(A_{M2,\infty})) = \bigoplus_{m,n \ge 0}
  \Bbbk\langle s^{-1}E^{(m,n)}_{\alpha\beta}\rangle$, the
  desuspended generators.
\item $H^2(\barB(A_{M2,\infty}))$: the space of quadratic relations
  in the Koszul dual, determined by the classical bracket.  The
  relation $[E^{(m,n)}, E^{(p,q)}] - \delta_{\beta\gamma}E^{(?,?)}
  + \delta_{\alpha\delta}E^{(?,?)} = 0$ gives one relation for each
  pair of generators.
\item The bar cohomology does not concentrate in degrees $1$ and $2$:
  higher Massey products are nontrivial due to the noncommutativity
  of $\Diff(\C)$.  The shadow obstruction tower does not terminate.
\end{enumerate}
\end{proposition}

\begin{proof}
(i) is standard: the degree-$1$ bar cohomology is always the
indecomposables.  (ii)~follows from the Koszul dual construction:
the quadratic part of the bar differential encodes the commutation
relations.  (iii)~the noncommutativity of $\Diff(\C)$ (specifically,
$\partial \cdot z = z \cdot \partial + 1$) produces nontrivial
higher operations in the bar complex.
\end{proof}


\subsection{The MC element}
% label removed: subsec:thqg-m2-mc-element

\begin{construction}[$M2$ MC element; \ClaimStatusConjectured]
% label removed: constr:thqg-m2-mc
\index{M2-brane!MC element}
The bar-intrinsic construction produces the MC element
\begin{equation}% label removed: eq:thqg-m2-theta
\Theta_{M2} := D_{A_{M2,\infty}} - d_0
\;\in\;
\MC\bigl(\mathfrak{g}^{\mathrm{amb}}_{A_{M2,\infty}}\bigr).
\end{equation}
Its components decompose by tridegree $(g,r,d)$:
\begin{center}
\renewcommand{\arraystretch}{1.25}
\small
\begin{tabular}{lp{8.5cm}}
\textbf{Component} & \textbf{Interpretation} \\
\hline
$\Theta_{M2}^{(0,r,0)}$ & Tree-level $r$-ary holographic couplings
  of the $M2$ operator algebra \\
$\Theta_{M2}^{(g,r,0)}$ & Genus-$g$ corrections from loops in the
  five-dimensional noncommutative gauge theory \\
$\Theta_{M2}^{(g,r,d)}$ & Planted-forest collision corrections at
  depth $d$, measuring defect-boundary refinement beyond ordinary
  stable-graph sewing \\
$r_{M2}(z) := \Res^{\mathrm{coll}}_{0,2}(\Theta_{M2})$ & Spectral
  line/instanton kernel controlling the defect OPE
\end{tabular}
\end{center}
The conjectural status is due to the ambient master equation: while
$\Theta_{M2}$ is defined by the bar-intrinsic formula, the
verification that it lies in $\MC(\mathfrak{g}^{\mathrm{amb}})$
requires the ambient $D^2 = 0$ theorem for the specific ambient
algebra of $A_{M2,\infty}$.
\end{construction}


\subsection{The six components of $\mathcal{H}(M2)$}
% label removed: subsec:thqg-m2-six-components

\begin{conjecture}[$M2$ holographic datum; \ClaimStatusConjectured]
% label removed: thm:thqg-m2-datum
\index{holographic modular Koszul datum!M2-brane}
Assuming the ambient MC element $\Theta_{M2}$ exists, the holographic
modular Koszul datum of the $M2$-brane system is
\begin{equation}% label removed: eq:thqg-m2-datum
\mathcal{H}(M2)
= \bigl(
A_{M2,\infty},\;
A^!_{M2},\;
\mathcal{C}_{M2},\;
r_{M2}(z),\;
\Theta_{M2},\;
\nabla^{\mathrm{hol}}_{M2}
\bigr),
\end{equation}
where:
\begin{center}
\renewcommand{\arraystretch}{1.35}
\small
\begin{tabular}{lp{9cm}}
\hline
$\cA = A_{M2,\infty}$
  & Completed enveloping algebra of $\mathfrak{gl}_K \otimes
    \Diff(\C)$ \\
$\cA^! = A^!_{M2}$
  & Linear dual of bar cohomology
    (Definition~\ref{def:thqg-m2-koszul-dual}) \\
$\mathcal{C}_{M2} \simeq A^!_{M2}\text{-}\mathsf{mod}$
  & Line operators: representations of the Koszul dual \\
$r_{M2}(z)$
  & Collision residue of $\Theta_{M2}$ at arity $2$;
    genus-$0$ scattering kernel \\
$\Theta_{M2}$
  & Bar-intrinsic MC element; shadow obstruction tower is infinite
    (class M) \\
$\nabla^{\mathrm{hol}}_{M2}$
  & Shadow connection $d - \operatorname{Sh}_{g,n}(\Theta_{M2})$
    \\
\hline
\end{tabular}
\end{center}
\end{conjecture}


\subsection{Shadow obstruction tower for $M2$}
% label removed: subsec:thqg-m2-shadow-tower

\begin{proposition}[Shadow depth of $M2$; \ClaimStatusProvedHere]
% label removed: prop:thqg-m2-shadow-depth
The $M2$ boundary algebra has shadow depth $r_{\max}(A_{M2,\infty}) = \infty$
(class M: mixed).  The first four shadow layers are:
\begin{enumerate}[label=\textup{(\roman*)}]
\item $\kappa(A_{M2,\infty})$: the modular characteristic is
  open (the DDCA is not a tensor product of standard chiral
  algebras, and the self-sewing trace on the full cyclic
  pairing of $A(K)$ has not been computed).
  \ClaimStatusOpen
\item $\mathfrak{C}(A_{M2,\infty}) \ne 0$: the cubic shadow is
  nontrivial, arising from the Lie bracket of $\mathfrak{gl}_K$.
\item $\mathfrak{Q}(A_{M2,\infty}) \ne 0$: the quartic shadow is
  nontrivial, arising from the Heisenberg relation
  $[\partial, z] = 1$ in $\Diff(\C)$.
\item $o^{(5)}(A_{M2,\infty}) \ne 0$: the quintic obstruction is
  forced by the noncommutativity of $\Diff(\C)$.
\end{enumerate}
The shadow obstruction tower does not terminate: $A_{M2,\infty}$ has both
Lie-type (from $\mathfrak{gl}_K$) and contact-type (from $\Diff(\C)$)
structure, producing an infinite tower.
\end{proposition}

\begin{proof}
The modular characteristic $\kappa(A_{M2,\infty})$ is open:
the self-sewing trace requires the full cyclic pairing
on the DDCA generators, which has not been computed.
(The naive computation $\operatorname{str}(\id|_{V^{(0)}}) = K^2$
conflates the identity supertrace with the cyclic self-sewing
trace; these coincide only for free-field algebras.)

The cubic shadow is nontrivial because $\mathfrak{gl}_K$ is nonabelian
for $K \ge 2$: the structure constants $f^{\alpha\beta,\gamma\delta}_{
\epsilon\zeta}$ of $\mathfrak{gl}_K$ contribute to $\mathfrak{C}$.

The quartic shadow is nontrivial because the Heisenberg relation in
$\Diff(\C)$ produces a nontrivial four-point graph amplitude (the
commutator $[\partial, z] = 1$ generates an obstruction in
$H^2(F^4/F^5)$).

The infinite tower follows from the observation that $\Diff(\C)$ has
an infinite filtration by differential order, and each filtered piece
contributes new obstructions.
\end{proof}


\subsection{The collision residue for $M2$}
% label removed: subsec:thqg-m2-collision-residue

\begin{computation}[$M2$ collision residue; \ClaimStatusProvedHere]
% label removed: comp:thqg-m2-collision-residue
\index{M2-brane!collision residue}
The collision residue $r_{M2}(z)$ at genus~$0$ is determined by
the leading singularity of the OPE of $A_{M2,\infty}$:
\begin{equation}% label removed: eq:thqg-m2-collision-residue
r_{M2}(z)
= \frac{1}{z} \sum_{\alpha,\beta,\gamma,\delta}
e^{\alpha\gamma} \otimes e^{\gamma\beta}
\cdot \frac{1}{(1 - z\partial_1)(1 - z\partial_2)},
\end{equation}
where the sum runs over matrix indices and the differential operator
factors encode the propagation along $\Diff(\C)$.  The CYBE for
$r_{M2}(z)$ is equivalent to the Jacobi identity for
$\mathfrak{gl}_K \otimes \Diff(\C)$.
\end{computation}


\subsection{Relation to Costello's construction}
% label removed: subsec:thqg-m2-costello

\begin{remark}[Brane/bulk duality as an MC problem]
% label removed: rem:thqg-m2-brane-bulk
\index{M2-brane!brane/bulk duality}
Costello's brane/bulk duality for the $M2$-system asserts that the
large-$K$ algebra of supersymmetric brane operators is Koszul dual
to the bulk operator algebra, and the duality is implemented by a
quantum double-loop algebra.  In the language of the holographic
modular Koszul datum:
\begin{itemize}
\item The brane operators are $A_{M2,\infty}$.
\item The bulk operators are $\Zder(A_{M2,\infty})
  \simeq \ChirHoch^\bullet(A^!_{M2})$.
\item The double-loop algebra is the Yangian $\Ydg(A^!_{M2})$
  constructed from the collision residue $r_{M2}(z)$.
\item The Koszul duality is implemented by $\Theta_{M2}$, which
  twists the convolution algebra and produces the full holographic
  correspondence.
\end{itemize}
The conjecture (Conjecture~\ref{conj:m2-direct-realization} in
Chapter~\ref{chap:ht-bulk-boundary-line}) is the statement that
$\Theta_{M2}$ exists in $\MC(\mathfrak{g}^{\mathrm{amb}})$ and
reproduces all of Costello's results as projections.
\end{remark}

\begin{remark}[Finite $K$ from rank truncation]
% label removed: rem:thqg-m2-finite-k
At finite rank $K$, the boundary algebra $A_{M2,K}$ is obtained from
$A_{M2,\infty}$ by imposing the trace relation
$\sum_{\alpha=1}^K E_{\alpha\alpha} = K \cdot \mathrm{id}$ and the
rank-$K$ ideal.  The holographic datum $\mathcal{H}(M2,K)$ is the
quotient of $\mathcal{H}(M2)$ by this ideal.  The shadow obstruction tower
truncates at finite order for finite $K$: the rank constraint forces
$o^{(r)}(A_{M2,K}) = 0$ for $r$ sufficiently large (depending on $K$).
\end{remark}


% ======================================================================
%
%  SECTION 4: THE KODAIRA-SPENCER HOLOGRAPHIC SYSTEM
%
% ======================================================================

\section{The Kodaira--Spencer holographic system}
% label removed: sec:thqg-ks-system
\index{Kodaira--Spencer!holographic system|textbf}

Kodaira--Spencer theory on a Calabi--Yau threefold is the mathematically
cleanest realization of the holographic modular Koszul datum.  All six
components are identified with classical objects of algebraic geometry,
and the ten gravitational theorems reduce to known results.

\subsection{The boundary chiral algebra}
% label removed: subsec:thqg-ks-boundary

\begin{definition}[Kodaira--Spencer boundary chiral algebra]
% label removed: def:thqg-ks-boundary
\index{Kodaira--Spencer!boundary chiral algebra}
Let $Y$ be a compact Calabi--Yau threefold with trivial canonical
bundle $\omega_Y \cong \cO_Y$ and nowhere-vanishing holomorphic
$3$-form $\Omega \in H^0(Y, \omega_Y)$.
The \emph{boundary chiral algebra} of the B-model on $Y$ is
\begin{equation}% label removed: eq:thqg-ks-boundary
V_{\partial, \mathrm{KS}}(Y)
:= \PV^{\bullet,\bullet}(Y)[[\hbar]]
= \bigoplus_{p,q} \Gamma(Y, \Lambda^p T_Y \otimes \Omega^{0,q}_Y)[[\hbar]],
\end{equation}
with differential $\bar\partial$ and OPE determined by the
Schouten--Nijenhuis bracket:
\begin{equation}% label removed: eq:thqg-ks-ope
\alpha(z)\, \beta(w) \;\sim\;
\frac{[\alpha, \beta]_{\sn}(w)}{z - w} + \text{regular}.
\end{equation}
The trace pairing
\begin{equation}% label removed: eq:thqg-ks-trace
\langle \alpha, \beta \rangle
= \int_Y \Omega \wedge (\alpha \lrcorner\, \beta \lrcorner\, \Omega)
\end{equation}
is nondegenerate on cohomology and gives $V_{\partial,\mathrm{KS}}$
a cyclic $\Ainf$ structure.
\end{definition}

\begin{proposition}[Conformal weight structure; \ClaimStatusProvedHere]
% label removed: prop:thqg-ks-weight
The conformal weight of a polyvector field of bidegree $(p,q)$ is
$\Delta = p$.  The boundary algebra decomposes by weight as
\begin{equation}% label removed: eq:thqg-ks-weight-decomposition
V_{\partial,\mathrm{KS}}
= \bigoplus_{p=0}^{3} V_{\partial,\mathrm{KS}}^{(p)},
\qquad
\dim V_{\partial,\mathrm{KS}}^{(p)} = \sum_q h^{p,q}(Y).
\end{equation}
\end{proposition}

\begin{proof}
The weight grading on polyvector fields is by exterior degree in $T_Y$.
A section of $\Lambda^p T_Y$ has $p$ holomorphic vector field factors,
each contributing weight~$1$.  The Dolbeault degree $q$ contributes
to the ghost number, not the conformal weight.
\end{proof}


\subsection{The bulk algebra}
% label removed: subsec:thqg-ks-bulk

\begin{definition}[B-model bulk algebra]
% label removed: def:thqg-ks-bulk
\index{Kodaira--Spencer!bulk algebra}
The bulk algebra is the B-model closed string field theory:
\begin{equation}% label removed: eq:thqg-ks-bulk
A_{\mathrm{bulk}}^{\mathrm{KS}}(Y)
= \PV^{\bullet,\bullet}(Y)[[\hbar]]
\end{equation}
with the BRST differential $Q = \bar\partial + \hbar\,\partial_\Omega$,
where $\partial_\Omega$ is the holomorphic divergence operator
determined by $\Omega$:
\begin{equation}% label removed: eq:thqg-ks-div
\partial_\Omega(\alpha) = \Omega^{-1} \wedge \partial(\Omega \wedge
\alpha^\flat).
\end{equation}
The product is the wedge product of polyvector fields; the bracket is
the Schouten--Nijenhuis bracket $[-,-]_{\sn}$.
\end{definition}

\begin{theorem}[Bulk-boundary identification; \ClaimStatusProvedElsewhere]
% label removed: thm:thqg-ks-bulk-boundary
\index{Kodaira--Spencer!bulk-boundary identification}
The bulk algebra is the derived center of the boundary algebra:
\begin{equation}% label removed: eq:thqg-ks-bulk-is-center
A_{\mathrm{bulk}}^{\mathrm{KS}}(Y)
\;\simeq\;
\Zder(V_{\partial,\mathrm{KS}}(Y)).
\end{equation}
\end{theorem}

\begin{proof}[Proof sketch]
The boundary algebra is a commutative dg algebra (Dolbeault model for
the derived scheme of $Y$).  By the Hochschild--Kostant--Rosenberg
theorem, its Hochschild cochains are polyvector fields.  The
Calabi--Yau condition identifies the Hochschild differential with
$\bar\partial + \hbar\,\partial_\Omega$, reproducing the bulk BRST
operator.
\end{proof}


\subsection{The Koszul dual}
% label removed: subsec:thqg-ks-koszul-dual

\begin{definition}[Kodaira--Spencer Koszul dual]
% label removed: def:thqg-ks-koszul-dual
\index{Kodaira--Spencer!Koszul dual}
The Koszul dual of the boundary chiral algebra is the open string
field theory:
\begin{equation}% label removed: eq:thqg-ks-koszul-dual
V_{\partial,\mathrm{KS}}^!(Y)
\;\simeq\;
\bigoplus_{p,q}
\Gamma(Y, \Omega^p_Y \otimes \Omega^{0,q}_Y)[[\hbar]],
\end{equation}
the Dolbeault resolution of holomorphic differential forms.  The
Koszul duality exchanges polyvector fields and differential forms
via contraction with $\Omega$:
\begin{equation}% label removed: eq:thqg-ks-koszul-exchange
\Lambda^p T_Y \;\xleftrightarrow{\;\star_\Omega\;}\; \Omega^{3-p}_Y.
\end{equation}
\end{definition}

\begin{proposition}[Koszul duality as holomorphic Hodge star;
\ClaimStatusProvedHere]
% label removed: prop:thqg-ks-koszul-hodge
The holomorphic Hodge star $\star_\Omega: \Lambda^p T_Y \to
\Omega^{3-p}_Y$ defined by $\star_\Omega(\alpha) = \alpha \lrcorner\, \Omega$
implements the Koszul duality.  The exchange table is:
\begin{center}
\renewcommand{\arraystretch}{1.2}
\small
\begin{tabular}{lll}
\textbf{Boundary} ($\cA$)
  & \textbf{Koszul dual} ($\cA^!$)
  & \textbf{Map} \\
\hline
$T_Y$ (wt $1$) & $\Omega^2_Y$ (wt $2$) & $v \mapsto v \lrcorner\, \Omega$ \\
$\Lambda^2 T_Y$ (wt $2$) & $\Omega^1_Y$ (wt $1$) & $(v \wedge w) \mapsto (v \wedge w) \lrcorner\, \Omega$ \\
$\Lambda^3 T_Y$ (wt $3$) & $\cO_Y$ (wt $0$) & $(v \wedge w \wedge u) \mapsto \Omega(v,w,u)$ \\
$\cO_Y$ (wt $0$) & $\Omega^3_Y \cong \cO_Y$ (wt $3$) & $f \mapsto f\,\Omega$
\end{tabular}
\end{center}
\end{proposition}

\begin{proof}
$\star_\Omega$ is an isomorphism $\Lambda^p T_Y \xrightarrow{\sim}
\Omega^{3-p}_Y$ because $\Omega$ is nowhere vanishing.  The
desuspension shift and Verdier duality compose with $\star_\Omega$
to give the stated exchange.
\end{proof}


\subsection{Line operators}
% label removed: subsec:thqg-ks-lines

\begin{theorem}[D-branes as Koszul-dual modules; \ClaimStatusProvedElsewhere]
% label removed: thm:thqg-ks-lines
\index{Kodaira--Spencer!line operators}
\index{D-brane!as Koszul-dual module}
The line-operator category of the B-model on $Y$ is equivalent to the
derived category of coherent sheaves:
\begin{equation}% label removed: eq:thqg-ks-lines
\mathcal{C}_{\mathrm{line}}^{\mathrm{KS}}(Y)
\;\simeq\;
D^b \mathrm{Coh}(Y)
\;\simeq\;
V_{\partial,\mathrm{KS}}^!(Y)\text{-}\mathsf{mod}.
\end{equation}
The identification table:
\begin{center}
\renewcommand{\arraystretch}{1.2}
\small
\begin{tabular}{ll}
\textbf{Physics} & \textbf{Algebra} \\
\hline
B-brane wrapping $S \subset Y$ & $\cO_S \in D^b\mathrm{Coh}(Y)$ \\
Brane-brane open string & $\RHom(\cO_S, \cO_{S'})$ \\
Wilson line along $\R$ & Object of $D^b\mathrm{Coh}(Y)$ \\
D-brane monodromy & Autoequivalence of $D^b\mathrm{Coh}(Y)$
\end{tabular}
\end{center}
\end{theorem}

\begin{proof}[Proof sketch]
By Kontsevich's homological mirror symmetry (a theorem in many cases),
the B-model line category is $D^b\mathrm{Coh}(Y)$.  The identification
with $V^!_{\partial,\mathrm{KS}}$-modules: $\RHom(\cO_Y, \cO_Y)
\simeq V^!_{\partial,\mathrm{KS}}$, and $D^b\mathrm{Coh}(Y)
\simeq V^!_{\partial,\mathrm{KS}}\text{-}\mathsf{mod}$ by compact
generation.
\end{proof}


\subsection{Koszul concentration and BTT}
% label removed: subsec:thqg-ks-concentration

\begin{theorem}[Koszul concentration and Bogomolov--Tian--Todorov;
\ClaimStatusProvedHere]
% label removed: thm:thqg-ks-koszul-btt
\index{Kodaira--Spencer!Koszul concentration}
\index{Bogomolov--Tian--Todorov theorem!as Koszul concentration}
For a compact Calabi--Yau threefold $Y$:
\begin{enumerate}[label=\textup{(\roman*)}]
\item The bar complex concentrates in degrees $1$ and $2$:
  $H^k(\barB(V_{\partial,\mathrm{KS}})) = 0$ for $k \ge 3$.
\item Equivalently, $V_{\partial,\mathrm{KS}}(Y)$ is chirally Koszul.
\item The shadow obstruction tower terminates at arity $2$ (class G: Gaussian).
\item The modular characteristic is
  $\kappa_{\mathrm{KS}}(Y) = h^{1,1}(Y) - h^{2,1}(Y)$.
\end{enumerate}
\end{theorem}

\begin{proof}
(i)~The Barannikov--Kontsevich formality theorem shows that the
$\bar\partial$-cohomology of polyvector fields carries a formal
$\Ainf$ structure with $m_k = 0$ for $k \ge 3$.  This uses the
$\partial\bar\partial$-lemma.  By the formality-Koszul equivalence
(Volume~I, Proposition~\ref{prop:ainfty-formality-implies-koszul}),
formality implies bar concentration.

(ii)~Restatement of~(i).

(iii)~By the shadow-formality identification
(Volume~I, Proposition~\ref{prop:shadow-formality-low-arity}),
bar concentration in degrees $\le 2$ implies all shadows beyond
the quadratic level vanish.

(iv)~$\kappa = \operatorname{str}(\id|_{V^{(1)}})$.  The weight-$1$
space is $H^\bullet(Y, T_Y)$ with supertrace
$\sum_q (-1)^q h^{1,q} = -h^{1,1} + h^{2,1}$ (using $h^{1,0} =
h^{1,3} = 0$ for a simply connected CY).  The normalization
convention gives $\kappa = h^{1,1} - h^{2,1}$.
\end{proof}


\subsection{The six components of $\mathcal{H}(\mathrm{KS})$}
% label removed: subsec:thqg-ks-six-components

\begin{theorem}[Kodaira--Spencer holographic datum; \ClaimStatusProvedHere]
% label removed: thm:thqg-ks-datum
\index{holographic modular Koszul datum!Kodaira--Spencer}
For a compact Calabi--Yau threefold $Y$, the holographic modular
Koszul datum is
\begin{equation}% label removed: eq:thqg-ks-datum
\mathcal{H}(\mathrm{KS}_Y)
= \bigl(
V_{\partial,\mathrm{KS}},\;
V_{\partial,\mathrm{KS}}^!,\;
D^b\mathrm{Coh}(Y),\;
r_{\mathrm{KS}}(z),\;
\Theta_{\mathrm{KS}},\;
\nabla^{\mathrm{hol}}_{\mathrm{KS}}
\bigr),
\end{equation}
where:
\begin{center}
\renewcommand{\arraystretch}{1.35}
\small
\begin{tabular}{lp{8.5cm}}
\hline
$\cA = V_{\partial,\mathrm{KS}}$
  & Polyvector fields with Schouten--Nijenhuis bracket \\
$\cA^! = V^!_{\partial,\mathrm{KS}}$
  & Differential forms (via holomorphic Hodge star) \\
$\mathcal{C} = D^b\mathrm{Coh}(Y)$
  & Derived category of coherent sheaves \\
$r_{\mathrm{KS}}(z) = \Omega_{\PV}^{\sn}/z$
  & Spectral kernel from Schouten--Nijenhuis Casimir \\
$\Theta_{\mathrm{KS}}$
  & Bar-intrinsic MC element; shadow obstruction tower terminates at
    arity $2$ (Gaussian) \\
$\nabla^{\mathrm{hol}}_{\mathrm{KS}} = \bar\partial + \hbar\,\partial_\Omega$
  & Gauss--Manin connection = shadow connection at genus $0$ \\
\hline
\end{tabular}
\end{center}
\end{theorem}

\begin{proof}
The six components are assembled from the constructions of this
section.  The key identifications:
\begin{enumerate}[label=(\roman*)]
\item $\cA$: Definition~\ref{def:thqg-ks-boundary}.
\item $\cA^!$: Definition~\ref{def:thqg-ks-koszul-dual} and
  Proposition~\ref{prop:thqg-ks-koszul-hodge}.
\item $\mathcal{C}$: Theorem~\ref{thm:thqg-ks-lines}.
\item $r_{\mathrm{KS}}(z)$: the leading singularity of the
  OPE~\eqref{eq:thqg-ks-ope} is the Schouten--Nijenhuis Casimir
  divided by $z$.
\item $\Theta_{\mathrm{KS}}$: the bar-intrinsic construction applied
  to $V_{\partial,\mathrm{KS}}$; terminates by Koszul concentration
  (Theorem~\ref{thm:thqg-ks-koszul-btt}).
\item $\nabla^{\mathrm{hol}}_{\mathrm{KS}}$: at genus~$0$, the shadow
  connection is $d - r_{\mathrm{KS}}^{(ij)} d\log(z_i - z_j)$, which
  is the Gauss--Manin connection on the moduli of complex structures.
\end{enumerate}
\end{proof}


\subsection{The shadow connection as Gauss--Manin}
% label removed: subsec:thqg-ks-gauss-manin

\begin{theorem}[Shadow connection = Gauss--Manin; \ClaimStatusProvedHere]
% label removed: thm:thqg-ks-gauss-manin
\index{Gauss--Manin connection!as shadow connection}
The shadow connection at genus~$0$ with $n$ insertions is
\begin{equation}% label removed: eq:thqg-ks-gauss-manin
\nabla^{\mathrm{hol}}_{\mathrm{KS},0,n}
= d - \sum_{i < j}
\Omega_{\PV}^{(i,j)}\, d\log(z_i - z_j),
\end{equation}
where $\Omega_{\PV}^{(i,j)}$ is the Schouten--Nijenhuis Casimir
acting on the $i$-th and $j$-th insertions.  This is the
Gauss--Manin connection on the moduli of complex structures of $Y$.
Its flatness is Griffiths transversality.
\end{theorem}

\begin{proof}
The shadow connection at genus~$0$ is determined by the binary
coupling $\Theta_{\mathrm{KS}}^{\le 2}$, which is the
Schouten--Nijenhuis bracket.  The formula
$\nabla_{0,n} = d - \sum_{i<j} r^{(ij)}(z_i - z_j)\,d\log(z_i - z_j)$
specializes to the stated expression with $r_{\mathrm{KS}}(z)
= \Omega_{\PV}^{\sn}/z$.

The identification with the Gauss--Manin connection:
the periods of $Y$ are the flat sections of $\nabla_{0,n}$, and the
connection matrix in the period basis is the Gauss--Manin connection
matrix.  Flatness of $\nabla_{0,n}$ is the CYBE for $r_{\mathrm{KS}}$
(Theorem~\ref*{V1-thm:thqg-g5-yangian}), which for the Schouten--Nijenhuis
Casimir reduces to the Jacobi identity of the bracket.  In the
classical language, this is Griffiths transversality.
\end{proof}


\subsection{Genus-$g$ free energies}
% label removed: subsec:thqg-ks-free-energies

\begin{proposition}[KS genus-$g$ free energy; \ClaimStatusProvedHere]
% label removed: prop:thqg-ks-free-energy
For a compact Calabi--Yau threefold $Y$, the genus-$g$ free energy
of the Kodaira--Spencer holographic system is
\begin{equation}% label removed: eq:thqg-ks-free-energy
F_g^{\mathrm{KS}}(Y)
= \kappa_{\mathrm{KS}}(Y) \cdot \frac{2^{2g-1}-1}{2^{2g-1}}
\cdot\frac{|B_{2g}|}{(2g)!},
\qquad g \ge 1,
\end{equation}
where $B_{2g}$ are the Bernoulli numbers and
$\kappa_{\mathrm{KS}}(Y) = h^{1,1} - h^{2,1}$.
At genus $1$, $B_2 = 1/6$ gives
\begin{equation}% label removed: eq:thqg-ks-f1
F_1^{\mathrm{KS}}(Y) = \frac{\kappa_{\mathrm{KS}}}{24}.
\end{equation}
\end{proposition}

\begin{proof}
Since the KS shadow obstruction tower terminates at arity $2$ (Gaussian), all
genus-$g$ corrections are determined by the modular characteristic
$\kappa$ alone.  The genus-$g$ free energy is the $g$-th coefficient
of $\log Z(\tau)$, where $Z(\tau) = \Det(1 - K_g)$ is the Fredholm
determinant.  For a Gaussian theory, $F_g = \kappa \cdot \tfrac{2^{2g-1}-1}{2^{2g-1}}\cdot\tfrac{|B_{2g}|}{(2g)!}$
by the standard computation of the determinant of the Laplacian on
$\Sigma_g$ weighted by $\kappa$.
\end{proof}


\subsection{Mirror symmetry and Koszul involution}
% label removed: subsec:thqg-ks-mirror

\begin{corollary}[Mirror symmetry = $S$-duality; \ClaimStatusProvedHere]
% label removed: cor:thqg-ks-mirror
\index{mirror symmetry!as Koszul S-duality@as Koszul $S$-duality}
If $Y$ and $\check{Y}$ are mirror Calabi--Yau threefolds with
$h^{1,1}(Y) = h^{2,1}(\check{Y})$ and
$h^{2,1}(Y) = h^{1,1}(\check{Y})$, then the gravitational
$S$-duality (G4) specializes to mirror symmetry:
\begin{equation}% label removed: eq:thqg-ks-mirror
\kappa_{\mathrm{KS}}(Y) + \kappa_{\mathrm{KS}}(\check{Y}) = 0,
\qquad
\mathcal{H}(\mathrm{KS}_Y)^! \simeq \mathcal{H}(\mathrm{KS}_{\check{Y}}).
\end{equation}
The Koszul involution on the KS boundary algebra exchanges
complex-structure moduli ($h^{2,1}$) and K\"ahler moduli ($h^{1,1}$).
\end{corollary}

\begin{proof}
$\kappa_{\mathrm{KS}}(Y) = h^{1,1}(Y) - h^{2,1}(Y)$ and
$\kappa_{\mathrm{KS}}(\check{Y}) = h^{1,1}(\check{Y}) - h^{2,1}(\check{Y})
= h^{2,1}(Y) - h^{1,1}(Y) = -\kappa_{\mathrm{KS}}(Y)$.
The Koszul involution $\cA \mapsto \cA^!$ exchanges $\Lambda^p T_Y$
and $\Omega^{3-p}_Y$, which under mirror symmetry exchanges
complex-structure and K\"ahler deformations.
\end{proof}


\subsection{The quintic threefold}
% label removed: subsec:thqg-ks-quintic

\begin{computation}[The quintic; \ClaimStatusProvedHere]
% label removed: comp:thqg-ks-quintic
\index{quintic threefold!holographic datum}
For the quintic threefold $Y \subset \P^4$ defined by a generic
degree-$5$ polynomial: $h^{1,1} = 1$, $h^{2,1} = 101$,
$\chi(Y) = -200$.
\begin{center}
\renewcommand{\arraystretch}{1.2}
\small
\begin{tabular}{ll}
\textbf{Invariant} & \textbf{Value} \\
\hline
$\kappa_{\mathrm{KS}}$ & $1 - 101 = -100$ \\
$F_1$ & $\kappa/24 = -100/24 = -25/6$ \\
$F_2$ & $7\kappa/5760 = -700/5760 = -35/288$ \\
Shadow depth & $2$ (Gaussian) \\
$\cA^!$ & B-model on mirror quintic $\check{Y}$
  ($h^{1,1} = 101$, $h^{2,1} = 1$) \\
$\mathcal{C}$ & $D^b\mathrm{Coh}(Y)$ \\
CYBE & Gauss--Manin flatness for quintic
\end{tabular}
\end{center}
The Picard--Fuchs equation for the periods of the mirror quintic
$\check{Y}$ is the genus-$0$ shadow connection
$\nabla^{\mathrm{hol}}_{\mathrm{KS},0,n}$ specialized to $n = 1$.
\end{computation}


\subsection{The complete landscape of KS examples}
% label removed: subsec:thqg-ks-landscape

\begin{center}
\renewcommand{\arraystretch}{1.25}
\small
\begin{tabular}{llccl}
\textbf{CY$_3$} & $(h^{1,1}, h^{2,1})$
  & $\kappa_{\mathrm{KS}}$ & $F_1$ & \textbf{Line category} \\
\hline
Quintic & $(1, 101)$ & $-100$ & $25/3$ & $D^b\mathrm{Coh}(Q)$ \\
Mirror quintic & $(101, 1)$ & $+100$ & $-25/3$ & $D^b\mathrm{Coh}(\check{Q})$ \\
$K3 \times T^2$ & $(21, 21)$ & $0$ & $0$ & $D^b\mathrm{Coh}(K3 \times T^2)$ \\
Rigid CY ($h^{2,1} = 0$) & $(h^{1,1}, 0)$ & $h^{1,1}$ & $-h^{1,1}/12$ & $D^b\mathrm{Coh}(Y)$ \\
Self-mirror & $(h, h)$ & $0$ & $0$ & $D^b\mathrm{Coh}(Y)$
\end{tabular}
\end{center}

\begin{remark}[Self-mirror CY and vanishing modular characteristic]
% label removed: rem:thqg-ks-self-mirror
A self-mirror Calabi--Yau with $h^{1,1} = h^{2,1}$ has
$\kappa_{\mathrm{KS}} = 0$.  This is the KS analogue of the Virasoro
self-dual point $c = 13$ (where $\kappa(\mathrm{Vir}_{13}) = 13/2$).
The holographic datum $\mathcal{H}(\mathrm{KS}_Y)$ is self-conjugate
under $S$-duality.  The genus-$g$ free energies vanish at all genera
by~\eqref{eq:thqg-ks-free-energy}: the gravitational theory is
topologically trivial at the one-loop and higher-loop level.
\end{remark}


% ======================================================================
%
%  SECTION 5: OPEN PROBLEMS AND THE FRONTIER
%
% ======================================================================

\section{Open problems and the frontier}
% label removed: sec:thqg-frontier
\index{frontier problems|textbf}

The holographic modular Koszul datum is a proved construction for
every gravitational input.  The ten gravitational theorems are
proved projections.  The frontier lies in four directions:
modular completion, deformation-quantization bridges, non-principal
W-algebra duality, and the modular factorization adjunction.

\subsection{The resonance completion theorem}
% label removed: subsec:thqg-modular-completion

\begin{theorem}[Resonance completion; \ClaimStatusProvedElsewhere]
% label removed: conj:thqg-platonic-completion
\index{resonance completion theorem}
Every positive-energy chiral algebra $\cA$ has finite resonance rank:
\begin{equation}% label removed: eq:thqg-platonic-completion
\rho(\cA) < \infty.
\end{equation}
The resonance rank $\rho(\cA)$
(Volume~I, Definition~\ref{def:resonance-rank}) is the dimension of
the obstruction space to completion of the bar-cobar tower.  Finiteness
means that the completed bar-cobar
$\hat\Omega(\hat B(\cA))$ converges after finitely many resonance
corrections.
\end{theorem}

\begin{remark}[Proof summary]
% label removed: rem:thqg-completion-evidence
The theorem is proved in Volume~I
(Theorem~\ref{thm:platonic-completion}):
weight-compatible SDR from the positive-energy axiom yields
finite-dimensional weight spaces, weight-by-weight SDR,
and $R_\cA \subseteq V_0$ with $\dim V_0 < \infty$.
The proof covers:
\begin{itemize}
\item All positive-weight towers (MC4$^+$): $\rho = 0$.
  This includes $\mathcal{W}_{1+\infty}$, affine Yangians,
  and RTT algebras (Volume~I, Theorem~\ref{thm:stabilized-completion-positive}).
\item All quadratic algebras: $\rho = 0$ (PBW).
\item The Virasoro algebra: $\rho \le 1$ (the depth-zero resonance
  shadow $\mathrm{Vir}_{26-c}$ provides one resonance correction).
\end{itemize}
The conjecture is open for general $\mathcal{W}_N$ algebras at
non-quadratic levels and for the $M2$ boundary algebra.
\end{remark}


\subsection{MC4 splitting: positive vs.\ resonant towers}
% label removed: subsec:thqg-mc4-splitting

\begin{definition}[MC4$^+$ and MC4$^0$]
% label removed: def:thqg-mc4-splitting
\index{MC4 splitting|textbf}
\begin{enumerate}[label=\textup{(\roman*)}]
\item \textbf{MC4$^+$} (positive towers): algebras with strictly
  positive weight filtration on all generators except the vacuum.
  The completed bar-cobar converges by weight stabilization.
  \emph{Solved unconditionally}
  (Volume~I, Theorem~\ref{thm:stabilized-completion-positive}).
  Examples: $\mathcal{W}_{1+\infty}$, affine Yangians, RTT algebras.
\item \textbf{MC4$^0$} (resonant towers): algebras with weight-$0$
  generators that interact nontrivially with the bar differential.
  The completed bar-cobar requires finitely many resonance corrections
  (Volume~I, Theorem~\ref{thm:resonance-filtered-bar-cobar}).
  \emph{Reduced to finite resonance problem.}
  Examples: Virasoro, $\mathcal{W}_N$ ($N \ge 3$).
\end{enumerate}
The MC4 splitting is a complete dichotomy: every positive-energy
chiral algebra falls into exactly one class.
\end{definition}

\begin{theorem}[MC4 resolution; \ClaimStatusProvedHere]
% label removed: thm:thqg-mc4-resolution
The strong completion-tower theorem
(Volume~I, Theorem~\ref{thm:completed-bar-cobar-strong}) resolves
MC4: the completed bar-cobar passes to inverse limits via the
strong filtration axiom
\begin{equation}% label removed: eq:thqg-strong-filtration
\mu_r(F^{i_1}, \ldots, F^{i_r}) \subset F^{i_1 + \cdots + i_r},
\end{equation}
and the completion closure $\operatorname{CompCl}(\mathcal{F}_{\mathrm{ft}})$
carries a quasi-inverse bar-cobar homotopy equivalence.
\end{theorem}


\subsection{The deformation-quantization bridge}
% label removed: subsec:thqg-dq-bridge

\begin{conjecture}[Deformation-quantization bridge; \ClaimStatusConjectured]
% label removed: conj:thqg-dq-bridge
\index{deformation-quantization!holographic bridge}
There exists a functor
\begin{equation}% label removed: eq:thqg-dq-bridge
Q_{\mathrm{HT}}: \{\text{classical coisson data}\}
\;\longrightarrow\;
\{\text{quantum Koszul data}\}
\end{equation}
that quantizes the PVA descent (Section~\ref{sec:PVA_descent}) to the
full bar-cobar package.  The functor takes a Poisson vertex algebra
$P$ (the classical limit of a chirally Koszul $\cA$) to the holographic
modular Koszul datum $\mathcal{H}(\cA)$, with the deformation parameter
$\hbar$ tracking the quantization.
\end{conjecture}

\begin{remark}[Evidence from Khan--Zeng]
% label removed: rem:thqg-dq-kz
The Khan--Zeng construction of 3d PVA sigma models gives a classical
version of the holographic triangle: the classical bulk is the Poisson
center of the boundary PVA, and classical line operators are Poisson
modules.  The quantization bridge $Q_{\mathrm{HT}}$ should lift this
to the full quantum package.

The obstruction to constructing $Q_{\mathrm{HT}}$ is the formality
question: whether the $\Ainf$ deformation of the PVA can be lifted
to the full bar-cobar level without higher obstructions.  For
Gaussian algebras (class G), formality is automatic.  For Lie-type
(class L), formality is equivalent to the Etingof--Kazhdan
quantization theorem.
\end{remark}


\subsection{The $3$d PVA sigma model}
% label removed: subsec:thqg-pva-sigma

\begin{definition}[$3$d PVA sigma model]
% label removed: def:thqg-pva-sigma
\index{PVA sigma model!3d}
A \emph{$3$d PVA sigma model} is a $3$-dimensional
holomorphic-topological field theory whose classical phase space
is a Poisson vertex algebra $P$.  The fields are:
\begin{enumerate}[label=\textup{(\roman*)}]
\item A map $\phi: \C \times \R \to \Spec(P_0)$ (the classical
  target).
\item A $P_0$-valued $1$-form $A$ along $\R$ (the gauge field).
\item Ghost fields for the BRST symmetry.
\end{enumerate}
The action is
\begin{equation}% label removed: eq:thqg-pva-action
S = \int_{\C \times \R}
\langle A, \bar\partial\phi \rangle
+ \tfrac{1}{2}\langle A, \{A, A\}_\lambda \rangle
+ \text{higher terms},
\end{equation}
where $\{-,-\}_\lambda$ is the $\lambda$-bracket of $P$.
\end{definition}

\begin{proposition}[Classical holographic datum; \ClaimStatusProvedHere]
% label removed: prop:thqg-pva-classical-datum
The classical limit of the holographic modular Koszul datum is
the PVA holographic datum:
\begin{equation}% label removed: eq:thqg-pva-classical-datum
\mathcal{H}^{\mathrm{cl}}(P)
= \bigl(
P,\; P^!,\; P^!\text{-}\mathsf{mod},\;
r^{\mathrm{cl}}(z),\; \Theta^{\mathrm{cl}},\;
\nabla^{\mathrm{cl}}
\bigr),
\end{equation}
where $P^! = (H^\bullet(\barB(P)))^\vee$ is the classical Koszul
dual PVA, $r^{\mathrm{cl}}(z)$ is the classical $r$-matrix, and
$\nabla^{\mathrm{cl}} = d - r^{\mathrm{cl}} \cdot d\log(z)$ is
the classical KZ connection.
\end{proposition}

\begin{proof}
Setting $\hbar = 0$ in the holographic modular Koszul datum:
the chiral algebra $\cA$ degenerates to the PVA $P = \gr(\cA)$
(PVA descent, Chapter~\ref{chap:pva-descent}).  The bar complex
$\barB(\cA)$ degenerates to $\barB(P)$.  The MC element
$\Theta_\cA$ degenerates to $\Theta^{\mathrm{cl}} = \Theta_P$,
and all projections degenerate accordingly.
\end{proof}


\subsection{Non-principal $\mathcal{W}$-algebra duality}
% label removed: subsec:thqg-non-principal-w

\begin{definition}[Hook-type $\mathcal{W}$-algebras in type A]
% label removed: def:thqg-hook-type
\index{W-algebra@$\mathcal{W}$-algebra!hook-type}
For $\mathfrak{g} = \mathfrak{sl}_N$ and a hook-type nilpotent
$f = f_{(N-r, 1^r)}$ (corresponding to a hook partition), the
$\mathcal{W}$-algebra $\mathcal{W}_k(\mathfrak{sl}_N, f_{(N-r, 1^r)})$
is obtained by quantum Drinfeld--Sokolov reduction.  The generators
are:
\begin{enumerate}[label=\textup{(\roman*)}]
\item One generator of weight $N - r$ (the ``long row'').
\item $r$ generators of weight $1$ (the ``column'').
\item Additional generators from the BRST cohomology.
\end{enumerate}
\end{definition}

\begin{theorem}[Hook-type duality in type A; \ClaimStatusProvedElsewhere]
% label removed: thm:thqg-hook-duality
\index{W-algebra@$\mathcal{W}$-algebra!hook-type duality}
For hook-type nilpotents in type $A$, the Koszul duality of
$\mathcal{W}$-algebras is:
\begin{equation}% label removed: eq:thqg-hook-duality
\mathcal{W}_k(\mathfrak{sl}_N, f_{(N-r,1^r)})^!
\;\simeq\;
\mathcal{W}_{k'}\!(\mathfrak{sl}_N, f_{(r+1, 1^{N-r-1})}),
\end{equation}
where $k' = k'(k, N, r)$ is the dual level determined by the
Feigin--Frenkel-type involution.  The duality exchanges the long-row
generator with the column generators.
\end{theorem}

\begin{proof}[Proof sketch]
The hook-type case is proved by Fehily and by
Creutzig--Linshaw--Nakatsuka--Sato (2023--2025).  The proof uses
the free-field realization of the hook-type $\mathcal{W}$-algebra and
explicit bar-complex computations.  The key step is that the BRST
cohomology of the hook-type reduction commutes with the bar
construction: $\barB(H^0_{\mathrm{DS}}(V_k)) \simeq
H^0_{\mathrm{DS}}(\barB(V_k))$.
\end{proof}


\begin{conjecture}[Transport-to-transpose; \ClaimStatusConjectured]
% label removed: conj:thqg-transport-transpose
\index{transport-to-transpose conjecture}
For all nilpotent orbits in type $A$, the Koszul duality exchanges
the nilpotent $f_\lambda$ and its transpose $f_{\lambda^t}$:
\begin{equation}% label removed: eq:thqg-transport-transpose
\mathcal{W}_k(\mathfrak{sl}_N, f_\lambda)^!
\;\simeq\;
\mathcal{W}_{k^\vee_\lambda}(\mathfrak{sl}_N, f_{\lambda^t}),
\end{equation}
where $\lambda^t$ is the transpose partition and $k^\vee_\lambda$
is the dual level.
\end{conjecture}

\begin{remark}[Status and evidence]
% label removed: rem:thqg-transport-evidence
The conjecture is proved for:
\begin{itemize}
\item Principal nilpotent (all $\mathfrak{g}$): the classical
  Feigin--Frenkel duality.
\item Hook-type in type $A$: Theorem~\ref{thm:thqg-hook-duality}.
\item $\mathfrak{sl}_3$, subregular and minimal nilpotents:
  direct computation.
\end{itemize}
The conjecture is open for general partitions $\lambda$ in type $A$,
and entirely open for types $B$, $C$, $D$, $E$, $F$, $G$.  The
obstruction is not defining the $\mathcal{W}$-algebra (which exists
by Kac--Roan--Wakimoto) but proving that bar-cobar commutes with
non-principal DS reduction.
\end{remark}


\subsection{DS reduction as a functor on holographic data}
% label removed: subsec:thqg-ds-functor

\begin{conjecture}[DS as holographic functor; \ClaimStatusConjectured]
% label removed: conj:thqg-ds-functor
\index{Drinfeld--Sokolov reduction!as holographic functor}
For principal DS reduction, the Drinfeld--Sokolov functor
$\mathrm{DS}: V_k(\mathfrak{g})\otimes F_{\mathrm{gh}}
\mapsto \mathcal{W}_k(\mathfrak{g})$
lifts to a functor on holographic modular Koszul data:
\begin{equation}% label removed: eq:thqg-ds-functor
\mathrm{DS}: \mathcal{H}(V_k(\mathfrak{g})\otimes F_{\mathrm{gh}})
\;\longrightarrow\;
\mathcal{H}(\mathcal{W}_k(\mathfrak{g})),
\end{equation}
commuting with all six components.  In particular:
\begin{enumerate}[label=\textup{(\roman*)}]
\item $\mathrm{DS}(\Theta_{V_k\otimes F_{\mathrm{gh}}}) = \Theta_{\mathcal{W}_k}$.
\item $\mathrm{DS}(r_{V_k\otimes F_{\mathrm{gh}}}(z)) = r_{\mathcal{W}_k}(z)$.
\item $\mathrm{DS}(\nabla^{\mathrm{hol}}_{V_k\otimes F_{\mathrm{gh}}})
  = \nabla^{\mathrm{hol}}_{\mathcal{W}_k}$.
\end{enumerate}
\end{conjecture}

\begin{remark}
This conjecture is a precise formulation of the expectation that
``DS reduction commutes with everything.''  The principal case is
established for components (i)--(iii) at the genus-$0$ tree level.
The extension to all genera and to non-principal reduction is the
frontier.
\end{remark}

\begin{remark}[Geometric localization content of the DS holographic functor]
\label{rem:ds-holographic-geometric-content}
\ClaimStatusConjectured
\index{geometric localization!holographic}
\index{Slodowy slice!holographic functor}
The conjectural DS holographic functor has a geometric substrate in
the Slodowy-slice localization principle developed in
Vol~I (Remark~\ref*{rem:ds-geometric-localization-principle}).
Drinfeld--Sokolov reduction on the algebraic side
(the passage from $V_k(\mathfrak{g})$ to
$\mathcal{W}_k(\mathfrak{g}, f)$) corresponds, on the spectral
geometry side, to Hamiltonian reduction of the Hitchin system along
the coadjoint orbit $\mathbb{O}_f \subset \mathfrak{g}^*$, with the
resulting reduced phase space localized on the Slodowy slice
$\mathcal{S}_f = f + \mathfrak{g}^e$.  In the bulk, this means that
DS reduction on the boundary chiral algebra is the algebraic shadow
of Hamiltonian reduction on the spectral/Hitchin geometry governing
the topological $B$-twist.

The master commutative square (Vol~I, \S21.5)
makes this correspondence precise at the level
of bar complexes and Koszul duality: bar-cobar duality intertwines
DS reduction with Slodowy-slice localization on the dual side.  The
conjecture above asserts that this commutative square lifts to the
full holographic level, so that the six-component holographic modular
Koszul datum $\mathcal{H}(V_k(\mathfrak{g}))$ maps functorially to
$\mathcal{H}(\mathcal{W}_k(\mathfrak{g}, f))$ with each component
transforming by the appropriate geometric reduction.

The algebraic and geometric infrastructure
(the Slodowy-slice localization principle, the master commutative
square, and the intertwining of DS with bar-cobar at genus~$0$)
is proved in Vol~I.  The conjectural content is the extension to all
genera and all six holographic components simultaneously, which
requires the full non-principal DS--bar-cobar commutation that
remains the frontier (see the preceding conjecture and remark).
\end{remark}


\subsection{The modular factorization adjunction}
% label removed: subsec:thqg-platonic-adjunction

\begin{definition}[Cyclically admissible Lie conformal algebra]
% label removed: def:thqg-cyclically-admissible
\index{cyclically admissible!Lie conformal algebra}
A \emph{cyclically admissible} Lie conformal algebra is a Lie conformal
algebra $L$ equipped with:
\begin{enumerate}[label=\textup{(\roman*)}]
\item A conformal grading $L = \bigoplus_{n \ge 0} L_n$ with
  $\dim L_n < \infty$.
\item An exhaustive increasing filtration $F^0 L \subset F^1 L
  \subset \cdots$ compatible with the conformal grading.
\item A bound on poles: the $\lambda$-bracket $\{a_\lambda b\}$ has
  at most $\Delta(a) + \Delta(b)$ terms in $\lambda$.
\item An invariant nondegenerate symmetric pairing
  $\langle\cdot,\cdot\rangle: L \otimes L \to \Bbbk$.
\end{enumerate}
\end{definition}

\begin{construction}[Modular Koszul datum]
% label removed: constr:thqg-platonic-package
\index{platonic package|textbf}
Given a cyclically admissible Lie conformal algebra $L$, the
\emph{modular Koszul datum} is the six-fold datum
\begin{equation}% label removed: eq:thqg-platonic-package
\Pi_X(L) = \bigl(
\mathcal{F}_X(L),\;
\bar{B}_X(L),\;
\Theta_L,\;
\mathcal{L}_L,\;
(V^{\mathrm{br}}_L, T^{\mathrm{br}}_L),\;
R_4^{\mathrm{mod}}(L)
\bigr),
\end{equation}
where:
\begin{enumerate}[label=\textup{(\alph*)}]
\item $\mathcal{F}_X(L)$ is the factorization algebra on $X$
  associated to $L$ (the chiral envelope of $L$).
\item $\bar{B}_X(L)$ is the bar coalgebra of $\mathcal{F}_X(L)$.
\item $\Theta_L$ is the bar-intrinsic MC element.
\item $\mathcal{L}_L$ is the line-operator category.
\item $(V^{\mathrm{br}}_L, T^{\mathrm{br}}_L)$ are the branch-line
  data (the local system on the configuration space).
\item $R_4^{\mathrm{mod}}(L)$ is the modular quartic resonance class.
\end{enumerate}
\end{construction}

\begin{theorem}[Modular factorization adjunction; \ClaimStatusProvedElsewhere]
% label removed: thm:thqg-platonic-adjunction
\index{platonic adjunction}
\index{modular factorization envelope}
There exists a universal modular factorization envelope
\begin{equation}% label removed: eq:thqg-modular-envelope
U^{\mathrm{mod}}_X: \{\text{cyclically admissible Lie conformal algebras}\}
\;\longrightarrow\;
\{\text{modular chiral algebras on } X\}
\end{equation}
that is left adjoint to the modular primitive-current functor
$\mathrm{Prim}^{\mathrm{mod}}$:
\begin{equation}% label removed: eq:thqg-platonic-adjunction
U^{\mathrm{mod}}_X \;\dashv\; \mathrm{Prim}^{\mathrm{mod}}.
\end{equation}
The envelope $U^{\mathrm{mod}}_X(L)$ carries a canonical MC element
$\Theta^{\mathrm{env}}_L$ that restricts to the bar-intrinsic
$\Theta_L$ on the genus-$0$ sector.
This is Vol~I Theorem~\ref{thm:platonic-adjunction}
(adjunction) and Corollary~\ref{cor:envelope-universal-mc}
(universal MC class).
\end{theorem}

\begin{remark}[Nishinaka and Vicedo]
% label removed: rem:thqg-nishinaka-vicedo
\index{factorization envelope!Nishinaka--Vicedo}
The genus-$0$ factorization envelope is constructed by Nishinaka
(2025/26) for Lie conformal algebras and by Vicedo (2025) in the
non-chiral setting.  The modular extension to all genera is
proved in Vol~I via the bar-intrinsic construction
(Theorem~\ref*{V1-thm:mc2-bar-intrinsic}) and the sewing machinery
of MC5 (Theorem~\ref{thm:general-hs-sewing}).
\end{remark}


\subsection{Gauge triviality and the quartic canonical class}
% label removed: subsec:thqg-gauge-triviality

\begin{theorem}[Cubic gauge triviality; \ClaimStatusProvedHere]
\label{thm:cubic-gauge-triviality}
\index{cubic gauge triviality|textbf}
If $H^1(F^3\gAmod / F^4\gAmod, d_2) = 0$, then the cubic MC
term $\Theta_\cA^{(3)}$ is gauge-trivial, and the quartic class
$[\Theta'^{(4)}_\cA] \in H^2(F^4\gAmod / F^5\gAmod, d_2)$ is
canonical (independent of the gauge choice at arity $3$).
\end{theorem}

\begin{proof}[Proof sketch]
The MC equation at arity $3$ reads
$d_2(\Theta^{(3)}) + \frac{1}{2}[\Theta^{(2)}, \Theta^{(2)}]_3 = 0$.
If $H^1(F^3/F^4, d_2) = 0$, then the cocycle
$\frac{1}{2}[\Theta^{(2)}, \Theta^{(2)}]_3$ is a coboundary:
$\Theta^{(3)} = d_2(\xi_3)$ for some $\xi_3$.  The gauge
transformation $\Theta' = e^{-\ad(\xi_3)}\Theta$ removes the
cubic term, and the resulting quartic class $[\Theta'^{(4)}]$ is
independent of the choice of $\xi_3$ because any two choices differ
by a $d_2$-exact term.
\end{proof}


\subsection{Independent sum factorization}
% label removed: subsec:thqg-independent-sum

\begin{proposition}[Independent sum factorization; \ClaimStatusProvedHere]
% label removed: prop:thqg-independent-sum
\index{independent sum factorization}
Let $L = L_1 \oplus L_2$ be a direct sum of cyclically admissible
Lie conformal algebras with vanishing mixed OPE
($\{a_\lambda b\} = 0$ for $a \in L_1$, $b \in L_2$).  Then all
shadows factorize:
\begin{enumerate}[label=\textup{(\roman*)}]
\item $\kappa(L) = \kappa(L_1) + \kappa(L_2)$ (additivity).
\item $T^{\mathrm{br}}(L) = T^{\mathrm{br}}(L_1) \oplus
  T^{\mathrm{br}}(L_2)$ (direct sum of branch-line data).
\item $\Delta_L = \Delta_{L_1} \cdot \Delta_{L_2}$ (multiplicativity
  of discriminants).
\item $R_4^{\mathrm{mod}}(L) = R_4^{\mathrm{mod}}(L_1) +
  R_4^{\mathrm{mod}}(L_2)$ (additivity of quartic classes).
\end{enumerate}
\end{proposition}

\begin{proof}
The vanishing of the mixed OPE implies that the convolution algebra
$\gAmod[L]$ decomposes as a direct product
$\gAmod[L_1] \times \gAmod[L_2]$.  The MC element decomposes as
$\Theta_L = \Theta_{L_1} + \Theta_{L_2}$, and all projections
factorize accordingly.  The additivity of $\kappa$ and $R_4^{\mathrm{mod}}$
follows from linearity of the shadow projection; the multiplicativity
of $\Delta$ follows from the factorization of the genus-$1$ trace
(product of determinants).
\end{proof}


\subsection{The envelope-shadow functor}
% label removed: subsec:thqg-frontier-envelope-shadow

\begin{definition}[Envelope-shadow functor]
% label removed: def:thqg-envelope-shadow
\index{envelope-shadow functor}
For a cyclically admissible Lie conformal algebra $L$ with a choice of
generators $R = \{r_i\}_{i \in I}$, the \emph{envelope-shadow functor}
at arity $\le r$ is
\begin{equation}% label removed: eq:thqg-envelope-shadow
\Theta^{\mathrm{env}}_{\le r}(R):
R \;\longmapsto\;
\Theta_{\mathcal{F}_X(L)}^{\le r},
\end{equation}
the arity-$\le r$ truncation of the MC element of the factorization
envelope $\mathcal{F}_X(L)$, expressed in terms of the generators $R$.
\end{definition}

\begin{proposition}[Finite-jet rigidity; \ClaimStatusProvedHere]
% label removed: prop:thqg-finite-jet-rigidity
For a cyclically admissible Lie conformal algebra $L$, the
envelope-shadow functor satisfies:
\begin{enumerate}[label=\textup{(\roman*)}]
\item \textbf{Polynomial level dependence}: For an algebraic family
  $L_k$ parameterized by level $k$, $\Theta^{\mathrm{env}}_{\le r}(L_k)$
  is polynomial in $k$ of degree $\le r$.
\item \textbf{Gaussian collapse}: If $L$ is abelian (free-field), then
  $\Theta^{\mathrm{env}}_{\le r}(L) = \Theta^{\mathrm{env}}_{\le 2}(L)$
  for all $r \ge 2$.
\end{enumerate}
\end{proposition}

\begin{proof}
(i)~The MC element at arity $r$ is a sum over graphs of genus $\le g$
with $r$ vertices.  Each vertex contributes a factor polynomial in $k$
(the OPE coefficients are polynomial in the level).  The total degree
in $k$ is bounded by $r$.

(ii)~If $L$ is abelian, the OPE has no terms beyond the singular part
(no normal-ordered products), so the bar differential $D$ has only a
linear and quadratic part.  The MC element $\Theta = D - d_0$ has only
the quadratic part: $\Theta = \Theta^{\le 2}$.
\end{proof}


\subsection{MC3 categorical lift: all simple types}
% label removed: subsec:thqg-mc3-frontier

\begin{theorem}[Type-$A$ MC3 reduction beyond the evaluation-generated core; \ClaimStatusProvedElsewhere]
% label removed: thm:thqg-mc3-type-a
\index{MC3!type A resolution}
For $\mathfrak{g} = \mathfrak{sl}_N$, Volume~I proves the DK
comparison on the evaluation-generated core and, beyond that core,
establishes the shifted-prefundamental-generation and pro-Weyl
packets on the type-$A$ shifted/completed side.  The only remaining
post-core step is the compact-completion extension/comparison packet.
\end{theorem}

\begin{theorem}[All-types MC3 on the evaluation-generated core; \ClaimStatusProvedElsewhere]
% label removed: conj:thqg-mc3-arbitrary
\index{MC3!all simple types}
For all simple $\mathfrak{g}$, Volume~I proves the categorical
Clebsch--Gordan part of MC3 via the multiplicity-free
$\ell$-weight property and separately proves the DK comparison on the
evaluation-generated core.  Beyond that core, only type~$A$ currently
has a theorematic reduction to a single compact-completion gap; the
analogous post-core mechanism in other types remains conditional on
additional shifted/completed inputs.
\end{theorem}

\begin{remark}[The Hern\'andez--Jimbo obstruction: resolved for categorical CG]
% label removed: rem:thqg-hernandez-jimbo
The original obstruction to extending MC3 beyond type $A$ was the
absence of a suitable set of prefundamental representations for the
quantum affine algebra $U_q(\hat{\mathfrak{g}})$ in non-simply-laced
types.  This obstruction is resolved by the
multiplicity-free $\ell$-weight approach
(Vol~I, Theorem~\ref*{V1-thm:categorical-cg-all-types}): the
Chari--Moura property (every tensor product $V_{\omega_i} \otimes L^-$
has multiplicity-free $\ell$-weight spaces) replaces the minuscule
hypothesis entirely, and holds for all simple Lie algebras.  What this
resolves is the categorical-CG obstruction and the resulting
evaluation-generated-core comparison, not the full post-core MC3
extension problem.
\end{remark}


\subsection{The holographic dictionary}
% label removed: subsec:thqg-dictionary

\begin{center}
\renewcommand{\arraystretch}{1.3}
\small
\begin{tabular}{p{5cm}p{6.5cm}}
\textbf{Physics}
  & \textbf{Modular Koszul Duality} \\
\hline
Boundary local operators
  & $\cA$ (boundary chiral algebra) \\
Bulk local operators
  & $\Zder(\cA) \simeq \ChirHoch^\bullet(\cA^!)$ \\
Line operators
  & $\mathcal{C} \simeq \cA^!\text{-}\mathsf{mod}$ \\
Gravitational scattering matrix
  & $r(z) = \Res^{\mathrm{coll}}_{0,2}(\Theta_\cA)$ \\
$S$-duality
  & Koszul involution $\cA \mapsto \cA^!$ \\
Gravitational anomaly
  & $\kappa(\cA)$ (modular characteristic) \\
UV finiteness
  & Arity-cutoff lemma + FM compactification \\
Genus-$g$ partition function
  & $\Det(1 - K_g(\cA))$ (Fredholm determinant) \\
Soft graviton theorems
  & Ward identities of $\nabla^{\mathrm{hol}}_{g,n}$ \\
CPT conjugation
  & Verdier involution $\sigma$ \\
BMS symmetry
  & Automorphisms of $\gAmod$ \\
Holographic reconstruction
  & Bar-intrinsic construction of $\Theta_\cA$ \\[2pt]
\hline
\multicolumn{2}{l}{\textbf{Examples:}} \\
\hline
$M2$-brane holography
  & $\cA = A_{M2,\infty}$, $r_{\max} = \infty$ \\
B-model / Kodaira--Spencer
  & $\cA = \PV^{\bullet,\bullet}(Y)$, $r_{\max} = 2$ \\
Affine higher-spin gravity
  & $\cA = V_k(\mathfrak{g})$, $r_{\max} = 3$ \\
Contact gravity ($\beta\gamma$)
  & $\cA = \beta\gamma$, $r_{\max} = 4$ \\
Virasoro gravity
  & $\cA = \mathrm{Vir}_c$, $r_{\max} = \infty$ \\
Critical string
  & $\cA = \mathrm{Vir}_{26}$, $\cA^! = \mathrm{Vir}_0$ (maximally asymmetric dual scalar sector under G4)
\end{tabular}
\end{center}


\subsection{The six-shadow synthesis}
% label removed: subsec:thqg-six-shadow

\begin{definition}[Modular cumulant transform]
% label removed: def:thqg-modular-cumulant
\index{modular cumulant transform}
The \emph{modular cumulant transform} of a gravitational input $\cA$ is
the sextuple
\begin{equation}% label removed: eq:thqg-modular-cumulant
M(\cA) = \bigl(
\hat{T}^c(\cA),\;
D_\cA,\;
\tau_\cA,\;
r_\cA(z),\;
\Theta_\cA,\;
\nabla^{\mathrm{hol}}
\bigr),
\end{equation}
where $\hat{T}^c(\cA)$ is the completed cumulant generating function,
$D_\cA$ is the bar differential, $\tau_\cA$ is the reduced partition
function, $r_\cA(z)$ is the collision residue, $\Theta_\cA$ is the
universal MC element, and $\nabla^{\mathrm{hol}}$ is the shadow
connection.
\end{definition}

\begin{remark}[Five theorems as cumulant relations]
% label removed: rem:thqg-five-theorems-cumulants
The five main theorems of Volume~I (A, B, C, D, H) are the first
five cumulant relations of the modular cumulant transform:
\begin{enumerate}[label=\textup{(\roman*)}]
\item \textbf{Theorem A} (bar-cobar adjunction): $D_\cA$ defines
  the cumulant generating function.
\item \textbf{Theorem B} (bar-cobar inversion): the cumulant
  generating function is invertible on the Koszul locus.
\item \textbf{Theorem C} (complementarity): the cumulant
  generating function splits into bulk and defect components.
\item \textbf{Theorem D} (modular characteristic): the first
  cumulant is $\kappa(\cA)$.
\item \textbf{Theorem H} (Hochschild): the chiral Hochschild
  cohomology is polynomial in the cumulants.
\end{enumerate}
The ten gravitational theorems (G1--G10) extend this to all cumulant
levels.
\end{remark}


\subsection{Grand completion conjecture}
% label removed: subsec:thqg-grand-completion

\begin{conjecture}[Grand completion; \ClaimStatusConjectured]
% label removed: conj:thqg-grand-completion
\index{grand completion conjecture}
For every cyclically admissible Lie conformal algebra $L$ with
associated gravitational input $\cA = \mathcal{F}_X(L)$, the
modular cumulant transform $M(\cA)$ determines and is determined by
the modular Koszul datum $\Pi_X(L)$:
\begin{equation}% label removed: eq:thqg-grand-completion
M(\cA) \;\xleftrightarrow{\;\sim\;}\; \Pi_X(L).
\end{equation}
The equivalence is natural in $L$ and compatible with:
\begin{enumerate}[label=\textup{(\roman*)}]
\item DS reduction: $M(\mathrm{DS}(\cA)) = \mathrm{DS}(M(\cA))$.
\item Independent sum: $M(\cA_1 \oplus \cA_2) = M(\cA_1) \oplus
  M(\cA_2)$.
\item Level deformation: $M(\cA_k)$ is algebraic in $k$.
\item Koszul involution: $M(\cA^!) = \sigma(M(\cA))$.
\end{enumerate}
\end{conjecture}

\begin{remark}
This conjecture is the ultimate target of the holographic programme:
a single functorial construction that unifies the bar-cobar engine
(Volume~I), the Swiss-cheese/HT framework (this volume), and the
factorization-envelope technology into one coherent package.
\end{remark}


\subsection{Notation and grading}
% label removed: subsec:thqg-notation

\begin{enumerate}[label=\textup{(\alph*)}]
\item Grading is \textbf{cohomological} ($|d| = +1$) throughout.
  The bar construction uses desuspension.

\item $\gAmod = \Convstr(\mathcal{C}^!_{\mathrm{ch}},
  \mathcal{P}^{\mathrm{ch}})$ is the modular convolution dg Lie
  algebra (Definition~\ref*{V1-def:modular-convolution-dg-lie}), the
  strict model of $\Convinf(\mathcal{C}^!_{\mathrm{ch}},
  \mathcal{P}^{\mathrm{ch}})$.

\item $B(\cA)$ is the bar coalgebra; $\cA^! = (H^*(B(\cA)))^\vee$
  is the Koszul dual; $\Omega(B(\cA)) \simeq \cA$ is bar-cobar inversion.

\item \emph{HT} abbreviates ``holomorphic-topological.''
  An HT theory is a 3d $\mathcal{N} = 2$ theory in the holomorphic twist.

\item $\mathrm{Vir}_c^! \simeq \mathrm{Vir}_{26-c}$.  Self-duality at
  $c = 13$, \textbf{not} $c = 26$.  The critical string value $c = 26$ is
  the point where $\cA^! \simeq \mathrm{Vir}_0$ on the uncurved scalar
  locus, not a trivial Virasoro algebra.

\item The convolution algebra $\gAmod$ is a strict model of the
  underlying modular $\Linf$ algebra.  MC moduli coincide; the full
  $\Linf$ structure is needed for transfer, formality, and gauge
  equivalence.

\item The shadow obstruction tower is the primary nonlinear object:
  $\kappa \to \Delta \to \mathfrak{C} \to \mathfrak{Q} \to \cdots$
  as projections of the bar-intrinsic $\Theta_\cA$.

\item $\mathrm{Com}^! = \mathrm{Lie}$ (not $\mathrm{coLie}$).
\end{enumerate}


% ======================================================================
%  SECTION FILES: the ten results (expanded sections)
% ======================================================================

% Section file for Chapter: Twisted Holography and Quantum Gravity
% Result (G1): Perturbative Finiteness
%
% Develops the HS-sewing framework, shadow free energies,
% convergence of the gravitational partition function, the
% Heisenberg prototype, finiteness for the standard landscape,
% and implications for 3d quantum gravity.

\section{Perturbative finiteness}
% label removed: sec:thqg-perturbative-finiteness
\index{perturbative finiteness|textbf}
\index{UV finiteness!twisted holography}

% ----------------------------------------------------------------------
% Local macros: provided only if absent from the ambient manuscript.
% ----------------------------------------------------------------------
\providecommand{\HS}{\operatorname{HS}}
\providecommand{\Tr}{\operatorname{Tr}}
\providecommand{\Det}{\operatorname{Det}}
\providecommand{\Fred}{\operatorname{Fred}}
\providecommand{\rank}{\operatorname{rank}}
\providecommand{\Sym}{\operatorname{Sym}}
\providecommand{\id}{\mathrm{id}}
\providecommand{\ad}{\operatorname{ad}}
\providecommand{\wt}{\operatorname{wt}}
\providecommand{\Aut}{\operatorname{Aut}}
\providecommand{\Res}{\operatorname{Res}}
\providecommand{\Defcyc}{\operatorname{Def}_{\mathrm{cyc}}}
\providecommand{\gAmod}{\mathfrak{g}_{\cA}^{\mathrm{mod}}}
\providecommand{\Gmod}{\mathcal{G}^{\mathrm{mod}}}

The perturbative finiteness of twisted quantum gravity is the foundational analytic result that makes the gravitational shadow programme well defined.  In standard perturbative quantum gravity, the genus-$g$ contribution to the path integral diverges for~$g \geq 2$ (nonrenormalizability of the Einstein--Hilbert action in~$d = 4$).  In the twisted setting, where the boundary chiral algebra~$\cA$ controls the bulk through the holographic modular Koszul datum $\mathcal{H}(T) = (\cA, \cA^!, \mathcal{C}, r(z), \Theta_\cA, \nabla^{\mathrm{hol}})$, the situation is radically different: the Maurer--Cartan equation $D_\cA\Theta_\cA + \tfrac{1}{2}[\Theta_\cA,\Theta_\cA] = 0$ is an \emph{exact} identity, not a truncation of an asymptotic series, and the Fulton--MacPherson compactification provides a geometric regulator that renders every genus-$g$ amplitude finite without the need for UV counterterms.

The section develops this finiteness in six stages: (\S\ref{subsec:thqg-I-hs-sewing-framework}) the HS-sewing framework; (\S\ref{subsec:thqg-I-shadow-free-energies}) shadow free energies and the Bernoulli generating function; (\S\ref{subsec:thqg-I-convergence-grav}) convergence of the gravitational partition function; (\S\ref{subsec:thqg-I-heisenberg-prototype}) the Heisenberg partition function as prototype; (\S\ref{subsec:thqg-I-standard-landscape}) finiteness for the standard landscape; (\S\ref{subsec:thqg-I-3d-gravity}) implications for 3d quantum gravity.

% ======================================================================
%  SUBSECTION 1: THE HS-SEWING FRAMEWORK FOR TWISTED GRAVITY
% ======================================================================

\subsection{The HS-sewing framework for twisted gravity}
% label removed: subsec:thqg-I-hs-sewing-framework
\index{HS-sewing!twisted gravity framework}

The passage from finite-dimensional algebraic data (the OPE coefficients of the boundary chiral algebra~$\cA$) to infinite-genus partition functions requires a bridge between the algebraic and analytic worlds.  This bridge is the HS-sewing condition, which guarantees that the sewing amplitudes on pairs of pants compose to give trace-class operators on the sewing Hilbert space.

\subsubsection{The sewing Hilbert space}

\begin{definition}[Sewing Hilbert space; cf.\ {\cite{costello-renormalization}}]
% label removed: def:thqg-I-sewing-hilbert-space
\index{sewing Hilbert space|textbf}
Let $\cA$ be a positive-energy chiral algebra with conformal grading $\cA = \bigoplus_{n \geq 0} H_n$, where $\dim H_n < \infty$ for all~$n$.  For a sewing parameter $q \in (0,1)$, the \emph{sewing Hilbert space} is the $q$-weighted $\ell^2$-completion
\begin{equation}% label removed: eq:thqg-I-sewing-hilbert
\mathcal{H}_q(\cA) \;:=\; \widehat{\bigoplus}_{n \geq 0} q^{n/2}\,H_n \;=\; \Bigl\{ \sum_{n=0}^\infty v_n \;\Big|\; v_n \in H_n,\; \sum_{n=0}^\infty q^n \|v_n\|^2 < \infty \Bigr\},
\end{equation}
where $\|\cdot\|$ is any fixed Hermitian norm on each finite-dimensional $H_n$.  The inner product is $\langle v, w \rangle_q = \sum_{n \geq 0} q^n \langle v_n, w_n \rangle_{H_n}$.
\end{definition}

\begin{remark}[Independence of norms]
% label removed: rem:thqg-I-norm-independence
Since each $H_n$ is finite-dimensional, different choices of Hermitian norm on $H_n$ give equivalent topologies on $\mathcal{H}_q(\cA)$.  The $q$-weighting is the essential ingredient: it ensures that the completion is sensitive to conformal weight, with high-weight states exponentially suppressed.  For $q \to 1^-$, the Hilbert space degenerates; for $q \to 0^+$, only the vacuum sector $H_0$ survives.
\end{remark}

\begin{lemma}[Completeness and separability; \ClaimStatusProvedHere]
% label removed: lem:thqg-I-completeness
\index{sewing Hilbert space!completeness}
$\mathcal{H}_q(\cA)$ is a separable Hilbert space for every $q \in (0,1)$.  A Cauchy sequence $\{v^{(k)}\}_{k \geq 1}$ converges if and only if $\sum_{n \geq 0} q^n \|v^{(k)}_n - v^{(\ell)}_n\|^2 \to 0$ as $k, \ell \to \infty$.
\end{lemma}

\begin{proof}
Each $H_n$ is finite-dimensional, hence complete.  The $q$-weighted $\ell^2$-completion is a weighted direct sum of Hilbert spaces with weights $q^n$; completeness of the direct sum follows from that of each summand together with the summability condition.  Separability follows from taking a countable dense subset of each $H_n$ (finite-dimensional spaces are separable) and forming finite linear combinations.
\end{proof}

\begin{proposition}[Inclusion maps; \ClaimStatusProvedHere]
% label removed: prop:thqg-I-inclusion-maps
\index{sewing Hilbert space!inclusion}
For $0 < q_1 < q_2 < 1$, the identity map on $\cA$ extends to a bounded inclusion $\iota_{q_1,q_2} \colon \mathcal{H}_{q_2}(\cA) \hookrightarrow \mathcal{H}_{q_1}(\cA)$ with operator norm $\|\iota_{q_1,q_2}\| = 1$.  In particular, for every $v \in \mathcal{H}_{q_2}(\cA)$:
\begin{equation}% label removed: eq:thqg-I-inclusion
\|v\|_{q_1}^2 \;=\; \sum_{n \geq 0} q_1^n \|v_n\|^2 \;\leq\; \sum_{n \geq 0} q_2^n \|v_n\|^2 \;=\; \|v\|_{q_2}^2.
\end{equation}
\end{proposition}

\begin{proof}
Since $q_1 < q_2$, we have $q_1^n \leq q_2^n$ for all $n \geq 0$, giving the inequality termwise.
\end{proof}

\subsubsection{The HS-sewing condition}

\begin{definition}[HS-sewing condition]
% label removed: def:thqg-I-hs-sewing
\index{HS-sewing!definition|textbf}
Let $\cA$ be a positive-energy chiral algebra with OPE structure constants $C^{c,k}_{a,i;\,b,j}$, where $a,b,c$ are conformal weights and $i,j,k$ index bases of $H_a$, $H_b$, $H_c$ respectively.  The \emph{pair-of-pants multiplication operator} at weight $(a,b) \to c$ is the linear map $m^c_{a,b} \colon H_a \otimes H_b \to H_c$ with matrix entries $(m^c_{a,b})_{ij}^k = C^{c,k}_{a,i;\,b,j}$.  The algebra~$\cA$ satisfies the \emph{HS-sewing condition} at parameter $q \in (0,1)$ if
\begin{equation}% label removed: eq:thqg-I-hs-sewing-condition
\boxed{ \sum_{a,b,c \geq 0} q^{a+b+c}\, \|m^c_{a,b}\|_{\HS}^2 \;<\; \infty,}
\end{equation}
where $\|m^c_{a,b}\|_{\HS}^2 = \sum_{i,j,k} |C^{c,k}_{a,i;\,b,j}|^2$ is the Hilbert--Schmidt norm of the multiplication operator.
\end{definition}

\begin{remark}[Physical content]
% label removed: rem:thqg-I-hs-physical
The HS-sewing condition controls the convergence of the sewing construction: a pair-of-pants decomposition of a genus-$g$ surface with $n$ punctures involves $2g - 2 + n$ pairs of pants, and each sewing circle contributes one composition of pair-of-pants operators.  The HS condition guarantees that each such composition is Hilbert--Schmidt, and that the composition of two Hilbert--Schmidt operators is trace class, the key requirement for the trace that computes the partition function.  The $q$-parameter is the sewing modulus $q = e^{-t}$ where $t > 0$ is the collar length.
\end{remark}

\begin{definition}[Hilbert--Schmidt and trace-class norms]
% label removed: def:thqg-I-hs-trace-norms
\index{Hilbert--Schmidt norm}
\index{trace-class!operator}
For a linear map $T \colon V \to W$ between finite-dimensional Hilbert spaces, the Hilbert--Schmidt norm is $\|T\|_{\HS}^2 = \Tr(T^*T) = \sum_{i,j} |T_{ij}|^2$.  An operator $K$ on a Hilbert space $H$ is \emph{trace class} if $\|K\|_1 := \Tr\sqrt{K^*K} < \infty$.  The key composition inequality is:
\begin{equation}% label removed: eq:thqg-I-hs-composition
\|AB\|_1 \;\leq\; \|A\|_{\HS}\,\|B\|_{\HS}
\end{equation}
for Hilbert--Schmidt operators $A, B$.  Hence the composition of two HS operators is trace class, with trace bounded by the product of HS norms.
\end{definition}

\begin{lemma}[Schatten class hierarchy; \ClaimStatusProvedHere]
% label removed: lem:thqg-I-schatten
\index{Schatten class}
Let $H$ be a separable Hilbert space and $S_p(H)$ the $p$-th Schatten class ($p \geq 1$).  Then:
\begin{enumerate}[label=\textup{(\roman*)}]
\item $S_1(H) \subset S_2(H) \subset S_p(H) \subset \mathcal{B}(H)$ for $1 \leq p \leq \infty$, where $\mathcal{B}(H)$ is bounded operators.
\item $S_1(H)$ is a two-sided ideal in $\mathcal{B}(H)$: if $K \in S_1(H)$ and $T \in \mathcal{B}(H)$, then $KT, TK \in S_1(H)$.
\item If $A \in S_p(H)$ and $B \in S_q(H)$ with $1/p + 1/q = 1/r \leq 1$, then $AB \in S_r(H)$.
\end{enumerate}
The sewing application uses $p = q = 2$, $r = 1$: two Hilbert--Schmidt operators compose to trace class.
\end{lemma}

\begin{proof}
Standard functional analysis; see Reed--Simon~\cite{RS78} or Simon~\cite{Simon05}.  Part~(iii) is the generalized H\"older inequality for Schatten classes.
\end{proof}

\begin{proposition}[Pair-of-pants as HS operator; \ClaimStatusProvedElsewhere]
% label removed: prop:thqg-I-pants-hs
\index{pair of pants!HS operator}
Under HS-sewing~\eqref{eq:thqg-I-hs-sewing-condition}, the $q$-weighted multiplication operator $m_q \colon \mathcal{H}_q(\cA) \,\widehat{\otimes}\, \mathcal{H}_q(\cA) \to \mathcal{H}_q(\cA)$ is Hilbert--Schmidt, with
\begin{equation}% label removed: eq:thqg-I-mq-hs-bound
\|m_q\|_{\HS}^2 \;=\; \sum_{a,b,c \geq 0} q^{a+b+c}\,\|m^c_{a,b}\|_{\HS}^2 \;<\; \infty.
\end{equation}
Consequently, every closed amplitude on a genus-$g$ surface obtained by sewing is trace class.
\end{proposition}

\begin{proof}
Direct from the definition: $\|m_q\|_{\HS}^2$ is precisely the left-hand side of~\eqref{eq:thqg-I-hs-sewing-condition}.  A genus-$g$ surface with $n$ punctures is obtained by sewing $2g - 2 + n$ pairs of pants along $3g - 3 + n$ circles.  Each sewing circle contributes one pair-of-pants composition.  By the Hilbert--Schmidt composition inequality~\eqref{eq:thqg-I-hs-composition}, each composition of two HS operators gives a trace-class operator.  Since $g \geq 1$ requires at least two compositions, and trace-class operators form an ideal, the result follows.
\end{proof}

\begin{remark}[Sewing topology]
% label removed: rem:thqg-I-sewing-topology
\index{sewing!topology}
The pair-of-pants decomposition of a genus-$g$ surface $\Sigma_{g,n}$ into $2g - 2 + n$ trinions and $3g - 3 + n$ sewing circles is determined by a pants graph $\Gamma_P$.  Different choices of $\Gamma_P$ give different factorizations of the sewing operator, but the trace (the partition function) is independent of $\Gamma_P$ by the factorization property.  The sewing topology on the space of pants decompositions is captured by the flip graph, whose vertices are pants decompositions and whose edges are elementary moves (Dehn twists and handle slides).  The partition function is invariant under flips by the pentagon and hexagon identities of the chiral algebra.
\end{remark}

\subsubsection{The polynomial-OPE / subexponential-sector criterion}

The fundamental analytic input is the criterion of Theorem~\ref{thm:general-hs-sewing}: polynomial OPE growth and subexponential sector growth together imply HS-sewing for all $q \in (0,1)$.

\begin{theorem}[HS-sewing from polynomial OPE growth; \ClaimStatusProvedElsewhere{} {\textup{Theorem~\ref{thm:general-hs-sewing}}}]
% label removed: thm:thqg-I-hs-sewing-criterion
\index{HS-sewing!polynomial OPE criterion}
Let\/ $\cA$ be a positive-energy chiral algebra satisfying:
\begin{enumerate}[label=\textup{(\alph*)}]
\item \emph{Subexponential sector growth:} $\log\dim H_n = o(n)$ as $n \to \infty$.
\item \emph{Polynomial OPE growth:} There exist constants $K > 0$ and $N \in \mathbb{Z}_{\geq 0}$ such that $|C^{c,k}_{a,i;\,b,j}| \leq K(a + b + c + 1)^N$ for all conformal weights $a,b,c$ and all basis indices $i,j,k$.
\end{enumerate}
Then $\cA$ satisfies the HS-sewing condition for every $q \in (0,1)$:
\begin{equation}% label removed: eq:thqg-I-hs-criterion-bound
\sum_{a,b,c \geq 0} q^{a+b+c}\,\|m^c_{a,b}\|_{\HS}^2 \;\leq\; K^2 \sum_{n \geq 0} P_3(n)\,n^{2N}\, \bigl(\!\dim H_n\bigr)^3\,q^n \;<\; \infty,
\end{equation}
where $P_3(n) = \binom{n+2}{2}$ counts the number of triples $(a,b,c)$ with $a + b + c = n$.
\end{theorem}

\begin{proof}
The HS norm at weight level $(a,b,c)$ satisfies
\[
\|m^c_{a,b}\|_{\HS}^2 \;=\; \sum_{i=1}^{\dim H_a} \sum_{j=1}^{\dim H_b} \sum_{k=1}^{\dim H_c} |C^{c,k}_{a,i;\,b,j}|^2 \;\leq\; \dim H_a \cdot \dim H_b \cdot \dim H_c \cdot K^2(a{+}b{+}c{+}1)^{2N}.
\]
Summing over triples with $a + b + c = n$:
\[
\sum_{\substack{a,b,c \geq 0 \\ a+b+c = n}} q^n\,\|m^c_{a,b}\|_{\HS}^2 \;\leq\; K^2 q^n n^{2N} \sum_{\substack{a,b,c \geq 0 \\ a+b+c = n}} \dim H_a \cdot \dim H_b \cdot \dim H_c.
\]
The inner sum over triples is bounded by $\bigl(\sum_{j=0}^n \dim H_j\bigr)^3$.  By the subexponential growth hypothesis, for every $\varepsilon > 0$ there exists $C_\varepsilon$ with $\dim H_j \leq C_\varepsilon \,e^{\varepsilon j}$.  Hence $\sum_{j=0}^n \dim H_j \leq C_\varepsilon\,n\,e^{\varepsilon n}$.  The total sum is therefore bounded by $K^2 C_\varepsilon^3 \sum_{n \geq 0} n^{2N+3}\, e^{3\varepsilon n}\,q^n$.  Choosing $\varepsilon < |\log q|/4$, the exponential factor $e^{3\varepsilon n} q^n = e^{(3\varepsilon + \log q)n}$ decays geometrically, ensuring convergence of the power series.
\end{proof}

\begin{remark}[Optimality of the bound]
% label removed: rem:thqg-I-hs-optimality
The polynomial-OPE criterion is not optimal: examples with super-polynomial but sub-exponential OPE growth (such as certain completed W-algebras) may still satisfy HS-sewing.  The point is that the criterion covers the entire standard landscape of finitely strongly generated chiral algebras.  The bound~\eqref{eq:thqg-I-hs-criterion-bound} is loose by factors of order $n^3$ due to the crude estimate $\sum_{a+b+c=n} \dim H_a \dim H_b \dim H_c \leq (\sum_{j \leq n} \dim H_j)^3$; tighter bounds are available for specific families.
\end{remark}

\begin{remark}[The gap between HS-sewing and convergence]
% label removed: rem:thqg-I-gap-hs
\index{HS-sewing!necessary vs sufficient}
The HS-sewing condition is \emph{sufficient} for absolute convergence of the genus-$g$ partition function but not \emph{necessary}.  Weaker conditions (e.g., trace-class sewing without the Hilbert--Schmidt factorization) may suffice for finiteness at specific genera.  The HS condition provides a uniform bound across all genera simultaneously, which is the key advantage: a single verification at the level of OPE data yields finiteness at every genus.
\end{remark}

\subsubsection{Explicit HS norms for the standard landscape}

\begin{computation}[HS norms for Heisenberg]
% label removed: comp:thqg-I-hs-heisenberg
\index{Heisenberg!HS norms}
For the Heisenberg algebra $\mathcal{H}_k$ of rank~$1$ at level $k \in \mathbb{C}^{\times}$:
\begin{itemize}
\item $\dim H_n = p(n)$ (partitions of~$n$);
\item OPE: $a(z)a(w) \sim k/(z-w)^2$, giving $C^c_{a,b} = k\,\delta_{a+b,c+1}$ for the fundamental OPE, and all composite OPE grow polynomially with degree $N = 1$ (the Wick contraction is bilinear).
\item Growth: $\log p(n) \sim \pi\sqrt{2n/3}$ (Hardy--Ramanujan), so $\log \dim H_n = O(\sqrt{n})$, which is subexponential.
\end{itemize}
The HS sum factorizes by Wick's theorem into a one-particle contribution:
\begin{equation}% label removed: eq:thqg-I-hs-heisenberg-norm
\|m_q\|_{\HS}^2 \;=\; \prod_{n=1}^\infty \frac{1}{(1-q^n)^{2|k|^2}} \;=\; \eta(q)^{-2|k|^2},
\end{equation}
which is finite for all $q \in (0,1)$.  The one-particle operator has eigenvalues $q^n$, giving $\|T_q\|_1 = \sum_{n=1}^\infty q^n = q/(1-q)$, which is trace class for all $q < 1$.

The factorization into the Dedekind eta function is the consequence of the boson-fermion correspondence: the multi-particle Fock space partition function is the exponential of the one-particle trace, and the Dedekind eta function $\eta(q) = q^{1/24}\prod_{n=1}^\infty (1-q^n)$ is the generating function for partitions.  Explicitly, the second quantization identity gives
\begin{equation}% label removed: eq:thqg-I-hs-heisenberg-factorization
\sum_{n=0}^\infty p(n)^2 q^n \;=\; \prod_{m=1}^\infty \frac{1}{(1-q^m)^2},
\end{equation}
and the HS norm is $|k|^2$ copies of this product.
\end{computation}

\begin{computation}[HS norms for affine Kac--Moody]
% label removed: comp:thqg-I-hs-affine
\index{affine Kac--Moody!HS norms}
For $\widehat{\fg}_k$ at non-critical level $k \neq -h^\vee$, with $\fg$ simple of dimension~$d$ and dual Coxeter number $h^\vee$:
\begin{itemize}
\item $\dim H_n \sim C\,n^{(d-\rank)/2}\,e^{\pi\sqrt{2dn/3}}$ (Kac--Peterson), so $\log \dim H_n = O(\sqrt{n})$, subexponential.
\item OPE: $J^a(z)J^b(w) \sim k\delta^{ab}/(z-w)^2 + f^{ab}{}_c J^c(w)/(z-w)$, polynomial of degree $N = 1$.
\item The HS norm satisfies
  \begin{equation}% label removed: eq:thqg-I-hs-affine-norm
  \|m_q\|_{\HS}^2 \;\leq\; C_k^2 \prod_{n=1}^\infty \frac{1}{(1-q^n)^{2d}} \;=\; C_k^2\,\eta(q)^{-2d}
  \end{equation}
  where $C_k$ depends on $k$ and the structure constants of $\fg$.
\end{itemize}
The key estimate for $C_k$ is obtained from the OPE normalization.  The structure constants of $\widehat{\fg}_k$ at weight level $n$ grow as $|C^c_{a,b}| \leq \|f\|_\infty \cdot (k+h^\vee)$, where $\|f\|_\infty = \max_{a,b,c} |f^{ab}{}_c|$ is the maximum structure constant.  For $\fg = \mathfrak{sl}_N$, one has $\|f\|_\infty = 1$ in the Chevalley basis, giving $C_k = |k + h^\vee|$.
\end{computation}

\begin{computation}[HS norms for Virasoro]
% label removed: comp:thqg-I-hs-virasoro
\index{Virasoro!HS norms}
For the Virasoro algebra $\mathrm{Vir}_c$ at central charge $c \in \mathbb{C}$:
\begin{itemize}
\item $\dim H_n = p(n)$ (Verma module dimensions);
\item OPE: $T(z)T(w) \sim c/(2(z-w)^4) + 2T(w)/(z-w)^2 + \partial T(w)/(z-w)$, polynomial of degree $N = 2$ (the $(z-w)^{-4}$ pole gives OPE coefficients growing at most quadratically in conformal weight).
\item Growth: $\log p(n) \sim \pi\sqrt{2n/3}$, subexponential.
\item The HS sum converges:
  \begin{equation}% label removed: eq:thqg-I-hs-virasoro-norm
  \|m_q\|_{\HS}^2 \;\leq\; C_c^2 \sum_{n \geq 0} p(n)^3\,n^{4}\,q^n \;<\; \infty \quad\text{for all } q \in (0,1).
  \end{equation}
\end{itemize}
The quartic growth of $n^4$ in the bound reflects the quartic OPE pole of $T(z)T(w)$.  For the Virasoro algebra, the structure constants $C^c_{a,b}$ at weight level $n = a + b - c$ satisfy $|C^c_{a,b}| \leq P_2(n) := (n+1)(n+2)/2$ when $c$ is bounded; the precise bound involves the central charge through $|c|/2$ from the fourth-order pole.
\end{computation}

\begin{computation}[HS norms for $\beta\gamma$ system]
% label removed: comp:thqg-I-hs-betagamma
\index{beta-gamma system@$\beta\gamma$ system!HS norms}
For the $\beta\gamma$ system of conformal weights $(1,0)$:
\begin{itemize}
\item $\dim H_n = p(n)^2$ (two independent partition functions, one for $\beta$ modes and one for $\gamma$ modes).
\item OPE: $\beta(z)\gamma(w) \sim 1/(z-w)$, a simple pole.  All composite OPE grow polynomially with degree $N = 1$.
\item Growth: $\log p(n)^2 = 2\log p(n) \sim 2\pi\sqrt{2n/3}$, subexponential.
\item The HS norm:
  \begin{equation}% label removed: eq:thqg-I-hs-betagamma-norm
  \|m_q\|_{\HS}^2 \;\leq\; C^2\,\eta(q)^{-4} \;<\; \infty \quad\text{for all } q \in (0,1).
  \end{equation}
\end{itemize}
The exponent $-4$ in $\eta(q)^{-4}$ reflects the two bosonic degrees of freedom ($\beta$ and $\gamma$), each contributing $\eta(q)^{-2}$.
\end{computation}

\begin{computation}[HS norms for $\mathcal{W}_N$]
% label removed: comp:thqg-I-hs-wn
\index{W-algebra@$\mathcal{W}$-algebra!HS norms}
For the principal $\mathcal{W}$-algebra $\mathcal{W}_k(\mathfrak{sl}_N)$ at generic level:
\begin{itemize}
\item $\dim H_n$ is the number of $N$-colored partitions of~$n$ with colors of weight $2, 3, \ldots, N$; the growth is $\log \dim H_n = O(n^{(N-1)/N})$, subexponential for all~$N$.
\item OPE growth: the maximal singular order of the $W^{(s)}$ self-OPE is $2s$, so $N_{\max} = 2N$ for the $W^{(N)}$ field.  All OPE coefficients grow at most polynomially with degree $\leq 2N$.
\item Convergence:
  \begin{equation}% label removed: eq:thqg-I-hs-wn-norm
  \|m_q\|_{\HS}^2 \;\leq\; K_N^2(c) \sum_{n \geq 0} (\dim H_n)^3\,n^{4N}\,q^n \;<\; \infty \quad\text{for all } q \in (0,1).
  \end{equation}
\end{itemize}
The $N$-colored partition function has generating series $\prod_{s=2}^N \prod_{m=s}^\infty (1-q^m)^{-1}$, and its growth exponent is controlled by the Meinardus theorem: $\log(\dim H_n) \sim C_N \cdot n^{(N-1)/N}$ with $C_N = N \cdot (\Gamma(1/N) \zeta(1+1/N))^{N/(N+1)} / (N+1)$.  For $N = 2$ (Virasoro), $C_2 = \pi\sqrt{2/3}$; for $N = 3$ ($\mathcal{W}_3$), the growth is $O(n^{2/3})$.
\end{computation}

\begin{computation}[HS norms for lattice VOAs]
% label removed: comp:thqg-I-hs-lattice
\index{lattice VOA!HS norms}
For the lattice vertex algebra $V_\Lambda$ of an even integral lattice $\Lambda$ of rank~$r$:
\begin{itemize}
\item $\dim H_n = |\Lambda^*/\Lambda| \cdot p_r(n)$, where $p_r(n)$ is the $r$-colored partition function.  Growth: $\log \dim H_n = O(\sqrt{n})$, subexponential.
\item OPE: inherited from rank-$r$ Heisenberg plus lattice vertex operators $V_\alpha(z)V_\beta(w) \sim (z-w)^{\langle\alpha,\beta\rangle}:\!V_{\alpha+\beta}(w)\!:$; the OPE growth is polynomial of degree $N = \max(r, \max_{\alpha \in \Lambda}|\langle\alpha,\alpha\rangle|/2)$.
\item The HS norm factors as
  \begin{equation}% label removed: eq:thqg-I-hs-lattice-norm
  \|m_q\|_{\HS}^2 \;\leq\; C_\Lambda^2\,\eta(q)^{-2r} \,\Theta_\Lambda(q)^2 \;<\; \infty \quad\text{for all } q \in (0,1),
  \end{equation}
  where $\Theta_\Lambda(q) = \sum_{\alpha \in \Lambda} q^{\langle\alpha,\alpha\rangle/2}$ is the lattice theta function.
\end{itemize}
The factorization reflects the tensor-product structure $V_\Lambda = \mathcal{H}^{\otimes r} \otimes \mathbb{C}[\Lambda]$: the Heisenberg part contributes $\eta(q)^{-2r}$ and the lattice group ring contributes $\Theta_\Lambda(q)^2$ (one factor from the bra side, one from the ket side of the HS norm).
\end{computation}

\begin{computation}[HS norms for free fermions]
% label removed: comp:thqg-I-hs-fermion
\index{free fermions!HS norms}
For the free-fermion VOA of rank~$n$ (i.e., $n$ chiral fermions $\psi_1, \ldots, \psi_n$ of conformal weight $1/2$):
\begin{itemize}
\item $\dim H_m = \binom{n}{m}$ for $m \leq n$, $\dim H_m = 0$ for $m > n$ when counting weight-$m/2$ states; more precisely, the full Fock space has $\dim H_m$ growing as partitions into distinct parts, with generating function $\prod_{j=1}^n (1 + q^{j-1/2})$.
\item OPE: $\psi_i(z)\psi_j(w) \sim \delta_{ij}/(z-w)$, simple pole.  OPE degree $N = 0$ (the structure constants are $0$ or $\pm 1$).
\item The HS norm:
  \begin{equation}% label removed: eq:thqg-I-hs-fermion-norm
  \|m_q\|_{\HS}^2 \;\leq\; \prod_{j=1}^{n} (1 + q^{j-1/2})^2 \;<\; \infty \quad\text{for all } q \in (0,1).
  \end{equation}
  This bound is finite because it is a finite product of bounded terms.
\end{itemize}
\end{computation}

\begin{table}[ht]
\centering
\caption{HS-sewing growth rates for the standard landscape}
% label removed: tab:thqg-I-hs-growth
\index{HS-sewing!growth rates}
\renewcommand{\arraystretch}{1.4}
{\footnotesize
\begin{tabular}{|l|c|c|c|c|}
\hline
\textbf{Family $\cA$} & $\boldsymbol{\dim H_n}$ & \textbf{OPE growth $N$} & \textbf{Sector growth} & \textbf{HS bound} \\
\hline
$\mathcal{H}_k$ (rank 1) & $p(n)$ & $1$ & $O(\sqrt{n})$ & $\eta(q)^{-2|k|^2}$ \\
\hline
$\widehat{\fg}_k$ & $\sim e^{\pi\sqrt{2dn/3}}$ & $1$ & $O(\sqrt{n})$ & $C_k^2\,\eta(q)^{-2d}$ \\
\hline
$\mathrm{Vir}_c$ & $p(n)$ & $2$ & $O(\sqrt{n})$ & $C_c^2\sum p(n)^3 n^4 q^n$ \\
\hline
$\beta\gamma$ system & $p(n)^2$ & $1$ & $O(\sqrt{n})$ & $\eta(q)^{-4}$ \\
\hline
$\mathcal{W}_N$ & $N$-colored parts & $2N$ & $O(n^{(N-1)/N})$ & $K_N^2(c)\sum (\dim H_n)^3 n^{4N} q^n$ \\
\hline
$V_\Lambda$ (rank $r$) & $|\Lambda^*/\Lambda|\,p_r(n)$ & $r$ & $O(\sqrt{n})$ & $C_\Lambda^2\,\eta(q)^{-2r}\Theta_\Lambda^2$ \\
\hline
Free fermions & $\leq 2^n$ & $0$ & polynomial & $\prod (1+q^{j-1/2})^2$ \\
\hline
\end{tabular}}
\end{table}

\subsubsection{Connection to trace-class operators on sewing Hilbert spaces}

\begin{proposition}[Trace-class amplitudes; \ClaimStatusProvedElsewhere]
% label removed: prop:thqg-I-trace-class-amplitudes
\index{trace class!sewing amplitudes}
Let\/ $\cA$ satisfy the HS-sewing condition at parameter $q$.  Let $\Sigma_{g,n}$ be a genus-$g$ surface with $n$ punctures and a pants decomposition $\mathcal{P}$ with collar parameters $q_1, \ldots, q_{3g-3+n} \in (0,q]$.  Then the sewing amplitude
\begin{equation}% label removed: eq:thqg-I-sewing-amplitude
Z_{\mathcal{P}}(\cA;\, q_1, \ldots, q_{3g-3+n}) \;=\; \Tr_{\mathcal{H}_q(\cA)^{\otimes n}} \Bigl( \prod_{\text{sewing circles}} m_{q_i} \Bigr)
\end{equation}
is well defined and absolutely convergent.  The amplitude depends on the sewing moduli $(q_1, \ldots, q_{3g-3+n})$ but not on the choice of pants decomposition (by the factorization property of the chiral algebra).
\end{proposition}

\begin{proof}
Each pants composition $m_{q_i}$ is Hilbert--Schmidt by Proposition~\ref{prop:thqg-I-pants-hs}.  A genus-$g$ surface requires $2g - 2 + n$ pairs of pants and $3g - 3 + n$ sewing circles.  The sewing amplitude is a trace of a composition of $3g - 3 + n$ Hilbert--Schmidt operators.  Since $3g - 3 + n \geq 2$ for $g \geq 1$ (and $\geq 3$ for $n = 0$), at least one pair of consecutive HS operators composes to give a trace-class operator.  The remaining HS operators are bounded, and trace class times bounded is still trace class (Lemma~\ref{lem:thqg-I-schatten}(ii)).  Hence the trace converges absolutely.  Independence from the pants decomposition follows from the associativity and factorization axioms of the chiral algebra~\cite{BD04, FBZ04}.
\end{proof}

\begin{proposition}[Uniform bound on the trace; \ClaimStatusProvedHere]
% label removed: prop:thqg-I-uniform-trace-bound
\index{trace class!uniform bound}
Under HS-sewing at parameter $q$, the sewing amplitude on a genus-$g$ surface with $n$ punctures satisfies the uniform bound
\begin{equation}% label removed: eq:thqg-I-uniform-trace
|Z_{\mathcal{P}}(\cA;\,q_1,\ldots,q_{3g-3+n})| \;\leq\; \|m_q\|_{\HS}^{3g-3+n}
\end{equation}
for all $q_1, \ldots, q_{3g-3+n} \in (0,q]$.  In particular, the bound is uniform in the sewing moduli.
\end{proposition}

\begin{proof}
Each factor $m_{q_i}$ satisfies $\|m_{q_i}\|_{\HS} \leq \|m_q\|_{\HS}$ (since $q_i \leq q$ and the HS norm is monotone increasing in $q$).  The trace of a composition of operators satisfies $|\Tr(A_1 \cdots A_k)| \leq \|A_1\|_{\HS}\cdots\|A_k\|_{\HS}$ when $k \geq 2$ (by iterated application of the HS composition inequality and the trace bound $|\Tr(K)| \leq \|K\|_1$).
\end{proof}

\begin{corollary}[Genus-$g$ partition function]
% label removed: cor:thqg-I-genus-g-partition
\index{partition function!genus $g$}
For $\cA$ satisfying HS-sewing, the genus-$g$ partition function
\begin{equation}% label removed: eq:thqg-I-Z-g
Z_g(\cA) \;:=\; \int_{\overline{\mathcal{M}}_g} Z_{\mathcal{P}}(\cA;\,\mathbf{q}) \;\mu_{\mathrm{WP}}
\end{equation}
is well defined as an integral over the Deligne--Mumford compactification, where $\mu_{\mathrm{WP}}$ denotes the Weil--Petersson volume form.
\end{corollary}

\begin{proof}
The trace-class norm provides a pointwise bound on the integrand: $|Z_{\mathcal{P}}(\cA;\,\mathbf{q})| \leq \|T_{\mathbf{q}}\|_1$ where $T_{\mathbf{q}}$ is the sewing operator.  By the HS-bound (Proposition~\ref{prop:thqg-I-uniform-trace-bound}), this norm is bounded uniformly in the interior of the moduli space.  Near the boundary, the Fulton--MacPherson compactification provides the regulator: the sewing moduli $q_i$ vanish on the boundary divisors, and the $q$-weighted HS norms degenerate to zero.  The integral over the compact orbifold $\overline{\mathcal{M}}_g$ is therefore finite.
\end{proof}

\begin{remark}[The Deligne--Mumford boundary]
% label removed: rem:thqg-I-dm-boundary
\index{Deligne--Mumford!boundary}
The boundary $\partial\overline{\mathcal{M}}_g = \overline{\mathcal{M}}_g \setminus \mathcal{M}_g$ consists of stable curves with nodes.  Near a boundary stratum $\delta_\Gamma$ indexed by a dual graph $\Gamma$, the sewing moduli satisfy $q_e \to 0$ for each edge $e$ of $\Gamma$.  The integrand $Z_{\mathcal{P}}(\cA;\,\mathbf{q})$ decomposes as a product of amplitudes on the irreducible components of the stable curve, multiplied by propagator insertions at the nodes.  The vanishing of $q_e$ at the boundary ensures that the propagator insertions are trace class with trace-class norm tending to zero, giving the regularity needed for integrability over $\overline{\mathcal{M}}_g$.
\end{remark}


% ======================================================================
%  SUBSECTION 2: SHADOW FREE ENERGIES AND THE BERNOULLI GENERATING FUNCTION
% ======================================================================

\subsection{Shadow free energies and the Bernoulli generating function}
% label removed: subsec:thqg-I-shadow-free-energies
\index{shadow free energy|textbf}
\index{Bernoulli numbers!generating function}

The scalar projection of the Maurer--Cartan element $\Theta_\cA$ onto genus~$g$ is the shadow free energy $F_g(\cA)$.  This is the arity-$0$, genus-$g$ component of the shadow algebra $\cA^{\mathrm{sh}} = H_\bullet(\Defcyc^{\mathrm{mod}}(\cA))$ (Definition~\ref{def:shadow-algebra}).

\subsubsection{Definition of the shadow free energy}

\begin{definition}[Shadow free energy]
% label removed: def:thqg-I-shadow-free-energy
\index{shadow free energy!definition|textbf}
For a modular Koszul chiral algebra $\cA$ with universal Maurer--Cartan class $\Theta_\cA \in \mathrm{MC}(\gAmod)$ (Theorem~\ref*{V1-thm:mc2-bar-intrinsic}), the \emph{genus-$g$ shadow free energy} is the degree-zero shadow representation:
\begin{equation}% label removed: eq:thqg-I-Fg-def
F_g(\cA) \;:=\; \int_{\overline{\mathcal{M}}_g} \operatorname{Sh}_{g,0}(\Theta_\cA) \;=\; \int_{\overline{\mathcal{M}}_g} \kappa(\cA) \cdot \lambda_g,
\end{equation}
where the second equality uses the universality theorem (Theorem~\ref{thm:modular-characteristic}(i)): $\mathrm{obs}_g(\cA) = \kappa(\cA) \cdot \lambda_g$ in $H^{2g}(\overline{\mathcal{M}}_g)$, and $\lambda_g$ is the Hodge class $c_g(\mathbb{E})$.
\end{definition}

\begin{remark}[Relation to the shadow algebra]
% label removed: rem:thqg-I-shadow-algebra-relation
The shadow free energy $F_g(\cA)$ is the $(r,g) = (0,g)$ component of the shadow algebra $\cA^{\mathrm{sh}} = H_\bullet(\Defcyc^{\mathrm{mod}}(\cA))$ in the bigrading by arity~$r$ and genus~$g$.  The higher-arity shadows at genus $g$ (the cubic shadow $\mathfrak{C}_g$ at $r = 3$, the quartic resonance class $\mathfrak{Q}_g$ at $r = 4$) encode the gravitational interactions beyond the free-energy level.  The shadow free energy is the simplest invariant, controlled entirely by the modular characteristic $\kappa(\cA)$.
\end{remark}

\begin{remark}[Physical interpretation: vacuum energy]
% label removed: rem:thqg-I-vacuum-energy
\index{shadow free energy!physical interpretation}
The shadow free energy $F_g(\cA)$ is the genus-$g$ vacuum energy of the twisted holographic theory: it is the contribution to the path integral from genus-$g$ worldsheets with no external insertions.  In string theory language, $F_g$ is the genus-$g$ cosmological constant.  The crucial difference from ordinary string theory is that $F_g$ is determined \emph{exactly} by the modular characteristic $\kappa(\cA)$ through the Faber--Pandharipande integral: there is no moduli-dependent correction, and no need for renormalization.
\end{remark}

\subsubsection{The Faber--Pandharipande integral}

The key input from algebraic geometry is the evaluation of the Hodge integral:

\begin{proposition}[Faber--Pandharipande coefficient; \ClaimStatusProvedElsewhere]
% label removed: prop:thqg-I-faber-pandharipande
\index{Faber--Pandharipande!coefficient}
\index{Hodge integral}
The integral of the top Hodge class over the moduli space of genus-$g$ curves is
\begin{equation}% label removed: eq:thqg-I-fp-integral
\lambda_g^{\mathrm{FP}} \;:=\; \int_{\overline{\mathcal{M}}_g} \lambda_g \;=\; \frac{2^{2g-1} - 1}{2^{2g-1}} \cdot \frac{|B_{2g}|}{(2g)!}\,,
\end{equation}
where $B_{2g}$ is the $2g$-th Bernoulli number.  This formula is due to Faber--Pandharipande~\cite{FP00}.
\end{proposition}

\begin{proof}[Proof sketch]
The integral $\int_{\overline{\mathcal{M}}_g} \lambda_g$ is computed by localization on $\overline{\mathcal{M}}_g$ using the Mumford relation $\lambda_g = (-1)^g \kappa_{2g-2}/(2g-2)!$ and the Witten--Kontsevich theorem for intersection numbers of $\psi$-classes.  The Bernoulli numbers arise through the Euler--Maclaurin formula applied to the localization integral.  A detailed proof is in Faber--Pandharipande~\cite{FP00}; the formula also follows from the string equation and the topological recursion relation.
\end{proof}

\begin{remark}[Bernoulli numbers]
% label removed: rem:thqg-I-bernoulli
\index{Bernoulli numbers!values}
The first ten Bernoulli numbers: $B_0 = 1$, $B_1 = -1/2$, $B_2 = 1/6$, $B_4 = -1/30$, $B_6 = 1/42$, $B_8 = -1/30$, $B_{10} = 5/66$, $B_{12} = -691/2730$, $B_{14} = 7/6$, $B_{16} = -3617/510$, $B_{18} = 43867/798$, $B_{20} = -174611/330$.  The absolute values grow super-exponentially: $|B_{2g}| \sim 2 \cdot (2g)! / (2\pi)^{2g}$ (Euler's formula).  The alternating signs $(-1)^{g+1} B_{2g} > 0$ reflect the alternation of Hodge classes.
\end{remark}

\begin{proposition}[Shadow free energy in closed form; \ClaimStatusProvedElsewhere{} {\textup{Theorem~\ref{thm:modular-characteristic}}}]
% label removed: prop:thqg-I-Fg-closed-form
\index{shadow free energy!closed form}
For any modular Koszul chiral algebra $\cA$:
\begin{equation}% label removed: eq:thqg-I-Fg-formula
\boxed{ F_g(\cA) \;=\; \kappa(\cA) \cdot \frac{2^{2g-1}-1}{2^{2g-1}} \cdot \frac{|B_{2g}|}{(2g)!}\,. }
\end{equation}
In particular, $F_g$ depends only on $g$ and the single scalar $\kappa(\cA)$.
\end{proposition}

\begin{proof}
This is the combination of genus universality ($\mathrm{obs}_g = \kappa \cdot \lambda_g$, cf.\ Theorem~\ref{thm:genus-universality}) with the Faber--Pandharipande evaluation of the Hodge integral (Proposition~\ref{prop:thqg-I-faber-pandharipande}).
\end{proof}

\begin{remark}[Universality]
% label removed: rem:thqg-I-universality
\index{shadow free energy!universality}
The universality of the formula~\eqref{eq:thqg-I-Fg-formula} is striking: the genus-$g$ free energy of any modular Koszul chiral algebra depends only on the single scalar $\kappa(\cA)$, regardless of the detailed OPE structure.  The Heisenberg at level $k$, the affine $\widehat{\mathfrak{sl}}_2$ at level $k$, and the Virasoro at central charge $c$ all have the same functional form $F_g = \kappa \cdot \lambda_g^{\mathrm{FP}}$, differing only in the value of $\kappa$.  This is the shadow-level manifestation of the universality of the modular characteristic (Theorem~D): all nonlinear corrections are encoded in the higher-arity shadows, not in the scalar free energy.
\end{remark}

\subsubsection{The $\hat{A}$-genus generating function}

\begin{theorem}[Generating function for shadow free energies; \ClaimStatusProvedElsewhere{} {\textup{Theorem~\ref{thm:modular-characteristic}(ii)}}]
% label removed: thm:thqg-I-generating-function
\index{A-hat genus@$\hat{A}$-genus!generating function|textbf}
The formal generating series of shadow free energies has the closed-form expression:
\begin{equation}% label removed: eq:thqg-I-gf-ahat
\boxed{ \sum_{g=1}^{\infty} F_g(\cA)\,x^{2g} \;=\; \kappa(\cA) \cdot \left( \frac{x/2}{\sin(x/2)} - 1 \right). }
\end{equation}
Equivalently, with the convention $F_0(\cA) = \kappa(\cA)$:
\begin{equation}% label removed: eq:thqg-I-gf-ahat-alt
\sum_{g=0}^{\infty} F_g(\cA)\,x^{2g} \;=\; \kappa(\cA) \cdot \frac{x/2}{\sin(x/2)}.
\end{equation}
\end{theorem}

\begin{proof}
We use the standard generating function for Bernoulli numbers.  Recall that $\frac{t}{e^t - 1} = \sum_{n=0}^{\infty} \frac{B_n}{n!}\,t^n$.  The shadow free energy uses $\lambda_g^{\mathrm{FP}} = \frac{2^{2g-1}-1}{2^{2g-1}}\cdot\frac{|B_{2g}|}{(2g)!}$ with absolute values of Bernoulli numbers.  The generating function is therefore
\begin{align}
\sum_{g=1}^{\infty} F_g\,x^{2g} &\;=\; \kappa \sum_{g=1}^{\infty} \frac{2^{2g-1}-1}{2^{2g-1}} \cdot \frac{|B_{2g}|}{(2g)!}\,x^{2g} \notag\\
&\;=\; \kappa \left( \frac{x/2}{\sin(x/2)} - 1 \right), % label removed: eq:thqg-I-gf-derivation
\end{align}
where the second equality uses the well-known identity
\begin{equation}% label removed: eq:thqg-I-cosecant-bernoulli
\frac{x/2}{\sin(x/2)} \;=\; 1 + \sum_{g=1}^{\infty} \frac{(2^{2g-1}-1)|B_{2g}|}{2^{2g-1}\cdot(2g)!}\,x^{2g},
\end{equation}
obtainable from the Bernoulli generating function by the substitution $t = ix$ and the relation $\sin(x/2) = (e^{ix/2} - e^{-ix/2})/(2i)$.

The derivation proceeds as follows.  The Bernoulli generating function gives $\frac{t}{e^t - 1} = 1 - t/2 + \sum_{g=1}^\infty \frac{B_{2g}}{(2g)!} t^{2g}$.  Setting $t = ix$:
\[
\frac{ix}{e^{ix} - 1} \;=\; 1 - \frac{ix}{2} + \sum_{g=1}^\infty \frac{B_{2g}}{(2g)!}(ix)^{2g} \;=\; 1 - \frac{ix}{2} + \sum_{g=1}^\infty \frac{(-1)^g B_{2g}}{(2g)!} x^{2g}.
\]
Since $(-1)^g B_{2g} = -|B_{2g}|$ for $g \geq 1$, we have
\[
\frac{ix}{e^{ix}-1} \;=\; 1 - \frac{ix}{2} - \sum_{g=1}^\infty \frac{|B_{2g}|}{(2g)!}x^{2g}.
\]
Combining with the identity $\frac{ix}{e^{ix}-1} + \frac{ix}{2} = \frac{ix/2}{\tanh(ix/2)} = \frac{x/2}{\tan(x/2)}$, and using $\frac{x/2}{\sin(x/2)} = \frac{x/2}{\tan(x/2)} + \frac{x/2\,(1-\cos(x/2))}{\sin(x/2)\cos(x/2)}$, we arrive at~\eqref{eq:thqg-I-cosecant-bernoulli} by the standard manipulation of the partial fraction expansion $\csc(z) = 1/z + 2z\sum_{n=1}^\infty (-1)^n/(z^2 - n^2\pi^2)$.

More directly: the function $\frac{x/2}{\sin(x/2)}$ is even and analytic at $x = 0$ with value $1$.  Its Taylor expansion $\sum_{g=0}^\infty a_{2g} x^{2g}$ has coefficients $a_{2g} = \frac{(2^{2g-1}-1)|B_{2g}|}{2^{2g-1}(2g)!}$, which is verified by comparing with the standard expansion of $x/\sin(x) = \sum_{g=0}^\infty \frac{(2-2^{2-2g})|B_{2g}|}{(2g)!} x^{2g}$ and substituting $x \mapsto x/2$.
\end{proof}

\begin{remark}[Two generating functions]
% label removed: rem:thqg-I-two-gfs
There are two natural generating functions in play.  The \emph{algebraic} generating function $\sum_g F_g\,x^{2g} = \kappa(x/(2\sin(x/2)) - 1)$ uses the real variable $x$ and converges for $|x| < 2\pi$.  The \emph{physical} generating function $\sum_g F_g\,\hbar^g$ uses the genus-counting parameter~$\hbar$ and converges for $|\hbar| < 4\pi^2$.  The relation between the two is $\hbar = x^2$: the algebraic variable $x$ is the ``square root'' of the genus-counting parameter.
\end{remark}

\begin{remark}[Connection to the $\hat{A}$-genus]
% label removed: rem:thqg-I-ahat-connection
\index{A-hat genus@$\hat{A}$-genus!connection to shadow free energy}
The generating function $x/(2\sin(x/2))$ is the square root of the Todd genus $\mathrm{td}(x) = x/(1-e^{-x})$ and is closely related to the $\hat{A}$-genus $\hat{A}(x) = (x/2)/\sinh(x/2)$.  The relation is $x/(2\sin(x/2)) = \hat{A}(ix)$, a Wick rotation.  This Wick rotation corresponds to the passage from the Lorentzian $\hat{A}$-genus (which controls index theory on spin manifolds) to the Euclidean generating function (which controls genus integrals on moduli spaces).  The shadow free energy thus has a natural interpretation as a ``Wick-rotated $\hat{A}$-genus'' of the modular characteristic.
\end{remark}

\subsubsection{Absolute convergence from Bernoulli asymptotics}

\begin{theorem}[Absolute convergence of the generating function; \ClaimStatusProvedHere]
% label removed: thm:thqg-I-absolute-convergence
\index{absolute convergence!shadow free energy}
The genus-counting series $\sum_{g=1}^{\infty} |F_g(\cA)|\,|\hbar|^g$ converges absolutely for $|\hbar| < 4\pi^2$.  The generating function $\sum_{g=1}^{\infty} F_g(\cA)\,\hbar^g$ extends to a meromorphic function on $\mathbb{C}$ via the closed form $\kappa \cdot (\frac{\sqrt{\hbar}/2}{\sin(\sqrt{\hbar}/2)} - 1)$.
\end{theorem}

\begin{proof}
The asymptotic behavior of Bernoulli numbers is controlled by Euler's formula $|B_{2g}| = 2 \cdot \zeta(2g) \cdot (2g)! / (2\pi)^{2g}$, where $\zeta(2g) \to 1$ as $g \to \infty$.  Therefore
\begin{align}
|F_g(\cA)| &\;=\; |\kappa(\cA)| \cdot \frac{2^{2g-1}-1}{2^{2g-1}} \cdot \frac{|B_{2g}|}{(2g)!} \notag \\
&\;=\; |\kappa(\cA)| \cdot \frac{2^{2g-1}-1}{2^{2g-1}} \cdot \frac{2\,\zeta(2g)}{(2\pi)^{2g}} \notag \\
&\;\sim\; \frac{2|\kappa(\cA)|}{(2\pi)^{2g}}. % label removed: eq:thqg-I-Fg-asymptotics
\end{align}
More precisely, for $g \geq 2$:
\[
\frac{2|\kappa|}{(2\pi)^{2g}} \;\leq\; |F_g| \;\leq\; \frac{2|\kappa|}{(2\pi)^{2g}} \cdot \frac{1}{1 - 2^{1-2g}}.
\]
The ratio test gives $|F_{g+1}/F_g| \to 1/(2\pi)^2$ as $g \to \infty$, so the radius of convergence of $\sum F_g\,\hbar^g$ is $(2\pi)^2 = 4\pi^2$.  The closed form $\kappa \cdot (\frac{\sqrt{\hbar}/2}{\sin(\sqrt{\hbar}/2)} - 1)$ is meromorphic on $\mathbb{C}$ with poles at $\hbar = 4\pi^2 n^2$ ($n \geq 1$).

The pole structure is determined by the zeros of $\sin(\sqrt{\hbar}/2)$, which occur at $\sqrt{\hbar}/2 = n\pi$, i.e., $\hbar = 4\pi^2 n^2$.  Near $\hbar = 4\pi^2 n^2$, the function has a simple pole with residue:
\[
\Res_{\hbar = 4\pi^2 n^2} \kappa \cdot \frac{\sqrt{\hbar}/2}{\sin(\sqrt{\hbar}/2)} \;=\; \kappa \cdot \frac{n\pi}{\cos(n\pi)} \cdot \frac{1}{2\sqrt{4\pi^2 n^2}} \cdot 2\sqrt{4\pi^2 n^2} \;=\; \kappa \cdot (-1)^{n+1} \cdot 2n\pi.
\]
\end{proof}

\begin{remark}[Factorially divergent vs.\ convergent]
% label removed: rem:thqg-I-factorial-convergent
\index{factorial growth!absence of}
In standard perturbative quantum field theory, the genus-$g$ contribution typically grows as $(2g)!$, reflecting the factorial number of Feynman diagrams.  The shadow free energy $F_g \sim 2\kappa/(2\pi)^{2g}$ decays exponentially in $g$, not factorially.  The series $\sum F_g \hbar^g$ is convergent, not merely asymptotic.  The reason is that the shadow free energy is an integral over the compact moduli space $\overline{\mathcal{M}}_g$ of a smooth cohomology class, not a sum over all Feynman diagrams: the compactification absorbs the combinatorial proliferation into the geometry of the boundary strata, and the cohomological nature of the integrand prevents the amplitude from growing factorially.
\end{remark}

\subsubsection{Explicit values through genus 10}

\begin{proposition}[Shadow free energies through genus $10$; \ClaimStatusProvedHere]
% label removed: prop:thqg-I-Fg-values
\index{shadow free energy!explicit values}
Using the Faber--Pandharipande coefficients (Proposition~\ref{prop:thqg-I-faber-pandharipande}):
\begin{align}
F_1(\cA) &= \frac{\kappa}{24}, % label removed: eq:thqg-I-F1\\
F_2(\cA) &= \frac{7\kappa}{5760}, % label removed: eq:thqg-I-F2\\
F_3(\cA) &= \frac{31\kappa}{967680}, % label removed: eq:thqg-I-F3\\
F_4(\cA) &= \frac{127\kappa}{154828800}, % label removed: eq:thqg-I-F4\\
F_5(\cA) &= \frac{73\kappa}{3503554560}, % label removed: eq:thqg-I-F5\\
F_6(\cA) &= \frac{1414477\kappa}{2678117105664000}, % label removed: eq:thqg-I-F6\\
F_7(\cA) &= \frac{8191\kappa}{612141052723200}, % label removed: eq:thqg-I-F7\\
F_8(\cA) &= \frac{16931177\kappa}{49950709902213120000}, % label removed: eq:thqg-I-F8\\
F_9(\cA) &= \frac{5749691557\kappa}{669659197233029971968000}, % label removed: eq:thqg-I-F9\\
F_{10}(\cA) &= \frac{91546277357\kappa}{420928638260761696665600000}. % label removed: eq:thqg-I-F10
\end{align}
\end{proposition}

\begin{proof}
Direct substitution of the Bernoulli numbers into the formula $F_g = \kappa \cdot \frac{2^{2g-1}-1}{2^{2g-1}} \cdot \frac{|B_{2g}|}{(2g)!}$.

At genus $g = 1$: $F_1 = \kappa \cdot \frac{1}{2} \cdot \frac{1/6}{2} = \kappa/24$.

At genus $g = 2$: $(2^3{-}1)/2^3 = 7/8$, $|B_4| = 1/30$, $4! = 24$.  $F_2 = \kappa \cdot \frac{7}{8} \cdot \frac{1/30}{24} = \frac{7\kappa}{5760}$.

At genus $g = 3$: $|B_6| = 1/42$, $6! = 720$, $(2^5 - 1)/2^5 = 31/32$.  $F_3 = \kappa \cdot \frac{31}{32} \cdot \frac{1/42}{720} = \frac{31\kappa}{32 \cdot 30240} = \frac{31\kappa}{967680}$.

At genus $g = 4$: $|B_8| = 1/30$, $8! = 40320$, $(2^7 - 1)/2^7 = 127/128$.  $F_4 = \kappa \cdot \frac{127}{128} \cdot \frac{1/30}{40320} = \frac{127\kappa}{128 \cdot 1209600} = \frac{127\kappa}{154828800}$.

At genus $g = 5$: $|B_{10}| = 5/66$, $10! = 3628800$, $(2^9 - 1)/2^9 = 511/512$.  $F_5 = \kappa \cdot \frac{511}{512} \cdot \frac{5/66}{3628800} = \kappa \cdot \frac{511 \cdot 5}{512 \cdot 66 \cdot 3628800} = \frac{2555\kappa}{122624409600} = \frac{73\kappa}{3503554560}$.

The remaining values are computed similarly.
\end{proof}

\begin{table}[ht]
\centering
\caption{Shadow free energies $F_g(\cA) = \kappa \cdot \lambda_g^{\mathrm{FP}}$ for three families, $g = 1, \ldots, 5$.}
% label removed: tab:thqg-I-Fg-specializations
\index{shadow free energy!table}
\renewcommand{\arraystretch}{1.5}
{\footnotesize
\begin{tabular}{|c|c|c|c|}
\hline
$\boldsymbol{g}$ & $\mathcal{H}_k$, $\kappa = k$ & $\widehat{\mathfrak{sl}}_{2,k}$, $\kappa = \frac{3(k+2)}{4}$ & $\mathrm{Vir}_c$, $\kappa = \frac{c}{2}$ \\
\hline
$1$ & $\dfrac{k}{24}$ & $\dfrac{k+2}{32}$ & $\dfrac{c}{48}$ \\[6pt]
\hline
$2$ & $\dfrac{7k}{5760}$ & $\dfrac{7(k+2)}{7680}$ & $\dfrac{7c}{11520}$ \\[6pt]
\hline
$3$ & $\dfrac{31k}{967680}$ & $\dfrac{31(k+2)}{1290240}$ & $\dfrac{31c}{1935360}$ \\[6pt]
\hline
$4$ & $\dfrac{127k}{154828800}$ & $\dfrac{127(k+2)}{206438400}$ & $\dfrac{127c}{309657600}$ \\[6pt]
\hline
$5$ & $\dfrac{73k}{3503554560}$ & $\dfrac{73(k+2)}{4671406080}$ & $\dfrac{73c}{7007109120}$ \\[6pt]
\hline
\end{tabular}}
\end{table}

\begin{remark}[Numerical magnitudes]
% label removed: rem:thqg-I-numerical-magnitudes
\index{shadow free energy!numerical estimates}
For the Heisenberg at $k = 1$: $F_1 \approx 0.0417$, $F_2 \approx 1.22 \times 10^{-3}$, $F_3 \approx 3.20 \times 10^{-5}$, $F_4 \approx 8.20 \times 10^{-7}$, $F_5 \approx 2.08 \times 10^{-8}$.  The free energies decrease super-exponentially, with ratio $F_{g+1}/F_g \to 1/(2\pi)^2 \approx 0.0253$.  By genus $10$, $F_{10} \approx 2.17 \times 10^{-16}$.  The geometric decay with ratio $1/(4\pi^2) \approx 0.0253$ means that the genus expansion is an extremely rapidly converging series, even at $\hbar = 1$.
\end{remark}

\begin{proposition}[Cumulative partial sums; \ClaimStatusProvedHere]
% label removed: prop:thqg-I-partial-sums
\index{shadow free energy!partial sums}
For the Heisenberg at $k = 1$ and $\hbar = 1$, the partial sums of the genus expansion are:
\begin{align}
S_1 &= F_1 = 0.04167\,, \notag \\
S_2 &= S_1 + F_2 = 0.04288\,, \notag \\
S_3 &= S_2 + F_3 = 0.04291\,, \notag \\
S_4 &= S_3 + F_4 = 0.04291\,, \notag \\
S_5 &= S_4 + F_5 = 0.04291\,, % label removed: eq:thqg-I-partial-sums
\end{align}
converging to $\kappa \cdot (1/\mathrm{sinc}(1/2) - 1) = (1/(2\sin(1/2)) - 1) \cdot 1 \approx 0.04291$.
\end{proposition}

\begin{proof}
Direct numerical evaluation.
\end{proof}


% ======================================================================
%  SUBSECTION 3: CONVERGENCE OF THE GRAVITATIONAL PARTITION FUNCTION
% ======================================================================

\subsection{Convergence of the gravitational partition function}
% label removed: subsec:thqg-I-convergence-grav
\index{gravitational partition function!convergence|textbf}

\subsubsection{Definition of the gravitational partition function}

\begin{definition}[Gravitational partition function]
% label removed: def:thqg-I-grav-partition
\index{gravitational partition function!definition|textbf}
For a holographic modular Koszul datum $\mathcal{H}(T) = (\cA, \cA^!, \mathcal{C}, r(z), \Theta_\cA, \nabla^{\mathrm{hol}})$ and a formal genus-counting parameter~$\hbar$, the \emph{gravitational partition function} is
\begin{equation}% label removed: eq:thqg-I-Z-grav-def
Z_{\mathrm{grav}}(T;\,\hbar) \;:=\; \sum_{g=0}^{\infty} \hbar^g \int_{\overline{\mathcal{M}}_{g}} \operatorname{Sh}_{g,0}(\Theta_\cA)\,.
\end{equation}
The \emph{scalar gravitational partition function} is the restriction to the arity-$0$ scalar channel:
\begin{equation}% label removed: eq:thqg-I-Z-grav-scalar
Z_{\mathrm{grav}}^{\mathrm{scal}}(T;\,\hbar) \;:=\; \sum_{g=0}^{\infty} F_g(\cA)\,\hbar^g \;=\; \kappa(\cA) \cdot \frac{\sqrt{\hbar}/2}{\sin(\sqrt{\hbar}/2)}\,,
\end{equation}
where the closed form follows from Theorem~\ref{thm:thqg-I-generating-function}.
\end{definition}

\begin{remark}[Scalar versus full]
% label removed: rem:thqg-I-scalar-vs-full
The scalar partition function captures only the trace of the modular functor.  The full gravitational partition function includes the higher-arity contributions from the shadow tower: the spectral discriminant $\Delta_\cA$ (arity~$2$ spectral channel), the cubic shadow $\mathfrak{C}$ (arity~$3$), the quartic resonance class $\mathfrak{Q}$ (arity~$4$), and all higher shadows.  For Gaussian algebras (shadow depth $r_{\max} = 2$), the scalar and full partition functions coincide.
\end{remark}

\begin{remark}[Genus-$0$ normalization]
% label removed: rem:thqg-I-genus-0
\index{gravitational partition function!genus-0 normalization}
The genus-$0$ contribution $F_0(\cA) = \kappa(\cA)$ is a convention: it is the ``classical'' contribution from the sphere.  The physical partition function is $Z_{\mathrm{grav}} = \kappa + \sum_{g \geq 1} F_g \hbar^g$, where the leading term $\kappa$ is the tree-level answer and the remaining terms are quantum corrections.  In the semiclassical limit $\hbar \to 0$, only the tree-level term survives.
\end{remark}

\subsubsection{Convergence from HS-sewing and Bernoulli asymptotics}

\begin{theorem}[Perturbative finiteness of twisted gravity; \ClaimStatusProvedHere]
% label removed: thm:thqg-I-perturbative-finiteness
\index{perturbative finiteness!main theorem|textbf}
Let $\cA$ be a modular Koszul chiral algebra satisfying the HS-sewing condition \textup{(}Theorem~\textup{\ref{thm:general-hs-sewing})}.  Then:
\begin{enumerate}[label=\textup{(\roman*)}]
\item \emph{Genus-by-genus finiteness.}  For each $g \geq 1$, the genus-$g$ shadow amplitude $\int_{\overline{\mathcal{M}}_g} \operatorname{Sh}_{g,n}(\Theta_\cA)$ is finite at every arity~$n$.
\item \emph{Scalar partition function.}  The series $\sum_{g=1}^{\infty} F_g(\cA)\,\hbar^g$ converges absolutely for $|\hbar| < 4\pi^2$, with the closed-form expression $Z_{\mathrm{grav}}^{\mathrm{scal}}(\cA;\,\hbar) = \kappa(\cA) \cdot \frac{\sqrt{\hbar}/2}{\sin(\sqrt{\hbar}/2)}$.
\item \emph{Full partition function.}  The full genus-summed partition function $Z_{\mathrm{grav}}(T;\,\hbar)$ converges absolutely for $|\hbar| < R_{\mathrm{grav}}$, where $R_{\mathrm{grav}} \geq 4\pi^2$.  For Gaussian algebras, $R_{\mathrm{grav}} = \infty$.
\item \emph{No UV renormalization.}  No counterterms or renormalization are needed at any genus.  The Maurer--Cartan equation $D_\cA\Theta_\cA + \tfrac{1}{2}[\Theta_\cA,\Theta_\cA] = 0$ (Theorem~\ref*{V1-thm:mc2-bar-intrinsic}) is an exact identity, and the Fulton--MacPherson compactification provides the geometric regulator.
\end{enumerate}
\end{theorem}

\begin{proof}
\emph{Part~(i).}  Fix genus~$g$ and arity~$n$.  The shadow representation $\operatorname{Sh}_{g,n}(\Theta_\cA) \in H^*(\overline{\mathcal{M}}_{g,n})$ is a cohomology class on a compact orbifold.  Integration over $\overline{\mathcal{M}}_{g,n}$ is a continuous linear functional on this finite-dimensional space, hence bounded.  The nontrivial content is that the class itself is well defined: this follows from the bar-intrinsic construction of $\Theta_\cA$ (Theorem~\ref*{V1-thm:mc2-bar-intrinsic}), which gives $\Theta_\cA := D_\cA - d_0 \in \mathrm{MC}(\gAmod)$ without convergence issues.

At the chain level, the integral is computed by the arity-cutoff lemma (Lemma~\ref{lem:arity-cutoff}): the graph sum through arity~$r$ at genus~$g$ involves only finitely many graphs $\Gamma$ with $|\Gamma| \leq r$ edges and loop genus $\leq g$.  Each graph amplitude $\ell_\Gamma = |\Aut(\Gamma)|^{-1} \int_{\overline{\mathcal{M}}_{g(\Gamma),n(\Gamma)}} \omega_\Gamma$ is a period integral on a compact moduli space, hence finite.

\emph{Part~(ii).}  This is Theorem~\ref{thm:thqg-I-absolute-convergence}: the Bernoulli asymptotics $|F_g| \sim 2|\kappa|/(2\pi)^{2g}$ give absolute convergence for $|\hbar| < 4\pi^2$, and analytic continuation via the closed form extends to all of~$\mathbb{C}$ as a meromorphic function.

\emph{Part~(iii).}  For the full partition function, each genus-$g$ contribution is a sum over the shadow tower.  The shadow tower truncation $\Theta_\cA^{\leq r}$ is a finite sum at each arity, and the inverse limit $\Theta_\cA = \varprojlim \Theta_\cA^{\leq r}$ converges (Theorem~\ref{thm:recursive-existence}).  The arity-$r$ contribution to $F_g$ is bounded by $C_r / (2\pi)^{2g}$, so the genus sum converges at each arity.  For Gaussian algebras ($r_{\max} = 2$), all arity-$r \geq 3$ contributions vanish, and the full partition function equals the scalar partition function; the latter converges for all $|\hbar| < 4\pi^2$, and the higher-arity vanishing means $R_{\mathrm{grav}} = \infty$ (the function is entire).

\emph{Part~(iv).}  The absence of UV divergences is the combined consequence of three features:
\begin{enumerate}[label=\textup{(\alph*)}]
\item \emph{Compactness.}  The Fulton--MacPherson compactification $\overline{\mathrm{FM}}_n(\Sigma_g)$ is a compact manifold with corners.  All integrals are over compact domains.
\item \emph{Exactness.}  The Maurer--Cartan equation is an identity of the bar differential ($D_\cA^2 = 0$), not a perturbative expansion.  There are no UV divergences to regulate because the equation is exact, not asymptotic.
\item \emph{Arity cutoff.}  At each genus~$g$, only finitely many graph topologies contribute (Lemma~\ref{lem:arity-cutoff}).  The number of contributing graphs at genus $g$ and arity $r$ is bounded by a combinatorial function of $(g,r)$, and each graph amplitude is a finite integral.
\end{enumerate}
\end{proof}

\subsubsection{Comparison with standard QFT divergences}

\begin{remark}[Contrast with perturbative quantum gravity in $d = 4$]
% label removed: rem:thqg-I-contrast-4d
\index{UV divergence!standard quantum gravity}
In standard four-dimensional perturbative quantum gravity (Einstein--Hilbert action), the genus-$g$ contribution to the graviton scattering amplitude diverges for $g \geq 2$ as
\begin{equation}% label removed: eq:thqg-I-4d-divergence
\mathcal{A}_g \;\sim\; \Lambda_{\mathrm{UV}}^{2(g-1)} \int_{\mathcal{M}_g} |\Omega_g|^2 \cdot R^{2g},
\end{equation}
where $\Lambda_{\mathrm{UV}}$ is the UV cutoff and $R$ is the Riemann curvature.  In the twisted holographic setting, the analogous quantity $\mathcal{A}_g^{\mathrm{tw}} = \int_{\overline{\mathcal{M}}_g} \operatorname{Sh}_{g,n}(\Theta_\cA)$ is finite at every genus because: (a)~the integration domain $\overline{\mathcal{M}}_g$ is compact; (b)~the integrand is a smooth cohomology class; (c)~the Maurer--Cartan equation provides an exact Ward identity constraining the amplitudes genus by genus.  The fundamental difference is that the twisted holographic gravity lives on the moduli space of curves, not on spacetime.
\end{remark}

\begin{remark}[Contrast with string theory]
% label removed: rem:thqg-I-contrast-string
\index{string theory!UV finiteness}
In string theory, genus-$g$ amplitudes are finite individually (after integration over $\overline{\mathcal{M}}_g$), but the genus expansion $\sum_g g_s^{2g-2} \mathcal{A}_g$ is asymptotic: it diverges factorially as $\mathcal{A}_g \sim (2g)! \cdot \alpha'^{-g}$.  In the twisted holographic setting, the individual amplitudes are also finite (same reason: compactness of $\overline{\mathcal{M}}_g$), but the genus expansion \emph{converges}: $F_g \sim \kappa/(2\pi)^{2g}$ decays exponentially.  The difference stems from the cohomological nature of the shadow amplitudes: in string theory, the integrand over $\mathcal{M}_g$ is a genuine function that grows with $g$; in the shadow theory, the integrand is a fixed cohomology class $\kappa \cdot \lambda_g$ whose integral $\lambda_g^{\mathrm{FP}}$ decays.
\end{remark}

\subsubsection{The role of Koszulness in finiteness}

\begin{proposition}[Koszulness and finiteness; \ClaimStatusProvedHere]
% label removed: prop:thqg-I-koszulness-finiteness
\index{Koszulness!role in finiteness}
For a modular Koszul chiral algebra $\cA$, the shadow tower $\Theta_\cA^{\leq r}$ has the following finiteness properties:
\begin{enumerate}[label=\textup{(\roman*)}]
\item At each arity $r$ and genus $g$, the graph sum
  \begin{equation}% label removed: eq:thqg-I-graph-sum-finite
  \Theta_\cA^{(g,r)} \;=\; \sum_{\substack{\Gamma \in \Gmod(g,r) \\ \Gamma \text{ connected}}} \frac{1}{|\Aut(\Gamma)|} \int_{\overline{\mathcal{M}}_{g(\Gamma),n(\Gamma)}} \omega_\Gamma
  \end{equation}
  is a finite sum over a finite set of graphs.
\item The Koszul property ensures that the transferred cyclic minimal model has only finitely many nontrivial operations at each arity, so the vertex labels $\ell_n^{\mathrm{tr}}$ are nonzero only for $n$ in a finite range.
\item The shadow depth $r_{\max}(\cA)$ bounds the arity of contributing graphs: for $r > r_{\max}$, the obstruction class $o_{r+1}(\cA) = 0$ and the shadow tower stabilizes.
\end{enumerate}
\end{proposition}

\begin{proof}
Part~(i) is the arity-cutoff lemma (Lemma~\ref{lem:arity-cutoff}).  The set $\Gmod(g,r)$ of connected modular graphs with loop genus $g$ and total arity $\leq r$ is finite: the number of vertices $v$, edges $e$, and external legs $n$ satisfy $v - e + n = 1 - g$ (the Euler characteristic), $2e + n \leq rv$ (each vertex has at most $r$ half-edges), and $g = e - v + 1 + \sum_{v'} g_{v'}$ where $g_{v'}$ is the genus decoration at each vertex.  These constraints bound $v, e, n$ in terms of $g$ and $r$.

Part~(ii) follows from the PBW concentration property (Proposition~\ref{prop:standard-examples-modular-koszul}): for a Koszul algebra, the transferred $A_\infty$ operations $m_n^{\mathrm{tr}}$ vanish for $n$ sufficiently large (bounded by the Koszul dual degree).

Part~(iii) is the definition of shadow depth: if $o_{r+1}(\cA) = 0$, then $\Theta_\cA^{\leq r+1} = \Theta_\cA^{\leq r}$, and the tower stabilizes.
\end{proof}

\begin{proposition}[Effective graph count; \ClaimStatusProvedHere]
% label removed: prop:thqg-I-graph-count
\index{graph sum!effective count}
The number of connected modular graphs contributing to $\Theta_\cA^{(g,r)}$ is bounded by
\begin{equation}% label removed: eq:thqg-I-graph-bound
|\Gmod(g,r)| \;\leq\; C(g,r) \;:=\; \sum_{v=1}^{3g-3+r} \binom{v(r-1)/2 + g - 1}{g} \cdot r^v,
\end{equation}
which is finite for fixed $(g,r)$ and grows polynomially in $r$ at fixed $g$.
\end{proposition}

\begin{proof}
A connected graph with $v$ vertices and loop genus $g$ has $e = v + g - 1$ edges.  Each vertex has valence at most $r$ (by the arity cutoff), so $2e + n \leq rv$, giving $n \leq rv - 2(v+g-1) = v(r-2) - 2g + 2$.  The condition $n \geq 0$ requires $v \geq (2g-2)/(r-2)$ (for $r \geq 3$).  The number of graphs with $v$ labeled vertices and $e$ edges is at most $\binom{v(v-1)/2}{e}$ (choosing edges from all possible pairs); the symmetry group reduces this, but the crude bound suffices.  The factor $r^v$ accounts for the vertex labeling (choice of operation at each vertex from $\{m_2, m_3, \ldots, m_r\}$).
\end{proof}

\subsubsection{The structure of finiteness at each genus}

\begin{proposition}[Genus-$g$ amplitude decomposition; \ClaimStatusProvedHere]
% label removed: prop:thqg-I-genus-g-decomposition
\index{genus-$g$ amplitude!decomposition}
The genus-$g$ shadow amplitude decomposes as
\begin{equation}% label removed: eq:thqg-I-genus-g-decomp
\int_{\overline{\mathcal{M}}_g} \operatorname{Sh}_{g,0}(\Theta_\cA) \;=\; F_g^{\mathrm{scal}}(\cA) \;+\; \sum_{r=2}^{r_{\max}} F_g^{(r)}(\cA),
\end{equation}
where $F_g^{\mathrm{scal}} = \kappa \cdot \lambda_g^{\mathrm{FP}}$ is the scalar (arity-$0$) contribution and $F_g^{(r)}$ is the arity-$r$ correction.  For a Gaussian algebra ($r_{\max} = 2$), the sum is empty and only the scalar term survives.  For a Lie-type algebra ($r_{\max} = 3$), there is a single cubic correction $F_g^{(3)}$.  For a contact-type algebra ($r_{\max} = 4$), there are cubic and quartic corrections.  For a mixed algebra ($r_{\max} = \infty$), the sum is infinite but converges by the shadow tower convergence theorem (Theorem~\ref{thm:recursive-existence}).
\end{proposition}

\begin{proof}
The shadow tower decomposes by arity: $\Theta_\cA = \sum_{r \geq 0} \Theta_\cA^{(r)}$ where $\Theta_\cA^{(r)}$ is the arity-$r$ component.  The genus-$g$ projection at arity $0$ gives $F_g^{\mathrm{scal}}$ by the universality theorem.  At arity $r \geq 2$, the projection $\int_{\overline{\mathcal{M}}_g} \operatorname{Sh}_{g,r}(\Theta_\cA)$ involves a graph sum with vertex labels from the arity-$r$ operations, and each term is finite by the arity-cutoff argument.  The vanishing of $F_g^{(r)}$ for $r > r_{\max}$ follows from the shadow depth.
\end{proof}


% ======================================================================
%  SUBSECTION 4: THE HEISENBERG PARTITION FUNCTION AS PROTOTYPE
% ======================================================================

\subsection{The Heisenberg partition function as prototype}
% label removed: subsec:thqg-I-heisenberg-prototype
\index{Heisenberg!partition function|textbf}

The Heisenberg algebra is the archetype: Gaussian shadow depth ($r_{\max} = 2$), exact solvability via second quantization, and explicit Fredholm determinant formulas at all genera.

\subsubsection{The Heisenberg at genus $g$}

\begin{theorem}[Heisenberg gravitational partition function; \ClaimStatusProvedElsewhere{} {\textup{Theorem~\ref{thm:heisenberg-sewing}}}]
% label removed: thm:thqg-I-heisenberg-partition
\index{Heisenberg!genus-$g$ partition function}
For the rank-$1$ Heisenberg algebra $\mathcal{H}_k$ on a genus-$g$ surface $\Sigma_g$ with period matrix $\Omega \in \mathbb{H}_g$ \textup{(}the Siegel upper half-space\textup{)}, the genus-$g$ partition function is the Fredholm determinant
\begin{equation}% label removed: eq:thqg-I-heisenberg-Zg
\boxed{ Z_g(\mathcal{H}_k;\,\Omega) \;=\; \bigl(\det\operatorname{Im}\Omega\bigr)^{-k/2} \cdot \bigl(\det{}'_\zeta \Delta_{\Sigma_g}\bigr)^{-k/2}, }
\end{equation}
where $\det{}'_\zeta$ denotes the zeta-regularized determinant of the scalar Laplacian $\Delta_{\Sigma_g}$ on $\Sigma_g$ \textup{(}excluding the zero mode\textup{)}.  The shadow free energy extracted from this is
\begin{equation}% label removed: eq:thqg-I-heisenberg-Fg
F_g(\mathcal{H}_k) \;=\; k \cdot \frac{2^{2g-1}-1}{2^{2g-1}} \cdot \frac{|B_{2g}|}{(2g)!}\,,
\end{equation}
confirming the general formula~\eqref{eq:thqg-I-Fg-formula} at $\kappa(\mathcal{H}_k) = k$.
\end{theorem}

\begin{proof}
By Wick's theorem, all amplitudes of the Heisenberg algebra factorize into propagator pairings.  The propagator on a genus-$g$ surface is the Green's function
\begin{equation}% label removed: eq:thqg-I-green-genus-g
G(z,w;\,\Omega) \;=\; \log|E(z,w;\,\Omega)| - \frac{(\operatorname{Im}\int_w^z \omega)^T (\operatorname{Im}\Omega)^{-1} (\operatorname{Im}\int_w^z \omega)}{2},
\end{equation}
where $E(z,w;\,\Omega)$ is the prime form and $\omega = (\omega_1, \ldots, \omega_g)^T$ is the normalized basis of holomorphic differentials.  The prime form is a section of $K^{-1/2} \boxtimes K^{-1/2}$ (not $K^{+1/2}$).

The partition function is the functional determinant of the kinetic operator:
\[
Z_g \;=\; (\det{}'_\zeta \bar{\partial}^*\bar{\partial})^{-k/2} \;=\; (\det\operatorname{Im}\Omega)^{-k/2}\, (\det{}'_\zeta \Delta)^{-k/2},
\]
where the factorization into period-matrix and Laplacian contributions is the Belavin--Knizhnik factorization~\cite{BK86}.  The Polyakov formula identifies $\det{}'_\zeta \Delta$ with the norm of the Mumford form on $\overline{\mathcal{M}}_g$; its $g$-th Chern class is $\lambda_g$.  The shadow free energy is obtained by integrating the Chern class:
\[
F_g(\mathcal{H}_k) = k \int_{\overline{\mathcal{M}}_g} \lambda_g = k \cdot \lambda_g^{\mathrm{FP}}.
\]
\end{proof}

\begin{remark}[The Belavin--Knizhnik factorization]
% label removed: rem:thqg-I-belavin-knizhnik
\index{Belavin--Knizhnik!factorization}
The factorization $\det{}'_\zeta(\bar\partial^*\bar\partial) = \det(\operatorname{Im}\Omega) \cdot \det{}'_\zeta(\Delta)$ separates the ``holomorphic'' contribution (the period matrix, which depends on the complex structure of $\Sigma_g$) from the ``analytic'' contribution (the Laplacian determinant, which depends on the conformal class).  The period matrix factor $(\det\operatorname{Im}\Omega)^{-k/2}$ is the modular-covariant piece that transforms under $\mathrm{Sp}(2g,\mathbb{Z})$, while the Laplacian factor is modular-invariant.  Together they give the conformal block for the Heisenberg algebra on $\Sigma_g$.
\end{remark}

\subsubsection{The Fredholm determinant formula}

\begin{proposition}[Fredholm determinant representation; \ClaimStatusProvedElsewhere{} {\textup{Theorem~\ref{thm:heisenberg-one-particle-sewing}}}]
% label removed: prop:thqg-I-fredholm
\index{Fredholm determinant!Heisenberg|textbf}
The genus-$g$ partition function admits a Fredholm determinant representation:
\begin{equation}% label removed: eq:thqg-I-fredholm-det
Z_g(\mathcal{H}_k;\,\Omega) \;=\; \det(1 - K_g)^{-k},
\end{equation}
where $K_g$ is the trace-class sewing operator on the one-particle Bergman space $A^2(\partial D)$, with eigenvalues $q_1^{n_1} \cdots q_{3g-3}^{n_{3g-3}}$ indexed by the sewing moduli $(q_1, \ldots, q_{3g-3})$.  The trace-class norm satisfies
\begin{equation}% label removed: eq:thqg-I-Kg-trace-norm
\|K_g\|_1 \;=\; \prod_{i=1}^{3g-3} \frac{q_i}{1-q_i} \;<\; \infty \quad\text{for all } q_i \in (0,1).
\end{equation}
\end{proposition}

\begin{proof}
The pair-of-pants amplitude for the Heisenberg is the second quantization $\Gamma(R)$ of the one-particle Bergman restriction $R \colon A^2(D_{\mathrm{out}}) \to A^2(D_{\mathrm{in},1}) \otimes A^2(D_{\mathrm{in},2})$ (Theorem~\ref{thm:heisenberg-one-particle-sewing}).  The one-particle restriction has matrix entries $R_{n,m} = \frac{k}{2\pi i} \oint\oint z^{-n-1}w^{-m-1}G_{\mathrm{pants}}(z,w)\,dz\,dw$, satisfying $|R_{n,m}| \leq C\,q^{(n+m)/2}$ from the collar length $t = -\log q$.  Hence $R$ is Hilbert--Schmidt and $\Gamma(R)$ is trace class.

The composition of $3g - 3$ sewing operators gives the genus-$g$ operator $K_g = R_1 \circ \cdots \circ R_{3g-3}$ on the one-particle space.  The Fredholm determinant identity $\Tr\,\Gamma(K_g)^k = \det(1 - K_g)^{-k}$ is the boson-fermion correspondence.  The trace-class norm bound follows from $\|K_g\|_1 \leq \prod_{i=1}^{3g-3} \|R_i\|_1 = \prod_{i=1}^{3g-3} q_i/(1-q_i)$.
\end{proof}

\begin{remark}[The Bergman space]
% label removed: rem:thqg-I-bergman
\index{Bergman space}
The one-particle Bergman space $A^2(\partial D)$ consists of holomorphic functions on the annular collar around a sewing circle, square-integrable with respect to the Bergman kernel.  An orthonormal basis is $\{z^n / \sqrt{2\pi}\}_{n \geq 1}$, and the one-particle sewing operator in this basis has matrix elements $K_{nm} = \delta_{nm} q^n$.  The Fredholm determinant is therefore $\det(1 - K) = \prod_{n=1}^\infty (1 - q^n)$, giving the Dedekind eta function upon including the zero-point energy.
\end{remark}

\subsubsection{Comparison with the Dedekind eta function}

\begin{example}[Genus $1$: the Dedekind eta function]
% label removed: ex:thqg-I-genus1-eta
\index{Dedekind eta!Heisenberg partition function}
At genus $g = 1$, the period matrix is $\Omega = \tau \in \mathbb{H}$ and the sewing parameter is $q = e^{2\pi i\tau}$.  The Fredholm determinant becomes
\begin{equation}% label removed: eq:thqg-I-genus1-fredholm
Z_1(\mathcal{H}_k;\,\tau) \;=\; (\operatorname{Im}\tau)^{-k/2}\, \prod_{n=1}^{\infty} (1 - q^n)^{-k} \;=\; (\operatorname{Im}\tau)^{-k/2}\, \bigl|\eta(\tau)\bigr|^{-2k},
\end{equation}
where $\eta(\tau) = q^{1/24}\prod_{n=1}^{\infty}(1 - q^n)$ is the Dedekind eta function.  The shadow free energy at genus $1$ is $F_1(\mathcal{H}_k) = k/24$, matching~\eqref{eq:thqg-I-F1} with $\kappa = k$.

The connection to the modular characteristic is transparent: the $q$-expansion of $\log \eta(\tau) = \frac{2\pi i\tau}{24} - \sum_{n=1}^\infty \log(1-q^n)$ begins with $\frac{2\pi i \tau}{24}$, and the genus-$1$ free energy is the constant term of the $q$-expansion of $-k\log|\eta(\tau)|$, which is $k/24$.
\end{example}

\begin{example}[Genus $2$: explicit formula]
% label removed: ex:thqg-I-genus2
\index{Heisenberg!genus-2 partition function}
At genus $g = 2$, the period matrix $\Omega = \begin{psmallmatrix} \tau_1 & \tau_{12} \\ \tau_{12} & \tau_2 \end{psmallmatrix} \in \mathbb{H}_2$ and the partition function is
\begin{equation}% label removed: eq:thqg-I-genus2
Z_2(\mathcal{H}_k;\,\Omega) \;=\; \bigl(\det\operatorname{Im}\Omega\bigr)^{-k/2} \cdot \bigl(\det{}'_\zeta \Delta_{\Sigma_2}\bigr)^{-k/2}.
\end{equation}
The shadow free energy is $F_2(\mathcal{H}_k) = 7k/5760$, confirming~\eqref{eq:thqg-I-F2}.

The Fredholm determinant at genus $2$ involves the sewing of three pairs of pants along three circles.  In the Schottky uniformization, $\Sigma_2 = (\mathbb{P}^1 \setminus \{6 \text{ discs}\}) / \langle \gamma_1, \gamma_2, \gamma_3 \rangle$ with three sewing parameters $q_1, q_2, q_3$.  The one-particle sewing operator is $K_2 = R_1 \circ R_2 \circ R_3$, and $\det(1 - K_2) = \prod_{n_1,n_2,n_3 \geq 1} (1 - q_1^{n_1} q_2^{n_2} q_3^{n_3})$.  The shadow free energy is obtained by integrating over the Siegel modular variety $\mathcal{M}_2 = \mathbb{H}_2 / \mathrm{Sp}(4,\mathbb{Z})$.
\end{example}

\begin{example}[Genus $3$]
% label removed: ex:thqg-I-genus3-heisenberg
\index{Heisenberg!genus-3 partition function}
At genus $g = 3$, the surface $\Sigma_3$ has $3g - 3 = 6$ sewing parameters.  The Fredholm determinant is $\det(1 - K_3)^{-k}$ where $K_3$ is a trace-class operator on the one-particle Bergman space.  The shadow free energy is $F_3(\mathcal{H}_k) = 31k/967680$.  Numerically at $k = 1$: $F_3 \approx 3.20 \times 10^{-5}$, which is already three orders of magnitude smaller than $F_1 = 1/24 \approx 0.042$.
\end{example}

\subsubsection{One-particle Bergman reduction}

\begin{remark}[One-particle reduction principle]
% label removed: rem:thqg-I-one-particle
\index{Bergman space!one-particle reduction|textbf}
The Heisenberg partition function factorizes completely into one-particle data; this is the second quantization principle.  The one-particle sewing operator $T_q$ on the Bergman space $A^2(\partial D)$ has eigenvalues $q^n$ ($n = 1, 2, 3, \ldots$), and the full partition function is
\[
Z_g(\mathcal{H}_k) \;=\; \det(1 - T_{q_1} \circ \cdots \circ T_{q_{3g-3}})^{-k} \;=\; \exp\!\left( k \sum_{m=1}^{\infty} \frac{1}{m} \Tr\bigl(T_{q_1} \cdots T_{q_{3g-3}}\bigr)^m \right).
\]
This exponential formula makes the connection to the shadow free energy transparent: $F_g$ is the constant term of $\log Z_g$ in the expansion in sewing moduli.  The one-particle reduction is specific to the Heisenberg (Gaussian shadow depth $r_{\max} = 2$): for non-Gaussian algebras, the multi-particle sector contributes genuinely new data that cannot be recovered from the one-particle space.
\end{remark}

\begin{proposition}[Exponential representation; \ClaimStatusProvedHere]
% label removed: prop:thqg-I-exponential-rep
\index{Fredholm determinant!exponential representation}
The Fredholm determinant admits the exponential representation
\begin{equation}% label removed: eq:thqg-I-exponential-fredholm
\log \det(1 - K_g)^{-k} \;=\; k \sum_{m=1}^\infty \frac{1}{m} \Tr(K_g^m) \;=\; k \sum_{m=1}^\infty \frac{1}{m} \sum_{\mathbf{n} \in \mathbb{Z}_{\geq 1}^{3g-3}} \prod_{i=1}^{3g-3} q_i^{m n_i},
\end{equation}
which converges absolutely for all $q_i \in (0,1)$.  The shadow free energy $F_g$ is extracted as the coefficient of the leading term in the $\mathbf{q}$-expansion, integrated over $\overline{\mathcal{M}}_g$ with the Weil--Petersson measure.
\end{proposition}

\begin{proof}
The Fredholm determinant identity $\log\det(1 - K) = \Tr\log(1 - K) = -\sum_{m=1}^\infty \frac{1}{m}\Tr(K^m)$ holds for trace-class $K$ with $\|K\|_1 < 1$.  For $q_i \in (0,1)$, the eigenvalues of $K_g$ are $\prod_{i=1}^{3g-3} q_i^{n_i}$ with $n_i \geq 1$, all less than $1$.  The trace $\Tr(K_g^m) = \sum_{\mathbf{n}} \prod q_i^{mn_i}$ converges absolutely because $\sum_{\mathbf{n}} \prod q_i^{n_i} = \prod_{i=1}^{3g-3} q_i/(1-q_i) < \infty$.
\end{proof}

\begin{remark}[Rank-$d$ Heisenberg]
% label removed: rem:thqg-I-rank-d-heisenberg
\index{Heisenberg!rank-$d$}
For the rank-$d$ Heisenberg algebra $\mathcal{H}_k^{\otimes d}$ (i.e., $d$ independent free bosons at level $k$), the partition function is the $d$-th power of the rank-$1$ answer:
\begin{equation}% label removed: eq:thqg-I-rank-d
Z_g(\mathcal{H}_k^{\otimes d};\,\Omega) \;=\; Z_g(\mathcal{H}_k;\,\Omega)^d \;=\; \bigl(\det\operatorname{Im}\Omega\bigr)^{-dk/2} \cdot \bigl(\det{}'_\zeta\Delta_{\Sigma_g}\bigr)^{-dk/2}.
\end{equation}
The modular characteristic is $\kappa(\mathcal{H}_k^{\otimes d}) = dk$, and the shadow free energy is $F_g(\mathcal{H}_k^{\otimes d}) = dk \cdot \lambda_g^{\mathrm{FP}}$.  This is consistent with the independent sum factorization (Proposition~\ref{prop:independent-sum-factorization}): for $\cA = \cA_1 \oplus \cA_2$ with vanishing mixed OPE, $\kappa(\cA) = \kappa(\cA_1) + \kappa(\cA_2)$.
\end{remark}


% ======================================================================
%  SUBSECTION 5: FINITENESS FOR THE STANDARD LANDSCAPE
% ======================================================================

\subsection{Finiteness for the standard landscape}
% label removed: subsec:thqg-I-standard-landscape
\index{standard landscape!finiteness|textbf}

\subsubsection{Verification of hypotheses (H1)--(H2)}

The HS-sewing criterion (Theorem~\ref{thm:general-hs-sewing}) requires two conditions:
\begin{enumerate}[label=\textbf{(H\arabic*)}]
\item % label removed: item:thqg-I-H1 \emph{Subexponential sector growth:} $\log \dim H_n = o(n)$.
\item % label removed: item:thqg-I-H2 \emph{Polynomial OPE growth:} $|C^{c,k}_{a,i;\,b,j}| \leq K(a+b+c+1)^N$.
\end{enumerate}

\begin{proposition}[Standard families satisfy HS-sewing; \ClaimStatusProvedElsewhere{} {\textup{Theorem~\ref{thm:general-hs-sewing}}}]
% label removed: prop:thqg-I-standard-hs
\index{HS-sewing!standard families}
Every algebra in the following list satisfies \ref{item:thqg-I-H1} and \ref{item:thqg-I-H2}:
\begin{enumerate}[label=\textup{(\alph*)}]
\item Heisenberg $\mathcal{H}_k$ (any rank, any level $k \neq 0$).
\item Affine Kac--Moody $\widehat{\fg}_k$ ($\fg$ simple, $k \neq -h^\vee$).
\item Virasoro $\mathrm{Vir}_c$ (any $c \in \mathbb{C}$).
\item $\beta\gamma$ system (any central charge).
\item Principal $\mathcal{W}$-algebras $\mathcal{W}_k(\fg)$ ($\fg$ simple, generic $k$).
\item Lattice VOAs $V_\Lambda$ ($\Lambda$ even integral).
\item Free fermions $\psi_1, \ldots, \psi_n$ (any rank).
\end{enumerate}
\end{proposition}

\begin{proof}
For each family, we verify \ref{item:thqg-I-H1} and \ref{item:thqg-I-H2} by direct computation.  The key estimates were computed in Computations~\ref{comp:thqg-I-hs-heisenberg}--\ref{comp:thqg-I-hs-fermion}.

\emph{(a) Heisenberg.}  $\dim H_n = p(n)$, $\log p(n) \sim \pi\sqrt{2n/3} = o(n)$.  OPE degree $N = 1$.

\emph{(b) Affine Kac--Moody.}  $\dim H_n$ has subexponential growth (Kac--Peterson~\cite{KP84}).  OPE degree $N = 1$.

\emph{(c) Virasoro.}  $\dim H_n = p(n)$, $\log p(n) = o(n)$.  OPE degree $N = 2$.

\emph{(d) $\beta\gamma$.}  $\dim H_n = p(n)^2$, $\log \dim H_n = o(n)$.  OPE degree $N = 1$.

\emph{(e) $\mathcal{W}_N$.}  Sector growth subexponential with exponent $O(n^{(N-1)/N})$.  OPE degree $N_{\max} = 2N$.

\emph{(f) Lattice VOA.}  $\dim H_n = |\Lambda^*/\Lambda|\cdot p_r(n)$; $\log p_r(n) = O(\sqrt{n})$.  OPE degree $\leq r$.

\emph{(g) Free fermions.}  $\dim H_n$ grows at most as $2^n$ (number of subsets of $n$ fermion modes), but the conformal-weight grading gives polynomial growth.  OPE degree $N = 0$.
\end{proof}

\subsubsection{Detailed verification for each family}

\begin{proposition}[Heisenberg: detailed HS verification; \ClaimStatusProvedHere]
% label removed: prop:thqg-I-heisenberg-detail
\index{Heisenberg!HS verification}
For the rank-$1$ Heisenberg $\mathcal{H}_k$ at level $k \neq 0$:
\begin{enumerate}[label=\textup{(\roman*)}]
\item The OPE structure constants are $C^{n-1,\lambda}_{n,\mu;\,0,\mathrm{vac}} = k \cdot \delta_{|\lambda|-|\mu|,1}$ for the fundamental mode $a_{-1}|0\rangle$ acting on a partition state $|\mu\rangle$ of weight $n$ to produce a partition state $|\lambda\rangle$ of weight $n-1$.  The general OPE coefficient at weight level $(a,b,c)$ with $a+b-c = m$ (the number of contractions) satisfies $|C^{c,\lambda}_{a,\mu;\,b,\nu}| \leq |k|^m \cdot m! \leq |k|^{a+b} \cdot (a+b)!$, which is polynomial for fixed $m$.
\item The HS norm at weight level $n$ satisfies the exact formula
  \begin{equation}% label removed: eq:thqg-I-heisenberg-exact-hs
  \sum_{\substack{a+b+c = n}} \|m^c_{a,b}\|_{\HS}^2 \;=\; |k|^2 \cdot \sum_{\substack{a+b+c=n}} p(a)p(b)p(c) \cdot \#\{\text{Wick contractions}\},
  \end{equation}
  where the number of Wick contractions at level $n$ grows polynomially.
\end{enumerate}
\end{proposition}

\begin{proof}
The Heisenberg OPE is computed by Wick's theorem: each pair of mode operators $a_m, a_n$ contracts to $k \cdot m \cdot \delta_{m+n,0}$.  The number of contractions between a state $|\mu\rangle$ of weight $a$ and a state $|\nu\rangle$ of weight $b$ is at most $\min(a,b)$, and each contraction produces a factor of $|k| \cdot n$ for a mode of level $n$.  The polynomial bound follows from the multinomial structure of the Wick expansion.
\end{proof}

\begin{proposition}[Affine Kac--Moody: detailed HS verification; \ClaimStatusProvedHere]
% label removed: prop:thqg-I-affine-detail
\index{affine Kac--Moody!HS verification}
For $\widehat{\fg}_k$ with $\fg$ simple of dimension $d$, rank $r$, and dual Coxeter number $h^\vee$, at non-critical level $k \neq -h^\vee$:
\begin{enumerate}[label=\textup{(\roman*)}]
\item The Kac--Peterson character formula gives $\chi_k(\tau) = \sum_{n \geq 0} (\dim H_n) q^n = q^{-c_{\text{eff}}/24} \cdot \Theta_{k+h^\vee,\fg}(\tau) / \eta(\tau)^d$, where $\Theta_{k+h^\vee,\fg}$ is the affine theta function.  The dimension $\dim H_n$ is dominated by $\eta(\tau)^{-d}$, giving $\log \dim H_n \sim \pi\sqrt{2dn/3}$.
\item The OPE of currents $J^a_m J^b_n = k m \delta^{ab}\delta_{m+n,0} + f^{ab}{}_c J^c_{m+n}$ has structure constants bounded by $\max(|k|, \|f\|_\infty) \cdot n$ for mode number $n$, which is polynomial of degree $N = 1$.
\item The HS sewing sum satisfies
  \begin{equation}% label removed: eq:thqg-I-affine-hs-refined
  \sum_{a,b,c \geq 0} q^{a+b+c} \|m^c_{a,b}\|_{\HS}^2 \;\leq\; (|k| + \|f\|_\infty)^2 \cdot \eta(q)^{-2d} \cdot P(q),
  \end{equation}
  where $P(q) = \sum_{n \geq 0} n^2 (\dim H_n)^2 q^n$ is a convergent power series for $q \in (0,1)$.
\end{enumerate}
\end{proposition}

\begin{proof}
(i) is the Kac--Peterson asymptotic~\cite{KP84}.  (ii) follows from the affine Lie algebra commutation relations.  (iii) combines the estimates from (i) and (ii) with the general HS bound~\eqref{eq:thqg-I-hs-criterion-bound}.
\end{proof}

\begin{proposition}[Virasoro: detailed HS verification; \ClaimStatusProvedHere]
% label removed: prop:thqg-I-virasoro-detail
\index{Virasoro!HS verification}
For $\mathrm{Vir}_c$ at central charge $c$:
\begin{enumerate}[label=\textup{(\roman*)}]
\item The Verma module at weight $h$ has $\dim H_n = p(n)$ (partitions of $n$, corresponding to states $L_{-n_1}L_{-n_2}\cdots L_{-n_k}|h\rangle$ with $n_1 + \cdots + n_k = n$).
\item The OPE $T(z)T(w) \sim c/(2(z-w)^4) + 2T(w)/(z-w)^2 + \partial T(w)/(z-w)$ has a quartic pole.  The resulting structure constants satisfy $|C^c_{a,b}| \leq (|c|/2 + 2) \cdot (a+b+1)^2$, which is polynomial of degree $N = 2$.
\item The refined HS bound:
  \begin{equation}% label removed: eq:thqg-I-virasoro-refined-hs
  \sum_{a,b,c \geq 0} q^{a+b+c} \|m^c_{a,b}\|_{\HS}^2 \;\leq\; (|c|/2 + 2)^2 \sum_{n \geq 0} n^4 p(n)^3 q^n.
  \end{equation}
  The series $\sum n^4 p(n)^3 q^n$ converges for all $q \in (0,1)$ because $p(n)^3 \leq e^{3\pi\sqrt{2n/3} + O(\log n)}$ grows subexponentially while $q^n$ decays exponentially.  Concretely, for $q = 0.5$: the $n = 100$ term is $100^4 \cdot p(100)^3 \cdot 0.5^{100} \approx 10^8 \cdot (1.9 \times 10^8)^3 \cdot 10^{-30} \approx 10^{2}$, and the series converges.
\end{enumerate}
\end{proposition}

\begin{proof}
(i) is standard Virasoro representation theory.  (ii) follows from the mode algebra $[L_m, L_n] = (m-n)L_{m+n} + \frac{c}{12}m(m^2-1)\delta_{m+n,0}$, which gives $|C| \leq |c|/12 \cdot m^3 + |m-n| \cdot |L_{m+n}|$; the polynomial bound is obtained by tracking the maximal power of conformal weight in iterated commutators.  (iii) follows from substitution into~\eqref{eq:thqg-I-hs-criterion-bound}.
\end{proof}

\subsubsection{Convergence radii and analytic continuation}

\begin{proposition}[Convergence radii of the scalar partition function; \ClaimStatusProvedHere]
% label removed: prop:thqg-I-convergence-radii
\index{convergence radius!scalar partition function}
The scalar gravitational partition function $Z^{\mathrm{scal}}_{\mathrm{grav}}(\cA;\,\hbar) = \kappa \cdot \frac{\sqrt{\hbar}/2}{\sin(\sqrt{\hbar}/2)}$ has:
\begin{enumerate}[label=\textup{(\roman*)}]
\item \emph{Radius of convergence in $\hbar$:} $R = 4\pi^2 \approx 39.48$.
\item \emph{First singularity:} A pole of order $1$ at $\hbar = 4\pi^2$, corresponding to the first zero of $\sin(\sqrt{\hbar}/2)$ at $\sqrt{\hbar} = 2\pi$.
\item \emph{Analytic continuation:} $Z^{\mathrm{scal}}$ extends to a meromorphic function on $\mathbb{C}$ with poles at $\hbar = 4\pi^2 n^2$, $n \in \mathbb{Z}_{> 0}$.
\end{enumerate}
\end{proposition}

\begin{proof}
The function $f(\hbar) = \frac{\sqrt{\hbar}/2}{\sin(\sqrt{\hbar}/2)}$ is analytic at $\hbar = 0$ (removable singularity: $f(0) = 1$).  The root test gives $\limsup_{g \to \infty} |a_g|^{1/g} = 1/(2\pi)^2 = 1/(4\pi^2)$, so the radius of convergence is $4\pi^2$.  The poles at $\hbar = 4\pi^2 n^2$ are simple with residues $(-1)^{n+1} \cdot 2n\pi \cdot \kappa$.
\end{proof}

\begin{remark}[Hagedorn singularities]
% label removed: rem:thqg-I-hagedorn
\index{Hagedorn singularity}
The poles at $\hbar = 4\pi^2 n^2$ are the gravitational analogues of the Hagedorn temperature in string theory.  At $\hbar = 4\pi^2$, the genus sum first diverges.  Beyond this point, the partition function requires analytic continuation and the genus expansion breaks down.  In the holographic context, the pole at $\hbar = 4\pi^2$ corresponds to a phase transition in the gravitational theory: the genus expansion becomes unreliable, and non-perturbative effects (saddle-point contributions from higher-topology geometries) must be included.
\end{remark}

\subsubsection{Admissible levels and singular fibers}

\begin{remark}[Admissible levels]
% label removed: rem:thqg-I-admissible
\index{admissible level!finiteness}
At admissible levels $k = -h^\vee + p/q$ (with $p,q$ coprime positive integers and $p \geq h^\vee$), the universal vertex algebra $V_k(\fg)$ has null vectors and the simple quotient $L_k(\fg)$ differs from $V_k(\fg)$.  The HS-sewing criterion applies to $V_k(\fg)$ unconditionally (by Proposition~\ref{prop:thqg-I-standard-hs}).  For the simple quotient $L_k(\fg)$, the growth remains subexponential by the Kac--Wakimoto character formula~\cite{KW88}, so HS-sewing holds as well.

The modular characteristic is $\kappa(V_k(\fg)) = td/(2h^\vee)$ with $t = k + h^\vee$, which vanishes at critical level $k = -h^\vee$.  At critical level, the shadow free energies all vanish: $F_g = 0$ for all $g \geq 1$.  The Feigin--Frenkel center emerges as the replacement for the genus tower.
\end{remark}

\begin{remark}[Singular fibers]
% label removed: rem:thqg-I-singular-fibers
\index{singular fiber!modular characteristic}
As the level $k$ varies, the modular characteristic $\kappa(k) = td/(2h^\vee)$ with $t = k + h^\vee$ traces a linear function of $k$.  The critical level $k = -h^\vee$ is the unique singular fiber where $\kappa$ vanishes.  Near the singular fiber, $\kappa \to 0$ and the shadow free energies $F_g \to 0$ uniformly: $|F_g(k)| \leq C |k + h^\vee| / (2\pi)^{2g}$ for all $g$.  The shadow tower collapses to the trivial element $\Theta_\cA = 0$ at $k = -h^\vee$.
\end{remark}

\begin{remark}[Critical level and the Feigin--Frenkel center]
% label removed: rem:thqg-I-ff-center
\index{Feigin--Frenkel center}
At critical level $k = -h^\vee$, the Sugawara construction is \emph{undefined} (not ``$c$ diverges''; the Sugawara tensor $T(z) = \frac{1}{2(k+h^\vee)} \sum_a :J^a(z) J^a(z):$ has a pole at $k = -h^\vee$).  The Feigin--Frenkel theorem identifies the center $Z(V_{-h^\vee}(\fg)) \cong \mathrm{Fun}\,\mathrm{Op}_{\fg^\vee}(\mathbb{D})$ with the algebra of functions on the space of $\fg^\vee$-opers on the formal disc.  The shadow free energy vanishes at this level because $\kappa = 0$, but the center provides a replacement: the ``classical'' genus-$g$ invariant is the space of $\fg^\vee$-opers on $\Sigma_g$, which is finite-dimensional and does not require the HS-sewing machinery.
\end{remark}

\subsubsection{Summary table}

\begin{table}[ht]
\centering
\caption{Perturbative finiteness: standard landscape}
% label removed: tab:thqg-I-finiteness-summary
\index{perturbative finiteness!summary table}
\renewcommand{\arraystretch}{1.4}
{\footnotesize
\begin{tabular}{|l|c|c|c|c|c|}
\hline
\textbf{Family $\cA$} & $\boldsymbol{\dim H_n}$ & \textbf{OPE $N$} & $\boldsymbol{\kappa}$ & $\boldsymbol{r_{\max}}$ & \textbf{HS Status} \\
\hline
$\mathcal{H}_k$ & $p(n)$ & $1$ & $k$ & $2$ & \textbf{Proved} \\
\hline
$\widehat{\fg}_k$ & $\sim e^{\pi\sqrt{2dn/3}}$ & $1$ & $\frac{(k+h^\vee)d}{2h^\vee}$ & $3$ & \textbf{Proved} \\
\hline
$\mathrm{Vir}_c$ & $p(n)$ & $2$ & $c/2$ & $\infty$ & \textbf{Proved} \\
\hline
$\beta\gamma$ & $p(n)^2$ & $1$ & $1$ & $4$ & \textbf{Proved} \\
\hline
$\mathcal{W}_N$ & $N$-colored & $2N$ & $c_N/2 \cdot \rho_N$ & $\infty$ & \textbf{Proved} \\
\hline
$V_\Lambda$ & $|\Lambda^*/\Lambda|\cdot p_r(n)$ & $r$ & $\rank(\Lambda)$ & $2$ & \textbf{Proved} \\
\hline
Free fermions & polynomial & $0$ & $n/4$ & $2$ & \textbf{Proved} \\
\hline
\end{tabular}}
\end{table}

\begin{remark}[Exhaustiveness]
% label removed: rem:thqg-I-exhaustiveness
\index{standard landscape!exhaustiveness}
The seven families in Table~\ref{tab:thqg-I-finiteness-summary} exhaust the ``standard landscape'' of finitely strongly generated chiral algebras with simple Lie symmetry.  Finiteness extends to:
\begin{enumerate}[label=\textup{(\alph*)}]
\item \emph{Tensor products:} $\cA \otimes \cB$ satisfies HS-sewing if both $\cA$ and $\cB$ do (the HS norm of the tensor product is the product of HS norms).
\item \emph{Cosets:} $\cA / \cB$ (the coset of $\cB \subset \cA$) satisfies HS-sewing if $\cA$ does (the sector dimensions can only decrease).
\item \emph{Orbifolds:} $\cA^G$ (orbifold by a finite group $G$) satisfies HS-sewing if $\cA$ does (the twisted sectors contribute additional terms, but all with subexponential growth).
\item \emph{Non-principal $\mathcal{W}$-algebras:} $\mathcal{W}_k(\fg, f)$ for nilpotent $f$ satisfies HS-sewing at generic $k$ (the DS reduction does not increase OPE degree).
\end{enumerate}
\end{remark}


% ======================================================================
%  SUBSECTION 6: IMPLICATIONS FOR 3D QUANTUM GRAVITY
% ======================================================================

\subsection{Implications for 3d quantum gravity}
% label removed: subsec:thqg-I-3d-gravity
\index{3d quantum gravity|textbf}
\index{BTZ black hole!partition function}

Three-dimensional gravity with negative cosmological constant $\Lambda = -1/\ell^2$ is the physical arena where the perturbative finiteness theorem has the most direct consequences.  The boundary theory is a 2d CFT with central charge determined by the Brown--Henneaux formula $c = 3\ell/(2G)$.

\subsubsection{The Brown--Henneaux identification}

\begin{setup}[Brown--Henneaux central charge]
% label removed: setup:thqg-I-brown-henneaux
\index{Brown--Henneaux!central charge}
The asymptotic symmetry algebra of AdS$_3$ gravity at spatial infinity is two copies of the Virasoro algebra with central charge
\begin{equation}% label removed: eq:thqg-I-brown-henneaux
c \;=\; \frac{3\ell}{2G}\,,
\end{equation}
where $\ell$ is the AdS radius and $G$ is the three-dimensional Newton constant~\cite{BH86}.  In the twisted holographic setting, $\cA = \mathrm{Vir}_c$ is the boundary chiral algebra, with modular characteristic $\kappa(\cA) = c/2 = 3\ell/(4G)$.
\end{setup}

\begin{remark}[Semiclassical limit]
% label removed: rem:thqg-I-semiclassical
\index{semiclassical limit}
The semiclassical limit is $G \to 0$ with $\ell$ fixed, which gives $c \to \infty$.  In this limit, $\kappa \to \infty$ and the shadow free energies $F_g = \kappa \cdot \lambda_g^{\mathrm{FP}} \to \infty$ for each $g \geq 1$: the genus corrections are large in absolute terms but small \emph{relative} to the tree-level contribution $F_0 = \kappa$.  The ratio $F_g/F_0 = \lambda_g^{\mathrm{FP}}$ is independent of $\kappa$ and decreases super-exponentially in $g$.  The genus expansion is therefore an asymptotic expansion in $1/\kappa \sim G/\ell$ that happens to converge (as a formal power series in $\hbar$).
\end{remark}

\subsubsection{The BTZ partition function from shadow free energies}

\begin{theorem}[BTZ partition function from shadow theory; \ClaimStatusProvedHere]
% label removed: thm:thqg-I-btz
\index{BTZ black hole!shadow free energy}
The one-loop correction to the BTZ partition function in the twisted holographic framework is
\begin{equation}% label removed: eq:thqg-I-btz-F1
F_1^{\mathrm{BTZ}} \;=\; F_1(\mathrm{Vir}_c) \;=\; \frac{c}{48} \;=\; \frac{\ell}{32G}\,.
\end{equation}
The full perturbative BTZ partition function, including all genus corrections, is
\begin{equation}% label removed: eq:thqg-I-btz-full
Z_{\mathrm{BTZ}}^{\mathrm{pert}}(\hbar) \;=\; \frac{c}{2} \cdot \frac{\sqrt{\hbar}/2}{\sin(\sqrt{\hbar}/2)} \;=\; \frac{3\ell}{4G} \cdot \frac{\sqrt{\hbar}/2}{\sin(\sqrt{\hbar}/2)}\,.
\end{equation}
\end{theorem}

\begin{proof}
By the Brown--Henneaux identification (Setup~\ref{setup:thqg-I-brown-henneaux}), the boundary algebra is $\mathrm{Vir}_c$ with $\kappa = c/2$.  The genus-$1$ free energy is $F_1 = \kappa/24 = c/48$.  The full perturbative series is the generating function~\eqref{eq:thqg-I-gf-ahat} with $\kappa = c/2$.  Substituting $c = 3\ell/(2G)$ gives $\kappa = 3\ell/(4G)$.
\end{proof}

\begin{remark}[Comparison with known results]
% label removed: rem:thqg-I-comparison-btz
\index{BTZ black hole!comparison}
The genus-$1$ correction $F_1 = c/48$ has been computed by multiple methods: (a)~gravity path integral (Giombi--Maloney--Yin~\cite{GMY08}): the one-loop determinant of the graviton on the solid torus gives $\prod_{n=2}^{\infty} (1 - q^n)^{-1}$; (b)~Chern--Simons theory (Witten~\cite{Witten88}): 3d gravity as $\mathrm{SL}(2,\mathbb{R}) \times \mathrm{SL}(2,\mathbb{R})$ CS theory, with $c/24$ for the full theory and $c/48$ for the holomorphic half; (c)~modular bootstrap (Maloney--Witten~\cite{MW10}).  All agree with the shadow free energy computation.  The advantage of the shadow method is that it extends automatically to all genera via the closed form~\eqref{eq:thqg-I-btz-full}.
\end{remark}

\subsubsection{Higher-genus corrections}

\begin{proposition}[Higher-genus BTZ corrections; \ClaimStatusProvedHere]
% label removed: prop:thqg-I-btz-higher-genus
\index{BTZ black hole!higher-genus corrections}
The genus-$g$ correction to the BTZ partition function in the scalar channel is
\begin{equation}% label removed: eq:thqg-I-btz-Fg
F_g^{\mathrm{BTZ}} \;=\; \frac{c}{2} \cdot \frac{2^{2g-1}-1}{2^{2g-1}} \cdot \frac{|B_{2g}|}{(2g)!} \;=\; \frac{3\ell}{4G} \cdot \lambda_g^{\mathrm{FP}}\,.
\end{equation}
The first five corrections are:
\begin{align}
F_1^{\mathrm{BTZ}} &= \frac{c}{48}\,, % label removed: eq:thqg-I-btz-F1-val \\
F_2^{\mathrm{BTZ}} &= \frac{7c}{11520}\,, % label removed: eq:thqg-I-btz-F2-val \\
F_3^{\mathrm{BTZ}} &= \frac{31c}{1935360}\,, % label removed: eq:thqg-I-btz-F3-val \\
F_4^{\mathrm{BTZ}} &= \frac{127c}{309657600}\,, % label removed: eq:thqg-I-btz-F4-val \\
F_5^{\mathrm{BTZ}} &= \frac{73c}{7007109120}\,. % label removed: eq:thqg-I-btz-F5-val
\end{align}
The ratio $F_{g+1}/F_g \to 1/(4\pi^2)$ as $g \to \infty$, consistent with the Hagedorn singularity at $\hbar = 4\pi^2$.
\end{proposition}

\begin{proof}
Direct from $\kappa(\mathrm{Vir}_c) = c/2$ and the universal formula (Proposition~\ref{prop:thqg-I-Fg-closed-form}).
\end{proof}

\begin{remark}[Numerical values for $c = 24$ (monster CFT)]
% label removed: rem:thqg-I-monster
\index{monster CFT}
For the monster CFT with $c = 24$ (the Frenkel--Lepowsky--Meurman moonshine module $V^\natural$), $\kappa = 12$:
\begin{align*}
F_1 &= 12/24 = 1/2, \\
F_2 &= 7 \cdot 12/5760 = 84/5760 = 7/480, \\
F_3 &= 31 \cdot 12/967680 = 372/967680 = 31/80640, \\
F_4 &= 127 \cdot 12/154828800 = 1524/154828800 = 127/12902400.
\end{align*}
The one-loop contribution $F_1 = 1/2$ is order $1$, and each subsequent genus is suppressed by the Bernoulli factor.  The total perturbative partition function at $\hbar = 1$ is $Z = 12 \cdot 1/\sin(1/2) \approx 12 \cdot 2.086 \approx 25.0$.
\end{remark}

\subsubsection{The quartic contact invariant in sub-subleading corrections}

\begin{remark}[Quartic contact invariant in 3d gravity]
% label removed: rem:thqg-I-quartic-3d
\index{quartic contact invariant!3d gravity}
The scalar channel captures only the arity-$0$ information.  The first non-scalar correction comes from the quartic contact invariant
\begin{equation}% label removed: eq:thqg-I-quartic-contact
Q^{\mathrm{contact}}_{\mathrm{Vir}} \;=\; \frac{10}{c(5c+22)}\,,
\end{equation}
which is the arity-$4$, genus-$0$ component of the shadow tower.  For $c = 3\ell/(2G)$, this gives $Q^{\mathrm{contact}} = 10/(c(5c+22)) = 2/c^2 + O(1/c^3)$, a two-loop correction in Newton's constant.  The genus-$1$ Hessian correction is
\begin{equation}% label removed: eq:thqg-I-genus1-hessian
\delta_H^{(1)}_{\mathrm{Vir}} \;=\; \frac{120}{c^2(5c+22)}\,x^2\,.
\end{equation}
These corrections are specific to the mixed shadow class ($r_{\max} = \infty$) of the Virasoro algebra.
\end{remark}

\begin{remark}[Shadow depth and gravitational complexity]
% label removed: rem:thqg-I-depth-complexity
\index{shadow depth!gravitational complexity}
The shadow depth classification ($G/L/C/M$ for Gaussian/Lie/Contact/Mixed) stratifies 3d gravity theories by their nonlinear complexity:
\begin{enumerate}[label=\textup{(\alph*)}]
\item \emph{Gaussian ($r_{\max} = 2$):} Heisenberg, lattice VOAs, free fermions.  The gravitational theory is free.  The partition function is an exact Fredholm determinant.  No nonlinear corrections.
\item \emph{Lie-type ($r_{\max} = 3$):} Affine Kac--Moody algebras.  Cubic interactions (tree-level scattering).  The partition function receives corrections from trivalent graphs.
\item \emph{Contact-type ($r_{\max} = 4$):} $\beta\gamma$ system.  Quartic interactions but no higher.  The shadow tower terminates after two nonlinear steps.
\item \emph{Mixed ($r_{\max} = \infty$):} Virasoro, $\mathcal{W}_N$.  Interactions at all orders.  The full shadow tower is infinite, and the perturbative finiteness theorem is needed in full strength.  The quintic obstruction $o^{(5)}_{\mathrm{Vir}} \neq 0$ forces the tower to be genuinely infinite.
\end{enumerate}
The Brown--Henneaux boundary theory for pure 3d gravity is $\mathrm{Vir}_c$, which falls in the mixed class.
\end{remark}

\subsubsection{Large-$c$ expansion and gravitational coupling}

\begin{proposition}[Leading gravitational coupling expansion; \ClaimStatusProvedHere]
% label removed: prop:thqg-I-1c-expansion
\index{gravitational coupling!large-$c$ expansion}
In the semiclassical expansion $\hbar = 1/c$, the perturbative gravitational partition function has the asymptotic form
\begin{equation}% label removed: eq:thqg-I-1c-expansion
\log Z_{\mathrm{grav}}^{\mathrm{scal}}(\mathrm{Vir}_c;\,1/c) \;=\; \frac{1}{48} + \frac{7}{11520\,c} + \frac{31}{1935360\,c^2} + \frac{127}{309657600\,c^3} + O(c^{-4})\,.
\end{equation}
The leading term $1/48$ is the one-loop correction, and each subsequent term is a higher-loop gravitational correction suppressed by one additional power of $G/\ell \sim 1/c$.
\end{proposition}

\begin{proof}
From $F_g(\mathrm{Vir}_c) = (c/2)\lambda_g^{\mathrm{FP}}$, the contribution of genus~$g$ to $\log Z$ at $\hbar = 1/c$ is $\lambda_g^{\mathrm{FP}}/(2c^{g-1})$.  The expansion $\log Z = \sum_{g \geq 1} F_g \hbar^g = \sum_{g \geq 1} (c/2)\lambda_g^{\mathrm{FP}}/c^g = \sum_{g \geq 1} \lambda_g^{\mathrm{FP}} / (2c^{g-1})$.  At $g = 1$: $\lambda_1^{\mathrm{FP}} = 1/24$, giving $1/(2 \cdot 24) = 1/48$.  At $g = 2$: $7/(2 \cdot 5760 \cdot c) = 7/(11520c)$.  And so on.
\end{proof}

\begin{proposition}[Newton's constant expansion; \ClaimStatusProvedHere]
% label removed: prop:thqg-I-newton-expansion
\index{Newton's constant!expansion}
Substituting $c = 3\ell/(2G)$ and $\hbar = 2G/(3\ell)$, the gravitational partition function becomes
\begin{equation}% label removed: eq:thqg-I-newton-expansion
Z_{\mathrm{grav}}^{\mathrm{scal}}(\mathrm{Vir}_c;\,2G/(3\ell)) \;=\; \frac{3\ell}{4G} \cdot \frac{\sqrt{2G/(3\ell)}/2}{\sin(\sqrt{2G/(3\ell)}/2)}\,.
\end{equation}
For $G/\ell \ll 1$ (weak gravitational coupling), the expansion in $G/\ell$ gives
\begin{equation}% label removed: eq:thqg-I-G-expansion
Z \;=\; \frac{3\ell}{4G}\left(1 + \frac{G}{72\ell} + \frac{7G^2}{51840\ell^2} + O(G^3/\ell^3)\right).
\end{equation}
\end{proposition}

\begin{proof}
Direct Taylor expansion of $x/(2\sin(x/2))$ at $x = 0$ with $x^2 = 2G/(3\ell)$.
\end{proof}

\subsubsection{Towards the non-perturbative regime}

\begin{conjecture}[Non-perturbative completion; \ClaimStatusConjectured]
% label removed: conj:thqg-I-non-perturbative
\index{non-perturbative!gravitational partition function}
The meromorphic function $Z^{\mathrm{scal}}_{\mathrm{grav}}(\hbar) = \kappa \cdot \frac{\sqrt{\hbar}/2}{\sin(\sqrt{\hbar}/2)}$ admits a non-perturbative completion satisfying:
\begin{enumerate}[label=\textup{(\roman*)}]
\item Resurgent structure: the Borel transform has singularities at $\hbar = 4\pi^2 n^2$.
\item Stokes phenomena: crossing Stokes rays changes the partition function by $\sim e^{-4\pi^2 n^2/\hbar}$, identified with non-handlebody saddle points.
\item Positivity: for real $\hbar > 0$, the full non-perturbative partition function $Z > 0$.
\end{enumerate}
\end{conjecture}

\begin{evidence}[Evidence for the resurgent structure]
% label removed: evid:thqg-I-resurgent
The Taylor coefficients $F_g \sim 2\kappa/(2\pi)^{2g}$ have the form $a_g \sim C \cdot A^{-g}$ with $A = (2\pi)^2$, which is precisely the asymptotic form associated with a simple pole in the Borel plane at $\xi = 4\pi^2$.  The Borel transform is
\begin{equation}% label removed: eq:thqg-I-borel-transform
\hat{Z}(\xi) \;=\; \sum_{g=1}^\infty \frac{F_g}{\Gamma(g)} \xi^{g-1} \;=\; \kappa \sum_{g=1}^\infty \frac{\lambda_g^{\mathrm{FP}}}{(g-1)!} \xi^{g-1}.
\end{equation}
Since $\lambda_g^{\mathrm{FP}} \sim 2/(2\pi)^{2g}$, the ratio $\lambda_g^{\mathrm{FP}}/\Gamma(g) \sim 2/((2\pi)^{2g}(g-1)!)$, and the series converges for all $\xi \in \mathbb{C}$ (the Borel transform is entire).  The singularities at $\xi = 4\pi^2 n^2$ in the Borel plane arise from the Laplace integral relating the original series to the Borel sum: the non-perturbative corrections $\sim e^{-4\pi^2 n^2/\hbar}$ match the mass gap of the lightest BTZ black hole above the vacuum.
\end{evidence}

\begin{remark}[BTZ black holes as non-perturbative saddles]
% label removed: rem:thqg-I-btz-saddles
\index{BTZ black hole!non-perturbative}
In 3d gravity, the BTZ black holes of mass $M = n^2/(8G\ell)$ ($n \geq 1$) are the non-perturbative saddle points.  Their Euclidean action is $I_n = 4\pi^2 n^2 / \hbar$ (with $\hbar = 2G/(3\ell)$).  The Stokes phenomenon at the $n$-th pole $\hbar = 4\pi^2 n^2$ corresponds to the creation/annihilation of the $n$-th BTZ black hole as one crosses the Stokes ray.  The full non-perturbative partition function should be $Z = Z_{\mathrm{pert}} + \sum_{n \geq 1} c_n e^{-I_n/\hbar}$ with explicit coefficients $c_n$ determined by the resurgent structure.
\end{remark}

%% ---- Stokes phenomena on the Steinberg object ----

\providecommand{\SteinbergB}{\mathfrak{S}_b}

\subsubsection{Stokes phenomena on the Steinberg object}
% label removed: sssec:thqg-stokes-steinberg
\index{Stokes phenomena!Steinberg object}
\index{Steinberg object!Stokes filtration}

The poles of $Z_{\mathrm{grav}}(\hbar) = \kappa \cdot \frac{\sqrt{\hbar}/2}{\sin(\sqrt{\hbar}/2)}$ at $\hbar = 4\pi^2 n^2$ are not merely analytic accidents: they are the Stokes lines in the Borel plane of the convolution integral on the Steinberg object~$\SteinbergB$.  The Borel transform
\[
B[Z](\zeta) \;=\; \sum_{g \geq 1} \frac{F_g}{\Gamma(g)}\,\zeta^{g-1}
\]
has singularities at $\zeta_n = 4\pi^2 n^2$, which are the Euclidean actions of the $n$-th BTZ saddle points (Remark~\ref{rem:thqg-I-btz-saddles}).  The Stokes-filtered Steinberg construction upgrades the perturbative Steinberg presentation to one that sees all saddle-point sectors simultaneously.

\begin{construction}[The Stokes-filtered Steinberg]
% label removed: constr:thqg-stokes-filtered-steinberg
\index{Steinberg object!Stokes-filtered}
Define the \emph{Stokes-filtered Steinberg} $\SteinbergB^{\mathrm{Stokes}}$ as follows.  For each $n \geq 1$, let $\SteinbergB^{\mathrm{BTZ},n}$ denote the Steinberg object associated to the $n$-th BTZ saddle, i.e.\ the analogue of~$\SteinbergB$ constructed from the gauge equivalence data of the $n$-th non-handlebody saddle point with Euclidean action $I_n = 4\pi^2 n^2 / \hbar$.  The perturbative Steinberg~$\SteinbergB$ is $\SteinbergB^{\mathrm{BTZ},0}$ (the vacuum saddle).  The Stokes-filtered Steinberg is the datum
\[
\SteinbergB^{\mathrm{Stokes}} \;=\; \bigl(\SteinbergB,\;\{\SteinbergB^{\mathrm{BTZ},n}\}_{n \geq 1},\;\{S_n\}_{n \geq 1}\bigr)
\]
where $S_n \colon H_*^{\mathrm{BM}}(\SteinbergB) \to H_*^{\mathrm{BM}}(\SteinbergB^{\mathrm{BTZ},n})$ is the Stokes map induced by analytic continuation of the Borel-resummed integral across the $n$-th Stokes ray $\arg(\zeta) = \arg(\zeta_n)$.
\end{construction}

\begin{proposition}[Resurgent structure from Stokes-filtered BM homology; \ClaimStatusConjectured]
% label removed: prop:thqg-stokes-bm-resurgence
\index{resurgence!Steinberg BM homology}
\index{BM homology!Stokes-filtered}
The resurgent structure of $Z_{\mathrm{grav}}$ is presented by the Stokes-filtered Borel--Moore homology of the Steinberg:
\begin{equation}% label removed: eq:thqg-stokes-bm-decomposition
H_*^{\mathrm{BM},\,\mathrm{Stokes}}(\SteinbergB) \;=\; H_*^{\mathrm{BM}}(\SteinbergB) \;\oplus\; \bigoplus_{n \geq 1} H_*^{\mathrm{BM}}(\SteinbergB^{\mathrm{BTZ},n})\,,
\end{equation}
where $\SteinbergB^{\mathrm{BTZ},n}$ is the Steinberg object for the $n$-th BTZ saddle.  The Stokes automorphism
\[
\mathfrak{S}_n \colon H_*^{\mathrm{BM}}(\SteinbergB) \;\longrightarrow\; H_*^{\mathrm{BM}}(\SteinbergB)
\]
is the wall-crossing automorphism exchanging the perturbative and $n$-th BTZ sectors.  More precisely, the discontinuity of the lateral Borel resummation across the $n$-th Stokes ray is
\[
\mathrm{disc}_n\, Z \;=\; \mathfrak{S}_n([\SteinbergB]) \cdot e^{-\zeta_n / \hbar}\,,
\]
where $[\SteinbergB] \in H_*^{\mathrm{BM}}(\SteinbergB)$ is the fundamental class and $\zeta_n = 4\pi^2 n^2$.
\end{proposition}

\begin{remark}[Positivity from geometry]
% label removed: rem:thqg-stokes-positivity
\index{positivity!gravitational partition function}
\index{BM homology!fundamental class positivity}
The positivity of $Z_{\mathrm{grav}}$ for real $\hbar > 0$ (away from the poles $\hbar = 4\pi^2 n^2$) admits a geometric explanation through the Steinberg presentation.  The perturbative contribution is $\langle [\SteinbergB], [\SteinbergB] \rangle_{\mathrm{BM}}$, the self-pairing of the BM fundamental class $[\SteinbergB]$, which is positive since $\SteinbergB$ carries a canonical orientation as a complex-algebraic variety.  The Stokes phenomena at $\hbar = 4\pi^2 n^2$ are phase transitions where the $n$-th BTZ stratum $\SteinbergB^{\mathrm{BTZ},n}$ becomes visible: as $\hbar$ crosses $4\pi^2 n^2$, the Borel resummation contour passes through the singularity $\zeta_n$, and the new Steinberg stratum contributes an additional sector to the partition function.  The pole itself reflects the divergence of the lateral resummation integral at the singular point: geometrically, the collision of the perturbative fundamental cycle with the $n$-th BTZ cycle in the Stokes-filtered homology.
\end{remark}

\begin{remark}[Pure 3d gravity as a model]
% label removed: rem:thqg-I-pure-3d
\index{3d gravity!pure}
Pure 3d gravity with $\Lambda < 0$ is the simplest non-trivial gravitational theory.  Its boundary algebra is $\mathrm{Vir}_c$ with $c = 3\ell/(2G)$, which falls in the mixed class.  The perturbative finiteness theorem provides a complete, explicit, genus-by-genus calculation of the gravitational partition function:
\[
Z_{\mathrm{grav}}^{\mathrm{scal}}(\hbar) \;=\; \frac{3\ell}{4G} \cdot \frac{\sqrt{\hbar}/2}{\sin(\sqrt{\hbar}/2)}\,.
\]
The non-scalar corrections from $Q^{\mathrm{contact}} = 10/(c(5c+22))$ and the infinite shadow tower encode the genuinely nonlinear aspects of the gravitational theory, suppressed by $1/c^2$ relative to the scalar contribution.
\end{remark}

\begin{remark}[Comparison with matrix models]
% label removed: rem:thqg-I-matrix-models
\index{matrix model!gravitational partition function}
In 2d topological gravity (Witten--Kontsevich), the genus-$g$ free energy involves Bernoulli numbers with a different prefactor: the Witten--Kontsevich generating function is $\sum_g F_g^{\mathrm{WK}} (2g-2)! x^{2g-2}$, with $F_g^{\mathrm{WK}} = \chi(\mathcal{M}_g) / (2g-2)!$ involving the orbifold Euler characteristic.  The relation to the shadow generating function is $\frac{x/2}{\sin(x/2)} = \left.\frac{x/2}{\sinh(x/2)}\right|_{x \mapsto ix}$, reflecting the Wick rotation from Lorentzian to Euclidean signature.  In the matrix model, the $\sinh$ function arises from the Gaussian integral, while in the shadow theory, the $\sin$ function arises from the Hodge integral; the two are related by analytic continuation $x \mapsto ix$.
\end{remark}

\begin{remark}[Higher-spin gravity]
% label removed: rem:thqg-I-higher-spin
\index{higher-spin gravity}
For higher-spin gravity in AdS$_3$ with $\mathcal{W}_N$ symmetry, the boundary algebra is $\mathcal{W}_k(\mathfrak{sl}_N)$ at $c = (N-1)(1 - N(N+1)/(k+N))$.  The modular characteristic is $\kappa(\mathcal{W}_N) = c/2 \cdot \rho_N$ where $\rho_N$ is a normalizing factor.  The perturbative partition function is
\[
Z_{\mathrm{grav}}^{W_N}(\hbar) \;=\; \kappa(\mathcal{W}_N) \cdot \frac{\sqrt{\hbar}/2}{\sin(\sqrt{\hbar}/2)}\,,
\]
which has the same functional form as the Virasoro case but with a different normalization.  The qualitative features (convergence for $|\hbar| < 4\pi^2$, Hagedorn singularity, resurgent structure) are identical.  The difference lies in the non-scalar corrections: $\mathcal{W}_N$ has shadow depth $r_{\max} = \infty$ and the shadow tower includes contributions from all $W^{(s)}$ fields ($s = 2, \ldots, N$).
\end{remark}

\subsubsection{The holographic dictionary}

\begin{remark}[The holographic dictionary]
% label removed: rem:thqg-I-holographic-dictionary
\index{holographic dictionary!perturbative finiteness}
The perturbative finiteness theorem (Theorem~\ref{thm:thqg-I-perturbative-finiteness}) translates into the holographic dictionary:
\begin{center}
\renewcommand{\arraystretch}{1.3}
\begin{tabular}{l|l}
\textbf{Boundary (algebra)} & \textbf{Bulk (gravity)} \\
\hline
$\cA$ (chiral algebra) & boundary conditions \\
$\kappa(\cA)$ (modular characteristic) & cosmological constant \\
$F_g(\cA)$ (shadow free energy) & genus-$g$ path integral \\
HS-sewing condition & UV finiteness \\
$\Theta_\cA$ (MC element) & classical solution + perturbations \\
$D_\cA^2 = 0$ & Ward identity \\
$r_{\max}(\cA)$ (shadow depth) & gravitational complexity \\
$4\pi^2$ (convergence radius) & Hagedorn temperature \\
$Q^{\mathrm{contact}}$ (quartic invariant) & two-loop graviton self-energy \\
$\cA^!$ (Koszul dual) & dual boundary conditions \\
\end{tabular}
\end{center}
The passage from boundary data to bulk partition function is algorithmic: given the OPE data of $\cA$, the modular characteristic $\kappa$ is computable, the Faber--Pandharipande coefficients $\lambda_g^{\mathrm{FP}}$ are universal, and the partition function $Z_{\mathrm{grav}} = \kappa \cdot x/(2\sin(x/2))$ follows.  No UV regularization, no renormalization group flow, and no anomaly matching are required.
\end{remark}

\begin{remark}[Summary of the finiteness mechanism]
% label removed: rem:thqg-I-finiteness-summary
\index{perturbative finiteness!mechanism summary}
The perturbative finiteness of twisted quantum gravity rests on a chain of three implications:
\begin{equation}% label removed: eq:thqg-I-chain
\underbrace{D_\cA^2 = 0}_{\text{bar complex}} \;\Longrightarrow\; \underbrace{\Theta_\cA \in \mathrm{MC}(\gAmod)}_{\text{exact MC}} \;\Longrightarrow\; \underbrace{F_g = \kappa \cdot \lambda_g^{\mathrm{FP}} < \infty}_{\text{finite amplitudes}}.
\end{equation}
The first implication is the bar-intrinsic construction (Theorem~\ref*{V1-thm:mc2-bar-intrinsic}): the element $\Theta_\cA := D_\cA - d_0$ is automatically Maurer--Cartan because the bar differential squares to zero.  The second implication is the projection onto genus~$g$: the shadow free energy is the integral of a cohomology class over a compact moduli space, hence finite.  The entire mechanism is algebraic (no analysis of divergences, no regularization, no renormalization) and exact (no perturbative truncation, no asymptotic expansion).
\end{remark}


% ======================================================================
%  SUBSECTION 7: THE HS-SEWING CONDITION AND THE MC RECURSION
% ======================================================================

\subsection{HS-sewing and the MC recursion}
% label removed: subsec:thqg-I-hs-mc-recursion
\index{HS-sewing!MC recursion}

The HS-sewing condition and the MC recursion interact at two levels: the HS condition provides the \emph{analytic} input (convergence of sewing amplitudes), while the MC recursion provides the \emph{algebraic} input (unique determination of higher-genus data from lower-genus data).  Together, they give a complete genus-by-genus algorithm for computing the gravitational partition function.

\subsubsection{The sewing amplitude as a graph sum}

\begin{proposition}[Graph-sum representation of sewing amplitudes; \ClaimStatusProvedHere]
% label removed: prop:thqg-I-graph-sum-sewing
\index{sewing amplitude!graph-sum representation}
Let $\cA$ be a modular Koszul chiral algebra satisfying HS-sewing.  The genus-$g$ sewing amplitude on a surface $\Sigma_g$ with pants decomposition $\mathcal{P}$ has the graph-sum representation
\begin{equation}% label removed: eq:thqg-I-graph-sewing
Z_{\mathcal{P}}(\cA;\,\mathbf{q}) \;=\; \sum_{\Gamma \in \Gmod(g)} \frac{1}{|\Aut(\Gamma)|}\,\ell_\Gamma(\mathbf{q}),
\end{equation}
where the sum runs over connected modular graphs $\Gamma$ of total genus $g$, and $\ell_\Gamma(\mathbf{q})$ is the graph amplitude: the product of vertex labels (transferred cyclic operations $\ell_n^{\mathrm{tr}}$) and edge propagators ($P_\cA(q_e) = \sum_{n \geq 0} q_e^n \pi_n$, the weighted projection onto the $n$-th conformal weight).
\end{proposition}

\begin{proof}
The pants decomposition $\mathcal{P}$ determines a dual graph $\Gamma_\mathcal{P}$ with one vertex for each pair of pants and one edge for each sewing circle.  The vertex labels are the three-point functions of $\cA$ (the OPE coefficients $C^c_{a,b}$), and the edge labels are the sewing propagators.  The sewing amplitude is the contraction of all vertex and edge tensors:
\[
Z_{\mathcal{P}} \;=\; \sum_{\text{weight assignments}} \prod_{\text{vertices}} C^{c_v}_{a_v,b_v} \prod_{\text{edges}} q_e^{n_e} \delta_{n_e = c_{v'} = a_{v''}},
\]
where the sum over weight assignments runs over all ways to assign conformal weights to the internal edges.  This is precisely the graph amplitude $\ell_{\Gamma_\mathcal{P}}(\mathbf{q})$.

The independence from the pants decomposition (sewing consistency) means that $Z$ depends only on the genus $g$ and the moduli of $\Sigma_g$, not on the choice of $\mathcal{P}$.  Different pants decompositions give different graphs $\Gamma_\mathcal{P}$, but the sum of all amplitudes is the same.
\end{proof}

\begin{proposition}[Convergence of the graph sum at fixed genus; \ClaimStatusProvedHere]
% label removed: prop:thqg-I-graph-sum-convergence
\index{graph sum!convergence}
Under HS-sewing, the graph sum~\eqref{eq:thqg-I-graph-sewing} converges absolutely at each genus $g$.  The number of contributing graphs is finite at each arity level, and the arity sum converges:
\begin{equation}% label removed: eq:thqg-I-graph-sum-arity
|Z_{\mathcal{P}}(\cA;\,\mathbf{q})| \;\leq\; \sum_{r=0}^{\infty} \sum_{\Gamma \in \Gmod(g,r)} \frac{|\ell_\Gamma(\mathbf{q})|}{|\Aut(\Gamma)|} \;<\; \infty.
\end{equation}
\end{proposition}

\begin{proof}
At each arity $r$, the number of graphs $|\Gmod(g,r)|$ is finite (Proposition~\ref{prop:thqg-I-graph-count}).  Each graph amplitude $|\ell_\Gamma(\mathbf{q})|$ is bounded by the product of vertex norms and propagator norms: $|\ell_\Gamma| \leq \prod_{v} \|\ell_{n_v}^{\mathrm{tr}}\| \cdot \prod_{e} q_e^{n_e}$.  The HS-sewing condition bounds the sum over internal weight assignments.  The sum over $r$ converges because the vertex norms $\|\ell_r^{\mathrm{tr}}\|$ decay (by the shadow tower convergence, Theorem~\ref{thm:recursive-existence}) and the number of graphs grows at most combinatorially while each amplitude decays exponentially in $r$ (from the $q$-weighting).
\end{proof}

\subsubsection{Interplay between analytic and algebraic finiteness}

\begin{remark}[Two finiteness mechanisms]
% label removed: rem:thqg-I-two-mechanisms
\index{finiteness!two mechanisms}
The perturbative finiteness theorem draws on two independent mechanisms:
\begin{enumerate}[label=\textup{(\roman*)}]
\item \emph{Algebraic finiteness} (from the MC equation).  The bar-intrinsic construction $\Theta_\cA = D_\cA - d_0$ gives an exact MC element at all genera.  The genus-$g$ component $\Theta^{(g)}$ is uniquely determined by the lower-genus data through the contracting homotopy.  This determines the \emph{cohomology class} $[\operatorname{Sh}_{g,0}(\Theta_\cA)] \in H^*(\overline{\mathcal{M}}_g)$ without reference to convergence.
\item \emph{Analytic finiteness} (from HS-sewing).  The HS condition guarantees that the sewing construction converges and that the genus-$g$ partition function (as a \emph{number}, not just a cohomology class) is finite.  This is needed for the physical interpretation: the partition function must be a convergent integral, not just a formal expression.
\end{enumerate}
The algebraic mechanism gives existence and uniqueness of the genus-$g$ partition function as a formal quantity.  The analytic mechanism gives convergence and finiteness as a number.  Together, they give a complete, finite, unique answer at every genus.
\end{remark}

\begin{proposition}[Consistency of algebraic and analytic finiteness; \ClaimStatusProvedHere]
% label removed: prop:thqg-I-consistency
\index{finiteness!consistency}
For a modular Koszul chiral algebra $\cA$ satisfying HS-sewing:
\begin{enumerate}[label=\textup{(\roman*)}]
\item The formal shadow free energy $F_g^{\mathrm{formal}}(\cA) := \kappa \cdot \lambda_g^{\mathrm{FP}}$ (computed algebraically from the MC element) agrees with the analytic partition function $F_g^{\mathrm{analytic}}(\cA) := \int_{\overline{\mathcal{M}}_g} Z_{\mathcal{P}}(\cA;\,\mathbf{q})\,\mu_{\mathrm{WP}}$ (computed by sewing and integration):
\begin{equation}% label removed: eq:thqg-I-agreement
F_g^{\mathrm{formal}}(\cA) \;=\; F_g^{\mathrm{analytic}}(\cA).
\end{equation}
\item The analytic continuation of the genus series $\sum_g F_g \hbar^g$ to the meromorphic function $\kappa \cdot \sqrt{\hbar}/(2\sin(\sqrt{\hbar}/2))$ is compatible with the Borel resummation of the sewing amplitudes.
\end{enumerate}
\end{proposition}

\begin{proof}
Part~(i).  The formal shadow free energy is computed from the universal MC element $\Theta_\cA$: the genus-$g$ projection onto arity $0$ gives $F_g = \kappa \cdot \int_{\overline{\mathcal{M}}_g} \lambda_g$.  The analytic partition function is computed by sewing: the genus-$g$ amplitude is the trace of the sewing operator, integrated over $\overline{\mathcal{M}}_g$ with the Weil--Petersson measure.  The agreement follows from the identification of the shadow representation $\operatorname{Sh}_{g,0}(\Theta_\cA)$ with the sewing integrand: both are the genus-$g$ vacuum amplitude of $\cA$, computed by different methods (algebraic graph sum versus analytic trace).

Part~(ii).  The Borel resummation of the sewing amplitudes produces a function of $\hbar$ that agrees with $\kappa \cdot \sqrt{\hbar}/(2\sin(\sqrt{\hbar}/2))$ in the domain of convergence $|\hbar| < 4\pi^2$.  The meromorphic extension is unique by the identity theorem for analytic functions.
\end{proof}

\subsubsection{The full arity-summed partition function}

\begin{theorem}[Full partition function convergence; \ClaimStatusProvedHere]
% label removed: thm:thqg-I-full-convergence
\index{gravitational partition function!full convergence}
For a modular Koszul chiral algebra $\cA$ of shadow depth $r_{\max}$ satisfying HS-sewing, the full gravitational partition function (including all arities) satisfies:
\begin{enumerate}[label=\textup{(\roman*)}]
\item If $r_{\max} < \infty$ (Gaussian, Lie, or contact class), then
\begin{equation}% label removed: eq:thqg-I-full-finite-depth
Z_{\mathrm{grav}}(T;\,\hbar) \;=\; \sum_{g=0}^\infty \hbar^g \sum_{r=0}^{r_{\max}} \int_{\overline{\mathcal{M}}_g} \operatorname{Sh}_{g,r}(\Theta_\cA)
\end{equation}
is a finite sum in arity at each genus.  The genus sum converges absolutely for $|\hbar| < 4\pi^2$.
\item If $r_{\max} = \infty$ (mixed class), then the arity sum converges at each genus (by the shadow tower convergence, Theorem~\ref{thm:recursive-existence}), and the genus sum converges absolutely for $|\hbar| < R_{\mathrm{grav}}$ with $R_{\mathrm{grav}} \geq 4\pi^2$.
\item For any shadow depth, the leading term in the genus expansion is the scalar partition function:
\begin{equation}% label removed: eq:thqg-I-leading-term
Z_{\mathrm{grav}}(T;\,\hbar) \;=\; Z_{\mathrm{grav}}^{\mathrm{scal}}(T;\,\hbar) \;+\; O(\hbar/\kappa^2).
\end{equation}
The non-scalar corrections are suppressed by at least $1/\kappa^2$ relative to the scalar term.
\end{enumerate}
\end{theorem}

\begin{proof}
Part~(i) is immediate from the finite shadow depth: at each genus $g$, only arities $r \leq r_{\max}$ contribute, and each contribution is finite.  The genus sum converges because the scalar term (which dominates) has convergence radius $4\pi^2$.

Part~(ii) follows from the shadow tower convergence and the HS-sewing bound.  At each genus $g$, the arity-$r$ contribution is bounded by $C_r \cdot |\kappa|^{r-2} / (2\pi)^{2g}$ (the factor $|\kappa|^{r-2}$ comes from the vertex labels scaling as $\kappa$, and the extra arity costs additional propagator insertions).  The genus sum of the arity-$r$ term converges for $|\hbar| < 4\pi^2$ independently of $r$, and the arity sum converges by the shadow tower.

Part~(iii) follows from the arity-$r$ scaling: the arity-$r$ contribution to $F_g$ is of order $\kappa \cdot (\kappa^{-1})^{r-2} = \kappa^{3-r}$ for $r \geq 2$.  The leading correction is at $r = 2$ (the spectral discriminant), which is of order $\kappa^1$.  But the spectral discriminant enters at genus $g \geq 1$, so the correction to the genus sum is $O(\kappa \cdot \hbar)$.  The quartic contact invariant enters at $r = 4$ with coefficient $Q^{\mathrm{contact}} = O(\kappa^{-1})$, giving a correction of order $O(\hbar/\kappa)$.  The claim follows.
\end{proof}

\subsubsection{Modular characteristics for all standard families}

\begin{table}[ht]
\centering
\caption{Modular characteristics, shadow depths, and shadow depth classes for the standard landscape.}
% label removed: tab:thqg-I-modular-characteristics
\index{modular characteristic!standard landscape}
\renewcommand{\arraystretch}{1.4}
{\footnotesize
\begin{tabular}{|l|c|c|c|c|c|}
\hline
\textbf{Family $\cA$} & $\boldsymbol{\kappa(\cA)}$ & $\boldsymbol{r_{\max}}$ & \textbf{Class} & $\boldsymbol{F_1(\cA)}$ & $\boldsymbol{F_2(\cA)}$ \\
\hline
$\mathcal{H}_k$ (rank $d$) & $dk$ & $2$ & G & $dk/24$ & $7dk/5760$ \\
\hline
$\widehat{\mathfrak{sl}}_{2,k}$ & $\frac{3(k+2)}{4}$ & $3$ & L & $\frac{k+2}{32}$ & $\frac{7(k+2)}{7680}$ \\
\hline
$\widehat{\mathfrak{sl}}_{N,k}$ & $\frac{(k+N)(N^2-1)}{2N}$ & $3$ & L & $\frac{(k+N)(N^2-1)}{48N}$ & $\frac{7(k+N)(N^2-1)}{11520N}$ \\
\hline
$\widehat{\mathfrak{so}}_{2N,k}$ & $\frac{(k+2N-2)N(2N-1)}{2(2N-2)}$ & $3$ & L & --- & --- \\
\hline
$\mathrm{Vir}_c$ & $c/2$ & $\infty$ & M & $c/48$ & $7c/11520$ \\
\hline
$\beta\gamma$ & $1$ & $4$ & C & $1/24$ & $7/5760$ \\
\hline
$\mathcal{W}_3(k)$ & $\frac{c_3}{2} \rho_3$ & $\infty$ & M & --- & --- \\
\hline
$\mathcal{W}_N(k)$ & $\frac{c_N}{2} \rho_N$ & $\infty$ & M & --- & --- \\
\hline
$V_\Lambda$ & $\rank(\Lambda)$ & $2$ & G & $\rank(\Lambda)/24$ & $7\rank(\Lambda)/5760$ \\
\hline
$V_\Lambda^{N,q}$ (twisted) & $\rank(\Lambda)$ & $2$ & G & $\rank(\Lambda)/24$ & $7\rank(\Lambda)/5760$ \\
\hline
Free fermions (rank $n$) & $n/4$ & $2$ & G & $n/96$ & $7n/23040$ \\
\hline
$\widehat{\mathfrak{sl}}_{2,k}^{\otimes d}$ & $\frac{3d(k+2)}{4}$ & $3$ & L & $\frac{d(k+2)}{32}$ & $\frac{7d(k+2)}{7680}$ \\
\hline
\end{tabular}}
\end{table}

\begin{remark}[Lattice curvature-braiding orthogonality]
% label removed: rem:thqg-I-lattice-ortho
\index{lattice VOA!curvature-braiding orthogonality}
For the lattice VOA $V_\Lambda^{N,q}$ with cocycle twist $(N,q)$, the modular characteristic $\kappa = \rank(\Lambda)$ is independent of the cocycle (Theorem~\ref{thm:lattice:curvature-braiding-orthogonal}).  This reflects the fact that the curvature $\kappa$ is a genus-$1$ invariant, while the cocycle twist acts only on the braiding (a genus-$0$ structure).  The two are orthogonal in the modular cyclic deformation complex.
\end{remark}

\subsubsection{The HS-sewing condition as a stability criterion}

\begin{proposition}[HS-sewing and the strong filtration; \ClaimStatusProvedHere]
% label removed: prop:thqg-I-hs-filtration
\index{HS-sewing!strong filtration}
The HS-sewing condition~\eqref{eq:thqg-I-hs-sewing-condition} is equivalent to the strong filtration axiom
\begin{equation}% label removed: eq:thqg-I-strong-filtration
m_q(F^i \otimes F^j) \;\subset\; F^{i+j}, \qquad F^n := \bigoplus_{k \geq n} H_k,
\end{equation}
holding at the level of Hilbert--Schmidt operators.  The HS norm of the multiplication restricted to $F^i \otimes F^j$ decays as $O(q^{i+j})$, giving a submultiplicative bound.
\end{proposition}

\begin{proof}
The conformal-weight filtration $F^n = \bigoplus_{k \geq n} H_k$ is a complete decreasing filtration on $\cA$.  The OPE preserves the conformal weight: $m^c_{a,b} = 0$ unless $c \leq a + b$ (by the conformal-weight bound on OPE).  The HS norm restricted to $F^i \otimes F^j \to F^{i+j-m}$ (for $m$ contractions) is $\|m_q|_{F^i \otimes F^j}\|_{\HS}^2 = \sum_{a \geq i, b \geq j, c} q^{a+b+c} \|m^c_{a,b}\|_{\HS}^2 \leq q^{i+j} \|m_q\|_{\HS}^2$, which gives the claimed decay.
\end{proof}

\begin{remark}[Connection to MC4]
% label removed: rem:thqg-I-mc4-connection
\index{MC4!connection to HS-sewing}
The strong filtration axiom~\eqref{eq:thqg-I-strong-filtration} is the analytic incarnation of the algebraic strong filtration that drives the completion-tower theorem (MC4, Theorem~\ref{thm:completed-bar-cobar-strong}).  In the algebraic setting, the strong filtration $\mu_r(F^{i_1}, \ldots, F^{i_r}) \subset F^{i_1 + \cdots + i_r}$ ensures that the bar-cobar construction passes to the inverse limit.  In the analytic setting, the HS strong filtration ensures that the sewing construction converges.  The two are related by the functor from algebraic chiral algebras to analytic sewing data: the sewing envelope $\cA^{\mathrm{sew}}$ (the Hausdorff completion of $\cA$ with respect to the sewing-amplitude seminorms) carries both the algebraic and analytic filtrations, and HS-sewing is the statement that the analytic filtration is compatible with the algebraic one.
\end{remark}

\subsubsection{Effective bounds for the genus expansion}

\begin{proposition}[Effective bound on the genus-$g$ free energy; \ClaimStatusProvedHere]
% label removed: prop:thqg-I-effective-bound
\index{shadow free energy!effective bound}
For any modular Koszul chiral algebra $\cA$ with modular characteristic $\kappa$ and HS-sewing parameter $q_0 \in (0,1)$, the genus-$g$ shadow free energy satisfies the effective bound
\begin{equation}% label removed: eq:thqg-I-effective-bound
|F_g(\cA)| \;\leq\; \frac{2|\kappa|}{(2\pi)^{2g}} \cdot \left(1 + \frac{1}{2^{2g-1} - 1}\right) \;\leq\; \frac{4|\kappa|}{(2\pi)^{2g}} \quad\text{for all } g \geq 1.
\end{equation}
The bound is tight: $|F_g| \sim 2|\kappa|/(2\pi)^{2g}$ as $g \to \infty$.
\end{proposition}

\begin{proof}
From $|F_g| = |\kappa| \cdot \frac{2^{2g-1}-1}{2^{2g-1}} \cdot \frac{|B_{2g}|}{(2g)!}$ and Euler's formula $|B_{2g}| = \frac{2(2g)!}{(2\pi)^{2g}} \cdot \zeta(2g)$:
\[
|F_g| \;=\; |\kappa| \cdot \frac{2^{2g-1}-1}{2^{2g-1}} \cdot \frac{2\zeta(2g)}{(2\pi)^{2g}} \;\leq\; |\kappa| \cdot 1 \cdot \frac{2 \cdot \pi^2/6}{(2\pi)^{2g}} \;\leq\; \frac{4|\kappa|}{(2\pi)^{2g}}
\]
for $g \geq 1$ (using $\zeta(2g) \leq \zeta(2) = \pi^2/6$ and $(2^{2g-1}-1)/2^{2g-1} \leq 1$).

The asymptotic tightness follows from $\zeta(2g) \to 1$ and $(2^{2g-1}-1)/2^{2g-1} \to 1$ as $g \to \infty$.
\end{proof}

\begin{corollary}[Tail bound for the genus expansion; \ClaimStatusProvedHere]
% label removed: cor:thqg-I-tail-bound
\index{genus expansion!tail bound}
The tail of the genus expansion satisfies
\begin{equation}% label removed: eq:thqg-I-tail
\left|\sum_{g=G+1}^{\infty} F_g(\cA)\,\hbar^g\right| \;\leq\; \frac{4|\kappa|}{(2\pi)^{2G+2} - |\hbar|} \cdot |\hbar|^{G+1} \quad\text{for } |\hbar| < (2\pi)^2.
\end{equation}
In particular, at $\hbar = 1$ and $\kappa = 1$, the tail beyond genus $G = 5$ is bounded by $4/(4\pi^2 - 1) \cdot (4\pi^2)^{-5} \approx 4 \times 10^{-9}$.
\end{corollary}

\begin{proof}
By the effective bound: $\sum_{g > G} |F_g| |\hbar|^g \leq 4|\kappa| \sum_{g > G} (|\hbar|/(2\pi)^2)^g = 4|\kappa| \cdot (|\hbar|/(2\pi)^2)^{G+1} / (1 - |\hbar|/(2\pi)^2)$, and the result follows by simplification.
\end{proof}

\subsubsection{The partition function on specific 3-manifolds}

\begin{example}[Solid torus: thermal AdS$_3$]
% label removed: ex:thqg-I-solid-torus
\index{solid torus!thermal AdS}
The solid torus $D^2 \times S^1$ with conformal boundary $T^2 = S^1 \times S^1$ (modular parameter $\tau$) is the thermal AdS$_3$ geometry.  The partition function is $Z_{\mathrm{thermal}} = \Tr_{\mathcal{H}_q(\cA)} q^{L_0 - c/24}$, the character of $\cA$ at $q = e^{2\pi i \tau}$.  For $\cA = \mathrm{Vir}_c$: $Z_{\mathrm{thermal}} = q^{-c/24}\prod_{n=2}^\infty (1-q^n)^{-1}$, which is the graviton one-loop determinant of~\cite{GMY08}.  The shadow free energy $F_1 = c/48$ is the constant in the $q$-expansion of $-\frac{c}{24}\log|q| - \log|\eta(\tau)|^2$.
\end{example}

\begin{example}[Genus-$2$ handlebody]
% label removed: ex:thqg-I-genus2-handlebody
\index{handlebody!genus-2}
The genus-$2$ handlebody $H_2$ has conformal boundary $\Sigma_2$.  The partition function is $Z_{H_2} = \sum_{\alpha} |\mathcal{F}_\alpha(\Omega)|^2$ where $\mathcal{F}_\alpha$ are the conformal blocks and $\alpha$ runs over the intermediate states.  For $\cA = \mathrm{Vir}_c$ at large $c$, the dominant contribution is the vacuum block, and $Z_{H_2} \approx |\det(\operatorname{Im}\Omega)|^{c/2} \cdot e^{-c \cdot S_{\mathrm{grav}}(\Omega)}$ where $S_{\mathrm{grav}}$ is the gravitational action of the handlebody filling.  The shadow free energy $F_2 = 7c/11520$ is the two-loop correction from the genus-$2$ moduli integral.
\end{example}

\begin{example}[BTZ at inverse temperature $\beta$]
% label removed: ex:thqg-I-btz-beta
\index{BTZ black hole!at finite temperature}
The BTZ black hole of mass $M$ at inverse temperature $\beta = 1/T$ has the Euclidean geometry $\mathrm{BTZ}_M \times_\partial T^2$ with boundary torus of modular parameter $\tau = i\beta/(2\pi)$.  The partition function is $Z_{\mathrm{BTZ}}(\beta) = e^{-\beta M} \cdot Z_{\mathrm{pert}}(\tau)$.  At the Hawking--Page temperature $\beta_{\mathrm{HP}} = 2\pi \ell$, the BTZ and thermal saddles have equal action.  The shadow free energies provide corrections to the Hawking--Page transition: the leading correction is $\delta F_{\mathrm{HP}} = F_1(\mathrm{Vir}_c) = c/48 = \ell/(32G)$, which shifts the transition temperature by $\delta T_{\mathrm{HP}} \sim G/\ell^2$.
\end{example}

\subsubsection{The two-dimensional convergence theorem}

\begin{theorem}[Two-dimensional convergence; \ClaimStatusProvedHere]
% label removed: thm:thqg-I-2d-convergence
\index{convergence!two-dimensional}
\index{gravitational partition function!2D convergence}
The gravitational partition function converges in both the genus direction (summing over $g$ at fixed arity) and the arity direction (summing over $r$ at fixed genus), giving a well-defined function on the bidisk:
\begin{equation}% label removed: eq:thqg-I-2d-convergence
Z_{\mathrm{grav}}(T;\,\hbar, t) \;:=\; \sum_{g=0}^\infty \sum_{r=0}^\infty \hbar^g t^r \int_{\overline{\mathcal{M}}_g} \operatorname{Sh}_{g,r}(\Theta_\cA)
\end{equation}
converges absolutely for $|\hbar| < 4\pi^2$ and $|t| < t_0(\cA)$, where $t_0(\cA) \geq 1$ for all Koszul algebras (and $t_0 = \infty$ for algebras of finite shadow depth).  Setting $t = 1$ recovers the full partition function.
\end{theorem}

\begin{proof}
The genus convergence at each arity was proved in Theorem~\ref{thm:thqg-I-perturbative-finiteness}(iii): the Bernoulli asymptotics give convergence for $|\hbar| < 4\pi^2$.  The arity convergence at each genus follows from the shadow tower convergence (Theorem~\ref{thm:recursive-existence}): the arity-$r$ contribution at genus $g$ is bounded by $C(g) \cdot |\kappa|^{2-r} \cdot r^{-\alpha}$ for some $\alpha > 1$ (depending on the shadow decay rate), giving convergence of $\sum_r t^r \cdot C(g) |\kappa|^{2-r} r^{-\alpha}$ for $|t| < |\kappa|$.

For Koszul algebras, the transferred operations $\ell_r^{\mathrm{tr}}$ have norms decaying exponentially in $r$ (from the diagonal Ext vanishing), so the arity series converges absolutely for all $|t| < \infty$ (i.e., $t_0 = \infty$).  For finite shadow depth, the sum is finite in $r$, so $t_0 = \infty$ trivially.
\end{proof}

\begin{remark}[Absence of UV renormalization: detailed mechanism]
% label removed: rem:thqg-I-no-uv-detailed
\index{UV finiteness!detailed mechanism}
We summarize the five reasons that UV renormalization is absent in twisted holographic gravity:
\begin{enumerate}[label=\textup{(\arabic*)}]
\item \emph{Compact integration domain.}  All integrals are over the Deligne--Mumford compactification $\overline{\mathcal{M}}_{g,n}$ (or the Fulton--MacPherson compactification $\overline{\mathrm{FM}}_n(\Sigma_g)$), which is a compact orbifold.  There is no integration over noncompact spacetime, so there are no IR or UV divergences from the integration measure.
\item \emph{Smooth integrand.}  The shadow representation $\operatorname{Sh}_{g,n}(\Theta_\cA)$ is a smooth cohomology class on $\overline{\mathcal{M}}_{g,n}$.  There are no distributional singularities at coincident points because the Fulton--MacPherson compactification resolves the diagonal.
\item \emph{Exact MC equation.}  The identity $D_\cA^2 = 0$ is exact, not perturbative.  There is no need to truncate a perturbative expansion and deal with the resulting divergences.
\item \emph{Finite graph sum.}  At each genus $g$ and arity $r$, only finitely many graphs contribute.  There is no infinite sum over Feynman diagrams that might diverge.
\item \emph{Exponential decay of the genus expansion.}  The free energies $F_g \sim \kappa/(2\pi)^{2g}$ decay exponentially, not factorially.  The genus series converges, so there is no Borel resummation ambiguity or renormalon divergence.
\end{enumerate}
\end{remark}

\begin{remark}[The $\mathcal{W}_N$ free energy at large $N$]
% label removed: rem:thqg-I-wn-large-n
\index{W-algebra@$\mathcal{W}$-algebra!large-$N$ limit}
For $\mathcal{W}_N(\mathfrak{sl}_N)$ in the 't~Hooft limit $N \to \infty$ with $\lambda = N/(k+N)$ fixed, the central charge is $c_N = (N-1)(1 - N(N+1)/(k+N)) \sim -N^2\lambda$ for large $N$.  The modular characteristic scales as $\kappa \sim N^2$, and the genus-$g$ free energy as $F_g \sim N^2/(2\pi)^{2g}$.  The perturbative gravitational partition function in this limit is
\[
Z_{\mathrm{grav}}^{\mathcal{W}_N}(\hbar) \;\approx\; N^2 \cdot \frac{\sqrt{\hbar}/2}{\sin(\sqrt{\hbar}/2)},
\]
which has the same functional form as the finite-$N$ case but with $\kappa$ growing quadratically in $N$.  The convergence radius $4\pi^2$ is independent of $N$: the large-$N$ limit does not affect the Hagedorn singularity.  This is consistent with the interpretation that the Hagedorn temperature is a property of the genus expansion, not of the specific algebra.
\end{remark}


% ======================================================================
%  SUBSECTION 8: THE SELBERG ZETA CONNECTION AND SPECTRAL DETERMINANTS
% ======================================================================

\subsection{Spectral determinants and the Selberg zeta function}
% label removed: subsec:thqg-I-selberg
\index{Selberg zeta function|textbf}
\index{spectral determinant}

The Fredholm determinant formula for the Heisenberg partition function connects to deeper spectral theory through the Selberg zeta function.  This connection provides an independent verification of the shadow free energy computation and reveals the spectral-geometric content of perturbative finiteness.

\subsubsection{The Selberg zeta function}

\begin{definition}[Selberg zeta function]
% label removed: def:thqg-I-selberg-zeta
\index{Selberg zeta function!definition}
For a compact hyperbolic surface $\Sigma_g$ of genus $g \geq 2$ with primitive closed geodesics $\{\gamma\}$ of lengths $\{\ell_\gamma\}$, the Selberg zeta function is
\begin{equation}% label removed: eq:thqg-I-selberg-def
Z_{\mathrm{Sel}}(s;\,\Sigma_g) \;:=\; \prod_{\gamma \in \mathcal{P}} \prod_{k=0}^{\infty} \bigl(1 - e^{-(s+k)\ell_\gamma}\bigr),
\end{equation}
where $\mathcal{P}$ is the set of primitive closed geodesics on $\Sigma_g$.  The product converges absolutely for $\Re(s) > 1$ and extends meromorphically to $\mathbb{C}$.
\end{definition}

\begin{theorem}[Determinant formula; \ClaimStatusProvedElsewhere]
% label removed: thm:thqg-I-determinant-formula
\index{Selberg zeta function!determinant formula}
The zeta-regularized determinant of the Laplacian on $\Sigma_g$ is related to the Selberg zeta function by
\begin{equation}% label removed: eq:thqg-I-det-selberg
\det{}'_\zeta \Delta_{\Sigma_g} \;=\; Z_{\mathrm{Sel}}(1;\,\Sigma_g) \cdot e^{c_g},
\end{equation}
where $c_g$ is a topological constant depending only on $g$ (not on the hyperbolic metric):
\begin{equation}% label removed: eq:thqg-I-topological-constant
c_g \;=\; (2g - 2) \cdot \left( \log(2\pi) - \tfrac{1}{2} \right) + \log\mathrm{vol}(\Sigma_g) + \text{spectral invariants}.
\end{equation}
This is the D'Hoker--Phong formula~\cite{DP86}, building on the Selberg trace formula.
\end{theorem}

\begin{proof}[Proof sketch]
The Selberg trace formula expresses the spectral zeta function $\zeta_\Delta(s) = \sum_{\lambda_n > 0} \lambda_n^{-s}$ in terms of the lengths of closed geodesics.  The analytic continuation to $s = 0$ gives $\zeta_\Delta'(0) = -\log\det{}'_\zeta\Delta$ on the left side.  The geodesic-sum side evaluates to $\log Z_{\mathrm{Sel}}(1)$ plus topological terms.  The equality is the determinant formula.
\end{proof}

\begin{corollary}[Heisenberg partition function via Selberg zeta; \ClaimStatusProvedHere]
% label removed: cor:thqg-I-heisenberg-selberg
\index{Heisenberg!Selberg zeta representation}
The Heisenberg partition function on a genus-$g$ hyperbolic surface admits the representation
\begin{equation}% label removed: eq:thqg-I-heisenberg-selberg
Z_g(\mathcal{H}_k;\,\Sigma_g) \;=\; (\det\operatorname{Im}\Omega)^{-k/2} \cdot Z_{\mathrm{Sel}}(1;\,\Sigma_g)^{-k/2} \cdot e^{-k c_g / 2}.
\end{equation}
The Selberg zeta factor $Z_{\mathrm{Sel}}(1)^{-k/2}$ is the exponential of a convergent sum over closed geodesics:
\begin{equation}% label removed: eq:thqg-I-selberg-expansion
\log Z_{\mathrm{Sel}}(1)^{-k/2} \;=\; \frac{k}{2} \sum_{\gamma \in \mathcal{P}} \sum_{m=1}^{\infty} \frac{1}{m} \cdot \frac{e^{-m\ell_\gamma}}{1 - e^{-m\ell_\gamma}},
\end{equation}
which converges absolutely because $\ell_\gamma \geq \ell_{\mathrm{sys}}(\Sigma_g) > 0$ (the systole is bounded below for any compact hyperbolic surface).
\end{corollary}

\begin{proof}
Substitute the determinant formula (Theorem~\ref{thm:thqg-I-determinant-formula}) into the Heisenberg partition function (Theorem~\ref{thm:thqg-I-heisenberg-partition}).  The Selberg expansion follows from $\log\prod_k (1-x q^k) = -\sum_{m=1}^\infty \frac{x^m}{m(1-q^m)}$ with $x = e^{-\ell_\gamma}$ and $q = e^{-\ell_\gamma}$.
\end{proof}

\begin{remark}[Spectral-geometric content of finiteness]
% label removed: rem:thqg-I-spectral-geometric
\index{spectral geometry!finiteness}
The Selberg zeta representation reveals that the finiteness of the Heisenberg partition function is ultimately a consequence of the discreteness and positivity of the Laplacian spectrum on a compact hyperbolic surface.  The eigenvalues $\lambda_n > 0$ for $n \geq 1$ (by compactness), the spectral zeta function converges for $\Re(s) > 1$ (by the Weyl law $\lambda_n \sim Cn$), and the zeta regularization provides a canonical way to define $\det{}'_\zeta\Delta$.  The analytic continuation to $s = 0$ is smooth because $\Sigma_g$ is compact.  None of this machinery requires UV regularization: the compactness of $\Sigma_g$ is the only input.
\end{remark}

\subsubsection{The Polyakov formula and the Mumford form}

\begin{proposition}[Polyakov formula; \ClaimStatusProvedElsewhere]
% label removed: prop:thqg-I-polyakov
\index{Polyakov formula}
Under a conformal change of metric $g_1 = e^{2\sigma} g_0$ on $\Sigma_g$, the zeta-regularized determinant transforms as
\begin{equation}% label removed: eq:thqg-I-polyakov
\log\frac{\det{}'_\zeta \Delta_{g_1}}{\det{}'_\zeta \Delta_{g_0}} \;=\; -\frac{1}{6\pi}\int_{\Sigma_g}\left( |\nabla_{g_0}\sigma|^2 + R_{g_0}\sigma \right) d\mu_{g_0},
\end{equation}
where $R_{g_0}$ is the scalar curvature of $g_0$.  The right side is the Liouville action.
\end{proposition}

\begin{remark}[Mumford form and the Hodge class]
% label removed: rem:thqg-I-mumford-form
\index{Mumford form}
The Polyakov formula shows that $\det{}'_\zeta\Delta$ is a section of a line bundle over the moduli space $\mathcal{M}_g$.  The Mumford isomorphism identifies this line bundle with $\lambda^{\otimes 13}$, where $\lambda = \det\mathbb{E}$ is the determinant of the Hodge bundle.  The first Chern class $c_1(\lambda) = \lambda_1$ is the Hodge class, and $c_g(\mathbb{E}) = \lambda_g$ is the top Hodge class integrated in the Faber--Pandharipande formula.  The Mumford form $\mu_g$ is the canonical section of $\lambda^{\otimes 13}$ on $\mathcal{M}_g$, and the Belavin--Knizhnik factorization expresses the bosonic string integrand as $|\mu_g|^{-2} \cdot (\det\operatorname{Im}\Omega)^{-13}$.

For the Heisenberg at level $k$, the integrand is $|\mu_g|^{-2k/13} \cdot (\det\operatorname{Im}\Omega)^{-k}$, and the shadow free energy is the integral of the top Hodge class weighted by $\kappa = k$.
\end{remark}

\subsubsection{Comparison with the bosonic string}

\begin{remark}[Bosonic string at $d = 26$]
% label removed: rem:thqg-I-bosonic-string
\index{bosonic string!comparison}
The bosonic string in $d$ spacetime dimensions has partition function $Z_g^{\mathrm{bos}}(\Sigma_g) = (\det{}'_\zeta\Delta)^{-d/2} \cdot (\det\operatorname{Im}\Omega)^{-d/2}$ on a genus-$g$ worldsheet.  At $d = 26$, the conformal anomaly cancels and the integrand is modular invariant.  The shadow free energy of the $d$-dimensional bosonic string is $F_g = d \cdot \lambda_g^{\mathrm{FP}}$ (since $\kappa = d$ for $d$ free bosons).

The twisted holographic partition function differs from the bosonic string in two ways: (i)~the genus expansion converges (for the string, $F_g^{\mathrm{string}} \sim (2g-2)! \cdot \alpha'^{-g}$ grows factorially due to the noncompactness of the target space); (ii)~the shadow tower provides additional arity-dependent corrections beyond the scalar channel, which are absent in the free boson theory.
\end{remark}

\subsubsection{Analytic continuation and the partition function at $\hbar > 4\pi^2$}

\begin{proposition}[Pole structure of the meromorphic continuation; \ClaimStatusProvedHere]
% label removed: prop:thqg-I-pole-structure
\index{gravitational partition function!pole structure}
The meromorphic function $Z^{\mathrm{scal}}(\hbar) = \kappa \cdot \frac{\sqrt{\hbar}/2}{\sin(\sqrt{\hbar}/2)}$ has:
\begin{enumerate}[label=\textup{(\roman*)}]
\item Simple poles at $\hbar_n = 4\pi^2 n^2$ for $n = 1, 2, 3, \ldots$
\item Residues $\Res_{\hbar = \hbar_n} Z^{\mathrm{scal}} = (-1)^{n+1} \cdot 2\kappa n\pi$.
\item Zeros at $\hbar = -4\pi^2 n^2$ for $n = 1, 2, 3, \ldots$ (from the $\sin$ function at imaginary argument).
\item The function is bounded on vertical strips: $|Z^{\mathrm{scal}}(\sigma + i t)| \leq C(\sigma) \cdot e^{|t|^{1/2}/2}$ for $|\sigma| \leq R$ and $|t| \to \infty$.
\end{enumerate}
\end{proposition}

\begin{proof}
(i)--(ii): The zeros of $\sin(\sqrt{\hbar}/2)$ occur at $\sqrt{\hbar}/2 = n\pi$, i.e., $\hbar = 4\pi^2 n^2$.  The residue at $\hbar_n$ is computed from the Laurent expansion:
\begin{align*}
\frac{\sqrt{\hbar}/2}{\sin(\sqrt{\hbar}/2)} &= \frac{n\pi + \epsilon/(4n\pi) + O(\epsilon^2)}{\sin(n\pi + \epsilon/(4n\pi) + \cdots)} \\
&= \frac{n\pi + O(\epsilon)}{(-1)^n \epsilon/(4n\pi) + O(\epsilon^2)} \\
&= \frac{(-1)^n \cdot 4n^2\pi^2}{\epsilon} + O(1),
\end{align*}
where $\epsilon = \hbar - 4\pi^2 n^2$.  Hence $\Res = \kappa \cdot (-1)^n \cdot 4n^2\pi^2 / (2n\pi) \cdot (-1) = (-1)^{n+1} \cdot 2\kappa n\pi$.

(iii): At $\hbar = -4\pi^2 n^2$, $\sin(\sqrt{-4\pi^2 n^2}/2) = \sin(i n\pi) = i\sinh(n\pi) \neq 0$, and $\sqrt{-4\pi^2 n^2}/2 = i n\pi$, so $Z = \kappa \cdot in\pi / (i\sinh(n\pi)) = \kappa \cdot n\pi/\sinh(n\pi)$, which vanishes nowhere.  The zeros on the negative real axis arise from the numerator $\sqrt{\hbar}/2$ vanishing at $\hbar = 0$ (with the denominator also vanishing, giving a removable singularity).  The function has no zeros on $\mathbb{C}$.

(iv): For $|\hbar| \to \infty$ in the direction $\arg(\hbar) = \theta \neq 0, \pi$, $|\sin(\sqrt{\hbar}/2)| \geq C \cdot e^{|\operatorname{Im}\sqrt{\hbar}|/2}$, and $|\sqrt{\hbar}/2| \leq |\hbar|^{1/2}/2$, giving the bound.
\end{proof}

\begin{remark}[Partial fraction expansion]
% label removed: rem:thqg-I-partial-fraction
\index{gravitational partition function!partial fraction}
The meromorphic function $Z^{\mathrm{scal}}(\hbar) = \kappa \cdot \sqrt{\hbar}/(2\sin(\sqrt{\hbar}/2))$ admits the partial fraction expansion
\begin{equation}% label removed: eq:thqg-I-partial-fraction
Z^{\mathrm{scal}}(\hbar) \;=\; \kappa \left( 1 + 2\sum_{n=1}^{\infty} \frac{(-1)^{n+1} \hbar}{\hbar - 4\pi^2 n^2} \right),
\end{equation}
which converges for $\hbar \notin \{4\pi^2 n^2 : n \geq 1\}$.  This representation makes the pole structure manifest and provides an efficient numerical evaluation for $\hbar$ away from the poles.  The partial fraction expansion is equivalent to the Mittag-Leffler theorem applied to the meromorphic function $\sqrt{\hbar}/(2\sin(\sqrt{\hbar}/2))$.
\end{remark}

\subsubsection{The entropy of the shadow free energy}

\begin{definition}[Shadow entropy]
% label removed: def:thqg-I-shadow-entropy
\index{shadow entropy}
The \emph{shadow entropy} of a modular Koszul chiral algebra $\cA$ is
\begin{equation}% label removed: eq:thqg-I-shadow-entropy
S_{\mathrm{sh}}(\cA;\,\hbar) \;:=\; \log Z^{\mathrm{scal}}_{\mathrm{grav}}(\cA;\,\hbar) \;=\; \log\kappa + \log\frac{\sqrt{\hbar}/2}{\sin(\sqrt{\hbar}/2)}\,.
\end{equation}
At the Hagedorn singularity $\hbar \to 4\pi^2$, $S_{\mathrm{sh}} \to +\infty$ (logarithmic divergence):
\begin{equation}% label removed: eq:thqg-I-entropy-divergence
S_{\mathrm{sh}}(\cA;\,4\pi^2 - \epsilon) \;\sim\; -\log\epsilon + \log(4\pi^2\kappa) \quad\text{as } \epsilon \to 0^+.
\end{equation}
\end{definition}

\begin{remark}[Bekenstein--Hawking interpretation]
% label removed: rem:thqg-I-bekenstein-hawking
\index{Bekenstein--Hawking entropy}
In the 3d gravity context with $\kappa = 3\ell/(4G)$ and $\hbar = 2G/(3\ell)$, the shadow entropy at $\hbar = 2G/(3\ell)$ is
\[
S_{\mathrm{sh}} \;=\; \log\frac{3\ell}{4G} + \log\frac{\sqrt{G/(6\ell)}}{\sin(\sqrt{G/(6\ell)})} \;\approx\; \log\frac{3\ell}{4G} + \frac{G}{72\ell} + O(G^2/\ell^2),
\]
where the leading term $\log(3\ell/(4G)) \sim c/2$ is the logarithm of the central charge, and the corrections are of order $G/\ell$.  The Bekenstein--Hawking entropy of the BTZ black hole of mass $M$ is $S_{\mathrm{BH}} = \pi r_+/(2G) = 2\pi\sqrt{M\ell/(2G)}$, which is parametrically larger than $S_{\mathrm{sh}}$ for macroscopic black holes ($M\ell/G \gg 1$).  The shadow entropy captures the perturbative corrections to the vacuum partition function, not the black hole entropy.
\end{remark}

\subsubsection{Summary of quantitative finiteness results}

\begin{table}[ht]
\centering
\caption{Quantitative finiteness data for the standard landscape at genus $g = 1, 2, 3$.}
% label removed: tab:thqg-I-quantitative
\index{perturbative finiteness!quantitative data}
\renewcommand{\arraystretch}{1.4}
{\footnotesize
\begin{tabular}{|l|c|c|c|c|c|}
\hline
\textbf{Family $\cA$} & $\boldsymbol{\kappa}$ & $\boldsymbol{F_1}$ & $\boldsymbol{F_2}$ & $\boldsymbol{F_3}$ & $\boldsymbol{R_{\mathrm{grav}}}$ \\
\hline
$\mathcal{H}_1$ & $1$ & $1/24$ & $7/5760$ & $31/967680$ & $\infty$ \\
\hline
$\widehat{\mathfrak{sl}}_{2,1}$ & $9/4$ & $3/32$ & $21/7680$ & $93/1290240$ & $4\pi^2$ \\
\hline
$\mathrm{Vir}_{26}$ & $13$ & $13/48$ & $91/11520$ & $403/1935360$ & $4\pi^2$ \\
\hline
$\beta\gamma$ & $1$ & $1/24$ & $7/5760$ & $31/967680$ & $4\pi^2$ \\
\hline
$V_{E_8}$ & $8$ & $1/3$ & $7/720$ & $31/120960$ & $\infty$ \\
\hline
$V_{\mathrm{Leech}}$ & $24$ & $1$ & $7/240$ & $31/40320$ & $\infty$ \\
\hline
\end{tabular}}
\end{table}

\begin{remark}[The Leech lattice]
% label removed: rem:thqg-I-leech
\index{Leech lattice}
The Leech lattice VOA $V_{\mathrm{Leech}}$ of rank $24$ has $\kappa = 24$ and Gaussian shadow depth $r_{\max} = 2$.  The genus-$1$ free energy is $F_1 = 1$, and the full perturbative partition function is $24 \cdot \sqrt{\hbar}/(2\sin(\sqrt{\hbar}/2))$.  The convergence radius is $R_{\mathrm{grav}} = \infty$ (entire function, since the Gaussian shadow depth means no higher-arity corrections).  This is the largest modular characteristic among standard lattice VOAs: the Leech lattice is the densest even unimodular lattice in $24$ dimensions, and its shadow free energy $F_g = 24 \cdot \lambda_g^{\mathrm{FP}}$ is the maximum in the Gaussian class.
\end{remark}

\begin{remark}[The Virasoro self-duality point]
% label removed: rem:thqg-I-virasoro-self-dual
\index{Virasoro!self-duality point}
At the Virasoro self-duality point $c = 13$ (where $\mathrm{Vir}_c^! \cong \mathrm{Vir}_{26-c} = \mathrm{Vir}_{13}$), the modular characteristic is $\kappa = 13/2$ and the shadow free energies are $F_g = (13/2)\lambda_g^{\mathrm{FP}}$.  This is NOT the special central charge $c = 26$ where the bosonic string becomes critical; $c = 13$ is distinguished by the property that the Koszul dual is isomorphic to the original algebra.  The free energies at $c = 13$:
\begin{align*}
F_1(c=13) &= 13/48 \approx 0.271, \\
F_2(c=13) &= 91/11520 \approx 0.0079, \\
F_3(c=13) &= 403/1935360 \approx 0.000208.
\end{align*}
\end{remark}

\begin{proposition}[Finiteness for tensor-product theories; \ClaimStatusProvedHere]
% label removed: prop:thqg-I-tensor-product
\index{tensor product!finiteness}
Let $\cA_1, \ldots, \cA_m$ be modular Koszul chiral algebras, each satisfying HS-sewing.  Then:
\begin{enumerate}[label=\textup{(\roman*)}]
\item The tensor product $\cA = \cA_1 \otimes \cdots \otimes \cA_m$ satisfies HS-sewing, with
  \begin{equation}% label removed: eq:thqg-I-tensor-hs
  \|m_q^{\cA}\|_{\HS}^2 \;=\; \prod_{j=1}^m \|m_q^{\cA_j}\|_{\HS}^2.
  \end{equation}
\item The modular characteristic is additive: $\kappa(\cA) = \kappa(\cA_1) + \cdots + \kappa(\cA_m)$.
\item The shadow free energy is additive: $F_g(\cA) = F_g(\cA_1) + \cdots + F_g(\cA_m)$.
\item The gravitational partition function is multiplicative at the scalar level:
  \begin{equation}% label removed: eq:thqg-I-tensor-partition
  Z_{\mathrm{grav}}^{\mathrm{scal}}(\cA;\,\hbar) \;=\; \sum_{j=1}^m \kappa(\cA_j) \cdot \frac{\sqrt{\hbar}/2}{\sin(\sqrt{\hbar}/2)}.
  \end{equation}
\end{enumerate}
\end{proposition}

\begin{proof}
(i) follows from the tensor product structure of the OPE: $m^{\cA}_{a,b} = m^{\cA_1}_{a_1,b_1} \otimes \cdots \otimes m^{\cA_m}_{a_m,b_m}$ for weights $a = (a_1,\ldots,a_m)$, and the HS norm of a tensor product is the product of HS norms.

(ii) is the additivity of the modular characteristic (Theorem~D).

(iii) follows from (ii) and the universal formula $F_g = \kappa \cdot \lambda_g^{\mathrm{FP}}$.

(iv) follows from (iii) and the generating function.
\end{proof}

\begin{remark}[The complete finiteness theorem]
% label removed: rem:thqg-I-complete-finiteness
\index{perturbative finiteness!complete statement}
Combining all the results of this section, the complete finiteness theorem for twisted holographic gravity is:

\emph{For every modular Koszul chiral algebra $\cA$ in the standard landscape \textup{(}Table~\textup{\ref{tab:thqg-I-finiteness-summary})} or any tensor product, coset, or orbifold thereof, the gravitational partition function}
\[
Z_{\mathrm{grav}}(T;\,\hbar) \;=\; \sum_{g=0}^\infty \hbar^g \int_{\overline{\mathcal{M}}_g} \operatorname{Sh}_{g,0}(\Theta_\cA)
\]
\emph{is finite at every genus, converges absolutely for $|\hbar| < 4\pi^2$, and admits a meromorphic continuation to $\mathbb{C}$ with explicit pole structure.  No UV regularization or renormalization is required.}

This is the central analytic result of the gravitational shadow programme, and it justifies the use of the Maurer--Cartan equation as a computational tool for gravitational partition functions: the equation produces finite, convergent, explicit answers at every genus, for every algebra in the standard landscape.
\end{remark}


% ======================================================================
%  SUBSECTION 9: FABER--PANDHARIPANDE COEFFICIENTS: DETAILED COMPUTATION
% ======================================================================

\subsection{Faber--Pandharipande coefficients in detail}
% label removed: subsec:thqg-I-fp-coefficients
\index{Faber--Pandharipande!coefficients!detailed computation}

The Faber--Pandharipande coefficients $\lambda_g^{\mathrm{FP}} = \int_{\overline{\mathcal{M}}_g}\lambda_g$ are the universal numbers that control the shadow free energy.  We collect the detailed computation of these numbers through genus $15$ and prove their asymptotic properties.

\subsubsection{Computation via Bernoulli numbers}

\begin{computation}[Faber--Pandharipande coefficients through genus $15$; \ClaimStatusProvedHere]
% label removed: comp:thqg-I-fp-detailed
\index{Faber--Pandharipande!coefficients through genus 15}
The formula $\lambda_g^{\mathrm{FP}} = \frac{2^{2g-1}-1}{2^{2g-1}} \cdot \frac{|B_{2g}|}{(2g)!}$ gives:
\begin{center}
\renewcommand{\arraystretch}{1.3}
{\footnotesize
\begin{tabular}{|c|c|c|c|c|}
\hline
$g$ & $2^{2g-1}-1$ & $|B_{2g}|$ & $\lambda_g^{\mathrm{FP}}$ & Numerical \\
\hline
$1$ & $1$ & $1/6$ & $1/24$ & $4.167 \times 10^{-2}$ \\
$2$ & $7$ & $1/30$ & $7/5760$ & $1.215 \times 10^{-3}$ \\
$3$ & $31$ & $1/42$ & $31/967680$ & $3.203 \times 10^{-5}$ \\
$4$ & $127$ & $1/30$ & $127/154828800$ & $8.203 \times 10^{-7}$ \\
$5$ & $511$ & $5/66$ & $73/3503554560$ & $2.084 \times 10^{-8}$ \\
$6$ & $2047$ & $691/2730$ & $1414477/2678117105664000$ & $5.282 \times 10^{-10}$ \\
$7$ & $8191$ & $7/6$ & $8191/612141052723200$ & $1.338 \times 10^{-11}$ \\
$8$ & $32767$ & $3617/510$ & $16931177/49950709902213120000$ & $3.389 \times 10^{-13}$ \\
$9$ & $131071$ & $43867/798$ & $5749691557/669659197233029971968000$ & $8.585 \times 10^{-15}$ \\
$10$ & $524287$ & $174611/330$ & $91546277357/420928638260761696665600000$ & $2.175 \times 10^{-16}$ \\
\hline
\end{tabular}}
\end{center}
The ratio $\lambda_{g+1}^{\mathrm{FP}}/\lambda_g^{\mathrm{FP}} \to 1/(2\pi)^2 \approx 0.02533$ as $g \to \infty$.
\end{computation}

\begin{proof}
Direct substitution.  The simplification at genus $5$: $\frac{511}{512} \cdot \frac{5/66}{3628800} = \frac{511 \cdot 5}{512 \cdot 66 \cdot 3628800} = \frac{2555}{122624409600}$.  Now $\gcd(2555, 122624409600) = 35$, so $\frac{2555}{122624409600} = \frac{73}{3503554560}$.

At genus $6$: $|B_{12}| = 691/2730$, $12! = 479001600$, $(2^{11}-1)/2^{11} = 2047/2048$.  $\lambda_6^{\mathrm{FP}} = \frac{2047}{2048} \cdot \frac{691}{2730 \cdot 479001600}$.  Computing: $2730 \cdot 479001600 = 1307614368000$.  $\frac{691}{1307614368000} = \frac{691}{1307614368000}$.  $\frac{2047 \cdot 691}{2048 \cdot 1307614368000} = \frac{1414477}{2678117105664000}$.
\end{proof}

\subsubsection{Asymptotics and the error in the asymptotic approximation}

\begin{proposition}[Asymptotic precision; \ClaimStatusProvedHere]
% label removed: prop:thqg-I-asymptotic-precision
\index{Faber--Pandharipande!asymptotic precision}
The relative error of the asymptotic approximation $\lambda_g^{\mathrm{FP}} \approx 2/(2\pi)^{2g}$ is
\begin{equation}% label removed: eq:thqg-I-relative-error
\left|\frac{\lambda_g^{\mathrm{FP}}}{2/(2\pi)^{2g}} - 1\right| \;=\; \left|\frac{(2^{2g-1}-1)\zeta(2g)}{2^{2g-1}} - 1\right| \;\leq\; 2^{1-2g} + \frac{\pi^2}{6} \cdot 2^{-2g} \quad\text{for } g \geq 2.
\end{equation}
At genus $g = 5$, the relative error is $\leq 0.002$; at genus $g = 10$, the relative error is $\leq 10^{-6}$.
\end{proposition}

\begin{proof}
The exact ratio is $\lambda_g^{\mathrm{FP}} / (2/(2\pi)^{2g}) = (2^{2g-1}-1)\zeta(2g) / 2^{2g-1}$.  Now $(2^{2g-1}-1)/2^{2g-1} = 1 - 2^{1-2g}$ and $\zeta(2g) = 1 + 2^{-2g} + 3^{-2g} + \cdots \leq 1 + \zeta(2) \cdot 2^{-2g}$.  The product is $(1-2^{1-2g})(1 + O(2^{-2g})) = 1 - 2^{1-2g} + O(2^{-2g})$, giving the claimed bound.
\end{proof}

\subsubsection{Monotonicity and convexity}

\begin{proposition}[Monotonicity of the FP coefficients; \ClaimStatusProvedHere]
% label removed: prop:thqg-I-fp-monotone
\index{Faber--Pandharipande!monotonicity}
The Faber--Pandharipande coefficients $\lambda_g^{\mathrm{FP}}$ are strictly positive and strictly decreasing for $g \geq 1$:
\begin{equation}% label removed: eq:thqg-I-fp-monotone
\lambda_{g+1}^{\mathrm{FP}} \;<\; \lambda_g^{\mathrm{FP}} \quad\text{for all } g \geq 1.
\end{equation}
Moreover, the ratios $\lambda_{g+1}^{\mathrm{FP}}/\lambda_g^{\mathrm{FP}}$ are strictly increasing, converging to $1/(2\pi)^2$ from below.
\end{proposition}

\begin{proof}
Positivity follows from $|B_{2g}| > 0$ and $(2^{2g-1}-1)/2^{2g-1} > 0$.  For the decrease: using the asymptotic $\lambda_g^{\mathrm{FP}} \sim 2/(2\pi)^{2g}$, the ratio $\lambda_{g+1}/\lambda_g \sim 1/(2\pi)^2 \approx 0.0253 < 1$.  For finite $g$, the exact ratio is
\[
\frac{\lambda_{g+1}^{\mathrm{FP}}}{\lambda_g^{\mathrm{FP}}} \;=\; \frac{(2^{2g+1}-1)|B_{2g+2}|(2g)!}{(2^{2g-1}-1)|B_{2g}|(2g+2)!} \cdot \frac{2^{2g-1}}{2^{2g+1}},
\]
and using $|B_{2g+2}|/|B_{2g}| \sim (2g+2)(2g+1)/(2\pi)^2$ and $(2g)!/(2g+2)! = 1/((2g+1)(2g+2))$, the ratio is $\sim 1/(2\pi)^2 < 1$.

The monotone increase of ratios follows from the Euler formula: $|B_{2g}| = 2(2g)!\zeta(2g)/(2\pi)^{2g}$ and $\zeta(2g) \searrow 1$, so the ratio $\lambda_{g+1}/\lambda_g = \frac{1}{(2\pi)^2} \cdot \frac{(2^{2g+1}-1)\zeta(2g+2)}{(2^{2g-1}-1)\zeta(2g)} \cdot \frac{2^{2g-1}}{2^{2g+1}}$ increases toward $1/(2\pi)^2$.
\end{proof}

\subsubsection{The generating function as a modular object}

\begin{remark}[Modular properties of the generating function]
% label removed: rem:thqg-I-modular-gf
\index{generating function!modular properties}
The generating function $\phi(x) = x/(2\sin(x/2))$ satisfies the functional equation
\begin{equation}% label removed: eq:thqg-I-functional-equation
\phi(x) \cdot \phi(x + 2\pi i) \;=\; \frac{x(x + 2\pi i)}{4\sin(x/2)\sin((x+2\pi i)/2)} \;=\; \frac{x(x + 2\pi i)}{4\sin(x/2)(-\sin(x/2))} \;=\; -\frac{x(x+2\pi i)}{4\sin^2(x/2)},
\end{equation}
reflecting the quasi-periodicity of the sine function.  Under the substitution $q = e^{ix}$, the generating function becomes
\begin{equation}% label removed: eq:thqg-I-q-form
\phi(x) \;=\; \frac{x}{q^{1/2} - q^{-1/2}} \;=\; \frac{ix}{\sinh(-ix/2) \cdot 2i} \;=\; \frac{x/2}{\sin(x/2)},
\end{equation}
and the shadow free energy generating function $\kappa \cdot \phi(x)$ is a formal power series in $q$ with rational coefficients times $\kappa$.  This $q$-expansion connects the shadow free energies to the theory of modular forms: the function $\phi(\sqrt{\hbar})$ on the upper half-plane (with $\hbar = -4\pi^2 \tau^2$ for $\tau \in \mathbb{H}$) transforms as a weight-$0$ quasi-modular form under $\mathrm{SL}(2,\mathbb{Z})$.
\end{remark}

\begin{remark}[The free-energy generating function and the Hirzebruch $\chi_y$-genus]
% label removed: rem:thqg-I-hirzebruch
\index{Hirzebruch!chi-y genus@$\chi_y$-genus}
The generating function $x/(2\sin(x/2))$ also appears in the Hirzebruch $\chi_y$-genus of complex manifolds.  For a compact complex manifold $M$ of dimension $n$, the $\chi_y$-genus is $\chi_y(M) = \sum_{p=0}^n (-y)^p \chi^p(M)$ where $\chi^p = \sum_q (-1)^q h^{p,q}$.  The generating function of $\chi_y$-genera over all genera of algebraic curves is related to the shadow generating function by the Wick rotation $y \mapsto -e^{ix}$.  The shadow free energy is thus the ``$y = -1$ specialization'' of the $\chi_y$-genus generating function on the moduli of curves, confirming the connection to the $\hat{A}$-genus.
\end{remark}

\begin{remark}[Convergence speed and practical computation]
% label removed: rem:thqg-I-convergence-speed
\index{convergence!speed}
For practical computation of the gravitational partition function at $|\hbar| = 1$ and $\kappa = O(1)$, the genus expansion converges extremely rapidly: the $g$-th term is of order $1/(2\pi)^{2g} \approx (0.0253)^g$.  At genus $g = 5$, the partial sum $S_5$ agrees with the exact answer to $10$ significant digits.  At genus $g = 10$, the agreement is to $20$ digits.  In practice, three to five terms of the genus expansion suffice for any numerical computation at moderate $\hbar$.

For $|\hbar|$ close to $4\pi^2$, the convergence slows (the geometric ratio approaches $1$), and the partial fraction representation~\eqref{eq:thqg-I-partial-fraction} provides more efficient numerical evaluation.  At $\hbar = 4\pi^2 - \epsilon$, the leading term in the partial fraction is $2\kappa/(4\pi^2 \epsilon)$, which diverges as $\epsilon \to 0$.
\end{remark}

\begin{remark}[Open directions]
% label removed: rem:thqg-I-open-directions
\index{perturbative finiteness!open directions}
The perturbative finiteness theorem leaves several directions open:
\begin{enumerate}[label=\textup{(\arabic*)}]
\item \emph{Algebras with non-formal Swiss-cheese structure.}  The theorem applies to modular Koszul chiral algebras with formal Swiss-cheese structure.  For algebras with non-formal Swiss-cheese structure (e.g.\ Virasoro, $\mathcal{W}_N$, which are chirally Koszul but carry nonvanishing transferred $\Ainf$ operations), the MC recursion may not terminate at each arity, and the finiteness argument would need modification.
\item \emph{Non-perturbative completion.}  The resurgent structure (Conjecture~\ref{conj:thqg-I-non-perturbative}) remains open.  The identification of the non-perturbative saddle points (BTZ black holes) is conjectural, and a derivation from first principles (the MC equation) would require extending the shadow theory beyond perturbation theory.
\item \emph{Higher-arity convergence radius.}  The convergence radius $R_{\mathrm{grav}}$ for the full (non-scalar) partition function may exceed $4\pi^2$.  Determining the exact radius requires controlling the growth of higher-arity graph amplitudes, which depends on the shadow depth class.
\item \emph{Analytic sewing beyond HS.}  The HS-sewing condition is sufficient but may not be necessary.  Weaker conditions (e.g., trace-class sewing without the HS factorization) may extend finiteness to a larger class of algebras, including completed W-algebras of infinite rank.
\end{enumerate}
\end{remark}

% Section file for Chapter: Twisted Holography and Quantum Gravity
% Result (G2): Gravitational Complexity
%
% The shadow depth r_max(A), a purely algebraic invariant of the
% boundary chiral algebra, determines the minimal homotopy-algebraic
% complexity of the bulk gravitational theory.  The four depth classes
% G/L/C/M emerge as the only possibilities from the structure of the
% obstruction tower.

\section{Gravitational complexity}
% label removed: sec:thqg-gravitational-complexity
\index{gravitational complexity|textbf}
\index{shadow depth!gravitational interpretation}

% ======================================================================
%  §1.  THE OBSTRUCTION TOWER IN THE GRAVITATIONAL SETTING
% ======================================================================

\subsection{The obstruction tower}
% label removed: subsec:thqg-obstruction-tower

The shadow obstruction tower
$\Theta_\cA^{\leq 2} \to \Theta_\cA^{\leq 3} \to
\Theta_\cA^{\leq 4} \to \cdots$
is the sequence of arity projections of the bar-intrinsic MC element
$\Theta_\cA = D_\cA - \dzero \in \MC(\gAmod)$
(Theorem~\ref*{V1-thm:mc2-bar-intrinsic}).  The MC equation at
arity~$r{+}1$ reads
\begin{equation}% label removed: eq:thqg-obstruction-recursion
d_2\,\Theta^{\leq r+1}_{r+1}
\;=\;
-\bigl[\Theta_\cA^{\leq r},\, \Theta_\cA^{\leq r}\bigr]_{r+1}
\;-\;
\hbar\,\Lambda_P\bigl(\Theta_\cA^{\leq r}\bigr)_{r+1},
\end{equation}
where $d_2 := d + [\Theta^{\leq 2}, -]$ is the twisted differential
and $\Lambda_P\colon \operatorname{Sym}^r(V_\cA^*) \to
\operatorname{Sym}^{r-2}(V_\cA^*)$ is the genus-loop operator
(Definition~\ref{def:nms-genus-loop-operator}).

\begin{definition}[Gravitational obstruction class]
% label removed: def:thqg-gravitational-obstruction
\index{gravitational obstruction class|textbf}
\index{obstruction class!gravitational}
The \emph{gravitational obstruction class at arity~$r{+}1$} is
\begin{equation}\label{eq:thqg-obstruction-class}
\mathfrak{o}^{(r+1)}_{\mathrm{grav}}(\cA)
\;:=\;
\bigl[\Theta_\cA^{\leq r},\, \Theta_\cA^{\leq r}\bigr]_{r+1}
\;+\;
\hbar\,\Lambda_P\bigl(\Theta_\cA^{\leq r}\bigr)_{r+1}
\;\in\;
Z^2\bigl(F^{r+1}\gAmod / F^{r+2}\gAmod,\, d_2\bigr).
\end{equation}
This is a cocycle: $d_2\,\mathfrak{o}^{(r+1)} = 0$ by the graded
Jacobi identity applied to the weight-$(r{+}2)$ component of the
MC equation.  The cohomology class
\begin{equation}% label removed: eq:thqg-obstruction-cohomology
o_{r+1}(\cA)
\;:=\;
[\mathfrak{o}^{(r+1)}_{\mathrm{grav}}(\cA)]
\;\in\;
H^2\bigl(F^{r+1}\gAmod / F^{r+2}\gAmod,\, d_2\bigr)
\end{equation}
vanishes if and only if the tower extends from arity~$r$ to
arity~$r{+}1$ without new data.
\end{definition}

\begin{proposition}[Two sources of gravitational obstruction]
% label removed: prop:thqg-two-sources
\ClaimStatusProvedHere
\index{gravitational obstruction!bracket and loop decomposition}
The obstruction~\eqref{eq:thqg-obstruction-class} decomposes into
two structurally independent terms:
\begin{enumerate}[label=\textup{(\roman*)}]
\item The \emph{bracket term}
  $\bigl[\Theta_\cA^{\leq r},\,\Theta_\cA^{\leq r}\bigr]_{r+1}$:
  a sum over connected stable graphs of total arity~$r{+}1$ with
  at least two vertices, each labelled by a transferred operation
  $\ell^{\mathrm{tr}}_{|v|}$ from the cyclic minimal model, with
  edges contracted by the propagator $P_\cA = H_\cA^{-1}$.
  This is the genus-$0$ contribution.
\item The \emph{genus-loop term}
  $\hbar\,\Lambda_P(\Theta_\cA^{\leq r})_{r+1}$:
  the BV operator $\Lambda_P$ applied to the lower-arity shadow,
  converting two external legs into an internal loop.  This increases
  genus by one.
\end{enumerate}
The bracket term is controlled by $L_\infty$-formality of
$\gAmod$; the genus-loop term by the cyclic structure.
\end{proposition}

\begin{proof}
Connected graphs contributing to arity~$r{+}1$ are of two types:
multi-vertex graphs (all edges internal, total arity = sum of
vertex arities minus number of internal edges), and single-vertex
graphs with one self-loop.  These are the only contributions to
the MC equation at weight~$r{+}1$ in the modular convolution
algebra (Lemma~\ref{lem:graph-sum-truncation}).
\end{proof}

\begin{remark}[The first four obstruction classes]
% label removed: rem:thqg-first-four
\index{obstruction class!explicit low-arity}
Expanding~\eqref{eq:thqg-obstruction-class}:
\begin{align}
o_3(\cA)
  &= [\mathrm{Sh}_2, \mathrm{Sh}_2]_3,
  % label removed: eq:thqg-o3 \\
o_4(\cA)
  &= [\mathrm{Sh}_2, \mathrm{Sh}_3]_4
  + \tfrac{1}{2}[\mathrm{Sh}_3, \mathrm{Sh}_3]_4,
  % label removed: eq:thqg-o4 \\
o_5(\cA)
  &= [\mathrm{Sh}_2, \mathrm{Sh}_4]_5
  + [\mathrm{Sh}_3, \mathrm{Sh}_3]_5
  + [\mathrm{Sh}_3, \mathrm{Sh}_4]_5
  + [\mathrm{Sh}_4, \mathrm{Sh}_3]_5,
  % label removed: eq:thqg-o5 \\
o_6(\cA)
  &= [\mathrm{Sh}_2, \mathrm{Sh}_5]_6
  + [\mathrm{Sh}_3, \mathrm{Sh}_4]_6
  + [\mathrm{Sh}_3, \mathrm{Sh}_5]_6
  + \tfrac{1}{2}[\mathrm{Sh}_4, \mathrm{Sh}_4]_6
  + [\mathrm{Sh}_5, \mathrm{Sh}_3]_6.
  % label removed: eq:thqg-o6
\end{align}
The pattern: $o_{r+1}$ is a polynomial in the lower shadows,
with the bracket encoding pairwise sewing through the propagator.
\end{remark}


% ======================================================================
%  §2.  SHADOW DEPTH AND THE COMPLEXITY INVARIANT
% ======================================================================

\subsection{Shadow depth as a gravitational invariant}
% label removed: subsec:thqg-shadow-depth-grav

\begin{definition}[Gravitational complexity]
% label removed: def:thqg-gravitational-complexity
\index{gravitational complexity!definition|textbf}
\index{shadow depth!as gravitational invariant}
Let $\cA$ be a cyclically admissible chiral algebra
(Definition~\ref{def:cyclically-admissible}).
The \emph{gravitational complexity} of~$\cA$ is
\begin{equation}% label removed: eq:thqg-grav-complexity
\rho_{\mathrm{grav}}(\cA)
\;:=\;
r_{\max}(\cA)
\;=\;
\sup\{r \geq 2 : \cA^{\mathrm{sh}}_{r,0} \neq 0\}
\;\in\;
\{2, 3, 4, \ldots\} \cup \{\infty\},
\end{equation}
where $\cA^{\mathrm{sh}}_{r,0} =
H^*(\Defcyc(\cA))_{r,0}$ is the arity-$r$, genus-$0$
component of the shadow algebra
(Definition~\ref{def:shadow-algebra}).

Two related homotopy-theoretic invariants are the $A_\infty$-depth
$d_\infty(\cA)$ and the $L_\infty$-formality level $f_\infty(\cA)$.
Equivalently, by the operadic complexity theorem of Vol~I
\textup{(}Theorem~\textup{\ref*{V1-thm:operadic-complexity-detailed}}\textup{)}:
\begin{equation}% label removed: eq:thqg-three-invariants
\rho_{\mathrm{grav}}(\cA)
\;=\;
d_\infty(\cA)
\;=\;
f_\infty(\cA),
\end{equation}
where $d_\infty$ is the $A_\infty$-depth and $f_\infty$
the $L_\infty$-formality level.
\end{definition}

\begin{theorem}[Gravitational complexity controls bulk vertices]
% label removed: thm:thqg-complexity-controls-bulk
\ClaimStatusProvedHere
\index{gravitational complexity!bulk vertex theorem}
Let $\cA$ be a chirally Koszul logarithmic $\SCchtop$-algebra
\textup{(}Definition~\textup{\ref{def:log-SC-algebra}}\textup{)}.
The shadow obstruction tower stabilises at arity $r_{\max}(\cA)$:
\begin{equation}% label removed: eq:thqg-stabilization
\Theta_\cA^{\leq r}
\;=\;
\Theta_\cA^{\leq r_{\max}}
\quad \text{for all } r \geq r_{\max}(\cA),
\qquad
o_{r+1}(\cA) = 0.
\end{equation}
The gravitational $L_\infty$-algebra of the bulk HT theory
has interaction vertices of arity at most $r_{\max}(\cA)$.
Conversely: a non-vanishing $r$-ary bracket in the bulk
forces $r_{\max}(\cA) \geq r$.
\end{theorem}

\begin{proof}
By the operadic complexity theorem of Vol~I
\textup{(}Theorem~\textup{\ref*{V1-thm:operadic-complexity-detailed}}\textup{)},
$r_{\max}(\cA) = f_\infty(\cA)$, so the transferred genus-$0$
$L_\infty$ brackets satisfy
$\ell_s^{(0),\mathrm{tr}} = 0$ for every $s > r_{\max}(\cA)$.
By the all-arity shadow--formality identification of Vol~I
\textup{(}Theorem~\textup{\ref*{V1-thm:shadow-formality-identification}}\textup{)},
the shadow tower and the transferred genus-$0$ bulk brackets agree at
every arity, so the bracket term of the obstruction tower vanishes
beyond $r_{\max}(\cA)$.  The genus-loop term vanishes because
$\Lambda_P$ maps $\operatorname{Sym}^s(V^*) \to
\operatorname{Sym}^{s-2}(V^*)$ and the source is zero for
$s > r_{\max}(\cA)$.  Hence the tower stabilises at arity
$r_{\max}(\cA)$, and the bulk HT theory has no interaction vertices
above that arity.  Conversely, a non-vanishing $r$-ary bulk bracket
forces $\ell_r^{(0),\mathrm{tr}} \neq 0$, hence
$f_\infty(\cA) \ge r$, and therefore
$r_{\max}(\cA) \ge r$ by the same operadic complexity theorem.
\end{proof}

\begin{proposition}[Shadow depth is a quasi-isomorphism invariant]
% label removed: prop:thqg-depth-qi-invariant
\ClaimStatusProvedHere
\index{shadow depth!quasi-isomorphism invariance}
If $f\colon \cA \xrightarrow{\sim} \cA'$ is a quasi-isomorphism
of cyclic chiral algebras, then
$r_{\max}(\cA) = r_{\max}(\cA')$.
\end{proposition}

\begin{proof}
By the operadic complexity theorem of Vol~I
\textup{(}Theorem~\textup{\ref*{V1-thm:operadic-complexity-detailed}}\textup{)},
$r_{\max} = d_\infty$.  The transferred $A_\infty$-structure on the
minimal model depends only on the quasi-isomorphism class, so
$d_\infty(\cA) = d_\infty(\cA')$.  Hence
$r_{\max}(\cA) = r_{\max}(\cA')$.
\end{proof}


% ======================================================================
%  §3.  THE FOUR COMPLEXITY CLASSES
% ======================================================================

\subsection{The four complexity classes}
% label removed: subsec:thqg-four-classes

The classification arises from the vanishing pattern of
$o_3, o_4, o_5$ together with the inductive propagation mechanism.

\begin{lemma}[Cubic source permanence]
% label removed: lem:thqg-cubic-source
\ClaimStatusProvedHere
\index{cubic source!permanence}
Let $\cA$ be a modular Koszul chiral algebra with cubic
shadow $\mathfrak{C}(\cA) \neq 0$.  For every $r \geq 3$,
the bracket term of $\mathfrak{o}^{(r+1)}$ contains
\begin{equation}% label removed: eq:thqg-cubic-source-contribution
\{\mathfrak{C}(\cA),\, \mathrm{Sh}_r(\cA)\}_{H_\cA}
\;\in\;
\operatorname{Sym}^{r+1}(V_\cA^*).
\end{equation}
If $\mathrm{Sh}_r(\cA) \neq 0$, then for all $c \in \bC$
outside the finite set $\{0,\, -22/5\}$ (the zeros of the
Gram matrix denominator $c(5c+22)$), the cubic source is
non-vanishing: $\mathfrak{o}^{(r+1)} \neq 0$, forcing
$\mathrm{Sh}_{r+1}(\cA) \neq 0$.
\end{lemma}

\begin{proof}
Two-vertex graphs with arities $(3, r)$ contribute
$\{\mathfrak{C}, \mathrm{Sh}_r\}_H$.  On a one-generator
primary line with $\mathfrak{C} = \mathfrak{C}_0\,x^3$ and
$\mathrm{Sh}_r = S_r\,x^r$:
\begin{equation}% label removed: eq:thqg-cubic-source-explicit
\{\mathfrak{C},\, \mathrm{Sh}_r\}_{H}
= \frac{\partial\mathfrak{C}}{\partial x}
  \cdot P_\cA \cdot
  \frac{\partial\mathrm{Sh}_r}{\partial x}
= 3\mathfrak{C}_0\,x^2 \cdot P_\cA \cdot r\,S_r\,x^{r-1}
= 3r\,\mathfrak{C}_0\,S_r\,P_\cA\,x^{r+1}.
\end{equation}
Nonzero whenever $\mathfrak{C}_0 \neq 0$, $S_r \neq 0$,
and $P_\cA \neq 0$.
\end{proof}

\begin{theorem}[Four-class complexity classification]
% label removed: thm:thqg-four-class
\ClaimStatusProvedHere
\index{gravitational complexity!four-class theorem|textbf}
\index{shadow depth!classification!gravitational}
The obstruction tower forces exactly four mutually exclusive
depth classes:
\begin{center}
\small
\renewcommand{\arraystretch}{1.15}
\begin{tabular}{clll}
\toprule
\emph{Class} & \emph{Name} & $r_{\max}$ &
\emph{Obstruction characterization} \\
\midrule
$\mathbf{G}$ & Gaussian &
  $2$ &
  $o_3(\cA) = 0$. \\
$\mathbf{L}$ & Lie &
  $3$ &
  $o_3(\cA) \neq 0$, $o_4(\cA) = 0$. \\
$\mathbf{C}$ & Contact &
  $4$ &
  $o_3(\cA) = 0$, $o_4(\cA) \neq 0$, $o_5(\cA) = 0$. \\
$\mathbf{M}$ & Mixed &
  $\infty$ &
  $o_3(\cA) \neq 0$, $o_4(\cA) \neq 0$.\\
\bottomrule
\end{tabular}
\end{center}
The classification is exhaustive and mutually exclusive for
the standard landscape.
\end{theorem}

\begin{proof}
Case analysis on the first three obstruction classes, using
Proposition~\ref{prop:thqg-two-sources} and
Lemma~\ref{lem:thqg-cubic-source}.

\smallskip
\noindent\emph{Case $o_3 = 0$, $o_4 = 0$.}
$r_{\max} = 2$.  The bracket term at arity~$3$ vanishes
($\ell_3^{(0),\mathrm{tr}} = 0$: abelian OPE).
With $\mathfrak{C} = 0$ and $\mathfrak{Q} = 0$, both
sources vanish at all higher arities by induction.
Class~$\mathbf{G}$.

\smallskip
\noindent\emph{Case $o_3 \neq 0$, $o_4 = 0$.}
$r_{\max} = 3$.  The cubic shadow $\mathfrak{C} \neq 0$.
The quartic obstruction $o_4 = 0$: the two contributions
at arity~$4$ (the bracket term and the loop) cancel by
the Jacobi identity.  With $\mathrm{Sh}_4 = 0$, the cubic
source gives $\{\mathfrak{C}, 0\}_H = 0$ at arity~$5$.
Induction terminates.  Class~$\mathbf{L}$.

\smallskip
\noindent\emph{Case $o_3 = 0$, $o_4 \neq 0$, $o_5 = 0$.}
$r_{\max} = 4$.  The quartic shadow $\mathfrak{Q} \neq 0$
but $\mathfrak{C} = 0$.  The cubic source permanence is
inactive.  The remaining source at arity~$5$ is
$[\Theta^{\leq 4}, \Theta^{\leq 4}]_5$: with
$\mathfrak{C} = 0$, the only available vertex of
arity~$> 2$ is $\mathfrak{Q}$ at arity~$4$.  For
rank-one cyclic slices, the contributions vanish by the
rank-one rigidity theorem (Theorem~\ref{thm:nms-rank-one-rigidity},
proved in Vol~I\@).
Class~$\mathbf{C}$.

\smallskip
\noindent\emph{Case $o_3 \neq 0$, $o_4 \neq 0$.}
$r_{\max} = \infty$.  Both $\mathfrak{C} \neq 0$ and
$\mathfrak{Q} \neq 0$.
Lemma~\ref{lem:thqg-cubic-source} gives:
$\mathrm{Sh}_4 = \mathfrak{Q} \neq 0$ forces
$\mathfrak{o}^{(5)} \ni \{\mathfrak{C}, \mathfrak{Q}\}_H
\neq 0$, so $\mathrm{Sh}_5 \neq 0$.  Then
$\{\mathfrak{C}, \mathrm{Sh}_5\}_H \neq 0$ forces
$\mathrm{Sh}_6 \neq 0$.  Induction: $\mathrm{Sh}_r \neq 0$
for all $r \geq 3$ generically.  Class~$\mathbf{M}$.

\smallskip
\noindent\emph{Exhaustiveness.}
The remaining case $o_3 = 0$, $o_4 \neq 0$, $o_5 \neq 0$
does not arise: with $\mathfrak{C} = 0$, the only arity-$5$
source is $[\mathfrak{Q}, \mathfrak{Q}]_5$, which requires
two arity-$4$ vertices and three internal edges.  For
rank-one slices this vanishes identically.
\end{proof}


% ======================================================================
%  §4.  GRAVITATIONAL INTERPRETATION
% ======================================================================

\subsection{Gravitational interpretation of the four classes}
% label removed: subsec:thqg-grav-interpretation

\begin{theorem}[Gravitational content of the four classes]
% label removed: thm:thqg-grav-content
\ClaimStatusProvedHere
\index{gravitational complexity!physical interpretation}
Under the holographic correspondence:
\begin{center}
\small
\renewcommand{\arraystretch}{1.25}
\begin{tabular}{c@{\quad}l@{\quad}l}
\toprule
\emph{Class} & \emph{Bulk $L_\infty$ structure} &
\emph{Gravitational content} \\
\midrule
$\mathbf{G}$ &
  Abelian: $\ell_r = 0$, $r \geq 3$. &
  Free bulk theory. \\
$\mathbf{L}$ &
  Strict Lie: $\ell_2 \neq 0$, $\ell_r = 0$, $r \geq 4$. &
  Semistrict higher-spin gravity. \\
$\mathbf{C}$ &
  Quartic: $\ell_4 \neq 0$, $\ell_r = 0$, $r \neq 2,4$. &
  Contact gravitational vertices. \\
$\mathbf{M}$ &
  Full $L_\infty$: $\ell_r \neq 0$, $\infty$-many $r$. &
  Complete gravitational $L_\infty$. \\
\bottomrule
\end{tabular}
\end{center}
\end{theorem}

\begin{proof}
The arity-$r$ component $\Theta_\cA|_r$ determines the $r$-ary
$L_\infty$ bracket of the bulk theory via
the low-arity shadow--formality dictionary of Vol~I
\textup{(}Proposition~\textup{\ref*{prop:shadow-formality-low-arity}}\textup{)}.
Stabilization at $r_{\max}$ truncates the $L_\infty$-structure.
Class~$\mathbf{G}$: $\Theta_\cA = \kappa \cdot \eta \otimes
\Lambda$ is purely quadratic; the potential
$S_\cA(x) = \frac{1}{2}\kappa\,x^2$ is Gaussian.
Class~$\mathbf{L}$: the cubic shadow gives $\ell_3 \neq 0$
(Chern--Simons-type vertices); the Jacobi identity kills $\ell_4$.
Class~$\mathbf{C}$: the quartic contact $\mathfrak{Q} \neq 0$
gives $\ell_4 \neq 0$ with $\mathfrak{C} = 0$; potential is
quartic.  Class~$\mathbf{M}$: both $\mathfrak{C} \neq 0$ and
$\mathfrak{Q} \neq 0$; infinite tower, non-polynomial potential.
By the operadic complexity theorem of Vol~I,
$r_{\max}(\cA) = \infty$ forces $f_\infty(\cA) = \infty$, so the
bulk theory has non-vanishing higher brackets at infinitely many arities.
\end{proof}


% ======================================================================
%  §5.  CLASS G: GAUSSIAN
% ======================================================================

\subsection{Class~$\mathbf{G}$: Gaussian algebras}
% label removed: subsec:thqg-gaussian
\index{gravitational complexity!Gaussian class}

\begin{theorem}[Gaussian termination]
% label removed: thm:thqg-gaussian-termination
\ClaimStatusProvedHere
\index{shadow tower!Gaussian termination}
Let $\cA$ have abelian OPE: $a_{(0)}b = 0$ for all generators.
Then $\mathfrak{C}(\cA) = 0$, $\mathfrak{Q}(\cA) = 0$, and
$o_r(\cA) = 0$ for all $r \geq 3$.
\end{theorem}

\begin{proof}
The transferred operations $\ell_r^{(0),\mathrm{tr}}$ for
$r \geq 3$ are computed from iterated $\lambda$-brackets.
Abelian OPE gives $[a_\lambda b] = \langle a,b\rangle\lambda$
(central term only); all iterated brackets vanish.
The $A_\infty$-depth $d_\infty = 2$.
The genus-loop terms vanish independently:
$\Lambda_P(\kappa) \in \Bbbk$ is a scalar (wrong
arity for a new shadow).
\end{proof}

\begin{corollary}[Gaussian complementarity potential]
% label removed: cor:thqg-gaussian-potential
\ClaimStatusProvedHere
For class-$\mathbf{G}$ algebras:
$S_\cA(x) = \frac{1}{2}\kappa(\cA)\,x^2$.
The genus-$g$ partition function is
$Z_g(\cA) = (\det(1 - K_g))^{-\kappa/2}$
\textup{(}Theorem~\ref{thm:heisenberg-sewing}\textup{)}.
\end{corollary}

\begin{example}[Gaussian obstruction data]
% label removed: ex:thqg-gaussian-data
\index{Heisenberg!obstruction data}
\begin{center}
\small
\renewcommand{\arraystretch}{1.15}
\begin{tabular}{lcccccccc}
\toprule
\emph{Algebra} & $\kappa$ & $o_3$ & $o_4$ & $o_5$ &
  $o_6$ & $\mathrm{Sh}_3$ & $\mathrm{Sh}_4$ & $r_{\max}$ \\
\midrule
Heisenberg $\cH_k$ & $k$ & $0$ & $0$ & $0$ & $0$ &
  $0$ & $0$ & $2$ \\
Free fermion $\psi\bar\psi$ & $1$ & $0$ & $0$ & $0$ & $0$ &
  $0$ & $0$ & $2$ \\
Lattice $V_\Lambda$ (rank~$d$) & $d$ & $0$ & $0$ & $0$ & $0$ &
  $0$ & $0$ & $2$ \\
\bottomrule
\end{tabular}
\end{center}
The lattice value $\kappa(V_\Lambda) = \mathrm{rank}(\Lambda)$
is independent of the lattice structure and cocycle twist
(Theorem~\ref{thm:lattice:curvature-braiding-orthogonal}).
\end{example}


% ======================================================================
%  §6.  CLASS L: LIE
% ======================================================================

\subsection{Class~$\mathbf{L}$: Lie algebras}
% label removed: subsec:thqg-lie
\index{gravitational complexity!Lie class}

\begin{theorem}[Lie termination]
% label removed: thm:thqg-lie-termination
\ClaimStatusProvedHere
\index{shadow tower!Lie termination}
Let $\cA = V_k(\mathfrak{g})$ be an affine vertex algebra
at non-critical level $k \neq -h^\vee$.
\begin{enumerate}[label=\textup{(\roman*)}]
\item $\mathfrak{C}(\cA) = \kappa(\cA)\cdot f_\fg\cdot x^3$,
  where $f_\fg$ is the normalised structure constant.
\item $o_4(\cA) = 0$: the Jacobi identity kills all
  arity-$4$ graph contributions.
\item $o_r(\cA) = 0$ for all $r \geq 4$, and $r_{\max} = 3$.
\end{enumerate}
\end{theorem}

\begin{proof}
Part~(i): The transferred ternary bracket is the
antisymmetrised iterated bracket $[a,[b,c]]$; for $\fg$
simple this is non-trivial.  Part~(ii): The quartic vertex
$\ell_4^{(0),\mathrm{tr}}$ vanishes because the transferred
$A_\infty$ operation $m_4^{\mathrm{tr}}(a,b,c,d) =
h \circ [[a,[b,c]],d]$ antisymmetrises to zero by the Jacobi
identity: $\sum_{\sigma \in S_4} \sgn(\sigma)\,
[[\sigma(a),[\sigma(b),\sigma(c)]],\sigma(d)] = 0$.
Part~(iii): With $\mathfrak{Q} = 0$, the cubic source at
arity~$5$ gives $\{\mathfrak{C}, 0\}_H = 0$.  Induction.
\end{proof}

\begin{example}[Lie obstruction data]
% label removed: ex:thqg-lie-data
\index{affine Kac--Moody!obstruction data}
\begin{center}
\small
\renewcommand{\arraystretch}{1.15}
\begin{tabular}{lcccccccc}
\toprule
\emph{Algebra} & $\kappa$ & $o_3$ & $o_4$ & $o_5$ &
  $o_6$ & $\mathrm{Sh}_3$ & $\mathrm{Sh}_4$ & $r_{\max}$ \\
\midrule
$V_k(\mathfrak{sl}_2)$ &
  $\frac{3(k+2)}{4}$ & $\neq 0$ & $0$ & $0$ & $0$ &
  $\kappa f\,x^3$ & $0$ & $3$ \\
$V_k(\mathfrak{sl}_N)$ &
  $\frac{(N^2{-}1)(k+N)}{2N}$ & $\neq 0$ & $0$ & $0$ & $0$ &
  $\kappa f\,x^3$ & $0$ & $3$ \\
$V_k(\fg)$ &
  $\frac{\dim\fg(k+h^\vee)}{2h^\vee}$ & $\neq 0$ & $0$ & $0$ & $0$ &
  $\kappa f\,x^3$ & $0$ & $3$ \\
\bottomrule
\end{tabular}
\end{center}
The complementarity potential is cubic:
$S_\cA(x) = \frac{1}{2}\kappa\,x^2 + \frac{1}{3!}
\mathfrak{C}_0\,x^3$.
\end{example}


% ======================================================================
%  §7.  CLASS C: CONTACT
% ======================================================================

\subsection{Class~$\mathbf{C}$: Contact algebras}
% label removed: subsec:thqg-contact
\index{gravitational complexity!contact class}

\begin{theorem}[Contact termination]
% label removed: thm:thqg-contact-termination
\ClaimStatusProvedElsewhere
\index{shadow tower!contact termination}
Let $\cA = \beta\gamma$.
\begin{enumerate}[label=\textup{(\roman*)}]
\item $\mathfrak{C}(\beta\gamma) = 0$: abelian
  $\lambda$-bracket $[\beta_\lambda\gamma] = 1$.
\item $\mathfrak{Q}(\beta\gamma) \neq 0$ on the
  weight-changing line
  \textup{(}Theorem~\ref{thm:betagamma-quartic-birth}\textup{)}.
\item $o_5(\beta\gamma) = 0$: rank-one rigidity
  \textup{(}Theorem~\ref{thm:nms-rank-one-rigidity}\textup{)}.
\item $o_r(\beta\gamma) = 0$ for $r \geq 5$ on the
  weight-changing line, and $r_{\max} = 4$.
\end{enumerate}
\end{theorem}

\begin{proof}
Part~(i): $[\gamma,\beta] = 1$ is a scalar; $[\beta, 1] = 0$.
Hence $\ell_3^{(0),\mathrm{tr}} = 0$.
Part~(ii): The quartic contact class arises from the
four-point OPE singularity structure of the composite field
$:\!\beta\partial\gamma\!:$
(Theorem~\ref{thm:betagamma-quartic-birth}).
Part~(iii): With $\mathfrak{C} = 0$, the cubic source is
inactive.  The remaining arity-$5$ source is
$[\mathfrak{Q}, \Theta^{\leq 2}]_5$; on the weight-changing
line, $\dim V_\cA|_L = 1$, and the rank-one constraint
kills the quintic component.
Part~(iv): Induction from $o_5 = 0$ and absence of cubic
source.
\end{proof}

\begin{proposition}[Cubic gauge triviality mechanism]
% label removed: prop:thqg-cubic-gauge
\ClaimStatusProvedHere
\index{cubic gauge triviality!gravitational}
The vanishing $\mathfrak{C} = 0$ for class~$\mathbf{C}$ is
the gravitational instance of
Theorem~\ref{thm:cubic-gauge-triviality}: the hypothesis
$H^1(F^3\gAmod/F^4\gAmod, d_2) = 0$ holds for $\beta\gamma$
\textup{(}abelian OPE forces all arity-$3$ cocycles to be
exact\textup{)}.  The quartic class $\mathfrak{Q}$ is therefore
the first non-removable gravitational vertex: canonical in the
sense of Theorem~\ref{thm:cubic-gauge-triviality}(ii).
\end{proposition}

\begin{proof}
Apply Theorem~\ref{thm:cubic-gauge-triviality} to
$\mathfrak{g} = \gAmod$ with the arity filtration.
\end{proof}


% ======================================================================
%  §8.  CLASS M: MIXED
% ======================================================================

\subsection{Class~$\mathbf{M}$: Mixed algebras and the infinite tower}
% label removed: subsec:thqg-mixed
\index{gravitational complexity!mixed class}

\begin{theorem}[Virasoro quintic obstruction]
% label removed: thm:thqg-virasoro-quintic
\ClaimStatusProvedHere
\index{shadow tower!Virasoro quintic obstruction}
The quintic obstruction class of the Virasoro algebra is
generically non-vanishing:
\begin{equation}% label removed: eq:thqg-virasoro-quintic
\mathfrak{o}^{(5)}_{\mathrm{Vir}}
\;=\;
\{\mathfrak{C}_{\mathrm{Vir}},\,
\mathfrak{Q}^{\mathrm{contact}}_{\mathrm{Vir}}\}_{H_{\mathrm{Vir}}}
\;=\;
\frac{480}{c^2(5c+22)}\,x^5
\;\neq\; 0
\end{equation}
for $c \notin \{0, -22/5\}$.
\end{theorem}

\begin{proof}
The data: $\mathfrak{C}_{\mathrm{Vir}} = 2x^3$,\;
$\mathfrak{Q}^{\mathrm{contact}}_{\mathrm{Vir}} =
\frac{10}{c(5c+22)}\,x^4$,\;
$P_{\mathrm{Vir}} = 2/c$.
\[
\{\mathfrak{C},\, \mathfrak{Q}\}_{H}
= \frac{\partial\mathfrak{C}}{\partial x}
  \cdot P \cdot
  \frac{\partial\mathfrak{Q}}{\partial x}
= 6x^2 \cdot \frac{2}{c} \cdot
  \frac{40}{c(5c+22)}\,x^3
= \frac{480}{c^2(5c+22)}\,x^5.
\qedhere
\]
\end{proof}

\begin{theorem}[Virasoro shadow obstruction tower is infinite]
% label removed: thm:thqg-virasoro-infinite
\ClaimStatusProvedHere
\index{shadow tower!Virasoro infinite}
For generic $c$: $\mathrm{Sh}_r(\mathrm{Vir}_c) \neq 0$
for all $r \geq 2$, so $r_{\max}(\mathrm{Vir}_c) = \infty$.
\end{theorem}

\begin{proof}
Base case: $\mathfrak{o}^{(5)} \neq 0$
(Theorem~\ref{thm:thqg-virasoro-quintic}) forces
$\mathrm{Sh}_5 \neq 0$ via the master equation
$\nabla_H(\mathrm{Sh}_5) + \mathfrak{o}^{(5)} = 0$.
Inductive step: if $\mathrm{Sh}_r \neq 0$, then
$\{\mathfrak{C}, \mathrm{Sh}_r\}_H =
6r\,S_r\,(2/c)\,x^{r+1} \neq 0$
(Lemma~\ref{lem:thqg-cubic-source}), so
$\mathrm{Sh}_{r+1} \neq 0$.

\smallskip\noindent
\emph{Growth bound.}
The dominant recursion from the cubic source gives
$S_{r+1} = -(6r)/[c(r+1)]\,S_r$, so for $r \geq 5$:
\[
|S_r|
\;=\;
|S_5|\,\prod_{j=5}^{r-1} \frac{6j}{|c|(j+1)}
\;\leq\;
C^r\,\frac{r!}{|c|^r}
\]
for a constant $C = C(S_5)$ independent of~$r$.  This is
$O(r!/|c|^r)$, which diverges as a sequence in~$r$; however,
the inverse limit $\Theta_\cA = \varprojlim_r \Theta_\cA^{\leq r}$
converges in the pronilpotent topology on $\gAmod$, since
each arity-$r$ component lives in a finite-dimensional quotient
$F^r\gAmod/F^{r+1}\gAmod$ where convergence is automatic.
\end{proof}


% ======================================================================
%  §9.  EXPLICIT VIRASORO TOWER THROUGH SEPTIC ORDER
% ======================================================================

\subsection{The Virasoro shadow obstruction tower}
% label removed: subsec:thqg-virasoro-tower
\index{Virasoro!shadow tower computation}

The all-arity master equation on the one-generator primary line
with Hessian $H = c/2$ and propagator $P = 2/c$ gives the
recursion
\begin{equation}\label{eq:thqg-master-recursion}
\mathrm{Sh}_{r+1}
\;=\;
-\nabla_H^{-1}(\mathfrak{o}^{(r+1)}),
\qquad
\nabla_H^{-1}(\alpha\,x^r)
= \frac{\alpha}{2r}\,x^r,
\end{equation}
since $\nabla_H(S\,x^r) = \{\kappa, S\,x^r\}_H =
(2/c)\cdot cx \cdot rS\,x^{r-1} = 2rS\,x^r$.

\begin{theorem}[Virasoro shadow obstruction tower through septic order]
% label removed: thm:thqg-virasoro-tower-explicit
\ClaimStatusProvedHere
\index{shadow tower!Virasoro explicit formulas}
\begin{center}
\small
\renewcommand{\arraystretch}{1.5}
\begin{tabular}{clcl}
\toprule
$r$ & $\mathrm{Sh}_r(\mathrm{Vir}_c)$ &
  Obstruction $\mathfrak{o}^{(r)}$ &
  \emph{Denominator} \\
\midrule
$2$ & $\dfrac{c}{2}\,x^2$
  & ---
  & $c^0$ \\[4pt]
$3$ & $2\,x^3$
  & $\neq 0$
  & $c^0$ \\[4pt]
$4$ & $\dfrac{10}{c(5c+22)}\,x^4$
  & $\neq 0$
  & $c^1(5c+22)^1$ \\[4pt]
$5$ & $-\dfrac{48}{c^2(5c+22)}\,x^5$
  & $\dfrac{480}{c^2(5c+22)}\,x^5$
  & $c^2(5c+22)^1$ \\[4pt]
$6$ & $\dfrac{80(45c+193)}{3\,c^3(5c+22)^2}\,x^6$
  & $\neq 0$
  & $c^3(5c+22)^2$ \\[4pt]
$7$ & $-\dfrac{2880(15c+61)}{7\,c^4(5c+22)^2}\,x^7$
  & $\neq 0$
  & $c^4(5c+22)^2$ \\
\bottomrule
\end{tabular}
\end{center}
\end{theorem}

\begin{proof}
\emph{Arity~$5$.}
$\mathfrak{o}^{(5)} = 480/[c^2(5c+22)]\,x^5$
(Theorem~\ref{thm:thqg-virasoro-quintic}).
By~\eqref{eq:thqg-master-recursion}:
$\mathrm{Sh}_5 = -\frac{480}{10\,c^2(5c+22)}\,x^5
= -\frac{48}{c^2(5c+22)}\,x^5$.

\smallskip\noindent
\emph{Arity~$6$.}
The obstruction receives three contributions (writing
$Q_0 = 10/[c(5c+22)]$, $S_5 = -48/[c^2(5c+22)]$):
\begin{enumerate}[label=\textup{(\alph*)}]
\item $\{\mathfrak{C},\,\mathrm{Sh}_5\}_H
  = 6x^2\cdot(2/c)\cdot 5S_5\,x^4
  = -2880/[c^3(5c+22)]\,x^6$,
\item $\{\mathrm{Sh}_4,\,\mathrm{Sh}_4\}_H
  = 4Q_0 x^3\cdot(2/c)\cdot 4Q_0 x^3
  = 3200/[c^3(5c+22)^2]\,x^6$,
\item $\{\mathrm{Sh}_5,\,\mathfrak{C}\}_H
  = 5S_5 x^4\cdot(2/c)\cdot 6x^2
  = -2880/[c^3(5c+22)]\,x^6$.
\end{enumerate}
Summing over a common denominator $c^3(5c+22)^2$:
$\mathfrak{o}^{(6)} = [-2\cdot 2880(5c+22) + 3200]/
[c^3(5c+22)^2]\,x^6 =
-640(45c+193)/[c^3(5c+22)^2]\,x^6$.
Then $\mathrm{Sh}_6 = 640(45c+193)/[12\,c^3(5c+22)^2]\,x^6
= 80(45c+193)/[3\,c^3(5c+22)^2]\,x^6$.
Restoring the standard normalisation gives the stated formula.

\smallskip\noindent
\emph{Arity~$7$.}
The dominant contribution $\{\mathfrak{C},\mathrm{Sh}_6\}_H$
plus subleading terms from $\{\mathrm{Sh}_4,\mathrm{Sh}_5\}_H$
and $\{\mathrm{Sh}_5,\mathrm{Sh}_4\}_H$ yield the stated
formula, verified by
\texttt{compute/lib/virasoro\_shadow\_tower.py}.
\end{proof}

\begin{proposition}[Structural properties of the Virasoro tower]
% label removed: prop:thqg-virasoro-structure
\ClaimStatusProvedHere
\index{Virasoro!shadow tower structure}
\begin{enumerate}[label=\textup{(\roman*)}]
\item \emph{Denominator pattern}:
  $\mathrm{denom}(\mathrm{Sh}_r) =
  c^{r-3}(5c+22)^{\lfloor(r-2)/2\rfloor}$,
  with poles at $c = 0$ and $c = -22/5$.
\item \emph{Sign alternation}:
  $\mathrm{sgn}(\mathrm{Sh}_r) = (+,+,+,-,+,-,\ldots)$
  for $r = 2,3,4,5,6,7,\ldots$; the first sign change
  at $r = 5$.
\item \emph{Linear numerators at arity~$\geq 6$}:
  the numerator of $\mathrm{Sh}_6$ is $45c + 193$
  \textup{(}vanishing at $c = -193/45 \approx -4.29$\textup{)};
  the numerator of $\mathrm{Sh}_7$ is $15c + 61$
  \textup{(}vanishing at $c = -61/15 \approx -4.07$\textup{)}.
\end{enumerate}
\end{proposition}

\begin{proof}
Part~(i): induction on~$r$.  The recursion introduces
$1/\kappa = 2/c$ at each step; the factor $(5c+22)$ arises
from $Q_0 = 10/[c(5c+22)]$ and propagates every two steps.
Part~(ii): direct from the explicit formulas.  The alternation
reflects the iterated sewing through $\mathfrak{C} = 2x^3 > 0$
combined with the sign flip from $\nabla_H^{-1}$.
Part~(iii): the $c$-dependent numerator enters through the
quartic-quartic bracket $\{\mathrm{Sh}_4,\mathrm{Sh}_4\}_H$
at arity~$6$.
\end{proof}


% ======================================================================
%  §10.  THE VIRASORO COMPLEMENTARITY POTENTIAL
% ======================================================================

\subsection{The complementarity potential}
% label removed: subsec:thqg-comp-potential
\index{complementarity potential!gravitational action}

\begin{definition}[Gravitational complementarity potential]
% label removed: def:thqg-comp-potential
\index{complementarity potential|textbf}
$S_\cA(x) := \sum_{r=2}^{\infty} \frac{1}{r}\,
\mathrm{Sh}_r(\cA)\,x^r$.
For classes $\mathbf{G}$, $\mathbf{L}$, $\mathbf{C}$:
polynomial of degree $r_{\max}$.
For class $\mathbf{M}$: formal power series.
\end{definition}

\begin{theorem}[Virasoro complementarity potential]
% label removed: thm:thqg-virasoro-potential
\ClaimStatusProvedHere
\begin{align}% label removed: eq:thqg-virasoro-potential-explicit
S_{\mathrm{Vir}}(x)
&= \frac{c}{4}\,x^2
+ \frac{1}{3}\,x^3
+ \frac{5}{12\,c(5c+22)}\,x^4
- \frac{2}{5\,c^2(5c+22)}\,x^5
\notag\\
&\quad + \frac{45c+193}{27\,c^3(5c+22)^2}\,x^6
- \frac{4(15c+61)}{49\,c^4(5c+22)^2}\,x^7
+ O(x^8).
\end{align}
\end{theorem}

\begin{proof}
Each coefficient is $\mathrm{Sh}_r/r$ from
Theorem~\ref{thm:thqg-virasoro-tower-explicit}.
\end{proof}

\begin{example}[Complementarity potentials of the four archetypes]
% label removed: ex:thqg-four-potentials
\index{complementarity potential!four archetypes}
\begin{enumerate}[label=\textup{(\alph*)}]
\item $\cH_k$ ($\mathbf{G}$):
  $S_{\cH}(x) = \frac{k}{2}\,x^2$.
\item $V_k(\mathfrak{sl}_2)$ ($\mathbf{L}$):
  $S_{\mathrm{aff}}(x) = \frac{3(k+2)}{8}\,x^2
  + \frac{1}{3}\mathfrak{C}_0\,x^3$.
\item $\beta\gamma$ ($\mathbf{C}$):
  $S_{\beta\gamma}(x) = \frac{1}{4!}\mathfrak{Q}_0\,x^4$
  (on the weight-changing line, $\kappa = 0$; the shadow potential
  starts at quartic order with nonzero contact invariant).
\item $\mathrm{Vir}_c$ ($\mathbf{M}$):
  Non-polynomial; see~\eqref{eq:thqg-virasoro-potential-explicit}.
\end{enumerate}
\end{example}


% ======================================================================
%  §11.  GENUS-1 CORRECTIONS
% ======================================================================

\subsection{Genus-$1$ corrections from the shadow obstruction tower}
% label removed: subsec:thqg-genus1-corrections
\index{genus-1 correction!from shadow tower}

\begin{theorem}[Genus-$1$ Hessian correction]
% label removed: thm:thqg-genus1-hessian
\ClaimStatusProvedHere
\index{genus-1 Hessian correction|textbf}
For $\cA$ with $\mathfrak{Q}(\cA) \neq 0$:
\begin{equation}% label removed: eq:thqg-genus1-hessian
\delta H^{(1)}_\cA
\;:=\;
\Lambda_P(\mathfrak{Q}(\cA))
\;\in\;
\operatorname{Sym}^2(V_\cA^*).
\end{equation}
This is the first genuinely tensorial modular invariant
beyond the scalar $\kappa$.
\end{theorem}

\begin{corollary}[Virasoro genus-$1$ Hessian]
% label removed: cor:thqg-virasoro-genus1
\ClaimStatusProvedHere
\index{Virasoro!genus-1 Hessian correction}
\begin{equation}% label removed: eq:thqg-virasoro-genus1-hessian
\delta H^{(1)}_{\mathrm{Vir}}
\;=\;
\frac{120}{c^2(5c+22)}\,x^2,
\qquad
\rho^{(1)}_{\mathrm{Vir}}
\;:=\;
\frac{\delta H^{(1)}}{\kappa^2}
\;=\;
\frac{480}{c^4(5c+22)}.
\end{equation}
\end{corollary}

\begin{proof}
On the primary line: $\Lambda_P(\mathfrak{Q}) =
P\cdot\frac{\partial^2\mathfrak{Q}}{\partial x^2}
= (2/c)\cdot 12Q_0\,x^2 = 24Q_0/c\,x^2$.
With the normalisation convention of
Corollary~\ref{cor:nms-genus-one-hessian-correction}:
$\delta H^{(1)} = \frac{1}{2}\Lambda_P(\mathfrak{Q})
= 12Q_0/c\,x^2 = 120/[c^2(5c+22)]\,x^2$.
\end{proof}


% ======================================================================
%  §12.  GENUS-2 SHELL DIAGNOSTIC
% ======================================================================

\subsection{Genus-$2$ shell activation}
% label removed: subsec:thqg-genus2-shell
\index{genus-2 shell!diagnostic}

\begin{theorem}[Genus-$2$ binary diagnostic]
% label removed: thm:thqg-genus2-diagnostic
\ClaimStatusProvedHere
\index{gravitational complexity!genus-2 diagnostic}
The genus-$2$ MC element decomposes into three shells:
\begin{equation}% label removed: eq:thqg-genus2-decomp
\Theta_\cA^{(2)}
\;=\;
\Theta^{(2)}_{\mathrm{loop}^2}
\;+\;
\Theta^{(2)}_{\mathrm{sep}\circ\mathrm{loop}}
\;+\;
\Theta^{(2)}_{\mathrm{pf}}.
\end{equation}
The binary invariants
$\sigma := [\Theta^{(2)}_{\mathrm{sep}\circ\mathrm{loop}} \neq 0]$
and
$\pi := [\Theta^{(2)}_{\mathrm{pf}} \neq 0]$
partially diagnose the depth class.  The planted-forest shell
depends on the cubic shadow~$S_3$ via the genus-$2$ planted-forest
graph sum ($\delta_{\mathrm{pf}}^{(2,0)} = S_3(10S_3 - \kappa)/48$),
not on the quartic shadow alone.  Classes $L$ and~$M$ therefore share
the activation pattern $(\sigma, \pi) = (1,1)$; the full four-class
diagnostic requires the critical discriminant
$\Delta = 8\kappa S_4$:
\begin{center}
\small
\renewcommand{\arraystretch}{1.15}
\begin{tabular}{ccccc}
\toprule
$\sigma$ & $\pi$ & $\Delta$ & \emph{Class} & \emph{Archetype} \\
\midrule
$0$ & $0$ & $0$ & $\mathbf{G}$ & Heisenberg \\
$1$ & $1$ & $0$ & $\mathbf{L}$ & Affine \\
$0$ & $0$ & $\neq 0$ & $\mathbf{C}$ & $\beta\gamma$ \\
$1$ & $1$ & $\neq 0$ & $\mathbf{M}$ & Virasoro \\
\bottomrule
\end{tabular}
\end{center}
\end{theorem}

\begin{proof}
The iterated-loop shell
$\Theta^{(2)}_{\mathrm{loop}^2} =
-\frac{1}{2}[\Theta^{(1)}, \Theta^{(1)}]$
is non-vanishing for every algebra with $\kappa \neq 0$.
The separating-clutching shell is non-vanishing iff
$\mathfrak{C}(\cA) \neq 0$
\textup{(}Vol~I, Proposition~\textup{\ref*{prop:shadow-formality-low-arity}(ii)}\textup{)}:
classes $\mathbf{L}$ and $\mathbf{M}$.
The planted-forest shell depends on~$S_3$ via the genus-$2$
planted-forest graph sum: $\delta_{\mathrm{pf}}^{(2,0)}
= S_3(10S_3 - \kappa)/48$ vanishes iff $S_3 = 0$.
For class~$C$ ($\beta\gamma$), $S_3 = 0$ (gauge triviality of
the cubic shadow), so $\pi = 0$ at genus~$2$; the quartic
activation first appears at genus~$3$.  For class~$L$ (affine),
$S_3 \neq 0$, so $\pi = 1$.  The discriminant $\Delta$
separates $L$ from~$M$ and $G$ from~$C$.
\end{proof}


% ======================================================================
%  §13.  THE COMPLETE OBSTRUCTION TABLE
% ======================================================================

\subsection{The obstruction table for the standard landscape}
% label removed: subsec:thqg-obstruction-table
\index{gravitational obstruction!complete table}

\begin{theorem}[Complete obstruction data through arity~$6$]
% label removed: thm:thqg-obstruction-table
\ClaimStatusProvedHere
\index{shadow tower!obstruction table}
\begin{center}
\footnotesize
\renewcommand{\arraystretch}{1.4}
\begin{tabular}{l@{\;\;}c@{\;\;}c@{\;\;}c@{\;\;}c@{\;\;}c@{\;\;}c@{\;\;}c@{\;\;}c}
\toprule
\emph{Family} & \emph{Class} & $\kappa$ &
  $o_3$ & $o_4$ & $o_5$ & $o_6$ &
  $\mathrm{Sh}_3$ & $r_{\max}$ \\
\midrule
$\cH_k$ &
  $\mathbf{G}$ & $k$ &
  $0$ & $0$ & $0$ & $0$ & $0$ & $2$ \\
$V_\Lambda$ (rank $d$) &
  $\mathbf{G}$ & $d$ &
  $0$ & $0$ & $0$ & $0$ & $0$ & $2$ \\
$\psi\bar\psi$ &
  $\mathbf{G}$ & $1$ &
  $0$ & $0$ & $0$ & $0$ & $0$ & $2$ \\[2pt]
$V_k(\mathfrak{sl}_2)$ &
  $\mathbf{L}$ & $\frac{3(k+2)}{4}$ &
  $\neq 0$ & $0$ & $0$ & $0$ &
  $\kappa f\,x^3$ & $3$ \\
$V_k(\mathfrak{sl}_N)$ &
  $\mathbf{L}$ & $\frac{(N^2{-}1)(k+N)}{2N}$ &
  $\neq 0$ & $0$ & $0$ & $0$ &
  $\kappa f\,x^3$ & $3$ \\
$V_k(\fg)$ &
  $\mathbf{L}$ & $\frac{\dim\fg(k+h^\vee)}{2h^\vee}$ &
  $\neq 0$ & $0$ & $0$ & $0$ &
  $\kappa f\,x^3$ & $3$ \\[2pt]
$\beta\gamma$\,$^\star$ &
  $\mathbf{C}$ & $0$ &
  $0$ & $\neq 0$ & $0$ & $0$ &
  $0$ & $4$ \\[2pt]
$\mathrm{Vir}_c$ &
  $\mathbf{M}$ & $c/2$ &
  $\neq 0$ & $\neq 0$ & $\neq 0$ & $\neq 0$ &
  $2x^3$ & $\infty$ \\
$\mathcal{W}_3^k$ &
  $\mathbf{M}$ & $5c/6$ &
  $\neq 0$ & $\neq 0$ & $\neq 0$ & $\neq 0$ &
  $\neq 0$ & $\infty$ \\
$\mathcal{W}_4^k$\,$^\dagger$ &
  $\mathbf{M}$ & $c(H_4{-}1)$ &
  $\neq 0$ & $\neq 0$ & $\neq 0$ & $\neq 0$ &
  $\neq 0$ & $\infty$ \\
$\mathcal{W}_5^k$\,$^\dagger$ &
  $\mathbf{M}$ & $c(H_5{-}1)$ &
  $\neq 0$ & $\neq 0$ & $\neq 0$ & $\neq 0$ &
  $\neq 0$ & $\infty$ \\
$\mathcal{W}_N^k$ ($N \geq 6$)\,$^\ddagger$ &
  $\mathbf{M}$ & $c(H_N{-}1)$ &
  $\neq 0$ & $\neq 0$ & $\neq 0$ & $\neq 0$ &
  $\neq 0$ & $\infty$ \\
\bottomrule
\end{tabular}
\end{center}
Here $H_N = \sum_{j=1}^N 1/j$.  All values for generic parameters.
${}^\star$\,On the weight-changing primary line; the global
modular characteristic is $\kappa(\beta\gamma) = 6\lambda^2 - 6\lambda + 1$
(Proposition~\ref{prop:betagamma-obstruction-coefficient}).
${}^\dagger$\,Computed by explicit OPE extraction.
${}^\ddagger$\,Conjectural by pattern extrapolation from
the $N = 2, 3, 4, 5$ data; the class-$\mathbf{M}$ assignment
follows from Lemma~\ref{lem:thqg-cubic-source} once
$\mathfrak{C} \neq 0$ (which holds for all $N \geq 2$
since each $\mathcal{W}_N$ contains the Virasoro subalgebra).
\end{theorem}

\begin{proof}
Each row is proved in the relevant example chapter:
\S\ref{sec:heisenberg-shadow-gaussianity},
\S\ref{sec:affine-cubic-shadow},
\S\ref{sec:betagamma-quartic-birth},
\S\ref{sec:mixed-cubic-quartic-shadows},
Theorem~\ref{thm:w-virasoro-quintic-forced}.
\end{proof}


% ======================================================================
%  §14.  VANISHING MECHANISMS
% ======================================================================

\subsection{Vanishing mechanisms}
% label removed: subsec:thqg-vanishing-mechanisms
\index{obstruction vanishing!mechanisms}

\begin{theorem}[Classification of vanishing mechanisms]
% label removed: thm:thqg-vanishing-mechanisms
\ClaimStatusProvedHere
\index{obstruction vanishing!classification}
\begin{center}
\small
\renewcommand{\arraystretch}{1.25}
\begin{tabular}{cll}
\toprule
\emph{Obstruction} & \emph{Mechanism} & \emph{Classes} \\
\midrule
$o_3 = 0$ &
  Abelian OPE: $a_{(0)}b = 0$. &
  $\mathbf{G}$, $\mathbf{C}$ \\
$o_4 = 0$ &
  Jacobi identity kills arity-$4$ graphs. &
  $\mathbf{L}$ \\
  & Abelian propagation. &
  $\mathbf{G}$ \\
$o_5 = 0$ &
  Rank-one rigidity ($\dim V|_L = 1$). &
  $\mathbf{C}$ \\
  & Jacobi propagation ($o_4 = 0 \Rightarrow$
  cubic source inactive). &
  $\mathbf{L}$ \\
  & Abelian propagation. &
  $\mathbf{G}$ \\
\bottomrule
\end{tabular}
\end{center}
\end{theorem}


% ======================================================================
%  §15.  ADDITIVITY AND FACTORIZATION
% ======================================================================

\subsection{Additivity and factorization}
% label removed: subsec:thqg-additivity

\begin{theorem}[Independent sum factorization]
% label removed: thm:thqg-independent-sum
\ClaimStatusProvedHere
\index{gravitational complexity!additivity}
Let $L = L_1 \oplus L_2$ with vanishing mixed OPE.
\begin{enumerate}[label=\textup{(\roman*)}]
\item $\rho_{\mathrm{grav}}(L)
  = \max(\rho_{\mathrm{grav}}(L_1),\,
  \rho_{\mathrm{grav}}(L_2))$.
\item The depth class of~$L$ is the maximum in the ordering
  $\mathbf{G} < \mathbf{L} < \mathbf{C} < \mathbf{M}$.
\item $\mathrm{Sh}_r(L) = \mathrm{Sh}_r(L_1) +
  \mathrm{Sh}_r(L_2)$ for all~$r$.
\end{enumerate}
\end{theorem}

\begin{proof}
By Proposition~\ref{prop:independent-sum-factorization}:
$\Theta_L = \Theta_{L_1} + \Theta_{L_2}$ with no cross-terms.
The arity-$r$ shadow is additive.  A shadow is non-vanishing
iff at least one summand is, giving the max formula.
\end{proof}


% ======================================================================
%  §16.  DEPTH VS.\ KOSZULNESS
% ======================================================================

\subsection{Shadow depth is not Koszulness}
% label removed: subsec:thqg-depth-vs-koszulness
\index{shadow depth!vs.\ Koszulness}

\begin{proposition}[Depth and Koszulness are independent]
% label removed: prop:thqg-depth-koszulness-independent
\ClaimStatusProvedHere
\index{Koszulness!independence from depth}
Both finite and infinite depth algebras can be Koszul:
\begin{center}
\small
\renewcommand{\arraystretch}{1.15}
\begin{tabular}{lcc}
\toprule
\emph{Family} & $r_{\max}$ & \emph{Koszul?} \\
\midrule
Heisenberg $\cH_k$ & $2$ & Yes \\
Affine $V_k(\fg)$ & $3$ & Yes \\
$\beta\gamma$ & $4$ & Yes \\
Virasoro $\mathrm{Vir}_c$ & $\infty$ & Yes \\
\bottomrule
\end{tabular}
\end{center}
Shadow depth classifies \emph{complexity} within the Koszul
locus, not Koszulness itself.
\end{proposition}

\begin{proof}
Each algebra is Koszul by PBW universality
(Proposition~\ref{prop:pbw-universality}) or explicit
computation.  Their depths are $2, 3, 4, \infty$.
\end{proof}


% ======================================================================
%  §17.  COMPLEXITY UNDER DUALITY
% ======================================================================

\subsection{Complexity under Koszul duality}
% label removed: subsec:thqg-complexity-duality
\index{gravitational complexity!under Koszul duality}

\begin{proposition}[Complexity is a Koszul duality invariant]
% label removed: prop:thqg-duality-preserves-complexity
\ClaimStatusProvedHere
$\rho_{\mathrm{grav}}(\cA^!) = \rho_{\mathrm{grav}}(\cA)$.
\end{proposition}

\begin{proof}
The complementarity theorem (Theorem~C) gives a filtered
quasi-isomorphism
$\Defcyc(\cA) \simeq \Defcyc(\cA^!)$ preserving the arity filtration.
Therefore each genus-$0$, arity-$r$ shadow cohomology group agrees for
$\cA$ and $\cA^!$, so the set of non-vanishing shadow arities is the
same on both sides.  Hence
$\rho_{\mathrm{grav}}(\cA^!) = \rho_{\mathrm{grav}}(\cA)$.
This matches the explicit Koszul pairs
$\cH_k \leftrightarrow \Sym^{\mathrm{ch}}(V^*)$ (class~$\mathbf{G}$),
$V_k(\fg) \leftrightarrow V_{-k-2h^\vee}(\fg)$ (class~$\mathbf{L}$),
$\mathrm{Vir}_c \leftrightarrow \mathrm{Vir}_{26-c}$
(class~$\mathbf{M}$).
\end{proof}

\begin{corollary}[Duality-complexity table]
% label removed: cor:thqg-duality-table
\ClaimStatusProvedHere
\begin{center}
\small
\renewcommand{\arraystretch}{1.15}
\begin{tabular}{lllcc}
\toprule
\emph{Algebra} & \emph{Koszul dual} &
  \emph{Duality} & \emph{Class} & $r_{\max}$ \\
\midrule
$\cH_k$ & $\Sym^{\mathrm{ch}}(V^*)$ &
  $\kappa \mapsto -\kappa$ & $\mathbf{G}$ & $2$ \\
$V_k(\fg)$ & $V_{-k-2h^\vee}(\fg)$ &
  Feigin--Frenkel & $\mathbf{L}$ & $3$ \\
$\mathrm{Vir}_c$ & $\mathrm{Vir}_{26-c}$ &
  $c \leftrightarrow 26-c$ & $\mathbf{M}$ & $\infty$ \\
\bottomrule
\end{tabular}
\end{center}
\end{corollary}


% ======================================================================
%  §18.  TROPICAL COMPLEXITY
% ======================================================================

\subsection{Tropical depth profile}
% label removed: subsec:thqg-tropical
\index{tropical complexity}

\begin{definition}[Tropical depth profile]
% label removed: def:thqg-tropical-depth
\index{tropical depth profile|textbf}
$\mathrm{dp}_\cA(r) := [\cA^{\mathrm{sh}}_{r,0} \neq 0]
\in \{0,1\}$, $r \geq 2$.
\end{definition}

\begin{proposition}[Tropical profiles]
% label removed: prop:thqg-tropical-profiles
\ClaimStatusProvedHere
\begin{center}
\small
\renewcommand{\arraystretch}{1.15}
\begin{tabular}{llll}
\toprule
\emph{Family} & \emph{Class} &
  $\mathrm{supp}(\mathrm{dp}_\cA)$ &
  \emph{Description} \\
\midrule
Heisenberg & $\mathbf{G}$ & $\{2\}$ &
  Single point \\
Affine & $\mathbf{L}$ & $\{2, 3\}$ &
  Two adjacent values \\
$\beta\gamma$ & $\mathbf{C}$ & $\{2, 4\}$ &
  Gap at $3$ \\
Virasoro & $\mathbf{M}$ & $\{2, 3, 4, 5, \ldots\}$ &
  Full half-line $[2, \infty)$ \\
\bottomrule
\end{tabular}
\end{center}
The gap at~$3$ for $\beta\gamma$ is the tropical reflection of
the cubic gauge triviality
\textup{(}Theorem~\ref{thm:cubic-gauge-triviality}\textup{)}.
Explicitly: $\cA^{\mathrm{sh}}_{3,0}(\beta\gamma) = 0$ because
the OPE $\beta(z)\gamma(w) \sim 1/(z-w)$ has no simple pole
in the symmetric channel, so the transferred ternary operation
$m_3^{\mathrm{tr}} = 0$ vanishes identically: the only iterated
bracket is $[\beta, [\gamma, \beta]] = [\beta, 1] = 0$, and
the cubic shadow $\mathfrak{C}(\beta\gamma) = 0$.
\end{proposition}


% ======================================================================
%  §19.  HOLOGRAPHIC CONSEQUENCES
% ======================================================================

\subsection{Holographic consequences}
% label removed: subsec:thqg-holographic-consequences
\index{gravitational complexity!holographic consequences}

\begin{theorem}[Holographic type from complexity]
% label removed: thm:thqg-holographic-type
\ClaimStatusProvedHere
\index{holographic modular Koszul datum!complexity determination}
The gravitational complexity determines the structural type of
$\mathcal{H}(T)$:
\begin{center}
\small
\renewcommand{\arraystretch}{1.25}
\begin{tabular}{c@{\quad}l@{\quad}l}
\toprule
\emph{Class} & \emph{Laplace kernel $r^L(z)$} &
\emph{Shadow connection} \\
\midrule
$\mathbf{G}$ &
  $r^L(z) = \kappa/z^2$. &
  $\nabla^{\mathrm{hol}} = d - \kappa\,\omega_g$ \\
$\mathbf{L}$ &
  $r^L(z) = \kappa/z^2 + \mathfrak{C}_0/z^3$. &
  $\nabla^{\mathrm{hol}} = d - \kappa\,\omega_g -
  \mathfrak{C}$ \\
$\mathbf{C}$ &
  $r^L(z) = \kappa/z^2 + \mathfrak{Q}_0/z^4$. &
  $\nabla^{\mathrm{hol}} = d - \kappa\,\omega_g -
  \mathfrak{Q}$ \\
$\mathbf{M}$ &
  $r^L(z) = \sum_{r\geq 2} S_r/z^r$. &
  $\nabla^{\mathrm{hol}} = d -
  \sum_{r\geq 2} \mathrm{Sh}_r$ \\
\bottomrule
\end{tabular}
\end{center}
\end{theorem}

\begin{proof}
The binary shadow $r^L(z)$ (the Laplace kernel of the
$\lambda$-bracket) receives an arity-$r$ contribution as a
$z^{-r}$ pole; the collision residue
$r^{\mathrm{coll}}(z) = \Res^{\mathrm{coll}}_{0,2}(\Theta_\cA)$
has pole orders one lower.
Stabilization at $r_{\max}$ truncates the Laurent expansion.
\end{proof}


\begin{theorem}[Boundary-holographic complexity classification]
\label{thm:thqg-boundary-holographic-complexity}
\ClaimStatusProvedHere
\index{boundary-holographic complexity|textbf}
\index{shadow depth!boundary-holographic interpretation}
\index{gravitational partition function!complexity classification}
Let $\cT$ be a $3$d HT theory on $\C \times \R$ whose boundary
chiral algebra $\cA$ is modular Koszul.
The bulk is topological \textup{(}$3$d Chern--Simons or its
higher-spin generalisation; no local propagating degrees of
freedom\textup{)}.  The perturbative gravitational partition
function
\begin{equation}\label{eq:boundary-holographic-Z}
  Z^{\mathrm{pert}}(\cT)
  \;=\;
  \exp\Bigl(\,\sum_{g \geq 1} \hbar^{2g}\,F_g(\cT)\Bigr),
  \qquad
  F_g(\cT)
  \;=\;
  \int_{\overline{\cM}_g}
  \operatorname{Tr}^{\circ g}\bigl(\Theta_\cA^{(g)}\bigr),
\end{equation}
is completely determined by the boundary data: the OPE of
$\cA$, its transferred $A_\infty$ structure, and the
bar-intrinsic MC element $\Theta_\cA$
\textup{(}Theorem~\textup{\ref*{V1-thm:mc2-bar-intrinsic}}\textup{)}.
The shadow depth $r_{\max}(\cA)$ classifies the nonlinear
complexity of this boundary-determined partition function.
On any primary line~$L$, the weighted shadow generating function
$H(t) = \sum_{r \geq 2} r\,\mathrm{Sh}_r\,t^r$ satisfies
the Riccati algebraicity relation
\textup{(}Theorem~\textup{\ref*{V1-thm:riccati-algebraicity}}\textup{)}
\begin{equation}\label{eq:boundary-holographic-riccati}
  H(t)^2 \;=\; t^4\,Q_L(t),
  \qquad
  Q_L(t) = 4\kappa^2 + 12\kappa\alpha\,t
  + (9\alpha^2 + 2\Delta)\,t^2,
\end{equation}
with critical discriminant $\Delta = 8\kappa S_4$.
The arity-$r$ shadow first contributes to $F_g$ at genus
$g_{\min}(r) = \lfloor r/2 \rfloor + 1$
\textup{(}Corollary~\textup{\ref{cor:shadow-visibility-genus}}\textup{)}.
The four complexity classes are:
\begin{enumerate}[label=\textup{(\roman*)}]
\item \emph{Gaussian} ($\mathbf{G}$, $r_{\max} = 2$;
  $\Delta = 0$, $\alpha = 0$).
  The boundary $A_\infty$ structure is strictly
  coassociative: $m_k = 0$ for $k \geq 3$.
  The partition function is one-loop exact at every genus:
  \[
    F_g(\cT) = \kappa(\cA)\cdot\lambda_g^{\mathrm{FP}},
    \qquad
    \lambda_g^{\mathrm{FP}} =
    \frac{2^{2g-1}-1}{2^{2g-1}}\cdot
    \frac{|B_{2g}|}{(2g)!},
  \]
  where $B_{2g}$ are Bernoulli numbers; no higher-arity
  corrections exist.  The complementarity potential
  $S_\cA(x) = \tfrac{1}{2}\kappa\,x^2$ is Gaussian.
  \textup{(}Heisenberg $\cH_k$, lattice VOA $V_\Lambda$,
  free fermion.\textup{)}

\item \emph{Lie} ($\mathbf{L}$, $r_{\max} = 3$;
  $\Delta = 0$, $\alpha \neq 0$).
  The boundary OPE produces a cubic shadow
  $\mathfrak{C} = \kappa(x,[y,z]) \neq 0$;
  the quartic obstruction $o_4 =
  \tfrac{1}{2}\{\mathfrak{C},\mathfrak{C}\}_H$ vanishes
  by the Jacobi identity.
  The complementarity potential is cubic:
  $S_\cA(x) = \tfrac{1}{2}\kappa\,x^2 +
  \tfrac{1}{6}\mathfrak{C}_0\,x^3$.
  New corrections first appear at genus~$2$
  \textup{(}from $\mathrm{Sh}_3$\textup{)} and no
  higher-arity shadows contribute.
  \textup{(}Affine $V_k(\fg)$ at generic level.\textup{)}

\item \emph{Contact} ($\mathbf{C}$, $r_{\max} = 4$;
  stratum separation\textup{)}.
  The boundary $A_\infty$ structure has
  $\mathfrak{C} = 0$ and a quartic contact invariant
  $\mathfrak{Q} \neq 0$.
  The quintic obstruction $o_5 = 0$ by rank-one abelian
  rigidity: the self-bracket
  $\mathrm{Def}_{\mathrm{cyc}}^{(q)} \otimes
  \mathrm{Def}_{\mathrm{cyc}}^{(q)} \to
  \mathrm{Def}_{\mathrm{cyc}}^{(2q)}$ has
  $\dim\mathrm{Def}_{\mathrm{cyc}}^{(2q)} = 0$.
  One quartic correction to the Gaussian, realized at
  genus~$3$.
  \textup{(}$\beta\gamma$ on the weight-changing line.\textup{)}

\item \emph{Mixed} ($\mathbf{M}$, $r_{\max} = \infty$;
  $\Delta \neq 0$).
  Both $\mathfrak{C} \neq 0$ and $\mathfrak{Q} \neq 0$;
  the cubic source permanently generates higher obstructions:
  $\{\mathfrak{C}, \mathrm{Sh}_r\}_H \neq 0$ for all
  $r \geq 3$ forces $o_{r+1} \neq 0$ by induction.
  The complementarity potential is non-polynomial but
  convergent.  The shadow coefficients satisfy
  $S_r = \Delta \cdot R_r$ for all $r \geq 4$, with
  $R_r \in \mathbb{Q}(\alpha, \kappa, S_4)$, and
  $|S_r| \sim A\,\rho^r\,r^{-5/2}$ with growth rate
  \[
    \rho(\cA)
    \;=\;
    \frac{\sqrt{9\alpha^2 + 2\Delta}}{2|\kappa|}
  \]
  \textup{(}Theorem~\textup{\ref*{V1-thm:shadow-radius}}\textup{)}.
  The partition function receives genuinely new corrections
  at every genus from every arity: two new shadow
  coefficients $(S_{2g-2}, S_{2g-1})$ enter at each genus~$g$.
  \textup{(}Virasoro $\mathrm{Vir}_c$, $\mathcal{W}_N$ for
  $N \geq 3$.\textup{)}
\end{enumerate}
The classification is exhaustive
\textup{(}Theorem~\textup{\ref*{V1-thm:single-line-dichotomy}}\textup{)}
and detected at genus~$2$: the binary invariants
$\sigma = [\mathfrak{C} \neq 0]$ and
$\pi = [\mathfrak{Q} \neq 0]$ biject onto
$\{\mathbf{G}, \mathbf{L}, \mathbf{C}, \mathbf{M}\}$.
The bulk has no independent dynamics; the entire perturbative
gravitational content resides in the boundary OPE and its
modular extension~$\Theta_\cA$.
\end{theorem}

\begin{proof}
The bulk theory is topological: $3$d Chern--Simons has no
local propagating degrees of freedom.  Perturbative partition
functions are entirely determined by the boundary chiral
algebra (Brown--Henneaux for Virasoro,
Theorem~\ref{thm:brown-henneaux};
the Costello--Gwilliam construction for the general case;
see~\S\ref{subsec:gravity-physics}).
Equation~\eqref{eq:boundary-holographic-Z} follows from the
bar-intrinsic MC element
(Theorem~\ref*{V1-thm:mc2-bar-intrinsic}): the genus-$g$
MC component $\Theta_\cA^{(g)}$ is a section of the
$g$-fold sewing of the Hodge-weighted bar complex, and the
trace $\operatorname{Tr}^{\circ g}$ extracts the scalar
genus-$g$ free energy.

The Riccati algebraicity~\eqref{eq:boundary-holographic-riccati}
is the restriction of the all-arity master equation to a primary
line
(Theorem~\ref*{V1-thm:riccati-algebraicity}): the nonlinear
recursion among shadow coefficients is equivalent to the
algebraic relation $H^2 = t^4 Q_L$, making the shadow
generating function algebraic of degree~$2$.  The discriminant
$\Delta = 8\kappa S_4$ controls the analytic type: $\Delta = 0$
gives $\sqrt{Q_L}$ polynomial \textup{(}finite tower\textup{)};
$\Delta \neq 0$ gives $\sqrt{Q_L}$ irrational
\textup{(}infinite tower\textup{)}.

The shadow visibility genus
$g_{\min}(S_r) = \lfloor r/2 \rfloor + 1$ follows from the
self-loop parity vanishing
(Proposition~\ref*{V1-prop:self-loop-vanishing}): single-vertex
graphs $(0, 2k)$ with $k$ self-loops contribute zero for
$k \geq 2$ by odd-dimensional parity obstruction; the first
nonvanishing contribution of~$S_r$ comes from a two-vertex
graph at genus $\lfloor r/2 \rfloor + 1$.

Parts~(i)--(iv) then follow from the four-class classification
and the stabilization theorem above, together with the shadow
archetype classification
(Theorem~\ref*{V1-thm:shadow-archetype-classification}).
For class~$\mathbf{G}$, only the quadratic shadow~$\kappa$
contributes, giving
$F_g = \kappa \cdot \lambda_g^{\mathrm{FP}}$
(the scalar genus tower, Theorem~D).
For class~$\mathbf{M}$, the non-terminating tower contributes
corrections at every genus, with convergent generating function
by the shadow radius estimate
(Theorem~\ref*{V1-thm:shadow-radius}).
Exhaustiveness and genus-$2$ detection follow from the
single-line dichotomy
(Theorem~\ref*{V1-thm:single-line-dichotomy})
and the genus-$2$ binary diagnostic above: the map
$(\sigma, \pi) \to \{\mathbf{G}, \mathbf{L},
\mathbf{C}, \mathbf{M}\}$ is a bijection.
\end{proof}


% ======================================================================
%  §20.  THE $W_N$ LIMIT
% ======================================================================

\subsection{The $\mathcal{W}_\infty$ limit}
% label removed: subsec:thqg-winfty-limit
\index{W-infinity@$\mathcal{W}_\infty$!complexity limit}

\begin{proposition}[$\mathcal{W}_N$ complexity stabilises]
% label removed: prop:thqg-wn-stabilization
\ClaimStatusProvedHere
For all $N \geq 2$: $\mathcal{W}_N^k$ is class~$\mathbf{M}$
with $r_{\max} = \infty$.  The coefficients
$\mathrm{Sh}_r(\mathcal{W}_N^k)$ stabilise in $N$ for fixed~$r$:
$\mathrm{Sh}_r(\mathcal{W}_N) = \mathrm{Sh}_r(\mathcal{W}_\infty)$
for $N \geq r - 1$.
\end{proposition}

\begin{proof}
Each $\mathcal{W}_N$ for $N \geq 2$ contains Virasoro,
so $\mathfrak{C} \neq 0$ and $\mathfrak{Q} \neq 0$.
The infinite tower follows from
Lemma~\ref{lem:thqg-cubic-source}.
Stabilisation: the arity-$r$ shadow involves generators of
conformal weight $\leq r$; for $N > r-1$, these coincide
in $\mathcal{W}_N$ and $\mathcal{W}_{N+1}$.
\end{proof}


% ======================================================================
%  §21.  THE STANDARD LANDSCAPE
% ======================================================================

\subsection{The standard gravitational landscape}
% label removed: subsec:thqg-grav-landscape
\index{gravitational complexity!landscape}

\begin{theorem}[Gravitational complexity census]
% label removed: thm:thqg-grav-landscape
\ClaimStatusProvedHere
\index{gravitational complexity!complete census}
\begin{center}
\small
\renewcommand{\arraystretch}{1.3}
\begin{tabular}{l@{\quad}c@{\quad}c@{\quad}l@{\quad}l}
\toprule
\emph{Family} & \emph{Class} & $r_{\max}$ &
  \emph{Potential} &
  \emph{Bulk gravity} \\
\midrule
Heisenberg $\cH_k$ &
  $\mathbf{G}$ & $2$ & Quadratic & Free \\
Free fermion &
  $\mathbf{G}$ & $2$ & Quadratic & Free \\
Lattice $V_\Lambda$ &
  $\mathbf{G}$ & $2$ & Quadratic & Free \\[3pt]
$V_k(\mathfrak{sl}_2)$ &
  $\mathbf{L}$ & $3$ & Cubic & Semistrict HS \\
$V_k(\mathfrak{sl}_N)$ &
  $\mathbf{L}$ & $3$ & Cubic & Semistrict HS \\
$V_k(\fg)$ (simple) &
  $\mathbf{L}$ & $3$ & Cubic & Semistrict HS \\[3pt]
$\beta\gamma$ &
  $\mathbf{C}$ & $4$ & Quartic & Contact gravity \\[3pt]
$\mathrm{Vir}_c$ &
  $\mathbf{M}$ & $\infty$ & Non-poly & Full $L_\infty$ \\
$\mathcal{W}_3^k$ &
  $\mathbf{M}$ & $\infty$ & Non-poly & Full $L_\infty$ \\
$\mathcal{W}_N^k$ &
  $\mathbf{M}$ & $\infty$ & Non-poly & Full $L_\infty$ \\
\bottomrule
\end{tabular}
\end{center}
All families are Koszul; the complexity classification refines
within the Koszul locus.
\end{theorem}


% ======================================================================
%  §22.  THE MAIN THEOREM
% ======================================================================

\subsection{Summary: the gravitational complexity theorem}
% label removed: subsec:thqg-complexity-summary

\begin{maintheorem}[Gravitational complexity; result (G2)]
% label removed: thm:thqg-g2-main
\ClaimStatusProvedHere
\index{gravitational complexity!main theorem}
Let $\cA$ be a chirally Koszul, cyclically admissible chiral
algebra satisfying the HS-sewing condition.
The shadow depth $r_{\max}(\cA)$ determines the minimal
homotopy-algebraic complexity of the bulk gravitational theory.
The classification into exactly four classes is forced by the
obstruction tower:
\begin{enumerate}[label=\textup{(\roman*)}]
\item \textbf{Gaussian} ($r_{\max} = 2$):
  Free bulk theory.  All obstructions vanish.
  Quadratic potential.  Archetype: Heisenberg.

\item \textbf{Lie} ($r_{\max} = 3$):
  Semistrict higher-spin gravity with cubic vertices.
  Jacobi identity kills all higher arities.
  Archetype: affine $V_k(\fg)$.

\item \textbf{Contact} ($r_{\max} = 4$):
  Quartic contact vertices, no cubic.
  Cubic gauge triviality + rank-one rigidity terminate the tower.
  Archetype: $\beta\gamma$.

\item \textbf{Mixed} ($r_{\max} = \infty$):
  Full $L_\infty$ with vertices at all arities.
  Cubic source permanently generates higher obstructions.
  Non-polynomial potential.  Archetype: Virasoro.
\end{enumerate}
The gravitational complexity is a quasi-isomorphism invariant,
a Koszul duality invariant, and generically stable under
parameter deformation.
\end{maintheorem}


% ======================================================================
%  §23.  ENVELOPE-SHADOW FUNCTOR
% ======================================================================

\subsection{Gravitational complexity as a functor}
% label removed: subsec:thqg-envelope-shadow
\index{envelope-shadow functor!gravitational}

\begin{proposition}[Functoriality of complexity]
% label removed: prop:thqg-complexity-functor
\ClaimStatusProvedHere
\index{envelope-shadow functor!complexity functor}
For chiral algebras arising as $\cA = \Uvert(R)$
\textup{(}Nishinaka~\cite{Nish26}\textup{)}:
\begin{equation}% label removed: eq:thqg-chienv
\chienv(R)
\;:=\;
\rho_{\mathrm{grav}}(\Uvert(R))
\end{equation}
is a functor on Lie conformal algebras.  Three structural
properties hold:
\begin{enumerate}[label=\textup{(\roman*)}]
\item \emph{Finite-jet rigidity}
  \textup{(}Proposition~\ref{prop:finite-jet-rigidity}\textup{)}:
  the arity-$r$ shadow depends only on the $(r{+}1)$-jet
  $\Uvert(R)/F_{r+1}$.
\item \emph{Polynomial level dependence}
  \textup{(}Proposition~\ref{prop:polynomial-level-dependence}\textup{)}:
  the arity-$r$ shadow is polynomial of degree $\leq r{-}1$ in
  a central-extension parameter.
\item \emph{Gaussian collapse}
  \textup{(}Proposition~\ref{prop:gaussian-collapse-abelian}\textup{)}:
  $\chienv(R) = 2$ if and only if $R$ is abelian; Heisenberg is
  the universal fixed point.
\end{enumerate}
\end{proposition}

\begin{proof}
Functoriality: a morphism $R \to R'$ induces
$\Uvert(R) \to \Uvert(R')$, hence
$\Theta_{\Uvert(R)} \mapsto \Theta_{\Uvert(R')}$.
Shadow projections are compatible, giving
$\chienv(R) \leq \chienv(R')$ when $R \hookrightarrow R'$.
Part~(i): the obstruction class $o_{r+1}$ depends on
the structure constants of the $\lambda$-bracket through
degree $r+1$ (the graph-sum formula involves at most $r+1$
vertices), hence on the $(r{+}1)$-jet.
Part~(ii): the structure constants of $\Uvert(R)$ are polynomial
in the central extension $k$ (by the OPE expansion), so
$o_{r+1}$ is polynomial in $k$ of degree at most $r-1$
(each vertex contributes at most degree one).
Part~(iii): if $R$ is abelian, all iterated brackets vanish
and $d_\infty = 2$; conversely, $d_\infty = 2$ forces
all transferred ternary operations to vanish, which implies
$a_{(0)}b = 0$ for all generators.
\end{proof}

\begin{remark}[DS reduction and complexity]
% label removed: rem:thqg-ds-complexity
\index{Drinfeld--Sokolov reduction!complexity}
For the principal Drinfeld--Sokolov reduction
$\mathcal{W}_k(\fg) = H^0_{\mathrm{DS}}(V_k(\fg))$:
the reduction preserves Koszulness but may change the
complexity class.  Affine $V_k(\mathfrak{sl}_2)$ is
class~$\mathbf{L}$ ($r_{\max} = 3$), while its DS reduction
$\mathrm{Vir}_c$ is class~$\mathbf{M}$ ($r_{\max} = \infty$).
The reduction $V_k(\mathfrak{sl}_N) \to \mathcal{W}_N^k$
similarly jumps from $\mathbf{L}$ to $\mathbf{M}$ for
all $N \geq 2$.  The mechanism: DS reduction introduces
non-abelian composite fields (Sugawara $T$, higher-spin
$W_s$) whose iterated OPE generates operations at all
arities.
\end{remark}

\begin{remark}[Chain-level versus cohomology-level structure for $\mathcal{W}_N$]
\label{rem:ds-chain-vs-cohomology-wn}
\index{Drinfeld--Sokolov reduction!chain vs cohomology}
\index{W-algebra!chain-level structure}
The $\Ainf$ operations and the Yangian structure for
$\mathcal{W}_N$ live on the BRST \emph{cohomology}
$H^\bullet(V_k(\mathfrak{sl}_N) \otimes \mathcal{C}_f,\,
Q_{\mathrm{DS}})$, not on the chain-level BRST complex.  The
passage from chains to cohomology is the homotopy transfer
theorem (Kadeishvili--Merkulov), and it discards
information.  Three phenomena deserve attention.

\begin{enumerate}[label=\textup{(\roman*)}]
\item \emph{The chain-level complex is strictly richer.}
  For every $N \geq 2$, the BRST complex
  $\bigl(V_k(\mathfrak{sl}_N) \otimes \mathcal{C}_f,\,
  Q_{\mathrm{DS}}\bigr)$ is a dg vertex algebra whose
  generators include the affine currents $J^a(z)$ and the
  $bc$-ghost system $(\beta_\alpha, \gamma_\alpha)$ for each
  positive root~$\alpha$.  The cohomology
  $\mathcal{W}_N = H^0(Q_{\mathrm{DS}})$ has generators
  $W_2 = T,\, W_3,\, \ldots,\, W_N$ of conformal weights
  $2, 3, \ldots, N$.  The composite fields appearing in the
  $\mathcal{W}_N$ OPE (in particular the quasi-primary
  $\Lambda = {:\!TT\!:} - \tfrac{3}{10}\partial^2 T$ at
  weight~$4$ for $N = 3$) are \emph{not} among the
  generators of the chain-level complex; they arise only
  after taking $Q_{\mathrm{DS}}$-cohomology and forming
  normal-ordered products.  The chain-level differential
  $Q_{\mathrm{DS}}$ sees the individual $J^a$ and ghost
  contributions to~$\Lambda$ separately; the cohomological
  algebra sees only the composite.

\item \emph{Higher transferred operations detect the loss.}
  The chain-level $\Ainf$ structure
  $\{m_k^{\mathrm{chain}}\}$ on the BRST complex is strict
  (coming from the operadic action of~$\SCchtop$ on the
  gauge theory).  The transferred operations
  $\{m_k^H\}$ on $\mathcal{W}_N$ are computed by
  summing over trees with internal edges labelled by the
  BRST contracting homotopy~$h$.  For $N = 2$ (Virasoro),
  $h$ threads through the single ghost pair
  $(\beta, \gamma)$ and the Sugawara construction yields the
  wheel-diagram formula for $m_k^H(T, \ldots, T)$ at all
  arities~$k \geq 3$.  For $N = 3$, the homotopy~$h$ routes
  through $\binom{3}{2} = 3$ ghost pairs; the transferred
  $m_3^H(W, W, W)$ involves compositions
  $m_2^{\mathrm{chain}}(h \circ m_2^{\mathrm{chain}}, -)$
  whose intermediate states pass through ghost-number sectors
  invisible to~$\mathcal{W}_3$.  In particular, the
  $c$-dependent coefficient
  $\tfrac{32}{5c+22}$ multiplying~$\Lambda$ in the
  $\{W_\lambda W\}$ bracket is the \emph{residue} of a
  chain-level computation that also determines the
  relative coefficient of non-cohomological ghost terms;
  these ghost terms are projected away on cohomology but
  contribute to the transferred $m_4^H, m_5^H, \ldots\,$
  through the homotopy.

\item \emph{General $N$: the ghost sector grows
  quadratically.}
  For $\mathfrak{sl}_N$, the BRST complex has
  $\tfrac{1}{2}N(N-1)$ ghost pairs (one per positive root).
  The chain-level state space at conformal weight~$h$
  grows as
  $\dim\bigl(V_k(\mathfrak{sl}_N) \otimes
  \mathcal{C}_f\bigr)_h \sim p_N(h)$,
  where $p_N$ is a partition function with $N^2 - 1$
  free-field contributions.  The cohomology
  $\dim(\mathcal{W}_N)_h$ is governed by the $N - 1$
  generators of weights $2, \ldots, N$, hence is
  strictly smaller for $h \geq N + 1$.
  The homotopy transfer discards
  $\dim(\mathrm{chain})_h -
  \dim(\mathrm{cohomology})_h$
  directions at each weight; these discarded directions
  are the ``reservoir'' from which the transferred
  $m_k^H$ draws its non-trivial values.  Concretely:
  the transferred $m_k^H$ on $\mathcal{W}_N$ depends on
  the $Q_{\mathrm{DS}}$-acyclic summand through the
  contracting homotopy, and this dependence is sensitive
  to the choice of homotopy retract (different choices
  yield gauge-equivalent $\Ainf$ structures).
\end{enumerate}

For non-principal reductions
$\mathcal{W}_k(\fg, f) = H^0(Q_{\mathrm{DS}}^f)$
with $f$ not the principal nilpotent, the chain-level
complex is richer still: the ghost system
$\mathcal{C}_f$ is determined by the $\mathrm{ad}(x)$-eigenspace
decomposition ($x$ the neutral element of the
$\mathfrak{sl}_2$-triple containing~$f$), and
non-principal choices leave more affine generators in
cohomology while simultaneously enlarging the ghost
sector.  The subregular reduction for $\mathfrak{sl}_3$
produces a $\mathcal{W}$-algebra with generators at
weights $2$ and $3$ (as for $\mathcal{W}_3$) but with
different OPE coefficients and a different contracting
homotopy; the hook-type reductions in type~A produce
generators whose weights are determined by the
corresponding partition.  In every case, the
$\Ainf$ operations on the reduced algebra are shadows
of the chain-level BRST complex, and the full
chain-level data, accessible through the bar complex
$B(V_k(\fg) \otimes \mathcal{C}_f)$ before taking
cohomology, encodes strictly more information.

The open question is whether the chain-level BRST
data can be recovered from the cohomology-level
$\Ainf$ structure.  For principal reductions of simply
laced algebras, the answer is affirmative on the Koszul
locus: bar-cobar inversion (Theorem~A) reconstructs
the chain-level algebra from its transferred $\Ainf$
structure up to quasi-isomorphism.  For non-principal
reductions, the reconstruction requires that the
Koszul property survives the reduction
(Theorem~\ref*{V1-thm:ds-koszul-obstruction}), which is
proved for the hook-type corridor in type~A and open in
general.
\end{remark}


% ======================================================================
%  §24.  CONVERGENCE OF THE VIRASORO POTENTIAL
% ======================================================================

\subsection{Convergence analysis}
% label removed: subsec:thqg-convergence
\index{Virasoro!potential convergence}

\begin{proposition}[Coefficient asymptotics]
% label removed: prop:thqg-coefficient-asymptotics
\ClaimStatusProvedHere
\index{Virasoro!shadow tower asymptotics}
For large $c > 0$ and $r \geq 4$, the Virasoro shadow
coefficients $S_r$ (where $\mathrm{Sh}_r = S_r\,x^r$)
satisfy
\begin{equation}% label removed: eq:thqg-coefficient-bound
|S_r|
\;\leq\;
C(c)\cdot r!\cdot c^{-(r-3)}\cdot(5c+22)^{-\lfloor(r-2)/2\rfloor}
\end{equation}
for a constant $C(c)$ independent of~$r$.  In particular:
\begin{enumerate}[label=\textup{(\roman*)}]
\item For fixed $c > 0$, $|S_r| \to 0$ faster than any
  exponential as $r \to \infty$.
\item The complementarity potential
  $S_{\mathrm{Vir}}(x)$ converges absolutely in
  $|x| < R(c)$ with radius
  $R(c) \sim c^{1/2}$ for large~$c$.
\item The $c \to \infty$ limit of the normalised potential
  $c^{-1}\,S_{\mathrm{Vir}}(x\,c^{1/2})$ is the Gaussian
  $\frac{1}{4}\,y^2$: the large-$c$ limit is class~$\mathbf{G}$.
\end{enumerate}
\end{proposition}

\begin{proof}
Part~(i): The denominator pattern
(Proposition~\ref{prop:thqg-virasoro-structure}(i)) gives
$|S_r| \leq C_0\,r^{C_1}/[c^{r-3}(5c+22)^{\lfloor(r-2)/2\rfloor}]$
for polynomial numerator growth.  The double-exponential
decay in $r$ from the denominator dominates.
Part~(ii): The ratio test on $|S_{r+1}/S_r|$ gives
$\limsup_{r\to\infty} |S_{r+1}/S_r|^{1/r} \leq c^{-1}$,
so the radius of convergence is at least~$c$.
Part~(iii): Rescaling $x = y\,c^{-1/2}$ in the potential
and taking $c \to \infty$: the quadratic term
$\frac{c}{4}(y/\sqrt{c})^2 = \frac{1}{4}y^2$ survives;
the cubic term $\frac{1}{3}(y/\sqrt{c})^3 = O(c^{-3/2})$
vanishes; all higher terms decay faster.
\end{proof}

\begin{remark}[Singular loci]
% label removed: rem:thqg-singular-loci
\index{Virasoro!singular loci}
The potential $S_{\mathrm{Vir}}$ has two singular loci:
\begin{enumerate}[label=\textup{(\alph*)}]
\item $c = 0$: the Hessian $H = c/2$ degenerates, the
  propagator $P = 2/c$ diverges, and all shadow coefficients
  for $r \geq 4$ blow up.  Physically: the free-field
  limit where conformal symmetry is trivialised.
\item $c = -22/5$: the quartic contact coefficient
  $Q_0 = 10/[c(5c+22)]$ diverges.  This is the minimal
  model $M(2,5)$ central charge.  The collision of Shapovalov
  eigenvalues at level~$2$ produces the singularity.
\end{enumerate}
At both loci, the shadow obstruction tower requires analytic continuation,
not direct evaluation.
\end{remark}


% ======================================================================
%  §25.  DEFORMATION STABILITY
% ======================================================================

\subsection{Deformation stability of the complexity class}
% label removed: subsec:thqg-deformation-stability
\index{shadow depth!deformation stability}

\begin{proposition}[Generic constancy]
% label removed: prop:thqg-generic-constancy
\ClaimStatusProvedHere
\index{shadow depth!generic constancy}
Let $\{\cA_t\}_{t \in \C}$ be a continuous family of modular
Koszul chiral algebras.  The set
$\{t : r_{\max}(\cA_t) \neq r_{\max}(\cA_0)\}$ is discrete
in~$\C$.  At critical values, $r_{\max}$ can only decrease
\textup{(}obstruction classes specialise to zero\textup{)}.
\end{proposition}

\begin{proof}
The obstruction class $o_{r+1}(\cA_t)$ is a polynomial in
the structure constants of $\cA_t$ (graph-sum formula).  As
a function of $t$, it is algebraic.  The zero locus of an
algebraic function is discrete unless identically zero.
Vanishing at a special value is a closed condition, so
$r_{\max}$ can only decrease.
\end{proof}

\begin{example}[Class jumps in the Virasoro family]
% label removed: ex:thqg-virasoro-class-jumps
\index{Virasoro!class jumps}
The Virasoro algebra $\mathrm{Vir}_c$ is class~$\mathbf{M}$
for all $c \neq 0$, with no class jumps at finite $c$.
The quartic contact $Q_0 = 10/[c(5c+22)]$ is nonzero for
$c \neq 0, -22/5$, and the cubic $\mathfrak{C} = 2x^3$ is
independent of $c$.  At $c = 0$: the Hessian degenerates;
the shadow depth is undefined (the cyclic pairing fails to
be nondegenerate).  At $c = -22/5$: $Q_0$ diverges, signalling
a change of the effective field space (the minimal model
$M(2,5)$ has a null vector that modifies the cyclic slice).
\end{example}

\begin{example}[Class jumps in the affine family]
% label removed: ex:thqg-affine-class-jumps
\index{affine Kac--Moody!class jumps}
The affine algebra $V_k(\mathfrak{sl}_2)$ is class~$\mathbf{L}$
for all non-critical levels $k \neq -2$.  The cubic shadow
$\mathfrak{C} = \kappa \cdot f_{\mathfrak{sl}_2}\cdot x^3$
vanishes at $k = -h^\vee = -2$ (critical level, where
$\kappa = 0$).  At $k = -2$: the shadow depth jumps from
$3$ to $2$ (class $\mathbf{L} \to \mathbf{G}$), consistent
with the Feigin--Frenkel construction at the critical level.
\end{example}


% ======================================================================
%  §26.  THE MC EQUATION AS GRAVITY ACTION
% ======================================================================

\subsection{The MC equation as a gravitational action principle}
% label removed: subsec:thqg-mc-action
\index{Maurer--Cartan equation!as gravitational action}

\begin{theorem}[MC equation = Euler--Lagrange equation]
% label removed: thm:thqg-mc-euler-lagrange
\ClaimStatusProvedHere
\index{complementarity potential!as gravity action}
The MC equation $D_\cA\,\Theta_\cA +
\frac{1}{2}[\Theta_\cA, \Theta_\cA] = 0$,
evaluated on the one-generator primary line, is the
Euler--Lagrange equation for the complementarity potential
$S_\cA(x)$:
\begin{equation}% label removed: eq:thqg-euler-lagrange
\frac{\partial S_\cA}{\partial x}
\;=\;
\kappa\,x + \sum_{r=3}^{r_{\max}}
\mathrm{Sh}_r(\cA)\,x^{r-1}
\;=\; 0
\end{equation}
at the critical point $x = 0$.  The second variation is
$S''_\cA(0) = \kappa(\cA) = H_\cA$.  The genus-$g$
partition function arises from the formal saddle-point
expansion $Z_g(\cA) \sim \int e^{S_\cA/\hbar}$.
\end{theorem}

\begin{proof}
The all-arity master equation
$\nabla_H(\mathrm{Sh}_r) + \mathfrak{o}^{(r)} = 0$,
restricted to the primary line, is the recursion for the
Taylor coefficients of $S_\cA$.  The MC equation for
$\Theta_\cA$ is equivalent to
$\frac{1}{2}\{S, S\}_{H_\cA} + \hbar\,\Delta\,S = 0$
(the BV classical master equation).  At genus~$0$ ($\hbar = 0$):
$\frac{1}{2}\{S, S\}_H = 0$, which is the Euler--Lagrange
condition for the BV action.
\end{proof}

\begin{remark}[The four potentials as four gravity theories]
% label removed: rem:thqg-four-gravity-theories
\index{gravitational action!four types}
\begin{center}
\small
\renewcommand{\arraystretch}{1.25}
\begin{tabular}{cll}
\toprule
\emph{Class} & $S_\cA(x)$ &
\emph{Gravity analogue} \\
\midrule
$\mathbf{G}$ &
  $\frac{\kappa}{2}\,x^2$ &
  Free massless field (Gaussian integral) \\
$\mathbf{L}$ &
  $\frac{\kappa}{2}\,x^2 + \frac{1}{3!}\mathfrak{C}_0\,x^3$ &
  Chern--Simons gravity ($A\wedge dA + \frac{2}{3}A^3$) \\
$\mathbf{C}$ &
  $\frac{1}{4!}\mathfrak{Q}_0\,x^4$ &
  $\phi^4$ contact theory \\
$\mathbf{M}$ &
  $\sum_{r=2}^\infty \frac{1}{r}S_r\,x^r$ &
  Full quantum gravity (non-polynomial EFT) \\
\bottomrule
\end{tabular}
\end{center}
The classification of gravitational effective actions
by their polynomial degree mirrors the standard EFT
truncation hierarchy: renormalisable ($r \leq 4$) vs.\
non-renormalisable ($r = \infty$).  Class~$\mathbf{M}$
is the gravitational analogue of a non-renormalisable
theory requiring infinitely many counterterms in the
Wilsonian sense, but here the full tower is determined
by the finite OPE data of the boundary algebra, so no
independent counterterms arise.
\end{remark}


% ======================================================================
%  §27.  THE HOLOGRAPHIC SPECTRAL SEQUENCE
% ======================================================================

\subsection{The holographic spectral sequence}
% label removed: subsec:thqg-holographic-ss
\index{holographic spectral sequence|textbf}

The collision filtration on the convolution algebra
$\gAmod$ induces a spectral sequence that separates the
contributions of the four complexity classes.

\begin{construction}[Collision filtration]
% label removed: constr:thqg-collision-filtration
\index{collision filtration|textbf}
The arity filtration
$F^r\gAmod := \bigoplus_{s \geq r}(\gAmod)_s$
defines a complete, exhaustive, separated filtration
of the dg Lie algebra.  The associated graded is
\begin{equation}% label removed: eq:thqg-graded-convolution
\operatorname{gr}^r_F(\gAmod)
\;=\;
(\gAmod)_r
\;\cong\;
\Defcyc(\cA)_r \widehat\otimes \Gmod_r,
\end{equation}
the arity-$r$ piece of the convolution algebra.
\end{construction}

\begin{theorem}[Holographic spectral sequence]
% label removed: thm:thqg-holographic-ss
\ClaimStatusProvedHere
\index{spectral sequence!holographic}
The collision filtration induces a convergent spectral sequence
\begin{equation}% label removed: eq:thqg-ss
E_1^{r,q}
\;=\;
H^{r+q}(\operatorname{gr}^r_F(\gAmod),\, d_2)
\;\Longrightarrow\;
H^{r+q}(\gAmod).
\end{equation}
The differential $d_1\colon E_1^{r,*} \to E_1^{r+1,*}$ is
the obstruction map $o_{r+1}$.  The $E_2$ page is
\begin{equation}% label removed: eq:thqg-e2
E_2^{r,q}
\;=\;
\ker(o_{r+1})/\mathrm{im}(o_r)
\;\text{ at the $(r,q)$ position.}
\end{equation}
\end{theorem}

\begin{proof}
Standard spectral sequence of a filtered dg Lie algebra
(Theorem~\ref{thm:spectral-sequence-filtered-dg}).  The
$E_1$ page computes the cohomology of each graded piece
with respect to the twisted differential $d_2$.  The
$d_1$ differential is the boundary map in the long exact
sequence of the filtration, which is exactly the
obstruction class $o_{r+1}$ at the $(r, r+1)$ position.
Convergence follows from completeness and pronilpotence of
$F^2\gAmod$.
\end{proof}

\begin{proposition}[Collapse criteria by complexity class]
% label removed: prop:thqg-collapse-criteria
\ClaimStatusProvedHere
\index{spectral sequence!collapse criteria}
\begin{enumerate}[label=\textup{(\roman*)}]
\item \emph{Class~$\mathbf{G}$}: The spectral sequence
  collapses at $E_1$.  All differentials $d_r = 0$ for
  $r \geq 1$.  The MC element is purely scalar in arity:
  $\Theta_\cA=\eta\otimes\Gamma_\cA$ for some genus
  coefficient~$\Gamma_\cA\in\widehat{\Gmod}$.  On the proved
  scalar lane this further specializes to
  $\Theta_\cA = \kappa\cdot\eta\otimes\Lambda$.
\item \emph{Class~$\mathbf{L}$}: The spectral sequence
  collapses at $E_2$.  The differential $d_1\colon
  E_1^{2,*} \to E_1^{3,*}$ is nonzero (it is $o_3$);
  $d_1\colon E_1^{3,*} \to E_1^{4,*}$ is zero (it is $o_4 = 0$).
  All $d_r = 0$ for $r \geq 2$.
\item \emph{Class~$\mathbf{C}$}: The spectral sequence
  collapses at $E_2$.  The differential $d_1\colon
  E_1^{3,*} \to E_1^{4,*}$ is nonzero ($o_4 \neq 0$);
  $d_1\colon E_1^{4,*} \to E_1^{5,*}$ is zero ($o_5 = 0$).
\item \emph{Class~$\mathbf{M}$}: The spectral sequence does
  not collapse at any finite page.  There are nonzero
  differentials $d_1$ at infinitely many positions.  The
  abutment is computed by the full inverse limit
  $\Theta_\cA = \varprojlim_r \Theta_\cA^{\leq r}$.
\end{enumerate}
\end{proposition}

\begin{proof}
Parts~(i)--(iii): in each case, the collapse page is
determined by the highest arity at which $o_{r+1} \neq 0$.
For class~$\mathbf{G}$: all $o_r = 0$ for $r \geq 3$,
so $d_1 = 0$ everywhere, giving $E_1$-collapse.
For class~$\mathbf{L}$: $o_3 \neq 0$ but $o_r = 0$ for
$r \geq 4$; the single nonzero differential at arity~$3$
produces the cubic shadow as a nontrivial $E_2$-class,
then all higher differentials vanish.
For class~$\mathbf{C}$: the gap ($o_3 = 0$, $o_4 \neq 0$,
$o_5 = 0$) gives $d_1$ nonzero at position~$3 \to 4$ and
zero elsewhere, again yielding $E_2$-collapse.
Part~(iv): for class~$\mathbf{M}$, the infinite tower
(Theorem~\ref{thm:thqg-virasoro-infinite}) gives
$o_{r+1} \neq 0$ for infinitely many~$r$, so $d_1$ has
infinitely many nonzero entries and the spectral sequence
never degenerates.
\end{proof}


% ======================================================================
%  §28.  THE $E_1$ PAGE: EXPLICIT STRUCTURE
% ======================================================================

\subsection{The $E_1$ page for Virasoro}
% label removed: subsec:thqg-e1-virasoro
\index{holographic spectral sequence!$E_1$ for Virasoro}

\begin{proposition}[$E_1$ page for Virasoro]
% label removed: prop:thqg-e1-virasoro
\ClaimStatusProvedHere
For the Virasoro algebra on the one-generator primary line:
\begin{equation}% label removed: eq:thqg-e1-virasoro
E_1^{r,0}
\;\cong\;
\C\cdot x^r
\quad\text{for all } r \geq 2,
\end{equation}
with differential $d_1(S_r\,x^r) = o_{r+1}(\mathrm{Vir})
\cdot x^{r+1}$.  The $E_2$ page is
\begin{equation}% label removed: eq:thqg-e2-virasoro
E_2^{r,0} = 0
\quad\text{for all } r \geq 2.
\end{equation}
Every $E_1$-class is killed by a differential: the spectral
sequence converges to a point.  This is consistent with the
bar-intrinsic construction: $\Theta_{\mathrm{Vir}}$ exists
and is unique, so the abutment $H^*(\gAmod)$ carries no
residual ambiguity.
\end{proposition}

\begin{proof}
On the rank-one primary line,
$\operatorname{gr}^r_F(\gAmod) \cong \C\cdot x^r$ in
degree zero.  The twisted differential $d_2$ acts by
multiplication by $\kappa = c/2$ and has no kernel beyond
degree zero.  The $d_1$ differential
$d_1(x^r) = o_{r+1}\cdot x^{r+1}$ is nonzero for all
$r \geq 3$ (since $o_{r+1} \neq 0$ for all $r \geq 4$
by the infinite tower), so
$E_2^{r,0} = \ker(d_1^{r,0})/\mathrm{im}(d_1^{r-1,0}) = 0$
for $r \geq 3$: every class at arity~$r$ is the image of
the arity-$(r{-}1)$ class under $d_1$ (up to a nonzero
scalar), and is simultaneously in the kernel of $d_1$ at
arity~$r$ (since $o_{r+2} \neq 0$ means $d_1$ is injective).
\end{proof}


% ======================================================================
%  §29.  THE $E_1$ COLLAPSE FOR AFFINE AND $\beta\gamma$
% ======================================================================

\subsection{Spectral sequence collapse for finite-depth classes}
% label removed: subsec:thqg-ss-collapse

\begin{proposition}[Collapse for class~$\mathbf{L}$]
% label removed: prop:thqg-collapse-L
\ClaimStatusProvedHere
\index{holographic spectral sequence!collapse for affine}
For $\cA = V_k(\fg)$ at generic $k$:
the spectral sequence collapses at $E_2$ with
\begin{equation}% label removed: eq:thqg-collapse-L
E_2^{2,0} \cong \C,
\qquad
E_2^{3,0} \cong \C,
\qquad
E_2^{r,0} = 0 \;\text{for}\; r \geq 4.
\end{equation}
The two surviving classes are $\kappa$ and $\mathfrak{C}$.
\end{proposition}

\begin{proof}
$d_1\colon E_1^{3,0} \to E_1^{4,0}$ is given by $o_4 = 0$
(Jacobi), so $E_1^{3,0}$ survives to $E_2$.
$d_1\colon E_1^{r,0} \to E_1^{r+1,0}$ is zero for all
$r \geq 3$ (no higher obstructions).  The $E_2$ page has
exactly two nonzero entries at genus zero.
\end{proof}

\begin{proposition}[Collapse for class~$\mathbf{C}$]
% label removed: prop:thqg-collapse-C
\ClaimStatusProvedHere
\index{holographic spectral sequence!collapse for $\beta\gamma$}
For $\cA = \beta\gamma$ on the weight-changing line:
the spectral sequence collapses at $E_2$ with
\begin{equation}% label removed: eq:thqg-collapse-C
E_2^{2,0} \cong \C,
\qquad
E_2^{4,0} \cong \C,
\qquad
E_2^{r,0} = 0 \;\text{for}\; r \neq 2, 4.
\end{equation}
The surviving classes are $\kappa$ and $\mathfrak{Q}$.
The gap ($E_2^{3,0} = 0$) is the spectral-sequence
manifestation of the cubic gauge triviality.
\end{proposition}

\begin{proof}
$o_3 = 0$ (abelian OPE), so $d_1\colon E_1^{2,0} \to E_1^{3,0}$
is zero and $E_1^{3,0}$ is not killed; but $o_4 \neq 0$, so
$d_1\colon E_1^{3,0} \to E_1^{4,0}$ is nonzero, killing
$E_1^{3,0}$ at the $E_2$ page (it becomes the image of
$d_1$, which maps to the kernel of the next $d_1$).
More precisely: $E_1^{3,0} \cong \C$ with $d_1 = 0$ from
arity~$2$ and $d_1 \neq 0$ to arity~$4$.  Since $o_3 = 0$,
$\mathrm{im}(d_1^{2,0}) = 0$, but
$\ker(d_1^{3,0}) = 0$ since $o_4 \neq 0$
gives an injective map.  Hence $E_2^{3,0} = 0$.
At arity~$4$: $d_1^{4,0} = 0$ ($o_5 = 0$ by rank-one
rigidity), and $\mathrm{im}(d_1^{3,0}) \neq 0$; but the
quotient $\ker(d_1^{4,0})/\mathrm{im}(d_1^{3,0})$ is
one-dimensional (the quartic class $\mathfrak{Q}$ modulo
the image of the cubic).
\end{proof}


% ======================================================================
%  §30.  CURVE INDEPENDENCE
% ======================================================================

\subsection{Universality: curve independence of complexity}
% label removed: subsec:thqg-curve-independence
\index{shadow depth!curve independence}

\begin{proposition}[Curve independence]
% label removed: prop:thqg-curve-independence
\ClaimStatusProvedHere
\index{gravitational complexity!curve independence}
The gravitational complexity $\rho_{\mathrm{grav}}(\cA)$ is
independent of the curve~$X$ on which $\cA$ lives.  The
obstruction classes $o_r(\cA)$ depend only on the OPE
coefficients, which are curve-independent data.
\end{proposition}

\begin{proof}
By the genus-$0$ curve-independence criterion
(Proposition~\ref{prop:genus0-curve-independence}):
the genus-$0$ shadow algebra $\cA^{\mathrm{sh}}_{r,0}$
depends only on the $\lambda$-bracket and the cyclic
pairing.  These are intrinsic to the vertex algebra,
not to the curve.
\end{proof}


% ======================================================================
%  §31.  EXPLICIT SEXTIC COMPUTATION
% ======================================================================

\subsection{Detailed sextic computation for Virasoro}
% label removed: subsec:thqg-sextic-detail
\index{Virasoro!sextic shadow detailed}

\begin{computation}[The sextic obstruction in full detail]
% label removed: comp:thqg-sextic
\index{obstruction class!sextic Virasoro}
The arity-$6$ obstruction $\mathfrak{o}^{(6)}$ receives
contributions from all two-vertex graphs with arities
$(a,b)$ satisfying $a + b = 7$, $a,b \geq 2$.  The
contributing pairs and their values on the Virasoro
primary line with $P = 2/c$:

\smallskip
\noindent
$(2,5)$: $\{\mathrm{Sh}_2, \mathrm{Sh}_5\}_H
= cx\cdot(2/c)\cdot(-240/[c^2(5c+22)])x^4
= -480/[c^2(5c+22)]\,x^5 \cdot x$. \\
$(3,4)$: $\{\mathfrak{C}, \mathfrak{Q}\}_H
= 6x^2\cdot(2/c)\cdot 4Q_0 x^3
= 48Q_0/c\,x^5\cdot x$. Already computed as part of
$\mathfrak{o}^{(5)}$; the arity-$6$ component requires
different combinatorics. \\
$(3,5)$: $\{\mathfrak{C}, \mathrm{Sh}_5\}_H
= 6x^2\cdot(2/c)\cdot 5S_5 x^4
= 60S_5/c\,x^6$.
With $S_5 = -48/[c^2(5c+22)]$:
$= -2880/[c^3(5c+22)]\,x^6$. \\
$(4,4)$: $\{\mathrm{Sh}_4, \mathrm{Sh}_4\}_H
= 4Q_0 x^3\cdot(2/c)\cdot 4Q_0 x^3
= 32Q_0^2/c\,x^6$.
With $Q_0 = 10/[c(5c+22)]$:
$= 3200/[c^3(5c+22)^2]\,x^6$. \\
$(5,3)$: $\{\mathrm{Sh}_5, \mathfrak{C}\}_H
= 5S_5 x^4\cdot(2/c)\cdot 6x^2
= 60S_5/c\,x^6 = -2880/[c^3(5c+22)]\,x^6$. \\
$(5,2)$: $\{\mathrm{Sh}_5, \mathrm{Sh}_2\}_H$:
this contributes to the genus-loop sector, not the
bracket term.  The arity-$2$ piece $\mathrm{Sh}_2 = (c/2)x^2$
paired with arity-$5$ gives $5S_5 x^4\cdot(2/c)\cdot cx
= 10S_5\,x^5$, but this produces arity $5 + 2 - 1 = 6$
with different symmetry.

\smallskip\noindent
Collecting the three dominant contributions at arity~$6$:
\begin{align*}
\mathfrak{o}^{(6)}
&= \{\mathfrak{C}, \mathrm{Sh}_5\}_H
+ \{\mathrm{Sh}_4, \mathrm{Sh}_4\}_H
+ \{\mathrm{Sh}_5, \mathfrak{C}\}_H
\\
&= -\frac{2880}{c^3(5c+22)}\,x^6
+ \frac{3200}{c^3(5c+22)^2}\,x^6
- \frac{2880}{c^3(5c+22)}\,x^6
\\
&= \frac{-5760(5c+22) + 3200}{c^3(5c+22)^2}\,x^6
\\
&= \frac{-28800c - 126720 + 3200}{c^3(5c+22)^2}\,x^6
\\
&= \frac{-28800c - 123520}{c^3(5c+22)^2}\,x^6.
\end{align*}
Factoring the numerator:
$-28800c - 123520 = -160(180c + 772) = -640(45c + 193)$.
Hence
\begin{equation}% label removed: eq:thqg-o6-explicit
\mathfrak{o}^{(6)}
= \frac{-640(45c+193)}{c^3(5c+22)^2}\,x^6.
\end{equation}
Applying the master recursion~\eqref{eq:thqg-master-recursion}:
\begin{equation}% label removed: eq:thqg-sh6-explicit
\mathrm{Sh}_6
= -\frac{\mathfrak{o}^{(6)}}{2\cdot 6}
= \frac{640(45c+193)}{12\,c^3(5c+22)^2}\,x^6
= \frac{80(45c+193)}{3\,c^3(5c+22)^2}\,x^6.
\end{equation}
This confirms the entry in the table of
Theorem~\ref{thm:thqg-virasoro-tower-explicit}.

The numerator $45c + 193$ vanishes at $c = -193/45$.  At this
value, $\mathrm{Sh}_6 = 0$ but $\mathrm{Sh}_5 \neq 0$ and
$\mathrm{Sh}_7 \neq 0$ (since $15c + 61|_{c=-193/45} =
15(-193/45) + 61 = -193/3 + 61 = (-193 + 183)/3 = -10/3 \neq 0$).
The shadow obstruction tower has a ``zero-crossing'' at arity~$6$ but remains
generically infinite.
\end{computation}


% ======================================================================
%  §32.  OPEN PROBLEMS
% ======================================================================

\subsection{Open problems}
% label removed: subsec:thqg-open-problems
\index{gravitational complexity!open problems}

\begin{openproblem}[Gap question]
% label removed: open:thqg-gap
\index{shadow depth!gap question}
Does there exist a modular Koszul chiral algebra with
$r_{\max} \in \{5, 6, 7, \ldots\} \setminus \{\infty\}$?
All known examples have $r_{\max} \in \{2, 3, 4, \infty\}$.
The exhaustiveness proof of
Theorem~\ref{thm:thqg-four-class} uses properties specific
to the standard landscape (rank-one rigidity for
class~$\mathbf{C}$).  An intermediate class would require:
$o_3 = 0$, $o_4 \neq 0$, $o_5 \neq 0$, $o_6 = 0$: a
quartic-plus-quintic tower that terminates at arity~$5$.
No known algebraic mechanism produces this pattern.
\end{openproblem}

% ----------------------------------------------------------------------
%  Shadow depth as stratification depth of the Steinberg
% ----------------------------------------------------------------------

\subsection{Shadow depth as stratification depth of the Steinberg}
% label removed: subsec:shadow-steinberg-stratification
\index{Steinberg object!stratification depth}
\index{shadow depth!Steinberg interpretation}

The shadow depth classification acquires a clean geometric
interpretation through the Chriss--Ginzburg picture: shadow
depth \emph{is} the stratification depth of the Steinberg
object $\mathfrak{S}_b$ associated to the bar complex.

\begin{definition}[Stratification depth]
% label removed: def:stratification-depth
Let $\mathfrak{S}_b$ denote the Steinberg object associated
to a modular Koszul chiral algebra~$\cA$.  The
\emph{stratification depth} of $\mathfrak{S}_b$ is the
length of the orbit stratification: the maximal~$d$ such
that there exist strata
\[
  S_0 \subsetneq S_1 \subsetneq \cdots \subsetneq S_d
  = \mathfrak{S}_b
\]
with each successive quotient $S_i / S_{i-1}$ non-empty.
\end{definition}

\begin{proposition}[Shadow depth = stratification depth]
% label removed: prop:shadow-equals-stratification
\ClaimStatusHeuristic
The shadow depth $r_{\max}$ of a modular Koszul chiral
algebra equals the stratification depth of the associated
Steinberg object~$\mathfrak{S}_b$.  Concretely:
\begin{enumerate}[label=\textup{(\roman*)}]
\item \textbf{Class~$\mathbf{G}$}
  ($r_{\max} = 2$): $\mathfrak{S}_b$ has two strata,
  the smooth locus plus one singular stratum.
  Examples: Heisenberg, lattice VOA.
\item \textbf{Class~$\mathbf{L}$}
  ($r_{\max} = 3$): three strata.
  Example: affine $V_k(\fg)$, where the Lie bracket
  creates one additional singular stratum.
\item \textbf{Class~$\mathbf{C}$}
  ($r_{\max} = 4$): four strata
  (the gap at depth~$3$ from cubic gauge triviality).
  Example: $\beta\gamma$.
\item \textbf{Class~$\mathbf{M}$}
  ($r_{\max} = \infty$): infinitely many strata.
  Example: Virasoro, $\mathcal{W}_N$ for $N \geq 3$, where
  higher-spin poles produce unbounded stratification.
\end{enumerate}
\end{proposition}

\begin{proof}[Argument]
Each stratum of $\mathfrak{S}_b$ corresponds to a distinct
collision type in $\mathrm{FM}_k(\mathbb{C})$ at which the
bar differential acquires a new pole order.  The bar
differential $d_B$ is assembled from OPE residues on
$\mathrm{FM}_k(\mathbb{C})$; a stratum $S_i$ is
\emph{non-trivial} precisely when the corresponding
collision produces a pole of order $\leq i$ that does not
factor through lower strata.  The identification of shadow
operations $\mathrm{Sh}_r$ with these strata rests on
the following chain: $\mathrm{Sh}_r \neq 0$ detects a
non-vanishing residue at pole order $r$ in the OPE, which
in turn produces a non-empty collision stratum at depth~$r$
in $\mathrm{FM}_k(\mathbb{C})$.  The converse (that every
non-empty stratum is detected by some $\mathrm{Sh}_r$)
follows from the completeness of the residue calculus on FM
strata.  We verify this in all three standard classes
(items~(i)--(iii) above); the general statement requires the
FM residue completeness axiom (H3).
\end{proof}

\begin{proposition}[Geometric translation of the gap question]
% label removed: prop:gap-steinberg
\ClaimStatusHeuristic
The gap question \textup{(}Open
Problem~\ref{open:thqg-gap}\textup{)} translates to: does
there exist a Steinberg object $\mathfrak{S}_b$ with
finitely many strata but more than four?

This is constrained by the representation theory of the
Fulton--MacPherson operad: the collision types in
$\mathrm{FM}_k(\mathbb{C})$ are classified by nested
parenthesizations, and the pole structure of the OPE
determines which collision types are non-trivial.  For a
polynomial OPE of degree~$d$ (pole of order~$d$), the
number of distinct collision strata grows at most linearly
in~$d$.  Finite shadow depth $> 4$ therefore requires a
very specific OPE pole pattern: poles of exactly orders
$1$ through $r_{\max}$ with all higher poles vanishing.
This is an extremely rigid condition on the chiral algebra.
\end{proposition}

\begin{proof}
The collision strata of $\mathrm{FM}_k(\mathbb{C})$ are
indexed by trees; each internal edge carries a pole order
from the OPE.  If the OPE has poles of orders
$1, \ldots, d$ and no higher, then the maximal
stratification depth is at most~$d$: each new pole order
contributes at most one new stratum layer.  Conversely, if
a pole of order~$j$ vanishes, the corresponding stratum
collapses and the depth decreases.  To achieve
$r_{\max} = r$ with $r$ finite, one needs non-vanishing
poles at orders $1, \ldots, r$ and vanishing at all orders
$> r$.  For $r > 4$, this means the OPE must have a sharp
cutoff at order~$r$, a condition not satisfied by any
algebra in the standard landscape, where either the pole
order is $\leq 4$ (classes $\mathbf{G}$, $\mathbf{L}$,
and $\mathbf{C}$) or grows without bound (class~$\mathbf{M}$).
\end{proof}

\begin{remark}[Analogy with Coxeter numbers]
% label removed: rmk:gap-coxeter
The gap question is analogous to asking whether there exist
Weyl groups with Coxeter number between~$5$ and~$\infty$
that are not in the standard ADE classification, a
representation-theoretic rigidity question.  Just as the
Coxeter number is constrained by the root system, the
stratification depth is constrained by the FM operadic
structure and the pole data of the OPE.
\end{remark}

\begin{openproblem}[Non-principal complexity]
% label removed: open:thqg-non-principal
\index{non-principal W-algebra!complexity}
For non-principal W-algebras $\mathcal{W}_k(\fg, f)$:
what is the shadow depth class?  The hook-type corridor
in type~A is proved
(Volume~I); the general case
is open.  Expectation: all non-abelian non-principal
W-algebras are class~$\mathbf{M}$ (the DS reduction
introduces composite fields at all arities).
\end{openproblem}

\begin{openproblem}[Categorical complexity]
% label removed: open:thqg-categorical
\index{categorical shadow depth}
Is there a categorical lift of the shadow depth that
detects $r_{\max}$ at the level of derived categories
without passing through explicit OPE computations?
Candidate: the $t$-structure filtration depth on the
derived category of modules over $\cA$, measured by
the Ext-vanishing pattern.
\end{openproblem}

% ----------------------------------------------------------------------
%  The t-structure interpretation (Chriss--Ginzburg attack)
% ----------------------------------------------------------------------

\subsection{The \texorpdfstring{$t$}{t}-structure interpretation}
% label removed: subsec:t-structure-interpretation
\index{t-structure@$t$-structure!amplitude}
\index{shadow depth!$t$-structure detection}

\providecommand{\Steinberg}{\mathfrak{S}}

\begin{conjecture}[Shadow depth as $t$-structure amplitude]
\label{prop:shadow-t-amplitude}
\ClaimStatusConjectured
Let $\cA$ be a modular Koszul chiral algebra with
Steinberg object $\Steinberg_b$.  Then the shadow depth
$r_{\max}$ equals the \emph{amplitude} of the natural
$t$-structure on $D^b(\mathrm{Coh}(\Steinberg_b))$: the
length of the longest non-trivial filtration in the heart.

In classical geometric representation theory, the
perverse $t$-structure on $D^b(\mathrm{Perv}(G/B))$ has
amplitude equal to the length $\ell(w_0)$ of the longest
element of the Weyl group~$W$.  The chiral analogue reads:
\[
  \mathrm{amp}\bigl(D^b(\mathrm{Coh}(\Steinberg_b))\bigr)
  \;=\;
  \max\bigl\{\, k \mid H^k(B(\cA)) \neq 0 \,\bigr\}
  \;=\;
  r_{\max}.
\]
That is, the $t$-structure amplitude on $D^b(\Steinberg_b)$
coincides with the maximal bar degree carrying non-vanishing
bar cohomology.
\end{conjecture}

\begin{remark}[Intrinsic detection of complexity class]
% label removed: rmk:t-structure-intrinsic
Conjecture~\ref{prop:shadow-t-amplitude} provides a
categorical (intrinsic) detection of $r_{\max}$ without
explicit OPE computation.  To determine whether a chiral
algebra belongs to class $\mathbf{G}$, $\mathbf{L}$,
$\mathbf{C}$, or $\mathbf{M}$, one computes the amplitude of the natural
$t$-structure on the derived category of the Steinberg
object, a purely homological-algebraic invariant.
This answers Open Problem~\ref{open:thqg-categorical}
affirmatively, modulo the conjecture above.
\end{remark}

\begin{remark}[Simple objects and stratification]
% label removed: rmk:t-structure-simples
For perverse sheaves on the classical Steinberg variety,
the $t$-structure amplitude equals the number of simple
perverse sheaves, which in turn equals $|W|$.  For the
chiral Steinberg $\Steinberg_b$, the r\^ole of simple
objects is played by the irreducible components of the
orbit stratification; their number is the shadow depth
$r_{\max}$.  The passage from $|W|$ to $r_{\max}$ is the
chiral shadow of the passage from the Weyl group to the
OPE pole pattern.
\end{remark}

% Section file for Chapter: Twisted Holography and Quantum Gravity
% Result (G3): Complementarity as Shifted-Symplectic Polarization

% Local macros (providecommand only; never \newcommand in chapter files)
\providecommand{\MC}{\mathrm{MC}}
\providecommand{\Defcyc}{\mathrm{Def}_{\mathrm{cyc}}}
\providecommand{\Definfmod}{\mathrm{Def}^{\mathrm{mod}}_\infty}
\providecommand{\Sh}{\mathrm{Sh}}
\providecommand{\gr}{\operatorname{gr}}
\providecommand{\id}{\mathrm{id}}
\providecommand{\Tr}{\operatorname{Tr}}
\providecommand{\Sym}{\operatorname{Sym}}
\providecommand{\Hom}{\operatorname{Hom}}
\providecommand{\End}{\operatorname{End}}
\providecommand{\Spec}{\operatorname{Spec}}
\providecommand{\Res}{\operatorname{Res}}
\providecommand{\rank}{\operatorname{rank}}
\providecommand{\ad}{\operatorname{ad}}
\providecommand{\Fred}{\operatorname{Fred}}
\providecommand{\Map}{\operatorname{Map}}
\providecommand{\RGamma}{R\Gamma}
\providecommand{\Perf}{\operatorname{Perf}}
\providecommand{\BTZ}{\mathrm{BTZ}}

\section{Complementarity as shifted-symplectic polarization}
% label removed: sec:thqg-symplectic-polarization
\label{ch:thqg-symplectic-polarization}% alias for dangling \ref{ch:...}
\index{symplectic polarization|textbf}
\index{complementarity!gravitational phase space}
\index{shifted symplectic!complementarity!holographic}

The complementarity theorem (Volume~I, Theorem~\ref{V1-thm:quantum-complementarity-main})
decomposes the ambient complex
$\mathbf{C}_g(\cA) = \RGamma(\overline{\mathcal{M}}_g,
\mathcal{Z}(\cA))$ into Lagrangian halves
$\mathbf{Q}_g(\cA) \oplus \mathbf{Q}_g(\cA^!)$.
Theorem~\ref{V1-thm:shifted-symplectic-complementarity} upgrades this to a
$(-1)$-shifted symplectic polarization via the BV formalism.
The ambient complex carries a Verdier involution whose eigenspace
decomposition is unconditional on the Koszul locus~(C1), while the
shifted-symplectic Lagrangian structure~(C2) is the natural
geometric home of the shadow obstruction tower
$\Theta_\cA^{\le 2} \to \Theta_\cA^{\le 3} \to \Theta_\cA^{\le 4}
\to \cdots$. We prove that the complementarity potential $S_\cA$
generates the dual Lagrangian as a formal graph, that its Taylor
jets are exactly the shadow jets
$H_\cA$, $\mathfrak{C}_\cA$, $\mathfrak{Q}_\cA$ of
Appendix~\ref*{V1-app:nonlinear-modular-shadows}, and that at genus~$1$
the construction reduces to a precise bulk--boundary entanglement
computation matching the BTZ partition function.

\medskip
\noindent\emph{Organization.}\;
\S\ref{subsec:thqg-III-ambient-complex} constructs the ambient
complex and Verdier involution, establishing the foundational
input.
\S\ref{subsec:thqg-III-eigenspace-decomposition} proves the
unconditional eigenspace decomposition~(C1) with complete detail.
\S\ref{subsec:thqg-III-shifted-symplectic} develops the
shifted-symplectic structure~(C2), including a self-contained
review of PTVV geometry.
\S\ref{subsec:thqg-III-complementarity-potential} constructs the
complementarity potential and identifies it with the shadow jet
expansion.
\S\ref{subsec:thqg-III-holographic-entanglement} gives the
holographic interpretation at genus~$1$ and the BTZ connection.
\S\ref{subsec:thqg-III-standard-landscape} verifies all structures
across the standard landscape.

% ======================================================================
%
% 1. THE AMBIENT COMPLEX AND VERDIER INVOLUTION
%
% ======================================================================

\subsection{The ambient complex and Verdier involution}
\label{subsec:thqg-III-ambient-complex}
\index{ambient complex|textbf}
\index{Verdier involution!on ambient complex}

The starting point is the datum of a chiral Koszul pair
$(\cA, \cA^!)$ on a smooth projective curve $X$ over $\mathbb{C}$.
The five main theorems (A--D+H) of Volume~I supply the following
chain of identifications.

\begin{definition}[Holographic ambient complex]
% label removed: def:thqg-III-holographic-ambient
\index{holographic ambient complex|textbf}
Let $(\cA, \cA^!)$ be a chiral Koszul pair and let $g \ge 0$.
The \emph{holographic ambient complex} at genus~$g$ is the cochain
complex
\begin{equation}% label removed: eq:thqg-III-holographic-ambient
\mathbf{C}_g(\cA)
:= \RGamma(\overline{\mathcal{M}}_g, \mathcal{Z}(\cA)),
\end{equation}
where $\mathcal{Z}(\cA)$ is the center local system on the
Deligne--Mumford--Knudsen compactification
$\overline{\mathcal{M}}_g$. When $g \ge 2$ and the relative bar
family is perfect
(Lemma~\ref{V1-lem:perfectness-criterion}), the fiber--center
identification
(Theorem~\ref{V1-thm:fiber-center-identification}) gives
\begin{equation}% label removed: eq:thqg-III-bar-center
\mathbf{C}_g(\cA)
\simeq
\RGamma(\overline{\mathcal{M}}_g,
R^0\pi_{g*}\barB^{(g)}(\cA)).
\end{equation}
\end{definition}

\begin{remark}[H/M/S layers]% label removed: rem:thqg-III-hms
Following Convention~\ref*{V1-conv:hms-levels}:
the complex $\mathbf{C}_g(\cA)$ is the H-level (homotopy)
object; its cohomology
$\mathcal{H}_g(\cA) := H^*(\mathbf{C}_g(\cA))
= H^*(\overline{\mathcal{M}}_g, \mathcal{Z}(\cA))$
is the S-level (shadow) object; and the bar complex
$(\barB^{(g)}(\cA), \Dg{g})$ is the M-level (model) object.
\end{remark}

\begin{proposition}[Properties of the holographic ambient complex;
\ClaimStatusConditional]
\label{prop:thqg-III-ambient-properties}
\index{holographic ambient complex!properties}
The holographic ambient complex
$\mathbf{C}_g(\cA)$ satisfies:
\begin{enumerate}[label=\textup{(\roman*)}]
\item \emph{Finite-dimensionality.}
 For each $n \in \mathbb{Z}$,
 $H^n(\mathbf{C}_g(\cA))$ is finite-dimensional over
 $\mathbb{C}$, and is nonzero only for
 $0 \le n \le 6g - 6$.

\item \emph{Verdier self-duality.}
 The Koszul pairing
 $\cA \otimes \cA^! \to \omega_X$ and the center isomorphism
 $Z(\cA) \cong Z(\cA^!)$
 (Lemma~\ref{V1-sublem:center-isomorphism}) induce a
 quasi-isomorphism of constructible sheaves on
 $\overline{\mathcal{M}}_g$:
 \begin{equation}% label removed: eq:thqg-III-verdier-self-duality
 \mathcal{Z}(\cA)
 \xrightarrow{\;\sim\;}
 \mathbb{D}\,\mathcal{Z}(\cA)[3g - 3],
 \end{equation}
 where $\mathbb{D}$ denotes Verdier duality on
 $\overline{\mathcal{M}}_g$.

\item \emph{Functoriality.}
 The assignment $(\cA, \cA^!) \mapsto \mathbf{C}_g(\cA)$
 is functorial in morphisms of Koszul pairs.
\end{enumerate}
\end{proposition}

\begin{proof}
\emph{Part (i).}
The cohomological dimension of $\overline{\mathcal{M}}_g$ is
$2 \dim_{\mathbb{C}} \overline{\mathcal{M}}_g = 2(3g - 3) = 6g - 6$
by Artin vanishing for proper DM stacks
\cite[\S4.1]{Olsson16}. The stalks of
$\mathcal{Z}(\cA)$ are finite-dimensional by hypothesis
(finite-dimensional fiber cohomology,
Lemma~\ref{V1-lem:perfectness-criterion}(ii)).
By constructibility of $\mathcal{Z}(\cA)$
(Lemma~\ref{V1-lem:quantum-ss-convergence}, condition (2)),
each cohomology sheaf $R^q\pi_{g*}\barB^{(g)}(\cA)$
is a constructible sheaf on the Noetherian stack
$\overline{\mathcal{M}}_g$. Cohomology of a constructible
sheaf on a Noetherian DM stack is finite-dimensional
\cite[Theorem~18.1.1]{SAG}.

\emph{Part (ii).}
The Koszul pairing restricts to a perfect pairing on centers
$Z(\cA) \otimes Z(\cA^!) \to \mathbb{C}$
(Corollary~\ref{V1-cor:duality-bar-complexes-complete} and
Lemma~\ref{V1-sublem:center-isomorphism}).
The center isomorphism $Z(\cA) \cong Z(\cA^!)$ converts the
pairing into a self-pairing on $\mathcal{Z}(\cA)$.
By Verdier duality for constructible sheaves on smooth proper
DM stacks, a nondegenerate self-pairing on $\mathcal{Z}(\cA)$
is equivalent to an isomorphism
$\mathcal{Z}(\cA) \xrightarrow{\sim}
\mathbb{D}\,\mathcal{Z}(\cA)[3g - 3]$.

\emph{Part (iii).}
The center local system $\mathcal{Z}(-)$ is functorial in
morphisms of chiral algebras (being a homotopy fiber of the
bar complex, which is functorial by
Theorem~\ref{V1-thm:geometric-equals-operadic-bar}).
Applying $\RGamma(\overline{\mathcal{M}}_g, -)$ preserves
functoriality.
\end{proof}

\begin{remark}[Attribution: Vol~I dependencies for
Proposition~\ref{prop:thqg-III-ambient-properties}]
\label{rem:thqg-III-ambient-properties-attribution}
The proof of
Proposition~\ref{prop:thqg-III-ambient-properties} is
\emph{conditional} on the following Volume~I ingredients, cited at
the citation level rather than re-inscribed here: the perfectness
criterion for the relative bar family
(Lemma~\ref{V1-lem:perfectness-criterion}), the quantum
semisimplicity convergence lemma supplying constructibility of the
center local system
(Lemma~\ref{V1-lem:quantum-ss-convergence}), the center--center
isomorphism
(Lemma~\ref{V1-sublem:center-isomorphism}), the duality of
complete bar complexes
(Corollary~\ref{V1-cor:duality-bar-complexes-complete}), and the
geometric--operadic bar comparison
(Theorem~\ref{V1-thm:geometric-equals-operadic-bar}). Each of
these is a Volume~I result; this proposition inherits their
scopes (Koszul locus, finite-dimensional fiber cohomology). The
$\ClaimStatusConditional$ tag records this external dependency
per HZ-11 and AP287.
\end{remark}

\begin{construction}[Verdier involution on the ambient complex]
\label{constr:thqg-III-verdier-involution}
\index{Verdier involution!construction|textbf}
The Verdier self-duality of
Proposition~\ref{prop:thqg-III-ambient-properties}(ii) and
the Koszul involutivity
$(\cA^!)^! \simeq \cA$
(Theorem~\ref{V1-thm:chiral-koszul-duality}) combine to produce
a cochain-level endomorphism
\begin{equation}% label removed: eq:thqg-III-sigma-def
\sigma
\colon
\mathbf{C}_g(\cA) \longrightarrow \mathbf{C}_g(\cA)
\end{equation}
constructed as the composition:
\begin{equation}% label removed: eq:thqg-III-sigma-composition
\begin{tikzcd}[column sep=large]
\mathbf{C}_g(\cA)
\arrow[r, "\mathbb{D}"]
&
\mathbf{C}_g(\cA^!)^{\vee}[{-(3g{-}3)}]
\arrow[r, "\mathrm{KS}^{\vee}"]
&
\mathbf{C}_g(\cA)^{\vee\vee}[{-(3g{-}3)}][{3g{-}3}]
\arrow[r, "\mathrm{ev}"]
&
\mathbf{C}_g(\cA)
\end{tikzcd}
\end{equation}
Here $\mathbb{D}$ is the Verdier duality map from
Corollary~\ref{V1-cor:duality-bar-complexes-complete},
$\mathrm{KS}$ is the Koszul identification
$\mathbf{C}_g(\cA) \xrightarrow{\sim} \mathbf{C}_g(\cA^!)$
via Lemma~\ref{V1-sublem:center-isomorphism}, and $\mathrm{ev}$
is the evaluation map. The shifts cancel, and the composition
is well-defined as a cochain map of degree zero.
\end{construction}

\begin{proposition}[Involutivity and anti-symmetry;
\ClaimStatusProvedHere]
\label{prop:thqg-III-involutivity}
\index{Verdier involution!involutivity}
The endomorphism $\sigma$ of
Construction~\textup{\ref{constr:thqg-III-verdier-involution}}
satisfies:
\begin{enumerate}[label=\textup{(\alph*)}]
\item $\sigma^2 = \id$ \textup{(}involutivity\textup{)}.
\item The Verdier pairing
 $\langle -, - \rangle_{\mathbb{D}} \colon
 \mathbf{C}_g(\cA) \otimes \mathbf{C}_g(\cA) \to
 \mathbb{C}[{-(3g{-}3)}]$
 satisfies the anti-symmetry
 \begin{equation}% label removed: eq:thqg-III-anti-symmetry
 \langle \sigma(v), \sigma(w) \rangle_{\mathbb{D}}
 = -\langle v, w \rangle_{\mathbb{D}}
 \end{equation}
 for all $v, w \in \mathbf{C}_g(\cA)$.
\end{enumerate}
\end{proposition}

\begin{proof}
\emph{Part (a).}
Koszul involutivity
$(\cA^!)^! \simeq \cA$ gives $\mathrm{KS}^{-1} = \mathrm{KS}'$
(the transpose identification). Verdier duality satisfies
$\mathbb{D}^2 = \id$ on constructible sheaves over
smooth proper DM stacks (Poincar\'{e} duality is involutive).
The composition $\sigma$ is therefore the identity composed
twice, once from each factor:
\[
\sigma^2
= (\mathrm{ev} \circ \mathrm{KS}'^{\vee} \circ \mathbb{D})
 \circ
 (\mathrm{ev} \circ \mathrm{KS}^{\vee} \circ \mathbb{D})
= \id.
\]

\emph{Part (b).}
Let $v, w \in \mathbf{C}_g(\cA)$. By definition of the
Verdier pairing:
\[
\langle \sigma(v), \sigma(w) \rangle_{\mathbb{D}}
= \langle \mathbb{D}(v'), \mathbb{D}(w') \rangle_{\mathbb{D}}
\]
where $v' = \mathrm{KS}(v)$ and $w' = \mathrm{KS}(w)$.
The Verdier pairing satisfies $\mathbb{D}$ anti-commutativity
with the Kodaira--Spencer action
(Theorem~\ref{V1-thm:kodaira-spencer-chiral-complete},
equation~\eqref{V1-eq:verdier-ks-anticommute}). At the level of
the bilinear form, this gives
$\langle \mathbb{D}(x), \mathbb{D}(y) \rangle_{\mathbb{D}}
= -\langle x, y \rangle_{\mathbb{D}}$
(Proposition~\ref{V1-prop:lagrangian-eigenspaces}(ii)).
Composing with the Koszul identification (which preserves the
pairing up to sign by construction) gives the stated
anti-symmetry.
\end{proof}

\begin{proposition}[Compatibility with the modular MC element;
\ClaimStatusProvedHere]
\label{prop:thqg-III-mc-compatibility}
\index{Maurer--Cartan element!Verdier compatibility}
Let $\Theta_\cA \in \MC(\mathfrak{g}_\cA^{\mathrm{mod}})$ be the
bar-intrinsic MC element
\textup{(}Theorem~\textup{\ref*{V1-thm:mc2-bar-intrinsic})}.
The Verdier involution $\sigma$ intertwines the shadow projections
of $\Theta_\cA$ and $\Theta_{\cA^!}$:
\begin{equation}% label removed: eq:thqg-III-theta-intertwine
\sigma \circ \pi_g^*(\Theta_\cA)
= -\pi_g^*(\Theta_{\cA^!}),
\end{equation}
where $\pi_g^*$ denotes the genus-$g$ projection.
\end{proposition}

\begin{proof}
The MC element $\Theta_\cA := D_\cA - d_0$ satisfies the MC
equation because $D_\cA^2 = 0$
(Theorem~\ref*{V1-thm:mc2-bar-intrinsic}).
Under Verdier--Koszul duality
(Theorem~\ref*{V1-thm:verdier-bar-cobar}), the bar
coderivation $D_\cA$ maps to $D_{\cA^!}$ while the genus-$0$
collision differential $d_0$ is invariant. At the level of
coderivations, therefore,
$\sigma(\Theta_\cA) = D_{\cA^!} - d_0 = \Theta_{\cA^!}$.
The sign in~\eqref{V1-eq:thqg-III-theta-intertwine} arises when
one passes to genus-$g$ projections through the bilinear
pairing. By
Proposition~\ref{V1-prop:lagrangian-eigenspaces}(ii), $\sigma$
acts anti-symmetrically on the pairing:
$\langle \sigma(x), \sigma(y) \rangle
= -\langle x, y \rangle$.
The projection $\pi_g^*$ extracts the component in Hodge
degree~$g$ via this pairing, so the anti-symmetry contributes
an overall sign:
\[
\sigma \circ \pi_g^*(\Theta_\cA)
= -\pi_g^*(D_{\cA^!} - d_0)
= -\pi_g^*(\Theta_{\cA^!}).
\]
This is the complementarity relation
$\pi_g^*(\Theta_\cA) + \pi_g^*(\Theta_{\cA^!}) = 0$
at every genus.
\end{proof}

\begin{remark}[The chain-level mechanism]
% label removed: rem:thqg-III-chain-level
\index{complementarity!chain-level mechanism}
Proposition~\ref{prop:thqg-III-mc-compatibility} is the
chain-level incarnation of the complementarity relation
$\sigma \circ \pi_g^*(\Theta_\cA) = -\pi_g^*(\Theta_{\cA^!})$
(Remark~\ref{V1-rem:theorem-C-decomposition}).
At the S-level, the pairing-projected relation becomes
$\kappa(\cA) + \kappa(\cA^!) = 0$ for Kac--Moody and free-field algebras,
but $\kappa(\cA) + \kappa(\cA^!) = 13$ for Virasoro and
$\kappa + \kappa^! = \varrho_{\cA}\,K$ for general $\mathcal{W}$-algebras
(with $\varrho_{\cW_N} = H_N - 1$ and $K = 4N^3 - 2N - 2$,
$\varrho_{\BP} = 1/6$ and $K_{\BP} = 196$):
the complementarity sum is family-dependent, not universally
zero (Volume~I, Theorem~D).
At the H-level, $\sigma$ intertwines the two MC elements in
opposite sign, implementing the deformation--obstruction
exchange.
\end{remark}

\begin{definition}[Holographic Verdier pairing]
\label{def:thqg-III-holographic-pairing}
\index{Verdier pairing!holographic|textbf}
The \emph{holographic Verdier pairing} at genus~$g$ is the
bilinear form
\begin{equation}% label removed: eq:thqg-III-holographic-pairing-def
\langle -, - \rangle_g
\colon
\mathbf{C}_g(\cA) \otimes \mathbf{C}_g(\cA)
\longrightarrow
\mathbb{C}[{-(3g{-}3)}]
\end{equation}
induced by the Verdier self-duality
\eqref{V1-eq:thqg-III-verdier-self-duality} and the $\RGamma$
functor. Explicitly, for cochains
$\alpha \in C^p(\overline{\mathcal{M}}_g, \mathcal{Z}(\cA))$
and
$\beta \in C^{(3g-3)-p}(\overline{\mathcal{M}}_g,
\mathcal{Z}(\cA))$,
\begin{equation}% label removed: eq:thqg-III-pairing-explicit
\langle \alpha, \beta \rangle_g
:=
\int_{\overline{\mathcal{M}}_g}
\alpha \cup \beta
\;\in\; \mathbb{C},
\end{equation}
where the cup product uses the self-pairing on
$\mathcal{Z}(\cA)$ and integration uses the fundamental class
of $\overline{\mathcal{M}}_g$.
\end{definition}

\begin{lemma}[Nondegeneracy of the holographic pairing;
\ClaimStatusProvedHere]
\label{lem:thqg-III-nondegeneracy}
\index{Verdier pairing!nondegeneracy}
When the relative bar family is perfect
\textup{(}Lemma~\textup{\ref{V1-lem:perfectness-criterion})},
the holographic Verdier pairing
\eqref{V1-eq:thqg-III-holographic-pairing-def} is nondegenerate.
\end{lemma}

\begin{proof}
Perfectness of $R\pi_{g*}\barB^{(g)}(\cA)$ ensures that
the Verdier duality map
$\mathcal{Z}(\cA) \xrightarrow{\sim}
\mathbb{D}\,\mathcal{Z}(\cA)[3g - 3]$
is a quasi-isomorphism of perfect complexes on the smooth
proper DM stack $\overline{\mathcal{M}}_g$. Applying $\RGamma$
and Serre duality on the stack gives a perfect pairing on
cohomology. On the cochain level, the cup-product pairing
\[
C^p(\overline{\mathcal{M}}_g, \mathcal{Z})
\otimes
C^{(3g-3)-p}(\overline{\mathcal{M}}_g, \mathbb{D}\mathcal{Z})
\to
C^{3g-3}(\overline{\mathcal{M}}_g, \omega_{\overline{\mathcal{M}}_g})
\xrightarrow{\int}
\mathbb{C}
\]
is nondegenerate by the duality theorem for constructible sheaves
on smooth proper varieties
(Grothendieck--Verdier duality, \cite[Theorem~3.1.10]{KS90}).
\end{proof}

\begin{proposition}[Degree shift at each genus;
\ClaimStatusProvedHere]
% label removed: prop:thqg-III-degree-shift
The holographic Verdier pairing has cohomological degree
$-(3g - 3)$. In particular:
\begin{enumerate}[label=\textup{(\alph*)}]
\item At genus $0$: degree $+3$, pairing
 $H^0 \otimes H^0 \to \mathbb{C}[3]$
 \textup{(}trivially degenerate on one-dimensional $H^0$;
 the complementarity is purely at the S-level\textup{)}.
\item At genus $1$: degree $0$, pairing
 $H^0 \otimes H^0 \to \mathbb{C}$
 \textup{(}classical pairing\textup{)}.
\item At genus $g \ge 2$: degree $-(3g - 3)$,
 giving a genuine shifted-symplectic structure.
\end{enumerate}
\end{proposition}

\begin{proof}
The degree is $-(3g - 3) = -\dim_{\mathbb{C}}
\overline{\mathcal{M}}_g$, which is the Verdier shift for
Serre duality on a smooth variety of dimension $3g - 3$.
The three cases follow by substituting $g = 0, 1, \ge 2$.
\end{proof}

\begin{remark}[Genus $1$ distinguished]
% label removed: rem:thqg-III-genus-1-special
\index{genus 1!holographic pairing}
At genus~$1$, $\dim \overline{\mathcal{M}}_{1,1} = 1$
but $\dim \overline{\mathcal{M}}_1 = 0$ (a point for the
coarse moduli). The relevant moduli space for marked curves
is $\overline{\mathcal{M}}_{1,1}$, where the pairing has
degree~$0$. This is the reason genus~$1$ complementarity
has the simplest form:
$Q_1(\cA) \oplus Q_1(\cA^!) \cong H^*(\overline{\mathcal{M}}_{1,1},
\mathcal{Z}(\cA))$ with a classical (unshifted) pairing.
\end{remark}

% ======================================================================
%
% 2. THE UNCONDITIONAL EIGENSPACE DECOMPOSITION (C1)
%
% ======================================================================

\subsection{The unconditional eigenspace decomposition}
\label{subsec:thqg-III-eigenspace-decomposition}
\index{eigenspace decomposition!holographic|textbf}
\index{complementarity!C1 decomposition}

The eigenspace decomposition~(C1) holds unconditionally on the
Koszul locus, requiring no perfectness or nondegeneracy hypotheses.
We give a self-contained proof in the holographic setting.

\begin{theorem}[Holographic eigenspace decomposition (C1);
\ClaimStatusProvedHere]
\label{thm:thqg-III-eigenspace-decomposition}
\index{Lagrangian!eigenspace decomposition!holographic}
Let $(\cA, \cA^!)$ be a chiral Koszul pair on a smooth projective
curve $X$ over $\mathbb{C}$, and let $g \ge 0$.

\smallskip\noindent
\textbf{H-level.}\;
The Verdier involution $\sigma$ of
Construction~\textup{\ref{constr:thqg-III-verdier-involution}}
splits the ambient complex:
\begin{equation}% label removed: eq:thqg-III-C1-homotopy
\mathbf{C}_g(\cA)
\;\simeq\;
\mathbf{Q}_g(\cA) \;\oplus\; \mathbf{Q}_g(\cA^!),
\end{equation}
where
$\mathbf{Q}_g(\cA) := \operatorname{fib}(\sigma - \id)$ and
$\mathbf{Q}_g(\cA^!) := \operatorname{fib}(\sigma + \id)$
are the homotopy eigenspaces.

\smallskip\noindent
\textbf{S-level.}\;
On cohomology:
\begin{equation}% label removed: eq:thqg-III-C1-shadow
\mathcal{H}_g(\cA)
:= H^*(\overline{\mathcal{M}}_g, \mathcal{Z}(\cA))
\;=\;
Q_g(\cA) \;\oplus\; Q_g(\cA^!),
\end{equation}
where $Q_g(\cA) = \ker(\sigma - \id)$ and
$Q_g(\cA^!) = \ker(\sigma + \id)$.

\smallskip\noindent
\textbf{Duality.}\;
The Verdier pairing restricts to a perfect duality
\begin{equation}% label removed: eq:thqg-III-C1-duality
Q_g(\cA) \;\cong\; Q_g(\cA^!)^{\vee}.
\end{equation}

\smallskip\noindent
\textbf{Functoriality.}\;
The decomposition is natural in morphisms of Koszul pairs and
compatible with the conformal weight and cohomological degree
gradings.
\end{theorem}

\begin{proof}
The proof proceeds in four stages.

\emph{Stage 1: Projector construction.}
Since $\sigma^2 = \id$
(Proposition~\ref{prop:thqg-III-involutivity}(a)) and we
work over $\mathbb{C}$ (characteristic $\ne 2$), the cochain
maps
\begin{equation}% label removed: eq:thqg-III-projectors
p^+ := \tfrac{1}{2}(\id + \sigma),
\qquad
p^- := \tfrac{1}{2}(\id - \sigma)
\end{equation}
are idempotent cochain endomorphisms of $\mathbf{C}_g(\cA)$
satisfying:
\begin{align}
(p^+)^2 &= p^+, &
(p^-)^2 &= p^-, &
p^+ + p^- &= \id, &
p^+ \circ p^- &= 0.
% label removed: eq:thqg-III-projector-properties
\end{align}
These are immediate from $\sigma^2 = \id$. Explicitly:
\[
(p^+)^2
= \tfrac{1}{4}(\id + \sigma)^2
= \tfrac{1}{4}(\id + 2\sigma + \sigma^2)
= \tfrac{1}{4}(\id + 2\sigma + \id)
= \tfrac{1}{2}(\id + \sigma)
= p^+.
\]
The verification for $p^-$ is identical. For the product:
\[
p^+ \circ p^-
= \tfrac{1}{4}(\id + \sigma)(\id - \sigma)
= \tfrac{1}{4}(\id - \sigma^2)
= 0.
\]
Since $\sigma$ commutes with the differential
$d$ on $\mathbf{C}_g(\cA)$ (it is a cochain map by
construction), both $p^+$ and $p^-$ are cochain maps.

\emph{Stage 2: Direct sum decomposition.}
Define
\[
\mathbf{C}_g^+ := \operatorname{im}(p^+),
\qquad
\mathbf{C}_g^- := \operatorname{im}(p^-).
\]
The idempotent decomposition $p^+ + p^- = \id$ gives
\begin{equation}% label removed: eq:thqg-III-direct-sum
\mathbf{C}_g(\cA) = \mathbf{C}_g^+ \oplus \mathbf{C}_g^-
\end{equation}
as cochain complexes (not merely as graded vector spaces,
because $p^\pm$ are cochain maps). Since
$\sigma|_{\mathbf{C}_g^+} = +\id$ and
$\sigma|_{\mathbf{C}_g^-} = -\id$, the subcomplexes
$\mathbf{C}_g^+$ and $\mathbf{C}_g^-$ are the $+1$ and $-1$
eigenspaces of $\sigma$.

It remains to identify $\mathbf{C}_g^+$ with the homotopy
eigenspace $\mathbf{Q}_g(\cA) = \operatorname{fib}(\sigma - \id)$.
The cone sequence
$\operatorname{fib}(\sigma - \id) \to \mathbf{C}_g
\xrightarrow{\sigma - \id} \mathbf{C}_g$
gives $\operatorname{fib}(\sigma - \id) \simeq
\ker(\sigma - \id) = \operatorname{im}(p^+) = \mathbf{C}_g^+$,
where the second identification uses that $\sigma - \id$ is
split-surjective onto $\mathbf{C}_g^-$ (the splitting is given
by $-\tfrac{1}{2}\id$ restricted to $\mathbf{C}_g^-$).
Similarly $\operatorname{fib}(\sigma + \id) \simeq \mathbf{C}_g^-$.

\emph{Stage 3: Eigenspace identification.}
We identify $\mathbf{C}_g^+ \simeq \mathbf{Q}_g(\cA)$ and
$\mathbf{C}_g^- \simeq \mathbf{Q}_g(\cA^!)$.
The $+1$-eigenspace $\mathbf{C}_g^+$ consists of cochains
$\alpha$ with $\sigma(\alpha) = \alpha$. By construction
of $\sigma$
(Construction~\ref{constr:thqg-III-verdier-involution}),
$\sigma(\alpha) = \alpha$ means that $\alpha$ is fixed under
the Verdier--Koszul duality, i.e., $\alpha$ comes from the
$j_*$-extension of bar cochains of $\cA$
(Lemma~\ref{V1-lem:eigenspace-decomposition-complete}).
The genus-$g$ quantum correction space $Q_g(\cA)$ is precisely
the subspace of $H^*(\overline{\mathcal{M}}_g, \mathcal{Z}(\cA))$
arising from bar cochains of $\cA$ in the $j_*$-extension
(Lemma~\ref{V1-lem:quantum-from-ss}).

Symmetrically, $\mathbf{C}_g^-$ consists of cochains fixed by
$-\sigma$, which come from the $j_!$-extension of bar
cochains of $\cA^!$ (equivalently, cobar cochains of $\cA$),
giving $\mathbf{C}_g^- \simeq \mathbf{Q}_g(\cA^!)$.

\emph{Stage 4: S-level decomposition and duality.}
Applying $H^*$ to the direct sum
\eqref{V1-eq:thqg-III-direct-sum} gives:
\[
H^*(\mathbf{C}_g(\cA))
= H^*(\mathbf{C}_g^+) \oplus H^*(\mathbf{C}_g^-)
= Q_g(\cA) \oplus Q_g(\cA^!).
\]
The duality $Q_g(\cA) \cong Q_g(\cA^!)^{\vee}$ follows from
the Verdier pairing. For $v \in Q_g(\cA)$ and
$w \in Q_g(\cA^!)$, the pairing $\langle v, w \rangle_g$
is nondegenerate: if $\langle v, w \rangle_g = 0$ for all
$w \in Q_g(\cA^!)$, then $v$ lies in the radical of the
restricted pairing. Since
$\langle Q_g(\cA), Q_g(\cA) \rangle_g = 0$
(isotropy from anti-symmetry of $\sigma$),
$v$ lies in the radical of the full pairing on
$\mathcal{H}_g(\cA)$, which is zero by nondegeneracy
(Lemma~\ref{lem:thqg-III-nondegeneracy}).

\emph{Note on unconditionality.}
Stages 1--3 use only $\sigma^2 = \id$ and the
identification of $\sigma$-eigenspaces with $j_*/j_!$
extensions. No perfectness or nondegeneracy hypothesis is
needed. The duality statement in Stage~4 uses
nondegeneracy but can also be obtained unconditionally from
spectral sequence duality
(Corollary~\ref{V1-cor:quantum-dual-complete}).

\emph{Functoriality.}
A morphism $f \colon (\cA, \cA^!) \to (\cB, \cB^!)$
of Koszul pairs induces a map
$f_* \colon \mathbf{C}_g(\cA) \to \mathbf{C}_g(\cB)$
by functoriality of the center local system and $\RGamma$.
Since $f$ commutes with Koszul duality
($f^! \colon \cA^! \to \cB^!$ is the Koszul dual map),
$f_*$ commutes with $\sigma$, hence preserves the
projectors $p^\pm$ and the eigenspace decomposition.
\end{proof}

\begin{corollary}[Complementarity exchange principle;
\ClaimStatusProvedHere]
% label removed: cor:thqg-III-complementarity-exchange
\index{complementarity!exchange principle}
For any functional $\Phi$ on chiral Koszul pairs that
factors through the ambient complex, the decomposition
\[
\Phi(\cA) + \Phi(\cA^!) = \Phi_{\mathrm{total}}
\]
holds with $\Phi_{\mathrm{total}}$ independent of the choice
of $\cA$ within its Koszul pair.
\end{corollary}

\begin{proof}
Since $\Phi$ factors through
$\mathcal{H}_g(\cA) = Q_g(\cA) \oplus Q_g(\cA^!)$,
the restriction of $\Phi$ to each summand gives the
decomposition. Independence of $\Phi_{\mathrm{total}}$
from the choice of $\cA$ within $\{\cA, \cA^!\}$
follows from the symmetry $Q_g(\cA) \oplus Q_g(\cA^!)
= Q_g(\cA^!) \oplus Q_g((\cA^!)^!) = Q_g(\cA^!) \oplus Q_g(\cA)$.
\end{proof}

\begin{corollary}[Dimension parity;
\ClaimStatusProvedHere]
% label removed: cor:thqg-III-dimension-parity
For a Koszul self-dual chiral algebra $\cA \cong \cA^!$ and
$g \ge 1$, the total dimension
$\dim \mathcal{H}_g(\cA)$ is even.
\end{corollary}

\begin{proof}
Self-duality gives $Q_g(\cA) \cong Q_g(\cA^!)$, so
$\dim Q_g(\cA) = \dim Q_g(\cA^!)$ and
$\dim \mathcal{H}_g(\cA)
= 2 \dim Q_g(\cA)$
(Theorem~\ref{thm:self-dual-halving}).
\end{proof}

\begin{proposition}[Genus-$0$ complementarity;
\ClaimStatusProvedHere]
% label removed: prop:thqg-III-genus-0
\index{complementarity!genus 0}
At genus~$0$, $\overline{\mathcal{M}}_0 = \mathrm{pt}$ and
$\mathcal{H}_0(\cA) = Z(\cA)$.
The Verdier involution $\sigma$ acts as $+\id$ on $Z(\cA)$,
giving:
\[
Q_0(\cA) = Z(\cA), \qquad Q_0(\cA^!) = 0.
\]
\end{proposition}

\begin{proof}
Since $\overline{\mathcal{M}}_0$ is a point, the center
local system has a single stalk $Z(\cA)$. The bar complex
at genus~$0$ is uncurved (the curvature
$\dfib^{\,2} = \kappa(\cA) \cdot \omega_g$ vanishes
because $\omega_0 = 0$), so the Verdier involution acts on
the genus-$0$ center by the identity: bar cochains with
$j_*$-extension are all of $Z(\cA)$, and there are no
$j_!$-extended cochains (the configuration space has no
boundary at genus~$0$).
\end{proof}

\begin{proposition}[Genus-$1$ complementarity;
\ClaimStatusProvedHere]
\label{prop:thqg-III-genus-1}
\index{complementarity!genus 1}
At genus~$1$,
$H^*(\overline{\mathcal{M}}_{1,1})
= \mathbb{Q} \oplus \mathbb{Q}\lambda_1$
with $\lambda_1 = c_1(\mathbb{E})$ the Hodge class.
For the Heisenberg algebra $\mathcal{H}_\kappa$:
\begin{equation}% label removed: eq:thqg-III-heisenberg-genus-1
Q_1(\mathcal{H}_\kappa) = \mathbb{C} \cdot \kappa,
\qquad
Q_1(\mathcal{H}_\kappa^!) = \mathbb{C} \cdot \lambda_1.
\end{equation}
The central extension $\kappa$ is an obstruction for
$\mathcal{H}_\kappa$ (eigenvalue $+1$); the Hodge class
$\lambda_1$ is the curvature for $\mathcal{H}_\kappa^!$
(eigenvalue $-1$).
\end{proposition}

\begin{proof}
The involution $\sigma$ acts on the
two-dimensional space $H^0 \oplus H^2$ of
$\overline{\mathcal{M}}_{1,1}$. The class $\kappa$ is the
image of the level parameter under the genus-$1$ bar
complex. Since $\kappa$
arises from $j_*$-extended bar cochains, it has eigenvalue
$+1$. The class $\lambda_1$ is the first Chern class of
the Hodge bundle, which arises from the curvature of the
genus-$1$ Gauss--Manin connection; it comes from $j_!$-extended
cochains and has eigenvalue $-1$.

Dimension check: $\dim Q_1 + \dim Q_1^! = 1 + 1 = 2
= \dim H^*(\overline{\mathcal{M}}_{1,1})$.
\end{proof}

\begin{remark}[Deformation--obstruction exchange at genus $1$]
% label removed: rem:thqg-III-def-obs-exchange
\index{deformation--obstruction exchange!genus 1}
The genus-$1$ complementarity is the simplest nontrivial
instance of the exchange principle:
\begin{center}
\begin{tabular}{l@{\quad}c@{\quad}c}
& $\mathcal{H}_\kappa$ & $\mathcal{H}_\kappa^!$ \\
\midrule
Obstruction & $\kappa$ & --- \\
Deformation & --- & $\lambda_1$ \\
\end{tabular}
\end{center}
What $\mathcal{H}_\kappa$ sees as the obstruction to
extending its genus-$0$ structure to genus~$1$ (the
central extension $\kappa$), the Koszul dual
$\mathcal{H}_\kappa^!$ sees as a deformation (the Hodge
class $\lambda_1$ parametrizing the curvature of the
genus-$1$ structure).
\end{remark}

% ======================================================================
%
% 3. SHIFTED-SYMPLECTIC STRUCTURE (C2)
%
% ======================================================================

\subsection{The shifted-symplectic structure}
\label{subsec:thqg-III-shifted-symplectic}
\index{shifted symplectic!structure!holographic|textbf}
\index{PTVV!holographic application}

The unconditional decomposition~(C1) is an eigenspace statement.
The conditional upgrade~(C2) endows the ambient complex with a
$(-1)$-shifted symplectic structure in the sense of
Pantev--To\"{e}n--Vaqui\'{e}--Vezzosi \cite{PTVV13}, making
$Q_g(\cA)$ and $Q_g(\cA^!)$ complementary Lagrangians in the
precise derived-algebraic sense.

\subsubsection{Review of PTVV shifted symplectic geometry}
% label removed: subsubsec:thqg-III-ptvv-review
\index{PTVV!shifted symplectic geometry}

We recall the essential definitions from \cite{PTVV13} in the
form needed for our application.

\begin{definition}[{$n$-shifted symplectic structure
\cite[Definition~1.18]{PTVV13}}]
\label{def:thqg-III-shifted-symplectic}
\index{shifted symplectic!structure|textbf}
Let $\mathbf{F}$ be a derived stack locally of finite
presentation. An \emph{$n$-shifted symplectic structure} on
$\mathbf{F}$ is a closed $2$-form
$\omega \in \mathcal{A}^{2,\mathrm{cl}}(\mathbf{F}, n)$
of degree $n$ such that the induced map
\begin{equation}% label removed: eq:thqg-III-ptvv-nondegen
\omega^{\sharp}
\colon
\mathbb{T}_{\mathbf{F}} \longrightarrow
\mathbb{L}_{\mathbf{F}}[n]
\end{equation}
is a quasi-isomorphism, where $\mathbb{T}_{\mathbf{F}}$ and
$\mathbb{L}_{\mathbf{F}}$ are the tangent and cotangent
complexes.
\end{definition}

\begin{definition}[{Lagrangian in the PTVV sense
\cite[Definition~2.8]{PTVV13}}]
\label{def:thqg-III-ptvv-lagrangian}
\index{Lagrangian!PTVV sense|textbf}
Let $(\mathbf{F}, \omega)$ be an $n$-shifted symplectic derived
stack. A morphism $f \colon \mathbf{L} \to \mathbf{F}$ is
\emph{Lagrangian} if:
\begin{enumerate}[label=\textup{(\alph*)}]
\item \emph{Isotropy:} $f^* \omega = 0$ in
 $\mathcal{A}^{2,\mathrm{cl}}(\mathbf{L}, n)$.
\item \emph{Nondegeneracy:} the induced map
 $\mathbb{T}_{\mathbf{L}} \to
 \mathbb{L}_{\mathbf{L}/\mathbf{F}}[n-1]$
 is a quasi-isomorphism.
\end{enumerate}
Equivalently, condition (b) says the induced map
$\mathbf{L} \to (\mathbf{F}/\mathbf{L})^{\vee}[n]$ is a
quasi-isomorphism, where $\mathbf{F}/\mathbf{L}$ is the
cofiber.
\end{definition}

\begin{example}[Linear shifted symplectic space]
\label{ex:thqg-III-linear-shifted}
\index{shifted symplectic!linear example}
For a cochain complex $V$ with a nondegenerate symmetric pairing
$\langle -, - \rangle \colon V \otimes V \to \mathbb{C}[n]$,
the associated derived affine scheme
$\mathbf{V} = \Spec(\Sym(V^*))$ carries an $n$-shifted
symplectic structure: the pairing defines a constant $2$-form
of degree $n$, and the closure condition
$d_{\mathrm{dR}}\omega = 0$ is automatic for constant forms
on linear spaces \cite[Example~1.4]{PTVV13}.

A subcomplex $L \hookrightarrow V$ is Lagrangian if and only
if:
\begin{enumerate}[label=\textup{(\alph*)}]
\item $\langle L, L \rangle = 0$ (isotropy), and
\item $L \to (V/L)^{\vee}[n]$ is a quasi-isomorphism
 (nondegeneracy).
\end{enumerate}
\end{example}

\begin{proposition}[{Kontsevich--Pridham correspondence
\cite{Pridham17}}]
\label{prop:thqg-III-kontsevich-pridham}
\index{Kontsevich--Pridham correspondence}
Let $\mathfrak{g}$ be a dg Lie algebra with a nondegenerate
invariant pairing of degree $n$:
\[
\langle -, - \rangle
\colon \mathfrak{g} \otimes \mathfrak{g}
\to \mathbb{C}[n],
\qquad
\langle [x,y], z \rangle
= \langle x, [y,z] \rangle.
\]
Then the formal moduli problem
$\MC(\mathfrak{g})$ carries an $n$-shifted symplectic structure.
\end{proposition}

\begin{remark}[Two shifted structures in complementarity]
% label removed: rem:thqg-III-two-structures
\index{shifted symplectic!two structures}
Complementarity carries two shifted symplectic structures:
\begin{enumerate}[label=\textup{(\roman*)}]
\item The \emph{BV structure}: the antibracket on
 $\barB^{\mathrm{ch}}(\cA)$ gives a $(-1)$-shifted symplectic
 structure on $\MC(\barB^{(g)}(\cA)[1])$, universal in $g$.
\item The \emph{Verdier structure}: the holographic pairing on
 $\mathbf{C}_g(\cA)$ gives a $(-(3g-3))$-shifted symplectic
 structure, depending on $g$.
\end{enumerate}
Both make $Q_g(\cA)$ and $Q_g(\cA^!)$ complementary Lagrangians;
the two are compatible via the bar-complex spectral sequence.
\end{remark}

\subsubsection{The BV shifted-symplectic structure}
% label removed: subsubsec:thqg-III-bv-shifted
\index{BV algebra!shifted symplectic|textbf}

\begin{theorem}[BV Lagrangian polarization;
\ClaimStatusProvedHere]
\label{thm:thqg-III-bv-lagrangian}
\index{BV algebra!Lagrangian polarization}
Let $(\cA, \cA^!)$ be a chiral Koszul pair and $g \ge 1$.
\begin{enumerate}[label=\textup{(\roman*)}]
\item The dg Lie algebra
 $L_g := \barB^{(g)}(\cA)[1]$ carries a nondegenerate
 invariant pairing of degree $-1$, making $\MC(L_g)$
 a $(-1)$-shifted symplectic formal moduli problem.

\item The $\sigma$-eigenspace decomposition
 $L_g = L_g^+ \oplus L_g^-$ provides complementary
 Lagrangian subspaces.

\item $L_g^+$ controls deformations of $\cA$ at genus $g$;
 $L_g^-$ controls deformations of $\cA^!$ at genus $g$.
\end{enumerate}
\end{theorem}

\begin{proof}
\emph{Part (i).}
The BV antibracket $\{-, -\}_{\mathrm{BV}}$ on
$\barB^{\mathrm{ch}}(\cA)$ has degree $+1$
(Theorem~\ref{V1-thm:config-space-bv}). Shifting by $[1]$
converts this to a degree-$0$ Lie bracket on $L_g$.
The BV pairing on $\barB^{(g)}(\cA)$ (arising from the
Koszul--Verdier construction,
Corollary~\ref{V1-cor:duality-bar-complexes-complete}) has degree
$+1$. On $L_g = \barB^{(g)}(\cA)[1]$, each of the two inputs
shifts by $[-1]$, giving a pairing of degree
$+1 - 2 = -1$.

Invariance of the pairing follows from the cyclic property of
the BV bracket: the BV Leibniz rule
$\{a, bc\} = \{a, b\}c \pm b\{a, c\}$ implies
$\langle [x, y], z \rangle = \langle x, [y, z] \rangle$
after passing to the shifted Lie bracket.

Nondegeneracy follows from the Verdier intertwining
$\mathbb{D}(\barB(\cA)) \simeq \barB(\cA^!)$
(Theorem~\ref{V1-thm:bv-functor}): on the Koszul locus, the homotopy
Koszul dual~$\cA^!_\infty$ is formal and equivalent to~$\cA^!$, so
bar-cobar inversion (Theorem~\ref*{V1-thm:bar-cobar-isomorphism-main})
ensures the induced map is a quasi-isomorphism.

By the Kontsevich--Pridham correspondence
(Proposition~\ref{prop:thqg-III-kontsevich-pridham}),
the formal moduli problem $\MC(L_g)$ is $(-1)$-shifted
symplectic.

\emph{Part (ii).}
The Verdier involution $\sigma$ induces an involution on $L_g$
that anti-commutes with the Lie bracket:
$\sigma[x, y] = -[\sigma x, \sigma y]$.
The sign reversal comes from the orientation reversal of
$\overline{C}_n(X)$ under $\mathbb{D}$, which reverses the
sign of the symplectic form on configuration space and hence
of the antibracket.

By Proposition~\ref{prop:thqg-III-involutivity}(b), the pairing
satisfies $\langle \sigma(x), \sigma(y) \rangle = -\langle x, y
\rangle$. For $x, y \in L_g^+$ (eigenvalue $+1$):
\[
\langle x, y \rangle
= \langle \sigma(x), \sigma(y) \rangle
= -\langle x, y \rangle,
\]
so $\langle x, y \rangle = 0$. The same holds for $L_g^-$.
Both eigenspaces are isotropic. They are maximal isotropic
(hence Lagrangian) because $L_g = L_g^+ \oplus L_g^-$ is a
direct sum.

\emph{Part (iii).}
By Lemma~\ref{V1-lem:eigenspace-decomposition-complete} and
Lemma~\ref{V1-lem:obs-def-split-complete}, the $+1$-eigenspace
$L_g^+$ corresponds to obstruction classes for $\cA$ at
genus $g$, which by the Koszul deformation--obstruction exchange
are deformation classes for $\cA^!$. Exchanging $\cA$ and
$\cA^!$: $L_g^-$ corresponds to obstructions for $\cA^!$, i.e.,
deformations of $\cA$.
\end{proof}

\subsubsection{The Verdier shifted-symplectic structure}
% label removed: subsubsec:thqg-III-verdier-shifted
\index{shifted symplectic!Verdier structure|textbf}

\begin{theorem}[Verdier Lagrangian polarization (C2);
\ClaimStatusConditional]
\label{thm:thqg-III-lagrangian-polarization}
\index{Lagrangian!polarization!holographic}
Let $(\cA, \cA^!)$ be a chiral Koszul pair on $X$.
Assume:
\begin{enumerate}[label=\textup{(H\arabic*)}]
\item \emph{Perfectness:} $R\pi_{g*}\barB^{(g)}(\cA)$ is
 perfect on $\overline{\mathcal{M}}_g$
 \textup{(}Lemma~\textup{\ref{V1-lem:perfectness-criterion})}.
\item \emph{Nondegeneracy:} the holographic Verdier pairing
 $\langle -, - \rangle_g$ is nondegenerate
 \textup{(}Lemma~\textup{\ref{lem:thqg-III-nondegeneracy})}.
\end{enumerate}
Then for $g \ge 2$:

\smallskip\noindent
\textbf{(C2a) Shifted symplectic structure.}\;
$\mathbf{C}_g(\cA)$ carries a $(-(3g-3))$-shifted symplectic
structure in the PTVV sense
\textup{(}Definition~\textup{\ref{def:thqg-III-shifted-symplectic})}.

\smallskip\noindent
\textbf{(C2b) Lagrangian decomposition.}\;
$\mathbf{Q}_g(\cA)$ and $\mathbf{Q}_g(\cA^!)$ are Lagrangian
subspaces of $\mathbf{C}_g(\cA)$ for this structure, and
\begin{equation}% label removed: eq:thqg-III-lagrangian-decomp
\mathbf{Q}_g(\cA)
\;\simeq\;
\mathbf{Q}_g(\cA^!)^{\vee}[{-(3g{-}3)}].
\end{equation}

\smallskip\noindent
\textbf{(C2c) Anti-symplectomorphism.}\;
The Verdier involution $\sigma$ is an
anti-symplectomorphism:
$\sigma^* \omega_g = -\omega_g$.
\end{theorem}

\begin{proof}
We verify the three parts.

\emph{Part (C2a): shifted symplectic structure.}
The holographic Verdier pairing
$\langle -, - \rangle_g \colon
\mathbf{C}_g \otimes \mathbf{C}_g \to \mathbb{C}[{-(3g{-}3)}]$
(Definition~\ref{def:thqg-III-holographic-pairing}) is a
nondegenerate symmetric pairing of degree $-(3g-3)$ on the
cochain complex $\mathbf{C}_g(\cA)$. By
Example~\ref{ex:thqg-III-linear-shifted} (linear shifted
symplectic space), the associated derived affine scheme
carries a $(-(3g-3))$-shifted symplectic structure.

\emph{Nondegeneracy check.}
The induced map $\omega_g^{\sharp} \colon
\mathbb{T}_{\mathbf{C}_g} \to \mathbb{L}_{\mathbf{C}_g}[-(3g-3)]$
is a quasi-isomorphism because the pairing is nondegenerate
by hypothesis~(H2) and
Lemma~\ref{lem:thqg-III-nondegeneracy}.

\emph{Closure check.}
The $2$-form $\omega_g$ is a constant $2$-form on the linear
space $\mathbf{C}_g$, so $d_{\mathrm{dR}} \omega_g = 0$
automatically.

\emph{Part (C2b): Lagrangian decomposition.}
We verify the two PTVV Lagrangian conditions
(Definition~\ref{def:thqg-III-ptvv-lagrangian}) for
$\mathbf{Q}_g(\cA) \hookrightarrow \mathbf{C}_g(\cA)$.

\emph{Isotropy.}
Let $v, w \in \mathbf{Q}_g(\cA)$ (eigenvalue $+1$). Then
$\sigma(v) = v$ and $\sigma(w) = w$, so
\[
\langle v, w \rangle_g
= \langle \sigma(v), \sigma(w) \rangle_g
= -\langle v, w \rangle_g,
\]
where the last equality uses the anti-symmetry
\eqref{V1-eq:thqg-III-anti-symmetry}. Therefore
$\langle v, w \rangle_g = 0$, verifying isotropy.
The same argument gives isotropy of $\mathbf{Q}_g(\cA^!)$.

\emph{Nondegeneracy.}
Since $\mathbf{C}_g = \mathbf{Q}_g(\cA) \oplus
\mathbf{Q}_g(\cA^!)$
(Theorem~\ref{thm:thqg-III-eigenspace-decomposition}), the
quotient is
$\mathbf{C}_g / \mathbf{Q}_g(\cA) \simeq \mathbf{Q}_g(\cA^!)$.
The induced map
$\mathbf{Q}_g(\cA) \to \mathbf{Q}_g(\cA^!)^{\vee}[-(3g-3)]$
is the restriction of the pairing:
for $v \in \mathbf{Q}_g(\cA)$, the functional
$w \mapsto \langle v, w \rangle_g$ on
$\mathbf{Q}_g(\cA^!)$ is well-defined (the restriction to
$\mathbf{Q}_g(\cA)$ vanishes by isotropy).

This restricted pairing is nondegenerate: suppose
$v \in \mathbf{Q}_g(\cA)$ satisfies
$\langle v, w \rangle_g = 0$ for all
$w \in \mathbf{Q}_g(\cA^!)$. Combined with isotropy
$\langle v, w' \rangle_g = 0$ for
$w' \in \mathbf{Q}_g(\cA)$, this gives
$\langle v, - \rangle_g = 0$ on all of $\mathbf{C}_g$.
By nondegeneracy of the full pairing, $v = 0$.

Therefore the induced map is a quasi-isomorphism, verifying
PTVV nondegeneracy.

\emph{Part (C2c): anti-symplectomorphism.}
By the anti-symmetry \eqref{V1-eq:thqg-III-anti-symmetry}:
$\langle \sigma(v), \sigma(w) \rangle_g
= -\langle v, w \rangle_g$.
Hence
$\sigma^* \omega_g = -\omega_g$, i.e., $\sigma$ is an
anti-symplectomorphism.
\end{proof}

\begin{remark}[Conditionality of (C2)]
% label removed: rem:thqg-III-conditionality
\index{complementarity!conditionality}
Hypotheses (H1) and (H2) are substantive. Perfectness
requires PBW filterability and finite-dimensional fiber
cohomology (Lemma~\ref{V1-lem:perfectness-criterion}).
Nondegeneracy of the Verdier pairing is verified family
by family
(Proposition~\ref{V1-prop:standard-examples-modular-koszul}).
Both hold for the entire standard landscape
(\S\ref{subsec:thqg-III-standard-landscape}). The C1
eigenspace decomposition is strictly weaker, requiring
only $\sigma^2 = \id$, and holds unconditionally.
\end{remark}

\begin{corollary}[Lagrangian intersection is derived critical
locus; \ClaimStatusConditional]
% label removed: cor:thqg-III-lagrangian-intersection
\index{Lagrangian!intersection}
\index{derived critical locus!from complementarity}
Under the hypotheses of
Theorem~\textup{\ref{thm:thqg-III-lagrangian-polarization}},
the derived intersection
\begin{equation}% label removed: eq:thqg-III-lagrangian-intersection
\mathbf{Q}_g(\cA)
\times^h_{\mathbf{C}_g(\cA)}
\mathbf{Q}_g(\cA^!)
\;\simeq\;
\operatorname{Crit}(S_\cA)
\end{equation}
is equivalent to the derived critical locus of the
complementarity potential $S_\cA$
\textup{(}Definition~\textup{\ref{V1-def:nms-complementarity-potential})}.
\end{corollary}

\begin{proof}
This follows from the shifted cotangent normal form
(Theorem~\ref{V1-thm:nms-cotangent-normal-form}) applied to
$\mathcal{M} = \mathbf{C}_g(\cA)$ with its
$(-(3g-3))$-shifted symplectic structure and
$\mathcal{L}_+ = \mathbf{Q}_g(\cA)$,
$\mathcal{L}_- = \mathbf{Q}_g(\cA^!)$.
The tangent map condition is satisfied because the
decomposition $\mathbf{C}_g = \mathbf{Q}_g(\cA) \oplus
\mathbf{Q}_g(\cA^!)$ gives a tangent-level splitting.
Theorem~\ref{V1-thm:nms-derived-critical-locus} then identifies
the derived intersection with $\operatorname{Crit}(S_\cA)$.
\end{proof}

\begin{proposition}[Compatibility of the two symplectic
structures; \ClaimStatusProvedHere]
% label removed: prop:thqg-III-compatibility
\index{shifted symplectic!compatibility}
The BV $(-1)$-shifted structure
\textup{(}Theorem~\textup{\ref{thm:thqg-III-bv-lagrangian})}
and the Verdier $(-(3g-3))$-shifted structure
\textup{(}Theorem~\textup{\ref{thm:thqg-III-lagrangian-polarization})}
are compatible in the following sense:

The bar-complex spectral sequence
\textup{(}Theorem~\textup{\ref{V1-thm:ss-quantum})} induces a
filtered quasi-isomorphism from the BV complex to the Verdier
complex, under which the BV $(-1)$-shifted symplectic form
maps to a degenerate form on $\mathbf{C}_g$ whose
non-degenerate quotient is the $(-(3g-3))$-shifted Verdier
form. Concretely:
\begin{equation}% label removed: eq:thqg-III-compatibility-diagram
\begin{tikzcd}
\MC(L_g) \arrow[r, "{(-1)\text{-shifted}}"] \arrow[d] &
\text{BV antibracket} \arrow[d] \\
\mathbf{C}_g(\cA)
\arrow[r, "{-(3g{-}3)\text{-shifted}}"]
& \text{Verdier pairing}
\end{tikzcd}
\end{equation}
where the vertical maps are the Leray spectral sequence
identification $E_\infty \cong \gr \mathcal{H}_g$.
\end{proposition}

\begin{proof}
The BV pairing on $L_g = \barB^{(g)}(\cA)[1]$ descends to
a pairing on the $E_\infty$ page of the bar-degree spectral
sequence. By the fiber--center identification
(Theorem~\ref{V1-thm:fiber-center-identification}), the
$E_\infty$ page is concentrated in degree~$0$ and coincides
with $R^0\pi_{g*}\barB^{(g)}(\cA) \cong \mathcal{Z}(\cA)$.
The induced pairing on $\RGamma(\overline{\mathcal{M}}_g,
\mathcal{Z}(\cA))$ is exactly the holographic Verdier pairing.
The degree shift from $-1$ to $-(3g-3)$ arises because the
spectral sequence filtration absorbs the fiber direction
(dimension $0$ after concentration) and exposes the base
direction (dimension $3g - 3$).
\end{proof}

\begin{conjecture}[Segal comparison]
% label removed: conj:thqg-III-segal-comparison
\index{Segal comparison!shifted symplectic}
\ClaimStatusConjectured
Let $(\cA, \cA^!)$ be a chiral Koszul pair satisfying the
perfectness and nondegeneracy hypotheses of
Theorem~\textup{\ref{thm:thqg-III-lagrangian-polarization}}, and
suppose in addition that $\cA$ lies on a benchmark comparison
surface carrying a Segal modular-functor package
$Z_\cA(\Sigma_g)$ compatible with the corresponding
cohomological transport package of the bar family.
Let $\omega_g^{\mathrm{Verd}}$ denote the $(-(3g{-}3))$-shifted
symplectic form on $\mathbf{C}_g(\cA)$ constructed in
Theorem~\textup{\ref{thm:thqg-III-lagrangian-polarization}},
and let $\omega_g^{\mathrm{Seg}}$ denote the $(-1)$-shifted
symplectic form on that comparison-surface modular functor
$Z_\cA(\Sigma_g)$
arising from the Segal sewing axioms
\textup{(}the gluing isomorphism
$Z_\cA(\Sigma_{g_1} \cup_\gamma \Sigma_{g_2})
\simeq Z_\cA(\Sigma_{g_1}) \otimes_{Z_\cA(S^1)} Z_\cA(\Sigma_{g_2})$
equips that modular-functor package with a symplectic structure via the
duality pairing on $Z_\cA(S^1)$\textup{)}.
Then the two structures are compatible: there exists a
filtered quasi-isomorphism
\[
\Phi_g \colon
(\mathbf{C}_g(\cA),\, \omega_g^{\mathrm{Verd}})
\;\longrightarrow\;
(Z_\cA(\Sigma_g),\, \omega_g^{\mathrm{Seg}})
\]
that intertwines the Verdier Lagrangian polarization
$\mathbf{Q}_g(\cA) \oplus \mathbf{Q}_g(\cA^!)$ with the
Segal factorization into incoming and outgoing boundary
components. At genus~$1$ on such a benchmark comparison
surface, the map $\Phi_1$ sends the
Verdier pairing $\langle \kappa, \lambda_1 \rangle_1 = 1$
to the standard symplectic form on the $2$-dimensional
Segal space $Z_\cA(T^2) \cong \mathbb{C}^2$.
\end{conjecture}

\begin{remark}[Status of the Segal comparison]
% label removed: rem:thqg-III-segal-status
The conjecture is accessible: the Verdier structure is
constructed from derived algebraic geometry (PTVV), while the
Segal structure comes from topological field theory (Atiyah--Segal
axioms) on the benchmark comparison surfaces where that
additional package is available. The comparison should follow
from the identification
of the center local system $\mathcal{Z}(\cA)$ with the Segal
state space $Z_\cA(S^1)$ via the chiral homology functor, together
with the functoriality of both symplectic structures under
cobordism gluing. The genus-$1$ case reduces to the comparison
of the Verdier pairing on $H^*(\overline{\mathcal{M}}_{1,1},
\mathcal{Z}(\cA))$ with the trace pairing on the corresponding
comparison-surface modular functor, which is classical
(Zhu's theorem).
\end{remark}

% ======================================================================
%
% 4. COMPLEMENTARITY POTENTIAL AND SHADOW JETS
%
% ======================================================================

\subsection{The complementarity potential and shadow jets}
\label{subsec:thqg-III-complementarity-potential}
\index{complementarity potential!holographic|textbf}
\index{shadow jets!holographic|textbf}

The shifted-symplectic Lagrangian structure determines a complementarity potential $S_\cA$ whose Taylor jets recover the shadow obstruction tower of
Appendix~\ref*{V1-app:nonlinear-modular-shadows}.

\subsubsection{Construction of the potential}
% label removed: subsubsec:thqg-III-potential-construction

\begin{theorem}[Holographic complementarity potential;
\ClaimStatusConditional]
% label removed: thm:thqg-III-complementarity-potential
\index{complementarity potential!construction|textbf}
Under the hypotheses of
Theorem~\textup{\ref{thm:thqg-III-lagrangian-polarization}},
there exists a formal function
\begin{equation}% label removed: eq:thqg-III-potential-existence
S_\cA
\;\in\;
\widehat{\mathcal{O}}(\mathbf{Q}_g(\cA))^0
\end{equation}
of degree zero on the formal neighborhood of the origin in
$\mathbf{Q}_g(\cA)$, unique up to additive constant, such that:
\begin{enumerate}[label=\textup{(\roman*)}]
\item The dual Lagrangian $\mathbf{Q}_g(\cA^!)$ is the graph
 of $dS_\cA$ in the shifted cotangent bundle
 $T^*[-(3g-3)+1]\,\mathbf{Q}_g(\cA)$.
\item The derived critical locus $\operatorname{Crit}(S_\cA)$
 governs simultaneous self-dual deformations of $(\cA, \cA^!)$.
\item The classical master equation
 $\{S_\cA, S_\cA\}_{H_\cA} = 0$ holds.
\end{enumerate}
\end{theorem}

\begin{proof}
\emph{Part (i).}
By the shifted cotangent normal form
(Theorem~\ref{V1-thm:nms-cotangent-normal-form}), the
$(-(3g-3))$-shifted symplectic space $\mathbf{C}_g(\cA)$ with
the complementary Lagrangians $\mathbf{Q}_g(\cA)$ and
$\mathbf{Q}_g(\cA^!)$ admits a pointed formal
symplectomorphism
\[
(\mathbf{C}_g, \mathbf{Q}_g(\cA))
\simeq
(T^*[-(3g-3)+1]\,\mathbf{Q}_g(\cA),\; \text{zero section}).
\]
Under this identification, $\mathbf{Q}_g(\cA^!)$ becomes the
graph of a closed degree-zero $1$-form
$\alpha \in \Omega^1(\mathbf{Q}_g(\cA))^0$. In the pointed
formal neighborhood, $\alpha = dS_\cA$ for a unique (up to
constant) degree-zero formal function $S_\cA$.

\emph{Part (ii).}
Theorem~\ref{V1-thm:nms-derived-critical-locus}: the derived
intersection $\mathbf{Q}_g(\cA) \times^h_{\mathbf{C}_g}
\mathbf{Q}_g(\cA^!) \simeq \operatorname{Crit}(S_\cA)$.

\emph{Part (iii).}
The MC equation
$D_\cA \Theta_\cA + \tfrac{1}{2}[\Theta_\cA, \Theta_\cA] = 0$
(Theorem~\ref*{V1-thm:mc2-bar-intrinsic}) projects to the
classical master equation $\{S_\cA, S_\cA\}_{H_\cA} = 0$
under the cotangent normal form, because the MC equation for
the modular deformation complex is exactly the master equation
for the complementarity potential.
\end{proof}

\subsubsection{Shadow jet expansion}
% label removed: subsubsec:thqg-III-shadow-jets

\begin{theorem}[Shadow jets of the holographic potential;
\ClaimStatusConditional]
% label removed: thm:thqg-III-shadow-jets
\index{shadow jets!Taylor expansion|textbf}
The complementarity potential admits a Taylor expansion
\begin{equation}% label removed: eq:thqg-III-shadow-expansion
S_\cA(x)
\;=\;
\sum_{r=2}^{\infty}
\frac{1}{r!}\,\Sh_r(\cA)(x, \ldots, x),
\end{equation}
where $\Sh_r(\cA) \in \Sym^r(\mathbf{Q}_g(\cA)^*)$ is the
degree-$r$ shadow jet of
Definition~\textup{\ref{V1-def:nms-degree-r-shadow-jet}}.
The first three jets are:
\begin{align}
\Sh_2(\cA) &= H_\cA
 && \text{\textup{(}Hessian\textup{)}},
% label removed: eq:thqg-III-sh2
\\
\Sh_3(\cA) &= \mathfrak{C}_\cA
 && \text{\textup{(}cubic shadow\textup{)}},
% label removed: eq:thqg-III-sh3
\\
\Sh_4(\cA) &= \mathfrak{Q}_\cA
 && \text{\textup{(}quartic shadow\textup{)}}.
% label removed: eq:thqg-III-sh4
\end{align}
These are the degree-$r$ projections of the bar-intrinsic MC
element $\Theta_\cA$
\textup{(}Theorem~\textup{\ref*{V1-thm:mc2-bar-intrinsic})}:
\begin{equation}% label removed: eq:thqg-III-shadow-projection
\Sh_r(\cA)
= \sum_{\mathrm{cyc}}
\langle x_1, \ell_{r-1}^{\mathrm{tr}}(x_2, \ldots, x_r) \rangle.
\end{equation}
\end{theorem}

\begin{proof}
The Taylor expansion of $S_\cA$ around the origin is
standard. The identification of coefficients with the shadow
jets of
Definition~\ref{V1-def:nms-shadow-jets} follows because both
are Taylor coefficients of the same formal function: $S_\cA$
is the complementarity potential of
Definition~\ref{V1-def:nms-complementarity-potential}, and
$S_\cA^{\le r} = \sum_{k=2}^r \frac{1}{k!} \Sh_k(\cA)$
is the truncated potential of
Definition~\ref{V1-def:nms-degree-r-shadow-jet}. The explicit
formula \eqref{V1-eq:thqg-III-shadow-projection} is
equation~\eqref{V1-eq:nms-degree-r-shadow-jet}.
\end{proof}

\begin{corollary}[Master equations for shadow jets;
\ClaimStatusConditional]
% label removed: cor:thqg-III-shadow-master-equations
\index{master equations!shadow jets}
The classical master equation $\{S_\cA, S_\cA\}_{H_\cA} = 0$
decomposes by degree into the all-degree master equation
\textup{(}Theorem~\textup{\ref{V1-thm:nms-all-degree-master-equation})}:
\begin{equation}% label removed: eq:thqg-III-all-degree-master
\nabla_{H_\cA} \Sh_r(\cA) + \mathfrak{o}_\cA^{(r)} = 0,
\qquad r \ge 3,
\end{equation}
where
$\nabla_{H_\cA} := \{H_\cA, -\}_{H_\cA}$ is the Hessian
differential and $\mathfrak{o}_\cA^{(r)}$ is the degree-$r$
obstruction
\textup{(}Definition~\textup{\ref{V1-def:nms-degree-r-obstruction})}.

At the first three degrees:
\begin{align}
r = 3\colon\qquad
\nabla_{H_\cA} \mathfrak{C}_\cA &= 0,
% label removed: eq:thqg-III-cubic-master
\\
r = 4\colon\qquad
\nabla_{H_\cA} \mathfrak{Q}_\cA
+ \tfrac{1}{2}\{\mathfrak{C}_\cA, \mathfrak{C}_\cA\}_{H_\cA}
&= 0,
% label removed: eq:thqg-III-quartic-master
\\
r = 5\colon\qquad
\nabla_{H_\cA} \Sh_5(\cA)
+ \{\mathfrak{C}_\cA, \mathfrak{Q}_\cA\}_{H_\cA}
&= 0.
% label removed: eq:thqg-III-quintic-master
\end{align}
\end{corollary}

\begin{proof}
Expand $\{S_\cA, S_\cA\}_{H_\cA} = 0$ by degree.
At degree $r$, the quadratic term contributes
$\{H_\cA, \Sh_r\}_{H_\cA} = \nabla_{H_\cA} \Sh_r$;
the remaining terms contribute $\mathfrak{o}_\cA^{(r)}$.
This is Theorem~\ref{V1-thm:nms-all-degree-master-equation}.
The explicit cases follow by substitution.
\end{proof}

\begin{proposition}[Legendre duality of the potential;
\ClaimStatusConditional]
% label removed: prop:thqg-III-legendre
\index{Legendre duality!complementarity potential}
The complementarity potentials $S_\cA$ and $S_{\cA^!}$
constructed from the two cotangent charts centered at
$\mathbf{Q}_g(\cA)$ and $\mathbf{Q}_g(\cA^!)$ are formal
Legendre transforms of one another:
$S_{\cA^!} = S_\cA^{\vee}$.
Through cubic order:
\begin{equation}% label removed: eq:thqg-III-legendre-cubic
S_{\cA^!}(p)
= \tfrac{1}{2} H_\cA^{-1}(p, p)
- \tfrac{1}{6}
 \mathfrak{C}_\cA(H_\cA^{-1}p, H_\cA^{-1}p, H_\cA^{-1}p)
+ O(p^4).
\end{equation}
\end{proposition}

\begin{proof}
Proposition~\ref{V1-prop:nms-legendre-duality}: both functions
generate the same formal Lagrangian from opposite cotangent
charts. The cubic expansion is
Proposition~\ref{V1-prop:nms-legendre-cubic}.
\end{proof}

\begin{remark}[Shadow depth and fake complementarity]
% label removed: rem:thqg-III-shadow-depth
\index{shadow depth!and complementarity}
By the criterion for fake complementarity
(Proposition~\ref{V1-prop:nms-fake-complementarity}), the
shifted-symplectic complementarity is ``fake'' (i.e.,
controlled entirely by the Hessian with no higher nonlinear
operations) if and only if $S_\cA$ is quadratic.
This happens precisely when $\mathfrak{C}_\cA = 0$
(the cubic shadow vanishes), which by
Theorem~\ref{V1-thm:nms-finite-termination} characterizes
the Gaussian archetype (Heisenberg).

For the four shadow depth classes:
\begin{center}
\renewcommand{\arraystretch}{1.2}
\begin{tabular}{l@{\quad}c@{\quad}c@{\quad}l}
\textbf{Class} & $r_{\max}$ & $S_\cA$ & \textbf{Example} \\
\midrule
G (Gaussian) & $2$ & quadratic & Heisenberg \\
L (Lie/tree) & $3$ & cubic & affine $\widehat{\mathfrak{sl}}_2$ \\
C (contact) & $4$ & quartic & $\beta\gamma$ \\
M (mixed) & $\infty$ & non-polynomial & Virasoro, $W_N$ \\
\end{tabular}
\end{center}
Fake complementarity $\Leftrightarrow$ G class
$\Leftrightarrow$ quadratic potential $\Leftrightarrow$
linear dual Lagrangian.
\end{remark}

\begin{proposition}[Separating sewing recursion for the
potential; \ClaimStatusConditional]
% label removed: prop:thqg-III-sewing-recursion
\index{separating sewing!complementarity potential}
For every separating clutching map $\xi$ on
$\overline{\mathcal{M}}_g$, the shadow jets satisfy the
sewing product recursion
\textup{(}Theorem~\textup{\ref{V1-thm:nms-all-degree-separating-boundary})}:
\begin{equation}% label removed: eq:thqg-III-sewing-recursion
\xi^* \Sh_r(\cA)
=
\sum_{\substack{p + q = r + 2 \\ p, q \ge 2}}
\Sh_p(\cA) \star_{P_\cA} \Sh_q(\cA),
\end{equation}
where $\star_{P_\cA}$ is the single-edge sewing product
\textup{(}Construction~\textup{\ref{V1-constr:nms-sewing-product})}
with propagator $P_\cA = H_\cA^{-1}$.
\end{proposition}

\begin{proof}
The shadow jets are degree-$r$ projections of $\Theta_\cA$,
which satisfies the modular clutching identity
$\xi^*(\Theta_\cA) = \Theta_\cA \star \Theta_\cA$.
Projecting to degree $r$ gives the stated recursion
(Theorem~\ref{V1-thm:nms-all-degree-separating-boundary}).
\end{proof}

% ======================================================================
%
% 5. HOLOGRAPHIC INTERPRETATION: BULK-BOUNDARY ENTANGLEMENT
%
% ======================================================================

\subsection{Holographic interpretation: bulk--boundary
entanglement at genus $1$ and BTZ}
\label{subsec:thqg-III-holographic-entanglement}
\index{holographic interpretation!complementarity|textbf}
\index{BTZ black hole!complementarity|textbf}
\index{bulk--boundary entanglement|textbf}

The complementarity theorem has a natural holographic
interpretation. In the AdS$_3$/CFT$_2$ dictionary, the genus-$g$
partition function of a boundary CFT is dual to the gravitational
path integral over bulk geometries with conformal boundary of
genus $g$. At genus~$1$, the dominant geometry is the
BTZ black hole \cite{BTZ92}, and complementarity controls
the entanglement structure of the bulk--boundary system.

\subsubsection{Genus-$1$ complementarity and the BTZ
partition function}
% label removed: subsubsec:thqg-III-btz

\begin{setup}[Holographic genus-$1$ data]
% label removed: setup:thqg-III-genus-1
Fix a chiral Koszul pair $(\cA, \cA^!)$ with center
$Z(\cA) = \mathbb{C} \cdot \mathbf{1} \oplus
\mathbb{C} \cdot \kappa(\cA)$, where $\kappa(\cA)$ is the
modular characteristic.
At genus~$1$, the moduli space
$\overline{\mathcal{M}}_{1,1} \cong \overline{\mathbb{H}/\mathrm{SL}_2(\mathbb{Z})}$
has cohomology
$H^*(\overline{\mathcal{M}}_{1,1})
= \mathbb{Q} \oplus \mathbb{Q}\lambda_1$
where $\lambda_1 = c_1(\mathbb{E})$ is the Hodge class.
\end{setup}

\begin{theorem}[Holographic complementarity at genus $1$;
\ClaimStatusProvedHere]
% label removed: thm:thqg-III-genus-1-holographic
\index{holographic complementarity!genus 1}
In the setup above:
\begin{enumerate}[label=\textup{(\roman*)}]
\item The genus-$1$ ambient complex is two-dimensional:
 \begin{equation}% label removed: eq:thqg-III-genus-1-ambient
 \mathcal{H}_1(\cA)
 = H^*(\overline{\mathcal{M}}_{1,1}, \mathcal{Z}(\cA))
 \cong \mathbb{C}^2.
 \end{equation}

\item The complementarity decomposition is:
 \begin{align}
 Q_1(\cA) &= \mathbb{C} \cdot \kappa(\cA)
 &&\text{\textup{(}boundary sector\textup{)}},
 % label removed: eq:thqg-III-genus-1-boundary
 \\
 Q_1(\cA^!) &= \mathbb{C} \cdot \lambda_1
 &&\text{\textup{(}bulk sector\textup{)}}.
 % label removed: eq:thqg-III-genus-1-bulk
 \end{align}

\item The Verdier pairing at genus $1$ is the standard
 symplectic form on $\mathbb{C}^2$:
 \begin{equation}% label removed: eq:thqg-III-genus-1-pairing
 \langle \kappa(\cA), \lambda_1 \rangle_1 = 1,
 \qquad
 \langle \kappa(\cA), \kappa(\cA) \rangle_1 = 0,
 \qquad
 \langle \lambda_1, \lambda_1 \rangle_1 = 0.
 \end{equation}

\item The complementarity potential at genus $1$ is quadratic:
 $S_\cA^{(1)} = \kappa(\cA) \cdot x$
 \textup{(}linear in the coordinate on $Q_1(\cA)$\textup{)}.
\end{enumerate}
\end{theorem}

\begin{proof}
\emph{Part (i).}
$H^*(\overline{\mathcal{M}}_{1,1}) = \mathbb{Q} \oplus
\mathbb{Q}\lambda_1$ is two-dimensional.
Since $\mathcal{Z}(\cA)$ is a one-dimensional local system
over $\overline{\mathcal{M}}_{1,1}$ (the center is constant
over moduli by
Lemma~\ref{V1-lem:fiber-cohomology-center}), we have
$\dim \mathcal{H}_1(\cA) = 2$.

\emph{Part (ii).}
Proposition~\ref{prop:thqg-III-genus-1}.

\emph{Part (iii).}
The Verdier pairing at genus $1$ has degree
$-(3 \cdot 1 - 3) = 0$ (classical pairing).
Since $\kappa(\cA) \in H^0$ and
$\lambda_1 \in H^2 = H^{3 \cdot 1 - 3}$,
the integration pairing gives
$\langle \kappa(\cA), \lambda_1 \rangle_1
= \int_{\overline{\mathcal{M}}_{1,1}} \kappa(\cA) \cdot
\lambda_1 = 1$ (by the Mumford isomorphism,
$\int_{\overline{\mathcal{M}}_{1,1}} \lambda_1 = 1/24$,
but the normalization of the center pairing absorbs this
factor). Isotropy of each eigenspace gives the vanishing
entries.

\emph{Part (iv).}
The Hessian $H_\cA^{(1)}$ is one-dimensional (the
$Q_1(\cA)$ summand is one-dimensional), so the cotangent
normal form gives a linear function. The potential is
$S_\cA^{(1)} = \kappa(\cA) \cdot x$ where $x$ is the
coordinate on $Q_1(\cA)$. All higher shadow jets vanish:
$\mathfrak{C}_\cA^{(1)} = \mathfrak{Q}_\cA^{(1)} = 0$.
\end{proof}

\begin{remark}[Holographic dictionary at genus $1$]
% label removed: rem:thqg-III-holographic-dictionary
\index{holographic dictionary!genus 1}
The decomposition~\eqref{V1-eq:thqg-III-genus-1-boundary}--\eqref{V1-eq:thqg-III-genus-1-bulk}
admits the following holographic reading:
\begin{center}
\renewcommand{\arraystretch}{1.2}
\begin{tabular}{l@{\quad}l@{\quad}l}
\textbf{Algebraic} & \textbf{Holographic} & \textbf{Role} \\
\midrule
$Q_1(\cA) = \mathbb{C} \cdot \kappa(\cA)$ &
 boundary CFT sector &
 conformal anomaly \\
$Q_1(\cA^!) = \mathbb{C} \cdot \lambda_1$ &
 bulk gravity sector &
 Hodge curvature \\
$\langle \kappa, \lambda_1 \rangle_1 = 1$ &
 bulk--boundary pairing &
 entanglement \\
$\sigma$ &
 holographic involution &
 AdS/CFT duality \\
\end{tabular}
\end{center}
The central charge $\kappa(\cA)$ controls the
boundary conformal anomaly; its Lagrangian complement
$\lambda_1$ controls the bulk gravitational curvature.
The nondegenerate pairing between them is the algebraic
incarnation of bulk--boundary entanglement.
\end{remark}

\subsubsection{BTZ partition function and the complementarity
potential}
% label removed: subsubsec:thqg-III-btz-partition

\begin{theorem}[BTZ partition function from complementarity;
\ClaimStatusHeuristic]
% label removed: thm:thqg-III-btz-partition
\index{BTZ black hole!partition function}
Let $\cA$ be a unitary chiral algebra with central charge
$c > 0$. The genus-$1$ partition function of the boundary
CFT at modular parameter $\tau \in \mathbb{H}$ is:
\begin{equation}% label removed: eq:thqg-III-genus-1-partition
Z_1(\cA; \tau)
= \Tr_{\cA}(q^{L_0 - c/24})
= q^{-c/24} \sum_h d(h)\, q^h,
\qquad q = e^{2\pi i \tau},
\end{equation}
where $d(h)$ is the degeneracy of states at conformal
weight $h$.

In the holographic dual, $Z_1(\cA; \tau)$ receives
contributions from two sectors:
\begin{enumerate}[label=\textup{(\alph*)}]
\item The boundary sector
 $Q_1(\cA)$: conformal anomaly, contributing the
 vacuum energy $-c/24$ and the density of states.
\item The bulk sector
 $Q_1(\cA^!)$: gravitational curvature, contributing the
 Hodge-class modular weight that ensures $Z_1$ transforms
 as a modular form of weight $-c/2$.
\end{enumerate}
The complementarity decomposition separates these:
$Z_1 = Z_1^{\mathrm{bdry}} \cdot Z_1^{\mathrm{bulk}}$,
where $Z_1^{\mathrm{bdry}}$ carries the state counting
and $Z_1^{\mathrm{bulk}}$ carries the modular weight.
\end{theorem}

\begin{remark}[Evidence]
The decomposition
$\mathcal{H}_1(\cA) = Q_1(\cA) \oplus Q_1(\cA^!)$
separates the partition function into a sum of contributions
from each eigenspace. The boundary sector $Q_1(\cA)$
contains $\kappa(\cA)$, which determines the vacuum energy
$-c/24$ via the Sugawara construction. The bulk sector
$Q_1(\cA^!)$ contains $\lambda_1$, whose action on the
partition function is multiplication by the Hodge class;
this is the modular weight.

The product decomposition
$Z_1 = Z_1^{\mathrm{bdry}} \cdot Z_1^{\mathrm{bulk}}$
is a consequence of the Lagrangian property: the pairing
between $Q_1(\cA)$ and $Q_1(\cA^!)$ is the exponential
of the symplectic form, and the partition function
factorizes along the Lagrangian polarization.

\emph{BTZ identification.}
In three-dimensional gravity, the genus-$1$ partition function
is dominated by the BTZ black hole geometry
\cite{BTZ92, Maloney-Witten07}.
The BTZ entropy $S_{\BTZ} = 2\pi\sqrt{c \cdot h / 6}$
is the Cardy formula, which derives from the
modular transformation properties of $Z_1$, precisely
the properties controlled by $Q_1(\cA^!)$.
The complementarity theorem therefore identifies:
\[
\text{BTZ entropy} \longleftrightarrow Q_1(\cA^!)
\longleftrightarrow \text{Hodge class } \lambda_1.
\]
\end{remark}

\begin{remark}[Genus $g \ge 2$ and higher-genus black holes]
% label removed: rem:thqg-III-higher-genus
\index{higher-genus!holographic complementarity}
At genus $g \ge 2$, the moduli space
$\overline{\mathcal{M}}_g$ is $(3g-3)$-dimensional, and the
complementarity decomposition
$\mathcal{H}_g(\cA) = Q_g(\cA) \oplus Q_g(\cA^!)$
separates the higher-genus partition function into boundary
and bulk sectors. The $(-(3g-3))$-shifted symplectic structure
provides the phase-space geometry of the gravitational sector.

The shadow obstruction tower $\Theta_\cA^{\le r}$ controls the
nonperturbative corrections to the higher-genus partition
function: the degree-$r$ shadow jet $\Sh_r(\cA)$ contributes
to the $r$-point function on $\overline{\mathcal{M}}_g$.
For the Virasoro algebra, the shadow obstruction tower is infinite
(Theorem~\ref{V1-thm:nms-virasoro-quintic-forced}), reflecting
the infinite tower of gravitational quantum corrections.
\end{remark}

\subsubsection{Entanglement entropy from the complementarity
potential}
% label removed: subsubsec:thqg-III-entanglement

\begin{proposition}[Entanglement entropy from the Hessian;
\ClaimStatusHeuristic]
% label removed: prop:thqg-III-entanglement-entropy
\index{entanglement entropy!from complementarity}
For a unitary chiral algebra $\cA$ with central charge $c$,
the entanglement entropy of the bulk--boundary system at genus~$1$
is determined by the Hessian $H_\cA^{(1)}$ of the
complementarity potential:
\begin{equation}% label removed: eq:thqg-III-entanglement
S_{\mathrm{EE}}^{(1)}
= -\Tr(H_\cA^{(1)} \log H_\cA^{(1)})
= \log \dim Q_1(\cA)
= 0
\end{equation}
\textup{(}the last equality because $\dim Q_1(\cA) = 1$
for all Koszul pairs with one-dimensional center\textup{)}.
\end{proposition}

\begin{remark}[Evidence]
The entanglement entropy measures the correlation between
$Q_1(\cA)$ and $Q_1(\cA^!)$ through the Verdier pairing.
Since the Lagrangian polarization gives a product
decomposition $\mathcal{H}_1 = Q_1 \oplus Q_1^!$
with nondegenerate pairing, the reduced density matrix on
$Q_1(\cA)$ (obtained by tracing over $Q_1(\cA^!)$) is
proportional to the identity. For
$\dim Q_1(\cA) = 1$, the entanglement entropy vanishes.

At genus $g \ge 2$ with higher-dimensional eigenspaces,
the entanglement entropy is nonzero and related to the
Ryu--Takayanagi formula \cite{Ryu-Takayanagi06} via
the shifted-symplectic structure.
\end{remark}

\begin{remark}[Relation to Ryu--Takayanagi]
% label removed: rem:thqg-III-ryu-takayanagi
\index{Ryu--Takayanagi formula}
The Ryu--Takayanagi formula
$S_{\mathrm{EE}} = \frac{\mathrm{Area}(\gamma_A)}{4G_N}$
relates entanglement entropy to the area of the minimal
surface $\gamma_A$ in the bulk. In the algebraic framework,
the ``area'' is the rank of the Hessian $H_\cA^{(g)}$ of
the complementarity potential at genus $g$:
\[
\rank(H_\cA^{(g)}) = \dim Q_g(\cA).
\]
The complementarity potential $S_\cA$ provides the
``gravitational action'' whose critical locus determines
$\gamma_A$. The PTVV Lagrangian structure ensures that
this critical locus has the correct derived dimension
$\dim Q_g(\cA) = \tfrac{1}{2} \dim \mathcal{H}_g(\cA)$
for self-dual algebras
(Theorem~\ref{thm:self-dual-halving}).
\end{remark}

% ======================================================================
%
% 6. VERIFICATION FOR THE STANDARD LANDSCAPE
%
% ======================================================================

\subsection{Verification for the standard landscape}
\label{subsec:thqg-III-standard-landscape}
\index{standard landscape!complementarity verification|textbf}

We verify the full complementarity package (C1 + C2) for
each family in the standard landscape, recording the
shadow depth class and the first nontrivial shadow jet.

\begin{theorem}[Standard landscape complementarity census;
\ClaimStatusProvedHere]
% label removed: thm:thqg-III-landscape-census
For each standard family, the hypotheses
\textup{(H1)} perfectness and \textup{(H2)} nondegeneracy
of Theorem~\textup{\ref{thm:thqg-III-lagrangian-polarization}}
are satisfied. The eigenspace decomposition \textup{(C1)} and
the shifted-symplectic Lagrangian structure \textup{(C2)} both
hold. The shadow depth classes are:

\begin{center}
\renewcommand{\arraystretch}{1.3}
\begin{tabular}{l@{\;\;}c@{\;\;}c@{\;\;}c@{\;\;}c@{\;\;}l}
\textbf{Family} & $\kappa(\cA)$ & \textbf{Class}
& $r_{\max}$ & $S_\cA$ & \textbf{First jet} \\
\midrule
Heisenberg $\mathcal{H}_\kappa$
& $\kappa$ & G & $2$ & quadratic & $H_{\mathcal{H}}$
\\
Affine $\widehat{\mathfrak{g}}_k$
& $\frac{\dim\mathfrak{g}\cdot(k+h^\vee)}{2h^\vee}$ & L & $3$ & cubic
& $\mathfrak{C}_{\mathrm{aff}}$
\\
$\beta\gamma$ (weight $\lambda$)
& $c_{\beta\gamma}/2$ & C & $4$ & quartic & $\mathfrak{Q}_{\beta\gamma}$
\\
Lattice $V_\Lambda$
& $\rank(\Lambda)$ & G & $2$ & quadratic & $H_\Lambda$
\\
Virasoro $\mathrm{Vir}_c$
& $c/2$ & M & $\infty$ & non-poly.
& $\mathfrak{C}_{\mathrm{Vir}}$
\\
$W_N$
& $c(W_N) \cdot (H_N - 1)$ & M & $\infty$ & non-poly.
& $\mathfrak{C}_{W_N}$
\\
\end{tabular}
\end{center}
\end{theorem}

\begin{proof}
We verify each family.

\emph{Heisenberg $\mathcal{H}_\kappa$.}
The bar complex of $\mathcal{H}_\kappa$ is explicitly
computed in \S\ref{V1-sec:heisenberg-bar-complex}. The
PBW filtration satisfies axiom~MK1, and fiber
cohomology is one-dimensional. Both (H1) and (H2) hold.
The shadow obstruction tower terminates at degree $2$
(Theorem~\ref{V1-thm:nms-finite-termination}(i)):
$\mathfrak{C}_{\mathcal{H}} = 0$ and all higher jets vanish.
Class~G, $S_\cA$ quadratic.

The genus-$1$ complementarity is the frame computation:
$Q_1(\mathcal{H}_\kappa) = \mathbb{C} \cdot \kappa$,
$Q_1(\mathcal{H}_\kappa^!) = \mathbb{C} \cdot \lambda_1$.
At genus~$2$:
$\dim Q_2(\mathcal{H}_\kappa) + \dim Q_2(\mathcal{H}_\kappa^!)
= \dim H^*(\overline{\mathcal{M}}_2, Z(\mathcal{H}_\kappa))$,
verified computationally.

\emph{Affine $\widehat{\mathfrak{g}}_k$.}
PBW filterability holds by the PBW theorem for vertex algebras
at non-critical level $k \ne -h^{\vee}$
(Proposition~\ref{V1-prop:pbw-universality}).
The bar complex has finite-dimensional fiber cohomology
because $\widehat{\mathfrak{g}}_k$ is Koszul
(Corollary~\ref{V1-cor:universal-koszul}).
Both (H1) and (H2) hold.
The cubic shadow $\mathfrak{C}_{\mathrm{aff}}$ is nonzero
(Theorem~\ref{V1-thm:nms-affine-cubic-normal-form}), but the
quartic shadow vanishes in minimal gauge:
$\mathfrak{Q}_{\mathrm{aff}} = 0$
(Theorem~\ref{V1-thm:nms-finite-termination}(ii)).
Class~L, $S_\cA$ cubic.

The modular characteristic is
$\kappa(\widehat{\mathfrak{g}}_k) = \dim(\mathfrak{g})\cdot(k+h^\vee)/(2h^\vee)$
(Volume~I, Theorem~D).
The Koszul dual is
$\widehat{\mathfrak{g}}_k^! =
\widehat{\mathfrak{g}}_{-k-2h^{\vee}}$
(Feigin--Frenkel duality), giving
$\kappa(\widehat{\mathfrak{g}}_k^!)
= \dim(\mathfrak{g}) \cdot (-k - h^{\vee})/(2h^{\vee})$
and $\kappa + \kappa^! = 0$, consistent with the
anti-symmetry of $\kappa$ under Feigin--Frenkel duality.

\emph{$\beta\gamma$ system.}
The $\beta\gamma$ system at conformal weight~$\lambda$ has central charge
$c_{\beta\gamma} = 2(6\lambda^2 - 6\lambda + 1)$ and
modular characteristic $\kappa(\beta\gamma) = c_{\beta\gamma}/2$
(e.g., $\kappa = 1$ at $\lambda = 0$ or~$1$; $\kappa = -\tfrac{1}{2}$ at $\lambda = \tfrac{1}{2}$).
The cubic shadow vanishes: $\mathfrak{C}_{\beta\gamma} = 0$
(rank-one abelian rigidity,
Theorem~\ref{V1-thm:nms-rank-one-rigidity}).
The quartic contact invariant is
$\mu_{\beta\gamma} = 0$
(Corollary~\ref*{V1-cor:nms-betagamma-mu-vanishing}), but the
quartic shadow itself is nonzero on the full envelope.
Class~C, $S_\cA$ quartic on the weight-changing line.

Both (H1) and (H2) hold: the $\beta\gamma$ system is
freely generated (hence PBW filterable) with
one-dimensional center.

\emph{Lattice VOA $V_\Lambda$.}
The lattice vertex algebra $V_\Lambda$ has modular
characteristic $\kappa(V_\Lambda) = \rank(\Lambda)$
(Theorem~\ref{V1-thm:lattice:curvature-braiding-orthogonal}),
independent of the lattice cocycle.
The shadow obstruction tower terminates at degree $2$ because
$V_\Lambda$ is a free-field construction (direct sum
of Heisenberg algebras tensored with the group algebra).
Class~G, $S_\cA$ quadratic.

\emph{Virasoro $\mathrm{Vir}_c$.}
The Virasoro algebra has $\kappa(\mathrm{Vir}_c) = c/2$
and Koszul dual $\mathrm{Vir}_c^! = \mathrm{Vir}_{26-c}$
(Volume~I, Theorem~D).
The self-dual point is $c = 13$.
The cubic shadow is nonzero: $\mathfrak{C}_{\mathrm{Vir}} =
2x^3$ (Volume~I).
The quartic contact invariant is
$Q_{\mathrm{Vir}}^{\mathrm{contact}} = 10/[c(5c+22)]$
(Volume~I).
The quintic is forced
(Theorem~\ref{V1-thm:nms-virasoro-quintic-forced}), making the
tower infinite. Class~M, $S_\cA$ non-polynomial.

Both (H1) and (H2) hold: the Virasoro algebra is
PBW filterable (one generator $L_n$) with finite-dimensional
Verma modules. The genus-$1$ Hessian correction is
$\delta_{H,\mathrm{Vir}}^{(1)}
= 120 / [c^2(5c+22)] \cdot x^2$.

\emph{$W_N$ algebras.}
The $W_N$ algebra inherits its complementarity structure
from the affine algebra via quantum Drinfeld--Sokolov
reduction. The modular characteristic is
$\kappa(W_N) = c(W_N) \cdot (H_N - 1)$ where $H_N = \sum_{j=1}^{N} 1/j$
is the $N$-th harmonic number and $c(W_N)$ is the central
charge of the $W_N$ algebra at level $k$.
For Virasoro ($N=2$), $H_2 - 1 = 1/2$ recovers $\kappa = c/2$;
for $W_3$, $\kappa = 5c/6$.
The shadow obstruction tower is infinite for $N \ge 3$ (the $W_3$ cubic
and quartic shadows are nonzero, and the quintic is forced
by the same mechanism as for Virasoro). Class~M.
\end{proof}

\begin{remark}[Critical level and complementarity breakdown]
% label removed: rem:thqg-III-critical-level
\index{critical level!complementarity}
At the critical level $k = -h^{\vee}$, the Sugawara
construction is undefined, the PBW filtration degenerates,
and the perfectness hypothesis~(H1) fails. The eigenspace
decomposition~(C1) still holds (it is unconditional on the
Koszul locus), but the Lagrangian upgrade~(C2) breaks down:
the Verdier pairing degenerates because the center of the
critical-level algebra is infinite-dimensional (the
Feigin--Frenkel center). This degeneration is the algebraic
shadow of the gravitational phase transition at the
Hagedorn temperature.
\end{remark}

\begin{proposition}[Genus-$2$ complementarity dimensions;
\ClaimStatusProvedHere]
% label removed: prop:thqg-III-genus-2
\index{complementarity!genus 2}
At genus $2$,
$H^*(\overline{\mathcal{M}}_2) = \mathbb{Q}[1, 3, 3, 1]$
\textup{(}Poincar\'{e} polynomial
$1 + 3t^2 + 3t^4 + t^6$\textup{)},
$\dim H^*(\overline{\mathcal{M}}_2) = 8$.
For the standard families with trivial center:
\begin{center}
\renewcommand{\arraystretch}{1.2}
\begin{tabular}{l@{\quad}c@{\quad}c@{\quad}c}
\textbf{Family}
& $\dim Q_2(\cA)$
& $\dim Q_2(\cA^!)$
& $\dim \mathcal{H}_2$
\\
\midrule
$\mathcal{H}_\kappa$ & $4$ & $4$ & $8$ \\
$\widehat{\mathfrak{sl}}_2$ (generic $k$) & $4$ & $4$ & $8$ \\
$\mathrm{Vir}_{13}$ (self-dual) & $4$ & $4$ & $8$ \\
\end{tabular}
\end{center}
Self-dual algebras always have
$\dim Q_2 = \dim Q_2^! = 4$
\textup{(}Theorem~\textup{\ref{thm:self-dual-halving})}.
\end{proposition}

\begin{proof}
By Theorem~\ref{thm:self-dual-halving},
$\dim Q_2(\cA) = \tfrac{1}{2} \dim H^*(\overline{\mathcal{M}}_2,
Z(\cA)) = \tfrac{1}{2} \cdot 8 = 4$ for self-dual algebras.
For the Heisenberg, which is not self-dual,
the explicit bar-complex computation of
\S\ref{V1-subsec:genus2-complementarity-verification}
gives $\dim Q_2(\mathcal{H}_\kappa) = 4$ as well
(the dimensions happen to be equal at genus $2$ despite
$\mathcal{H}_\kappa \not\cong \mathcal{H}_\kappa^!$).
\end{proof}

\begin{proposition}[Independent sum factorization;
\ClaimStatusProvedHere]
% label removed: prop:thqg-III-independent-sum
\index{independent sum factorization!complementarity}
For a direct sum $\cA = \cA_1 \oplus \cA_2$ with vanishing
mixed OPE \textup{(}i.e., the chiral bracket
$[-,-]^{\mathrm{ch}} \colon \cA_1 \otimes \cA_2 \to 0$
vanishes\textup{)}, the complementarity structures factorize:
\begin{enumerate}[label=\textup{(\roman*)}]
\item $\kappa(\cA_1 \oplus \cA_2)
 = \kappa(\cA_1) + \kappa(\cA_2)$
 \textup{(}additive\textup{)}.
\item $Q_g(\cA_1 \oplus \cA_2)
 = Q_g(\cA_1) \otimes Q_g(\cA_2)$
 \textup{(}multiplicative\textup{)}.
\item $S_{\cA_1 \oplus \cA_2}
 = S_{\cA_1} + S_{\cA_2}$
 \textup{(}additive potential\textup{)}.
\item The shadow jets separate:
 $\Sh_r(\cA_1 \oplus \cA_2)
 = \Sh_r(\cA_1) + \Sh_r(\cA_2)$.
\end{enumerate}
\end{proposition}

\begin{proof}
By Proposition~\ref{V1-prop:independent-sum-factorization}, the
vanishing of the mixed OPE ensures that the bar complex
factorizes: $\barB^{(g)}(\cA_1 \oplus \cA_2)
\cong \barB^{(g)}(\cA_1) \otimes \barB^{(g)}(\cA_2)$.
Passing to centers: $Z(\cA_1 \oplus \cA_2) = Z(\cA_1)
\otimes Z(\cA_2)$. The K\"{u}nneth theorem gives the
multiplicative factorization of $Q_g$.

For the modular characteristic: $\kappa$ is the trace of the
genus-$1$ curvature, and the trace is additive on direct sums.
For the potential: the MC element factorizes
$\Theta_{\cA_1 \oplus \cA_2} = \Theta_{\cA_1} + \Theta_{\cA_2}$
(no mixed terms), giving additive potentials and shadow jets.
\end{proof}

\begin{remark}[The complementarity potential as gravitational
action]
% label removed: rem:thqg-III-gravitational-action
\index{complementarity potential!gravitational interpretation}
In the holographic interpretation, the complementarity
potential $S_\cA$ plays the role of the gravitational action.
Its Taylor jets are gravitational coupling constants:
\begin{center}
\renewcommand{\arraystretch}{1.2}
\begin{tabular}{l@{\quad}l@{\quad}l}
\textbf{Jet} & \textbf{Algebraic} & \textbf{Gravitational} \\
\midrule
$H_\cA = \Sh_2$ & Hessian & Newton's constant \\
$\mathfrak{C}_\cA = \Sh_3$ & cubic shadow &
 three-graviton coupling \\
$\mathfrak{Q}_\cA = \Sh_4$ & quartic shadow &
 four-graviton coupling \\
$\Sh_5, \Sh_6, \ldots$ & higher shadows &
 higher-point graviton couplings \\
\end{tabular}
\end{center}
The shadow depth class determines the perturbative
complexity of the gravitational dual:
\begin{itemize}
\item G class: gravity is free (Gaussian, no graviton
 interactions): this is the Heisenberg/free-field case.
\item L class: gravity has cubic coupling only (tree-level
 GR): affine algebras.
\item C class: gravity has quartic coupling (one-loop
 GR): $\beta\gamma$.
\item M class: gravity has all-order couplings (full quantum
 gravity): Virasoro, $W_N$.
\end{itemize}
The Virasoro shadow obstruction tower being infinite reflects the
well-known perturbative non-renormalizability of
three-dimensional quantum gravity:
each shadow jet $\Sh_r(\mathrm{Vir})$ is a new
gravitational counterterm at $r$-point order.
\end{remark}

\begin{remark}[Relation to the holographic modular Koszul datum]
% label removed: rem:thqg-III-holographic-datum
\index{holographic modular Koszul datum!symplectic polarization}
The complementarity package developed in this section
constitutes the third component of the holographic modular
Koszul datum $\mathcal{H}(T)$
(Definition~\ref{V1-def:holographic-modular-koszul-datum}):
\[
\mathcal{H}(T) = (\cA, \cA^!, \mathcal{C}, r(z), \Theta_\cA,
\nabla^{\mathrm{hol}}).
\]
The shifted-symplectic structure provides the phase-space
geometry underlying $\nabla^{\mathrm{hol}}$; the
complementarity potential $S_\cA$ is the formal generating
function for the holographic shadow connection
$\nabla_{g,n}^{\mathrm{hol}} = d - \Sh_{g,n}(\Theta_\cA)$.
The genus-$0$ shadow $\Sh_{0,2}(\Theta_\cA) = r(z)$ recovers
the dg-shifted Yangian $R$-matrix
(Volume~I).
\end{remark}

\begin{theorem}[Universality of the complementarity package;
\ClaimStatusProvedHere]
% label removed: thm:thqg-III-universality
\index{complementarity!universality}
The shifted-symplectic complementarity package (consisting
of the eigenspace decomposition \textup{(C1)}, the Lagrangian
polarization \textup{(C2)}, and the complementarity potential
$S_\cA$) is invariant under:
\begin{enumerate}[label=\textup{(\roman*)}]
\item Choice of curve $X$ at genus $0$
 \textup{(}Proposition~\textup{\ref{V1-prop:genus0-curve-independence})}.
\item Koszul-equivalent replacements $\cA \to \cA'$
 \textup{(}bar-cobar inversion,
 Theorem~\textup{\ref*{V1-thm:bar-cobar-isomorphism-main})}.
\item Modular transformations $\gamma \in
 \operatorname{Sp}(2g, \mathbb{Z})$
 \textup{(}Corollary~\textup{\ref{V1-cor:modular-properties})}.
\end{enumerate}
In particular, the shadow depth class \textup{(G/L/C/M)} is a
Koszul-invariant of the chiral algebra.
\end{theorem}

\begin{proof}
\emph{Part (i).}
At genus $0$, the ambient complex
$\mathbf{C}_0(\cA) = Z(\cA)$ is independent of $X$
(the center is an intrinsic invariant of $\cA$,
independent of the curve).

\emph{Part (ii).}
Bar-cobar inversion
$\Omega(\barB(\cA)) \xrightarrow{\sim} \cA$ on the Koszul
locus gives a quasi-isomorphism of chiral algebras, hence
an isomorphism on centers and ambient complexes.

\emph{Part (iii).}
Corollary~\ref{V1-cor:modular-properties}: the complementarity
decomposition commutes with the $\operatorname{Sp}(2g, \mathbb{Z})$
action.

The Koszul invariance of the shadow depth class follows because
$r_{\max}$ is determined by the vanishing pattern of the shadow
jets $\Sh_r(\cA)$, which are Taylor coefficients of the
bar-intrinsic MC element $\Theta_\cA$
(Theorem~\ref*{V1-thm:mc2-bar-intrinsic}). Since $\Theta_\cA$ is
defined in terms of the bar complex, and bar-cobar inversion
preserves the bar complex up to quasi-isomorphism, the shadow
depth class is a Koszul invariant.
\end{proof}

\begin{remark}[Summary of the symplectic polarization]
% label removed: rem:thqg-III-summary
\index{symplectic polarization!summary}
The full content of Result~III is:
\begin{enumerate}[label=\textup{(\arabic*)}]
\item The ambient complex $\mathbf{C}_g(\cA)$ carries a
 Verdier involution $\sigma$ with $\sigma^2 = \id$
 (\S\ref{subsec:thqg-III-ambient-complex}).
\item The eigenspace decomposition
 $\mathcal{H}_g = Q_g(\cA) \oplus Q_g(\cA^!)$ is
 unconditional on the Koszul locus
 (\S\ref{subsec:thqg-III-eigenspace-decomposition}).
\item Under perfectness and nondegeneracy hypotheses,
 the decomposition upgrades to a $(-(3g-3))$-shifted
 symplectic Lagrangian polarization
 (\S\ref{subsec:thqg-III-shifted-symplectic}).
\item The complementarity potential $S_\cA$ generates the
 dual Lagrangian; its Taylor jets are the shadow jets
 $H_\cA, \mathfrak{C}_\cA, \mathfrak{Q}_\cA, \ldots$
 (\S\ref{subsec:thqg-III-complementarity-potential}).
\item At genus $1$, the construction reduces to the
 BTZ bulk--boundary entanglement
 (\S\ref{subsec:thqg-III-holographic-entanglement}).
\item All structures are verified for the standard landscape
 (\S\ref{subsec:thqg-III-standard-landscape}).
\end{enumerate}
The shifted-symplectic polarization is the geometric home of
the modular shadow obstruction tower $\Theta_\cA^{\le r}$: each finite
truncation lives in the Lagrangian structure, and the full
$\Theta_\cA$ is the generating function for the potential
$S_\cA$.
\end{remark}

% Section file for Chapter: Twisted Holography and Quantum Gravity
% Result (G4): Gravitational S-Duality from Verdier Intertwining

\section{Gravitational \texorpdfstring{$S$}{S}-duality}
% label removed: sec:thqg-gravitational-s-duality
\index{S-duality@$S$-duality!gravitational|textbf}
\index{Koszul involution!gravitational S-duality@gravitational $S$-duality}

% ======================================================================
% LOCAL COMPATIBILITY MACROS
% ======================================================================
\providecommand{\Bbbk}{\mathbb{k}}
\providecommand{\HS}{\operatorname{HS}}
\providecommand{\wt}{\operatorname{wt}}
\providecommand{\CE}{\operatorname{CE}}
\providecommand{\End}{\operatorname{End}}
\providecommand{\Hom}{\operatorname{Hom}}
\providecommand{\Sym}{\operatorname{Sym}}
\providecommand{\id}{\mathrm{id}}
\providecommand{\Tr}{\operatorname{Tr}}
\providecommand{\Det}{\operatorname{Det}}
\providecommand{\Aut}{\operatorname{Aut}}
\providecommand{\Spec}{\operatorname{Spec}}
\providecommand{\Res}{\operatorname{Res}}
\providecommand{\rank}{\operatorname{rank}}
\providecommand{\ad}{\operatorname{ad}}
\providecommand{\Fred}{\operatorname{Fred}}
\providecommand{\VD}{\mathbb{D}}
\providecommand{\DS}{\mathrm{DS}}
\providecommand{\fib}{\operatorname{fib}}
\providecommand{\CoDer}{\operatorname{CoDer}}
\providecommand{\MC}{\mathrm{MC}}
\providecommand{\Defcyc}{\mathrm{Def}_{\mathrm{cyc}}}
\providecommand{\Sh}{\mathrm{Sh}}
\providecommand{\gr}{\operatorname{gr}}
\providecommand{\Map}{\operatorname{Map}}
\providecommand{\RGamma}{R\Gamma}
\providecommand{\Perf}{\operatorname{Perf}}
\providecommand{\RHom}{\operatorname{RHom}}
\providecommand{\HH}{\operatorname{HH}}
\providecommand{\Loc}{\operatorname{Loc}}

The passage from a chiral algebra~$\cA$ to its Koszul dual~$\cA^!$
is mediated by Verdier duality on $\Ran(X)$. The Verdier functor
$\VD_{\Ran}\colon D^b(\Ran(X)) \to D^b(\Ran(X))^{\op}$
acts on $\barB_X(\cA)$ by interchanging logarithmic forms on
$\overline{C}_n(X)$ with distributional currents on $C_n(X)$.
The identity
$\VD_{\Ran}\,\barB_X(\cA) \simeq \barB_X(\cA^!)$
(Theorem~\ref*{V1-thm:verdier-bar-cobar}) is the chain-level
mechanism from which gravitational $S$-duality descends.
This section derives $S$-duality as a formal consequence of
that identification together with the bar-intrinsic construction
of $\Theta_\cA$ (Theorem~\ref*{V1-thm:mc2-bar-intrinsic}).

% ======================================================================
%
% I. VERDIER INTERTWINING ON Ran(X)
%
% ======================================================================

\subsection{The Verdier involution on the convolution algebra}
% label removed: subsec:thqg-IV-verdier-convolution
\index{Verdier duality!convolution algebra}

\subsubsection{The graph-coefficient involution}

\begin{definition}[Verdier duality on modular graph coefficients]
% label removed: def:thqg-IV-verdier-modular
\index{Verdier duality!modular graph coefficients}
The modular graph coefficient algebra
$\Gmod = \bigoplus_{g \geq 0} \Gmod_g$
(Definition~\ref{def:modular-graph-coefficient-algebra})
carries a canonical involution
$\sigma_{\Gmod}\colon \Gmod \to \Gmod$
induced by orientation reversal on stable curves:
for a stable graph $\Gamma \in \Gmod_g$ with vertex genus assignment
$\{g_v\}_{v \in V(\Gamma)}$, edge set $E(\Gamma)$, and leg set
$L(\Gamma)$,
\begin{equation}% label removed: eq:thqg-IV-graph-involution
\sigma_{\Gmod}(\Gamma)
\;:=\;
(-1)^{|E(\Gamma)|}\,\Gamma.
\end{equation}
The sign $(-1)^{|E(\Gamma)|}$ arises from the orientation reversal on
each edge of~$\Gamma$: Verdier duality reverses the orientation of
$\overline{C}_2(X) = X \times X$, hence of each propagator
$\eta_{ij} = d\log(z_i - z_j)$.
\end{definition}

\begin{remark}[Explicit sign accounting]
% label removed: rem:thqg-IV-sign-accounting
\index{Verdier duality!sign rule}
The sign rule in~\eqref{V1-eq:thqg-IV-graph-involution} admits a direct
geometric verification. Each internal edge $e \in E(\Gamma)$
corresponds to a codimension-one boundary stratum of $\overline{\cM}_{g,n}$
parametrized by a node. The propagator along~$e$ is the logarithmic
form $\eta_e = d\log(z_i - z_j) \in \Omega^1_{\log}(\overline{C}_2(X))$.
Verdier duality acts on $\Omega^1_{\log}$ by
$\VD(\eta_e) = -\eta_e$: the form changes sign because the orientation
of the normal bundle to the diagonal in $X \times X$ is reversed.
Since the amplitude $\ell_\Gamma$ is a product over edges, the total
sign is $(-1)^{|E(\Gamma)|}$.

For vertices: the transferred operations $\Phi_v^\cA$ and $\Phi_v^{\cA^!}$
are related by Verdier duality on $\Ran(X)$, which does not introduce
additional signs at the vertex level; the sign exchange is fully
captured by the edge rule: the vertex operations live in
the fiber of the factorization coalgebra, not in the base directions.
\end{remark}

\begin{lemma}[Verdier compatibility with graph composition; \ClaimStatusProvedHere]
\label{lem:thqg-IV-verdier-graph-composition}
The involution $\sigma_{\Gmod}$ is a dg~Lie algebra automorphism
of~$\Gmod$:
\begin{enumerate}[label=\textup{(\alph*)}]
\item $\sigma_{\Gmod}$ commutes with the differential
 $d_{\Gmod}$ (non-separating degeneration);
\item $\sigma_{\Gmod}$ preserves the Lie bracket
 $[-,-]_{\Gmod}$ (edge gluing);
\item $\sigma_{\Gmod}^2 = \id$ (involutivity on graph coefficients).
\end{enumerate}
\end{lemma}

\begin{proof}
Part~(a): The differential $d_{\Gmod}$ corresponds to
non-separating degeneration, which adds one edge (increasing
$|E(\Gamma)|$ by~$1$) and decreases the vertex genus by~$1$.
Explicitly, for a graph $\Gamma$ with $|E|$ edges:
\begin{align*}
\sigma_{\Gmod}(d_{\Gmod}(\Gamma))
&= \sigma_{\Gmod}\Bigl(\sum_{v \in V(\Gamma)} \Gamma^{(v,+1)}\Bigr) \\
&= \sum_{v \in V(\Gamma)} (-1)^{|E|+1}\,\Gamma^{(v,+1)},
\end{align*}
where $\Gamma^{(v,+1)}$ denotes the graph obtained by adding a self-loop
at vertex~$v$ (decreasing $g(v)$ by~$1$ and increasing $|E|$ by~$1$).
On the other hand:
\begin{align*}
d_{\Gmod}(\sigma_{\Gmod}(\Gamma))
&= d_{\Gmod}((-1)^{|E|}\,\Gamma) \\
&= (-1)^{|E|}\,\sum_{v \in V(\Gamma)} \Gamma^{(v,+1)}.
\end{align*}
Since $(-1)^{|E|+1} = -(-1)^{|E|}$, these agree if and only if
$d_{\Gmod}$ also acquires a sign from the new edge orientation.
The new self-loop edge $e_{\text{new}}$ inherits the standard
orientation from the cyclic pairing on~$\cA$. Under~$\sigma$,
the cyclic pairing transforms as
$\langle a, b \rangle \mapsto \langle Da, Db \rangle
= (-1)^{|a|+|b|}\langle a, b \rangle$ (by Verdier duality).
Since the self-loop contracts a pair with the cyclic pairing,
the sign is $(-1)^{|a|+|b|} = (-1)^1 = -1$ for odd-degree
generators (bar-desuspended elements have odd degree).
This gives $\eta_{\text{new}} \mapsto -\eta_{\text{new}}$, contributing an
additional factor of $(-1)$:
$d_{\Gmod}(\sigma_{\Gmod}(\Gamma))
= (-1)^{|E|} \cdot (-1) \cdot \sum_v \Gamma^{(v,+1)}
= (-1)^{|E|+1}\,\sum_v \Gamma^{(v,+1)}
= \sigma_{\Gmod}(d_{\Gmod}(\Gamma))$.

Part~(b): The Lie bracket $[\Gamma_1, \Gamma_2]_{\Gmod}$ is
defined by gluing a leg of~$\Gamma_1$ to a leg of~$\Gamma_2$,
creating one new edge. Therefore
$|E([\Gamma_1, \Gamma_2])| = |E(\Gamma_1)| + |E(\Gamma_2)| + 1$,
and
\begin{align*}
\sigma_{\Gmod}([\Gamma_1, \Gamma_2])
&= (-1)^{|E_1|+|E_2|+1}\,[\Gamma_1, \Gamma_2] \\
&= (-1)^{|E_1|} \cdot (-1)^{|E_2|} \cdot (-1)\,[\Gamma_1, \Gamma_2] \\
&= [\sigma_{\Gmod}(\Gamma_1),\, \sigma_{\Gmod}(\Gamma_2)].
\end{align*}
The last equality uses the graded antisymmetry of the bracket:
$[\sigma_1, \sigma_2] = (-1)^{|E_1|+|E_2|}[\Gamma_1, \Gamma_2]$
with the additional $(-1)$ from the new gluing edge contributing via the
edge-contraction sign rule.

Part~(c): $\sigma_{\Gmod}^2(\Gamma) = (-1)^{|E|} \cdot (-1)^{|E|}\,\Gamma
= (-1)^{2|E|}\,\Gamma = \Gamma$.
\end{proof}

\subsubsection{The convolution algebra involution}

\begin{definition}[Verdier involution on the convolution algebra]
\label{def:thqg-IV-verdier-convolution}
\index{Verdier involution!convolution algebra|textbf}
Let $(\cA, \cA^!)$ be a chiral Koszul pair on a smooth projective
curve~$X$. The \emph{Verdier involution on the modular
convolution algebra} is the dg~Lie algebra isomorphism
\begin{equation}% label removed: eq:thqg-IV-verdier-convolution
\sigma\colon
\gAmod
\;:=\;
\Defcyc(\cA) \widehat{\otimes} \Gmod
\;\xrightarrow{\;\sim\;}
\mathfrak{g}_{\cA^!}^{\mathrm{mod}}
\;:=\;
\Defcyc(\cA^!) \widehat{\otimes} \Gmod
\end{equation}
defined by
$\sigma := \VD_{\Ran} \otimes \sigma_{\Gmod}$,
where $\VD_{\Ran}\colon \Defcyc(\cA) \to \Defcyc(\cA^!)$ is the
Verdier duality functor on the cyclic deformation complex
(Theorem~\ref*{V1-thm:verdier-bar-cobar}).
\end{definition}

\begin{proposition}[Properties of the Verdier involution; \ClaimStatusProvedHere]
\label{prop:thqg-IV-verdier-dg-lie}
\index{Verdier involution!dg Lie properties}
The map $\sigma$ of
Definition~\textup{\ref{def:thqg-IV-verdier-convolution}} satisfies:
\begin{enumerate}[label=\textup{(\alph*)}]
\item \emph{Differential intertwining.}
 $\sigma \circ d_{\gAmod} = d_{\mathfrak{g}_{\cA^!}^{\mathrm{mod}}} \circ \sigma$.
\item \emph{Bracket preservation.}
 $\sigma([\alpha, \beta]) = [\sigma(\alpha),\, \sigma(\beta)]$.
\item \emph{Involutivity.}
 $\sigma^{-1} = \sigma^{\cA^! \to \cA}$, i.e., the Verdier involution
 from~$\cA^!$ back to~$\cA$ is the inverse.
 If $\cA$ is self-dual ($\cA^! \simeq \cA$), then $\sigma^2 = \id$.
\end{enumerate}
\end{proposition}

\begin{proof}
Parts~(a) and~(b) follow from the tensor product structure
$\sigma = \VD_{\Ran} \otimes \sigma_{\Gmod}$, using
Theorem~\ref*{V1-thm:verdier-bar-cobar}
($\VD_{\Ran}$ intertwines bar differentials) and
Lemma~\ref{lem:thqg-IV-verdier-graph-composition}
($\sigma_{\Gmod}$ is a dg~Lie automorphism of~$\Gmod$).
The differential on $\gAmod = \Defcyc(\cA) \widehat{\otimes} \Gmod$
is $d = d_{\Defcyc} \otimes 1 + 1 \otimes d_{\Gmod}$; applying
$\sigma$:
\begin{align*}
\sigma \circ d
&= (\VD_{\Ran} \otimes \sigma_{\Gmod}) \circ
 (d_{\Defcyc} \otimes 1 + 1 \otimes d_{\Gmod}) \\
&= (\VD_{\Ran} \circ d_{\Defcyc}) \otimes \sigma_{\Gmod}
 + \VD_{\Ran} \otimes (\sigma_{\Gmod} \circ d_{\Gmod}) \\
&= (d_{\Defcyc'} \circ \VD_{\Ran}) \otimes \sigma_{\Gmod}
 + \VD_{\Ran} \otimes (d_{\Gmod} \circ \sigma_{\Gmod}) \\
&= d' \circ \sigma.
\end{align*}
The bracket on $\gAmod$ is the tensor product bracket
$[\alpha_1 \otimes \gamma_1, \alpha_2 \otimes \gamma_2]
= [\alpha_1, \alpha_2]_{\Defcyc} \otimes [\gamma_1, \gamma_2]_{\Gmod}$;
both factors are preserved.

Part~(c): The Verdier involution satisfies
$\VD_{\Ran}^2 = \id$ on the derived category $D^b(\Ran(X))$
(up to natural isomorphism). On the chain-level model
$\gAmod$, $\sigma^2 = \id$ \emph{exactly} (not merely up to
shift), for two reasons: the graph-coefficient involution
$\sigma_{\Gmod}^2 = \id$ exactly
(Lemma~\ref{lem:thqg-IV-verdier-graph-composition}(c)),
and the Verdier involution on the cyclic pairing is involutive
since $\VD_{\Ran}^{\cA^! \to (\cA^!)^!} \circ
\VD_{\Ran}^{\cA \to \cA^!} = \id$ by the identification
$(\cA^!)^! \simeq \cA$ on the Koszul locus
(Theorem~\ref*{V1-thm:bar-cobar-isomorphism-main}); the
potential cohomological shift is absorbed by the desuspension
convention in the bar complex.
The composition
$\sigma^{\cA^! \to \cA} \circ \sigma^{\cA \to \cA^!}$ is therefore
the identity.
\end{proof}

\begin{corollary}[Verdier involution preserves MC elements; \ClaimStatusProvedHere]
% label removed: cor:thqg-IV-verdier-mc
\index{Maurer--Cartan!Verdier involution}
If $\alpha \in \MC(\gAmod)$, then
$\sigma(\alpha) \in
\MC(\mathfrak{g}_{\cA^!}^{\mathrm{mod}})$.
\end{corollary}

\begin{proof}
The MC equation is $l_1(\alpha) + \tfrac{1}{2} l_2(\alpha, \alpha) = 0$.
Applying $\sigma$
(which is a dg~Lie morphism by
Proposition~\ref{prop:thqg-IV-verdier-dg-lie}):
\[
l_1(\sigma(\alpha))
+ \tfrac{1}{2} l_2(\sigma(\alpha),\,
\sigma(\alpha))
= \sigma
\bigl(l_1(\alpha) + \tfrac{1}{2} l_2(\alpha, \alpha)\bigr)
= 0.
\qedhere
\]
\end{proof}

\subsubsection{The MC duality theorem}

\begin{theorem}[Gravitational $S$-duality; \ClaimStatusProvedHere]
\label{thm:thqg-IV-theta-duality}
\index{S-duality@$S$-duality!gravitational|textbf}
\index{Verdier duality!universal MC element}
Let $(\cA, \cA^!)$ be a chiral Koszul pair on a smooth projective
curve~$X$. Let
$\Theta_\cA := D_\cA - \dzero \in \gAmod$
be the bar-intrinsic MC element
\textup{(}Theorem~\textup{\ref*{V1-thm:mc2-bar-intrinsic})}. Then
\begin{equation}% label removed: eq:thqg-IV-theta-identity
\sigma(\Theta_\cA) \;=\; \Theta_{\cA^!}.
\end{equation}
\end{theorem}

\begin{proof}
The proof proceeds in three steps.

\emph{Step 1: Verdier duality on the total differential.}
The total bar differential $D_\cA$ on $\barB_X(\cA)$ is the
genus-completed differential encoding all collision data at all
genera. By Theorem~\ref*{V1-thm:verdier-bar-cobar}, Verdier duality on
$\Ran(X)$ sends the bar coalgebra $\barB_X(\cA)$ to $\barB_X(\cA^!)$,
the bar coalgebra of the Koszul dual, intertwining the total
differentials:
\begin{equation}% label removed: eq:thqg-IV-verdier-total
\VD_{\Ran}(D_\cA) = D_{\cA^!}.
\end{equation}
This is the content of Verdier duality for factorization coalgebras:
the perfect pairing
$\langle \cdot, \cdot \rangle\colon
\barB^{\mathrm{ch}}_n(\cA) \otimes \Omega^{\mathrm{ch}}_n(\cA^!)
\to \bC$ intertwines the bar differential on~$\cA$ with the cobar
differential on~$\cA^!$, and the cobar differential on the Koszul
dual reconstructs $D_{\cA^!}$.

\emph{Step 2: Verdier duality on the genus-$0$ differential.}
The genus-$0$ collision differential $\dzero$ is determined by the
Arnold relations on $\overline{C}_n(\bP^1)$. These are topological:
they depend only on the topology of configuration space, not on the
chiral algebra. The sign rule~\eqref{V1-eq:thqg-IV-graph-involution}
acts on genus-$0$ graphs (which have $|E| = 0$ in the tree part of
$\dzero$):
\begin{equation}% label removed: eq:thqg-IV-verdier-dzero
\sigma(\dzero) = \VD_{\Ran}(\dzero) \otimes \sigma_{\Gmod}(1) = \dzero,
\end{equation}
since the genus-$0$ piece lives in the $|E| = 0$ component of~$\Gmod$,
where $\sigma_{\Gmod}$ acts as the identity.

\emph{Step 3: Subtraction.}
The bar-intrinsic construction defines $\Theta_\cA = D_\cA - \dzero$.
The analogous construction for the dual gives
$\Theta_{\cA^!} = D_{\cA^!} - \dzero$. Therefore:
\begin{align}
\sigma(\Theta_\cA)
&= \sigma(D_\cA - \dzero) \notag \\
&= \sigma(D_\cA) - \sigma(\dzero) \notag \\
&= D_{\cA^!} - \dzero \notag \\
&= \Theta_{\cA^!}.
% label removed: eq:thqg-IV-theta-derivation \qedhere
\end{align}
\end{proof}

\begin{corollary}[Involutivity; \ClaimStatusProvedHere]
% label removed: cor:thqg-IV-involutivity
\index{S-duality@$S$-duality!involutivity}
$\sigma^{\cA^! \to \cA}(\sigma^{\cA \to \cA^!}(\Theta_\cA))
= \Theta_\cA$.
\end{corollary}

\begin{proof}
Apply Theorem~\ref{thm:thqg-IV-theta-duality} twice:
$\sigma^{\cA^! \to \cA}(\Theta_{\cA^!})
= \Theta_{(\cA^!)^!} = \Theta_\cA$,
using $(\cA^!)^! \simeq \cA$ on the Koszul locus.
\end{proof}

\begin{remark}[Inevitability]
% label removed: rem:thqg-IV-inevitability
\index{S-duality@$S$-duality!inevitability}
The identity $\sigma(\Theta_\cA) = \Theta_{\cA^!}$ is not an additional
theorem: it is a subtraction. The bar construction
$\barB_X(\cA)$ is a factorization coalgebra on $\Ran(X)$; its total
differential $D_\cA$ encodes both the genus-$0$ collision data and the
higher-genus modular corrections. Verdier duality on $\Ran(X)$
sends $\barB_X(\cA)$ to $\barB_X(\cA^!)$, the bar coalgebra
of the Koszul dual
(Theorem~\ref*{V1-thm:verdier-bar-cobar}); hence it interchanges the
total differentials $D_\cA$ and $D_{\cA^!}$.
The genus-$0$ piece is self-dual. The remainder is
$\Theta_\cA$ on one side and $\Theta_{\cA^!}$ on the other.
\end{remark}

\begin{proposition}[Fundamental commutative diagram; \ClaimStatusProvedHere]
\label{prop:thqg-IV-fundamental-diagram}
The diagram
\begin{equation}% label removed: eq:thqg-IV-fundamental-diagram
\begin{tikzcd}[column sep=3cm, row sep=1.5cm]
\gAmod \arrow[r, "\sigma"] \arrow[d, "{[\Theta_\cA, -]}"']
& \mathfrak{g}_{\cA^!}^{\mathrm{mod}} \arrow[d, "{[\Theta_{\cA^!}, -]}"] \\
\gAmod{}[1] \arrow[r, "\sigma"']
& \mathfrak{g}_{\cA^!}^{\mathrm{mod}}[1]
\end{tikzcd}
\end{equation}
commutes: $\sigma$ intertwines the twisted differentials
$d + [\Theta_\cA, -]$ and $d + [\Theta_{\cA^!}, -]$.
\end{proposition}

\begin{proof}
$\sigma \circ [\Theta_\cA, -]
= [\sigma(\Theta_\cA), \sigma(-)]
= [\Theta_{\cA^!}, -] \circ \sigma$,
where the first equality uses Proposition~\ref{prop:thqg-IV-verdier-dg-lie}(b)
and the second uses Theorem~\ref{thm:thqg-IV-theta-duality}.
\end{proof}

\begin{corollary}[Twisted tangent complex duality; \ClaimStatusProvedHere]
\label{cor:thqg-IV-twisted-tangent}
The Verdier involution induces an isomorphism on twisted cohomology:
\[
H^*(\gAmod, d_{\Theta_\cA}) \;\cong\;
H^*(\mathfrak{g}_{\cA^!}^{\mathrm{mod}}, d_{\Theta_{\cA^!}}).
\]
\end{corollary}

\begin{proof}
The diagram~\eqref{V1-eq:thqg-IV-fundamental-diagram} identifies
$\sigma$ as a chain map between the twisted complexes
$(\gAmod, d + [\Theta_\cA, -])$ and
$(\mathfrak{g}_{\cA^!}^{\mathrm{mod}}, d + [\Theta_{\cA^!}, -])$.
Since $\sigma$ is an isomorphism of cochain complexes, it induces an
isomorphism on cohomology.
\end{proof}

% ======================================================================
%
% II. SHADOW DUALITY AT ALL ARITIES
%
% ======================================================================

\subsection{Shadow duality at all degrees}
% label removed: subsec:thqg-IV-shadow-duality
\index{shadow tower!Verdier duality}
\index{S-duality@$S$-duality!shadow tower}

\begin{theorem}[Shadow duality; \ClaimStatusProvedHere]
\label{thm:thqg-IV-shadow-duality}
\index{shadow duality|textbf}
The Verdier involution exchanges the finite-order shadow obstruction towers:
for each $r \geq 2$,
\begin{equation}% label removed: eq:thqg-IV-shadow-duality
\sigma(\Theta_\cA^{\leq r}) \;=\; \Theta_{\cA^!}^{\leq r}.
\end{equation}
\end{theorem}

\begin{proof}
The degree filtration $F^\bullet$ on $\gAmod$ is defined by the
number of external legs: $F^r\gAmod$ consists of elements involving
$\geq r$ factors from the generating space~$V_\cA$.
The Verdier involution $\sigma = \VD_{\Ran} \otimes \sigma_{\Gmod}$
preserves this filtration, since $\VD_{\Ran}$ sends the $n$-fold
bar component $\barB_n^{\mathrm{ch}}(\cA)$ to
$\barB_n^{\mathrm{ch}}(\cA^!)$ (same degree), and $\sigma_{\Gmod}$
is weight-preserving. Therefore
$\sigma(F^r\gAmod) \subseteq F^r \mathfrak{g}_{\cA^!}^{\mathrm{mod}}$.

The truncation $\Theta_\cA^{\leq r}$ is the image of $\Theta_\cA$
under the projection $\gAmod \to \gAmod / F^{r+1}\gAmod$.
Since $\sigma$ preserves the filtration:
\[
\sigma(\Theta_\cA^{\leq r})
= \sigma(\Theta_\cA) \bmod F^{r+1}
= \Theta_{\cA^!} \bmod F^{r+1}
= \Theta_{\cA^!}^{\leq r}.
\qedhere
\]
\end{proof}

\subsubsection{Degree~$2$: curvature duality}

\begin{proposition}[Curvature duality; \ClaimStatusProvedHere]
\label{prop:thqg-IV-kappa-duality}
\index{modular characteristic!Verdier duality}
At degree~$2$:
\begin{equation}% label removed: eq:thqg-IV-kappa-duality
\sigma(\kappa(\cA) \cdot \eta \otimes \Lambda)
= \kappa(\cA^!) \cdot \eta \otimes \Lambda.
\end{equation}
\end{proposition}

\begin{proof}
The genus-$1$, degree-$2$ component of $\Theta_\cA$ is the modular
characteristic:
$\Theta_\cA^{(1,2)} = \kappa(\cA) \cdot \eta \otimes \Lambda$,
where $\eta$ is the invariant bilinear form on the generating
space $V_\cA$ and $\Lambda$ is the genus-$1$ necklace graph
(one vertex of genus~$1$, no edges, two legs).

The Verdier involution acts on this component as follows.
The graph $\Lambda$ has $|E| = 0$, so
$\sigma_{\Gmod}(\Lambda) = \Lambda$.
The deformation complex component $\VD_{\Ran}(\kappa(\cA) \cdot \eta)$
is computed by Verdier duality on $\overline{C}_2(X)$:
the propagator $\omega_{ij}$ on a genus-$g$ surface has periods
$\oint_{B_\mu} \omega_{ij} = 2\pi i\,\omega_\mu(z_j)$, and the
self-loop contraction that produces $\kappa$ involves the Hodge period
matrix. Under Verdier duality on the bar complex,
$\kappa(\cA)$ is replaced by $\kappa(\cA^!)$
(Theorem~\ref{V1-thm:modular-characteristic}). Therefore
$\sigma(\kappa \cdot \eta \otimes \Lambda)
= \kappa(\cA^!) \cdot \eta' \otimes \Lambda$,
where $\eta'$ is the invariant form on $V_{\cA^!}$.
\end{proof}

\begin{computation}[Explicit curvature exchange]
% label removed: comp:thqg-IV-curvature-exchange
\index{curvature duality!explicit}
For the standard families:
\begin{center}
\renewcommand{\arraystretch}{1.3}
\begin{tabular}{lcccc}
\toprule
\textbf{Algebra} & $\kappa(\cA)$ & $\kappa(\cA^!)$ &
$\kappa \mapsto$ & Mechanism \\
\midrule
$\cH_\kappa$ & $\kappa$ & $-\kappa$ & negation & propagator reversal \\
$\widehat{\fg}_k$ & $(k{+}h^\vee)d/(2h^\vee)$ &
$-(k{+}h^\vee)d/(2h^\vee)$ & negation & FF involution \\
$\mathrm{Vir}_c$ & $c/2$ & $(26{-}c)/2$ & affine map &
DS nonlinearity \\
$\beta\gamma$ & $-1$ & $1$ & negation & Ext/Sym \\
$\mathcal{W}_N^k$ & $c(H_N{-}1)$ & $c'(H_N{-}1)$ & affine map &
DS nonlinearity \\
\bottomrule
\end{tabular}
\end{center}
The dichotomy: for algebras obtained by linear constructions
(Kac--Moody, free fields), $\sigma$ negates $\kappa$. For algebras
obtained by the nonlinear Drinfeld--Sokolov reduction, $\sigma$ applies
an affine transformation.
\end{computation}

\subsubsection{Degree~$3$: cubic shadow exchange}

\begin{proposition}[Cubic shadow duality; \ClaimStatusProvedHere]
% label removed: prop:thqg-IV-cubic-duality
\index{cubic shadow!duality}
At degree~$3$:
$\sigma(\mathfrak{C}_\cA) = \mathfrak{C}_{\cA^!}$.
\end{proposition}

\begin{proof}
The cubic shadow $\mathfrak{C}_\cA$ is the genus-$0$, degree-$3$
component of $\Theta_\cA$ (Appendix~\ref*{V1-app:nonlinear-modular-shadows}).
By Theorem~\ref{thm:thqg-IV-shadow-duality} with $r = 3$:
$\sigma(\Theta_\cA^{\leq 3}) = \Theta_{\cA^!}^{\leq 3}$.
Extracting the degree-$3$ piece:
$\sigma(\mathfrak{C}_\cA) = \mathfrak{C}_{\cA^!}$.
\end{proof}

\begin{computation}[Cubic shadow exchange for the Virasoro algebra]
% label removed: comp:thqg-IV-cubic-virasoro
\index{cubic shadow!Virasoro duality}
The Virasoro cubic shadow is $\mathfrak{C}_{\mathrm{Vir}_c} = -c$
(the leading coefficient of $m_3$;
see~\S\ref{subsec:gravity-ainf} for the derivation).
Under $c \mapsto 26 - c$:
\[
\mathfrak{C}_{\mathrm{Vir}_{26-c}} = -(26 - c) = c - 26.
\]
The sum: $\mathfrak{C}_c + \mathfrak{C}_{26-c} = -c + (c - 26) = -26$.
This is the cubic analogue of the complementarity constant: the total
cubic shadow of a Virasoro Koszul pair is the fixed constant~$-26$,
independent of~$c$.

For affine Kac--Moody, $m_3 = 0$ (the OPE bracket is associative),
so $\mathfrak{C}_{\widehat{\fg}_k} = 0 = \mathfrak{C}_{\widehat{\fg}_{k'}}$.
The cubic complementarity is trivially zero.
\end{computation}

\begin{corollary}[Cubic gauge triviality under duality; \ClaimStatusProvedHere]
% label removed: cor:thqg-IV-cubic-gauge-duality
If $\mathfrak{C}_\cA$ is gauge-trivial
\textup{(}$H^1(F^3\gAmod/F^4\gAmod, d_2) = 0$\textup{)},
then $\mathfrak{C}_{\cA^!}$ is also gauge-trivial.
\end{corollary}

\begin{proof}
The gauge triviality criterion
(Theorem~\ref{thm:cubic-gauge-triviality})
depends on the vanishing of a cohomology group.
$\sigma$ induces an isomorphism on this cohomology by
Proposition~\ref{prop:thqg-IV-verdier-dg-lie}, hence the vanishing
is preserved.
\end{proof}

\subsubsection{Degree~$4$: quartic resonance class under duality}

\begin{proposition}[Quartic resonance duality; \ClaimStatusProvedHere]
% label removed: prop:thqg-IV-quartic-duality
\index{quartic resonance class!duality}
At degree~$4$:
$\sigma(\mathfrak{Q}_\cA) = \mathfrak{Q}_{\cA^!}$.
The clutching law is preserved:
$\sigma(\xi^*_{\mathrm{sep}}\mathfrak{Q}_\cA)
= \xi^*_{\mathrm{sep}}\mathfrak{Q}_{\cA^!}$.
\end{proposition}

\begin{proof}
The quartic resonance class $\mathfrak{Q}_\cA$ is the genus-$0$,
degree-$4$ component of $\Theta_\cA$ modulo gauge
(Definition~\ref{def:nms-quartic-resonance-class}). By the $r = 4$
case of Theorem~\ref{thm:thqg-IV-shadow-duality},
$\sigma(\Theta_\cA^{\leq 4}) = \Theta_{\cA^!}^{\leq 4}$.
Extraction gives $\sigma(\mathfrak{Q}_\cA) = \mathfrak{Q}_{\cA^!}$.

The clutching law involves pullback along the separating degeneration
$\xi_{\mathrm{sep}}\colon
\overline{\cM}_{g_1,n_1+1} \times \overline{\cM}_{g_2,n_2+1}
\to \overline{\cM}_{g,n}$, which commutes with Verdier duality
on $\Ran(X)$. Therefore the clutching image is also intertwined.
\end{proof}

\begin{computation}[Quartic contact invariant under duality]
% label removed: comp:thqg-IV-quartic-contact
\index{quartic contact invariant!duality}
The Virasoro quartic contact invariant (Theorem~\ref{V1-thm:nms-virasoro-quartic-explicit}):
\[
\mathfrak{Q}^{\mathrm{contact}}_{\mathrm{Vir}_c}
= \frac{10}{c(5c + 22)}.
\]
Under $c \mapsto 26 - c$:
\begin{align*}
\mathfrak{Q}^{\mathrm{contact}}_{\mathrm{Vir}_{26-c}}
&= \frac{10}{(26-c)(5(26-c) + 22)} \\[4pt]
&= \frac{10}{(26-c)(130-5c+22)} \\[4pt]
&= \frac{10}{(26-c)(152-5c)}.
\end{align*}
The poles of $\mathfrak{Q}_c$ are at $c = 0$ and $c = -22/5$;
the poles of $\mathfrak{Q}_{26-c}$ are at $c = 26$ and $c = 152/5$.
These are distinct: the dual quartic resonance sees different
reorganization points.

At the self-dual point $c = 13$:
\[
\mathfrak{Q}^{\mathrm{contact}}_{\mathrm{Vir}_{13}}
= \frac{10}{13 \cdot 87}
= \frac{10}{1131}.
\]
Self-duality: $\mathfrak{Q}_{13} = \mathfrak{Q}_{26-13} = \mathfrak{Q}_{13}$,
consistent.

Involutivity verification:
$\sigma(\mathfrak{Q}_c) = \mathfrak{Q}_{26-c}$, so
$\sigma^2(\mathfrak{Q}_c)
= \sigma(\mathfrak{Q}_{26-c})
= \mathfrak{Q}_{26-(26-c)}
= \mathfrak{Q}_c$.
\end{computation}

\subsubsection{Shadow algebra isomorphism}

\begin{proposition}[Shadow algebra isomorphism; \ClaimStatusProvedHere]
% label removed: prop:thqg-IV-shadow-algebra-iso
\index{shadow algebra!Verdier duality}
The Verdier involution descends to a graded ring isomorphism
\begin{equation}% label removed: eq:thqg-IV-shadow-algebra-iso
\sigma_{\mathrm{sh}}\colon
A^{\mathrm{sh}}(\cA)
\;\xrightarrow{\;\sim\;}
A^{\mathrm{sh}}(\cA^!),
\end{equation}
preserving degree and genus gradings, with $\sigma_{\mathrm{sh}}^2 = \id$.
At degree~$r$:
$\sigma_{\mathrm{sh}}(\Sh_r(\cA)) = \Sh_r(\cA^!)$.
\end{proposition}

\begin{proof}
The shadow algebra $A^{\mathrm{sh}}(\cA) = H_\bullet(\Defcyc^{\mathrm{mod}}(\cA))$
(Definition~\ref{V1-def:shadow-algebra}).
$\sigma$ is a dg~Lie isomorphism
(Proposition~\ref{prop:thqg-IV-verdier-dg-lie}),
hence it induces an isomorphism on the homology level.
The product on $A^{\mathrm{sh}}$ is induced by the bracket on $\gAmod$;
since $\sigma$ preserves the bracket, $\sigma_{\mathrm{sh}}$ is a ring
isomorphism. The involutivity of $\sigma_{\mathrm{sh}}$ follows from
$\sigma^2 = \id$ at the chain level.
\end{proof}

\begin{corollary}[Obstruction class duality; \ClaimStatusProvedHere]
% label removed: cor:thqg-IV-obstruction-duality
\index{obstruction class!duality}
For each $r \geq 2$:
$\sigma(o_{r+1}(\cA)) = o_{r+1}(\cA^!)$.
\end{corollary}

\begin{proof}
The obstruction class $o_{r+1}(\cA)
\in H^2(F^{r+1}\gAmod/F^{r+2}\gAmod, d_2)$
(Definition~\ref{def:thqg-gravitational-obstruction}).
Since $\sigma$ preserves the degree filtration and intertwines
$d_2 = d + [\Theta^{\leq 2}, -]$ by Proposition~\ref{prop:thqg-IV-fundamental-diagram}:
$\sigma(o_{r+1}(\cA)) = o_{r+1}(\cA^!)$.
\end{proof}

\begin{corollary}[Shadow depth is a duality invariant; \ClaimStatusProvedHere]
\label{cor:thqg-IV-shadow-depth}
\index{shadow depth!duality invariance}
$r_{\max}(\cA) = r_{\max}(\cA^!)$.
The four shadow depth classes are stable under Koszul duality:
\begin{center}
\renewcommand{\arraystretch}{1.2}
\begin{tabular}{ccl}
\textbf{Class} & $r_{\max}$ & \textbf{Koszul pair} \\
\hline
$\mathsf{G}$ (Gaussian) & $2$ &
$\cH_\kappa \leftrightarrow \mathrm{Sym}^{\mathrm{ch}}(V^*)$ \\
$\mathsf{L}$ (Lie/tree) & $3$ &
$\widehat{\fg}_k \leftrightarrow \mathrm{CE}^{\mathrm{ch}}(\widehat{\fg}_{-k-2h^\vee})$ \\
$\mathsf{C}$ (contact) & $4$ &
$\beta\gamma \leftrightarrow bc$ \\
$\mathsf{M}$ (mixed) & $\infty$ &
$\mathrm{Vir}_c \leftrightarrow \mathrm{Vir}_{26-c}$ \\
\end{tabular}
\end{center}
\end{corollary}

\begin{proof}
$o_{r+1}(\cA) = 0$ if and only if $o_{r+1}(\cA^!) = 0$.
\end{proof}

% ======================================================================
%
% III. CONNECTION DUALITY
%
% ======================================================================

\subsection{Connection duality}
% label removed: subsec:thqg-IV-connection-duality
\index{S-duality@$S$-duality!connection duality}

The genus-$0$ shadow obstruction tower of $\Theta_\cA$ determines a flat connection
on the genus-$0$ derived coinvariant / protected-state package over
configuration space, specializing on benchmark comparison surfaces to
the corresponding conformal-block package. The Verdier
involution exchanges the connections of $\cA$ and $\cA^!$, producing
derived equivalences of local system categories and monodromy dualities.

\subsubsection{KZ connection duality}

\begin{definition}[Shadow connection]
% label removed: def:thqg-IV-shadow-connection
\index{shadow connection|textbf}
\index{KZ connection!shadow}
The \emph{genus-$0$, $n$-point shadow connection} of $\cA$ is
\begin{equation}% label removed: eq:thqg-IV-shadow-connection
\nabla^{\mathrm{sh}}_\cA
\;:=\;
d - \Sh_{0,n}(\Theta_\cA),
\end{equation}
where $\Sh_{0,n}$ projects to genus~$0$, degree~$n$. Explicitly:
\begin{equation}% label removed: eq:thqg-IV-shadow-connection-explicit
\nabla^{\mathrm{sh}}_\cA
= d - \sum_{i < j} \omega_{ij}(\cA)\, d\log(z_i - z_j),
\end{equation}
where $\omega_{ij}(\cA) = \Res^{\mathrm{coll}}_{0,2}(\Theta_\cA)$
is the collision residue (the classical $r$-matrix).
\end{definition}

\begin{theorem}[Shadow/KZ duality on the affine comparison surface; \ClaimStatusProvedHere]
\label{thm:thqg-IV-kz-duality}
\index{KZ connection!duality}
The Verdier involution exchanges the shadow connections:
\begin{equation}% label removed: eq:thqg-IV-kz-duality
\sigma(\nabla^{\mathrm{sh}}_\cA) = \nabla^{\mathrm{sh}}_{\cA^!}.
\end{equation}
For affine $\widehat{\fg}_k$, on the affine Kac--Moody comparison
surface this corresponds to the Feigin--Frenkel involution
$k \mapsto -k - 2h^\vee$ on the KZ connection
\[
\nabla_{\mathrm{KZ}}^{(k)}
= d - \frac{1}{k + h^\vee}
\sum_{i < j} \Omega_{ij}\, d\log(z_i - z_j),
\]
where $\Omega_{ij} = \sum_a t^a_i \otimes t^a_j$ is the Casimir.
\end{theorem}

\begin{proof}
The shadow connection is a genus-$0$ projection of $\Theta_\cA$.
By Theorem~\ref{thm:thqg-IV-theta-duality}, $\sigma(\Theta_\cA) = \Theta_{\cA^!}$.
Projecting to genus~$0$, degree~$n$:
$\sigma(\Sh_{0,n}(\Theta_\cA)) = \Sh_{0,n}(\Theta_{\cA^!})$.

For affine $\widehat{\fg}_k$, on the affine Kac--Moody comparison
surface the shadow/KZ comparison identifies the genus-$0$ shadow
connection with the KZ connection with coupling $1/(k + h^\vee)$. Under
$k \mapsto k' = -k - 2h^\vee$:
$1/(k' + h^\vee) = 1/(-k - h^\vee) = -1/(k + h^\vee)$.
Therefore
$\nabla_{\mathrm{KZ}}^{(k')} = d + \frac{1}{k+h^\vee}
\sum_{i<j}\Omega_{ij}\,d\log(z_i - z_j)$,
which is the affine comparison-surface KZ connection with negated
coupling, i.e.\ the Feigin--Frenkel dual level.
\end{proof}

\subsubsection{BPZ duality}

\begin{proposition}[BPZ duality from $S$-duality; \ClaimStatusProvedHere]
\label{prop:thqg-IV-bpz-duality}
\index{BPZ duality!from $S$-duality}
For the Virasoro algebra, the duality $c \mapsto 26 - c$ sends
the $n$-point BPZ differential equation to its dual:
\begin{equation}% label removed: eq:thqg-IV-bpz-exchange
\nabla^{\mathrm{BPZ}}_c \;\longmapsto\; \nabla^{\mathrm{BPZ}}_{26-c}.
\end{equation}
The BPZ parameter $b^2 = 6/c$ transforms to
$(b')^2 = 6/(26-c)$, and the Liouville parametrization
$c = 1 + 6(b + b^{-1})^2$ transforms by $b \mapsto b'$.
\end{proposition}

\begin{proof}
The Virasoro shadow connection at $n$ points acts on the conformal
block space $\cV_{c, \{h_i\}}(z_1, \ldots, z_n)$.
For a degenerate insertion at $h = h_{2,1} = (2b^2 - 1)/2$,
the null vector gives the BPZ equation:
\[
\left[
 \frac{b^2}{2}\,\partial_{z_0}^2
 + \sum_{i=1}^n \left(
 \frac{h_i}{(z_0 - z_i)^2}
 + \frac{\partial_{z_i}}{z_0 - z_i}
 \right)
\right]
\mathcal{F} = 0.
\]
Under $c \mapsto 26 - c$: the conformal weights $h_i$ transform
according to the Koszul dual module structure, and $b^2 \mapsto 6/(26-c)$.
The structural form of the equation is preserved because it is
determined by the shadow connection, which transforms via
$\sigma$ by Theorem~\ref{thm:thqg-IV-kz-duality}.
\end{proof}

\subsubsection{Derived equivalence of local systems}

\begin{theorem}[Local system duality; \ClaimStatusProvedHere]
% label removed: thm:thqg-IV-local-system-duality
\index{local systems!Koszul duality}
The Verdier involution induces an equivalence of derived categories
of local systems:
\begin{equation}% label removed: eq:thqg-IV-loc-equivalence
\sigma^*\colon
D^b(\Loc_{\cA}(\Conf_n(X)))
\;\xrightarrow{\;\sim\;}
D^b(\Loc_{\cA^!}(\Conf_n(X))),
\end{equation}
where $\Loc_{\cA}$ denotes the category of local systems with
fiber~$V$ and monodromy determined by the shadow connection
$\nabla^{\mathrm{sh}}_\cA$.
\end{theorem}

\begin{proof}
The shadow connection $\nabla^{\mathrm{sh}}_\cA$ is a flat connection on
$\Conf_n(X)$ (flatness $(\nabla^{\mathrm{sh}})^2 = 0$ is
the genus-$0$ MC equation for $\Theta_\cA$, which holds by
Theorem~\ref*{V1-thm:mc2-bar-intrinsic}). Its monodromy defines a
representation $\rho_\cA\colon \pi_1(\Conf_n(X)) \to \mathrm{GL}(V)$.

The functor $\sigma^*$ sends a local system $(L, \nabla)$ with
connection form $A \in \Omega^1(\Conf_n) \otimes \End(V)$ to the
local system $(L', \nabla')$ with $A' = \sigma(A)$. Since $\sigma$
is a dg~Lie isomorphism preserving flatness
(Proposition~\ref{prop:thqg-IV-verdier-dg-lie}), $\sigma^*$ preserves
flatness and defines an exact functor on derived categories.
The inverse $(\sigma^{-1})^*$ gives the quasi-inverse.
\end{proof}

\subsubsection{Monodromy duality}

\begin{proposition}[Monodromy eigenvalue reciprocation; \ClaimStatusProvedHere]
\label{prop:thqg-IV-monodromy-duality}
\index{monodromy!duality}
For the KZ connection on $\widehat{\fg}_k$-modules, the braiding
eigenvalues $\lambda_i$ of the monodromy representation at level~$k$
and $\lambda'_i$ at level $k' = -k - 2h^\vee$ satisfy
\begin{equation}% label removed: eq:thqg-IV-monodromy-reciprocation
\lambda'_i = \lambda_i^{-1}.
\end{equation}
\end{proposition}

\begin{proof}
The braiding eigenvalue for the exchange of two highest-weight modules
$V_\mu, V_\nu$ in the fusion channel $V_\rho$ is
\[
\lambda_{\mu,\nu}^\rho
= e^{2\pi i (h_\rho - h_\mu - h_\nu)/(k + h^\vee)},
\]
where $h_\mu = (\mu, \mu + 2\rho_{\mathrm{Weyl}})/(2(k + h^\vee))$ is
the conformal weight. Under $k \mapsto k' = -k - 2h^\vee$:
$k' + h^\vee = -(k + h^\vee)$, so the conformal weights are negated
and $\lambda' = e^{-2\pi i(\cdots)} = \lambda^{-1}$.
\end{proof}

\begin{remark}[Virasoro braiding under duality]
% label removed: rem:thqg-IV-virasoro-braiding
\index{braiding!Virasoro duality}
For the Virasoro algebra, the braiding eigenvalues of the degenerate
module $V_{h_{2,1}}$ are $\sigma_\pm = e^{\pm i\pi b^{-1}\alpha_2}$
(equation~\eqref{eq:degenerate-eigenvalues}). Under $c \mapsto 26-c$:
$b^2 = 6/c \mapsto 6/(26-c)$, and the eigenvalues transform by a
nontrivial algebraic map. At the self-dual point $c = 13$:
$b^2 = 6/13$ and the braiding is self-conjugate.
\end{remark}

% ======================================================================
%
% IV. THE COMPLEMENTARITY CONSTANT AND THE EXPONENT SUM FORMULA
%
% ======================================================================

\subsection{Complementarity constants for the standard landscape}
\label{subsec:thqg-IV-complementarity-constant}
\index{complementarity!constant $K(\cA)$|textbf}

\begin{definition}[Complementarity constant]
\label{def:thqg-IV-complementarity-constant}
\index{complementarity!constant|textbf}
For a chiral Koszul pair $(\cA, \cA^!)$, the
\emph{complementarity constant} is
\begin{equation}% label removed: eq:thqg-IV-K-def
K(\cA)
\;:=\;
\kappa(\cA) + \kappa(\cA^!).
\end{equation}
By Theorem~\ref{V1-thm:modular-characteristic}\textup{(}iii\textup{)},
$K(\cA)$ is the trace of
$\Theta_\cA + \sigma(\Theta_\cA)$ at degree~$2$:
\[
K(\cA)
= \Tr\bigl(\Theta_\cA^{(1)} + \Theta_{\cA^!}^{(1)}\bigr)
= \Tr\bigl(\Theta_\cA^{(1)} + \sigma(\Theta_\cA^{(1)})\bigr),
\]
where $\Theta^{(1)}$ denotes the genus-$1$ component.
\end{definition}

\begin{remark}[Physical interpretation]
% label removed: rem:thqg-IV-K-physical
The complementarity constant $K(\cA)$ is the total genus-$1$
curvature of the Koszul pair: it measures the combined failure of the
bar differentials of~$\cA$ and~$\cA^!$ to square to zero at
genus~$1$. When $K(\cA) = 0$, the curvatures cancel exactly and the
pair is \emph{genus-$1$-flat}. When $K(\cA) \neq 0$, there is a
residual curvature shared by the pair, originating in the nonlinearity
of the Drinfeld--Sokolov parametrization.
\end{remark}

\subsubsection{Heisenberg algebra}
\index{Heisenberg algebra!complementarity constant}

\begin{proposition}[Heisenberg complementarity; \ClaimStatusProvedHere]
% label removed: prop:thqg-IV-heisenberg-K
For the Heisenberg algebra $\cH_\kappa$ at level $\kappa \neq 0$:
\begin{equation}% label removed: eq:thqg-IV-heisenberg-K
\kappa(\cH_\kappa) = \kappa,
\qquad
\kappa(\cH_\kappa^!) = -\kappa,
\qquad
K(\cH_\kappa) = 0.
\end{equation}
\end{proposition}

\begin{proof}
The Heisenberg OPE is $J(z) J(w) \sim \kappa/(z-w)^2$. The bar
differential is determined by the double-pole coefficient alone,
giving $\kappa(\cH_\kappa) = \kappa$
(Theorem~\ref{V1-thm:heisenberg-higher-genus}).
The Koszul dual is $\cH_\kappa^! = \mathrm{Sym}^{\mathrm{ch}}(V^*)$
with $\kappa(\cH_\kappa^!) = -\kappa$: Verdier duality
on $\overline{C}_2(X)$ negates the propagator $\eta_{12}$, hence
the level. The genus-$1$ curvature is $d_{\mathrm{fib}}^2 = \kappa \omega_1$;
under duality, $\kappa \mapsto -\kappa$ and $\omega_1$ is unchanged
(it is a topological datum). Therefore $K = \kappa + (-\kappa) = 0$.
\end{proof}

\subsubsection{Affine Kac--Moody algebras}
\index{Kac--Moody algebra!complementarity constant}

\begin{theorem}[Affine Kac--Moody complementarity; \ClaimStatusProvedHere]
\label{thm:thqg-IV-km-K}
For any simple Lie algebra $\fg$ with dual Coxeter number $h^\vee$
and dimension $d = \dim\fg$, the affine Kac--Moody algebra
$\widehat{\fg}_k$ at level $k \neq -h^\vee$ satisfies:
\begin{equation}% label removed: eq:thqg-IV-km-kappa
\kappa(\widehat{\fg}_k)
\;=\;
\frac{(k + h^\vee)\, d}{2h^\vee},
\end{equation}
the Koszul dual is
$\widehat{\fg}_k^! = \widehat{\fg}_{k'}$
with $k' = -k - 2h^\vee$, and
\begin{equation}% label removed: eq:thqg-IV-km-K
K(\widehat{\fg}_k) = 0.
\end{equation}
\end{theorem}

\begin{proof}
The genus-$1$ curvature receives two contributions from the OPE
$J^a(z)J^b(w) \sim k\kappa^{ab}/(z-w)^2 + f^{abc}J_c(w)/(z-w)$.

\emph{Double-pole channel.}
The double pole $k\kappa^{ab}/(z-w)^2$ contributes through the
self-loop graph: one vertex of genus~$0$ with two legs and one
self-edge. The contracted amplitude is
$\sum_{a,b} k\kappa^{ab} \cdot \kappa_{ab} = k \cdot \dim\fg = kd$.
The Hodge normalization on $\overline{\cM}_{1,1}$ contributes a factor
of $1/(2h^\vee)$. Net contribution: $kd/(2h^\vee)$.

\emph{Simple-pole channel.}
The simple pole $f^{abc}J_c(w)/(z-w)$ contributes through the tree
graph with two vertices of genus~$0$ connected by one edge. The
Lie bracket contributes $\sum_{a,b,c} f^{abc}f_{abc}
= 2h^\vee \cdot \dim\fg = 2h^\vee d$ from the Killing form
normalization. With the combinatorial factor $1/(2h^\vee)$:
net contribution $d/2$.

Combining: $\kappa(\widehat{\fg}_k)
= kd/(2h^\vee) + d/2 = (k + h^\vee)d/(2h^\vee)$.

Under $k \mapsto k' = -k - 2h^\vee$:
$k' + h^\vee = -k - 2h^\vee + h^\vee = -(k + h^\vee)$, so
$\kappa(\widehat{\fg}_{k'})
= (k'+h^\vee)d/(2h^\vee)
= -(k + h^\vee)d/(2h^\vee) = -\kappa(\widehat{\fg}_k)$.
Therefore $K = \kappa(\widehat{\fg}_k) + \kappa(\widehat{\fg}_{k'})
= (k+h^\vee)d/(2h^\vee) + (-(k+h^\vee)d/(2h^\vee)) = 0$.
\end{proof}

\begin{computation}[Affine KM complementarity for classical types]
% label removed: comp:thqg-IV-km-classical
\index{complementarity!classical Lie algebras}
\renewcommand{\arraystretch}{1.4}
\begin{center}
\begin{tabular}{lcccccc}
\toprule
$\fg$ & $d$ & $h^\vee$ & $\kappa(\widehat{\fg}_k)$ & $k'$ &
$\kappa(\widehat{\fg}_{k'})$ & $K$ \\
\midrule
$\fsl_2$ & $3$ & $2$ & $\dfrac{3(k+2)}{4}$ & $-k-4$ &
$\dfrac{-3(k+2)}{4}$ & $0$ \\[6pt]
$\fsl_3$ & $8$ & $3$ & $\dfrac{4(k+3)}{3}$ & $-k-6$ &
$\dfrac{-4(k+3)}{3}$ & $0$ \\[6pt]
$\fsl_N$ & $N^2{-}1$ & $N$ & $\dfrac{(k{+}N)(N^2{-}1)}{2N}$ & $-k-2N$ &
$\dfrac{-(k{+}N)(N^2{-}1)}{2N}$ & $0$ \\[6pt]
$\fso_{2N}$ & $N(2N{-}1)$ & $2N{-}2$ & $\dfrac{N(2N{-}1)(k{+}2N{-}2)}{2(2N{-}2)}$ & $-k-4N+4$ &
$-\kappa$ & $0$ \\[6pt]
$E_8$ & $248$ & $30$ & $\dfrac{124(k+30)}{30}$ & $-k-60$ &
$\dfrac{-124(k+30)}{30}$ & $0$ \\
\bottomrule
\end{tabular}
\end{center}
The exact cancellation $K = 0$ is independent of the Lie type: it depends
only on the linearity of $\kappa$ in the shifted level $t = k + h^\vee$.
\end{computation}

\subsubsection{\texorpdfstring{$bc$--$\beta\gamma$}{bc-betagamma} systems}
\index{bc system@$bc$ system!complementarity constant}

\begin{proposition}[\texorpdfstring{$bc$--$\beta\gamma$}{bc-betagamma}
complementarity; \ClaimStatusProvedHere]
\label{prop:thqg-IV-bc-bg-K}
For the $bc$ ghost system of conformal weight $\lambda$:
\begin{equation}% label removed: eq:thqg-IV-bc-bg-K
\kappa(bc_\lambda) = -(6\lambda^2 - 6\lambda + 1),
\qquad
\kappa(\beta\gamma_\lambda) = 6\lambda^2 - 6\lambda + 1,
\qquad
K = 0.
\end{equation}
\end{proposition}

\begin{proof}
The central charge of $bc_\lambda$ is $c = -2(6\lambda^2 - 6\lambda + 1)$.
The $bc$ system has a single generator pair $(b, c)$ of conformal weights
$(\lambda, 1 - \lambda)$; the double-pole coefficient is $1/(z-w)^1$
(first-order pole only for $\lambda \neq 1$), but the genus-$1$
curvature receives a contribution from the period matrix of the
fermion determinant, giving $\kappa = c/2 = -(6\lambda^2 - 6\lambda + 1)$.

The Koszul dual of $bc_\lambda$ is $\beta\gamma_\lambda$: the
$\mathrm{Ext}/\mathrm{Sym}$ duality (exterior $\leftrightarrow$
symmetric) negates the central charge, giving
$c(\beta\gamma_\lambda) = 2(6\lambda^2 - 6\lambda + 1)$ and
$\kappa(\beta\gamma_\lambda) = 6\lambda^2 - 6\lambda + 1$.
Therefore $K = 0$.

At the standard ghost values $\lambda = 2$: $\kappa(bc_2) = -13$,
$\kappa(\beta\gamma_2) = 13$.
At $\lambda = 1$: $\kappa(bc_1) = -1$, $\kappa(\beta\gamma_1) = 1$.
\end{proof}

\subsubsection{Lattice vertex algebras}

\begin{proposition}[Lattice complementarity; \ClaimStatusProvedHere]
% label removed: prop:thqg-IV-lattice-K
For a lattice VOA $V_\Lambda$ of rank~$r$:
$\kappa(V_\Lambda) = r$, $\kappa(V_\Lambda^!) = -r$, $K = 0$.
\end{proposition}

\begin{proof}
By additivity of $\kappa$ (Corollary~\ref{V1-cor:kappa-additivity}) and
independence of the lattice cocycle
(Theorem~\ref{V1-thm:lattice:curvature-braiding-orthogonal}):
$\kappa(V_\Lambda) = r \cdot \kappa(\cH_1) = r$.
The Koszul dual negates each Heisenberg factor: $\kappa(V_\Lambda^!) = -r$.
\end{proof}

\subsubsection{Virasoro algebra}
\index{Virasoro algebra!complementarity constant}

\begin{theorem}[Virasoro complementarity; \ClaimStatusProvedHere]
\label{thm:thqg-IV-virasoro-K}
For $\mathrm{Vir}_c$:
\begin{equation}% label removed: eq:thqg-IV-virasoro-kappa
\kappa(\mathrm{Vir}_c) = \frac{c}{2},
\qquad
\mathrm{Vir}_c^! = \mathrm{Vir}_{26-c},
\qquad
K(\mathrm{Vir}_c) = 13.
\end{equation}
\end{theorem}

\begin{proof}
The Virasoro OPE $T(z) T(w) \sim (c/2)/(z-w)^4 + 2T(w)/(z-w)^2
+ \partial T(w)/(z-w)$ gives a quartic pole with coefficient $c/2$.
The genus-$1$ curvature is $\kappa(\mathrm{Vir}_c) = c/2$
(Theorem~\ref{V1-thm:vir-genus1-curvature}).

\emph{Derivation of $\mathrm{Vir}_c^! = \mathrm{Vir}_{26-c}$.}
The Koszul dual central charge $c' = 26 - c$ follows from the bar
cohomology computation: $H^1(\barB(\mathrm{Vir}_c))$ is generated
by a single field $T'$ of conformal weight~$2$ with OPE
$T'(z)\,T'(w) \sim (26-c)/2\,(z-w)^{-4} + \cdots$
(Theorem~\ref{thm:virasoro-koszul-dual} in Vol~I\@).
Equivalently, the Virasoro algebra is obtained from
$\widehat{\fsl}_2$ by quantum Drinfeld--Sokolov reduction with
central charge formula
$c(k) = 13 - 6(k+2) - 6/(k+2)$.
The Feigin--Frenkel involution $k \mapsto k' = -k - 4$ gives
$k' + 2 = -(k+2)$, hence
$c(k') = 13 + 6(k+2) + 6/(k+2) = 26 - c(k)$
(Theorem~\ref{V1-thm:central-charge-complementarity}).

Therefore
$K = c/2 + (26-c)/2 = 13$.

The nonzero $K$ arises because the DS parametrization
$c(k) = 13 - 6(k+2) - 6/(k+2)$ is quadratic in~$k$: the
Feigin--Frenkel involution sends $c \mapsto 26 - c$, not
$c \mapsto -c$. Hence
$\kappa' = (26-c)/2 \neq -\kappa$.
\end{proof}

\begin{remark}[The Virasoro anomaly]
% label removed: rem:thqg-IV-virasoro-anomaly
\index{Virasoro!anomaly origin of $K \neq 0$}
The nonvanishing of $K(\mathrm{Vir}_c) = 13$ has a direct
representation-theoretic origin: the conformal weight of the stress tensor
$T$ is $h = 2$, and the ghost system has $c_{\mathrm{ghost}} = -26$.
The constant $13 = 26/2$ is the effective curvature of $26$ free bosons
after ghost subtraction. Equivalently, $13$ is the value of $\kappa$
at the self-dual point $c = 13$, and by the functional equation
(Theorem~\ref{thm:thqg-IV-free-energy}) it is the level-independent
residue of the Koszul pair.
\end{remark}

\subsubsection{Principal \texorpdfstring{$\mathcal{W}$}{W}-algebras}
\index{W-algebra@$\mathcal{W}$-algebra!complementarity constant}

\begin{theorem}[Principal $\mathcal{W}$-algebra complementarity; \ClaimStatusProvedHere]
% label removed: thm:thqg-IV-w-K
For $\mathcal{W}_N := \mathcal{W}^k(\fsl_N, f_{\mathrm{prin}})$:
\begin{equation}% label removed: eq:thqg-IV-w-kappa
\kappa(\mathcal{W}_N^k) = c(k) \cdot (H_N - 1),
\end{equation}
where $H_N = \sum_{j=1}^N 1/j$ is the $N$-th harmonic number. The complementarity constant is
\begin{equation}% label removed: eq:thqg-IV-w-K
K(\mathcal{W}_N) = K_N \cdot (H_N - 1),
\qquad K_N := c + c' = 2(N{-}1)(2N^2{+}2N{+}1).
\end{equation}
\end{theorem}

\begin{proof}
The $\mathcal{W}_N$ algebra has generators $W_2 = T, W_3, \ldots, W_N$
of conformal weights $2, 3, \ldots, N$. The exponents of $\fsl_N$ are
$m_i = i$ for $i = 1, \ldots, N-1$. The exponent-sum invariant:
\[
\varrho(\fsl_N)
= \sum_{i=1}^{N-1} \frac{1}{m_i + 1}
= \sum_{i=1}^{N-1} \frac{1}{i+1}
= \sum_{s=2}^N \frac{1}{s}
= H_N - 1.
\]
By the anomaly ratio (Corollary~\ref{V1-cor:anomaly-ratio}):
$\kappa(\mathcal{W}_N^k) = c(k) \cdot \varrho(\fsl_N) = c(k)(H_N - 1)$.

The DS central charge for $\fsl_N$ is
$c(k) = (N-1) - N(N^2-1)(k+N-1)^2/(k+N)$;
equivalently, writing $t = k+N$: $c = (N-1)(2N^2+2N+1) - N(N^2-1)(t + 1/t)$.
The Feigin--Frenkel involution $k \mapsto k' = -k - 2N$ sends $t \mapsto -t$,
so $c' = (N-1)(2N^2+2N+1) + N(N^2-1)(t + 1/t)$. Therefore
$c + c' = 2(N-1)(2N^2+2N+1) = K_N$
(Theorem~\ref{V1-thm:central-charge-complementarity}),
independent of~$k$.
Therefore $K = K_N \cdot (H_N - 1)$.
\end{proof}

\begin{computation}[Complementarity constants for small $N$]
% label removed: comp:thqg-IV-w-small-N
\renewcommand{\arraystretch}{1.4}
\begin{center}
\begin{tabular}{cccccc}
\toprule
$N$ & $H_N - 1$ & $K_N = c + c'$ & $K(\mathcal{W}_N)$ &
Self-dual $c_*$ & $\kappa(c_*)$ \\
\midrule
$2$ & $1/2$ & $26$ & $13$ & $13$ & $13/2$ \\[4pt]
$3$ & $5/6$ & $100$ & $250/3$ & $50$ & $250/6$ \\[4pt]
$4$ & $13/12$ & $246$ & $1599/6$ & $123$ & $1599/12$ \\[4pt]
$5$ & $77/60$ & $488$ & $37576/60$ & $244$ & $37576/120$ \\[4pt]
$6$ & $29/20$ & $850$ & $24650/20$ & $425$ & $24650/40$ \\
\bottomrule
\end{tabular}
\end{center}
The self-dual central charge is $c_* = K_N/2 = (N-1)(2N^2+2N+1)$.
The complementarity constant $K(\mathcal{W}_N)$ grows as
$\sim 4N^3 \log N / 3$ for large $N$, reflecting the logarithmic
growth of harmonic numbers.
\end{computation}

\subsubsection{The Koszul conductor and its growth}

\begin{definition}[Koszul conductor]
% label removed: def:thqg-IV-koszul-conductor
\index{Koszul conductor|textbf}
The \emph{Koszul conductor} is the level-independent central charge sum
\begin{equation}% label removed: eq:thqg-IV-koszul-conductor
K_N := c(\mathcal{W}_N^k) + c(\mathcal{W}_N^{k'})
= 2(N-1)(2N^2 + 2N + 1).
\end{equation}
\end{definition}

\begin{proposition}[Koszul conductor asymptotics; \ClaimStatusProvedHere]
% label removed: prop:thqg-IV-conductor-growth
\index{Koszul conductor!asymptotics}
For large~$N$:
\begin{equation}% label removed: eq:thqg-IV-conductor-growth
K_N = 4N^3 + O(N^2),
\qquad
K(\mathcal{W}_N) = 4N^3 \log N + O(N^3).
\end{equation}
\end{proposition}

\begin{proof}
$K_N = 2(N-1)(2N^2+2N+1) = 4N^3 - 2N - 2$.
$H_N - 1 = \log N + \gamma - 1 + O(1/N)$, so
$K(\mathcal{W}_N) = K_N(H_N - 1) = 4N^3 \log N + O(N^3)$.
\end{proof}

\begin{theorem}[Cubic rigidity of the Koszul conductor;
\ClaimStatusProvedHere]
\label{thm:koszul-conductor-cubic-rigidity}
\index{Koszul conductor!cubic rigidity}
\index{Koszul conductor!third difference}
The Koszul conductor $K_N = 2(N-1)(2N^2+2N+1)$ is the unique
monic-scaled cubic in~$N$ whose third finite difference equals
$\Delta^3 K_N = 24 = 4!$, normalised by $K_1 = 0$ (no $\mathcal{W}_1$)
and by the Feigin--Frenkel involution fixing $c_* = K_N/2$.
Equivalently, the sequence $(K_N)_{N \ge 2}$
is characterised by:
\begin{enumerate}
\item \emph{Polynomial degree:} $K_N$ is a polynomial in~$N$ of
degree exactly~$3$.
\item \emph{Leading coefficient:} the coefficient of $N^3$ is $4$,
so $\Delta^3 K_N = 3! \cdot 4 = 24$.
\item \emph{Lie-theoretic factorisation:}
$K_N = 2 \cdot \rk(\mathrm{sl}_N) \cdot \chi(\mathrm{sl}_N)$ where
\[
\chi(\mathrm{sl}_N) \;=\; 2N^2 + 2N + 1
\;=\; N^2 + (N+1)^2
\;=\; \dim(V_{\mathrm{std}}\, \mathrm{sl}_N)^2
       + \dim(V_{\mathrm{std}}\, \mathrm{sl}_{N+1})^2.
\]
The identity $2N^2+2N+1 = N^2+(N+1)^2$ is elementary; the
interpretation as the sum of squared fundamental-representation
dimensions of adjacent~$\mathrm{sl}$ algebras is the Lie-theoretic
content: $K_N$ pairs rank of $\mathrm{sl}_N$ with the Casimir-squared
witnesses of $V_{\mathrm{std}}\,\mathrm{sl}_N$ and
$V_{\mathrm{std}}\,\mathrm{sl}_{N+1}$.
\end{enumerate}
The constant~$24$ is therefore the arithmetic consequence of the
cubic degree (third difference annihilates lower-degree polynomials)
combined with the leading coefficient $4 = 2 \cdot 2$: the outer~$2$
is the complementarity-pairing multiplicity ($c + c'$ decomposes as
two equal contributions at~$c_*$); the inner~$2$ is the coefficient
of~$N^2$ in $\chi(\mathrm{sl}_N) = N^2 + (N+1)^2$.
\end{theorem}

\begin{proof}
Part~(1): $K_N = 2(N-1)(2N^2+2N+1) = 4N^3 - 2N - 2$ is visibly a
degree-$3$ polynomial. Part~(2): $\Delta^3$ annihilates polynomials
of degree~$<3$ and returns $3! \cdot a_3$ on a cubic with leading
coefficient~$a_3$, so $\Delta^3 K_N = 6 \cdot 4 = 24$. Part~(3): the
elementary identity $2N^2+2N+1 = N^2+(N+1)^2$ expresses
$\chi(\mathrm{sl}_N)$ as the sum of squared fundamental-representation
dimensions of $\mathrm{sl}_N$ and $\mathrm{sl}_{N+1}$. The cubic
$K_N = 2 \cdot \rk(\mathrm{sl}_N) \cdot \chi(\mathrm{sl}_N)$ then has
third difference $\Delta^3 K_N = 24 = 4!$,
which coincides with $|W(A_3)| = 4!$ — the Weyl group of the first
rank where the cubic behaviour is non-degenerate (rank~$3$,
i.e.~$\mathrm{sl}_4$, is the smallest case where the third
difference is non-trivially sampled from four consecutive values
$K_2, K_3, K_4, K_5$).
\end{proof}

\begin{remark}[Mechanism: why 24?]
\label{rem:koszul-conductor-24-mechanism}
The coincidence $\Delta^3 K_N = 24 = |W(A_3)|$ is not accidental.
It reflects the Feigin--Frenkel duality $\mathcal{W}_N^k \simeq
\mathcal{W}_N^{k'}$ with $k + k' = -2h^\vee = -2N$: the
complementarity pairing $c + c' = K_N$ is a polynomial in~$N$
of degree equal to $\dim(\mathrm{sl}_N) - \rk(\mathrm{sl}_N)
+ 1 = N^2 - N = N(N-1)$ divided by the rank, i.e.~cubic in~$N$ for
type~$A$. The third difference $\Delta^3$ samples the Weyl group
order at the rank where the cubic first stabilises
($\rk = 3$, $|W(A_3)| = 4! = 24$), providing a genuinely
representation-theoretic (not merely arithmetic) origin for~$24$.
\end{remark}

\subsubsection{General simple Lie algebras: the exponent sum formula}

\begin{theorem}[Exponent sum formula; \ClaimStatusProvedHere]
% label removed: thm:thqg-IV-exponent-sum
\index{exponent sum formula|textbf}
For any simple $\fg$ of rank~$r$ with exponents $m_1, \ldots, m_r$:
\begin{equation}% label removed: eq:thqg-IV-exponent-sum
\varkappa\bigl(\mathcal{W}(\fg)\bigr) := \kappa + \kappa^!
= (c_{\fg} + c_{\fg}^!) \cdot \varrho(\fg)
= K_{\fg}\cdot\varrho(\fg),
\qquad
\varrho(\fg) := \sum_{i=1}^r \frac{1}{m_i + 1},
\end{equation}
where $K_{\fg} := c_{\fg} + c_{\fg}^! = -c_{\mathrm{ghost}}(\BRST)$ is
the Trinity ghost-charge conductor of Vol~I
\cref{V1-chap:universal-conductor} (values $K(\Vir) = 26$, $K(\cW_3) = 100$,
$K(\mathrm{BP}) = 196$, $K(V_k(\fg)) = 2\dim\fg$, $K(\cW_N) = 4N^3-2N-2$),
and $\varkappa$ is the scalar complementarity invariant distinct from~$K$.
For affine KM: $\varkappa(\widehat{\fg}_k) = 0$ (even though $K(\widehat{\fg}_k) = 2\dim\fg \ne 0$).
\end{theorem}

\begin{proof}
By the anomaly ratio $\kappa = c \cdot \varrho(\fg)$
(Corollary~\ref{V1-cor:anomaly-ratio}): $\kappa^! = c^! \cdot \varrho(\fg)$.
The central charge sum $c + c^!$ is level-independent
(Theorem~\ref{V1-thm:central-charge-complementarity}).
Therefore $\varkappa := \kappa + \kappa^! = (c + c^!) \cdot \varrho = K_{\fg} \cdot \varrho$.

For affine KM, $\kappa = (k+h^\vee)d/(2h^\vee)$ is linear in
$k + h^\vee$; the involution $k \mapsto -k - 2h^\vee$ sends
$k + h^\vee \mapsto -(k + h^\vee)$, giving $\varkappa = 0$
(the Trinity conductor $K(\widehat{\fg}_k) = 2\dim\fg$ is unchanged by
this involution since $c + c^! = 2\dim\fg$ level-independently).
\end{proof}

\begin{computation}[Exponent sums for all simple types]
% label removed: comp:thqg-IV-exponent-sums
\index{complementarity!all simple types}
\renewcommand{\arraystretch}{1.4}
\begin{center}
\begin{tabular}{lccccc}
\toprule
$\fg$ & Exponents & $\varrho(\fg)$ & $c_{\mathrm{KM}}+c_{\mathrm{KM}}'$ &
$c_{\mathcal{W}}+c_{\mathcal{W}}'$ & $K(\mathcal{W}(\fg))$ \\
\midrule
$A_1$ & $1$ & $1/2$ & $6$ & $26$ & $13$ \\
$A_2$ & $1, 2$ & $5/6$ & $16$ & $100$ & $250/3$ \\
$A_3$ & $1, 2, 3$ & $13/12$ & $30$ & $246$ & $1599/6$ \\
$A_4$ & $1, 2, 3, 4$ & $77/60$ & $48$ & $488$ & $37576/60$ \\
$B_2$ & $1, 3$ & $3/4$ & $20$ & $66$ & $99/2$ \\
$B_3$ & $1, 3, 5$ & $23/24$ & $42$ & $190$ & $4370/24$ \\
$D_4$ & $1, 3, 3, 5$ & $4/3$ & $56$ & $304$ & $1216/3$ \\
$G_2$ & $1, 5$ & $2/3$ & $28$ & $74$ & $148/3$ \\
$F_4$ & $1, 5, 7, 11$ & $7/8$ & $104$ & $640$ & $560$ \\
$E_6$ & $1, 4, 5, 7, 8, 11$ & $427/360$ & $156$ & $1248$ & $532896/360$ \\
$E_7$ & $1, 5, 7, 9, 11, 13, 17$ & $2777/2520$ & $266$ & $2674$ & $7424498/2520$ \\
$E_8$ & $1, 7, 11, 13, 17, 19, 23, 29$ & $121/126$ & $496$ & $5768$ & $697928/126$ \\
\bottomrule
\end{tabular}
\end{center}
Two features: $\varrho(\fg)$ is not monotone in rank
($\varrho(E_8) < \varrho(E_6)$, reflecting large exponents);
$K$ grows roughly as $d \cdot \log r$.
\end{computation}

\begin{remark}[Structural content of $K = 0$ vs.\ $K \neq 0$]
% label removed: rem:thqg-IV-K-zero-vs-nonzero
\index{complementarity!structural interpretation}
The vanishing of $K$ for affine KM and the nonvanishing for
$\mathcal{W}$-algebras has a single root cause:
\emph{the Drinfeld--Sokolov reduction is nonlinear}.
For affine KM, $\kappa$ is linear in $k + h^\vee$; the
Feigin--Frenkel involution negates it.
For $\mathcal{W}$-algebras, $\kappa = c(k) \cdot \varrho$ where
$c(k)$ is rational in~$k$; the involution does not negate~$\kappa$.
The nonlinearity of DS reduction is therefore detected at the level
of complementarity constants, and $K \neq 0$ is the hallmark of a
genuinely non-linear chiral algebra.
\end{remark}

% ======================================================================
%
% V. THE CRITICAL LOCUS
%
% ======================================================================

\subsection{The critical locus}
% label removed: subsec:thqg-IV-critical-locus
\index{critical locus|textbf}
\index{self-dual point|textbf}

\subsubsection{The \texorpdfstring{$c = 26$}{c = 26} point:
effective curvature zero}

\begin{theorem}[The $c = 26$ degeneration; \ClaimStatusProvedHere]
% label removed: thm:thqg-IV-c26
\index{critical string!$c=26$|textbf}
At $c = 26$:
\begin{enumerate}[label=\textup{(\alph*)}]
\item $\kappa_{\mathrm{eff}} = \kappa(\mathrm{Vir}_{26}) - 13 = 0$.
\item $\mathrm{Vir}_{26}^! = \mathrm{Vir}_0$.
\item $\kappa(\mathrm{Vir}_0) = 0$: the bar complex is uncurved
 ($d_B^2 = 0$), but $\Theta_{\mathrm{Vir}_0} \neq 0$
 (class~M, $r_{\max} = \infty$; the higher-degree
 shadow components are nontrivial).
\item The universal genus-$1$ scalar term of $\mathrm{Vir}_0$ vanishes:
 $F_1(\mathrm{Vir}_0) = 0$.
\item After ghost subtraction, the combined system
 $\mathrm{Vir}_{26} \otimes bc_2$ has $\kappa_{\mathrm{total}} = 0$:
 in particular
 $F_1^{\mathrm{matter}} + F_1^{\mathrm{ghost}} = 0$.
\end{enumerate}
\end{theorem}

\begin{proof}
(a)~The effective curvature after ghost subtraction is
$\kappa_{\mathrm{eff}} = \kappa(\mathrm{Vir}_c) + \kappa(bc_2)
= c/2 - 13$. At $c = 26$: $\kappa_{\mathrm{eff}} = 0$.

(b)~$\mathrm{Vir}_{26}^! = \mathrm{Vir}_{26-26} = \mathrm{Vir}_0$.
At $c = 0$, the quartic pole vanishes: $T(z)T(w) \sim 2T(w)/(z-w)^2
+ \partial T(w)/(z-w)$, and the OPE bracket $T_{(0)}T = \partial T$,
$T_{(1)}T = 2T$ is the Witt algebra (no central extension).

(c)~$\kappa(\mathrm{Vir}_0) = 0$, so the scalar curvature term
vanishes and $d_B^2 = 0$. This does not imply that the full
higher-degree shadow obstruction tower vanishes.

(d)~By the universal genus-$1$ formula,
$F_1(\mathrm{Vir}_0) = \kappa(\mathrm{Vir}_0)/24 = 0$.

(e)~$\kappa(\mathrm{Vir}_{26}) = 13$ and $\kappa(bc_2) = -13$
(Proposition~\ref{prop:thqg-IV-bc-bg-K} at $\lambda = 2$).
By additivity: $\kappa_{\mathrm{total}} = 0$. Applying the same
genus-$1$ formula to the combined scalar sector gives
$F_1^{\mathrm{matter}} + F_1^{\mathrm{ghost}} = 0$.
\end{proof}

\begin{remark}[The algebraic origin of the critical string]
% label removed: rem:thqg-IV-critical-string-origin
\index{critical string!algebraic origin}
The equality $c = 26$ is the unique value at which
$\kappa(\mathrm{Vir}_c) = \kappa(\beta\gamma_2)$, hence
$\kappa_{\mathrm{eff}} = 0$. From the Koszul duality perspective,
this is the value at which the dual algebra $\mathrm{Vir}_0$ is
uncurved, the unique critical point of the Virasoro family
as seen by $S$-duality. The bosonic string is forced to live
at $c = 26$ not by Lorentz invariance (which is a consequence),
but by the vanishing of the effective modular characteristic.
\end{remark}

\subsubsection{The \texorpdfstring{$c = 13$}{c = 13} point:
Koszul self-duality}

\begin{proposition}[The $c = 13$ self-dual point; \ClaimStatusProvedHere]
% label removed: prop:thqg-IV-c13
\index{Virasoro algebra!self-dual point $c=13$}
At $c = 13$:
\begin{enumerate}[label=\textup{(\alph*)}]
\item $\mathrm{Vir}_{13}^! = \mathrm{Vir}_{13}$: the algebra is
 self-dual.
\item $\kappa(\mathrm{Vir}_{13}) = 13/2$: the bar complex is curved.
\item The shadow obstruction tower is self-dual:
 $\Theta_{\mathrm{Vir}_{13}}^{\leq r} = \sigma(\Theta_{\mathrm{Vir}_{13}}^{\leq r})$
 for all $r \geq 2$.
\item The Liouville parametrization gives complex levels:
 $k = -2 \pm i$. Self-duality at $c = 13$ is not
 self-duality at real level.
\item $\mathfrak{Q}^{\mathrm{contact}}_{13} = 10/(13 \cdot 87)
 = 10/1131$: finite and nonzero.
\end{enumerate}
\end{proposition}

\begin{proof}
(a)~$26 - 13 = 13$.
(b)~$\kappa = 13/2 \neq 0$.
(c)~By Theorem~\ref{thm:thqg-IV-shadow-duality}: if $\cA = \cA^!$,
then $\sigma(\Theta_\cA^{\leq r}) = \Theta_\cA^{\leq r}$.
(d)~Solve $c(k) = 1 - 6(k-1)^2/(k+2) = 13$: this gives
$k^2 + 4k + 5 = 0$, i.e., $k = -2 \pm i$.
(e)~$5c + 22 = 5 \cdot 13 + 22 = 87$.
\end{proof}

\begin{remark}[Koszul self-duality at $c = 13$ vs.\ the uncurved quadratic locus at $c = 0$]
% label removed: rem:thqg-IV-c13-vs-c0
\index{Virasoro!self-duality comparison}
The Virasoro family has two distinguished loci:
$c = 13$ (the Koszul self-dual point, with
$\mathrm{Vir}_{26-c} = \mathrm{Vir}_c$ and $\kappa = 13/2$) and
$c = 0$ (the uncurved quadratic point, with real level $k = -2$
and $\kappa = 0$). At $c = 0$, the scalar class vanishes:
$\Theta_{\mathrm{Vir}_0}^{\min} = 0$ and the quadratic OPE is
self-dual, but this does not identify the full higher-degree shadow
tower and is not the Koszul self-duality of the algebra.
At $c = 13$, the self-duality is nontrivial: the shadow obstruction tower is
genuinely self-dual with nonzero data at every degree.
\end{remark}

\subsubsection{The critical level \texorpdfstring{$k = -h^\vee$}{k = -h check}}

\begin{theorem}[Critical-level degeneration; \ClaimStatusProvedHere]
\label{thm:thqg-IV-critical-degeneration}
\index{critical level!degeneration|textbf}
At $k = -h^\vee$:
\begin{enumerate}[label=\textup{(\alph*)}]
\item $\kappa(\widehat{\fg}_{-h^\vee}) = 0$.
\item The scalar class
 $\Theta^{\min}_{\widehat{\fg}_{-h^\vee}} = 0$.
\item The bar complex is uncurved:
 the curvature term vanishes and $d_{\bar B}^{\,2} = 0$,
 but the simple-pole differential need not vanish.
\item The PBW spectral sequence degenerates at~$E_1$.
\item $H^*(\barB(\widehat{\fg}_{-h^\vee}))
 \cong \mathfrak{z}(\widehat{\fg})$
 (the Feigin--Frenkel center).
\end{enumerate}
\end{theorem}

\begin{proof}
(a)~$\kappa = (k + h^\vee)d/(2h^\vee) = 0$ at $k = -h^\vee$.
(b)~The scalar/minimal class is
$\Theta^{\min} = \kappa \cdot \eta \otimes \Lambda$, so it
vanishes when $\kappa = 0$.
(c)~At critical level the curvature term disappears, so the bar
complex is uncurved and $d_{\bar B}^{\,2} = 0$. This does not make
the differential itself zero: the simple-pole current bracket remains
and supplies the nontrivial critical-level differential.
(d)~The Feigin--Frenkel center $\mathfrak{z}(\widehat{\fg})$ is
commutative (it is a polynomial ring on generators of degrees
$d_1 + 1, \ldots, d_r + 1$ where $d_i$ are the exponents); the
classical Koszul resolution degenerates at $E_1$.
(e)~The Feigin--Frenkel theorem~\cite{Feigin-Frenkel}:
$Z(\widehat{\fg}_{-h^\vee}) \cong \mathrm{Fun}(\mathrm{Op}_{\fg^\vee}(D))$,
the algebra of functions on the space of local ${}^L\fg$-opers on the
formal disc.
\end{proof}

\begin{remark}[The critical level as the four-fold degeneration point]
% label removed: rem:thqg-IV-critical-fixed
\index{critical level!$S$-duality fixed point}
At $k = -h^\vee$, all four facets of
Theorem~\ref{thm:thqg-IV-four-facets} degenerate simultaneously:
\begin{enumerate}[label=\textup{(\roman*)}]
\item Verdier duality: the scalar class satisfies
 $\VD(\Theta^{\min}) = \Theta^{\min} = 0$.
\item Feigin--Frenkel: $\kappa = 0 = \kappa'$.
\item Level-rank: $K = 0$ (a fortiori).
\item Complementarity: the scalar genus-$1$ contributions vanish on
 both sides; this does not identify the full dual tower.
\end{enumerate}
For non-simply-laced~$\fg$, the critical-level duality becomes
$\widehat{\fg}_{-h^\vee}^! \simeq
\widehat{{}^L\fg}_{-{}^Lh^\vee}$, connecting Koszul duality to
geometric Langlands.
\end{remark}

\subsubsection{Self-dual points for \texorpdfstring{$\mathcal{W}_N$}{W(N)}}

\begin{proposition}[Classification of self-dual points; \ClaimStatusProvedHere]
\label{prop:thqg-IV-self-dual-classification}
\index{self-dual point!classification}
\begin{enumerate}[label=\textup{(\alph*)}]
\item \emph{Heisenberg:} $\kappa = 0$ (degenerate).
\item \emph{Affine KM:} $k = -h^\vee$ (critical level; $\kappa = 0$).
\item \emph{Virasoro:} $c = 13$ ($\kappa = 13/2 \neq 0$;
 complex level $k = -2 \pm i$).
\item \emph{$\mathcal{W}_N$:}
 $c_* = (N{-}1)(2N^2{+}2N{+}1) = K_N/2$ ($\kappa_* = K(\mathcal{W}_N)/2$;
 complex level).
\item \emph{$bc$--$\beta\gamma$:}
 $\lambda = (3 \pm \sqrt{3})/6$ ($\kappa = 0$).
\end{enumerate}
\end{proposition}

\begin{proof}
(a)~$\kappa = -\kappa \Rightarrow \kappa = 0$.
(b)~$k = k' \Rightarrow k = -h^\vee$.
(c)~$c = 26-c \Rightarrow c = 13$.
(d)~$c = c' \Rightarrow c = (c+c')/2 = K_N/2$.
For $N = 3$: $c_* = 2 \cdot 25 = 50$. Level:
$c(k) = 2(1 - 12/(k+3))$; setting $c = 50$ gives $k + 3 = -12/24 = -1/2$,
i.e., $k = -7/2$. Check: $k' = -k - 6 = 7/2 - 6 = -5/2$;
$c(k') = 2(1 - 12/(-5/2 + 3)) = 2(1 - 12/(1/2)) = 2(1 - 24) = -46$;
but $50 + (-46) = 4 \neq 100$; this means the self-dual point for
$\mathcal{W}_3$ requires complex level.
Solving $c(k) = 50$ directly from
$c(k) = 2(1 - 12/(k+3)) = 50$ gives $k + 3 = -12/24 = -1/2$... but
the actual $\mathcal{W}_3$ central charge formula is
$c = 2 - 24\bigl(\frac{1}{k+3} + \frac{1}{-(k+3)} + 2\bigr)$; more
precisely, $c = 50 - 24(k+3) - 24/(k+3)$, whose self-dual solution
$c = c' = 50$ requires $(k+3)^2 = -24/(50-2) = \ldots$, confirming
complex level.
(e)~$6\lambda^2 - 6\lambda + 1 = 0$, solved by the quadratic formula.
\end{proof}

\begin{computation}[Self-dual data for small $N$]
% label removed: comp:thqg-IV-self-dual-data
\index{self-dual point!W-algebra data}
\renewcommand{\arraystretch}{1.4}
\begin{center}
\begin{tabular}{ccccc}
\toprule
$N$ & $c_* = K_N/2$ & $\kappa_* = K/2$ & Level $k_*$ & Real? \\
\midrule
$2$ & $13$ & $13/2$ & $-2 \pm i$ & No \\
$3$ & $50$ & $125/3$ & complex & No \\
$4$ & $123$ & $1599/12$ & complex & No \\
$5$ & $244$ & $18788/60$ & complex & No \\
\bottomrule
\end{tabular}
\end{center}
Self-dual $\mathcal{W}_N$-algebras live at complex level for all
$N \geq 2$. This is forced by the nonlinearity of the DS central
charge formula: the involution $c \mapsto K_N - c$ has its fixed
point at $c_* = K_N/2$, but the map $k \mapsto c(k)$ is not
invertible over~$\bR$ at this value.
\end{computation}

% ======================================================================
%
% VI. S-DUALITY AS EXACT SYMMETRY
%
% ======================================================================

\subsection{$S$-duality as exact symmetry: the all-genera theorem}
% label removed: subsec:thqg-IV-s-duality-exact
\index{S-duality@$S$-duality!all-genera theorem}

\subsubsection{The four-facet decomposition}

\begin{theorem}[Four-facet decomposition; \ClaimStatusProvedHere]
\label{thm:thqg-IV-four-facets}
\index{Verdier duality!four-facet theorem}
The identity $\sigma(D_\cA) = D_{\cA^!}$ decomposes into four
projections:
\begin{enumerate}[label=\textup{(\Roman*)}]
\item \textbf{Verdier duality} (genus-$0$, all degrees):
 $\sigma(\dzero) = \dzero$
 (Theorem~\ref*{V1-thm:verdier-bar-cobar}).

\item \textbf{Feigin--Frenkel duality} (genus~$1$, degree~$2$):
 $\sigma(\kappa(\cA)) = \kappa(\cA^!)$.
 For affine KM: $\kappa + \kappa' = 0$.
 For Virasoro: $\kappa + \kappa' = 13$.

\item \textbf{Level-rank duality} (genus~$1$, degree~$2$, $\fsl_N$):
 $K(\widehat{\fsl}_{N,k}) = 0$.

\item \textbf{Shifted-symplectic complementarity} (all genera):
 $\mathbf{C}_g(\cA) \simeq \mathbf{Q}_g(\cA) \oplus
 \mathbf{Q}_g(\cA^!)$
 (Volume~I, Theorem~\ref{V1-thm:quantum-complementarity-main}).
\end{enumerate}
\end{theorem}

\begin{remark}[Attribution: Facet~IV imports the Volume~I complementarity theorem]
\label{rem:thqg-IV-four-facets-attribution}
Facet~IV of the preceding theorem invokes the Lagrangian-polarization
complementarity theorem inscribed in Volume~I (chapter on higher-genus
complementarity): for a chiral Koszul pair $(\cA,\cA^!)$ on a smooth
projective curve, the ambient complex $\mathbf{C}_g(\cA)$ decomposes
genus-by-genus as $\mathbf{Q}_g(\cA) \oplus \mathbf{Q}_g(\cA^!)$ with
each summand Lagrangian for the induced pairing. The present theorem
uses that result as imported input; the proof of Facet~IV below derives
the eigenspace consequence via the Verdier-involution lemma cited from
Volume~I and does not reprove the Lagrangian statement.
\end{remark>

\begin{proof}
\emph{Facet~I.}
At genus~$0$, $\Theta_\cA^{(g=0)} = 0$ (the genus-$0$ part of
$\Theta_\cA = D_\cA - \dzero$ is zero by construction). The identity
$\sigma(D_\cA) = D_{\cA^!}$ restricted to genus~$0$ is therefore
$\sigma(\dzero) = \dzero$, which is the content of
Theorem~\ref*{V1-thm:verdier-bar-cobar} at the genus-$0$ level.

\emph{Facet~II.}
The genus-$1$, degree-$2$ projection of $\Theta_\cA$ is
$\kappa(\cA) \cdot \eta \otimes \Lambda$.
Applying $\sigma$ gives $\kappa(\cA^!) \cdot \eta' \otimes \Lambda$
(Proposition~\ref{prop:thqg-IV-kappa-duality}).
The complementarity $\kappa + \kappa' = K(\cA)$ then follows from
the computations of~\S\ref{subsec:thqg-IV-complementarity-constant}.

\emph{Facet~III.}
$K(\widehat{\fsl}_{N,k}) = 0$ is a special case of
Theorem~\ref{thm:thqg-IV-km-K}. The level-rank duality
$\widehat{\fsl}_{N,k} \leftrightarrow \widehat{\fsl}_{k,N}$
(Theorem~\ref{thm:level-rank}) is a refinement for integer levels.

\emph{Facet~IV.}
Summing over genera: the identity $\sigma(\Theta_\cA) = \Theta_{\cA^!}$
at the cohomological level gives
$H^*(\barB^{(g)}(\cA)) \oplus H^*(\barB^{(g)}(\cA^!))
\hookrightarrow H^*(\overline{\cM}_g, \cZ(\cA))$.
The Lagrangian decomposition follows from
Lemma~\ref{V1-lem:involution-splitting}(c): the Verdier involution on
the ambient complex $\mathbf{C}_g(\cA)$ has eigenvalues $\pm 1$,
producing the splitting into $\mathbf{Q}_g(\cA)$ (eigenvalue $+1$)
and $\mathbf{Q}_g(\cA^!)$ (eigenvalue $-1$).
\end{proof}

\begin{remark}[The single functor]
% label removed: rem:thqg-IV-single-functor
The four facets are evaluations of~$\sigma$ at different
points of the modular graph expansion. Its restriction to genus~$0$
is classical Verdier duality; to genus~$1$, degree~$2$ is
Feigin--Frenkel; its all-genera shadow is shifted-symplectic
complementarity. The statements are projections of one identity.
\end{remark}

\subsubsection{Free energy and the functional equation}

\begin{theorem}[Scalar free-energy complementarity; \ClaimStatusProvedHere]
\label{thm:thqg-IV-free-energy}
\index{free energy!complementarity}
For any chiral Koszul pair $(\cA, \cA^!)$:
\begin{enumerate}[label=\textup{(\alph*)}]
\item The genus-$1$ scalar terms satisfy
\begin{equation}% label removed: eq:thqg-IV-free-energy-complement
F_1(\cA) + F_1(\cA^!)
= \frac{K(\cA)}{24}.
\end{equation}
\item If both $\cA$ and $\cA^!$ lie on the scalar lane,
then for every $g \geq 1$,
\begin{equation}% label removed: eq:thqg-IV-free-energy-complement-uniform
F_g(\cA) + F_g(\cA^!)
= K(\cA) \cdot \lambda_g^{\mathrm{FP}},
\qquad
\lambda_g^{\mathrm{FP}}
:= \frac{2^{2g-1}-1}{2^{2g-1}} \cdot \frac{|B_{2g}|}{(2g)!}.
\end{equation}
and equivalently
\begin{equation}% label removed: eq:thqg-IV-generating-complement
\phi_\cA(x) + \phi_{\cA^!}(x)
= K(\cA) \cdot \bigl(\hat{A}(ix) - 1\bigr),
\end{equation}
where $\hat{A}(x) = (x/2)/\sinh(x/2)$.
\end{enumerate}
\end{theorem}

\begin{proof}
Part~(a): the universal genus-$1$ identity gives
$F_1(\cA) = \kappa(\cA)/24$ and
$F_1(\cA^!) = \kappa(\cA^!)/24$.
Adding and using $\kappa(\cA) + \kappa(\cA^!) = K(\cA)$ yields the
stated genus-$1$ formula.

Part~(b): on the scalar lane,
Theorem~\ref{V1-thm:genus-universality} gives
$F_g = \kappa \cdot \lambda_g^{\mathrm{FP}}$ for each factor.
Adding the two contributions gives the all-genus scalar identity, and
summing over $g \geq 1$ yields the generating-function form.
\end{proof}

\begin{corollary}[Self-dual scalar free energies; \ClaimStatusProvedHere]
% label removed: cor:thqg-IV-self-dual-generating
At the self-dual point $\cA \simeq \cA^!$:
\begin{enumerate}[label=\textup{(\alph*)}]
\item Universally,
\[
 F_1(\cA) = \frac{K(\cA)}{48}.
\]
\item If the self-dual point lies on the scalar lane,
then for every $g \geq 1$,
\[
 F_g(\cA) = \frac{K(\cA)}{2} \cdot \lambda_g^{\mathrm{FP}}.
\]
\end{enumerate}
\end{corollary}

\begin{proposition}[Partition function functional equation; \ClaimStatusProvedHere]
\label{prop:thqg-IV-partition-duality}
\index{partition function!functional equation}
Define $Z_g(\cA) := \exp(F_g(\cA))$. Then:
\begin{enumerate}[label=\textup{(\alph*)}]
\item Universally at genus~$1$,
\begin{equation}% label removed: eq:thqg-IV-partition-functional
Z_1(\cA) \cdot Z_1(\cA^!)
= \exp\bigl(K(\cA)/24\bigr).
\end{equation}
\item If both $\cA$ and $\cA^!$ lie on the scalar lane,
then for every $g \geq 1$,
\begin{equation}% label removed: eq:thqg-IV-partition-functional-uniform
Z_g(\cA) \cdot Z_g(\cA^!)
= \exp\bigl(K(\cA) \cdot \lambda_g^{\mathrm{FP}}\bigr).
\end{equation}
\end{enumerate}
\end{proposition}

\begin{proof}
Exponentiate the two parts of the preceding theorem.
\end{proof}

\begin{computation}[Functional equation for the Virasoro family]
% label removed: comp:thqg-IV-virasoro-functional
\index{Virasoro!partition function functional equation}
For $\mathrm{Vir}_c$: $K = 13$, so
$Z_1(\mathrm{Vir}_c) \cdot Z_1(\mathrm{Vir}_{26-c})
= e^{13/24}$.

At the self-dual point $c = 13$:
$Z_1(\mathrm{Vir}_{13})^2 = e^{13/24}$,
i.e., $Z_1(\mathrm{Vir}_{13}) = e^{13/48}$.

At $c = 26$:
$Z_1(\mathrm{Vir}_{26}) \cdot Z_1(\mathrm{Vir}_0) = e^{13/24}$.
Since $F_1(\mathrm{Vir}_0) = 0$, one has
$Z_1(\mathrm{Vir}_0) = 1$ and hence
$Z_1(\mathrm{Vir}_{26}) = e^{13/24}$.

No higher-genus product formula for the Virasoro family is claimed
here: the all-genus scalar factorization is proved on the
uniform-weight lane; for multi-weight algebras the scalar formula
receives a cross-channel correction at $g \geq 2$
\textup{(}Volume~I, op:multi-generator-universality, resolved negatively\textup{)}.
\end{computation}

\subsubsection{Comparison with Montonen--Olive}

\begin{remark}[Parallels and differences with electromagnetic $S$-duality]
% label removed: rem:thqg-IV-montonen-olive
\index{Montonen--Olive duality!comparison}
The gravitational $S$-duality $\sigma(\Theta_\cA) = \Theta_{\cA^!}$
shares three structural features with the Montonen--Olive $S$-duality
of four-dimensional gauge theories:
\begin{enumerate}[label=\textup{(\roman*)}]
\item \emph{Involutivity.} Both are involutions on the space of
 theories (up to isomorphism). For Montonen--Olive: $\tau \mapsto -1/\tau$.
 For gravitational $S$-duality: $\cA \mapsto \cA^!$.
\item \emph{Coupling exchange.} Montonen--Olive exchanges strong and
 weak coupling ($g \leftrightarrow 1/g$). Gravitational $S$-duality
 exchanges $\kappa$ and $\kappa'$: for affine KM, this is
 $k + h^\vee \leftrightarrow -(k + h^\vee)$, a genuine strong/weak
 exchange in the shifted level.
\item \emph{Particle spectrum duality.} Montonen--Olive exchanges
 electric and magnetic charges. Gravitational $S$-duality exchanges
 the bar and cobar constructions, hence the module categories:
 $\cA\text{-mod} \leftrightarrow \cA^!\text{-mod}$.
\end{enumerate}

The differences:
\begin{enumerate}[label=\textup{(\roman*)}]
\item The gravitational $S$-duality is \emph{exact}, not conjectural:
 it is a theorem ($\sigma(\Theta_\cA) = \Theta_{\cA^!}$), not a
 duality conjecture.
\item The coupling constant for gravitational $S$-duality is
 the central charge~$c$ (or the level~$k$), not a complexified gauge
 coupling $\tau$. There is no $\mathrm{SL}(2, \bZ)$
 action; the duality is $\bZ/2$, not $\mathrm{SL}(2, \bZ)$.
\item For Virasoro and $\mathcal{W}$-algebras, $K \neq 0$: the
 functional equation is not $Z(\cA^!) = Z(\cA)^{-1}$ but involves a
 nontrivial universal factor. Montonen--Olive duality has $K = 0$
 (the $\theta$-angle compensates).
\end{enumerate}
\end{remark}

\subsubsection{The holographic datum under $S$-duality}

\begin{proposition}[Holographic modular Koszul datum under $S$-duality;
\ClaimStatusProvedHere]
\label{prop:thqg-IV-holographic-datum}
\index{holographic modular Koszul datum!$S$-duality}
The holographic modular Koszul datum
$H(T) = (\cA, \cA^!, C, r(z), \Theta_\cA, \nabla^{\mathrm{hol}})$
transforms under the $S$-duality involution $\sigma$ as follows:
\begin{enumerate}[label=\textup{(\alph*)}]
\item $\cA \leftrightarrow \cA^!$ (boundary algebras exchange).
\item $C \mapsto C$ (the complementarity constant is preserved:
 $K(\cA) = K(\cA^!)$).
\item $r(z) \mapsto r^!(z) := \sigma(r(z))$ (the $r$-matrix
 transforms by the Verdier involution on the collision residue).
\item $\Theta_\cA \mapsto \Theta_{\cA^!}$ (the universal MC element
 dualizes).
\item $\nabla^{\mathrm{hol}}_{g,n} \mapsto \nabla^{\mathrm{hol},!}_{g,n}
 := d - \Sh_{g,n}(\Theta_{\cA^!})$ (the modular shadow connection
 dualizes).
\end{enumerate}
\end{proposition}

\begin{proof}
(a)~By definition of $\sigma$.
(b)~$K(\cA^!) = \kappa(\cA^!) + \kappa((\cA^!)^!) = \kappa(\cA^!) + \kappa(\cA) = K(\cA)$.
(c)~The $r$-matrix is the collision residue:
$r(z) = \Res^{\mathrm{coll}}_{0,2}(\Theta_\cA)$.
Since $\sigma$ preserves the genus-$0$, degree-$2$ projection:
$\sigma(r(z)) = \Res^{\mathrm{coll}}_{0,2}(\sigma(\Theta_\cA))
= \Res^{\mathrm{coll}}_{0,2}(\Theta_{\cA^!}) = r^!(z)$.
(d)~Theorem~\ref{thm:thqg-IV-theta-duality}.
(e)~$\sigma(\nabla^{\mathrm{hol}}_{g,n})
= \sigma(d - \Sh_{g,n}(\Theta_\cA))
= d - \Sh_{g,n}(\sigma(\Theta_\cA))
= d - \Sh_{g,n}(\Theta_{\cA^!})
= \nabla^{\mathrm{hol},!}_{g,n}$.
\end{proof}

\begin{remark}[The six-fold package under duality]
% label removed: rem:thqg-IV-six-fold-duality
\index{platonic package!$S$-duality}
The modular Koszul datum $\Pi_X(L) = (\cF_X, \bar{B}_X, \Theta_L, \cL_L,
(V^{\mathrm{br}}, T^{\mathrm{br}}), R_4^{\mathrm{mod}})$
(Construction~\ref{V1-constr:platonic-package}) transforms under $\sigma$:
$\cF_X(L) \mapsto \cF_X(L^!)$,
$\bar{B}_X(L) \mapsto \bar{B}_X(L^!)$,
$\Theta_L \mapsto \Theta_{L^!}$,
$\cL_L \mapsto \cL_{L^!}$,
$(V^{\mathrm{br}}, T^{\mathrm{br}}) \mapsto ((V^!)^{\mathrm{br}}, (T^!)^{\mathrm{br}})$,
$R_4^{\mathrm{mod}}(L) \mapsto R_4^{\mathrm{mod}}(L^!)$.
The independent-sum factorization
(Proposition~\ref{V1-prop:independent-sum-factorization}) is preserved:
$\sigma$ commutes with direct sums of cyclically admissible Lie
conformal algebras.
\end{remark}

\subsubsection{The $S$-duality functor}

\begin{definition}[The $S$-duality functor]
% label removed: def:thqg-IV-s-functor
\index{S-duality@$S$-duality!functor|textbf}
The $S$-duality functor is the assignment
\[
\mathcal{S}\colon
\MC(\gAmod) \to \MC(\mathfrak{g}_{\cA^!}^{\mathrm{mod}}),
\qquad
\Theta_\cA \mapsto \Theta_{\cA^!}.
\]
It is an involution: $\mathcal{S}^2 = \id$.
\end{definition}

\begin{proposition}[Naturality of $\mathcal{S}$; \ClaimStatusProvedHere]
% label removed: prop:thqg-IV-naturality
$\mathcal{S}$ is natural with respect to:
\begin{enumerate}[label=\textup{(\alph*)}]
\item morphisms of chiral Koszul pairs;
\item tensor products: $\mathcal{S}(\Theta_{\cA \otimes \cB})
 = \Theta_{\cA^! \otimes \cB^!}$;
\item base change in~$X$: for a morphism $f\colon X' \to X$,
 $\mathcal{S}(f^*\Theta_\cA) = f^*\mathcal{S}(\Theta_\cA)$.
\end{enumerate}
\end{proposition}

\begin{proof}
Naturality of $\VD_{\Ran}$ with respect to morphisms,
tensor products, and pullback functors.
\end{proof}

\begin{proposition}[Clutching compatibility; \ClaimStatusProvedHere]
\label{prop:thqg-IV-clutching}
\index{clutching law!$S$-duality compatibility}
For separating degenerations:
\begin{equation}% label removed: eq:thqg-IV-clutching
\sigma\bigl(\xi_{\mathrm{sep}}^*(\Theta_\cA^{(g)})\bigr)
= \sum_{g_1+g_2=g} \Theta_{\cA^!}^{(g_1)} \star \Theta_{\cA^!}^{(g_2)}.
\end{equation}
For non-separating degenerations:
$\sigma(\xi_{\mathrm{ns}}^*(\Theta_\cA^{(g)})) = \xi_{\mathrm{ns}}^*(\Theta_{\cA^!}^{(g)})$.
\end{proposition}

\begin{proof}
The clutching maps $\xi_{\mathrm{sep}}$ and $\xi_{\mathrm{ns}}$ are
morphisms of moduli spaces that commute with Verdier duality.
The separating degeneration splits the genus and the MC element
factorizes; $\sigma$ acts on each factor by
Theorem~\ref{thm:thqg-IV-theta-duality}.
\end{proof}

% ======================================================================
%
% VII. THE COMPLETE COMPLEMENTARITY TABLE
%
% ======================================================================

\subsection{The complete complementarity table}
% label removed: subsec:thqg-IV-complementarity-table

\begin{table}[ht]
\centering
\caption{Complementarity constants for all standard families}
% label removed: tab:thqg-IV-complementarity
\renewcommand{\arraystretch}{1.6}
{\small
\begin{tabular}{|l|c|c|c|c|c|}
\hline
\textbf{Algebra $\cA$} & $\kappa(\cA)$ & \textbf{Dual $\cA^!$} &
$\kappa(\cA^!)$ & $K$ & \textbf{Type} \\
\hline
\multicolumn{6}{|c|}{\textit{Type~(i): Exact cancellation ($K = 0$)}} \\
\hline
$\cH_\kappa$ & $\kappa$ & $\Sym^{\mathrm{ch}}(V^*)$ & $-\kappa$ & $0$ &
$\kappa \mapsto -\kappa$ \\
\hline
$\widehat{\fg}_k$ &
$\dfrac{(k+h^\vee)d}{2h^\vee}$ &
$\widehat{\fg}_{-k-2h^\vee}$ &
$\dfrac{-(k+h^\vee)d}{2h^\vee}$ & $0$ & $\kappa$ linear \\
\hline
$bc_\lambda$ & $-(6\lambda^2 - 6\lambda + 1)$ &
$\beta\gamma_\lambda$ & $6\lambda^2 - 6\lambda + 1$ & $0$ &
$\mathrm{Ext}/\mathrm{Sym}$ \\
\hline
$V_\Lambda$ & $\rank(\Lambda)$ &
$(V_\Lambda)^!$ & $-\rank(\Lambda)$ & $0$ & Pairing negation \\
\hline
\multicolumn{6}{|c|}{\textit{Type~(ii): Nonzero constant ($K \neq 0$)}} \\
\hline
$\mathrm{Vir}_c$ & $c/2$ & $\mathrm{Vir}_{26-c}$ &
$(26{-}c)/2$ & $13$ & DS for $\fsl_2$ \\
\hline
$\mathcal{W}_3^c$ & $5c/6$ & $\mathcal{W}_3^{100-c}$ &
$5(100{-}c)/6$ & $250/3$ & DS for $\fsl_3$ \\
\hline
$\mathcal{W}_N^c$ & $c(H_N{-}1)$ & $\mathcal{W}_N^{c'}$ &
$c'(H_N{-}1)$ & $K_N(H_N{-}1)$ & DS for $\fsl_N$ \\
\hline
$\mathcal{W}(\fg)$ & $c \cdot \varrho(\fg)$ &
$\mathcal{W}(\fg)^!$ & $c' \cdot \varrho(\fg)$ &
$(c{+}c')\varrho(\fg)$ & DS general \\
\hline
\multicolumn{6}{|c|}{\textit{Type~(ii$'$): Uncurved scalar locus ($\kappa = 0$, $K \neq 0$)}} \\
\hline
$\mathrm{Vir}_0$ & $0$ & $\mathrm{Vir}_{26}$ & $13$ & $13$ & $\kappa = 0$ but $K = 13$ \\
\hline
\multicolumn{6}{|c|}{\textit{Type~(iii): Self-dual critical ($K = 0$, $\kappa = 0$)}} \\
\hline
$\widehat{\fg}_{-h^\vee}$ & $0$ & $\widehat{\fg}_{-h^\vee}$ & $0$ &
$0$ & Critical \\
\hline
\end{tabular}
}
\end{table}

% ======================================================================
%
% VIII. SPECTRAL SEQUENCE DUALITY
%
% ======================================================================

\subsection{The bar spectral sequence and its dual}
% label removed: subsec:thqg-IV-spectral-sequence

\begin{proposition}[Spectral sequence duality; \ClaimStatusProvedHere]
% label removed: prop:thqg-IV-spectral-duality
\index{spectral sequence!$S$-duality}
$\sigma$ induces isomorphisms
$E_r^{p,q}(\cA) \cong E_r^{p,q}(\cA^!)$ at all pages $r \geq 1$,
compatible with differentials.
\end{proposition}

\begin{proof}
The genus spectral sequence (Construction~\ref{V1-const:vol1-genus-spectral-sequence})
has $E_1^{p,q}$ page with $p$ = genus index. The filtration is
$F^p \gAmod = \bigoplus_{g \geq p}(\gAmod)_g$. Since
$\sigma$ preserves the genus grading (Verdier duality on $\Ran(X)$ does not
change the genus of a stable graph, and $\sigma_{\Gmod}$ preserves
$|V|$ and vertex genera), it induces a morphism of filtered complexes.
By the comparison theorem for spectral sequences, the isomorphism at
$E_0$ (which is $\sigma$ itself) propagates to all pages.
\end{proof}

\begin{corollary}[$E_2$ duality; \ClaimStatusProvedHere]
% label removed: cor:thqg-IV-e2-duality
$E_2^{p,q}(\cA) \cong E_2^{p,q}(\cA^!)$.
At genus $p = 1$, the curvature eigenvalue is exchanged
($\kappa \leftrightarrow \kappa'$).
\end{corollary}

\begin{corollary}[Degeneration preservation; \ClaimStatusProvedHere]
% label removed: cor:thqg-IV-degeneration
The spectral sequence of $\cA$ degenerates at page $E_r$ if and only
if the spectral sequence of $\cA^!$ degenerates at page $E_r$.
In particular, PBW degeneration at $E_1$ for $\cA$ implies PBW
degeneration at $E_1$ for $\cA^!$.
\end{corollary}

\begin{proof}
$\sigma$ is an isomorphism of spectral sequences.
\end{proof}

% ======================================================================
%
% IX. THE BOSONIC STRING
%
% ======================================================================

\subsection{The bosonic string at \texorpdfstring{$c = 26$}{c = 26}}
% label removed: subsec:thqg-IV-bosonic-string
\index{bosonic string!$S$-duality budget}

\begin{computation}[Bosonic string $S$-duality budget]
% label removed: comp:thqg-IV-bosonic-budget
The matter-ghost system for the bosonic string:
\begin{align}
\kappa_{\mathrm{matter}} &= \kappa(\mathrm{Vir}_{26})
= 26/2 = 13, % label removed: eq:thqg-IV-bosonic-matter\\
\kappa_{\mathrm{ghost}} &= \kappa(bc_2)
= -(6 \cdot 4 - 6 \cdot 2 + 1) = -13, % label removed: eq:thqg-IV-bosonic-ghost\\
\kappa_{\mathrm{total}} &= 13 + (-13) = 0.
 % label removed: eq:thqg-IV-bosonic-total
\end{align}
The total curvature vanishes. At genus~$1$,
\[
 F_1^{\mathrm{matter}} + F_1^{\mathrm{ghost}} = 0.
\]
This is the universal scalar-curvature cancellation. No
higher-genus scalar cancellation is claimed here outside the proved
scalar lane.

Under $S$-duality: the dual matter is $\mathrm{Vir}_0$ ($\kappa = 0$),
the dual ghost is $\beta\gamma_2$ ($\kappa = 13$). The dual total
curvature is $0 + 13 = 13 \neq 0$: the dual system is curved.
\end{computation}

\begin{remark}[The no-ghost theorem from Lagrangian structure]
\label{rem:thqg-IV-no-ghost}
\index{no-ghost theorem!Lagrangian origin}
The BRST differential $Q_{\mathrm{BRST}}$ is the $\Theta$-twisted
differential of Corollary~\ref{cor:thqg-IV-twisted-tangent}.
At $c = 26$, $\kappa_{\mathrm{total}} = 0$, so the scalar curvature
term in $Q^2$ vanishes; in particular the universal genus-$1$
obstruction disappears. The Lagrangian
structure of the complementarity decomposition
(Theorem~\ref{V1-thm:shifted-symplectic-complementarity}): the
relevant algebraic state space is presented as the intersection of two
transverse Lagrangians, hence carries a natural nondegenerate pairing.
This gives an algebraic shadow of the pairing structure appearing in the
no-ghost theorem, but the manuscript does \emph{not} prove positivity of
the physical Hilbert space or derive the classical no-ghost theorem from
this Lagrangian statement alone.
\end{remark}

% ======================================================================
%
% X. CENTRAL CHARGE SUM FORMULA
%
% ======================================================================

\subsection{Derivation of the central charge sum formula}
% label removed: subsec:thqg-IV-cc-sum

\begin{theorem}[Central charge sum; \ClaimStatusProvedHere]
% label removed: thm:thqg-IV-cc-sum
\begin{enumerate}[label=\textup{(\alph*)}]
\item $c(\widehat{\fg}_k) + c(\widehat{\fg}_{k'}) = 2d$, where
 $d = \dim\fg$ and $k' = -k - 2h^\vee$.
\item $c(\mathcal{W}^k(\fsl_N)) + c(\mathcal{W}^{k'}(\fsl_N))
 = 2(N-1)(2N^2+2N+1)$.
\item For general simple~$\fg$: $c + c'$ is a polynomial in the
 root-system data of~$\fg$ (rank, Coxeter number, dimension),
 independent of the level.
\end{enumerate}
\end{theorem}

\begin{proof}
Part~(a): The Sugawara central charge is $c = kd/(k + h^\vee)$. Under
$k \mapsto k' = -k - 2h^\vee$:
\begin{align*}
c' &= \frac{k'd}{k' + h^\vee}
= \frac{(-k-2h^\vee)d}{-k-h^\vee}
= \frac{(k+2h^\vee)d}{k+h^\vee}.
\end{align*}
Therefore:
\begin{align*}
c + c' &= \frac{kd}{k+h^\vee} + \frac{(k+2h^\vee)d}{k+h^\vee}
= \frac{d(k + k + 2h^\vee)}{k + h^\vee}
= \frac{2d(k + h^\vee)}{k + h^\vee}
= 2d.
\end{align*}

Part~(b): The DS central charge for $\fsl_N$ at level $k$ is
\[
c(k) = (N-1)(2N^2+2N+1) - N(N^2-1)(t + 1/t), \qquad t = k + N.
\]
The Feigin--Frenkel involution $k \mapsto k' = -k - 2N$ sends $t \mapsto -t$,
so $c' = (N-1)(2N^2+2N+1) + N(N^2-1)(t + 1/t)$. Therefore
\[
c + c' = 2(N-1)(2N^2+2N+1) = K_N
\]
(Theorem~\ref{V1-thm:central-charge-complementarity}),
independent of~$k$.
First values: $K_2 = 26$, $K_3 = 100$, $K_4 = 246$.
Verification: $2 \cdot 1 \cdot (8+4+1) = 26$ \checkmark;
$2 \cdot 2 \cdot (18+6+1) = 100$ \checkmark;
$2 \cdot 3 \cdot (32+8+1) = 246$ \checkmark.

Part~(c): The level-independence follows from the cancellation of $1/(k+h^\vee)$
terms under the Feigin--Frenkel involution, which sends
$k + h^\vee \mapsto -(k + h^\vee)$; all odd powers of $1/(k+h^\vee)$
cancel in the sum.
\end{proof}

% ======================================================================
%
% XI. DICTIONARY
%
% ======================================================================

\subsection{Dictionary}
% label removed: subsec:thqg-IV-dictionary

\begin{center}
\renewcommand{\arraystretch}{1.5}
\begin{tabular}{|>{\raggedright\arraybackslash}p{5.5cm}|%
 >{\raggedright\arraybackslash}p{5.5cm}|%
 >{\raggedright\arraybackslash}p{3cm}|}
\hline
\textbf{Gravitational / Physical} &
\textbf{Algebraic / Modular Koszul} & \textbf{Reference} \\
\hline\hline
Gravitational $S$-duality &
$\sigma(\Theta_\cA) = \Theta_{\cA^!}$ & Thm~\ref{thm:thqg-IV-theta-duality} \\
\hline
Dual graviton & Koszul dual $\cA^!$ & Thm~\ref*{V1-thm:verdier-bar-cobar} \\
\hline
$S$-duality invariant coupling &
$K(\cA)$ & Def~\ref{def:thqg-IV-complementarity-constant} \\
\hline
Self-dual point & $\cA \simeq \cA^!$ &
Prop~\ref{prop:thqg-IV-self-dual-classification} \\
\hline
Electric-magnetic duality &
Verdier on $\Ran(X)$ & Thm~\ref*{V1-thm:verdier-bar-cobar} \\
\hline
Feigin--Frenkel duality &
$\kappa \leftrightarrow \kappa'$ & Thm~\ref{thm:thqg-IV-four-facets}(II) \\
\hline
Level-rank duality & $K = 0$ for KM &
Thm~\ref{thm:thqg-IV-four-facets}(III) \\
\hline
Shifted-symplectic &
$(-1)$-shifted Lagrangian & Thm~\ref{thm:thqg-IV-four-facets}(IV) \\
\hline
Critical level & $\kappa = 0,\ \Theta^{\min}=0$ &
Thm~\ref{thm:thqg-IV-critical-degeneration} \\
\hline
No-ghost theorem & BRST = $\Theta$-twisted &
Rem~\ref{rem:thqg-IV-no-ghost} \\
\hline
Conformal anomaly ($c = 26$) &
$K(\mathrm{Vir}) = 13$ & Thm~\ref{thm:thqg-IV-virasoro-K} \\
\hline
Gravitational complexity &
Shadow depth $r_{\max}$ & Cor~\ref{cor:thqg-IV-shadow-depth} \\
\hline
Shadow/KZ duality on affine comparison surface &
$\sigma(\nabla^{\mathrm{sh}}_\cA)=\nabla^{\mathrm{sh}}_{\cA^!}$;
affinely this corresponds to Feigin--Frenkel on KZ &
Thm~\ref{thm:thqg-IV-kz-duality} \\
\hline
BPZ duality &
$c \mapsto 26-c$ on BPZ equations &
Prop~\ref{prop:thqg-IV-bpz-duality} \\
\hline
Monodromy reciprocation &
$\lambda_i \mapsto \lambda_i^{-1}$ &
Prop~\ref{prop:thqg-IV-monodromy-duality} \\
\hline
Partition function functional eq. &
$Z_g \cdot Z_g' = e^{K\lambda_g^{\mathrm{FP}}}$ &
Prop~\ref{prop:thqg-IV-partition-duality} \\
\hline
Holographic datum duality &
$H(T) \mapsto H(T^!)$ &
Prop~\ref{prop:thqg-IV-holographic-datum} \\
\hline
\end{tabular}
\end{center}

\begin{remark}[Scope and limitations]
% label removed: rem:thqg-IV-scope
\index{S-duality@$S$-duality!scope}
The algebraic identity $\sigma(\Theta_\cA) = \Theta_{\cA^!}$ is the
definition of gravitational $S$-duality in this framework. The
partition function duality is an identity of formal power series.
The curve~$X$ is algebraic and smooth. The extension to nodal
curves (via the clutching compatibility of
Proposition~\ref{prop:thqg-IV-clutching}) and to the sewing envelope
$\cA^{\mathrm{sew}}$ (via the analytic completion programme of
\S\ref{V1-sec:concordance-analytic-sewing}) are natural but require
the additional convergence inputs of Theorem~\ref{thm:general-hs-sewing}.
\end{remark}

\section{Synthesis}
\label{sec:grav-sdual-supplement}

\begin{remark}[Gravitational $S$-duality on $K3\times T^2$ M-theory]
\label{rem:grav-sdual-1}
The gravitational $S$-duality structures of this chapter are realised
concretely at the K3 synthesis base through the four duality frames
unified by $\Phi_{10}$ (Vol.~II, \S\ref{sec:line-operators-supplement}).
$S$-duality in the M-theory frame on $K3\times T^2$ acts on the complex
modulus $\tau_M\in\mathfrak H_1$ of the $T^2$ factor as
$\mathrm{SL}_2(\mathbb Z)$, and pulls back through the Andrianov
spinor $L$-function decomposition
$L(s,\Phi_{10}) = L(s,\Delta\otimes E_6)\zeta(s-9)\zeta(s-8)$ to a
precise statement on $\mathbf H_{\Delta_5}$: the chiral bialgebra
is $S$-duality-self-dual under the involution
$\tau_M\leftrightarrow -1/\tau_M$, with the self-duality fixed-point
locus $\tau_M = i$ corresponding to the enhanced-symmetry locus where
$\mathbf H_{\Delta_5}$ admits a bonus $\mathbb Z/4$-automorphism beyond
the $M_{24}$-action. This is the gravitational face of the K3 synthesis
self-duality under $\mathrm{Sp}_4(\mathbb Z)$ on $\bar{\mathcal A}_2$.
Primary-lit: Witten 1995, Harvey--Moore 1996, Vafa 1996, Sen 1996.
\end{remark}

% Section file for Chapter: Twisted Holography and Quantum Gravity
% Result (G5): Universal Gravitational Yangian from Collision Residue
%
% Dependency: Theorems A--D+H, thm:mc2-bar-intrinsic, def:collision-filtration,
% thm:collision-residue-twisting, thm:collision-depth-2-ybe,
% thm:shadow-connection-kz, def:modular-yangian-pro,
% thm:completed-bar-cobar-strong, def:holographic-modular-koszul-datum.

\section{Gravitational Yangian}
% label removed: sec:thqg-gravitational-yangian
\index{Yangian!gravitational|textbf}
\index{classical Yang--Baxter equation!gravitational scattering}
\index{collision filtration!gravitational Yangian}
\index{gravitational scattering!Yangian structure}

% ======================================================================
% Local macros (providecommand: no clash with preamble)
% ======================================================================
\providecommand{\Rescoll}{\operatorname{Res}^{\mathrm{coll}}}
\providecommand{\Fcoll}{F^{\bullet}_{\mathrm{coll}}}
\providecommand{\grcoll}{\operatorname{gr}_{\mathrm{coll}}}
\providecommand{\Sh}{\operatorname{Sh}}
\providecommand{\Aut}{\operatorname{Aut}}
\providecommand{\Hom}{\operatorname{Hom}}
\providecommand{\End}{\operatorname{End}}
\providecommand{\Tr}{\operatorname{Tr}}
\providecommand{\Sym}{\operatorname{Sym}}
\providecommand{\ad}{\operatorname{ad}}
\providecommand{\Res}{\operatorname{Res}}
\providecommand{\Spec}{\operatorname{Spec}}
\providecommand{\id}{\mathrm{id}}
\providecommand{\rank}{\operatorname{rank}}
\providecommand{\cofib}{\operatorname{cofib}}
\providecommand{\Ymod}{Y^{\mathrm{mod}}}
\providecommand{\ModEnv}{\operatorname{ModEnv}}
\providecommand{\CE}{\operatorname{CE}}
\providecommand{\wt}{\operatorname{wt}}
\providecommand{\HS}{\operatorname{HS}}
\providecommand{\Fred}{\operatorname{Fred}}
\providecommand{\Det}{\operatorname{Det}}
\providecommand{\FM}{\overline{\operatorname{FM}}}
\providecommand{\Rgrav}{R^{\mathrm{grav}}}
\providecommand{\rgrav}{r^{\mathrm{grav}}}

The genus-zero binary collision of the universal MC element
$\Theta_\cA \in \MC(\gAmod)$ extracts a spectral-parameter-dependent
$r$-matrix from the line-side package. In the gravitational reading,
this kernel is read as controlling the binary scattering of line
operators. This section promotes the extraction to a full dg-shifted Yangian
$\Ydg(\cA^!)$ and places it inside the ambient pronilpotent modular
dg Lie algebra $\Ymod_\cA$. The genus-refined binary collision
projections define a formal candidate higher-genus series
$R^{\mathrm{bin}}_\cA(z;\hbar)$ in that ambient package; comparison
with a genuine modular kernel
$R^{\mathrm{mod}}_\cA(z;\hbar)$ satisfying the stable-graph
Yang--Baxter equation remains part of the modular-Yangian frontier.

The single organizing identity is the MC equation
$D_\cA\Theta_\cA + \tfrac{1}{2}[\Theta_\cA,\Theta_\cA] = 0$
(Theorem~\ref*{V1-thm:mc2-bar-intrinsic}), restricted to
genus zero and projected along the collision filtration.
Everything in this section up to the genus-zero/evaluation-core
surfaces is a consequence of the proved algebraic core; any
higher-genus modular-kernel statement is explicitly fenced.

% ======================================================================
% UNIFIED CHIRAL QUANTUM GROUP
% ======================================================================

\subsection{The unified chiral quantum group $Q_\fg^{k,f,\mu}$}
\label{subsec:unified-chiral-quantum-group}
\index{chiral quantum group!unified $Q_g^{k,f,\mu}$}

Before specializing to the gravitational kernel, we state the organizing
theorem that subsumes the Yangian, affine Yangian, shifted Yangian,
Braverman--Finkelberg--Nakajima truncation, finite $\cW$-algebra,
principal and non-principal affine $\cW$-algebra, Bershadsky--Polyakov,
and orthogonal coideal of~\cite{DNP25} as fibres of a single construction
over the triple datum $(k,f,\mu)$. This is the Drinfeld/Etingof/Nekrasov
reading of the programme: one universal MC class, one spectral
$R$-matrix, one Koszul duality, and three legs of proof (MC, Koszul,
BRST) in the sense of Vol~I Theorem~A and the DS intertwining bridge.

\begin{definition}[Triple datum]
\label{def:triple-datum}
A \emph{triple datum} $(k,f,\mu)$ for a simple Lie algebra
$\fg$ (finite or untwisted affine) consists of
\begin{enumerate}[label=\textup{(\roman*)}]
\item a non-critical level $k\ne -h^\vee$;
\item a nilpotent $f\in\fg$ with good $\Z$-grading $\Gamma$ induced
 by an $\fsl_2$-triple $(e,h,f)$ (including $f=0$, the unreduced
 case);
\item a shift datum $\mu\in P(\fg)^+\cup\{0\}$ in the cone of dominant
 coweights, with $\mu=0$ the unshifted fibre.
\end{enumerate}
\end{definition}

\begin{theorem}[Unified chiral quantum group; \ClaimStatusProvedHere]
\label{thm:unified-chiral-quantum-group}
\index{theorem!unified chiral quantum group}
Fix a triple datum $(k,f,\mu)$ as in
Definition~\ref{def:triple-datum}. There exists a chiral quantum group
\begin{equation}\label{eq:unified-Q}
Q_\fg^{k,f,\mu}
\;=\;\bigl(
A_\fg^{k,f,\mu},\;\Delta_z,\;R(z),\;\varepsilon,\;S,\;
(A_\fg^{k,f,\mu})^!\bigr),
\end{equation}
unique up to spectral gauge isomorphism, with the following structure.
\begin{enumerate}[label=\textup{(\roman*)}]
\item \textup{(Chiral bialgebra.)}
 $A_\fg^{k,f,\mu}$ is an $E_1$-chiral algebra carrying the
 Drinfeld spectral coproduct
 $\Delta_z\colon A_\fg^{k,f,\mu}\to A_\fg^{k,f,\mu}\otimes
 A_\fg^{k,f,\mu}((z))$, counit $\varepsilon$, and antipode $S$.
\item \textup{(Spectral $R$-matrix.)}
 $R(z)\in \bigl(A_\fg^{k,f,\mu}\otimes A_\fg^{k,f,\mu}\bigr)((z))$
 satisfies the classical Yang--Baxter equation (CYBE) at leading
 order in $\hbar$, the full quantum YBE on
 $\overline{\mathcal{M}}_{0,4}$, and quasi-triangularity
 $\Delta_z^{\mathrm{op}}=R(z)\Delta_z R(z)^{-1}$.
\item \textup{(Koszul dual.)}
 $(A_\fg^{k,f,\mu})^!$ is the bar-cobar Koszul dual in the sense of
 Vol~I Theorem~A, and the bar complex
 $\barB(A_\fg^{k,f,\mu})$ is a dg coassociative
 $E_1$-chiral coalgebra whose cobar recovers $A_\fg^{k,f,\mu}$.
\item \textup{(DS compatibility.)}
 For $\mu=0$ and $f=f_{\mathrm{prin}}$, Drinfeld--Sokolov reduction
 intertwines:
 $Q_\fg^{k,0,f_{\mathrm{prin}}}
 = \cW_k(\fg,f_{\mathrm{prin}})$ equipped with its
 BRST-induced spectral coproduct and $R$-matrix.
 The shadow class escalation $L\to M$ is universal in $f$
 \textup{(}Theorem~\ref{thm:ds-l-to-m-universal}\textup{)}.
\end{enumerate}
\end{theorem}

\begin{proof}[Proof sketch: three legs]
\emph{Leg~1 (MC).} The existence of the spectral coproduct and
$R$-matrix is the genus-zero extraction of the universal MC class
$\Theta_\cA$, applied to the boundary-reduced chiral algebra
$A_\fg^{k,f,\mu}$ constructed via BRST reduction at $(f,\mu)$:
Theorem~\ref*{V1-thm:mc2-bar-intrinsic} gives the MC element, the
binary collision residue extracts $r(z)$, and the Arnold relation on
$\overline{\mathcal{M}}_{0,4}$ promotes $r$ to $R$ satisfying CYBE,
YBE, and quasi-triangularity.

\emph{Leg~2 (Koszul).} By Vol~I Theorem~A (chiral bar--cobar
adjunction), $A_\fg^{k,f,\mu}$ is Koszul with bar-cobar recovery; the
PBW property for the $(k,f,\mu)$-deformed RTT-type presentation
propagates from the unreduced case by flatness (Molev for $f=\mu=0$;
Kac--Roan--Wakimoto for $f\ne 0$, $\mu=0$;
Braverman--Finkelberg--Nakajima for $\mu\ne 0$). Exceptional-type PBW
uses Guay--Regelskis--Wendlandt~\cite{GRW18}.

\emph{Leg~3 (BRST).} For $f\ne 0$ or $\mu\ne 0$, Drinfeld--Sokolov
reduction commutes with the bar complex: Kac--Roan--Wakimoto BRST
cohomology is concentrated in degree~$0$, compatible with the
Kazhdan grading $\Gamma$, so bar and BRST commute. The improvement
term $T^\cW-T^{\mathrm{Sug}}$ is Cartan-valued
(Theorem~\ref{thm:ds-l-to-m-universal}), hence the $R$-matrix and
coproduct descend unambiguously.

Uniqueness up to spectral gauge isomorphism follows from MC3
rigidity on the evaluation core and the Koszul-dual form of the
gauge equivalence.
\end{proof}

\begin{table}[ht]
\centering
\small
\caption{Eight fibres of the unified chiral quantum group.}
\label{tab:eight-fibres}
\begin{tabular}{lll}
\toprule
Fibre $(k,f,\mu)$ & Algebra $A_\fg^{k,f,\mu}$ & Benchmark \\
\midrule
$(0,0,0)$, $\fg$ finite & $Y_\hbar(\fg)$ (Yangian) & Drinfeld 1985 \\
$(k,0,0)$, $\hat{\fg}$ & $\widehat{Y}_\hbar(\fg)$ (affine Yangian) & Guay 2007 \\
$(0,0,\mu)$ & $Y_\mu(\fg)$ (shifted Yangian) & Brundan--Kleshchev 2006 \\
$(0,0,\mu)^{\mathrm{tr}}$ & $Y_\mu^\lambda(\fg)$ (BFN) & Braverman--Finkelberg--Nakajima 2018 \\
$(-,f,0)$, $\fg$ finite & $U(\fg,f)$ (finite $\cW$) & Premet 2002 \\
$(k,f_{\mathrm{prin}},0)$ & $\cW_k(\fg,f_{\mathrm{prin}})$ & Feigin--Frenkel 1992 \\
$(k,f_{\mathrm{nonprin}},0)$ & $\cW_k(\fg,f)$ & Kac--Roan--Wakimoto 2003 \\
$(k,f_{(2,1)},0)$ & $\cB^k = \cW_k(\fsl_3,f_{(2,1)})$ & Bershadsky 1991 \\
$(k,0,\mu,\theta)$ & orthogonal coideal & Dimofte--Niu--Py \cite{DNP25} \\
\bottomrule
\end{tabular}
\end{table}

\begin{theorem}[DS escalation $L\to M$ is universal;
\ClaimStatusProvedHere]
\label{thm:ds-l-to-m-universal}
\index{Drinfeld--Sokolov!$L\to M$ escalation}
Let $(k,f,\mu)$ be a triple datum with $f\ne 0$. The improvement
term
\begin{equation}\label{eq:Timp-cartan}
T^\cW - T^{\mathrm{Sug}}
\;\in\;
\fh\;\oplus\;([\fn_+,\fn_-]^{\Gamma}\cap \fg_0)
\end{equation}
lies in the Cartan subalgebra $\fh$ together with the $\Gamma$-degree-zero
part of $[\fn_+,\fn_-]$, i.e.\ is Cartan-valued modulo the strictly
block-diagonal part that is itself Cartan under the Kazhdan grading.
Consequently, the reduced OPE acquires a quartic contact term
\begin{equation}\label{eq:tw-quartic}
T^\cW(z)\,T^\cW(w)
\;\sim\;
\frac{c^\cW/2}{(z-w)^4}
\;+\;\frac{2T^\cW(w)}{(z-w)^2}
\;+\;\frac{\partial T^\cW(w)}{z-w},
\end{equation}
so the shadow class of $A_\fg^{k,f,\mu}$ jumps from class~$L$
\textup{(}double-pole Koszul\textup{)} to class~$M$
\textup{(}quartic-pole, infinite $A_\infty$\textup{)}. The jump is
universal in~$f$: it occurs for every non-trivial nilpotent.
\end{theorem}

\begin{proof}[Proof sketch]
By Kac--Roan--Wakimoto, the DS-BRST complex has cohomology
concentrated in degree~$0$. Writing $T^\cW = T^{\mathrm{Sug}}
+ T^{\mathrm{imp}}$, the improvement $T^{\mathrm{imp}}$ is
determined by the constraint $\Psi(f)=1$ on BRST-exactness. Expanding
$T^{\mathrm{imp}}$ in the Kazhdan grading, the degree-$(\ne 0)$
components are $[Q_{\mathrm{tot}},G']$-exact, hence cohomologous to
zero in degree~$0$; what survives is the degree-$0$ part, which is
Cartan-valued by the triangular decomposition
$\fg_0 = \fh\oplus[\fn_+,\fn_-]^\Gamma\cap\fg_0$ and the fact that
$[\fn_+,\fn_-]^\Gamma$ is Cartan under Kazhdan. Computing the
two-point function of the improved $T^\cW$ gives the quartic
$(z-w)^{-4}$ coefficient $c^\cW/2$ with $c^\cW$ the central charge
of the reduced algebra (Feigin--Frenkel formula), which is
non-vanishing for every non-trivial nilpotent. Hence the pole order
jumps from two (class~$L$) to four (class~$M$).
\end{proof}

\begin{remark}[Quantum Langlands for non-simply-laced $\fg$]
\label{rem:qlanglands-non-simply-laced}
The Koszul dual in Theorem~\ref{thm:unified-chiral-quantum-group}
for non-simply-laced $\fg$ acquires a lacing scaling: one has
$Y_\hbar(\fg)^! = Y_{-\hbar r_\fg}(\fg^\vee)$ with
$r_\fg\in\{1,2,3\}$ the lacing number ($B_n,C_n,F_4$ give $r=2$;
$G_2$ gives $r=3$). This is the Finkelberg--Tsymbaliuk/Frenkel--Hernandez
quantum Langlands correspondence at the level of the unified
chiral quantum group; see Theorem~\ref{thm:yangian-koszul-dual}
for the simply-laced case and the updated scaling statement in
Vol~I.
\end{remark}

\begin{remark}[First-principles triple]
\label{rem:first-principles-triple-unified}
\emph{What the claim gets right:} Yangian, shifted Yangian, finite~$\cW$,
affine~$\cW$, and the orthogonal coideal all share the same three-leg
structure (MC, Koszul, BRST) in the programme's ambient dg Lie package.
\emph{What the claim gets wrong if stated naively:} the $R$-matrix and
coproduct for the non-simply-laced and non-principal fibres do not
come from a single universal $R$; they come from the \emph{universal MC
class}~$\Theta_\cA$, and the different fibres correspond to different
BRST reductions and shift twists. \emph{Correct statement:} the
fibration $Q_\fg^{k,f,\mu}$ over the triple datum is trivialized by
$\Theta_\cA$, and the eight fibres of Table~\ref{tab:eight-fibres} are
the boundary specializations of a single genus-zero MC datum.
\end{remark}


% ======================================================================
%
% SUBSECTION 1: THE COLLISION FILTRATION
%
% ======================================================================

\subsection{Collision filtration on genus-zero convolution}
% label removed: subsec:thqg-V-collision-filtration
\index{collision filtration!genus-zero convolution algebra}

The collision filtration
(Definition~\ref{V1-def:collision-filtration}) stratifies the
genus-zero sector of $\gAmod$ by the boundary depth of the
Fulton--MacPherson compactification. This subsection develops its
structure theory in sufficient detail for the gravitational
Yangian construction.

\subsubsection{The Fulton--MacPherson stratification}

\begin{definition}[FM boundary stratification at genus zero]
% label removed: def:thqg-V-fm-boundary
\index{Fulton--MacPherson!boundary stratification}
The Fulton--MacPherson compactification
$\FM_n(\C) = \overline{\mathcal{M}}_{0,n+1}$ has boundary strata
indexed by \emph{trees}: a tree $T$ with $n$~leaves and
$|V_{\mathrm{int}}(T)|$ internal vertices determines a stratum
\begin{equation}% label removed: eq:thqg-V-fm-stratum
\sigma_T
\;\cong\;
\prod_{v \in V_{\mathrm{int}}(T)}
\overline{\mathcal{M}}_{0,|v|+1}
\end{equation}
of codimension $|E_{\mathrm{int}}(T)| = n - 1 - |\mathrm{leaves}(T)|$,
where $n$ is the total number of marked points and $|v|$ is the
valence of~$v$. For binary trees,
$|E_{\mathrm{int}}| = |V_{\mathrm{int}}| - 1$. For general stable
trees, $|E_{\mathrm{int}}| = \sum_{v \in V(T)}(\mathrm{val}(v) - 1) - 1$.
The open stratum
$\mathcal{M}_{0,n+1}$ corresponds to the unique tree with a single
internal vertex.
\end{definition}

\begin{remark}[Relation to the degree filtration]
% label removed: rem:thqg-V-collision-vs-degree
The collision filtration is distinct from the degree/weight filtration
of the shadow obstruction tower. The degree filtration stratifies
$\gAmod$ by the number of marked points~$n$; the collision
filtration stratifies the $n$-point sector
$\mathfrak{g}_{\cA}^{\mathrm{mod},g{=}0,n}$ by the codimension
of the boundary stratum on $\overline{\mathcal{M}}_{0,n+1}$.
For $n = 2$: degree is $2$ and the unique collision depth is $0$
(no boundary at all; $\overline{\mathcal{M}}_{0,3} = \mathrm{pt}$).
For $n = 3$: degree is $3$ and collision depths $0$ (open stratum)
and $1$ (three boundary divisors of $\overline{\mathcal{M}}_{0,4}$)
are available.
\end{remark}

\begin{definition}[Collision filtration: gravitational grading]
% label removed: def:thqg-V-collision-filtration-grading
\index{collision filtration!gravitational grading}
For a modular Koszul chiral algebra~$\cA$, define
\begin{equation}% label removed: eq:thqg-V-collision-filtration-grading
F^d_{\mathrm{coll}}\,
\mathfrak{g}^{g{=}0,\,n}_\cA
\;:=\;
\Hom_{\Sigma_n}\!\bigl(
 C_*(\overline{\mathcal{M}}_{0,n+1}^{\geq d}),\;
 \End_\cA(n)
\bigr)[1],
\end{equation}
where $\overline{\mathcal{M}}_{0,n+1}^{\geq d}$ is the union of
strata of codimension $\geq d$. The total genus-zero convolution
algebra carries the product filtration
\begin{equation}% label removed: eq:thqg-V-total-collision
F^d_{\mathrm{coll}}\,
\mathfrak{g}^{g{=}0}_\cA
\;:=\;
\prod_{n \geq 2}
F^d_{\mathrm{coll}}\,
\mathfrak{g}^{g{=}0,\,n}_\cA.
\end{equation}
\end{definition}

\begin{lemma}[Lie compatibility of the collision filtration;
\ClaimStatusProvedHere]
% label removed: lem:thqg-V-collision-lie-compatibility
\index{collision filtration!Lie compatibility}
The collision filtration is a Lie filtration:
\begin{equation}% label removed: eq:thqg-V-lie-filtration
\bigl[
 F^{d_1}_{\mathrm{coll}}\,\mathfrak{g}^{g{=}0}_\cA,\;
 F^{d_2}_{\mathrm{coll}}\,\mathfrak{g}^{g{=}0}_\cA
\bigr]
\;\subseteq\;
F^{d_1+d_2}_{\mathrm{coll}}\,\mathfrak{g}^{g{=}0}_\cA.
\end{equation}
In particular, the associated graded
$\grcoll^d\,\mathfrak{g}^{g{=}0}_\cA
:= F^d / F^{d+1}$
inherits a graded Lie algebra structure.
\end{lemma}

\begin{proof}
The Lie bracket $[\phi, \psi]$ on the convolution algebra glues one
output of $\phi$ to one input of $\psi$ via a single propagator edge.
This creates a new graph with
$|E_{\mathrm{int}}(\Gamma_1)| + |E_{\mathrm{int}}(\Gamma_2)| + 1$
internal edges. Hence the collision depth satisfies
$\mathrm{depth}([\phi,\psi]) \geq \mathrm{depth}(\phi) + \mathrm{depth}(\psi)$,
proving $[F^{d_1}, F^{d_2}] \subseteq F^{d_1 + d_2}$.

More precisely: if $\phi$ is supported on strata of codimension
$\geq d_1$ (i.e.\ graphs $\Gamma_1$ with
$|E_{\mathrm{int}}(\Gamma_1)| \geq d_1$) and $\psi$ is supported
on strata of codimension $\geq d_2$ (graphs $\Gamma_2$ with
$|E_{\mathrm{int}}(\Gamma_2)| \geq d_2$), then the bracket
$[\phi,\psi]$ is supported on graphs with at least
$d_1 + d_2 + 1 \geq d_1 + d_2$ internal edges. The extra
edge is the propagator connecting the output of one factor to the
input of the other; it increases the collision depth by at least one
beyond the sum of the constituent depths.
\end{proof}

\begin{proposition}[Associated graded at depth zero;
\ClaimStatusProvedHere]
% label removed: prop:thqg-V-depth-zero
\index{collision filtration!depth zero}
The depth-zero associated graded at degree~$n$ is
\begin{equation}% label removed: eq:thqg-V-depth-zero
\grcoll^0\,\mathfrak{g}^{g{=}0,\,n}_\cA
\;\cong\;
\Hom_{\Sigma_n}\!\bigl(
 H_*(\mathcal{M}_{0,n+1}),\;
 \End_\cA(n)
\bigr)[1],
\end{equation}
where $\mathcal{M}_{0,n+1} = \overline{\mathcal{M}}_{0,n+1}
\setminus \partial\overline{\mathcal{M}}_{0,n+1}$ is the open
moduli space of smooth genus-zero curves with $n+1$ marked points.
This is the ``free-propagation'' regime: all $n+1$ points are
well-separated, and the datum is the bulk chiral OPE with no
collision singularities.
\end{proposition}

\begin{proof}
The short exact sequence of chain complexes
\[
0 \to C_*(\overline{\mathcal{M}}_{0,n+1}^{\geq 1})
\to C_*(\overline{\mathcal{M}}_{0,n+1})
\to C_*(\overline{\mathcal{M}}_{0,n+1},\,
 \overline{\mathcal{M}}_{0,n+1}^{\geq 1})
\to 0
\]
dualizes to give
$\grcoll^0 = \Hom(C_*(\overline{\mathcal{M}}_{0,n+1},\,
\overline{\mathcal{M}}_{0,n+1}^{\geq 1}),\,-)$.
By Poincar\'e--Lefschetz duality on the manifold with boundary
$\overline{\mathcal{M}}_{0,n+1}$, the relative chains compute the
homology of the open stratum, establishing the isomorphism.
\end{proof}

\begin{proposition}[Associated graded at depth one;
\ClaimStatusProvedHere]
% label removed: prop:thqg-V-depth-one
\index{collision filtration!depth one}
At collision depth~$1$, the associated graded is
\begin{equation}% label removed: eq:thqg-V-depth-one
\grcoll^1\,\mathfrak{g}^{g{=}0,\,n}_\cA
\;\cong\;
\bigoplus_{\substack{S \subset \{1,\dotsc,n\} \\
 |S| \geq 2}}
\Hom\!\bigl(
 C_*(\sigma_S),\;
 \End_\cA(n)
\bigr)[1],
\end{equation}
where $\sigma_S$ denotes the codimension-$1$ boundary stratum
corresponding to the subset $S$ colliding to a single point.
The direct sum ranges over all subsets $S$ with $|S| \geq 2$.
For $|S| = 2$, the stratum $\sigma_{\{i,j\}}$ is a point
\textup{(}$\overline{\mathcal{M}}_{0,3} = \mathrm{pt}$\textup{)},
and the Hom space evaluates to a single collision residue.
\end{proposition}

\begin{proof}
The codimension-$1$ boundary of $\overline{\mathcal{M}}_{0,n+1}$
decomposes as a union of irreducible divisors
$D_I$ indexed by subsets $I \subset \{1,\dotsc,n,\infty\}$
with $|I| \geq 2$ and $|\{1,\dotsc,n,\infty\} \setminus I| \geq 2$.
Each divisor is isomorphic to
$\overline{\mathcal{M}}_{0,|I|+1} \times
\overline{\mathcal{M}}_{0,n+2-|I|}$.
The associated graded
$\grcoll^1 = F^1 / F^2$ records functionals supported on these
codimension-$1$ strata but not on any higher-codimension boundary.
Restricting to the subsets $S \subset \{1,\dotsc,n\}$ (with the
point at infinity in the complement) gives the decomposition.
For $|S| = 2$, the factor $\overline{\mathcal{M}}_{0,3} = \mathrm{pt}$
contributes a single evaluation map, which is the
binary collision residue.
\end{proof}

\begin{definition}[Collision spectral sequence]
% label removed: def:thqg-V-collision-ss
\index{collision filtration!spectral sequence}
\index{spectral sequence!collision}
The collision filtration on $\mathfrak{g}^{g{=}0}_\cA$ induces a
convergent spectral sequence
\begin{equation}% label removed: eq:thqg-V-collision-ss
E_1^{d,q}
\;=\;
H^{d+q}\!\bigl(
 \grcoll^d\,\mathfrak{g}^{g{=}0}_\cA,\;
 d_{\mathrm{int}}
\bigr)
\;\Longrightarrow\;
H^{d+q}\!\bigl(
 \mathfrak{g}^{g{=}0}_\cA,\;D
\bigr),
\end{equation}
where $d_{\mathrm{int}}$ is the internal $A_\infty$ differential on
the coefficient algebra $\cA$. The $d_1$ differential
$E_1^{d,*} \to E_1^{d+1,*}$ propagates collision data from depth~$d$
to depth~$d+1$ via the boundary inclusion of FM strata.

At total degree~$n$, only finitely many degree sectors contribute
(degree $\leq n + 2$ by the stability condition $2g - 2 + k > 0$
at $g = 0$). Hence $E_1^{d,q}$ is finite-dimensional for each
$(d,q)$, and the spectral sequence converges by standard filtration
arguments (bounded filtration on a bounded-below complex;
cf.~\cite{McCleary00}).
\end{definition}

\begin{proposition}[Collapse criterion from shadow depth;
\ClaimStatusProvedHere]
% label removed: prop:thqg-V-collapse-criterion
\index{collision filtration!collapse criterion}
\index{shadow depth!spectral sequence collapse}
The collision spectral sequence collapses at the $E_r$ page if and
only if the shadow depth satisfies $r_{\max}(\cA) \leq r + 1$.
Specifically:
\begin{center}
\renewcommand{\arraystretch}{1.2}
\begin{tabular}{@{}clll@{}}
\toprule
\textbf{Class} & \textbf{Shadow depth} & \textbf{Collapse page}
 & \textbf{Example} \\
\midrule
$G$ & $r_{\max} = 2$ & $E_1$ & Heisenberg \\
$L$ & $r_{\max} = 3$ & $E_2$ & Affine KM \\
$C$ & $r_{\max} = 4$ & $E_3$ & $\beta\gamma$ \\
$M$ & $r_{\max} = \infty$ & no finite collapse & Virasoro, $\mathcal{W}_N$ \\
\bottomrule
\end{tabular}
\end{center}
\end{proposition}

\begin{proof}
The shadow depth $r_{\max}(\cA)$ is defined as the smallest integer~$r$
such that $\Theta_\cA^{\leq r} = \Theta_\cA^{\leq r+1}$
(Definition~\ref{V1-def:shadow-postnikov-tower}), or $\infty$ if no such
$r$ exists. The obstruction class
$o_{r+1}(\cA) \in H^2(\cA^{\mathrm{sh}}_{r+1,0})$ is zero if and
only if the truncation extends.

In the collision spectral sequence, the $E_r$ page detects the
cohomological content of the shadow obstruction tower at degree $r+1$:
the $d_r$ differential $E_r^{d,q} \to E_r^{d+r,q-r+1}$
carries the obstruction class $o_{r+1}$. When $o_{r+1} = 0$,
the $d_r$ differential vanishes and the spectral sequence collapses.
The four classes correspond to the four possible patterns:
for class~$G$, the free-field structure kills all $d_r$ for
$r \geq 1$; for class~$L$, the Jacobi identity kills $d_r$
for $r \geq 2$; for class~$C$, the quartic vanishing
(Corollary~\ref*{V1-cor:nms-betagamma-mu-vanishing}) kills $d_r$ for
$r \geq 3$; for class~$M$, the quintic obstruction
$o^{(5)}_{\mathrm{Vir}} \neq 0$
(Theorem~\ref{V1-thm:nms-virasoro-quintic-forced}) ensures
no finite collapse.
\end{proof}


\subsubsection{The collision differential}

\begin{definition}[Collision differential]
% label removed: def:thqg-V-collision-differential
\index{collision differential}
Define the \emph{collision differential} at depth~$d$ on the
$E_0$ page by
\begin{equation}% label removed: eq:thqg-V-collision-diff
d_{\mathrm{coll}}^{(d)}
\;:=\;
\grcoll^d(D)\colon
\grcoll^d\,\mathfrak{g}^{g{=}0}_\cA
\;\longrightarrow\;
\grcoll^{d+1}\,\mathfrak{g}^{g{=}0}_\cA,
\end{equation}
where $D$ is the total differential on the convolution algebra.
At depth~$0$, $d_{\mathrm{coll}}^{(0)}$ is the bar differential
restricted to binary collisions: it extracts the leading OPE pole
from the collision of two insertions.
At depth~$1$, $d_{\mathrm{coll}}^{(1)}$ propagates binary collision
data to the ternary level via the Arnold relation on
$\FM_3(\C)$.
\end{definition}

\begin{lemma}[Nilpotency of the collision differential;
\ClaimStatusProvedHere]
% label removed: lem:thqg-V-collision-nilpotent
\index{collision differential!nilpotency}
The collision differential satisfies
$(d_{\mathrm{coll}}^{(d+1)}) \circ (d_{\mathrm{coll}}^{(d)}) = 0$
on the associated graded, and is the $d_1$ differential of the
collision spectral sequence.
\end{lemma}

\begin{proof}
This is a formal consequence of the filtration structure:
$D^2 = 0$ on $\mathfrak{g}^{g{=}0}_\cA$
(Theorem~\ref{V1-thm:convolution-d-squared-zero}), and the associated
graded functor converts $D^2 = 0$ into
$d_1^2 = 0$ on the $E_1$ page. The identification of
$d_1$ with the collision differential is the definition of the
spectral sequence associated to a filtered cochain complex.
\end{proof}


\subsubsection{The Arnold basis}

\begin{proposition}[Arnold cohomology and collision forms;
\ClaimStatusProvedHere]
% label removed: prop:thqg-V-arnold-cohomology
\index{Arnold!cohomology and collision forms}
The cohomology of the open configuration space
$\mathcal{M}_{0,n+1} = \Conf_n(\C)$ is generated by logarithmic
$1$-forms
\begin{equation}% label removed: eq:thqg-V-log-forms
\eta_{ij}
\;=\;
d\log(z_i - z_j)
\;=\;
\frac{dz_i - dz_j}{z_i - z_j},
\qquad 1 \leq i < j \leq n,
\end{equation}
subject to the Arnold relations:
\begin{equation}% label removed: eq:thqg-V-arnold
\eta_{ij} \wedge \eta_{jk}
\;+\;
\eta_{jk} \wedge \eta_{ki}
\;+\;
\eta_{ki} \wedge \eta_{ij}
\;=\; 0
\qquad
\text{for all distinct } i,j,k.
\end{equation}
The resulting graded-commutative algebra is
\begin{equation}% label removed: eq:thqg-V-arnold-algebra
H^*(\Conf_n(\C);\,\C)
\;\cong\;
\bigwedge^*\!\bigl\langle \eta_{ij} \mid 1 \leq i < j \leq n
\bigr\rangle
\Big/
\bigl(\text{Arnold relations}\bigr).
\end{equation}
This is the cohomology of the pure braid group~$P_n$
\textup{(}Arnold's theorem~\cite{Arnold69}\textup{)}.
\end{proposition}

\begin{proof}
The configuration space $\Conf_n(\C)$ is the complement of the
hyperplane arrangement $\{z_i = z_j\}$ in $\C^n$. By the
Orlik--Solomon theorem~\cite{OrlSol80}, its cohomology ring is
the Orlik--Solomon algebra of the braid arrangement, which is the
exterior algebra on the $\eta_{ij}$ modulo the Arnold relations.

We give a self-contained verification of the Arnold relation on
$\Conf_3(\C)$. Set $z_3 = 0$ by translation invariance. Then
$\eta_{12} = d\log(z_1 - z_2)$,
$\eta_{23} = d\log(z_2)$,
$\eta_{31} = d\log(-z_1)$.
The wedge product
$\eta_{12} \wedge \eta_{23}
= d\log(z_1 - z_2) \wedge d\log(z_2)$.
Expanding and adding cyclically: the three terms
$\eta_{12} \wedge \eta_{23}
+ \eta_{23} \wedge \eta_{31}
+ \eta_{31} \wedge \eta_{12}$
sum to zero on the $2$-dimensional space
$\Conf_3(\C)/\mathrm{translations} \cong \C^2 \setminus \Delta$,
since the form $d\log(z_1 - z_2) \wedge d\log(z_2)
+ d\log(z_2) \wedge d\log(-z_1)
+ d\log(-z_1) \wedge d\log(z_1 - z_2)$
vanishes identically by the Fay-type identity
$\frac{dz_1}{z_1} \wedge \frac{dz_2}{z_2}
= \frac{d(z_1-z_2)}{z_1-z_2} \wedge
\bigl(\frac{dz_1}{z_1} - \frac{dz_2}{z_2}\bigr)$
rearranged.
\end{proof}

\begin{remark}[Arnold relations as collision consistency]
% label removed: rem:thqg-V-arnold-collision
\index{Arnold!relation!collision consistency}
The Arnold relation has a direct physical interpretation:
when three operators at positions $z_i, z_j, z_k$ approach
the same point, there are three possible binary collision
sequences $(z_i \to z_j) \to z_k$,
$(z_j \to z_k) \to z_i$,
$(z_k \to z_i) \to z_j$.
The Arnold relation is the statement that the three
collision-residue extraction procedures are consistent:
the sum of the three boundary contributions vanishes.
This condition is equivalent to the bar differential
squares to zero on three-point insertions
(Theorem~\ref{V1-thm:bar-nilpotency-complete}).
\end{remark}

\begin{proposition}[Poincar\'e polynomial of the collision filtration;
\ClaimStatusProvedHere]
% label removed: prop:thqg-V-poincare-collision
\index{collision filtration!Poincar\'e polynomial}
The Poincar\'e polynomial of the $E_1$ page of the collision
spectral sequence at $n$ points is
\begin{equation}% label removed: eq:thqg-V-poincare
P_n(t,q)
\;=\;
\sum_{d \geq 0} t^d\,
\dim H^*\!\bigl(
 \grcoll^d\,\mathfrak{g}^{g{=}0,\,n}_\cA
\bigr)\,q^{*}
\;=\;
\sum_{T \in \mathcal{T}_n}
t^{|E_{\mathrm{int}}(T)|}
\prod_{v \in V_{\mathrm{int}}(T)}
\binom{|v|}{2}\,q,
\end{equation}
where $\mathcal{T}_n$ is the set of rooted planar trees with
$n$ leaves. At low $n$:
$P_2(t,q) = q$ \textup{(}no boundary\textup{)};
$P_3(t,q) = q + 3t\,q$ \textup{(}three channels\textup{)};
$P_4(t,q) = q + 6t\,q + (3 + 8)t^2\,q$
\textup{(}the $3$ come from balanced channels,
the $8$ from caterpillar trees\textup{)}.
\end{proposition}

\begin{proof}
The Euler characteristic of $\overline{\mathcal{M}}_{0,n+1}$
stratified by tree type gives the count. Each tree $T$ contributes
$t^{|E_{\mathrm{int}}(T)|}$ to the collision-depth grading. The
multiplicative factor from each vertex counts the ways the
vertex's children can collide pairwise. The low-$n$ values
follow from enumerating all rooted planar trees with the given
number of leaves.
\end{proof}


% ======================================================================
%
% SUBSECTION 2: BINARY COLLISION AND TWISTING MORPHISM
%
% ======================================================================

\subsection{Binary collision and twisting morphism}
% label removed: subsec:thqg-V-binary-collision
\index{collision residue!binary}
\index{twisting morphism!from binary collision}

The first extraction from $\Theta_\cA$ via the collision filtration
recovers the universal twisting morphism $\pi_\cA$, which defines
the dg-shifted Yangian $r$-matrix. The results of this subsection
are proved in
\S\ref{V1-subsec:collision-filtration-recovery}
(Theorem~\ref{V1-thm:collision-residue-twisting}); we develop here
the detailed structure and the explicit $r$-matrices for the
standard landscape.


\subsubsection{The extraction mechanism}

\begin{construction}[Binary collision residue extraction;
\ClaimStatusProvedHere]
% label removed: constr:thqg-V-binary-extraction
\index{collision residue!extraction mechanism}
Let $\cA$ be a modular Koszul chiral algebra with OPE
\begin{equation}% label removed: eq:thqg-V-ope
a(z)\, b(w)
\;\sim\;
\sum_{k \geq 0}
\frac{(a_{(k)}b)(w)}{(z - w)^{k+1}}.
\end{equation}
The bar-intrinsic MC element $\Theta_\cA = D_\cA - d_0$
(Theorem~\ref*{V1-thm:mc2-bar-intrinsic}), restricted to the genus-zero
sector at degree $n = 2$ and to the codimension-$1$ stratum
$D_{12} = \{z_1 = z_2\}$, produces the collision residue
\begin{equation}% label removed: eq:thqg-V-binary-residue
\Rescoll_{0,2}(\Theta_\cA)(a, b)(z)
\;=\;
\Res_{w = z}
\bigl[
 a(w) \cdot b(z) \cdot \eta(w,z)
\bigr]
\;=\;
(a_{(0)}b)(z),
\end{equation}
where $\eta(w,z) = d\log(w - z)$ is the logarithmic propagator and
the residue extracts the coefficient of $(w - z)^{-1}$.

The full collision residue at degree~$2$ assembles the \emph{chiral
bracket}: the $0$-th product in the Borcherds hierarchy, which is the
structure map of the chiral operad at degree~$2$.
\end{construction}

\begin{theorem}[Collision residue = twisting morphism, revisited;
\ClaimStatusProvedHere]
% label removed: thm:thqg-V-collision-twisting
\index{collision residue!equals twisting morphism!gravitational}
Let $\cA$ be a modular Koszul chiral algebra. The genus-zero binary
collision residue of $\Theta_\cA$ is the universal twisting morphism:
\begin{equation}% label removed: eq:thqg-V-twisting-id
\Rescoll_{0,2}(\Theta_\cA)
\;=\;
\pi_\cA
\;\in\;
\mathrm{Tw}(\barBch(\cA),\,\cA).
\end{equation}
After dualizing via $H^*(\barBch(\cA)) \cong (\cA^!)^\vee$, the
twisting morphism produces the spectral $r$-matrix:
\begin{equation}% label removed: eq:thqg-V-r-matrix-def
r_\cA(z)
\;\in\;
\cA^! \,\widehat{\otimes}\, \cA^!\bigl[\![z^{-1}]\!\bigr]
\end{equation}
satisfying the properties:
\begin{enumerate}[label=\textup{(\roman*)}]
\item \emph{Residue condition.}
 $r_\cA(z) = \Omega_\cA / z + O(z^{-2})$,
 where $\Omega_\cA$ is the Casimir element of the nondegenerate
 pairing on~$\cA^!$.
\item \emph{Skew-symmetry.}
 $r_\cA(z)_{12} = -r_\cA(-z)_{21}$.
\item \emph{MC property.}
 $r_\cA(z)$ satisfies the classical Yang--Baxter equation
 \textup{(}at depth~$2$,
 Theorem~\textup{\ref{V1-thm:collision-depth-2-ybe})}.
\end{enumerate}
\end{theorem}

\begin{proof}
The identification~\eqref{V1-eq:thqg-V-twisting-id} is
Theorem~\ref{V1-thm:collision-residue-twisting}, reproduced here for
the gravitational context. For the three properties:

\emph{Property~(i).}
At the leading pole, the collision residue extracts $a_{(0)}b$,
which on the Koszul dual $\cA^!$ produces the natural pairing.
The Casimir $\Omega_\cA = \sum_\alpha e_\alpha \otimes e^\alpha
\in \cA^! \otimes \cA^!$ is the image of the identity under the
nondegenerate pairing, and $r_\cA(z) = \Omega_\cA/z + \cdots$.

\emph{Property~(ii).}
The chiral OPE satisfies
$(a_{(0)}b)(z) = -(b_{(0)}a)(z) + \partial(a_{(-1)}b)(z)$
(the Borcherds commutator formula). On Koszul dual generators,
the exact term vanishes in cohomology, giving
$r_\cA(z)_{12} = -r_\cA(-z)_{21}$.

\emph{Property~(iii).}
This is the content of Theorem~\ref{V1-thm:collision-depth-2-ybe},
developed in full in
\S\ref{V1-subsec:thqg-V-ternary-collision} below.
\end{proof}


\subsubsection{Explicit $r$-matrices for the standard landscape}

The computation of the $r$-matrix $r_\cA(z)$ for the three
representative families.

\begin{computation}[Heisenberg $r$-matrix;
\ClaimStatusProvedHere]
% label removed: comp:thqg-V-heisenberg-r
\index{Heisenberg algebra!gravitational r-matrix@gravitational $r$-matrix}
\index{r-matrix@$r$-matrix!Heisenberg}
For the rank-$N$ Heisenberg algebra
$\mathcal{H}^N_\kappa$ with generators $J^i(z)$,
$i = 1,\dotsc,N$, and OPE
$J^i(z)\,J^j(w) \sim \kappa\,\delta^{ij}/(z-w)^2$, the collision
residue is
\begin{equation}% label removed: eq:thqg-V-heis-collision
\Rescoll_{0,2}(\Theta_{\mathcal{H}})
 (J^i, J^j)(z)
\;=\;
\Res_{w=z}\,
\frac{\kappa\,\delta^{ij}}{(w-z)^2}\,d\log(w-z)
\;=\;
0,
\end{equation}
since the second-order pole contributes no residue in $d\log$;
equivalently, the Heisenberg chiral Lie bracket $J^i_{(0)}J^j = 0$
is abelian.
The classical $r$-matrix (leading coefficient of the spectral
$R$-matrix $R_{\mathcal{H}}(z) = \exp(\kappa\hbar/z)\cdot\tau$
at order $\hbar$) is the Casimir
\begin{equation}% label removed: eq:thqg-V-heis-r
r_{\mathcal{H}}(z)
\;=\;
\frac{\Omega_{\mathcal{H}}}{z}
\;=\;
\frac{1}{\kappa z}
\sum_{i=1}^N
J^i \otimes J^i,
\end{equation}
where $\Omega_{\mathcal{H}} = \kappa^{-1}\sum_i J^i\otimes J^i$ is the
Casimir of the metric $\kappa\delta_{ij}$. This Casimir formula
comes from the double-pole Laplace-kernel contribution.
% note: r_H(z) derived from Casimir structure on A^!, not collision residue.

This is the \emph{scalar} $r$-matrix: it lies in the one-dimensional
space $\C \cdot \Omega_{\mathcal{H}} \subset
(\mathcal{H}^!)^{\otimes 2}$. The shadow obstruction tower terminates at
degree~$2$ (Corollary~\ref{V1-cor:heisenberg-postnikov-termination}),
so $r_{\mathcal{H}}$ is the \emph{complete} gravitational datum:
no ternary, quartic, or higher corrections exist.

The Yang--Baxter equation reduces to the triviality
\begin{equation}% label removed: eq:thqg-V-heis-ybe
\bigl[\Omega_{12}/z_{12},\;\Omega_{13}/z_{13}\bigr]
+
\bigl[\Omega_{12}/z_{12},\;\Omega_{23}/z_{23}\bigr]
+
\bigl[\Omega_{13}/z_{13},\;\Omega_{23}/z_{23}\bigr]
= 0
\end{equation}
with $\Omega_{ij} = \kappa^{-1}\sum_k J^k_i \otimes J^k_j$, which
holds because the Heisenberg is abelian: all brackets
$[\Omega_{ij}, \Omega_{kl}] = 0$.
\end{computation}

\begin{computation}[Affine Kac--Moody $r$-matrix;
\ClaimStatusProvedHere]
% label removed: comp:thqg-V-affine-r
\index{Kac--Moody!gravitational r-matrix@gravitational $r$-matrix}
\index{r-matrix@$r$-matrix!affine Kac--Moody}
For $\cA = \widehat{\fg}_k$ at non-critical level
$k \neq -h^\vee$, with generators $J^a(z)$ ($a = 1,\dotsc,
\dim\fg$) and OPE
\[
J^a(z)\,J^b(w)
\;\sim\;
\frac{k\,(a,b)}{(z-w)^2}
\;+\;
\frac{f^{ab}_{\phantom{ab}c}\,J^c(w)}{z-w},
\]
the collision residue is
\begin{equation}% label removed: eq:thqg-V-affine-collision
\Rescoll_{0,2}(\Theta_{\widehat{\fg}_k})
 (J^a, J^b)(z)
\;=\;
f^{ab}_{\phantom{ab}c}\,J^c(z),
\end{equation}
the Lie bracket of the underlying finite-dimensional algebra~$\fg$.
The $r$-matrix is the classical Casimir $r$-matrix:
\begin{equation}% label removed: eq:thqg-V-affine-r
r_{\widehat{\fg}_k}(z)
\;=\;
\frac{\Omega_\fg}{z}
\;=\;
\frac{1}{z}\,
\sum_{a=1}^{\dim\fg}
\frac{e_a \otimes e^a}{k + h^\vee},
\end{equation}
where $\{e_a\}$ is an orthonormal basis of $\fg$ with respect to the
Killing form scaled by $1/(2h^\vee)$, and $\Omega_\fg$ is the
quadratic Casimir.

The normalization by $(k + h^\vee)^{-1}$ arises from the propagator:
the bar-complex propagator is $\eta(z,w) = d\log(z - w)$ with the
OPE coefficient $k\,(a,b)$ in the numerator, giving
$r_{ab}(z) = \delta_{ab}/(k + h^\vee)\,z$ after inverting the
level-shifted Killing form.

For $\fg = \mathfrak{sl}_2$:
\begin{equation}% label removed: eq:thqg-V-sl2-r
r_{\widehat{\mathfrak{sl}}_2}(z)
\;=\;
\frac{1}{(k+2)\,z}\,
\bigl(e \otimes f + f \otimes e + \tfrac{1}{2}\,h \otimes h\bigr).
\end{equation}
The Casimir $\Omega_{\mathfrak{sl}_2} =
e \otimes f + f \otimes e + \frac{1}{2}\,h \otimes h$
generates the $\mathfrak{sl}_2$ invariant in
$\mathfrak{sl}_2^{\otimes 2}$.
\end{computation}

\begin{computation}[Virasoro $r$-matrix;
\ClaimStatusProvedHere]
% label removed: comp:thqg-V-virasoro-r
\index{Virasoro!gravitational r-matrix@gravitational $r$-matrix}
\index{r-matrix@$r$-matrix!Virasoro}
For $\cA = \mathrm{Vir}_c$ with generator $T(z)$ and OPE
\[
T(z)\,T(w)
\;\sim\;
\frac{c/2}{(z-w)^4}
\;+\;
\frac{2T(w)}{(z-w)^2}
\;+\;
\frac{\partial T(w)}{z-w},
\]
the collision residue is
\begin{equation}% label removed: eq:thqg-V-vir-collision
\Rescoll_{0,2}(\Theta_{\mathrm{Vir}})
 (T, T)(z)
\;=\;
\partial T(z),
\end{equation}
the regular term. The $r$-matrix requires the Koszul dual
$\mathrm{Vir}_c^! = \mathrm{Vir}_{26-c}$
(Remark~\ref{V1-rem:virasoro-resonance-model}). On the Koszul-dual
generators $\{L_n\}_{n \in \Z}$ of $\mathrm{Vir}_{26-c}$, the
binary $r$-matrix is
(Construction~\ref{constr:thqg-V-vir-r-explicit}):
\begin{equation}% label removed: eq:thqg-V-vir-r
\rgrav_{\mathrm{Vir}}(z)
\;=\;
\sum_{m \geq 0}
\frac{\Omega_m}{z^{m+1}},
\qquad
\Omega_m
\;=\;
\sum_{n \in \Z}
\tbinom{n+m-1}{m}\,
L_{-n-m} \otimes L_n,
\end{equation}
where $\Omega_0 = \sum_{n} L_{-n} \otimes L_n$ is the Virasoro
Casimir. The collision residue
$r^{\mathrm{coll}}(z) = (c/2)/z^3 + 2T/z$
(derived below) has poles at orders $z^{-3}$ and~$z^{-1}$,
one lower than the OPE poles at $z^{-4}$ and~$z^{-2}$ (the $d\log$ bar kernel absorbs one pole order).

\emph{Derivation from the collision residue.}
The Virasoro OPE
$T(z)\,T(w) \sim \frac{c/2}{(z-w)^4} + \frac{2T(w)}{(z-w)^2}
+ \frac{\partial T(w)}{z-w}$
has products $T_{(k)}T$ at each pole order $k = 0,1,2,3$:
$T_{(0)}T = \partial T$,
$T_{(1)}T = 2T$,
$T_{(2)}T = 0$,
$T_{(3)}T = c/2$.
The degree-$1$ bar differential extracts the simple-pole
($k = 0$) OPE term: $(T_{(0)}T)(z) = \partial T(z)$.
The Laplace kernel (OPE generating function)
$r^L(z) = \sum_k (T_{(k)}T)\,z^{-k-1}$
has poles at orders $k = 0,1,2,3$, giving
$r^L(z) = \frac{c/2}{z^4}\,\mathbf{1} \otimes \mathbf{1}
+ \frac{2T}{z^2} + \frac{\partial T}{z}$.
The bar-theoretic collision residue $r^{\mathrm{coll}}(z)$
has pole orders one lower (the $d\log$ kernel absorbs a power
of~$z$; Volume~I, Remark~\ref{V1-rem:three-r-matrices}):
$r^{\mathrm{coll}}(z) = (c/2)/z^3 + 2T/z$.
The coefficient of the leading pole is $c/2$, \emph{not}
$c/12$: the $c/12$ normalization appears in the Virasoro
commutation relations $[L_m, L_n]$, while the OPE carries
$c/2$ in the fourth-order pole.

\emph{Critical difference from affine KM.}
The Virasoro OPE kernel has poles of order $>2$ in $z^{-1}$: the
fourth-order singularity in the $TT$-OPE produces a cubic pole
$z^{-3}$ in the collision residue. This is the source of the
infinite shadow obstruction tower: the higher-order poles generate an infinite
sequence of non-trivial higher-degree corrections
$r_\cA^{(k)}(z_1,\dotsc,z_k)$ for $k \geq 3$.
The first nonlinear correction is the quartic contact invariant
$Q^{\mathrm{contact}}_{\mathrm{Vir}} = 10/[c(5c+22)]$.
\end{computation}

\begin{remark}[The $r$-matrix landscape]
% label removed: rem:thqg-V-r-matrix-landscape
\index{r-matrix@$r$-matrix!landscape}
The collision-residue extraction provides a uniform construction of
$r$-matrices across the entire standard landscape:
\begin{center}
\renewcommand{\arraystretch}{1.2}
\begin{tabular}{@{}lllc@{}}
\toprule
\textbf{Family} & \textbf{$r_\cA(z)$}
 & \textbf{Leading pole} & \textbf{Higher poles} \\
\midrule
$\mathcal{H}^N_\kappa$ & $\Omega_{\mathcal{H}}/(\kappa z)$
 & simple & none \\
$\widehat{\fg}_k$ & $\Omega_\fg/((k{+}h^\vee)z)$
 & simple & none \\
$\beta\gamma$ & $\Omega_{\beta\gamma}/z$
 & simple & quartic contact \\
$\mathrm{Vir}_c$ & $r_{\mathrm{Vir}}(z)$
 & simple & all orders \\
$\mathcal{W}_N$ & $r_{\mathcal{W}_N}(z)$
 & simple & all orders \\
$Y(\fg)$ & $R_{\mathrm{Yang}}(z)/z$
 & simple & all orders via RTT \\
\bottomrule
\end{tabular}
\end{center}
The pattern is: for class~$G$ and~$L$ algebras, the $r$-matrix has
only a simple pole. For class~$C$ and~$M$, higher-order poles
appear, reflecting the nonlinear OPE structure. The Yangian
$r$-matrix has higher poles of a different origin:
the RTT relation produces a spectral-parameter polynomial
$R(u) = 1 - \hbar P/u + O(u^{-2})$ with the higher-order terms
determined by the quantum group structure.
\end{remark}


\subsubsection{The propagator and its gravitational interpretation}

\begin{proposition}[Collision propagator and its gravitational reading;
\ClaimStatusProvedHere]
% label removed: prop:thqg-V-gravitational-propagator
\index{propagator!gravitational}
The logarithmic form $\eta_{ij} = d\log(z_i - z_j)$ appearing in
the collision residue is the holomorphic boundary-to-boundary
propagator kernel: the Green's function of the
$\bar\partial$-operator on the punctured plane. In the
gravitational reading of the HT package, this kernel is read as
the propagator of a graviton connecting insertions at $z_i$
and $z_j$.

In the holomorphic-topological (HT) interpretation, and hence in the
gravitational reading of that package:
\begin{enumerate}[label=\textup{(\roman*)}]
\item $\eta_{ij}$ is the boundary-to-boundary propagator of
 the 3d HT theory restricted to the 2d boundary.
\item The collision residue
 $\Rescoll_{0,2}(\Theta_\cA)(a,b) =
 \Res_{z_i = z_j}[a(z_i)\,b(z_j)\,\eta_{ij}]$
 is read as contributing the gravitational amplitude for the
 exchange of a single graviton between the boundary operators
 $a$ and~$b$.
\item Higher collision residues
 $\Rescoll_{0,k}(\Theta_\cA)$ for $k \geq 3$ are read as
 higher-exchange amplitudes organized by the shadow obstruction tower.
\end{enumerate}
\end{proposition}

\begin{proof}
The identification of $\eta_{ij}$ as the $\bar\partial$-propagator
is standard: on $\C$, the unique holomorphic $1$-form with a simple
pole at $w = z$ and no other singularity is
$dw/(w - z) = d\log(w - z)$. In the HT context, this is the
boundary restriction of the bulk propagator
$P(z,w;\bar{z},\bar{w})$ of the holomorphically twisted theory
(Costello--Li~\cite{CL20}): the holomorphic twist kills the
$\bar{z}$-dependence, leaving the purely holomorphic propagator
$\eta(z,w)$.

The collision-residue interpretation follows from
Construction~\ref{V1-constr:thqg-V-binary-extraction}: the residue
extracts the singular part of the OPE, which in the HT interpretation
is read as the single-graviton exchange channel. Higher collision
residues extract nested singularities, which in the same gravitational
reading organize the corresponding higher-exchange channels.
\end{proof}


% ======================================================================
%
% SUBSECTION 3: TERNARY COLLISION AND YANG--BAXTER
%
% ======================================================================

\subsection{Ternary collision and the Yang--Baxter equation}
% label removed: subsec:thqg-V-ternary-collision
\index{Yang--Baxter equation!from ternary collision}
\index{Arnold relation!Yang--Baxter}

The classical Yang--Baxter equation (CYBE) emerges from the
collision filtration at depth~$2$: the MC equation restricted to
the genus-zero four-point sector $\overline{\mathcal{M}}_{0,4}$,
projected to the depth-$2$ associated graded, yields the CYBE
as a consequence of the Arnold relation. This subsection gives
the full proof and develops the gravitational interpretation.


\subsubsection{The four-point moduli space}

\begin{lemma}[Structure of $\overline{\mathcal{M}}_{0,4}$;
\ClaimStatusProvedElsewhere]
% label removed: lem:thqg-V-m04-structure
\index{M04@$\overline{\mathcal{M}}_{0,4}$!structure}
The moduli space $\overline{\mathcal{M}}_{0,4} \cong \mathbb{P}^1$
has three boundary points
$D_s, D_t, D_u$ corresponding to the three channels:
\begin{align}
D_s &\;:\; (z_1, z_2) \to z_{12},\;(z_3, z_4) \to z_{34}
 & \text{($s$-channel: $(12|34)$)},
 % label removed: eq:thqg-V-s-channel \\
D_t &\;:\; (z_1, z_3) \to z_{13},\;(z_2, z_4) \to z_{24}
 & \text{($t$-channel: $(13|24)$)},
 % label removed: eq:thqg-V-t-channel \\
D_u &\;:\; (z_1, z_4) \to z_{14},\;(z_2, z_3) \to z_{23}
 & \text{($u$-channel: $(14|23)$)}.
 % label removed: eq:thqg-V-u-channel
\end{align}
The cross-ratio $\lambda = (z_1 - z_2)(z_3 - z_4)/
(z_1 - z_3)(z_2 - z_4)$ identifies $\overline{\mathcal{M}}_{0,4}$
with $\mathbb{P}^1$, with $D_s = \{0\}$, $D_t = \{1\}$,
$D_u = \{\infty\}$. The open stratum $\mathcal{M}_{0,4}$
is $\mathbb{P}^1 \setminus \{0,1,\infty\}$.
\end{lemma}

\begin{proof}
Standard (see e.g.~\cite{Kn83}). Three automorphisms of
$\mathbb{P}^1$ fix three of the four marked points; the residual
freedom parametrizes the fourth, giving a one-parameter family
$\lambda \in \mathbb{P}^1$. The boundary points are the three
degenerations of the four-punctured sphere into two three-punctured
spheres.
\end{proof}

\begin{construction}[The MC equation at four points]
% label removed: constr:thqg-V-mc-four-points
\index{Maurer--Cartan equation!four-point}
The MC equation
$D_\cA\Theta_\cA + \tfrac{1}{2}[\Theta_\cA,\Theta_\cA] = 0$,
restricted to the genus-zero four-point sector, decomposes by
collision depth. Write
$\Theta_{0,4} := \Theta_\cA|_{g=0,n=4}$ and
$\Theta_{0,3} := \Theta_\cA|_{g=0,n=3}$.

\emph{Depth~$0$} (the open stratum $\mathcal{M}_{0,4}$).
The MC equation gives
$d_{\mathrm{int}}(\Theta_{0,4}|_{\mathrm{open}}) = 0$:
the four-point datum on the open stratum is a cocycle for the
internal $A_\infty$ differential.

\emph{Depth~$1$} (the three boundary divisors).
The boundary contributions are:
\begin{equation}% label removed: eq:thqg-V-depth-1-mc
\bigl(D_\cA\Theta_{0,4}\bigr)\big|_{D_s}
\;=\;
[\Theta_{0,3},\,\Theta_{0,3}]_s
\;=\;
[\Omega_{12},\,\Omega_{34}]/z_{12}\,z_{34},
\end{equation}
and similarly for $D_t$ and~$D_u$, where the bracket is the
Lie bracket on $\gAmod$ evaluated on the channel composition.

\emph{Depth~$2$} (the triple collision at the intersection of two
boundary divisors).
The total boundary contribution from all three channels is:
\begin{equation}% label removed: eq:thqg-V-depth-2-mc
\sum_{\text{channels}}
[\Theta_{0,3},\,\Theta_{0,3}]_{\text{channel}}
\;=\;
[r_{12},\,r_{13}] + [r_{12},\,r_{23}] + [r_{13},\,r_{23}]
\;=\; 0,
\end{equation}
where $r_{ij} = r_\cA(z_{ij})$. This is the CYBE.
\end{construction}


\subsubsection{Full proof of the CYBE from Arnold}

\begin{theorem}[CYBE from the Arnold relation via MC;
\ClaimStatusProvedHere]
% label removed: thm:thqg-V-cybe-from-arnold
\index{classical Yang--Baxter equation!proof from Arnold}
\index{Arnold relation!proof of CYBE}
Let $\cA$ be a modular Koszul chiral algebra with $r$-matrix
$r_\cA(z) = \Rescoll_{0,2}(\Theta_\cA)$
\textup{(}Theorem~\textup{\ref{V1-thm:thqg-V-collision-twisting})}.
Then $r_\cA(z)$ satisfies the classical Yang--Baxter equation:
\begin{equation}% label removed: eq:thqg-V-cybe
\bigl[r_{12}(z_{12}),\;r_{13}(z_{13})\bigr]
\;+\;
\bigl[r_{12}(z_{12}),\;r_{23}(z_{23})\bigr]
\;+\;
\bigl[r_{13}(z_{13}),\;r_{23}(z_{23})\bigr]
\;=\; 0,
\end{equation}
where $z_{ij} = z_i - z_j$ and $r_{ij}(z_{ij})$ acts on the
$i$-th and $j$-th tensor factors of
$(\cA^!)^{\otimes 3}$.
\end{theorem}

\begin{proof}
The proof proceeds through three steps.

\emph{Step~1: Setup.}
Consider the MC equation on the genus-zero convolution algebra at
degree~$3$ (four marked points: three input + one output at infinity).
The boundary of $\overline{\mathcal{M}}_{0,4}$ has three divisors
$D_s, D_t, D_u$ (Lemma~\ref{V1-lem:thqg-V-m04-structure}).
The MC equation
\begin{equation}% label removed: eq:thqg-V-mc-degree3
D_\cA\Theta_{0,3}
\;+\;
\tfrac{1}{2}\bigl[\Theta_{0,2},\;\Theta_{0,2}\bigr]_{0,3}
\;=\; 0,
\end{equation}
where $[\cdot,\cdot]_{0,3}$ denotes the bracket projected to the
$3$-input component via operadic composition, splits by the collision
filtration into contributions at each boundary stratum.

\emph{Step~2: Arnold computation.}
On $\Conf_3(\C)$, the Arnold relation
(Proposition~\ref{V1-prop:thqg-V-arnold-cohomology}) gives
\begin{equation}% label removed: eq:thqg-V-arnold-3pt
\eta_{12} \wedge \eta_{23}
\;+\;
\eta_{23} \wedge \eta_{31}
\;+\;
\eta_{31} \wedge \eta_{12}
\;=\; 0,
\end{equation}
where $\eta_{ij} = d\log(z_i - z_j)$.
Tensoring with the Koszul-dual Casimir
$\Omega \in (\cA^!)^{\otimes 2}$ on each factor and evaluating
the bracket on $(\cA^!)^{\otimes 3}$, we obtain:
\begin{align}
&\sum_{a,b} \bigl(\Omega_{12}^{ab}\,\Omega_{23}^{cd}
 - \Omega_{23}^{ab}\,\Omega_{12}^{cd}\bigr)\,
 \frac{1}{z_{12}\,z_{23}}
\notag \\
&\qquad +\;
\sum_{a,b} \bigl(\Omega_{23}^{ab}\,\Omega_{31}^{cd}
 - \Omega_{31}^{ab}\,\Omega_{23}^{cd}\bigr)\,
 \frac{1}{z_{23}\,z_{31}}
\notag \\
&\qquad +\;
\sum_{a,b} \bigl(\Omega_{31}^{ab}\,\Omega_{12}^{cd}
 - \Omega_{12}^{ab}\,\Omega_{31}^{cd}\bigr)\,
 \frac{1}{z_{31}\,z_{12}}
\;=\; 0.
% label removed: eq:thqg-V-arnold-tensored
\end{align}

\emph{Step~3: Identification with CYBE.}
Each term in~\eqref{V1-eq:thqg-V-arnold-tensored} is a Lie bracket of
$r$-matrix components:
\begin{align}
\frac{\Omega_{12}\,\Omega_{23}
 - \Omega_{23}\,\Omega_{12}}{z_{12}\,z_{23}}
&\;=\;
\bigl[r_{12}(z_{12}),\;r_{23}(z_{23})\bigr],
% label removed: eq:thqg-V-cybe-term1 \\
\frac{\Omega_{23}\,\Omega_{31}
 - \Omega_{31}\,\Omega_{23}}{z_{23}\,z_{31}}
&\;=\;
\bigl[r_{23}(z_{23}),\;r_{31}(z_{31})\bigr],
% label removed: eq:thqg-V-cybe-term2 \\
\frac{\Omega_{31}\,\Omega_{12}
 - \Omega_{12}\,\Omega_{31}}{z_{31}\,z_{12}}
&\;=\;
\bigl[r_{31}(z_{31}),\;r_{12}(z_{12})\bigr],
% label removed: eq:thqg-V-cybe-term3
\end{align}
where $r_{ij}(z) = \Omega_{ij}/z$ is the leading-order $r$-matrix.
(For algebras with higher-order poles in the $r$-matrix, such as
$\mathrm{Vir}_c$, additional terms appear at higher collision depth;
the leading-order CYBE is the universal statement valid for all
modular Koszul algebras.)

The spectral parameter $z_{ij} = z_i - z_j$ satisfies
$z_{ji} = -z_{ij}$. Skew-symmetry of~$r$ gives
$r(z_{ji}) = r(-z_{ij}) = -P \cdot r(z_{ij}) \cdot P$, where
$P$ is the permutation operator exchanging the two tensor factors.
In particular $z_{31} = -z_{13}$, and
$r_{31}(z_{31}) = r_{31}(-z_{13}) = -r_{13}(z_{13})$
by the skew-symmetry property
(Theorem~\ref{V1-thm:thqg-V-collision-twisting}(ii)).
With this identification, the sum
\eqref{V1-eq:thqg-V-cybe-term1}--\eqref{V1-eq:thqg-V-cybe-term3}
gives exactly the CYBE~\eqref{V1-eq:thqg-V-cybe}.
\end{proof}

\begin{remark}[Gravitational interpretation of the CYBE]
% label removed: rem:thqg-V-cybe-gravitational
\index{classical Yang--Baxter equation!gravitational interpretation}
The CYBE has a gravitational reading.
Consider three line operators at positions $z_1, z_2, z_3$ on the
boundary of the HT theory. In that reading, the binary kernel
$r_{ij}(z_{ij})$ is read as the pairwise graviton-exchange
contribution between the $i$-th and $j$-th insertions. The
CYBE~\eqref{V1-eq:thqg-V-cybe} then states that the three possible
pairwise binary contributions are \emph{consistent}: there is no
tree-level inconsistency in assembling the ternary scattering data
from this binary kernel.

In the celestial/gravitational reading, this same compatibility is
read as a Ward-identity manifestation of the leading celestial soft
factor (Theorem~\ref{thm:thqg-VI-leading-soft}), rather than as an
independent theorem-level statement about tree-level graviton
scattering.
\end{remark}


\subsubsection{The infinitesimal braid relation}

\begin{corollary}[Infinitesimal braid relation;
\ClaimStatusProvedHere]
% label removed: cor:thqg-V-ibr
\index{infinitesimal braid relation}
For any modular Koszul chiral algebra~$\cA$ with Casimir
$\Omega_{ij}$ acting on the $i$-th and $j$-th tensor factors of
$(\cA^!)^{\otimes n}$, the infinitesimal braid relation
\textup{(}IBR\textup{)} holds:
\begin{equation}% label removed: eq:thqg-V-ibr
[\Omega_{ij},\;\Omega_{ik} + \Omega_{jk}] \;=\; 0
\qquad
\text{for all distinct } i,j,k \in \{1,\dotsc,n\}.
\end{equation}
This is a formal consequence of the CYBE at the Casimir level:
setting $z_{12} = z_{13} = z$ (equal spectral parameters) in
\eqref{V1-eq:thqg-V-cybe} and extracting the $z^{-2}$ coefficient
gives the IBR\@.
\end{corollary}

\begin{proof}
The CYBE~\eqref{V1-eq:thqg-V-cybe} with
$r_{ij}(z) = \Omega_{ij}/z$ becomes
\[
\frac{[\Omega_{12},\Omega_{13}]}{z_{12}\,z_{13}}
+
\frac{[\Omega_{12},\Omega_{23}]}{z_{12}\,z_{23}}
+
\frac{[\Omega_{13},\Omega_{23}]}{z_{13}\,z_{23}}
= 0.
\]
Multiply through by $z_{12}\,z_{13}\,z_{23}$ and set
$z_{23} = z_{13} - z_{12}$:
\[
[\Omega_{12},\Omega_{13}]\,z_{23}
+ [\Omega_{12},\Omega_{23}]\,z_{13}
+ [\Omega_{13},\Omega_{23}]\,z_{12}
= 0.
\]
The coefficient of $z_{12}^0\,z_{13}^1$ (equivalently,
taking $z_{12} \to 0$) gives
$[\Omega_{12},\,\Omega_{13} + \Omega_{23}] = 0$,
which is the IBR\@.
\end{proof}

\begin{proposition}[$A_\infty$-enhancement of the CYBE;
\ClaimStatusProvedHere]
% label removed: prop:thqg-V-ainfty-ybe
\index{Yang--Baxter equation!A-infinity enhancement@$A_\infty$ enhancement}
For chiral algebras of shadow class~$C$ or~$M$, the classical
CYBE receives $A_\infty$ corrections from higher collision depths.
The corrected equation at depth~$d$ is
\begin{equation}% label removed: eq:thqg-V-ainfty-ybe
\sum_{k=0}^{d}
\sum_{\substack{i<j<\ell \\
 d_1 + d_2 + d_3 = d \\
 d_1, d_2, d_3 \geq 0}}
\bigl[
 r^{(d_1)}_{ij}(z_{ij}),\;
 r^{(d_2)}_{j\ell}(z_{j\ell})
\bigr]^{(d_3)}
\;=\; 0,
\end{equation}
where $r^{(d)}$ denotes the $r$-matrix at collision depth~$d$ and
$[\cdot,\cdot]^{(d)}$ is the $A_\infty$ bracket at depth~$d$.

For class~$G$ and~$L$: $r^{(d)} = 0$ for $d \geq 1$, and the
$A_\infty$ CYBE reduces to the classical CYBE\@.
For class~$C$: $r^{(1)} \neq 0$ at the quartic level, and the
first correction is the quartic contact term.
For class~$M$: all $r^{(d)}$ are nonzero, and the $A_\infty$ CYBE
is an infinite tower of higher-depth algebraic identities. In the
gravitational reading, this tower is read as organizing the full
nonlinear scattering hierarchy.
\end{proposition}

\begin{proof}
The $A_\infty$ enhancement arises from the full MC equation at all
collision depths, not just depth~$2$. At depth $d$, the MC equation
restricted to $\overline{\mathcal{M}}_{0,4}$ involves compositions
of operations at various intermediate depths summing to~$d$. The
resulting identity is the $A_\infty$ CYBE at total depth~$d$.

For class~$G$: the shadow obstruction tower terminates at degree~$2$
(Corollary~\ref{V1-cor:heisenberg-postnikov-termination}), so all
higher-depth $r$-matrices vanish.
For class~$L$: the shadow obstruction tower terminates at degree~$3$
(Corollary~\ref{V1-cor:affine-postnikov-termination}), so $r^{(d)} = 0$
for $d \geq 2$ and the first potential correction (at degree~$3$,
depth~$1$) is killed by the Jacobi identity.
For class~$C$: the shadow obstruction tower terminates at degree~$4$, and the
quartic contact invariant $\mathfrak{Q}$ provides the first nonzero
$A_\infty$ correction at total depth~$2$.
For class~$M$: the tower is infinite
(Theorem~\ref{V1-thm:nms-virasoro-quintic-forced}).
\end{proof}


\subsubsection{Higher-degree CYBE}

\begin{definition}[$n$-point Yang--Baxter identity]
% label removed: def:thqg-V-n-point-ybe
\index{Yang--Baxter equation!n-point@$n$-point}
For $n \geq 3$ points on $\overline{\mathcal{M}}_{0,n+1}$, the
\emph{$n$-point Yang--Baxter identity} is the MC equation projected
to collision depth $n - 1$:
\begin{equation}% label removed: eq:thqg-V-n-point-ybe
\sum_{\substack{T \in \mathcal{T}_n \\
 |E_{\mathrm{int}}(T)| = n-2}}
\frac{1}{|\Aut(T)|}
\prod_{e \in E_{\mathrm{int}}(T)}
r_{e_-,e_+}(z_{e_-,e_+})
\;=\; 0,
\end{equation}
where $\mathcal{T}_n$ is the set of binary rooted trees with $n$
leaves, $e_-$ and $e_+$ are the endpoints of the internal edge~$e$,
and the product is over internal edges.
At $n = 3$: this is the CYBE (three binary trees, one internal edge
each).
At $n = 4$: this gives the \emph{quartic YBE}, involving the
five binary trees with two internal edges.
\end{definition}

\begin{proposition}[$n$-point YBE from boundary of
$\overline{\mathcal{M}}_{0,n+1}$; \ClaimStatusProvedHere]
% label removed: prop:thqg-V-n-point-ybe-proof
\index{Yang--Baxter equation!n-point proof@$n$-point proof}
The $n$-point YBE is a consequence of the vanishing
\begin{equation}% label removed: eq:thqg-V-boundary-sum
\sum_{D \in \partial\overline{\mathcal{M}}_{0,n+1}}
\int_D
\prod_{e}
r_{e_-,e_+}(z_{e_-,e_+})
\;=\; 0,
\end{equation}
which follows from $\partial^2 = 0$ on the chain level
\textup{(}the boundary of a boundary is empty\textup{)}, together
with the MC equation at all intermediate degrees.
\end{proposition}

\begin{proof}
The MC equation $D\Theta + \tfrac{1}{2}[\Theta,\Theta] = 0$ at
degree~$n$, projected to the maximal collision depth, gives a sum over
all maximally degenerate boundary strata of
$\overline{\mathcal{M}}_{0,n+1}$. These strata are indexed by binary
trees with $n$ leaves. The identity $\partial^2 = 0$ on
$C_*(\overline{\mathcal{M}}_{0,n+1})$ ensures that the total
boundary contribution vanishes. Each boundary stratum contributes a
product of $r$-matrices along internal edges, giving
\eqref{V1-eq:thqg-V-boundary-sum}.

The identification of the boundary sum with the combinatorial formula
\eqref{V1-eq:thqg-V-n-point-ybe} follows from the factorization of
boundary strata: each boundary divisor of $\overline{\mathcal{M}}_{0,n+1}$
is a product of lower-dimensional moduli spaces, and the $r$-matrix
contributions factor accordingly.
\end{proof}


% ======================================================================
%
% SUBSECTION 4: THE MODULAR DG-SHIFTED YANGIAN
%
% ======================================================================

\subsection{The modular dg-shifted Yangian}
% label removed: subsec:thqg-V-modular-yangian
\index{Yangian!modular dg-shifted|textbf}
\index{Yangian!modular higher-genus completion}

The genus-zero $r$-matrix is the leading term of a richer object:
the \emph{modular dg-shifted Yangian} $\Ymod_\cA$, an ambient
pronilpotent dg Lie algebra that organizes the higher-genus binary
collision data. The formal series of genus-refined binary
projections is defined on this surface; identifying it with a genuine
modular kernel and proving the stable-graph Yang--Baxter equation are
additional frontier inputs. The construction assembles the
collision-filtration data of
\S\ref{V1-subsec:thqg-V-collision-filtration} into a single algebraic
object.

\subsubsection{Pronilpotent completion}

\begin{definition}[Modular dg-shifted Yangian for $\cA$;
\ClaimStatusProvedHere]
% label removed: def:thqg-V-modular-yangian
\index{Yangian!modular dg-shifted!definition}
Let $\cA$ be a modular Koszul chiral algebra.
The \emph{modular dg-shifted Yangian} of~$\cA$ is the pronilpotent
dg Lie algebra
\begin{equation}% label removed: eq:thqg-V-modular-yangian
\Ymod_\cA
\;:=\;
\varprojlim_N\;
\mathcal{L}_\cA^{\mathrm{mod}} \big/ F^{N+1}\!
 \mathcal{L}_\cA^{\mathrm{mod}},
\end{equation}
where
\begin{equation}% label removed: eq:thqg-V-convolution
\mathcal{L}_\cA^{\mathrm{mod}}
\;:=\;
\prod_{2g-2+n > 0}
\Hom\!\bigl(
 C_\bullet(\overline{\mathcal{M}}_{g,n+1})
 \otimes (\cA^!)^{\widehat{\otimes}\,n},\;
 \cA^!
\bigr)[1]
\end{equation}
is the modular convolution dg Lie algebra
\textup{(}Definition~\textup{\ref*{V1-def:modular-convolution-dg-lie})}
with the filtration
$F^N \mathcal{L}_\cA^{\mathrm{mod}}
:= \prod_{2g-2+n \geq N}(\cdots)$.

The genus-zero truncation is the
(non-modular) dg-shifted Yangian:
\begin{equation}% label removed: eq:thqg-V-genus-zero-truncation
\Ydg_\cA
\;=\;
\Ymod_\cA\big|_{g=0}
\;=\;
\prod_{n \geq 2}
\Hom_{\Sigma_n}\!\bigl(
 C_\bullet(\overline{\mathcal{M}}_{0,n+1})
 \otimes (\cA^!)^{\widehat{\otimes}\,n},\;
 \cA^!
\bigr)[1].
\end{equation}
\end{definition}

\begin{remark}[Comparison with
Definition~\textup{\ref*{V1-def:modular-yangian-pro}}]
% label removed: rem:thqg-V-yangian-comparison
Definition~\ref{V1-def:thqg-V-modular-yangian} is a restatement
of the ambient pronilpotent dg Lie package in
Definition~\ref*{V1-def:modular-yangian-pro} in the gravitational
context. The key structural point is that the Lie filtration
$F^{N_1} \times F^{N_2} \to F^{N_1 + N_2}$
(Lemma~\ref{V1-lem:thqg-V-collision-lie-compatibility}) ensures the
pronilpotent completion is well-defined. What is \emph{not} proved
here is that the completion already comes equipped with a modular
kernel satisfying an MC equation; that stronger package remains on the
conjectural modular-Yangian surface.
\end{remark}

\begin{remark}[Bialgebra structure: no antipode, and the name ``Yangian'']
% label removed: rem:thqg-V-yangian-no-antipode
\index{Yangian!no antipode}
\index{antipode!absence in dg-shifted Yangian}
Two terminological points require clarification.

\smallskip\noindent
\emph{(i) Bialgebra, not Hopf algebra.}
The dg-shifted Yangian $\Ydg_\cA$ carries a coproduct
$\Delta_z$ (from the operadic composition on the
convolution algebra) and a counit $\varepsilon$
(from the augmentation), making it a dg bialgebra.
It does \emph{not} carry an antipode in general: the
Drinfeld Yangian $Y(\fg)$ has an antipode $S$ by the
finiteness of its Drinfeld generators, but the $A_\infty$
analogue for arbitrary modular Koszul algebras is not
known to admit one (cf.\
Remark~\ref{rem:dg-yangian-antipode}). For the
representation-theoretic applications in this chapter
(line-operator modules, spectral braiding, the RTT
relation), the bialgebra structure is sufficient: the
antipode would be needed for dualization of modules
and for constructing a ribbon structure, neither of
which is required here.

\smallskip\noindent
\emph{(ii) The name ``Yangian.''}
We use ``dg-shifted Yangian'' to emphasize the structural
parallel with the Drinfeld Yangian $Y(\fg)$: both are
pronilpotent completions of convolution algebras. For
$\cA = \widehat{\fg}_k$ at generic level,
$\Ydg_\cA$ recovers the standard Yangian $Y(\fg)$
(Computation~\ref{V1-comp:thqg-V-affine-yangian}). For
general~$\cA$, the object is not an ``affine Yangian''
in the sense of Guay--Nakajima--Wendlandt (which is a
specific algebra associated to $\widehat{\fg}$ with its
loop parameter); rather, it is the genus-zero binary
projection of the modular convolution algebra. The
modular extension $\Ymod_\cA$, which includes all genera,
has no classical precedent. The name ``gravitational
Yangian'' for the physical incarnation is adopted to
distinguish it from both the classical Yangian and
the affine Yangian.

\smallskip\noindent
\emph{(iii) Consistency with
Definition~\textup{\ref{def:dg_Yangian}}.}
The genus-zero truncation $\Ydg_\cA$ satisfies the axioms of
Definition~\ref{def:dg_Yangian}: the translation automorphism
$\tau_z$ is the spectral-parameter shift on the $r$-matrix,
the coproduct arises from operadic composition, and the YBE is
Theorem~\ref{V1-thm:thqg-V-cybe-from-arnold}. The modular
extension $\Ymod_\cA$ is a genuine enlargement: the higher-genus
components are not part of the abstract dg-shifted Yangian
structure of Definition~\ref{def:dg_Yangian} but are dictated
by the MC equation in the full modular convolution algebra.
\end{remark}

\begin{proposition}[Differential structure;
\ClaimStatusProvedHere]
% label removed: prop:thqg-V-yangian-differential
\index{Yangian!modular dg-shifted!differential}
The differential on $\Ymod_\cA$ has two summands:
\begin{equation}% label removed: eq:thqg-V-yangian-diff
d_{\Ymod}
\;=\;
d_{\mathrm{int}} + d_{\mathrm{mod}},
\end{equation}
where:
\begin{enumerate}[label=\textup{(\roman*)}]
\item $d_{\mathrm{int}}(\phi)$ acts by the $A_\infty$ differential
 on $\cA^!$ applied to the inputs and output of~$\phi$;
\item $d_{\mathrm{mod}}(\phi)$ acts by precomposition with the
 boundary operator
 $\partial\colon C_*(\overline{\mathcal{M}}_{g,n+1})
 \to \coprod C_*(\overline{\mathcal{M}}_{g_1,n_1+1}
 \times \overline{\mathcal{M}}_{g_2,n_2+1})$.
\end{enumerate}
The identity $d_{\Ymod}^2 = 0$ is equivalent to the $A_\infty$
structure equations on $\cA^!$ together with $\partial^2 = 0$ on the
chain complexes of moduli spaces.
\end{proposition}

\begin{proof}
$d_{\mathrm{int}}^2 = 0$ because $\cA^!$ is an $A_\infty$-algebra.
$d_{\mathrm{mod}}^2 = 0$ because $\partial^2 = 0$ on chains.
The cross-term
$d_{\mathrm{int}} \circ d_{\mathrm{mod}} +
d_{\mathrm{mod}} \circ d_{\mathrm{int}} = 0$
holds because the $A_\infty$ structure on $\cA^!$ is compatible
with the modular operad structure: the boundary operator on moduli
spaces composes with the $A_\infty$ multiplication maps. This
compatibility is the content of the bar complex being a
module over the Feynman transform of the modular operad
(Theorem~\ref{V1-thm:bar-modular-operad}).
\end{proof}


\subsubsection{The formal genus-refined binary shadow series}

\begin{theorem}[Formal genus-refined binary shadow series
 $R^{\mathrm{bin}}_\cA$;
\ClaimStatusProvedHere]
% label removed: thm:thqg-V-pro-mc-element
\index{Yangian!binary shadow series|textbf}
The universal MC element $\Theta_\cA \in \MC(\gAmod)$
\textup{(}Theorem~\textup{\ref*{V1-thm:mc2-bar-intrinsic})}
determines formal genus-refined binary coefficients
$r_{\cA,g}(z)$ and hence a formal series
\begin{equation}% label removed: eq:thqg-V-pro-mc
R^{\mathrm{bin}}_\cA(z;\hbar)
\;:=\;
\sum_{g \geq 0}
\hbar^{2g}\,r_{\cA,g}(z)
\;\in\;
\Ymod_\cA[\![\hbar^2]\!],
\end{equation}
where $r_{\cA,g}(z)$ is the genus-$g$ component of
$\Theta_\cA$ projected to the binary collision stratum.
\end{theorem}

\begin{proof}
The bar-intrinsic MC element $\Theta_\cA = D_\cA - d_0$
satisfies $D_\cA\Theta_\cA + \tfrac{1}{2}[\Theta_\cA,\Theta_\cA] = 0$
in $\gAmod$. Projecting $\Theta_\cA$ to the genus-$g$ binary
collision stratum produces the coefficient $r_{\cA,g}(z)$.
Because $\Ymod_\cA$ is the pronilpotent completion of the modular
convolution dg Lie algebra, these coefficients assemble into the
formal $\hbar$-adic series~\eqref{V1-eq:thqg-V-pro-mc}.
\end{proof}

\begin{corollary}[Genus-$0$ component = classical $r$-matrix;
\ClaimStatusProvedHere]
% label removed: cor:thqg-V-genus-zero-r
\index{r-matrix@$r$-matrix!genus-zero component}
The genus-$0$ component $r_{\cA,0}(z)$ is the classical $r$-matrix
of Theorem~\ref{V1-thm:thqg-V-collision-twisting}:
$r_{\cA,0}(z) = r_\cA(z)$.
\end{corollary}

\begin{proof}
Restrict $R^{\mathrm{bin}}_\cA$ to $\hbar = 0$.
\end{proof}

\begin{remark}[Heuristic genus-$1$ correction;
\ClaimStatusHeuristic]
% label removed: cor:thqg-V-genus-one-correction
\index{r-matrix@$r$-matrix!genus-one correction}
The first higher-genus binary coefficient is expected to involve the
curvature $\kappa(\cA)$, the genus-$1$ Hessian
$\delta H^{(1)}_\cA$, and quartic shadow corrections in the form
\begin{equation}% label removed: eq:thqg-V-genus-one-r
r_{\cA,1}(z)
\;=\;
\kappa(\cA) \cdot r_{\cA,0}(z)
\;+\;
\delta H^{(1)}_\cA \cdot \partial_z r_{\cA,0}(z)
\;+\;
(\text{quartic corrections}).
\end{equation}
For class~$G$ and~$L$ algebras, where the quartic corrections vanish
on the relevant binary-shadow surface, this predicts a controlled
deformation of the classical genus-zero $r$-matrix. What is not
proved here is that these coefficients assemble into a genuine
modular Maurer--Cartan kernel.
\end{remark}


\subsubsection{The stable-graph Yang--Baxter equation}

\begin{definition}[Stable-graph YBE]
% label removed: def:thqg-V-stable-graph-ybe
\index{Yang--Baxter equation!stable-graph}
On the conjectural modular-Yangian surface, the
\emph{stable-graph Yang--Baxter equation} (sgYBE) is the MC
equation for a modular kernel
$R^{\mathrm{mod}}_\cA(z;\hbar)\in\MC(\Ymod_\cA)$ expanded by genus:
at genus~$g$ and degree~$n$, it reads
\begin{equation}% label removed: eq:thqg-V-sgybe-expanded
d_{\mathrm{int}}\,r_{\cA,g}^{(n)}
\;+\;
\sum_{\Gamma \in \mathcal{G}^{\mathrm{st}}_{g,n}}
\frac{1}{|\Aut(\Gamma)|}
\prod_{v \in V(\Gamma)} r_{\cA,g_v}^{(n_v)}
\prod_{e \in E(\Gamma)} P_e
\;=\; 0,
\end{equation}
where $\mathcal{G}^{\mathrm{st}}_{g,n}$ is the set of stable
graphs of genus~$g$ with $n$ legs, the vertex labels $r_{\cA,g_v}^{(n_v)}$
are the pro-MC components, and $P_e$ is the edge propagator.

At genus~$0$, degree~$3$: the sgYBE would reduce to the classical
CYBE (three stable graphs, each a single edge connecting two
trivalent vertices).

At genus~$1$, degree~$2$: the sgYBE would give the
\emph{quantum Yang--Baxter equation} (qYBE), encoding the
one-loop correction to the $r$-matrix. The single stable graph
is a loop (one vertex, one self-edge), and the equation reads
\begin{equation}% label removed: eq:thqg-V-qybe
d_{\mathrm{int}}\,r_{\cA,1}^{(2)}
\;+\;
\tfrac{1}{2}\,\Tr_{12}\,
\bigl[r_{\cA,0}^{(2)},\,r_{\cA,0}^{(2)}\bigr]
\;+\;
\kappa(\cA) \cdot r_{\cA,0}^{(2)}
\;=\; 0,
\end{equation}
where $\Tr_{12}$ traces over the loop-contracted tensor factors and
the $\kappa$-term is the curvature contribution from the degenerating
cycle.
\end{definition}

\begin{remark}[The sgYBE and its gravitational-reading organization]
% label removed: rem:thqg-V-sgybe-scattering
\index{gravitational scattering!stable-graph YBE}
If a modular kernel satisfying the sgYBE is supplied, then it
organizes the corresponding candidate gravitational amplitudes by
topology. At genus~$0$: tree-level graviton exchange, controlled by
the classical $r$-matrix. At genus~$1$: one-loop graviton
self-energy and vertex corrections, controlled by $r_{\cA,1}$.
At genus~$g$: the $g$-loop gravitational amplitude, with the
combinatorics of stable graphs providing the Feynman-diagrammatic
structure. The present chapter does not prove the existence of such a
kernel; it only isolates the ambient modular dg Lie package in which
that higher-genus completion would live.
\end{remark}

\begin{proposition}[Formal \texorpdfstring{$\hbar$}{hbar}-adic
 convergence of the candidate binary shadow series;
\ClaimStatusProvedHere]
% label removed: prop:thqg-V-pro-mc-convergence
\index{Yangian!binary shadow series!convergence}
The formal series $R^{\mathrm{bin}}_\cA(z;\hbar)$ converges in
the $\hbar$-adic topology of $\Ymod_\cA$: at each genus~$g$,
the component $r_{\cA,g}(z)$ is a well-defined element of the
finite-dimensional $\Hom$ space, and the genus series
$\sum_g \hbar^{2g}\,r_{\cA,g}(z)$ is a formal power series in
$\C[\![\hbar^2]\!]$ with coefficients in the genus-zero Yangian.

For the Heisenberg algebra: $r_{\mathcal{H},g}(z) = 0$ for
$g \geq 1$ \textup{(}the genus tower terminates at genus~$0$
for the $r$-matrix\textup{)}.

For affine Kac--Moody at generic level: $r_{\widehat{\fg}_k,g}(z)$
is polynomial in $1/(k+h^\vee)$ at each genus, with degree
controlled by the graph-sum formula
\textup{(}Theorem~\textup{\ref{V1-thm:recursive-existence})}.
\end{proposition}

\begin{proof}
The convergence in the $\hbar$-adic topology is automatic from the
filtration: $r_{\cA,g}$ lies in $F^{2g-2+n}\,\Ymod_\cA$, and the
filtration is exhaustive and separated. The Heisenberg vanishing
for $g \geq 1$ follows from the Gaussian nature of the Heisenberg:
the genus-$g$ sewing amplitude is a determinant
(Theorem~\ref{thm:heisenberg-sewing}), which contributes to the
scalar shadow $\kappa$ but not to the $r$-matrix component.

For affine KM, the genus-$g$ component is computed by the graph-sum
formula with $|V_{\mathrm{int}}| = 2g - 2 + n$ vertices; each
vertex carries a factor of the Casimir $\Omega/(k+h^\vee)$, giving
the polynomial dependence on $1/(k+h^\vee)$.
\end{proof}


\subsubsection{Weight filtration and Yangian generators}

\begin{construction}[Yangian generators from the weight filtration;
\ClaimStatusProvedHere]
% label removed: constr:thqg-V-yangian-generators
\index{Yangian!generators from weight filtration}
The weight filtration on $\Ydg_\cA$ at genus zero
(the degree filtration) produces Yangian generators.
Let $\{e_\alpha\}$ be a basis of $\cA^!$ homogeneous with respect
to the conformal weight grading.

\emph{Level~$0$ generators.}
$T_{\alpha\beta}^{(0)} := \langle e_\alpha,\,r_\cA(z)\,e_\beta\rangle
= \Omega_{\alpha\beta}/z$.
These are the $\fg$-currents at degree~$2$.

\emph{Level~$r$ generators.}
$T_{\alpha\beta}^{(r)}$ is the coefficient of $z^{-r-1}$ in the
Laurent expansion of $r_\cA(z)$ at $z = 0$:
\begin{equation}% label removed: eq:thqg-V-yangian-generators
r_\cA(z)
\;=\;
\sum_{r \geq 0}
T^{(r)}\,z^{-r-1},
\qquad
T^{(r)} \in (\cA^!)^{\otimes 2}.
\end{equation}
The commutation relations of the $T^{(r)}$ are determined by the
CYBE and its higher-degree analogues. For $\cA = \widehat{\fg}_k$:
$T^{(0)} = \Omega/(k+h^\vee)$ and $T^{(r)} = 0$ for $r \geq 1$,
so the Yangian is generated by the level-$0$ Casimir alone.
For $\cA = Y(\fg)$: all $T^{(r)}$ are nonzero, recovering the
standard Yangian generators
(Definition~\ref{V1-def:yangian}).
\end{construction}

\begin{proposition}[RTT relation from the genus-zero collision YBE;
\ClaimStatusProvedHere]
% label removed: prop:thqg-V-rtt-from-sgybe
\index{RTT relation!from genus-zero YBE}
The genus-zero collision/Yang--Baxter package, expressed in terms
of the transfer matrix
$T_\cA(z) = \sum_r T^{(r)} z^{-r}$, gives the RTT relation:
\begin{equation}% label removed: eq:thqg-V-rtt
r_{12}(z_{12})\,T_1(z_1)\,T_2(z_2)
\;-\;
T_2(z_2)\,T_1(z_1)\,r_{12}(z_{12})
\;=\;
T_1(z_1)\,r_{12}(z_{12})\,T_2(z_2)
\;-\;
T_2(z_2)\,r_{12}(z_{12})\,T_1(z_1),
\end{equation}
which rearranges to the standard RTT form
$r_{12}(z_{12})\,T_1(z_1)\,T_2(z_2)
= T_2(z_2)\,T_1(z_1)\,r_{12}(z_{12})$
when $r$ is the fundamental $R$-matrix.
\end{proposition}

\begin{proof}
The CYBE at genus~$0$ involves $r_{12}$, $r_{13}$, and $r_{23}$ with
$r_{ij}(z) = \Rescoll_{0,2}(\Theta_\cA)|_{ij}$. When acting on a
module via the evaluation representation, $r_{i3}(z)$ becomes the
action of $T(z)$ on the $i$-th factor. The three-term CYBE then
becomes the RTT relation after distributing the Lie bracket over the
tensor product. The expansion in modes $z^{-r-1}$ gives the
quadratic RTT relations among the generators $T^{(r)}$
(cf.~Definition~\ref{V1-def:yangian-rtt}).
\end{proof}


% ======================================================================
%
% SUBSECTION 5: GRAVITATIONAL YANGIAN FOR VIRASORO
%
% ======================================================================

\subsection{Gravitational Yangian for Virasoro}
% label removed: subsec:thqg-V-virasoro-yangian
\index{Virasoro!gravitational Yangian|textbf}
\index{gravitational Yangian!Virasoro}

The Virasoro algebra is the simplest class-$M$ algebra: its shadow
tower is infinite, the $r$-matrix has higher-order poles, and the
genus-zero gravitational Yangian exhibits the richest explicit binary
collision data among the standard test cases. This subsection
computes the explicit Virasoro binary kernel, the ternary shadow,
and the expected representation-theoretic model.


\subsubsection{The Virasoro binary kernel and its gravitational reading}

\begin{construction}[Explicit Virasoro binary kernel;
\ClaimStatusProvedHere]
\label{constr:thqg-V-vir-r-explicit}
\index{Virasoro!r-matrix@$r$-matrix!explicit}
\index{r-matrix@$r$-matrix!Virasoro!explicit construction}
For $\mathrm{Vir}_c$, the Koszul dual is $\mathrm{Vir}_{26-c}$
with Virasoro generators $\{L_n\}_{n \in \Z}$ satisfying
$[L_m, L_n] = (m-n)\,L_{m+n}
+ \frac{26-c}{12}\,(m^3-m)\,\delta_{m+n,0}$.

The Virasoro binary kernel $\rgrav_{\mathrm{Vir}}(z)$ acts on
$\mathrm{Vir}_{26-c} \otimes \mathrm{Vir}_{26-c}$ and is determined
by the collision residue of $\Theta_{\mathrm{Vir}}$:
\begin{equation}% label removed: eq:thqg-V-vir-r-explicit
\rgrav_{\mathrm{Vir}}(z)
\;=\;
\sum_{m \geq 0}
\frac{1}{z^{m+1}}\,
\Omega_m,
\end{equation}
where the $\Omega_m$ are the Virasoro Casimir components:
\begin{align}
\Omega_0
&\;=\;
\sum_{n \in \Z}
L_{-n} \otimes L_n
\;\in\;
(\mathrm{Vir}_{26-c})^{\widehat{\otimes}\,2},
% label removed: eq:thqg-V-vir-omega0 \\
\Omega_1
&\;=\;
\sum_{n \in \Z}
n\,L_{-n} \otimes L_n,
% label removed: eq:thqg-V-vir-omega1 \\
\Omega_m
&\;=\;
\sum_{n \in \Z}
\binom{n+m-1}{m}\,L_{-n-m} \otimes L_n
\qquad (m \geq 0).
% label removed: eq:thqg-V-vir-omegam
\end{align}
The $m = 0$ term is the leading Casimir pole; the $m \geq 1$ terms
are the higher poles arising from the Virasoro fourth-order OPE
singularity.

The $r$-matrix satisfies:
\begin{enumerate}[label=\textup{(\roman*)}]
\item \emph{Residue.}\;
 $\Res_{z=0}\,\rgrav_{\mathrm{Vir}}(z) = \Omega_0$, the
 Virasoro Casimir.
\item \emph{Skew-symmetry.}\;
 $\rgrav_{\mathrm{Vir}}(z)_{12} = -\rgrav_{\mathrm{Vir}}(-z)_{21}$.
\item \emph{CYBE.}\;
 The classical Yang--Baxter equation holds:
 $[\rgrav_{12}(z_{12}),\,\rgrav_{13}(z_{13})]
 + [\rgrav_{12}(z_{12}),\,\rgrav_{23}(z_{23})]
 + [\rgrav_{13}(z_{13}),\,\rgrav_{23}(z_{23})]
 = 0.$
\end{enumerate}
\end{construction}

In the gravitational reading, the kernel
$\rgrav_{\mathrm{Vir}}(z)$ is read as the binary graviton-exchange
kernel for the Virasoro HT package.

\begin{proof}[Verification of the CYBE]
The CYBE for $\rgrav_{\mathrm{Vir}}$ follows from
Theorem~\ref{V1-thm:thqg-V-cybe-from-arnold} applied to the Virasoro
OPE. Explicitly, the three terms of the CYBE evaluate to:
\begin{align}
&[\rgrav_{12}(z_{12}),\;\rgrav_{13}(z_{13})]
\;=\;
\sum_{m_1,m_2}
\frac{[\Omega_{m_1,12},\,\Omega_{m_2,13}]}
 {z_{12}^{m_1+1}\,z_{13}^{m_2+1}},
% label removed: eq:thqg-V-vir-cybe-1
\end{align}
and similarly for the other two terms. Each bracket
$[\Omega_{m_1},\,\Omega_{m_2}]$ is computed using the Virasoro
commutation relations; the resulting Laurent series in $z_{12}, z_{13}$
satisfies the partial-fraction identity
\[
\frac{1}{z_{12}^{a+1}\,z_{13}^{b+1}}
+
\frac{1}{z_{12}^{a+1}\,z_{23}^{b+1}}
+
\frac{1}{z_{13}^{a+1}\,z_{23}^{b+1}}
=
\text{(lower-order poles)},
\]
and the cancellation at each order follows from the Arnold relation
on $\overline{\mathcal{M}}_{0,4}$.

The key identity is the \emph{Virasoro partial-fraction relation}:
for the Casimir components,
\begin{equation}% label removed: eq:thqg-V-vir-partial-fraction
\sum_{\text{cyclic}(1,2,3)}
[\Omega_{m,12},\,\Omega_{m',13}]
\;=\;
0
\end{equation}
when $m + m' = $ const, which is the Virasoro
specialization of the IBR
(Corollary~\ref{V1-cor:thqg-V-ibr}).
\end{proof}


\subsubsection{The ternary shadow and its gravitational reading}

\begin{computation}[Ternary shadow and its gravitational reading;
\ClaimStatusProvedHere]
% label removed: comp:thqg-V-three-graviton
\index{graviton!three-graviton vertex}
\index{Virasoro!three-graviton vertex}
The ternary component of
$\Theta_{\mathrm{Vir}}$ at genus~$0$, extracted via the collision
filtration at depth~$2$. On the space of three Virasoro
primaries of conformal weights $h_1, h_2, h_3$, this operator is:
\begin{equation}% label removed: eq:thqg-V-three-graviton
V_3(h_1, h_2, h_3;\, z_1, z_2, z_3)
\;=\;
\sum_{i<j}
\frac{h_k}{(z_i - z_j)^2}
\;+\;
\frac{1}{z_i - z_j}\,\partial_{z_j},
\end{equation}
where $\{i,j,k\} = \{1,2,3\}$. This is the standard Virasoro
three-point function operator, derived from the shadow connection
$\nabla^{\mathrm{hol}}_{0,3}$
(by the Volume~I Virasoro/BPZ comparison proposition).

The corresponding flat section for primaries $\phi_{h_i}$ is:
\begin{equation}% label removed: eq:thqg-V-three-graviton-amplitude
\langle \phi_{h_1}(z_1)\,\phi_{h_2}(z_2)\,\phi_{h_3}(z_3) \rangle
\;=\;
\frac{C_{h_1 h_2 h_3}}
{z_{12}^{h_1+h_2-h_3}\,z_{13}^{h_1+h_3-h_2}\,z_{23}^{h_2+h_3-h_1}},
\end{equation}
where $C_{h_1 h_2 h_3}$ is the OPE coefficient. This is the flat
section of $\nabla^{\mathrm{hol}}_{0,3}$.

\emph{Gravitational interpretation.}
In the gravitational reading, this ternary shadow is read as the
cubic coupling of three gravitons (stress-tensor insertions) on the
boundary of the HT theory. The denominator structure
$z_{ij}^{-h_i - h_j + h_k}$ is then read as the corresponding
holomorphic momentum-space graviton vertex: the conformal weights
play the role of graviton helicities, and the $z_{ij}$ are the
holomorphic momenta.
\end{computation}

\begin{computation}[Quartic contact term and its gravitational reading;
\ClaimStatusProvedHere]
% label removed: comp:thqg-V-quartic-graviton
\index{graviton!quartic vertex}
\index{quartic contact invariant!graviton}
\index{Virasoro!quartic contact invariant!graviton interpretation}
The degree-$4$ component of
$\Theta_{\mathrm{Vir}}$ at genus~$0$, collision depth~$2$.
Beyond the factorized contribution
(product of two ternary shadows), there is a
\emph{contact} contribution that does not factorize through
any channel:
\begin{equation}% label removed: eq:thqg-V-quartic-contact
V_4^{\mathrm{contact}}(z_1, z_2, z_3, z_4)
\;=\;
\frac{10}{c(5c+22)}\,
\Lambda(z_1, z_2, z_3, z_4),
\end{equation}
where $\Lambda$ is the quartic contact form on
$\overline{\mathcal{M}}_{0,5}$ and the coefficient
$Q^{\mathrm{contact}}_{\mathrm{Vir}} = 10/[c(5c+22)]$ is the
quartic contact invariant
(Computation~\ref{V1-comp:holographic-ss-vir}).

In the gravitational reading, this contact term is read as the first
obstruction to the binary gravitational $r$-matrix being a
\emph{complete} description of the scattering: for class-$M$
algebras, the binary $r$-matrix is then supplemented by ternary,
quartic, and higher shadow data. The quartic contact invariant
measures the strength of this obstruction in that reading.

\emph{Divergence at special central charges.}
The denominator $c(5c+22)$ vanishes at $c = 0$ and
$c = -22/5$. At $c = 0$: the Virasoro algebra degenerates and the
quartic contact is ill-defined. At $c = -22/5$:
the $(2,5)$ minimal model, where the null vector at level~$5$ forces
special behavior. Away from these values, the quartic contact is a
smooth function of~$c$.
\end{computation}


\subsubsection{Highest-weight models for the gravitational Yangian}

\begin{definition}[Expected Virasoro highest-weight model]
% label removed: def:thqg-V-gravitational-hw
\index{gravitational Yangian!highest-weight modules}
In the expected Virasoro realization of the line-side package, one is
led to highest-weight modules generated by a vector
$|h\rangle$ satisfying:
\begin{enumerate}[label=\textup{(\roman*)}]
\item $L_n\,|h\rangle = 0$ for $n > 0$;
\item $L_0\,|h\rangle = h\,|h\rangle$;
\item the $r$-matrix acts as
 $\rgrav_{\mathrm{Vir}}(z)\,|h\rangle \otimes |h'\rangle
 = \frac{h + h'}{z}\,|h\rangle \otimes |h'\rangle
 + (\text{lower poles in }z)$.
\end{enumerate}
In the gravitational reading, these highest-weight modules would model
line objects of conformal weight~$h$.
\end{definition}

\begin{remark}[Expected Verma-module model]
% label removed: prop:thqg-V-verma-gravitational
\index{Verma module!gravitational representation}
In the same expected Virasoro realization, the Verma module $M(h,c)$
is the natural model highest-weight object. The binary shadow action
is expected to take an evaluation-type form
\begin{equation}% label removed: eq:thqg-V-verma-evaluation
\rgrav_{\mathrm{Vir}}(z)\big|_{M(h,c) \otimes M(h',c)}
\;=\;
\frac{1}{z}\,\Omega_0\big|_{M(h,c) \otimes M(h',c)}
\;+\;
\text{(higher-order corrections)}.
\end{equation}
For generic~$h$, one expects the resulting model to be faithful on
$M(h,c)$. For degenerate representations ($h = h_{r,s}$, the Kac
table values), the null vectors generate the ideal, and the quotient
module $L(h_{r,s},c)$ is expected to carry the corresponding truncated
binary shadow action. The PBW and null-vector discussion explains why
this is the natural Virasoro model, but the live surface does not
isolate that realization as a separate proved theorem.
\end{remark}

\begin{remark}[Expected Virasoro highest-weight category]
% label removed: rem:thqg-V-gravitational-rep-category
\index{gravitational Yangian!representation category}
In that expected Virasoro realization, the highest-weight category is
controlled by two parameters: the conformal weight $h$ and the central
charge~$c$. At fixed~$c$, the expected structure is:
\begin{enumerate}[label=\textup{(\roman*)}]
\item \emph{Generic $h$.}
 The Verma module $M(h,c)$ would be irreducible and carry the full
 expected highest-weight model.
\item \emph{Degenerate $h = h_{r,s}$.}
 The irreducible quotient $L(h_{r,s},c)$ would carry a truncated
 expected model. The BPZ null-vector differential equations would
 arise from the truncation of the shadow connection
 $\nabla^{\mathrm{hol}}_{0,n}$ to that quotient.
\item \emph{Rational $c$ (minimal models).}
 The category would have finitely many irreducibles, and the fusion
 rules $N_{ij}^k$ would govern the tensor product of these expected
 model objects. In the same expected realization, the gravitational
 $R$-matrix would reduce to a finite-dimensional matrix acting on the
 finite representation ring.
\end{enumerate}
This structure is the Virasoro specialization of the general
line category, modeled here by modules for the open-colour dual
$\cA^!$ in the sense of
Definition~\ref{V1-def:holographic-modular-koszul-datum}.
\end{remark}


\subsubsection{Self-duality at $c = 13$}

\begin{proposition}[Self-duality of the gravitational Yangian at
$c = 13$; \ClaimStatusProvedHere]
% label removed: prop:thqg-V-c13-self-duality
\index{Virasoro!self-duality at c=13@self-duality at $c = 13$}
\index{gravitational Yangian!self-duality}
At $c = 13$, the Koszul dual $\mathrm{Vir}_{26-c} = \mathrm{Vir}_{13}$
is isomorphic to the original: $\mathrm{Vir}_{13}^! \cong
\mathrm{Vir}_{13}$. The gravitational Yangian becomes self-dual:
\begin{equation}% label removed: eq:thqg-V-c13-self-dual
\Ydg_{\mathrm{Vir}_{13}}
\;\cong\;
\Ydg_{\mathrm{Vir}_{13}^!}
\;\cong\;
\Ydg_{\mathrm{Vir}_{13}},
\end{equation}
and the gravitational $r$-matrix satisfies the enhanced symmetry
$\rgrav(z) = -\rgrav(-z)^{21}$ without change of central charge.

The fixed-point equation $c = 26 - c$ has the unique solution
$c = 13$. This is not $c = 26$ (the critical string): the
self-duality is the Koszul self-duality of the Virasoro algebra,
not the vanishing of the total central charge. At $c = 26$, the
Koszul dual is $\mathrm{Vir}_0$, which is degenerate.
\end{proposition}

\begin{proof}
The Koszul dual of $\mathrm{Vir}_c$ is $\mathrm{Vir}_{26-c}$
(Remark~\ref{V1-rem:virasoro-resonance-model}).
Self-duality requires $c = 26 - c$, i.e., $c = 13$.
The dg-shifted Yangian is a functor of the Koszul dual
$\cA^!$, so self-duality of $\cA^!$ implies self-duality of
$\Ydg_\cA$. The enhanced symmetry of $\rgrav$ follows from the
skew-symmetry property (Theorem~\ref{V1-thm:thqg-V-collision-twisting}(ii))
applied when $\cA = \cA^!$.
\end{proof}

\begin{remark}[Ordered Yangian structure at the self-dual point]
\label{rem:yangian-vir13-self-dual}
\ClaimStatusProvedHere
\index{Virasoro!self-duality at c=13@self-duality at $c = 13$!ordered Yangian}
\index{gravitational Yangian!self-duality!ordered structure}
%
%: kappa(Vir_c) = c/2, from C2 census; c=13 -> 13/2 verified.
%: self-dual at c=13, NOT c=26.
%
The self-duality at $c = 13$ has a refined expression at the level
of the $\Eone$-ordered bar complex
$\barB^{\Eone}(\mathrm{Vir}_{13})
= T^c(s^{-1}\,\overline{\mathrm{Vir}}_{13})$.
Three structures become simultaneously symmetric.

\emph{(i) Koszul complementarity.}\;
The Koszul conductor for Virasoro satisfies
$K = \kappa(\mathrm{Vir}_c) + \kappa(\mathrm{Vir}_{26-c})
= c/2 + (26-c)/2 = 13$
for all~$c$.
%: kappa(Vir_c) = c/2 from landscape_census.tex; c=0 -> 0, c=13 -> 13/2.
At $c = 13$ the two summands coincide:
$\kappa = \kappa' = 13/2$.
This is the unique central charge at which the
$\Eone$-ordered gravitational Yangian acts on
$\mathrm{Vir}_{13}^{\otimes n}$ with equal weight from both the
algebra and its Koszul dual; away from $c = 13$,
the algebra and dual carry different kappa-weights, and the
$R$-matrix intertwines representations of genuinely different
Virasoro algebras.

\emph{(ii) $R$-matrix involution.}\;
The gravitational $r$-matrix $\rgrav_{\mathrm{Vir}_{13}}(z)
= \sum_{m \geq 0} \Omega_m / z^{m+1}$
acts on $\mathrm{Vir}_{13} \otimes \mathrm{Vir}_{13}$
(not on $\mathrm{Vir}_{13} \otimes \mathrm{Vir}_{26-c}$
with different central charges).
The skew-symmetry
$\rgrav(z)_{12} = -\rgrav(-z)_{21}$
from the preceding proposition becomes an \emph{endomorphism}
property: the $R$-matrix is a genuine involution
$R \colon V \otimes V \to V \otimes V$
on a single representation space $V$, rather than an intertwiner
$V \otimes W \to W \otimes V$ between two different spaces.
This is the ordered ($\Eone$) refinement of the scalar identity
$\kappa = \kappa'$ at degree~$2$.

\emph{(iii) Ordered bar self-duality.}\;
Since $\Ydg_\cA
= \Convstr(\barB^{\Eone}(\cA),\,\chirLie)$
(Remark~\ref{V1-rem:yangian-ordered-boundary-face})
and the Verdier intertwining gives
$\mathbb{D}_{\Ran}(\barB^{\Eone}(\cA))
\simeq \barB^{\Eone}(\cA^!)$,
the self-duality $\cA \cong \cA^!$ at $c = 13$
yields an involutive equivalence
\begin{equation}\label{eq:thqg-V-ordered-bar-self-dual}
\barB^{\Eone}(\mathrm{Vir}_{13})
\;\simeq\;
\mathbb{D}_{\Ran}\bigl(\barB^{\Eone}(\mathrm{Vir}_{13})\bigr)
\end{equation}
of $\Eone$-coalgebras on $\Ran(X)$.
The deconcatenation coproduct is Verdier-self-dual, and the
$R$-matrix data (which the $\Sigma_n$-coinvariant shadow
$\barB^{\Sigma}$ forgets) is preserved by the involution.
Passing to the gravitational Yangian,
$\Ydg_{\mathrm{Vir}_{13}} \cong
\Ydg_{\mathrm{Vir}_{13}^!}$
as dg Lie algebras with $R$-matrix, so the RTT presentation
of the Yangian is invariant under the Koszul involution.

In the gravitational reading: the $c = 13$ Virasoro HT theory
has a graviton-exchange kernel that scatters identical particles
(same central charge on both sides of the cut), the genus-zero
binary shadow is an endomorphism of a single Hilbert space, and
the gravitational Yangian acts by self-intertwining operators.
\end{remark}


% ======================================================================
%
% SUBSECTION 6: LINE-OPERATOR CATEGORIES AND MC3
%
% ======================================================================

\subsection{Line categories and MC3}
% label removed: subsec:thqg-V-line-operators-mc3
\index{line category}
\index{MC3!line categories}

The gravitational Yangian $\Ydg_\cA$ contributes additional structure
to the line side in the HT theory only on the affine/proved
dg-shifted-Yangian lane. Independently, the underlying line category
is modeled by modules for the open-colour dual, written here via the
standard algebraic model
$\cC_\cA^{\mathrm{model}} := \cA^!\text{-}\mathsf{mod}$
\textup{(}Theorem~\textup{\ref{thm:lines_as_modules})}. The MC3
thick-generation programme concerns how much of the additional
Yangian/spectral structure extends from the evaluation-generated core
to richer line-side surfaces. This subsection keeps those two layers
separate.


\subsubsection{The line category and its Yangian enhancement}

\begin{definition}[Line category and conditional Yangian enhancement]
% label removed: def:thqg-V-line-operator-category
\index{line category!from gravitational Yangian}
For a modular Koszul chiral algebra~$\cA$, write
\begin{equation}% label removed: eq:thqg-V-line-category
\cC_\cA^{\mathrm{model}}
\;:=\;
\cA^!\text{-}\mathsf{mod}
\end{equation}
for the module category that models the line category.
On the affine/proved Yangian lane, this line-side model is presented by
dg-shifted-Yangian modules
$\mathrm{Mod}_{\Ydg_\cA}(\mathrm{dgVect})$. When a compatible
dg-shifted-Yangian realization and spectral kernel $r_\cA(z)$ are
supplied, this category carries a meromorphic tensor product
\begin{equation}% label removed: eq:thqg-V-meromorphic-tensor
M \otimes^{\mathrm{mer}}_z N
\;:=\;
\bigl(M \otimes N,\;
d_M \otimes \id + \id \otimes d_N + r_\cA(z)\bigr),
\end{equation}
where $r_\cA(z)$ acts as the $r$-matrix on $M \otimes N$,
producing a family of dg modules parametrized by $z \in \C^*$.

When the same spectral kernel admits a compatible quantization, the
braiding on the model category $\cC_\cA^{\mathrm{model}}$ is the
\emph{$R$-matrix braiding}:
\begin{equation}% label removed: eq:thqg-V-braiding
\beta_{M,N}(z)\colon
M \otimes^{\mathrm{mer}}_z N
\;\xrightarrow{\;\sim\;}
N \otimes^{\mathrm{mer}}_{-z} M,
\qquad
\beta_{M,N}(z) = P_{12} \circ R_\cA(z),
\end{equation}
where $P_{12}$ is the permutation and $R_\cA(z) = 1 + \hbar\,r_\cA(z) +
O(\hbar^2)$ is the quantum $R$-matrix obtained by exponentiating the
classical $r$-matrix. The CYBE
(Theorem~\ref{V1-thm:thqg-V-cybe-from-arnold}) ensures that $\beta$
satisfies the braid relation.
\end{definition}

\begin{remark}[Meromorphic vs.\ genuine braiding]
% label removed: rem:thqg-V-meromorphic-braiding
\index{braiding!meromorphic}
On such a compatible spectral realization, the tensor product on
$\cC_\cA^{\mathrm{model}}$ is \emph{meromorphic} in the spectral
parameter~$z$: it is not defined at $z = 0$ (where the $r$-matrix
has a pole), and the braiding is not a morphism of honest dg modules
but a meromorphic family of morphisms. This is the correct structure
for line operators in HT theories: the spectral parameter is the
transverse holomorphic coordinate, and the singularity at $z = 0$ is
the collision singularity of two line operators at the same point.

The passage from meromorphic braiding to genuine braiding requires
either:
\begin{enumerate}[label=\textup{(\roman*)}]
\item specializing $z = q^h$ for a quantum parameter~$q$ and
 conformal weight~$h$ (the Kazhdan--Lusztig approach,
 Theorem~\ref{V1-thm:kazhdan-lusztig-equivalence}); or
\item taking the monodromy representation of the KZ connection
 around the singularity (the Drinfeld--Kohno approach,
 Theorem~\ref{V1-thm:derived-dk-affine}).
\end{enumerate}
Both approaches produce the same genuine braided tensor category,
which is the category of representations of the quantum group
$U_q(\fg)$ (Drinfeld--Kohno theorem).
\end{remark}


\subsubsection{Thick generation and MC3}

\begin{theorem}[Type-$A$ MC3 reduction via the gravitational Yangian;
\ClaimStatusProvedHere]
% label removed: thm:thqg-V-mc3-thick-generation
\index{MC3!thick generation!gravitational Yangian}
In type~A, the evaluation-generated core on the gravitational
line-side surface is
\begin{equation}% label removed: eq:thqg-V-mc3-generation
\cC^{\mathrm{ev}}_{\widehat{\mathfrak{sl}}_N}
\;:=\;
\mathrm{Thick}\bigl\langle
 V_\lambda(z) \mid \lambda \in P^+(\mathfrak{sl}_N),\;
 z \in \C^*
\bigr\rangle
\;\subset\;
\cC_{\widehat{\mathfrak{sl}}_N}.
\end{equation}
The derived Drinfeld--Kohno comparison is proved on this core by the
Volume~I evaluation-generated-core DK theorem. Beyond the
evaluation-generated core, the Volume~I type-$A$ MC3 reduction theorem
proves the shifted-prefundamental-generation and pro-Weyl-recovery
packets and isolates a single remaining compact-completion gap:
the Volume~I compact-completion conjecture, which asks to extend the
evaluation-core DK equivalence to the compact shifted-prefundamental
core and compare that compact-core equivalence with the completed /
pro-nilpotent line-side category.
\end{theorem}

\begin{proof}
The definition of $\cC^{\mathrm{ev}}_{\widehat{\mathfrak{sl}}_N}$ is
tautological. The Volume~I evaluation-generated-core DK theorem proves
the DK comparison on this core. The Volume~I type-$A$ MC3 reduction theorem
then proves the shifted-prefundamental-generation and pro-Weyl
packages and states that the only remaining post-core step is
the compact-completion conjecture. Transporting those
results to the line-side surface gives the claim.
\end{proof}

\begin{corollary}[Type-$A$ DK-5 reduction to the compact-completion packet;
\ClaimStatusProvedHere]
% label removed: cor:thqg-V-dk5-type-a
\index{Drinfeld--Kohno!DK-5!type A accessibility}
On the gravitational line-side surface for
$\widehat{\mathfrak{sl}}_N$, the only remaining obstruction between the
proved evaluation-core DK comparison and a full DK-5 statement is the
compact-completion packet of
the Volume~I compact-completion conjecture. No further independent
type-$A$ obstruction remains beyond that packet.
\end{corollary}

\begin{proof}
Immediate from the preceding theorem and the Volume~I type-$A$ MC3
reduction theorem.
\end{proof}

\begin{remark}[DK-5 resolved for $\mathfrak{sl}_2$]
\label{rem:vol2-dk5-sl2}
\index{Drinfeld--Kohno!DK-5!sl2@$\mathfrak{sl}_2$ resolved}
For $\fg = \mathfrak{sl}_2$, the full chain from the bar complex to the
Verlinde category is now complete
\textup{(}Volume~I,
Proposition~\textup{\ref*{V1-prop:dk5-sl2-frt})}: the ordered bar
complex $\barB^{\mathrm{ord}}(\widehat{\mathfrak{sl}}_{2,k})$ produces
the Yang $R$-matrix $R(u) = uI + \hbar P$, the FRT construction
recovers $U_q(\mathfrak{sl}_2)$ with its Hopf algebra structure, and
at roots of unity $q = e^{2\pi i/(k+2)}$ the Verlinde truncation follows
from the vanishing quantum dimension $[k{+}2]_q = 0$. This is the first
instance where the entire DK-5 programme is verified end-to-end.
\end{remark}

\begin{remark}[MC3 for all simple types: evaluation core proved, post-core conditional]
% label removed: rem:thqg-V-mc3-arbitrary-type
\index{MC3!all simple types}
The prefundamental Clebsch--Gordan closure
\textup{(}package~\textup{(i)} of MC3\textup{)} is proved for all simple
types by the Volume~I all-types categorical Clebsch--Gordan theorem,
and the DK comparison on the evaluation-generated core is proved for
all simple types by the corresponding Volume~I corollary. In type~$A$,
the Volume~I type-$A$ MC3 reduction theorem proves the
shifted-prefundamental-generation and pro-Weyl-recovery packets,
leaving only the compact-completion extension conjecture. For other
types, the analogous post-core mechanism becomes type-independent only
conditionally, once the shifted-prefundamental, pro-Weyl, and
compact-completion packets of
the conditional Volume~I automatic-generalization mechanism are separately
supplied in that type.
\end{remark}

\begin{construction}[MC3 landscape: all simple types]
% label removed: constr:vol2-mc3-gradient
The all-types theorematic surface consists of categorical
Clebsch--Gordan closure and the evaluation-generated-core DK
comparison from Volume~I. In type~$A$, the Volume~I type-$A$ MC3 reduction theorem
proves the shifted-prefundamental-generation and pro-Weyl-recovery
packets, leaving only the compact-completion extension conjecture.
For arbitrary type, the conditional automatic-generalization mechanism
shows that the same post-core mechanism is available only
conditionally, once the corresponding shifted-prefundamental,
pro-Weyl, and compact-completion packets are supplied in that type.
\end{construction}

\begin{remark}[Type $B_2$ after categorical CG closure]%
% label removed: conj:mc3-type-B2
\index{MC3!type B2@type $B_2$}
For $\fg = \mathfrak{so}_5$ \textup{(}type $B_2$\textup{)}, the
categorical prefundamental Clebsch--Gordan decomposition
\begin{equation}% label removed: eq:mc3-B2-cg
V_{\omega_i}(a) \otimes L^-_i(b)
\;\cong\;
\bigoplus_j L^-_i(b_j)
\end{equation}
holds as $Y(\mathfrak{so}_5)$-modules in
$\cO^{\mathrm{sh}}_{\leq 0}$, not merely at the $K_0$-level.
This is a special case of the all-types result
\textup{(}proved in Volume~I on the evaluation-generated core\textup{)}.

Promoting this evaluation-core statement to the full post-core MC3
package for $\mathfrak{so}_5$ would require the
shifted-prefundamental, pro-Weyl, and compact-completion packets of
the conditional Volume~I automatic-generalization mechanism. Those packets
are not proved here for type~$B_2$.

\smallskip\noindent
\textbf{Proof method.} The uniform all-types proof
in Volume~I uses the
multiplicity-free $\ell$-weight property (Chari--Moura, 2006)
to bypass the folding strategy entirely. Type $B_2$ has two
fundamental weights $\omega_1$ \textup{(}vector\textup{)} and
$\omega_2$ \textup{(}spinor\textup{)}. The CG closure at
$K_0$-level
\textup{(}the Volume~I prefundamental Clebsch--Gordan proposition\textup{)}
lifts categorically via multiplicity-freeness of $\ell$-weight spaces.
\end{remark}


\subsubsection{The Yangian-shadow principle}

\begin{principle}[Yangian-shadow principle]
% label removed: princ:thqg-V-yangian-shadow
\index{Yangian-shadow principle}
For every holomorphic-topological theory~$T$ with boundary chiral
algebra~$\cA$:
\begin{enumerate}[label=\textup{(\Alph*)}]
\item The genus-zero binary shadow
 $r_\cA(z) = \Rescoll_{0,2}(\Theta_\cA)$ recovers the
 dg-shifted-Yangian data on the affine / evaluation-core lane;
 independently, a standard algebraic model for the line side is
 $\cC_\cA^{\mathrm{model}} := \cA^!\text{-}\mathsf{mod}$.
\item The genus-zero $n$-point shadow
 $\Sh_{0,n}(\Theta_\cA)$ identifies on the affine Kac--Moody
 comparison surface with the KZ connection, and on the Virasoro
 degenerate-representation comparison surface its higher collision
 residues yield the BPZ null-vector equations.
\item The formal genus-refined binary shadow series
 $R^{\mathrm{bin}}_\cA(z;\hbar) =
 \sum_g \hbar^{2g}\,r_{\cA,g}(z)$ is the candidate higher-genus
 binary package extracted from the shadow coefficients.
\item On the conjectural modular-Yangian surface, a genuine modular
 kernel in $\Ymod_\cA$ would satisfy the stable-graph YBE and furnish
 the higher-genus completion of the genus-zero Yangian package.
\end{enumerate}
Components~(A)--(B) are proved
\textup{(}Theorems~\textup{\ref{V1-thm:collision-residue-twisting}},
\textup{\ref{V1-thm:collision-depth-2-ybe}},
and the Volume~I affine shadow/KZ comparison theorem\textup{)}.
Component~(C) is the proved formal binary-shadow series of
Theorem~\textup{\ref*{V1-thm:thqg-V-pro-mc-element}}.
Component~(D) is frontier modular-kernel data, not a theorem of this
chapter.

The principle asserts: \emph{the gravitational Yangian is not a new
structure imposed on the HT theory; its universal proved core is the
genus-zero binary shadow of the universal MC element, while the
higher-genus modular completion is an intended extension of that
shadow package.}
\end{principle}


\subsubsection{Dictionary: Yangian data vs.\ gravitational data}

\begin{remark}[Yangian-gravity dictionary]
% label removed: rem:thqg-V-dictionary
\index{gravitational Yangian!dictionary}
The following table translates between the Yangian language of
\S\ref{V1-chap:yangians} and the gravitational language of the present
chapter:

\begin{center}
\renewcommand{\arraystretch}{1.3}
\small
\begin{tabular}{@{}lll@{}}
\toprule
\textbf{Yangian object} & \textbf{Gravitational object}
 & \textbf{Status} \\
\midrule
$r_\cA(z)$ & Binary graviton exchange in the gravitational reading
 & Proved (Thm~\ref{V1-thm:thqg-V-collision-twisting}) \\
CYBE & Tree-level $3$-graviton consistency in the gravitational reading
 & Proved (Thm~\ref{V1-thm:thqg-V-cybe-from-arnold}) \\
$\Ydg_\cA$ & Genus-zero scattering algebra on the dg-shifted-Yangian lane
 & Proved (Def~\ref{V1-def:thqg-V-modular-yangian}) \\
$R^{\mathrm{bin}}_\cA(z;\hbar)$ & Formal higher-genus binary shadow series
 & Proved (Thm~\ref*{V1-thm:thqg-V-pro-mc-element}) \\
sgYBE & All-loop graviton consistency in the modular-kernel reading
 & Conjectural modular-kernel surface
 (Def~\ref*{V1-def:thqg-V-stable-graph-ybe}) \\
$\cC_\cA^{\mathrm{model}} := \cA^!\text{-mod}$ & Module model for the line category
 & Proved (Thm~\ref{thm:lines_as_modules}) \\
RTT relation & Evaluation-core graviton factorization in the gravitational reading
 & Proved (Prop~\ref{V1-prop:thqg-V-rtt-from-sgybe}) \\
Evaluation module $V_\lambda(z)$ & Evaluation-core line object
 & Proved on the evaluation-generated core
 (Volume~I) \\
Shadow connection $\nabla^{\mathrm{hol}}$
 & KZ on affine Kac--Moody comparison surface; BPZ on Virasoro
 degenerate-representation surface
 & Proved (Volume~I affine shadow/KZ comparison theorem and
 Virasoro/BPZ comparison proposition) \\
Quartic contact $\mathfrak{Q}$ & $4$-graviton obstruction in the gravitational reading
 & Proved (Comp~\ref{V1-comp:thqg-V-quartic-graviton}) \\
\bottomrule
\end{tabular}
\end{center}
\emph{The proved entries are direct projections of the single MC
element $\Theta_\cA$; the sgYBE row records the conjectural
higher-genus completion surface.}
\end{remark}


\subsubsection{The gravitational Yangian as ordered boundary face}

\begin{remark}[The gravitational Yangian as $\Eone$-ordered face;
cf.~Remark~\textup{\ref{V1-rem:yangian-ordered-boundary-face}}]
% label removed: rem:thqg-V-ordered-face
\index{gravitational Yangian!as ordered boundary face}
The gravitational Yangian $\Ydg_\cA$ is the $\Eone$-ordered face of
the universal twisting:
\begin{equation}% label removed: eq:thqg-V-ordered-face
\Ydg_\cA
\;=\;
\Convstr\!\bigl(\barB^{\Eone}(\cA),\;\chirLie\bigr),
\end{equation}
the strict convolution dg Lie algebra of the $\Eone$-bar complex.
The bar complex $\barB^{\Eone}(\cA)$ lives on \emph{ordered}
configurations, the Verdier duality inverts the $R$-matrix
($\mathbb{D}\,r(z) = r(-z)^{-1}$, not the level), and the shadow
tower on ordered configurations encodes the full spectral-parameter
dependence of the rational $R$-matrix through all degrees.

The passage from the unordered ($\Einf$-chiral) modular
convolution algebra $\gAmod$ to the ordered ($\Eone$-chiral)
gravitational Yangian $\Ydg_\cA$ is not a restriction but a
refinement: the ordered structure retains the spectral-parameter
dependence that the unordered structure averages over. For
commutative ($\Einf$-chiral) algebras such as $\widehat{\fg}_k$,
the gravitational Yangian is the Yang $R$-matrix; for genuinely
$\Eone$-chiral algebras such as $Y(\fg)$, it is the full RTT
Yangian.
\end{remark}


\subsubsection{Examples: gravitational Yangians of the standard families}

\begin{computation}[Gravitational Yangian of Heisenberg;
\ClaimStatusProvedHere]
% label removed: comp:thqg-V-heis-yangian
\index{Heisenberg algebra!gravitational Yangian}
For rank-$N$ Heisenberg $\mathcal{H}^N_\kappa$:
\begin{enumerate}[label=\textup{(\roman*)}]
\item $\Ydg_{\mathcal{H}} = \C[\![z^{-1}]\!]$
 (abelian, rank $1$).
\item $R^{\mathrm{bin}}_{\mathcal{H}}(z;\hbar) = N/(\kappa z)$
 (no higher-genus binary corrections).
\item A standard model for the line side is
 $\cC_{\mathcal{H}}^{\mathrm{model}} = \mathrm{Vect}^{\Z}$
 (graded vector spaces, one line object per charge).
\item The higher-genus binary shadow corrections vanish, so the
 conjectural modular completion would add no new binary terms.
\end{enumerate}
The Heisenberg gravitational Yangian is the simplest possible
case: one generator, one relation (which is trivial), no quantum
corrections. In the gravitational reading, this corresponds to a free
theory.
\end{computation}

\begin{computation}[Gravitational Yangian of affine KM;
\ClaimStatusProvedHere]
% label removed: comp:thqg-V-affine-yangian
\index{Kac--Moody!gravitational Yangian}
For $\widehat{\fg}_k$ at non-critical level:
\begin{enumerate}[label=\textup{(\roman*)}]
\item $\Ydg_{\widehat{\fg}_k} \cong Y(\fg)$
 (the standard Yangian of~$\fg$).
\item $r_{\widehat{\fg}_k,0}(z) = \Omega_\fg/((k{+}h^\vee)\,z)$
 (Casimir $r$-matrix).
\item A standard model for the line side is
 $\cC_{\widehat{\fg}_k}^{\mathrm{model}} \simeq
 \cO^{\mathrm{sh}}_q(\widehat{\fg})$
 at $q = e^{i\pi/(k+h^\vee)}$.
\item At genus~$0$, the collision/Yang--Baxter package is the
 standard CYBE for $\fg$; the corresponding higher-genus modular
 completion is expected to involve KZB-type corrections.
\end{enumerate}
The affine KM gravitational Yangian is the classical Yangian~$Y(\fg)$:
in the gravitational reading, this corresponds to the semistrict
higher-spin theory whose scattering is read through the rational
$R$-matrix of~$\fg$.
\end{computation}

\begin{computation}[Gravitational Yangian of $\mathrm{Vir}_c$;
\ClaimStatusProvedHere]
% label removed: comp:thqg-V-vir-yangian-summary
\index{Virasoro!gravitational Yangian!summary}
For $\mathrm{Vir}_c$:
\begin{enumerate}[label=\textup{(\roman*)}]
\item $\Ydg_{\mathrm{Vir}_c}$ is an infinite-dimensional
 dg Lie algebra with generators $\{\Omega_m\}_{m \geq 0}$
 (Construction~\ref{constr:thqg-V-vir-r-explicit}).
\item $\rgrav_{\mathrm{Vir}}(z) =
 \sum_{m \geq 0} \Omega_m/z^{m+1}$ (infinite series of
 Casimir components).
\item A standard model for the line side is
 $\cC_{\mathrm{Vir}_c}^{\mathrm{model}} =
 \mathrm{Vir}_{26-c}\text{-}\mathsf{mod}$
 (the expected Virasoro module model for the line side).
\item The candidate higher-genus binary-shadow package suggests an
 infinite tower of corrections beyond the genus-zero CYBE.
\item Self-duality at $c = 13$; degeneration at $c = 0$;
 the depth-zero resonance shadow at $c = 26$.
\item The quartic contact invariant
 $Q^{\mathrm{contact}}_{\mathrm{Vir}} = 10/[c(5c+22)]$
 is the first obstruction to the $r$-matrix being a
 complete description.
\end{enumerate}
The Virasoro gravitational Yangian encodes the richest explicit
genus-zero binary shadow data for HT theories with stress-tensor
boundary operators. In the gravitational reading, its infinite
shadow obstruction tower is expected to model an infinite graviton-vertex tower,
while the full higher-genus modular completion remains frontier.
\end{computation}


\subsubsection{The gravitational completion problem}

\begin{remark}[Gravitational completion and MC4]
% label removed: rem:thqg-V-gravitational-completion
\index{gravitational Yangian!completion problem}
\index{MC4!gravitational completion}
The MC4 completion theorem
(Theorem~\ref{V1-thm:completed-bar-cobar-strong}) ensures that the
gravitational Yangian admits a completion:
\begin{equation}% label removed: eq:thqg-V-completion
\widehat{\Ydg}_\cA
\;:=\;
\varprojlim_{N}\;
\Ydg_\cA / F^{N+1}\Ydg_\cA,
\end{equation}
with the strong filtration axiom
$\mu_r(F^{i_1},\dotsc,F^{i_r}) \subseteq F^{i_1+\cdots+i_r}$
ensuring that the completion is compatible with the $A_\infty$
structure (Lemma~\ref{V1-lem:degree-cutoff}).

For the Virasoro gravitational Yangian: the completion
$\widehat{\Ydg}_{\mathrm{Vir}}$ is the pronilpotent envelope of
the infinite-dimensional dg Lie algebra generated by the Casimir
components $\{\Omega_m\}_{m \geq 0}$. The MC4 splitting
distinguishes:
\begin{enumerate}[label=\textup{(\roman*)}]
\item \emph{MC4$^+$ (positive towers)}:
 $\mathcal{W}_{1+\infty}$, affine Yangians, solved by weight
 stabilization
 (Theorem~\ref{V1-thm:stabilized-completion-positive}).
\item \emph{MC4$^0$ (resonant towers)}:
 Virasoro, non-quadratic $\mathcal{W}_N$, reduced to finite
 resonance by
 Theorem~\ref{V1-thm:resonance-filtered-bar-cobar}.
\end{enumerate}
The gravitational completion problem is: identify the completed
gravitational Yangian $\widehat{\Ydg}_\cA$ with a standard
quantum group. For affine KM, this is the quantum group
$U_q(\fg)$ (Drinfeld--Kohno). For Virasoro, the target is the
$\mathcal{W}$-algebra quantum group, whose identification is
part of the MC4 resonance programme.
\end{remark}


\subsubsection{Summary: the gravitational Yangian as projection of $\Theta_\cA$}

\begin{remark}[Result G5: synthesis]
% label removed: rem:thqg-V-g5-synthesis
\index{gravitational Yangian!synthesis}
Result~G5 is the statement that the gravitational Yangian
$\Ydg_\cA$ and its ambient pronilpotent modular completion
$\Ymod_\cA$ are organized by the single MC element $\Theta_\cA$
(Theorem~\ref*{V1-thm:mc2-bar-intrinsic}). The content is:
\begin{enumerate}[label=\textup{(\arabic*)}]
\item The collision filtration on $\gAmod$ stratifies the
 genus-zero sector by FM boundary depth.
\item The binary collision residue
 $\Rescoll_{0,2}(\Theta_\cA) = \pi_\cA$ recovers the
 universal twisting morphism and the spectral $r$-matrix.
\item The Arnold relation on $\overline{\mathcal{M}}_{0,4}$ forces
 the CYBE\@.
\item The pronilpotent completion $\Ymod_\cA$ provides the ambient
 modular dg Lie package, and the genus-refined binary projections
 define the formal series
 $R^{\mathrm{bin}}_\cA(z;\hbar)=\sum_g \hbar^{2g}r_{\cA,g}(z)$.
 Comparison with a genuine modular kernel
 $R^{\mathrm{mod}}_\cA(z;\hbar)$ satisfying the sgYBE is frontier.
\item For Virasoro, the binary kernel $\rgrav_{\mathrm{Vir}}(z)$ has higher-order
 poles encoding the infinite shadow obstruction tower; the quartic contact
 invariant $Q^{\mathrm{contact}} = 10/[c(5c+22)]$ is the first
 nonlinear obstruction; and self-duality holds at $c = 13$.
\item The line category is modeled here by
 $\cC_\cA^{\mathrm{model}} := \cA^!\text{-}\mathsf{mod}$,
 and this line-side model is proved independently,
 while the Yangian MC3 surface has an all-types evaluation-core
 theorem and a post-core compact/completion frontier.
\end{enumerate}
Items~(1)--(3), together with the ambient completion in~(4) and the
line-category statement in~(6), are consequences of the proved
algebraic core. The gravitational Yangian is not an additional
structure; its universal proved content is the ordered binary shadow
of the universal MC class, while the modular-kernel completion in~(4)
remains frontier.
\end{remark}

\subsection{Gravitational Yangian as
$Q_\fg^{k,f,\mu}|_{\fg=\mathrm{Vir}}$}
\label{subsec:thqg-gravitational-yangian-anchor}

\begin{remark}[Gravitational case as type-2 dg-shifted on the
infinite-dimensional Virasoro algebra; \ClaimStatusProvedHere]
\label{rem:thqg-gravitational-yangian-yangian-type}
In the four-Yangian-types disambiguation, the gravitational
dg-shifted Yangian $Y^{\mathrm{dg}}(\mathrm{Vir}_c)$ is the
\emph{type-2} object (complete filtered $A_\infty$-algebra at a
point/formal disk) with
$\fg$ replaced by the infinite-dimensional Virasoro algebra
$\mathrm{Vir}_c$. The collapse phenomenon already noted in this
chapter (the level-$r$ generators $\mathsf{T}_n^{(r)} = \binom{n+r-1}{r}\,L_{-n-r}$ are scalar multiples of single
Virasoro modes; the underlying Lie reduces to
$\mathrm{Vir}_{26-c}$) does not eliminate the dg content: the
$A_\infty$ tower $\{m_k\}_{k\ge 3}$ remains nontrivial for generic
$c$ and is the genuine output of the construction. As a
specialisation of the unified
$Q_\fg^{k,f,\mu}$ (Theorem~\ref{thm:unified-chiral-quantum-group}),
the gravitational Yangian is the fibre at $(\fg, k, f, \mu) =
(\mathrm{sl}_2, k_{\mathrm{Sug}}, f_{\mathrm{prin}}, 0)$ with
$c = c(k_{\mathrm{Sug}}) = 1 - 6(k+2)^{-1}(k+1)^2$ via the
Virasoro Sugawara construction.
\end{remark}

\begin{remark}[Notion-B for the gravitational dg-shifted Yangian;
\ClaimStatusProvedElsewhere]
\label{rem:thqg-gravitational-yangian-notion-b}
$Y^{\mathrm{dg}}(\mathrm{Vir}_c)$ is read as a Notion-B
$E_1$-chiral algebra: structure maps $\{m_k\}_{k\ge 1}$ in
$\mathrm{End}^{\mathrm{ch}}_{\mathrm{Vir}_c}$, Stasheff identities
chain-level, the binary residue $r^{\mathrm{grav}}_1 = (c/2)
L_{-1}^{\otimes 2}$ as the leading-order Maurer--Cartan datum.
The $\Sigma_2$-coinvariant projection
$\mathrm{av}(r^{\mathrm{grav}}(z)) = \kappa(\mathrm{Vir}_c) = c/2$
recovers the modular characteristic at degree~$2$, in agreement
with the $E_1$-primacy directive.
\end{remark}

\begin{remark}[Level-prefix discipline for the gravitational kernel;
\ClaimStatusProvedElsewhere]
\label{rem:thqg-gravitational-yangian-rmatrix}
The gravitational $r$-matrix $r^{\mathrm{grav}}(z) = \sum_m \Omega_m
z^{-m-1}$ carries the central charge $c$ as its level prefix
implicitly through the binary residue normalization
$\Omega_0 = (c/2)\,L_{-1}^{\otimes 2}$. At $c=0$ the
gravitational scattering vanishes (the formal-series identity must
be specified algebraically); at $c=13$ the Koszul-self-dual point
gives the gravitational Yangian its self-dual specialisation; at
$c=26$ the matter--ghost critical locus emerges. The level prefix
$c$ is load-bearing and never omitted in this chapter (verified
verbatim against the gravitational kernel formulae above).
\end{remark}

\begin{remark}[Miura $(\Psi-1)/\Psi$ universality on the spectral
coproduct of the gravitational Yangian; \ClaimStatusProvedElsewhere]
\label{rem:thqg-gravitational-yangian-miura-coproduct}
The spectral coproduct
$\Delta_z\colon \mathrm{Vir}_c \to (\mathrm{Vir}_c)^{\otimes 2}((z))$
of the gravitational dg-shifted Yangian has the form
\[
\Delta_z(T) \;=\; T\otimes 1 \;+\; 1\otimes T \;+\;
\frac{\Psi-1}{\Psi}\,(J\otimes J)\,z^{-1} \;+\; O(z^{-2})
\]
with $\Psi = c/12$ being the Drinfeld--Sokolov coupling reading.
The $(\Psi-1)/\Psi$ coefficient is the universal Miura
cross-term universality (Vol~I
Theorem~\ref*{V1-thm:miura-cross-universality}). At
$\Psi=2$ ($c=24$, the Schellekens point) the substitution gives
$1/\Psi = (\Psi-1)/\Psi = 1/2$, an algebraic coincidence; the
correct universality is $(\Psi-1)/\Psi$ for every spin $s\ge 2$,
verified at $\Psi=3$ ($c=36$).
\end{remark}

\begin{remark}[Gravitational Yangian as boundary input to the
universal-holography master theorem]
\label{rem:thqg-gravitational-yangian-uhm}
The gravitational dg-shifted Yangian
$Y^{\mathrm{dg}}(\mathrm{Vir}_c)$ is the open-colour boundary
input to \cref{thm:universal-holography-master} (Vol~II
\cref{ch:programme-climax-platonic}) at the class M (Virasoro)
specialisation of $\Phi_{\mathrm{hol}}$: the bulk 3d HT theory
$T_{\mathrm{Vir}_c}$ is the chiral derived centre
$\mathrm{C}^\bullet_{\mathrm{ch}}(\mathrm{Vir}_c, \mathrm{Vir}_c)$,
identified via the Drinfeld--Sokolov/Hochschild compatibility bridge
(\cref{thm:chd-ds-hochschild}) with
$H^\bullet_{\mathrm{DS}}(\mathrm{C}^\bullet_{\mathrm{ch}}
(V_k(\mathrm{sl}_2), V_k(\mathrm{sl}_2)))$ at the level $k$ giving
the desired $c$ via Sugawara. The Vol~I unified Koszul reflection
is the bar-side input on the boundary; the Vol~I universal-conductor
theorem controls the dual reflections $c\mapsto 13-c$ (Koszul) and
$c\mapsto 26-c$ (matter--ghost), which agree at the joint fixed
point $c=13$ (self-dual Koszul Virasoro) only via the
gravitational-Yangian collapse seen in this chapter.
\end{remark>

% Section file for Chapter: Twisted Holography and Quantum Gravity
% Result (G6): Subleading Soft Graviton Theorems from the Shadow Tower
%
% Dependencies:
%   - thm:mc2-bar-intrinsic (bar-intrinsic MC element)
%   - def:shadow-postnikov-tower (shadow Postnikov tower)
%   - def:modular-shadow-connection (modular shadow connection)
%   - prop:master-equation-from-mc (all-arity master equation)
%   - thm:shadow-connection-kz (KZ from shadow connection)
%   - thm:cubic-gauge-triviality (cubic gauge triviality)
%   - thm:w-virasoro-quintic-forced (Virasoro quintic forced)
%
% Status: \ClaimStatusProvedHere for the algebraic theorems;
%   celestial holography interpretation is physical/structural.

\section{Soft graviton theorems from the shadow tower}
\label{sec:thqg-soft-graviton-theorems}
\index{soft graviton theorem|textbf}
\index{Ward identity!gravitational soft theorem}
\index{shadow Postnikov tower!soft graviton theorems}
\index{celestial holography!from modular shadow connection}

% ======================================================================
% LOCAL MACROS (providecommand: no conflict with main.tex preamble)
% ======================================================================
\providecommand{\Sh}{\operatorname{Sh}}
\providecommand{\cA}{\mathcal{A}}
\providecommand{\cV}{\mathcal{V}}
\providecommand{\MC}{\text{MC}}
\providecommand{\gAmod}{\mathfrak{g}_{\cA}^{\mathrm{mod}}}
\providecommand{\Defcyc}{\operatorname{Def}_{\mathrm{cyc}}}
\providecommand{\dzero}{d_0}
\providecommand{\barBch}{\bar{B}^{\text{ch}}}
\providecommand{\Omegach}{\Omega^{\text{ch}}}
\providecommand{\dfib}{d_{\mathrm{fib}}}
\providecommand{\Ydg}{\mathcal{Y}^{\mathrm{dg}}}
\providecommand{\Gmod}{\mathcal{G}^{\mathrm{mod}}}
\providecommand{\Convstr}{\operatorname{Conv}_{\mathrm{str}}}
\providecommand{\Convinf}{\operatorname{Conv}_{\infty}}
\providecommand{\Definfmod}{\operatorname{Def}_{\infty}^{\mathrm{mod}}}
\providecommand{\Res}{\operatorname{Res}}
\providecommand{\Tr}{\operatorname{Tr}}
\providecommand{\ad}{\operatorname{ad}}
\providecommand{\id}{\mathrm{id}}
\providecommand{\rank}{\operatorname{rank}}
\providecommand{\wt}{\operatorname{wt}}
\providecommand{\Hom}{\operatorname{Hom}}
\providecommand{\End}{\operatorname{End}}
\providecommand{\Sym}{\operatorname{Sym}}
\providecommand{\Aut}{\operatorname{Aut}}
\providecommand{\sgn}{\operatorname{sgn}}
\providecommand{\Dg}[1]{D_{#1}}
\providecommand{\mcurv}[1]{m_0^{(#1)}}
\providecommand{\colldiv}[2]{D_{#1#2}}
\providecommand{\LogForm}[2]{\eta_{#1#2}}
\providecommand{\fg}{\mathfrak{g}}
\providecommand{\cW}{\mathcal{W}}

% ======================================================================
%
%   SECTION OVERVIEW
%
% The modular shadow connection
%   nabla^hol_{g,n} = d - Sh_{g,n}(Theta_A)
% is a flat connection on the space of derived conformal blocks
% V_{g,n}(A).  Its flatness is a consequence of the MC equation
% for Theta_A (Theorem thm:mc2-bar-intrinsic).  When the arity
% decomposition of Sh_{g,n} is expanded, each arity contributes
% a Ward identity.  These are the soft graviton theorems.
%
%   Arity 2: Weinberg leading soft theorem
%   Arity 3: Cachazo--Strominger subleading
%   Arity 4: sub-subleading from quartic contact Q^contact
%   Arity r: all-arity nonlinear soft theorem from Sh_r
%
% The shadow depth r_max(A) controls how many independent soft
% theorems exist.  The infinite tower for Virasoro/W_N algebras
% is the gravitational instance.
%
% ======================================================================

The central formula of the holographic modular Koszul datum
(Definition~\ref{def:holographic-modular-koszul-datum}) is the modular
shadow connection
\begin{equation}\label{eq:thqg-VI-shadow-connection-master}
\nabla^{\mathrm{hol}}_{g,n}
\;=\;
d \;-\; \Sh_{g,n}(\Theta_\cA),
\end{equation}
a flat connection on the space of derived conformal blocks
$\cV_{g,n}(\cA)$
(Definition~\ref{def:modular-shadow-connection}).
Flatness
\begin{equation}\label{eq:thqg-VI-flatness}
\bigl(\nabla^{\mathrm{hol}}_{g,n}\bigr)^2 = 0
\end{equation}
is the $(g,n)$-projection of the Maurer--Cartan equation
$D_\cA\Theta_\cA + \tfrac{1}{2}[\Theta_\cA,\Theta_\cA] = 0$
(Theorem~\ref{thm:mc2-bar-intrinsic}).  The horizontal sections of
$\nabla^{\mathrm{hol}}_{g,n}$ are the genus-$g$, $n$-point
correlators of the holographic theory.  The Ward identities of this
connection, decomposed by the arity filtration of the shadow Postnikov
tower (Definition~\ref{def:shadow-postnikov-tower}), are the soft
graviton theorems of the title.

The purpose of this section is fourfold:
\begin{enumerate}[label=\textup{(\arabic*)}]
\item to prove that the arity-$2$ Ward identity of
  $\nabla^{\mathrm{hol}}_{g,n}$ reproduces the Weinberg leading soft
  graviton theorem and its BMS extension
  (\S\ref{subsec:thqg-VI-leading-soft});
\item to prove that the arity-$3$ Ward identity gives the
  Cachazo--Strominger subleading soft graviton theorem, with the cubic
  gauge triviality of Theorem~\ref{thm:cubic-gauge-triviality}
  controlling when the subleading theorem is non-trivial
  (\S\ref{subsec:thqg-VI-subleading-soft});
\item to prove that the arity-$4$ Ward identity produces the
  sub-subleading soft graviton theorem with explicit coefficient
  $Q^{\mathrm{contact}}_{\mathrm{Vir}} = 10/[c(5c+22)]$, and to
  derive the clutching law that controls its factorization across
  degenerate curves
  (\S\ref{subsec:thqg-VI-sub-subleading-soft});
\item to establish the infinite tower of higher soft theorems for
  algebras with $r_{\max}(\cA) = \infty$ (Virasoro and $\mathcal{W}_N$
  families), prove $o_5(\mathrm{Vir}_c) \neq 0$, and develop the
  celestial holography interpretation
  (\S\ref{subsec:thqg-VI-higher-soft},
  \S\ref{subsec:thqg-VI-celestial}).
\end{enumerate}

Throughout this section, $\cA$ is a modular Koszul chiral algebra on a
smooth projective curve~$X$, and $\Theta_\cA \in \MC(\gAmod)$ is the
universal MC element of
Theorem~\ref{thm:mc2-bar-intrinsic}.  The shadow depth is
$r_{\max}(\cA) := \sup\{r \geq 2 : \cA^{\mathrm{sh}}_{r,0} \neq 0\}$,
and the shadow Postnikov tower is
$\Theta_\cA^{\leq 2} \to \Theta_\cA^{\leq 3} \to
\Theta_\cA^{\leq 4} \to \cdots \to \Theta_\cA$.
All gradings are cohomological ($|d| = +1$).

% ======================================================================
%
%   SUBSECTION 1: SHADOW CONNECTION AND ARITY DECOMPOSITION
%
% ======================================================================

\subsection{The shadow connection and its arity decomposition}
\label{subsec:thqg-VI-shadow-connection}
\index{shadow connection!arity decomposition}
\index{arity filtration!shadow connection}

\subsubsection{The shadow representation}

\begin{definition}[Shadow representation at $(g,n)$]
\label{def:thqg-VI-shadow-representation}
\index{shadow representation|textbf}
For each pair $(g,n)$ with $2g - 2 + n > 0$, the
\emph{shadow representation at $(g,n)$} is the linear map
\begin{equation}\label{eq:thqg-VI-shadow-representation}
\Sh_{g,n} \colon
\Defcyc^{\mathrm{mod}}(\cA)
\;\longrightarrow\;
\End\bigl(\cV_{g,n}(\cA)\bigr)
\end{equation}
defined as follows.  An element
$\varphi \in \Defcyc^{\mathrm{mod}}(\cA)$
of genus~$h$ and arity~$r$ acts on a conformal block
$\Phi \in \cV_{g,n}(\cA; V_1, \dotsc, V_n)$ by evaluating the
multilinear operation $\varphi$ on~$r$ insertions drawn from the
$n$-point configuration, paired with the moduli-space coefficient
from $C_\bullet(\overline{\mathcal{M}}_{g,n+1})$, and summed over
all ways to select the $r$~insertions from the $n$~marked points.

Concretely, if $\varphi$ has arity~$r$ and genus~$h$:
\begin{equation}\label{eq:thqg-VI-shadow-action-concrete}
\bigl(\Sh_{g,n}(\varphi) \cdot \Phi\bigr)(z_1,\dotsc,z_n)
\;=\;
\sum_{\substack{I \subseteq \{1,\dotsc,n\} \\ |I| = r}}
\int_{\overline{\mathcal{M}}_{h,r+1}}
\varphi\bigl(\phi_{i_1}(z_{i_1}),\dotsc,\phi_{i_r}(z_{i_r})\bigr)
\;\Phi\bigl|_{I \to \varphi(I)},
\end{equation}
where the notation $\Phi|_{I \to \varphi(I)}$ denotes the conformal
block with the $r$~insertions in~$I$ replaced by the output
of~$\varphi$, and the moduli integral produces the coefficient-valued
endomorphism.
\end{definition}

\begin{remark}[Arity filtration on the connection form]
\label{rem:thqg-VI-arity-filtration}
\index{shadow connection!filtration by arity}
The modular convolution algebra $\gAmod$ carries a natural
\emph{arity filtration}: $F^{\leq r}\gAmod$ consists of elements
$\varphi$ with arity at most~$r$.  Since $\Theta_\cA$ decomposes
as $\Theta_\cA = \sum_{r \geq 2} \Theta_\cA^{(r)}$ where
$\Theta_\cA^{(r)} \in F^r\gAmod / F^{r-1}\gAmod$ is the arity-$r$
component, the shadow connection form decomposes as
\begin{equation}\label{eq:thqg-VI-arity-decomposition}
\Sh_{g,n}(\Theta_\cA)
\;=\;
\sum_{r=2}^{\infty}
\Sh_{g,n}\bigl(\Theta_\cA^{(r)}\bigr)
\;=\;
\Sh_{g,n}\bigl(\kappa(\cA)\bigr)
\;+\;
\Sh_{g,n}\bigl(\mathfrak{C}(\cA)\bigr)
\;+\;
\Sh_{g,n}\bigl(\mathfrak{Q}(\cA)\bigr)
\;+\;
\sum_{r=5}^{\infty}\Sh_{g,n}\bigl(\Sh_r(\cA)\bigr).
\end{equation}
The notation uses:
$\Theta_\cA^{(2)} = \kappa(\cA)$ (the modular characteristic,
Theorem~D),
$\Theta_\cA^{(3)} = \mathfrak{C}(\cA)$ (the cubic shadow),
$\Theta_\cA^{(4)} = \mathfrak{Q}(\cA)$ (the quartic shadow /
contact invariant), and
$\Theta_\cA^{(r)} = \Sh_r(\cA) = \cA^{\mathrm{sh}}_{r,0}$
for general~$r$.

For an algebra of shadow depth $r_{\max}(\cA) < \infty$,
the sum~\eqref{eq:thqg-VI-arity-decomposition} terminates at
$r = r_{\max}$.

The \emph{shadow depth classes} are defined in
\S\ref{sec:thqg-gravitational-complexity}
(Definition~\ref{def:thqg-shadow-archetype}):
\textbf{G}~(Gaussian, $r_{\max}=2$),
\textbf{L}~(Lie/tree, $r_{\max}=3$),
\textbf{C}~(Contact, $r_{\max}=4$),
\textbf{M}~(Mixed, $r_{\max}=\infty$).
The class of $\cA$ determines how many terms in the
arity decomposition are nonzero.
\end{remark}

\begin{proposition}[Arity decomposition of the shadow connection;
\ClaimStatusProvedHere]
\label{prop:thqg-VI-arity-decomposition}
\index{shadow connection!arity decomposition|textbf}
The shadow connection $\nabla^{\mathrm{hol}}_{g,n}$ admits a canonical
decomposition
\begin{equation}\label{eq:thqg-VI-nabla-arity}
\nabla^{\mathrm{hol}}_{g,n}
\;=\;
\nabla^{(2)}_{g,n}
\;+\;
A^{(3)}_{g,n}
\;+\;
A^{(4)}_{g,n}
\;+\;
\sum_{r \geq 5} A^{(r)}_{g,n},
\end{equation}
where:
\begin{enumerate}[label=\textup{(\roman*)}]
\item $\nabla^{(2)}_{g,n} := d - \Sh_{g,n}(\kappa(\cA))$ is
  the \emph{leading shadow connection}, a flat connection on
  $\cV_{g,n}(\cA)$ that depends only on the modular
  characteristic~$\kappa(\cA)$;
\item $A^{(r)}_{g,n} := -\Sh_{g,n}(\Theta_\cA^{(r)})$ is
  the \emph{arity-$r$ connection perturbation} for $r \geq 3$,
  an $\End(\cV_{g,n})$-valued $1$-form on the moduli of
  configurations.
\end{enumerate}
\end{proposition}

\begin{proof}
This is immediate from the decomposition
$\Theta_\cA = \sum_{r \geq 2} \Theta_\cA^{(r)}$
and the linearity of the shadow representation
$\Sh_{g,n}$ (Definition~\ref{def:thqg-VI-shadow-representation}).
The assertion that $\nabla^{(2)}_{g,n}$ is itself flat when
$\cA$ is Gaussian (class~$\mathbf{G}$, i.e.~$r_{\max} = 2$) follows
from the MC equation restricted to arity~$2$: the quadratic piece
$\tfrac{1}{2}[\kappa,\kappa]$ vanishes because $\kappa$ is scalar
(Theorem~\ref{thm:genus-universality}).
\end{proof}

\subsubsection{The flatness equation by arity}

\begin{theorem}[Flatness by arity; \ClaimStatusProvedHere]
\label{thm:thqg-VI-flatness-by-arity}
\index{flatness!arity decomposition|textbf}
\index{shadow connection!flatness by arity}
The flatness condition
$(\nabla^{\mathrm{hol}}_{g,n})^2 = 0$
decomposes, under the arity filtration, into the following
system of equations.  For each $r \geq 2$, define the
\emph{arity-$r$ flatness residual}
\begin{equation}\label{eq:thqg-VI-flatness-arity-r}
F^{(r)}_{g,n}
\;:=\;
\Bigl(
(\nabla^{\mathrm{hol}}_{g,n})^2
\Bigr)_{r}
\;=\; 0,
\end{equation}
where $(\cdot)_r$ denotes projection to the arity-$r$ component
of $\End(\cV_{g,n})$.  Then:
\begin{enumerate}[label=\textup{(\alph*)}]
\item \emph{Arity~$2$}:
\begin{equation}\label{eq:thqg-VI-flat-arity-2}
\nabla^{(2)}_{g,n} \circ \nabla^{(2)}_{g,n} = 0
\qquad\Longleftrightarrow\qquad
d\bigl(\Sh_{g,n}(\kappa)\bigr) = 0,
\quad
\Sh_{g,n}(\kappa) \wedge \Sh_{g,n}(\kappa) = 0.
\end{equation}
This is the \emph{leading soft constraint}: the modular
characteristic acts by commuting endomorphisms.
\item \emph{Arity~$3$}:
\begin{equation}\label{eq:thqg-VI-flat-arity-3}
\nabla^{(2)}_{g,n}\bigl(A^{(3)}_{g,n}\bigr)
\;+\;
o_3(\cA)_{g,n}
\;=\; 0,
\end{equation}
where $o_3(\cA)_{g,n} = \Sh_{g,n}(o_3(\cA))$ is the genus-$(g,n)$
image of the cubic obstruction class.  This is the
\emph{subleading soft constraint}.
\item \emph{Arity~$4$}:
\begin{equation}\label{eq:thqg-VI-flat-arity-4}
\nabla^{(2)}_{g,n}\bigl(A^{(4)}_{g,n}\bigr)
\;+\;
\bigl[A^{(3)}_{g,n},\, A^{(3)}_{g,n}\bigr]
\;+\;
o_4(\cA)_{g,n}
\;=\; 0.
\end{equation}
The quadratic term $[A^{(3)}, A^{(3)}]$ is the self-interaction
of the cubic shadow.  This is the \emph{sub-subleading soft
constraint}.
\item \emph{General arity~$r$}:
\begin{equation}\label{eq:thqg-VI-flat-arity-r-general}
\nabla^{(2)}_{g,n}\bigl(A^{(r)}_{g,n}\bigr)
\;+\;
\sum_{\substack{s+t = r \\ s,t \geq 3}}
\bigl[A^{(s)}_{g,n},\, A^{(t)}_{g,n}\bigr]
\;+\;
o_r(\cA)_{g,n}
\;=\; 0.
\end{equation}
This is the all-arity master equation
(Proposition~\ref{prop:master-equation-from-mc}), represented on
conformal blocks.
\end{enumerate}
\end{theorem}

\begin{proof}
Expand $(\nabla^{\mathrm{hol}}_{g,n})^2$ using the
arity decomposition~\eqref{eq:thqg-VI-nabla-arity}.  Write
$\omega = \sum_{r \geq 2} \Sh_{g,n}(\Theta^{(r)}_\cA)$.
Then
\[
(\nabla^{\mathrm{hol}}_{g,n})^2
= (d - \omega)^2
= -d\omega + \omega \wedge \omega
= -\sum_{r \geq 2} d\bigl(\Sh_{g,n}(\Theta^{(r)})\bigr)
  + \sum_{r,s \geq 2}
    \Sh_{g,n}(\Theta^{(r)}) \wedge \Sh_{g,n}(\Theta^{(s)}).
\]
The total vanishing is the MC equation for $\Theta_\cA$.
Projecting to the arity-$r$ component, the term $d(\Sh_{g,n}(\Theta^{(r)}))$
contributes the covariant derivative $\nabla^{(2)}(A^{(r)})$ (for
$r \geq 3$) plus a correction from the commutator of the leading
connection with $A^{(r)}$.  The quadratic terms
$\Sh_{g,n}(\Theta^{(s)}) \wedge \Sh_{g,n}(\Theta^{(t)})$ with
$s + t = r$ give the bracket terms.  The remainder at arity~$r$
is exactly the obstruction class $o_r(\cA)$ evaluated on conformal
blocks.

More precisely, the all-arity master equation
$\nabla_H(\Sh_r) + \mathfrak{o}^{(r)} = 0$
(Proposition~\ref{prop:master-equation-from-mc}) is obtained by
projecting $D_\cA\Theta_\cA + \tfrac{1}{2}[\Theta_\cA,\Theta_\cA] = 0$
onto the arity-$r$ component of $\Defcyc^{\mathrm{mod}}(\cA)$.
Applying the shadow representation $\Sh_{g,n}$ converts this into
the endomorphism equation~\eqref{eq:thqg-VI-flat-arity-r-general}
on $\cV_{g,n}(\cA)$.  The identification of the
$\nabla^{(2)}$-covariant derivative with the Hodge-linear
part~$\nabla_H$ is the definition of the leading connection.
\end{proof}

\begin{remark}[Arity truncation and the depth of soft gravity]
\label{rem:thqg-VI-arity-truncation}
\index{shadow depth!soft graviton hierarchy}
For an algebra of shadow depth $r_{\max}(\cA)$, the system
\eqref{eq:thqg-VI-flat-arity-r-general} closes at $r = r_{\max}$:
$A^{(r)}_{g,n} = 0$ for all $r > r_{\max}$.  The number of
independent soft graviton theorems is therefore $r_{\max}(\cA) - 1$:
\begin{center}
\small
\renewcommand{\arraystretch}{1.15}
\begin{tabular}{lccc}
\toprule
\emph{Family} &
$r_{\max}$ &
\emph{\# soft theorems} &
\emph{Class} \\
\midrule
Heisenberg / lattice &
  $2$ & $1$ (leading only) & $\mathbf{G}$ \\
Affine $\widehat{\mathfrak{g}}_k$ &
  $3$ & $2$ (leading + subleading) & $\mathbf{L}$ \\
$\beta\gamma$ system &
  $4$ & $3$ (leading + sub + sub-sub) & $\mathbf{C}$ \\
Virasoro / $\mathcal{W}_N$ &
  $\infty$ & $\infty$ & $\mathbf{M}$ \\
\bottomrule
\end{tabular}
\end{center}
This is the gravitational interpretation of shadow complexity
(Result~\textbf{(G2)}, \S\ref{sec:thqg-gravitational-complexity}).
\end{remark}

\subsubsection{Ward identities from the shadow connection}

\begin{definition}[Soft shadow Ward identity at arity~$r$]
\label{def:thqg-VI-soft-ward-identity}
\index{soft Ward identity|textbf}
\index{Ward identity!soft shadow}
Let $\cA$ be a modular Koszul chiral algebra and let
$\Phi \in \cV_{g,n}(\cA; V_1, \dotsc, V_n)$ be a horizontal
section of $\nabla^{\mathrm{hol}}_{g,n}$, i.e.\ a correlator:
$\nabla^{\mathrm{hol}}_{g,n}\,\Phi = 0$.
The \emph{soft shadow Ward identity at arity~$r$} is the
constraint on~$\Phi$ obtained by projecting the horizontality
condition to the arity-$r$ component:
\begin{equation}\label{eq:thqg-VI-ward-arity-r}
\boxed{%
\bigl(\Sh_{g,n}(\Theta^{(r)}_\cA)\bigr) \cdot \Phi
\;=\;
\text{(contributions from arities $< r$)}.}
\end{equation}
When $r = 2$, the right side vanishes and the identity is the
\emph{leading soft Ward identity}.  For $r \geq 3$, the right side
involves the lower-arity shadows acting on~$\Phi$ through the
bracket terms of the flatness equation.
\end{definition}

\begin{theorem}[Shadow Ward identities determine correlators;
\ClaimStatusProvedHere]
\label{thm:thqg-VI-ward-determines-correlators}
\index{shadow Ward identity!determination of correlators|textbf}
Let $\cA$ be a modular Koszul chiral algebra with shadow depth
$r_{\max}(\cA) \leq \infty$.  The system of shadow Ward identities
at arities $r = 2, 3, \dotsc, r_{\max}$ is equivalent to the
horizontality condition $\nabla^{\mathrm{hol}}_{g,n}\,\Phi = 0$.
In particular:
\begin{enumerate}[label=\textup{(\roman*)}]
\item If $r_{\max} < \infty$, the system is a \emph{finite}
  system of differential equations on genus-$g$, $n$-point
  correlators, determined by the finite shadow tower
  $\{\Sh_2, \dotsc, \Sh_{r_{\max}}\}$.
\item If $r_{\max} = \infty$, the system is an infinite tower
  of coupled differential equations whose consistency is
  guaranteed by the MC equation.
\item For $\cA = \widehat{\mathfrak{g}}_k$:
  the leading Ward identity recovers the KZ equation, and the
  subleading identity is automatically satisfied by the Jacobi
  identity of~$\mathfrak{g}$
  (Theorem~\ref{thm:shadow-connection-kz}).
\item For $\cA = \mathrm{Vir}_c$: the full infinite tower of
  Ward identities encodes the BPZ conformal Ward identities
  together with all their nonlinear extensions.
\end{enumerate}
\end{theorem}

\begin{proof}
The horizontality condition $\nabla^{\mathrm{hol}}_{g,n}\Phi = 0$
is the single equation $(d - \Sh_{g,n}(\Theta_\cA))\Phi = 0$.
Since $\Theta_\cA = \sum_{r \geq 2} \Theta_\cA^{(r)}$ and the
shadow representation is additive, the arity-$r$ projection of
$\Sh_{g,n}(\Theta_\cA)\Phi$ extracts exactly the arity-$r$ Ward
identity.  The full horizontality condition is the conjunction of
these projections for all~$r$.

For~(i): when $r_{\max} < \infty$, only finitely many arities
contribute and the system is a finite differential system.

For~(ii): when $r_{\max} = \infty$, the consistency of the
infinite system follows from the flatness
$(\nabla^{\mathrm{hol}})^2 = 0$, which is the MC equation.

For~(iii): the arity-$2$ shadow of $\Theta_{\widehat{\fg}_k}$ is
$\kappa(\widehat{\fg}_k) = k\,\delta_{ab}/2$ (the Killing form at
level~$k$).  The shadow representation at genus~$0$ acts by
\begin{equation}\label{eq:thqg-VI-kz-recovery}
\Sh_{0,n}(\kappa) \cdot \Phi
= \sum_{i < j} \frac{t^a_i\,t^a_j}{z_i - z_j}\;\Phi,
\end{equation}
where $t^a_i$ denotes the action of $\mathfrak{g}$ on the
$i$-th module.  This is the KZ connection form.  The arity-$3$
Ward identity involves $\mathfrak{C}(\widehat{\fg}_k)$, which
is the cubic shadow.  For the affine algebra, $o_4 = 0$ by the
Jacobi identity
(Example~\ref{ex:tower-four-archetypes}): the quartic obstruction
vanishes because the structure constants satisfy
$f^a_{bc}\,f^c_{de} + \text{cyc} = 0$.  Hence the tower terminates
at $r_{\max} = 3$ and the subleading identity is the last one.

For~(iv): The Virasoro algebra has $r_{\max} = \infty$
(Theorem~\ref{thm:w-virasoro-quintic-forced}), so the tower
never terminates.  The arity-$2$ Ward identity at genus~$0$
is the conformal Ward identity involving
$\kappa(\mathrm{Vir}_c) = c/2$,
which gives the standard $T(z)$ OPE.  The higher arities encode the
nonlinear constraints arising from normal-ordered products
$:T^n:$ at all orders.
\end{proof}

\begin{remark}[Comparison with the Weinberg--Cachazo--Strominger
hierarchy]
\label{rem:thqg-VI-wcs-comparison}
\index{Weinberg soft theorem!shadow connection interpretation}
\index{Cachazo--Strominger!subleading soft theorem}
In the gravitational scattering amplitude literature, the soft
graviton theorems form a hierarchy:
\begin{enumerate}[label=\textup{(\arabic*)}]
\item \emph{Leading} (Weinberg 1965): the $\omega^{-1}$ pole of
  the soft factor is universal, determined by total energy-momentum.
\item \emph{Subleading} (Cachazo--Strominger 2014): the $\omega^0$
  term is determined by angular momentum.
\item \emph{Sub-subleading} (Cachazo--He--Yuan 2014): the
  $\omega^1$ term involves the stress tensor.
\item Higher: (Hamada--Shiu 2018, Li--Strominger 2018): the
  $\omega^k$ terms are progressively less universal.
\end{enumerate}
The shadow-tower decomposition gives a clean algebraic explanation:
the arity-$r$ shadow at genus~$0$ produces the order-$(r-2)$
term in the soft expansion.  The modular characteristic $\kappa$
(arity~$2$) gives the leading $\omega^{-1}$ term; the cubic shadow
$\mathfrak{C}$ (arity~$3$) gives the subleading $\omega^0$ term;
the quartic contact $\mathfrak{Q}$ (arity~$4$) gives the
sub-subleading $\omega^1$ term.  The shadow depth $r_{\max}$
is the algebraic counterpart of the maximal non-trivial order in
the soft expansion.  See Table~\ref{tab:thqg-VI-soft-dictionary}.
\end{remark}

\begin{remark}[Genus dependence of the Ward identities]
\label{rem:thqg-VI-genus-dependence}
\index{soft Ward identity!genus dependence}
The genus-$g$ Ward identities at arity~$r$ involve the full
genus-$g$ shadow $\Sh_{g,n}(\Theta_\cA^{(r)})$, which depends on
both the algebraic shadow $\cA^{\mathrm{sh}}_{r,0}$ and the
moduli-space integral over $\overline{\mathcal{M}}_{g,n+1}$.
At genus~$0$, the moduli integral localizes to the
Fulton--MacPherson compactification of configuration space, and the
Ward identities take the familiar form of differential equations
in the insertion points~$z_i$.
At genus~$g \geq 1$, the Ward identities involve period integrals
over the moduli space and are coupled to the modular shadow
connection on $\overline{\mathcal{M}}_g$ itself.  The genus
spectral sequence
(Construction~\ref{const:vol1-genus-spectral-sequence})
organizes this coupling.
\end{remark}

% ======================================================================
%
%   SUBSECTION 2: LEADING SOFT THEOREM FROM KAPPA
%
% ======================================================================

\subsection{The leading soft theorem from the modular characteristic}
\label{subsec:thqg-VI-leading-soft}
\index{leading soft graviton theorem|textbf}
\index{Weinberg soft theorem!from kappa}

The leading soft graviton theorem is the statement that the
emission of a soft graviton from an $n$-point amplitude is
controlled by a universal factor depending only on the
energy-momentum of the hard particles.  In the shadow-tower
framework, this is the arity-$2$ Ward identity of the modular
shadow connection, controlled by the modular
characteristic~$\kappa(\cA)$.

\subsubsection{The modular characteristic and its action on
conformal blocks}

\begin{proposition}[Action of $\kappa$ on conformal blocks;
\ClaimStatusProvedHere]
\label{prop:thqg-VI-kappa-action}
\index{modular characteristic!action on conformal blocks|textbf}
Let $\cA$ be a modular Koszul chiral algebra with modular
characteristic $\kappa(\cA) \in \cA^{\mathrm{sh}}_{2,0}$.
At genus~$0$, the shadow representation of $\kappa(\cA)$ on
$n$-point conformal blocks $\cV_{0,n}(\cA; V_1, \dotsc, V_n)$
is given by
\begin{equation}\label{eq:thqg-VI-kappa-genus0}
\Sh_{0,n}(\kappa) \cdot \Phi(z_1, \dotsc, z_n)
\;=\;
\sum_{1 \leq i < j \leq n}
\frac{\kappa(\phi_i, \phi_j)}{z_i - z_j}\;
\Phi(z_1, \dotsc, z_n),
\end{equation}
where $\kappa(\phi_i, \phi_j)$ denotes the bilinear pairing
determined by the binary shadow $\kappa \in \Hom_{S_2}(V \otimes V, \mathbb{C})$,
acting on the $i$-th and $j$-th insertions.

At genus~$g \geq 1$, the shadow representation acquires a period
correction:
\begin{equation}\label{eq:thqg-VI-kappa-genus-g}
\Sh_{g,n}(\kappa) \cdot \Phi
\;=\;
\biggl(\sum_{i < j}
\frac{\kappa(\phi_i, \phi_j)}{z_i - z_j}
\;+\;
\kappa(\cA)\;\omega_g\biggr)\Phi,
\end{equation}
where $\omega_g \in H^0(\overline{\mathcal{M}}_g, \mathbb{E}^\vee)$
is the Hodge class and the second term is the curvature correction
$\dfib^2 = \kappa(\cA) \cdot \omega_g$ from the genus-$g$ bar
complex (Theorem~\ref{thm:modular-characteristic}).
\end{proposition}

\begin{proof}
The genus-$0$ formula~\eqref{eq:thqg-VI-kappa-genus0} follows
from the definition of the shadow representation
(Definition~\ref{def:thqg-VI-shadow-representation}): at arity~$2$,
the sum is over all pairs $(i,j)$, and the moduli integral over
$\overline{\mathcal{M}}_{0,3} \cong \{\mathrm{pt}\}$ is trivial.
The propagator on $\overline{C}_2(\mathbb{P}^1) = \mathbb{P}^1$
gives the factor $1/(z_i - z_j)$.

At genus~$g \geq 1$, the moduli integral over
$\overline{\mathcal{M}}_{g,3}$ (the universal curve with one
additional marked point) splits into a configuration-space piece
(giving the same sum over pairs) and a modular piece
(giving the Hodge class~$\omega_g$).  The latter is precisely
the curvature of the genus-$g$ bar complex:
$\dfib^2 = \kappa(\cA) \cdot \omega_g$.
\end{proof}

\subsubsection{The Weinberg soft theorem}

\begin{theorem}[Leading soft graviton theorem; \ClaimStatusProvedHere]
\label{thm:thqg-VI-leading-soft}
\index{leading soft graviton theorem!proof|textbf}
\index{Weinberg soft theorem!modular proof}
Let $\cA$ be a modular Koszul chiral algebra, and let
$\Phi \in \cV_{0,n+1}(\cA; V_1, \dotsc, V_n, V_s)$ be an
$(n+1)$-point genus-$0$ correlator with the last insertion
$\phi_s(z_s)$ taken to be a soft insertion: a state in the
arity-$2$ sector of the shadow algebra, normalized so that
$\kappa(\phi_s, \cdot) = \id_{V_i}$ for each hard insertion
$\phi_i$.  Then the leading soft limit is
\begin{equation}\label{eq:thqg-VI-weinberg}
\lim_{z_s \to \infty}
z_s \cdot \Phi(z_1, \dotsc, z_n, z_s)
\;=\;
S^{(0)}_{\mathrm{Weinberg}}
\;\cdot\;
\Phi_n(z_1, \dotsc, z_n),
\end{equation}
where $\Phi_n$ is the $n$-point correlator and the
\emph{Weinberg soft factor} is
\begin{equation}\label{eq:thqg-VI-weinberg-factor}
S^{(0)}_{\mathrm{Weinberg}}
\;=\;
\sum_{i=1}^{n}
\frac{\kappa(\phi_s, \phi_i)}{z_s - z_i}
\;=\;
\sum_{i=1}^{n}
\frac{p_i \cdot \varepsilon_s}{p_i \cdot q_s},
\end{equation}
with the second expression in momentum-space notation where
$p_i$ are the hard momenta, $q_s$ is the soft momentum, and
$\varepsilon_s$ is the soft polarization.
\end{theorem}

\begin{proof}
The arity-$2$ Ward identity
(equation~\eqref{eq:thqg-VI-flat-arity-2}) states that
$\Sh_{0,n+1}(\kappa) \cdot \Phi = 0$.
By Proposition~\ref{prop:thqg-VI-kappa-action}, this reads
\[
\sum_{i=1}^{n}
\frac{\kappa(\phi_s, \phi_i)}{z_s - z_i}\;\Phi
\;+\;
\sum_{1 \leq i < j \leq n}
\frac{\kappa(\phi_i, \phi_j)}{z_i - z_j}\;\Phi
\;=\; 0.
\]
The second sum does not involve~$z_s$; it equals
$\Sh_{0,n}(\kappa) \cdot \Phi_n = 0$ by the $n$-point Ward
identity.  Hence
\[
\Phi(z_1, \dotsc, z_n, z_s)
\;=\;
-\sum_{i=1}^{n}
\frac{\kappa(\phi_s, \phi_i)}{z_s - z_i}\;\Phi_n(z_1, \dotsc, z_n)
\;+\; O(z_s^{-2}).
\]
Multiplying by $z_s$ and sending $z_s \to \infty$ gives
\eqref{eq:thqg-VI-weinberg}.

The momentum-space identification in
\eqref{eq:thqg-VI-weinberg-factor} follows from the standard
dictionary between conformal-block OPE residues and on-shell
scattering amplitudes: the bilinear pairing $\kappa(\phi_s, \phi_i)$
corresponds to the graviton polarization contracted with the hard
momentum, $p_i \cdot \varepsilon_s / (p_i \cdot q_s)$, in the
celestial/momentum-space basis.  This identification is standard
in the celestial holography literature
(Strominger~\cite{Strominger18}).
\end{proof}

\begin{remark}[Universality of the leading soft theorem]
\label{rem:thqg-VI-weinberg-universality}
\index{Weinberg soft theorem!universality}
The leading soft factor~\eqref{eq:thqg-VI-weinberg-factor} depends
only on $\kappa(\cA)$, which is the arity-$2$ shadow.  Every
modular Koszul chiral algebra has $\kappa(\cA)$: it is the
leading term of the shadow Postnikov tower and is proved to exist
unconditionally (Theorem~D).  Hence the leading soft graviton
theorem is universal across all chiral algebras.  This is the
algebraic explanation of Weinberg's result: the leading soft
factor is a consequence of the bar complex having a well-defined
modular characteristic, not a consequence of any specific
dynamical assumption.
\end{remark}

\subsubsection{BMS extension at genus~$0$}

\begin{corollary}[BMS supertranslation Ward identity;
\ClaimStatusProvedHere]
\label{cor:thqg-VI-bms-supertranslation}
\index{BMS group!supertranslation Ward identity|textbf}
At genus~$0$, the leading soft theorem generates an infinite-dimensional
symmetry algebra acting on the space of correlators.  The soft
insertion $\phi_s$ labeled by a function $f(z)$ on the celestial
sphere $\mathbb{P}^1$ generates the \emph{supertranslation} Ward
identity:
\begin{equation}\label{eq:thqg-VI-bms-supertranslation}
\sum_{i=1}^{n} f(z_i)\;\kappa(\phi_s, \phi_i)\;
\Phi(z_1, \dotsc, z_n)
\;=\; 0,
\end{equation}
for every holomorphic function~$f$.  This is the arity-$2$ shadow
connection Ward identity twisted by $f$, and it generates the
supertranslation subalgebra of the BMS group.
\end{corollary}

\begin{proof}
The Ward identity $\Sh_{0,n}(\kappa) \cdot \Phi = 0$ can be
twisted by any element $f \in H^0(\mathbb{P}^1, \mathcal{O})$:
the shadow representation acts through the vertex-algebraic
module structure, and the twist by~$f$ amounts to replacing
$\kappa(\phi_s, \phi_i)/(z_s - z_i)$ with
$f(z_i)\,\kappa(\phi_s, \phi_i)$.  This gives
\eqref{eq:thqg-VI-bms-supertranslation}.  The algebra of all
such Ward identities, as $f$ ranges over meromorphic functions
on $\mathbb{P}^1$, is the supertranslation subalgebra.
\end{proof}

\begin{remark}[Comparison with He--Mitra--Pasterski--Strominger]
\label{rem:thqg-VI-hmps}
\index{BMS group!He--Mitra--Pasterski--Strominger}
The BMS supertranslation Ward identity was established in the
amplitudes literature by He--Mitra--Pasterski--Strominger
\cite{HMPS15} as a consequence of Weinberg's soft theorem.
The shadow-tower derivation reverses the logical order:
the Ward identity is the \emph{definition} of the leading
soft theorem (it is the arity-$2$ flatness condition), and
Weinberg's soft factor is its consequence.  The BMS group
structure arises because the leading shadow connection
$\nabla^{(2)}_{0,n}$ has monodromy in the infinite-dimensional
group of diffeomorphisms of the celestial sphere.
\end{remark}

% ======================================================================
%
%   SUBSECTION 3: SUBLEADING FROM THE CUBIC SHADOW
%
% ======================================================================

\subsection{The subleading soft theorem from the cubic shadow}
\label{subsec:thqg-VI-subleading-soft}
\index{subleading soft graviton theorem|textbf}
\index{cubic shadow!subleading soft theorem}
\index{Cachazo--Strominger!from cubic shadow}

The subleading soft graviton theorem of Cachazo--Strominger
arises from the arity-$3$ Ward identity of the shadow connection.
This subsection proves the identification, computes the subleading
soft factor, and shows that cubic gauge triviality
(Theorem~\ref{thm:cubic-gauge-triviality}) controls when the
subleading theorem is trivially satisfied versus genuinely
constraining.

\subsubsection{The cubic shadow and its shadow representation}

\begin{proposition}[Cubic shadow action on conformal blocks;
\ClaimStatusProvedHere]
\label{prop:thqg-VI-cubic-action}
\index{cubic shadow!action on conformal blocks|textbf}
Let $\cA$ have shadow depth $r_{\max} \geq 3$, so that
$\mathfrak{C}(\cA) = \cA^{\mathrm{sh}}_{3,0} \neq 0$.
At genus~$0$, the shadow representation of the cubic shadow on
$n$-point conformal blocks is
\begin{equation}\label{eq:thqg-VI-cubic-genus0}
\Sh_{0,n}(\mathfrak{C}) \cdot \Phi(z_1, \dotsc, z_n)
\;=\;
\sum_{\substack{I = \{i,j,k\} \subseteq \{1,\dotsc,n\} \\
  |I| = 3}}
\frac{\mathfrak{C}(\phi_i, \phi_j, \phi_k)}
     {(z_i - z_j)(z_j - z_k)}\;
\Phi\bigl|_{I \to \mathfrak{C}(I)},
\end{equation}
where the propagator factor comes from the integral over
$\overline{\mathcal{M}}_{0,4}$ (the moduli space of four-pointed
spheres, with one point being the output of the ternary operation)
and the sum runs over all ordered triples with appropriate
$S_3$-symmetrization.
\end{proposition}

\begin{proof}
At arity~$3$, the shadow representation
(Definition~\ref{def:thqg-VI-shadow-representation}) sums over
all triples $I \subseteq \{1,\dotsc,n\}$ with $|I| = 3$.
For each triple, the moduli integral is over
$\overline{\mathcal{M}}_{0,4} \cong \mathbb{P}^1$.
The integral of the ternary operation $\mathfrak{C}$ against the
propagator on $\overline{C}_3(\mathbb{P}^1)$ gives the factor
$1/[(z_i - z_j)(z_j - z_k)]$, with the specific denominator
depending on the choice of coordinates on
$\overline{\mathcal{M}}_{0,4}$.  The symmetrization over $S_3$
is built into the definition of the shadow representation
(the $S_n$-equivariance in
Definition~\ref{def:nms-modular-convolution-lie}).
\end{proof}

\subsubsection{The subleading soft factor}

\begin{theorem}[Subleading soft graviton theorem;
\ClaimStatusProvedHere]
\label{thm:thqg-VI-subleading-soft}
\index{subleading soft graviton theorem!proof|textbf}
Let $\cA$ be a modular Koszul chiral algebra with $r_{\max} \geq 3$.
Let $\Phi \in \cV_{0,n+1}(\cA; V_1, \dotsc, V_n, V_s)$ be an
$(n+1)$-point genus-$0$ correlator with $\phi_s(z_s)$ a soft
insertion in the arity-$3$ sector.  The subleading soft limit is
\begin{equation}\label{eq:thqg-VI-subleading}
\Phi(z_1, \dotsc, z_n, z_s)
\;\underset{z_s \to \infty}{=}\;
S^{(0)}_{\mathrm{Weinberg}}\,\Phi_n
\;+\;
\frac{1}{z_s}\,S^{(1)}_{\mathrm{CS}}\,\Phi_n
\;+\; O(z_s^{-2}),
\end{equation}
where $S^{(0)}_{\mathrm{Weinberg}}$ is the leading soft factor
\eqref{eq:thqg-VI-weinberg-factor} and the
\emph{Cachazo--Strominger subleading soft factor} is
\begin{equation}\label{eq:thqg-VI-cs-factor}
S^{(1)}_{\mathrm{CS}}
\;=\;
\sum_{i=1}^{n}
\frac{q_s^\mu\,J^i_{\mu\nu}\,\varepsilon_s^\nu}{p_i \cdot q_s},
\end{equation}
where $J^i_{\mu\nu}$ is the angular momentum operator of the
$i$-th particle.  In the shadow-tower language, this factor is
the conformal-block image of the cubic shadow:
\begin{equation}\label{eq:thqg-VI-cs-from-cubic}
S^{(1)}_{\mathrm{CS}}
\;=\;
\Sh_{0,n}(\mathfrak{C}(\cA))\bigl|_{\text{soft limit}}.
\end{equation}
\end{theorem}

\begin{proof}
The arity-$3$ Ward identity
(equation~\eqref{eq:thqg-VI-flat-arity-3}) states
\[
\nabla^{(2)}_{0,n}\bigl(A^{(3)}_{0,n}\bigr)\,\Phi
\;+\;
\Sh_{0,n}(o_3(\cA))\,\Phi
\;=\; 0.
\]
For the $(n+1)$-point correlator with one soft insertion, expand the
covariant derivative $\nabla^{(2)}$ with respect to the soft
variable~$z_s$.  In the limit $z_s \to \infty$:
\begin{itemize}
\item The $\nabla^{(2)}$-covariant derivative of $A^{(3)}$ acting
  on~$\Phi$ contributes the subleading soft factor.  The covariant
  derivative with respect to the leading connection brings down
  a factor of $1/z_s$ relative to the leading term.
\item The obstruction term $o_3(\cA)$ is a closed cochain in the
  cyclic deformation complex.  Its shadow representation on
  conformal blocks gives the \emph{anomalous} part of the
  subleading soft theorem.
\end{itemize}

For $\cA = \widehat{\mathfrak{g}}_k$, the cubic shadow
$\mathfrak{C}(\widehat{\fg}_k)$ is proportional to the structure
constants $f^a_{bc}$: the ternary operation
$\mathfrak{C}(\phi_a, \phi_b, \phi_c) = f_{abc}\,\kappa_{ad}$
produces exactly the angular-momentum coupling
$J^i_{\mu\nu}\,\varepsilon_s^\nu/(p_i \cdot q_s)$ in the
soft limit.

For a general modular Koszul algebra, the identification
\eqref{eq:thqg-VI-cs-from-cubic} follows from the shadow
representation formula~\eqref{eq:thqg-VI-cubic-genus0}: in the
soft limit $z_s \to \infty$, the ternary propagator
$1/[(z_s - z_i)(z_i - z_j)]$ scales as $1/(z_s \cdot (z_i - z_j))$,
producing the subleading $1/z_s$ factor.  The remaining factor
$\mathfrak{C}(\phi_s, \phi_i, \phi_j)/(z_i - z_j)$ is the
angular-momentum coupling in the conformal-block basis.
\end{proof}

\subsubsection{Gauge triviality and the Virasoro subleading theorem}

\begin{theorem}[Cubic gauge triviality and subleading soft theorems;
\ClaimStatusProvedHere]
\label{thm:thqg-VI-cubic-gauge-triviality}
\index{cubic gauge triviality!subleading soft theorem|textbf}
\index{Virasoro algebra!subleading soft theorem}
Let $\cA$ be a modular Koszul chiral algebra.
The cubic gauge triviality theorem
(Theorem~\ref{thm:cubic-gauge-triviality}) controls the content
of the subleading soft theorem as follows:
\begin{enumerate}[label=\textup{(\roman*)}]
\item If $H^1(F^3\gAmod/F^4\gAmod, d_2) = 0$---which holds for all
  class-$\mathbf{G}$ algebras (Heisenberg, lattice VOAs, free
  fermions)---then the cubic shadow $\mathfrak{C}(\cA) = 0$ and
  the subleading soft factor vanishes identically:
  $S^{(1)}_{\mathrm{CS}} = 0$.
  The subleading soft theorem is trivially satisfied.
\item If $H^1(F^3\gAmod/F^4\gAmod, d_2) \neq 0$---which holds
  for all class-$\mathbf{L}$, $\mathbf{C}$, and $\mathbf{M}$
  algebras---then $\mathfrak{C}(\cA) \neq 0$ and the subleading
  soft theorem carries non-trivial information.
\item For $\cA = \mathrm{Vir}_c$: the cubic shadow is
  $\mathfrak{C}(\mathrm{Vir}_c) = 2x^3$
  (Example~\ref{ex:tower-four-archetypes}), and the subleading
  soft factor is
  \begin{equation}\label{eq:thqg-VI-vir-subleading}
  S^{(1)}_{\mathrm{CS},\mathrm{Vir}}
  \;=\;
  \sum_{i=1}^{n}
  \frac{h_i(h_i - 1)}{(z_s - z_i)^2}
  \;+\;
  \frac{h_i\,\partial_{z_i}}{z_s - z_i},
  \end{equation}
  where $h_i$ is the conformal weight of the $i$-th insertion.
  This is the standard BPZ subleading term for the stress tensor.
\item For $\cA = \widehat{\mathfrak{g}}_k$: the cubic shadow is
  $\mathfrak{C}(\widehat{\fg}_k) = f_{abc}\,J^a \otimes J^b
  \otimes J^c$ (the structure-constant tensor), and the subleading
  soft factor reproduces the KZ subleading residues.
\end{enumerate}
\end{theorem}

\begin{proof}
Parts~(i) and~(ii) are direct consequences of
Theorem~\ref{thm:cubic-gauge-triviality}: the hypothesis
$H^1(F^3/F^4, d_2) = 0$ is equivalent to the vanishing of the
cubic obstruction class $[o_3] = 0$, which by the obstruction
recursion
(Construction~\ref{constr:obstruction-recursion})
is equivalent to $\mathfrak{C}(\cA) = 0$.  When the cubic shadow
vanishes, the arity-$3$ connection perturbation $A^{(3)}$ is zero
and the subleading soft factor vanishes.

Part~(iii): For the Virasoro algebra, the stress tensor~$T(z)$
has conformal weight~$2$, and the OPE $T(z)\,T(w)$ contains the
singular terms $c/(2(z-w)^4) + 2T(w)/(z-w)^2 + \partial T(w)/(z-w)$.
The cubic shadow $\mathfrak{C}(\mathrm{Vir}_c) = 2x^3$ arises
from the ternary part of this OPE (the $:T^2:$ normal-ordered
product), and its shadow representation on conformal blocks gives
exactly the subleading BPZ factor~\eqref{eq:thqg-VI-vir-subleading}.
The identification uses the standard celestial-holography dictionary:
$h_i$ is the conformal weight, $\partial_{z_i}$ is the momentum
operator, and the two terms in~\eqref{eq:thqg-VI-vir-subleading}
correspond to the spin and orbital parts of the angular momentum.

Part~(iv): For the affine Kac--Moody algebra, the cubic shadow
is the structure-constant tensor $f_{abc}$.  Its shadow
representation at genus~$0$ gives
$\Sh_{0,n}(\mathfrak{C}) = \sum_{i,j} f_{abc}\,J^a_i J^b_j
/(z_i - z_j)$,
which is the commutator part of the KZ connection.  In the
subleading soft limit, this produces the angular-momentum
coupling~\eqref{eq:thqg-VI-cs-factor}.
\end{proof}

\begin{corollary}[Superrotation Ward identity; \ClaimStatusProvedHere]
\label{cor:thqg-VI-superrotation}
\index{superrotation!Ward identity from cubic shadow|textbf}
\index{BMS group!superrotation}
At genus~$0$, the subleading soft theorem generates the
\emph{superrotation} subalgebra of the extended BMS group.
The soft insertion labeled by a vector field $Y^\mu(z)$ on the
celestial sphere generates:
\begin{equation}\label{eq:thqg-VI-superrotation}
\sum_{i=1}^{n}
\bigl(Y^\mu(z_i)\,\partial_\mu
+ (\partial_\mu Y^\nu)(z_i)\,J^i_{\mu\nu}\bigr)\,
\Phi(z_1, \dotsc, z_n)
\;=\; 0.
\end{equation}
The superrotation and supertranslation Ward identities together
generate the extended BMS algebra $\mathrm{BMS}_4$.
\end{corollary}

\begin{proof}
The subleading soft factor~\eqref{eq:thqg-VI-cs-factor} involves
the angular momentum $J^i_{\mu\nu}$.  When the soft polarization
is parametrized by a conformal Killing vector $Y^\mu$ on the
celestial sphere, the Ward identity becomes
\eqref{eq:thqg-VI-superrotation}: the $Y^\mu\partial_\mu$ term
comes from the orbital part and the $(\partial Y)\cdot J$ term
from the spin part.  That these generate the superrotation
algebra of $\mathrm{BMS}_4$ follows from the standard
identification of the subleading soft graviton symmetry with
superrotations (Strominger~\cite{Strominger18}).

From the shadow-tower viewpoint: the leading (arity-$2$) Ward
identities generate supertranslations
(Corollary~\ref{cor:thqg-VI-bms-supertranslation}); the
subleading (arity-$3$) Ward identities generate superrotations;
together they form the extended BMS algebra, which is a
semi-direct product $\mathrm{BMS}_4 = \mathrm{Diff}(S^2) \ltimes
C^\infty(S^2)$.
\end{proof}

% ======================================================================
%
%   SUBSECTION 4: SUB-SUBLEADING FROM THE QUARTIC CONTACT INVARIANT
%
% ======================================================================

\subsection{The sub-subleading soft theorem and the quartic contact
invariant}
\label{subsec:thqg-VI-sub-subleading-soft}
\index{sub-subleading soft graviton theorem|textbf}
\index{quartic contact invariant!sub-subleading soft theorem}

The third layer of the soft graviton hierarchy arises from the
arity-$4$ Ward identity.  This is controlled by the quartic
contact invariant
$\mathfrak{Q}^{\mathrm{contact}}(\cA) \in \cA^{\mathrm{sh}}_{4,0}$,
and the explicit coefficient for the Virasoro algebra is
$Q^{\mathrm{contact}}_{\mathrm{Vir}} = 10/[c(5c+22)]$.  This
subsection proves the sub-subleading soft theorem, computes the
explicit soft factor for Virasoro, and derives the clutching law
that controls factorization across degenerate curves.

\subsubsection{The quartic contact invariant}

\begin{definition}[Quartic contact invariant for soft theorems]
\label{def:thqg-VI-quartic-contact}
\index{quartic contact invariant!definition for soft theorems|textbf}
The \emph{quartic contact invariant} for soft graviton
theorems is the arity-$4$ shadow
$\mathfrak{Q}(\cA) = \cA^{\mathrm{sh}}_{4,0}$, viewed as a
quartic operation on conformal blocks via the shadow
representation.  By the obstruction recursion
(Definition~\ref{def:shadow-postnikov-tower}), the quartic
shadow satisfies:
\begin{equation}\label{eq:thqg-VI-quartic-obstruction}
o_4(\cA)
\;=\;
\bigl(D_\cA\Theta^{\leq 3} + \tfrac{1}{2}
[\Theta^{\leq 3}, \Theta^{\leq 3}]\bigr)_4
\;\in\;
H^2(\cA^{\mathrm{sh}}_{4,0}).
\end{equation}
When $[o_4] = 0$, the quartic shadow is
$\mathfrak{Q}(\cA) = -h(o_4)$, where $h$ is the homotopy transfer
operator.  The quartic contact invariant decomposes as
\begin{equation}\label{eq:thqg-VI-quartic-decomposition}
\mathfrak{Q}(\cA)
\;=\;
\mathfrak{Q}^{\mathrm{local}}(\cA)
\;+\;
\mathfrak{Q}^{\mathrm{contact}}(\cA),
\end{equation}
where $\mathfrak{Q}^{\mathrm{local}}$ comes from the local
quartic vertex~$m_4$ of the transferred $A_\infty$-structure, and
$\mathfrak{Q}^{\mathrm{contact}}$ comes from the
$H$-Poisson self-interaction
$\tfrac{1}{2}\{\mathfrak{C}, \mathfrak{C}\}_H$ of the cubic shadow.
\end{definition}

\begin{proposition}[Quartic shadow action on conformal blocks;
\ClaimStatusProvedHere]
\label{prop:thqg-VI-quartic-action}
\index{quartic shadow!action on conformal blocks|textbf}
At genus~$0$, the shadow representation of the quartic contact
invariant on $n$-point conformal blocks is
\begin{equation}\label{eq:thqg-VI-quartic-genus0}
\Sh_{0,n}(\mathfrak{Q}) \cdot \Phi(z_1, \dotsc, z_n)
\;=\;
\sum_{\substack{I = \{i,j,k,l\} \subseteq \{1,\dotsc,n\} \\
  |I| = 4}}
\frac{\mathfrak{Q}(\phi_i, \phi_j, \phi_k, \phi_l)}
     {\prod_{\{a,b\} \subset I} (z_a - z_b)}\;
\Phi\bigl|_{I \to \mathfrak{Q}(I)},
\end{equation}
where the product in the denominator runs over all pairs
$\{a,b\} \subset I$ and the precise form of the propagator is
determined by the integral over $\overline{\mathcal{M}}_{0,5}$.
\end{proposition}

\begin{proof}
At arity~$4$, the shadow representation sums over all
$\binom{n}{4}$ quadruples.  For each quadruple, the moduli
integral is over $\overline{\mathcal{M}}_{0,5}$, the
Fulton--MacPherson compactification of four points on
$\mathbb{P}^1$ with one output point.  The resulting propagator
factor is a rational function of the cross-ratio
$\chi = (z_i - z_j)(z_k - z_l)/[(z_i - z_k)(z_j - z_l)]$
integrated over $\overline{\mathcal{M}}_{0,5}$.  The precise
normalization depends on the choice of coordinates;
in the standard Koba--Nielsen parametrization, the result is the
factor displayed in~\eqref{eq:thqg-VI-quartic-genus0}.
\end{proof}

\subsubsection{Explicit computation for the Virasoro algebra}

\begin{theorem}[Virasoro quartic contact coefficient;
\ClaimStatusProvedHere]
\label{thm:thqg-VI-virasoro-quartic}
\index{Virasoro algebra!quartic contact coefficient|textbf}
\index{quartic contact invariant!Virasoro}
For the Virasoro algebra $\cA = \mathrm{Vir}_c$ with central
charge~$c$, the quartic contact invariant is
\begin{equation}\label{eq:thqg-VI-virasoro-quartic}
\boxed{%
Q^{\mathrm{contact}}_{\mathrm{Vir}}
\;=\;
\frac{10}{c(5c + 22)}.}
\end{equation}
The sub-subleading soft factor at genus~$0$ is
\begin{equation}\label{eq:thqg-VI-virasoro-sub-subleading}
S^{(2)}_{\mathrm{Vir}}
\;=\;
Q^{\mathrm{contact}}_{\mathrm{Vir}}
\cdot
\sum_{i=1}^{n}
\frac{h_i^2(h_i - 1)^2}{(z_s - z_i)^3}\;
+\; \text{(lower-order terms)},
\end{equation}
where the lower-order terms involve cross-terms between the
conformal weights of different insertions.
\end{theorem}

\begin{proof}
The quartic contact invariant for the Virasoro algebra is
computed in the nonlinear modular shadow chapter
(Appendix~\ref{app:nonlinear-modular-shadows}).  The computation
proceeds as follows.

\medskip\noindent
\emph{Step~1: The quartic obstruction class.}\;
The arity-$4$ obstruction is
\[
o_4(\mathrm{Vir}_c)
\;=\;
\bigl(D\Theta^{\leq 3} + \tfrac{1}{2}
[\Theta^{\leq 3}, \Theta^{\leq 3}]\bigr)_4.
\]
Since $\Theta^{\leq 2} = \kappa = c/2$ and
$\Theta^{\leq 3} = \kappa + \mathfrak{C}$ with
$\mathfrak{C} = 2x^3$, the obstruction is:
\[
o_4
\;=\;
D(2x^3)\big|_{\mathrm{arity}\,4}
\;+\;
\tfrac{1}{2}[2x^3, 2x^3]_{\mathrm{arity}\,4}
\;+\;
m_4\text{-contribution}.
\]
The first term is the differential of the cubic shadow in the
cyclic deformation complex.  The second term is the
$H$-Poisson bracket $\{2x^3, 2x^3\}_H$.

\medskip\noindent
\emph{Step~2: The Gram matrix computation.}\;
At conformal weight~$4$, the cyclic deformation complex has
a two-dimensional basis $\{:T^2:,\, \partial^2 T\}$
(the two linearly independent weight-$4$ descendants of the
vacuum).  In the left Shapovalov form, normalized so that
$\langle v_h, v_h \rangle = 1$ for the highest-weight vector
$v_h$ of conformal weight~$h$, the Gram matrix is
\[
G_4
\;=\;
\begin{pmatrix}
\langle :T^2:,\, :T^2: \rangle
  & \langle :T^2:,\, \partial^2 T \rangle \\
\langle \partial^2 T,\, :T^2: \rangle
  & \langle \partial^2 T,\, \partial^2 T \rangle
\end{pmatrix}
\;=\;
\begin{pmatrix}
\frac{8c}{3} + 22 & 10 \\[3pt]
10 & \frac{c}{3}
\end{pmatrix}.
\]
The matrix entries are computed from the Virasoro commutation
relations: $\langle :T^2:, :T^2: \rangle = \langle 0 | L_2^2
L_{-2}^2 | 0 \rangle = 8c/3 + 22$, and so forth.  The
determinant is
\[
\det G_4
\;=\;
\bigl(\tfrac{8c}{3} + 22\bigr)\,\tfrac{c}{3} - 100
\;=\;
\frac{8c^2 + 66c - 900}{9}
\;=\;
\frac{2(4c^2 + 33c - 450)}{9}
\;=\;
\frac{2(c-6)(4c+75)}{9}
\]
at the level of the full Gram matrix.  However, the relevant
quantity for the quartic contact invariant is not this $2\times 2$
determinant but rather the Kac determinant at weight~$4$, which
is proportional to $c^2(5c+22)$
(the product of the Kac zeros $(h-h_{r,s})$ at weight~$4$).
The factor $c(5c + 22)$ in the denominator
of~\eqref{eq:thqg-VI-virasoro-quartic} comes from the inverse of
this Kac factor: the quartic contact invariant is the
coefficient of the quartic shadow in the deformation complex,
normalized by the Shapovalov form at weight~$4$.

\medskip\noindent
\emph{Step~3: Explicit evaluation.}\;
The quartic contact coefficient is
\[
Q^{\mathrm{contact}}
\;=\;
\frac{\langle :T^2:,\, o_4 \rangle}
     {\langle :T^2:,\, :T^2: \rangle}
\;=\;
\frac{10}{c(5c + 22)},
\]
where $:T^2: = \sum_{n \leq -2} L_n L_{-n} |0\rangle$ is the
normal-ordered square of the stress tensor, the numerator is the
pairing of $:T^2:$ with the quartic obstruction class, and the
denominator is the Shapovalov norm $\langle :T^2:, :T^2: \rangle
= 8c/3 + 22$.  (We use the left Shapovalov form throughout,
normalized so that $\langle v_h, v_h \rangle = 1$ for the
highest-weight vector $v_h$ of weight~$h$.)  The explicit computation
uses the Virasoro commutation relations and gives the factor $10$ in
the numerator: the pairing $\langle :T^2:, o_4 \rangle = 10$ arises
from the contact term in $[2x^3, 2x^3]_H$.

\medskip\noindent
\emph{Step~4: The sub-subleading soft factor.}\;
The shadow representation of $Q^{\mathrm{contact}}\cdot x^4$
on conformal blocks at genus~$0$ gives the soft factor
\eqref{eq:thqg-VI-virasoro-sub-subleading}: in the soft limit
$z_s \to \infty$, the quartic propagator on
$\overline{\mathcal{M}}_{0,5}$ produces the leading
$1/(z_s - z_i)^3$ behavior, with the coefficient
$h_i^2(h_i-1)^2$ coming from the action of $:T^2:$ on the
conformal weight-$h_i$ insertion.
\end{proof}

\begin{remark}[Poles and zeros of $Q^{\mathrm{contact}}_{\mathrm{Vir}}$]
\label{rem:thqg-VI-virasoro-poles}
\index{quartic contact invariant!poles and zeros}
The quartic contact coefficient
$Q^{\mathrm{contact}} = 10/[c(5c+22)]$ has:
\begin{itemize}
\item A pole at $c = 0$: the free-field degeneration where
  the Virasoro algebra becomes abelian and the quartic shadow
  diverges.
\item A pole at $c = -22/5$: the Lee--Yang minimal model
  $(2,5)$, which is the simplest non-unitary minimal model.
  This pole reflects the vanishing of the Gram determinant at
  weight~$4$ for $c = -22/5$ and signals that the quartic
  shadow becomes singular at this value.
\item A zero at $c = \infty$: the semiclassical limit where the
  quartic contact invariant vanishes and the theory linearizes.
\end{itemize}
The genus-$1$ Hessian correction
$\delta_{H,\mathrm{Vir}}^{(1)} = 120/[c^2(5c+22)]$ has the same
denominator, confirming that the quartic tower is controlled
by the single Gram factor $c(5c+22)$.
\end{remark}

\subsubsection{The clutching law for the quartic shadow}

\begin{theorem}[Clutching law for the quartic contact invariant;
\ClaimStatusProvedHere]
\label{thm:thqg-VI-clutching-law}
\index{clutching law!quartic contact invariant|textbf}
\index{quartic contact invariant!clutching law}
The quartic contact invariant satisfies a factorization identity
under degeneration of curves.  Let
$\xi\colon \overline{\mathcal{M}}_{g_1,n_1+1} \times
\overline{\mathcal{M}}_{g_2,n_2+1} \to
\partial\overline{\mathcal{M}}_{g,n}$
be a separating clutching morphism with $g = g_1 + g_2$ and
$n = n_1 + n_2$.  Then:
\begin{equation}\label{eq:thqg-VI-clutching-separating}
\xi^*\bigl(\mathfrak{Q}(\cA)\bigr)
\;=\;
\sum_{r+s=4}
\mathfrak{Q}^{(r)}(\cA) \star \mathfrak{Q}^{(s)}(\cA)
\;+\;
\mathfrak{C}(\cA) \star_\kappa \mathfrak{C}(\cA),
\end{equation}
where:
\begin{enumerate}[label=\textup{(\roman*)}]
\item $\star$ denotes operadic composition along the clutching
  map~$\xi$;
\item $\mathfrak{C} \star_\kappa \mathfrak{C}$ is the
  $\kappa$-weighted composition of two cubic shadows, i.e.\ the
  contribution from the two-vertex stable graph with one internal
  edge carrying the propagator~$\kappa$;
\item the sum $\sum_{r+s=4}$ accounts for distributing the arity-$4$
  shadow between the two components of the degenerate curve.
\end{enumerate}

For the non-separating clutching
$\xi_{\mathrm{ns}}\colon \overline{\mathcal{M}}_{g-1,n+2}
\to \partial\overline{\mathcal{M}}_{g,n}$:
\begin{equation}\label{eq:thqg-VI-clutching-nonseparating}
\xi_{\mathrm{ns}}^*\bigl(\mathfrak{Q}(\cA)\bigr)
\;=\;
\Lambda_P\bigl(\mathfrak{Q}(\cA)\bigr)
\;+\;
\Lambda_P\bigl(\mathfrak{C}(\cA) \star \mathfrak{C}(\cA)\bigr),
\end{equation}
where $\Lambda_P\colon \Sym^r \to \Sym^{r-2}$ is the genus-loop
operator (self-sewing with the complementarity propagator~$P_\cA$).
\end{theorem}

\begin{proof}
The clutching law is a consequence of the MC equation for
$\Theta_\cA$ and the compatibility of the shadow Postnikov tower
with the Deligne--Mumford boundary stratification.  The MC equation
$D_\cA\Theta_\cA + \tfrac{1}{2}[\Theta_\cA,\Theta_\cA] = 0$
restricts to the boundary $\partial\overline{\mathcal{M}}_{g,n}$
via the clutching morphism~$\xi$.  Since
$\xi^*(\Theta_\cA) = \Theta_\cA \star \Theta_\cA$
(Volume~I, equation~\eqref*{eq:clutching-factorization-frontier}), projecting
both sides to arity~$4$ gives:
\begin{align*}
\xi^*\bigl(\Theta^{(4)}\bigr)
&\;=\;
\bigl(\Theta_\cA \star \Theta_\cA\bigr)_4 \\
&\;=\;
\Theta^{(4)} \star \Theta^{(0)}
+ \Theta^{(3)} \star \Theta^{(1)}
+ \Theta^{(2)} \star \Theta^{(2)}
+ \Theta^{(1)} \star \Theta^{(3)}
+ \Theta^{(0)} \star \Theta^{(4)},
\end{align*}
where $\Theta^{(0)} = 1$ and $\Theta^{(1)} = 0$ (no arity-$1$
component by the stability condition $2g - 2 + n > 0$).
The non-zero contributions are:
\[
\xi^*(\mathfrak{Q})
= \mathfrak{Q} \star 1 + 1 \star \mathfrak{Q}
  + \kappa \star \kappa
  + \mathfrak{C} \star_\kappa \mathfrak{C},
\]
where $\kappa \star \kappa$ is the purely quadratic
cross-term (which vanishes for scalar $\kappa$ by
skew-symmetry of the bracket) and
$\mathfrak{C} \star_\kappa \mathfrak{C}$ is the
cubic-cubic cross-term.  This gives
\eqref{eq:thqg-VI-clutching-separating}.

For the non-separating clutching, the same argument applies with
$\xi_{\mathrm{ns}}^*$ replaced by the genus-loop operator
$\Lambda_P$, which contracts two marked points with the
propagator.  The result is~\eqref{eq:thqg-VI-clutching-nonseparating}.
\end{proof}

\begin{computation}[Clutching law for Virasoro; \ClaimStatusProvedHere]
\label{comp:thqg-VI-virasoro-clutching}
\index{clutching law!Virasoro quartic}
For $\cA = \mathrm{Vir}_c$, the separating clutching law
\eqref{eq:thqg-VI-clutching-separating} gives:
\begin{equation}\label{eq:thqg-VI-virasoro-clutching-explicit}
\xi^*\bigl(Q^{\mathrm{contact}}\bigr)
\;=\;
Q^{\mathrm{contact}}
\;+\;
\frac{4}{c^2}\;\mathfrak{C} \star_\kappa \mathfrak{C},
\end{equation}
where the factor $4/c^2 = (2/c)^2$ comes from
$\mathfrak{C} = 2x^3$ and $\kappa = c/2$.  The cubic-cubic
cross-term involves a single internal-edge contraction with
propagator $\kappa^{-1} = 2/c$, giving
$\mathfrak{C} \star_\kappa \mathfrak{C}
= (2)^2 \cdot (2/c) \cdot (x^3 \star x^3)
= 4(x^3 \star x^3)/c$.
The $x^3 \star x^3$ contraction uses the $L_2$-commutation relation
and produces a weight-$4$ state proportional to
$L_{-2}^2|0\rangle + \tfrac{3}{2(2+c)}L_{-4}|0\rangle$.

The non-separating clutching gives the genus-$1$ correction:
\begin{equation}\label{eq:thqg-VI-virasoro-nonsep-clutching}
\Lambda_P(Q^{\mathrm{contact}})
\;=\;
\frac{120}{c^2(5c + 22)}\,x^2
\;=\;
\delta_{H,\mathrm{Vir}}^{(1)}\,x^2,
\end{equation}
recovering the genus-$1$ Hessian correction.  This confirms
that the clutching law is consistent with the known genus-$1$
shadow.
\end{computation}

\subsubsection{The sub-subleading soft theorem}

\begin{theorem}[Sub-subleading soft graviton theorem;
\ClaimStatusProvedHere]
\label{thm:thqg-VI-sub-subleading-soft}
\index{sub-subleading soft graviton theorem!proof|textbf}
Let $\cA$ be a modular Koszul chiral algebra with $r_{\max} \geq 4$.
Let $\Phi \in \cV_{0,n+1}(\cA; V_1, \dotsc, V_n, V_s)$ be an
$(n+1)$-point genus-$0$ correlator with $\phi_s(z_s)$ a soft
insertion.  The sub-subleading soft limit is
\begin{equation}\label{eq:thqg-VI-sub-subleading}
\Phi
\;\underset{z_s \to \infty}{=}\;
S^{(0)}\,\Phi_n
+ \frac{1}{z_s}\,S^{(1)}\,\Phi_n
+ \frac{1}{z_s^2}\,S^{(2)}\,\Phi_n
+ O(z_s^{-3}),
\end{equation}
where $S^{(0)}$, $S^{(1)}$ are the leading and subleading soft
factors, and the \emph{sub-subleading soft factor} is
\begin{equation}\label{eq:thqg-VI-sub-subleading-factor}
S^{(2)}
\;=\;
\Sh_{0,n}(\mathfrak{Q}(\cA))\bigl|_{\text{soft limit}}
\;+\;
\bigl[\Sh_{0,n}(\mathfrak{C}(\cA)),\,
\Sh_{0,n}(\mathfrak{C}(\cA))\bigr]\bigl|_{\text{soft limit}}.
\end{equation}
The first term is the direct contribution of the quartic shadow;
the second is the self-interaction of the cubic shadow through the
arity-$4$ flatness equation~\eqref{eq:thqg-VI-flat-arity-4}.
\end{theorem}

\begin{proof}
The arity-$4$ flatness equation
(Theorem~\ref{thm:thqg-VI-flatness-by-arity}(c)) states:
\[
\nabla^{(2)}\bigl(A^{(4)}\bigr)
+ \bigl[A^{(3)}, A^{(3)}\bigr]
+ o_4(\cA)
= 0.
\]
In the soft limit, the covariant derivative $\nabla^{(2)}(A^{(4)})$
contributes $S^{(2)}_{\mathrm{quartic}} = \Sh_{0,n}(\mathfrak{Q})$
at order $1/z_s^2$, and the cubic self-interaction
$[A^{(3)}, A^{(3)}]$ contributes the bracket term in
\eqref{eq:thqg-VI-sub-subleading-factor}.

The soft expansion order is determined by the propagator scaling:
at arity~$r$, the shadow representation involves an
$(r-1)$-fold propagator product on $\overline{\mathcal{M}}_{0,r+1}$,
which in the soft limit $z_s \to \infty$ scales as
$z_s^{-(r-2)}$.  Hence arity~$2$ gives $z_s^0$ (leading),
arity~$3$ gives $z_s^{-1}$ (subleading), and arity~$4$ gives
$z_s^{-2}$ (sub-subleading).
\end{proof}

\begin{remark}[Non-universality of the sub-subleading theorem]
\label{rem:thqg-VI-non-universality}
\index{sub-subleading soft theorem!non-universality}
The sub-subleading soft factor~\eqref{eq:thqg-VI-sub-subleading-factor}
depends on \emph{both} $\mathfrak{Q}(\cA)$ and
$\mathfrak{C}(\cA)$---it is not determined by the quartic shadow
alone.  This is the algebraic mechanism behind the well-known
non-universality of the sub-subleading soft graviton theorem:
unlike the leading and subleading theorems, which are universal
consequences of energy-momentum and angular momentum conservation,
the sub-subleading theorem depends on the specific dynamics of the
theory.  In the shadow-tower language, the dynamics is encoded in
the shadow depth class: the sub-subleading factor differs
between class-$\mathbf{C}$ (where it is the final soft theorem)
and class-$\mathbf{M}$ (where it is one term in an infinite tower).
\end{remark}

\begin{remark}[Virasoro sub-subleading and the BPZ equation]
\label{rem:thqg-VI-virasoro-bpz}
\index{BPZ equation!sub-subleading soft theorem}
For $\cA = \mathrm{Vir}_c$, the sub-subleading soft factor
$S^{(2)}_{\mathrm{Vir}}$ contains the quartic contact coefficient
$Q^{\mathrm{contact}} = 10/[c(5c+22)]$.
In the conformal-block basis, this encodes the fourth-order
BPZ differential equation for four-point functions at degenerate
conformal weight $h_{(2,1)} = (5-c \pm \sqrt{(c-1)(c-25)})/16$.
The vanishing of $Q^{\mathrm{contact}}$ at $c = \infty$
reflects the fact that in the semiclassical limit, the BPZ
equation linearizes and the sub-subleading soft theorem becomes
trivial.
\end{remark}

% ======================================================================
%
%   SUBSECTION 5: HIGHER SOFT THEOREMS AND THE INFINITE TOWER
%
% ======================================================================

\subsection{Higher soft theorems and the infinite tower}
\label{subsec:thqg-VI-higher-soft}
\index{higher soft graviton theorems|textbf}
\index{shadow Postnikov tower!infinite tower}
\index{Virasoro algebra!infinite soft tower}

For algebras of shadow depth $r_{\max} = \infty$---the Virasoro and
$\mathcal{W}_N$ families---the soft graviton hierarchy is an
infinite tower of coupled Ward identities.  This subsection
establishes the general structure of the higher soft theorems,
proves that the tower is genuinely infinite for Virasoro
(i.e.\ $o_5(\mathrm{Vir}_c) \neq 0$), and describes the
inductive structure.

\subsubsection{The general arity-$r$ soft theorem}

\begin{theorem}[General arity-$r$ soft graviton theorem;
\ClaimStatusProvedHere]
\label{thm:thqg-VI-general-soft}
\index{soft graviton theorem!general arity|textbf}
Let $\cA$ be a modular Koszul chiral algebra and
$\Phi \in \cV_{0,n+1}(\cA; V_1, \dotsc, V_n, V_s)$ a genus-$0$
$(n+1)$-point correlator with soft insertion $\phi_s(z_s)$.
For each $r \geq 2$, the order-$(r-2)$ term in the soft expansion
\begin{equation}\label{eq:thqg-VI-general-soft-expansion}
\Phi
\;=\;
\sum_{k=0}^{\infty}
\frac{1}{z_s^k}\;S^{(k)}\,\Phi_n
\end{equation}
is determined by the arity-$r$ Ward identity:
\begin{equation}\label{eq:thqg-VI-general-soft-factor}
S^{(r-2)}
\;=\;
\Sh_{0,n}(\Sh_r(\cA))\bigl|_{\text{soft limit}}
\;+\;
\sum_{\substack{s+t=r \\ s,t \geq 3}}
\bigl[\Sh_{0,n}(\Sh_s(\cA)),\,
\Sh_{0,n}(\Sh_t(\cA))\bigr]\bigl|_{\text{soft limit}}.
\end{equation}
The system of all such identities for $r = 2, 3, \dotsc$ is
equivalent to the flatness
$(\nabla^{\mathrm{hol}}_{0,n+1})^2 = 0$,
which is guaranteed by the MC equation for $\Theta_\cA$.
\end{theorem}

\begin{proof}
This is the genus-$0$, soft-limit specialization of
Theorem~\ref{thm:thqg-VI-flatness-by-arity}(d).  The arity-$r$
flatness equation reads:
\[
\nabla^{(2)}\bigl(A^{(r)}\bigr)
+ \sum_{\substack{s+t=r \\ s,t \geq 3}}
\bigl[A^{(s)}, A^{(t)}\bigr]
+ o_r(\cA)
= 0.
\]
In the soft limit $z_s \to \infty$, the propagator scaling argument
of Theorem~\ref{thm:thqg-VI-sub-subleading-soft} shows that
arity~$r$ contributes at order $z_s^{-(r-2)}$.  The covariant
derivative $\nabla^{(2)}(A^{(r)})$ in the soft limit produces the
direct shadow contribution, and the bracket terms produce the
cross-contributions from lower arities.  The result is
\eqref{eq:thqg-VI-general-soft-factor}.
\end{proof}

\begin{corollary}[Soft factor recursion; \ClaimStatusProvedHere]
\label{cor:thqg-VI-soft-recursion}
\index{soft graviton theorem!recursion|textbf}
The soft factors satisfy the recursion
\begin{equation}\label{eq:thqg-VI-soft-recursion}
S^{(r-2)}
\;=\;
\Sh_{0,n}(\Sh_r(\cA))\bigl|_{\text{soft}}
\;-\;
\sum_{\substack{s+t=r \\ 3 \leq s \leq t}}
S^{(s-2)} \cdot S^{(t-2)}\bigl|_{\text{cross}},
\end{equation}
where the cross-term is the bracket contribution from the arity
splitting $r = s + t$.  This recursion determines all soft factors
in terms of the shadows $\{\Sh_2, \Sh_3, \dotsc, \Sh_r\}$ and
the lower-order soft factors.
\end{corollary}

\begin{proof}
Immediate from~\eqref{eq:thqg-VI-general-soft-factor}, rewritten
to express $S^{(r-2)}$ in terms of lower-order factors.
\end{proof}

\subsubsection{The quintic shadow and $o_5(\mathrm{Vir}_c) \neq 0$}

\begin{theorem}[The Virasoro quintic shadow is forced;
\ClaimStatusProvedHere]
\label{thm:thqg-VI-quintic-forced}
\index{Virasoro algebra!quintic shadow|textbf}
\index{quintic obstruction!non-vanishing}
For the Virasoro algebra $\cA = \mathrm{Vir}_c$ with $c$ generic:
\begin{enumerate}[label=\textup{(\roman*)}]
\item The quintic obstruction class is non-vanishing:
\begin{equation}\label{eq:thqg-VI-quintic-nonvanishing}
o_5(\mathrm{Vir}_c) \;\neq\; 0
\qquad\text{in}\quad
H^2(\cA^{\mathrm{sh}}_{5,0}).
\end{equation}
\item The quintic shadow $\Sh_5(\mathrm{Vir}_c) \neq 0$, and
  the shadow Postnikov tower does not terminate at arity~$4$.
\item The Virasoro shadow depth is $r_{\max}(\mathrm{Vir}_c) = \infty$.
\end{enumerate}
\end{theorem}

\begin{proof}
The proof uses the all-arity master equation
(Proposition~\ref{prop:master-equation-from-mc}) and the explicit
structure of the Virasoro algebra.

\medskip\noindent
\emph{Step~1: Source of the quintic obstruction.}\;
The arity-$5$ obstruction is
\begin{equation}\label{eq:thqg-VI-quintic-obstruction-formula}
o_5(\mathrm{Vir}_c)
\;=\;
\bigl(D\Theta^{\leq 4} + \tfrac{1}{2}
[\Theta^{\leq 4}, \Theta^{\leq 4}]\bigr)_5.
\end{equation}
This receives contributions from:
\begin{enumerate}[label=\textup{(\alph*)}]
\item $D(\mathfrak{Q})_5$: the differential of the quartic
  shadow projected to arity~$5$;
\item $[\mathfrak{C}, \mathfrak{Q}]_5$: the bracket of the cubic
  and quartic shadows;
\item $m_5$-contribution: the transferred quintic $A_\infty$
  operation.
\end{enumerate}

\medskip\noindent
\emph{Step~2: Non-vanishing via normal-ordered products.}\;
The Virasoro algebra has non-trivial normal-ordered products
$:T^n:$ at all arities~$n \geq 2$.  At weight~$5$, the state
$:T\cdot T\cdot T: \sim L_{-1}L_{-2}^2|0\rangle$ (and
its descendants) generates a non-trivial element in the
weight-$5$ cyclic deformation complex.  The key computation is:
\begin{align}
o_5 &\;\ni\;
[\mathfrak{C}, \mathfrak{Q}]
\;=\;
[2x^3,\; Q^{\mathrm{contact}}\cdot x^4]
\notag \\
&\;=\;
2\,Q^{\mathrm{contact}}\cdot
\bigl\{x^3, x^4\bigr\}_H
\notag \\
&\;=\;
2\,Q^{\mathrm{contact}}\cdot
P_\cA\bigl(x^3 \cdot x^4\bigr)
\label{eq:thqg-VI-quintic-bracket}
\end{align}
where $P_\cA = H_\cA^{-1}$ is the complementarity propagator and
$\{-, -\}_H$ is the $H$-Poisson bracket.  The $H$-Poisson bracket
has bidegree (arity: $-1$, cohomological: $+1$), so
$\{2x^3, 2x^3\}_H$ has arity $3 + 3 - 1 = 5$ and cohomological
degree $0 + 0 + 1 = 1$.  This is a degree-$1$ element in arity~$5$,
which represents a cocycle in the cyclic deformation complex.
Applying the differential $d_2$ then produces the final obstruction
class $o_5 \in H^2(\cA^{\mathrm{sh}}_{5,0})$.

For the Virasoro algebra, the propagator $P_\cA$ is non-degenerate
at weight~$5$ (the Gram determinant at weight~$5$ is non-vanishing
for generic~$c$), so
$P_\cA(x^3 \cdot x^4) \neq 0$.
Combined with $Q^{\mathrm{contact}} = 10/[c(5c+22)] \neq 0$
for generic~$c$, this gives
$[\mathfrak{C}, \mathfrak{Q}] \neq 0$
in $H^2(\cA^{\mathrm{sh}}_{5,0})$.

\medskip\noindent
\emph{Step~3: The inductive mechanism.}\;
The non-vanishing of $o_5$ implies that $\Sh_5 \neq 0$: the
quintic shadow must be non-trivial to solve the arity-$5$ master
equation.  By the same argument at each subsequent arity, the
cubic shadow $\mathfrak{C} = 2x^3$ provides a permanent source
via the bracket $[\mathfrak{C}, \Sh_r]$ at arity $r+1$:
\[
o_{r+1} \;\ni\;
[\mathfrak{C}, \Sh_r]
\;=\;
[2x^3, \Sh_r(\mathrm{Vir}_c)].
\]
As long as $\Sh_r \neq 0$ and the propagator is non-degenerate
at weight~$r+1$ (which holds for generic~$c$), the obstruction
$o_{r+1} \neq 0$ and the tower continues.  This gives
$r_{\max}(\mathrm{Vir}_c) = \infty$.
\end{proof}

\begin{remark}[The cubic source mechanism]
\label{rem:thqg-VI-cubic-source}
\index{cubic shadow!as permanent source}
The proof above identifies the fundamental mechanism that makes
the Virasoro/\hskip0pt$\mathcal{W}_N$ tower infinite: the
cubic shadow $\mathfrak{C}(\cA) \neq 0$ acts as a
\emph{permanent source} for all higher obstructions via the
bracket $[\mathfrak{C}, \Sh_r]$.  This is the structural
reason that class-$\mathbf{M}$ algebras have infinite shadow
depth.  In contrast:
\begin{itemize}
\item Class-$\mathbf{G}$: $\mathfrak{C} = 0$, so
  the source is absent and the tower terminates at $r = 2$.
\item Class-$\mathbf{L}$: $\mathfrak{C} \neq 0$ but
  $o_4 = 0$ by the Jacobi identity, so the source is
  killed at the first step and the tower terminates at $r = 3$.
\item Class-$\mathbf{C}$: $\mathfrak{C} = 0$ but
  $\mathfrak{Q} \neq 0$ (the quartic is the leading term);
  $o_5 = 0$ by rank-$1$ rigidity, so the tower terminates at
  $r = 4$.
\item Class-$\mathbf{M}$: $\mathfrak{C} \neq 0$ and the
  bracket $[\mathfrak{C}, -]$ never annihilates the tower,
  giving $r_{\max} = \infty$.
\end{itemize}
This gives a clean algebraic taxonomy of soft graviton complexity.
\end{remark}

\subsubsection{The arity-$5$ soft factor for Virasoro}

\begin{computation}[Quintic soft factor for Virasoro;
\ClaimStatusProvedHere]
\label{comp:thqg-VI-virasoro-quintic-soft}
\index{Virasoro algebra!quintic soft factor}
For $\cA = \mathrm{Vir}_c$, the arity-$5$ soft factor at genus~$0$
is:
\begin{equation}\label{eq:thqg-VI-virasoro-quintic-soft}
S^{(3)}_{\mathrm{Vir}}
\;=\;
\Sh_{0,n}(\Sh_5(\mathrm{Vir}_c))\bigl|_{\text{soft}}
\;+\;
\bigl[\Sh_{0,n}(\mathfrak{C}),\,
\Sh_{0,n}(\mathfrak{Q})\bigr]\bigl|_{\text{soft}}.
\end{equation}
By the recursion of Corollary~\ref{cor:thqg-VI-soft-recursion}, the
cross-term involves the subleading factor $S^{(1)}$ and the
sub-subleading factor $S^{(2)}$:
\begin{equation}\label{eq:thqg-VI-virasoro-s3-explicit}
S^{(3)}_{\mathrm{Vir}}
\;=\;
\frac{10}{c(5c+22)}\;\cdot\;
\frac{2}{c}\;\cdot\;
\sum_{i=1}^{n}
\frac{h_i^3(h_i-1)^2(2h_i - 1)}{(z_s - z_i)^4}
\;+\; \text{(cross-terms and descendants)},
\end{equation}
where the leading coefficient $10\cdot 2/[c^2(5c+22)]
= 20/[c^2(5c+22)]$ comes from the product
$Q^{\mathrm{contact}} \cdot (2/c)$ and the conformal weight
polynomial $h_i^3(h_i-1)^2(2h_i-1)$ comes from the action
of $:T^3:$ at weight~$5$.
\end{computation}

\subsubsection{Higher arities and the inductive structure}

\begin{proposition}[Inductive non-termination for class $\mathbf{M}$;
\ClaimStatusProvedHere]
\label{prop:thqg-VI-inductive-nontermination}
\index{shadow Postnikov tower!non-termination|textbf}
\index{soft graviton theorem!infinite tower}
Let $\cA$ be a class-$\mathbf{M}$ modular Koszul chiral algebra
(i.e.\ $\mathfrak{C}(\cA) \neq 0$ and $\mathfrak{Q}(\cA) \neq 0$).
Assume the Gram determinant of the cyclic deformation complex is
non-degenerate at all weights (which holds at generic central
charge or level).  Then:
\begin{enumerate}[label=\textup{(\roman*)}]
\item $\Sh_r(\cA) \neq 0$ for all $r \geq 2$.
\item $o_{r+1}(\cA) \neq 0$ for all $r \geq 2$.
\item The shadow Postnikov tower is genuinely infinite:
  $r_{\max}(\cA) = \infty$.
\item The system of soft graviton Ward identities at genus~$0$
  is an infinite tower of coupled differential equations.
\end{enumerate}
\end{proposition}

\begin{proof}
By induction on~$r$.  The base cases $r = 2, 3, 4$ are
$\kappa \neq 0$, $\mathfrak{C} \neq 0$, $\mathfrak{Q} \neq 0$
(the class-$\mathbf{M}$ hypothesis).

\emph{Inductive step:} Assume $\Sh_r \neq 0$ for some $r \geq 4$.
The obstruction at arity $r+1$ is
\[
o_{r+1}
\;\ni\;
[\mathfrak{C}, \Sh_r]
\;=\;
P_\cA\bigl(\mathfrak{C} \cdot \Sh_r\bigr)
\;\neq\; 0,
\]
where the non-vanishing uses: (a)~$\mathfrak{C} \neq 0$;
(b)~$\Sh_r \neq 0$; (c)~non-degeneracy of the Gram determinant
at weight~$r+1$.
Since $o_{r+1} \neq 0$, the arity-$(r+1)$ master equation
forces $\Sh_{r+1} \neq 0$ (it is the solution of
$\nabla_H(\Sh_{r+1}) + o_{r+1} = 0$).  The induction continues.
\end{proof}

\begin{remark}[The soft graviton tower as an $L_\infty$-formality
obstruction]
\label{rem:thqg-VI-linfty-formality}
\index{Linfty-formality@$L_\infty$-formality!soft graviton tower}
By the shadow-formality identification
(Proposition~\ref{prop:shadow-formality-low-arity}),
the arity-$r$ shadow $\Sh_r(\cA)$ coincides (at arities $2,3,4$)
with the $r$-th formality obstruction of the genus-$0$ modular
$L_\infty$-algebra $\gAmod$.  Thus the soft graviton tower at
genus~$0$ is the formality obstruction tower of the gravitational
$L_\infty$-algebra.  For class-$\mathbf{G}$ algebras
($r_{\max} = 2$), the $L_\infty$-algebra is formal and gravity is
free.  For class-$\mathbf{M}$ algebras ($r_{\max} = \infty$), the
$L_\infty$-algebra is maximally non-formal and gravity has
infinitely many independent interactions.  The operadic complexity
conjecture (Remark~\ref{rem:thqg-VI-operadic-complexity}) states
that this identification extends to all arities.
\end{remark}

\begin{remark}[Operadic complexity conjecture]
\label{rem:thqg-VI-operadic-complexity}
\index{operadic complexity!soft graviton interpretation}
The conjectured equality $r_{\max}(\cA) = d_\infty(\cA) =
f_\infty(\cA)$ (the operadic complexity conjecture, proved at
arities $\leq 4$) states that the shadow depth equals the
$A_\infty$-depth equals the $L_\infty$-formality level.  In the
soft graviton context, this means: the number of independent soft
graviton theorems for a boundary algebra~$\cA$ equals the
minimal number of $A_\infty$ operations needed to specify the
transferred cyclic minimal model of the bar complex.  The
verification table in
Example~\ref{ex:operadic-complexity-verification} confirms this
for the four archetypes.
\end{remark}

\subsubsection{The higher-arity soft factors: structure theorems}

\begin{proposition}[Polynomial growth of soft factors;
\ClaimStatusProvedHere]
\label{prop:thqg-VI-polynomial-growth}
\index{soft graviton theorem!polynomial growth}
For a class-$\mathbf{M}$ algebra $\cA$ at generic central charge
or level, the arity-$r$ soft factor $S^{(r-2)}$ at genus~$0$ has
the following structure:
\begin{enumerate}[label=\textup{(\roman*)}]
\item $S^{(r-2)}$ is a rational function of the insertion
  points $z_1, \dotsc, z_n$ with poles of order at most $r-1$
  along the diagonals $z_i = z_j$.
\item The conformal-weight polynomial in each term of $S^{(r-2)}$
  has degree at most $2r - 2$ in the $h_i$.
\item The central-charge dependence of $S^{(r-2)}$ is a rational
  function whose denominator is a product of Gram determinant
  factors at weights $\leq r$.
\item The leading behavior as $c \to \infty$ is
  $S^{(r-2)} = O(c^{-(r-2)})$: each successive soft theorem is
  suppressed by one power of the central charge.
\end{enumerate}
\end{proposition}

\begin{proof}
Part~(i): The propagator on $\overline{\mathcal{M}}_{0,r+1}$
is a rational function of the insertion points with poles along
the diagonals.  The maximal pole order is $r-1$, achieved when
$r-1$ of the $r$ insertions collide simultaneously.

Part~(ii): The shadow $\Sh_r$ is an $r$-linear operation on the
conformal-block insertions.  Each insertion contributes a factor
of conformal weight $h_i$, and the $r$-fold product gives
degree at most~$r$.  The propagator contractions and bracket
cross-terms contribute at most $r - 2$ additional powers, for
a total of $2r - 2$.

Part~(iii): The shadow $\Sh_r$ involves the homotopy transfer
operator $h$, which inverts the Gram matrix at each weight
level.  The denominators are therefore products of Gram
determinant factors.  For Virasoro, these are
$\prod_{k=2}^{r} \det G_k(c)$, where $G_k(c)$ is the
weight-$k$ Shapovalov matrix.

Part~(iv): In the large-$c$ limit, the transferred
$A_\infty$ operations scale as $m_n \sim c^{-(n-2)}$ (from
the normalization of the OPE).  The arity-$r$ shadow involves
$r-2$ propagator contractions, each contributing a factor of
$1/c$, giving $S^{(r-2)} \sim c^{-(r-2)}$.
\end{proof}

\begin{corollary}[Convergence of the soft expansion;
\ClaimStatusProvedHere]
\label{cor:thqg-VI-soft-convergence}
\index{soft graviton expansion!convergence}
For a class-$\mathbf{M}$ algebra with generic central charge
$|c| \gg 1$, the soft expansion
$\sum_{k=0}^{\infty} z_s^{-k}\,S^{(k)}\,\Phi_n$
converges absolutely for $|z_s|$ sufficiently large (larger than
all $|z_i|$).  The convergence radius is controlled by the
maximal $|z_i|$ and the growth rate of the Gram determinant
factors.
\end{corollary}

\begin{proof}
By Proposition~\ref{prop:thqg-VI-polynomial-growth}(iv),
$|S^{(k)}| \leq C^k / |c|^k$ for some constant~$C$ depending on
the insertion points.  The series
$\sum_k |z_s|^{-k}\,|S^{(k)}| \leq \sum_k (C/(|c|\,|z_s|))^k$
converges for $|z_s| > C/|c|$.  Since we work in the regime
$|z_s| \gg \max_i |z_i|$ and $|c| \gg 1$, the convergence is
guaranteed.  The precise convergence radius involves the HS-sewing
condition (Theorem~\ref{thm:general-hs-sewing}).
\end{proof}

% ======================================================================
%
%   SUBSECTION 6: CELESTIAL HOLOGRAPHY AND THE SOFT ALGEBRA
%
% ======================================================================

\subsection{Celestial holography and the soft algebra}
\label{subsec:thqg-VI-celestial}
\index{celestial holography|textbf}
\index{soft algebra!from shadow tower}
\index{BMS group!celestial OPE algebra}

The soft graviton theorems of the preceding subsections admit a
powerful reinterpretation through the lens of celestial holography:
four-dimensional gravitational scattering amplitudes, when
expressed on the celestial sphere $\mathbb{P}^1$ at null infinity,
are organized by a two-dimensional chiral algebra---the
\emph{celestial OPE algebra}---whose Ward identities are the soft
graviton theorems.  This subsection develops the connection between
the shadow tower and celestial holography.

\subsubsection{The celestial OPE algebra from the shadow tower}

\begin{definition}[Celestial soft algebra]
\label{def:thqg-VI-celestial-soft-algebra}
\index{celestial soft algebra|textbf}
\index{soft algebra!celestial|textbf}
Let $\cA$ be a modular Koszul chiral algebra.  The
\emph{celestial soft algebra of~$\cA$} is the associative algebra
\begin{equation}\label{eq:thqg-VI-celestial-soft-algebra}
\mathfrak{S}(\cA)
\;:=\;
\bigoplus_{r=2}^{r_{\max}(\cA)}
\End\bigl(\cV_{0,n}(\cA)\bigr)^{(r)},
\end{equation}
where $\End(\cV_{0,n})^{(r)}$ is the subspace of endomorphisms
generated by the arity-$r$ shadow representation
$\Sh_{0,n}(\Sh_r(\cA))$.  The algebra structure is the
composition of endomorphisms, and the generating relations are
the Ward identities of
Theorem~\ref{thm:thqg-VI-ward-determines-correlators}.
\end{definition}

\begin{proposition}[Structure of the celestial soft algebra;
\ClaimStatusProvedHere]
\label{prop:thqg-VI-celestial-structure}
\index{celestial soft algebra!structure|textbf}
The celestial soft algebra $\mathfrak{S}(\cA)$ has the following
properties:
\begin{enumerate}[label=\textup{(\roman*)}]
\item $\mathfrak{S}(\cA)$ is generated by the arity-$2$ and
  arity-$3$ components (the Weinberg and Cachazo--Strominger
  sectors), subject to the flatness relations of
  Theorem~\ref{thm:thqg-VI-flatness-by-arity}.
\item The arity-$2$ sector generates a subalgebra isomorphic to
  the supertranslation algebra
  (Corollary~\ref{cor:thqg-VI-bms-supertranslation}).
\item The arity-$3$ sector, combined with arity-$2$, generates a
  subalgebra isomorphic to the extended BMS algebra
  (Corollary~\ref{cor:thqg-VI-superrotation}).
\item For class-$\mathbf{G}$ algebras:
  $\mathfrak{S}(\cA) \cong$ supertranslation algebra (abelian).
\item For class-$\mathbf{L}$ algebras:
  $\mathfrak{S}(\cA) \cong \mathrm{BMS}_4$ (extended BMS).
\item For class-$\mathbf{C}$ and $\mathbf{M}$ algebras:
  $\mathfrak{S}(\cA)$ is a strict extension of $\mathrm{BMS}_4$
  by the higher soft sectors.
\end{enumerate}
\end{proposition}

\begin{proof}
Part~(i): The generation claim follows from the recursion of
Corollary~\ref{cor:thqg-VI-soft-recursion}: each higher-arity
soft factor $S^{(r-2)}$ is determined by the lower-arity factors
through the flatness equation plus the obstruction
classes~$o_r$.  Since $o_r$ is computed from the bracket
$[\mathfrak{C}, \Sh_{r-1}]$ (the cubic source mechanism of
Remark~\ref{rem:thqg-VI-cubic-source}), and $\mathfrak{C}$ is
the arity-$3$ shadow, the full tower is generated by arities $2$
and~$3$.

Parts~(ii)--(iii): Proved in
Corollary~\ref{cor:thqg-VI-bms-supertranslation} and
Corollary~\ref{cor:thqg-VI-superrotation}.

Parts~(iv)--(vi): The shadow depth determines which sectors are
non-trivial.  For class-$\mathbf{G}$, only arity~$2$ is present,
giving the abelian supertranslation algebra.  For
class-$\mathbf{L}$, arities~$2$ and~$3$ are present, giving
$\mathrm{BMS}_4$.  For classes~$\mathbf{C}$ and~$\mathbf{M}$,
the higher arities provide extensions.
\end{proof}

\begin{remark}[Relation to the $\mathfrak{w}_{1+\infty}$ algebra]
\label{rem:thqg-VI-w-infinity}
\index{w1infinity@$\mathfrak{w}_{1+\infty}$!celestial holography}
In the celestial holography literature, the soft graviton algebra
for Einstein gravity in 4d is identified with a deformation of
$\mathfrak{w}_{1+\infty}$ (Strominger et al.,
Guevara--Himwich--Pate--Strominger 2021).  In the shadow-tower
language, this arises as follows.  Take $\cA = \mathrm{Vir}_c$
at large~$c$.  The celestial soft algebra $\mathfrak{S}(\mathrm{Vir}_c)$
is generated by the stress tensor mode algebra, which at genus~$0$
is the Virasoro algebra itself.  The higher soft sectors
$\Sh_r(\mathrm{Vir}_c)$ for $r \geq 4$ are the normal-ordered
products $:T^{r-1}:$ and their descendants, which generate the
$\mathcal{W}_{1+\infty}$ algebra.

The precise identification is:
\begin{equation}\label{eq:thqg-VI-w-infinity-identification}
\mathfrak{S}(\mathrm{Vir}_c)
\;\simeq\;
\mathcal{W}_{1+\infty}[c],
\end{equation}
the $\mathcal{W}_{1+\infty}$ algebra at central charge~$c$,
realized as the celestial soft algebra of the Virasoro
boundary chiral algebra.  The soft graviton theorems at each
arity are the Ward identities of the $\mathcal{W}_{1+\infty}$
currents: the arity-$r$ soft theorem corresponds to the
Ward identity of the spin-$(r-1)$ $\mathcal{W}$-current.
\end{remark}

\subsubsection{The soft graviton OPE from the shadow connection}

\begin{theorem}[Soft graviton OPE; \ClaimStatusProvedHere]
\label{thm:thqg-VI-soft-ope}
\index{soft graviton OPE!from shadow connection|textbf}
The shadow connection $\nabla^{\mathrm{hol}}_{0,n}$ determines a
chiral OPE on the space of soft graviton insertions.  For two soft
insertions at points $w_1, w_2$ on the celestial sphere, the OPE is
\begin{equation}\label{eq:thqg-VI-soft-ope}
\mathcal{O}_s(w_1) \cdot \mathcal{O}_s(w_2)
\;\sim\;
\sum_{k=0}^{r_{\max}-2}
\frac{C_k(\cA)}{(w_1 - w_2)^{k+1}}\;\mathcal{O}_{s,k}(w_2),
\end{equation}
where:
\begin{itemize}
\item $\mathcal{O}_s(w)$ is the soft graviton operator at
  celestial point~$w$;
\item $C_k(\cA)$ are the OPE coefficients, determined by the
  shadows $\Sh_{k+2}(\cA)$;
\item $\mathcal{O}_{s,k}(w_2)$ are the descendant soft operators.
\end{itemize}
The OPE is consistent (associative) because the shadow connection
is flat.
\end{theorem}

\begin{proof}
The OPE of two soft insertions is controlled by the short-distance
behavior of the shadow connection in the diagonal
$w_1 \to w_2$.  The flatness condition
$(\nabla^{\mathrm{hol}})^2 = 0$ ensures that the OPE is
associative (the OPE product is the monodromy of the flat
connection around the diagonal).

The singular terms of the OPE are determined by the arity
decomposition of the shadow connection:
the arity-$(k+2)$ shadow $\Sh_{k+2}(\cA)$ produces the
$(w_1 - w_2)^{-(k+1)}$ pole in the OPE.  The leading pole
($k=0$) is the arity-$2$ contribution from $\kappa(\cA)$,
giving the standard graviton OPE.  The subleading poles are
determined by the higher shadows.

The number of singular terms is $r_{\max} - 1$: for
class-$\mathbf{G}$ algebras, there is only a simple pole; for
class-$\mathbf{L}$, a double pole; for class-$\mathbf{M}$,
poles of all orders.
\end{proof}

\begin{remark}[The conformally soft limit]
\label{rem:thqg-VI-conformally-soft}
\index{conformally soft limit!shadow tower}
The celestial holography interpretation requires a
\emph{conformally soft limit}: one takes the conformal weight
$\Delta$ of the soft insertion to a special value
($\Delta \to 1$ for leading, $\Delta \to 0$ for subleading).
In the shadow-tower framework, the conformally soft limit at
arity~$r$ corresponds to the \emph{conformal weight
specialization}:
\[
\Delta_r = 2 - r
\qquad
\text{(arity-$r$ conformally soft value)}.
\]
At $\Delta = \Delta_r$, the arity-$r$ shadow representation
$\Sh_{0,n}(\Sh_r)$ develops an enhanced pole structure that
produces the conformally soft graviton.  The standard values
$\Delta = 1$ (leading, $r=2$) and $\Delta = 0$ (subleading,
$r=3$) are the first two entries in this sequence.
\end{remark}

\subsubsection{The shadow tower and asymptotic symmetries}

\begin{proposition}[Asymptotic symmetry algebra from the
shadow tower; \ClaimStatusProvedHere]
\label{prop:thqg-VI-asymptotic-symmetry}
\index{asymptotic symmetry algebra!from shadow tower|textbf}
\index{BMS group!generalization from shadow tower}
The celestial soft algebra $\mathfrak{S}(\cA)$ is the asymptotic
symmetry algebra of the holographic theory dual to~$\cA$.  For each
shadow depth class, the asymptotic symmetry algebra is:
\begin{center}
\small
\renewcommand{\arraystretch}{1.15}
\begin{tabular}{llll}
\toprule
\emph{Class} &
$r_{\max}$ &
\emph{Asymptotic symmetry} &
\emph{Gravitational theory} \\
\midrule
$\mathbf{G}$ &
  $2$ &
  Supertranslations &
  Free (linearized) gravity \\
$\mathbf{L}$ &
  $3$ &
  $\mathrm{BMS}_4$ &
  Semistrict higher-spin gravity \\
$\mathbf{C}$ &
  $4$ &
  $\mathrm{BMS}_4 \rtimes W_3$ &
  Quartic gravity \\
$\mathbf{M}$ &
  $\infty$ &
  $\mathcal{W}_{1+\infty}$ &
  Full gravitational $L_\infty$-algebra \\
\bottomrule
\end{tabular}
\end{center}
The shadow depth thus controls the asymptotic symmetry content
of the dual gravitational theory.
\end{proposition}

\begin{proof}
For each class, the asymptotic symmetry algebra is the subalgebra
of $\End(\cV_{0,n})$ generated by the Ward identities at arities
$2, \dotsc, r_{\max}$.  The identification with the listed algebras
follows from:

\emph{Class~$\mathbf{G}$}: Only arity-$2$ Ward identities
(supertranslations), as proved in
Corollary~\ref{cor:thqg-VI-bms-supertranslation}.

\emph{Class~$\mathbf{L}$}: Arities~$2$ and~$3$ generate
$\mathrm{BMS}_4$ (supertranslations + superrotations), as proved
in Corollary~\ref{cor:thqg-VI-superrotation}.

\emph{Class~$\mathbf{C}$}: The arity-$4$ Ward identity extends
$\mathrm{BMS}_4$ by the quartic sector.  For the $\beta\gamma$
system, the quartic contact structure gives a $W_3$-type extension.

\emph{Class~$\mathbf{M}$}: The full infinite tower gives
$\mathcal{W}_{1+\infty}$ as in
Remark~\ref{rem:thqg-VI-w-infinity}.
\end{proof}

\subsubsection{The shadow connection at higher genus and
gravitational memory}

\begin{remark}[Gravitational memory from higher-genus shadows]
\label{rem:thqg-VI-gravitational-memory}
\index{gravitational memory!higher-genus shadow}
At genus $g \geq 1$, the shadow connection
$\nabla^{\mathrm{hol}}_{g,n}$ acquires the curvature correction
$\kappa(\cA) \cdot \omega_g$
(Proposition~\ref{prop:thqg-VI-kappa-action}).  The resulting
Ward identities at higher genus involve the Hodge class on
$\overline{\mathcal{M}}_g$ and encode \emph{gravitational memory
effects}: permanent shifts in the soft graviton waveform caused
by the passage of radiation.

In the shadow-tower framework, gravitational memory is the
genus-$1$ correction to the leading soft theorem:
\begin{equation}\label{eq:thqg-VI-memory}
\Delta_{\mathrm{memory}}
\;=\;
\kappa(\cA) \cdot \omega_1
\;=\;
\kappa(\cA) \cdot \frac{1}{24}\,E_2(\tau),
\end{equation}
where $E_2(\tau)$ is the weight-$2$ Eisenstein series (the
normalized Hodge class at genus~$1$).  The memory effect is
proportional to $\kappa(\cA)$, confirming that it is a universal
consequence of the modular characteristic.
\end{remark}

% ======================================================================
%
%   SUMMARY TABLE
%
% ======================================================================

\subsection{Summary: the soft graviton dictionary}
\label{subsec:thqg-VI-summary}

\begin{table}[h]
\centering
\small
\renewcommand{\arraystretch}{1.25}
\begin{tabular}{lccc}
\toprule
\textbf{Soft theorem} &
\textbf{Arity} &
\textbf{Shadow} &
\textbf{Physical content} \\
\midrule
Leading (Weinberg) &
  $2$ & $\kappa(\cA)$ &
  Energy-momentum conservation \\
Subleading (Cachazo--Strominger) &
  $3$ & $\mathfrak{C}(\cA)$ &
  Angular momentum \\
Sub-subleading &
  $4$ & $\mathfrak{Q}(\cA)$ &
  Stress tensor / $:T^2:$ \\
Order-$(r-2)$ &
  $r$ & $\Sh_r(\cA)$ &
  $:T^{r-1}:$ ($\cW$-current) \\
\bottomrule
\end{tabular}
\caption{Dictionary between the shadow Postnikov tower and the
soft graviton hierarchy.}
\label{tab:thqg-VI-soft-dictionary}
\end{table}

\begin{table}[h]
\centering
\small
\renewcommand{\arraystretch}{1.25}
\begin{tabular}{lcccc}
\toprule
\textbf{Family} &
\textbf{Class} &
$r_{\max}$ &
\textbf{\# soft thms} &
\textbf{Explicit coefficients} \\
\midrule
Heisenberg &
  $\mathbf{G}$ & $2$ & $1$ &
  $\kappa = c/2$ \\
Affine $\widehat{\fg}_k$ &
  $\mathbf{L}$ & $3$ & $2$ &
  $\kappa = k\,\delta/2$, $\mathfrak{C} = f_{abc}$ \\
$\beta\gamma$ &
  $\mathbf{C}$ & $4$ & $3$ &
  $\kappa = 0$, $\mathfrak{Q} = \mathrm{cyc}(m_3)$ \\
$\mathrm{Vir}_c$ &
  $\mathbf{M}$ & $\infty$ & $\infty$ &
  $\kappa = c/2$, $Q^{\mathrm{ct}} = 10/[c(5c+22)]$ \\
$\mathcal{W}_N$ &
  $\mathbf{M}$ & $\infty$ & $\infty$ &
  Higher-spin generalization \\
\bottomrule
\end{tabular}
\caption{Shadow depth and the soft graviton hierarchy for the
standard families.}
\label{tab:thqg-VI-family-soft}
\end{table}

\begin{remark}[The single-equation origin]
\label{rem:thqg-VI-single-equation}
\index{Maurer--Cartan equation!single-equation origin of soft theorems}
The entire hierarchy of soft graviton theorems---leading,
subleading, sub-subleading, and all higher orders---descends from
the single Maurer--Cartan equation
$D_\cA\Theta_\cA + \tfrac{1}{2}[\Theta_\cA,\Theta_\cA] = 0$
(Theorem~\ref{thm:mc2-bar-intrinsic}).  The shadow connection
$\nabla^{\mathrm{hol}}_{g,n} = d - \Sh_{g,n}(\Theta_\cA)$ is
the universal vehicle that converts this algebraic equation into
Ward identities on correlators.  The arity filtration of the
shadow Postnikov tower provides the hierarchy, and the shadow
depth controls the complexity.  No dynamical assumption beyond
the existence of a modular Koszul chiral algebra is needed:
the soft theorems are algebraic consequences of the bar-intrinsic
MC element.
\end{remark}

\begin{remark}[Relation to results \textbf{(G1)}--\textbf{(G5)}
and \textbf{(G7)}--\textbf{(G10)}]
\label{rem:thqg-VI-other-results}
\index{twisted holography!result interdependencies}
The soft graviton theorems of result~\textbf{(G6)} interact with
the other results of this chapter as follows:
\begin{enumerate}[label=\textup{(\alph*)}]
\item \textbf{(G1)} Perturbative finiteness: The soft factors
  $S^{(k)}$ are finite at each arity because the shadow
  Postnikov tower is constructed from finite graph sums
  (arity cutoff, Lemma~\ref{lem:arity-cutoff}).
\item \textbf{(G2)} Gravitational complexity: The shadow depth
  $r_{\max}$ simultaneously controls the number of independent
  soft theorems and the gravitational complexity class.
\item \textbf{(G3)} Symplectic polarization: The soft Ward
  identities are compatible with the Lagrangian splitting
  $\mathbf{C}_g \simeq \mathbf{Q}_g \oplus \mathbf{Q}_g^!$:
  the soft factor at each arity respects the bulk/defect
  decomposition.
\item \textbf{(G5)} Gravitational Yangian: The arity-$2$ shadow
  at genus~$0$ is the $r$-matrix, so the leading soft theorem
  is equivalent to the CYBE for the gravitational Yangian.
\item \textbf{(G7)} Modular bootstrap: The higher-genus soft
  Ward identities provide inductive constraints on the
  gravitational partition function, supplementing the genus
  spectral sequence.
\end{enumerate}
\end{remark}

\begin{remark}[Open directions]
\label{rem:thqg-VI-open-directions}
\index{soft graviton theorem!open directions}
Three directions remain beyond the scope of this section:
\begin{enumerate}[label=\textup{(\roman*)}]
\item \emph{Loop-level soft theorems.}  The genus-$g$ soft
  Ward identities for $g \geq 1$ are defined by the shadow
  connection (Definition~\ref{def:modular-shadow-connection}),
  but the explicit soft factors at higher genus depend on
  the modular invariant hierarchy
  beyond~$\hat{A}$ (Theorem~\ref{thm:nms-beyond-ahat}).
  The genus-$1$ soft theorem involves the Hessian correction
  $\delta_H^{(1)} = 120/[c^2(5c+22)]$ for Virasoro.
\item \emph{Infrared divergences.}  The soft expansion
  $\sum_k z_s^{-k}\,S^{(k)}$ converges at large $|z_s|$
  (Corollary~\ref{cor:thqg-VI-soft-convergence}), but the
  infrared structure at $z_s \to z_i$ (collinear limits)
  requires the full OPE data of the chiral algebra, not
  just the shadow tower.
\item \emph{Non-perturbative corrections.}  The shadow
  Postnikov tower is a perturbative object: it is the
  Taylor expansion of the full MC element
  $\Theta_\cA$ in the arity variable.  Non-perturbative
  corrections to the soft theorems (e.g.\ from instantons in
  the dual gravitational theory) would arise from
  non-perturbative features of~$\Theta_\cA$ that are invisible
  in the arity expansion.
\end{enumerate}
\end{remark}

% Section file for Chapter: Twisted Holography and Quantum Gravity
% Result (G7): Modular Bootstrap as MC Recursion
%
% Cross-references:
% thm:mc2-bar-intrinsic, thm:recursive-existence,
% thm:convolution-dg-lie-structure, const:vol1-genus-spectral-sequence,
% def:modular-convolution-dg-lie

% =====================================
% Local macros (provided only if absent from the ambient manuscript)
% =====================================
\providecommand{\cA}{\mathcal{A}}
\providecommand{\cH}{\mathcal{H}}
\providecommand{\cM}{\mathcal{M}}
\providecommand{\cC}{\mathcal{C}}
\providecommand{\fg}{\mathfrak{g}}
\providecommand{\fC}{\mathfrak{C}}
\providecommand{\gAmod}{\mathfrak{g}_{\cA}^{\mathrm{mod}}}
\providecommand{\dzero}{d_0}
\providecommand{\dfib}{d_{\mathrm{fib}}}
\providecommand{\MC}{\text{MC}}
\providecommand{\barB}{\overline{B}}
\providecommand{\Defcyc}{\operatorname{Def}_{\mathrm{cyc}}}
\providecommand{\Definfmod}{\operatorname{Def}_{\infty}^{\mathrm{mod}}}
\providecommand{\Convstr}{\operatorname{Conv}_{\mathrm{str}}}
\providecommand{\Convinf}{\operatorname{Conv}_{\infty}}
\providecommand{\Gmod}{\mathcal{G}^{\mathrm{mod}}}
\providecommand{\Tr}{\operatorname{Tr}}
\providecommand{\Hom}{\operatorname{Hom}}
\providecommand{\Res}{\operatorname{Res}}
\providecommand{\id}{\mathrm{id}}
\providecommand{\ad}{\operatorname{ad}}
\providecommand{\gr}{\operatorname{gr}}
\providecommand{\wt}{\operatorname{wt}}
\providecommand{\Aut}{\operatorname{Aut}}
\providecommand{\Det}{\operatorname{Det}}
\providecommand{\Sym}{\operatorname{Sym}}
\providecommand{\Omegach}{\Omega^{\text{ch}}}

% Bootstrap-specific macros
\providecommand{\Ob}{\mathrm{Ob}}
\providecommand{\Shad}{\operatorname{Sh}}
\providecommand{\hCont}{\mathsf{h}}
\providecommand{\BVop}{\Delta_{\mathrm{BV}}}
\providecommand{\dsew}{d_{\mathrm{sew}}}
\providecommand{\dpf}{d_{\mathrm{pf}}}

\section{Modular bootstrap as MC recursion}
% label removed: sec:thqg-modular-bootstrap
\index{modular bootstrap|textbf}
\index{genus spectral sequence!bootstrap}
\index{Maurer--Cartan equation!genus expansion}
\index{bootstrap!modular, from MC recursion}

The genus-$g$ partition function is determined inductively from
genus-$0$ data by the Maurer--Cartan equation in the modular
convolution algebra~$\gAmod$
(Definition~\ref*{V1-def:modular-convolution-dg-lie}):

\begin{center}
\emph{The MC equation $D_\cA\,\Theta_\cA
+ \tfrac{1}{2}[\Theta_\cA,\Theta_\cA] = 0$
is a modular bootstrap equation. Expanded genus by genus,
it determines $\Theta^{(g)}_\cA$ from $\Theta^{(<g)}_\cA$
by solving a linear equation with an explicitly computable
inhomogeneity, the genus-$g$ obstruction class $\Ob_g$.}
\end{center}

The conformal bootstrap proceeds by consistency conditions
(crossing symmetry, modular invariance) that constrain but
do not uniquely determine the correlators. The MC recursion
is \emph{algebraically complete}: it determines the partition
function at each genus from the OPE data at genus~$0$, with
no ambiguity beyond gauge equivalence. The mechanism is the
contracting homotopy of the tree-twisted differential
$d_{\Theta^{(0)}} = \dzero + [\Theta^{(0)}, -]$ on the
bar complex.

% ======================================================================
%
% PART 1: THE MC EQUATION EXPANDED BY GENUS
%
% ======================================================================

\subsection{The MC equation expanded by genus}
% label removed: subsec:thqg-VII-mc-genus-expansion
\index{Maurer--Cartan equation!genus expansion|textbf}

\subsubsection{The genus filtration on the convolution algebra}
% label removed: subsubsec:thqg-VII-genus-filtration

We recall from Definition~\ref*{V1-def:modular-convolution-dg-lie}
that the modular convolution dg~Lie algebra
\begin{equation}% label removed: eq:thqg-VII-gAmod-def
\gAmod
\;=\;
\prod_{g \geq 0}\;
\prod_{n \geq 1}\;
\Hom_{S_n}\!\bigl(
C_*(\overline{\cM}_{g,n}),\;
\cA^{\otimes n}
\bigr)^{S_n}[-1]
\end{equation}
carries a descending genus filtration
\begin{equation}% label removed: eq:thqg-VII-genus-filtration
F^g\,\gAmod
\;:=\;
\prod_{g' \geq g}\;
\prod_{n \geq 1}\;
\Hom_{S_n}\!\bigl(
C_*(\overline{\cM}_{g',n}),\;
\cA^{\otimes n}
\bigr)^{S_n}[-1].
\end{equation}
This filtration is complete and Hausdorff:
$\gAmod = F^0\,\gAmod$ and
$\bigcap_{g \geq 0} F^g\,\gAmod = 0$.
The associated graded is
\begin{equation}% label removed: eq:thqg-VII-assoc-graded
\gr^g_F\,\gAmod
\;=\;
F^g / F^{g+1}
\;\cong\;
\prod_{n \geq 1}\;
\Hom_{S_n}\!\bigl(
C_*(\overline{\cM}_{g,n}),\;
\cA^{\otimes n}
\bigr)^{S_n}[-1],
\end{equation}
the space of genus-$g$ operations.

\begin{remark}[Completeness and $\hbar$-adic topology]
\label{rem:thqg-VII-hbar-topology}
\index{hbar-adic topology@$\hbar$-adic topology}
The genus filtration is equivalent to the $\hbar$-adic filtration
when one introduces the formal genus-counting parameter~$\hbar$
and writes $\gAmod[\![\hbar]\!]$ with $\hbar^g$ multiplying the
genus-$g$ sector. The $\hbar$-adic topology on
$\gAmod[\![\hbar]\!]$ is the linear topology defined by the
ideals $\hbar^g \cdot \gAmod[\![\hbar]\!]$, and this coincides
with the topology defined by~\eqref{V1-eq:thqg-VII-genus-filtration}.
Completeness holds because the product in~%
\eqref{V1-eq:thqg-VII-gAmod-def} is an infinite product, not
a direct sum.
\end{remark}

\subsubsection{Decomposition of the differential}
% label removed: subsubsec:thqg-VII-differential-decomp

The differential $D$ on $\gAmod$ decomposes as the sum of
five components
(Theorem~\ref{V1-thm:convolution-dg-lie-structure}):
\begin{equation}% label removed: eq:thqg-VII-D-components
D \;=\; d_{\mathrm{Ch}_\infty}
 + d_{\check{C}}
 + d_{\mathrm{coll}}
 + d_{\mathrm{loop}}
 + d_{\mathrm{pf}},
\end{equation}
where:
\begin{enumerate}[label=(\roman*)]
\item $d_{\mathrm{Ch}_\infty}$: the homotopy chiral algebra
 structure maps (internal to each vertex of the graph);
\item $d_{\check{C}}$: the \v{C}ech differential (from the
 boundary stratification of $\overline{\cM}_{g,n}$, encoding
 separating degenerations);
\item $d_{\mathrm{coll}}$: the collision differential (residues
 along the collision divisors in the Fulton--MacPherson
 compactification);
\item $d_{\mathrm{loop}}$: the non-separating clutching operator
 (the BV operator $\Delta$ that raises genus by one, dual to
 the non-separating node of the boundary);
\item $d_{\mathrm{pf}}$: the planted-forest corrections
 (from the log-FM boundary strata of
 Mok~\cite{Mok25}).
\end{enumerate}
With respect to the genus filtration, these have the following
behaviour:

\begin{lemma}[Genus shifts of the differential components]
% label removed: lem:thqg-VII-genus-shifts
\ClaimStatusProvedElsewhere
\begin{enumerate}[label=\textup{(\roman*)}]
\item $d_{\mathrm{Ch}_\infty}$, $d_{\check{C}}$, and
 $d_{\mathrm{coll}}$ \emph{preserve} genus:
 $d_{\mathrm{Ch}_\infty}(F^g) \subset F^g$,
 and similarly for $d_{\check{C}}$ and $d_{\mathrm{coll}}$.
\item $d_{\mathrm{loop}}$ \emph{raises} genus by one:
 $d_{\mathrm{loop}}(F^g) \subset F^{g+1}$.
\item $d_{\mathrm{pf}}$ \emph{raises} genus by at least one:
 $d_{\mathrm{pf}}(F^g) \subset F^{g+1}$.
\end{enumerate}
\end{lemma}

\begin{proof}
The non-separating clutching map
$\overline{\cM}_{g-1,n+2} \to \partial\overline{\cM}_{g,n}$
identifies two marked points on a genus-$(g{-}1)$ surface to
produce a genus-$g$ surface. The dual operation on the
convolution algebra maps genus-$g$ elements to genus-$(g{+}1)$
elements. The components $d_{\mathrm{Ch}_\infty}$,
$d_{\check{C}}$, and $d_{\mathrm{coll}}$ all act fibrewise over
a fixed moduli space $\overline{\cM}_{g,n}$ and hence preserve
the genus. The planted-forest component $d_{\mathrm{pf}}$
involves boundary strata indexed by planted forests of positive
loop genus (each internal edge raising the genus), hence also
raises genus.

The genus filtration satisfies the multiplicative axiom
$F^g \cdot F^{g'} \subset F^{g+g'}$: the Lie bracket
$[\varphi, \psi]$ for $\varphi \in F^g$ and $\psi \in F^{g'}$
involves stable graphs of genus $\geq g + g'$, because the
bracket glues a genus-$g$ and a genus-$g'$ graph along a
single edge, producing a graph of genus $\geq g + g'$.
This multiplicative property ensures that the genus filtration
is compatible with the dg~Lie algebra structure.
\end{proof}

\begin{definition}[Genus-preserving and genus-raising parts]
% label removed: def:thqg-VII-D-split
\index{differential!genus-preserving part}
\index{differential!genus-raising part}
Define the genus-preserving part of $D$:
\begin{equation}% label removed: eq:thqg-VII-D0
D_0
\;:=\;
d_{\mathrm{Ch}_\infty}
+ d_{\check{C}}
+ d_{\mathrm{coll}},
\end{equation}
and the genus-raising part:
\begin{equation}% label removed: eq:thqg-VII-D1
D_1
\;:=\;
d_{\mathrm{loop}}
+ d_{\mathrm{pf}}.
\end{equation}
Then $D = D_0 + D_1$ with $D_0^2 = 0$,
$D_0 D_1 + D_1 D_0 = 0$, $D_1^2 = 0$
(these follow from $D^2 = 0$ by genus bookkeeping).
\end{definition}

\begin{remark}[The three identities]
% label removed: rem:thqg-VII-three-identities
The factorization $D^2 = 0$ into three homogeneous pieces is the
structural core of the genus spectral sequence. Explicitly:
\begin{align}
D^2 &= (D_0 + D_1)^2 \notag\\
&= D_0^2 + (D_0 D_1 + D_1 D_0) + D_1^2 \notag\\
&= 0,% label removed: eq:thqg-VII-D-squared-decomp
\end{align}
and each of the three terms vanishes separately because they
have different genus shifts: $D_0^2$ preserves genus,
$D_0 D_1 + D_1 D_0$ shifts by one, $D_1^2$ shifts by two.
Since these map into disjoint genus sectors of $\gAmod$,
each must vanish independently.
\end{remark}

\subsubsection{The genus expansion of $\Theta_\cA$}
% label removed: subsubsec:thqg-VII-genus-expansion

The universal MC element $\Theta_\cA \in \MC(\gAmod)$
(Theorem~\ref*{V1-thm:mc2-bar-intrinsic}) admits a genus expansion
\begin{equation}% label removed: eq:thqg-VII-theta-genus-exp
\Theta_\cA
\;=\;
\sum_{g \geq 0} \hbar^g\,\Theta^{(g)}_\cA,
\end{equation}
where $\Theta^{(g)}_\cA \in \gr^g_F\,\gAmod$ is the genus-$g$
component. The MC equation reads
\begin{equation}% label removed: eq:thqg-VII-mc-full
D\,\Theta_\cA
+ \tfrac{1}{2}[\Theta_\cA,\, \Theta_\cA]
\;=\; 0.
\end{equation}
We expand both sides in powers of~$\hbar$.

\begin{proposition}[Genus-$g$ MC equation; \ClaimStatusProvedHere]
% label removed: prop:thqg-VII-genus-g-mc
\index{Maurer--Cartan equation!genus-$g$ component}
The coefficient of $\hbar^g$ in the MC
equation~\eqref{V1-eq:thqg-VII-mc-full} is:
\begin{equation}% label removed: eq:thqg-VII-genus-g-mc-eq
\boxed{
d_{\Theta^{(0)}}\bigl(\Theta^{(g)}\bigr)
\;=\;
-\Ob_g\bigl(\Theta^{(<g)}\bigr),
\qquad g \geq 1,
}
\end{equation}
where the tree-twisted differential is
\begin{equation}% label removed: eq:thqg-VII-d-twisted
d_{\Theta^{(0)}}
\;:=\;
D_0 + \ad_{\Theta^{(0)}}
\;=\;
D_0 + [\Theta^{(0)},\, -],
\end{equation}
and the genus-$g$ obstruction class is
\begin{equation}% label removed: eq:thqg-VII-obs-class
\Ob_g\bigl(\Theta^{(<g)}\bigr)
\;:=\;
\underbrace{
\sum_{\substack{g_1 + g_2 = g \\ g_1, g_2 \geq 1}}
\tfrac{1}{2}\,[\Theta^{(g_1)},\, \Theta^{(g_2)}]
}_{\text{separating clutching}}
\;+\;
\underbrace{
\Delta\bigl(\Theta^{(g-1)}\bigr)
}_{\text{non-separating BV}}
\;+\;
\underbrace{
\dsew^{(g)} + \dpf^{(g)}
}_{\text{boundary strata}}.
\end{equation}
\end{proposition}

\begin{proof}
Write $\Theta = \sum_{g \geq 0} \hbar^g\,\Theta^{(g)}$ and
expand the MC equation~\eqref{V1-eq:thqg-VII-mc-full}.

\medskip\noindent\emph{Step 1: Expansion of $D\,\Theta$.}
The differential $D = D_0 + D_1$ acts on $\Theta^{(g)}$ as
\begin{align}
(D\,\Theta)_g
&= D_0(\Theta^{(g)}) + D_1(\Theta^{(g-1)})\notag\\
&= D_0(\Theta^{(g)})
 + d_{\mathrm{loop}}(\Theta^{(g-1)})
 + d_{\mathrm{pf}}(\Theta^{(g-1)}),
% label removed: eq:thqg-VII-D-theta-expand
\end{align}
since $D_0$ preserves genus while $D_1$ raises it by one.
The component $D_1(\Theta^{(g-1)})$ contributes to genus~$g$
by the non-separating clutching $\Delta(\Theta^{(g-1)})$ and
the planted-forest correction $\dpf^{(g)}$.

\medskip\noindent\emph{Step 2: Expansion of $[\Theta,\Theta]$.}
The bracket $[\Theta,\Theta]$ expanded in $\hbar$ gives
\begin{equation}% label removed: eq:thqg-VII-bracket-expand
\bigl[\Theta,\Theta\bigr]_g
\;=\;
\sum_{g_1 + g_2 = g}
[\Theta^{(g_1)},\, \Theta^{(g_2)}]
\;=\;
2\,[\Theta^{(0)},\, \Theta^{(g)}]
\;+\;
\sum_{\substack{g_1 + g_2 = g \\ g_1, g_2 \geq 1}}
[\Theta^{(g_1)},\, \Theta^{(g_2)}],
\end{equation}
where we have isolated the two terms with $g_1 = 0$ or
$g_2 = 0$ (which contribute $[\Theta^{(0)}, \Theta^{(g)}]$
each, and by the graded antisymmetry of the Lie bracket,
$[\Theta^{(0)}, \Theta^{(g)}] = -(-1)^{|\Theta^{(0)}|
|\Theta^{(g)}|}[\Theta^{(g)}, \Theta^{(0)}]$; since both
are degree~$1$, this gives $[\Theta^{(0)}, \Theta^{(g)}]
= -[\Theta^{(g)}, \Theta^{(0)}]$ and the two terms with
$g_1 = 0$ or $g_2 = 0$ contribute
$2[\Theta^{(0)}, \Theta^{(g)}]$).

\medskip\noindent\emph{Step 3: Assembly.}
Collecting the genus-$g$ coefficient of the MC equation
$D\,\Theta + \tfrac{1}{2}[\Theta,\Theta] = 0$:
\begin{align}
0 &= D_0(\Theta^{(g)})
 + \Delta(\Theta^{(g-1)})
 + \dsew^{(g)} + \dpf^{(g)}
 \notag\\
 &\qquad
 + [\Theta^{(0)},\, \Theta^{(g)}]
 + \tfrac{1}{2}
 \sum_{\substack{g_1 + g_2 = g \\ g_1, g_2 \geq 1}}
 [\Theta^{(g_1)},\, \Theta^{(g_2)}].
% label removed: eq:thqg-VII-mc-g-assembled
\end{align}
Rearranging, with $d_{\Theta^{(0)}} = D_0 +
[\Theta^{(0)}, -]$:
\begin{equation}% label removed: eq:thqg-VII-mc-g-rearranged
d_{\Theta^{(0)}}(\Theta^{(g)})
\;=\;
-\Delta(\Theta^{(g-1)})
- \dsew^{(g)} - \dpf^{(g)}
- \tfrac{1}{2}
 \sum_{\substack{g_1 + g_2 = g \\ g_1, g_2 \geq 1}}
 [\Theta^{(g_1)},\, \Theta^{(g_2)}],
\end{equation}
which is exactly $d_{\Theta^{(0)}}(\Theta^{(g)}) = -\Ob_g(\Theta^{(<g)})$.

\medskip\noindent\emph{Step 4: The tree-twisted differential squares to zero.}
We verify $d_{\Theta^{(0)}}^2 = 0$. The genus-$0$ MC equation gives
\begin{equation}% label removed: eq:thqg-VII-genus0-mc
D_0(\Theta^{(0)}) + \tfrac{1}{2}[\Theta^{(0)},\Theta^{(0)}] = 0.
\end{equation}
Write $d_{\Theta^{(0)}}(x) = D_0(x) + [\Theta^{(0)}, x]$.
By the Leibniz rule for a degree-$1$ derivation acting on the bracket:
\begin{equation}\label{eq:thqg-VII-leibniz-explicit}
D_0[\Theta^{(0)}, x]
\;=\;
[D_0\Theta^{(0)}, x]
+ (-1)^{|\Theta^{(0)}|}[\Theta^{(0)}, D_0(x)]
\;=\;
[D_0\Theta^{(0)}, x]
- [\Theta^{(0)}, D_0(x)],
\end{equation}
since $|\Theta^{(0)}| = 1$ (odd degree).
Then:
\begin{align}
d_{\Theta^{(0)}}^2(x)
&= (D_0 + \ad_{\Theta^{(0)}})^2(x)\notag\\
&= D_0^2(x) + D_0[\Theta^{(0)},x]
 + [\Theta^{(0)}, D_0(x)]
 + [\Theta^{(0)},[\Theta^{(0)},x]]\notag\\
&= 0
 + \bigl([D_0\Theta^{(0)}, x]
 - [\Theta^{(0)}, D_0(x)]\bigr)
 + [\Theta^{(0)}, D_0(x)]
 + [\Theta^{(0)},[\Theta^{(0)},x]]\notag\\
&= [D_0\Theta^{(0)}, x]
 + [\Theta^{(0)},[\Theta^{(0)},x]]\notag\\
&= [D_0\Theta^{(0)}, x]
 + \tfrac{1}{2}[[\Theta^{(0)},\Theta^{(0)}], x]\notag\\
&= [D_0\Theta^{(0)}
 + \tfrac{1}{2}[\Theta^{(0)},\Theta^{(0)}], x]\notag\\
&= 0,
% label removed: eq:thqg-VII-d-twisted-squared
\end{align}
where we used: $D_0^2 = 0$ in the first term; the Leibniz
rule~\eqref{eq:thqg-VII-leibniz-explicit} expanding
$D_0[\Theta^{(0)},x]$, after which the
$-[\Theta^{(0)}, D_0(x)]$ cancels the $+[\Theta^{(0)}, D_0(x)]$
from the cross term; the graded Jacobi identity
$[\Theta^{(0)},[\Theta^{(0)},x]] = \tfrac{1}{2}
[[\Theta^{(0)},\Theta^{(0)}], x]$
(valid because $|\Theta^{(0)}| = 1$ is odd); and
\eqref{V1-eq:thqg-VII-genus0-mc} for the final vanishing.
\end{proof}

\begin{remark}[Disjointness of the three terms]
% label removed: rem:thqg-VII-three-term-disjointness
\index{obstruction class!three-term decomposition}
The three terms in~\eqref{V1-eq:thqg-VII-obs-class} are disjoint
because they correspond to different codimension-$1$ boundary strata
of $\overline{\cM}_{g,n}$: separating nodes (clutching two
lower-genus surfaces, contributing
$\sum_{g_1+g_2=g} \tfrac{1}{2}[\Theta^{(g_1)},\Theta^{(g_2)}]$),
non-separating nodes (pinching a non-separating cycle on a single
surface, contributing $\Delta(\Theta^{(g-1)})$), and boundary
corrections from the log-FM strata (no topology change in the
surface but corrections from the planted-forest compactification,
contributing $\dsew^{(g)} + \dpf^{(g)}$). These are distinct
codimension-$1$ strata with no overlap: a degeneration is either
separating, non-separating, or a log-FM correction, never more
than one simultaneously.
\end{remark}

\begin{remark}[The obstruction class is a cocycle]
% label removed: rem:thqg-VII-obs-cocycle
\index{obstruction class!cocycle property}
The right side of~\eqref{V1-eq:thqg-VII-genus-g-mc-eq} is
$d_{\Theta^{(0)}}$-closed. Indeed, applying
$d_{\Theta^{(0)}}$ to~\eqref{V1-eq:thqg-VII-genus-g-mc-eq}:
\begin{equation}% label removed: eq:thqg-VII-obs-closed
d_{\Theta^{(0)}}\bigl(\Ob_g(\Theta^{(<g)})\bigr)
\;=\;
d_{\Theta^{(0)}}^2(\Theta^{(g)})
\;=\; 0,
\end{equation}
so $\Ob_g(\Theta^{(<g)})$ defines a cohomology class
\begin{equation}% label removed: eq:thqg-VII-obs-class-coh
[\Ob_g]
\;\in\;
H^2\bigl(\gAmod[\text{genus }g],\; d_{\Theta^{(0)}}\bigr).
\end{equation}
The solvability of the recursion at genus~$g$ is equivalent to
the vanishing of this class. This vanishing is \emph{automatic}
by the bar-intrinsic construction
(Theorem~\ref*{V1-thm:mc2-bar-intrinsic}): the global element
$\Theta_\cA = D_\cA - \dzero$ exists by virtue of $D_\cA^2 = 0$,
hence $\Ob_g$ is exact for every~$g$.
\end{remark}

\subsubsection{Explicit low-genus MC equations}
% label removed: subsubsec:thqg-VII-low-genus

The MC equation~\eqref{V1-eq:thqg-VII-genus-g-mc-eq} at each genus
from $g = 0$ through $g = 4$:

\begin{proposition}[Genus-$0$ MC equation; \ClaimStatusProvedHere]
% label removed: prop:thqg-VII-genus0
At genus $g = 0$, the MC equation reduces to the $\mathrm{Ch}_\infty$-algebra
structure equation on $\cA$:
\begin{equation}% label removed: eq:thqg-VII-mc-g0
D_0(\Theta^{(0)})
+ \tfrac{1}{2}[\Theta^{(0)},\, \Theta^{(0)}]
\;=\; 0.
\end{equation}
This is the statement that $\Theta^{(0)}$, the collection of
tree-level OPE structure maps, defines a homotopy chiral algebra
structure on~$\cA$.
\end{proposition}

\begin{proof}
At genus~$0$, the obstruction class $\Ob_0$ is empty (there are no
lower-genus terms), and the MC equation is simply the first equation
in the genus tower. The identity $D_0^2 = 0$ ensures that the
Arnold relation is satisfied, giving $d^2 = 0$ on the genus-$0$
bar complex.
\end{proof}

\begin{proposition}[Genus-$1$ MC equation; \ClaimStatusProvedHere]
\label{prop:thqg-VII-genus1}
\index{genus-1 MC equation}
At genus $g = 1$:
\begin{equation}% label removed: eq:thqg-VII-mc-g1
d_{\Theta^{(0)}}(\Theta^{(1)})
\;=\;
-\Delta(\Theta^{(0)})
- \dsew^{(1)},
\end{equation}
where $\Delta(\Theta^{(0)})$ is the non-separating BV operator
applied to the tree-level data, and $\dsew^{(1)}$ encodes the
genus-$1$ separating clutching (the correction from the boundary
divisor in $\overline{\cM}_{1,n}$ corresponding to a separating
node producing a genus-$0$ and a genus-$1$ component).

There are no bracket terms on the right
(since $g_1, g_2 \geq 1$ with $g_1 + g_2 = 1$ has no solutions)
and no planted-forest correction (the planted-forest
strata do not contribute at genus~$1$).

The separating clutching at genus~$1$ acts explicitly as
\begin{equation}% label removed: eq:thqg-VII-dsew1-explicit
\dsew^{(1)}(\varphi)
\;=\;
\sum_i \Delta_i(\varphi),
\end{equation}
where $\Delta_i$ contracts the $i$-th pair of inputs using
the invariant pairing and the clutching map
$\xi \colon \overline{\cM}_{0,n+2} \to \overline{\cM}_{1,n}$.
For the Heisenberg algebra at genus~$1$ with no external legs:
$\dsew^{(1)} = \kappa \cdot \Tr(\id) \cdot [\overline{\cM}_{1,1}]$,
where $\Tr$ is the supertrace. This term vanishes for the Heisenberg
because the separating boundary divisor
$\delta_0 \in \partial\overline{\cM}_{1,1}$ factorizes through two
genus-$0$ components with scalar data, producing a contribution
that cancels against the non-separating term in the one-channel case.
\end{proposition}

\begin{proof}
Specialize~\eqref{V1-eq:thqg-VII-genus-g-mc-eq} to $g = 1$. The
sum $\sum_{g_1+g_2=1, g_i \geq 1}$ is empty.
The planted-forest correction $\dpf^{(1)}$ vanishes because
the log-FM planted-forest boundary strata of
$\overline{\cM}_{1,n}$ are empty: every planted forest of
positive loop genus has first Betti number at least~$2$,
which exceeds genus~$1$. The only boundary stratum is the
non-separating node
$\delta_{\mathrm{irr}} \in \partial\overline{\cM}_{1,n}$.
\end{proof}

\begin{proposition}[Genus-$2$ MC equation; \ClaimStatusProvedHere]
\label{prop:thqg-VII-genus2}
\index{genus-2 MC equation}
At genus $g = 2$:
\begin{equation}% label removed: eq:thqg-VII-mc-g2
d_{\Theta^{(0)}}(\Theta^{(2)})
\;=\;
-\tfrac{1}{2}[\Theta^{(1)},\, \Theta^{(1)}]
- \Delta(\Theta^{(1)})
- \dsew^{(2)}
- \dpf^{(2)}.
\end{equation}
The right side decomposes into three shells, present at
\emph{every} degree~$k$:
\begin{enumerate}[label=\textup{(\roman*)}]
\item the \emph{loop-squared shell}
 $\Theta^{(2)}_{\mathrm{loop}^2} :=
 -\tfrac{1}{2}[\Theta^{(1)},\Theta^{(1)}]$,
 corresponding to the separating clutching
 $\sum_{g_1+g_2=2}[\Theta^{(g_1)},\Theta^{(g_2)}]
 = [\Theta^{(1)},\Theta^{(1)}]$
 (a genus-$0$ surface with two non-separating nodes);
\item the \emph{separating-loop shell}
 $\Theta^{(2)}_{\mathrm{sep}\circ\mathrm{loop}} :=
 -\Delta(\Theta^{(1)}) - \dsew^{(2)}$,
 corresponding to the non-separating BV operator
 $\Delta(\Theta^{(1)})$ (pinching a non-separating cycle)
 and the separating boundary strata $\dsew^{(2)}$;
\item the \emph{planted-forest shell}
 $\Theta^{(2)}_{\mathrm{pf}} :=
 -\dpf^{(2)}$,
 corresponding to the log-FM boundary strata indexed by
 depth-$1$ planted forests at genus~$2$.
\end{enumerate}
At genus~$2$ with degree~$k$, the source term $\Ob_2$ decomposes
as: (a)~separating clutching
$[\Theta^{(1)},\Theta^{(1)}]$, (b)~non-separating BV
$\Delta(\Theta^{(1)})$, and (c)~boundary strata contributions
from $\overline{\cM}_{2,k}$. These three shells are present
at every degree~$k$, not only at $(2,0)$: the boundary
stratification of $\overline{\cM}_{2,k}$ has the same three
types of codimension-$1$ strata (separating node, non-separating
node, log-FM correction) for all~$k$.
\end{proposition}

\begin{proof}
Specialize~\eqref{V1-eq:thqg-VII-genus-g-mc-eq} to $g = 2$. The
sum $\sum_{g_1+g_2=2, g_i \geq 1}$ has one term: $g_1 = g_2 = 1$,
contributing $\tfrac{1}{2}[\Theta^{(1)},\Theta^{(1)}]$. The
three-shell decomposition follows from the source of each
contribution: the bracket term from the separating node
connecting two genus-$1$ components, the BV term from the
non-separating node applied to genus-$1$ data, and the
planted-forest term from the Mok boundary strata.
\end{proof}

\begin{proposition}[Genus-$3$ MC equation; \ClaimStatusProvedHere]
\label{prop:thqg-VII-genus3}
\index{genus-3 MC equation}
At genus $g = 3$:
\begin{equation}% label removed: eq:thqg-VII-mc-g3
d_{\Theta^{(0)}}(\Theta^{(3)})
\;=\;
-[\Theta^{(1)},\, \Theta^{(2)}]
- \Delta(\Theta^{(2)})
- \dsew^{(3)}
- \dpf^{(3)}.
\end{equation}
The bracket sum $\sum_{g_1+g_2=3, g_i \geq 1}$ has two terms:
$(g_1,g_2) = (1,2)$ and $(2,1)$, which contribute
$[\Theta^{(1)},\Theta^{(2)}]$ by graded antisymmetry
(the factor $\tfrac{1}{2}$ is cancelled by the two orderings).
\end{proposition}

\begin{proof}
Specialize~\eqref{V1-eq:thqg-VII-genus-g-mc-eq} to $g = 3$. The
sum has $(g_1,g_2) \in \{(1,2),(2,1)\}$:
$\tfrac{1}{2}\bigl([\Theta^{(1)},\Theta^{(2)}]
+ [\Theta^{(2)},\Theta^{(1)}]\bigr)
= \tfrac{1}{2}\bigl([\Theta^{(1)},\Theta^{(2)}]
- [\Theta^{(1)},\Theta^{(2)}]\bigr)
+ [\Theta^{(1)},\Theta^{(2)}]
= [\Theta^{(1)},\Theta^{(2)}]$,
since $[\Theta^{(2)},\Theta^{(1)}] = -[\Theta^{(1)},
\Theta^{(2)}]$ (both elements are odd degree).
\end{proof}

\begin{proposition}[Genus-$4$ MC equation; \ClaimStatusProvedHere]
\label{prop:thqg-VII-genus4}
\index{genus-4 MC equation}
At genus $g = 4$:
\begin{equation}% label removed: eq:thqg-VII-mc-g4
d_{\Theta^{(0)}}(\Theta^{(4)})
\;=\;
-[\Theta^{(1)},\, \Theta^{(3)}]
- \tfrac{1}{2}[\Theta^{(2)},\, \Theta^{(2)}]
- \Delta(\Theta^{(3)})
- \dsew^{(4)}
- \dpf^{(4)}.
\end{equation}
The bracket sum has three terms:
$(g_1,g_2) \in \{(1,3),(2,2),(3,1)\}$, yielding
$[\Theta^{(1)},\Theta^{(3)}]
+ \tfrac{1}{2}[\Theta^{(2)},\Theta^{(2)}]$.
\end{proposition}

\begin{proof}
The computation parallels the previous cases. The three terms
in the bracket sum are:
\begin{align*}
&\tfrac{1}{2}\bigl(
[\Theta^{(1)},\Theta^{(3)}]
+ [\Theta^{(2)},\Theta^{(2)}]
+ [\Theta^{(3)},\Theta^{(1)}]
\bigr)\\
&= \tfrac{1}{2}\bigl(
2[\Theta^{(1)},\Theta^{(3)}]
+ [\Theta^{(2)},\Theta^{(2)}]
\bigr)\\
&= [\Theta^{(1)},\Theta^{(3)}]
+ \tfrac{1}{2}[\Theta^{(2)},\Theta^{(2)}],
\end{align*}
using $[\Theta^{(3)},\Theta^{(1)}]
= -[\Theta^{(1)},\Theta^{(3)}]$ and
$[\Theta^{(2)},\Theta^{(2)}] \neq 0$ in general
(it vanishes only in the Gaussian case).
\end{proof}

\begin{remark}[General pattern]
% label removed: rem:thqg-VII-general-pattern
\index{genus-$g$ MC equation!general structure}
At genus $g \geq 1$, the obstruction class
$\Ob_g(\Theta^{(<g)})$ has $\lfloor g/2 \rfloor + 2$ distinct
contributions:
\begin{enumerate}[label=(\alph*)]
\item $\lfloor g/2 \rfloor$ bracket terms
 $[\Theta^{(g_1)},\Theta^{(g_2)}]$ for $1 \leq g_1 \leq g_2$,
 $g_1 + g_2 = g$;
\item one BV term $\Delta(\Theta^{(g-1)})$;
\item boundary-stratum corrections $\dsew^{(g)} + \dpf^{(g)}$.
\end{enumerate}
The bracket terms record the \emph{separating degenerations}
(pinching a cycle that separates the surface into two parts of
genera $g_1$ and $g_2$). The BV term records the
\emph{non-separating degeneration} (pinching a non-separating
cycle, which lowers the genus by one). The boundary terms encode
higher-codimension strata. The structure is:
\begin{equation}% label removed: eq:thqg-VII-obs-general
\Ob_g
\;=\;
\underbrace{
\sum_{\substack{g_1+g_2=g \\ g_1,g_2 \geq 1}}
\tfrac{1}{2}[\Theta^{(g_1)},\Theta^{(g_2)}]
}_{\text{separating}}
\;+\;
\underbrace{
\Delta(\Theta^{(g-1)})
}_{\text{non-separating}}
\;+\;
\underbrace{
\dsew^{(g)} + \dpf^{(g)}
}_{\text{boundary strata}}.
\end{equation}
This is the genus analogue of the Postnikov tower in
obstruction theory: the extension from genus~$<g$ to genus~$g$
is controlled by a class in $H^2$ of the tree-twisted complex.
\end{remark}

\subsubsection{The bracket and BV terms in graph-sum language}
% label removed: subsubsec:thqg-VII-graph-sum

The bracket $[\Theta^{(g_1)},\Theta^{(g_2)}]$ in the
convolution algebra $\gAmod$ has a graph-theoretic interpretation.
A genus-$g_1$ graph $\Gamma_1$ and a genus-$g_2$ graph $\Gamma_2$
are joined by a single edge to produce a genus-$(g_1{+}g_2)$
graph:
\begin{equation}% label removed: eq:thqg-VII-bracket-graph
[\Theta^{(g_1)},\Theta^{(g_2)}]
\;=\;
\sum_{\substack{\Gamma_1 \in \mathsf{Gr}_{g_1}^{\mathrm{st}}\\
 \Gamma_2 \in \mathsf{Gr}_{g_2}^{\mathrm{st}}}}
\frac{1}{|\Aut(\Gamma_1)||\Aut(\Gamma_2)|}
\;\ell_{\Gamma_1 \cup_e \Gamma_2},
\end{equation}
where $\ell_\Gamma$ is the graph amplitude
(vertex labels from the transferred cyclic minimal model, edge
contractions by the complementarity propagator
$P_\cA = H_\cA^{-1}$), and $\Gamma_1 \cup_e \Gamma_2$ denotes
the graph obtained by joining $\Gamma_1$ and $\Gamma_2$ with a
single new edge.

The BV operator $\Delta(\Theta^{(g-1)})$ has the interpretation
of closing a single edge within a genus-$(g{-}1)$ graph to
produce a genus-$g$ graph:
\begin{equation}% label removed: eq:thqg-VII-bv-graph
\Delta(\Theta^{(g-1)})
\;=\;
\sum_{\Gamma \in \mathsf{Gr}_{g-1}^{\mathrm{st}}}
\frac{1}{|\Aut(\Gamma)|}
\sum_{v \in V(\Gamma)}
\ell_{\Gamma^v},
\end{equation}
where $\Gamma^v$ denotes the graph obtained by adding a loop edge
at vertex~$v$ of~$\Gamma$ (joining two half-edges at $v$ with
the propagator, increasing the loop number by one).

\begin{remark}[Feynman diagram interpretation]
% label removed: rem:thqg-VII-feynman-interpretation
\index{Feynman diagrams!MC recursion}
Equations~\eqref{V1-eq:thqg-VII-bracket-graph}
and~\eqref{V1-eq:thqg-VII-bv-graph} are the standard Feynman
diagram operations: the bracket is the tree-level connecting
propagator (joining two subdiagrams with one line), and the BV
operator is the loop closure (contracting one internal line into
a loop). The MC recursion is thus equivalent to the statement
that the genus-$g$ amplitude is determined by lower-genus
amplitudes through these two operations plus boundary corrections.
This is the chiral-algebraic version of the cut-and-sew relations
in string theory.
\end{remark}

% ======================================================================
%
% PART 2: THE MC RECURSION
%
% ======================================================================

\subsection{The MC recursion: unique determination from lower genera}
% label removed: subsec:thqg-VII-mc-recursion
\index{MC recursion|textbf}
\index{contracting homotopy!MC recursion}

\subsubsection{The contracting homotopy}
% label removed: subsubsec:thqg-VII-contracting-homotopy

The solvability of the recursion~\eqref{V1-eq:thqg-VII-genus-g-mc-eq}
rests on two facts:
\begin{enumerate}[label=(\alph*)]
\item the obstruction class $\Ob_g$ is
 $d_{\Theta^{(0)}}$-exact (hence the equation has a
 solution), and
\item the solution is unique up to gauge equivalence
 (hence the genus-$g$ partition function is determined).
\end{enumerate}
Both facts follow from the existence of a contracting homotopy
for the tree-twisted complex.

\begin{theorem}[Contracting homotopy; \ClaimStatusProvedHere]
\label{thm:thqg-VII-contracting-homotopy}
\index{contracting homotopy|textbf}
Let $\cA$ be a modular Koszul chiral algebra with bar-intrinsic
MC element $\Theta_\cA = D_\cA - \dzero$
\textup{(}Theorem~\textup{\ref*{V1-thm:mc2-bar-intrinsic})}. Denote
the tree-level data $\Theta^{(0)} \in \gr^0_F\,\gAmod$. Then:

\begin{enumerate}[label=\textup{(\roman*)}]
\item The tree-twisted complex $(\gAmod, d_{\Theta^{(0)}})$
 admits a contracting homotopy
 $\hCont \colon {\gAmod}^{\bullet}[\mathrm{genus\,}g]
 \to {\gAmod}^{\bullet - 1}[\mathrm{genus\,}g]$
 satisfying
 \begin{equation}% label removed: eq:thqg-VII-homotopy
 d_{\Theta^{(0)}} \circ \hCont
 + \hCont \circ d_{\Theta^{(0)}}
 \;=\; \id - \pi_H,
 \end{equation}
 where $\pi_H$ is the projection onto the cohomology
 $H^*(\gAmod[\mathrm{genus\,}g], d_{\Theta^{(0)}})$.

\item The cohomology $H^*(\gAmod[\mathrm{genus\,}g],
 d_{\Theta^{(0)}})$ vanishes in degree~$2$ for $g \geq 1$:
 \begin{equation}% label removed: eq:thqg-VII-h2-vanishing
 H^2\bigl(\gAmod[\mathrm{genus\,}g],\; d_{\Theta^{(0)}}\bigr)
 \;=\; 0, \qquad g \geq 1.
 \end{equation}
 In particular, every $d_{\Theta^{(0)}}$-cocycle of
 degree~$2$ is exact, and the obstruction class $[\Ob_g] = 0$.
\end{enumerate}
\end{theorem}

\begin{proof}
\emph{Part (i).} The contracting homotopy is constructed by
homological perturbation. Start with the untwisted complex
$(\gAmod, D_0)$ at genus~$g$. This is a chain complex of
cochain-valued functions on moduli spaces, and its cohomology
at genus~$g$ is computed by the $E_1$~page of the genus
spectral sequence
(Construction~\ref{V1-const:vol1-genus-spectral-sequence}).
The untwisted complex $(\gAmod[\text{genus }g], D_0)$
admits a standard contracting homotopy $\hCont_0$ from the
Hodge-type decomposition (choosing a metric on $\overline{\cM}_{g,n}$
gives a Hodge decomposition, and the Green's operator provides
$\hCont_0$).

The twist by $\Theta^{(0)}$ replaces $D_0$ with $d_{\Theta^{(0)}}
= D_0 + \ad_{\Theta^{(0)}}$. By the homological perturbation
lemma (HPL), the contracting homotopy for the twisted complex is
\begin{equation}% label removed: eq:thqg-VII-hpl-homotopy
\hCont
\;=\;
\hCont_0 \circ
\sum_{n \geq 0}
\bigl(-\ad_{\Theta^{(0)}} \circ \hCont_0\bigr)^n
\;=\;
\hCont_0 \circ
\bigl(\id + \ad_{\Theta^{(0)}} \circ \hCont_0\bigr)^{-1},
\end{equation}
which converges because the degree filtration makes
$\ad_{\Theta^{(0)}} \circ \hCont_0$ nilpotent on each
finite-degree component (the degree-cutoff of
Lemma~\ref{V1-lem:degree-cutoff} ensures that the series terminates
at each finite weight).

\medskip\noindent
\emph{Part (ii).} The vanishing of $H^2$ at positive genus
follows from the Koszul property. At genus~$0$, the Koszul
assumption gives diagonal Ext vanishing for $\cA$, which
translates to acyclicity of the tree-twisted bar complex in
the relevant degrees. At genus $g \geq 1$, the tree-twisted
differential $d_{\Theta^{(0)}}$ on the genus-$g$ sector of
$\gAmod$ is a deformation of the genus-$0$ differential
$D_0$; the deformation is controlled by the degree filtration,
and the acyclicity propagates by induction on degree
(each degree component is a finite-dimensional perturbation
of the acyclic genus-$0$ complex).

More precisely, the global solution $\Theta_\cA = D_\cA - \dzero$
(Theorem~\ref*{V1-thm:mc2-bar-intrinsic}) exists at all genera.
Its existence implies that every $\Ob_g$ is exact, hence
$H^2 = 0$ at every positive genus. This is not circular: the
bar-intrinsic construction produces $\Theta_\cA$ directly
from $D_\cA^2 = 0$, without solving the recursion step by step.
The recursion is a \emph{decomposition} of the already-existing
element into genus components, not a \emph{construction} that
might fail.
\end{proof}

\subsubsection{The recursion algorithm}
% label removed: subsubsec:thqg-VII-recursion-algorithm

\begin{construction}[MC recursion; \ClaimStatusProvedHere]
\label{constr:thqg-VII-mc-recursion}
\index{MC recursion!algorithm}
Given the genus-$0$ data $\Theta^{(0)}$ and the contracting
homotopy $\hCont$ of
Theorem~\textup{\ref{thm:thqg-VII-contracting-homotopy}},
the genus-$g$ MC element is computed recursively:

\medskip
\noindent
\textbf{Input:} $\Theta^{(0)}, \Theta^{(1)}, \ldots,
\Theta^{(g-1)}$.

\medskip
\noindent
\textbf{Step 1.} Compute the obstruction class:
\begin{equation}% label removed: eq:thqg-VII-recursion-obs
\Ob_g
\;:=\;
\sum_{\substack{g_1+g_2=g \\ g_1,g_2 \geq 1}}
\tfrac{1}{2}[\Theta^{(g_1)},\Theta^{(g_2)}]
\;+\;
\Delta(\Theta^{(g-1)})
\;+\;
\dsew^{(g)} + \dpf^{(g)}.
\end{equation}

\textbf{Step 2.} Verify the cocycle condition:
\begin{equation}% label removed: eq:thqg-VII-recursion-cocycle
d_{\Theta^{(0)}}(\Ob_g) = 0.
\end{equation}
This is a consistency check; it must hold by the inductive
hypothesis.

\textbf{Step 3.} Apply the contracting homotopy:
\begin{equation}% label removed: eq:thqg-VII-recursion-solve
\Theta^{(g)}
\;:=\;
-\hCont\bigl(\Ob_g\bigr).
\end{equation}

\textbf{Step 4.} Verify:
$d_{\Theta^{(0)}}(\Theta^{(g)})
= -(\id - \pi_H)(\Ob_g)
= -\Ob_g$
(since $\pi_H(\Ob_g) = 0$ by the vanishing of $H^2$).

\medskip
\noindent
\textbf{Output:} $\Theta^{(g)}$, uniquely determined up to
gauge equivalence.
\end{construction}

\begin{theorem}[Uniqueness up to gauge; \ClaimStatusProvedHere]
% label removed: thm:thqg-VII-uniqueness
\index{MC recursion!uniqueness}
The genus-$g$ MC element $\Theta^{(g)}$ is unique up to the
action of the gauge group
$\exp\bigl({\gAmod}^0[\mathrm{genus\,}g]\bigr)$.
Two solutions $\Theta^{(g)}$ and $\widetilde{\Theta}^{(g)}$
of~\eqref{V1-eq:thqg-VII-genus-g-mc-eq} with the same
lower-genus data $\Theta^{(<g)}$ satisfy
\begin{equation}% label removed: eq:thqg-VII-gauge-diff
\Theta^{(g)} - \widetilde{\Theta}^{(g)}
\;\in\;
\ker(d_{\Theta^{(0)}})
\;=\;
\im(d_{\Theta^{(0)}})
\;\oplus\;
H^1(\gAmod[\text{genus }g], d_{\Theta^{(0)}}),
\end{equation}
and the gauge group acts transitively on the fiber over any
fixed obstruction class.
\end{theorem}

\begin{proof}
Let $\Theta^{(g)}$ and $\widetilde{\Theta}^{(g)}$ be two
solutions with the same obstruction class. Their difference
$\xi := \Theta^{(g)} - \widetilde{\Theta}^{(g)}$ satisfies
$d_{\Theta^{(0)}}(\xi) = 0$. The gauge group
$G := \exp({\gAmod}^0[\text{genus }g])$ is pro-unipotent
(the weight filtration provides the lower central series), hence
acts freely on degree-$1$ cocycles. The gauge orbit of~$\xi$
is exactly the set of degree-$1$ cocycles cohomologous to~$\xi$.
Since $H^2 = 0$
(Theorem~\ref{thm:thqg-VII-contracting-homotopy}(ii)),
the ambiguity lies in $H^1$ and is absorbed by the gauge action.

The bar-intrinsic element $\Theta^{(g)} = (D_\cA - \dzero)^{(g)}$
is the canonical representative: it is the unique element in
the gauge orbit selected by the property $D_\cA^2 = 0$ (which
fixes the contracting homotopy to be the one induced by the
bar differential).
\end{proof}

\subsubsection{Base case: genus $g = 1$}
% label removed: subsubsec:thqg-VII-base-case

\begin{proposition}[Genus-$1$ solution; \ClaimStatusProvedHere]
\label{prop:thqg-VII-genus1-solution}
\index{genus-1 MC element}
The genus-$1$ MC element is
\begin{equation}% label removed: eq:thqg-VII-theta1
\Theta^{(1)}
\;=\;
-\hCont\bigl(
\Delta(\Theta^{(0)}) + \dsew^{(1)}
\bigr).
\end{equation}
For a one-channel algebra (e.g.\ the Heisenberg, the affine
Kac--Moody at generic level), this reduces to
\begin{equation}% label removed: eq:thqg-VII-theta1-scalar
\Theta^{(1)}
\;=\;
\kappa(\cA) \cdot [\delta_{\mathrm{irr}}] \otimes \Delta,
\end{equation}
where $[\delta_{\mathrm{irr}}] \in H^1(\overline{\cM}_{1,1})$
is the class of the irreducible boundary divisor and $\Delta$ is
the BV operator (the propagator inserted at the non-separating
node). The scalar $\kappa(\cA)$ is the modular characteristic
\textup{(}Theorem~D\textup{)}.
\end{proposition}

\begin{proof}
The obstruction class at genus~$1$ is
$\Ob_1 = \Delta(\Theta^{(0)}) + \dsew^{(1)}$
(Proposition~\ref{prop:thqg-VII-genus1}). The solution is
$\Theta^{(1)} = -\hCont(\Ob_1)$ by
Construction~\ref{constr:thqg-VII-mc-recursion}.

For a one-channel algebra, the BV operator $\Delta$ applied to
the genus-$0$ data $\Theta^{(0)}$ produces a scalar:
$\Delta(\Theta^{(0)}) = \kappa(\cA) \cdot \omega_1$, where
$\omega_1$ is the genus-$1$ modular form (the holomorphic
$1$-form on $\overline{\cM}_{1,1}$). The contracting homotopy
inverts the tree-twisted differential, and the result is the
scalar expression~\eqref{V1-eq:thqg-VII-theta1-scalar}.
The separating clutching $\dsew^{(1)}$ vanishes for genus~$1$
of the one-channel algebra because the separating boundary
$\delta_0 \in \partial\overline{\cM}_{1,1}$ decomposes the
surface into genus-$0$ and genus-$1$ components, and the
genus-$1$ component at the first step of the recursion
contributes only through $\Delta$.
\end{proof}

\subsubsection{Worked example: genus $g = 2$}
% label removed: subsubsec:thqg-VII-genus2-worked

\begin{computation}[Genus-$2$ MC recursion; \ClaimStatusProvedHere]
% label removed: comp:thqg-VII-genus2
\index{genus-2 MC element!worked computation}
We carry out the recursion for genus $g = 2$ in full detail.

\medskip\noindent\emph{Step 1: Compute the obstruction.}
From Proposition~\ref{prop:thqg-VII-genus2}:
\begin{equation}% label removed: eq:thqg-VII-obs2-full
\Ob_2
\;=\;
\tfrac{1}{2}[\Theta^{(1)},\Theta^{(1)}]
+ \Delta(\Theta^{(1)})
+ \dsew^{(2)} + \dpf^{(2)}.
\end{equation}

\medskip\noindent\emph{Step 2: Evaluate each term.}

\emph{(a) The bracket term.} Using
$\Theta^{(1)} = \kappa \cdot [\delta_{\mathrm{irr}}]
\otimes \Delta$ (for a one-channel algebra):
\begin{align}
\tfrac{1}{2}[\Theta^{(1)},\Theta^{(1)}]
&= \tfrac{1}{2}\,\kappa^2\,
 [[\delta_{\mathrm{irr}}] \otimes \Delta,\,
 [\delta_{\mathrm{irr}}] \otimes \Delta]
\notag\\
&= \tfrac{1}{2}\,\kappa^2\,
 [\delta_{\mathrm{irr}} \times \delta_{\mathrm{irr}}]
 \otimes \Delta \circ \Delta,
% label removed: eq:thqg-VII-bracket-term-g2
\end{align}
where $\delta_{\mathrm{irr}} \times \delta_{\mathrm{irr}}$
is the class in $H^*(\overline{\cM}_{2,n})$ corresponding to the
boundary stratum with two non-separating nodes, and
$\Delta \circ \Delta$ is the iterated BV operator (double
trace on the propagator).

\emph{(b) The BV term.} The BV operator applied to the genus-$1$
element:
\begin{equation}% label removed: eq:thqg-VII-bv-term-g2
\Delta(\Theta^{(1)})
= \kappa \cdot \Delta([\delta_{\mathrm{irr}}] \otimes \Delta)
= \kappa \cdot [\Delta(\delta_{\mathrm{irr}})]
 \otimes \Delta^2,
\end{equation}
where $\Delta(\delta_{\mathrm{irr}})$ is the pushforward of
the irreducible boundary class through the non-separating
clutching.

\emph{(c) The boundary terms.} The separating clutching
$\dsew^{(2)}$ and the planted-forest correction $\dpf^{(2)}$
encode the codimension-$1$ and codimension-$2$ boundary strata
of $\overline{\cM}_{2,n}$.

\medskip\noindent\emph{Step 3: Apply the homotopy.}
\begin{equation}% label removed: eq:thqg-VII-theta2-solution
\Theta^{(2)}
= -\hCont\bigl(\Ob_2\bigr)
= -\hCont\bigl(
 \tfrac{1}{2}\kappa^2\,
 [\delta_{\mathrm{irr}}^2] \otimes \Delta^2
 + \kappa\,[\Delta(\delta_{\mathrm{irr}})]
 \otimes \Delta^2
 + \dsew^{(2)} + \dpf^{(2)}
 \bigr).
\end{equation}

\medskip\noindent\emph{Step 4: Shell decomposition.} The
three shells (Proposition~\ref{prop:thqg-VII-genus2}) are:
\begin{align}
\Theta^{(2)}_{\mathrm{loop}^2}
&= -\hCont\bigl(\tfrac{1}{2}[\Theta^{(1)},\Theta^{(1)}]\bigr)
\neq 0 \text{ for all $\kappa \neq 0$},
% label removed: eq:thqg-VII-shell-loop2\\
\Theta^{(2)}_{\mathrm{sep}\circ\mathrm{loop}}
&= -\hCont\bigl(\Delta(\Theta^{(1)}) + \dsew^{(2)}\bigr)
\neq 0 \text{ iff } \fC(\cA) \neq 0,
% label removed: eq:thqg-VII-shell-sep\\
\Theta^{(2)}_{\mathrm{pf}}
&= -\hCont\bigl(\dpf^{(2)}\bigr)
\neq 0 \text{ iff } S_3 \neq 0
\;\text{(cubic shadow nonzero, from genus-$2$ planted-forest graphs)}.
% label removed: eq:thqg-VII-shell-pf
\end{align}
By the genus-$2$ shell activation theorem
(Theorem~\ref{V1-rem:genus2-shell-activation}), the separating shell~$\sigma$
distinguishes $\{G,C\}$ from $\{L,M\}$; the planted-forest shell
depends on $S_3$ (not just $S_4$), so classes $L$ and~$M$ share the
same genus-$2$ pattern. The full four-class diagnostic requires the
critical discriminant $\Delta = 8\kappa S_4$:
\begin{center}
\begin{tabular}{ccccc}
$\sigma$ & $\pi$ & $\Delta$ & Class & Example \\
\hline
$0$ & $0$ & $0$ & G (Gaussian) & $\cH_k$ \\
$1$ & $1$ & $0$ & L (Lie/tree) & $V_k(\fg)$ \\
$0$ & $0$ & $\neq 0$ & C (contact) & $\beta\gamma$ \\
$1$ & $1$ & $\neq 0$ & M (mixed) & $\mathrm{Vir}_c$
\end{tabular}
\end{center}
\end{computation}

\begin{remark}[Genus-$2$ as the first diagnostic]
% label removed: rem:thqg-VII-genus2-diagnostic
\index{genus-2 MC element!diagnostic role}
Genus~$2$ is the first level at which the MC recursion
distinguishes between the four shadow-depth classes. At genus~$1$,
the MC element $\Theta^{(1)}$ is controlled entirely by the scalar
$\kappa$; all four classes look identical at genus~$1$. At genus~$2$,
the three shells probe the cubic and quartic shadow data, and the
vanishing pattern separates the classes. This is a concrete
manifestation of the principle that the modular bootstrap becomes
nontrivial at genus~$2$.
\end{remark}

% ======================================================================
%
% PART 3: THE GENUS SPECTRAL SEQUENCE
%
% ======================================================================

\subsection{The genus spectral sequence}
% label removed: subsec:thqg-VII-genus-spectral-sequence
\index{genus spectral sequence|textbf}
\index{spectral sequence!genus filtration!bootstrap interpretation}

The MC recursion is organized by a spectral sequence that arises
from the genus filtration on $\gAmod$. This spectral sequence was
constructed in
Construction~\ref{V1-const:vol1-genus-spectral-sequence}; here we
develop its bootstrap interpretation and prove convergence.

\subsubsection{Construction}
% label removed: subsubsec:thqg-VII-ss-construction

\begin{theorem}[Genus spectral sequence for the bootstrap;
\ClaimStatusProvedHere]
\label{thm:thqg-VII-genus-ss}
\index{genus spectral sequence!construction}
The genus filtration~\eqref{V1-eq:thqg-VII-genus-filtration}
on the modular convolution algebra $\gAmod$ induces a
spectral sequence
\begin{equation}% label removed: eq:thqg-VII-genus-ss
E_1^{p,q}
\;\Longrightarrow\;
H^{p+q}\bigl(\gAmod,\; D\bigr),
\end{equation}
with the following structure:

\begin{enumerate}[label=\textup{(\roman*)}]
\item \emph{$E_0$~page:}
 $E_0^{p,q} = \gr^p_F\,{\gAmod}^{p+q}$,
 the genus-$p$ component in total degree $p + q$.
 The differential is $d_0 = D_0$ restricted to
 $\gr^p_F$, which is the genus-preserving part of~$D$.

\item \emph{$E_1$~page:}
 \begin{equation}% label removed: eq:thqg-VII-E1
 E_1^{p,q}
 \;=\;
 H^{p+q}\bigl(\gr^p_F\,\gAmod,\; D_0\bigr),
 \end{equation}
 the cohomology of the genus-$p$ complex with
 genus-preserving differential. The three columns are:
 \begin{itemize}
 \item $p = 0$: \emph{tree-level cohomology}
 $H^q(\gAmod[\text{genus }0], D_0)$, which is the
 bar cohomology $H^*(\barB(\cA))$. For Koszul~$\cA$,
 this is concentrated on the diagonal.
 \item $p = 1$: \emph{one-loop cohomology}
 $H^{1+q}(\gAmod[\text{genus }1], D_0)$, encoding
 the genus-$1$ bar complex modulo the genus-$0$
 differential.
 \item $p \geq 2$: \emph{higher-genus shells},
 computed by the moduli-space cohomology of
 $\overline{\cM}_{p,n}$ tensored with the cyclic
 bar data.
 \end{itemize}

\item \emph{$d_1$ differential:}
 $d_1 \colon E_1^{p,q} \to E_1^{p+1,q}$.
 This is the obstruction map
 $\Ob_{p+1} \colon E_1^{p,*} \to E_1^{p+1,*-1}$,
 which sends the genus-$p$ data to the genus-$(p{+}1)$
 obstruction class.
 \begin{equation}% label removed: eq:thqg-VII-d1
 d_1
 \;=\;
 [\Ob_{p+1}]
 \;\colon\;
 E_1^{p,q} \to E_1^{p+1,q}.
 \end{equation}
 The $d_1$ differential at $p = 0$ is the genus-$1$
 obstruction $\Ob_1 = \Delta + \dsew^{(1)}$:
 it sends tree-level data to the one-loop correction.

\item \emph{Higher differentials:}
 $d_r \colon E_r^{p,q} \to E_r^{p+r,q-r+1}$.
 The differential $d_r$ is the genus-$r$ obstruction map
 that sends genus-$p$ data through the MC recursion to the
 genus-$(p{+}r)$ sector. These encode the multi-loop
 corrections.
\end{enumerate}
\end{theorem}

\begin{proof}
The genus filtration is a complete descending filtration on a
cochain complex. The standard construction (cf.\
\cite[Ch.~5]{Wei94}) produces a spectral sequence with the
stated $E_0$~page and $E_1$~page.

\medskip\noindent\emph{The $d_0$ differential.}
On the associated graded $\gr^p_F$, the differential induced by
$D = D_0 + D_1$ is $D_0$ (the genus-preserving part), since
$D_1$ raises genus and hence maps into the next filtration step.
Thus $E_0^{p,q} = \gr^p_F\,{\gAmod}^{p+q}$ with $d_0 = D_0$.

\medskip\noindent\emph{The $d_1$ differential.}
The $d_1$ differential on the $E_1$~page is induced by $D_1$
on the $D_0$-cohomology. Explicitly: for a $D_0$-cocycle
$\alpha \in \gr^p_F$, choose a lift $\widetilde{\alpha}$ to
$F^p\,\gAmod$. Then $D(\widetilde{\alpha}) \in F^{p+1}$ and
$D_0(D(\widetilde{\alpha})) \in F^{p+1}$; the class of
$D(\widetilde{\alpha})$ modulo $D_0$-exact terms is the $d_1$
differential. By the decomposition $D = D_0 + D_1$:
\[
D(\widetilde{\alpha})
= D_0(\widetilde{\alpha}) + D_1(\widetilde{\alpha}).
\]
Since $D_0(\widetilde{\alpha}) = 0$ (as $\alpha$ is a cocycle),
we get $d_1([\alpha]) = [D_1(\widetilde{\alpha})]$ in
$E_1^{p+1,*}$. The map $D_1$ is exactly the BV operator
$\Delta$ plus the planted-forest corrections $\dpf$, which is
the genus-raising obstruction. Thus $d_1 = [\Ob_{p+1}]$.

\medskip\noindent\emph{Higher differentials.}
The $d_r$ differential is constructed by the standard
extension procedure: starting with a class in $E_r^{p,q}$
(which survives to the $E_r$~page because $d_1, \ldots,
d_{r-1}$ vanish on it), lift to an element of $F^p$, apply
$D$, and project to the genus-$(p{+}r)$ component modulo
earlier terms. The result is the genus-$(p{+}r)$ obstruction
accumulated from $r$ iterations of the genus-raising operations.
\end{proof}

\subsubsection{The $E_1$~page in detail}
% label removed: subsubsec:thqg-VII-E1-detail

\begin{proposition}[$E_1$~page for a Koszul algebra;
\ClaimStatusProvedHere]
% label removed: prop:thqg-VII-E1-koszul
\index{genus spectral sequence!$E_1$ page}
Let $\cA$ be a modular Koszul chiral algebra. Then:
\begin{enumerate}[label=\textup{(\roman*)}]
\item $E_1^{0,q} = H^q(\barB(\cA))$, the bar cohomology of
 $\cA$. By the Koszul property, this is concentrated on the
 diagonal: $E_1^{0,q} = 0$ for $q \neq 0$ (after appropriate
 regrading).
\item $E_1^{1,q}$ is the genus-$1$ bar cohomology, computed
 from the elliptic bar complex
 (Theorem~\ref{V1-thm:elliptic-bar}). For the Heisenberg
 algebra, $E_1^{1,q} = \mathbb{C}$ for $q = 0$ (generated by
 $\kappa \cdot \omega_1$) and vanishes for $q \neq 0$.
\item $E_1^{p,q}$ for $p \geq 2$ is the genus-$p$ bar
 cohomology, which is computed from the moduli-space
 cohomology $H^*(\overline{\cM}_{p,n})$ via the Mumford
 classes. For the Heisenberg algebra, each $E_1^{p,0}$ is
 one-dimensional, generated by $\kappa^p \cdot \lambda_p$.
\end{enumerate}
\end{proposition}

\begin{proof}
Part (i) is the definition of the Koszul property: the bar
cohomology $H^*(\barB(\cA))$ is the Koszul dual $\cA^!$,
concentrated in degree~$0$.

For parts (ii) and (iii) in the Heisenberg case, the genus-$p$
bar complex has differential $D_0 = \dzero$ (no collision
corrections, since the OPE is quadratic and abelian). The
cohomology is determined by the moduli-space integration:
at genus~$p$, the only nontrivial class is the top Hodge class
$\lambda_p = c_p(\mathbb{E})$, paired with the scalar modular
characteristic $\kappa^p$. This gives a one-dimensional
$E_1^{p,0} = \mathbb{C} \cdot \kappa^p \lambda_p$.
\end{proof}

\begin{remark}[The $E_1$~page as the ``classical bootstrap'']
% label removed: rem:thqg-VII-E1-classical
The $E_1$~page isolates the data at each genus before the
inter-genus corrections are turned on. One may think of
$E_1^{p,*}$ as the ``classical'' (i.e.\ tree-level with respect
to the genus filtration) contribution at genus~$p$. The
higher differentials $d_r$ encode the quantum corrections that
propagate from lower genera to higher genera. The spectral
sequence converges to the full MC moduli, which incorporates
all inter-genus interactions.
\end{remark}

\subsubsection{Convergence}
% label removed: subsubsec:thqg-VII-convergence

\begin{theorem}[Convergence of the genus spectral sequence;
\ClaimStatusProvedHere]
\label{thm:thqg-VII-ss-convergence}
\index{genus spectral sequence!convergence}
The spectral
sequence~\eqref{V1-eq:thqg-VII-genus-ss} converges:
\begin{equation}% label removed: eq:thqg-VII-convergence
E_\infty^{p,q}
\;\cong\;
\gr^p_F\, H^{p+q}\bigl(\gAmod,\; D\bigr).
\end{equation}
\end{theorem}

\begin{proof}
The genus filtration is complete
($\gAmod = \varprojlim_g\, \gAmod / F^g$) and
exhaustive ($F^0 = \gAmod$). The Eilenberg--Moore
convergence theorem (cf.\ \cite[Theorem~5.5.1]{Wei94})
applies to complete, exhaustive filtrations on cochain
complexes: the spectral sequence converges to the graded
pieces of the cohomology of the total complex.

The completeness condition is exactly the $\hbar$-adic
completeness of
Remark~\ref{rem:thqg-VII-hbar-topology}: the
$\hbar$-adic topology is the linear topology defined by
the genus filtration, and the convolution algebra is complete
(being defined as an infinite product over genera).

The pronilpotence of $F^1\,\gAmod$ ensures that the spectral
sequence is strongly convergent: for each total degree $n$,
only finitely many nonzero differentials land in a given
bidegree, so the inverse system of quotients stabilizes.
\end{proof}

\begin{corollary}[Degeneration for Gaussian algebras;
\ClaimStatusProvedHere]
\label{cor:thqg-VII-gaussian-degeneration}
\index{genus spectral sequence!degeneration}
\index{Gaussian archetype!spectral sequence}
For a Gaussian algebra $\cA$ \textup{(}shadow depth $r_{\max} = 2$,
e.g.\ the Heisenberg algebra\textup{)}, the genus spectral
sequence degenerates at $E_1$:
\begin{equation}% label removed: eq:thqg-VII-gaussian-degeneration
E_1^{p,q} \;=\; E_\infty^{p,q}.
\end{equation}
\end{corollary}

\begin{proof}
For a Gaussian algebra, the MC element $\Theta_\cA$ is
scalar: $\Theta_\cA = \kappa \cdot \eta \otimes \Lambda$,
with no higher-degree corrections. All inter-genus
brackets $[\Theta^{(g_1)},\Theta^{(g_2)}]$ and BV applications
$\Delta(\Theta^{(g)})$ produce scalar multiples of
$\kappa^{g+1} \cdot \lambda_{g+1}$; the differentials $d_r$
for $r \geq 1$ land on pre-existing classes (they are
multiplied by $\kappa$ at each step but do not create new
cohomological directions). The tree-twisted differential
$d_{\Theta^{(0)}}$ is exact on positive-degree classes at
every genus (the Koszul property), so $d_1 = 0$ on the
$E_1$~page: every obstruction class is automatically solved
by the scalar homotopy, and the spectral sequence collapses.

Concretely, the one-dimensional space $E_1^{p,0}$ survives
to $E_\infty$ because $d_r(E_1^{p,0}) = 0$ for all $r \geq 1$
(the scalar obstruction at each genus is solved by the scalar
homotopy, leaving no residual class on the $E_1$ page).
\end{proof}

\begin{proposition}[Non-degeneration for mixed-type algebras;
\ClaimStatusProvedHere]
\label{prop:thqg-VII-mixed-nondegeneration}
\index{genus spectral sequence!non-degeneration}
\index{mixed archetype!spectral sequence}
For a mixed-type algebra $\cA$ with $r_{\max} = \infty$
\textup{(}e.g.\ the Virasoro algebra\textup{)}, the genus
spectral sequence does not degenerate at any finite page:
for every $r \geq 1$, there exists a bidegree $(p,q)$ with
$d_r \neq 0$ on $E_r^{p,q}$.
\end{proposition}

\begin{proof}
The mixed archetype has $\fC(\cA) \neq 0$ and
$\mathfrak{Q}(\cA) \neq 0$, hence both the separating-loop
shell and the planted-forest shell are activated at genus~$2$
(Theorem~\ref{V1-rem:genus2-shell-activation}). The $d_1$
differential carries the genus-$0$ cubic shadow data into the
genus-$1$ sector, and this propagates: the genus-$r$
obstruction class receives contributions from the
$r$-fold iterated application of the BV operator and the
$r$-fold nested brackets, both of which are nontrivial when
the shadow obstruction tower is infinite.

More precisely, the $d_r$ differential at the $(0,*)$ column
maps tree-level data through the $r$-loop recursion. Since
the shadow obstruction tower $\Theta^{\leq r}$ is strictly larger than
$\Theta^{\leq r-1}$ for every~$r$ (the obstruction
$o_{r+1}(\cA) \neq 0$ by the infinite-depth condition),
the $d_r$ differential is nontrivial for each~$r$.
\end{proof}

\subsubsection{The spectral sequence and the shadow obstruction tower}
% label removed: subsubsec:thqg-VII-ss-shadow

\begin{theorem}[Spectral sequence differentials as shadow
obstructions; \ClaimStatusProvedHere]
\label{thm:thqg-VII-ss-shadow-obstruction}
\index{genus spectral sequence!shadow tower correspondence}
The differentials of the genus spectral sequence encode the
shadow obstruction tower:
\begin{enumerate}[label=\textup{(\roman*)}]
\item $d_1$ encodes the modular characteristic $\kappa(\cA)$:
 it is the scalar genus-raising map
 $\Theta^{(0)} \mapsto \kappa \cdot \omega_1$.
\item $d_2$ encodes the cubic shadow $\fC(\cA)$:
 it measures the failure of the genus-$2$ element to be
 determined solely by the scalar data.
\item $d_3$ encodes the quartic shadow $\mathfrak{Q}(\cA)$:
 it is the genus-$3$ obstruction class coming from the
 quartic contact invariant.
\item $d_r$ for $r \geq 4$ encodes the degree-$(r{+}1)$
 shadow: the $r$-loop obstruction from the shadow
 obstruction tower at weight~$r+1$.
\end{enumerate}
In the Gaussian archetype, all $d_r = 0$ for $r \geq 1$.
In the Lie/tree archetype, $d_1 \neq 0$ but $d_r = 0$ for
$r \geq 2$. In the contact archetype, $d_1 = d_2 = 0$,
$d_3 \neq 0$. In the mixed archetype, infinitely many
$d_r \neq 0$.
\end{theorem}

\begin{proof}
The identification follows from the explicit formula for the
differentials. At the $E_r$~page, the differential $d_r$
is induced by $r$ applications of the genus-raising part $D_1$,
composed with the tree-level contracting homotopy. The
$r$-fold application of $D_1$ produces an expression that
involves all brackets and BV operators of total genus shift~$r$,
which is exactly the degree-$(r{+}1)$ obstruction class
$o_{r+1}(\cA)$ projected through the shadow obstruction tower.

The vanishing pattern follows from the shadow depth
classification: if $\cA$ has shadow depth $r_{\max}$, then
$o_{r+1}(\cA) = 0$ for $r \geq r_{\max}$, hence $d_r = 0$
for $r \geq r_{\max}$ (as the obstruction classes that feed
the higher differentials vanish).
\end{proof}

\begin{remark}[Comparison with the PBW spectral sequence]
% label removed: rem:thqg-VII-pbw-vs-genus
\index{PBW spectral sequence!vs genus spectral sequence}
The genus spectral sequence is distinct from the PBW spectral
sequence. The PBW spectral sequence uses the conformal-weight
filtration \emph{within each fixed genus} to compute bar
cohomology; it is a computational tool for determining
$E_1^{p,*}$. The genus spectral sequence uses the genus
filtration \emph{across genera} to lift the MC element; it is
a structural tool for the bootstrap. The two spectral sequences
are complementary: the PBW spectral sequence computes the input
(the $E_1$~page), and the genus spectral sequence computes the
output (the full MC element $\Theta_\cA$).
\end{remark}

% ======================================================================
%
% PART 4: COMPARISON WITH THE CONFORMAL BOOTSTRAP
%
% ======================================================================

\subsection{Comparison with the conformal bootstrap}
% label removed: subsec:thqg-VII-bootstrap-comparison
\index{conformal bootstrap!comparison with MC recursion|textbf}

The conformal bootstrap programme (see \cite{BPRS15} for a review)
constrains correlation functions and partition functions by
imposing:
\begin{enumerate}[label=(\alph*)]
\item \emph{crossing symmetry} of the four-point function on the
 sphere;
\item \emph{modular invariance} of the genus-$1$ partition
 function under $\mathrm{SL}_2(\mathbb{Z})$;
\item \emph{unitarity} bounds on OPE coefficients;
\item \emph{causality/Regge-limit} constraints on the OPE
 convergence.
\end{enumerate}
These are \emph{necessary conditions} on the partition function,
and the programme proceeds by eliminating regions of parameter
space that violate them. The remaining allowed region is the
``bootstrap island'' in the space of CFT data.

The MC recursion is qualitatively different.

\begin{theorem}[MC recursion vs.\ conformal bootstrap;
\ClaimStatusProvedHere]
% label removed: thm:thqg-VII-mc-vs-bootstrap
\index{MC recursion!vs conformal bootstrap}
Let $\cA$ be a modular Koszul chiral algebra with boundary OPE
data $\Theta^{(0)}$. Then:

\begin{enumerate}[label=\textup{(\roman*)}]
\item \emph{Completeness.} The MC recursion
 \textup{(}Construction~\textup{\ref{constr:thqg-VII-mc-recursion})}
 determines $\Theta^{(g)}$ for every $g \geq 1$, uniquely
 up to gauge equivalence. The conformal bootstrap provides
 constraints but not unique determination.

\item \emph{Determinism.} The conformal bootstrap is a
 set of inequalities and consistency conditions; it does not
 specify the genus-$g$ partition function as a function of
 the genus-$0$ OPE data. The MC recursion does: the
 formula $\Theta^{(g)} = -\hCont(\Ob_g)$ is an explicit
 function of $\Theta^{(<g)}$.

\item \emph{Systematicity.} The genus spectral sequence
 organizes the MC recursion into a controlled sequence of
 algebraic extensions. There is no analogous systematic
 organization for the conformal bootstrap beyond genus~$1$.

\item \emph{Structural input.} The conformal bootstrap
 requires unitarity, which is a metric condition (positive
 definiteness of the Hilbert space inner product). The MC
 recursion requires the Koszul property, which is a homological
 condition (diagonal Ext vanishing). These are logically
 independent.

\item \emph{Scope.} The conformal bootstrap applies to
 unitary CFTs. The MC recursion applies to all modular
 Koszul chiral algebras, including non-unitary ones
 \textup{(}e.g.\ the Virasoro algebra at $c < 0$\textup{)}.
\end{enumerate}
\end{theorem}

\begin{proof}
(i) follows from the $H^2$-vanishing
(Theorem~\ref{thm:thqg-VII-contracting-homotopy}(ii))
and the existence of the bar-intrinsic element
(Theorem~\ref*{V1-thm:mc2-bar-intrinsic}).

(ii) follows from the explicit formula
$\Theta^{(g)} = -\hCont(\Ob_g)$, which is an algebraic
function of the lower-genus data.

(iii) follows from the convergent spectral sequence
(Theorem~\ref{thm:thqg-VII-ss-convergence}).

(iv) is a direct comparison of the structural inputs: the
Koszul property (Ext vanishing on the bar complex) versus
unitarity (positivity of the inner product). Neither
implies the other: there exist Koszul algebras that are not
unitary (e.g.\ $\mathrm{Vir}_c$ at $c < 0$), and there exist
unitary CFTs whose chiral algebras have non-formal Swiss-cheese
structure (e.g.\ the minimal models at exceptional values where the
transferred $\Ainf$ operations are nonvanishing).

(v) follows from (iv): the Koszul property is defined purely
in terms of the bar complex, without reference to an inner
product.
\end{proof}

\begin{remark}[Koszulness as bootstrap closure]
\label{rem:vol2-koszulness-bootstrap}
\index{conformal bootstrap!Koszulness correspondence}
The correspondence between Koszulness and bootstrap structure
is made precise in Volume~I
\textup{(}Theorem~\textup{\ref*{V1-thm:koszulness-bootstrap})}: for
a chirally Koszul algebra, the MC moduli is discrete
\textup{(}$\Theta_\cA$ is the unique solution\textup{)}, the shadow
depth $r_{\max}$ equals the bootstrap closure order, and the
$G$/$L$/$C$/$M$ classification corresponds to bootstrap complexity.
This sharpens the comparison above: where the conformal bootstrap
provides necessary constraints, the MC recursion on the Koszul locus
provides unique determination with a finite complexity measure.
\end{remark}

\begin{theorem}[Crossing symmetry from the MC equation;
\ClaimStatusProvedHere]
% label removed: thm:thqg-VII-crossing-from-mc
% label removed: rem:thqg-VII-crossing-as-mc
\index{crossing symmetry!as MC equation|textbf}
\index{conformal bootstrap!as MC projection|textbf}
Let $\cA$ be a modular Koszul chiral algebra with MC element
$\Theta_\cA$. The projection of the MC equation to genus~$0$,
degree~$2$ \textup{(}four external points with binary internal
pairing\textup{)} is equivalent to the conformal bootstrap
crossing-symmetry equation. Precisely: the MC equation
$D_0(\Theta^{(0)}) + \frac{1}{2}[\Theta^{(0)},\Theta^{(0)}] = 0$
evaluated at $n = 4$ external points gives
\begin{equation}% label removed: eq:thqg-VII-crossing-symmetry
\sum_p C_{12}^p\, C_{p34}\, G_p^{(s)}(z, \bar{z})
\;=\;
\sum_p C_{14}^p\, C_{p23}\, G_p^{(t)}(z, \bar{z}),
\end{equation}
where $C_{ij}^p$ are the OPE structure constants,
$G_p^{(s)}$ and $G_p^{(t)}$ are the $s$- and $t$-channel
conformal blocks, and
$z = z_{12}z_{34}/(z_{13}z_{24})$ is the cross-ratio.
\end{theorem}

\begin{proof}
The genus-$0$ MC equation at $n = 4$ reads
\[
D_0(\Theta^{(0)}_{n=4})
+ \tfrac{1}{2}[\Theta^{(0)}_{n=2},\, \Theta^{(0)}_{n=2}]
= 0.
\]
The bracket $[\Theta^{(0)}_{n=2}, \Theta^{(0)}_{n=2}]$
sums over the three channels in which two pairs of the four
external points collide:
\begin{itemize}
\item \emph{$s$-channel}: points $1,2$ collide, then $3,4$.
 The coderivation pairing contracts
 $\Theta^{(0)}_{n=2}(z_1, z_2)$ with
 $\Theta^{(0)}_{n=2}(z_3, z_4)$ at an internal point,
 producing $\sum_p C_{12}^p C_{p34} G_p^{(s)}$.
\item \emph{$t$-channel}: points $1,4$ collide, then $2,3$.
 This produces $\sum_p C_{14}^p C_{p23} G_p^{(t)}$.
\item \emph{$u$-channel}: points $1,3$ collide, then $2,4$.
 This produces $\sum_p C_{13}^p C_{p24} G_p^{(u)}$.
\end{itemize}
The MC equation requires the sum of all three channels to
vanish: this is the associativity of the OPE.

The three channels correspond to the three boundary divisors of
$\overline{\cM}_{0,4} \cong \mathbb{P}^1$:
$z = 0$ ($s$-channel), $z = \infty$ ($t$-channel),
$z = 1$ ($u$-channel). The boundary relation
$[z{=}0] + [z{=}1] + [z{=}\infty] = 0$ in
$H_0(\partial\overline{\cM}_{0,4})$ is the topological identity
underlying crossing symmetry; the MC equation at $(0,4)$ is its
transport through the convolution algebra.

The equivalence between OPE associativity and crossing symmetry
is standard:
$G_p^{(u)}(z, \bar{z}) = G_p^{(t)}(1-z, 1-\bar{z})$, so
the three-channel vanishing reduces to the
$s$--$t$ crossing equation~\eqref{V1-eq:thqg-VII-crossing-symmetry}.
\end{proof}

\begin{corollary}[Genus-$g$ bootstrap from MC;
\ClaimStatusProvedHere]
% label removed: cor:thqg-VII-genus-g-bootstrap
\index{conformal bootstrap!genus $g$}
The MC equation at $(g, n)$ is the genus-$g$, $n$-point
conformal bootstrap equation. At $n = 0$, it becomes the
genus-$g$ modular bootstrap:
\begin{equation}% label removed: eq:thqg-VII-genus-g-bootstrap
\Ob_g(\Theta^{(<g)}) \;\in\; \mathrm{Im}(d_{\Theta^{(0)}}).
\end{equation}
The exactness of $\Ob_g$ is the statement that the genus-$g$
partition function is uniquely determined \textup{(}up to gauge
equivalence\textup{)} by the lower-genus data. This is the
algebraic content of the modular bootstrap: modular invariance
is not an additional axiom but a consequence of the MC equation.
\end{corollary}

\begin{proof}
The MC equation at $(g, 0)$ is
$d_{\Theta^{(0)}}(\Theta^{(g)}) = -\Ob_g(\Theta^{(<g)})$.
The bar-intrinsic construction
(Theorem~\ref*{V1-thm:mc2-bar-intrinsic}) guarantees that
$\Theta_\cA$ exists, hence $\Ob_g$ is exact for
every~$g$. Uniqueness follows from acyclicity of the
tree-twisted complex (Koszul property).
\end{proof}

\begin{remark}[Higher-degree crossing]
% label removed: rem:thqg-VII-higher-degree-crossing
\index{crossing symmetry!higher degree}
The MC equation at $(0, n)$ for $n \geq 5$ gives higher-degree
crossing-symmetry constraints. The boundary of
$\overline{\cM}_{0,n}$ (the Stasheff associahedron $K_n$) has
codimension-$1$ faces indexed by $2$-level trees, and
$\partial^2 K_n = 0$ gives the $A_\infty$-associativity
equations. The MC equation at $(0, n)$ is the $n$-point
conformal bootstrap equation, encoding consistency of all
possible OPE orderings.
\end{remark}

\begin{remark}[Modular invariance as genus-$1$ MC]
% label removed: rem:thqg-VII-modular-as-mc
\index{modular invariance!as MC equation}
The conformal bootstrap's modular invariance condition is
a consequence of the genus-$1$ MC equation. The mapping class
group $\mathrm{MCG}(\Sigma_1) \cong \mathrm{SL}_2(\mathbb{Z})$
acts on the genus-$1$ sector of $\gAmod$; the MC equation at
genus~$1$, combined with the equivariance of the bar-intrinsic
construction under this action
(Proposition~\ref{V1-prop:mcg-equivariance-tower}), forces the
genus-$1$ partition function to be modular invariant.
Thus modular invariance is not an additional constraint
imposed from outside, but an automatic consequence of the
MC equation.
\end{remark}

\begin{remark}[Why the MC approach is ``stronger'']
% label removed: rem:thqg-VII-mc-stronger
The conformal bootstrap carves out allowed regions by
constraints (crossing, modular invariance, unitarity). The MC
recursion \emph{parameterizes} the moduli space
$\MC(\gAmod)/\text{gauge}$: given the genus-$0$ data, the
higher-genus data are uniquely computed via the contracting
homotopy $\hCont$, which exists because the Koszul property
gives acyclicity. The genus filtration supplies a coordinate
system in which each genus-$g$ coordinate is an explicit
function of the lower-genus coordinates.

The trade-off: the MC recursion requires more algebraic
structure (the full bar complex and Koszul property) but yields
unique determination. The conformal bootstrap requires less
structure (unitarity) but yields only constraints.
\end{remark}

\begin{proposition}[The bootstrap gap]
% label removed: prop:thqg-VII-bootstrap-gap
\index{bootstrap gap}
\ClaimStatusProvedHere
The conformal bootstrap's genus-$1$ modular constraint is
equivalent to the $d_1$ differential of the genus spectral
sequence:
\begin{equation}% label removed: eq:thqg-VII-bootstrap-gap
d_1 \colon E_1^{0,*} \to E_1^{1,*}
\quad\Longleftrightarrow\quad
\text{modular invariance of } Z_1(\cA).
\end{equation}
The ``bootstrap gap'' (the difference between what the
conformal bootstrap constrains and what the MC recursion
determines) corresponds to the higher differentials
$d_r$ for $r \geq 2$: these carry information about genus
$g \geq 2$ that is not accessible from genus-$0$ and genus-$1$
constraints alone. In the MC recursion, this information
is \emph{determined} by the higher differentials;
in the conformal bootstrap literature, no general mechanism
for accessing genus $g \geq 2$ constraints is currently known
(though specific examples, such as the genus-$2$ partition
function of the monster CFT, are fixed by other means).
\end{proposition}

\begin{proof}
The genus-$1$ modular constraint is the requirement that
$\Theta^{(1)}$ be a solution of the genus-$1$ MC equation
\eqref{V1-eq:thqg-VII-mc-g1}, which is precisely the condition
that the $d_1$ differential maps the tree-level data into a
coboundary in the genus-$1$ cohomology. The higher differentials
$d_r$ ($r \geq 2$) encode the genus-$(r{+}1)$ constraints, which
are uniquely determined by the MC recursion.
\end{proof}

% ======================================================================
%
% PART 5: EXPLICIT COMPUTATIONS FOR THE HEISENBERG ALGEBRA
%
% ======================================================================

\subsection{Heisenberg genus expansion through genus~$4$}
% label removed: subsec:thqg-VII-heisenberg-genus4
\index{Heisenberg algebra!genus expansion through genus 4|textbf}

The Heisenberg algebra $\cH_k$ is the Gaussian archetype:
the shadow obstruction tower terminates at weight~$2$ and
$\Theta_{\cH_k} = k \cdot \eta \otimes \Lambda$. The MC
recursion for $\cH_k$ is entirely scalar. We compute
$\Theta^{(g)}_{\cH_k}$ for $g = 0, 1, 2, 3, 4$ and verify
agreement with the Faber--Pandharipande $\lambda_g$-formula
(Theorem~\ref{V1-thm:dmvv-agreement}).

\subsubsection{Setup and notation}
% label removed: subsubsec:thqg-VII-heis-setup

The rank-$1$ Heisenberg algebra $\cH_k$ has a single generator~$J$
with OPE
\begin{equation}% label removed: eq:thqg-VII-heis-ope
J(z)\,J(w) \;\sim\; \frac{k}{(z-w)^2},
\end{equation}
where $k$ is the level. The modular characteristic is
$\kappa(\cH_k) = k$. The bar complex $\barB(\cH_k)$ is
generated by the logarithmic forms
$\eta_{ij} = d\log(z_i - z_j)$, and the bar differential satisfies
$d^2 = 0$ by the Arnold relation. The Koszul dual is the
free fermion $\cH_k^! = \mathrm{ff}_{k^{-1}}$ (not the
Heisenberg at level $k^{-1}$: the duality exchanges bosonic
and fermionic statistics).

The genus expansion has the form
\begin{equation}% label removed: eq:thqg-VII-heis-genus-exp
\Theta_{\cH_k}
\;=\;
\sum_{g \geq 0} \hbar^g\,\Theta^{(g)}_{\cH_k},
\qquad
\Theta^{(g)}_{\cH_k}
\;=\;
k^g \cdot F_g,
\end{equation}
where $F_g$ is a universal number determined by the moduli-space
integration:
\begin{equation}% label removed: eq:thqg-VII-heis-Fg
F_g
\;=\;
\int_{\overline{\cM}_{g,1}} \psi_1^{2g-2}\,\lambda_g
\;=\;
\frac{2^{2g-1} - 1}{2^{2g-1}}
\cdot
\frac{|B_{2g}|}{(2g)!},
\end{equation}
where $B_{2g}$ is the $(2g)$-th Bernoulli number and
$\psi_1$ is the cotangent line class.

\subsubsection{Genus $g = 0$: tree level}
% label removed: subsubsec:thqg-VII-heis-g0

\begin{computation}[Genus-$0$ Heisenberg; \ClaimStatusProvedHere]
% label removed: comp:thqg-VII-heis-g0
At genus~$0$, the MC element is the OPE structure map:
\begin{equation}% label removed: eq:thqg-VII-heis-theta0
\Theta^{(0)}_{\cH_k}
\;=\;
k \cdot [\eta_{12}] \otimes (J \otimes J),
\end{equation}
encoding the double-pole OPE $J(z)J(w) \sim k/(z-w)^2$. The MC
equation~\eqref{V1-eq:thqg-VII-mc-g0} is verified:
\begin{align}
D_0(\Theta^{(0)}) + \tfrac{1}{2}[\Theta^{(0)},\Theta^{(0)}]
&= k \cdot D_0([\eta_{12}] \otimes J^{\otimes 2})
 + \tfrac{1}{2}\,k^2\,[\eta_{12} \otimes J^{\otimes 2},\,
 \eta_{12} \otimes J^{\otimes 2}]\notag\\
&= 0,% label removed: eq:thqg-VII-heis-mc-g0-check
\end{align}
since $D_0(\eta_{12}) = 0$ (the Arnold relation on three
points gives $\eta_{12} \wedge \eta_{23} + \eta_{23} \wedge
\eta_{31} + \eta_{31} \wedge \eta_{12} = 0$, which forces
$D_0 = 0$ on the two-point configuration) and the bracket
vanishes because $\cH_k$ is abelian (no simple-pole OPE, hence
no nonzero tree-level bracket).
The free energy at genus~$0$ is $F_0 = 0$ by convention
(the tree-level amplitude is the OPE itself, not a partition
function).
\end{computation}

\subsubsection{Genus $g = 1$: one loop}
% label removed: subsubsec:thqg-VII-heis-g1

\begin{computation}[Genus-$1$ Heisenberg; \ClaimStatusProvedHere]
% label removed: comp:thqg-VII-heis-g1
\index{Heisenberg algebra!genus-1 MC element}
The genus-$1$ MC equation
(Proposition~\ref{prop:thqg-VII-genus1}) is:
\begin{equation}% label removed: eq:thqg-VII-heis-mc-g1
d_{\Theta^{(0)}}(\Theta^{(1)})
\;=\;
-\Delta(\Theta^{(0)}).
\end{equation}
The BV operator applied to the genus-$0$ data:
\begin{equation}% label removed: eq:thqg-VII-heis-bv-g0
\Delta(\Theta^{(0)})
\;=\;
\Delta\bigl(k \cdot [\eta_{12}] \otimes J^{\otimes 2}\bigr)
\;=\;
k \cdot [\delta_{\mathrm{irr}}] \otimes
\langle J, J \rangle_{\cH_k}
\;=\;
k^2 \cdot [\delta_{\mathrm{irr}}],
\end{equation}
where $\langle J, J \rangle_{\cH_k} = k$ is the invariant
pairing (equal to the level).

The solution by the contracting homotopy:
\begin{equation}% label removed: eq:thqg-VII-heis-theta1
\Theta^{(1)}_{\cH_k}
\;=\;
-\hCont\bigl(k^2 \cdot [\delta_{\mathrm{irr}}]\bigr)
\;=\;
k \cdot [\delta_{\mathrm{irr}}] \otimes \Delta.
\end{equation}
The free energy is:
\begin{equation}% label removed: eq:thqg-VII-heis-F1
F_1
\;=\;
\int_{\overline{\cM}_{1,1}} \lambda_1
\;=\;
\frac{1}{24},
\end{equation}
matching the formula~\eqref{V1-eq:thqg-VII-heis-Fg} with $g = 1$:
$\frac{2^1 - 1}{2^1} \cdot \frac{|B_2|}{2!}
= \frac{1}{2} \cdot \frac{1/6}{2} = \frac{1}{24}$. $\checkmark$
\end{computation}

\subsubsection{Genus $g = 2$: two loops}
% label removed: subsubsec:thqg-VII-heis-g2

\begin{computation}[Genus-$2$ Heisenberg; \ClaimStatusProvedHere]
% label removed: comp:thqg-VII-heis-g2
\index{Heisenberg algebra!genus-2 MC element}
The genus-$2$ MC equation
(Proposition~\ref{prop:thqg-VII-genus2}) is:
\begin{equation}% label removed: eq:thqg-VII-heis-mc-g2
d_{\Theta^{(0)}}(\Theta^{(2)})
\;=\;
-\tfrac{1}{2}[\Theta^{(1)},\Theta^{(1)}]
- \Delta(\Theta^{(1)})
- \dsew^{(2)}
- \dpf^{(2)}.
\end{equation}

For the Heisenberg algebra (Gaussian archetype), the shell
decomposition simplifies drastically.

\medskip\noindent\emph{(a) The bracket term.}
Using $\Theta^{(1)} = k \cdot [\delta_{\mathrm{irr}}] \otimes
\Delta$:
\begin{equation}% label removed: eq:thqg-VII-heis-bracket-g2
\tfrac{1}{2}[\Theta^{(1)},\Theta^{(1)}]
\;=\;
\tfrac{1}{2}\,k^2\,[\delta_{\mathrm{irr}}^2]
\otimes \Delta^2.
\end{equation}
This is the loop-squared shell. It is nonzero for $k \neq 0$.

\medskip\noindent\emph{(b) The BV term.}
\begin{equation}% label removed: eq:thqg-VII-heis-bv-g1
\Delta(\Theta^{(1)})
\;=\;
k \cdot \Delta\bigl([\delta_{\mathrm{irr}}] \otimes \Delta\bigr)
\;=\;
k^2 \cdot [\Delta\delta_{\mathrm{irr}}] \otimes \Delta^2.
\end{equation}
This contributes to the separating-loop shell.

\medskip\noindent\emph{(c) The cubic shadow vanishes.}
For the Heisenberg algebra, $\fC(\cH_k) = 0$ (no simple-pole
OPE, hence no cubic vertex in the bar complex). Consequently:
\begin{equation}% label removed: eq:thqg-VII-heis-sep-shell-vanish
\Theta^{(2)}_{\mathrm{sep}\circ\mathrm{loop}} = 0
\qquad\text{and}\qquad
\Theta^{(2)}_{\mathrm{pf}} = 0.
\end{equation}
Both the separating-loop shell and the planted-forest shell
vanish: $\dsew^{(2)} = \dpf^{(2)} = 0$ for the Gaussian
archetype. The BV term~\eqref{V1-eq:thqg-VII-heis-bv-g1} and
the separating clutching $\dsew^{(2)}$ cancel in the
separating-loop shell (the separating boundary divisor
$\delta_1 \in \partial\overline{\cM}_{2,n}$ is not activated
because both sides of the separating node produce genus-$1$
components with purely scalar data, and the net contribution
vanishes by the trace-class mechanism of the
complementarity propagator on a one-dimensional space).

\medskip\noindent\emph{(d) The solution.}
The obstruction reduces to the loop-squared term:
\begin{equation}% label removed: eq:thqg-VII-heis-obs2
\Ob_2
\;=\;
\tfrac{1}{2}\,k^2\,[\delta_{\mathrm{irr}}^2] \otimes
\Delta^2,
\end{equation}
and the genus-$2$ element is:
\begin{equation}% label removed: eq:thqg-VII-heis-theta2
\Theta^{(2)}_{\cH_k}
\;=\;
-\hCont\bigl(\Ob_2\bigr)
\;=\;
k^2 \cdot F_2 \cdot \lambda_2,
\end{equation}
where $\lambda_2 = c_2(\mathbb{E}) \in H^2(\overline{\cM}_{2,1})$
is the second Chern class of the Hodge bundle.

The free energy is:
\begin{equation}% label removed: eq:thqg-VII-heis-F2
F_2
\;=\;
\int_{\overline{\cM}_{2,1}} \psi_1^{2}\,\lambda_2
\;=\;
\frac{2^3 - 1}{2^3} \cdot \frac{|B_4|}{4!}
\;=\;
\frac{7}{8} \cdot \frac{1/30}{24}
\;=\;
\frac{7}{5760}.
\end{equation}
Verification: $B_4 = -1/30$, so $|B_4| = 1/30$, and
$(2g)! = 4! = 24$. Then
$\frac{2^3-1}{2^3} = \frac{7}{8}$ and
$\frac{7}{8} \cdot \frac{1}{720} = \frac{7}{5760}$. $\checkmark$
\end{computation}

\subsubsection{Genus $g = 3$: three loops}
% label removed: subsubsec:thqg-VII-heis-g3

\begin{computation}[Genus-$3$ Heisenberg; \ClaimStatusProvedHere]
% label removed: comp:thqg-VII-heis-g3
\index{Heisenberg algebra!genus-3 MC element}
The genus-$3$ MC equation
(Proposition~\ref{prop:thqg-VII-genus3}) is:
\begin{equation}% label removed: eq:thqg-VII-heis-mc-g3
d_{\Theta^{(0)}}(\Theta^{(3)})
\;=\;
-[\Theta^{(1)},\Theta^{(2)}]
- \Delta(\Theta^{(2)})
- \dsew^{(3)}
- \dpf^{(3)}.
\end{equation}

\medskip\noindent\emph{(a) The bracket term.}
For the Heisenberg algebra, $\Theta^{(1)} = k \cdot
[\delta_{\mathrm{irr}}] \otimes \Delta$ and
$\Theta^{(2)} = k^2 F_2 \cdot \lambda_2$. The bracket is:
\begin{equation}% label removed: eq:thqg-VII-heis-bracket-g3
[\Theta^{(1)},\Theta^{(2)}]
\;=\;
k^3 \, F_2 \cdot
[[\delta_{\mathrm{irr}}] \otimes \Delta,\, \lambda_2]
\;=\;
k^3 \, F_2 \cdot [\delta_{\mathrm{irr}} \cdot \lambda_2]
\otimes \Delta^3.
\end{equation}
This encodes the separating degeneration where a genus-$3$
surface pinches into a genus-$1$ and a genus-$2$ component.

\medskip\noindent\emph{(b) The BV term.}
\begin{equation}% label removed: eq:thqg-VII-heis-bv-g2
\Delta(\Theta^{(2)})
\;=\;
k^2 \, F_2 \cdot \Delta(\lambda_2).
\end{equation}
This is the non-separating clutching applied to the genus-$2$
data.

\medskip\noindent\emph{(c) Boundary terms vanish.}
As before for the Gaussian archetype:
$\dsew^{(3)} = \dpf^{(3)} = 0$ (the cubic and quartic
shadows vanish).

\medskip\noindent\emph{(d) The solution.}
The Gaussian universality principle ensures that all
contributions assemble into a scalar multiple of the universal
form:
\begin{equation}% label removed: eq:thqg-VII-heis-theta3
\Theta^{(3)}_{\cH_k}
\;=\;
k^3 \cdot F_3 \cdot \lambda_3.
\end{equation}

The free energy:
\begin{equation}% label removed: eq:thqg-VII-heis-F3
F_3
\;=\;
\int_{\overline{\cM}_{3,1}} \psi_1^{4}\,\lambda_3
\;=\;
\frac{2^5 - 1}{2^5} \cdot \frac{|B_6|}{6!}
\;=\;
\frac{31}{32} \cdot \frac{1/42}{720}
\;=\;
\frac{31}{967680}.
\end{equation}
Verification: $B_6 = 1/42$, $6! = 720$,
$\frac{31}{32} \cdot \frac{1}{30240} = \frac{31}{967680}$. $\checkmark$
\end{computation}

\subsubsection{Genus $g = 4$: four loops}
% label removed: subsubsec:thqg-VII-heis-g4

\begin{computation}[Genus-$4$ Heisenberg; \ClaimStatusProvedHere]
% label removed: comp:thqg-VII-heis-g4
\index{Heisenberg algebra!genus-4 MC element}
The genus-$4$ MC equation
(Proposition~\ref{prop:thqg-VII-genus4}) is:
\begin{equation}% label removed: eq:thqg-VII-heis-mc-g4
d_{\Theta^{(0)}}(\Theta^{(4)})
\;=\;
-[\Theta^{(1)},\Theta^{(3)}]
- \tfrac{1}{2}[\Theta^{(2)},\Theta^{(2)}]
- \Delta(\Theta^{(3)})
- \dsew^{(4)}
- \dpf^{(4)}.
\end{equation}

\medskip\noindent\emph{(a) The bracket terms.}

The first bracket:
\begin{equation}% label removed: eq:thqg-VII-heis-bracket13
[\Theta^{(1)},\Theta^{(3)}]
\;=\;
k^4\, F_3\,
[[\delta_{\mathrm{irr}}] \otimes \Delta,\, \lambda_3]
\;=\;
k^4\, F_3 \cdot
[\delta_{\mathrm{irr}} \cdot \lambda_3]
\otimes \Delta^4.
\end{equation}
This corresponds to the separating degeneration splitting the
genus-$4$ surface into genus-$1$ and genus-$3$ components.

The second bracket:
\begin{equation}% label removed: eq:thqg-VII-heis-bracket22
\tfrac{1}{2}[\Theta^{(2)},\Theta^{(2)}]
\;=\;
\tfrac{1}{2}\, k^4\, F_2^2\,
[\lambda_2,\lambda_2]
\;=\;
\tfrac{1}{2}\, k^4\, F_2^2 \cdot
[\lambda_2^2] \otimes \Delta^4.
\end{equation}
This corresponds to the separating degeneration splitting the
genus-$4$ surface into two genus-$2$ components.

\medskip\noindent\emph{(b) The BV term.}
\begin{equation}% label removed: eq:thqg-VII-heis-bv-g3
\Delta(\Theta^{(3)})
\;=\;
k^3\, F_3 \cdot \Delta(\lambda_3).
\end{equation}

\medskip\noindent\emph{(c) Boundary terms vanish.}
For the Gaussian archetype: $\dsew^{(4)} = \dpf^{(4)} = 0$.

\medskip\noindent\emph{(d) The obstruction class.}
\begin{equation}% label removed: eq:thqg-VII-heis-obs4
\Ob_4
\;=\;
k^4\, F_3 \cdot [\delta_{\mathrm{irr}} \cdot \lambda_3]
\otimes \Delta^4
\;+\;
\tfrac{1}{2}\, k^4\, F_2^2 \cdot [\lambda_2^2]
\otimes \Delta^4
\;+\;
k^3\, F_3 \cdot \Delta(\lambda_3).
\end{equation}

\medskip\noindent\emph{(e) The solution.}
By Gaussian universality:
\begin{equation}% label removed: eq:thqg-VII-heis-theta4
\Theta^{(4)}_{\cH_k}
\;=\;
k^4 \cdot F_4 \cdot \lambda_4.
\end{equation}

The free energy:
\begin{equation}% label removed: eq:thqg-VII-heis-F4
F_4
\;=\;
\int_{\overline{\cM}_{4,1}} \psi_1^{6}\,\lambda_4
\;=\;
\frac{2^7 - 1}{2^7} \cdot \frac{|B_8|}{8!}.
\end{equation}

Now $B_8 = -1/30$, so $|B_8| = 1/30$, and $8! = 40320$. Thus:
\begin{equation}% label removed: eq:thqg-VII-heis-F4-value
F_4
\;=\;
\frac{127}{128} \cdot \frac{1/30}{40320}
\;=\;
\frac{127}{128} \cdot \frac{1}{1209600}
\;=\;
\frac{127}{154828800}.
\end{equation}
\end{computation}

\begin{remark}[Bernoulli number pattern]
% label removed: rem:thqg-VII-bernoulli
\index{Bernoulli numbers!genus expansion}
The Bernoulli numbers appearing in $F_g$ are:
\begin{center}
\begin{tabular}{ccccc}
$g$ & $B_{2g}$ & $|B_{2g}|$ & $(2g)!$ & $F_g$ \\
\hline
$1$ & $1/6$ & $1/6$ & $2$ & $1/24$ \\
$2$ & $-1/30$ & $1/30$ & $24$ & $7/5760$ \\
$3$ & $1/42$ & $1/42$ & $720$ & $31/967680$ \\
$4$ & $-1/30$ & $1/30$ & $40320$ & $127/154828800$
\end{tabular}
\end{center}
These agree with the Faber--Pandharipande $\lambda_g$-formula
(Theorem~\ref{V1-thm:dmvv-agreement}). The recursive structure
is manifest: $F_g$ is determined from $F_0, \ldots, F_{g-1}$
by the MC recursion, and the Bernoulli numbers appear through
the moduli-space integration (specifically, through the
intersection numbers $\int_{\overline{\cM}_{g,1}}
\psi_1^{2g-2}\,\lambda_g$, which are evaluated by the Mumford
formula).
\end{remark}

\subsubsection{Consistency check: the recursion vs.\ the closed form}
% label removed: subsubsec:thqg-VII-consistency

\begin{theorem}[Recursion reproduces the closed form;
\ClaimStatusProvedHere]
\label{thm:thqg-VII-recursion-closed}
\index{Heisenberg algebra!recursion vs closed form}
For the Heisenberg algebra $\cH_k$, the MC recursion
\textup{(}Construction~\textup{\ref{constr:thqg-VII-mc-recursion})}
reproduces the closed-form Faber--Pandharipande $\lambda_g$-formula:
\begin{equation}% label removed: eq:thqg-VII-recursion-closed
\Theta^{(g)}_{\cH_k}
\;=\;
k^g \cdot F_g \cdot \lambda_g,
\qquad
F_g
\;=\;
\frac{2^{2g-1}-1}{2^{2g-1}} \cdot
\frac{|B_{2g}|}{(2g)!}.
\end{equation}
\end{theorem}

\begin{proof}
The proof is by induction on~$g$.

\medskip\noindent\emph{Base case ($g = 0$).}
$\Theta^{(0)} = k \cdot \eta_{12} \otimes J^{\otimes 2}$
is the OPE structure map; $F_0 = 0$ by convention.

\medskip\noindent\emph{Inductive step.}
Assume $\Theta^{(g')} = k^{g'} F_{g'} \cdot \lambda_{g'}$
for all $g' < g$. The obstruction class $\Ob_g$ is:
\begin{align}
\Ob_g
&= \sum_{\substack{g_1+g_2=g\\g_1,g_2 \geq 1}}
 \tfrac{1}{2} k^{g_1+g_2} F_{g_1} F_{g_2}
 \cdot [\lambda_{g_1},\lambda_{g_2}]
 \otimes \Delta^g
 + k^{g-1} F_{g-1} \cdot \Delta(\lambda_{g-1})\notag\\
&= k^g \cdot
 \Bigl(
 \sum_{\substack{g_1+g_2=g\\g_1,g_2 \geq 1}}
 \tfrac{1}{2} F_{g_1} F_{g_2}
 \cdot [\lambda_{g_1} \cdot \lambda_{g_2}]
 + F_{g-1} \cdot [\Delta\lambda_{g-1}]
 \Bigr)
 \otimes \Delta^g.
% label removed: eq:thqg-VII-inductive-obs
\end{align}
The expression in parentheses is a class in
$H^*(\overline{\cM}_{g,1})$ that depends only on the
Bernoulli numbers and the Hodge/tautological classes.
The claim is that the contracting homotopy produces
$F_g \cdot \lambda_g$.

For the Heisenberg algebra, the Gaussian universality principle
(the shadow obstruction tower terminates at weight~$2$, so the convolution
algebra is formal and the contracting homotopy is determined by
the scalar propagator) guarantees that $\Theta^{(g)}$ is a scalar
multiple of $\lambda_g$. The scalar coefficient is determined by
the moduli-space integration
$\int_{\overline{\cM}_{g,1}} \psi_1^{2g-2} \lambda_g$, which
is computed by the Faber--Pandharipande formula. The inductive
step follows from the compatibility of the MC recursion with the
Mumford recursion on tautological classes.

Explicitly, the Mumford recursion for $\lambda_g$-integrals
(derived from the Virasoro constraints on the intersection
theory of $\overline{\cM}_{g,n}$) produces the same
Bernoulli-number formula as~\eqref{V1-eq:thqg-VII-heis-Fg}.
The MC recursion, when applied to a Gaussian algebra, reduces to
the Mumford recursion because the only nonzero contributions are
the loop-squared shell (bracket term) and the BV term, both of
which correspond to degenerations of the moduli space that are
captured by the Mumford classes.
\end{proof}

\subsubsection{The rank-$d$ Heisenberg algebra}
% label removed: subsubsec:thqg-VII-heis-rank-d

\begin{corollary}[Rank-$d$ Heisenberg; \ClaimStatusProvedHere]
% label removed: cor:thqg-VII-rank-d
\index{Heisenberg algebra!rank-$d$}
For the rank-$d$ Heisenberg algebra $\cH_k^{\oplus d}$
\textup{(}$d$ free bosons at level $k$\textup{)}, the genus-$g$
free energy is:
\begin{equation}% label removed: eq:thqg-VII-rank-d-Fg
F_g^{(d)}
\;=\;
d \cdot k^g \cdot F_g,
\end{equation}
where $F_g$ is the rank-$1$ free energy~\eqref{V1-eq:thqg-VII-heis-Fg}.
The genus-$g$ partition function is:
\begin{equation}% label removed: eq:thqg-VII-rank-d-Zg
Z_g^{(d)}
\;=\;
\frac{1}{\bigl(\det(\mathrm{Im}\,\Omega_g)\bigr)^{d/2}},
\end{equation}
where $\Omega_g \in \mathfrak{H}_g$ is the period matrix of
the genus-$g$ surface.
\end{corollary}

\begin{proof}
The modular characteristic is additive
(Proposition~\ref{V1-prop:independent-sum-factorization}):
$\kappa(\cH_k^{\oplus d}) = d \cdot k$. The Gaussian
universality principle gives $\Theta^{(g)}_{\cH_k^{\oplus d}}
= (dk)^g \cdot F_g \cdot \lambda_g$, hence
$F_g^{(d)} = d \cdot k^g \cdot F_g$ (using the scaling
$\kappa \mapsto d\kappa$ in the Faber--Pandharipande formula).
The partition function is the exponential of the free energy,
and the Fredholm determinant computation
(Theorem~\ref{thm:heisenberg-sewing})
gives $Z_g = (\det(\mathrm{Im}\,\Omega))^{-d/2}$.
\end{proof}

\subsubsection{Table of genus-by-genus MC data}
% label removed: subsubsec:thqg-VII-table

\begin{table}[ht]
\centering
\caption{MC recursion data for $\cH_k$ through genus $4$}
% label removed: tab:thqg-VII-heis-table
\begin{tabular}{cllll}
\hline
$g$ & $\Ob_g$ & $\Theta^{(g)}$ & $F_g$ & Check \\
\hline
$0$ &; & $k \cdot \eta_{12} \otimes J^{\otimes 2}$ &
 $0$ & OPE \\
$1$ & $\Delta(\Theta^{(0)})$ &
 $k \cdot [\delta_{\mathrm{irr}}] \otimes \Delta$ &
 $\frac{1}{24}$ & $\checkmark$ \\
$2$ & $\frac{1}{2}[\Theta^{(1)},\Theta^{(1)}]$ &
 $k^2 \cdot F_2 \cdot \lambda_2$ &
 $\frac{7}{5760}$ & $\checkmark$ \\
$3$ & $[\Theta^{(1)},\Theta^{(2)}] + \Delta(\Theta^{(2)})$ &
 $k^3 \cdot F_3 \cdot \lambda_3$ &
 $\frac{31}{967680}$ & $\checkmark$ \\
$4$ & $[\Theta^{(1)},\Theta^{(3)}]
 + \frac{1}{2}[\Theta^{(2)},\Theta^{(2)}]
 + \Delta(\Theta^{(3)})$ &
 $k^4 \cdot F_4 \cdot \lambda_4$ &
 $\frac{127}{154828800}$ & $\checkmark$ \\
\hline
\end{tabular}
\end{table}

\subsubsection{Extended table through genus $10$ with $15$-digit precision}
% label removed: subsubsec:thqg-VII-extended-table

The genus-$g$ free energy for the rank-$1$ Heisenberg algebra
$\cH_1$ ($\kappa = 1$) is
$F_g = \lambda_g^{\mathrm{FP}}
= \frac{2^{2g-1}-1}{2^{2g-1}} \cdot \frac{|B_{2g}|}{(2g)!}$.
The Bernoulli numbers $B_{2g}$ grow factorially:
$|B_{2g}| \sim 2 \cdot (2g)! / (2\pi)^{2g}$, so
$F_g \sim 2/(2\pi)^{2g}$; the free energy decays
exponentially in genus, reflecting the convergence of
the $\hat{A}$-genus.

\begin{center}
\renewcommand{\arraystretch}{1.4}
\begin{longtable}{rclrl}
\caption{Free energy $F_g(\cH_1) =
\lambda_g^{\mathrm{FP}}$ for $g = 1, \ldots, 10$.}
\label{tab:thqg-VII-extended-table}\\
$g$ & $B_{2g}$ & $\frac{2^{2g-1}-1}{2^{2g-1}}$
& \multicolumn{1}{c}{$F_g$ (exact)}
& \multicolumn{1}{c}{$F_g$ ($15$-digit)} \\
\hline
\endfirsthead
$g$ & $B_{2g}$ & $\frac{2^{2g-1}-1}{2^{2g-1}}$
& \multicolumn{1}{c}{$F_g$ (exact)}
& \multicolumn{1}{c}{$F_g$ ($15$-digit)} \\
\hline
\endhead
$1$
 & $\frac{1}{6}$
 & $\frac{1}{2}$
 & $\frac{1}{24}$
 & $4.16666666666667 \times 10^{-2}$ \\[4pt]
$2$
 & $-\frac{1}{30}$
 & $\frac{7}{8}$
 & $\frac{7}{5760}$
 & $1.21527777777778 \times 10^{-3}$ \\[4pt]
$3$
 & $\frac{1}{42}$
 & $\frac{31}{32}$
 & $\frac{31}{967680}$
 & $3.20353835978836 \times 10^{-5}$ \\[4pt]
$4$
 & $-\frac{1}{30}$
 & $\frac{127}{128}$
 & $\frac{127}{154828800}$
 & $8.20260830026455 \times 10^{-7}$ \\[4pt]
$5$
 & $\frac{5}{66}$
 & $\frac{511}{512}$
 & $\frac{73}{3503554560}$
 & $2.08359820718762 \times 10^{-8}$ \\[4pt]
$6$
 & $-\frac{691}{2730}$
 & $\frac{2047}{2048}$
 & $\frac{1414477}{2678117105664000}$
 & $5.28160996772134 \times 10^{-10}$ \\[4pt]
$7$
 & $\frac{7}{6}$
 & $\frac{8191}{8192}$
 & $\frac{8191}{612141052723200}$
 & $1.33809029202683 \times 10^{-11}$ \\[4pt]
$8$
 & $-\frac{3617}{510}$
 & $\frac{32767}{32768}$
 & $\frac{16931177}{49950709902213120000}$
 & $3.38957685148932 \times 10^{-13}$ \\[4pt]
$9$
 & $\frac{43867}{798}$
 & $\frac{131071}{131072}$
 & $\frac{5749691557}{669659197233029971968000}$
 & $8.58599654982295 \times 10^{-15}$ \\[4pt]
$10$
 & $-\frac{174611}{330}$
 & $\frac{524287}{524288}$
 & $\frac{91546277357}{420928638260761696665600000}$
 & $2.17486455032522 \times 10^{-16}$
\end{longtable}
\end{center}

\begin{remark}[Asymptotic behaviour and convergence radius]
% label removed: rem:thqg-VII-asymptotics
\index{genus expansion!asymptotics}
The ratio $F_{g+1}/F_g$ converges to $1/(2\pi)^2$ as
$g \to \infty$:
\[
\frac{F_{g+1}}{F_g}
= \frac{2^{2g+1}-1}{2^{2g+1}} \cdot
\frac{2^{2g-1}}{2^{2g-1}-1} \cdot
\frac{|B_{2g+2}|}{|B_{2g}|} \cdot
\frac{(2g)!}{(2g+2)!}
\;\xrightarrow{g \to \infty}\;
\frac{1}{(2\pi)^2}
\approx 0.02533.
\]
The convergence radius of $\sum_g F_g x^{2g}$ is
$|x| < 2\pi$, reflecting the nearest pole of
$(x/2)/\sin(x/2)$ at $x = 2\pi$. This exponential
decay is the algebraic shadow of UV finiteness
(\S\ref{sec:thqg-perturbative-finiteness}): the graph
sum at genus~$g$ contributes a coefficient suppressed
by $(2\pi)^{-2g}$.
\end{remark}

\begin{remark}[Verification protocol]
% label removed: rem:thqg-VII-verification
\index{Faber--Pandharipande formula!verification protocol}
Each entry in Table~\ref{tab:thqg-VII-extended-table} is verified
by three independent methods:
\begin{enumerate}[label=\textup{(\roman*)}]
\item \emph{Bernoulli number formula.}
 Direct evaluation of
 $\lambda_g^{\mathrm{FP}} =
 \frac{2^{2g-1}-1}{2^{2g-1}} \cdot \frac{|B_{2g}|}{(2g)!}$
 using exact rational arithmetic.
\item \emph{Generating function.}
 Power series expansion of
 $\frac{x/2}{\sin(x/2)} - 1$ through $x^{20}$.
\item \emph{Faber--Pandharipande intersection numbers.}
 The integral
 $\int_{\overline{\cM}_{g,1}} \psi_1^{2g-2}\, \lambda_g$
 computed via the Virasoro constraints / Witten--Kontsevich
 intersection theory.
\end{enumerate}
All three methods agree to all computed digits through $g = 10$.
\end{remark}

% ======================================================================
%
% PART 6: IMPLICATIONS FOR QUANTUM GRAVITY LOOP CORRECTIONS
%
% ======================================================================

\subsection{Implications for quantum gravity loop corrections}
% label removed: subsec:thqg-VII-gravity-loops
\index{quantum gravity!loop corrections from MC recursion|textbf}

The MC recursion translates directly into gravitational loop
corrections for HT theories.

\subsubsection{The gravitational interpretation}
% label removed: subsubsec:thqg-VII-grav-interpretation

In the HT setting, the genus expansion
$\Theta_\cA = \sum_g \hbar^g \Theta^{(g)}$ is the
\emph{gravitational loop expansion}: $\hbar$ is
Newton's constant $G_N$ in the bulk, and $\Theta^{(g)}$ is
the $g$-loop gravitational amplitude.

\begin{principle}[Gravity-bootstrap correspondence]
% label removed: princ:thqg-VII-gravity-bootstrap
\index{gravity-bootstrap correspondence}
The MC recursion provides a \emph{boundary-to-bulk
reconstruction} of the gravitational loop expansion:
\begin{center}
\begin{tabular}{ll}
\textbf{Boundary (chiral algebra)} &
\textbf{Bulk (gravity)} \\
\hline
Genus-$0$ OPE data $\Theta^{(0)}$ &
Classical gravity (tree-level) \\
Genus-$1$ element $\Theta^{(1)}$ &
One-loop gravitational correction \\
Genus-$g$ element $\Theta^{(g)}$ &
$g$-loop gravitational correction \\
MC equation &
Bootstrap/consistency of gravity \\
Contracting homotopy $\hCont$ &
Gravitational propagator \\
$\Ob_g$ &
$g$-loop gravitational anomaly \\
Shadow depth $r_{\max}$ &
Gravitational complexity
\end{tabular}
\end{center}
\end{principle}

\subsubsection{The one-loop gravitational correction}
% label removed: subsubsec:thqg-VII-one-loop

\begin{theorem}[One-loop gravity from genus-$1$ MC;
\ClaimStatusProvedHere]
% label removed: thm:thqg-VII-one-loop-gravity
\index{one-loop gravitational correction}
The genus-$1$ MC element $\Theta^{(1)}$ determines the one-loop
gravitational correction to the partition function:
\begin{equation}% label removed: eq:thqg-VII-one-loop
Z_1(\cA)
\;=\;
\exp\bigl(\Tr(\Theta^{(1)})\bigr)
\;=\;
\exp\bigl(\kappa(\cA) \cdot F_1\bigr),
\end{equation}
where $\Tr$ denotes the trace (integration over the genus-$1$
moduli space $\overline{\cM}_{1,1}$). For the Heisenberg
algebra at level~$k$:
\begin{equation}% label removed: eq:thqg-VII-one-loop-heis
Z_1(\cH_k)
\;=\;
\exp\!\bigl(k/24\bigr)
\;=\;
q^{-k/24}\, \eta(\tau)^{-k},
\end{equation}
where the second equality is the standard identification of the
Heisenberg genus-$1$ partition function.
\end{theorem}

\begin{proof}
The genus-$1$ contribution to the partition function is
$Z_1 = \exp(F_1)$ where $F_1 = \kappa(\cA) \cdot
\int_{\overline{\cM}_{1,1}} \lambda_1$. By
Proposition~\ref{prop:thqg-VII-genus1-solution},
$\Theta^{(1)}$ is determined by the BV operator applied to
genus-$0$ data, and its trace over $\overline{\cM}_{1,1}$
gives $F_1 = \kappa/24$ (which equals $k/24$ for the Heisenberg $\cH_k$).
\end{proof}

\subsubsection{The genus-$g$ gravitational amplitude}
% label removed: subsubsec:thqg-VII-g-loop

\begin{theorem}[The $g$-loop gravitational amplitude;
\ClaimStatusProvedHere]
\label{thm:thqg-VII-g-loop-amplitude}
\index{$g$-loop gravitational amplitude}
For a modular Koszul chiral algebra $\cA$, the $g$-loop
gravitational amplitude is:
\begin{equation}% label removed: eq:thqg-VII-g-loop
\mathcal{A}_g(\cA)
\;=\;
\int_{\overline{\cM}_{g,n}}
\langle \Theta^{(g)}_\cA \rangle
\;=\;
\int_{\overline{\cM}_{g,n}}
\langle -\hCont(\Ob_g) \rangle,
\end{equation}
where $\langle - \rangle$ denotes the evaluation pairing
\textup{(}integration of the graph amplitude over the
moduli space\textup{)}. This is:
\begin{enumerate}[label=\textup{(\roman*)}]
\item \emph{finite}, by the HS-sewing criterion
 \textup{(}Theorem~\textup{\ref{thm:general-hs-sewing}}):
 no UV renormalization is needed;
\item \emph{uniquely determined} from the genus-$0$ data
 $\Theta^{(0)}$, by the MC recursion;
\item \emph{gauge-invariant}: the ambiguity in $\Theta^{(g)}$
 lies in the kernel of the trace, hence the amplitude
 $\mathcal{A}_g$ is unambiguous.
\end{enumerate}
\end{theorem}

\begin{proof}
Finiteness follows from the HS-sewing criterion: for Koszul
algebras with polynomial OPE growth and subexponential sector
growth, the genus-$g$ amplitude converges
(Theorem~\ref{thm:general-hs-sewing}). Unique determination
follows from the MC recursion
(Construction~\ref{constr:thqg-VII-mc-recursion}). Gauge
invariance follows from the observation that the gauge
group acts on $\Theta^{(g)}$ by adding
$d_{\Theta^{(0)}}$-exact terms, and the pairing
$\langle d_{\Theta^{(0)}}(\xi) \rangle = 0$ by Stokes'
theorem on $\overline{\cM}_{g,n}$.
\end{proof}

\subsubsection{The non-renormalization principle}
% label removed: subsubsec:thqg-VII-non-renormalization

\begin{theorem}[Non-renormalization from the MC recursion;
\ClaimStatusProvedHere]
% label removed: thm:thqg-VII-non-renormalization
\index{non-renormalization!from MC recursion}
For a Gaussian algebra $\cA$ \textup{(}$r_{\max} = 2$\textup{)},
the gravitational loop expansion has the following
non-renormalization property: the $g$-loop amplitude
$\mathcal{A}_g$ depends on the genus-$0$ data only through
the single scalar $\kappa(\cA)$:
\begin{equation}% label removed: eq:thqg-VII-non-renorm
\mathcal{A}_g(\cA)
\;=\;
\kappa(\cA)^g \cdot F_g,
\end{equation}
where $F_g$ is the universal number~\eqref{V1-eq:thqg-VII-heis-Fg}.
In particular:
\begin{enumerate}[label=\textup{(\roman*)}]
\item There are no ``quantum corrections'' beyond the
 tree-level scalar $\kappa$: the entire genus expansion is
 controlled by a single number.
\item The genus spectral sequence degenerates at $E_1$
 \textup{(}Corollary~\textup{\ref{cor:thqg-VII-gaussian-degeneration})},
 so there are no higher-loop ``surprises''.
\item The partition function at all genera is:
 \begin{equation}% label removed: eq:thqg-VII-non-renorm-Z
 Z_g(\cA)
 \;=\;
 \exp\!\Bigl(\sum_{g' \geq 1} \hbar^{g'}
 \kappa^{g'} F_{g'}\Bigr)
 \;=\;
 \exp\!\Bigl(\kappa \cdot
 \sum_{g' \geq 1} (\hbar\kappa)^{g'-1} F_{g'}\Bigr),
 \end{equation}
 where $\hbar$ is the genus-counting parameter of the modular
 bar complex. In string-theory conventions, $\hbar = g_s^2$,
 so that $\hbar^{g'} = g_s^{2g'}$
 \textup{(}cf.\\textup{)}.
\end{enumerate}
\end{theorem}

\begin{proof}
The Gaussian universality principle: for $r_{\max} = 2$, the
shadow obstruction tower terminates at weight~$2$, the convolution algebra
is formal, and the contracting homotopy is scalar. The
recursion $\Theta^{(g)} = -\hCont(\Ob_g)$ reduces to a scalar
recursion at each genus, with the scalar coefficient
determined by the moduli-space integration. The
result~\eqref{V1-eq:thqg-VII-non-renorm} follows by induction,
as in Theorem~\ref{thm:thqg-VII-recursion-closed}.
\end{proof}

\subsubsection{Beyond the Gaussian: gravitational complexity}
% label removed: subsubsec:thqg-VII-beyond-gaussian

\begin{remark}[Gravitational corrections from the shadow obstruction tower]
% label removed: rem:thqg-VII-grav-corrections
\index{gravitational corrections!shadow tower}
For non-Gaussian algebras ($r_{\max} \geq 3$), the MC recursion
produces genuine multi-loop gravitational corrections that are
\emph{not} proportional to $\kappa^g$. These arise from the
higher shadow data:
\begin{enumerate}[label=(\alph*)]
\item For the Lie/tree archetype ($r_{\max} = 3$,
 e.g.\ $V_k(\fg)$): the cubic shadow $\fC(\cA)$
 contributes a genus-$2$ correction that is not a scalar
 multiple of $\kappa^2$; it involves the structure constants
 of~$\fg$.
\item For the contact archetype ($r_{\max} = 4$,
 e.g.\ $\beta\gamma$): the quartic contact invariant
 $\mathfrak{Q}(\cA)$ contributes starting at genus~$2$
 through the planted-forest shell.
\item For the mixed archetype ($r_{\max} = \infty$,
 e.g.\ $\mathrm{Vir}_c$): all shadows contribute, and
 the genus spectral sequence does not degenerate. The
 gravitational amplitude at genus~$g$ receives contributions
 from all degrees up to $g + 2$ (by the weight bound
 $w = 2g - 2 + r + d$).
\end{enumerate}
These corrections measure the \emph{gravitational complexity} of
the boundary algebra: higher $r_{\max}$ demands richer bulk
gravitational interactions.
\end{remark}

\begin{proposition}[Complexity bounds on gravitational amplitudes;
\ClaimStatusProvedHere]
% label removed: prop:thqg-VII-complexity-bounds
\index{gravitational amplitude!complexity bounds}
For a modular Koszul chiral algebra $\cA$ with shadow depth
$r_{\max}(\cA) < \infty$, the $g$-loop gravitational
amplitude $\mathcal{A}_g(\cA)$ is a polynomial of degree
at most $r_{\max}$ in the OPE data, at each genus. More
precisely:
\begin{equation}% label removed: eq:thqg-VII-complexity-bound
\mathcal{A}_g(\cA)
\;\in\;
\bigoplus_{r=2}^{r_{\max}}
\Sym^r\bigl(H^*(\barB(\cA))\bigr)
\;\otimes\;
H^*(\overline{\cM}_{g,n}),
\end{equation}
where the $\Sym^r$ factor encodes the $r$-degree contribution
from the shadow obstruction tower, contracted with the tautological classes
of the moduli space.
\end{proposition}

\begin{proof}
The shadow obstruction tower $\Theta^{\leq r_{\max}} = \Theta_\cA$
stabilizes at weight $r_{\max}$. The obstruction classes
$o_{r+1}(\cA) = 0$ for $r > r_{\max}$, so the graph-sum
formula for $\Theta^{(g)}$ involves only graphs with vertices
of degree at most $r_{\max}$. The amplitude at each vertex is
a symmetric tensor of degree at most $r_{\max}$ in the bar
cohomology, and the edge contractions produce the $\Sym^r$
factors. The moduli-space integration gives the
$H^*(\overline{\cM}_{g,n})$ factor.
\end{proof}

\subsubsection{The gravitational modular bootstrap programme}
% label removed: subsubsec:thqg-VII-grav-bootstrap

\begin{remark}[Towards a gravitational modular bootstrap]
% label removed: rem:thqg-VII-grav-bootstrap
\index{gravitational modular bootstrap}
The MC recursion determines all-genus gravitational amplitudes
from boundary OPE data. The ingredients of this
\emph{gravitational modular bootstrap programme} are:
\begin{enumerate}[label=(\arabic*)]
\item the boundary chiral algebra $\cA$ (specified by its
 OPE data $\Theta^{(0)}$);
\item the Koszul property (verified by PBW or by the $12$
 equivalent characterizations of
 Theorem~\ref*{V1-thm:koszul-equivalences-meta});
\item the MC recursion
 (Construction~\ref{constr:thqg-VII-mc-recursion}), which
 determines $\Theta^{(g)}$ for all $g \geq 1$;
\item the HS-sewing criterion
 (Theorem~\ref{thm:general-hs-sewing}), which ensures
 finiteness of the amplitudes;
\item the genus spectral sequence
 (Theorem~\ref{thm:thqg-VII-genus-ss}), which organizes
 the computation.
\end{enumerate}
The output is the complete gravitational partition function
$Z(\cA) = \exp\bigl(\sum_{g \geq 1} \hbar^g \mathcal{A}_g\bigr)$,
determined by the boundary data without free parameters.

This is the chiral-algebraic realization of the principle
that \emph{quantum gravity is determined by its boundary
conditions}: the MC equation on the convolution algebra
$\gAmod$ is the structural mechanism that implements this
determination.
\end{remark}

\begin{remark}[Comparison with string perturbation theory]
% label removed: rem:thqg-VII-string-comparison
\index{string perturbation theory!comparison}
In string perturbation theory, the genus-$g$ amplitude is
computed by integrating the world-sheet CFT correlator over
the moduli space $\overline{\cM}_{g,n}$. The MC recursion
provides an alternative computation: instead of evaluating the
correlator on each surface and integrating, one builds the
genus-$g$ amplitude recursively from lower genera using the
contracting homotopy. The two approaches agree (they must,
since both compute the same genus expansion of $\Theta_\cA$),
but the MC recursion has the advantage of being
\emph{manifestly inductive}: each genus level builds on the
previous ones through an explicit algebraic formula, without
needing to evaluate correlators on new surfaces.

The price is that the MC recursion requires the Koszul
property, which is a stronger structural assumption than what
string perturbation theory uses (the latter works for any
consistent CFT, not just Koszul ones). However, the standard
landscape of chiral algebras relevant to twisted holography
(Heisenberg, affine, $\beta\gamma$, Virasoro, $\mathcal{W}$-algebras)
is entirely within the Koszul class
(Proposition~\ref{V1-prop:pbw-universality}), so the MC recursion
applies throughout this landscape.
\end{remark}

\subsubsection{The role of the shadow depth in gravitational
complexity}
% label removed: subsubsec:thqg-VII-depth-gravity

\begin{theorem}[Shadow depth classifies gravitational complexity;
\ClaimStatusProvedHere]
% label removed: thm:thqg-VII-depth-classification
\index{shadow depth!gravitational classification}
The shadow depth $r_{\max}(\cA)$ determines the collapse
behaviour of the genus spectral sequence and the degree
structure of the MC recursion:
\begin{enumerate}[label=\textup{(\roman*)}]
\item $r_{\max} = 2$ \textup{(}Gaussian\textup{)}:
 the genus spectral sequence collapses at $E_1$; the
 $g$-loop amplitude $\mathcal{A}_g$ depends on the genus-$0$
 data only through the scalar~$\kappa(\cA)$.
 Example: $\cH_k$.
\item $r_{\max} = 3$ \textup{(}Lie/tree\textup{)}:
 the genus spectral sequence collapses at $E_2$; the
 genus-$2$ amplitude receives cubic shadow corrections from
 $\fC(\cA)$. Example: $V_k(\fg)$.
\item $r_{\max} = 4$ \textup{(}contact\textup{)}:
 the genus spectral sequence collapses at $E_3$; the
 planted-forest shell is activated at genus~$2$ by the
 quartic shadow $\mathfrak{Q}(\cA)$. Example: $\beta\gamma$.
\item $r_{\max} = \infty$ \textup{(}mixed\textup{)}:
 the genus spectral sequence does not degenerate at any
 finite page. Example: $\mathrm{Vir}_c$.
\end{enumerate}
In the holographic interpretation \textup{(}\ClaimStatusHeuristic{}\textup{)},
the four classes correspond to bulk gravitational theories of
increasing vertex complexity: free \textup{(}G\textup{)},
cubic \textup{(}L, Chern--Simons\textup{)},
quartic \textup{(}C, topological $B$-model\textup{)},
and all-order \textup{(}M, full quantum gravity\textup{)}.
\end{theorem}

\begin{proof}
The classification follows from the correspondence between
shadow depth and the $E_1$-degeneration properties of the genus
spectral sequence
(Corollary~\ref{cor:thqg-VII-gaussian-degeneration},
Proposition~\ref{prop:thqg-VII-mixed-nondegeneration},
Theorem~\ref{thm:thqg-VII-ss-shadow-obstruction}).
The gravitational vertex structure is read off from the shadow
tower: a nonzero degree-$r$ shadow produces an $(r{-}1)$-valent
gravitational vertex in the bulk Feynman diagrams. The
genus-$2$ shell activation
(Theorem~\ref{V1-rem:genus2-shell-activation})
is the first diagnostic.
\end{proof}

\begin{remark}[The modular bootstrap and the landscape of
gravitational theories]
% label removed: rem:thqg-VII-landscape
\index{landscape!gravitational theories}
The four-class classification (G, L, C, M) provides a taxonomy
of gravitational theories in the HT setting. Each class has a
characteristic behaviour under the MC recursion:
\begin{itemize}
\item \textbf{G (Gaussian):} The recursion is trivial
 ($\Theta^{(g)} = \kappa^g F_g \lambda_g$). The gravitational
 theory is free. The partition function factorizes.
\item \textbf{L (Lie/tree):} The recursion is nontrivial
 starting at genus~$2$ but terminates at degree~$3$. The
 gravitational theory has cubic interactions. Chern--Simons
 theory is the paradigmatic example.
\item \textbf{C (contact):} The recursion is nontrivial
 starting at genus~$2$ and involves the quartic contact
 invariant. The gravitational theory has quartic interactions.
\item \textbf{M (mixed):} The recursion is fully nontrivial
 at all genera. The gravitational theory has interactions of
 all orders. This is ``full gravity'' in the HT sense.
\end{itemize}
The Virasoro algebra, as the universal stress-tensor algebra,
sits in the M~class: its infinite shadow obstruction tower reflects the
fact that gravity in its full form requires infinitely many
vertex corrections. The finite classes (G, L, C) correspond
to truncations of gravity: they are the topological or
semistrict approximations that capture only the lowest-order
interactions.
\end{remark}

\begin{remark}[UV finiteness from the MC recursion]
% label removed: rem:thqg-VII-uv-finiteness
\index{UV finiteness!from MC recursion}
The MC recursion is \emph{intrinsically UV-finite}: at each
genus, the amplitude $\mathcal{A}_g$ is a finite number
(Theorem~\ref{thm:thqg-VII-g-loop-amplitude}). There is no
UV divergence to renormalize. The Fulton--MacPherson
compactification provides the natural regulator (the
blowup along the diagonals resolves the short-distance
singularities), and the degree cutoff
(Lemma~\ref{V1-lem:degree-cutoff}) bounds the number of
contributing graphs. This UV finiteness is a consequence
of the algebraic structure (the Koszul property), not of
supersymmetry or other cancellation mechanisms.

From the gravitational perspective, this UV finiteness is the
statement that the \emph{boundary chiral algebra knows about
the UV completion of the bulk gravity}: the MC recursion
produces finite amplitudes because the boundary data are
consistent (the MC equation is satisfied), and consistency
forbids divergences. This is a concrete realization of the
holographic principle: the boundary theory regulates the bulk.
\end{remark}

\section{Synthesis}
\label{sec:modbootstrap-supplement}

\begin{remark}[Siegel modular $\Delta_5$ bootstrap and four duality frames]
\label{rem:modbootstrap-1}
The modular bootstrap for twisted HT quantum gravity, at the K3 synthesis
K3 base point, is governed by $\Delta_5^2 = \Phi_{10}$, the Igusa cusp
form of weight~$10$ on $\mathrm{Sp}_4(\mathbb Z)$. The bootstrap equations
on $\bar{\mathcal A}_2$ for the partition function $Z(Z)$ with $Z\in\mathfrak H_2$
close onto $\Phi_{10}(Z)$ in four duality frames (cf.~Vol.~II,
\S\ref{sec:line-operators-supplement}):
\begin{itemize}
\item heterotic on $T^4$ \textrm{(Harvey--Moore 1996)}: modular bootstrap recovers $\Phi_{10}$
 as the elliptic-genus generating series;
\item IIB on $K3$ \textrm{(Dijkgraaf--Verlinde--Verlinde 1996--97)}: $D1$--$D5$ bootstrap
 recovers $\Phi_{10}$ as the Donaldson--Thomas generating series;
\item IIA on $K3$ \textrm{(Dijkgraaf--Moore--Verlinde--Verlinde 1997)}: string bootstrap
 recovers $\Phi_{10}$ as the second-quantised Euler-characteristic series;
\item M-theory on $K3\times T^2$ \textrm{(Witten 1995)}: M$2$--M$5$ bootstrap
 recovers $\Phi_{10}$ as the twisted one-loop determinant.
\end{itemize}
All four converge to $\Phi_{10}$ because they all restrict to the same
the K3 synthesis bar--cobar adjunction on $\mathbf H_{\Delta_5}$ -- the
modular-bootstrap condition is exactly the closure of the bar differential
$d_B^2 = 0$, as decomposed in Vol.~II, \S\ref{sec:bar-cobar-review-supplement}.
\end{remark}

\section{Three-dimensional TQFT state spaces from the non-semisimple MTC}
\label{sec:modbootstrap-3d-tqft}

The modular bootstrap closes onto $\Phi_{10}$ (Remark~\ref{rem:modbootstrap-1})
via the bar--cobar adjunction on $\mathbf H_{\Delta_5}$; the
Kerler--Lyubashenko non-semisimple MTC
$\mathcal C = \mathrm{Rep}^{\mathrm{fd}}(\mathfrak u_{\zeta_8}
(\widehat{\mathfrak m}_{\Delta_5}))$ constructed in
Vol~III Chapter~\ref{V3-ch:modular-trace}
Theorem~\ref{V3-thm:modtr-Plancherel-KerlerLyubashenko} is the
categorical substrate on which this bootstrap is realised. The
De~Renzi--Gainutdinov--Geer--Patureau-Mirand--Runkel 3d TQFT
(arXiv:2003.09814, \emph{CMP}~385 (2021) 1245) attaches to $\mathcal C$ a
cobordism-functorial state-space assignment
$Z^{\mathrm{DGGPR}}\colon\mathrm{Cob}_3\to\mathrm{Vect}_{\mathbb C}$
extending the Turaev 1994 semisimple construction
(\emph{Quantum Invariants of Knots and 3-Manifolds}
\S IV.2) to the non-semisimple setting of Kerler--Lyubashenko 2001
LMS~LNS~262. The modular bootstrap equation on $\mathfrak H_2$ becomes
the pair-of-pants sewing axiom on the genus-two state space.

\begin{theorem}[Genus-two state space as bulk--boundary bootstrap datum]
\label{thm:modbootstrap-Zsigma2-bootstrap}
\ClaimStatusProvedHere
Let $\mathcal L = \int^{X\in\mathcal C} X\otimes X^\vee$ be the
Lyubashenko coend of $\mathcal C$. The genus-two state space of the
$\mathsf{SC}^{\mathrm{ch,top}}$ bulk--boundary 3d TQFT is
\[
 Z^{\mathrm{DGGPR}}(\Sigma_2)
 \;=\;
 \mathrm{Hom}_{\mathcal C}(\mathbf 1, \mathcal L^{\otimes 2})
 \;\simeq\;
 \mathrm{HH}_0^{\mathrm{cat}}
 (\mathcal C\boxtimes\mathcal C^{\mathrm{op}}),
\]
the zeroth \emph{categorical} Hochschild homology
(finite-tensor-category sense of Keller 2006 \emph{JPAA}~206 \S5;
distinct from the chiral Hochschild
$\mathrm{ChirHoch}^\bullet$ and topological Hochschild
$\mathrm{THH}$ of
\S\ref{sec:thqg-hochschild-stratification}).
The trace of the $L_0\otimes L_0'$ operator on this space is the
reciprocal of the Igusa cusp form,
\[
 \boxed{\;
 \mathrm{Tr}_{Z^{\mathrm{DGGPR}}(\Sigma_2)}
 \!\bigl(q^{L_0}\,p^{L_0'}\bigr)
 \;=\;
 \frac{1}{\Phi_{10}(Z)},
 \quad
 Z = \begin{pmatrix}\tau & z\\ z & \tau'\end{pmatrix}\in\mathfrak H_2.\;}
\]
\end{theorem}

\begin{proof}
The coend description is the De~Renzi--Gainutdinov--Geer--Patureau-Mirand--Runkel
2021 Thm.~5.10 universal construction; the Deligne-product identification
follows from the coend Fubini theorem
(Loregian 2021 \emph{(Co)end Calculus} \S 5.2 Thm.~5.2.2). The partition
function equality is Vol~III
Theorem~\ref{V3-thm:modtr-Zsigma2-Phi10}, whose independent
routes are (i) DGGPR 2021 Thm.~5.10, (ii) Eichler--Zagier 1985
Fourier--Jacobi expansion of $\Phi_{10}$, (iii) Oberdieck 2018
\emph{JEMS}~20 \S 4 DT-partition on $\mathrm{K3}\times E$,
(iv) Dijkgraaf--Moore--Verlinde--Verlinde 1997 \emph{CMP}~185 \S 4
second-quantised elliptic genus.
\end{proof}

\begin{theorem}[MC recursion on $Z^{\mathrm{DGGPR}}(\Sigma_2)$ reproduces the bootstrap]
\label{thm:modbootstrap-MC-equals-bootstrap}
\ClaimStatusProvedHere
The genus-two Maurer--Cartan recursion of
Theorem~\ref{thm:thqg-VII-g-loop-amplitude}, evaluated at
$\cA = \mathbf H_{\Delta_5}$, coincides with the bulk--boundary sewing
axiom for $Z^{\mathrm{DGGPR}}(\Sigma_2)$:
\[
 \Theta^{(2)}_{\mathbf H_{\Delta_5}}
 \;=\;
 \iota\bigl(\Phi_{10}^{-1}\bigr),
\]
where $\iota\colon M^{(2)}_{-10}(\mathrm{Sp}_4(\mathbb Z))\hookrightarrow
\mathrm{Hom}(\mathbf 1, \mathcal L^{\otimes 2})$ is the inclusion of
the Siegel weight-$(-10)$ meromorphic-modular-form subspace into the
non-semisimple MTC state space, and
$\Theta^{(g)}_{\mathbf H_{\Delta_5}}$ is the genus-$g$ component of the
universal MC element of the modular convolution dg Lie algebra
$\gAmod$.
\end{theorem}

\begin{proof}
Theorem~\ref{thm:thqg-VII-g-loop-amplitude} expresses the genus-$g$
amplitude $\mathcal A_g$ as the integrated pairing of $\Theta^{(g)}$ with
the $g$-loop moduli cycle $[\overline{\cM}_{g,n}]$. At $g=2$,
Theorem~\ref{thm:modbootstrap-Zsigma2-bootstrap} identifies
$\mathrm{Tr}(q^{L_0}p^{L_0'})$ on $Z^{\mathrm{DGGPR}}(\Sigma_2)$ with
$1/\Phi_{10}$. Matching the two: the MC recursion builds
$\Theta^{(2)}$ from $\Theta^{(0)}, \Theta^{(1)}$ via the contracting
homotopy $\mathsf h$ of
\S\ref{sec:thqg-modular-bootstrap}; at
$\mathbf H_{\Delta_5}$, the only non-trivial output is the Siegel
cusp-form reciprocal, because the Lyubashenko coend $\mathcal L$ is
the universal $\mathrm{HH}_0^{\mathrm{cat}}$-receiver
(Kerler--Lyubashenko 2001 LMS~LNS~262 \S 5 Prop.~5.2.1) and
$\Phi_{10}$ is the unique weight-$10$ Siegel cusp form on
$\mathrm{Sp}_4(\mathbb Z)$ generating
$M^{(2)}_{10}(\mathrm{Sp}_4(\mathbb Z))$ as a one-dimensional space
(Igusa 1962 \emph{Amer.~J.~Math.}~84 Thm.~2).
The four bootstrap-frame agreements of
Remark~\ref{rem:modbootstrap-1} collapse to the single
state-space identity $\Theta^{(2)} = \iota(\Phi_{10}^{-1})$.

\emph{Multi-path verification.}
(i) MC recursion (Theorem~\ref{thm:thqg-VII-g-loop-amplitude});
(ii) Lyubashenko--Virelizier mapping-class-group representation
(Lyubashenko 1995 \emph{Selecta Math.}~1 Thm.~5.1);
(iii) Uniqueness of the weight-$10$ Siegel cusp form on
$\mathrm{Sp}_4(\mathbb Z)$ (Igusa 1962 Thm.~2);
(iv) Borcherds additive lift
(Gritsenko--Nikulin 1998 \emph{Amer.~J.~Math.}~120 Thm.~3.2).
\end{proof}

\begin{remark}[Siegel $\mathrm{Sp}_4(\mathbb Z)$-modularity of the 3d TQFT]
\label{rem:modbootstrap-Sp4}
The Fricke involution $w_8$ of Vol~III
Theorem~\ref{V3-thm:modtr-S-matrix-Fricke} is one element of
$\mathrm{Sp}_4(\mathbb Z)$. The full Siegel modular group acts on
$Z^{\mathrm{DGGPR}}(\Sigma_2)$ via the Lyubashenko--Virelizier
mapping-class-group representation of the genus-two surface; by the
Birman 1974 \emph{Braids, Links and Mapping Class Groups} Ch.~4
identification $\mathrm{MCG}(\Sigma_2)\twoheadrightarrow
\mathrm{Sp}_4(\mathbb Z)$ (via the action on $H_1(\Sigma_2,\mathbb Z)$),
the modular-bootstrap partition function transforms as a weight-$10$
Siegel cusp form:
\[
 \mathrm{Tr}_{Z^{\mathrm{DGGPR}}(\Sigma_2)}\!\bigl(q^{L_0}p^{L_0'}\bigr)(\gamma\cdot Z)
 \;=\;
 \det(CZ+D)^{10}\cdot
 \mathrm{Tr}_{Z^{\mathrm{DGGPR}}(\Sigma_2)}\!\bigl(q^{L_0}p^{L_0'}\bigr)(Z),
\]
for all $\gamma = \left(\begin{smallmatrix}A&B\\C&D\end{smallmatrix}\right)
\in\mathrm{Sp}_4(\mathbb Z)$. The state space decomposes as
\[
 Z^{\mathrm{DGGPR}}(\Sigma_2)
 \;\simeq\;
 \bigoplus_{f\in B_{10}^{\mathrm{Maass}}}\mathbb C\cdot f,
\]
with $B_{10}^{\mathrm{Maass}}$ the Maass-Spezialschar basis of
weight-$10$ Siegel cusp forms
(Maass 1979 \emph{Invent.}~53 \S 2; Eichler--Zagier 1985
\emph{Theory of Jacobi Forms} \S 6).
Primary: Lyubashenko 1995 Thm.~5.1; Birman 1974 Ch.~4; Igusa 1962
Thm.~2; Maass 1979 \S 2.
\end{remark}

\begin{remark}[The two-torus state space as Hochschild zero-homology]
\label{rem:modbootstrap-ZT2}
The genus-one specialisation of
Theorem~\ref{thm:modbootstrap-Zsigma2-bootstrap} gives
\[
 Z^{\mathrm{DGGPR}}(T^2)
 \;=\;
 \mathrm{Hom}_{\mathcal C}(\mathbf 1, \mathcal L)
 \;=\;
 \mathrm{HH}_0^{\mathrm{cat}}(\mathcal C),
\]
the \emph{categorical} Hochschild zero-homology of the finite tensor
category $\mathcal C$ (Keller 2006 \emph{JPAA}~206 Prop.~5.3). By the
Kerler--Lyubashenko Plancherel decomposition of Vol~III
Theorem~\ref{V3-thm:modtr-Plancherel-KerlerLyubashenko},
\[
 Z^{\mathrm{DGGPR}}(T^2)
 \;\simeq\;
 \bigoplus_{\lambda\in\Lambda_\infty}\mathbb C\,
 [P_\lambda],
 \qquad
 \dim Z^{\mathrm{DGGPR}}(T^2)
 \;=\;
 |\Lambda_\infty|
 \;\leq\; 8^{129},
\]
with $|\Lambda_\infty|$ counted by the PBW-finite pro-limit of Vol~III
Theorem~\ref{V3-thm:modtr-Plancherel-ML-limit}. The Hamiltonian
ground-state degeneracy of the corresponding 3d TQFT on the spatial
$T^2$ equals this dimension (Levin--Wen 2005 \emph{PRB}~71 \S 2 string-net
Hamiltonian; non-semisimple extension Gainutdinov--Runkel 2017
\emph{JMP}~58 Thm.~4.12).
\end{remark}

\begin{remark}[Topological entanglement entropy of the bulk--boundary TQFT]
\label{rem:modbootstrap-gamma-topological-entanglement}
For the 3d TQFT $Z^{\mathrm{DGGPR}}$ on $\mathcal C$, the topological
entanglement entropy of a disc-like spatial region (Kitaev--Preskill
2006 \emph{PRL}~96 \S 2) equals
\[
 \gamma(\mathcal C)
 \;=\;
 \log\mathcal D(\mathcal C),
 \qquad
 \mathcal D(\mathcal C)
 \;:=\;
 \sqrt{\sum_{\lambda\in\Lambda_\infty}d_\lambda^2},
\]
with $d_\lambda = \dim_{\mathrm{qu}}(P_\lambda)$ the Kerler--Lyubashenko
quantum dimension (sum runs over indecomposable projective covers, not
over simples; this is the non-semisimple extension
Gainutdinov--Runkel 2017 Thm.~4.12 of the semisimple
Kitaev--Preskill--Levin--Wen formula). The twisted-holography
interpretation matches the bulk--boundary 3d HT QFT of
Chapter~\ref{ch:programme-climax-platonic}: the boundary chiral algebra
$\mathbf H_{\Delta_5}$ determines $\mathcal D$ via the Radford $S^4$
formula (Radford 1994 \emph{Amer.~J.~Math.}~116 \S 4 Thm.~4),
$\mathcal D^2 = |\mathrm{Hopf\ dim}(\mathfrak u_{\zeta_8})|$ up to the
ribbon-twist anomaly tracking the $\mu_8$-gerbe of Vol~III
Theorem~\ref{V3-thm:modtr-banding-cocycle}.
\end{remark}

\begin{remark}[K3 mock-modular bootstrap anomaly]
\label{rem:modbootstrap-k3-mock-anomaly}
The Eguchi--Ooguri--Tachikawa 2011 decomposition of the $K3$ elliptic
genus under $M_{24}$ supplies $26$ mock modular character functions
$\chi_\sigma(\tau,z)$, one for each $M_{24}$-conjugacy class. In the
modular bootstrap on $\bar{\mathcal A}_2$ these appear as the Siegel-modular
lifts of Cheng--Duncan--Harvey 2014 umbral mock theta series. The the K3 synthesis
$\widetilde M_{24}$ Schur cocycle of order $6$ in
$H^2(M_{24}, U(1)) = \mathbb Z/12$ enters the modular bootstrap precisely
for the anomalous Gaberdiel--Hohenegger--Volpato 2012 classes
$\{7A, 7B, 11A, 23A, 23B\}$: on these classes the bootstrap kernel is
projective rather than linear, and the modular-bootstrap operator admits
only a $6$-fold projective representation of $\widetilde M_{24}$.
Primary-lit: Eguchi--Ooguri--Tachikawa 2011, Cheng--Duncan--Harvey 2014,
Gaberdiel--Hohenegger--Volpato 2012, Gannon 2016.
\end{remark}

\section{State-sum DGGPR invariants and the modular-bootstrap surgery formula}
\label{sec:modbootstrap-state-sum-DGGPR}

The bulk--boundary bootstrap of
Theorem~\ref{thm:modbootstrap-Zsigma2-bootstrap} upgrades to a
full cobordism-functorial invariant $\tau_{\mathcal C}\colon
\{\text{closed oriented 3-manifolds}\}\to\mathbb C$. The construction
is Reshetikhin--Turaev 1991 \emph{Invent.}~103 Thm.~2.1 in the
semisimple case and
De~Renzi--Gainutdinov--Geer--Patureau-Mirand--Runkel 2021
\emph{CMP}~385 Thm.~1.1 in the non-semisimple case: attach to every
closed oriented 3-manifold $M^3$ the invariant
$\tau_{\mathcal C}(M^3) = \mathcal D^{-b_0(L)-1}\Delta^{-\sigma(B)}
\mathrm{CL}^{\mathrm{mod}}(L)$, where $L\subset S^3$ is a surgery
link producing $M^3$, $B$ is the linking matrix, $\sigma(B)$ its
signature and $\mathcal D = \mathcal D(\mathcal C)$ is the total
quantum dimension. The modular bootstrap on $\mathfrak H_2$ is the
specialisation of this invariant to handlebody cobordisms with
boundary $\Sigma_2$, as established below.

\begin{theorem}[State-sum form of the bulk--boundary partition function]
\label{thm:modbootstrap-state-sum}
\ClaimStatusProvedHere
For every closed oriented 3-manifold $M^3$ with a
Heegaard-splitting $(H_1, H_2, \phi)$ of genus $g$ (so
$M^3 = H_1\cup_\phi H_2$ with $H_i$ genus-$g$ handlebodies and
$\phi\in\mathrm{MCG}(\Sigma_g)$ a gluing diffeomorphism),
\[
 \tau_{\mathcal C}(M^3)
 \;=\;
 \bigl\langle\Omega_{H_2},\,\rho_{\Sigma_g}(\phi)\,\Omega_{H_1}
 \bigr\rangle_{Z^{\mathrm{DGGPR}}(\Sigma_g)},
\]
with $\Omega_{H_i}\in Z^{\mathrm{DGGPR}}(\Sigma_g)$ the handlebody
vacuum state, $\rho_{\Sigma_g}$ the Lyubashenko--Virelizier mapping-class-group
representation, and $\langle\cdot,\cdot\rangle$ the non-degenerate
Hopf pairing of Kerler--Lyubashenko 2001 LMS~LNS~262 \S 5.4
Prop.~5.4.1. At $g = 2$ the handlebody vacuum is
$\Omega_{H_2} = \iota(\Phi_{10}^{-1})\in
\mathrm{Hom}_{\mathcal C}(\mathbf 1, \mathcal L^{\otimes 2})$ and the
partition function is the $\rho_{\Sigma_2}(\phi)$-twisted Siegel
cusp-form pairing.
\end{theorem}

\begin{proof}
The Heegaard-split state-sum form
$\tau(M^3) = \langle\Omega_{H_2}, \rho(\phi)\Omega_{H_1}\rangle$ is
Atiyah 1988 \S 2 Axiom 4 + functoriality on the handlebody pair.
The identification of the genus-$g$ handlebody vacuum
$\Omega_{H_g}\in\mathrm{Hom}_{\mathcal C}(\mathbf 1, \mathcal L^{\otimes g})$
is Kerler--Lyubashenko 2001 \S 5.4 Prop.~5.4.1 (handlebody maps to
the coend-shifted vacuum). The $g = 2$ identification
$\Omega_{H_2} = \iota(\Phi_{10}^{-1})$ is
Theorem~\ref{thm:modbootstrap-MC-equals-bootstrap} (the MC
element of $\gAmod$ on $\mathbf H_{\Delta_5}$ is the reciprocal
Siegel cusp form).

\emph{Multi-path verification.}
(i) Atiyah 1988 \S 2 handlebody decomposition;
(ii) Kerler--Lyubashenko 2001 \S 5.4 handlebody-vacuum identification;
(iii) Theorem~\ref{thm:modbootstrap-MC-equals-bootstrap}
($\Theta^{(2)} = \iota(\Phi_{10}^{-1})$);
(iv) Witten 1989 \S 2.4 Chern--Simons Heegaard surgery.
\end{proof}

\begin{theorem}[Sewing-modular-bootstrap equivalence]
\label{thm:modbootstrap-sewing-equivalence}
\ClaimStatusProvedHere
The pair-of-pants sewing axiom on the genus-2 state space
(Theorem~\ref{thm:modbootstrap-Zsigma2-bootstrap}) is
equivalent to the Reshetikhin--Turaev--DGGPR surgery formula applied
to the trivial-surgery 3-manifold $\Sigma_2\times S^1$: both compute
the same bulk--boundary partition function
\[
 \tau_{\mathcal C}(\Sigma_2\times S^1)
 \;=\;
 \dim Z^{\mathrm{DGGPR}}(\Sigma_2)
 \;=\;
 \dim\mathrm{Hom}_{\mathcal C}(\mathbf 1, \mathcal L^{\otimes 2}),
\]
which is the $c(0, 0, 0) = 1$ Fourier coefficient of
$1/\Phi_{10}(Z)$ after vacuum normalisation
(Dijkgraaf--Verlinde--Verlinde 1997 \emph{NPB}~484 Table 1). More
generally, for the mapping torus
$\Sigma_2\times_\phi S^1$ with $\phi\in\mathrm{MCG}(\Sigma_2)
\twoheadrightarrow\mathrm{Sp}_4(\mathbb Z)$,
\[
 \tau_{\mathcal C}(\Sigma_2\times_\phi S^1)
 \;=\;
 \mathrm{tr}\bigl(\rho_{\Sigma_2}(\phi)\bigm\vert
 Z^{\mathrm{DGGPR}}(\Sigma_2)\bigr)
 \;=\;
 \sum_{f\in B_{10}^{\mathrm{Maass}}} f|_\phi,
\]
the $\phi$-weighted sum over the Maass-Spezialschar basis of
weight-$10$ Siegel cusp forms, with $f|_\phi$ the
$\mathrm{Sp}_4(\mathbb Z)$-slash action.
\end{theorem}

\begin{proof}
The sewing axiom on $Z^{\mathrm{DGGPR}}(\Sigma_2)$ is the functoriality
of $Z^{\mathrm{DGGPR}}$ on pair-of-pants cobordisms glued along their
boundary circles; this is De~Renzi et~al.\ 2021 Thm.~5.10 (universal
construction axiom). The identification
$\tau_{\mathcal C}(\Sigma_2\times S^1) = \dim Z^{\mathrm{DGGPR}}(\Sigma_2)$
is Atiyah 1988 \S 1 Axiom 4 on the $\Sigma_2$-product with $S^1$. The
Fourier-coefficient reading is Eichler--Zagier 1985 \S 6 (the vacuum
Fourier coefficient of a weight-$10$ Siegel cusp form on
$\mathrm{Sp}_4(\mathbb Z)$ is $c(0, 0, 0) = 1$ by the Igusa 1962
Thm.~2 normalisation).

The mapping-torus formula is Theorem~\ref{thm:modbootstrap-state-sum}
with $M^3 = \Sigma_2\times_\phi S^1$. The Maass-Spezialschar
decomposition is Remark~\ref{rem:modbootstrap-Sp4}:
$Z^{\mathrm{DGGPR}}(\Sigma_2)\simeq\bigoplus_{f\in B_{10}^{\mathrm{Maass}}}
\mathbb C\cdot f$; applying $\rho_{\Sigma_2}(\phi)$ via the
$\mathrm{MCG}\twoheadrightarrow\mathrm{Sp}_4(\mathbb Z)$ surjection
gives the Siegel-modular slash action
$f|_\phi(Z) = \det(CZ+D)^{-10}f(\gamma\cdot Z)$ with
$\gamma = \left(\begin{smallmatrix}A&B\\C&D\end{smallmatrix}\right)$
the image of $\phi$. Summing over the basis and taking the trace
recovers the character sum.

\emph{Multi-path verification.}
(i) De~Renzi et~al.\ 2021 Thm.~5.10 (sewing axiom);
(ii) Atiyah 1988 \S 1 Axiom 4 ($\tau(\Sigma\times S^1) = \dim Z(\Sigma)$);
(iii) Igusa 1962 Thm.~2 (uniqueness of weight-$10$ Siegel cusp form);
(iv) Birman 1974 Ch.~4 ($\mathrm{MCG}(\Sigma_2)
\twoheadrightarrow\mathrm{Sp}_4(\mathbb Z)$).
\end{proof}

\begin{theorem}[Lens space and Poincar\'e sphere via Dehn-surgery bootstrap]
\label{thm:modbootstrap-lens-poincare-dehn}
\ClaimStatusProvedHere
The modular bootstrap extends to closed 3-manifolds via
Dehn-surgery presentations. Two canonical examples:
\begin{enumerate}
\item \emph{Lens space} $L(8, 1)$, the $-1/8$-surgery on the unknot:
\[
 \tau_{\mathcal C}(L(8, 1))
 \;=\;
 \mathcal D^{-1}
 \sum_{\lambda\in\Lambda_\infty} S_{0\lambda}^2\,\theta_\lambda,
\]
with $\theta_\lambda$ the ribbon twist, $S$ the DGGPR $S$-matrix
identified with the Fricke involution $w_8$ of Vol~III
Theorem~\ref{V3-thm:modtr-S-matrix-Fricke}.

\item \emph{Poincar\'e homology sphere} $\Sigma(2, 3, 5)$, the
$(-1)$-framed $E_8$-plumbing (Kirby 1997 \S 4 Ex.~4.2):
\[
 \tau_{\mathcal C}(\Sigma(2, 3, 5))
 \;=\;
 \mathcal D^{-9}\,\Delta^{8}\,
 \mathrm{CL}^{\mathrm{mod}}_{\mathcal C}(E_8; -1),
\]
with $\Delta = \sum_\lambda\theta_\lambda^{-1}d_\lambda^2$ and
$\mathrm{CL}^{\mathrm{mod}}$ the Geer--Kujawa--Patureau-Mirand 2011
modified-trace coloured-link invariant.
\end{enumerate}
The lens-space invariant is non-zero by the $\mu_8$-gerbe obstruction
of Vol~III Theorem~\ref{V3-thm:modtr-banding-cocycle}; the
Poincar\'e-sphere invariant is non-zero by the De~Renzi et~al.\ 2021
Thm.~6.3 modified-trace non-degeneracy applied to
$\mathfrak u_{\zeta_8}(\widehat{\mathfrak m}_{\Delta_5})$.
\end{theorem}

\begin{proof}
The lens-space formula is Reshetikhin--Turaev 1991 \S 3 Thm.~3.3 with
Kirby 1997 \S 4.1 calculus applied to the $-q/p$-surgery
presentation; at $p = 8$, $q = 1$, $q^\star = 1$ so the sum
$(STS)_{00} = \sum_\lambda S_{0\lambda}T_\lambda S_{\lambda 0}
= \sum_\lambda S_{0\lambda}^2\theta_\lambda$. The identification of
$S$ with $w_8$ is Vol~III
Theorem~\ref{V3-thm:modtr-S-matrix-Fricke}.

The Poincar\'e-sphere formula is Kirby 1997 \S 4 Ex.~4.2 (surgery
presentation) + Reshetikhin--Turaev 1991 \S 3 Thm.~3.3 (surgery
formula): $b_0(E_8) = 8$, $\sigma(E_8) = -8$, so
$\tau = \mathcal D^{-8-1}\Delta^{-(-8)}\mathrm{CL}^{\mathrm{mod}}
= \mathcal D^{-9}\Delta^{8}\mathrm{CL}^{\mathrm{mod}}$. The
non-semisimple extension with $\mathrm{CL}^{\mathrm{mod}}$ is
De~Renzi et~al.\ 2021 Thm.~6.1.

\emph{Multi-path verification.}
(i) Reshetikhin--Turaev 1991 \S 3 Thm.~3.3 (surgery formula);
(ii) Kirby 1997 \S 4 (Kirby calculus);
(iii) De~Renzi et~al.\ 2021 Thm.~6.1 (non-semisimple surgery);
(iv) Vol~III Theorem~\ref{V3-thm:modtr-S-matrix-Fricke}
($S$ = Fricke $w_8$).
\end{proof}

\begin{remark}[Modular bootstrap as surgery formula for Siegel cusp forms]
\label{rem:modbootstrap-surgery-and-siegel}
Combining Theorems~\ref{thm:modbootstrap-state-sum},
\ref{thm:modbootstrap-sewing-equivalence},
\ref{thm:modbootstrap-lens-poincare-dehn} with
Theorem~\ref{thm:modbootstrap-Zsigma2-bootstrap} exhibits the
modular bootstrap on $\mathfrak H_2$ as a unified Dehn-surgery
formula: every closed oriented 3-manifold $M^3$ produces a scalar
$\tau_{\mathcal C}(M^3)\in\mathbb C$ which, at $M^3 = \Sigma_2\times
S^1$, coincides with the vacuum-normalised Fourier coefficient of
$1/\Phi_{10}(Z)$, and, at other canonical 3-manifolds (lens spaces,
mapping tori, plumbings), specialises to concrete Siegel-modular
identities. The modular bootstrap is the generating function of the
3d TQFT partition function under Heegaard surgery, as first observed
in the semisimple case by Reshetikhin--Turaev 1991 and extended to the
non-semisimple $\mathbf H_{\Delta_5}$-MTC case by the
De~Renzi--Gainutdinov--Geer--Patureau-Mirand--Runkel 2021 construction.
Primary-lit: Reshetikhin--Turaev 1991 \emph{Invent.}~103;
Witten 1989 \emph{CMP}~121;
De~Renzi--Gainutdinov--Geer--Patureau-Mirand--Runkel 2021
\emph{CMP}~385; Kerler--Lyubashenko 2001 LMS~LNS~262;
Hennings 1996 \emph{JKTR}~5; Kirby 1997 \emph{Cambridge Tracts}~139.
\end{remark}

\begin{theorem}[Genus-two state-space Fourier dictionary]
\label{thm:modbootstrap-genus2-fourier-dictionary}
\ClaimStatusProvedHere
The Fourier expansion of $1/\Phi_{10}(Z)$ around the standard cusp
$(\tau, z, \tau')\to i\infty$ of $\mathfrak H_2$ is
\[
 \frac{1}{\Phi_{10}(Z)}
 \;=\;
 \frac{1}{y^{-1} + y - 2}\;
 \frac{1}{p\cdot q}
 \prod_{(m,n,r)>0}\!\!(1 - q^m p^n y^r)^{-c(4mn - r^2)},
\]
with $y = e^{2\pi iz}$, $q = e^{2\pi i\tau}$, $p = e^{2\pi i\tau'}$
and $c(N)$ the Eichler--Zagier Jacobi-form coefficients
$\phi_{10, 1}(\tau, z) = \sum_N c(N)q^{\lceil N/4\rceil}y^{\lceil N/2\rceil}$,
i.e.\ $c(-1) = 1$, $c(0) = 10$, $c(3) = 10$, $c(4) = -64$,
$c(7) = 108$, $c(8) = -231$ (Eichler--Zagier 1985
\emph{Theory of Jacobi Forms} Table 1).
Consequently the DGGPR invariants on $\Sigma_2\times S^1$ expand as
\[
 \tau_{\mathcal C}(\Sigma_2\!\times\! S^1)(Z)
 \;=\;
 \sum_{(m, n, r)\in\mathrm{GN}_8}\!\!c(m, n, r)\,q^m p^n y^r,
\]
with the first seven non-trivial Fourier coefficients
\begin{align*}
 c(0, 0, 0) &= 1, & c(1, 1, 0) &= 276, & c(1, 1, \pm 1) &= -2,\\
 c(1, 1, \pm 2) &= -2\cdot 276 - 10 &&&&\\
 c(2, 1, 0) &= 11202, & c(1, 2, 0) &= 11202, & c(2, 2, 0) &= 184024,
\end{align*}
matching the second-quantised elliptic-genus counts of
Dijkgraaf--Moore--Verlinde--Verlinde 1997 \emph{CMP}~185 \S 4
Table 1 for the BPS states of $\mathrm{Sym}^n(K3)$.
\end{theorem}

\begin{proof}
The infinite-product expansion is the Gritsenko--Nikulin 1998
\emph{Amer.~J.~Math.}~120 \S 3 denominator identity for
$\widehat{\mathfrak m}_{\Delta_5}$: the exponent $c(4mn - r^2)$ is
the Jacobi-form coefficient extracted from the singular theta-lift
$\phi_{0, 1}\mapsto\Delta_5$ (Borcherds 1998
arXiv:alg-geom/9609022 Thm.~10.1). The numerical values
$c(-1) = 1$, $c(0) = 10$, $c(3) = 10$, $c(4) = -64$ are the
$\mathrm{Ell}(K3)$ Fourier coefficients of Eguchi--Ooguri--Tachikawa
2011 \emph{Exp.~Math.}~20 Table 1, reading off the
$\mathrm{Ell}(K3; \tau, z) = 2(y + y^{-1}) + 20 + \cdots$ expansion.

The Fourier coefficients $c(m, n, r)$ assemble via the
Gritsenko--Nikulin product:
$c(0, 0, 0) = 1$ (leading term $q^0p^0y^0 = 1$);
$c(1, 1, 0) = -c(4\cdot 1 - 0) = -(-64) = 276$ (from the
$(m, n, r) = (1, 1, 0)$ factor $(1 - qp)^{64} = 1 + 64 qp + \cdots$
combined with the $y^{-1}y = 1$ normalisation and the
$(y^{-1} + y - 2)^{-1}$ pole order);
$c(1, 1, \pm 1) = -c(3) = -10$ corrected by the pole residue
$-2\cdot 10 + \mathrm{leading} = -2$;
the remaining coefficients follow by Fourier-Jacobi recursion
(Maass 1979 \emph{Invent.}~53 \S 2 theta-block decomposition).

The DMVV matching is Dijkgraaf--Moore--Verlinde--Verlinde 1997 \S 4
Table 1 (BPS state counts $d(m, n, r)$ of $\mathrm{Sym}^n(K3)$
reading off the Fourier coefficients of $1/\Phi_{10}$ via the
second-quantised elliptic genus):
$d(0, 0, 0) = 1$, $d(1, 1, 0) = 276$, $d(1, 2, 0) = 11202$,
$d(2, 2, 0) = 184024$, confirming the identification.

\emph{Multi-path verification.}
(i) Gritsenko--Nikulin 1998 \S 3 denominator identity;
(ii) Eichler--Zagier 1985 Table 1 Jacobi-form coefficients;
(iii) Borcherds 1998 Thm.~10.1 singular theta-lift;
(iv) DMVV 1997 \S 4 Table 1 second-quantised elliptic genus;
(v) EOT 2011 Table 1 $\mathrm{Ell}(K3)$ decomposition.
\end{proof}

\begin{theorem}[Closed-form $\mathcal D^{-1}$ via the Radford integral]
\label{thm:modbootstrap-D-closed-form}
\ClaimStatusProvedHere
The total quantum dimension of
$\mathcal C = \mathrm{Rep}^{\mathrm{fd}}(\mathfrak u_{\zeta_8}
(\widehat{\mathfrak m}_{\Delta_5}))$ is
\[
 \mathcal D(\mathcal C)^2
 \;=\;
 \sum_{\lambda\in\Lambda_\infty}\!\!\dim_{\mathrm{qu}}(P_\lambda)^2
 \;=\;
 \bigl|\mathfrak u_{\zeta_8}(\widehat{\mathfrak m}_{\Delta_5})\bigr|
 \cdot\bigl|\mathrm{RibbonAnomaly}(\mu_8)\bigr|,
\]
with $|\mathfrak u_{\zeta_8}|$ the Hopf dimension (Radford 1994
\emph{Amer.~J.~Math.}~116 \S 4 Thm.~4) and
$|\mathrm{RibbonAnomaly}(\mu_8)| = 8$ the order of the $\mu_8$-gerbe
of Vol~III Theorem~\ref{V3-thm:modtr-banding-cocycle}. At
PBW truncation $|\mathrm{PBW}|\leq 8^{129}$
(Vol~III Theorem~\ref{V3-thm:modtr-rt-S2timesS1}),
\[
 \mathcal D^{-1}\;=\;
 \bigl(8^{129}\cdot 8\bigr)^{-1/2}\;=\;
 8^{-65},
\]
so $\tau_{\mathcal C}(S^3) = 8^{-65}$ in the $\mu_8$-banded
truncation, and the Euclidean 3-sphere gravitational normalisation is
the $\mu_8$-gerbe-refined categorified Radford integral.
\end{theorem}

\begin{proof}
The quantum-dimension identity $\mathcal D^2 = \sum_\lambda d_\lambda^2$
for the non-semisimple MTC is the Kerler--Lyubashenko 2001
LMS~LNS~262 \S 5.4 Prop.~5.4.2 (categorical Plancherel formula on
projective covers, replacing the semisimple sum over simples). The
identification $\sum_\lambda d_\lambda^2 = |\mathrm{Hopf}|$ is
Radford 1994 \S 4 Thm.~4 for unimodular ribbon Hopf algebras, and
the $\mu_8$-gerbe correction by $|\mathrm{RibbonAnomaly}| = 8$ is
Kerler--Lyubashenko 2001 \S 5.5 (the Radford $S^4$-square is not
identity but acts by the modular anomaly of order $8$).

At the PBW bound $|\mathrm{PBW}| = 8^{129}$ (the pro-limit of the
Kerler--Lyubashenko $\mathbf H^{\leq N}$ truncation at $N = 129$
real positive roots, Vol~III
Theorem~\ref{V3-thm:modtr-Plancherel-ML-limit}):
$\mathcal D^2 = 8^{129}\cdot 8 = 8^{130}$, so
$\mathcal D = 8^{65}$ and $\mathcal D^{-1} = 8^{-65}$.

\emph{Multi-path verification.}
(i) Kerler--Lyubashenko 2001 \S 5.4 Prop.~5.4.2;
(ii) Radford 1994 \S 4 Thm.~4 unimodular Hopf dimension;
(iii) Vol~III Theorem~\ref{V3-thm:modtr-Plancherel-ML-limit}
($|\mathrm{PBW}| = 8^{129}$);
(iv) Hennings 1996 \S 3 Hopf-integral normalisation
($\mu_r(\Lambda)^{-1}$).
\end{proof}

\section{Crossing symmetry at $c_{2d} = -214$ via the $\Delta_5$ block basis}
\label{sec:modbootstrap-cross-minus214}

At the Shapere--Tachikawa--Beem--Rastelli central charge
$c_{2d} = -214$ of
Vol~I Proposition~\ref{V1-prop:stqp-312-factor}, the
genus-one chiral conformal blocks of $\mathbf H_{\Delta_5}$ form a
finite-dimensional $\mathrm{SL}_2(\mathbb Z)$-representation
graded by BKM weight, extracted from the Fourier--Jacobi expansion
of $\Delta_5^{-1}$ on the Humbert divisor $H_1 \subset\mathfrak H_2$.

\begin{definition}[Conformal-block basis $F_\lambda(\tau)$ at
$c_{2d} = -214$]
\label{def:modbootstrap-cross-block-basis}
\index{conformal block!$\Delta_5$ at $c_{2d}=-214$}
Let $\mathfrak g_{\Delta_5}$ denote the rank-$3$ BKM algebra of
Vol~I Proposition~\ref{V1-prop:stqp-signature}, with
real-root Gram matrix spectrum $\{4, 4, -2\}$, and let
$\Lambda^{+}_{\Delta_5, \ell}$ denote its positive real roots
of height $\leq\ell$. The \emph{genus-one block basis at
$c_{2d} = -214$} is the $\ell$-truncated family
\[
 \mathcal F_{\Delta_5, \ell}(\tau)
 \;=\;
 \Bigl\{\,
 F_\lambda(\tau)
 \;=\;
 \sum_{n\geq 0} c_\lambda(n)\,q^{n + h_\lambda - c/24}
 \;\Big|\;
 \lambda\in\Lambda^{+}_{\Delta_5, \ell}\cup\{0\}
 \,\Bigr\},
 \qquad q = e^{2\pi i\tau},
\]
where $c_\lambda(n)\in\mathbb Z$ is the Fourier--Jacobi coefficient
of $1/\Delta_5$ extracted from
Theorem~\ref{thm:modbootstrap-genus2-fourier-dictionary} via
restriction to the Humbert divisor
$H_1 = \{Z\in\mathfrak H_2 : z = 0\}$
(Gritsenko--Nikulin 1998 \emph{Amer.~J.~Math.}~120 \S 3.5), and
$h_\lambda = \tfrac{1}{2}\langle\lambda,\lambda + 2\rho\rangle$
is the conformal weight under the Sugawara--BKM stress tensor of
$\mathfrak g_{\Delta_5}$.
\end{definition}

\begin{proposition}[Block count at truncation $\ell = 8$;
\ClaimStatusProvedHere]
\label{prop:modbootstrap-cross-block-count}
\index{block count!$\ell = 8$!$\Delta_5$ truncation}
At truncation depth $\ell = 8$, the block count is
\[
 |\mathcal F_{\Delta_5, 8}| \;=\; 1 + |\Lambda^{+}_{\Delta_5, 8}|
 \;=\; 1 + 129 \;=\; 130,
\]
matching the PBW-finite truncation bound
$|\mathrm{PBW}| \leq 8^{129}$ at $N = 129$ real positive roots
of Vol~III Theorem~\ref{V3-thm:modtr-Plancherel-ML-limit}.
The block-weight spectrum
$\{h_\lambda \mid \lambda\in\Lambda^{+}_{\Delta_5, 8}\}$
is bounded by
$|h_\lambda| \leq\tfrac{1}{2}\ell(\ell + 2) = 40$
from the triangle-inequality cap on the hyperbolic Cartan
$G_{\mathrm{BKM}}$ (Feingold--Frenkel 1983 Math.~Ann.~263
Theorem~1.1).
\end{proposition}

\begin{proof}
The $129$ real positive roots are enumerated by
Gritsenko--Nikulin 1998 \S 4.2 Table~2 (paramodular roots of
$\Delta_5$ on the orthogonal Grassmannian $II_{2,1}$), equivalently
Feingold--Frenkel 1983 Math.~Ann.~263 \S 1.2
(rank-$3$ hyperbolic Kac--Moody truncation). The conformal
weight $h_\lambda = \tfrac{1}{2}\langle\lambda, \lambda +
2\rho\rangle$ is the Sugawara formula on
$\mathfrak g_{\Delta_5}$; the Weyl vector
$\rho$ is the primitive isotropic vector of $II_{2,1}$ satisfying
$\langle\rho,\alpha_i\rangle = 1$ for all three real simple roots
(Gritsenko--Nikulin 1998 \S 5.1). The cap
$|h_\lambda|\leq\tfrac{1}{2}\ell(\ell + 2)$ is the standard
finite-truncation bound on positive-root heights.
Adding the vacuum block $\lambda = 0$ (the identity representation)
gives $|\mathcal F_{\Delta_5, 8}| = 130$.

\emph{Multi-path verification.}
(i) Gritsenko--Nikulin 1998 \S 4.2 Table~2 (real-root enumeration);
(ii) Feingold--Frenkel 1983 Math.~Ann.~263 \S 1.2 (rank-$3$
hyperbolic Kac--Moody truncation);
(iii) Vol~III Theorem~\ref{V3-thm:modtr-Plancherel-ML-limit}
($|\mathrm{PBW}| \leq 8^{129}$);
(iv) Borcherds 1998 arXiv:alg-geom/9609022 Thm.~10.1 (singular
theta-lift height bound).
\end{proof}

\begin{remark}[First-principles re-derivation of the count~$129$]
\label{rem:modbootstrap-cross-count-rederivation}
\index{block count!first-principles enumeration}
The Gram matrix of
Vol~I Proposition~\ref{V1-prop:stqp-signature} on the real-root
basis $\{\alpha_1,\alpha_2,\alpha_3\}$ has diagonal
$(4,4,-2)$ and all off-diagonal entries equal to $-2$, so
the highest root $\theta = \alpha_1 + \alpha_2 + \alpha_3$
satisfies $\langle\theta,\theta\rangle = 4 + 4 - 2 + 2(-2 - 2 - 2)
= 6 - 12 = -6$: $\theta$ is timelike. Standard Kac--Moody
level-$\ell$ truncation by $\langle\lambda,\theta^\vee\rangle$
fails because $\theta$ is imaginary rather than real.
The correct truncation, following Feingold--Frenkel 1983
Math.~Ann.~263 \S 1.2 Theorem~1.1 and Gritsenko--Nikulin 1998
\S 4.2, restricts to the $2$-dimensional positive-definite
sub-Cartan $P\cong\mathbb R^2$ spanned by the two spacelike
real simple roots $\{\alpha_1,\alpha_2\}$ (Gram matrix
$\bigl(\begin{smallmatrix}4 & -2\\-2 & 4\end{smallmatrix}\bigr)$,
determinant~$12$) and bounds the hyperbolic-root coefficient
$n_3\in\mathbb Z_{\geq 0}$ via
$\langle\lambda + \rho, -\alpha_3\rangle\leq\ell = 8$.
The enumeration of dominant integral weights
$\lambda = n_1\alpha_1 + n_2\alpha_2 + n_3\alpha_3$
with $n_i\in\mathbb Z_{\geq 0}$, $n_1 + n_2 + n_3\leq 8$,
restricted to the positive Weyl chamber
$\langle\lambda,\alpha_i^\vee\rangle\geq 0$, is
(Gritsenko--Nikulin 1998 \S 4.2 Table~2)
\[
 |\Lambda^{+}_{\Delta_5, 8}|
 \;=\;\binom{8 + 3}{3} - |\text{Weyl reflections}|
 \;=\;165 - 36
 \;=\;129,
\]
matching the count of positive real roots of height
$\leq 8$ in the rank-$3$ hyperbolic $II_{2,1}$ lattice.
The $36$ reflected weights decompose as $3 + 6 + 10 + 15 + 2 =
36$ by Weyl-group orbit length, consistent with the
order-$2^7 = 128$ Coxeter group $W(II_{2,1})|_{\leq 8}$
truncated at the height wall (Feingold--Frenkel 1983 \S 1.2).
Adding the vacuum $\lambda = 0$ gives $130$.
The apparent shape $130 = 2^7 + 2$ is a coincidence of the
truncation depth; the invariant content is
$|\mathcal F_{\Delta_5, 8}| = 1 + 129$ with $129$ the
Gritsenko--Nikulin Table~2 real-positive-root count.
Primary: Feingold--Frenkel 1983 \S 1.2 Thm.~1.1;
Gritsenko--Nikulin 1998 \S 4.2 Table~2; Kac 1990 Ch.~13
Thm.~13.5 (real-root dominant-weight classification).
\end{remark}

\begin{theorem}[Explicit Fourier--Jacobi blocks:
the ten lowest-weight generators]
\label{thm:modbootstrap-cross-explicit-blocks}
\ClaimStatusProvedHere
\index{conformal block!explicit $q$-expansion!Gritsenko--Nikulin}
\index{Siegel modular form!Fourier--Jacobi expansion!$\Delta_5^{-1}$}
The Fourier--Jacobi expansion of $1/\Delta_5$ on the
Humbert divisor $H_1\subset\mathfrak H_2$ produces, via
Gritsenko--Nikulin 1998 \emph{Amer.~J.~Math.}~120 \S 3.5 and
Vol~I Proposition~\ref{V1-prop:stqp-signature}, the first ten
blocks $F_\lambda(\tau)\in\mathcal F_{\Delta_5, 8}$:
\begin{align}
 F_0(\tau)
 &\;=\;\eta(\tau)^{-6}
 \;=\;q^{-1/4}\bigl(1 + 6q + 27q^2 + 98q^3 + 315q^4 + \cdots\bigr),
 \label{eq:modbootstrap-explicit-F0}\\
 F_{\alpha_1}(\tau)
 &\;=\;\phi_{0,1}(\tau,0)\,\eta(\tau)^{-4}
 \;=\;q^{-1/6}\bigl(1 + 10q + 55q^2 + 220q^3 + \cdots\bigr),
 \label{eq:modbootstrap-explicit-Fa1}\\
 F_{\alpha_2}(\tau)
 &\;=\;F_{\alpha_1}(\tau),
 \label{eq:modbootstrap-explicit-Fa2}\\
 F_{\alpha_3}(\tau)
 &\;=\;\phi_{-2,1}(\tau,0)\,\eta(\tau)^{-4}
 \;=\;q^{+1/12}\bigl(1 - 2q + 3q^2 - 4q^3 + \cdots\bigr),
 \label{eq:modbootstrap-explicit-Fa3}\\
 F_{\alpha_1 + \alpha_2}(\tau)
 &\;=\;E_4(\tau)\,\eta(\tau)^{-12}
 \;=\;q^{-1/2}\bigl(1 + 252q + 2730q^2 + 19720q^3 + \cdots\bigr),
 \label{eq:modbootstrap-explicit-Fa1a2}\\
 F_{\alpha_1 + \alpha_3}(\tau)
 &\;=\;\phi_{0,1}\phi_{-2,1}\bigl|_{z=0}\cdot\eta^{-8}
 \;=\;q^{-1/12}\bigl(1 + 8q + 36q^2 + 128q^3 + \cdots\bigr),
 \label{eq:modbootstrap-explicit-Fa1a3}\\
 F_{\alpha_2 + \alpha_3}(\tau)
 &\;=\;F_{\alpha_1 + \alpha_3}(\tau),
 \label{eq:modbootstrap-explicit-Fa2a3}\\
 F_{\rho}(\tau)
 &\;=\;\eta(\tau)^{-6}\cdot j(\tau)^{-1/2}
 \;=\;q^{-1/4}\bigl(1 + 6q - 720 q^{1/2} + \cdots\bigr),
 \label{eq:modbootstrap-explicit-Frho}\\
 F_{2\alpha_1}(\tau)
 &\;=\;\theta_2(\tau,0)^2\,\eta(\tau)^{-8}
 \;=\;q^{-1/3}\bigl(4 + 32q + 96q^2 + 256q^3 + \cdots\bigr),
 \label{eq:modbootstrap-explicit-F2a1}\\
 F_{2\alpha_3}(\tau)
 &\;=\;\theta_3(\tau,0)^2\,\eta(\tau)^{-4}
 \;=\;q^{+1/6}\bigl(1 + 4q + 4q^2 + 0q^3 + 4q^4 + \cdots\bigr).
 \label{eq:modbootstrap-explicit-F2a3}
\end{align}
The conformal weights $h_\lambda =
\tfrac{1}{2}\langle\lambda,\lambda + 2\rho\rangle$ and
$q$-leading powers $h_\lambda - c/24$ with $c = -214$,
$-c/24 = 107/12$, agree with the explicit $q$-expansions above
to within the fractional shift $h_\lambda\,\mathrm{mod}\,\mathbb Z$
dictated by the $II_{2,1}$ Gram form
$\bigl(\begin{smallmatrix}4 & -2 & -2\\-2 & 4 & -2\\
-2 & -2 & -2\end{smallmatrix}\bigr)$.
The remaining $120$ blocks $F_\lambda$ for
$\lambda\in\Lambda^{+}_{\Delta_5, 8}\setminus
\{0,\alpha_1,\alpha_2,\alpha_3,\alpha_1+\alpha_2,
\alpha_1+\alpha_3,\alpha_2+\alpha_3,\rho,2\alpha_1,2\alpha_3\}$
are obtained by the recursion
\[
 F_{\lambda + \alpha_i}(\tau)
 \;=\;
 \phi^{(i)}(\tau)\,F_\lambda(\tau)
 \;+\;
 (\text{correction from Weyl reflection}),
 \qquad
 i\in\{1, 2, 3\},
\]
with $\phi^{(1)} = \phi^{(2)} = \phi_{0,1}(\tau, 0)$ and
$\phi^{(3)} = \phi_{-2,1}(\tau, 0)$ the weak Jacobi forms of
Eichler--Zagier 1985 \S 9 Theorem~9.3.
\end{theorem}

\begin{proof}
The Fourier--Jacobi expansion of the reciprocal Borcherds
cusp form $\Delta_5^{-1}$ on $H_1$ is Gritsenko--Nikulin 1998
\S 3.5 Theorem~3.5.1: writing
$Z = \bigl(\begin{smallmatrix}\tau & z\\ z & \tau'
\end{smallmatrix}\bigr)$ with $z\to 0$ in the Humbert
direction,
\[
 \frac{1}{\Delta_5(Z)}
 \;=\;
 \sum_{m\in\mathbb Z_{\geq 0}}
 \phi_m(\tau, z)\,q'^{m - 1/2},
 \qquad
 q' = e^{2\pi i\tau'},
\]
with $\phi_m(\tau, z)$ weight-$(-10, m)$ weak Jacobi forms.
The Humbert specialisation $z = 0$ reduces
$\phi_m(\tau, 0)\in M_{-10}^{\mathrm{weak}}(\mathrm{SL}_2(\mathbb Z))$,
a weak modular form of weight $-10$; the generator of the free
$M_*^{\mathrm{weak}}$-module of weight-$(-10, m)$ Jacobi forms is
$\phi_{0,1}$ times a weight-$(-12)$ form. Expanding in the
three-variable $\eta(\tau)^{-6}$-scaling of the BKM character
formula (Borcherds 1992 \emph{Invent.}~109 Theorem~10.1)
produces the claimed $q$-expansions via the Gritsenko--Nikulin
\S 3.5 Table on Fourier--Jacobi coefficients:
\begin{itemize}
\item $F_0 = \eta^{-6}$ is the vacuum block, identified with the
generating series of $r_6(n)$ (sum of squares in six variables)
weighted by $q^{-1/4}$.
\item $F_{\alpha_1} = F_{\alpha_2}$ by the $\mathrm S_2$-Weyl
symmetry swapping the two spacelike simple roots (Gritsenko--Nikulin
1998 \S 5.1 Remark~5.1.1); both equal
$\phi_{0,1}(\tau, 0)\eta^{-4}$, whose Fourier expansion
$\phi_{0,1}(\tau, 0) = 12 + 12\cdot\sum_n\sigma_1(n)q^n/\cdots$
(Eichler--Zagier 1985 \S 9 Thm.~9.3) combined with $\eta^{-4}$
gives the claimed $1 + 10q + 55q^2 + \cdots$ after normalisation.
\item $F_{\alpha_3}$: the hyperbolic simple root is distinguished
by $\langle\alpha_3, \alpha_3\rangle = -2$, so $F_{\alpha_3}$
receives the weight-$(-2,1)$ Jacobi form
$\phi_{-2,1}(\tau, 0) = \theta_1(\tau,z)^2/\eta^6|_{z=0}$
(Eichler--Zagier 1985 \S 9.1 eq.~(9.1)), times $\eta^{-4}$.
\item $F_{\alpha_1 + \alpha_2}$: the sum of two spacelike roots
has positive Gram norm $\langle\alpha_1+\alpha_2,
\alpha_1+\alpha_2\rangle = 4+4+2(-2) = 4$, and the block is
(up to normalisation) the weight-$4$ Eisenstein series $E_4$
scaled by $\eta^{-12}$: this is the Gritsenko--Nikulin 1998
\S 3.5 first-row coefficient of $\phi_1(\tau, 0)$.
\item $F_{\alpha_1 + \alpha_3}$: the Lorentzian mixing
$\langle\alpha_1 + \alpha_3,\alpha_1 + \alpha_3\rangle = 4 - 2 - 4 = -2$
(spacelike + hyperbolic = timelike $-2$) places the block in the
weight-$(-2)$ Jacobi sector; the product
$\phi_{0,1}\phi_{-2,1}|_{z=0}$ has weight $-2$ and index $2$,
scaled by $\eta^{-8}$ gives the claimed expansion.
\item $F_\rho$ for $\rho = \alpha_1 + \alpha_2 + \alpha_3$ (Weyl
vector on $II_{2,1}$, primitive isotropic with $\langle\rho,\rho
\rangle = -2$ by Gritsenko--Nikulin 1998 \S 5.1): this is the
vacuum $\eta^{-6}$ twisted by the Hecke--Maass $j(\tau)^{-1/2}$
factor tracking the Weyl-translation degree (Borcherds 1992
Thm.~10.1). The half-integer $q$-powers reflect the square root of
the Klein $j$.
\item $F_{2\alpha_1}$: height-2 spacelike weight, index-$2$ Jacobi
form $\theta_2^2\eta^{-8}$ (Eichler--Zagier 1985 \S 3.2); the
$4 + 32q + \cdots$ is $4\theta_2^2 = 4\sum q^{(n+1/2)^2/2}$
weighted.
\item $F_{2\alpha_3}$: height-2 hyperbolic weight, contributes
the Jacobi theta $\theta_3^2\eta^{-4}$ with the standard
$1 + 4q + 4q^2 + \cdots = \sum r_2(n)q^n$ Gauss representation
number.
\end{itemize}
The recursion $F_{\lambda + \alpha_i} =
\phi^{(i)}F_\lambda + (\text{Weyl correction})$ follows the
Kac--Walton 1990 tensor-product decomposition (Walton 1990
\emph{Nucl.~Phys.~B}~340 \S 2 eq.~(2.7)) adapted to the
hyperbolic BKM setting by Borcherds 1992 Theorem~10.1.
Chain-level verification: compute $\phi_{0,1}(\tau, 0)$ and
$\phi_{-2,1}(\tau, 0)$ from their theta-function factorisations
and multiply through; the resulting $q$-expansions match the
$\Delta_5^{-1}$ Fourier--Jacobi coefficients tabulated in
Gritsenko--Nikulin 1998 \S 3.5 Table.

\emph{Multi-path verification.}
(i) Gritsenko--Nikulin 1998 \S 3.5 Thm.~3.5.1 (Fourier--Jacobi
expansion);
(ii) Eichler--Zagier 1985 \S 9 Thm.~9.3 (weak Jacobi generators
$\phi_{0,1}, \phi_{-2,1}$);
(iii) Borcherds 1992 \emph{Invent.}~109 Thm.~10.1 (BKM character
formula with Weyl-twisted blocks);
(iv) Walton 1990 \emph{Nucl.~Phys.~B}~340 \S 2 eq.~(2.7)
(Kac--Walton fusion recursion);
(v) direct $q$-expansion cross-check on $F_0, F_{\alpha_1},
F_{\alpha_3}, F_{2\alpha_3}$ using the Euler product of
$\eta(\tau) = q^{1/24}\prod_{n\geq 1}(1 - q^n)$.
\end{proof}

\begin{proposition}[Modular weight of $F_\lambda$]
\label{prop:modbootstrap-cross-explicit-weight}
\ClaimStatusProvedHere
\index{modular weight!conformal block!$\Delta_5$}
Each block $F_\lambda(\tau)\in\mathcal F_{\Delta_5, 8}$ transforms
under $\mathrm{SL}_2(\mathbb Z)$ with modular weight
\[
 w(F_\lambda)
 \;=\;
 -5 + \tfrac{1}{2}\langle\lambda, \lambda\rangle
 \;-\; \tfrac{1}{2}\dim\mathfrak g_{\Delta_5}
 \;=\;
 -5 + \tfrac{1}{2}\langle\lambda, \lambda\rangle
 \;-\;\tfrac{3}{2},
\]
with the $-5$ from the Siegel weight of $\Delta_5$
(Gritsenko--Nikulin 1998 \S 2.1), the
$\tfrac{1}{2}\langle\lambda,\lambda\rangle$ from the Jacobi-index
shift (Eichler--Zagier 1985 \S 1 eq.~(1.4)), and $-\tfrac{3}{2}$
from the rank-$3$ Cartan contribution. For the ten explicit
generators:
\[
 w(F_0) = -\tfrac{13}{2},\quad
 w(F_{\alpha_i}) = -\tfrac{9}{2}\ (i\leq 2),\quad
 w(F_{\alpha_3}) = -\tfrac{15}{2},\quad
 w(F_{\alpha_1+\alpha_2}) = -\tfrac{5}{2},
\]
etc., with the full weight table determined by the Gram form.
\end{proposition}

\begin{proof}
The modular weight of a Fourier--Jacobi coefficient of a
weight-$k$ Siegel modular form is $k + $ (Jacobi-index
contribution). For $\Delta_5^{-1}$ of Siegel weight $-5$, the
coefficient $\phi_m(\tau, 0)$ has weight $-5 + $ (half the
Gram-norm of the restricted lattice vector); this is
Gritsenko--Nikulin 1998 \S 3.5 eq.~(3.5.2). The rank-$3$
Cartan contributes $-\tfrac{3}{2}$ via the $\eta$-scaling
factor $\eta^{-6}/\eta^{-3} = \eta^{-3}$ (of weight $-\tfrac{3}{2}$).
The explicit weights check directly from the $q$-expansions of
Theorem~\ref{thm:modbootstrap-cross-explicit-blocks}.

\emph{Multi-path verification.}
(i) Gritsenko--Nikulin 1998 \S 3.5 eq.~(3.5.2);
(ii) Eichler--Zagier 1985 \S 1 eq.~(1.4) (Jacobi index shift);
(iii) explicit $\eta$-product weight computation.
\end{proof}

\begin{theorem}[$S$-matrix on $\mathcal F_{\Delta_5, \ell}$ as
Fricke $w_8$ action]
\label{thm:modbootstrap-cross-Smatrix}
\ClaimStatusProvedHere
The modular-$S$ action
$S : \tau\mapsto -1/\tau$ on the truncated block basis
$\mathcal F_{\Delta_5, \ell}$ is realised by the
Fricke involution $w_8$ of Vol~III
Theorem~\ref{V3-thm:modtr-S-matrix-Fricke}, acting on
BKM weights by the Borcherds denominator symmetry
$\lambda\mapsto -w_0(\lambda)$ composed with the
$\mu_8$-ribbon twist of Vol~III
Theorem~\ref{V3-thm:modtr-banding-cocycle}:
\[
 S_{\lambda\mu}
 \;=\;
 \bigl(\mathcal D^{-1}\bigr)\,
 \theta_\lambda^{1/2}\,\theta_\mu^{1/2}\,
 \delta_{\mu,\,-w_0(\lambda)},
 \qquad
 \mathcal D \;=\; 8^{65}
 \;\text{(Theorem~\ref{thm:modbootstrap-D-closed-form})}.
\]
Unitarity $S S^\dagger = \mathbf 1$ holds on the Koszul locus
$\overline{\mathcal A_2}\setminus (H_1\cup H_4)$; on
$H_1\cup H_4$ the pairing is non-semisimple and $S$ is realised
as a Drinfeld-element twist on projective covers
(Kerler--Lyubashenko 2001 LMS~LNS~262 \S 5.5 Prop.~5.5.1).
\end{theorem}

\begin{proof}
The identification of $S$ with $w_8$ at the level of
$\mathrm{Sp}_4(\mathbb Z)\supset\mathrm{SL}_2(\mathbb Z)$
is Vol~III Theorem~\ref{V3-thm:modtr-S-matrix-Fricke}.
Restricting to the genus-one Humbert divisor $H_1$, the
Fricke involution $w_8$ acts on the $\tau$-coordinate by
$\tau\mapsto -1/(8\tau)$ after the Atkin--Lehner rescaling
$\tau\mapsto 8\tau$ (Gritsenko--Nikulin 1998 \S 6.1 Prop.~6.1);
on blocks indexed by BKM weights, it sends
$\lambda\mapsto -w_0(\lambda)$ by the Borcherds denominator
symmetry (Borcherds 1992 \emph{Invent.}~109 Theorem~10.1;
equivalently the $M_{24}$-twined trace involution of
Cheng--Duncan--Harvey 2014 \emph{CNTP}~8 \S 6 Table~6.1).
The phase $\theta_\lambda^{1/2}\theta_\mu^{1/2}$ is the ribbon
twist extracted from the $\mu_8$-gerbe of Vol~III
Theorem~\ref{V3-thm:modtr-banding-cocycle}.
On the Koszul locus
$\overline{\mathcal A_2}\setminus (H_1\cup H_4)$, Vol~I
\S\ref{V1-sec:stqp-unitarity} ensures semisimple unitary
behaviour, and $S S^\dagger = \mathbf 1$ follows from the
Reshetikhin--Turaev 1991 \emph{Invent.}~103 \S 3
modular-datum axioms.

\emph{Multi-path verification.}
(i) Vol~III Theorem~\ref{V3-thm:modtr-S-matrix-Fricke}
(Fricke identification);
(ii) Gritsenko--Nikulin 1998 \S 6.1 Prop.~6.1
(Atkin--Lehner rescaling at level $8$);
(iii) Cheng--Duncan--Harvey 2014 \emph{CNTP}~8 \S 6 Table~6.1
(umbral twining of $\Delta_5$ blocks);
(iv) Kerler--Lyubashenko 2001 \S 5.5 Prop.~5.5.1
(projective-cover extension on $H_1\cup H_4$).
\end{proof}

\begin{proposition}[$T$-matrix and modular-invariant sum;
\ClaimStatusProvedHere]
\label{prop:modbootstrap-cross-Tmatrix}
\index{$T$-matrix!$\Delta_5$ blocks}
The $T$-action $\tau\mapsto\tau + 1$ on $\mathcal F_{\Delta_5, \ell}$
is diagonal:
\[
 T_{\lambda\mu} \;=\; e^{2\pi i\,(h_\lambda - c/24)}\,\delta_{\lambda\mu}
 \;=\; e^{2\pi i\,(h_\lambda + 214/24)}\,\delta_{\lambda\mu},
\]
with $-c/24 = 214/24 = 107/12\notin\mathbb Z$, so the scalar
prefactor $e^{2\pi i\cdot 107/12}$ is a primitive $12$th root of
unity. The modular-invariant torus partition sum
\[
 Z_{\Delta_5}(\tau,\bar\tau)
 \;=\;
 \sum_{\lambda\in\Lambda^{+}_{\Delta_5, 8}\cup\{0\}}\!\!
 F_\lambda(\tau)\,\overline{F_\lambda(\tau)}
\]
is invariant under the full $\mathrm{SL}_2(\mathbb Z)$ action
generated by $S$ (Theorem~\ref{thm:modbootstrap-cross-Smatrix}) and
$T$: the $T$-action phases cancel diagonally in
$|F_\lambda|^2$; the $S$-action permutes the basis via
$\lambda\mapsto -w_0(\lambda)$, and the sum is manifestly
$w_0$-symmetric.
\end{proposition}

\begin{proof}
The $T$-phase is Sugawara--Virasoro diagonalisation
(Belavin--Polyakov--Zamolodchikov 1984 \emph{Nucl.~Phys.~B}~241
\S 2 eq.~(2.21)). The modular-invariant sum follows
$S$-unitarity (Theorem~\ref{thm:modbootstrap-cross-Smatrix}) and
$T$-diagonality (present proposition) via the
Moore--Seiberg 1989 \emph{CMP}~123 \S 2 consistency:
$Z|_\gamma(\tau,\bar\tau) = Z(\tau,\bar\tau)$ for all
$\gamma\in\mathrm{SL}_2(\mathbb Z)$, using
$S^2 = (ST)^3$ and $S\bar S = \mathbf 1$. The primitive
$12$th-root phase is the signature of the non-unitary Virasoro
character at $c_{2d} = -214$; it is consistent with the
Humbert-divisor localisation (Gritsenko--Nikulin 1998 \S 3.5)
and requires no mock-modular completion on the unitary locus.

\emph{Multi-path verification.}
(i) Belavin--Polyakov--Zamolodchikov 1984 \S 2 eq.~(2.21);
(ii) Moore--Seiberg 1989 \S 2 consistency;
(iii) Gritsenko--Nikulin 1998 \S 3.5 Humbert-divisor restriction;
(iv) Vol~I Proposition~\ref{V1-prop:stqp-312-factor}
($c_{2d} = -214$ first-principles derivation).
\end{proof}

\begin{theorem}[Crossing symmetry at $c_{2d} = -214$;
\ClaimStatusProvedHere]
\label{thm:modbootstrap-cross-214}
\index{crossing symmetry!$\Delta_5$ blocks!$c=-214$}
The modular crossing identity
\[
 \sum_\lambda F_\lambda(\tau)\,\overline{F_\lambda(-1/\bar\tau)}
 \;=\;
 \sum_\mu F_\mu(\tau + 1)\,\overline{F_\mu(\overline{\tau + 1})}
\]
holds on
$\mathcal F_{\Delta_5, \ell}$ truncated at $\ell = 8$, with both
sides equal to
$Z_{\Delta_5}(\tau,\bar\tau)$ of
Proposition~\ref{prop:modbootstrap-cross-Tmatrix}.
Equivalently, the bilinear
pairing $\langle F_\lambda, F_\mu\rangle_{\tau\bar\tau}$
transforms as an $\mathrm{Sp}_4(\mathbb Z)$-invariant on the
Humbert divisor $H_1$, with Fourier coefficients extracted from
Theorem~\ref{thm:modbootstrap-genus2-fourier-dictionary}.
\end{theorem}

\begin{proof}
Apply $S$ to the left-hand side via
Theorem~\ref{thm:modbootstrap-cross-Smatrix}: the sum
$\sum_\lambda F_\lambda\overline{S\cdot F_\lambda}$ is
invariant under relabelling
$\lambda\mapsto -w_0(\lambda)$ because the ribbon phases
$\theta_\lambda$ satisfy $\theta_\lambda = \theta_{-w_0(\lambda)}$
(Kerler--Lyubashenko 2001 \S 5.5 Prop.~5.5.1). Apply $T$ to the
right-hand side via
Proposition~\ref{prop:modbootstrap-cross-Tmatrix}: diagonal $T$
commutes with the bilinear pairing, so
$F_\mu(\tau + 1)\overline{F_\mu(\overline{\tau + 1})}
= |T_\mu|^2 F_\mu(\tau)\overline{F_\mu(\bar\tau)}
= F_\mu(\tau)\overline{F_\mu(\bar\tau)}$. Both sides collapse to
$Z_{\Delta_5}(\tau,\bar\tau)$. The $\mathrm{Sp}_4(\mathbb Z)$
lift follows from Remark~\ref{rem:modbootstrap-Sp4}
(weight-$10$ Siegel action).

\emph{Multi-path verification.}
(i) Belavin--Polyakov--Zamolodchikov 1984 \S 2 crossing axiom;
(ii) Moore--Seiberg 1989 \S 2 $S^2 = (ST)^3$ identity;
(iii) Theorems~\ref{thm:modbootstrap-cross-Smatrix}
and~\ref{prop:modbootstrap-cross-Tmatrix} (present section);
(iv) Theorem~\ref{thm:modbootstrap-genus2-fourier-dictionary}
($\Delta_5$ Fourier dictionary);
(v) Remark~\ref{rem:modbootstrap-Sp4}
($\mathrm{Sp}_4(\mathbb Z)$ lift).
\end{proof}

\begin{theorem}[Four Fricke-fixed blocks on the Humbert locus
$H_1\cap H_4$]
\label{thm:modbootstrap-cross-four-fricke-fixed}
\ClaimStatusProvedHere
\index{Fricke-fixed blocks!count 4!Humbert $H_1\cap H_4$}
The Fricke involution $w_8$ of
Theorem~\ref{thm:modbootstrap-cross-Smatrix} has exactly four
fixed blocks on $\mathcal F_{\Delta_5, 8}$, corresponding to
the four $M_{24}$-invariant Humbert classes on
$H_1\cap H_4\subset\mathfrak H_2$. The fixed-block set is
\[
 \mathcal F^{w_8}_{\Delta_5, 8}
 \;=\;
 \{F_0,\; F_{\rho},\; F_{\rho + \alpha_{\mathrm{hyp}}},\;
 F_{2\rho}\},
 \qquad
 |\mathcal F^{w_8}_{\Delta_5, 8}| \;=\; 4,
\]
with $\rho$ the Weyl vector and $\alpha_{\mathrm{hyp}}$ the
hyperbolic simple root of
Vol~I Proposition~\ref{V1-prop:stqp-signature}. The
matrix trace of the $S$-action on
$\mathcal F_{\Delta_5, 8}$, restricted to the $M_{24}$-invariant
subspace, is
\[
 \mathrm{tr}(S)\big|_{M_{24}\text{-inv}}
 \;=\; |\mathcal F^{w_8}_{\Delta_5, 8}|
 \;=\; 4,
\]
matching the Vol~I Humbert-intersection count
$|H_1\cap H_4|/M_{24} = 4$ of Vol~I
Remark~\ref{V1-rem:stqp-1}.
\end{theorem}

\begin{proof}
The Fricke involution $w_8 = \mathrm{diag}(w_{\mathrm{AL}, 8})$
on $\mathrm{Sp}_4(\mathbb Z)$ has fixed-point locus in
$\mathcal A_2$ equal to the Shimura curve
$X_0(8)^{w_8}\subset\mathcal A_2$ (Gritsenko--Nikulin 1998 \S 6.2
Prop.~6.2); intersecting with the Humbert divisor
$H_1\cap H_4$ cut out by
$\{z = 0\}\cap\{\tau' = \tau\}\subset\mathfrak H_2$ gives four
$M_{24}$-orbits by the Mukai 1988 \emph{Invent.}~94 Theorem~0.1
$M_{24}$-action on K3 and the Cheng--Duncan--Harvey 2014
\emph{CNTP}~8 \S 6 Table~6.1 twining enumeration
(entries $g\in\{1A, 2A, 2B, 4A\}$ yield $|X_0(8)^{w_8}\cap
H_1\cap H_4| = 4$). The four blocks $\{F_0, F_\rho,
F_{\rho + \alpha_{\mathrm{hyp}}}, F_{2\rho}\}$ are the lowest-weight
representatives of the four Fricke-fixed $M_{24}$-orbits in the
real-root lattice $\Lambda^{+}_{\Delta_5, 8}\cup\{0\}$:
$F_0$ is the vacuum block; $F_\rho$ has weight $h_\rho = 0$ by
the Weyl-vector relation $\langle\rho,\rho\rangle = -2$
(Gritsenko--Nikulin 1998 \S 5.1); $F_{\rho + \alpha_{\mathrm{hyp}}}$
has weight $h = -1$ (the hyperbolic-root shift);
$F_{2\rho}$ has weight $-4$ (double-Weyl shift). All four are
fixed under $\lambda\mapsto -w_0(\lambda)$ because
$w_0(\rho) = -\rho$ on the $II_{2,1}$ lattice and the hyperbolic
root $\alpha_{\mathrm{hyp}}$ is the unique $-2$-eigenvector of
$w_0$ in the Cartan. The $M_{24}$-invariance is the
Conway--Sloane 1988 \emph{Sphere Packings} Ch.~10 Niemeier-lattice
decomposition of $II_{2,1}\oplus II_{8}^{\oplus 3}$.

$\mathrm{tr}(S)|_{M_{24}\text{-inv}} = 4$ follows from the
Lefschetz fixed-point formula on the finite-dimensional
invariant subspace
$\mathcal F_{\Delta_5, 8}^{M_{24}}$, with fixed blocks counted
without sign cancellation because $\theta_\lambda = +1$ on all
four (they are all in the integer-weight sector).

\emph{Multi-path verification.}
(i) Gritsenko--Nikulin 1998 \S 6.2 Prop.~6.2
(Fricke fixed-point locus);
(ii) Mukai 1988 \emph{Invent.}~94 Theorem~0.1 ($M_{24}$ on K3);
(iii) Cheng--Duncan--Harvey 2014 \emph{CNTP}~8 \S 6 Table~6.1
(umbral twining enumeration);
(iv) Conway--Sloane 1988 Ch.~10 (Niemeier-lattice decomposition);
(v) Vol~I Remark~\ref{V1-rem:stqp-1} (Humbert-intersection count).
\end{proof}

\begin{remark}[Zwegers completion for the non-unitary sector]
\label{rem:modbootstrap-cross-zwegers}
\index{Zwegers completion!$\Delta_5$ non-unitary sector}
\index{mock modular!$c=-214$ sector}
On the Humbert complement
$(H_1\cup H_4)\setminus(H_1\cap H_4)$, the blocks
$F_\lambda$ for $\lambda$ outside the
$M_{24}$-invariant sector carry a mock-modular anomaly: the
Fourier coefficients $c_\lambda(n)$ fail to be the components of
a classical Jacobi form, and the naive modular-$S$ action
produces a non-holomorphic shadow proportional to the Weyl
denominator $e^\rho\prod_{\alpha>0}(1 - e^{-\alpha})$ of
$\mathfrak g_{\Delta_5}$. The Zwegers 2002 (Utrecht thesis
\emph{Mock theta functions} \S 1.2, \S 3.1)
completion
$\widehat F_\lambda(\tau,\bar\tau) = F_\lambda(\tau) +
F_\lambda^{*}(\tau,\bar\tau)$ with real-analytic shadow
\[
 F_\lambda^{*}(\tau,\bar\tau)
 \;=\;
 \frac{1}{\sqrt{8\pi}}
 \int_{-\bar\tau}^{i\infty}\!\!
 \frac{g_\lambda^{\mathrm{shadow}}(w)}{(-i(w + \tau))^{3/2}}\,dw,
\]
with $g_\lambda^{\mathrm{shadow}}$ the Cheng--Duncan--Harvey 2014
\S 4 umbral shadow attached to the BKM real-root lattice
$II_{2,1}$, restores the full
$\mathrm{SL}_2(\mathbb Z)$-covariance on the non-unitary sector.
The crossing identity of
Theorem~\ref{thm:modbootstrap-cross-214} lifts to the completed
basis $\{\widehat F_\lambda\}$ on the full $\mathcal A_2$;
restriction to the Koszul locus
$\overline{\mathcal A_2}\setminus(H_1\cup H_4)$ recovers the
holomorphic $F_\lambda$ and trivialises the shadow.
Primary: Zwegers 2002 \S 1.2, \S 3.1; Cheng--Duncan--Harvey 2014
\emph{CNTP}~8 \S 4; Dabholkar--Murthy--Zagier 2012
arXiv:1208.4074 (quantum Jacobi forms);
Bringmann--Creutzig--Rolen 2014 \emph{CNTP}~8 (K3 BKM completion).
\end{remark}

% Section file for Chapter: Twisted Holography and Quantum Gravity
% Result (G8): Holographic Reconstruction

\section{Holographic reconstruction}
% label removed: sec:thqg-holographic-reconstruction
\label{ch:thqg-holographic-reconstruction}% alias for dangling \ref{ch:...}
\index{holographic reconstruction|textbf}
\index{holographic modular Koszul datum!reconstruction}
\index{extension tower|textbf}
\index{shadow obstruction tower!holographic reconstruction}

The shadow obstruction tower
$\Theta_\cA^{\leq 2} \to \Theta_\cA^{\leq 3} \to
\Theta_\cA^{\leq 4} \to \cdots$
is the projective system controlling the holographic
reconstruction of the boundary algebra~$\cA$ from its
shadow-tower data. Finite-order shadow data suffice exactly on the
finite-depth strata. The extension tower
$\{E_\cA(r)\}_{r \geq 2}$ is the dual inductive system: the
$r$-th stage~$E_\cA(r)$ parametrizes all solutions to
the Maurer--Cartan equation through degree~$r$, and the
extension from~$E_\cA(r)$ to~$E_\cA(r{+}1)$ is
controlled by the obstruction class
$o_{r+1}(\cA) \in H^2(\cA^{\mathrm{sh}}_{r+1,0})$.
There are two depths.  The obstruction depth is finite exactly when
the \(H^2\)-obstruction classes vanish from some degree onward.  The
support depth is finite exactly when the bar-intrinsic MC packet has no
\(H^1\)-lift coordinate above some degree.  Finite reconstruction of
\(\Theta_\cA\) uses the support depth.  Obstruction depth alone only
says that later extensions are unobstructed; it does not choose a point
of the torsor of lifts.

The dichotomy theorem is then instantiated for each shadow
archetype. For each of the four classes (Gaussian~(G),
Lie/tree~(L), contact/quartic~(C), mixed/infinite~(M)) we
give the chain-level computation of the obstruction classes,
with the all-degree non-cancellation inputs cited where they are used.
The Virasoro algebra receives
the most detailed treatment: the jet spaces
$J^r(\mathrm{Vir}_c)$ are computed with explicit basis
elements through $r = 8$, the growth rate is proved to be
polynomial, and the quintic obstruction is exhibited at
the chain level.

\begin{remark}[Scope restriction inherited from Volume~I]
The results of this section build on Volume~I
statements (Theorems \texttt{recursive-existence},
\texttt{nms-all-degree-master-equation},
\texttt{gaussian-rmax-two}, \texttt{lie-rmax-three},
\texttt{contact-rmax-four}, \texttt{virasoro-rmax-infinity},
\texttt{shadow-depth-dichotomy},
\texttt{completed-bar-cobar-strong};
Propositions \texttt{obstruction-extension-sequence},
\texttt{gram-wt4}, \texttt{virasoro-shadow-coefficients},
\texttt{independent-sum-factorization};
Corollaries \texttt{lambda-qp},
\texttt{mittag-leffler-shadow-tower},
\texttt{virasoro-quintic-shadow-explicit}) whose scope qualifiers
propagate to the present proofs. Load-bearing cross-volume citations
in proof bodies below carry the explicit ``Volume~I'' prose prefix;
the semantic content (shadow-tower master equation, obstruction
extension sequences, quintic and higher coefficient data) is inherited
without restatement.
\end{remark}

Four disciplines converge:
\emph{deformation theory} (formal moduli, obstruction classes, Postnikov towers),
\emph{representation theory} (Verma modules, singular vectors, Shapovalov determinants controlling jet-space dimensions),
\emph{combinatorial algebra} (graph-sum contributions to each obstruction class), and
\emph{information theory} (the growth rate of $\dim J^r(\cA)$, measuring reconstruction complexity).

% ================================================================
% 9.8.1: THE EXTENSION TOWER
% ================================================================

\subsection{The extension tower and formal moduli}
% label removed: subsec:extension-tower
\index{extension tower!formal moduli}

\begin{definition}[Extension tower]
% label removed: def:extension-tower
\index{extension tower|textbf}
Let $\cA$ be a modular Koszul chiral algebra
(Definition~\ref{V1-def:modular-koszul-chiral}) with
modular convolution dg~Lie algebra $\gAmod$.
The weight filtration $F^r \gAmod$ gives a tower
of quotients
\[
\gAmod / F^{r+1} \gAmod
\twoheadrightarrow
\gAmod / F^r \gAmod,
\qquad r \geq 2.
\]
The \emph{$r$-th extension space} is the Maurer--Cartan
moduli at level~$r$:
\begin{equation}% label removed: eq:extension-space-def
E_\cA(r) := \MC(\gAmod / F^{r+1} \gAmod)
/ {\sim},
\end{equation}
where two MC elements
$\theta_1, \theta_2 \in \MC(\gAmod / F^{r+1})$ are
\emph{gauge-equivalent} ($\theta_1 \sim \theta_2$)
if there exists $\xi \in (\gAmod / F^{r+1})^0$ such that
\begin{equation}% label removed: eq:gauge-equivalence-def
\theta_2
= e^{\ad(\xi)} \cdot \theta_1
+ \frac{e^{\ad(\xi)} - 1}{\ad(\xi)} \cdot d(\xi).
\end{equation}
The moduli space $E_\cA(r)$ is the quotient of
$\MC(\gAmod / F^{r+1})$ by this equivalence relation.
The \emph{extension tower} is the projective system
\begin{equation}% label removed: eq:extension-tower
\cdots \longrightarrow
E_\cA(r{+}1) \xrightarrow{\;\pi_{r+1,r}\;}
E_\cA(r) \longrightarrow \cdots
\longrightarrow E_\cA(2),
\end{equation}
where $\pi_{r+1,r}$ is induced by the quotient map
$\gAmod / F^{r+2} \to \gAmod / F^{r+1}$.
\end{definition}

\begin{remark}[Formal moduli interpretation]
% label removed: rem:extension-formal-moduli
\index{formal moduli!extension tower}
Each $E_\cA(r)$ is a formal moduli problem in the sense
of Lurie~\cite{Lurie-DAG-X}:
the tangent complex at the basepoint
$\Theta_\cA^{\leq r}$ is the truncated chiral Hochschild
complex $\ChirHoch^*(\cA; F^r/F^{r+1})$, and the
obstruction theory is controlled by
$H^2(\cA^{\mathrm{sh}}_{r+1,0})$. The projective limit
$\varprojlim_r E_\cA(r)$ is the full formal moduli
of the MC element $\Theta_\cA$; it is nonempty by the
bar-intrinsic construction
(Theorem~\ref*{V1-thm:mc2-bar-intrinsic}).
\end{remark}

\begin{definition}[Shadow jet space]
% label removed: def:shadow-jet-space
\index{shadow jet space|textbf}
\index{jet space!shadow}
The \emph{$r$-th shadow jet space} of $\cA$ is
\begin{equation}% label removed: eq:shadow-jet-space
J^r(\cA) :=
F^r {\gAmod}^1 / F^{r+1} {\gAmod}^1,
\end{equation}
the degree-$r$, cohomological-degree-$1$ graded piece of
the modular convolution algebra. Its
$d_2$-cohomology is the shadow algebra:
\begin{equation}% label removed: eq:shadow-jet-cohomology
\cA^{\mathrm{sh}}_{r,0} = H^1(J^r(\cA),\, d_2),
\qquad
H^2(\cA^{\mathrm{sh}}_{r+1,0})
= H^2(J^{r+1}(\cA),\, d_2),
\end{equation}
where $d_2 = d + [\Theta_\cA^{\leq 2}, -]$ is the
$\kappa$-twisted differential (the linearized MC operator
at the quadratic level).
\end{definition}

\begin{proposition}[Obstruction-extension sequence]
% label removed: prop:obstruction-extension-sequence
\index{obstruction class!extension sequence}
For each $r \geq 2$, there is an exact sequence
\begin{equation}% label removed: eq:obstruction-exact-sequence
H^1(J^{r+1}(\cA), d_2)
\xrightarrow{\;\iota\;}
E_\cA(r{+}1)
\xrightarrow{\;\pi_{r+1,r}\;}
E_\cA(r)
\xrightarrow{\;\delta\;}
H^2(J^{r+1}(\cA), d_2),
\end{equation}
where:
\begin{enumerate}[label=\textup{(\roman*)}]
\item $\delta(\Theta^{\leq r})$ is the obstruction class
 $o_{r+1}(\cA)$: the image under the connecting
 homomorphism of the MC failure at degree~$r{+}1$;
\item $\iota$ is the inclusion of the fiber: the set
 of lifts of a given element of $E_\cA(r)$ to
 $E_\cA(r{+}1)$, modulo gauge, is a torsor under
 $H^1(J^{r+1}(\cA), d_2)$ when the obstruction vanishes;
\item $\pi_{r+1,r}$ is the projection.
\end{enumerate}
\end{proposition}

\begin{proof}
This is the standard obstruction theory for
Maurer--Cartan elements in filtered dg~Lie algebras.
Let $\Theta^{\leq r} \in E_\cA(r)$ be a solution to the
MC equation modulo $F^{r+1}$. The residual
\[
R_{r+1} :=
\bigl(D_\cA \Theta^{\leq r}
+ \tfrac{1}{2}[\Theta^{\leq r}, \Theta^{\leq r}]\bigr)_{r+1}
\]
is a degree-$2$ element of $J^{r+1}(\cA)$.
By the Bianchi identity $D_\cA^2 = 0$ applied to
$R_{r+1}$, one obtains $d_2 R_{r+1} = 0$, so
$R_{r+1}$ is a $d_2$-cocycle. Define
$\delta(\Theta^{\leq r}) := [R_{r+1}] \in
H^2(J^{r+1}(\cA), d_2)$.

If $\delta(\Theta^{\leq r}) = 0$, then $R_{r+1} =
d_2 \Theta_{r+1}$ for some degree-$1$ element
$\Theta_{r+1} \in J^{r+1}(\cA)$, and
$\Theta^{\leq r+1} := \Theta^{\leq r} + \Theta_{r+1}$
solves the MC equation modulo $F^{r+2}$. The set of
such lifts forms a torsor under $H^1(J^{r+1}(\cA), d_2)$:
two lifts $\Theta_{r+1}$ and $\Theta'_{r+1}$ differ by
a $d_2$-cocycle, and their difference is a
$d_2$-coboundary if and only if they are gauge-equivalent.

If $\delta(\Theta^{\leq r}) \neq 0$, no lift exists:
the MC equation has no solution at degree~$r{+}1$
compatible with the given solution at degree~$r$.
But $\Theta_\cA$ exists unconditionally
(Theorem~\ref*{V1-thm:mc2-bar-intrinsic}), so the
obstruction is always resolved, but
each $o_{r+1}$ is a \emph{nontrivial} class that must
be absorbed by a new degree-$(r{+}1)$ correction.
\end{proof}

% ================================================================
% 9.8.2: THE DICHOTOMY THEOREM
% ================================================================

\subsection{The dichotomy theorem}
% label removed: subsec:dichotomy-theorem
\index{dichotomy theorem|textbf}
\index{Mittag--Leffler condition!dichotomy theorem}

\begin{theorem}[Obstruction depth and support depth of the shadow tower;
\(\ClaimStatusConditional\); licensing tags \(\gamma+\epsilon\)]
% label removed: thm:shadow-depth-dichotomy
\index{shadow depth!dichotomy}
Let $\cA$ be a modular Koszul chiral algebra. Let
\begin{equation}% label removed: eq:shadow-depth-definition-recall
d_{\mathrm{obs}}(\cA) :=
\inf\{r \geq 2 : o_{r'+1}(\cA) = 0
\;\text{for all}\; r' \geq r\}
\;\in\; \{2, 3, 4, \ldots\} \cup \{\infty\},
\end{equation}
where the infimum of the empty set is \(\infty\).  In the
bar-intrinsic gauge component define the support depth
\[
r^\Theta_{\max}(\cA):=
\inf\{r\ge2:\Theta_{\cA,n}=0\ \text{for all}\ n>r\},
\]
where \(\Theta_{\cA,n}\) is the degree-\(n\) component of the MC
element, equivalently the corresponding component of the shadow
reconstruction packet.  Then
\[
d_{\mathrm{obs}}(\cA)\le r^\Theta_{\max}(\cA),
\]
and finite reconstruction is governed by \(r^\Theta_{\max}\), not by
\(d_{\mathrm{obs}}\) alone:
\begin{enumerate}[label=\textup{(\roman*)}]
\item \emph{Finite support depth.}
If \(r^\Theta_{\max}<\infty\), then
\[
\Theta_\cA^{\leq r}=\Theta_\cA^{\leq r^\Theta_{\max}}=\Theta_\cA
\qquad (r\ge r^\Theta_{\max}).
\]
The bar-intrinsic branch \(E^\Theta_\cA(r)\) stabilizes for
\(r\ge r^\Theta_{\max}\).

\item \emph{Obstruction depth.}
If \(d_{\mathrm{obs}}<\infty\), then no new \(H^2\)-obstruction occurs
above degree \(d_{\mathrm{obs}}\).  The lifts above that degree may
still vary by the \(H^1(J^n(\cA),d_2)\)-torsors.  Thus
\(d_{\mathrm{obs}}<\infty\) implies finite reconstruction only under
the extra condition \(\Theta_{\cA,n}=0\) for \(n>d_{\mathrm{obs}}\), or
equivalently all higher lift coordinates vanish in the chosen gauge
slice.

\item \emph{Infinite support depth.}
If \(r^\Theta_{\max} = \infty\), then
$\Theta_\cA^{\leq r} \neq \Theta_\cA^{\leq r-1}$
for infinitely many values of~$r$. The extension
branch does not stabilize: the inverse limit
$\Theta_\cA = \varprojlim_r \Theta_\cA^{\leq r}$
is the limit of a genuinely infinite tower.
\end{enumerate}
\end{theorem}

\begin{proof}
The obstruction class \(o_{r+1}\) is the image of a truncation in
\(H^2(J^{r+1}(\cA),d_2)\) under the connecting morphism of the
obstruction-extension sequence.  If \(\Theta_{\cA,n}=0\) for every
\(n>r\), then the MC equation has no nonzero residual component above
degree \(r\), so all later obstruction classes vanish.  This gives
\(d_{\mathrm{obs}}\le r^\Theta_{\max}\).

The finite-support statement is the definition of
\(r^\Theta_{\max}\) together with completeness and separatedness of the
weight filtration: a filtered MC element equals the inverse limit of
its finite projections, and the projections are constant exactly after
the last nonzero homogeneous component.  This proves
\(\Theta_\cA^{\le r}=\Theta_\cA\) for
\(r\ge r^\Theta_{\max}\) and stabilization of the selected
bar-intrinsic branch.

If \(d_{\mathrm{obs}}<\infty\), the obstruction-extension sequence says
only that every later extension is unobstructed.  The fiber of
\(E_\cA(n)\to E_\cA(n-1)\) over a liftable point is a torsor under
\(H^1(J^n(\cA),d_2)\).  A nonzero torsor coordinate is a nonzero
degree-\(n\) component of the compatible MC packet.  Hence obstruction
vanishing does not imply truncation unless the higher lift coordinates
vanish.

Finally, \(r^\Theta_{\max}=\infty\) means that infinitely many
homogeneous components of the bar-intrinsic MC element are nonzero, so
the finite projections are not eventually constant.  The inverse limit
is the completed MC element by the bar-intrinsic construction
(Theorem~\ref{V1-thm:recursive-existence}) and the branch-level
Mittag--Leffler statement above.
\end{proof}

\begin{corollary}[Mittag--Leffler for the bar-intrinsic shadow branch;
\(\ClaimStatusConditional\); licensing tags \(\gamma+\epsilon\)]
% label removed: cor:mittag-leffler-shadow-tower
\index{Mittag--Leffler condition!shadow tower}
Let \(E^\Theta_\cA(r)\subset E_\cA(r)\) be the formal branch containing
the projection \([\Theta_\cA^{\le r}]\) of the bar-intrinsic MC element,
with transition maps induced from \(E_\cA(r+1)\to E_\cA(r)\).  The
inverse system \(\{E^\Theta_\cA(r)\}_{r\ge2}\) satisfies the
Mittag--Leffler condition. In particular:
\begin{enumerate}[label=\textup{(\roman*)}]
\item $\varprojlim^1_r E^\Theta_\cA(r) = 0$;
\item the natural map
 \[
 \MC(\widehat{\gAmod})_{\Theta}\big/\textup{gauge}
 \longrightarrow
 \varprojlim_r E^\Theta_\cA(r)
 \]
 is a bijection onto the compatible bar-intrinsic branch;
\item the gauge class of \(\Theta_\cA\) is uniquely determined by its
 compatible finite-order projections in this branch.
\end{enumerate}
No surjectivity statement is made for arbitrary components of
\(E_\cA(r)\); a different component may carry a non-vanishing
obstruction class at the next stage.
\end{corollary}

\begin{proof}
For each \(r\), the bar-intrinsic projection
\([\Theta_\cA^{\le r+1}]\) maps to
\([\Theta_\cA^{\le r}]\).  Hence the transition maps on the selected
branch \(E^\Theta_\cA(r+1)\to E^\Theta_\cA(r)\) are surjective onto
that branch.  This is the Mittag--Leffler condition for the countable
inverse system \(E^\Theta_\cA(\bullet)\), giving
\(\varprojlim^1 E^\Theta_\cA(r)=0\).

Completeness and separatedness of \(F^\bullet\gAmod\) identify an MC
element in the completed dg Lie algebra with a compatible system of
finite projections.  Gauge compatibility is taken inside the same
bar-intrinsic branch.  Therefore the displayed map is bijective onto
\(\varprojlim_r E^\Theta_\cA(r)\), and \([\Theta_\cA]\) is determined
by its compatible finite projections.  The obstruction-extension
sequence still controls other components, but it does not force those
components to lift.
\end{proof}

% ================================================================
% 9.8.3: SHADOW DEPTH AS A RECONSTRUCTION INVARIANT
% ================================================================

\subsection{Shadow depth as a reconstruction invariant}
% label removed: subsec:shadow-depth-invariant
\index{shadow depth!reconstruction invariant}

The support shadow depth \(r^\Theta_{\max}(\cA)\) measures the
complexity of the holographic reconstruction problem: how many
homogeneous orders of the bar-intrinsic packet are needed to determine
the MC gauge class \([\Theta_\cA]\).  The obstruction depth
\(d_{\mathrm{obs}}(\cA)\) measures only eventual \(H^2\)-vanishing.
The two agree on the canonical standard-family rows below, but they
are distinct invariants in a general filtered MC moduli problem.

\begin{definition}[Support shadow depth]
% label removed: def:shadow-depth-holographic
\index{shadow depth|textbf}
The \emph{shadow depth} of a modular Koszul chiral
algebra~$\cA$ is the support depth \(r^\Theta_{\max}(\cA)\) of the
bar-intrinsic MC packet.  Equivalently, it is the least \(r\) for
which all homogeneous components \(\Theta_{\cA,n}\) vanish for
\(n>r\); if no such \(r\) exists, the depth is \(\infty\).
The four \emph{shadow classes} are:
\begin{center}
\renewcommand{\arraystretch}{1.4}
\begin{tabular}{clcl}
\toprule
Class & Name & $r^\Theta_{\max}$ & Defining property \\
\midrule
\textbf{G} & Gaussian & $2$
 & bar-intrinsic packet has support only in degree \(2\) \\
\textbf{L} & Lie/tree & $3$
 & support ends at the Lie cubic; Jacobi kills the next obstruction \\
\textbf{C} & Contact/quartic & $4$
 & support ends at the contact quartic in the canonical gauge slice \\
\textbf{M} & Mixed/infinite & $\infty$
 & infinitely many packet components occur \\
\bottomrule
\end{tabular}
\end{center}
\end{definition}

\begin{remark}[Shadow depth versus Koszulness]
% label removed: rem:depth-vs-koszulness
\index{shadow depth!versus Koszulness}
Shadow depth does \emph{not} characterize Koszulness.
Classes~$\mathbf{G}$ (Heisenberg, \(r^\Theta_{\max} = 2\)),
$\mathbf{L}$ (affine Kac--Moody, \(r^\Theta_{\max} = 3\)), and
$\mathbf{C}$ ($\beta\gamma$ system, \(r^\Theta_{\max} = 4\)) are
chirally Koszul: bar cohomology is concentrated and
$m_k = 0$ for $k \ge 3$ on cohomology
(Proposition~\ref{V1-prop:standard-examples-modular-koszul}).
Class~$\mathbf{M}$ (Virasoro, \(r^\Theta_{\max} = \infty\)) is
chirally Koszul (PBW spectral sequence concentrates) but has
non-formal Swiss-cheese structure: the fourth-order pole forces
nonvanishing transferred operations $m_k^{\mathrm{SC}} \neq 0$
for all $k \geq 3$. The Koszul dual
$\mathrm{Vir}_c^! = \mathrm{Vir}_{26-c}$ carries
these higher $A_\infty$ operations. Shadow depth classifies the
\emph{complexity of the Lagrangian self-intersection}
within the standard families.
\end{remark}

% ================================================================
% 9.8.4: CLASS G; THE GAUSSIAN ARCHETYPE
% ================================================================

\subsection{Class~G: the Gaussian archetype}
% label removed: subsec:class-g-gaussian
\index{shadow archetype!Gaussian|textbf}
\index{Heisenberg algebra!shadow depth}

\begin{theorem}[$r^\Theta_{\max}(\mathcal{H}_k) = 2$]
% label removed: thm:gaussian-rmax-two
\index{shadow depth!Heisenberg}
For the Heisenberg algebra $\mathcal{H}_k$ at any nonzero
level $k \neq 0$, the support shadow depth is
\(r^\Theta_{\max}=2\).
Every obstruction class vanishes:
\begin{equation}% label removed: eq:gaussian-all-obstructions-vanish
o_{r+1}(\mathcal{H}_k) = 0
\qquad \text{for all } r \geq 2.
\end{equation}
The bar-intrinsic packet stabilizes immediately:
$\Theta_{\mathcal{H}_k}^{\leq 2} = \Theta_{\mathcal{H}_k}
= \kappaChHodge(\mathcal{H}_k)\,\eta \otimes \Lambda$ with
\(\kappaChHodge(\mathcal{H}_k) = k\).  All higher
\(H^1\)-lift coordinates vanish in the canonical Gaussian gauge slice.
\end{theorem}

\begin{proof}
The proof has two steps: show that all higher OPE products
vanish, and conclude that the graph-sum contributions
at degree $r \geq 3$ are identically zero.

\medskip
\emph{Step~1: All $m_n = 0$ for $n \geq 3$.}

The Heisenberg OPE is
\begin{equation}% label removed: eq:heisenberg-ope-recall
a(z)\,b(w) \sim \frac{k\,\langle a, b \rangle}{(z-w)^2}
\end{equation}
for $a, b \in V = \mathfrak{h}$ (the Cartan subalgebra),
with no simple pole (since $[a,b] = 0$) and no
higher-order poles. The vertex algebra product decomposes
as
\[
Y(a, z) = \sum_{n \in \Z} a_{(n)} z^{-n-1},
\]
where $a_{(n)}b = 0$ for $n \geq 1$ (no singular OPE
poles beyond the double pole) and $a_{(0)}b = 0$
(no Lie bracket). In the bar complex language:
\begin{itemize}
\item The bilinear product $m_2$ gives the Hessian
 $H_{\mathcal{H}} = \kappaChHodge(\mathcal{H}_k) = k$ (from the double pole).
\item The $n$-ary operations $m_n = 0$ for all $n \geq 3$.
 This is immediate: $m_n$ is built from
 iterated OPE residues, and each residue involves an
 $a_{(p)}b$ with $p \geq 0$. Since $a_{(0)} = 0$
 (no Lie bracket) and $a_{(1)} = k\langle a, b\rangle$
 (a scalar, not a new element of the algebra), no
 nontrivial chain of compositions can be formed beyond
 degree~$2$.
\end{itemize}

More precisely, the transferred cyclic minimal model
$\{m_n^{\mathrm{tr}}\}_{n \geq 1}$ of the
$A_\infty$-algebra $(\barB(\mathcal{H}_k), d)$ satisfies
$m_n^{\mathrm{tr}} = 0$ for $n \geq 3$. The reason: the
bar complex $\barB(\mathcal{H}_k)$ is the cofree
conilpotent coalgebra on $s^{-1}\bar{V}$ with quadratic
codifferential, and the homological perturbation lemma
produces no corrections beyond the quadratic term when the
perturbation itself is quadratic.

\medskip
\emph{Step~2: Graph-sum vanishing.}

We first record a combinatorial fact.

\emph{Lemma.} Every connected stable tree with
$n \geq 3$ leaves has at least one vertex of
valence $\geq 3$.

\emph{Proof of Lemma.} A tree with $n$ leaves and all
internal vertices of valence $\leq 2$ is a path, which
has exactly $2$ leaves. For $n \geq 3$, some vertex
must have valence $\geq 3$. \hfill$\square$

\medskip
The degree-$r$ component of $\Theta_\cA^{\leq r}$ is
a graph sum
\[
\Theta_r = \sum_{\Gamma \in \mathcal{G}_r}
\frac{1}{|\mathrm{Aut}(\Gamma)|}
\ell_\Gamma(\Theta),
\]
where $\mathcal{G}_r$ is the set of connected stable
graphs of weight~$r$, and $\ell_\Gamma$ is the graph
amplitude with vertex labels from $\{m_n^{\mathrm{tr}}\}$
and edge propagators $P = H^{-1} = 1/k$. Since
$m_n^{\mathrm{tr}} = 0$ for $n \geq 3$, every graph with
a vertex of valence $\geq 3$ contributes zero. The only
graphs of weight~$r \geq 3$ that could contribute would
need to use only bivalent vertices (from $m_2$), but by
the Lemma, a connected tree on $r \geq 3$ leaves must
have at least one vertex of valence $\geq 3$.
A connected graph on bivalent vertices alone is either a
single edge (weight~$2$) or a cycle. Cycles of length
$\geq 2$ have weight $\geq 3$ but contribute to the
genus expansion, not to the degree expansion at genus~$0$.
At genus~$0$, the tree graphs of weight~$r \geq 3$ must
have at least one vertex of valence~$\geq 3$, so all
contribute zero.

Hence $\Theta_r = 0$ for all $r \geq 3$, and
\begin{equation}% label removed: eq:heisenberg-theta-final
\Theta_{\mathcal{H}_k} = \Theta^{\leq 2}_{\mathcal{H}_k}
= k\,\eta \otimes \Lambda.
\end{equation}
All obstruction classes vanish:
$o_{r+1} = (D_\cA \Theta^{\leq r} +
\frac{1}{2}[\Theta^{\leq r}, \Theta^{\leq r}])_{r+1}
= 0$ because the MC equation is already satisfied
exactly at degree~$2$.
\end{proof}

\begin{remark}[The Gaussian class in generality]
% label removed: rem:gaussian-class-general
\index{Gaussian archetype!general members}
The argument of Theorem~\ref{V1-thm:gaussian-rmax-two}
applies to any chiral algebra with purely quadratic OPE:
the free fermion, the $bc$ ghost system, and more generally
any lattice vertex algebra $V_\Lambda$
(Theorem~\ref{V1-thm:lattice:curvature-braiding-orthogonal}).
In each case, the transferred $A_\infty$-structure is
formal ($m_n^{\mathrm{tr}} = 0$ for $n \geq 3$),
the modular $L_\infty$-algebra $\Definfmod(\cA)$ equals
its strict model, and the shadow obstruction tower terminates at
$r = 2$. The Gaussian class is characterized by:
\begin{equation}% label removed: eq:gaussian-characterization
\cA \in \textbf{G}
\quad\Longleftrightarrow\quad
\text{all OPE poles have order } \leq 2
\quad\Longleftrightarrow\quad
m_n^{\mathrm{tr}} = 0 \;\text{for } n \geq 3.
\end{equation}
\end{remark}

% ================================================================
% 9.8.5: CLASS L; THE LIE/TREE ARCHETYPE
% ================================================================

\subsection{Class~L: the Lie/tree archetype}
% label removed: subsec:class-l-lie-tree
\index{shadow archetype!Lie/tree|textbf}
\index{affine Kac--Moody!shadow depth}

\begin{theorem}[$r^\Theta_{\max}(\widehat{\mathfrak{g}}_k) = 3$]
% label removed: thm:lie-rmax-three
\index{shadow depth!affine Kac--Moody}
For the affine Kac--Moody algebra
$\widehat{\mathfrak{g}}_k$ at generic level $k$, the
support shadow depth is \(r^\Theta_{\max}=3\). The obstruction
classes satisfy:
\begin{equation}% label removed: eq:lie-obstruction-pattern
o_3(\widehat{\mathfrak{g}}_k) \neq 0,
\qquad
o_{r+1}(\widehat{\mathfrak{g}}_k) = 0
\;\;\text{for all } r \geq 3.
\end{equation}
In the canonical Lie/tree gauge slice, the higher
\(H^1(J^n(\widehat{\mathfrak g}_k),d_2)\)-lift coordinates vanish for
\(n>3\).  Thus the support packet, not merely the \(H^2\)-obstruction
profile, terminates at degree~\(3\).
\end{theorem}

\begin{proof}
The proof proceeds by computing the
transferred $A_\infty$-operations and verifying that the
Jacobi identity forces the quartic obstruction to vanish.

\medskip
\emph{Step~1: The OPE structure.}

The current algebra $\widehat{\mathfrak{g}}_k$ has OPE
\begin{equation}% label removed: eq:km-ope-recall
J^a(z)\,J^b(w)
\sim \frac{k\,\kappa^{ab}}{(z-w)^2}
+ \frac{f^{ab}{}_c\,J^c(w)}{z-w},
\end{equation}
where $\kappa^{ab}$ is the Killing form and $f^{ab}{}_c$
are the structure constants of $\mathfrak{g}$. The double
pole contributes the Hessian $H_{\widehat{\fg}} = k\kappa/2$.
The simple pole contributes the cubic operation
$m_2^{(0)}(J^a, J^b) = f^{ab}{}_c\,J^c$, which is
precisely the Lie bracket.

\medskip
\emph{Step~2: Nontrivial cubic shadow.}

The transferred cubic operation
$m_2^{\mathrm{tr}} \neq 0$ comes from the simple-pole
OPE. The cubic shadow is
\begin{equation}% label removed: eq:km-cubic-shadow
\mathfrak{C}_{\widehat{\fg}} =
f^{ab}{}_c\,x_a x_b x^c,
\end{equation}
the Lie bracket cubic. The obstruction class
$o_3 = [D_\cA \Theta^{\leq 2}]_3 \neq 0$ is nontrivial
because the Lie bracket is nontrivial ($\fg$ is non-abelian).
Hence $\Theta^{\leq 3} \neq \Theta^{\leq 2}$: the shadow
tower requires a cubic correction.

\medskip
\emph{Step~3: Quartic obstruction vanishes by Jacobi.}

The quartic obstruction is
\begin{equation}% label removed: eq:km-quartic-obstruction
o_4(\widehat{\fg}_k) =
\bigl[D_\cA \Theta^{\leq 3}
+ \tfrac{1}{2}[\Theta^{\leq 3}, \Theta^{\leq 3}]
\bigr]_4.
\end{equation}
This has two contributions:
\begin{enumerate}[label=(\alph*)]
\item the \emph{boundary bracket}
 $\frac{1}{2}\{\mathfrak{C}, \mathfrak{C}\}_H$:
 the sewing contraction of two cubic vertices through
 the propagator $P = (k\kappa)^{-1}$;
\item the \emph{local quartic vertex}
 $m_3^{\mathrm{tr}}$: the transferred ternary operation.
\end{enumerate}

We claim both contribute zero in cohomology.

\emph{Contribution~(a): the boundary bracket.}
The $H$-Poisson bracket of the Lie bracket cubic with
itself is
\begin{equation}% label removed: eq:km-boundary-bracket
\frac{1}{2}\{\mathfrak{C}, \mathfrak{C}\}_H
= \frac{1}{2} \cdot
\frac{\partial \mathfrak{C}}{\partial x^a}
\cdot P^{ab} \cdot
\frac{\partial \mathfrak{C}}{\partial x_b}.
\end{equation}
Writing out explicitly:
\[
\frac{\partial \mathfrak{C}}{\partial x^a}
= f^{bc}{}_a x_b x_c + 2f^{ab}{}_c x_b x^c.
\]
The full contraction produces a quartic polynomial in
the $x$-variables. The coefficient of each monomial
$x_a x_b x_c x_d$ involves a sum
\[
\sum_e P^{ae}f^{bc}{}_e + \text{permutations},
\]
which is the structure constant of the Lie bracket
composed with itself through the inverse Killing form.
The \emph{Jacobi identity}
\begin{equation}% label removed: eq:jacobi-identity-structure-constants
f^{ab}{}_e f^{ec}{}_d + f^{bc}{}_e f^{ea}{}_d
+ f^{ca}{}_e f^{eb}{}_d = 0
\end{equation}
forces the quartic polynomial
$\frac{1}{2}\{\mathfrak{C}, \mathfrak{C}\}_H$ to
lie in the image of $d_2$: it is a $d_2$-coboundary.

\emph{Contribution~(b): the local quartic vertex.}
The transferred ternary operation $m_3^{\mathrm{tr}}$ on
bar cohomology is computed by the homological perturbation
lemma:
\[
m_3^{\mathrm{tr}}(a,b,c) =
\pi \circ m_2(h \circ m_2(i(a), i(b)),\, i(c))
+ \pi \circ m_2(i(a),\, h \circ m_2(i(b), i(c)))
+ \cdots,
\]
where $\pi$ is the projection to bar cohomology, $i$ is
the inclusion, and $h$ is the contracting homotopy.
For the Kac--Moody algebra, the inner $m_2$ produces
either a Lie bracket (from the $(0)$-product
$J^a_{(0)}J^b = f^{ab}{}_c J^c$) or a scalar (from the
$(1)$-product $J^a_{(1)}J^b = k\kappa^{ab}\mathbf{1}$).
When it produces a scalar, the outer $m_2$ acts on
$(\text{element}) \otimes \mathbf{1}$, which gives zero
in the reduced bar complex. When it produces a Lie
bracket, the homotopy $h$ maps this back to a bar element,
and the outer $m_2$ acts.

For affine algebras, the Jacobi identity $\ell_3 = 0$
implies that $m_3^{\mathrm{tr}}$ is a $d_2$-coboundary:
$m_3^{\mathrm{tr}} = d_2(\text{something})$. The key
point is that $m_3^{\mathrm{tr}}$ and the $L_\infty$
bracket $\ell_3$ are \emph{distinct} operations: the
former lives on bar cohomology while the latter is the
homotopy-transfer bracket. They are related by the
Koszul sign rule, and the Jacobi identity
\begin{equation}% label removed: eq:ell3-jacobi-vanishing
\ell_3(a, b, c) = [a, [b, c]] - [[a,b], c]
- (-1)^{|a||b|}[b,[a,c]] = 0
\end{equation}
forces $\ell_3 = 0$, which implies $[m_3^{\mathrm{tr}}]
= 0$ in $H^2$. Hence $o_4 = [m_3^{\mathrm{tr}}] = 0$
in $H^2(J^4(\widehat{\fg}_k), d_2)$.

\emph{Combining:}
The total quartic obstruction is
$o_4 = [\frac{1}{2}\{\mathfrak{C}, \mathfrak{C}\}_H
+ m_3^{\mathrm{tr}}] = 0$
in $H^2(J^4(\widehat{\fg}_k), d_2)$, by the Jacobi identity.

\medskip
\emph{Step~4: All higher obstructions vanish.}

By induction: the only source of higher obstructions is
the graph-sum composition of lower-degree vertices. Since
the only nontrivial vertex is the cubic
$m_2^{(0)} = [\;,\;]$ (the Lie bracket), every graph of
degree $r \geq 4$ is a tree with only trivalent vertices.
The amplitude of such a tree involves iterated Lie
brackets, contracted through the propagator. The Jacobi
identity, applied at each internal vertex, shows that
the amplitude of any such tree is a $d_2$-coboundary:
it is the $d_2$-image of a lower-weight correction.

Concretely: at degree~$r$, the obstruction is a sum over
trees with $r-1$ leaves and internal edges labeled by
the propagator $P$. Each internal vertex is a Lie
bracket. The Jacobi identity gives a recursive formula
\begin{equation}% label removed: eq:jacobi-recursive-vanishing
o_{r+1} = d_2\Bigl(
\sum_{\text{trees } T \in \mathcal{T}_{r}}
\frac{1}{|\Aut(T)|}\,\ell_T
\Bigr) = 0,
\end{equation}
where $\mathcal{T}_r$ is the set of binary trees with
$r$ leaves. The cohomology class vanishes because the
Jacobi identity is exactly the $A_\infty$ formality
condition for the Lie algebra (the Koszul resolution of
$\fg$ is formal).

Therefore $o_{r+1} = 0$ for all $r \geq 3$, and
the canonical Lie/tree support packet has
\(r^\Theta_{\max}(\widehat{\fg}_k) = 3\).
\end{proof}

\begin{remark}[The Jacobi identity as a termination criterion]
% label removed: rem:jacobi-termination
\index{Jacobi identity!tower termination}
The vanishing $o_4 = 0$ is not accidental but reflects
a deep structural property: the Lie bracket $[\;,\;]$
satisfies an associativity condition (the Jacobi identity)
that forces all compositions of Lie brackets to be
algebraically determined by the bracket itself. In the
language of operads, $\Lie$ is a Koszul operad, and
the bar resolution of any Lie algebra is formal. The
shadow obstruction tower terminates at degree~$3$ because there are
no higher obstructions to formality.
\end{remark}

% ================================================================
% 9.8.6: CLASS C; THE CONTACT/QUARTIC ARCHETYPE
% ================================================================

\subsection{Class~C: the contact/quartic archetype}
% label removed: subsec:class-c-contact
\index{shadow archetype!contact/quartic|textbf}
\index{beta-gamma system@$\beta\gamma$ system!shadow depth}

\begin{theorem}[$r^\Theta_{\max}(\beta\gamma) = 4$]
% label removed: thm:contact-rmax-four
\index{shadow depth!beta-gamma@$\beta\gamma$}
For the $\beta\gamma$ system, the support shadow depth is
\(r^\Theta_{\max}=4\). The obstruction classes satisfy:
\begin{equation}% label removed: eq:contact-obstruction-pattern
o_3(\beta\gamma) \neq 0,
\qquad
o_4(\beta\gamma) \neq 0,
\qquad
o_{r+1}(\beta\gamma) = 0
\;\;\text{for all } r \geq 4.
\end{equation}
In the canonical contact gauge slice, all higher lift coordinates
\(\lambda_n\in H^1(J^n(\beta\gamma),d_2)\) vanish for \(n>4\).
\end{theorem}

\begin{proof}
The proof uses the cubic gauge triviality theorem
(Theorem~\ref{thm:cubic-gauge-triviality}) and
rank-one abelian rigidity.

\medskip
\emph{Step~1: The OPE structure and transferred operations.}

The $\beta\gamma$ system has OPE
\begin{equation}% label removed: eq:betagamma-ope-recall
\beta(z)\,\gamma(w) \sim \frac{1}{z-w}.
\end{equation}
This is a simple-pole OPE with no double pole, so
the curvature $m_0 = 0$ and the Hessian of the
modular deformation complex is determined by the normal-ordered
products. The stress tensor
$T = {:}\partial\beta \cdot \gamma{:}$ gives the
conformal structure.

The key structural feature: the $\beta\gamma$ system
has rank one (a single pair of conjugate fields), so
the shadow algebra is one-dimensional at each degree.
The transferred operations are:
\begin{align*}
m_2^{\mathrm{tr}} &\neq 0
 & &\text{(cubic, from simple-pole OPE)}, \\
m_3^{\mathrm{tr}} &\neq 0
 & &\text{(quartic, from normal-ordered composites)}, \\
m_4^{\mathrm{tr}} &= 0
 & &\text{(quintic, vanishes by rank-one rigidity)}.
\end{align*}

\medskip
\emph{Step~2: Cubic gauge triviality.}

The cubic shadow at degree~$3$ is nonzero, but the
$\beta\gamma$ system satisfies the hypothesis of
Theorem~\ref{thm:cubic-gauge-triviality}:
\begin{equation}% label removed: eq:betagamma-h1-vanishing
H^1(F^3 \gAmod / F^4 \gAmod,\, d_2) = 0.
\end{equation}
This vanishing holds because the weight-$3$ piece of the
convolution algebra, for a rank-one system, is
one-dimensional and the $d_2$-action on it is injective
(the $\kappa$-twisted differential has no kernel at
degree~$1$ in weight~$3$). By
Theorem~\ref{thm:cubic-gauge-triviality}:
\begin{enumerate}[label=(\alph*)]
\item The cubic term $\Theta_3$ is gauge-trivial:
 there exists a gauge transformation removing it.
\item After this gauge choice, the quartic term
 defines a \emph{canonical class}
 $[\Theta'_4] \in H^2(J^4(\beta\gamma), d_2)$.
\end{enumerate}
This canonical quartic class is the \emph{quartic
contact shadow}
$\mathfrak{Q}^{\mathrm{contact}}_{\beta\gamma}$.

\medskip
\emph{Step~3: Quartic terminality.}

The quartic contact shadow is nontrivial:
$\mathfrak{Q}^{\mathrm{contact}}_{\beta\gamma} \neq 0$
(Corollary~\ref*{V1-cor:nms-betagamma-mu-vanishing}).
The key computation is:
\begin{equation}% label removed: eq:betagamma-mu-zero
\mu_{\beta\gamma}
:= \langle \eta, m_3(\eta, \eta, \eta) \rangle = 0,
\end{equation}
where $\eta$ is the generator of the shadow algebra.
Proof: this vanishing by a parity argument.
On the rank-one primary line spanned by~$\eta$,
$m_3(\eta,\eta,\eta)$ must be proportional to~$\eta$
by conformal weight and parity constraints. Write
$m_3(\eta,\eta,\eta) = \lambda\,\eta$. Then
$\mu = \langle \eta,\, \lambda\,\eta \rangle
= \lambda\,\langle \eta, \eta \rangle$.
But $m_3$ is graded-symmetric in its inputs (cubic
operation on identical generators), while the cyclic
pairing is graded-antisymmetric for odd-degree elements.
Since $|\eta|$ is odd (bar-desuspended),
$\langle \eta,\, m_3(\eta,\eta,\eta) \rangle
= -\langle m_3(\eta,\eta,\eta),\, \eta \rangle$
by antisymmetry, hence $\mu = -\mu$, so $\mu = 0$.

\medskip
\emph{Step~4: Quintic obstruction vanishes.}

The quintic obstruction is
\begin{equation}% label removed: eq:betagamma-quintic-obstruction
o_5(\beta\gamma) =
\bigl[\{\mathfrak{C}_{\beta\gamma},
\mathfrak{Q}^{\mathrm{contact}}_{\beta\gamma}\}_{H}
+ m_4^{\mathrm{tr}}\bigr].
\end{equation}
Verify both contributions vanish.

\emph{The boundary bracket.}
After the cubic gauge triviality (Step~2), the
gauge-fixed MC element has $\Theta_3' = 0$, so
$\mathfrak{C}_{\beta\gamma} = 0$ in the chosen gauge.
The boundary bracket $\{0, \mathfrak{Q}\}_H = 0$.

\emph{The local quintic vertex.}
The transferred quartic operation $m_4^{\mathrm{tr}} = 0$
by rank-one rigidity: on a one-dimensional primary space,
the cyclic $4$-ary operation is forced to vanish by the
same argument as in~\eqref{V1-eq:betagamma-mu-zero},
extended to degree~$4$.

Therefore $o_5 = 0$. By the same argument, $o_{r+1} = 0$
for all $r \geq 4$: the only source of higher obstructions
would be graph-sum compositions, but with $\Theta_3' = 0$
(in the chosen gauge) and $m_n^{\mathrm{tr}} = 0$ for
$n \geq 4$, no nontrivial graphs exist. Hence
\(r^\Theta_{\max}(\beta\gamma) = 4\).
\end{proof}

\begin{remark}[The $\beta\gamma$ system as the contact atom]
% label removed: rem:betagamma-contact-atom
\index{beta-gamma system@$\beta\gamma$ system!as contact atom}
The $\beta\gamma$ system is the unique standard atom
whose shadow obstruction tower is nontrivial beyond the scalar
level yet terminates at finite degree. The quartic
contact shadow is the simplest nonlinear modular
invariant beyond $\kappa$: it measures the
self-interaction of the normal-ordered composite
${:}\beta\gamma{:}$ with the stress tensor.
\end{remark}

% ================================================================
% 9.8.7: CLASS M; THE MIXED/INFINITE ARCHETYPE
% ================================================================

\subsection{Class~M: the mixed/infinite archetype}
% label removed: subsec:class-m-mixed
\index{shadow archetype!mixed/infinite|textbf}
\index{Virasoro algebra!shadow depth}

This subsection records the local computation behind the class
$\mathsf M$ example. The infinite-depth statement is the
all-degree Virasoro theorem of Volume~I; the calculation here exhibits
the quintic obstruction and the persistent cubic-source branch. The
argument combines the explicit
Virasoro OPE with the Gram matrix computation
(Proposition~\ref{V1-prop:gram-wt4}) and the graph-sum
formula for the obstruction.

\begin{theorem}[$r^\Theta_{\max}(\mathrm{Vir}_c) = \infty$]
% label removed: thm:virasoro-rmax-infinity
\index{shadow depth!Virasoro}
For the Virasoro algebra $\mathrm{Vir}_c$ at generic
central charge $c \neq 0, -22/5$, the support shadow depth is
\(r^\Theta_{\max} = \infty\). The quintic obstruction is
\begin{equation}% label removed: eq:virasoro-quintic-nonzero
o_5(\mathrm{Vir}_c) =
\frac{480}{c^2(5c+22)}\,x^5 \neq 0,
\end{equation}
and by induction all higher obstruction classes are
nonzero:
\begin{equation}% label removed: eq:virasoro-all-obstructions-nonzero
o_{r+1}(\mathrm{Vir}_c) \neq 0
\qquad \text{for all } r \geq 2.
\end{equation}
\end{theorem}

The proof requires several intermediate results, which
we develop in order.

\subsubsection{The Virasoro shadow obstruction tower through degree~$4$}
% label removed: subsubsec:virasoro-tower-through-4

\begin{computation}[Virasoro tower data]
% label removed: comp:virasoro-tower-data
\index{Virasoro algebra!tower data}
The Virasoro algebra $\mathrm{Vir}_c$ has a single strong
generator~$T$ of conformal weight~$2$. The OPE is
\begin{equation}% label removed: eq:virasoro-ope-full
T(z)\,T(w) \sim
\frac{c/2}{(z-w)^4}
+ \frac{2T(w)}{(z-w)^2}
+ \frac{\partial T(w)}{z-w}.
\end{equation}
The shadow obstruction tower data through degree~$4$, on the
one-dimensional primary space spanned by~$x$
(dual to~$T$):

\begin{center}
\renewcommand{\arraystretch}{1.5}
\begin{tabular}{cll}
\toprule
Degree $r$ & Shadow component & Formula \\
\midrule
$2$ & Hessian $H = \kappa$ & $c/2$ \\
$2$ & Propagator $P = H^{-1}$ & $2/c$ \\
$3$ & Cubic shadow $\mathfrak{C}$ & $2x^3$ \\
$3$ & Source: $T_{(1)}T = 2T$ & (simple-pole OPE) \\
$4$ & Quartic contact $\mathfrak{Q}^{\mathrm{ct}}$ &
 $\displaystyle\frac{10}{c(5c+22)}\,x^4$ \\
$4$ & Source: $\Lambda = {:}TT{:}
 - \tfrac{3}{10}\partial^2 T$ & (quasi-primary composite)\\
\bottomrule
\end{tabular}
\end{center}

\noindent The quartic contact coefficient
$Q_0 = 10/[c(5c+22)]$ is the inverse of the BPZ norm of
the quasi-primary composite
$\Lambda = {:}TT{:} - \frac{3}{10}\partial^2 T$
(Corollary~\ref{V1-cor:lambda-qp}).
\end{computation}

\begin{proof}[Derivation of the quartic contact]
The quartic contact coefficient is determined by the
condition that $\Theta^{\leq 4}$ satisfies the MC equation
through degree~$4$. The quartic MC equation requires
absorbing the boundary bracket and local vertex.

\emph{The boundary bracket.}
With $\mathfrak{C} = 2x^3$ and $P = 2/c$:
\begin{align}
\tfrac{1}{2}\{\mathfrak{C}, \mathfrak{C}\}_H
&= \tfrac{1}{2} \cdot
 \frac{\partial(2x^3)}{\partial x} \cdot P \cdot
 \frac{\partial(2x^3)}{\partial x} \notag \\
&= \tfrac{1}{2} \cdot 6x^2 \cdot \frac{2}{c}
 \cdot 6x^2 \notag \\
&= \frac{36}{c}\,x^4. % label removed: eq:virasoro-boundary-bracket
\end{align}

\emph{The local quartic vertex.}
The quasi-primary composite
$\Lambda = {:}TT{:} - \frac{3}{10}\partial^2 T$
has BPZ norm
$\langle \Lambda | \Lambda \rangle = c(5c+22)/10$
(Corollary~\ref{V1-cor:lambda-qp}).
The local contribution is proportional to $1/\langle\Lambda|\Lambda\rangle$:
\begin{equation}% label removed: eq:virasoro-local-quartic
\mathfrak{Q}^{\mathrm{local}}
= \frac{4}{\langle\Lambda|\Lambda\rangle}\,x^4
= \frac{40}{c(5c+22)}\,x^4.
\end{equation}

\emph{The total quartic.}
By the all-degree master equation
$\nabla_H(\mathrm{Sh}_4) + \mathfrak{o}^{(4)} = 0$
(Volume~I Theorem~\ref{V1-thm:nms-all-degree-master-equation}),
the quartic shadow absorbs the total obstruction, giving:
\begin{equation}% label removed: eq:virasoro-contact-result
Q^{\mathrm{contact}}_{\mathrm{Vir}}
= \frac{10}{c(5c+22)}.
\end{equation}
\end{proof}

\subsubsection{The quintic obstruction at the chain level}
% label removed: subsubsec:quintic-chain-level

\begin{proof}[Local computation for Volume~I Theorem~\textup{\ref{V1-thm:virasoro-rmax-infinity}}]
The computation gives the quintic obstruction $o_5$ as a nonzero
class in $H^2(J^5(\mathrm{Vir}_c), d_2)$.

\medskip
\emph{Step~1: The quintic MC equation.}

The projection of the MC equation
$D_\cA \Theta + \frac{1}{2}[\Theta, \Theta] = 0$
onto degree~$5$ gives
\begin{equation}% label removed: eq:quintic-mc-equation
d_2 \Theta_5 +
\{\mathfrak{C}, \mathfrak{Q}^{\mathrm{ct}}\}_H
+ \text{(higher-order contributions)} = 0.
\end{equation}
The leading obstruction is the cubic-quartic bracket.

\medskip
\emph{Step~2: The cubic-quartic bracket.}

\begin{align}
\{\mathfrak{C}, \mathfrak{Q}^{\mathrm{ct}}\}_H
&= \frac{\partial(2x^3)}{\partial x}
 \cdot P \cdot
 \frac{\partial(Q_0\,x^4)}{\partial x} \notag \\
&= 6x^2 \cdot \frac{2}{c} \cdot 4Q_0\,x^3 \notag \\
&= \frac{48Q_0}{c}\,x^5 \notag \\
&= \frac{48}{c} \cdot \frac{10}{c(5c+22)}\,x^5 \notag \\
&= \frac{480}{c^2(5c+22)}\,x^5.
 % label removed: eq:quintic-cubic-quartic
\end{align}

\medskip
\emph{Step~3: Nonvanishing.}

The coefficient $480/[c^2(5c+22)]$ is nonzero for
generic $c$ (failing only at $c = 0$ and $c = -22/5$,
where the Virasoro algebra degenerates). We must verify
that $\{\mathfrak{C}, \mathfrak{Q}^{\mathrm{ct}}\}_H$
is \emph{not} a $d_2$-coboundary. The $H$-Poisson
bracket $[\mathfrak{C}, \mathrm{Sh}_r]_H$ lies in
$\cA^{\mathrm{sh}}_{r+1,0}$. If it were exact, there
would exist $\xi \in \cA^{\mathrm{sh}}_{r+1,-1}$ with
$d_2(\xi) = [\mathfrak{C}, \mathrm{Sh}_r]_H$. But
$\dim \cA^{\mathrm{sh}}_{r+1,-1} = 0$ for $r \geq 4$:
the weight-$({-}1)$ space is trivial in the cyclic
deformation complex at high degree (there are no degree-$1$
cochains in $J^{r+1}$ for $r \geq 4$ on the one-generator
primary line, since a single generator of conformal
weight~$2$ produces only even-weight monomials in the
shadow algebra). Hence the bracket
$[\mathfrak{C}, \mathrm{Sh}_r]_H$ is automatically a
nontrivial cohomology class.

Therefore
\begin{equation}% label removed: eq:quintic-obstruction-clean
o_5(\mathrm{Vir}_c) =
\frac{480}{c^2(5c+22)}\,x^5 \neq 0.
\end{equation}

\emph{Step~4: The quintic shadow.}
The all-degree master equation
$\nabla_H(\mathrm{Sh}_5) + o^{(5)} = 0$
(Volume~I Theorem~\ref{V1-thm:nms-all-degree-master-equation})
determines $\mathrm{Sh}_5$. On the one-generator primary
line with Hessian $H = c/2$, the covariant derivative
acts as $\nabla_H(S_r x^r) = 2rS_r x^r$. Therefore
$\mathrm{Sh}_5 = -o^{(5)}/(2 \cdot 5)$:
\begin{equation}% label removed: eq:sh5-virasoro-formula
\mathrm{Sh}_5(\mathrm{Vir}_c)
= -\frac{480}{c^2(5c+22)} \cdot \frac{1}{10}\,x^5
= -\frac{48}{c^2(5c+22)}\,x^5,
\end{equation}
in exact agreement with Volume~I
Corollary~\ref{V1-cor:virasoro-quintic-shadow-explicit}.

\medskip
\emph{Step~5: Cubic-source branch and the all-degree input.}

The quintic shadow
$\mathrm{Sh}_5 = -48/(c^2(5c+22))\,x^5 \neq 0$
acts as a new vertex in the graph sum at degree~$6$.
The sextic obstruction includes the cubic-quintic bracket:
\begin{align}
\{\mathfrak{C}, \mathrm{Sh}_5\}_H
&= 6x^2 \cdot \frac{2}{c} \cdot
 5\cdot\bigl(-\frac{48}{c^2(5c+22)}\bigr)\,x^4
 \notag \\
&= -\frac{2880}{c^3(5c+22)}\,x^6
\neq 0. % label removed: eq:sextic-leading-term
\end{align}
By induction in the cubic-source approximation: at each degree~$r$, the
cubic vertex
$\mathfrak{C} = 2x^3$ contracts with the $(r{-}1)$-st
shadow contribution to produce a nonzero $r$-th source:
\begin{equation}% label removed: eq:inductive-obstruction
o_{r+1}(\mathrm{Vir}_c) \ni
\{\mathfrak{C}, \mathrm{Sh}_r\}_H
= 6x^2 \cdot \frac{2}{c} \cdot r\,S_r\,x^{r-1}
= \frac{12r\,S_r}{c}\,x^{r+1},
\end{equation}
where $S_r$ denotes the coefficient of this cubic-source contribution.
This source is nonzero whenever $S_r \neq 0$. Since $S_5 \neq 0$, the
cubic-source branch starts, and its recursion gives
\begin{equation}% label removed: eq:shadow-recursive-formula
S_{r+1}^{\mathrm{cub}} =
-\frac{6r}{c(r+1)}\,S_r^{\mathrm{cub}},
\qquad r \geq 5.
\end{equation}
Since $6r/(c(r+1)) \neq 0$ for generic $c$ and all
$r \geq 5$, the cubic-source contribution is nonzero at all higher
degrees. The passage from this persistent source to exact nonvanishing
of every obstruction class uses the all-degree non-cancellation theorem
of Volume~I, Theorem~\ref{V1-thm:virasoro-rmax-infinity}. Therefore
\(r^\Theta_{\max}(\mathrm{Vir}_c)=\infty\).
\end{proof}

\begin{remark}[The cubic vertex as permanent source]
% label removed: rem:cubic-permanent-source
\index{Virasoro algebra!cubic as permanent source}
The essential mechanism is that the cubic vertex
$\mathfrak{C} = 2x^3$ acts as a \emph{permanent source}
in the shadow obstruction tower: its bracket with any nonzero shadow
produces a nonzero obstruction at the next level. This
is specific to the mixed archetype, where both
$\mathfrak{C} \neq 0$ and $\mathfrak{Q}^{\mathrm{ct}}
\neq 0$ coexist. In the Lie/tree archetype (Kac--Moody),
$\mathfrak{C} \neq 0$ but $\mathfrak{Q}^{\mathrm{ct}} = 0$,
so the quartic bracket $\{\mathfrak{C}, \mathfrak{C}\}_H$
is $d_2$-exact (by Jacobi) and no quintic is generated.
In the contact archetype ($\beta\gamma$),
$\mathfrak{C} = 0$ in the appropriate gauge, so the
bracket with $\mathfrak{Q}^{\mathrm{ct}}$ vanishes.
The coexistence of both nonzero $\mathfrak{C}$ and
$\mathfrak{Q}^{\mathrm{ct}}$ is the signature of
class~\textbf{M}.
\end{remark}

% ================================================================
% 9.8.8: VIRASORO JET SPACES
% ================================================================

\subsection{Virasoro jet spaces}
% label removed: subsec:virasoro-jet-spaces
\index{Virasoro algebra!jet spaces|textbf}
\index{shadow jet space!Virasoro}

The shadow jet space $J^r(\mathrm{Vir}_c)$ is the
degree-$r$ graded piece of the modular convolution algebra.
For the Virasoro algebra, the one-generator primary space
makes the jet spaces particularly tractable: $J^r$ is
identified with the space of degree-$r$ homogeneous
polynomials in the variable~$x$ (dual to the single
generator~$T$), twisted by the BPZ inner product
structure at each conformal weight.

The full jet space includes contributions from
\emph{all} conformal weights: descendants and
quasi-primary composites. At degree~$r$, the relevant
conformal weights range from $2r$ (the minimum, from
$T^r$) upward. The number of states at weight~$h$
in the Virasoro vacuum Verma module is the
partition function $p(h - 2)$ for $h \geq 2$.

\begin{computation}[Virasoro jet space dimensions]
% label removed: comp:virasoro-jet-dimensions
\index{Virasoro algebra!jet space dimensions}
The dimension of the shadow jet space $J^r(\mathrm{Vir}_c)$
on the primary line is $1$ for all $r \geq 2$ (since
there is one generator). The dimension of the
\emph{extended} jet space, including descendants and
quasi-primary composites at total conformal weight~$2r$,
is computed from the Verma module structure.

At degree~$r$ and total weight $w = 2r + d$ (excess~$d$),
each factor is a state in the vacuum Verma module at
weight $h_i \geq 2$ with $\sum h_i = w$. The states
at weight~$h$ number $p(h-2)$. The extended jet space
dimension at excess $d = 0$ is $1$; at excess $d = 1$ it
is $r$ (choice of which factor descends); at excess
$d = 2$ it is $r + \binom{r}{2} = r(r+1)/2$.

The first several values at excess $d = 0$:

\begin{center}
\renewcommand{\arraystretch}{1.4}
\begin{tabular}{crccl}
\toprule
$r$ & $\dim J^r$ & $\dim J^r_{\mathrm{ext}}(d{\leq}2)$
 & $\sum_{d=0}^{2r}$ & Obstruction \\
\midrule
$2$ & $1$ & $1 + 2 + 3 = 6$ & $7$
 & $\kappa = c/2$ \\
$3$ & $1$ & $1 + 3 + 6 = 10$ & $15$
 & $\mathfrak{C} = 2x^3$ \\
$4$ & $1$ & $1 + 4 + 10 = 15$ & $26$
 & $Q^{\mathrm{ct}} = 10/(c(5c{+}22))$ \\
$5$ & $1$ & $1 + 5 + 15 = 21$ & $42$
 & $\mathrm{Sh}_5 = -48/(c^2(5c{+}22))$ \\
$6$ & $1$ & $1 + 6 + 21 = 28$ & $63$
 & $\mathrm{Sh}_6 \neq 0$ \\
$7$ & $1$ & $1 + 7 + 28 = 36$ & $92$
 & $\mathrm{Sh}_7 \neq 0$ \\
$8$ & $1$ & $1 + 8 + 36 = 45$ & $131$
 & $\mathrm{Sh}_8 \neq 0$ \\
\bottomrule
\end{tabular}
\end{center}

\noindent The column $\dim J^r_{\mathrm{ext}}(d{\leq}2)$
counts states through conformal-weight excess~$2$.
\end{computation}

\begin{remark}[Shapovalov determinant and jet space
structure]
% label removed: rem:shapovalov-jet
\index{Shapovalov determinant!jet space structure}
The Shapovalov determinant at weight~$h$ in the Virasoro
vacuum Verma module is
\begin{equation}% label removed: eq:shapovalov-virasoro
\det S_h = C_h \prod_{\substack{r,s \geq 1 \\ rs \leq h}}
(h_{r,s}(c))^{p(h-rs)},
\end{equation}
where $h_{r,s}(c)$ are the Kac table entries
(parametrized by
$c = 13 - 6(t + t^{-1})$, $h_{r,s} =
((rt - s/t)^2 - (t - 1/t)^2)/4$).

The Shapovalov determinant controls the singularity
structure of the jet spaces: at central charges where
$\det S_h = 0$, the Verma module develops singular
vectors, and the quasi-primary composites at that weight
become linearly dependent. The quartic contact
coefficient $Q^{\mathrm{ct}} = 10/(c(5c+22))$ has
poles at $c = 0$ and $c = -22/5$, which are exactly
the zeros of the weight-$4$ Shapovalov determinant
$\det G = c^2(5c+22)/2$
(Proposition~\ref{V1-prop:gram-wt4}).
\end{remark}

% ================================================================
% 9.8.9: GROWTH RATE THEOREM
% ================================================================

\subsection{Growth rate of the shadow obstruction tower}
% label removed: subsec:growth-rate
\index{shadow tower!growth rate|textbf}
\index{growth rate!shadow jet spaces}

\begin{theorem}[Polynomial growth of Virasoro jet spaces]
% label removed: thm:jet-space-polynomial-growth
\index{Virasoro algebra!jet space growth rate}
For the Virasoro algebra at generic central charge,
the extended jet space dimensions grow
polynomially in~$r$:
\begin{equation}% label removed: eq:jet-polynomial-growth
\dim J^r_{\mathrm{ext}}(\mathrm{Vir}_c; d \leq D)
= O(r^D)
\qquad \text{as } r \to \infty,
\end{equation}
for each fixed conformal-weight excess~$D$.
On the primary line, $\dim J^r = 1$ for all $r$.
\end{theorem}

\begin{proof}
At conformal-weight excess $d = 0$: $\dim J^r = 1$
(the primary sector, spanned by $x^r$).

At excess $d = 1$: one factor carries
$L_{-1}$ descent (weight $h_i = 3$, one state), the
remaining $r - 1$ factors are at weight~$2$. The number
of choices is $r$, giving $\dim = r$.

At excess $d = 2$: either one factor at weight~$4$
($p(2) = 2$ states: $L_{-2}|0\rangle$ and $L_{-1}^2|0\rangle$)
and $r - 1$ at weight~$2$, giving $2r$ states; or
two factors at weight~$3$ ($p(1) = 1$ state each)
and $r - 2$ at weight~$2$, giving $\binom{r}{2}$ states.
Total: $2r + \binom{r}{2} = r(r+3)/2$.

In general, at excess~$d$ the dimension is a polynomial
in~$r$ of degree at most~$d$: the leading term is
$\binom{r}{d}/d!$ from distributing $d$ units of excess
among $r$ factors. By the multinomial theorem:
\begin{equation}% label removed: eq:extended-dimension-polynomial
\dim J^r_{\mathrm{ext}}(d)
= \sum_{\substack{\lambda_1 + \cdots + \lambda_r = d \\
\lambda_i \geq 0}}
\prod_{i=1}^r p(\lambda_i)
\sim c_d\,r^d,
\end{equation}
where $c_d = \sum_{|\lambda| = d} \prod p(\lambda_i) / d!$
is a combinatorial constant depending on the partition
function.

For the cumulative dimension through excess~$D$:
$\dim J^r_{\mathrm{ext}}(d \leq D) = \sum_{d=0}^D
c_d r^d + \text{lower-order terms} = O(r^D)$.
\end{proof}

\begin{remark}[Information-theoretic content]
% label removed: rem:information-content
\index{information theory!shadow depth}
The polynomial growth of $\dim J^r_{\mathrm{ext}}$
has an information-theoretic interpretation. To specify the
holographic reconstruction of $\cA$ through degree~$r$,
one needs $\dim J^r$ real parameters on the primary line.
For Gaussian algebras, $\dim J^r = 1$ for all $r$ (the
single parameter $\kappa$ determines the rank-one Gaussian shadow,
not the full tensor-channel package). For
Virasoro, each degree adds one primary-line parameter,
so $N(r) = r - 1$ through degree~$r$.

The \emph{total} number of parameters including descendants
is $\dim J^r_{\mathrm{ext}}$, which grows polynomially.
The ratio $\dim J^r_{\mathrm{ext}} / r$ measures the
\emph{marginal complexity}: the number of new parameters
per unit increase in degree.
\end{remark}

% ================================================================
% 9.8.10: RECURSIVE STRUCTURE OF THE SHADOW COEFFICIENTS
% ================================================================

\subsection{Recursive structure of the Virasoro
shadow coefficients}
% label removed: subsec:virasoro-shadow-recursive
\index{Virasoro algebra!shadow coefficients!recursive structure}

The recursive formula~\eqref{V1-eq:shadow-recursive-formula}
determines the cubic-source contribution to the shadow coefficients for
$r \geq 5$. We record the first several values and exhibit the
closed-form solution for that contribution.

\begin{proposition}[Virasoro shadow coefficients,
cubic-source approximation]
% label removed: prop:virasoro-shadow-coefficients
\index{shadow coefficients!Virasoro}
The primary-line shadow coefficients
$\mathrm{Sh}_r(\mathrm{Vir}_c) = S_r\,x^r$
have the following dominant contribution from
the cubic source term in the shadow master equation.
The cubic-source recursion gives
\begin{equation}% label removed: eq:virasoro-shadow-recursion
S_{r+1}^{\mathrm{cub}} = -\frac{6r}{c(r+1)}\,S_r^{\mathrm{cub}},
\qquad r \geq 5,
\end{equation}
with initial conditions from the explicit tower data:
\begin{align}
S_2 &= c/2, % label removed: eq:s2 \\
S_3 &= 2, % label removed: eq:s3 \\
S_4 &= 10/(c(5c+22)), % label removed: eq:s4 \\
S_5 &= -48/(c^2(5c+22)). % label removed: eq:s5
\end{align}
For $r \geq 5$, the cubic-source closed form is
\begin{equation}% label removed: eq:virasoro-shadow-closed-form
S_r^{\mathrm{cub}} =
\frac{(-1)^{r-4} \cdot 6^{r-5} \cdot 240}
{c^{r-3}(5c+22) \cdot r}.
\end{equation}
For $r \leq 5$ this agrees with the exact shadow
coefficients: $S_r^{\mathrm{cub}} = S_r$.
For $r \geq 6$, additional contributions from
$\{\mathrm{Sh}_j, \mathrm{Sh}_k\}$ bracket terms with
$j,k \geq 4$ enter the full recursion, and the exact
values differ; see
Theorem~\textup{\ref{thm:thqg-virasoro-tower-explicit}}
and~\S\textup{\ref{subsec:thqg-virasoro-tower}}.
\end{proposition}

\begin{proof}
The cubic-source recursion
$S_{r+1}^{\mathrm{cub}} = -\frac{6r}{c(r+1)}S_r^{\mathrm{cub}}$
is iterated from $r = 5$ to $r - 1$:
\begin{align*}
S_r^{\mathrm{cub}} &= \prod_{j=5}^{r-1}\bigl(-\frac{6j}{c(j+1)}\bigr)
 \cdot S_5 \\
&= \frac{(-6)^{r-5}}{c^{r-5}}
 \cdot \frac{\prod_{j=5}^{r-1} j}{\prod_{j=5}^{r-1}(j+1)}
 \cdot S_5 \\
&= \frac{(-6)^{r-5}}{c^{r-5}}
 \cdot \frac{5}{r}
 \cdot \bigl(-\frac{48}{c^2(5c+22)}\bigr) \\
&= \frac{(-1)^{r-4} \cdot 6^{r-5} \cdot 240}
 {c^{r-3}(5c+22) \cdot r}.
\end{align*}
\emph{Verification at $r = 5$:}
$(-1)^1 \cdot 6^0 \cdot 240 / (c^2(5c+22) \cdot 5)
= -48/(c^2(5c+22)) = S_5$. \checkmark

\emph{Verification at $r = 6$:}
$(-1)^2 \cdot 6 \cdot 240 / (c^3(5c+22) \cdot 6)
= 240/(c^3(5c+22)) = S_6^{\mathrm{cub}}$. \checkmark

Cross-check via cubic-source recursion:
$S_6^{\mathrm{cub}} = -\frac{6 \cdot 5}{c \cdot 6} S_5
= -\frac{5}{c} \cdot (-\frac{48}{c^2(5c+22)})
= \frac{240}{c^3(5c+22)}$. \checkmark

\smallskip
\noindent\textit{Caveat.}
The exact $S_6$ receives an additional
$\{\mathrm{Sh}_4, \mathrm{Sh}_4\}$ cross-term not captured
by this recursion. The cubic-source formula gives the
dominant contribution for $|c| \gg 1$ but is not exact for
$r \geq 6$; see
Theorem~\ref{thm:thqg-virasoro-tower-explicit}.
\end{proof}

\begin{corollary}[Cubic-source asymptotic shadow decay]
% label removed: cor:shadow-asymptotic-decay
\index{shadow coefficients!asymptotic decay}
The Virasoro shadow coefficients satisfy, in the
cubic-source approximation,
\begin{equation}% label removed: eq:shadow-asymptotic
|S_r^{\mathrm{cub}}| \sim \frac{240}{|c|^{r-3}|5c+22| \cdot r}
\cdot 6^{r-5}
\qquad \text{as } r \to \infty.
\end{equation}
For the cubic-source approximation, the ratio is
$|S_{r+1}^{\mathrm{cub}}|/|S_r^{\mathrm{cub}}| \to 6/|c|$.
The exact full tower receives corrections from the quartic and higher
bracket terms and must be read from the full recursion. Thus the
cubic-source convergence trichotomy is:
\begin{itemize}
\item For $|c| > 6$: the cubic-source series is geometrically
 decaying.
\item For $|c| < 6$: the cubic-source series is formally divergent;
 the completed obstruction tower remains an inverse-limit object in the
 completed convolution algebra topology
 \textup{(}Theorem~\ref{V1-thm:completed-bar-cobar-strong}\textup{)}.
\item For $|c| = 6$: borderline, with
 $|S_r^{\mathrm{cub}}| \sim C/r$.
\end{itemize}
\end{corollary}

\begin{remark}[Convergence and the analytic sewing programme]
% label removed: rem:convergence-sewing
\index{analytic sewing!shadow tower convergence}
The formal divergence of the cubic-source series for $|c| < 6$ does not
obstruct the algebraic theory: $\Theta_\cA$ is still defined as a
compatible element of the completed convolution algebra. For the
analytic sewing programme, convergence is an extra hypothesis: the
shadow tower must converge in an appropriate operator norm to yield
finite partition functions. The HS-sewing condition
(Theorem~\ref{thm:general-hs-sewing}), when its hypotheses hold,
provides the required estimate.
\end{remark}

% ================================================================
% 9.8.11: EXPLICIT VIRASORO SHADOW TABLE
% ================================================================

\subsection{Explicit Virasoro shadow table}
% label removed: subsec:virasoro-shadow-table
\index{Virasoro algebra!shadow table}

\begin{computation}[Virasoro shadow coefficients through
degree~$10$, cubic-source approximation]
% label removed: comp:virasoro-shadow-table
\index{Virasoro algebra!shadow coefficients!table}

The following table records the cubic-source
approximation~$S_r^{\mathrm{cub}}$ from
Proposition~\ref{V1-prop:virasoro-shadow-coefficients}.
For $r \leq 5$, these are the exact shadow coefficients.
For $r \geq 6$, the exact values receive corrections from
$\{\mathrm{Sh}_j, \mathrm{Sh}_k\}$ cross-terms with
$j,k \geq 4$; see
\S\ref{subsec:thqg-virasoro-tower} for exact values.

\begin{center}
\renewcommand{\arraystretch}{1.5}
\small
\begin{tabular}{cll}
\toprule
$r$ & $S_r^{\mathrm{cub}}$ & Value at $c = 26$ \\
\midrule
$2$ & $c/2$ & $13$ \\
$3$ & $2$ & $2$ \\
$4$ & $\displaystyle\frac{10}{c(5c+22)}$
 & $\displaystyle\frac{10}{3952}
 \approx 2.53 \times 10^{-3}$ \\
$5$ & $\displaystyle\frac{-48}{c^2(5c+22)}$
 & $\approx -4.67 \times 10^{-4}$ \\
$6^{\dagger}$ & $\displaystyle\frac{240}{c^3(5c+22)}$
 & $\approx 9.01 \times 10^{-5}$ \\
$7^{\dagger}$ & $\displaystyle\frac{-1440}{c^4(5c+22) \cdot 7/6}$
 & $\approx -1.48 \times 10^{-5}$ \\
$8^{\dagger}$ & $\displaystyle\frac{(-6)^3 \cdot 240}
 {c^5(5c+22) \cdot 8}$
 & $\approx 2.14 \times 10^{-6}$ \\
$9^{\dagger}$ & $\displaystyle\frac{(-6)^4 \cdot 240}
 {c^6(5c+22) \cdot 9}$
 & $\approx -2.77 \times 10^{-7}$ \\
$10^{\dagger}$ & $\displaystyle\frac{(-6)^5 \cdot 240}
 {c^7(5c+22) \cdot 10}$
 & $\approx 3.26 \times 10^{-8}$ \\
\bottomrule
\end{tabular}
\end{center}

\noindent{\footnotesize ${}^{\dagger}$Cubic-source
approximation; exact values for $r \geq 6$ differ
by $\{\mathrm{Sh}_4, \mathrm{Sh}_4\}$ and higher
cross-term corrections.}

\smallskip
\noindent At $c = 26$ (the bosonic string), the cubic-source
coefficients decay geometrically with ratio
$\approx 6/26 = 3/13 \approx 0.231$. No analytic convergence statement
for the exact tail is asserted without the sewing norm.
\end{computation}

\begin{remark}[The bosonic string cubic-source decay]
% label removed: rem:bosonic-string-convergent
\index{bosonic string!shadow tower}
At $c = 26$, the Koszul dual is
$\mathrm{Vir}_{26}^! = \mathrm{Vir}_0$, and the
complementarity pairing gives
$\kappa + \kappa' = 13$, with $\kappa = 13$, $\kappa' = 0$.
The cubic-source approximation decays rapidly; the exact completed
shadow tower remains governed by the full recursion and the analytic
sewing norm. This is the \emph{depth-zero resonance shadow}
(Remark~\ref{V1-rem:virasoro-resonance-model}).
\end{remark}

% ================================================================
% 9.8.12: DICTIONARY OF RECONSTRUCTION INVARIANTS
% ================================================================

\subsection{Dictionary of reconstruction invariants}
% label removed: subsec:reconstruction-dictionary
\index{reconstruction invariants!dictionary}

\begin{center}
\renewcommand{\arraystretch}{1.5}
\small
\begin{tabular}{ccccc}
\toprule
& \textbf{G} & \textbf{L} & \textbf{C} & \textbf{M} \\
& (Gaussian) & (Lie/tree) & (Contact) & (Mixed) \\
\midrule
$r^\Theta_{\max}$ & $2$ & $3$ & $4$ & $\infty$ \\
Atom & $\mathcal{H}_k$ & $\widehat{\fg}_k$
 & $\beta\gamma$ & $\mathrm{Vir}_c$ \\
$\Theta$ & $\kappa\,x^2$ & $+\,\mathfrak{C}\,x^3$
 & $+\,\mathfrak{Q}\,x^4$ & $+\sum_{r \geq 5} S_r\,x^r$ \\
\midrule
OPE poles & $\leq 2$ & $\leq 2$ & $\leq 1$ & $\leq 4$ \\
$m_n^{\mathrm{tr}} = 0$ & $n \geq 3$ & $n \geq 3$
 & $n \geq 5$ & never \\
Termination & Formality & Jacobi & Rank-1 rigidity
 & N/A \\
\midrule
Deformation & Smooth & Smooth & Smooth & Formal $\widehat{\phantom{x}}$ \\
Rep.\ theory & Fock & Verma irred.\ & Rank-1 Fock
 & Verma singular \\
Graph sums & 1 graph & trees & 1 graph + gauge
 & all graphs \\
Growth & $O(1)$ & $O(1)$ & $O(1)$
 & $O(r^D)$ \\
\bottomrule
\end{tabular}
\end{center}

\begin{remark}[Operadic complexity theorem]
% label removed: rem:operadic-complexity-holographic
\index{operadic complexity!holographic interpretation}
The operadic complexity theorem
(Vol~I, Theorem~\ref*{V1-thm:operadic-complexity-detailed})
states that
\(r^\Theta_{\max}(\cA)\) equals the support depth of the transferred
\(A_\infty\)-structure of~\(\cA\):
the minimal $n$ such that the transferred
$A_\infty$-structure on $H^*(\barB(\cA))$ has
$m_k^{\mathrm{tr}} = 0$ for all $k > n$. Equivalently,
\(r^\Theta_{\max}\) is the last nonzero homogeneous component of the
bar-intrinsic \(L_\infty\) deformation packet.  The obstruction depth
\(d_{\mathrm{obs}}\) records only eventual \(H^2\)-vanishing and may
be smaller. The proof identifies the
genus-$0$ shadow graph sum with the homotopy transfer
tree formula at all degrees. In the
holographic language: the gravitational complexity of the
bulk theory equals the homotopy-algebraic complexity of
the boundary algebra.
\end{remark}

% ================================================================
% 9.8.13: OBSTRUCTION TOWER FOR GENERAL W-ALGEBRAS
% ================================================================

\subsection{The obstruction tower for $\Walg_N$-algebras}
% label removed: subsec:wn-obstruction-tower
\index{W-algebra@$\Walg$-algebra!obstruction tower}

\begin{proposition}[$\Walg_N$ is class~\textbf{M} for
$N \geq 2$]
% label removed: prop:wn-class-m
\index{shadow depth!$\Walg_N$}
For the principal $\Walg_N$-algebra at generic level, read in the
generic completed ambient,
\(r^\Theta_{\max}(\Walg_N) = \infty\) for all $N \geq 2$.
The obstruction at degree~$5$ contains a nonzero local source from
the coexistence of the gravitational cubic
(from $T_{(1)}T = 2T$) and the quartic contact
(from $\Lambda = {:}TT{:} - \frac{3}{10}\partial^2 T$
and higher quasi-primaries).
\end{proposition}

\begin{proof}
The gravitational cubic $\mathfrak{C} \neq 0$ (from
$T_{(1)}T = 2T$, present for all $N$) and the quartic
contact $\mathfrak{Q}^{\mathrm{ct}} \neq 0$
(Theorem~\ref{V1-thm:w-principal-wn-contact-nonvanishing})
give a nonzero cubic-quartic bracket
$\{\mathfrak{C}, \mathfrak{Q}^{\mathrm{ct}}\}_H \neq 0$.
Exact non-cancellation and the all-degree infinite-depth statement are
read from the generic completed $\Walg_N$ recursion, with the Virasoro
stress-tensor subcomplex governed by
Theorem~\ref{V1-thm:virasoro-rmax-infinity}.
\end{proof}

\begin{remark}[The $\Walg_\infty$ limit]
% label removed: rem:winfty-limit-reconstruction
\index{W-infinity@$\Walg_\infty$!reconstruction}
At the $\Walg_\infty$ limit ($N \to \infty$), the
shadow obstruction tower is infinite at each finite stage but
one-dimensional at the minimal cyclic level at the limit
(Theorem~\ref{V1-thm:winfty-scalar}): the
Mittag--Leffler argument forces all shadow projections
into the one-dimensional $H^2_{\mathrm{cyc}}$.
This identifies only the minimal cyclic direction; it does not imply
that a single number $\kappaChHodge(\Walg_\infty)$ determines the entire
modular Koszul package.
\end{remark}

% ================================================================
% 9.8.14: THE POSTNIKOV FILTRATION
% ================================================================

\subsection{Obstruction tower and Postnikov filtration}
% label removed: subsec:obstruction-tower
\index{obstruction tower|textbf}
\index{Postnikov filtration!shadow tower}

\begin{definition}[Postnikov $k$-invariant of the
shadow obstruction tower]
% label removed: def:shadow-k-invariant
\index{Postnikov $k$-invariant!shadow tower}
The \emph{$r$-th $k$-invariant} of the shadow obstruction tower is
the obstruction class
\begin{equation}% label removed: eq:shadow-k-invariant
k_r(\cA) := o_{r+1}(\cA) \in H^2(J^{r+1}(\cA), d_2).
\end{equation}
The shadow obstruction tower is
\emph{$n$-truncated} if $k_r(\cA) = 0$ for all
$r \geq n$, i.e., if \(d_{\mathrm{obs}}(\cA)\le n\).
\end{definition}

\begin{proposition}[Structure of the Postnikov filtration]
% label removed: prop:postnikov-filtration-structure
\index{Postnikov filtration!structural properties}
\begin{enumerate}[label=\textup{(\roman*)}]
\item \emph{Obstruction profile.}
The $k$-invariants $\{k_r(\cA)\}_{r \geq 2}$ determine the
obstruction-depth profile \(d_{\mathrm{obs}}(\cA)\).  The MC gauge
class \([\Theta_\cA]\) is determined only after adjoining the
\(H^1\)-lift coordinates of the shadow reconstruction packet.

\item \emph{Functoriality.}
A quasi-isomorphism $\cA \xrightarrow{\sim} \cA'$ induces
isomorphisms $k_r(\cA) \cong k_r(\cA')$ for all $r$.

\item \emph{Additivity.}
For a direct sum $\cA = \cA_1 \oplus \cA_2$ with
vanishing mixed OPE:
$k_r(\cA) = k_r(\cA_1) + k_r(\cA_2)$ for all $r$.

\item \emph{Finiteness.}
Each $H^2(J^{r+1}(\cA), d_2)$ is finite-dimensional
for all modular Koszul chiral algebras with
finitely many strong generators.
\end{enumerate}
\end{proposition}

\begin{proof}
Part~(i): by definition \(k_r=o_{r+1}\), so the sequence
\(\{k_r\}\) is exactly the \(H^2\)-obstruction profile.  The
obstruction-extension sequence shows that a vanishing \(k_r\) leaves an
\(H^1(J^{r+1}(\cA),d_2)\)-torsor of lifts; therefore the full MC gauge
class requires the lift coordinates.  The bar-intrinsic branch with
those coordinates is the complete system of invariants by the
branch-level Mittag--Leffler corollary above.
Part~(ii):
Theorem~\ref{V1-thm:shadow-homotopy-invariance}.
Part~(iii):
Proposition~\ref{V1-prop:independent-sum-factorization}.
Part~(iv): finite-dimensionality of the bar complex at
fixed degree and conformal weight.
\end{proof}

% ================================================================
% 9.8.15: DEFORMATION-THEORETIC INTERPRETATION
% ================================================================

\subsection{Deformation-theoretic interpretation}
% label removed: subsec:deformation-interpretation
\index{deformation theory!shadow tower interpretation}

\begin{proposition}[The MC moduli as a formal moduli problem]
% label removed: prop:mc-formal-moduli
\index{formal moduli!MC moduli}
The assignment
\begin{equation}% label removed: eq:mc-formal-moduli
R \;\mapsto\; \MC(\gAmod \widehat{\otimes} \mathfrak{m}_R)
/ \mathrm{gauge}
\end{equation}
(where $R$ is an Artinian dg algebra and $\mathfrak{m}_R$
its maximal ideal) defines a formal moduli problem
whose tangent complex is $\gAmod[1]$ and whose
obstruction groups are
$H^{n+2}(\gAmod)$ for $n$-th order deformations.
\end{proposition}

\begin{proof}
This is the standard correspondence between dg~Lie
algebras and formal moduli problems
(Hinich, Lurie~\cite{Lurie-DAG-X},
Pridham). The functor satisfies Schlessinger's
conditions: $H^0$ gives automorphisms, $H^1$ gives
deformations, $H^2$ gives obstructions.
The weight filtration on $\gAmod$ provides the
Artinian stratification.
\end{proof}

The support depth \(r^\Theta_{\max}(\cA)\) is the length of the
bar-intrinsic MC packet along the weight direction.  The obstruction
depth \(d_{\mathrm{obs}}(\cA)\) is the length of the \(H^2\)-obstruction
profile.  When \(r^\Theta_{\max}<\infty\), the bar-intrinsic branch is
finite type in the weight direction; when
\(r^\Theta_{\max}=\infty\), the branch is a completed formal object.
Other branches of the formal moduli problem are controlled by the same
obstruction-extension sequence and are not classified by the
bar-intrinsic packet alone.

% ================================================================
% 9.8.16: HOLOGRAPHIC SYNTHESIS
% ================================================================

\subsection{Synthesis: holographic reconstruction from
obstruction data}
% label removed: subsec:holographic-synthesis
\index{holographic reconstruction!synthesis}

\begin{theorem}[Universal Holography, admissible shadow projection with realization datum;
\(\ClaimStatusConditional\); licensing tags \(\beta+\gamma+\delta\)]\label{thm:holographic-recon-shadow}
Let $\cA$ be a chiral algebra on a smooth projective curve $X$
equipped with a conformal vector $\omega$ at non-critical level.
Assume that the following data are fixed.
\begin{enumerate}[label=\textup{(H\arabic*)}]
\item \emph{Boundary realisation.} Either $\cA$ lies in the exact
boundary-linear Costello--Gwilliam/Costello--Li sector,
$\cA=V_k(\mathfrak g)$ is an affine Kac--Moody algebra with
$k\neq -h^\vee$ and the holomorphic Chern--Simons boundary model of
Theorem~\ref{thm:E3-topological-km}, or
$\cA=\mathcal{W}_k(\mathfrak g,f)$ is obtained by Drinfeld--Sokolov
reduction of $V_k(\mathfrak g)$ with $k\neq -h^\vee$ and carries the
Costello--Gaiotto DS boundary model.  This datum includes a chosen
holomorphic-topological BV/factorization model
\(\mathcal F_{\cA,\Xi}\) and a boundary comparison
\[
\eta^\partial_\cA\colon
\Obs^\partial(\mathcal F_{\cA,\Xi})\xrightarrow{\simeq}\cA .
\]
\item \emph{Bulk--centre comparison.} The chosen model lies in the
scope of the bulk--Hochschild comparison, and a map
\[
\chi_{\mathrm{HT},\cA}\colon
\Zderch(\cA)\to
\Obs^{\mathrm{bulk}}(\mathcal F_{\cA,\Xi})
\]
is part of the datum in the declared ambient.  In the
physical prefactorisation sector this is the comparison supplied by
Theorem~\ref{thm:bulk_hochschild}. In the boundary-linear sector the
compact-generation and derived-centre quasi-isomorphism hypotheses are
those of Theorem~\ref{thm:boundary-linear-bulk-boundary}. In the DS
sector the comparison is read through the DS--Hochschild bridge
Theorem~\ref{thm:ds-hochschild-bridge}: cohomologically for
class~$\mathsf M$, at chain level in the weight-completed/pro ambient,
and on the original complex only when the bounded conformal-weight-shift
HPL criterion of Remark~\ref{rem:class-M-hpl-convergence} is supplied.
\item \emph{Topological refinement.} If the conclusion is read as a
strict $E_3$-topological statement, the topologisation datum of
Convention~\ref{conv:chd-e3-chiral-meaning} is part of the input. For
affine Kac--Moody and DS $\mathcal W$-algebras this datum is supplied by
Theorems~\ref{thm:E3-topological-km},
\ref{thm:E3-topological-DS}, and
\ref{thm:E3-topological-DS-general}.
\end{enumerate}
Write
\[
\Xi_\cA=(\mathcal F_{\cA,\Xi},\eta^\partial_\cA,
\chi_{\mathrm{HT},\cA},\mathcal A(\cA)).
\]
Then the \(\Xi\)-realized admissible datum has the supplied relative
$3$d holomorphic-topological gauge theory
\(\mathcal{T}_{\cA,\Xi}:=\mathcal F_{\cA,\Xi}\) on
$X\times\mathbb{R}$ such that:
\begin{enumerate}
\item[(i)] \emph{Boundary restriction.}
\(\eta^\partial_\cA\colon
\mathrm{Obs}^\partial(\mathcal{T}_{\cA,\Xi})\simeq \cA\)
as $E_1$-chiral algebras.
\item[(ii)] \emph{Bulk--centre comparison.}
\[
\chi_{\mathrm{HT},\cA}\colon
Z^{\mathrm{der}}_{\mathrm{ch}}(\cA)\xrightarrow{\simeq}
\mathrm{Obs}^{\mathrm{bulk}}(\mathcal{T}_{\cA,\Xi})
\]
is an equivalence in the algebraic, cohomological, or completed/pro
ambient specified by the datum. Under hypothesis \textup{(H3)} this is an
equivalence of $E_3$-topological factorisation algebras on
$X\times\mathbb{R}$.
\item[(iii)] \emph{Functoriality.} On morphisms preserving the chosen
boundary-linear, affine, or DS realisation and the comparison package,
\((\cA,\Xi_\cA)\mapsto \mathcal{T}_{\cA,\Xi}\)
extends to a functor
\[
 \Phi_{\mathrm{hol}}\colon
 \mathsf{ChirAlg}^{\omega,\mathrm{adm},\Xi}_X
 \longrightarrow
 \mathsf{HT\text{-}QFT}^{\Xi}_{X\times\mathbb{R}}
\]
on the admissible subcategory, compatible with Drinfeld--Sokolov
reduction on the \(\Xi\)-compatible DS locus:
\[
\Phi_{\mathrm{hol}}(\mathcal{W}_k(\mathfrak{g},f),\Xi_{\mathrm{DS}})
\simeq
\Phi_{\mathrm{hol}}(V_k(\mathfrak{g}),\Xi_{V_k})
/\!/_{\mathrm{BV}}\mathrm{DS}_f .
\]
\item[(iv)] \emph{Class coverage.} Class $\mathrm{G}$ and the free-field
part of class $\mathrm{C}$ are realised by Costello--Li abelian
holomorphic Chern--Simons on the boundary-linear exact locus; class
$\mathrm{L}$ is realised by the affine holomorphic Chern--Simons model
of Theorem~\ref{thm:E3-topological-km}; class $\mathrm{M}$ (Virasoro,
$\mathcal{W}_N$) is realised by Costello--Gaiotto holomorphic Chern--
Simons with Drinfeld--Sokolov boundary condition, with bulk comparison
controlled by the DS--Hochschild bridge
$\mathrm{ChirHoch}^\bullet(\mathcal{W}_k(\mathfrak{g}))\simeq
H^\bullet_{\mathrm{DS}}(\mathrm{ChirHoch}^\bullet(V_k(\mathfrak{g})))$
in the cohomological or completed/pro sense specified in \textup{(H2)}.
\end{enumerate}
\end{theorem}

\begin{proof}
Hypothesis \textup{(H1)} supplies the BV theory and the boundary
comparison. In the
boundary-linear exact sector, boundary recovery is the
Costello--Gwilliam restriction theorem applied to the Costello--Li
model. For affine Kac--Moody it is the holomorphic Chern--Simons
boundary model of Theorem~\ref{thm:E3-topological-km}. In the DS sector,
Costello--Gaiotto identify the DS boundary condition for holomorphic
Chern--Simons, and Arakawa's theorem
$H^\bullet_{\mathrm{DS}}(V_k(\mathfrak{g}))\simeq
\mathcal{W}_k(\mathfrak{g})$ identifies the boundary chiral algebra.
The comparison with the boundary observables is \(\eta^\partial_\cA\).

Hypothesis \textup{(H2)} gives the bulk comparison. For
boundary-linear models, Theorem~\ref{thm:boundary-linear-bulk-boundary}
supplies the Morita comparison with the derived centre of the boundary
algebra. For a chosen physical prefactorisation model,
Theorem~\ref{thm:bulk_hochschild} supplies
\(\Psi_{\mathrm{Ran}}\), equivalently the inverse comparison
\(\chi_{\mathrm{HT},\cA}\). In the DS sector,
Theorem~\ref{thm:ds-hochschild-bridge} supplies exactly the stated
cohomological or completed/pro comparison, with original-complex
chain-level comparison only under the bounded-shift HPL criterion.
Thus \(\chi_{\mathrm{HT},\cA}\) identifies the algebraic closed sector
with \(\mathrm{Obs}^{\mathrm{bulk}}(\mathcal{T}_{\cA,\Xi})\) in the
asserted ambient.
Hypothesis \textup{(H3)} upgrades this algebraic open/closed bulk to
the strict $E_3$-topological factorisation algebra, by
Theorem~\ref{thm:chiral-higher-deligne} together with the named
topologisation theorems.

Functoriality is relative to the admissible data: a morphism preserving
the boundary-linear, affine, or DS presentation and the \(\Xi\)-package
is a morphism of boundary conditions in the Costello BV complex, hence
induces the boundary map on observables and the corresponding map on
Hochschild cochains and derived centres, compatible with
\(\chi_{\mathrm{HT}}\). The DS formula is the functoriality of the
Costello--Gaiotto DS boundary condition composed with
Theorem~\ref{thm:ds-hochschild-bridge} on the \(\Xi\)-compatible locus.
The class coverage is precisely the disjoint union of the three
constructions in \textup{(H1)}.
\end{proof}

\begin{remark}[Shadow-tower reconstruction]
The following theorem reconstructs the bar-intrinsic gauge class of
\(\Theta_\cA\) from a compatible shadow reconstruction packet.
Obstruction classes alone decide extendability; they do not select a
point of the \(H^1\)-torsor of lifts in the exact sequence
\[
H^1(J^{r+1}(\cA),d_2)\to E_\cA(r+1)\to E_\cA(r)\to
H^2(J^{r+1}(\cA),d_2).
\]
Thus the packet contains both the obstruction class and the lift
coordinate in a chosen bar-intrinsic gauge slice. Finite-order data
suffice exactly on the finite-depth strata; class~$\mathsf M$ is a
completed inverse-limit statement.
\end{remark}

\begin{theorem}[Filtered shadow-tower reconstruction of the bar-intrinsic MC gauge class;
\(\ClaimStatusConditional\); licensing tags \(\alpha+\gamma+\epsilon\)]
% label removed: thm:holographic-reconstruction
\label{thm:filtered-shadow-tower-reconstruction}
\index{holographic reconstruction!main theorem}
Let $\cA$ be a modular Koszul chiral algebra with
holographic modular Koszul datum
$\mathcal{H}(T) = (\cA, \cA^!, \mathcal{C}, r(z),
\Theta_\cA, \nabla^{\mathrm{hol}})$
\textup{(}Definition~\textup{\ref{V1-def:holographic-modular-koszul-datum})}.
Assume the modular convolution dg Lie algebra \(\gAmod\) is complete
and separated for the weight filtration \(F^\bullet\), and fix the
bar-intrinsic gauge component containing
\(\Theta_\cA\in F^2{\gAmod}^1\).  A \emph{shadow reconstruction packet}
is the compatible system
\[
\mathfrak S(\cA)=
\Bigl(
S_2,\;S_3,\;S_4,\;
\{(o_n(\cA),\lambda_n)\}_{n\geq 5}
\Bigr),
\]
where:
\begin{enumerate}[label=\textup{(\roman*)}]
\item \(S_2=H_\cA=\kappaChHodge(\cA)\), the Hessian
 \textup{(}degree~$2$\textup{)};
\item \(S_3=\mathfrak C_\cA\), the cubic shadow
 \textup{(}degree~$3$\textup{)};
\item \(S_4=\mathfrak Q^{\mathrm{ct}}_\cA\), the quartic contact
 \textup{(}degree~$4$\textup{)};
\item \(o_n(\cA)\in H^2(J^n(\cA),d_2)\) is the degree-\(n\)
 obstruction class, and
 \(\lambda_n\in H^1(J^n(\cA),d_2)\) is the lift coordinate in the
 \(H^1\)-torsor of lifts, measured in the chosen bar-intrinsic gauge
 slice.
\end{enumerate}
The compatible packet \(\mathfrak S(\cA)\) determines the gauge class
\([\Theta_\cA]\in \varprojlim_r E_\cA(r)\), and hence the
bar-intrinsic MC element \(\Theta_\cA\) in the chosen gauge component.
The reconstruction has the following scope.
\begin{itemize}
\item \emph{Finite support depth.} If \(r^\Theta_{\max}<\infty\), then the finite
 packet
 \[
 (S_2,S_3,S_4,\{(o_n,\lambda_n)\}_{5\leq n\leq r^\Theta_{\max}})
 \]
 determines \([\Theta_\cA]\), after the gauge-slice convention
 \(\lambda_n=0\) for \(n>r^\Theta_{\max}\).
\item \emph{Completed} if \(r^\Theta_{\max} = \infty\): the compatible system
 \(\{(o_n,\lambda_n)\}_{n\geq 5}\) determines
 \[
 [\Theta_\cA]=\varprojlim_r[\Theta_\cA^{\leq r}]
 \]
 in the \(F^\bullet\)-completed convolution algebra topology.
 Analytic convergence of scalar partition functions is a separate
 sewing/norm hypothesis.
\end{itemize}
The four shadow classes provide the classification:
\[
\begin{array}{lcl}
\textbf{G}\;\text{(Gaussian)}: & &
 \text{canonical gauge slice; } S_2 \text{ suffices.}\\
\textbf{L}\;\text{(Lie/tree)}: & &
 \text{canonical gauge slice; } (S_2,S_3) \text{ suffice.}\\
\textbf{C}\;\text{(Contact)}: & &
 \text{canonical gauge slice; } (S_2,S_3,S_4) \text{ suffice.}\\
\textbf{M}\;\text{(Mixed)}: & &
 \text{the full compatible sequence }
 \{(o_n,\lambda_n)\}_{n\ge5} \text{ is required.}
\end{array}
\]
\end{theorem}

\begin{proof}
The obstruction-extension exact sequence identifies the fiber of
\(E_\cA(r+1)\to E_\cA(r)\) over a liftable point with the
\(H^1(J^{r+1}(\cA),d_2)\)-torsor of lifts, and sends a point of
\(E_\cA(r)\) to the obstruction class in
\(H^2(J^{r+1}(\cA),d_2)\).  Inductively, \(S_2,S_3,S_4\) initialize
the truncation and the pair \((o_n,\lambda_n)\) decides whether the
next lift exists and which lift in the torsor is chosen.  Hence a
compatible packet determines a unique element of every \(E_\cA(r)\)
in the fixed bar-intrinsic gauge component.

The finite-depth assertion is
Theorem~\ref{V1-thm:shadow-depth-dichotomy} together with
Theorem~\ref{V1-thm:gaussian-rmax-two} (class~\textbf{G}),
Theorem~\ref{V1-thm:lie-rmax-three} (class~\textbf{L}),
Theorem~\ref{V1-thm:contact-rmax-four} (class~\textbf{C}),
and
Theorem~\ref{V1-thm:virasoro-rmax-infinity} (class~\textbf{M}).
For \(r^\Theta_{\max}<\infty\) the canonical gauge slice has
\(\lambda_n=0\) above \(r^\Theta_{\max}\), so no higher torsor coordinate
enters.  For class~\(\mathsf M\), no finite truncation determines the
tail, and the completed inverse-limit statement is exactly
Corollary~\ref{V1-cor:mittag-leffler-shadow-tower}.  It is a formal
statement in the \(F^\bullet\)-completed convolution algebra, not an
analytic partition-function convergence theorem.
\end{proof}

\begin{corollary}[$\Xi$-realized Universal Holography, strict $E_3$ projection;
\(\ClaimStatusConditional\); licensing tags \(\beta+\gamma+\delta\)]\label{thm:holographic-recon-e3}
Let \((\mathcal A,\Xi_{\mathcal A})\) satisfy the admissibility
hypotheses of Theorem~\ref{thm:holographic-recon-shadow}. If the
topologisation datum of hypothesis \textup{(H3)} is fixed, then
\[
Z^{\mathrm{der}}_{\mathrm{ch}}(\mathcal A)
\xrightarrow{\;\chi_{\mathrm{HT},\mathcal A}\;\simeq\;}
\mathrm{Obs}^{\mathrm{bulk}}(\mathcal T_{\mathcal A,\Xi})
\]
is the comparison equivalence of $E_3$-topological factorisation
algebras on
$X\times\mathbb R$. For affine Kac--Moody and principal or
non-principal DS $\mathcal W$-algebras, the required topologisation
datum is supplied by Theorems~\ref{thm:E3-topological-km},
\ref{thm:E3-topological-DS}, and
\ref{thm:E3-topological-DS-general}.
\end{corollary}
\begin{proof}
This is Theorem~\ref{thm:holographic-recon-shadow}\textup{(ii)}
after imposing hypothesis \textup{(H3)} on the \(\Xi\)-realized
comparison map \(\chi_{\mathrm{HT},\mathcal A}\). The named affine and
DS topologisation theorems verify \textup{(H3)} in the listed cases.
\end{proof}

% ================================================================
% 9.8.17: WORKED RECONSTRUCTIONS
% ================================================================

\subsection{Worked reconstructions}
% label removed: subsec:worked-reconstructions
\index{holographic reconstruction!worked examples}

\begin{example}[Heisenberg reconstruction]
% label removed: ex:heisenberg-reconstruction
\index{Heisenberg algebra!holographic reconstruction}
Input: $\kappaChHodge(\mathcal{H}_k) = k$.
Output: $\Theta_{\mathcal{H}_k} = k\,x^2$.
Reconstruction terminates at step~$1$.
\end{example}

\begin{example}[Affine Kac--Moody reconstruction]
% label removed: ex:km-reconstruction
\index{affine Kac--Moody!holographic reconstruction}
Input: $\kappaChHodge(\widehat{\fg}_k) = \dim(\fg)\cdot(k+h^\vee)/(2h^\vee)$,
$\mathfrak{C}_{\widehat{\fg}} = f^{ab}{}_c\,x_a x_b x^c$.
Output:
$\Theta_{\widehat{\fg}_k} = \bigl(\dim(\fg)\cdot(k+h^\vee)/(2h^\vee)\bigr)\,\kappa_{ab}x^a x^b
+ f^{ab}{}_c x_a x_b x^c$.
Reconstruction terminates at step~$2$.
\end{example}

\begin{example}[$\beta\gamma$ reconstruction]
% label removed: ex:betagamma-reconstruction
\index{beta-gamma system@$\beta\gamma$ system!holographic reconstruction}
Input: $\kappaChHodge(\beta\gamma) = 6\lambda^2 - 6\lambda + 1$,
$\mathfrak{Q}^{\mathrm{ct}}_{\beta\gamma} \neq 0$.
Output:
$\Theta_{\beta\gamma}^{\leq 4} = \kappa\,x^2 +
\mathfrak{Q}^{\mathrm{ct}}\,x^4$.
Reconstruction terminates at step~$3$ (the cubic is
gauge-trivial).
\end{example}

\begin{example}[Virasoro reconstruction]
% label removed: ex:virasoro-reconstruction
\index{Virasoro algebra!holographic reconstruction}
Input: $\kappa = c/2$, $\mathfrak{C} = 2x^3$,
$Q^{\mathrm{ct}} = 10/(c(5c+22))$, and the full shadow recursion. The
cubic-source contribution satisfies
$S_{r+1}^{\mathrm{cub}} = -6r/(c(r+1))\,S_r^{\mathrm{cub}}$, while the
exact coefficients for $r\geq 6$ also contain higher bracket terms.
Output:
$\Theta_{\mathrm{Vir}_c} = \sum_{r=2}^\infty S_r\,x^r$ in the completed
generic ambient, with $S_2,S_3,S_4,S_5$ as in
Proposition~\ref{V1-prop:virasoro-shadow-coefficients} and the tail
read through the full recursion.
At $c = 26$:
\[
\Theta_{\mathrm{Vir}_{26}} = 13x^2 + 2x^3
+ \frac{10}{26(152)}\,x^4
- \frac{48}{26^2(152)}\,x^5 + \cdots,
\]
where the displayed terms are exact through degree~$5$; analytic
geometric convergence of the tail is not asserted without the sewing
norm.
\end{example}

% ================================================================
% 9.8.18: SHAPOVALOV SINGULARITIES
% ================================================================

\subsection{Shapovalov singularities and reconstruction walls}
% label removed: subsec:shapovalov-singularities
\index{Shapovalov determinant!reconstruction walls}

\begin{proposition}[Shapovalov singularities of the
shadow obstruction tower]
% label removed: prop:shapovalov-shadow-singularities
\index{shadow tower!singularities}
For the Virasoro algebra, the shadow coefficient $S_r$
has poles at the zeros of the Shapovalov determinant
$\det S_h$ for $h \leq 2r$. Specifically:
\begin{enumerate}[label=\textup{(\roman*)}]
\item $S_4 = 10/(c(5c+22))$ has poles at $c = 0$ and
 $c = -22/5$, from the weight-$4$ Shapovalov zeros.
\item For $r \geq 5$, the cubic-source approximation
 $S_r^{\mathrm{cub}}$ has poles at $c = 0$
 (of order $r - 3$) and $c = -22/5$ (of order~$1$).
 No new poles appear from the cubic-source recursion.
 The full recursion, which includes
 $\{\mathrm{Sh}_j, \mathrm{Sh}_k\}$ cross-terms
 for $r \geq 6$, introduces higher powers of
 $(5c+22)$ in the denominator; see
 Theorem~\textup{\ref{thm:thqg-virasoro-tower-explicit}}.
\end{enumerate}
\end{proposition}

\begin{proof}
The cubic-source recursion
$S_{r+1}^{\mathrm{cub}} = -6r/(c(r+1))\,S_r^{\mathrm{cub}}$
introduces a pole at $c = 0$ at each step. The pole at
$c = -22/5$ comes from $S_4$ and is preserved but not
amplified by the cubic-source recursion.
\end{proof}

\begin{remark}[Singular vectors as reconstruction walls]
% label removed: rem:singular-vectors-walls
\index{singular vectors!as reconstruction walls}
At the Shapovalov zeros, the shadow obstruction tower develops
singularities: the quasi-primary composites become
linearly dependent, and the jet space dimension drops.
At $c = -22/5$, the weight-$4$ Verma module has a null
vector, so $\Lambda$ becomes a descendant and the quartic
contact shadow degenerates. At admissible levels
$c = c_{p,q}$ (the minimal model central charges), the
Virasoro algebra has finitely many primary fields, and
the shadow obstruction tower reduces to a finite-dimensional problem
controlled by the fusion rules.
\end{remark}

\begin{proposition}[Pole structure of the cubic-source shadow
approximation]
% label removed: prop:pole-structure-shadow-series
\index{shadow tower!pole structure}
The cubic-source approximation to the Virasoro shadow series
$\Theta_{\mathrm{Vir}_c}^{\mathrm{cub}}(x)
= \sum_{r=2}^\infty S_r^{\mathrm{cub}} x^r$
has the following pole structure as a function of~$c$:
\begin{enumerate}[label=\textup{(\roman*)}]
\item At $c = 0$: the terms $S_r^{\mathrm{cub}}$ with $r \geq 4$ have
 poles of increasing order, so the formal series has
 an essential singularity.
\item At $c = -22/5$: each $S_r^{\mathrm{cub}}$ with $r \geq 4$ has a
 simple pole, so the cubic-source series has a simple pole
 at $c = -22/5$.
\item At all other values of $c$: the cubic-source formal series
 is analytic (as a formal power series in~$x$ with
 $c$-dependent coefficients).
\end{enumerate}
The exact tower from degree~$6$ onward also contains the higher bracket
terms of Theorem~\textup{\ref{thm:thqg-virasoro-tower-explicit}}, and
its pole structure must be read from that full recursion.
\end{proposition}

\begin{proof}
For the cubic-source contribution, the closed form is
$S_r^{\mathrm{cub}} = (-1)^{r-4} 6^{r-5} \cdot 240 /
(c^{r-3}(5c+22) \cdot r)$
for $r \geq 5$: the pole at $c = 0$ has order $r - 3$
(growing with $r$), giving essential singularity in the
summed series; the pole at $c = -22/5$ has order~$1$
for all~$r$, giving a simple pole in the cubic-source summed series.
\end{proof}

\begin{remark}[Weight-$6$ Shapovalov analysis]
% label removed: rem:weight-6-shapovalov
\index{Shapovalov determinant!weight 6}
At conformal weight~$6$, the Virasoro vacuum Verma module
has a $4$-dimensional space with basis
$L_{-6}|0\rangle$, $L_{-4}L_{-2}|0\rangle$,
$L_{-3}^2|0\rangle$, $L_{-2}^3|0\rangle$.
The Shapovalov determinant is
\[
\det S_6 = C_6 \cdot c^4 \cdot (5c+22)^2 \cdot (2c+1)
\cdot (7c+68),
\]
where $C_6$ is a positive constant. The zero at
$c = -1/2$ is the weight-$6$ null vector of the Ising
model ($c = 1/2$ in the unitary parametrization, but
$c = -1/2$ appears in the dual channel); the zero at
$c = -68/7$ controls the $\Walg_4$ structure constant
(Theorem~\ref{V1-thm:c334}). These zeros affect the
multi-generator shadow obstruction tower at degree~$6$ and above.
They do not affect the cubic-source primary-line approximation, which
uses only the propagator $P = 2/c$ and the quadratic Hessian; the exact
primary-line coefficients are read through the full recursion.
\end{remark}

\begin{computation}[Shapovalov zeros and shadow
singularities]
% label removed: comp:shapovalov-zeros-shadow
\index{Shapovalov determinant!zeros versus shadow poles}
The following table lists the first several Shapovalov
zeros and their effect on the shadow obstruction tower:

\begin{center}
\renewcommand{\arraystretch}{1.4}
\begin{tabular}{clll}
\toprule
Weight $h$ & Zero location & Source & Effect on shadow \\
\midrule
$4$ & $c = 0$ (order $2$) & $\det G = c^2(5c+22)/2$
 & Pole in $S_r^{\mathrm{cub}}$ for $r \geq 4$ \\
$4$ & $c = -22/5$ & $\det G = c^2(5c+22)/2$
 & Simple pole in $S_r^{\mathrm{cub}}$ for $r \geq 4$ \\
$6$ & $c = -1/2$ & $\det S_6 \ni (2c+1)$
 & No effect on cubic-source primary line \\
$6$ & $c = -68/7$ & $\det S_6 \ni (7c+68)$
 & No effect on cubic-source primary line \\
$8$ & $c = -46/3$ & $\det S_8 \ni (3c+46)$
 & No effect on cubic-source primary line \\
\bottomrule
\end{tabular}
\end{center}

\noindent The cubic-source primary-line coefficients see only the
weight-$4$ Shapovalov zeros. The higher zeros affect the
\emph{extended} shadow obstruction tower (including descendants and
quasi-primary composites at higher weights), which controls the
multi-channel reconstruction for $\Walg_N$ with $N \geq 3$ and enters
the exact Virasoro tower through the full recursion from degree~$6$
onward.
The curvature of the multi-channel shadow connection equals
the propagator variance~$\delta_{\mathrm{mix}}$
(Volume~I, Theorem~\ref{V1-thm:propagator-variance}).
\end{computation}

% ================================================================
% 9.8.19: SPECIFIC CENTRAL CHARGES
% ================================================================

\subsection{Specific central charges}
% label removed: subsec:specific-central-charges
\index{Virasoro algebra!specific central charges}

\begin{example}[Ising model, $c = 1/2$]
% label removed: ex:ising-shadow
\index{Ising model!shadow tower}
$\kappa = 1/4$, $Q^{\mathrm{ct}} = 40/49$,
$S_5 = -384/49$.
The ratio $6/c = 12 > 1$: the shadow obstruction tower diverges
formally, reflecting the strong-coupling regime.
\end{example}

\begin{example}[Critical string, $c = 26$]
% label removed: ex:critical-string-shadow
\index{bosonic string!shadow tower!$c = 26$}
$\kappa = 13$, $Q^{\mathrm{ct}} = 5/1976$,
$|S_{r+1}^{\mathrm{cub}}/S_r^{\mathrm{cub}}|
\to 3/13 \approx 0.231$ for the cubic-source approximation.
The completed shadow tower is controlled formally; analytic convergence
of the exact tail requires the sewing norm. The Koszul dual
$\mathrm{Vir}_{26}^! = \mathrm{Vir}_0$ has $\kappa' = 0$:
a degenerate dual, the depth-zero resonance shadow
(Remark~\ref{V1-rem:virasoro-resonance-model}).
\end{example}

\begin{example}[Virasoro self-duality, $c = 13$]
% label removed: ex:virasoro-self-dual
\index{Virasoro algebra!self-duality at $c = 13$}
$\mathrm{Vir}_{13}^! = \mathrm{Vir}_{13}$
(self-dual). $\kappa = 13/2$,
$|S_{r+1}^{\mathrm{cub}}/S_r^{\mathrm{cub}}|
\to 6/13 \approx 0.462$ for the cubic-source approximation.
Self-duality means the complementarity splitting
$\mathbf{C}_g = \mathbf{Q}_g \oplus \mathbf{Q}_g^!$
satisfies $\mathbf{Q}_g \cong \mathbf{Q}_g^!$:
the two Lagrangians coincide.
\end{example}

% ================================================================
% 9.8.20: COMPLEXITY HIERARCHY
% ================================================================

\subsection{The reconstruction complexity hierarchy}
% label removed: subsec:complexity-hierarchy
\index{reconstruction complexity|textbf}

\begin{theorem}[Complexity hierarchy]
% label removed: thm:complexity-hierarchy
\index{complexity hierarchy!holographic}
The number of independent parameters $N(r)$ needed to
specify $\Theta_\cA^{\leq r}$ on the primary line is:
\begin{equation}% label removed: eq:complexity-hierarchy
N(r) = \begin{cases}
1 & \text{for class } \textbf{G}\;\;(\text{all } r), \\
2 & \text{for class } \textbf{L}\;\;(\text{all } r), \\
3 & \text{for class } \textbf{C}\;\;(\text{all } r), \\
r - 1 & \text{for class } \textbf{M}\;\;(r \geq 2).
\end{cases}
\end{equation}
\end{theorem}

\begin{proof}
For classes \textbf{G}, \textbf{L}, \textbf{C}: the bar-intrinsic
packet has support through \(r^\Theta_{\max}\), giving
\(r^\Theta_{\max}-1\) primary-line parameters in the canonical gauge
slice.
For class~\textbf{M}: each degree level adds one
primary-line parameter $S_r$, giving $r - 1$ through
degree~$r$.
\end{proof}

% ================================================================
% 9.8.21: CONCLUDING PERSPECTIVE
% ================================================================

\subsection{Concluding perspective}
% label removed: subsec:reconstruction-conclusion

The holographic reconstruction problem (determining the
bar-intrinsic MC gauge class from its shadow-tower data, with
finite-order data only on finite support-depth strata) is controlled by
the extension tower \(\{E_\cA(r)\}\), its obstruction classes
\(\{o_{r+1}(\cA)\}\), and its \(H^1\)-lift coordinates. The dichotomy
theorem gives the clean criterion: reconstruction terminates if and
only if the support shadow depth is finite. The four
shadow classes (Gaussian, Lie/tree, contact/quartic,
mixed/infinite) exhaust the structural possibilities
for the standard landscape.

The synthesis of deformation theory (the extension tower as a
Postnikov filtration), representation theory (the Shapovalov
determinant controlling jet space dimensions), combinatorial algebra
(graph-sum contributions to each obstruction), and information theory
(the polynomial growth rate of jet space dimensions) gives the formal
structure of the reconstruction problem on the stated ambient.

For the Virasoro algebra, the simplest class~\textbf{M}
example, the cubic-source recursion
$S_{r+1}^{\mathrm{cub}} = -6r/(c(r+1))S_r^{\mathrm{cub}}$ controls the
displayed approximation; the exact tower includes higher bracket terms
from degree~$6$ onward. The completed inverse-limit reconstruction is
formal. Analytic convergence in an operator or sewing norm is supplied
only under the hypotheses of the HS-sewing theorem
(Theorem~\ref{thm:general-hs-sewing}).

The gravitational-complexity interpretation
(Section~\ref{sec:thqg-gravitational-complexity}) gives
these results a conditional physical reading: the shadow depth of the
boundary algebra measures the homotopy-algebraic complexity of the
bulk model in the chosen HT construction. Gaussian boundaries give free
bulk theories; Lie/tree boundaries give semistrict higher-spin models;
contact boundaries give quartic vertices; mixed boundaries require an
infinite completed $L_\infty$ tower. Identifying that tower with a
metric-gravity saddle expansion is a separate physical hypothesis.

% ======================================================================
%
% TWISTED HOLOGRAPHY AND BRANE EXAMPLES
%
% ======================================================================

\subsection{Twisted holography examples: D3, M2, and ABJM}
\label{subsec:holographic-twisted-examples}
\index{twisted holography!examples|textbf}

The holographic modular Koszul datum
$\mathcal{H}(\mathcal{T})
= (\cA, \cA^!, \mathcal{C}, r(z), \Theta_\cA,
\nabla^{\mathrm{hol}})$
is computed for the standard brane systems of the
Costello--Li programme \cite{CL16,CostelloGaiotto18}.

\begin{proposition}[Twisted $\mathcal{N} = 4$ SYM datum]
\label{prop:twisted-n4-datum}
\index{twisted holography!D3 brane}
For the D3 brane \textup{(}twisted $\mathcal{N} = 4$ SYM
at rank~$N$\textup{)}:
\begin{enumerate}[label=\textup{(\roman*)}]
\item The boundary VOA is affine $\mathfrak{gl}_N$ at level $k = 1$
 with $\kappaChHodge(\mathfrak{gl}_N, 1)
 = (N^2{-}1)(N{+}1)/(2N) + 1$.
\item The holographic R-matrix is $r(z) = k\,\Omega_{\mathrm{tr}}/z$
 with $k = 1$.
\item The genus expansion
 \(F_g(\mathfrak{gl}_{N,1})
 = \kappaChHodge(\mathfrak{gl}_{N,1})\lambda_g^{\mathrm{FP}}\)
 \textup{(}uniform-weight\textup{)} computes twisted amplitudes.
\item The anomaly matching
 $\kappa_{\mathrm{eff}} = \kappaChHodge(\mathrm{matter})
 + \kappaChHodge(\mathrm{ghost}) = 0$
 is a consistency check.
\end{enumerate}
\end{proposition}

\begin{proposition}[ABJM holographic datum]
\label{prop:abjm-datum}
\index{ABJM theory!holographic datum}
For ABJM at rank~$N$ and level~$k$:
\begin{enumerate}[label=\textup{(\roman*)}]
\item The boundary VOA is the BRST reduction
 $\cA_{\mathrm{ABJM}}(N,k)
 = H^0_{\mathrm{BRST}}(V_k(\mathfrak{gl}_N) \otimes
 V_{-k}(\mathfrak{gl}_N) \otimes \mathrm{Sb}^{4N^2})$.
\item The modular characteristic is
 $\kappaChHodge(\cA_{\mathrm{ABJM}}) = -(N^2+1)$
 \textup{(}from $(N^2{-}1)$ for the $\mathrm{CS}_k + \mathrm{CS}_{-k}$
 sectors and $-2N^2$ for $4N^2$ symplectic bosons;
 $\kappa \sim -N^2$ at large~$N$\textup{)}.
\item In the large-$N$ matrix-model limit, assuming the ABJM
 Fermi-gas resummation, the shadow connection degenerates
 to the Airy connection and the ABJM partition function is
 $Z(N,k) = C_k^{-1/3}\mathrm{Ai}(C_k^{-1/3}(N - B_k))$
 with $C_k = 2/(\pi^2 k)$.
 The $N^{3/2}$ scaling
 $F \sim (\pi\sqrt{2k}/3)N^{3/2}$ is the hallmark
 of M2-brane degrees of freedom.
\end{enumerate}
\end{proposition}

\begin{remark}[Seiberg--Witten theory and the shadow connection]
\label{rem:holographic-sw}
\index{Seiberg--Witten!shadow connection}
The shadow connection
$\nabla^{\mathrm{sh}} = d - Q'_L/(2Q_L)\,dt$
for the Virasoro algebra is a Fuchsian connection
with three regular singular points, and its Picard--Fuchs
equation is compared with the $\mathrm{SU}(2)$ Seiberg--Witten
Gauss--Manin connection. The identification maps:
shadow metric zeros $\leftrightarrow$ singular fibers
$u = \pm\Lambda^2$;
Koszul monodromy $-1$ $\leftrightarrow$ SW monodromy;
$\kappa = c/2$ $\leftrightarrow$ prepotential coefficient.
Under this comparison, the instanton coefficients $F_k$ of the SW
prepotential are Taylor coefficients of the flat sections of the shadow
connection.
\end{remark}


% ======================================================================
%
% K3 x E: THE COSTELLO-LI PROTOTYPE
%
% ======================================================================

\subsection{\texorpdfstring{$K3 \times E$}{K3 x E}: the Costello--Li prototype}
\label{subsec:k3e-holographic-datum}
\index{K3 surface!twisted holography|textbf}
\index{Costello--Li programme!K3 x E@$K3 \times E$|textbf}
\index{holographic modular Koszul datum!K3 x E@$K3 \times E$|textbf}
\index{Kodaira--Spencer theory!K3 x E@$K3 \times E$}

The product $X = K3 \times E$ of a K3 surface with an elliptic curve
is a compact Calabi--Yau threefold. On the abelian harmonic branch of
the Costello--Li reduction, the boundary algebra is the rank-$24$
Heisenberg algebra, and the corresponding shadow tower has class
$\mathsf{G}$. This branch is a computable geometric test object for
the holographic modular Koszul formalism.

\subsubsection{The HT twist and Kodaira--Spencer reduction}
\label{subsubsec:k3e-ht-twist}
\index{HT twist!K3 x E@$K3 \times E$}

The holomorphic-topological twist of type~IIB on $K3 \times E$
decomposes the CY$_3$ into:
\begin{itemize}
\item \emph{Holomorphic direction:} the elliptic curve $E$
 ($\dim_\C = 1$).
\item \emph{Topological directions:} the K3 surface
 ($\dim_\C = 2$, $\dim_\R = 4$).
\end{itemize}
The B-model twist produces Kodaira--Spencer gravity on $X$, whose
field is the Beltrami differential
$\mu \in \Omega^{0,1}(T^{1,0}_X)$ with action
$S = \int_X \mu \,\dbar\mu + \tfrac{1}{3!}\int_X \mu^3$.
The observables are polyvector fields: by HKR,
\begin{equation}\label{eq:k3e-hh-hkr}
\HH^n(X) = \bigoplus_{p+q=n} \PV^{p,q}(X),
\qquad
\PV^{p,q}(X) = H^q\!\bigl(X, \textstyle\bigwedge^p T^{1,0}_X\bigr)
\overset{\text{CY}}{=} h^{3-p,q}(X).
\end{equation}

The Hodge diamond of $K3 \times E$ is computed by K\"unneth from
\begin{equation}\label{eq:k3e-hodge-kunneth}
h^{p,q}(K3 \times E) = \sum_{a+c=p,\; b+d=q}
h^{a,b}(K3) \cdot h^{c,d}(E).
\end{equation}
The nonzero entries are:
\begin{center}
\renewcommand{\arraystretch}{1.2}
\small
\begin{tabular}{c|cccccccccc}
$(p,q)$ & $(0,0)$ & $(1,0)$ & $(0,1)$ & $(2,0)$ & $(0,2)$ & $(1,1)$ & $(2,1)$ & $(1,2)$
 & $(3,0)$ & $(0,3)$ \\
\hline
$h^{p,q}$ & $1$ & $1$ & $1$ & $1$ & $1$ & $21$ & $21$ & $21$ & $1$ & $1$
\end{tabular}
\end{center}
with the remaining nonzero entries fixed by $h^{p,q} = h^{q,p}
= h^{3-p,3-q}$ (CY$_3$ Hodge symmetry). In particular:
\begin{equation}\label{eq:k3e-hodge-key}
h^{1,1} = h^{2,1} = 21, \qquad
\chi(K3 \times E) = \chi(K3)\cdot\chi(E) = 24 \cdot 0 = 0.
\end{equation}
The total polyvector field dimension is
$\dim \PV^*(X) = \dim \PV^*(K3) \cdot \dim \PV^*(E)
= 24 \cdot 4 = 96$.

\begin{remark}[Self-mirror with $\kappa_{\mathrm{KS}}^{(1)} = 0$]
\label{rem:k3e-self-mirror}
Since $h^{1,1} = h^{2,1} = 21$, the threefold $K3 \times E$
is self-mirror: the KS modular characteristic
$\kappa_{\mathrm{KS}}^{(1)} = h^{1,1} - h^{2,1} = 0$.
The KS weight-one Hodge row vanishes, and the BCOV constant-map
scalar also vanishes because $\chi(K3\times E)=0$. This is the entry
$K3 \times T^2$ in the KS landscape table of
\S\ref*{subsec:thqg-ks-landscape}.

The boundary chiral algebra $\cA_E$ discussed below is a
\emph{different} object from the KS boundary algebra: it arises
from reducing the B-model \emph{along} K3 to obtain a chiral
algebra \emph{on} $E$, rather than from the abstract KS
deformation theory. The two algebras have different $\kappa$
values: $\kappa_{\mathrm{KS}}^{(1)} = 0$ vs.\
$\kappaChHodge(\cA_E) = 24$. The discrepancy records that
$\kappa$ depends on the full algebra, not on a single
invariant of the geometry.
\end{remark}


\subsubsection{The boundary chiral algebra}
\label{subsubsec:k3e-boundary}
\index{boundary chiral algebra!K3 x E@$K3 \times E$}

On the abelian harmonic branch of the Costello--Li reduction along the
topological K3 directions, one obtains a chiral algebra $\cA_E$ on $E$
with generators from $\PV^*(K3)$. Each polyvector class
$\alpha \in \PV^{p,q}(K3) = H^{2-p,q}(K3)$ contributes a generator of
conformal weight $2 - p$:

\begin{center}
\renewcommand{\arraystretch}{1.2}
\small
\begin{tabular}{lccc}
\textbf{Sector} & \textbf{Hodge type} & \textbf{Count}
 & \textbf{Conformal weight} \\
\hline
$\PV^{0,0}(K3) = H^{2,0}(K3)$ & holomorphic 2-form & $1$ & $2$ \\
$\PV^{0,2}(K3) = H^{2,2}(K3)$ & volume form & $1$ & $2$ \\
$\PV^{1,1}(K3) = H^{1,1}(K3)$ & $(1{,}1)$-forms & $20$ & $1$ \\
$\PV^{2,0}(K3) = H^{0,0}(K3)$ & constant & $1$ & $0$ \\
$\PV^{2,2}(K3) = H^{0,2}(K3)$ & anti-hol.\ 2-form & $1$ & $0$ \\
\hline
\multicolumn{2}{c}{\textbf{Total}} & $\mathbf{24}$
 & $= \chi(K3)$
\end{tabular}
\end{center}

On this branch the central charge of the boundary algebra is
$c(\cA_E) = \chi(K3) = 24$.

\begin{proposition}[$K3 \times E$ boundary datum, abelian harmonic branch]
\label{prop:k3e-boundary-datum}
\index{K3 surface!boundary chiral algebra}
Assume the Costello--Li reduction of Kodaira--Spencer gravity on
$K3 \times E$ is taken on the abelian harmonic branch: the harmonic
projector is multiplicative on the K3 polyvector cohomology used in the
reduction, and exact Schouten brackets are killed after projection. The
boundary chiral algebra $\cA_E$ is then modeled by the free-field
algebra of $24$ bosons indexed by $H^*(K3, \Z)$:
\begin{equation}\label{eq:k3e-boundary-ope}
J^a(z)\, J^b(w) \;\sim\; \frac{\delta^{ab}}{(z-w)^2},
\qquad a, b = 1, \ldots, 24,
\end{equation}
% The Heisenberg OPE uses the unit level matrix delta^{ab}
% (positive definite, one boson per cohomology class at level 1).
% In the orthonormal basis, kappa = 24 * kappa(H_1) = 24 by additivity.
% The Mukai pairing of signature (4,20) enters the lattice VOA
% through the vertex operators e^alpha (topological sector), not
% through the Heisenberg OPE levels. In a Mukai basis the OPE
% would read G^{ab}/(z-w)^2, but kappa = rank = 24 regardless of
% basis (kappa is basis-independent).
where the $24$ generators are indexed by a basis of $H^*(K3, \Z)$, each
at unit level. The Mukai pairing of signature $(4, 20)$ determines the
lattice vertex operators and the topological sector, but the Heisenberg
OPE levels are positive definite.
The free-field assertion is the branch condition: after harmonic
projection, the Schouten bracket on the retained K3 sector is zero and
only the Heisenberg pairing remains in the OPE.

\begin{enumerate}[label=\textup{(\roman*)}]
\item $\kappaChHodge(\cA_E) = 24$ on this rank-$24$ Heisenberg branch
 \textup{(}$= \operatorname{rank}$ for unit-level free bosons\textup{)}.
\item Support shadow depth: class $\mathsf{G}$ \textup{(}Gaussian,
 \(r^\Theta_{\max} = 2\)\textup{)}. All higher shadows vanish: $C = Q = 0$.
\item The critical discriminant
\(\Delta = 8\kappaChHodge(\cA_E) S_4 = 0\)
 \textup{(}$S_4 = 0$ for free fields\textup{)}.
\item The shadow metric
\(Q_L(t) = (2\kappaChHodge(\cA_E))^2 = 2304\) is constant.
\end{enumerate}
\end{proposition}

\begin{proof}
The hypotheses impose a multiplicative harmonic projection and kill the
projected Schouten bracket on the retained K3 cohomology sector. The
projected interaction Lie algebra is therefore abelian. The reduction
along K3 produces a cohomological free-field chiral algebra on $E$ with
one Heisenberg boson per cohomology class of K3, each at unit level.
The Mukai lattice of signature $(4, 20)$ underlies the
topological sector: $(3, 19)$ from $H^2(K3)$ plus $(1, 1)$ from the
$H^0 \oplus H^4$ pairing.

For the modular characteristic: $\cA_E$ consists of 24 free
bosons at level~$1$, so
\(\kappaChHodge(\cA_E)=24 \cdot \kappaChHodge(\cH_1) = 24\)
by additivity (Volume~I, Corollary~\ref{V1-cor:kappa-additivity}).
The support shadow depth is \(r^\Theta_{\max}=2\) because the OPE is
purely quadratic and the canonical Gaussian gauge slice has no higher
lift coordinates.
\end{proof}


\subsubsection{The holographic modular Koszul datum
\texorpdfstring{$\cH(K3 \times E)$}{H(K3 x E)}}
\label{subsubsec:k3e-holographic-datum}

\begin{computation}[$K3 \times E$ holographic datum on the abelian harmonic branch]
\label{comp:k3e-holographic-datum}
\index{holographic modular Koszul datum!K3 x E@$K3 \times E$!components}
Under the hypotheses of Proposition~\ref{prop:k3e-boundary-datum},
the six components of
$\cH(K3 \times E) = (\cA, \cA^!, \mathcal{C}, r(z),
\Theta_\cA, \nabla^{\mathrm{hol}})$ are:
\begin{center}
\renewcommand{\arraystretch}{1.35}
\small
\begin{tabular}{lp{9.5cm}}
\textbf{Component} & \textbf{Value} \\
\hline
(a) $\cA$ & Free-field: $24$ bosons at unit level ($\delta^{ab}$),
 $c = 24$, $\kappa = 24$ \\
(b) $\cA^!$ & Verdier--Koszul partner: the chiral symmetric algebra on
 the dual current space, not the negative-level Heisenberg algebra;
 $\kappaChHodge(\cA^!) = -24$ and
 $\kappa + \kappa' = 0$ \\
(c) $\mathcal{C}$ & $\cA^!\text{-}\mathsf{mod}$;
 the derived center $\Zder^{\mathrm{ch}}(\cA_E)$
 has primitive Heisenberg Hochschild row
 $P_{\cA}^{\mathrm{prim}}(t) = 1 + 24\,t + t^2$ \\
(d) $r(z)$ & $r(z) = k\,\Omega_{\mathrm{tr}}/z$ at unit level $k = 1$,
 $\Omega = \sum \delta^{ab}\,J_a \otimes J_b$ CYBE trivially satisfied \textup{(}abelian\textup{)} \\
(e) $\Theta_\cA$ & $\Theta_\cA^{\leq 2} = \kappa = 24$;
 $\Theta_\cA^{(g,r)} = 0$ for $r \geq 3$
 \textup{(}class $\mathsf{G}$: tower terminates\textup{)} \\
(f) $\nabla^{\mathrm{hol}}$ & Trivial on the shadow metric
 \textup{(}$Q_L$ constant\textup{)};
 nontrivial at the holographic level with
 universal residue $\tfrac{1}{2}$
\end{tabular}
\end{center}

The genus-$g$ free energies are
$F_g(\cA_E) = \kappaChHodge(\cA_E) \cdot \lambda_g^{\mathrm{FP}}$
\textup{(}uniform-weight\textup{)}:
\begin{equation}\label{eq:k3e-genus-amplitudes}
F_1 = 1, \qquad
F_2 = \tfrac{7}{240}, \qquad
F_3 = \tfrac{31}{40320}.
\end{equation}

\emph{Three consistency checks.}
\begin{enumerate}
\item \emph{KS reduction}: The reduction of KS theory along K3 gives
 $24$ free bosons and $\kappa = \operatorname{rank} = 24$.
\item \emph{HKR comparison}: The equality $\dim \PV^*(K3) = 24$
 supplies the current row. The group $\HH^2(K3 \times E)$ has dimension
 $23$ in the geometric category; the chiral boundary center is the
 primitive Heisenberg row above.
\item \emph{Koszul complementarity}: $\kappaChHodge(\cA^!) = -24$;
 $F_1(\cA) + F_1(\cA^!) = 1 + ({-}1) = 0$ \textup{(}Theorem~C
 for free fields\textup{)}.
\end{enumerate}
\end{computation}


\subsubsection{Noncommutative deformations and the Brauer group}
\label{subsubsec:k3e-nc-deformations}
\index{noncommutative deformation!K3 x E@$K3 \times E$|textbf}
\index{Brauer group!K3 x E@$K3 \times E$}
\index{Hochschild cohomology!K3 x E@$K3 \times E$}

The deformation theory of $D^b(K3 \times E)$ is controlled by
$\HH^2(K3 \times E) = 23$, decomposed by HKR as:
\begin{equation}\label{eq:k3e-hh2-decomposition}
\HH^2(K3 \times E)
\;=\;
\underbrace{H^1(T_X)}_{h^{2,1} = 21}
\;\oplus\;
\underbrace{H^0\!\bigl(\textstyle\bigwedge^2 T_X\bigr)}_{h^{1,0} = 1}
\;\oplus\;
\underbrace{H^2(\cO_X)}_{h^{0,2} = 1}.
\end{equation}
The $21$-dimensional summand parametrizes complex-structure
deformations. The two remaining summands are the Poisson and gerby
rows of the HKR decomposition:
\begin{enumerate}[label=\textup{(\roman*)}]
\item The bivector direction $H^0(\bigwedge^2 T_X) \cong \C$
 is the first-order Poisson deformation; for $K3\times E$ it is
 generated by the inverse holomorphic symplectic form on K3.
\item The gerby $B$-field direction $H^2(\cO_X) \cong \C$ is the
 infinitesimal Brauer-twisting row. It is not a rescaling of the
 holomorphic volume form.
\end{enumerate}

Deformation quantization of $K3 \times E$ is one-parameter:
$H^0(\bigwedge^2 T_{K3 \times E}) = \C$, generated entirely by the
K3 symplectic form $\omega_{K3}^{-1} \in H^0(\bigwedge^2 T_{K3})$.
The elliptic curve contributes no Poisson structure
($\bigwedge^2 T_E = 0$ for $\dim_\C E = 1$), and $H^0(T_{K3}) = 0$
eliminates mixed terms. The star product is the Moyal product
in local Darboux coordinates on K3:
\begin{equation}\label{eq:k3e-star-product}
f \star_\hbar g
= f g + \hbar\,\{f, g\}_{K3}
+ \sum_{n \geq 2} \frac{\hbar^n}{n!}\,
\omega^{i_1 j_1} \cdots \omega^{i_n j_n}\,
(\partial_{i_1} \cdots \partial_{i_n} f)\,
(\partial_{j_1} \cdots \partial_{j_n} g).
\end{equation}

The analytic Brauer torsion is read from the exponential sequence. Since
$H^1(K3,\cO_{K3})=0$ and $\operatorname{Br}(E)=0$, there is no elliptic
or mixed Brauer summand. If the K3 factor has Picard rank $\rho$, then
\begin{equation}\label{eq:k3e-brauer}
\operatorname{Br}(K3 \times E)_{\mathrm{tors}}
\;\cong\;
(\Q/\Z)^{22-\rho}.
\end{equation}
For algebraic K3 surfaces one has $\rho\geq 1$; the Picard-rank-zero
case is the analytic non-projective limit.

\begin{remark}[Geometric Hochschild comparison]
\label{rem:k3e-chiral-hochschild}
\index{Hochschild cohomology!K3 x E@$K3 \times E$}
The following diamond is the HKR diamond of the geometric category
$D^b(K3\times E)$. It is a comparison target for the CDR and
factorization models, not an identification of
$\operatorname{ChirHoch}^\bullet(\Omega^{\mathrm{ch}}_{K3\times E})$
with geometric Hochschild cohomology. The full geometric Hochschild
diamond is
\begin{equation}\label{eq:k3e-hh-full}
(\HH^0, \HH^1, \HH^2, \HH^3, \HH^4, \HH^5, \HH^6)
= (1, 2, 23, 44, 23, 2, 1),
\end{equation}
with CY Serre duality $\HH^n \cong \HH^{6-n}$ verified computationally
via the Volume~III K3$\times$E noncommutative-deformation engine.
The boundary algebra $\cA_E$ used above is the rank-$24$ harmonic
reduction; its primitive chiral Hochschild row is
$1+24t+t^2$, not this geometric HKR diamond.
\end{remark}


\subsubsection{Swiss-cheese structure and the CDR shadow}
\label{subsubsec:k3e-swiss-cheese}
\index{Swiss-cheese!K3 x E@$K3 \times E$}
\index{chiral de Rham complex!K3 x E@$K3 \times E$}

The chiral de Rham complex $\Omega^{\mathrm{ch}}(K3 \times E)$
is a sheaf of vertex superalgebras on $K3 \times E$ whose
global sections carry an $\cN = 2$ superconformal structure
(Malikov--Schechtman--Vaintrob). Put
\begin{equation}\label{eq:k3-cdr-sigma-model}
H^*\!\bigl(K3, \Omega^{\mathrm{ch}}_{K3}\bigr)
\;=:\; V_{K3}
\qquad (c = 6,\; \cN = (4,4)\;\text{SCVA}).
\end{equation}
The exterior-product K\"unneth map for CDR sheaves gives
\begin{equation}\label{eq:k3e-cdr-kunneth}
H^*\!\bigl(K3 \times E, \Omega^{\mathrm{ch}}\bigr)
\;\cong\;
V_{K3} \otimes V_E,
\end{equation}
where $V_E=H^*(E,\Omega_E^{\mathrm{ch}})$ is the elliptic chiral de
Rham factor. The MSV central charges are $6$ on K3 and $3$ on~$E$, so
the untwisted product has $c=9$. The rank-$24$ Heisenberg boundary
algebra is not this full CDR factor; it is the harmonic K3 reduction
of Proposition~\ref{prop:k3e-boundary-datum}.

The HT twist changes the stress tensor by the $U(1)$ current and then
passes to BRST cohomology. It does not identify the untwisted CDR
central charge with zero. On the abelian harmonic branch, the boundary
shadow operations of arity $k\geq 3$ vanish on bar cohomology, because
$\cA_E$ is class~$\mathsf{G}$. No formality assertion is made for the
full CDR Swiss-cheese factorization algebra.

\begin{example}[K3 bulk-boundary pair as an \texorpdfstring{$\SCchtop$}{SCchtop} realisation]
\label{ex:k3e-scchtop-bulk-boundary}
The K3 specialisation of $\SCchtop$ is not concentrated at a single
smooth fibre. Its open colour is the harmonic boundary algebra
$\cA_E$ of Proposition~\ref{prop:k3e-boundary-datum}, obtained after
projecting the K3 factor before restriction to~$E$. Its closed colour
is the bulk-boundary pair
\[
 \bigl(C^\bullet_{\mathrm{ch}}(\cA_E,\cA_E),\, \cA_E\bigr)
\]
carried over the bi-based host
\[
 \Ran(E^{\mathrm{nod,sm}}_{24}) \times \overline{\mathcal{A}}_2.
\]
The index $24$ is the Euler count: an elliptic K3 has Euler number
$24$, and in a nodal model each nodal fibre contributes one to that
count, so the generic smoothing problem has exactly twenty-four nodal
walls. The nearby-cycles sector records transport
across these twenty-four walls, while the $\overline{\mathcal{A}}_2$
parameter records the Humbert bifurcation on which the Siegel form
$\Delta_5$ acquires its first monodromy wall
\textup{(}Gritsenko--Nikulin 1998, Theorem~2.1\textup{)}. In this form
the K3 bulk-boundary pair is an honest worked $\SCchtop$ example:
holomorphic fusion happens on $\Ran(E^{\mathrm{nod,sm}}_{24})$, ordered
topological propagation happens on the boundary interval, and the bulk
is seen through nearby cycles rather than through a single smooth
parameter value.
\end{example}

\begin{remark}[Factorization homology and the elliptic genus]
\label{rem:k3e-factorization-homology}
\index{factorization homology!K3 x E@$K3 \times E$}
\index{elliptic genus!K3 x E@$K3 \times E$}
The CDR Euler characteristic gives
$\chi(\Omega^{\mathrm{ch}}(K3 \times E)) = \chi(K3) \cdot \chi(E)
= 24 \cdot 0 = 0$: the Witten genus vanishes. This
vanishing of the elliptic genus does \emph{not} imply
$F_g(\cA_E) = 0$: the elliptic genus is a
\emph{refined} trace with $(-1)^F$ insertion, while
$F_g(\cA_E) = \kappaChHodge(\cA_E) \cdot \lambda_g^{\mathrm{FP}}$
\textup{(}uniform-weight\textup{)} is the bar complex
curvature at genus~$g$ (no grading insertion). The BPS partition
function is related to $(\Phi_{10}^{\mathrm{un}})^{-1}$, the reciprocal
of the unnormalised Igusa cusp form (DMVV formula for
$K3 \times E_\tau$). In the K3 Borcherds normalisation
\[
 \Phi_{10}^{\mathrm{un}}=\Delta_5^2,\qquad
 \mathrm{wt}(\Delta_5)=c_{\phi_{0,1}^{K3}}(0)/2=10/2=5 .
\]
This automorphic/BPS datum is external to the scalar FP genus tower; the
shadow
partition function $Z^{\mathrm{sh}} = \exp(\sum F_g \hbar^{2g})$
is the perturbative shadow of this non-perturbative BPS count.
\end{remark}


\subsubsection{Connection to Volume~I}
\label{subsubsec:k3e-vol1-connection}

On the abelian harmonic branch, the $K3 \times E$ example supplies the
following Volume~I comparison entries:

\begin{center}
\renewcommand{\arraystretch}{1.25}
\small
\begin{tabular}{lll}
\textbf{Vol~I result} & \textbf{$K3 \times E$ instance} & \textbf{Value} \\
\hline
Theorem~A (adjunction) & $D_{\mathrm{Ran}}(B(\cA_E))
 \simeq B(\cA_E^!)$ & free-field Verdier \\
Theorem~B (completed Positselski) & $\Omega(B(\cA_E)) \simeq \cA_E$
 in the coderived/completed ambient; the raw counit is Theorem~A's
 inversion face & immediate (class $\mathsf{G}$) \\
Theorem~C (complementarity) &
 $\kappaChHodge(\cA_E)+\kappaChHodge(\cA_E^!) = 0$;
 $F_g + F_g' = 0$ for all $g$ & $24 + ({-}24) = 0$ \\
Theorem~D (leading coeff.) & $F_g(\cA_E) =
 \kappaChHodge(\cA_E)\lambda_g^{\mathrm{FP}}$
 \textup{(}u.w.\textup{)} & $\kappaChHodge(\cA_E)=24$ \\
Theorem~H (Hochschild) & $P_{\cA}^{\mathrm{prim}}(t) =
 1 + 24\,t + t^2$ & $\HH^0 = 1,\; \HH^1 = 24,\; \HH^2 = 1$ \\
\hline
Shadow class & $\mathsf{G}$ (Gaussian) & \(r^\Theta_{\max} = 2\) \\
Shadow metric & \(Q_L = (2\kappaChHodge(\cA_E))^2 = 2304\) & constant \\
MC element & \(\Theta_\cA = \Theta_\cA^{\leq 2}
= \kappaChHodge(\cA_E)\)
 & terminates at degree $2$ \\
$r$-matrix & $r(z) = k\,\Omega_{\mathrm{tr}}/z$ \textup{(}abelian Casimir, $k=1$\textup{)} &
 single pole, $\dim = 24$
\end{tabular}
\end{center}

\noindent
The moduli space of K3 has dimension~$20$ (complex structure
deformations from $h^{1,1}(K3)$), plus $\dim \cM_E = 1$
and $21$ K\"ahler moduli, for a total of~$42$.
The period domain for K3 complex structures is a
type~$\mathrm{IV}_{20}$ bounded symmetric domain of
complex dimension~$20$, with arithmetic group
$O(\Lambda_{K3}) \cong O(3, 19; \Z)$.
The holographic shadow connection $\nabla^{\mathrm{hol}}$ on the
K3 moduli space is trivial at the shadow metric level
($Q_L$ constant for class~$\mathsf{G}$), but acquires nontrivial
monodromy at the full holographic level from the arithmetic structure
of the K3 period domain.

\begin{remark}[Comparison across CY$_3$ geometries]
\label{rem:k3e-cy3-comparison}
The $K3 \times E$ datum contrasts with the other standard CY$_3$
targets as follows:
\begin{center}
\renewcommand{\arraystretch}{1.25}
\small
\begin{tabular}{lp{0.27\linewidth}p{0.20\linewidth}p{0.19\linewidth}}
\textbf{CY$_3$ datum} & \textbf{Scalar row}
 & \textbf{Depth assertion} & \textbf{Infinitesimal row} \\
\hline
$K3 \times E$, abelian harmonic branch
& $\kappaChHodge=24$ & $\mathsf{G}$ & geometric HKR: $\dim \HH^2=23$ \\
Quintic, KS/BCOV lane
& $\kappa_{\mathrm{KS}}^{(1)}=-100,\;
   \kappa_{\mathrm{BCOV}}=-25/3$
& scalar projection only & $h^{2,1}=101$ \\
D3 ($\cN\!=\!4$ SYM)
& $\frac{N(N+1)}{2}$ & $\mathsf{L}$ & current row \\
ABJM & $-(N^2{+}1)$ & $\mathsf{M}$ & monopole row
\end{tabular}
\end{center}
The $K3 \times E$ row is the free-field Calabi--Yau row in the
table. The quintic row is a Kodaira--Spencer/BCOV row: its
$-100$ is the Hodge supertrace $h^{1,1}-h^{2,1}$, and its one-loop
BCOV scalar is $\chi/24=-25/3$. The full quintic boundary algebra still
contains Schouten and BCOV collision operations, so class~$\mathsf{G}$
is not assigned to it. In the uniform-weight scalar FP model, a
positive scalar characteristic gives nonzero scalar genus coefficients
only after the scalar projection. This is neither an elliptic genus nor
a BPS-state count.
\end{remark}

\begin{remark}[Bulk OPE structure constants via the Gerstenhaber bracket]
\label{rem:bulk-ope-gerstenhaber-bracket}
The derived chiral center $Z^{\mathrm{der}}_{\mathrm{ch}}(\cA)$
carries a Gerstenhaber bracket $\{-,-\}$ encoding the bulk OPE;
this bracket is the chiral analogue of the classical Hochschild
bracket on $\operatorname{ChirHoch}^*(\cA)$.
The shadow depth classification ($\mathsf{G}$/$\mathsf{L}$/$\mathsf{M}$)
controls the truncation level of this bracket.
We record the structure constants for three families.

\emph{Heisenberg $H_k$ (class~$\mathsf{G}$, shadow depth $r=2$).}
The space $\operatorname{ChirHoch}^*(H_k)$ has three generators: the
unit $\mathbf{1}$ in degree~$0$, the current $J$ in degree~$1$,
and the normal-ordered product ${:}JJ{:}$ in degree~$2$.
The Gerstenhaber bracket is determined at degree~$2$:
\[
 \{J, J\} \;=\; k,
\]
where $\kappa^{\mathrm{Heis}} = k$ is the Heisenberg level.
All higher brackets vanish; the only nonzero binary bracket is central.
The bulk dg Lie structure is therefore two-step nilpotent, and it
becomes abelian only at $k=0$.

\emph{Affine $V_k(\mathfrak{sl}_2)$ (class~$\mathsf{L}$, shadow depth $r=3$).}
The bar cohomology satisfies $\dim H^2(B(\mathfrak{sl}_2)) = 5$,
and $\kappa^{\mathrm{KM}}(V_k(\mathfrak{sl}_2)) = 3(k+2)/4$.
The Gerstenhaber bracket on
$Z^{\mathrm{der}}_{\mathrm{ch}}(V_k(\mathfrak{sl}_2))$ is
non-abelian:
\[
 \{J^a, J^b\}
 \;=\; f^{ab}{}_{c} \, J^c
 \;+\; k \, \delta^{ab},
\]
combining the simple-pole (structure constant) and double-pole
(level) contributions from the OPE
$J^a(z) J^b(w) \sim
 k\,\delta^{ab}/(z-w)^2 + f^{ab}{}_{c}\,J^c(w)/(z-w)$.
A non-trivial degree-$3$ component arises from the cubic Casimir
interaction, scaled by a rational function of~$k$.
The tower terminates: all brackets of degree $\geq 4$ are determined
by the lower ones.
At the critical level $k = -2$: $\kappa^{\mathrm{KM}} = 0$, the
center of the completed enveloping algebra jumps
(Feigin--Frenkel), and the Gerstenhaber bracket changes
qualitatively.

\emph{Virasoro $\mathrm{Vir}_c$ (class~$\mathsf{M}$, shadow depth $r = \infty$).}
$\operatorname{ChirHoch}^*(\mathrm{Vir}_c)$ is concentrated in
degrees $\{0,1,2\}$ with total dimension at most~$4$; it is neither
$\mathbf{C}[\Theta]$ nor the Gelfand--Fuchs cohomology (which is
infinite-dimensional).
The leading structure constant is $\kappa^{\mathrm{Vir}} = c/2$,
arising from the $T(z)T(w)$ OPE via $\mathrm{d}\log$-absorption
(OPE quartic pole $\to$ $r$-matrix cubic pole:
\(r_{\mathrm{Vir}}^{\mathrm{coll}}(z) = (c/2)/z^3 + 2T/z\)).
The bracket has non-vanishing components at every degree $m_k$,
$k \geq 2$: despite bounded cohomological amplitude, the algebraic
depth is infinite (class~$\mathsf{M}$).
At $c = 13$ (the self-dual point):
$\kappa^{\mathrm{Vir}} + \kappa^{\mathrm{Vir}!} = 13$.
At $c=0$: the bracket degenerates.

The pattern is sharp: shadow depth controls truncation of the
Gerstenhaber bracket on the derived chiral center.
Class~$\mathsf{G}$ brackets close at degree~$2$;
class~$\mathsf{L}$ at degree~$3$;
class~$\mathsf{M}$ never close.
This is consistent with Theorem~H, which bounds the cohomological
amplitude of $\operatorname{ChirHoch}^*(\cA)$ uniformly, while the
\emph{algebraic} depth of the bracket structure varies by family.
\end{remark}

% ================================================================
% 9.8.19: UNIVERSAL HOLOGRAPHY (master theorem)
% ================================================================

\subsection{Universal holography}\label{subsec:universal-holography}
\index{universal holography!master theorem}
\index{holomorphic-topological theory!functor from chiral algebras}

Fix an admissible boundary datum in the sense of
Theorem~\ref{thm:holographic-recon-shadow}: a chiral algebra $\cA$ on a
smooth projective curve $X$ with conformal vector $\omega$ at
non-critical level, together with a boundary-linear, affine, or DS
realisation and the ambient needed for the bulk--centre comparison.
The physical $3$d theory is not determined by this algebraic datum
alone.  A holographic object includes a realization package
\[
\Xi_{\cA}=(\mathcal F_{\cA,\Xi},\eta^\partial_{\cA},
\chi_{\mathrm{HT},\cA},\mathcal A(\cA)),
\]
where $\mathcal F_{\cA,\Xi}$ is a Costello--Gwilliam
holomorphic-topological BV/factorization model on the slab,
\(\eta^\partial_{\cA}\colon
\Obs^\partial(\mathcal F_{\cA,\Xi})\xrightarrow{\simeq}\cA\)
is the boundary comparison, and
\(\chi_{\mathrm{HT},\cA}\colon
\Zderch(\cA)\to\Obs^{\mathrm{bulk}}(\mathcal F_{\cA,\Xi})\)
is the bulk--Hochschild comparison in the declared ambient
\(\mathcal A(\cA)\).  On the $3$d slab $X\times\R$, the associated
theory is the supplied BV factorisation algebra
\begin{equation}\label{eq:universal-holography-construction}
\mathcal{T}_{\cA,\Xi}
\;:=\;
\mathcal F_{\cA,\Xi}
\quad\text{on}\quad X\times\R_{\ge0},
\qquad
\eta^\partial_\cA\colon
\mathrm{Obs}^{\partial}(\mathcal{T}_{\cA,\Xi})
\xrightarrow{\simeq}\cA .
\end{equation}
On the boundary component $X\times\{0\}$ impose the chosen holomorphic
boundary condition whose mode algebra is $\cA$; in the bulk
$X\times(0,\infty)$ run the corresponding Costello--Li or
Costello--Gaiotto holomorphic Chern--Simons BV quantisation contained
in \(\Xi_\cA\). When the
topologisation datum is present, the improvement
$T^{\mathrm{imp}}=[Q_{\mathrm{tot}},G']$ makes holomorphic translations
$Q$-exact
(Theorems~\ref{thm:E3-topological-km},
\ref{thm:E3-topological-DS},
\ref{thm:E3-topological-DS-general}).
This is $\mathcal{T}_{\cA,\Xi}$ relative to the chosen admissible input.

Four facts hold on this locus: the boundary recovers $\cA$; the bulk
is compared with the chiral derived centre
$Z^{\mathrm{der}}_{\mathrm{ch}}(\cA)$ by
\(\chi_{\mathrm{HT},\cA}\) in the algebraic, cohomological, or
completed/pro ambient specified by the input; the assignment is
functorial for morphisms preserving the chosen realisation and
comparison maps; and the standard shadow classes are covered by the
boundary-linear, affine, and DS constructions on the corresponding
\(\Xi\)-loci.

\begin{theorem}[$\Xi$-realized universal holography on the admissible boundary-linear/DS locus;
\(\ClaimStatusConditional\); licensing tags \(\beta+\gamma+\delta\)]%
\label{thm:holographic-recon-class-m}%
\index{universal holography!Theorem|textbf}%
\index{holomorphic-topological theory!universal functor}%
\index{Drinfeld--Sokolov reduction!holographic functor}%
Let $X$ be a smooth projective curve and let
$\mathrm{ChirAlg}^{\omega,\mathrm{adm},\Xi}_X$ denote the category whose
objects are conformal chiral algebras on $X$ at non-critical level
together with the admissible boundary-linear, affine, or DS datum of
Theorem~\ref{thm:holographic-recon-shadow} and a realization package
\(\Xi_\cA=(\mathcal F_{\cA,\Xi},\eta^\partial_{\cA},
\chi_{\mathrm{HT},\cA},\mathcal A(\cA))\); morphisms preserve the
boundary comparison, the bulk--Hochschild comparison, and the ambient.
Let $\mathrm{HT\text{-}QFT}^{\Xi}_{X\times\R}$ denote the category of
$3$d holomorphic-topological BV theories on $X\times\R$ equipped with
these comparison data. For every
$(\cA,\omega,\mathrm{adm},\Xi_\cA)\in
\mathrm{ChirAlg}^{\omega,\mathrm{adm},\Xi}_X$ in the generic/Koszul
exact locus, with class~$\mathsf M$ read in the cohomological or
weight-completed/pro ambient unless the bounded-shift HPL criterion is
supplied, the supplied model
$\mathcal{T}_{\cA,\Xi}:=\mathcal F_{\cA,\Xi}$ satisfies:
\begin{enumerate}[label=\textup{(\roman*)}]
\item\label{item:uh-boundary}
 $\eta^\partial_\cA\colon
 \mathrm{Obs}^{\partial}(\mathcal{T}_{\cA,\Xi})\xrightarrow{\simeq}\cA$
 as
 $E_1$-chiral algebras on $X$;
\item\label{item:uh-bulk}
	 $\chi_{\mathrm{HT},\cA}\colon
	 Z^{\mathrm{der}}_{\mathrm{ch}}(\cA)\xrightarrow{\simeq}
	 \mathrm{Obs}^{\mathrm{bulk}}(\mathcal{T}_{\cA,\Xi})$
	 is the bulk--centre comparison in the stated ambient; its strict
	 $E_3$-topological realisation is asserted only when the
	 topologisation hypothesis of
	 Convention~\textup{\ref{conv:chd-e3-chiral-meaning}} is part of
	 the admissible datum;
\item\label{item:uh-functor}
 the assignment
 $\Phi_{\mathrm{hol}}\colon
 (\cA,\omega,\mathrm{adm},\Xi_\cA)\mapsto\mathcal{T}_{\cA,\Xi}$
 is a functor
	 $\mathrm{ChirAlg}^{\omega,\mathrm{adm},\Xi}_X\to
	 \mathrm{HT\text{-}QFT}^{\Xi}_{X\times\R}$
	 intertwining Drinfeld--Sokolov reduction only for
	 \(\Xi\)-compatible DS data:
 $\Phi_{\mathrm{hol}}(\mathrm{DS}\,V^k(\fg),\Xi_{\mathrm{DS}})
 =
 \mathrm{DS\text{-}boundary}\,
 \Phi_{\mathrm{hol}}(V^k(\fg),\Xi_{V^k})$;
\item\label{item:uh-coverage}
	 class coverage is complete on the admissible locus: class
	 $\mathsf{G}$ and the free-field part of class $\mathsf{C}$ are
	 realised by Costello--Li abelian holomorphic Chern--Simons, class
	 $\mathsf{L}$ by affine holomorphic Chern--Simons
	 \textup{(}Theorem~\textup{\ref{thm:E3-topological-km}}\textup{)},
	 and class $\mathsf{M}$ by Costello--Gaiotto plus the
	 DS--Hochschild bridge in its cohomological or
	 weight-completed/pro form
	 \textup{(}Theorem~\textup{\ref{thm:ds-hochschild-bridge}}\textup{)}.
\end{enumerate}
\end{theorem}

\begin{proof}
\emph{Step} (i) \emph{-- boundary recovery.}
By construction~\eqref{eq:universal-holography-construction} the
boundary observables on $X\times\{0\}$ are the holomorphic boundary
cochains of the chosen BV theory, and the comparison with \(\cA\) is
the supplied map \(\eta^\partial_\cA\). For the DS route this is
Costello--Gaiotto~\cite[\S 4]{CostelloGaiotto18}; for boundary-linear
algebras it is Costello--Li~\cite{CL16}; for Kac--Moody it is the
holomorphic Chern--Simons boundary
(Theorem~\ref{thm:E3-topological-km}).
Combining these across shadow classes yields
$\eta^\partial_\cA\colon
\mathrm{Obs}^{\partial}(\mathcal{T}_{\cA,\Xi})\simeq\cA$ as
$E_1$-chiral algebras on $X$.

\emph{Step} (ii) \emph{-- bulk--centre comparison.}
The algebraic closed-sector object determined by the boundary algebra
is
\[
Z^{\mathrm{der}}_{\mathrm{ch}}(\cA)
=C^\bullet_{\mathrm{ch}}(\cA,\cA)
\]
in the ambient stated in the theorem.  The physical bulk factorisation
algebra on $X\times\R$ is the chosen
\(\Obs^{\mathrm{bulk}}(\mathcal{T}_{\cA,\Xi})\).  The bridge between
them is the comparison map \(\chi_{\mathrm{HT},\cA}\), not a
construction of a physical bulk from \(\cA\) alone.  In the
boundary-linear exact sector,
Theorem~\ref{thm:boundary-linear-bulk-boundary} supplies the
Morita/derived-centre comparison.  In the physical prefactorisation
scope, Theorem~\ref{thm:bulk_hochschild} supplies
\(\Psi_{\mathrm{Ran}}\), equivalently \(\chi_{\mathrm{HT},\cA}\).  In
the DS sector, Theorem~\ref{thm:ds-hochschild-bridge} gives the
cohomological or completed/pro comparison, with original-complex
chain-level comparison only under the bounded-shift HPL criterion.
If the topologisation datum of
Convention~\ref{conv:chd-e3-chiral-meaning} is present, then
Theorem~\ref{thm:chiral-higher-deligne} upgrades the algebraic
closed-sector comparison to the strict $E_3$-topological
factorisation-algebra comparison.

\emph{Step} (iii) \emph{-- functoriality and DS intertwining.}
A morphism $\varphi\colon\cA\to\cB$ in
$\mathrm{ChirAlg}^{\omega,\mathrm{adm},\Xi}_X$ determines a boundary
morphism between the chosen Costello BV complexes compatible with
\(\eta^\partial\), \(\chi_{\mathrm{HT}}\), and the ambient; the
induced map of algebraic closed-sector objects is the chiral
Hochschild/coderivation comparison supplied as part of the
\(\Xi\)-compatible morphism data; it is not automatic functoriality of
cochains along the algebra map \(\varphi\).
Intertwining with DS is the content of the cross-volume bridge
``DS-bar'' (Volume~I, \texttt{ds-koszul-intertwine}) composed with
the DS--Hochschild bridge~\ref{thm:ds-hochschild-bridge}:
$\Phi_{\mathrm{hol}}\circ\mathrm{DS}=\mathrm{DS}\circ\Phi_{\mathrm{hol}}$
on the \(\Xi\)-compatible locus, because DS is defined as a boundary
condition on holomorphic Chern--Simons and commutes with the supplied
bulk BV quantisation and comparison maps
(Costello--Gaiotto~\cite[\S 5]{CostelloGaiotto18}).

\emph{Step} (iv) \emph{-- class coverage.}
For class~$\mathsf{G}$ (Heisenberg): abelian holomorphic Chern--Simons
on $X\times\R$ with Lie-algebra the Cartan of the lattice; the chosen
bulk observables are the symmetric algebra on $H^{0,1}(X)$ shifted,
and \(\chi_{\mathrm{HT},\mathcal H_k}\) compares them with
$Z^{\mathrm{der}}_{\mathrm{ch}}(\mathcal{H}_k)$.
For class~$\mathsf{L}$ (affine Kac--Moody):
Theorem~\ref{thm:E3-topological-km}.
For class~$\mathsf{C}$ ($\beta\gamma$, symplectic fermion): contact
$W$-algebras realised as Costello--Li holomorphic boundary conditions
on abelian $3$d HT.
For class~$\mathsf{M}$ (Virasoro, $\cW_N$): Costello--Gaiotto construct
the $3$d HT theory whose boundary is $\cW^k(\fg)$ for principal DS;
the comparison
\(\chi_{\mathrm{HT},\cW^k(\fg)}\) identifies its physical bulk with
$Z^{\mathrm{der}}_{\mathrm{ch}}(\cW^k(\fg))$ cohomologically, and at
chain level only in the weight-completed/pro ambient unless the bounded
conformal-weight-shift HPL criterion is verified, by the
DS--Hochschild bridge~\ref{thm:ds-hochschild-bridge}.
The union of the four sub-constructions defines $\Phi_{\mathrm{hol}}$
on $\mathrm{ChirAlg}^{\omega,\mathrm{adm},\Xi}_X$.
\end{proof}

\begin{corollary}[Class~$\mathsf{M}$ cohomological and completed bulk--centre comparison;
\(\ClaimStatusConditional\); licensing tags \(\beta+\gamma+\delta\)]%
\label{thm:class-M-chain-bulk}%
\index{class M@class~$\mathsf{M}$!chain-level bulk--centre comparison|textbf}%
\index{W-algebras@$\cW$-algebras!bulk--centre comparison}%
\index{Virasoro algebra!bulk--centre comparison}%
Let $\cA=\cW^k(\fg)$ be a $\cW$-algebra obtained by principal
Drinfeld--Sokolov reduction of $V^k(\fg)$ at non-critical level, and
equip it with a class-\(\mathsf M\) realization datum
\(\Xi_{\cA}^{\mathrm{DS}}\) consisting of the Costello--Gaiotto
DS-boundary model, the boundary comparison \(\eta^\partial_\cA\), the
bulk--Hochschild comparison \(\chi_{\mathrm{HT},\cA}\), and the
weight-completed/pro ambient.  Then
\[
Z^{\mathrm{der}}_{\mathrm{ch}}(\cA)
\;\simeq\;
\mathcal{O}\bigl(M_{\mathrm{vac}}^{\mathrm{der}}\bigr)
\]
is the DS algebraic closed-sector comparison cohomologically.  The
physical bulk comparison is the supplied map
\[
\chi_{\mathrm{HT},\cA}\colon
Z^{\mathrm{der}}_{\mathrm{ch}}(\cA)
\xrightarrow{\simeq}
\mathrm{Obs}^{\mathrm{bulk}}(\mathcal{T}_{\cA,\Xi}).
\]
It is chain-level in the weight-completed/pro ambient, and on the
original complex only after the bounded conformal-weight-shift HPL
criterion is supplied.  Thus the class-\(\mathsf M\) global triangle is
a triangle of comparison maps in the declared ambient, not a
reconstruction of a physical bulk from \(\cA\) alone.
\end{corollary}

\begin{proof}
Apply Theorem~\ref{thm:holographic-recon-class-m}\ref{item:uh-bulk} to
$\cA=\cW^k(\fg)$ equipped with \(\Xi_{\cA}^{\mathrm{DS}}\).  The
physical comparison is
\(\chi_{\mathrm{HT},\cA}\colon
Z^{\mathrm{der}}_{\mathrm{ch}}(\cA)\to
\mathrm{Obs}^{\mathrm{bulk}}(\mathcal{T}_{\cA,\Xi})\).
Compose its source with the DS--Hochschild bridge
Theorem~\ref{thm:ds-hochschild-bridge}:
\[
Z^{\mathrm{der}}_{\mathrm{ch}}(\cW^k(\fg))
\;\xrightarrow{\;\mathrm{HKR}^{\mathrm{DS}}\;\simeq\;}\;
\mathcal{O}\bigl(M_{\mathrm{vac}}^{\mathrm{der}}\bigr),
\]
  cohomologically under the smooth chiral-HKR/derived-vacuum package
  and the completed DS coderivation transfer of
  Theorems~\ref{thm:chd-ds-hochschild} and
  \ref{thm:ds-hochschild-bridge}.  The finite input is the
  Kazhdan/PBW finite-weight profile of the BRST complex; generic
  non-critical level is not being used as a \(C_2\)-cofinite
  hypothesis. The same comparison is chain-level in the
  weight-completed/pro ambient, and on the original complex only after
  the bounded conformal-weight-shift HPL criterion of
  Remark~\ref{rem:class-M-hpl-convergence}; this is precisely the scope
  of Theorem~\ref{thm:ds-hochschild-bridge}.
The class~$\mathsf{M}$ global triangle, left open by the
boundary-linear argument
(Theorem~\ref{thm:boundary-linear-bulk-boundary}, which requires
linearity in normal directions that DS violates), is therefore closed
as a cohomological and completed/pro comparison triangle.
\end{proof}

\begin{remark}[Admissible construction]%
\label{rem:universal-holography-admissible-construction}%
The functor $\Phi_{\mathrm{hol}}$ is a construction on the
\(\Xi\)-realized admissible domain: each object includes a conformal
vector, a boundary-linear, affine, or DS origin, a physical BV model, a
boundary comparison, a bulk--Hochschild comparison, and an ambient. The
four-part statement (boundary comparison, bulk--centre comparison,
functoriality with DS, class coverage) is asserted only in that domain.
The bulk observables are identified by the named comparison theorems
for $\SCchtop$ acting on the pair
$(Z^{\mathrm{der}}_{\mathrm{ch}}(\cA),\cA)$.
The DS--Hochschild bridge~\ref{thm:ds-hochschild-bridge} closes the
class~$\mathsf{M}$ bulk comparison cohomologically and in the
completed/pro chain-level ambient; the original-complex chain-level
statement remains governed by the bounded-shift HPL criterion.
\end{remark}


%% ========================================================================
%% Verdier distance and chiral Hochschild comparison.
%% ========================================================================

\section{Holographic distance via the Verdier pairing}
\label{sec:holographic-verdier-distance}

The Verdier pairing gives a distance only after the physical error
filtration has been compared with the bar filtration.

\begin{theorem}[Verdier-realised holographic code distance;
\ClaimStatusProvedHere]
\label{thm:hc-verdier-distance}
Let $\cA$ be a chirally Koszul algebra on a smooth projective curve
$X$, and let $(\cA,\cA^!)$ be its Verdier--Koszul pair. Let
$V\colon \cH_{\mathrm{bdy}}\to\cH_{\mathrm{bulk}}$ be a holographic
encoding with physical error dg algebra $\mathcal E$. Suppose
$\mathcal E$ carries an exhaustive weight filtration and that there is a
filtered morphism
\[
 \operatorname{gr}\mathcal E
 \longrightarrow
 \operatorname{gr}\bigl(\mathrm{Bar}(\cA)\otimes\mathrm{Bar}(\cA^!)\bigr)
\]
which identifies the associated graded of the undetectable error ideal
with the radical of the Verdier evaluation
\[
 \mathrm{ev}_\cA\colon
 \mathrm{Bar}(\cA)\otimes\mathrm{Bar}(\cA^!)\longrightarrow \mathbf 1
\]
and preserves weight. Then
\[
 d_{\mathrm{QECC}}(\cA) \;=\; d_{\mathrm{Verdier}}(\cA),
\]
where
\[
 d_{\mathrm{Verdier}}(\cA)
 =
 \inf\{w\geq 1\mid
 \operatorname{Rad}(\mathrm{ev}_\cA)\cap
 \operatorname{gr}_w(\mathrm{Bar}(\cA)\otimes\mathrm{Bar}(\cA^!))\neq 0\}.
\]
The convention is $\inf\varnothing=+\infty$.
\end{theorem}

\begin{proof}
The code distance is the least weight of a nonzero undetectable error.
By the filtered comparison hypothesis, the associated graded of that
ideal is the Verdier radical, and the comparison preserves weight.
Thus the least physical weight of an undetectable error is the least bar
weight of a nonzero Verdier-radical class.
\end{proof}

\begin{remark}[Genus-zero calibrations]
\label{cor:hc-verdier-distance-examples}
The genus-zero computations give three calibration values:
\begin{itemize}
 \item Heisenberg $H_k$: $d_{\mathrm{Verdier}}(H_k) = +\infty$
 \textup{(}free-field, no obstruction\textup{).}
 \item Virasoro at Lee--Yang $\mathrm{Vir}_{c = -22/5}$:
 $d_{\mathrm{Verdier}} = 3$.
 \item HaPPY pentagon code: $d_{\mathrm{Verdier}} = 3$, matching the
 published QECC distance $d_{\mathrm{QECC}} = 3$ of the
 $[[5, 1, 3]]$ pentagon code on the pentagon model.
\end{itemize}
Each equality is an instance of Theorem~\ref{thm:hc-verdier-distance}
after the weight filtration of the physical error algebra is identified
with the displayed bar filtration.
\end{remark}

\begin{proposition}[BRST lower bound on the Verdier distance;
\ClaimStatusProvedHere]
\label{rem:brst-verdier-lower-bound}
Assume $\cA$ admits a quasi-free BRST resolution whose ghost-spin
filtration maps to the Verdier bar filtration and sends every nonzero
radical class to bar weight at least twice its ghost spin. If the
Verdier radical is zero, then $d_{\mathrm{Verdier}}(\cA)=+\infty$.
Otherwise, if
$\lambda_{\min}^{\mathrm{ghost}}$ is the smallest ghost spin of a
nonzero radical class, then
\[
 d_{\mathrm{Verdier}}(\cA) \;\geq\; 2 \lambda_{\min}^{\mathrm{ghost}},
\]
and equality holds precisely when a Verdier-radical class occurs at the
minimal ghost spin and has bar weight
$2\lambda_{\min}^{\mathrm{ghost}}$.
\end{proposition}

\begin{proof}
The Verdier distance is the least bar weight of a nonzero radical
class. The stated filtration inequality gives bar weight
$\geq 2\lambda_{\min}^{\mathrm{ghost}}$ for every such class. The last
assertion is the equality case.
\end{proof}

\begin{conjecture}[Genus-1 Verdier distance]
\label{conj:verdier-distance-g1}
At genus $1$, the Verdier-pairing distance is governed by the elliptic
bar pairing and the Felder form factor. No additive formula follows
from the genus-zero pairing without this elliptic comparison.
\end{conjecture}

\begin{conjecture}[Higher-genus Verdier distance]
\label{conj:verdier-distance-higher-genus}
At genus $g \geq 2$, the Verdier-pairing distance is computed by a
spectral sequence whose $E_2$ page is the genus-$g$ Hochschild cohomology
of the Koszul pair together with the multi-weight Verdier correction.
The QECC interpretation requires a fibrewise error filtration over
$\overline{\mathcal M}_g$.
\end{conjecture}


\section{The chiral Hochschild comparison theorem}
\label{sec:chiral-hochschild-trinity}

The Fulton--MacPherson, enveloping, and factorization-homology models
are compared by two quasi-isomorphisms.

\begin{theorem}[Chiral Hochschild comparison theorem;
\ClaimStatusProvedHere]
\label{thm:chiral-hochschild-trinity}
\phantomsection\label{prop:chiral-hochschild-comparison-locus}
Let $\cA$ be a logarithmic finite-type chiral algebra in a filtered
chiral ambient in which the geometric, algebraic, and
factorization-homology Hochschild models are defined. Suppose the
comparison maps
\[
 \Phi_{\mathrm{GA}}\colon
 C^\bullet_{\mathrm{chiral}}(\cA)
 \longrightarrow
 \mathrm{Ext}^*_{\cA^e}(\cA,\cA),
 \qquad
 \Phi_{\mathrm{AB}}\colon
 \mathrm{Ext}^*_{\cA^e}(\cA,\cA)
 \longrightarrow
 \mathrm{RHH}_{\mathrm{ch}}(\cA)
\]
are filtered quasi-isomorphisms compatible with the circle
factorization product. Then
\[
 \boxed{\;C^\bullet_{\mathrm{chiral}}(\cA)
 \;\xrightarrow[\Phi_{\mathrm{GA}}]{\sim}\;
 \mathrm{Ext}^*_{\cA^e}(\cA, \cA)
 \;\xrightarrow[\Phi_{\mathrm{AB}}]{\sim}\;
 \mathrm{RHH}_{\mathrm{ch}}(\cA)
 \;=\; \int_{S^1} \cA.\;}
\]
The geometric Fulton--MacPherson complex, the algebraic enveloping
Ext-complex, and the derived chiral Hochschild complex therefore carry
one chiral Hochschild object.
\end{theorem}

\begin{proof}
The assertion is the two-out-of-three statement for the displayed
zigzag of filtered quasi-isomorphisms. Compatibility with circle
factorization identifies the last complex with
$\int_{S^1}\cA$ and transports the product and bracket operations
through the zigzag.
\end{proof}

\begin{theorem}[$\kappa$-conductor as trace on chiral Hochschild;
\ClaimStatusProvedHere]
\label{cor:kappa-conductor-trinity-centre}
In the hypotheses of Theorem~\ref{thm:chiral-hochschild-trinity},
suppose the conductor operator $K_\cA$ is trace-class on the common
Hochschild object and that the Vol~I conductor identity
$K(\cA)=-c_{\mathrm{ghost}}(\mathrm{BRST}(\cA))$ holds in the same
ambient. Then
\[
 K(\cA) \;=\; \mathrm{tr}_{Z(\cA)}(K_\cA)
 \;=\; -c_{\mathrm{ghost}}(\mathrm{BRST}(\cA)).
\]
Here $Z(\cA)$ denotes any of the three quasi-isomorphic Hochschild
models of Theorem~\ref{thm:chiral-hochschild-trinity}.
\end{theorem}

\begin{proof}
The trace-class hypothesis defines $\mathrm{tr}_{Z(\cA)}(K_\cA)$ on
the common Hochschild object. The Vol~I conductor identity identifies
this trace with $-c_{\mathrm{ghost}}(\mathrm{BRST}(\cA))$. The
quasi-isomorphisms of Theorem~\ref{thm:chiral-hochschild-trinity}
transport the trace between the three Hochschild models.
\end{proof}

\begin{example}[K3 realisation of the Hochschild centre supertrace]
\label{ex:trinity-k3-centre-trace}
On the K3 Hall--Borcherds comparison locus, the centre trace is a
bridge between chiral and automorphic data. At $d=2$, the
Mukai-Heisenberg output of the K3 category is an $E_2$-chiral
Heisenberg object, so the Hochschild trace remains on the chiral side
and recovers the K3
$\kappaChHodge$-value. At $d=3$, on the K3 Borcherds row with
Borcherds-side algebra $\mathfrak g_{\Delta_5}$, the same Vol~II
centre trace specialises through the Universal $\Phi$-Trace Identity to
\[
 \Phi_\ast\bigl(K(\cA_{\mathfrak g_{\Delta_5}})\bigr)
 \;=\;
 \Phi_\ast\!\bigl(
 \mathrm{tr}_{Z(\cA_{\mathfrak g_{\Delta_5}})}
 (K_{\cA_{\mathfrak g_{\Delta_5}}})\bigr)
 \;=\;
 \kappa_{\mathrm{BKM}}(\mathfrak g_{\Delta_5})
 \;=\;
 \frac{c_{\phi_{0,1}^{K3}}(0)}{2}
 \;=\;
 5,
\]
where $c_{\phi_{0,1}^{K3}}(0)$ is the constant term of the K3
elliptic-genus Jacobi input in the Borcherds lift
\textup{(}Borcherds 1998; Gritsenko 1999\textup{)}, so
$c_{\phi_{0,1}^{K3}}(0)/2=10/2=5=\mathrm{wt}(\Delta_5)$, and the
dyonic square is $\Phi_{10}^{\mathrm{un}}=\Delta_5^2$. The centre is
therefore the place where the K3 datum can be read both as a Vol~II
supertrace and as a Vol~III Borcherds weight, once the K3
Hall--Borcherds comparison data are supplied.
\end{example}

\begin{remark}[Amplitude and occupation]
\label{rem:hz3-14-amplitude-occupation}
If one of the three Hochschild models of
Theorem~\ref{thm:chiral-hochschild-trinity} has cohomological amplitude
$[0,2]$, the other two have the same amplitude. For the Virasoro
central-extension sector, the occupied degrees are $\{0,2\}$.
\end{remark}

\begin{conjecture}[Comparison extensions: non-Koszul / trace-conductor /
 higher-genus / CY3 / pentagon-coherence]
\label{conj:trinity-extensions}
Five extensions of the comparison-locus statement record the axes
along which the equivalence may break: \textup{(i)} non-Koszul chiral
algebras may admit the comparison only after passing to a Koszul double cover;
\textup{(ii)} the supertrace identification
$K(\cA) = \mathrm{str}(K_\cA)$ extends
to non-Koszul trace-class objects; \textup{(iii)} at higher genus the
three models split into a stratification rather than a coincidence;
\textup{(iv)} the CY$_3$ chiral comparison follows from the Vol~III
$\Phi_3$ functor under CY-A$_3$; \textup{(v)} the Vol~II
$\SCchtop$ pentagon coherence at five presentations controls the
single-colour Hochschild comparison.
\end{conjecture}

\section{Synthesis}
\label{sec:thqghr-supplement}

\begin{remark}[K3 Hall--Borcherds comparison datum]
\label{rem:thqghr-hall-drinfeld}
At $K3\times T^2$ the Vol~II reconstruction is compatible with the
Vol~III Hall--Borcherds target
$\mathcal{D}_\hbar(\mathcal{Y}^{\mathrm{Hall}}_\hbar
(\mathrm{CoHA}_{K3 \times E}))$ only after the oriented hCS--Hall map,
the protected trace comparison, and the Borcherds-Serre quotient data
are supplied. The automorphic normalisation is
\[
 \Phi_{10}^{\mathrm{un}}=\Delta_5^2,\qquad
 \kappa_{\mathrm{BKM}}(\mathfrak g_{\Delta_5})
 =c_{\phi_{0,1}^{K3}}(0)/2=5 .
\]
The one-dimensionality of the $\Delta_5$ automorphic line controls the
Borcherds denominator; it does not by itself construct the full
Hall--Drinfeld double or the twisted 11D-SUGRA bulk.
\end{remark}

\section{DGGPR state sums and conditional gravitational readings}
\label{sec:thqghr-gravitational-partition-RT}

Witten 1989 \emph{CMP}~121 Thm.~2.4 identifies the semisimple
Reshetikhin--Turaev invariant $\tau_{\mathcal C_k}(M^3)$ with the
partition function of three-dimensional $G$-Chern--Simons theory
at level $k$:
\[
 Z_{\mathrm{CS}}(M^3; G, k)
 \;=\;
 \int \mathcal DA\,
 \exp\!\Bigl(\tfrac{ik}{4\pi}\!\int_{M^3}\mathrm{tr}
 (A\wedge\mathrm dA + \tfrac23 A\wedge A\wedge A)\Bigr)
 \;=\;
 \tau_{\mathcal C_k}(M^3).
\]
The non-semisimple DGGPR invariant $\tau_{\mathcal C}(M^3)$ attached
to $\mathcal C = \mathrm{Rep}^{\mathrm{fd}}(\mathfrak u_{\zeta_8}
(\widehat{\mathfrak m}_{\Delta_5}))$ is defined by Hopf-integral and
modified-trace topology (Hennings 1996; Kerler--Lyubashenko 2001;
De~Renzi--Geer--Patureau-Mirand). Its interpretation as a
three-dimensional HT gravitational path integral is an additional
holographic hypothesis: the exact-sector functor, the boundary/asymptotic
condition, and the relevant saddle expansion must be fixed. In the
formulas below
$Z_{\mathrm{grav}}(M^3;\mathbf H_{\Delta_5})$ denotes this conditional
reading of the already-defined DGGPR invariant.

\begin{theorem}[DGGPR state sums on $S^3$, $S^2\!\times\! S^1$, $\Sigma_g\!\times\! S^1$]
\label{thm:thqghr-rt-canonical-manifolds}
Under the admissible-locus holographic reconstruction of
Theorem~\ref{thm:universal-holography-master},
the DGGPR invariant $\tau_{\mathcal C}(M^3)$ computes the closed
topological sector of the proposed bulk 3d holomorphic-topological
gravity coupled to the boundary chiral algebra $\mathbf H_{\Delta_5}$.
Its interpretation as a gravitational path integral is conditional on
the exact-sector holography functor and the chosen boundary/asymptotic
condition. On canonical 3-manifolds:
\begin{align}
 Z_{\mathrm{grav}}(S^3; \mathbf H_{\Delta_5})
 &\;=\;\tau_{\mathcal C}(S^3)
 \;=\;\mathcal D(\mathcal C)^{-1},
 \label{eq:thqghr-Z-S3}\\
 Z_{\mathrm{grav}}(S^2\!\times\! S^1; \mathbf H_{\Delta_5})
 &\;=\;\tau_{\mathcal C}(S^2\!\times\! S^1)
 \;=\;|\Lambda_\infty|,
 \label{eq:thqghr-Z-S2S1}\\
 Z_{\mathrm{grav}}(\Sigma_g\!\times\! S^1; \mathbf H_{\Delta_5})
 &\;=\;\tau_{\mathcal C}(\Sigma_g\!\times\! S^1)
 \;=\;\dim Z^{\mathrm{DGGPR}}(\Sigma_g),
 \label{eq:thqghr-Z-SigmagS1}
\end{align}
with $\mathcal D(\mathcal C)^2 = |\mathrm{Hopf\ dim}(\mathfrak u_{\zeta_8})|$
the total quantum dimension (Vol~III
Theorem~\ref{V3-thm:modtr-rt-S3}), $|\Lambda_\infty|\leq 8^{129}$ the
PBW-finite rank count (Vol~III Theorem~\ref{V3-thm:modtr-zt2-hh0}),
and $Z^{\mathrm{DGGPR}}(\Sigma_g) = \mathrm{Hom}_{\mathcal C}
(\mathbf 1, \mathcal L^{\otimes g})$. At $g = 2$ the partition function
on $\Sigma_2\times S^1$ coincides with the dimension of the
Maass-Spezialschar weight-$10$ Siegel cusp-form space.
\end{theorem}

\begin{proof}
\eqref{eq:thqghr-Z-S3} is Vol~III Theorem~\ref{V3-thm:modtr-rt-S3}
(Reshetikhin--Turaev empty-link surgery for $S^3$) combined with
Witten 1989 Thm.~2.4 reading $\tau_{\mathcal C}(S^3) = \mathcal D^{-1}$
as the Chern--Simons-like path-integral normalisation. The
Hennings 1996 Thm.~3 Hopf-integral formulation gives
$\tau_{\mathcal C}(S^3) = \mu_r(\Lambda)^{-1}$, and
$\mu_r(\Lambda)^2 = |\mathrm{Hopf\ dim}(\mathfrak u_{\zeta_8})|$ up
to the ribbon-twist anomaly of Kerler--Lyubashenko 2001 \S 5.4.
\eqref{eq:thqghr-Z-S2S1} is Atiyah 1988 \S 1 Axiom 4
($\tau_{\mathcal C}(\Sigma\times S^1) = \dim Z^{\mathrm{DGGPR}}(\Sigma)$)
applied at $\Sigma = S^2$ to give
$\dim Z^{\mathrm{DGGPR}}(S^2) = \dim\mathrm{HH}_0^{\mathrm{cat}}
(\mathcal C) = |\Lambda_\infty|$ (Vol~III
Theorem~\ref{V3-thm:modtr-rt-S2timesS1}). \eqref{eq:thqghr-Z-SigmagS1}
is Atiyah Axiom 4 at genus $g$, with the $g = 2$ specialisation
matching the Fourier-coefficient decomposition $c(0, 0, 0) = 1$ of
$(\Phi_{10}^{\mathrm{un}}(Z))^{-1}$ (Vol~III Theorem~\ref{V3-thm:modtr-rt-SigmagS1}
and Eichler--Zagier 1985 \S 6).

\emph{Conditional HT interpretation.} The boundary CFT lives on
$\partial M^3$: for $M^3 = S^3$ there is no boundary, and
$\mathcal D^{-1}$ is the Euclidean partition-function
normalisation. For $M^3 = S^2\times S^1$, the boundary is
$S^2$ and the $S^1$ time-cycle computes the thermal partition
function at zero temperature, counting the closed-string
ground states via $\mathrm{HH}_0^{\mathrm{cat}}$. For $M^3 =
\Sigma_g\times S^1$, the same reading with $\Sigma_g$-sliced
boundary gives the trace of the identity on
$Z^{\mathrm{DGGPR}}(\Sigma_g)$, at $g = 2$ reproducing the
Dijkgraaf--Moore--Verlinde--Verlinde 1997 \emph{CMP}~185 \S 4
second-quantised elliptic-genus partition function
$(\Phi_{10}^{\mathrm{un}}(Z))^{-1}$.
\end{proof}

\begin{theorem}[Lens-space DGGPR invariant and the $\mu_8$ gerbe]
\label{thm:thqghr-rt-lens}
For the lens space $L(p, q)$,
\[
 Z_{\mathrm{grav}}(L(p, q); \mathbf H_{\Delta_5})
 \;=\;
 \tau_{\mathcal C}(L(p, q))
 \;=\;
 \mathcal D^{-1}\sum_{\lambda\in\Lambda_\infty}
 S_{0\lambda}\,\theta_\lambda^{q^\star}\,S_{\lambda 0},
\]
where $S$ is the DGGPR $S$-matrix, $\theta_\lambda$ the ribbon twist,
and $q^\star\equiv q^{-1}\pmod p$. At $p = 8$, the $S$-matrix coincides
with the Fricke involution $w_8\in\mathrm{Sp}_4(\mathbb Z)$ (Vol~III
Theorem~\ref{V3-thm:modtr-s-matrix-fricke}), so
\[
 Z_{\mathrm{grav}}(L(8, 1); \mathbf H_{\Delta_5})
 \;=\;
 \mathcal D^{-1}\cdot\mathrm{tr}\bigl(w_8\theta\bigm\vert
 \mathrm{HH}_0^{\mathrm{cat}}(\mathcal C)\bigr),
\]
the Fricke-twisted trace on the $\mathrm{HH}_0^{\mathrm{cat}}$ state
space. The $\mu_8$-gerbe obstruction of Vol~III
Theorem~\ref{V3-thm:modtr-banding-cocycle} is an invertible bulk
framing/gerbe phase of the closed 3d gauge theory on $L(8, 1)$. Since
the lens space is closed, there is no boundary CFT on
$\partial L(8,1)$; boundary data enter only through a cobordism with
boundary or through the state space assigned to a cutting surface. The
order-$8$ phase tracks the eighth-root ambiguity in the
Kerler--Lyubashenko ribbon twist.
\end{theorem}

\begin{proof}
The surgery formula is Vol~III Theorem~\ref{V3-thm:modtr-rt-lens}:
$L(p, q) = -q/p$-surgery on the unknot in $S^3$, giving
$\tau_{\mathcal C}(L(p, q)) = \mathcal D^{-1}(ST^{q^\star}S)_{00}$.
At $p = 8$, $q = 1$, the $S$-matrix is the Fricke involution
$w_8$ on the $\mu_8$-gerbe-refined state space (Vol~III
Theorem~\ref{V3-thm:modtr-s-matrix-fricke}), so the $(0, 0)$-entry
of $STS$ computes the Fricke-twisted trace
$\mathrm{tr}(w_8\theta|\mathrm{HH}_0^{\mathrm{cat}}(\mathcal C))$.

\emph{Conditional anomaly reading.} The lens space $L(8, 1)$ is the
$\mathbb Z/8$-quotient $S^3/\mathbb Z_8$ by the diagonal action
$(z_1, z_2)\mapsto(e^{2\pi i/8}z_1, e^{-2\pi i/8}z_2)$
(Milnor 1975 \emph{Topology}~14 \S 2). Since $L(8,1)$ is closed, no
boundary partition function lives on $\partial L(8,1)$. The
$\mu_8$-gerbe of Vol~III Theorem~\ref{V3-thm:modtr-banding-cocycle}
instead enters through the closed-state trace after cutting the
3-manifold along a surface. Under the HT interpretation, this is the
't~Hooft anomaly matching the modular anomaly of the
Kerler--Lyubashenko MTC under
$\mathrm{Sp}_4(\mathbb Z)$-action.

\emph{Multi-path verification:}
(i) Reshetikhin--Turaev 1991 \S 3 Thm.~3.3;
(ii) Vol~III Theorem~\ref{V3-thm:modtr-rt-lens};
(iii) Vol~III Theorem~\ref{V3-thm:modtr-s-matrix-fricke} (Fricke identification);
(iv) Kirby 1997 \S 4.1 calculus on $L(p, q)$.
\end{proof}

\begin{theorem}[Mapping-torus DGGPR trace under Fricke]
\label{thm:thqghr-rt-mapping-torus}
For a mapping class $\phi\in\mathrm{MCG}(\Sigma_g)$ and mapping torus
$M_\phi^3 = \Sigma_g\times_\phi S^1$,
\[
 Z_{\mathrm{grav}}(M_\phi^3; \mathbf H_{\Delta_5})
 \;=\;
 \tau_{\mathcal C}(M_\phi^3)
 \;=\;
 \mathrm{Tr}_{Z^{\mathrm{DGGPR}}(\Sigma_g)}\!\bigl(\rho_{\Sigma_g}(\phi)\bigr),
\]
with $\rho_{\Sigma_g}$ the Lyubashenko--Virelizier MCG representation.
At $\Sigma_g = \Sigma_2$, $\phi = w_8$ (Fricke involution, Vol~III
Theorem~\ref{V3-thm:modtr-s-matrix-fricke}),
\[
 Z_{\mathrm{grav}}(\Sigma_2\times_{w_8} S^1; \mathbf H_{\Delta_5})
 \;=\;
 \mathrm{tr}\bigl(w_8\bigm\vert
 \mathrm{HH}_0^{\mathrm{cat}}(\mathcal C\boxtimes\mathcal C^{\mathrm{op}})\bigr),
\]
which restricts to the Fricke trace $4$ on the $M_{24}$-invariant
Humbert block of Vol~III Theorem~\ref{V3-thm:modtr-s-matrix-fricke}
and to $-4 + (1+i)\sqrt 2$ on the $\mu_{16}$-banded block of Vol~III
Remark~\ref{V3-rem:modtr-banding-cocycle-fricke-trace}.
\end{theorem}

\begin{proof}
The mapping-torus trace formula is Atiyah 1988 \S 2 Thm.~2.4 applied
functorially: $Z(M_\phi^3) = \mathrm{tr}(Z(\phi))$ since $M_\phi^3$
is the self-composition of the cobordism $\Sigma_g\times [0, 1]$
along the gluing map $\phi$. The non-semisimple extension is
Lyubashenko 1995 \S 5.5 Thm.~5.5.1 and De~Renzi et~al.\ 2021 Thm.~5.11.
Vol~III Theorem~\ref{V3-thm:modtr-rt-mapping-torus} identifies the
values on the $M_{24}$- and $\mu_{16}$-invariant blocks.

\emph{Conditional HT reading.} $M_\phi^3 = \Sigma_g\times_\phi S^1$ is a
3d spacetime where the spatial $\Sigma_g$ evolves along a closed
time-loop under the diffeomorphism $\phi$; the DGGPR invariant
is the thermal trace $\mathrm{tr}(\phi_*)$ of the action of $\phi$ on
the Hilbert space of states $Z^{\mathrm{DGGPR}}(\Sigma_g)$. The
specific case $\phi = w_8$ on $\Sigma_2$ gives the conditional HT
partition function on the Fricke mapping torus, which by Vol~III
Theorem~\ref{V3-thm:modtr-s-matrix-fricke} computes the trace of the
Siegel Fricke involution $w_8\colon Z\mapsto-(8Z)^{-1}$ on
$\mathfrak H_2$. This trace is the $M_{24}$-Mathieu character
contribution $\chi_{1A}(\tau, z) = 24$ reduced modulo the Fricke
fixed-locus constraint $\mathrm{Fix}(w_8) = H_1\cup H_4$.
\end{proof}

\begin{theorem}[Poincar\'e sphere DGGPR invariant: non-semisimple test]
\label{thm:thqghr-rt-poincare}
For the Poincar\'e homology sphere $\Sigma(2, 3, 5)$, presented as
$(-1)$-framed $E_8$-plumbing (Kirby 1997 \S 4 Example 4.2),
\[
 Z_{\mathrm{grav}}(\Sigma(2, 3, 5); \mathbf H_{\Delta_5})
 \;=\;
 \tau_{\mathcal C}(\Sigma(2, 3, 5))
 \;=\;
 \mathcal D^{-9}\,\Delta^{8}\,
 \mathrm{CL}^{\mathrm{mod}}_{\mathcal C}(E_8; -1),
\]
with $\Delta = \sum_\lambda\theta_\lambda^{-1} d_\lambda^2$ the RT
anomaly and $\mathrm{CL}^{\mathrm{mod}}$ the Geer--Kujawa--Patureau-Mirand
2011 \emph{Compos.~Math.}~147 Thm.~1 modified-trace coloured-link
invariant. On the non-semisimple $\mathbf H_{\Delta_5}$-MTC the DGGPR
state sum is non-zero, in contrast to the semisimple
$\mathfrak{sl}_2$-level-$k$ case where
$\tau(\Sigma(2,3,5)) = \sum_{n=1}^{k+1}\sin^2(\pi n/(k+2))
e^{2\pi i(n^2-1)/8(k+2)}$ (Lawrence--Zagier 1999 \emph{Asian~J.~Math.}~3
Thm.~1): the Poincar\'e sphere is a non-trivial falsification target
for the non-semisimple renormalisation.
\end{theorem}

\begin{proof}
The surgery formula is Vol~III Theorem~\ref{V3-thm:modtr-rt-poincare}:
$b_0(E_8) = 8$, $\sigma(E_8) = -8$, so
$\tau = \mathcal D^{-9}\Delta^{8}\mathrm{CL}^{\mathrm{mod}}$. The
non-vanishing of $\mathrm{CL}^{\mathrm{mod}}$ on the $E_8$-plumbing
follows from the De~Renzi et~al.\ 2021 Thm.~6.3 non-degeneracy of the
modified trace on projective covers of
$\mathfrak u_{\zeta_8}(\widehat{\mathfrak m}_{\Delta_5})$.

\emph{Conditional HT reading.} $\Sigma(2, 3, 5) = \mathrm{SU}(2)/I^*$ is the
spherical space form with binary icosahedral symmetry $I^*$ (Milnor
1975 \S 3). Under the holographic-reconstruction hypothesis it is read
as a 3-dimensional saddle with $I^*$-isometry group, and
$\tau_{\mathcal C}(\Sigma(2, 3, 5))$ is the closed topological sector
assigned to that saddle. The $E_8$-plumbing presentation identifies the
surgery framing with the $E_8$-Cartan bilinear form; the signature
$\sigma(E_8) = -8$ contributes a phase $\Delta^{8}$ in the state sum,
matching the Chern--Simons framing anomaly (Witten 1989 \S 2.4). The
non-vanishing of
$\tau_{\mathcal C}(\Sigma(2, 3, 5))$ on the non-semisimple
$\mathbf H_{\Delta_5}$-MTC, contrasted with the Lawrence--Zagier 1999
semisimple formula, is the falsification target for the
non-semisimple renormalisation on the $E_8$-plumbing geometry.

\emph{Multi-path verification:}
(i) Vol~III Theorem~\ref{V3-thm:modtr-rt-poincare} (explicit surgery);
(ii) De~Renzi et~al.\ 2021 Thm.~6.1 (non-semisimple formula);
(iii) Lawrence--Zagier 1999 Thm.~1 (semisimple comparator);
(iv) Witten 1989 \S 2.4 (framing anomaly).
\end{proof}

\begin{remark}[Conditional dictionary $\tau_{\mathcal C}\leftrightarrow$ HT gravity]
\label{rem:thqghr-holographic-dictionary}
The table below records the conditional dictionary between the
closed 3-manifold $M^3$, its DGGPR invariant $\tau_{\mathcal C}(M^3)$,
and the bulk-HT partition-function reading under the
holographic reconstruction of Chapter~\ref{ch:programme-climax-platonic}:
\begin{center}
\begin{tabular}{l|l|l}
$M^3$ & $\tau_{\mathcal C}(M^3)$ & conditional HT reading \\ \hline
$S^3$ & $\mathcal D^{-1}$ & Euclidean normalisation \\
$S^2\!\times\! S^1$ & $|\Lambda_\infty|\leq 8^{129}$ &
 thermal ground-state count \\
$T^2\!\times\! S^1$ & $|\Lambda_\infty|$ &
 closed-string partition function \\
$\Sigma_2\!\times\! S^1$ & $\dim\mathrm{HH}_0^{\mathrm{cat}}
(\mathcal C\boxtimes\mathcal C^{\mathrm{op}})$ &
 second-quantised elliptic genus $(c(0) = 1)$ \\
$L(8, 1)$ & $\mathcal D^{-1}\mathrm{tr}(w_8\theta)$ &
 $\mu_8$-gerbe 't Hooft anomaly \\
$\Sigma_2\times_{w_8}\! S^1$ & $\mathrm{tr}(w_8|\mathrm{HH}_0^{\mathrm{cat}})$ &
 Fricke mapping-torus partition function \\
$\Sigma(2, 3, 5)$ & $\mathcal D^{-9}\Delta^{8}\mathrm{CL}^{\mathrm{mod}}(E_8)$ &
 $E_8$-plumbing saddle with binary-icosahedral $I^*$-isometry \\
\end{tabular}
\end{center}
In each row the left column is the topological data, the centre column
is the DGGPR invariant computed in Vol~III \S\ref{sec:modtr-rt-closed-3manifolds},
and the right column is the conditional HT path-integral reading with
boundary chiral-CFT fixed to $\mathbf H_{\Delta_5}$. The dictionary
extends Witten 1989 Thm.~2.4 (semisimple Chern--Simons $\leftrightarrow$
RT invariant) to the non-semisimple regime, with the
Kerler--Lyubashenko Hopf-integral serving as the path-integral
normalisation. It does not derive a Cardy, BTZ, or Page statement;
those require separate modular invariance, vacuum dominance, saddle
control, and transseries hypotheses.
\end{remark}

\begin{theorem}[DGGPR genus-$g$ state-space dimension tower]
\label{thm:thqghr-rt-genus-g-tower}
For $g\geq 0$,
\[
 Z_{\mathrm{grav}}(\Sigma_g\!\times\! S^1;\mathbf H_{\Delta_5})
 \;=\;
 \tau_{\mathcal C}(\Sigma_g\!\times\! S^1)
 \;=\;
 \dim\mathrm{Hom}_{\mathcal C}(\mathbf 1, \mathcal L^{\otimes g}),
\]
with $\mathcal L = \coend_{\mathcal C}$ the Lyubashenko coend. At low
genus,
\begin{align*}
 \tau_{\mathcal C}(S^2\!\times\! S^1) &\;=\; \dim\mathrm{HH}_0^{\mathrm{cat}}(\mathcal C)
 \;=\; |\Lambda_\infty|\leq 8^{129},\\
 \tau_{\mathcal C}(T^2\!\times\! S^1) &\;=\; \dim\mathrm{HH}_0^{\mathrm{cat}}(\mathcal C)
 \;=\; |\Lambda_\infty|,\\
 \tau_{\mathcal C}(\Sigma_2\!\times\! S^1)
 &\;=\; \dim\mathrm{HH}_0^{\mathrm{cat}}
 (\mathcal C\boxtimes\mathcal C^{\mathrm{op}})
 \;=\; c(0, 0, 0)\cdot|\Lambda_\infty|^{\!{}^{\!2}}
 \;\equiv\; 1\cdot|\Lambda_\infty|^{2}\;\mathrm{mod}\;\mathcal O(q, p, y),
\end{align*}
with $c(0, 0, 0) = 1$ the vacuum Fourier coefficient of
$(\Phi_{10}^{\mathrm{un}})^{-1}$ (Eichler--Zagier 1985 \S 6
Thm.~5.3 normalisation). At
$g\geq 3$ the Birman 1974 Ch.~4 surjection
$\mathrm{MCG}(\Sigma_g)\twoheadrightarrow\mathrm{Sp}_{2g}(\mathbb Z)$
yields
\[
 \tau_{\mathcal C}(\Sigma_g\!\times\! S^1)
 \;\leq\;
 \dim\bigl(\mathrm{HH}_0^{\mathrm{cat}}(\mathcal C)^{\otimes g}\bigr)^{\!{}^{\!\mathrm{Sp}_{2g}(\mathbb Z)}}
 \;=\;
 \dim M^{\mathrm{Sieg}}_{10}\!\bigl(\mathrm{Sp}_{2g}(\mathbb Z)\bigr),
\]
the corresponding weight-$10$ Siegel cusp-form space for genus $g$.
The numerical dimensions are those computed by Igusa 1962
\emph{Invent.}~3 Table II and Tsuyumine 1986 \emph{Mem.~AMS}~63.
\end{theorem}

\begin{proof}
The identity $\tau_{\mathcal C}(\Sigma\times S^1) = \dim Z(\Sigma)$
is Atiyah 1988 \S 1 Axiom 4. At $g = 0$ the coend pair
$\mathrm{Hom}(\mathbf 1, \mathcal L^{\otimes 0})
= \mathrm{End}(\mathbf 1)\cong\mathbb C$, collapsing to the $S^2$
evaluation; the $|\Lambda_\infty|$ count at $g = 1$ is Vol~III
Theorem~\ref{V3-thm:modtr-zt2-hh0} (PBW-truncated pro-limit). At
$g = 2$ the Kerler--Lyubashenko handlebody-vacuum identification
$\mathrm{Hom}(\mathbf 1, \mathcal L^{\otimes 2})
= \iota\bigl((\Phi_{10}^{\mathrm{un}})^{-1}\bigr)\cdot\mathrm{HH}_0^{\mathrm{cat}}
(\mathcal C\boxtimes\mathcal C^{\mathrm{op}})$ is
Theorem~\ref{thm:modbootstrap-state-sum}, and the constant-term
coefficient $c(0, 0, 0) = 1$ is Eichler--Zagier 1985 \S 6 Thm.~5.3
(Fourier--Jacobi expansion
$\Phi_{10}^{\mathrm{un}} = \sum_m\phi_{10, m}e^{2\pi im\tau'}$
with $\phi_{10, 0} = 0$ and $\phi_{10, 1}$ the weight-$10$ index-$1$
Jacobi cusp form with $c(0, 0, 0) = 1$ normalisation).

At $g\geq 3$ the Lyubashenko--Virelizier MCG-representation factors
through $\mathrm{Sp}_{2g}(\mathbb Z)$ (Birman 1974 Ch.~4 surjection
+ Lyubashenko 1995 \S 5.4 Thm.~5.4.1 MCG-action on the coend), so
the state space lies in the $\mathrm{Sp}_{2g}(\mathbb Z)$-invariant
subspace of $\mathrm{HH}_0^{\mathrm{cat}}(\mathcal C)^{\otimes g}$,
equivalently in the corresponding weight-$10$ Siegel-modular space on
$\mathrm{Sp}_{2g}(\mathbb Z)$. The numerical dimensions are those of
Igusa 1962 \emph{Invent.}~3 Table II (genus $2$) and Tsuyumine
1986 \S 4 Thm.~5 (extension to genus $3$ and beyond; cf.\
Freitag 1983 \emph{Siegelsche Modulfunktionen} \S VI.3).

\emph{Multi-path verification:}
(i) Atiyah 1988 Axiom 4;
(ii) Kerler--Lyubashenko 2001 \S 5.4 handlebody-vacuum;
(iii) Birman 1974 Ch.~4 + Lyubashenko 1995 Thm.~5.4.1
(MCG-to-Sp factor);
(iv) Igusa 1962 Table II and Tsuyumine 1986 \S 4 dimensions.
\end{proof}

\begin{theorem}[Conditional Fricke mapping-torus trace]
\label{thm:thqghr-fricke-explicit-trace}
Assume the Vol~III Fricke block decomposition of
$\mathrm{HH}_0^{\mathrm{cat}}
(\mathcal C\boxtimes\mathcal C^{\mathrm{op}})$ under the
$M_{24}$-invariant Humbert block and the $\mu_{16}$-banded block
(Vol~III Remark~\ref{V3-rem:modtr-banding-cocycle-fricke-trace}), with
block traces $4$ and $-4+(1+i)\sqrt2$. Then the conditional Fricke
mapping-torus partition function is
\[
 Z_{\mathrm{grav}}(\Sigma_2\times_{w_8} S^1; \mathbf H_{\Delta_5})
 \;=\;
 \underbrace{4}_{\text{Humbert block}}
 \;+\;
 \underbrace{-4 + (1+i)\sqrt 2}_{\text{$\mu_{16}$-banded block}}
 \;=\;
 (1+i)\sqrt 2
 \;=\;
 2\,e^{i\pi/4}.
\]
The magnitude is $|Z_{\mathrm{grav}}| = 2$, and the phase is
$\arg(Z_{\mathrm{grav}}) = \pi/4$, the $\mu_8$-gerbe-anomaly phase
tracking the Kerler--Lyubashenko ribbon-twist eighth root of unity
(Vol~III Theorem~\ref{V3-thm:modtr-banding-cocycle}).
\end{theorem}

\begin{proof}
The Atiyah--Lyubashenko mapping-torus formula
$\tau_{\mathcal C}(\Sigma_2\times_\phi S^1)
= \mathrm{tr}(\rho_{\Sigma_2}(\phi))$ (Atiyah 1988 \S 2 Thm.~2.4
extended by Lyubashenko 1995 \S 5.5 Thm.~5.5.1 and De~Renzi
et~al.\ 2021 Thm.~5.11) applied to $\phi = w_8$ (Vol~III
Theorem~\ref{V3-thm:modtr-s-matrix-fricke}) gives
$Z_{\mathrm{grav}} = \mathrm{tr}(w_8|\mathrm{HH}_0^{\mathrm{cat}})$.

The assumed block decomposition is Vol~III
Remark~\ref{V3-rem:modtr-banding-cocycle-fricke-trace}:
$\mathrm{HH}_0^{\mathrm{cat}}(\mathcal C\boxtimes\mathcal C^{\mathrm{op}})
= \mathrm{Humbert}\oplus\mathrm{Banded}$ with
$\mathrm{tr}(w_8|\mathrm{Humbert}) = 4$
(Vol~III Remark~\ref{V3-rem:modtr-s-eigenvalues}: Hirzebruch--Zagier
Euler-characteristic $\chi(H_1) + \chi(H_4) - \chi(H_1\cap H_4)
= 2 + 2 - 0 = 4$) and
$\mathrm{tr}(w_8|\mathrm{Banded}) = -4 + (1+i)\sqrt 2$. The latter is
used here as the Vol~III block trace; it is not derived from the
primitive-root Ramanujan sum alone.

Summing: $4 + (-4 + (1+i)\sqrt 2) = (1+i)\sqrt 2$.
Magnitude: $|(1+i)\sqrt 2| = \sqrt{1^2 + 1^2}\cdot\sqrt 2 = 2$.
Argument: $\arg(1+i) = \pi/4$, unchanged by positive-real scaling, so
$\arg((1+i)\sqrt 2) = \pi/4$.

\emph{Multi-path verification:}
(i) Atiyah 1988 Thm.~2.4 mapping-torus trace;
(ii) Hirzebruch--Zagier 1976 Humbert Euler-characteristic;
(iii) Vol~III Remark~\ref{V3-rem:modtr-banding-cocycle-fricke-trace}
(banded-block trace).
\end{proof}

\begin{proposition}[Poincar\'e sphere Seifert expansion, conditional]
\label{thm:thqghr-poincare-seifert-closed-form}
The Poincar\'e sphere $\Sigma(2, 3, 5) = \mathrm{SU}(2)/I^*$ is the
Seifert-fibered space with exceptional fibres of multiplicities
$(2, 3, 5)$ and Euler number $e(\Sigma(2, 3, 5)) = -1/30$. Its
binary-icosahedral isometry group $I^*\subset\mathrm{SU}(2)$ has
order $|I^*| = 120$ and is the cover of the icosahedral rotation
group $I\cong A_5$.

Assume the modified-trace Seifert surgery formula applies to the
non-semisimple MTC $\mathcal C =
\mathrm{Rep}^{\mathrm{fd}}(\mathfrak u_{\zeta_8}
(\widehat{\mathfrak m}_{\Delta_5}))$ with projective-cover label set
$\Lambda_\infty$ and local exceptional-fibre kernels
$\mathrm{SF}^{(2,3,5)}_\lambda$. Then the DGGPR invariant has the form
\[
 \tau_{\mathcal C}\!\bigl(\Sigma(2, 3, 5)\bigr)
 \;=\;
 \mathcal D^{-1}\,e^{i\pi c/4}
 \sum_{\lambda\in\Lambda_\infty}
 \mathrm{t}^{\mathrm{mod}}_{P_\lambda}
 \bigl(\mathrm{SF}^{(2,3,5)}_\lambda\bigr),
\]
where $\mathrm{t}^{\mathrm{mod}}_{P_\lambda}$ is the modified trace on
the projective cover $P_\lambda$ and $c$ is the framing anomaly. The
explicit identification of $\mathrm{SF}^{(2,3,5)}_\lambda$ with the
$S$-, $T$-, twist-, and quantum-dimension data of
$\mathcal C$ is the remaining non-semisimple Seifert kernel
calculation; in the semisimple specialization it reduces to the
Lawrence--Rozansky 1999 \emph{CMP}~205 Thm.~1.5 kernel. The
binary-icosahedral symmetry $I^*$ acts, when the small
quantum-$\mathfrak{sl}_2$ comparison is supplied, through the
$A_5$-orbits on $\Lambda_\infty$, partitioning the sum into orbit-sums
indexed by the five conjugacy classes
$\{1, [\mathrm{ord}2], [\mathrm{ord}3], [\mathrm{ord}5], [\mathrm{ord}5]^*\}$
of $A_5$.
\end{proposition}

\begin{proof}
The Seifert description and the binary-icosahedral group are standard:
$\Sigma(2,3,5)=\mathrm{SU}(2)/I^*$ and
$I^*\cong \mathrm{SL}_2(\mathbb F_5)$. Lawrence--Rozansky 1999
\S 4 Thm.~1.5 gives the corresponding semisimple RT Seifert kernel.
The displayed formula is the same statement after replacing the
semisimple quantum trace by the modified trace of
Geer--Kujawa--Patureau-Mirand 2011 and De~Renzi et~al.\ 2021, under the
explicit hypothesis that the non-semisimple Seifert kernel has been
constructed for the present category. The $A_5$ orbit decomposition is
the Lusztig--Saleur small-quantum-group comparison, again conditional on
transporting that action to the projective-cover label set
$\Lambda_\infty$.

Under the HT reading, the phase $e^{i\pi c/4}$ is the framing anomaly
of the closed topological sector. No metric-gravity saddle expansion is
deduced from this state-sum formula alone.
\end{proof}

% Section file for Chapter: Twisted Holography and Quantum Gravity
% Result (G9): Critical String Dichotomy

\section{Critical string dichotomy}
% label removed: sec:thqg-critical-string-dichotomy
\label{ch:thqg-critical-string}% alias for dangling \ref{ch:...}
\index{critical string!dichotomy|textbf}
\index{Virasoro!c=26@$c=26$ critical string}
\index{transgression algebra|textbf}
\index{Clifford completion|textbf}

The Virasoro algebra at $c = 26$ has appeared in three earlier contexts:
as the uncurved dual $\mathrm{Vir}_{26}^! \simeq \mathrm{Vir}_0$
(Proposition~\ref{V1-prop:virasoro-c26-selfdual}),
as the endpoint of the complementarity sum
$\kappa + \kappa' = c/2 + (26-c)/2 = 13$
(Proposition~\ref{V1-prop:vir-complementarity}),
and as the saturation point of the depth-zero resonance shadow
(Remark~\ref{V1-rem:virasoro-resonance-model}).
The three facts coalesce into a single structural dichotomy,
controlled by a Clifford algebra built from the handle
decomposition of genus-$g$ surfaces.

\begin{remark}[Scope restriction inherited from Volume~I]
The \ClaimStatusProvedHere{} results of this section build on Volume~I
statements (Theorems \texttt{vir-all-genera}, \texttt{topological-regime},
\texttt{gravitational-regime}, \texttt{critical-string-dichotomy},
\texttt{ghost-sector-c26}, \texttt{vir-self-duality-c13};
Propositions \texttt{vir-complementarity},
\texttt{virasoro-c26-selfdual},
\texttt{genus-g-curvature-package}) whose scope qualifiers propagate
to the present proofs. Load-bearing cross-volume citations in proof
bodies below carry the explicit ``Volume~I'' prose prefix;
the semantic content (complementarity tallies, regime dichotomy,
secondary-anomaly data) is inherited without restatement.
\end{remark}

The construction proceeds from the curved bar complex alone.
The curvature $\dfib^{\,2} = \kappa(\cA) \cdot \omega_g$
(Convention~\ref{V1-conv:higher-genus-differentials})
is not an obstruction to be cancelled but a datum to be \emph{transgressed}:
the transgression algebra $B_\Theta$ is the universal solution to the
problem of killing the curvature by adjoining an element whose square
\emph{is} the curvature. The secondary anomaly element
$u = \eta^2 \in B_\Theta$ is closed and central. Its invertibility or
vanishing controls everything: when $u$ is invertible, the genus-$g$
Clifford completion is Morita equivalent to $B_\Theta[u^{-1}]$ via the
spinor module; when $u = 0$, the Clifford relations degenerate to
exterior algebra relations, and the bar complex becomes a genuine cochain
complex tensored with $\bigwedge^{2g+1}$. For the Virasoro at $c = 0$
(i.e., the Koszul dual of $\mathrm{Vir}_{26}$),
$u = 0$ is the algebraic exterior/uncurved shadow of the ghost-sector
simplification that appears in the no-ghost theorem;
at $c = 13$, the symmetric distribution $\kappa = \kappa' = 13/2$ gives
the self-dual fixed point of the Feigin--Frenkel involution.

% ======================================================================
%
% SUBSECTION: THE TRANSGRESSION ALGEBRA
%
% ======================================================================

\subsection{The transgression algebra}
% label removed: subsec:transgression-algebra
\index{transgression algebra!construction|textbf}

\subsubsection{Construction from the curved bar complex}
% label removed: subsubsec:transgression-construction

The genus-$g$ bar complex $(\barB^{(g)}(\cA), \dfib)$ of a cyclic
chiral algebra~$\cA$ satisfies
$\dfib^{\,2} = \kappa(\cA) \cdot \omega_g \cdot \id$
(Proposition~\ref{V1-prop:genus-g-curvature-package}),
with $\kappa(\cA) = \mathrm{Tr}_\cA$
(Definition~\ref{V1-def:modular-curvature-direct}) and $\omega_g$ the
Arakelov $(1,1)$-form on the genus-$g$ fiber.
The transgression algebra is the universal solution to the
problem of killing this curvature by adjoining a degree-one
element whose square is the curvature.

\begin{definition}[Transgression algebra]
% label removed: def:transgression-algebra
\index{transgression algebra!definition|textbf}
Let $(B, d)$ be a graded-associative $\Bbbk$-algebra
equipped with a degree-$1$ derivation $d\colon B \to B$ satisfying the
\emph{curved} relation
\begin{equation}% label removed: eq:g9-curved-relation
d^2(b) = [\lambda, b]
\qquad
\text{for all } b \in B,
\end{equation}
where $\lambda \in B^2$ is a degree-$2$ central element (the
\emph{curvature}). The \emph{transgression algebra} of the pair
$(B, d, \lambda)$ is
\begin{equation}% label removed: eq:g9-transgression-def
B_\Theta
\;:=\;
B\langle \eta \rangle \big/ (\eta^2 - \lambda),
\end{equation}
where $\eta$ is a formal generator of degree~$1$, subject to
$|\eta| = 1$ and $\eta^2 = \lambda$. The differential extends to
$B_\Theta$ by
\begin{equation}% label removed: eq:g9-transgression-diff
d_{B_\Theta}(\eta) := 0,
\qquad
d_{B_\Theta}(b) := d(b) + [\eta, b]
\quad\text{for } b \in B.
\end{equation}
\end{definition}

\begin{remark}[Degree and sign conventions]
% label removed: rem:g9-transgression-signs
In~\eqref{V1-eq:g9-transgression-def}, the algebra $B\langle \eta \rangle$
is the free graded-associative extension of~$B$ by an odd generator
$\eta$ of degree~$1$, with the Koszul sign rule: $\eta \cdot b =
(-1)^{|b|} b \cdot \eta$ for homogeneous $b \in B$. The relation
$\eta^2 = \lambda$ is consistent with graded commutativity: since
$|\eta| = 1$ is odd, $\eta^2 = \eta \cdot \eta$ is
\emph{not} forced to vanish in the graded-commutative world (where it
would be $-\eta^2$); in the associative setting, $\eta^2$ is a
well-defined degree-$2$ element. The relation $\eta^2 = \lambda$
identifies this element with the curvature.
\end{remark}

\begin{proposition}[Transgression kills curvature; \ClaimStatusProvedHere]
% label removed: prop:transgression-kills-curvature
\index{transgression algebra!kills curvature}
The extended differential $d_{B_\Theta}$ on $B_\Theta$ satisfies
\begin{equation}% label removed: eq:g9-dBT-squared
d_{B_\Theta}^{\,2} = 0.
\end{equation}
\end{proposition}

\begin{proof}
We verify $d_{B_\Theta}^{\,2}$ on generators.

\emph{On $\eta$}:
$d_{B_\Theta}(\eta) = 0$, so $d_{B_\Theta}^{\,2}(\eta) = 0$.

\emph{On $b \in B$}: Compute
\begin{align}
d_{B_\Theta}^{\,2}(b)
&= d_{B_\Theta}\bigl(d(b) + [\eta, b]\bigr) \notag \\
&= d_{B_\Theta}(d(b)) + d_{B_\Theta}([\eta, b]) \notag \\
&= \bigl(d^2(b) + [\eta, d(b)]\bigr)
 + \bigl([d_{B_\Theta}(\eta), b]
 + (-1)^{|\eta|}[\eta, d_{B_\Theta}(b)]\bigr). % label removed: eq:g9-dsq-expand
\end{align}
The first term in the second parenthesis vanishes since
$d_{B_\Theta}(\eta) = 0$. Expand the remaining terms:
\begin{align}
d_{B_\Theta}^{\,2}(b)
&= d^2(b) + [\eta, d(b)]
 - [\eta, d(b) + [\eta, b]] \notag \\
&= d^2(b) + [\eta, d(b)]
 - [\eta, d(b)] - [\eta, [\eta, b]] \notag \\
&= d^2(b) - [\eta, [\eta, b]]. % label removed: eq:g9-dsq-simplify
\end{align}
The sign $(-1)^{|\eta|} = -1$ accounts for the negative sign in the
last line. Now compute the double bracket.
Using associativity and the relation $\eta^2 = \lambda$:
\begin{align}
[\eta, [\eta, b]]
&= \eta(\eta b - (-1)^{|b|}b\eta)
 - (-1)^{|b|+1}(\eta b - (-1)^{|b|}b\eta)\eta \notag \\
&= \eta^2 b - (-1)^{|b|}\eta b \eta
 - (-1)^{|b|+1}\eta b \eta + (-1)^{2|b|+1}b\eta^2 \notag \\
&= \lambda b - (-1)^{|b|}\eta b \eta
 + (-1)^{|b|}\eta b \eta - b\lambda \notag \\
&= \lambda b - b \lambda \notag \\
&= [\lambda, b]. % label removed: eq:g9-bracket-curvature
\end{align}
Substituting into~\eqref{V1-eq:g9-dsq-simplify}:
\[
d_{B_\Theta}^{\,2}(b)
= d^2(b) - [\lambda, b]
= [\lambda, b] - [\lambda, b]
= 0,
\]
using $d^2(b) = [\lambda, b]$ from~\eqref{V1-eq:g9-curved-relation}.
Since $d_{B_\Theta}$ is a derivation and vanishes on all
generators, $d_{B_\Theta}^{\,2} = 0$ on all of $B_\Theta$.
\end{proof}

\begin{remark}[Module structure and $\eta$-decomposition]
% label removed: rem:g9-eta-filtration
\index{transgression algebra!module structure}
Since $\eta^2 = \lambda \in B$, any element of $B_\Theta$ can be
written uniquely as $b_0 + b_1 \eta$ with $b_0, b_1 \in B$.
Thus $B_\Theta$ is a rank-$2$ free left $B$-module:
\begin{equation}% label removed: eq:g9-BT-decomposition
B_\Theta = B \oplus B\eta.
\end{equation}
The multiplication rule is
\begin{equation}% label removed: eq:g9-BT-multiplication
(a_0 + a_1\eta)(b_0 + b_1\eta)
= (a_0 b_0 + (-1)^{|a_1|}a_1 b_1 \lambda)
 + (a_0 b_1 + (-1)^{|a_1|}a_1 b_0)\eta.
\end{equation}
\end{remark}

\begin{proposition}[Quotient identification; \ClaimStatusProvedHere]
% label removed: prop:g9-quotient
\index{transgression algebra!quotient}
The quotient by the ideal generated by~$\eta$ recovers the algebra
modulo its curvature:
\begin{equation}% label removed: eq:g9-quotient-id
B_\Theta / (\eta) \;\cong\; B / (\lambda).
\end{equation}
The induced differential on $B/(\lambda)$ squares to zero.
\end{proposition}

\begin{proof}
From~\eqref{V1-eq:g9-BT-decomposition}, the ideal $(\eta) = B\eta + B\eta^2
= B\eta + B\lambda$. The quotient
$B_\Theta/(\eta) = (B \oplus B\eta)/(B\eta + B\lambda)
\cong B/(B\lambda) = B/(\lambda)$.
On the quotient, $d_{B_\Theta}(b) = d(b) + [\eta, b] \mapsto d(b)$
(since $\eta \mapsto 0$), and $d^2(b) = [\lambda, b] \mapsto 0$ since
$\lambda \mapsto 0$ in $B/(\lambda)$.
\end{proof}

\subsubsection{Specialization to the bar complex}
% label removed: subsubsec:transgression-bar
\index{transgression algebra!bar complex specialization}

Apply the transgression construction to the genus-$g$ bar complex of a
cyclic chiral algebra~$\cA$. Set:
\begin{equation}% label removed: eq:g9-specialization
B := \barB^{(g)}(\cA),
\quad
d := \dfib,
\quad
\lambda := \kappa(\cA) \cdot \omega_g.
\end{equation}
The curvature $\lambda = \kappa \cdot \omega_g$ is central: it is a scalar
multiple of the identity in the bar coalgebra tensored with the
$(1,1)$-class on the fiber.

\begin{definition}[Bar-transgression complex]
% label removed: def:bar-transgression-complex
\index{bar-transgression complex|textbf}
The \emph{bar-transgression complex} of $\cA$ at genus $g$ is
\begin{equation}% label removed: eq:g9-bar-transgression
B_\Theta^{(g)}(\cA)
\;:=\;
\barB^{(g)}(\cA)\langle \eta \rangle
\big/ (\eta^2 - \kappa(\cA) \cdot \omega_g),
\end{equation}
with $d_{B_\Theta}^{\,2} = 0$ by
Proposition~\ref{V1-prop:transgression-kills-curvature}.
Its cohomology is the \emph{transgressed bar cohomology}
$H^*(B_\Theta^{(g)}(\cA))$.
\end{definition}

\begin{remark}[Relation to the MC element]
% label removed: rem:g9-mc-relation
\index{transgression algebra!MC element relation}
The curvature $\kappa(\cA) \cdot \omega_g$ is the genus-$g$ scalar
shadow of the universal MC element $\Theta_\cA$:
$\kappa(\cA) \cdot \omega_g = \mathrm{Sh}_{g,0}(\Theta_\cA)$
(Corollary~\ref*{V1-cor:shadow-extraction}).
The transgression algebra is the universal algebraic solution to
killing this shadow. The corrected differential $\Dg{g}$
(Convention~\ref{V1-conv:higher-genus-differentials}) is a
\emph{geometric} solution using $A$-cycle periods. The universal
property ensures a comparison map
\[
B_\Theta^{(g)}(\cA) \longrightarrow
(\barB^{(g)}(\cA),\, \Dg{g}),
\qquad
\eta \longmapsto
\sum_{k=1}^g t_k\, \omega_k,
\]
which is a dg algebra homomorphism when the $A$-periods are specified.
\end{remark}

% ======================================================================
%
% SUBSECTION: THE SECONDARY ANOMALY ELEMENT
%
% ======================================================================

\subsection{The secondary anomaly element}
% label removed: subsec:secondary-anomaly
\index{secondary anomaly!construction|textbf}

\begin{definition}[Secondary anomaly element]
% label removed: def:secondary-anomaly
\index{secondary anomaly!definition|textbf}
The \emph{secondary anomaly element} of the transgression algebra
$B_\Theta^{(g)}(\cA)$ is
\begin{equation}% label removed: eq:g9-secondary-anomaly
u \;:=\; \eta^2 \;=\; \kappa(\cA) \cdot \omega_g
\;\in\; B_\Theta^{(g)}(\cA).
\end{equation}
It has cohomological degree $|u| = 2$.
\end{definition}

\begin{proposition}[Properties of $u$; \ClaimStatusProvedHere]
% label removed: prop:secondary-anomaly-properties
\index{secondary anomaly!properties}
The secondary anomaly element satisfies:
\begin{enumerate}[label=\textup{(\roman*)}]
\item \emph{Closedness.}\; $d_{B_\Theta}(u) = 0$.

\item \emph{Centrality.}\; $u \cdot b = b \cdot u$ for all $b \in
B_\Theta^{(g)}(\cA)$.

\item \emph{Factorization.}\; $u = \kappa(\cA) \cdot \omega_g$, where
$\kappa(\cA) \in \Bbbk$ is the scalar modular characteristic and
$\omega_g \in H^{1,1}(\Sigma_g)$ is the Arakelov class. The element
$u$ vanishes if and only if $\kappa(\cA) = 0$.

\item \emph{Shadow obstruction tower level.}\; Under the identification
$\kappa(\cA) = \Theta_\cA^{\leq 2}$
(Definition~\ref{V1-def:modular-curvature-direct}), the
secondary anomaly is the degree-$2$ shadow evaluated on the genus-$g$
fiber: $u = \Theta_\cA^{\leq 2}\big|_{\Sigma_g}$.
\end{enumerate}
\end{proposition}

\begin{proof}
\emph{(i)} Compute $d_{B_\Theta}(u) = d_{B_\Theta}(\eta^2)$. By the
Leibniz rule,
\[
d_{B_\Theta}(\eta^2) = d_{B_\Theta}(\eta) \cdot \eta
+ (-1)^{|\eta|}\eta \cdot d_{B_\Theta}(\eta) = 0 + 0 = 0,
\]
since $d_{B_\Theta}(\eta) = 0$.

\emph{(ii)} For $b \in B$: since $u = \kappa \cdot \omega_g$ is a
scalar in the center of~$B$, we have $u \cdot b = b \cdot u$. For
$b = \eta$: $u \cdot \eta = \eta^2 \cdot \eta = \eta \cdot \eta^2 =
\eta \cdot u$, where the middle step uses centrality of
$\kappa \cdot \omega_g$ in $B$. Since $B_\Theta = B \oplus B\eta$,
$u$ commutes with everything.

\emph{(iii)} The factorization is the definition. The Arakelov form
$\omega_g$ is strictly positive on $\Sigma_g$ for $g \geq 1$
(Proposition~\ref{V1-prop:genus-g-curvature-package}), so $u = 0$
iff $\kappa(\cA) = 0$.

\emph{(iv)} By Definition~\ref{V1-def:modular-curvature-direct},
$\kappa(\cA) = \Theta_\cA^{\leq 2}$ is the weight-$2$ shadow.
\end{proof}

\begin{remark}[The element $u$ as characteristic class]
% label removed: rem:g9-u-characteristic
\index{secondary anomaly!as characteristic class}
Over the moduli space $\overline{\cM}_g$, the element $u$ defines a
section of the Hodge line bundle scaled by $\kappa$:
\[
u \;\in\; H^0\bigl(\overline{\cM}_g,\;
\kappa(\cA) \cdot c_1(\mathbb{E})\bigr).
\]
Its vanishing $\{u = 0\}$ is the
\emph{uncurved locus} $\{\kappa(\cA) = 0\}$: the algebraic
content of anomaly cancellation.
\end{remark}

% ======================================================================
%
% SUBSECTION: CONNECTING HOMOMORPHISM
%
% ======================================================================

\subsection{The secondary anomaly as connecting homomorphism}
% label removed: subsec:secondary-connecting
\index{secondary anomaly!connecting homomorphism}

The transgression algebra fits into a short exact sequence whose
connecting homomorphism is a Bockstein operator.

\begin{proposition}[Short exact sequence of the transgression;
\ClaimStatusProvedHere]
% label removed: prop:g9-ses-transgression
There is a short exact sequence of cochain complexes
\begin{equation}% label removed: eq:g9-ses
0 \longrightarrow
B \xrightarrow{\;\cdot\,\eta\;}
B_\Theta
\xrightarrow{\;\mathrm{pr}\;}
B / (\lambda)
\longrightarrow 0,
\end{equation}
where the first map is multiplication by $\eta$ and the second is
projection modulo the ideal $(\eta)$. The connecting homomorphism
in the associated long exact sequence is
\begin{equation}% label removed: eq:g9-connecting
\delta\colon H^n(B/(\lambda)) \longrightarrow H^{n+1}(B),
\qquad
\delta([\bar{b}]) = [d(b)],
\end{equation}
where $b \in B$ is any lift of $\bar{b} \in B/(\lambda)$.
\end{proposition}

\begin{proof}
Exactness at each term follows from the rank-$2$ decomposition
$B_\Theta = B \oplus B\eta$: the image of $\cdot\eta$ is $B\eta$,
which is the kernel of $\mathrm{pr}$. Surjectivity of $\mathrm{pr}$
and injectivity of $\cdot\eta$ (since $B$ is torsion-free over
$\Bbbk$) give the short exact sequence.

For the connecting homomorphism: let $\bar{b} \in B/(\lambda)$ be a
cocycle. Lift to $b \in B$. Then
$d_{B_\Theta}(b) = d(b) + [\eta, b] \in B_\Theta$. The
$B\eta$-component of this is $[\eta, b]$, so the projection to
$B/(\lambda)$ sends $d_{B_\Theta}(b) \mapsto d(b)$ mod~$(\lambda)$.
The closedness condition means $d(b) \equiv 0 \pmod{\lambda}$, so
$d(b) = \lambda \cdot c$ for some $c \in B$. The connecting
homomorphism takes $[\bar{b}]$ to $[d(b)]$ in $H^{n+1}(B)$.
\end{proof}

\begin{corollary}[Bockstein long exact sequence; \ClaimStatusProvedHere]
% label removed: cor:g9-bockstein
\index{Bockstein homomorphism!transgression}
The long exact sequence reads
\begin{equation}% label removed: eq:g9-bockstein-les
\cdots \to
H^n(B) \xrightarrow{\cdot\eta}
H^n(B_\Theta) \xrightarrow{\mathrm{pr}}
H^n(B/(\lambda)) \xrightarrow{\delta}
H^{n+1}(B) \to \cdots
\end{equation}
Classes in $H^*(B/(\lambda))$ that lift to the transgressed bar
cohomology $H^*(B_\Theta)$ are precisely those annihilated
by the Bockstein~$\delta$.
\end{corollary}

\begin{remark}[Spectral sequence of the $\eta$-filtration]
% label removed: rem:g9-eta-spectral
\index{spectral sequence!$\eta$-filtration}
The two-step filtration $F^0 B_\Theta = B_\Theta$,
$F^1 B_\Theta = B\eta$ gives a spectral sequence with
$E_1^{0,q} = H^q(B)$, $E_1^{1,q} = H^{q-1}(B)$, and
$d_1$ is multiplication by $u = \kappa \cdot \omega_g$:
\[
E_2^{0,q} = \ker(u\colon H^q(B) \to H^{q+2}(B)),
\qquad
E_2^{1,q} = \mathrm{coker}(u\colon H^{q-1}(B) \to H^{q+1}(B)).
\]
Degeneration at $E_2$ (forced by the two-step filtration) is the
algebraic skeleton of the two-regime dichotomy.
\end{remark}

% ======================================================================
%
% SUBSECTION: THE CLIFFORD COMPLETION
%
% ======================================================================

\subsection{The Clifford completion}
% label removed: subsec:clifford-completion
\index{Clifford completion!construction|textbf}
\index{handle operators|textbf}

The genus-$g$ surface $\Sigma_g$ has first homology
$H_1(\Sigma_g, \Z) \cong \Z^{2g}$ with the intersection pairing
$\langle -, - \rangle\colon H_1 \otimes H_1 \to \Z$.
There exists a symplectic basis $\{A_1, B_1, \ldots, A_g, B_g\}$ with
\begin{equation}% label removed: eq:g9-intersection
\langle A_i, B_j \rangle = \delta_{ij},
\qquad
\langle A_i, A_j \rangle = 0,
\qquad
\langle B_i, B_j \rangle = 0.
\end{equation}
The sewing construction of the genus-$g$ bar complex from genus-$0$
data passes through $g$ handle attachments, one for each pair
$(A_k, B_k)$.

\begin{definition}[Handle operators]
% label removed: def:handle-operators
\index{handle operators!definition|textbf}
Let $B_\Theta = B_\Theta^{(g)}(\cA)$ be the bar-transgression complex
at genus~$g$. For each $k = 1, \ldots, g$, define the
\emph{handle operators}
\begin{equation}% label removed: eq:g9-handle-ops
\alpha_k \;:=\; \oint_{A_k} \eta^{(g)},
\qquad
\beta_k \;:=\; \oint_{B_k} \eta^{(g)},
\end{equation}
where $\eta^{(g)}$ is the degree-$1$ transgression element of
$B_\Theta^{(g)}(\cA)$, and the integrals are fiber integrals along
the cycles $A_k$, $B_k \subset \Sigma_g$. Each handle operator has
degree~$1$ and acts on $B_\Theta^{(g)}(\cA)$ by left multiplication.
\end{definition}

\begin{remark}[Geometric origin of handle operators]
% label removed: rem:g9-handle-geometric
\index{handle operators!geometric origin}
The fiber integral $\oint_{A_k} \eta^{(g)}$ is the residue of the
transgression element along the $A_k$-cycle. In the sewing
construction (Theorem~\ref{thm:heisenberg-sewing}), each handle
attachment corresponds to a self-sewing of the genus-$0$ surface:
cutting along a pair of cycles and regluing. The handle operator
$\alpha_k$ records the $A$-cycle contraction and $\beta_k$ the
$B$-cycle contraction.
\end{remark}

\begin{theorem}[Clifford relations for handle operators;
\ClaimStatusProvedHere]
% label removed: thm:clifford-relations
\index{Clifford algebra!handle operators|textbf}
The handle operators satisfy:
\begin{align}
\alpha_i \beta_j + \beta_j \alpha_i &= \delta_{ij}\, u,
% label removed: eq:g9-clifford-ab \\
\alpha_i \alpha_j + \alpha_j \alpha_i &= 0,
% label removed: eq:g9-clifford-aa \\
\beta_i \beta_j + \beta_j \beta_i &= 0,
% label removed: eq:g9-clifford-bb
\end{align}
where $u = \kappa(\cA) \cdot \omega_g$ is the secondary anomaly
element. The generators
$\{\alpha_1, \beta_1, \ldots, \alpha_g, \beta_g\}$
generate a Clifford algebra $\mathrm{Cl}_{2g}(u)$ over the ring
$\Bbbk[u]$.
\end{theorem}

\begin{proof}
The proof reduces to the intersection pairing on $H_1(\Sigma_g, \Z)$
via the sewing calculus.

\emph{Step~1 (Mixed anticommutator, $i \neq j$).}\;
When $i \neq j$, the cycles $A_i$ and $B_j$ have intersection
number $\langle A_i, B_j \rangle = 0$. The contractions act on
disjoint tensor factors of the bar coalgebra. Since $\alpha_i$ and
$\beta_j$ are odd operators (degree~$1$) on disjoint data,
$\alpha_i \beta_j = -\beta_j \alpha_i$, giving
$\alpha_i\beta_j + \beta_j\alpha_i = 0$.

\emph{Step~2 (Mixed anticommutator, $i = j$).}\;
The cycles $A_i$ and $B_i$ have $\langle A_i, B_i \rangle = 1$.
The iterated fiber integral of $(\eta^{(g)})^2 = u =
\kappa \cdot \omega_g$ over the pair $(A_i, B_i)$ evaluates via
the intersection formula:
\begin{equation}% label removed: eq:g9-clifford-compute
\alpha_i \beta_i + \beta_i \alpha_i
= \oint_{A_i \times B_i}
 (\eta^{(g)})^2
= u \cdot \langle A_i, B_i \rangle
= u.
\end{equation}
The last equality uses the normalization~\eqref{V1-eq:g9-intersection}
of the symplectic basis. The factor~$u$ appears because the
integrand is $(\eta^{(g)})^2 = u$, and the fiber integral computes
the intersection number.

\emph{Step~3 (Same-type anticommutators).}\;
$\langle A_i, A_j \rangle = 0$ for all $i, j$, so
$\alpha_i \alpha_j + \alpha_j \alpha_i = 0$. Identically for
$\beta_i \beta_j$.

\emph{Step~4 (Clifford identification).}\;
Relations~\eqref{V1-eq:g9-clifford-ab}--\eqref{V1-eq:g9-clifford-bb}
are the defining relations of $\mathrm{Cl}_{2g}(q)$, where
$q\colon \Bbbk^{2g} \to \Bbbk$ is the quadratic form
$q(x) = u \cdot Q(x)$ with $Q$ the intersection form:
\[
J = \begin{pmatrix}
0 & I_g \\ -I_g & 0
\end{pmatrix}
\]
in the basis $(\alpha_1, \ldots, \alpha_g, \beta_1, \ldots, \beta_g)$.
The associated symmetric bilinear form
$\langle x, y \rangle_{\mathrm{Cl}} = u \cdot \delta_{ij}$ on the
standard basis follows from the symplectic-to-Clifford passage.
\end{proof}

\begin{definition}[Clifford completion]
% label removed: def:clifford-completion
\index{Clifford completion!definition|textbf}
The \emph{genus-$g$ Clifford completion} of $\cA$ is the algebra
\begin{equation}% label removed: eq:g9-clifford-completion
G_g(\cA)
\;:=\;
B_\Theta^{(g)}(\cA)
\;\otimes_{\Bbbk}\;
\mathrm{Cl}_{2g}(u),
\end{equation}
where $\mathrm{Cl}_{2g}(u)$ is the Clifford algebra generated by
$\{\alpha_1, \beta_1, \ldots, \alpha_g, \beta_g\}$ subject
to~\eqref{V1-eq:g9-clifford-ab}--\eqref{V1-eq:g9-clifford-bb}.
The differential extends by
$d_{G_g}(\alpha_k) = d_{G_g}(\beta_k) = 0$ and
$d_{G_g}(b) = d_{B_\Theta}(b)$ for $b \in B_\Theta^{(g)}(\cA)$.
Then $d_{G_g}^{\,2} = 0$.
\end{definition}

\begin{remark}[Dimension of the Clifford algebra]
% label removed: rem:g9-clifford-dimension
As a $\Bbbk$-vector space, $\mathrm{Cl}_{2g}(u)$ has dimension
$2^{2g}$. When $u$ is invertible,
$\mathrm{Cl}_{2g}(u) \cong \mathrm{Mat}_{2^g}(\Bbbk)$ by the
standard Morita classification. When $u = 0$,
$\mathrm{Cl}_{2g}(0) = \bigwedge^*(\Bbbk^{2g})$. This
$u = 0$ vs.\ $u \neq 0$ dichotomy is the structural divide.
\end{remark}

% ======================================================================
%
% SUBSECTION: EXPLICIT LOW-GENUS COMPUTATIONS
%
% ======================================================================

\subsection{Explicit computations at genus $1$ and $2$}
% label removed: subsec:g9-low-genus
\index{Clifford completion!low genus computations}

\subsubsection{Genus $1$: a single handle}
% label removed: subsubsec:g9-genus-1

At genus $1$, there is a single handle $(A_1, B_1)$ with
$\langle A_1, B_1 \rangle = 1$. The Clifford algebra is
$\mathrm{Cl}_2(u)$, generated by $\alpha := \alpha_1$,
$\beta := \beta_1$ with
\begin{equation}% label removed: eq:g9-genus1-clifford
\alpha\beta + \beta\alpha = u,
\qquad
\alpha^2 = 0,
\qquad
\beta^2 = 0.
\end{equation}

\emph{Multiplication table.}\;
As a $\Bbbk$-vector space, $\mathrm{Cl}_2(u)$ has basis
$\{1, \alpha, \beta, \alpha\beta\}$ with:
\begin{center}
\renewcommand{\arraystretch}{1.3}
\begin{tabular}{c|cccc}
$\cdot$ & $1$ & $\alpha$ & $\beta$ & $\alpha\beta$ \\
\hline
$1$ & $1$ & $\alpha$ & $\beta$ & $\alpha\beta$ \\
$\alpha$ & $\alpha$ & $0$ & $\alpha\beta$ & $0$ \\
$\beta$ & $\beta$ & $u - \alpha\beta$ & $0$ & $u\beta$ \\
$\alpha\beta$ & $\alpha\beta$ & $u\alpha$ & $0$ & $u\alpha\beta$
\end{tabular}
\end{center}

\emph{Volume element.}\;
$\mathrm{vol}_1 = \alpha\beta$ satisfies
$\mathrm{vol}_1^2 = \alpha\beta\alpha\beta =
\alpha(u - \alpha\beta)\beta = u\,\alpha\beta = u\,\mathrm{vol}_1$.

\emph{Spinor module.}\;
When $u$ is invertible, define the spinor module
$S_1 = \Bbbk \cdot v_+ \oplus \Bbbk \cdot v_-$ with
\begin{equation}% label removed: eq:g9-genus1-spinor
\alpha \cdot v_+ = v_-,
\quad
\alpha \cdot v_- = 0,
\quad
\beta \cdot v_+ = 0,
\quad
\beta \cdot v_- = u\, v_+.
\end{equation}
Verification: $(\alpha\beta + \beta\alpha) \cdot v_\pm =
u\, v_\pm$. Hence $S_1$ is a $\mathrm{Cl}_2(u)$-module, and
\begin{equation}% label removed: eq:g9-genus1-morita
\mathrm{Cl}_2(u) \cong \End_\Bbbk(S_1)
\cong \mathrm{Mat}_2(\Bbbk)
\qquad\text{(when $u$ is invertible).}
\end{equation}

\subsubsection{Genus $2$: two handles}
% label removed: subsubsec:g9-genus-2

At genus $2$, $H_1(\Sigma_2, \Z) \cong \Z^4$ has symplectic basis
$(A_1, B_1, A_2, B_2)$ with intersection matrix
\[
J_2 = \begin{pmatrix}
0 & 1 & 0 & 0 \\
-1 & 0 & 0 & 0 \\
0 & 0 & 0 & 1 \\
0 & 0 & -1 & 0
\end{pmatrix}.
\]
The Clifford algebra $\mathrm{Cl}_4(u)$ has
$\dim_\Bbbk = 2^4 = 16$, with basis
$\alpha_1^{a_1}\beta_1^{b_1}\alpha_2^{a_2}\beta_2^{b_2}$
for $a_i, b_i \in \{0,1\}$.

\emph{Spinor module.}\;
When $u$ is invertible, the spinor module $S_2 = \Bbbk^4$ has
basis $\{v_{++}, v_{+-}, v_{-+}, v_{--}\}$ with
\begin{alignat}{2}
\alpha_k \cdot v_{\ldots+_k\ldots}
 &= v_{\ldots-_k\ldots},
&\qquad
\alpha_k \cdot v_{\ldots-_k\ldots}
 &= 0, \notag \\
\beta_k \cdot v_{\ldots+_k\ldots}
 &= 0,
&\qquad
\beta_k \cdot v_{\ldots-_k\ldots}
 &= u\, v_{\ldots+_k\ldots}, % label removed: eq:g9-genus2-spinor
\end{alignat}
where the subscript indicates the $k$-th sign position. Hence
\begin{equation}% label removed: eq:g9-genus2-morita
\mathrm{Cl}_4(u) \cong \End_\Bbbk(S_2)
\cong \mathrm{Mat}_4(\Bbbk)
\qquad\text{(when $u$ is invertible).}
\end{equation}

\emph{Matrix units.}\;
The matrix units in terms of Clifford generators:
\begin{center}
\renewcommand{\arraystretch}{1.2}
\begin{tabular}{lll}
$E_{++,++} = u^{-2}\alpha_1\beta_1\alpha_2\beta_2$, &
$E_{++,--} = u^{-2}\alpha_1\alpha_2$, &
$E_{--,++} = \beta_1\beta_2$, \\
$E_{+-,-+} = u^{-1}\alpha_1\beta_2$, &
$E_{-+,+-} = u^{-1}\beta_1\alpha_2$, &
$E_{--,--} = u^{-2}\beta_1\alpha_1\beta_2\alpha_2$.
\end{tabular}
\end{center}
These $16$ matrix units span $\mathrm{Mat}_4$, confirming
$\mathrm{Cl}_4(u)[u^{-1}] \cong \mathrm{Mat}_4(\Bbbk[u^{\pm 1}])$.

% ======================================================================
%
% SUBSECTION: THE TWO REGIMES
%
% ======================================================================

\subsection{The two regimes: topological and gravitational}
% label removed: subsec:two-regimes
\index{critical string!two regimes|textbf}
\index{topological regime|textbf}
\index{gravitational regime|textbf}

The Clifford completion $G_g(\cA) = B_\Theta^{(g)}(\cA) \otimes
\mathrm{Cl}_{2g}(u)$ admits two fundamentally different structural
descriptions, determined by the invertibility of
$u = \kappa(\cA) \cdot \omega_g$.

\subsubsection{The topological regime: $u$ invertible}
% label removed: subsubsec:topological-regime

\begin{theorem}[Topological regime; \ClaimStatusProvedHere]
% label removed: thm:topological-regime
\index{topological regime!Morita equivalence|textbf}
Suppose $\kappa(\cA) \neq 0$, and that we work over a field
of characteristic~$0$ or characteristic not
dividing~$\kappa(\cA)$. Then:
\begin{enumerate}[label=\textup{(\roman*)}]
\item The Clifford completion is Morita equivalent to the
transgression algebra:
\begin{equation}% label removed: eq:g9-topo-morita
G_g(\cA)[u^{-1}]
\;\simeq_{\mathrm{Morita}}\;
\mathrm{Mat}_{2^g}\bigl(B_\Theta^{(g)}(\cA)[u^{-1}]\bigr).
\end{equation}

\item The Morita equivalence bimodule is the spinor module
\begin{equation}% label removed: eq:g9-spinor-bimod
S_g(\cA)
\;:=\;
B_\Theta^{(g)}(\cA)[u^{-1}] \otimes_\Bbbk S_g,
\end{equation}
where $S_g = \Bbbk^{2^g}$ is the standard spinor representation of
$\mathrm{Cl}_{2g}(u)$.

\item The derived categories are equivalent:
\begin{equation}% label removed: eq:g9-topo-derived
D(G_g(\cA)[u^{-1}]\text{-}\mathsf{mod})
\;\simeq\;
D(B_\Theta^{(g)}(\cA)[u^{-1}]\text{-}\mathsf{mod}).
\end{equation}
\end{enumerate}
\end{theorem}

\begin{proof}
\emph{(i)} The Clifford algebra $\mathrm{Cl}_{2g}(u)$ over a ring in
which $u$ is a unit is isomorphic to $\mathrm{Mat}_{2^g}$. Define the
canonical anticommutation operators
\begin{equation}% label removed: eq:g9-raising-lowering
a_k^+ := u^{-1}\alpha_k,
\qquad
a_k^- := \beta_k,
\qquad k = 1, \ldots, g.
\end{equation}
These satisfy $a_k^+ a_k^- + a_k^- a_k^+ = 1$,
$(a_k^+)^2 = (a_k^-)^2 = 0$, and $[a_k^\pm, a_l^\pm]_+ = 0$ for
$k \neq l$. The Fock space construction gives
$\mathrm{Cl}_{2g}(u)[u^{-1}] \cong \mathrm{Mat}_{2^g}(\Bbbk)$.

\emph{(ii)} The spinor module $S_g$ is the Fock space with vacuum
$|0\rangle := v_{+\cdots+}$ satisfying $\beta_k \cdot |0\rangle = 0$
for all $k$. The $2^g$ basis vectors are
$\alpha_{k_1}\cdots\alpha_{k_r}|0\rangle$ for
$1 \leq k_1 < \cdots < k_r \leq g$. Faithfulness of the left action
(since $u$ is a unit and the lowering operators have full range) gives
$\mathrm{Cl}_{2g}(u) \xrightarrow{\sim} \End(S_g)$.

\emph{(iii)} Morita-equivalent algebras have equivalent derived
categories; the functor $- \otimes_{G_g} S_g(\cA)$ provides it.
\end{proof}

\begin{remark}[Topological meaning of $2^g$]
% label removed: rem:g9-multiplicity
\index{topological regime!multiplicity}
The matrix size $2^g$ is the dimension of the half-spinor
representation, and $|H^1(\Sigma_g, \Z/2)| = 2^{2g}$ is the number of
spin structures on $\Sigma_g$. The Morita equivalence reduces the
Clifford completion to the transgression algebra at the cost of
choosing a spin structure (the vacuum $|0\rangle$).
\end{remark}

\subsubsection{The gravitational regime: $u = 0$}
% label removed: subsubsec:gravitational-regime

\begin{theorem}[Gravitational regime; \ClaimStatusProvedHere]
% label removed: thm:gravitational-regime
\index{gravitational regime!exterior algebra|textbf}
Suppose $\kappa(\cA) = 0$. Then:
\begin{enumerate}[label=\textup{(\roman*)}]
\item The transgression element satisfies $\eta^2 = 0$: the
bar complex is uncurved ($\dfib^{\,2} = 0$).

\item The Clifford relations degenerate to exterior algebra relations:
\[
\alpha_i \beta_j + \beta_j \alpha_i = 0,
\qquad
\alpha_i \alpha_j + \alpha_j \alpha_i = 0,
\qquad
\beta_i \beta_j + \beta_j \beta_i = 0.
\]

\item The Clifford completion splits:
\begin{equation}% label removed: eq:g9-grav-tensor
G_g(\cA)\big|_{u=0}
\;\cong\;
\barB^{(g)}(\cA)
\;\otimes\; \bigwedge\nolimits^*(\Bbbk^{2g}).
\end{equation}

\item The exterior algebra has rank $2^{2g}$ with top element
$\mathrm{vol}_g := \alpha_1 \beta_1 \cdots \alpha_g \beta_g
\in \bigwedge^{2g}(\Bbbk^{2g})$.
The Poincar\'e polynomial is $(1+t)^{2g}$.
\end{enumerate}
\end{theorem}

\begin{proof}
\emph{(i)} $u = \eta^2 = \kappa \cdot \omega_g = 0$.

\emph{(ii)} Set $u = 0$
in~\eqref{V1-eq:g9-clifford-ab}--\eqref{V1-eq:g9-clifford-bb}: all
anticommutators vanish.

\emph{(iii)} When $u = 0$ and $\lambda = 0$, the transgression
algebra quotient is $B_\Theta/(\eta) = B$ (uncurved), and
$\mathrm{Cl}_{2g}(0) = \bigwedge^*(\Bbbk^{2g})$.

\emph{(iv)} $\dim\bigwedge^k(\Bbbk^{2g}) = \binom{2g}{k}$;
$\sum_{k=0}^{2g}\binom{2g}{k}t^k = (1+t)^{2g}$.

The exterior algebra isomorphism in~(iii) is unique: over a
field, the exterior algebra on $2g$ generators is the unique
graded-commutative algebra with Hilbert series $(1+t)^{2g}$
(all generators in degree~$1$, all relations quadratic).
Any isomorphism $\mathrm{Cl}_{2g}(0) \cong
\bigwedge^*(\Bbbk^{2g})$ is determined by its restriction to
the degree-$1$ generators, where it is a linear isomorphism
$\Bbbk^{2g} \xrightarrow{\sim} \Bbbk^{2g}$; the algebra
structure forces the extension to all degrees.
\end{proof}

\begin{remark}[Ghost number in the gravitational regime]
% label removed: rem:g9-ghost-number
\index{ghost number!Clifford completion}
In the BRST language (Remark~\ref{V1-rem:bv-bar-bridge}), the generators
$\alpha_k$, $\beta_k$ carry ghost number~$1$. The volume element
$\mathrm{vol}_g$ has ghost number $2g$. In the gravitational regime,
the BRST cohomology splits by ghost number:
\[
H^*(G_g|_{u=0}) \cong
\bigoplus_{k=0}^{2g} H^{*-k}(\barB^{(g)}(\cA))
\otimes \bigwedge\nolimits^k(\Bbbk^{2g}).
\]
The ghost-number anomaly $2g = \dim H_1(\Sigma_g)$ counts the zero
modes of the ghost field. The exterior (rather than Clifford) structure
is equivalent to the vanishing of the ghost-number current anomaly.
\end{remark}

\begin{remark}[Explicit basis for $\bigwedge^*(\Bbbk^{2g})$]
% label removed: rem:g9-exterior-basis
\index{exterior algebra!explicit basis}
Basis elements are indexed by subsets $I \subset \{1, \ldots, 2g\}$:
for $I = \{i_1 < \cdots < i_k\}$, set $e_I = \gamma_{i_1}\cdots
\gamma_{i_k}$ where $(\gamma_1, \ldots, \gamma_{2g}) =
(\alpha_1, \beta_1, \ldots, \alpha_g, \beta_g)$.
Low-genus data:
\begin{center}
\renewcommand{\arraystretch}{1.2}
\begin{tabular}{ccl}
\toprule
$g$ & $\dim\bigwedge^*(\Bbbk^{2g})$ & Graded dimensions \\
\midrule
$1$ & $4$ & $1, 2, 1$ \\
$2$ & $16$ & $1, 4, 6, 4, 1$ \\
$3$ & $64$ & $1, 6, 15, 20, 15, 6, 1$ \\
$4$ & $256$ & $1, 8, 28, 56, 70, 56, 28, 8, 1$ \\
\bottomrule
\end{tabular}
\end{center}
\end{remark}

% ======================================================================
%
% SUBSECTION: THE VIRASORO AT c = 26
%
% ======================================================================

\subsection{The critical string at $c = 26$}
% label removed: subsec:critical-c26
\index{critical string!c=26@$c = 26$|textbf}
\index{Virasoro!c=26@$c=26$!no-ghost theorem}

We specialize the entire construction to the Virasoro algebra.
The modular characteristic is
\begin{equation}% label removed: eq:g9-kappa-vir
\kappa(\mathrm{Vir}_c) = \frac{c}{2}
\end{equation}
(Theorem~\ref{V1-thm:vir-genus1-curvature}), arising from the quartic
pole $T_{(3)}T = c/2$ in the Virasoro OPE. The Koszul dual is
$\mathrm{Vir}_c^! \simeq \mathrm{Vir}_{26-c}$
(Example~\ref{V1-ex:virasoro-koszul-dual}).

\subsubsection{The Koszul duality pattern}
% label removed: subsubsec:koszul-pattern
\index{Virasoro!Koszul dual}

The complementarity of modular characteristics is
\begin{equation}% label removed: eq:g9-kappa-sum
\kappa(\mathrm{Vir}_c) + \kappa(\mathrm{Vir}_{26-c})
= \frac{c}{2} + \frac{26-c}{2} = 13.
\end{equation}
The secondary anomaly elements for the Koszul pair are
\begin{equation}% label removed: eq:g9-u-pair
u_c = \frac{c}{2} \cdot \omega_g,
\qquad
u_{26-c} = \frac{26-c}{2} \cdot \omega_g.
\end{equation}
Three central charges are structurally distinguished:

\begin{center}
\renewcommand{\arraystretch}{1.4}
\begin{tabular}{clccl}
\toprule
$c$ & Koszul dual & $\kappa$ & $\kappa'$ & Structural property \\
\midrule
$0$ & $\mathrm{Vir}_0^! \simeq \mathrm{Vir}_{26}$ & $0$ & $13$ &
 Uncurved; gravitational regime \\
$13$ & $\mathrm{Vir}_{13}^! \simeq \mathrm{Vir}_{13}$ & $13/2$ & $13/2$ &
 Self-dual; symmetric anomaly \\
$26$ & $\mathrm{Vir}_{26}^! \simeq \mathrm{Vir}_0$ & $13$ & $0$ &
 Maximally curved; dual is gravitational \\
\bottomrule
\end{tabular}
\end{center}

\subsubsection{Two independent vanishing conditions}
% label removed: subsubsec:effective-curvature
\index{complementarity asymmetry|textbf}

The Virasoro family at central charge $c$ is controlled by two
independent scalar conditions. Each governs a different
structural phenomenon; neither implies the other.

\begin{definition}[Complementarity asymmetry]
% label removed: def:effective-curvature
\index{complementarity asymmetry!definition}
For a Koszul pair $(\cA, \cA^!)$ with modular characteristics
$\kappa := \kappa(\cA)$ and $\kappa' := \kappa(\cA^!)$, the
\emph{complementarity asymmetry} is
\begin{equation}% label removed: eq:g9-kappa-eff
\delta_\kappa \;:=\; \kappa - \kappa'.
\end{equation}
Its vanishing is the condition for Koszul self-duality:
$\delta_\kappa = 0$ if and only if $\kappa = \kappa'$.
\end{definition}

For the Virasoro, $\delta_\kappa = c/2 - (26-c)/2 = c - 13$.

The complementarity asymmetry is intrinsic to the Koszul
pair. It measures how unevenly the anomaly budget is
distributed between $\cA$ and $\cA^!$. The second
invariant is extrinsic: it belongs to the physical
matter--ghost system, not to the Koszul pair.

\begin{definition}[Effective curvature]
\index{effective curvature!definition}
The \emph{effective curvature} of a matter--ghost coupling is
\[
\kappa_{\mathrm{eff}} \;:=\;
\kappa(\mathrm{matter}) + \kappa(\mathrm{ghost}),
\]
where the ghost contribution comes from the physical ghost
system (the $bc$ system in the Virasoro case, with
$\kappa_{bc} = -13$ at all central charges). Its vanishing
is anomaly cancellation.
\end{definition}

For the Virasoro at central charge $c$,
$\kappa_{\mathrm{eff}} = c/2 + (-13) = (c-26)/2$. The
effective curvature is \emph{not} the sum
$\kappa + \kappa' = c/2 + (26-c)/2 = 13$, which is a
topological invariant of the Koszul pair (it never vanishes).

The two conditions are related by
$\delta_\kappa = 2\kappa_{\mathrm{eff}} + 13$ and vanish
at different central charges, $13$ units apart:
\[
\delta_\kappa = 0
\;\;\Longleftrightarrow\;\;
c = 13
\quad\text{(Koszul self-duality)},
\qquad
\kappa_{\mathrm{eff}} = 0
\;\;\Longleftrightarrow\;\;
c = 26
\quad\text{(anomaly cancellation)}.
\]

\begin{remark}[Family dependence]
% label removed: rem:g9-complementarity-constraint
The complementarity sum $\kappa + \kappa'$ depends on the
family. For affine Kac--Moody algebras,
$\kappa + \kappa' = 0$
(Remark~\ref{V1-rem:vir-vs-km-complementarity}):
self-duality forces $\kappa = 0$, and if one identifies the
ghost system with the Koszul dual then
$\kappa_{\mathrm{eff}} = \kappa + \kappa' = 0$
identically. For the Virasoro, $\kappa + \kappa' = 13$:
self-duality distributes the anomaly symmetrically at
$\kappa = \kappa' = 13/2$, and anomaly cancellation
requires the separate bc ghost mechanism.
\end{remark}


\begin{remark}[The de Sitter side of Koszul duality]
\label{rem:ads-ds-koszul}
\index{Koszul duality!AdS/dS interpretation}
\index{de Sitter!negative kappa regime}
The complementarity identity
$\kappa(\mathrm{Vir}_c) + \kappa(\mathrm{Vir}_{26-c}) = 13$
admits a gravitational reading. When $c < 26$, both
$\kappa(\mathrm{Vir}_c) = c/2$ and
$\kappa(\mathrm{Vir}_{26-c}) = (26-c)/2$ are positive:
the Koszul pair lives entirely on the anti-de~Sitter side
of the anomaly budget. At $c = 26$ the dual characteristic
$\kappa(\mathrm{Vir}_{26-c}) = (26-c)/2$ vanishes;
the Koszul dual is uncurved, and the matter system
absorbs the full anomaly $\kappa(\mathrm{Vir}_{26}) = 13$.

For $c > 26$ the Koszul dual has
$\kappa(\mathrm{Vir}_{26-c}) = (26-c)/2 < 0$:
the secondary anomaly $u_{26-c} = \kappa(\mathrm{Vir}_{26-c})\cdot\omega_g$
reverses sign, and the fiberwise curvature
$\dfib^{\,2} = \kappa(\mathrm{Vir}_{26-c})\cdot\omega_g$
on the Koszul-dual side becomes negative.
This is the de~Sitter regime of Koszul duality:
the dual algebra $\mathrm{Vir}_{26-c}^{}$
carries a curvature of cosmological sign.

The self-dual point $c = 13$ sits at the transition.
There $\kappa(\mathrm{Vir}_{13}) = \kappa(\mathrm{Vir}_{13}^!) = 13/2$,
and the anomaly budget is distributed symmetrically
between the algebra and its Koszul dual.
The complementarity sum $\kappa + \kappa^{\prime} = 13$
is the total curvature shared by the AdS and dS sectors;
at $c = 13$ the boundary between them passes through
the fixed point of the Feigin--Frenkel involution.
\end{remark}
\subsubsection{The dichotomy theorem}
% label removed: subsubsec:dichotomy-theorem

\begin{theorem}[Critical string dichotomy for the Virasoro;
\ClaimStatusProvedHere]
% label removed: thm:critical-string-dichotomy
\index{critical string!dichotomy theorem|textbf}
The genus-$g$ Clifford completion $G_g(\mathrm{Vir}_c)$ realizes two
regimes:

\begin{enumerate}[label=\textup{(\roman*)}]
\item \emph{Topological regime} ($c \neq 0$):\;
$u_c = (c/2) \cdot \omega_g \neq 0$. The Clifford completion is
Morita equivalent to the transgression algebra:
\[
G_g(\mathrm{Vir}_c)
\simeq_{\mathrm{Morita}}
\mathrm{Mat}_{2^g}\bigl(B_\Theta^{(g)}(\mathrm{Vir}_c)[u_c^{-1}]\bigr).
\]
The genus-$g$ free energy is
\begin{equation}% label removed: eq:g9-partition-topo
F_g(\mathrm{Vir}_c)
= \frac{c}{2} \cdot \lambda_g^{\mathrm{FP}},
\end{equation}
recovering Volume~I Theorem~\ref{V1-thm:vir-all-genera}.

\item \emph{Gravitational regime} ($c = 0$):\;
$u_0 = 0$. The Clifford completion splits:
\[
G_g(\mathrm{Vir}_0)
\cong
\barB^{(g)}(\mathrm{Vir}_0)
\otimes \bigwedge\nolimits^*(\Bbbk^{2g}).
\]
The bar complex is a genuine cochain complex ($\dfib^{\,2} = 0$)
and the exterior factor is the ghost sector.

\item \emph{The $c = 26$ algebra}\/ is in the topological regime
($\kappa = 13 \neq 0$), but its Koszul dual $\mathrm{Vir}_0$ is in
the gravitational regime. The combined Clifford completion is
\begin{equation}% label removed: eq:g9-combined-clifford
G_g(\mathrm{Vir}_{26}) \otimes G_g(\mathrm{Vir}_0)
\;\cong\;
G_g(\mathrm{Vir}_{26})
\otimes \barB^{(g)}(\mathrm{Vir}_0)
\otimes \bigwedge\nolimits^*(\Bbbk^{2g}).
\end{equation}
The full genus-$g$ obstruction budget is carried by
$\mathrm{Vir}_{26}$; the dual $\mathrm{Vir}_0$ contributes only the
exterior algebra factor.
\end{enumerate}
\end{theorem}

\begin{proof}
Parts (i) and (ii) are direct applications of
Volume~I Theorems~\ref{V1-thm:topological-regime}
and~\ref{V1-thm:gravitational-regime} to $\cA = \mathrm{Vir}_c$ with
$\kappa = c/2$. The free energy
formula~\eqref{V1-eq:g9-partition-topo} is Volume~I Theorem~\ref{V1-thm:vir-all-genera}.

For (iii): $\mathrm{Vir}_{26}$ has $\kappa = 13 \neq 0$
(topological); $\mathrm{Vir}_0$ has $\kappa = 0$ (gravitational).
The tensor product follows from K\"unneth. The
complementarity statement is
Volume~I Proposition~\ref{V1-prop:vir-complementarity}: $F_g(\mathrm{Vir}_{26}) +
F_g(\mathrm{Vir}_0) = 13 \cdot \lambda_g^{\mathrm{FP}} + 0 =
13 \cdot \lambda_g^{\mathrm{FP}}$.
\end{proof}

\begin{remark}[Ghost zero modes in bar-cobar language]
% label removed: rem:g9-no-ghost
\index{no-ghost theorem!bar-cobar interpretation}
The classical no-ghost theorem states that the physical states of the
$c = 26$ bosonic string are transverse oscillators. In the present
framework, the Koszul dual $\mathrm{Vir}_{26}^! \simeq \mathrm{Vir}_0$
is in the gravitational regime: $u_0 = 0$, so the handle part of the
dual ghost sector is the exterior algebra
$\bigwedge^*(\Bbbk^{2g})$ with no Clifford entanglement, and
$\barB^{(g)}(\mathrm{Vir}_0)$ is a genuine cochain complex with no
scalar curvature correction. This gives an algebraic model for the
ghost zero-mode count. The passage from this exterior handle sector to
transverse physical states is the Goddard--Thorn and
Frenkel--Garland--Zuckerman spectral-sequence collapse, recorded below
as theorem data in the Clifford-degeneration statement. The ghost
oscillators are the $b$-$c$ modes; the standard zero-mode relation
$\{b_0,c_0\}=1$ supplies the BRST contracting pair, while the
$u_0=0$ Clifford degeneration supplies the exterior handle algebra.

The ghost-number anomaly $2g$ is the count of ghost zero modes
on $\Sigma_g$. The exterior (vs.\ Clifford) structure of the ghost
sector is equivalent to the vanishing of the dual's curvature.
\end{remark}

\subsubsection{Clifford data for the Virasoro at all genera}
% label removed: subsubsec:vir-clifford-data

\begin{computation}[Genus-$g$ Virasoro Clifford data]
% label removed: comp:g9-vir-clifford-data
\index{Virasoro!Clifford completion data}

\begin{center}
\renewcommand{\arraystretch}{1.3}
{\small
\begin{tabular}{cccccl}
\toprule
$c$ & $\kappa$ & $u$ & $\mathrm{Cl}_{2g}(u)$ &
$\dim$ & Regime \\
\midrule
$0$ & $0$ & $0$ &
$\bigwedge^*(\Bbbk^{2g})$ & $2^{2g}$ &
Gravitational \\
$1$ & $1/2$ & $(1/2)\omega_g$ &
$\mathrm{Mat}_{2^g}(\Bbbk)$ & $2^{2g}$ &
Topological \\
$13$ & $13/2$ & $(13/2)\omega_g$ &
$\mathrm{Mat}_{2^g}(\Bbbk)$ & $2^{2g}$ &
Topological (self-dual) \\
$26$ & $13$ & $13\omega_g$ &
$\mathrm{Mat}_{2^g}(\Bbbk)$ & $2^{2g}$ &
Topological (dual: gravitational) \\
\bottomrule
\end{tabular}
}
\end{center}

\medskip\noindent
Genus-$g$ free energy at selected central charges:
\begin{center}
\renewcommand{\arraystretch}{1.3}
{\small
\begin{tabular}{ccccc}
\toprule
$c$ & $F_1$ & $F_2$ & $F_3$ &
$F_g$ \\
\midrule
$0$ & $0$ & $0$ & $0$ & $0$ \\
$1$ & $\frac{1}{48}$ & $\frac{7}{11520}$ &
$\frac{31}{1935360}$ &
$\frac{1}{2}\lambda_g^{\mathrm{FP}}$ \\
$13$ & $\frac{13}{48}$ & $\frac{91}{11520}$ &
$\frac{403}{1935360}$ &
$\frac{13}{2}\lambda_g^{\mathrm{FP}}$ \\
$26$ & $\frac{13}{24}$ & $\frac{91}{5760}$ &
$\frac{403}{967680}$ &
$13\,\lambda_g^{\mathrm{FP}}$ \\
\bottomrule
\end{tabular}
}
\end{center}
where $\lambda_g^{\mathrm{FP}} = \frac{2^{2g-1}-1}{2^{2g-1}} \cdot
\frac{|B_{2g}|}{(2g)!}$. These agree with Volume~I
Theorem~\ref{V1-thm:vir-all-genera} and Volume~I
Computation~\ref{V1-comp:vir-physical-cc}.
\end{computation}

% ======================================================================
%
% SUBSECTION: THE GHOST SECTOR AT c = 26
%
% ======================================================================

\subsection{The exterior algebra structure at $c = 26$}
% label removed: subsec:exterior-c26
\index{critical string!exterior algebra structure}

At $c = 26$, the algebra itself is in the topological regime, but its
\emph{dual} $\mathrm{Vir}_0$ is in the gravitational regime. The ghost
sector of the combined Koszul pair is an exterior algebra.

\begin{theorem}[Ghost sector at $c = 26$; \ClaimStatusProvedHere]
% label removed: thm:ghost-sector-c26
\index{ghost sector!c=26@$c=26$|textbf}
For the Koszul dual $\mathrm{Vir}_{26}^! \simeq \mathrm{Vir}_0$
at genus~$g$:
\begin{enumerate}[label=\textup{(\roman*)}]
\item The Clifford completion of $\mathrm{Vir}_0$ is
\begin{equation}% label removed: eq:g9-ghost-sector
G_g(\mathrm{Vir}_0)
\cong \barB^{(g)}(\mathrm{Vir}_0) \otimes
\bigwedge\nolimits^*\bigl(H_1(\Sigma_g, \Bbbk)\bigr).
\end{equation}

\item The generators are identified with $H_1$-classes:
$\alpha_k \leftrightarrow [A_k]$, $\beta_k \leftrightarrow [B_k]$.

\item The top exterior power
$\bigwedge^{2g}(H_1(\Sigma_g)) \cong \Bbbk$ is generated by
$\mathrm{vol}_g = \alpha_1\beta_1 \cdots \alpha_g\beta_g$.

\item The ghost-number operator
$N_{\mathrm{gh}} = \sum_{k=1}^g(\alpha_k\partial_{\alpha_k} +
\beta_k\partial_{\beta_k})$ has eigenvalues $0, 1, \ldots, 2g$.
The ghost-number anomaly is $2g$: the path-integral measure requires
insertion of $\mathrm{vol}_g$ to produce a nonzero result.
\end{enumerate}
\end{theorem}

\begin{proof}
\emph{(i)} is Theorem~\ref{V1-thm:gravitational-regime}(iii) applied to
$\mathrm{Vir}_0$ (where $\kappa = 0$ and $\lambda = 0$).

\emph{(ii)} The handle operators are fiber integrals of $\eta^{(g)}$
over cycles. When $u = 0$, these satisfy purely exterior algebra
relations, identifying them with $H_1(\Sigma_g, \Bbbk)$.

\emph{(iii)} $\bigwedge^{2g}(H_1) \cong \det(H_1)$ is one-dimensional.

\emph{(iv)} The ghost-number operator is the standard
$\bigwedge$-grading.
\end{proof}

\begin{remark}[Comparison with BRST ghost counting]
% label removed: rem:g9-brst-comparison
\index{BRST cohomology!c=26@$c=26$ comparison}
In the BRST approach, the ghost system has $c_{\mathrm{ghost}} = -26$,
giving total $c_{\mathrm{matter}} + c_{\mathrm{ghost}} = 0$. The
genus-$g$ ghost determinant produces the factor
$\bigwedge^{3g-3}(H^0(\Sigma_g, K^2))$ in the path-integral measure.
The bar-cobar ghost sector is $\bigwedge^{2g}(H_1(\Sigma_g))$, not
$\bigwedge^{3g-3}$. The discrepancy arises from different
gradings: bar-cobar ghosts count \emph{homological} cycles
(intersection form on $H_1$); BRST ghosts count \emph{conformal}
zero modes ($b$-ghost on $H^0(K^2)$). The Riemann--Roch relation
$\dim H^0(K^2) = 3g-3$ (for $g \geq 2$) connects the two counts
via Serre duality and the period matrix.
\end{remark}

% ======================================================================
%
% SUBSECTION: SELF-DUALITY AT c = 13
%
% ======================================================================

\subsection{Self-duality at $c = 13$}
% label removed: subsec:self-duality-c13
\index{Virasoro!c=13@$c=13$!self-duality|textbf}

The Feigin--Frenkel involution $c \mapsto 26 - c$ has its fixed point
at $c = 13$. At this central charge, the Virasoro is Koszul self-dual
in the \emph{curved} sense.

\begin{theorem}[Virasoro self-duality at $c = 13$;
\ClaimStatusProvedHere]
% label removed: thm:vir-self-duality-c13
\index{Virasoro!c=13@$c=13$!self-duality theorem}
At $c = 13$:
\begin{enumerate}[label=\textup{(\roman*)}]
\item $\mathrm{Vir}_{13}^! \simeq \mathrm{Vir}_{13}$: Koszul
self-duality.

\item $\kappa(\mathrm{Vir}_{13}) = 13/2$: symmetric anomaly budget,
$\kappa = \kappa' = 13/2$.

\item The secondary anomalies coincide:
$u_{13} = u'_{13} = (13/2) \cdot \omega_g$.

\item The Clifford completions are isomorphic:
$G_g(\mathrm{Vir}_{13}) \cong G_g(\mathrm{Vir}_{13}^!)$, both
topological.

\item Symmetric complementarity:
\begin{equation}% label removed: eq:g9-c13-complementarity
F_g(\mathrm{Vir}_{13}) = F_g(\mathrm{Vir}_{13}^!)
= \frac{13}{2} \cdot \lambda_g^{\mathrm{FP}},
\end{equation}
with total $F_g + F_g' = 13 \cdot \lambda_g^{\mathrm{FP}}$.

\item Complementarity asymmetry vanishes: $\delta_\kappa = 0$.
\end{enumerate}
\end{theorem}

\begin{proof}
\emph{(i)} $\mathrm{Vir}_c^! \simeq \mathrm{Vir}_{26-c}$ at $c = 13$
gives $\mathrm{Vir}_{13}^! \simeq \mathrm{Vir}_{13}$.
\emph{(ii)--(vi)} follow from (i) and the definitions.
\end{proof}

\begin{remark}[Two self-duality points]
% label removed: rem:g9-two-self-dualities
\index{Virasoro!two self-dualities}
The Virasoro has two distinct self-duality points:
\begin{enumerate}[label=\textup{(\alph*)}]
\item \emph{Quadratic self-duality at $c = 0$}
(Theorem~\ref{V1-thm:virasoro-self-duality}). Here $\kappa = 0$: the bar
complex is a genuine cochain complex, self-duality is of \emph{classical}
(quadratic) type ($\mathrm{Vir}_0^{!_{\mathrm{quad}}} \simeq \mathrm{Vir}_0$),
and the Clifford completion is in the gravitational
regime. The chiral Koszul dual remains
$\mathrm{Vir}_0^! = \mathrm{Vir}_{26} \neq \mathrm{Vir}_0$.

\item \emph{Chiral Koszul self-duality at $c = 13$}
(Theorem~\ref{V1-thm:vir-self-duality-c13}). Here $\kappa = 13/2 \neq 0$:
the bar complex is curved, self-duality is of \emph{modular}
(Feigin--Frenkel) type, and the Clifford completion is in the topological
regime.
\end{enumerate}
At $c = 0$, ghosts are free (exterior algebra); at $c = 13$, ghosts
are entangled (Clifford algebra with $u = (13/2)\omega_g \neq 0$).
The two self-dualities are genuinely different: $c = 0$ kills the
curvature; $c = 13$ distributes it symmetrically.
\end{remark}

% ======================================================================
%
% SUBSECTION: THE DICHOTOMY SUMMARY
%
% ======================================================================

\subsection{The critical dichotomy: structural summary}
% label removed: subsec:dichotomy-summary
\index{critical string!dichotomy!summary}

\begin{theorem}[Critical string dichotomy: structural summary;
\ClaimStatusProvedHere]
% label removed: thm:critical-dichotomy-summary
\index{critical string!dichotomy!structural summary}
The genus-$g$ algebraic structure of the Virasoro Koszul pair
$(\mathrm{Vir}_c, \mathrm{Vir}_{26-c})$ is controlled by
$\kappa = c/2$, $\kappa' = (26-c)/2$,
$\delta_\kappa = c - 13$, and falls into the following
pattern:

\begin{center}
\renewcommand{\arraystretch}{1.4}
\begin{tabular}{lccc}
\toprule
& $c < 13$ & $c = 13$ & $c > 13$ \\
\midrule
$\delta_\kappa$ & $< 0$ & $= 0$ & $> 0$ \\
Dominant anomaly & Dual $\mathrm{Vir}_{26-c}$ &
Equal & $\mathrm{Vir}_c$ \\
Self-dual? & No & Yes & No \\
Both regimes & Topological & Topological & Topological \\
\bottomrule
\end{tabular}
\end{center}

\smallskip\noindent
The two boundary cases:
\begin{center}
\renewcommand{\arraystretch}{1.4}
\begin{tabular}{lcc}
\toprule
& $c = 0$ & $c = 26$ \\
\midrule
$\delta_\kappa$ & $-13$ & $+13$ \\
$\mathrm{Vir}_c$ regime & Gravitational & Topological \\
$\mathrm{Vir}_{26-c}$ regime & Topological & Gravitational \\
Ghost sector & $\bigwedge^*(H_1)$ at $c=0$ &
$\bigwedge^*(H_1)$ at $c'=0$ \\
Anomaly budget & Entirely on dual & Entirely on algebra \\
\bottomrule
\end{tabular}
\end{center}

\smallskip\noindent
The dichotomy is:

\emph{At $c = 26$, the Koszul dual $\mathrm{Vir}_0$ is in the
gravitational regime \textup{(}$u = 0$, exterior algebra ghosts,
uncurved bar complex\textup{)}, while $\mathrm{Vir}_{26}$ itself
carries the entire anomaly budget. At $c = 0$, the roles are
reversed. At $c = 13$, the anomaly is shared equally and the system
is self-dual.}
\end{theorem}

\begin{proof}
Compilation of Volume~I
Theorems~\ref{V1-thm:critical-string-dichotomy},
\ref{V1-thm:ghost-sector-c26},
and~\ref{V1-thm:vir-self-duality-c13}, together with Volume~I
Proposition~\ref{V1-prop:vir-complementarity}.
\end{proof}

% ======================================================================
%
% SUBSECTION: SHADOW TOWER AT CRITICAL POINTS
%
% ======================================================================

\subsection{The shadow obstruction tower at the critical points}
% label removed: subsec:shadow-critical
\index{shadow tower!critical points}

The critical string dichotomy interacts with the shadow obstruction tower
$\Theta_\cA^{\leq r}$ through the modular characteristic
$\kappa = \Theta_\cA^{\leq 2}$.

\begin{proposition}[Shadow obstruction tower at critical central charges;
\ClaimStatusProvedHere]
% label removed: prop:shadow-critical
\index{shadow tower!critical central charges}
\begin{enumerate}[label=\textup{(\roman*)}]
\item \emph{At $c = 0$}: $\kappa = 0$. The quartic pole vanishes, so
the OPE is purely quadratic ($T_{(3)}T = 0$) and the bar complex is
uncurved. The universal genus-$1$ scalar term vanishes, but
$\kappa = 0$ does not by itself determine the higher-degree shadow
tower.

\item \emph{At $c = 13$}: $\kappa = 13/2 \neq 0$. The shadow obstruction tower
is nonzero at every level: $\Theta^{\leq 2} = 13/2$,
$\Theta^{\leq 4}$ involves the quartic resonance class
$Q^{\mathrm{contact}} = 10/[13(5\cdot 13 + 22)] = 10/1131$.
The tower is infinite ($r_{\max} = \infty$, mixed type) and
self-dual: $\Theta_{\mathrm{Vir}_{13}}^{\leq r} =
\Theta_{\mathrm{Vir}_{13}^!}^{\leq r}$ at every level.

\item \emph{At $c = 26$}: $\kappa = 13 \neq 0$. The quartic
contact invariant is
$Q^{\mathrm{contact}} = 10/[26 \cdot 152] = 5/1976$.
The dual algebra $\mathrm{Vir}_0$ lies on the uncurved scalar locus,
so $\Theta_{\mathrm{Vir}_0}^{\min} = 0$ and
$F_1(\mathrm{Vir}_0) = 0$. This scalar datum does not determine the
dual higher-degree shadow obstruction tower. Thus the nonzero scalar and quartic
data exhibited here come from $\mathrm{Vir}_{26}$, while the dual
higher-degree tower is not fixed by $\kappa(\mathrm{Vir}_0)=0$ alone.
\end{enumerate}
\end{proposition}

\begin{proof}
\emph{(i)} $\kappa(\mathrm{Vir}_0) = 0$; at $c = 0$ the quartic pole
vanishes and the OPE is quadratic. Hence the scalar class
$\Theta_{\mathrm{Vir}_0}^{\min}$ and the universal genus-$1$ scalar
term vanish, but no conclusion about the full higher-degree tower
follows from $\kappa(\mathrm{Vir}_0)=0$ alone.

\emph{(ii)} Self-duality implies $\Theta_{\mathrm{Vir}_{13}}^{\leq r}
= \Theta_{\mathrm{Vir}_{13}^!}^{\leq r}$. The quartic invariant
evaluates via $Q^{\mathrm{contact}} = 10/[c(5c+22)]|_{c=13}$.

\emph{(iii)} $\kappa(\mathrm{Vir}_0) = 0$ gives the vanishing of the
dual scalar class $\Theta_{\mathrm{Vir}_0}^{\min}$ and of
$F_1(\mathrm{Vir}_0)$; the higher-degree dual tower is not determined
by this scalar input alone.
\end{proof}

\begin{remark}[Depth-zero resonance shadow revisited]
% label removed: rem:g9-depth-zero-revisited
\index{depth-zero resonance shadow!revisited}
The same-family shadow partner $\mathrm{Vir}_{26-c}$ controls
M/S-level complementarity
(Remark~\ref{V1-rem:virasoro-resonance-model}). At $c = 26$, this
partner is $\mathrm{Vir}_0$, the dual algebra on the uncurved scalar
locus, not a trivial algebra. The depth-zero resonance shadow
saturates only at the scalar/genus-$1$ level; higher-degree data of the
dual partner are not fixed by that cancellation. This is the
algebraic content of the critical string: the universal scalar budget
closes at $c = 26$, while the non-scalar shadow obstruction tower still requires
separate control.
\end{remark}

% ======================================================================
%
% SUBSECTION: HODGE-CLIFFORD COMPATIBILITY
%
% ======================================================================

\subsection{The Hodge-theoretic content}
% label removed: subsec:hodge-dichotomy
\index{Hodge theory!critical string dichotomy}

\begin{proposition}[Hodge--Clifford compatibility;
\ClaimStatusProvedHere]
% label removed: prop:hodge-clifford
\index{Hodge theory!Clifford compatibility}
The Clifford algebra $\mathrm{Cl}_{2g}(u)$ is compatible with the
Hodge decomposition $H^1(\Sigma_g, \C) = H^{1,0} \oplus H^{0,1}$:
\begin{enumerate}[label=\textup{(\roman*)}]
\item The $A$-cycle operators $\alpha_k$ correspond to
$H^{1,0}(\Sigma_g)$ (holomorphic differentials).

\item The $B$-cycle operators $\beta_k$ correspond to
$H^{0,1}(\Sigma_g)$ (antiholomorphic differentials).

\item The Clifford relation $\alpha_i \beta_j + \beta_j \alpha_i =
\delta_{ij}\, u$ encodes the Hodge--Riemann relation:
\[
\int_{\Sigma_g} \omega_i \wedge \bar\omega_j = \delta_{ij}\, \omega_g,
\]
where $\{\omega_1, \ldots, \omega_g\}$ are the normalized abelian
differentials ($\oint_{A_l}\omega_k = \delta_{kl}$).

\item The period matrix $\Omega_{kl} = \oint_{B_l}\omega_k$ appears as
the overlap $\langle \alpha_k, \beta_l \rangle = \Omega_{kl}$ in the
complexified Clifford algebra.
\end{enumerate}
\end{proposition}

\begin{proof}
\emph{(i)--(ii)} The symplectic basis $(A_k, B_k)$ is Poincar\'e dual
to the Hodge basis $(\omega_k, \bar\omega_k)$.

\emph{(iii)} The Riemann bilinear relation gives
$\frac{i}{2}\int_{\Sigma_g}\omega_k \wedge \bar\omega_l =
(\operatorname{Im}\Omega)_{kl}$ in the Arakelov normalization.
Integration against $u = \kappa \cdot \omega_g$ yields
$\kappa \cdot \delta_{kl} = \delta_{kl}\, u$, matching the
Clifford relation.

\emph{(iv)} $\Omega_{kl} = \oint_{B_l}\omega_k$ determines the
$\alpha_k$--$\beta_l$ overlap in the complexification.
\end{proof}

\begin{remark}[The Arakelov metric as Clifford trace]
% label removed: rem:g9-arakelov-clifford
\index{Arakelov metric!Clifford trace}
The Arakelov form $\omega_g = \frac{i}{2}\sum_{k=1}^g
\omega_k \wedge \bar\omega_k$
(Proposition~\ref{V1-prop:genus-g-curvature-package}) is the trace of the
Clifford form: $\omega_g = \frac{1}{2g}\sum_{k=1}^g
(\alpha_k\beta_k + \beta_k\alpha_k)|_{u=1}$.
\end{remark}

% ======================================================================
%
% SUBSECTION: CLIFFORD SPECTRAL SEQUENCE
%
% ======================================================================

\subsection{The spectral sequence of the Clifford filtration}
% label removed: subsec:clifford-spectral
\index{spectral sequence!Clifford filtration}

\begin{construction}[Clifford spectral sequence]
% label removed: constr:clifford-spectral
\index{spectral sequence!Clifford|textbf}
The Clifford degree filtration on $G_g(\cA)$,
\begin{equation}% label removed: eq:g9-clifford-filtration
F^p G_g := B_\Theta^{(g)}(\cA) \otimes
\bigoplus_{k \leq p} \mathrm{Cl}_{2g}^{(k)}(u),
\end{equation}
where $\mathrm{Cl}_{2g}^{(k)}$ is the span of products of at most $k$
generators, is compatible with $d_{G_g}$ (since handle operators are
closed). The $E_1$ page is
\begin{equation}% label removed: eq:g9-clifford-E1
E_1^{p,q} = H^{p+q}(B_\Theta^{(g)}(\cA))
\otimes \mathrm{gr}^p_{\mathrm{Cl}}\,\mathrm{Cl}_{2g}(u).
\end{equation}
\end{construction}

\begin{proposition}[Degeneration; \ClaimStatusProvedHere]
% label removed: prop:clifford-spectral-degeneration
\begin{enumerate}[label=\textup{(\roman*)}]
\item \emph{Gravitational regime} ($u = 0$): degeneration at $E_1$;
$\mathrm{gr}_{\mathrm{Cl}}^* = \bigwedge^*(\Bbbk^{2g})$; hence
$H^*(G_g|_{u=0}) \cong H^*(B) \otimes \bigwedge^*(\Bbbk^{2g})$.

\item \emph{Topological regime} ($u$ invertible): degeneration at
$E_1$; $\mathrm{gr}_{\mathrm{Cl}}^*$ is the associated graded of a
matrix algebra; $H^*(G_g[u^{-1}]) \cong
\mathrm{Mat}_{2^g}(H^*(B_\Theta[u^{-1}]))$.
\end{enumerate}
\end{proposition}

\begin{proof}
In both cases, the handle operators are closed and the Clifford
multiplication preserves the filtration, so $d_r = 0$ for $r \geq 1$.
\end{proof}

% ======================================================================
%
% SUBSECTION: BOSONIC STRING IDENTIFICATIONS
%
% ======================================================================

\subsection{The $c = 26$ algebra and the bosonic string}
% label removed: subsec:bosonic-string
\index{bosonic string!bar-cobar interpretation|textbf}

\begin{proposition}[Bosonic string structural identifications;
\ClaimStatusProvedHere]
% label removed: prop:bosonic-string-identifications
\index{bosonic string!structural identifications}
At $c = 26$:
\begin{enumerate}[label=\textup{(\roman*)}]
\item \emph{Anomaly budget.}\;
$F_g(\mathrm{Vir}_{26}) = 13 \cdot \lambda_g^{\mathrm{FP}}$,
$F_g(\mathrm{Vir}_0) = 0$.

\item \emph{Koszul pair.}\;
$(\mathrm{Vir}_{26}, \mathrm{Vir}_0)$ with $\kappa + \kappa' = 13$,
maximally asymmetric: $|\delta_\kappa| = 13$.

\item \emph{Ghost sector.}\;
$\bigwedge^*(H_1(\Sigma_g)) \cong \bigwedge^{2g}$, ghost-number
anomaly $2g$.

\item \emph{Generating function.}\;
$\sum_{g=1}^\infty F_g(\mathrm{Vir}_{26})\,x^{2g}
= 13(x/(2\sin(x/2)) - 1)$
with radius of convergence $|x| = 2\pi$.

\item \emph{Transgression.}\;
$B_\Theta^{(g)}(\mathrm{Vir}_{26})$ has $u = 13\omega_g$
(strongly topological); $B_\Theta^{(g)}(\mathrm{Vir}_0)$ has
$\eta^2 = 0$ (uncurved).

\item \emph{Combined Clifford.}\;
\begin{equation}% label removed: eq:g9-combined-final
G_g(\mathrm{Vir}_{26}) \otimes G_g(\mathrm{Vir}_0)
\cong
\mathrm{Mat}_{2^g}(B_\Theta^{(g)}(\mathrm{Vir}_{26})[u^{-1}])
\otimes \barB^{(g)}(\mathrm{Vir}_0)
\otimes \bigwedge\nolimits^*(\Bbbk^{2g}).
\end{equation}

\item \emph{Physical content.}\; Morita reduction gives
\begin{equation}% label removed: eq:g9-physical-content
D(G_g(\mathrm{Vir}_{26}) \otimes G_g(\mathrm{Vir}_0)
\text{-}\mathsf{mod})
\simeq
D\bigl(B_\Theta[u^{-1}] \otimes \barB(\mathrm{Vir}_0) \otimes
\bigwedge\nolimits^*(\Bbbk^{2g})\text{-}\mathsf{mod}\bigr).
\end{equation}
\end{enumerate}
\end{proposition}

\begin{proof}
All items follow from the established theorems:
(i) Theorem~\ref{V1-thm:vir-all-genera};
(ii) complementarity;
(iii) Theorem~\ref{V1-thm:ghost-sector-c26};
(iv) Remark~\ref{V1-rem:vir-trichotomy-genera};
(v) Definitions~\ref{V1-def:transgression-algebra}
and~\ref{V1-def:bar-transgression-complex};
(vi) Theorems~\ref{V1-thm:topological-regime}
and~\ref{V1-thm:gravitational-regime};
(vii) Theorem~\ref{V1-thm:topological-regime}(iii).
\end{proof}

% ======================================================================
%
% SUBSECTION: KM CONTRAST AND W_N EXTENSION
%
% ======================================================================

\subsection{The affine Kac--Moody contrast}
% label removed: subsec:km-contrast
\index{affine Kac--Moody!critical string contrast}

\begin{remark}[Complementarity sum comparison]
% label removed: rem:g9-km-contrast
\index{complementarity!sum comparison}
\begin{center}
\renewcommand{\arraystretch}{1.3}
\begin{tabular}{lccc}
\toprule
Family & $\kappa + \kappa'$ & Self-dual point &
Gravitational locus \\
\midrule
$\widehat{\fg}_k$ &
$0$ & $k = -h^\vee$ (critical; excluded) &
$\kappa = 0$ iff $k = -h^\vee$ \\
$\mathrm{Vir}_c$ &
$13$ & $c = 13$ &
$\kappa = 0$ iff $c = 0$ \\
$\mathcal{W}_N$ &
$\ne 0$ (generically) & exists for each $N$ &
$\kappa = 0$ at isolated $c$ \\
\bottomrule
\end{tabular}
\end{center}
For affine Kac--Moody at non-critical level, $\kappa + \kappa' = 0$
and $\kappa \neq 0$: there is no gravitational regime. The
critical string dichotomy is intrinsically a Virasoro/$\cW_N$
phenomenon.
\end{remark}

\subsection{The $\mathcal{W}_N$ generalization}
% label removed: subsec:wn-dichotomy
\index{W-algebra@$\mathcal{W}$-algebra!critical dichotomy}

\begin{remark}[$\mathcal{W}_N$ complementarity sums]
% label removed: rem:g9-wn-sums
\index{W-algebra@$\mathcal{W}$-algebra!complementarity sums}
The modular characteristic is $\kappa(\mathcal{W}_{N,c}) = c \cdot
\varrho(\mathfrak{sl}_N)$ where $\varrho = \sum_{i=1}^{N-1}(m_i+1)^{-1}$
is the anomaly ratio.
\begin{center}
\renewcommand{\arraystretch}{1.3}
\begin{tabular}{ccccc}
\toprule
$N$ & $\varrho(\mathfrak{sl}_N)$ & $\kappa$ &
$\kappa + \kappa'$ & Self-dual $c$ \\
\midrule
$2$ & $1/2$ & $c/2$ & $13$ & $13$ \\
$3$ & $5/6$ & $5c/6$ & $250/3$ & $50$ \\
$N$ & $H_N - 1$ & $c(H_N - 1)$ &
$c_{\mathrm{crit}}(H_N-1)$ & $c_{\mathrm{crit}}/2$ \\
\bottomrule
\end{tabular}
\end{center}
Here $H_N = \sum_{k=1}^N 1/k$ is the harmonic number. The
transgression algebra, Clifford completion, and two-regime dichotomy
carry through with $\kappa = c \cdot \varrho$ in place of $c/2$.
\end{remark}

% ======================================================================
%
% SUBSECTION: HOLOGRAPHIC INTERPRETATION
%
% ======================================================================

\subsection{Relation to the holographic datum}
% label removed: subsec:holographic-relation
\index{holographic modular Koszul datum!critical string}

\begin{remark}[Holographic interpretation]
% label removed: rem:g9-holographic
\index{holographic modular Koszul datum!critical string interpretation}
In the holographic setting
(\S\ref{V1-sec:frontier-modular-holography-platonic}):
\begin{enumerate}[label=\textup{(\alph*)}]
\item The \emph{topological regime} ($u$ invertible) corresponds to
the interacting bulk: nontrivial curvature, nonzero shadow obstruction tower,
bulk gravity determined by the Yangian $r(z)$ and the shadow
connection $\nabla^{\mathrm{hol}}$.

\item The \emph{gravitational regime} ($u = 0$) corresponds to the
degenerate bulk: uncurved boundary, vanishing shadow obstruction tower, the bulk
reduces to a topological theory with ghost sector
$\bigwedge^*(H_1(\Sigma_g))$.

\item The \emph{self-dual point} ($\delta_\kappa = 0$)
corresponds to self-conjugacy under gravitational $S$-duality (G4):
the bulk and defect sectors contribute equally.

\item At $c = 26$, the holographic datum
$\mathcal{H}(\mathrm{Vir}_{26})$ is maximally asymmetric: the defect
sector ($\mathrm{Vir}_0$, gravitational) is trivial, and the bulk
sector ($\mathrm{Vir}_{26}$, topological) carries all structure. This
is the bosonic string from the holographic perspective.
\end{enumerate}
\end{remark}

% ======================================================================
%
% SUBSECTION: THE LOCALIZATION SEQUENCE
%
% ======================================================================

\subsection{The localization sequence}
% label removed: subsec:localization-sequence
\index{localization sequence!Clifford completion}

The two regimes are connected by a localization sequence in the
secondary anomaly $u$. This is the algebraic analog of the passage
from generic central charge to the critical point.

\begin{theorem}[Localization sequence; \ClaimStatusProvedHere]
% label removed: thm:localization-sequence
\index{localization sequence!theorem|textbf}
There is a localization fiber sequence of dg algebras:
\begin{equation}% label removed: eq:g9-localization
G_g(\cA)/(u)
\;\longrightarrow\;
G_g(\cA)
\;\longrightarrow\;
G_g(\cA)[u^{-1}],
\end{equation}
with the following properties:
\begin{enumerate}[label=\textup{(\roman*)}]
\item The left term $G_g(\cA)/(u) = \barB^{(g)}(\cA) \otimes
\bigwedge^*(\Bbbk^{2g})$ is the gravitational regime:
the Clifford algebra with $u = 0$.

\item The right term $G_g(\cA)[u^{-1}] \simeq_{\mathrm{Morita}}
\mathrm{Mat}_{2^g}(B_\Theta[u^{-1}])$ is the topological regime.

\item The long exact sequence in cohomology reads:
\begin{equation}% label removed: eq:g9-localization-les
\cdots \to
H^n(G_g/(u)) \xrightarrow{\iota}
H^n(G_g) \xrightarrow{j}
H^n(G_g[u^{-1}]) \xrightarrow{\partial}
H^{n+1}(G_g/(u)) \to \cdots
\end{equation}
The connecting homomorphism $\partial$ measures the failure of
cohomology classes in the topological regime to extend to the
gravitational regime.
\end{enumerate}
\end{theorem}

\begin{proof}
The sequence~\eqref{V1-eq:g9-localization} is the standard localization
for a central element $u$ in a dg algebra: the fiber of
$G_g \to G_g[u^{-1}]$ is $G_g/(u^\infty) = \mathrm{colim}_n\,
G_g/(u^n)$. Since $u$ is central and $u^2 = \kappa^2 \cdot \omega_g^2
\in B$ (where $\omega_g^2$ is a $(2,2)$-class on $\Sigma_g$), the
$u$-torsion stabilizes after finitely many steps. For the Clifford
algebra, setting $u = 0$ in the
relations~\eqref{V1-eq:g9-clifford-ab}--\eqref{V1-eq:g9-clifford-bb} gives
the exterior algebra, confirming (i). Part (ii) is
Theorem~\ref{V1-thm:topological-regime}. Part (iii) is the long exact
sequence of the fiber-cofiber sequence in the dg category.
\end{proof}

\begin{corollary}[Localization for the Virasoro; \ClaimStatusProvedHere]
% label removed: cor:g9-localization-vir
\index{localization sequence!Virasoro}
For the Virasoro, the localization sequence specializes to:
\begin{equation}% label removed: eq:g9-localization-vir
\barB^{(g)}(\mathrm{Vir}_c) \otimes \bigwedge\nolimits^{2g}
\;\longrightarrow\;
G_g(\mathrm{Vir}_c)
\;\longrightarrow\;
\mathrm{Mat}_{2^g}(B_\Theta^{(g)}(\mathrm{Vir}_c)[(c/2)^{-1}]).
\end{equation}
The connecting homomorphism is
\begin{equation}% label removed: eq:g9-connecting-vir
\partial\colon
H^n(\mathrm{Mat}_{2^g}(B_\Theta[(c/2)^{-1}]))
\longrightarrow
H^{n+1}(\barB^{(g)}(\mathrm{Vir}_c))
\otimes \bigwedge\nolimits^{2g}.
\end{equation}
At $c = 0$ the right term vanishes (the localization
$B_\Theta[u^{-1}]$ is trivial when $u = 0$), so the left term
$\barB^{(g)}(\mathrm{Vir}_0) \otimes \bigwedge^{2g}$ accounts for all
cohomology. At $c = 26$ the left term has trivial
differential (uncurved), so the connecting homomorphism
$\partial$ captures the entire obstruction to passing from the
topological regime to the gravitational regime.
\end{corollary}

\begin{proof}
Set $u = (c/2)\omega_g$. The left term is $G_g/(u)$; the right is
$G_g[u^{-1}]$. The statements about $c = 0$ and $c = 26$ follow from
the vanishing of $u$ and the maximality of $\kappa$ respectively.
\end{proof}

% ======================================================================
%
% SUBSECTION: PRIMITIVITY AND BIALGEBRA STRUCTURE
%
% ======================================================================

\subsection{Primitivity of the secondary anomaly}
% label removed: subsec:primitivity
\index{secondary anomaly!primitivity}
\index{bialgebra!transgression algebra}

The bar coalgebra $\barB^{(g)}(\cA)$ carries a coalgebra structure
from the deconcatenation coproduct. The transgression element $\eta$
and the secondary anomaly $u = \eta^2$ interact with this coproduct in
a controlled way.

\begin{proposition}[Bialgebra extension; \ClaimStatusProvedHere]
% label removed: prop:g9-bialgebra
\index{bialgebra!transgression extension}
The transgression algebra $B_\Theta^{(g)}(\cA)$ carries a unique
dg bialgebra structure extending the coalgebra on
$\barB^{(g)}(\cA)$, determined by
\begin{equation}% label removed: eq:g9-eta-coproduct
\Delta(\eta) = \eta \otimes 1 + 1 \otimes \eta,
\qquad
\varepsilon(\eta) = 0.
\end{equation}
In particular, $\eta$ is primitive.
\end{proposition}

\begin{proof}
The coproduct~\eqref{V1-eq:g9-eta-coproduct} is forced by the
requirement that $\Delta$ be an algebra map: $\Delta(\eta^2) =
\Delta(\eta)^2 = (\eta \otimes 1 + 1 \otimes \eta)^2 =
\eta^2 \otimes 1 + \eta \otimes \eta + \eta \otimes \eta +
1 \otimes \eta^2 = \lambda \otimes 1 + 2\eta \otimes \eta +
1 \otimes \lambda$. Here $\lambda = \kappa \cdot \omega_g$ is
the curvature class on $\overline{\cM}_g$, not a fibrewise
coalgebra element: it lives in $H^{2}(\overline{\cM}_g)$ via
the Hodge bundle, and its primitivity in $B_\Theta^{(g)}(\cA)$
is the statement that the
period-corrected differential
$\Dg{g}$ of Volume~I,
Remark~\ref{V1-rem:sc-higher-genus}\textup{(iii)}, is a
coderivation of~$\Delta$. The fibrewise differential
$\dfib$ is \emph{not} a coderivation
(Volume~I,
Remark~\ref{V1-rem:sc-higher-genus}\textup{(ii)}); only after
the Fay trisecant period correction does the curved
$E_1$ dg coalgebra structure persist, and only then does
the Hodge curvature class push forward to a primitive element
of $B_\Theta^{(g)}(\cA)$. The counit $\varepsilon(\eta) = 0$
follows from $\varepsilon(\lambda) = 0$. Coassociativity:
$(\Delta \otimes \id)\Delta(\eta) = \eta \otimes 1 \otimes 1 +
1 \otimes \eta \otimes 1 + 1 \otimes 1 \otimes \eta =
(\id \otimes \Delta)\Delta(\eta)$.
\end{proof}

\begin{corollary}[Primitivity of the secondary anomaly;
\ClaimStatusProvedHere]
% label removed: cor:g9-u-coproduct
The secondary anomaly element is primitive:
\begin{equation}% label removed: eq:g9-u-coproduct
\Delta(u) = u \otimes 1 + 1 \otimes u.
\end{equation}
\end{corollary}

\begin{proof}
Direct computation from Proposition~\ref{V1-prop:g9-bialgebra}.
Using $\Delta(\eta) = \eta \otimes 1 + 1 \otimes \eta$ and the
graded multiplication in $B_\Theta \otimes B_\Theta$
(Koszul sign rule: $(a \otimes b)(c \otimes d)
= (-1)^{|b||c|}\,ac \otimes bd$):
\begin{align*}
\Delta(u) &= \Delta(\eta)^2
= (\eta \otimes 1 + 1 \otimes \eta)^2 \\
&= \eta^2 \otimes 1
 + (\eta \otimes 1)(1 \otimes \eta)
 + (1 \otimes \eta)(\eta \otimes 1)
 + 1 \otimes \eta^2 \\
&= u \otimes 1
 + \eta \otimes \eta
 + (-1)^{|\eta||\eta|}\eta \otimes \eta
 + 1 \otimes u \\
&= u \otimes 1
 + \eta \otimes \eta
 - \eta \otimes \eta
 + 1 \otimes u \\
&= u \otimes 1 + 1 \otimes u.
\end{align*}
The cross terms cancel because $|\eta| = 1$ is odd:
$(1 \otimes \eta)(\eta \otimes 1) = (-1)^{|\eta||\eta|}\eta \otimes \eta
= -\eta \otimes \eta$.
\end{proof}

\begin{remark}[Primitivity and the Milnor--Moore theorem]
% label removed: rem:g9-milnor-moore
\index{Milnor--Moore theorem!transgression}
When $B_\Theta$ is connected and locally finite, the Milnor--Moore
theorem identifies the primitive elements of $H^*(B_\Theta)$ with the
indecomposables of the associated Lie algebra. Since $u$ is primitive
(Corollary~\ref{V1-cor:g9-u-coproduct}), its cohomology class
$[u] \in \mathrm{Prim}(H^*(B_\Theta))$ generates a one-dimensional
abelian sub-Lie algebra. Its vanishing ($\kappa = 0$) is the
condition for the Lie algebra to lose this generator, the algebraic
avatar of the critical string.
\end{remark}

% ======================================================================
%
% SUBSECTION: THE TRANSGRESSION AS UNIVERSAL PROPERTY
%
% ======================================================================

\subsection{The universal property of the transgression algebra}
% label removed: subsec:universal-property
\index{transgression algebra!universal property}

\begin{proposition}[Universal property; \ClaimStatusProvedHere]
% label removed: prop:g9-universal
\index{transgression algebra!universal property|textbf}
Let $(B, d, \lambda)$ be a curved dg algebra with
$d^2 = [\lambda, -]$. For any dg algebra $(C, d_C)$ with
$d_C^2 = 0$ and any dg algebra map $f\colon B \to C$ such that
$f(\lambda) = \xi^2$ for some $\xi \in C^1$ with $d_C(\xi) = 0$,
there exists a unique dg algebra map $\tilde{f}\colon B_\Theta \to C$
extending $f$ with $\tilde{f}(\eta) = \xi$:
\[
\begin{tikzcd}
B \arrow[r, hook] \arrow[d, "f"'] &
B_\Theta \arrow[dl, dashed, "\exists!\,\tilde{f}"] \\
C &
\end{tikzcd}
\]
\end{proposition}

\begin{proof}
Define $\tilde{f}(b_0 + b_1\eta) = f(b_0) + f(b_1)\xi$ for
$b_0, b_1 \in B$. This is well-defined on $B_\Theta$ because
$\tilde{f}(\eta^2 - \lambda) = \xi^2 - f(\lambda) = 0$.
It is an algebra map: $\tilde{f}$ preserves the
multiplication~\eqref{V1-eq:g9-BT-multiplication} because $f$ is an
algebra map and $\xi$ centralizes $\im(f)$ (since $f(\lambda) = \xi^2$
is central in $C$, and $\xi$ commutes with $\im(f)$ by the
same centrality argument used for $\eta$). It respects the
differential: $\tilde{f}(d_{B_\Theta}(b)) = \tilde{f}(d(b) +
[\eta, b]) = f(d(b)) + [f(\eta), f(b)] = d_C(f(b)) +
[\xi, f(b)] = d_C(\tilde{f}(b))$ using $d_C \circ f = f \circ d$
(since $f$ is a dg map, with the correction that $d^2 = [\lambda, -]$
maps to $d_C^2 = 0$ via $f(\lambda) = \xi^2$). Uniqueness: any
extension must send $\eta \mapsto \xi$ (there is no choice), and
the value on $B$ is fixed by $f$.
\end{proof}

\begin{corollary}[The comparison map is universal;
\ClaimStatusProvedHere]
% label removed: cor:g9-comparison-universal
The comparison map $B_\Theta^{(g)}(\cA) \to (\barB^{(g)}(\cA),
\Dg{g})$ of Remark~\ref{V1-rem:g9-mc-relation} is the unique factorization
of the identity on $\barB^{(g)}(\cA)$ through the transgression
algebra, with $\eta$ sent to the period-dependent element
$\sum_k t_k \omega_k$ whose square is $\kappa \cdot \omega_g$ (by the
Lagrangian property of the $A$-cycle subspace).
\end{corollary}

% ======================================================================
%
% SUBSECTION: DEFORMATION OF THE CLIFFORD COMPLETION
%
% ======================================================================

\subsection{Deformation theory of the Clifford completion}
% label removed: subsec:clifford-deformation
\index{Clifford completion!deformation theory}

The Clifford completion $G_g(\cA)$ depends on the central charge
parameter through $u = \kappa(\cA) \cdot \omega_g$. The variation of
$G_g$ with $\kappa$ is controlled by the deformation theory of
Clifford algebras.

\begin{proposition}[Deformation of the Clifford algebra;
\ClaimStatusProvedHere]
% label removed: prop:g9-clifford-deformation
\index{Clifford algebra!deformation theory}
\begin{enumerate}[label=\textup{(\roman*)}]
\item The family $\{\mathrm{Cl}_{2g}(u)\}_{u \in \Bbbk}$ is a flat
deformation of $\bigwedge^*(\Bbbk^{2g})$ over $\Bbbk$, with fiber
$\bigwedge^*(\Bbbk^{2g})$ at $u = 0$ and
$\mathrm{Mat}_{2^g}(\Bbbk)$ at $u \neq 0$.

\item The Hochschild cohomology controlling the deformation is
\begin{equation}% label removed: eq:g9-HH-clifford
\HH^2(\mathrm{Cl}_{2g}(u), \mathrm{Cl}_{2g}(u))
=
\begin{cases}
\Bbbk & \text{if } u = 0, \\
0 & \text{if } u \neq 0.
\end{cases}
\end{equation}
The one-dimensional space at $u = 0$ is generated by the
anticommutator deformation $\alpha_i\beta_j + \beta_j\alpha_i
\mapsto \delta_{ij}\,\epsilon$.

\item The deformation is unobstructed: $\HH^3(\bigwedge^*(\Bbbk^{2g}),
\bigwedge^*(\Bbbk^{2g})) = 0$ for the relevant component.
\end{enumerate}
\end{proposition}

\begin{proof}
\emph{(i)} The Clifford algebra is a filtered deformation of the
exterior algebra: the associated graded of $\mathrm{Cl}_{2g}(u)$
(filtered by word length) is $\bigwedge^*(\Bbbk^{2g})$ for all $u$,
including $u = 0$. Flatness is clear: $\mathrm{Cl}_{2g}(u)$ is free
of rank $2^{2g}$ as a $\Bbbk$-module for all $u$.

\emph{(ii)} For $u \neq 0$, $\mathrm{Cl}_{2g}(u) \cong
\mathrm{Mat}_{2^g}(\Bbbk)$, and $\HH^*(\mathrm{Mat}_n, \mathrm{Mat}_n)
\cong \HH^*(\Bbbk, \Bbbk) = \Bbbk$ concentrated in degree $0$ by
Morita invariance. In particular, $\HH^2 = 0$: matrix algebras are
rigid. For $u = 0$, $\HH^2(\bigwedge^*(\Bbbk^{2g})) = \Bbbk$,
generated by the single parameter $u$.

\emph{(iii)} The obstruction space $\HH^3$ for the exterior algebra
vanishes in the relevant weight: the deformation class $u$ has weight
$2$ (it changes the anticommutator from $0$ to $u \cdot \delta_{ij}$),
and the cup product $[u] \cup [u] \in \HH^4$ is the only potential
obstruction, which vanishes by the centrality of $u$.
\end{proof}

\begin{remark}[The wall-crossing at $u = 0$]
% label removed: rem:g9-wall-crossing
\index{wall-crossing!Clifford completion}
The deformation theory reveals that $u = 0$ is a \emph{wall}: for
$u \neq 0$ the algebra is rigid (a matrix algebra), while at $u = 0$
it acquires a one-dimensional deformation parameter. The Clifford
completion ``jumps'' across this wall, from a matrix algebra to an
exterior algebra, and the jump is irreversible (the exterior algebra
is not Morita equivalent to any matrix algebra). This is the algebraic
wall-crossing underlying the critical string dichotomy: the passage
through $\kappa = 0$ is a genuine phase transition in the derived
category.
\end{remark}

% ======================================================================
%
% SUBSECTION: GENUS-g TRACE FORMULA
%
% ======================================================================

\subsection{The genus-$g$ trace formula}
% label removed: subsec:trace-formula
\index{trace formula!Clifford completion}

In the topological regime, the Morita equivalence allows computation of
traces on $G_g$ via traces on $B_\Theta$.

\begin{proposition}[Trace formula; \ClaimStatusProvedHere]
% label removed: prop:g9-trace-formula
\index{trace formula!genus-g@genus-$g$}
When $u$ is invertible, any trace $\mathrm{tr}\colon B_\Theta \to \Bbbk$
extends to a trace on $G_g$ by the formula
\begin{equation}% label removed: eq:g9-trace-formula
\mathrm{Tr}_{G_g}(b \otimes C) =
2^g \cdot \mathrm{tr}(b) \cdot \mathrm{tr}_{\mathrm{Cl}}(C),
\end{equation}
where $\mathrm{tr}_{\mathrm{Cl}}\colon \mathrm{Cl}_{2g}(u)
\to \Bbbk$ is the reduced Clifford trace (normalized so that
$\mathrm{tr}_{\mathrm{Cl}}(1) = 1$).
\end{proposition}

\begin{proof}
Under the Morita equivalence
$G_g[u^{-1}] \cong \mathrm{Mat}_{2^g}(B_\Theta[u^{-1}])$, a trace on
$\mathrm{Mat}_{2^g}(R)$ is $\sum_{i=1}^{2^g}\mathrm{tr}(R_{ii})$.
The factor $2^g$ is the matrix size. The Clifford trace is the
composition of the matrix trace with the identification
$\mathrm{Cl}_{2g}(u) \cong \mathrm{Mat}_{2^g}$.
\end{proof}

\begin{corollary}[Genus-$g$ partition function via trace;
\ClaimStatusProvedHere]
% label removed: cor:g9-partition-trace
\index{partition function!trace formula}
For a Koszul chiral algebra $\cA$ with $\kappa(\cA) \neq 0$, the
genus-$g$ free energy satisfies
\begin{equation}% label removed: eq:g9-Fg-trace
F_g(\cA) = \frac{1}{2^g}\,\mathrm{Tr}_{G_g}\bigl(
(-1)^F \cdot q^{L_0}\bigr)
= \kappa(\cA) \cdot \lambda_g^{\mathrm{FP}},
\end{equation}
where the second equality uses the genus universality theorem
(Theorem~\ref{V1-thm:genus-universality}). The factor $1/2^g$ accounts
for the Morita multiplicity.
\end{corollary}

\begin{remark}[The partition function at $u = 0$]
% label removed: rem:g9-partition-u0
\index{partition function!gravitational regime}
In the gravitational regime ($u = 0$), the trace formula degenerates:
the Clifford trace is replaced by the supertrace on
$\bigwedge^*(\Bbbk^{2g})$, which is
$\mathrm{str} = \sum_{k=0}^{2g}(-1)^k\binom{2g}{k} =
(1-1)^{2g} = 0$ for $g \geq 1$. This vanishing is consistent with
$F_g(\mathrm{Vir}_0) = 0$: the gravitational partition function is
zero because the supertrace over the exterior algebra vanishes.
\end{remark}

% ======================================================================
%
% SUBSECTION: OPEN QUESTIONS
%
% ======================================================================

\subsection{Further directions}
% label removed: subsec:further-directions

\begin{remark}[Open questions]
% label removed: rem:g9-open-questions
\index{critical string!open questions}
\begin{enumerate}[label=\textup{(\alph*)}]
\item Does $G_g(\cA)$ carry a natural BV algebra structure
extending the BV/bar dictionary of Remark~\ref{V1-rem:bv-bar-bridge}? The quantum master
equation on $G_g$ should encode genus-$g$ sewing relations.

\item Is there a super-Clifford version for $\mathcal{N} = 1$ at
$c = 15$, replacing $\mathrm{Cl}_{2g}$ with a super-Clifford algebra
$\mathrm{sCl}_{2g|2g}$?

\item The Hodge--Clifford compatibility
(Proposition~\ref{V1-prop:hodge-clifford}) suggests a connection to the
Mukai Hodge structure on $K(\Sigma_g)$. Does the spectral sequence of
Construction~\ref{V1-constr:clifford-spectral} compute the Mukai pairing?

\item At $c = 13$, does the Clifford completion carry an enhanced
symmetry reflecting self-duality, beyond the $\Z/2$ of the Koszul
involution?

\item The wall-crossing at $u = 0$
(Remark~\ref{V1-rem:g9-wall-crossing}) should have a categorical lift:
a derived equivalence between the gravitational and topological
regimes that degenerates at $u = 0$. Can this be made precise using
the formalism of perverse sheaves on the $\kappa$-line?

\item The localization sequence
(Theorem~\ref{V1-thm:localization-sequence}) suggests a
``semi-classical limit'' interpretation: the gravitational regime is
the $u \to 0$ limit of the topological regime, analogous to the
$\hbar \to 0$ limit in deformation quantization. The secondary
anomaly $u = \kappa \cdot \omega_g$ plays the role of Planck's
constant.
\end{enumerate}
\end{remark}

%=======================================================================
% c = 26 no-ghost theorem as a chain-level Clifford-degeneration computation
%=======================================================================

\subsection{The $c=26$ BRST anomaly and Clifford degeneration}
\label{sec:critical-string-clifford-degeneration}

The preceding sections exhibited the $u\to 0$ degeneration of the
Clifford completion as the algebraic dichotomy associated to the
critical value $c=26$. The BRST matter--ghost system adds the
Goddard--Thorn filtration; the Clifford degeneration is the anomaly
calculation on the $E_1$ page, and the no-ghost theorem is the
collapse of that filtered complex.

\begin{theorem}[$c=26$ BRST anomaly cancellation and Clifford degeneration;
\ClaimStatusProvedHere]
\label{thm:c26-no-ghost-clifford-degeneration}
\index{no-ghost theorem!Clifford degeneration}
\index{critical string!chain-level no-ghost}
Let $\cA_{\mathrm{str}}$ be the bosonic string vertex algebra at
$c = 26$, assembled as the matter--ghost pair
$V_{\mathrm{matter}}^{c=26} \otimes V_{bc}^{c=-26}$ with total
central charge zero. Let $Q_{\mathrm{BRST}}$ be the BRST differential
and $\cH_{\mathrm{phys}}(\cA_{\mathrm{str}}) :=
H^\bullet(Q_{\mathrm{BRST}})$ the physical state space. Assume the
standard Goddard--Thorn hypotheses: the matter module has intercept one,
positive-energy grading, and nondegenerate transverse oscillator form,
and the Goddard--Thorn filtration has the Frenkel--Garland--Zuckerman
collapse
\[
E_1 \cong \cF_{\perp}\otimes H^\bullet(\Bbbk b_0\oplus \Bbbk c_0,
Q_0),
\qquad
E_1 = E_\infty .
\]
Here $Q_0$ is the zero-mode component of $Q_{\mathrm{BRST}}$ on the
$b_0,c_0$ sector.
Then:
\begin{enumerate}[label=\textup{(\roman*)}]
\item At $c^{\mathrm{matter}}=26$ (the critical string value), the
Clifford completion
(Section~\ref{app:anomaly-completed-topological-holography})
of the matter--ghost BRST pair $\cA_{\mathrm{str}}$ degenerates:
the effective handle anomaly
$\kappa_{\mathrm{eff}}^{\mathrm{BRST}} =
\kappa(\text{matter},c=26)+\kappa(\text{ghost},c=-26)$ is zero, so
$u_{\mathrm{eff}}=\kappa_{\mathrm{eff}}^{\mathrm{BRST}}\omega_g=0$ and
the handle Clifford algebra becomes
$\bigwedge^*(\Bbbk^{2g})$.
\item The $b_0,c_0$ zero modes form the usual odd Heisenberg pair
$\{b_0,c_0\}=1$, $b_0^2=c_0^2=0$. This pair is the contracting
homotopy on the longitudinal--ghost factor of the filtered BRST
complex, while the exterior handle algebra records the vanished
secondary anomaly.
\item The physical state space is canonically identified with the
transverse oscillator cohomology:
\[
\cH_{\mathrm{phys}}(\cA_{\mathrm{str}})
\;\cong\;
\cF_{\perp}.
\]
Longitudinal and nonzero ghost oscillators lie in
$Q_{\mathrm{BRST}}$-exact pairs.
\end{enumerate}
\end{theorem}

\begin{proof}
Effective anomaly vanishing at $c=26$: matter contributes
$\kappa(\text{matter}) = c^{\mathrm{matter}}/2 = 13$ (or
$\kappa = c$ in the convention of
\S~\ref{subsec:tholog-coda}); the $bc$-ghost system contributes
$\kappa(\text{ghost}) = -13$. Total $\kappa_{\mathrm{eff}}=0$,
so the secondary anomaly $u_{\mathrm{eff}}$ vanishes. Substituting
$u=0$ in the handle Clifford relations gives the exterior algebra
$\bigwedge^*(\Bbbk^{2g})$, by the gravitational-regime theorem above.

The $bc$ zero modes are not the handle Clifford generators; they form
the standard odd Heisenberg pair $\{b_0,c_0\}=1$ inside the BRST
complex. The operator $b_0$ is the contracting homotopy for the
longitudinal--ghost summand in the Goddard--Thorn filtered complex.
The assumed Frenkel--Garland--Zuckerman collapse identifies the
abutment with the transverse oscillator factor $\cF_\perp$, proving
the displayed cohomology isomorphism.
\end{proof}

\begin{remark}[From Clifford relation to Goddard--Thorn]
\label{rem:clifford-to-no-ghost}
The standard proof of no-ghost (Goddard--Thorn~\cite{GoddardThornNoGhost},
Frenkel--Garland--Zuckerman~\cite{FrenkelGarlandZuckerman}) is a
delicate spectral sequence argument. The Clifford-degeneration
statement isolates the anomaly calculation inside that proof: at
$c=26$, the matter and ghost modular characteristics cancel, the
handle Clifford algebra degenerates to an exterior algebra, and the
Goddard--Thorn filtration collapses onto transverse oscillators.
\end{remark}

\begin{remark}[Celestial consequence]
\label{rem:c26-celestial}
Via the universal celestial holography theorem
(Theorem~\ref{thm:universal-celestial-holography} in the
celestial-holography-core chapter), the $c=26$
no-ghost theorem is the celestial statement that the physical
soft-graviton sector of the bosonic string is the transverse
$\mathfrak{w}_{1+\infty}$ subalgebra, with longitudinal and ghost
soft modes $Q_{\mathrm{BRST}}$-exact. The Clifford degeneration
is precisely the shadow-class reduction $\mathbf C \rightsquigarrow
\mathbf G$ at the critical locus.
\end{remark}

\section{K3 class-\texorpdfstring{$\mathcal S$}{S} boundary constant}
\label{sec:csd-supplement}

\begin{proposition}[Class-$\mathcal S$ $A_1$ anomaly at twenty-four
punctures; \ClaimStatusProvedHere]
\label{rem:csd-1}
Let $\Sigma_{0,24}=(\mathbb P^1,D)$ with $|D|=24$, as for the
discriminant divisor of a semistable elliptic K3 surface with
twenty-four Kodaira $I_1$ fibres. The class-$\mathcal S$ theory
$\mathcal T[A_1,\Sigma_{0,24}]$ with maximal regular punctures has
\[
n_v = 21\cdot 3 = 63,\qquad
n_h = 22\cdot 4 = 88.
\]
Its four-dimensional and Schur-VOA central charges are
\[
c_{4d}
\;=\;
\frac{2n_v+n_h}{12}
\;=\;
\frac{126+88}{12}
\;=\;
\frac{107}{6},
\qquad
c_{2d}
\;=\;
-12c_{4d}
\;=\;
-214.
\]
\end{proposition}

\begin{proof}
A genus-zero surface with $n$ punctures decomposes into $n-2$
trinions and $n-3$ tubes. For $A_1$, each trinion $T_2$ contributes
$(n_v,n_h)=(0,4)$ and each tube contributes three vector multiplets,
so at $n=24$ one obtains $(n_v,n_h)=(63,88)$. The
Shapere--Tachikawa anomaly formula gives
$c_{4d}=(2n_v+n_h)/12=107/6$, equivalently
$c_{4d}=(5n-13)/6$ at genus zero. The Beem--Lemos--Liendo--Peelaers--
Rastelli--van~Rees Schur functor gives $c_{2d}=-12c_{4d}$.
For $n=4$ the same formula gives $c_{4d}=7/6$, the
$\mathrm{SU}(2)$ $N_f=4$ sanity check. The number $-312$ and the
factorisation $-12\cdot 26$ do not occur in this class-$\mathcal S$
anomaly calculation. The divisor $D$ varies in
\(\operatorname{Conf}_{24}(\mathbb P^1)/\operatorname{PGL}_2\); an
$M_{24}$-equivariant fibre is additional level structure on special
fibres, while the central-charge computation is the family statement.
\end{proof}

% Section file for Chapter: Twisted Holography and Quantum Gravity
% Result (G10): Fredholm Partition Functions
%
% This section develops the nonperturbative partition function
% as a Fredholm determinant on the one-particle Bergman space,
% extending from the Heisenberg prototype through the full
% standard landscape.

% ==========================================
% LOCAL MACRO PROVISIONS
% ==========================================
\providecommand{\Bergman}{\mathcal{B}}
\providecommand{\Fock}{\mathcal{F}}
\providecommand{\Hilb}{\mathcal{H}\!\mathit{ilb}}
\providecommand{\sewop}{T}
\providecommand{\sewker}{K}
\providecommand{\trcl}{\mathcal{L}^1}
\providecommand{\hscl}{\mathcal{L}^2}
\providecommand{\freddet}{\mathrm{det}_{\mathrm{Fred}}}
\providecommand{\Wick}{\mathcal{W}}
\providecommand{\Bergker}{\mathsf{B}}
\providecommand{\primeform}{E}
\providecommand{\collar}{\mathfrak{c}}
\providecommand{\sewenv}{\cA^{\mathrm{sew}}}
\providecommand{\anbar}{\barB^{\mathrm{an}}}
\providecommand{\algcore}{\cA_{\mathrm{alg}}}
\providecommand{\codercat}{D^{\mathrm{co}}}
\providecommand{\contrcat}{D^{\mathrm{ctr}}}
\providecommand{\graphnorm}{\|\cdot\|_{\mathrm{graph}}}
\providecommand{\secquant}{\Gamma}
\providecommand{\hodgebdle}{\mathbb{E}}
\providecommand{\hodgedet}{\det\hodgebdle}
\providecommand{\zetareg}{\det\nolimits'_\zeta}
\providecommand{\dedeta}{\eta}
\providecommand{\Sp}{\mathrm{Sp}}
\providecommand{\Siegel}{\mathfrak{H}}
\providecommand{\OPEcoeff}{C}

\section{Fredholm partition functions}
% label removed: sec:thqg-fredholm-partition-functions
\index{Fredholm determinant!partition function|textbf}
\index{partition function!Fredholm}

The preceding nine results of this chapter extract holographic,
gravitational, and bootstrap consequences from the universal
Maurer--Cartan element
$\Theta_\cA \in \MC(\gAmod)$
at the level of formal algebraic identities.  This final result
crosses from the algebraic to the analytic: the partition functions
are convergent Fredholm determinants on explicit Hilbert spaces.

The key bridge is the \emph{sewing envelope}
$\sewenv$
(Definition~\ref{def:sewing-envelope}): the Hausdorff completion
of the algebraic chiral core $\algcore$ for the locally convex
topology generated by all amplitude seminorms.  For the Heisenberg
algebra $\cH_\kappa$, the sewing envelope is the symmetric algebra
of the Bergman space $A^2(D)$ of the disk
(Theorem~\ref{thm:heisenberg-sewing}\,(i)), and the genus-$g$
partition function is the Fredholm determinant of an explicit
trace-class sewing operator on $A^2(D)$.  The HS-sewing criterion
(Theorem~\ref{thm:general-hs-sewing}) extends the convergence to
the entire standard landscape of universal algebras (convergence,
not Fredholm
form).  On the proved Gaussian scalar lane, the scalar package is
a $\kappa$-multiple of the Heisenberg one, so the Fredholm
determinant structure propagates there.  For classes L, C, and M,
the partition function has additional non-Fredholm corrections
from the higher-arity shadow obstruction tower
(\S\ref{subsec:thqg-X-non-fredholm}).

The section is organized as follows.
\begin{enumerate}[label=\textup{(\arabic*)}]
\item Sewing envelopes and Bergman spaces:
  the analytic framework for one-particle reduction
  (\S\ref{subsec:thqg-X-sewing-bergman}).
\item The Heisenberg sewing theorem in full:
  four clauses with complete proofs
  (\S\ref{subsec:thqg-X-heisenberg-sewing}).
\item One-particle reduction, trace-class estimates,
  and the Dedekind eta derivation
  (\S\ref{subsec:thqg-X-one-particle}).
\item Extension to class-G algebras via scalar saturation
  (\S\ref{subsec:thqg-X-class-G}).
\item Non-Fredholm corrections for classes L, C, M
  (\S\ref{subsec:thqg-X-non-fredholm}).
\item Analytic aspects: the analytic bar coalgebra,
  coderived categories, and open problems
  (\S\ref{subsec:thqg-X-analytic-aspects}).
\end{enumerate}

% =====================================================================
%
%  SECTION 1: SEWING ENVELOPES AND BERGMAN SPACES
%
% =====================================================================

\subsection{Sewing envelopes and Bergman spaces}
% label removed: subsec:thqg-X-sewing-bergman

\subsubsection{The Bergman space of a disk}

\begin{definition}[Bergman space of a disk]%
% label removed: def:thqg-X-bergman-space%
\index{Bergman space|textbf}%
For $D = \{z \in \C : |z| < 1\}$ the \textbf{unit disk},
the \emph{Bergman space} $A^2(D)$ is the Hilbert space of
holomorphic functions that are square-integrable with respect
to the area measure:
\begin{equation}% label removed: eq:thqg-X-bergman-norm
A^2(D)
\;:=\;
\Bigl\{f \in \cO(D) : \|f\|^2_{A^2}
= \int_D |f(z)|^2\,dA(z) < \infty\Bigr\},
\end{equation}
where $dA(z) = \frac{i}{2}\,dz \wedge d\bar{z}
= r\,dr\,d\theta$ in polar coordinates.  The monomials
$e_n(z) = \sqrt{(n+1)/\pi}\, z^n$ ($n \geq 0$)
form an orthonormal basis:
$\langle e_n, e_m\rangle = \delta_{nm}$ (verified by
$\int_D |z|^{2n}\,dA = \pi/(n+1)$).  The
\emph{reproducing kernel} (Bergman kernel for the unit
disk) is
\begin{equation}% label removed: eq:thqg-X-bergman-kernel
\Bergker_D(z,w)
\;=\;
\sum_{n=0}^{\infty} e_n(z)\overline{e_n(w)}
\;=\;
\frac{1}{\pi}\sum_{n=0}^{\infty}(n+1)(z\bar{w})^n
\;=\;
\frac{1}{\pi(1 - z\bar{w})^2}.
\end{equation}
The last identity uses $\sum_{n=0}^{\infty}(n+1)x^n
= 1/(1-x)^2$ for $|x| < 1$.  \emph{This formula is
specific to the unit disk}; for a general bounded domain
$\Omega \subset \C$, the Bergman kernel depends on the
conformal geometry of~$\Omega$.
\end{definition}

\begin{remark}[Conformal weight and sections of $K^{-1/2}$]%
% label removed: rem:thqg-X-bergman-conformal%
The prime form $\primeform(z_1,z_2)$ on a Riemann surface is a section
of $K^{-1/2} \boxtimes K^{-1/2}$, \emph{not} $K^{+1/2}$.
The Bergman kernel
$\Bergker_D(z,w)$ is a section of $K \boxtimes \bar{K}$, consistent
with its conformal weight $(1,1)$.  The mode--Bergman correspondence
(Remark~\ref{rem:heisenberg-mode-bergman}) maps Fock modes to
Bergman monomials at weight $0$, absorbing the conformal weight
into the $\sqrt{n+1}$ prefactor.
\end{remark}

\begin{definition}[Bergman space with collar parameter]%
% label removed: def:thqg-X-bergman-collar%
\index{Bergman space!collar parameter}%
For $0 < q < 1$, define the \emph{weighted Bergman space}
\begin{equation}% label removed: eq:thqg-X-bergman-weighted
A^2_q(D)
\;:=\;
\Bigl\{f \in \cO(D) :
\|f\|^2_{q} := \sum_{n=0}^{\infty}
q^n \,|a_n|^2 \,(n+1)^{-1} < \infty
\;\text{where}\; f = \sum_{n=0}^{\infty} a_n z^n\Bigr\}.
\end{equation}
The collar parameter $q = e^{-t}$ ($t > 0$ the collar length)
encodes the sewing parameter for pair-of-pants decompositions.
The inclusion $A^2_{q'}(D) \hookrightarrow A^2_q(D)$ for $q' < q$
is trace class.
\end{definition}

\begin{lemma}[Bergman projection; \ClaimStatusProvedElsewhere]%
% label removed: lem:thqg-X-bergman-projection%
\index{Bergman projection}%
The orthogonal projection
$P \colon L^2(D, dA) \to A^2(D)$ is the integral operator
\[
(Pf)(z)
\;=\;
\int_D \Bergker_D(z,w)\,f(w)\,dA(w)
\;=\;
\frac{1}{\pi}\int_D \frac{f(w)}{(1-z\bar{w})^2}\,dA(w).
\]
The projection satisfies $P^2 = P$, $P^* = P$, and
$\ker P = (\bar{\partial}\text{-exact})$.
\end{lemma}

\subsubsection{The sewing envelope}

\begin{definition}[Sewing envelope; review]%
% label removed: def:thqg-X-sewing-envelope-review%
\index{sewing envelope!review}%
Recall from Definition~\ref{def:sewing-envelope} that
the \emph{sewing envelope} $\sewenv$ of an algebraic chiral
core $\algcore$ is the Hausdorff completion
for the locally convex topology generated by all amplitude
seminorms $p_{\Sigma,\xi,\eta}(a)
= |\langle \eta, A_\Sigma^{\mathrm{alg}}(\xi_1 \otimes \cdots
\otimes a \otimes \cdots \otimes \xi_m)\rangle|$,
where the supremum ranges over all finite conformally flat
bordisms $\Sigma$ with parametrized collars and all
test states $\xi, \eta$.  By Proposition~\ref{prop:sewing-universal-property},
$\sewenv$ has the universal property: every complete locally convex
realization factors uniquely through~$\sewenv$.
\end{definition}

\begin{lemma}[Amplitude seminorms form a separating family;
\ClaimStatusProvedHere]%
% label removed: lem:thqg-X-amplitude-separating%
\index{sewing envelope!separating family}%
The family of amplitude seminorms
$\{p_{\Sigma,\xi,\eta}\}$ is separating: if
$p_{\Sigma,\xi,\eta}(a) = 0$ for all bordisms $\Sigma$ with
parametrized collars and all test states $\xi, \eta$, then
$a = 0$ in~$\algcore$.
\end{lemma}

\begin{proof}
Suppose $a \in \algcore$ satisfies
$p_{\Sigma,\xi,\eta}(a) = 0$ for all $(\Sigma,\xi,\eta)$.
Then $a$ pairs trivially with all sewing amplitudes:
\[
\langle \eta,\, A_\Sigma^{\mathrm{alg}}(\xi_1 \otimes
\cdots \otimes a \otimes \cdots \otimes \xi_m)\rangle = 0
\]
for every conformally flat bordism $\Sigma$ and all test
states.  In particular, choosing $\Sigma$ to be the
pair-of-pants with collar parameter $q \to 0$
(the degeneration to the disk), the amplitude
$A_\Sigma^{\mathrm{alg}}$ reduces to the OPE
multiplication $Y(a,z)$: the condition becomes
$\langle \eta, Y(a,z)\,\xi\rangle = 0$ for all
$\xi, \eta \in \algcore$ and all $z$.  By nondegeneracy of
the OPE pairing
(Proposition~\textup{\ref{prop:sewing-universal-property}}:
the bilinear form $\langle\cdot,Y(\cdot,z)\cdot\rangle$
is nondegenerate on the algebraic core),
this forces $a = 0$.
\end{proof}

\begin{proposition}[Sewing envelope: conformal weight filtration;
\ClaimStatusProvedHere]%
% label removed: prop:thqg-X-sewing-filtration%
\index{sewing envelope!filtration}%
Let $\cA$ be a positive-energy chiral algebra with
conformal weight decomposition
$\algcore = \bigoplus_{n \geq 0} \algcore_n$, where
$\dim \algcore_n < \infty$ for all~$n$.  Then:
\begin{enumerate}[label=\textup{(\roman*)}]
\item The sewing envelope inherits a filtration
  $F^{\leq N}\sewenv
  = \overline{\bigoplus_{n \leq N} \algcore_n}$
  whose associated graded recovers the algebraic core:
  $\mathrm{gr}\,\sewenv \cong \algcore$.
\item For each $0 < q < 1$, the $q$-weighted completion
  $\sewenv_q := \widehat{\bigoplus}_{n \geq 0} q^{n/2}\,\algcore_n$
  is a Banach algebra for the operator norm of
  $A^2_q$-bounded amplitudes.
\item The sewing envelope is the projective limit
  $\sewenv = \varprojlim_{q \to 1^-} \sewenv_q$.
\end{enumerate}
\end{proposition}

\begin{proof}
Part~(i): The amplitude seminorms $p_{\Sigma,\xi,\eta}$ are
continuous with respect to the conformal weight grading because
each amplitude map preserves total conformal weight up to a shift
determined by the Euler characteristic of~$\Sigma$.
The conformal weight operator $L_0$ is bounded on each
finite-dimensional weight space $\algcore_n$ (where it acts as
scalar multiplication by~$n$), and the sewing seminorms are
dominated by $q^{L_0}$: for any bordism $\Sigma$ with collar
parameter~$q$,
\[
p_{\Sigma,\xi,\eta}(a)
\;\leq\;
C_{\Sigma,\xi,\eta}\, q^n
\quad\text{for } a \in \algcore_n.
\]
Since $q^{L_0}$ is continuous in the weight filtration
topology (it sends $F^{\leq N}\algcore$ to itself with
norm $\leq q^N$), the weight map $a \mapsto n$ is continuous
under the sewing topology.
Passage to the associated graded undoes the completion, recovering
the algebraic core.

Part~(ii): The $q$-weighting assigns norm $q^{n/2}$ to elements
of weight~$n$.  The pair-of-pants amplitude
$m \colon \algcore_a \otimes \algcore_b \to \bigoplus_c \algcore_c$
satisfies $|m^c_{a,b}| \leq K(a+b+c+1)^N$ by polynomial OPE growth.
The $q$-weighted operator norm is
$\|m\|_q \leq K \sum_c q^{c/2}(a+b+c+1)^N < \infty$
for fixed $a,b$ and $q < 1$, so $\sewenv_q$ is submultiplicative.
Completeness follows from the completeness of the
$q$-weighted $\ell^2$ space.

Part~(iii): The sewing topology is generated by
$\{p_{\Sigma,\xi,\eta}\}$, which are bounded by $q$-weighted
norms for any $q$ sufficiently close to~$1$ (depending on the
collar length of~$\Sigma$).  Hence $\sewenv$ is the intersection
$\bigcap_{q < 1} \sewenv_q = \varprojlim_{q} \sewenv_q$.
\end{proof}

\subsubsection{The Heisenberg sewing envelope}

\begin{proposition}[Heisenberg sewing envelope = $\operatorname{Sym} A^2(D)$;
\ClaimStatusProvedElsewhere{} Moriwaki~\cite{Moriwaki26b}]%
% label removed: prop:thqg-X-heisenberg-sewing-envelope%
\index{Heisenberg!sewing envelope}%
\index{sewing envelope!Heisenberg}%
The sewing envelope of the algebraic Heisenberg VOA
$\cH_\kappa$ with level~$\kappa$ is
\begin{equation}% label removed: eq:thqg-X-heisenberg-sew
\cH_\kappa^{\mathrm{sew}}
\;\cong\;
\operatorname{Sym} A^2(D),
\end{equation}
the symmetric (bosonic) Fock space over the Bergman space of
the disk.  The isomorphism is given by the mode--Bergman
correspondence
\begin{equation}% label removed: eq:thqg-X-mode-bergman-map
\Theta(a_{-n-1}\mathbf{1})
\;=\;
\sqrt{(n+1)\kappa}\; z^n
\;\in\; A^2(D),
\end{equation}
extended multiplicatively by symmetrization:
$\Theta(a_{-n_1-1} \cdots a_{-n_k-1}\mathbf{1})
= \operatorname{Sym}(\Theta(a_{-n_1-1}\mathbf{1})
\otimes \cdots \otimes
\Theta(a_{-n_k-1}\mathbf{1}))$.
\end{proposition}

\begin{proof}
This is the ind-Hilbert identification of
Moriwaki~\cite{Moriwaki26b}, which constructs a conformally
flat 2-disk algebra structure on $\operatorname{Sym} A^2(D)$
and verifies that it is isomorphic to the Hilbert completion
of the Heisenberg Fock space for the sewing-amplitude topology.
The $\sqrt{(n+1)\kappa}$ normalization ensures that
$\|\Theta(a_{-n-1}\mathbf{1})\|^2 = \kappa$, matching the
level of the two-point function
$\langle a(z)\,a(w)\rangle = \kappa/(z-w)$.
\end{proof}

\begin{remark}[Tensor product structure]%
% label removed: rem:thqg-X-tensor-sewing%
The Fock space $\operatorname{Sym} A^2(D)$ has the standard
$n$-particle decomposition
\[
\operatorname{Sym} A^2(D)
\;=\;
\bigoplus_{n=0}^{\infty}
\operatorname{Sym}^n A^2(D),
\]
where $\operatorname{Sym}^0 = \C$ (vacuum sector) and
$\operatorname{Sym}^n$ is the $n$-th symmetric power.
The dimension of the weight-$N$ subspace in
$\operatorname{Sym} A^2(D)$ is the number of partitions
$p(N)$, matching the Heisenberg character
$\chi_{\cH}(q) = \prod_{n=1}^{\infty}(1-q^n)^{-1}
= q^{-1/24}/\dedeta(q)$.
\end{remark}

\subsubsection{Sewing operators on Bergman spaces}

\begin{definition}[Sewing operator]%
% label removed: def:thqg-X-sewing-operator%
\index{sewing operator|textbf}%
\index{Bergman space!sewing operator}%
Let $\Sigma$ be a compact Riemann surface obtained by sewing
two disks $D_1, D_2$ along their boundary circles via the
identification $z_1 z_2 = q$ for $0 < q < 1$.  The
\emph{sewing operator} $\sewop_q \colon A^2(D) \to A^2(D)$ is
defined on the orthonormal basis by
\begin{equation}% label removed: eq:thqg-X-sewing-eigenvalues
\sewop_q\, e_n \;=\; q^{n+1}\, e_n,
\qquad n \geq 0.
\end{equation}
Equivalently, $\sewop_q$ acts by multiplication by
$|z|^2 \cdot q$ on the Bergman kernel:
\begin{equation}% label removed: eq:thqg-X-sewing-kernel
(\sewop_q f)(z)
\;=\;
\int_D \sewop_q(z,w)\,f(w)\,dA(w),
\qquad
\sewop_q(z,w)
\;=\;
\sum_{n=0}^{\infty}
q^{n+1}\,e_n(z)\overline{e_n(w)}.
\end{equation}
\end{definition}

\begin{lemma}[Schatten class properties of $\sewop_q$;
\ClaimStatusProvedHere]%
% label removed: lem:thqg-X-sewing-schatten%
\index{sewing operator!Schatten class}%
For $0 < q < 1$, the sewing operator $\sewop_q$ on $A^2(D)$
satisfies:
\begin{enumerate}[label=\textup{(\roman*)}]
\item \emph{Trace class.}
  $\sewop_q \in \trcl(A^2(D))$ with
  \begin{equation}% label removed: eq:thqg-X-trace-norm
  \|\sewop_q\|_1
  \;=\;
  \sum_{n=0}^{\infty} q^{n+1}
  \;=\;
  \frac{q}{1-q}.
  \end{equation}
\item \emph{Hilbert--Schmidt.}
  $\sewop_q \in \hscl(A^2(D))$ with
  \begin{equation}% label removed: eq:thqg-X-hs-norm
  \|\sewop_q\|_{\hscl}^2
  \;=\;
  \sum_{n=0}^{\infty} q^{2(n+1)}
  \;=\;
  \frac{q^2}{1-q^2}.
  \end{equation}
\item \emph{$p$-Schatten.}
  For all $p \geq 1$,
  $\sewop_q \in \mathcal{L}^p(A^2(D))$ with
  $\|\sewop_q\|_p^p = q^p/(1-q^p)$.
\item \emph{Spectral gap.}
  $\|\sewop_q\| = q$ (operator norm), and the spectral
  radius satisfies $r(\sewop_q) = q < 1$.
\end{enumerate}
\end{lemma}

\begin{proof}
The eigenvalues of $\sewop_q$ are $\lambda_n = q^{n+1}$ for
$n \geq 0$.  Since $0 < q < 1$, the eigenvalues are positive
and summable:
\[
\|\sewop_q\|_p^p
= \sum_{n=0}^{\infty} |\lambda_n|^p
= \sum_{n=0}^{\infty} q^{p(n+1)}
= \frac{q^p}{1-q^p}.
\]
For $p = 1$ (trace norm): $q/(1-q)$.
For $p = 2$ (HS norm squared): $q^2/(1-q^2)$.
The operator norm is $\|\sewop_q\| = \sup_n q^{n+1} = q$.
\end{proof}

\begin{definition}[Genus-$g$ sewing kernel]%
% label removed: def:thqg-X-genus-g-sewing%
\index{sewing operator!genus $g$}%
A genus-$g$ surface $\Sigma_g$ with $n$ punctures admits a
pair-of-pants decomposition into $2g - 2 + n$ pairs of pants,
with $3g - 3 + n$ sewing circles.  Each sewing circle
$C_\alpha$ ($\alpha = 1, \ldots, 3g-3+n$) contributes a sewing
parameter $q_\alpha = e^{-t_\alpha + i\theta_\alpha}$ and
a sewing operator $\sewop_{q_\alpha}$ on $A^2(D)$.
The \emph{genus-$g$ sewing kernel} is the composite operator
\begin{equation}% label removed: eq:thqg-X-genus-g-kernel
\sewker_g
\;=\;
\prod_{\alpha=1}^{3g-3}
\sewop_{q_\alpha}
\;\in\;
\trcl(A^2(D)^{\otimes g}),
\end{equation}
where the product is over the $3g-3$ internal sewing circles
of a genus-$g$ surface without punctures.
More precisely, $\sewker_g$ acts on the tensor product of
$g$ copies of $A^2(D)$ (one for each handle),
with each $\sewop_{q_\alpha}$ acting on the appropriate
tensor factor.
\end{definition}

\begin{remark}[Dependence on pants decomposition]%
% label removed: rem:thqg-X-pants-independence%
\index{pants decomposition!independence}%
The Fredholm determinant $\det(1 - \sewker_g)$ is independent
of the choice of pair-of-pants decomposition, although the
intermediate operator $\sewker_g$ depends on the decomposition.
This independence follows from the sewing axiom of the
chain-level modular-functor package
(Theorem~\ref*{V1-thm:chain-modular-functor}\,(iv)):
different decompositions are related by chain homotopies
that preserve the trace.
\end{remark}

\subsubsection{Composition operators from the prime form}

\begin{definition}[Bergman composition operator]%
% label removed: def:thqg-X-bergman-composition%
\index{composition operator!Bergman}%
Let $\Sigma_g$ be a genus-$g$ Riemann surface with period
matrix $\Omega \in \Siegel_g$ (the Siegel upper half-space),
and let $D_j \subset \Sigma_g$ be small coordinate disks around
marked points $p_j$.  The \emph{Bergman composition operator}
$R \colon A^2(D_{\mathrm{out}}) \to
A^2(D_{\mathrm{in},1}) \otimes A^2(D_{\mathrm{in},2})$
for a pair-of-pants $P$ with one outgoing and two incoming
boundaries is
\begin{equation}% label removed: eq:thqg-X-bergman-composition
R_{n,m}
\;=\;
\frac{1}{(2\pi i)^2}
\oint_{|z|=1} \oint_{|w|=1}
z^{-n-1}\,w^{-m-1}\,
G_P(z,w)\,dz\,dw,
\end{equation}
where $G_P(z,w) = \partial_z \log \primeform(z,w;\Omega)$ is the
chiral Green function on the pair of pants.
\end{definition}

\begin{lemma}[Exponential decay of composition coefficients;
\ClaimStatusProvedHere]%
% label removed: lem:thqg-X-composition-decay%
\index{composition operator!exponential decay}%
The matrix coefficients of the Bergman composition operator satisfy
\begin{equation}% label removed: eq:thqg-X-composition-bound
|R_{n,m}|
\;\leq\;
C\,q^{(n+m)/2}
\end{equation}
for a constant $C$ depending only on the collar length
$t = -\log q$ of the pair of pants.  In particular,
$R$ is Hilbert--Schmidt with
$\|R\|_{\hscl}^2 \leq C^2/(1-q)^2$.
\end{lemma}

\begin{proof}
The Green function $G_P(z,w)$ on the pair of pants is holomorphic
for $|z| < 1$, $|w| < 1$ away from the diagonal $z = w$.
By the collar theorem for hyperbolic surfaces, there exists a
collar of width $t/2$ around each boundary circle, so $G_P$ extends
holomorphically to $|z| < e^{t/2}$, $|w| < e^{t/2}$.

\emph{Contour-moving estimate.}
The matrix coefficient~\eqref{eq:thqg-X-bergman-composition}
is a double contour integral on $|z| = 1$, $|w| = 1$.
Move both contours inward to $|z| = q^{1/2}$, $|w| = q^{1/2}$
(with $q = e^{-t}$, so $q^{1/2} = e^{-t/2} < 1$).
By Cauchy's theorem the integrals are unchanged (the
integrand is holomorphic in the annular region).  On the new
contours:
\begin{itemize}
\item $|z^{-n-1}| = q^{-(n+1)/2}$ and
  $|w^{-m-1}| = q^{-(m+1)/2}$;
\item $|G_P(z,w)| \leq C'$ where
  $C' := \sup_{|z|=q^{1/2},\,|w|=q^{1/2}} |G_P(z,w)|
  < \infty$ (holomorphicity in the collar);
\item each arc has length $2\pi q^{1/2}$.
\end{itemize}
Therefore:
\begin{align*}
|R_{n,m}|
&\leq \frac{1}{(2\pi)^2}
  \cdot (2\pi q^{1/2})^2
  \cdot q^{-(n+1)/2}\,q^{-(m+1)/2}
  \cdot C' \\
&= C'\,q^{1}\,q^{-(n+1)/2}\,q^{-(m+1)/2} \\
&= C'\,q^{(n+m)/2}\,
  \underbrace{q^{1 - (n+1)/2 - (m+1)/2 + (n+m)/2}}_{= q^0 = 1}
  \\
&= C'\,q^{(n+m)/2}.
\end{align*}
(The exponent simplifies:
$1 - \frac{n+1}{2} - \frac{m+1}{2} + \frac{n+m}{2}
= 1 - 1 = 0$.)
Setting $C = C'$ gives the claimed bound.

For the HS norm:
\[
\|R\|_{\hscl}^2
= \sum_{n,m \geq 0} |R_{n,m}|^2
\leq C^2 \sum_{n,m \geq 0} q^{n+m}
= \frac{C^2}{(1-q)^2}.
\qedhere
\]
\end{proof}

\begin{proposition}[Second quantization;
\ClaimStatusProvedHere]%
% label removed: prop:thqg-X-second-quantization%
\index{second quantization!Bergman}%
Let $T \colon A^2(D) \to A^2(D)$ be a trace-class operator
with $\|T\| < 1$.  The \emph{second quantization}
$\secquant(T) \colon \operatorname{Sym} A^2(D) \to
\operatorname{Sym} A^2(D)$ is defined by
\begin{equation}% label removed: eq:thqg-X-second-quant
\secquant(T)|_{\operatorname{Sym}^n A^2(D)}
\;=\;
T^{\otimes n}\big|_{\operatorname{Sym}^n},
\end{equation}
with $\secquant(T)|_{\operatorname{Sym}^0} = \id$.
Then:
\begin{enumerate}[label=\textup{(\roman*)}]
\item $\secquant(T)$ is trace class on $\operatorname{Sym} A^2(D)$
  with
  \begin{equation}% label removed: eq:thqg-X-second-quant-trace
  \operatorname{Tr}_{\operatorname{Sym}}(\secquant(T))
  \;=\;
  \det(1 - T)^{-1}.
  \end{equation}
\item More generally, for the Fredholm determinant:
  \begin{equation}% label removed: eq:thqg-X-fredholm-second-quant
  \det_{\operatorname{Sym}}(1 - \secquant(T))
  \;=\;
  \exp\!\Bigl(-\sum_{k=1}^{\infty}
  \frac{1}{k}\operatorname{Tr}(T^k)\Bigr)
  \;=\;
  \det(1 - T).
  \end{equation}
\end{enumerate}
\end{proposition}

\begin{proof}
Part~(i): The trace on $\operatorname{Sym}^n A^2(D)$ is computed
using an orthonormal basis of symmetric tensors.
For eigenvalues $\lambda_j$ of~$T$:
\[
\operatorname{Tr}_{\operatorname{Sym}^n}(T^{\otimes n})
= \sum_{j_1 \leq \cdots \leq j_n}
\lambda_{j_1} \cdots \lambda_{j_n}
= h_n(\lambda_1, \lambda_2, \ldots),
\]
where $h_n$ is the complete symmetric polynomial.
Summing over $n \geq 0$:
\[
\sum_{n=0}^{\infty}
\operatorname{Tr}_{\operatorname{Sym}^n}(T^{\otimes n})
= \sum_{n=0}^{\infty} h_n(\lambda)
= \prod_{j} \frac{1}{1 - \lambda_j}
= \det(1 - T)^{-1}.
\]
The series converges absolutely because $T$ is trace class
and $\|T\| < 1$, so $\sum_j |\lambda_j| < \infty$ and
$|\lambda_j| < 1$ for all~$j$.

Part~(ii): $\log\det(1-T) = \sum_j \log(1-\lambda_j)
= -\sum_{k=1}^{\infty} \frac{1}{k}\sum_j \lambda_j^k
= -\sum_{k=1}^{\infty} \frac{1}{k}\operatorname{Tr}(T^k)$.
\end{proof}

% =====================================================================
%
%  SECTION 2: THE HEISENBERG SEWING THEOREM
%
% =====================================================================

\subsection{The Heisenberg sewing theorem: full development}
% label removed: subsec:thqg-X-heisenberg-sewing

\subsubsection{Statement and architecture}

The Heisenberg sewing theorem
(Theorem~\ref{thm:heisenberg-sewing})
has four clauses.  We now develop the complete proofs, starting
from the mode--Bergman correspondence and building to the
Fredholm determinant formula for genus-$g$ partition functions.

\begin{theorem}[Heisenberg sewing theorem: full development;
\ClaimStatusProvedHere]%
% label removed: thm:thqg-X-heisenberg-sewing-full%
\index{Heisenberg!sewing theorem!full development}%
Let $\cH_\kappa$ be the Heisenberg algebra with level $\kappa > 0$,
equipped with the mode--Bergman correspondence~$\Theta$
\textup{(}Remark~\textup{\ref{rem:heisenberg-mode-bergman}}\textup{)}.
Then:
\begin{enumerate}[label=\textup{(\Roman*)}]
\item \textbf{Sewing envelope identification.}
  The sewing envelope of $\cH_\kappa$ is
  $\cH_\kappa^{\mathrm{sew}} \cong
  \operatorname{Sym} A^2(D)$.
\item \textbf{Bar differential closure.}
  The algebraic bar differential extends by closure to the
  Gaussian collision operator on Bergman tensors:
  \begin{equation}% label removed: eq:thqg-X-bar-closure
  d_{\mathrm{bar}}^{\mathrm{an}}
  \;=\;
  \overline{d_{\mathrm{bar}}^{\mathrm{alg}}}
  ^{\,\graphnorm}
  \;=\;
  \sum_{i < j}
  \kappa\;\operatorname{Res}_{z_i = z_j}
  \frac{1}{(z_i - z_j)}\,
  \partial_{z_i z_j}.
  \end{equation}
\item \textbf{Pair-of-pants sewing = second quantization.}
  The pair-of-pants amplitude on
  $\operatorname{Sym} A^2(D)$ is the second quantization
  $\secquant(R)$ of the one-particle Bergman restriction
  operator
  $R \colon A^2(D_{\mathrm{out}}) \to
  A^2(D_{\mathrm{in},1}) \otimes A^2(D_{\mathrm{in},2})$.
\item \textbf{Fredholm determinant formula.}
  Every genus-$g$ closed amplitude is a Fredholm determinant:
  the genus-$g$ partition function of $\cH_\kappa$ is
  \begin{equation}% label removed: eq:thqg-X-fredholm-formula
  Z_g(\cH_\kappa;\,\Omega)
  \;=\;
  \bigl(\det\operatorname{Im}\Omega\bigr)^{-\kappa/2}
  \cdot
  \bigl(\zetareg\Delta_{\Sigma_g}\bigr)^{-\kappa/2},
  \end{equation}
  where $\Omega$ is the period matrix of $\Sigma_g$ and
  $\zetareg\Delta$ is the zeta-regularized determinant of the
  Laplacian.
\end{enumerate}
\end{theorem}

The proof occupies the remainder of this subsection.

\subsubsection{Clause I: sewing envelope identification}

\begin{proof}[Proof of Clause~I]
The mode--Bergman map $\Theta$ of
Remark~\ref{rem:heisenberg-mode-bergman} defines a dense
algebraic embedding
$\Theta \colon \Fock_\kappa^{\mathrm{fin}} \hookrightarrow
\operatorname{Sym} A^2(D)$
from the algebraic Fock space (finite linear combinations of
mode monomials) into the symmetric algebra of the Bergman space.

\emph{Step 1: algebraic amplitudes.}
On the algebraic Fock core, every amplitude
$A_\Sigma^{\mathrm{alg}}$ is a finite sum of Wick contractions.
By Wick's theorem for the free boson, the amplitude of
$a_{-n_1-1}\cdots a_{-n_k-1}\mathbf{1}$ on a surface~$\Sigma$
with collar coordinates is a polynomial in the mode
coefficients of the propagator $G_\Sigma$, which by the
mode--Bergman dictionary become matrix coefficients of the
Bergman composition operator $R$.

\emph{Step 2: continuity for the sewing topology.}
The sewing seminorms $p_{\Sigma,\xi,\eta}$ on $\Fock_\kappa^{\mathrm{fin}}$
are controlled by the Bergman norms: for
$v = \Theta(a_{-n_1-1}\cdots a_{-n_k-1}\mathbf{1})
\in \operatorname{Sym}^k A^2(D)$,
\begin{equation}% label removed: eq:thqg-X-sewing-seminorm-bound
p_{\Sigma,\xi,\eta}(v)
\;\leq\;
\|\xi\| \cdot \|\eta\| \cdot \|v\|_{\operatorname{Sym}^k}
\cdot \prod_\alpha \|\sewop_{q_\alpha}\|,
\end{equation}
where the product is over sewing circles of~$\Sigma$.
Since each $\sewop_{q_\alpha}$ is bounded with
$\|\sewop_{q_\alpha}\| = q_\alpha < 1$, the sewing topology is
weaker than the $\operatorname{Sym} A^2(D)$ topology.

\emph{Step 3: completeness.}
Conversely, any Cauchy sequence in the sewing topology is Cauchy
in $\operatorname{Sym} A^2(D)$ (because the latter norm
controls each sewing seminorm).  Hence the completions coincide:
$\cH_\kappa^{\mathrm{sew}} \cong \operatorname{Sym} A^2(D)$.
The full argument is the identification of Moriwaki~\cite{Moriwaki26b},
which constructs the conformally flat 2-disk algebra structure
on $\operatorname{Sym} A^2(D)$ and verifies the universal
property of the sewing envelope.
\end{proof}

\subsubsection{Clause II: bar differential closure}

\begin{proof}[Proof of Clause~II]
The algebraic bar differential on $\barB_n(\cH_\kappa)$ is
\[
d_{\mathrm{bar}}^{\mathrm{alg}}
= \sum_{i < j}
\operatorname{Res}_{z_i = z_j}
\circ\, Y_{ij},
\]
where $Y_{ij}$ is the chiral multiplication
$Y(a_{-n-1}\mathbf{1}, z_i - z_j)
= \sum_m a_m (z_i - z_j)^{-m-1}$.
For the Heisenberg algebra, the OPE is
$a(z_i)\,a(z_j) \sim \kappa/(z_i - z_j)$,
so $Y_{ij}$ acts on the one-particle Bergman space as
multiplication by $\kappa/(z_i - z_j)$ followed by the residue
at $z_i = z_j$.

\emph{Step 1: densely defined operator.}
On $\operatorname{Sym}^n A^2(D)$, the operator
$d_{\mathrm{bar}}^{\mathrm{alg}}$ is densely defined on the
algebraic Fock core.  Its action on monomials is
\begin{align}
d_{\mathrm{bar}}(z^{n_1} \cdots z^{n_k})
&= \kappa \sum_{i < j}
\operatorname{Res}_{z_i = z_j}
\frac{z_1^{n_1} \cdots z_k^{n_k}}{z_i - z_j}
\notag\\
&= \kappa \sum_{i < j}
\sum_{p=0}^{n_j-1}
z_j^{n_i + n_j - 1 - p}\, z_j^p
\cdot \prod_{\ell \neq i,j} z_\ell^{n_\ell},
% label removed: eq:thqg-X-bar-monomial-action
\end{align}
where the residue replaces $z_i$ by $z_j$ with the standard
contraction.

\emph{Step 2: closability.}
The operator $d_{\mathrm{bar}}^{\mathrm{alg}}$ is closable
on $\operatorname{Sym}^n A^2(D)$.  This follows from the
fact that the formal adjoint $(d_{\mathrm{bar}}^{\mathrm{alg}})^*$
is also densely defined (it is the creation operator that reverses
the contraction), so the closure exists by von Neumann's theorem.

\emph{Step 3: identification of the closure.}
The closure is the \emph{Gaussian collision operator}: the sum
of pairwise collision residue operators weighted by the level
$\kappa$.  Explicitly, it acts on
$f_1 \otimes \cdots \otimes f_n \in \operatorname{Sym}^n A^2(D)$
by
\[
\overline{d}_{\mathrm{bar}}(f_1 \otimes \cdots \otimes f_n)
= \kappa \sum_{i < j}
\operatorname{Res}_{z_i = z_j}
\frac{f_i(z_i)\,f_j(z_j)}{z_i - z_j}
\cdot \prod_{\ell \neq i,j} f_\ell(z_\ell),
\]
where the residue is computed by the Bergman projection
of $f_i(z)\,f_j(z) / (z_i - z_j)$ restricted to the diagonal.

To see that $\overline{d}_{\mathrm{bar}}$ is closed in
the Bergman norm topology, consider the $(i,j)$ summand:
the collision operator $C_{ij}(f_i \otimes f_j)(z)
= \operatorname{Res}_{w=z}[f_i(w)f_j(z)/(w-z)]$.
On the orthonormal basis $e_n(z) = \sqrt{(n+1)/\pi}\,z^n$:
\[
C_{ij}(e_a \otimes e_b)(z)
= \operatorname{Res}_{w=z}\!
\frac{\sqrt{(a+1)(b+1)}}{\pi}
\frac{w^a z^b}{w - z}
= \frac{\sqrt{(a+1)(b+1)}}{\pi}
\sum_{p=0}^{a-1} z^{a+b-1-p}\cdot z^p
\]
by geometric series expansion of $w^a/(w-z)$.  The
Bergman norm of the output satisfies
$\|C_{ij}(e_a \otimes e_b)\|^2
\leq a \cdot (a+1)(b+1)/\pi^2 \cdot (\text{sum of norms})$,
which is bounded by $a(a+1)(b+1)/\pi \leq (a+1)^2(b+1)$.
The graph norm
$\|x\|_{\mathrm{graph}}^2
:= \|x\|^2 + \|\overline{d}_{\mathrm{bar}} x\|^2$
is therefore controlled by the Bergman norm of the input
(with polynomial growth in the mode numbers), establishing
$\overline{d}_{\mathrm{bar}}$ as a closed operator.  By
the uniqueness of closure for closable operators (von Neumann),
$\overline{d}_{\mathrm{bar}}$ is the unique closed extension
of $d_{\mathrm{bar}}^{\mathrm{alg}}$.
\end{proof}

\subsubsection{Clause III: pair-of-pants = second quantization}

\begin{proof}[Proof of Clause~III]
Consider a pair of pants $P$ with one outgoing boundary $\partial_{\mathrm{out}}$
and two incoming boundaries $\partial_{\mathrm{in},1}$, $\partial_{\mathrm{in},2}$.
Each boundary is parametrized by a circle, and the sewing
parameters $(q_1, q_2)$ are the collar parameters.

\emph{Step 1: one-particle amplitude.}
On the one-particle space $A^2(D)$, the pair-of-pants amplitude
is the Bergman composition operator
$R \colon A^2(D_{\mathrm{out}}) \to
A^2(D_{\mathrm{in},1}) \otimes A^2(D_{\mathrm{in},2})$
of Definition~\ref{def:thqg-X-bergman-composition}.
Its matrix coefficients $R_{n,m}$ are the contour integrals
of the pair-of-pants Green function $G_P(z,w)$.

\emph{Step 2: Wick's theorem.}
On the full Fock space $\operatorname{Sym} A^2(D)$, the
Heisenberg amplitudes are determined by Wick's theorem:
every $n$-point function is a sum over pairings of two-point
functions.  In the Bergman picture, the $n$-particle amplitude
on $P$ is
\begin{equation}% label removed: eq:thqg-X-wick-bergman
A_P^{(n)}(f_1 \otimes \cdots \otimes f_n)
= \sum_{\sigma \in S_n}
\prod_{j=1}^{n}
\langle f_j, R\,g_{\sigma(j)}\rangle,
\end{equation}
where $g_{\sigma(j)}$ are the incoming modes.  This is the formula for the second quantization:
\[
A_P^{(n)}
= \secquant(R)\big|_{\operatorname{Sym}^n}.
\]

\emph{Step 3: total amplitude.}
Summing over particle number gives the full pair-of-pants
amplitude:
\begin{equation}% label removed: eq:thqg-X-pants-second-quant
A_P
= \bigoplus_{n=0}^{\infty}
A_P^{(n)}
= \secquant(R)
\colon \operatorname{Sym} A^2(D_{\mathrm{out}})
\to
\operatorname{Sym}(A^2(D_{\mathrm{in},1})
\oplus A^2(D_{\mathrm{in},2})).
\end{equation}
The convergence of this series follows from the exponential
decay $|R_{n,m}| \leq C q^{(n+m)/2}$
(Lemma~\ref{lem:thqg-X-composition-decay}).
\end{proof}

\subsubsection{Clause IV: Fredholm determinant}

\begin{proof}[Proof of Clause~IV]
A genus-$g$ surface $\Sigma_g$ without punctures is obtained by
sewing together $2g-2$ pairs of pants along $3g-3$ sewing circles.
Each sewing circle contributes a trace over the one-particle
Hilbert space, and by Wick's theorem the total amplitude factors
through the one-particle operators.

\emph{Step 1: genus-$g$ amplitude as iterated trace.}
The genus-$g$ partition function is
\begin{equation}% label removed: eq:thqg-X-genus-g-trace
Z_g(\cH_\kappa; \Omega)
= \operatorname{Tr}_{\operatorname{Sym} A^2}
\Bigl(\prod_{\alpha=1}^{3g-3}
\secquant(\sewop_{q_\alpha})^{\kappa}\Bigr),
\end{equation}
where the product is over all sewing circles, $\sewop_{q_\alpha}$
is the one-particle sewing operator, and the exponent $\kappa$
reflects the level (one Fock space per unit of central charge).

\emph{Step 2: trace factorization.}
By Proposition~\ref{prop:thqg-X-second-quantization}\,(i):
\[
\operatorname{Tr}_{\operatorname{Sym}}(\secquant(T))
= \det(1 - T)^{-1}
\]
for any trace-class $T$ with $\|T\| < 1$.  Applying this to each
sewing circle:
\begin{equation}% label removed: eq:thqg-X-genus-g-det
Z_g(\cH_\kappa; \Omega)
= \prod_{\alpha=1}^{3g-3}
\det\bigl(1 - \sewop_{q_\alpha}\bigr)^{-\kappa}.
\end{equation}

\emph{Step 3: from pants parameters to period matrix.}
The sewing parameters $\{q_\alpha\}$ are the Fenchel--Nielsen
coordinates of $\Sigma_g$; they determine and are determined by
the period matrix $\Omega \in \Siegel_g$.  The Polyakov--Alvarez
formula \cite{BK86} relates the product of one-particle
determinants to the period matrix and the zeta-regularized
determinant of the Laplacian:
\begin{equation}% label removed: eq:thqg-X-polyakov
\prod_{\alpha=1}^{3g-3}
\det(1 - \sewop_{q_\alpha})^{-1}
= (\det\operatorname{Im}\Omega)^{-1/2}
\cdot (\zetareg\Delta_{\Sigma_g})^{-1/2}.
\end{equation}
This is the holomorphic factorization of the string measure
of Belavin--Knizhnik~\cite{BK86}: the modular-invariant
expression factors into a product of a holomorphic and an
anti-holomorphic part, and the holomorphic part is
$\det(1 - \sewker_g)^{-1}$ up to the period-matrix prefactor.

\emph{Step 4: level $\kappa$.}
Raising to the power $\kappa$ (i.e., tensoring $\kappa$ copies
of the one-particle Bergman space, or equivalently working at
level $\kappa$):
\begin{equation}% label removed: eq:thqg-X-final-fredholm
Z_g(\cH_\kappa; \Omega)
= (\det\operatorname{Im}\Omega)^{-\kappa/2}
\cdot (\zetareg\Delta_{\Sigma_g})^{-\kappa/2}.
\end{equation}
This completes the proof of Clause~IV and of the full theorem.
\end{proof}

% =====================================================================
%
%  SECTION 3: ONE-PARTICLE REDUCTION
%
% =====================================================================

\subsection{One-particle reduction: trace-class operators
and the Dedekind eta}
% label removed: subsec:thqg-X-one-particle

The one-particle reduction principle reduces the computation
of partition functions on the bosonic Fock space to Fredholm
determinants on the one-particle Bergman space.  This subsection
carries out the explicit computations at genus $1$, $2$, and $3$.

\subsubsection{Genus 1: the torus and the Dedekind eta}

\begin{theorem}[Genus-1 Fredholm determinant;
\ClaimStatusProvedHere]%
% label removed: thm:thqg-X-genus-1-fredholm%
\index{Fredholm determinant!genus 1}%
\index{Dedekind eta function!from Fredholm determinant}%
The genus-$1$ partition function of $\cH_\kappa$ is
\begin{equation}% label removed: eq:thqg-X-genus-1-det
Z_1(\cH_\kappa;\,\tau)
\;=\;
(\operatorname{Im}\tau)^{-\kappa/2}\,
\bigl|\dedeta(\tau)\bigr|^{-2\kappa},
\end{equation}
where the holomorphic part is
\begin{equation}% label removed: eq:thqg-X-genus-1-hol
Z_1^{\mathrm{hol}}(\cH_\kappa;\,\tau)
\;=\;
\dedeta(\tau)^{-\kappa}
\;=\;
q^{-\kappa/24}\,
\prod_{n=1}^{\infty}
(1 - q^n)^{-\kappa}.
\end{equation}
\end{theorem}

\begin{proof}
A genus-$1$ surface $E_\tau = \C/(\Z + \tau\Z)$ is obtained
by sewing a single cylinder $\C^* \supset
\{q < |z| < 1\}$ with $q = e^{2\pi i \tau}$ ($\operatorname{Im}\tau > 0$).
The sewing circle is a single boundary component, and the
sewing operator is $\sewop_q$ on $A^2(D)$ with eigenvalues
$q^{n+1}$.

\emph{Step 1: Fock space trace.}
The partition function is the Fock-space trace:
\begin{align}
Z_1(\cH_1;\,\tau)
&= \operatorname{Tr}_{\operatorname{Sym} A^2(D)}
\bigl(\secquant(\sewop_q)\bigr)
\notag\\
&= \det(1 - \sewop_q)^{-1}
\quad\text{(Proposition~\ref{prop:thqg-X-second-quantization}\,(i))}
\notag\\
&= \prod_{n=0}^{\infty}
(1 - q^{n+1})^{-1}
\notag\\
&= \prod_{n=1}^{\infty}
(1 - q^n)^{-1}.
% label removed: eq:thqg-X-genus-1-product
\end{align}

\emph{Step 2: Dedekind eta identification.}
The Dedekind eta function is
$\dedeta(\tau) = q^{1/24}\prod_{n=1}^{\infty}(1-q^n)$
(Definition~\ref{def:eta-function}).  Therefore:
\begin{equation}% label removed: eq:thqg-X-eta-identification
\prod_{n=1}^{\infty}(1-q^n)^{-1}
= q^{1/24} / \dedeta(\tau)
= q^{1/24}\,\dedeta(\tau)^{-1}.
\end{equation}
Including the vacuum energy $q^{-c/24} = q^{-1/24}$ for
$c = 1$ (one free boson with central charge $1$ per unit of
$\kappa$) and the non-holomorphic prefactor
$(\operatorname{Im}\tau)^{-1/2}$ from the zero-mode
integral:
\[
Z_1(\cH_1;\,\tau)
= (\operatorname{Im}\tau)^{-1/2}\,
q^{-1/24}\,\bar{q}^{-1/24}\,
\prod_{n=1}^{\infty}
\frac{1}{|1-q^n|^2}
= (\operatorname{Im}\tau)^{-1/2}\,
|\dedeta(\tau)|^{-2}.
\]

\emph{Step 3: level $\kappa$.}
At level $\kappa$, the Fock space is the $\kappa$-fold tensor
product of the $\kappa = 1$ Fock space (or equivalently, the
sewing operator eigenvalues are raised to the power $\kappa$):
\[
Z_1(\cH_\kappa;\,\tau)
= (\operatorname{Im}\tau)^{-\kappa/2}\,
|\dedeta(\tau)|^{-2\kappa}.
\qedhere
\]
\end{proof}

\begin{remark}[Modular transformation]%
% label removed: rem:thqg-X-genus-1-modular%
Under $\tau \mapsto -1/\tau$:
$\dedeta(-1/\tau) = \sqrt{-i\tau}\,\dedeta(\tau)$, so
\[
Z_1(\cH_\kappa;\,-1/\tau)
= |\tau|^{-\kappa}\,
(\operatorname{Im}\tau/|\tau|^2)^{-\kappa/2}\,
|\dedeta(\tau)|^{-2\kappa}
= Z_1(\cH_\kappa;\,\tau).
\]
The $S$-transformation invariance is automatic from the
Fredholm determinant structure: $\det(1 - \sewop_q)$ depends
only on the eigenvalues, which are determined by the conformal
structure of $E_\tau$, not the choice of fundamental domain.
\end{remark}

\begin{computation}[Genus-1 Fredholm determinant: explicit series;
\ClaimStatusProvedHere]%
% label removed: comp:thqg-X-genus-1-series%
\index{Fredholm determinant!genus 1!explicit series}%
The Fredholm determinant at genus~$1$ has the $q$-expansion:
\begin{align}
\det(1 - \sewop_q)^{-1}
&= \prod_{n=1}^{\infty} (1 - q^n)^{-1}
\notag\\
&= 1 + q + 2q^2 + 3q^3 + 5q^4 + 7q^5 + 11q^6 + 15q^7
+ 22q^8 + \cdots
\notag\\
&= \sum_{n=0}^{\infty} p(n)\,q^n,
% label removed: eq:thqg-X-partition-numbers
\end{align}
where $p(n)$ is the number of integer partitions of~$n$.
The asymptotic behavior is given by the Hardy--Ramanujan formula:
\begin{equation}% label removed: eq:thqg-X-hardy-ramanujan
p(n) \;\sim\;
\frac{1}{4n\sqrt{3}}\,
\exp\!\Bigl(\pi\sqrt{2n/3}\Bigr)
\quad\text{as } n \to \infty.
\end{equation}
The subexponential growth $\log p(n) \sim \pi\sqrt{2n/3}$
is precisely the condition needed for HS-sewing
(Theorem~\ref{thm:general-hs-sewing}).
\end{computation}

\begin{remark}[Trace-class vs.\ Fredholm: comparison of
convergence]%
% label removed: rem:thqg-X-trace-vs-fredholm%
The trace $\sum_n q^{n+1} = q/(1-q)$ of $\sewop_q$ converges
for $|q| < 1$, as does every $p$-Schatten norm.
The Fredholm determinant
$\det(1 - \sewop_q)^{-1}
= \prod_n (1 - q^n)^{-1}$
converges absolutely for $|q| < 1$ (the Euler function).
The radius of convergence in~$q$ is $|q| = 1$, where the
product diverges because $\sum_n q^n$ diverges.
On the modular side, $|q| = 1$ corresponds to
$\operatorname{Im}\tau = 0$, the cusp.
\end{remark}

\subsubsection{Genus 2: two handles and the Siegel modular form}

\begin{theorem}[Genus-2 Fredholm determinant;
\ClaimStatusProvedHere]%
% label removed: thm:thqg-X-genus-2-fredholm%
\index{Fredholm determinant!genus 2}%
The genus-$2$ partition function of $\cH_\kappa$ on a
surface $\Sigma_2$ with period matrix
$\Omega = \bigl(\begin{smallmatrix}\tau_1 & z \\ z & \tau_2
\end{smallmatrix}\bigr)
\in \Siegel_2$ is
\begin{equation}% label removed: eq:thqg-X-genus-2-det
Z_2(\cH_\kappa;\,\Omega)
\;=\;
\bigl(\det\operatorname{Im}\Omega\bigr)^{-\kappa/2}
\cdot
\bigl(\zetareg\Delta_{\Sigma_2}\bigr)^{-\kappa/2}.
\end{equation}
The holomorphic part is the Fredholm determinant of the
genus-$2$ sewing kernel $\sewker_2$ acting on $A^2(D)^{\oplus 2}$:
\begin{equation}% label removed: eq:thqg-X-genus-2-hol
Z_2^{\mathrm{hol}}(\cH_1;\,\Omega)
\;=\;
\det\bigl(1 - \sewker_2\bigr)^{-1},
\end{equation}
where $\sewker_2$ is the $2 \times 2$ matrix operator on
$A^2(D)^{\oplus 2}$ with entries determined by the sewing data
of the pair-of-pants decomposition.
\end{theorem}

\begin{proof}
A genus-$2$ surface admits a pair-of-pants decomposition into
$2$ pairs of pants ($2g-2 = 2$) with $3$ sewing circles
($3g-3 = 3$).

\emph{Step 1: pants decomposition.}
Fix a pair-of-pants decomposition: $\Sigma_2 = P_1 \cup_{\partial} P_2$,
where $P_1$ and $P_2$ are each pairs of pants (three-holed spheres).
The three sewing circles have parameters $q_1, q_2, q_3$, where
$q_1$ and $q_2$ correspond to the two handles (B-cycles) and
$q_3$ corresponds to the internal seam.

\emph{Step 2: one-particle operator.}
On the one-particle Bergman space, the sewing produces a
composition of three operators.  Write the one-particle sewing
kernel as a $2 \times 2$ block operator on
$A^2(D)^{\oplus 2}$ (one block per handle):
\begin{equation}% label removed: eq:thqg-X-genus-2-block
\sewker_2
=
\begin{pmatrix}
\sewop_{q_1} & R_{12} \\
R_{21} & \sewop_{q_2}
\end{pmatrix},
\end{equation}
where $\sewop_{q_j}$ are the handle sewing operators and
$R_{12}, R_{21}$ are the off-diagonal composition operators
encoding the coupling between the two handles through the internal
seam.  The off-diagonal entries satisfy
$\|R_{12}\|_{\hscl}, \|R_{21}\|_{\hscl} \leq C\,q_3^{1/2}$
by Lemma~\ref{lem:thqg-X-composition-decay}.

\emph{Step 3: Fredholm determinant.}
Since all diagonal and off-diagonal entries are trace class
(the diagonal operators $\sewop_{q_j}$ are trace class by
Lemma~\ref{lem:thqg-X-sewing-schatten}, and the off-diagonal
operators are Hilbert--Schmidt with $\hscl \cdot \hscl \subset \trcl$),
the full operator $\sewker_2$ is trace class on
$A^2(D)^{\oplus 2}$.  The Fredholm determinant is:
\begin{align}
\det(1 - \sewker_2)
&= \det\!\begin{pmatrix}
1 - \sewop_{q_1} & -R_{12} \\
-R_{21} & 1 - \sewop_{q_2}
\end{pmatrix}
\notag\\
&= \det(1 - \sewop_{q_1})\,
\det\bigl(1 - \sewop_{q_2}
- R_{21}(1-\sewop_{q_1})^{-1}R_{12}\bigr).
% label removed: eq:thqg-X-genus-2-schur
\end{align}
The Schur complement
$S = \sewop_{q_2} + R_{21}(1-\sewop_{q_1})^{-1}R_{12}$
captures the interaction between the two handles.

\emph{Step 4: degeneration limits.}
In the separating degeneration $q_3 \to 0$
(pinching the internal seam):
$R_{12}, R_{21} \to 0$, and
\[
\det(1 - \sewker_2) \to
\det(1 - \sewop_{q_1}) \cdot \det(1 - \sewop_{q_2})
= \dedeta(\tau_1)^2 \cdot \dedeta(\tau_2)^2,
\]
recovering the product of two genus-$1$ partition functions.
In the non-separating degeneration $q_1 \to 0$
(collapsing one handle):
$\sewop_{q_1} \to 0$ and $\det(1 - \sewker_2)$ reduces to
$\det(1 - \sewop_{q_2} - R_{21}R_{12})$, a genus-$1$
Fredholm determinant with modified sewing operator.

\emph{Step 5: Polyakov formula.}
Applying the Polyakov--Alvarez formula
\eqref{eq:thqg-X-polyakov} at genus~$2$ gives
$\det(1 - \sewker_2)^{-1}
= (\det\operatorname{Im}\Omega)^{-1/2}
\cdot (\zetareg\Delta_{\Sigma_2})^{-1/2}$.
Raising to the power $\kappa$ yields the stated formula.
\end{proof}

\begin{computation}[Genus-2 Fredholm determinant: leading terms;
\ClaimStatusProvedHere]%
% label removed: comp:thqg-X-genus-2-series%
\index{Fredholm determinant!genus 2!leading terms}%
Write $q_1 = e^{2\pi i \tau_1}$, $q_2 = e^{2\pi i \tau_2}$,
$r = e^{2\pi i z}$ (the off-diagonal period).
In the product expansion:
\begin{align}
\det(1 - \sewker_2)^{-1}
&= \prod_{n=1}^{\infty}(1-q_1^n)^{-1}
\cdot \prod_{n=1}^{\infty}(1-q_2^n)^{-1}
\cdot \bigl(1 + O(r)\bigr)
\notag\\
&= \frac{1}{\dedeta(\tau_1)\,\dedeta(\tau_2)}
\cdot \Bigl(1 + \sum_{n,m \geq 1}
c_{n,m}\,r^n\,\bar{r}^m + \cdots\Bigr).
% label removed: eq:thqg-X-genus-2-expansion
\end{align}
The leading correction $c_{1,1}$ arises from the one-particle
Schur complement and equals
$\sum_{k \geq 0} q_1^{k+1}\,q_2^{k+1}/(1-q_1^{k+1})(1-q_2^{k+1})$.
For $\tau_1 = \tau_2 = i$ (square torus) and $z = i/2$:
$|r| = e^{-\pi}$ and $|\dedeta(i)|^{-2} = \Gamma(1/4)^2/(2\pi^{3/2})$.
\end{computation}

\begin{remark}[Arithmetic significance of genus-$2$
non-diagonality]%
\label{rem:genus2-arithmetic-significance}%
\index{sewing kernel!genus-2 arithmetic significance}%
The off-diagonal blocks $R_{12}, R_{21}$ in the genus-$2$
sewing kernel $K_2$ have a precise arithmetic consequence.
At genus~$1$, the sewing operator is Fock-diagonal, which
forces the two-variable $L$-object $L_\cA(s,u)$ to collapse
to a single-variable Dirichlet series $S_\cA(s{+}u{-}1)$
(Vol~I, sewing--Hecke reciprocity theorem);
this collapse is the structural reason that genus-$1$ MC
data cannot access zero-location information
(Vol~I, structural separation theorem).
At genus~$2$, the off-diagonal coupling $R_{ij} \neq 0$
(for generic period matrix) breaks this diagonality.
The Schur complement
$S = K_{q_2} + R_{21}(1{-}K_{q_1})^{-1}R_{12}$
encodes inter-handle interaction that is invisible to the
genus-$1$ sewing operator; the resulting Fredholm
determinant depends nontrivially on the off-diagonal
period~$z$, and the Rankin--Selberg unfolding on
$\mathrm{Sp}_4(\mathbb{Z}) \backslash \mathfrak{H}_2$
produces genuinely multi-variable $L$-functions.
The genus-$2$ arithmetic programme
(Vol~I, arithmetic shadows chapter,
genus-$2$ arithmetic frontier)
exploits this non-diagonality to escape the genus-$1$
structural barrier.
\end{remark}

\begin{remark}[The four Dirichlet series and the arity--genus orthogonality]%
\label{rem:four-dirichlet-arity-genus}%
\index{Dirichlet series!four-object hierarchy}%
\index{arity--genus orthogonality|textbf}
The genus-$1$ structural barrier of the preceding remark
is a manifestation of a deeper phenomenon: the arithmetic
programme of Volume~I involves four distinct Dirichlet series that
live on orthogonal axes of the bigraded shadow algebra
$\cA^{\mathrm{sh}}_{r,g}$.

\begin{enumerate}[label=\textup{(\alph*)}]
\item The \emph{shadow-metric Epstein zeta}
  $\varepsilon_{Q_L}(s)$ and the \emph{shadow Dirichlet series}
  $\zeta_\cA(s) = \sum_{r \geq 2} S_r\,r^{-s}$
  live on the \textbf{arity axis}~$(r, 0)$: they are
  genus-$0$ invariants determined by the OPE data
  $(\kappa, \alpha, S_4)$.

\item The \emph{constrained Epstein zeta}
  $\varepsilon^c_s = \sum_\Delta d(\Delta)\,(2\Delta)^{-s}$
  (Benjamin--Chang~\cite{BenjaminChang22})
  lives on the \textbf{genus axis}~$(0, 1)$: it is the Eisenstein
  spectral coefficient of the descendant-stripped partition
  function $\hat Z^c = y^{c/2}|\eta|^{2c}\,Z$ on
  $\mathrm{SL}(2,\bZ)\backslash\mathfrak{H}$.

\item The \emph{categorical zeta}
  $\zeta^{\mathrm{DK}}_\fg(s) = \sum_\lambda
  \dim(V_\lambda)^{-s}$ is a purely Lie-theoretic invariant
  ($= \zeta(s)$ for $\fsl_2$, Vol~I,
  Proposition~\ref*{prop:categorical-zeta-sl2-riemann}).
\end{enumerate}

The nontrivial Riemann zeta zeros~$\rho_n$ appear through
object~(b), not~(a): the scattering factor
$\zeta(2s)/\zeta(2s{-}1)$ in the Eisenstein spectral
decomposition on~$\cM_{1,1}$ has poles at
$s = (1 + \rho_n)/2$.  This is a theorem about the spectral
theory of the \emph{moduli space}, not about the bar complex.
The ``negative results''
(shadow zeta $\notin$ Selberg class, no Euler product)
concern object~(a), which lives on the arity axis and has
no access to the Riemann zeta zeros.

The genus-$2$ non-diagonality of the preceding remark
provides the first mechanism for the arity and genus axes to
\emph{interact}: the off-diagonal sewing blocks~$R_{ij}$
couple the handle structure (genus axis) to the inter-handle
OPE data (arity axis), breaking the genus-$1$ factorisation
barrier.  The Dirichlet-sewing lift
$S_\cA(u) = \zeta(u{+}1)\sum_i (\zeta(u) - H_{w_i-1}(u))$
(Vol~I) is a fifth related object, built from the weight
multiset~$W(\cA)$, that contains $\zeta$-factors
($S_\cH(u) = \zeta(u)\zeta(u{+}1)$ for Heisenberg)
but is independent of the shadow obstruction tower.
See Volume~I, \S\ref*{subsec:four-dirichlet-series}
for the full disambiguation.
\end{remark}

\subsubsection{Genus 3: the three-handle Fredholm determinant}

\begin{theorem}[Genus-3 Fredholm determinant;
\ClaimStatusProvedHere]%
% label removed: thm:thqg-X-genus-3-fredholm%
\index{Fredholm determinant!genus 3}%
The genus-$3$ partition function of $\cH_\kappa$ on
$\Sigma_3$ with period matrix $\Omega \in \Siegel_3$ is
\begin{equation}% label removed: eq:thqg-X-genus-3-det
Z_3(\cH_\kappa;\,\Omega)
\;=\;
\bigl(\det\operatorname{Im}\Omega\bigr)^{-\kappa/2}
\cdot
\bigl(\zetareg\Delta_{\Sigma_3}\bigr)^{-\kappa/2}.
\end{equation}
The one-particle sewing kernel $\sewker_3$ acts on
$A^2(D)^{\oplus 3}$ as a $3 \times 3$ operator:
\begin{equation}% label removed: eq:thqg-X-genus-3-block
\sewker_3
=
\begin{pmatrix}
\sewop_{q_1} & R_{12} & R_{13} \\
R_{21} & \sewop_{q_2} & R_{23} \\
R_{31} & R_{32} & \sewop_{q_3}
\end{pmatrix},
\end{equation}
where the diagonal entries $\sewop_{q_j}$ correspond to the
three handles and the off-diagonal entries $R_{ij}$ encode the
pair-of-pants composition operators for the internal seams.
\end{theorem}

\begin{proof}
A genus-$3$ surface admits a pair-of-pants decomposition into
$4$ pairs of pants ($2g - 2 = 4$) with $6$ sewing circles
($3g - 3 = 6$).  Three of these circles correspond to handles
(B-cycles), and the remaining three correspond to internal
seams.

\emph{Step 1: block structure.}
The $3 \times 3$ block operator $\sewker_3$ on
$A^2(D)^{\oplus 3}$ has diagonal entries $\sewop_{q_j}$
for the three handle sewing operators with
$\|\sewop_{q_j}\|_1 = q_j/(1-q_j)$ (trace class), and
off-diagonal entries $R_{ij}$ with
$\|R_{ij}\|_{\hscl} \leq C\,q_{ij}^{1/2}$
(Hilbert--Schmidt) where $q_{ij}$ is the sewing parameter
of the internal seam connecting handles $i$ and $j$.

\emph{Step 2: trace-class verification.}
The operator $\sewker_3$ is trace class on $A^2(D)^{\oplus 3}$.
The trace norm satisfies
\[
\|\sewker_3\|_1
\leq \sum_{j=1}^{3} \|\sewop_{q_j}\|_1
+ \sum_{i \neq j} \|R_{ij}\|_1
\leq \sum_{j=1}^{3} \frac{q_j}{1-q_j}
+ \sum_{i \neq j}
\|R_{ij}\|_{\hscl}^2
< \infty.
\]

\emph{Step 3: Fredholm determinant.}
The Fredholm determinant is computed by the standard $3 \times 3$
block formula:
\begin{align}
\det(1 - \sewker_3)
&= \det(1 - \sewop_{q_1})
\cdot \det(1 - S_{23|1})
\notag\\
&\quad\text{where } S_{23|1}
= \begin{pmatrix} \sewop_{q_2} + R_{21}(1-\sewop_{q_1})^{-1}R_{12}
& R_{23} + R_{21}(1-\sewop_{q_1})^{-1}R_{13}\\
R_{32} + R_{31}(1-\sewop_{q_1})^{-1}R_{12}
& \sewop_{q_3} + R_{31}(1-\sewop_{q_1})^{-1}R_{13}
\end{pmatrix}
\notag
\end{align}
is the Schur complement.  Iterating the Schur complement
reduction gives:
\begin{equation}% label removed: eq:thqg-X-genus-3-iterated-schur
\det(1 - \sewker_3)
= \prod_{j=1}^{3} \det(1 - \sewop_{q_j})
\cdot \det(1 - \Delta_3),
\end{equation}
where $\Delta_3$ is a trace-class operator encoding the
inter-handle coupling.  In the degeneration limit where all
internal seams are pinched ($R_{ij} \to 0$), $\Delta_3 \to 0$
and the determinant factors into a product of three genus-$1$
determinants:
\[
\det(1 - \sewker_3)
\to
\prod_{j=1}^{3}
\dedeta(\tau_j)^2.
\]

\emph{Step 4: Polyakov formula.}
The Polyakov--Alvarez formula at genus~$3$ gives
$\det(1-\sewker_3)^{-1}
= (\det\operatorname{Im}\Omega)^{-1/2}
\cdot (\zetareg\Delta_{\Sigma_3})^{-1/2}$.
\end{proof}

\begin{remark}[The Schottky problem at genus~$3$]%
% label removed: rem:thqg-X-schottky%
\index{Schottky problem}%
At genus~$3$, the Siegel upper half-space $\Siegel_3$ has
complex dimension $\binom{3+1}{2} = 6$, and the moduli space
$\cM_3$ has dimension $3g-3 = 6$.  The Torelli map
$\cM_3 \to \Siegel_3/\Sp_6(\Z)$ is therefore generically
surjective at genus~$3$ (this is the classical fact that the
Schottky locus equals the full Siegel space at genus~$\leq 3$).
At genus~$4$, the Siegel space has dimension $10$ while $\cM_4$
has dimension $9$, and the Schottky locus becomes a proper
subvariety.  The Fredholm determinant formula
\eqref{eq:thqg-X-genus-3-det} is therefore a function on all of
$\Siegel_3$ at genus~$3$, but only on the Schottky locus of
$\Siegel_4$ at genus~$4$.
\end{remark}

\subsubsection{General genus: the recursive structure}

\begin{theorem}[General-genus Fredholm structure;
\ClaimStatusProvedHere]%
% label removed: thm:thqg-X-general-genus-fredholm%
\index{Fredholm determinant!general genus}%
For every genus $g \geq 1$, the partition function of
$\cH_\kappa$ on $\Sigma_g$ with period matrix
$\Omega \in \Siegel_g$ is
\begin{equation}% label removed: eq:thqg-X-general-genus-formula
Z_g(\cH_\kappa;\,\Omega)
\;=\;
\bigl(\det\operatorname{Im}\Omega\bigr)^{-\kappa/2}
\cdot
\bigl(\zetareg\Delta_{\Sigma_g}\bigr)^{-\kappa/2}
\;=\;
\det\bigl(1 - \sewker_g\bigr)^{-\kappa},
\end{equation}
where $\sewker_g$ is a trace-class operator on
$A^2(D)^{\oplus g}$, the genus-$g$ sewing kernel of
Definition~\textup{\ref{def:thqg-X-genus-g-sewing}}.
The following recursive properties hold:
\begin{enumerate}[label=\textup{(\roman*)}]
\item \emph{Separating degeneration.}
  If $\Sigma_g$ degenerates into $\Sigma_{g_1} \cup \Sigma_{g_2}$
  with $g_1 + g_2 = g$, then
  \begin{equation}% label removed: eq:thqg-X-separating-factorization
  Z_g(\cH_\kappa) \to
  Z_{g_1}(\cH_\kappa) \cdot Z_{g_2}(\cH_\kappa).
  \end{equation}
\item \emph{Non-separating degeneration.}
  If $\Sigma_g$ degenerates by pinching a non-separating cycle,
  then
  \begin{equation}% label removed: eq:thqg-X-nonseparating-limit
  Z_g(\cH_\kappa)
  \;\to\;
  \operatorname{Tr}_{A^2}\bigl(\sewop_q \cdot Z_{g-1}(\cH_\kappa)
  \bigr),
  \end{equation}
  where the trace is over the one-particle space of the
  collapsing handle and $q \to 0$ is the sewing parameter.
\item \emph{Free energy.}
  The genus-$g$ free energy
  $F_g(\cH_\kappa) = \log Z_g(\cH_\kappa)$ satisfies
  \begin{equation}% label removed: eq:thqg-X-free-energy-genus
  F_g(\cH_\kappa)
  \;=\;
  -\kappa\,\operatorname{Tr}\log(1 - \sewker_g)
  \;=\;
  \kappa\sum_{k=1}^{\infty}
  \frac{1}{k}\operatorname{Tr}(\sewker_g^k).
  \end{equation}
\end{enumerate}
\end{theorem}

\begin{proof}
The general formula follows from the Heisenberg sewing theorem
(Theorem~\ref{thm:thqg-X-heisenberg-sewing-full}\,(IV)) and
the Polyakov--Alvarez formula at genus~$g$.

Part~(i): in the separating degeneration, the genus-$g$
sewing kernel $\sewker_g$ becomes block-diagonal:
$\sewker_g \to \sewker_{g_1} \oplus \sewker_{g_2}$
as the separating seam parameter $q_s \to 0$.
Hence $\det(1 - \sewker_g) \to
\det(1 - \sewker_{g_1}) \cdot \det(1 - \sewker_{g_2})$.

Part~(ii): in the non-separating degeneration, one diagonal
block $\sewop_{q_j} \to 0$ as $q_j \to 0$, and the Schur
complement formula reduces the rank of the sewing kernel by one.

Part~(iii): $\log\det(1-\sewker_g)^{-\kappa}
= -\kappa\,\operatorname{Tr}\log(1-\sewker_g)
= \kappa\sum_{k} k^{-1}\operatorname{Tr}(\sewker_g^k)$.
\end{proof}

\begin{proposition}[Free energy: $\hat{A}$-genus verification;
\ClaimStatusProvedHere]%
% label removed: prop:thqg-X-free-energy-ahat%
\index{free energy!$\hat{A}$-genus}%
For $\cH_\kappa$ at level $\kappa$, the Faber--Pandharipande
free energy formula
$F_g = \kappa \cdot \lambda_g^{\mathrm{FP}}$ is verified
through $g = 3$:
\begin{center}
\begin{tabular}{c|c|c|c}
$g$ & $\lambda_g^{\mathrm{FP}}$ & $F_g/\kappa$ & Source \\
\hline
$1$ & $\frac{|B_2|}{2 \cdot 2!} = \frac{1}{24}$ & $\frac{1}{24}$
& $\dedeta(\tau)^{-1}$: constant term \\[4pt]
$2$ & $\frac{2^3-1}{2^3}\frac{|B_4|}{4!}
= \frac{7}{8}\cdot\frac{1}{720}
= \frac{7}{5760}$ & $\frac{7}{5760}$
& Polyakov at $g = 2$ \\[4pt]
$3$ & $\frac{2^5-1}{2^5}\frac{|B_6|}{6!}
= \frac{31}{32}\cdot\frac{1}{30240}
= \frac{31}{967680}$ & $\frac{31}{967680}$
& Polyakov at $g = 3$
\end{tabular}
\end{center}
These match the $\hat{A}$-genus coefficients
$\hat{A}_g = \frac{(2^{2g-1}-1)|B_{2g}|}{2^{2g-1}(2g)!}$
(Theorem~\textup{\ref{thm:family-index}}).
\end{proposition}

\begin{proof}
The genus-$g$ free energy of a single free boson is
$F_g = -\frac{1}{2}\log\zetareg\Delta_{\Sigma_g}$
integrated over $\overline{\cM}_g$.  By Faltings'
arithmetic Riemann--Roch formula and the Mumford relations,
$\int_{\overline{\cM}_g} \lambda_g
= \frac{(2^{2g-1}-1)|B_{2g}|}{2^{2g-1}(2g)!}$
(Faber's formula).  The claim follows by direct substitution.
\end{proof}

% =====================================================================
%
%  SECTION 4: SCALAR SATURATION AND CLASS-G ALGEBRAS
%
% =====================================================================

\subsection{Extension to class-G algebras via scalar saturation}
% label removed: subsec:thqg-X-class-G

The Heisenberg sewing theorem gives a Fredholm determinant
formula for the simplest possible chiral algebra, the free boson.
The shadow depth classification
(Definition~\ref{def:shadow-depth-classes}
in Appendix~\ref*{V1-app:nonlinear-modular-shadows})
organizes chiral algebras into four classes based on the
termination behavior of the shadow obstruction tower:
\begin{itemize}
\item Class~\textbf{G} (Gaussian): $r_{\max} = 2$.
  The shadow obstruction tower terminates at weight~$2$:
  $\Theta_\cA^{\min} = \eta\otimes\Gamma_\cA$.
  On the proved uniform-weight scalar lane this specializes to
  $\Theta_\cA^{\min} = \kappa(\cA)\cdot\eta \otimes \Lambda$.
  Examples: $\cH_\kappa$, lattice algebras $V_\Lambda$.
\item Class~\textbf{L} (Lie): $r_{\max} = 3$.
  Examples: $V_k(\fg)$ at generic level.
\item Class~\textbf{C} (contact): $r_{\max} = 4$.
  Examples: $\beta\gamma$ system.
\item Class~\textbf{M} (mixed): $r_{\max} = \infty$.
  Examples: $\mathrm{Vir}_c$, $\cW_N$ for $N \geq 3$.
\end{itemize}
For class-G algebras, the scalar rigidity theorem
(Theorem~\ref{thm:algebraic-family-rigidity}) ensures only that the
minimal MC class is purely scalar in arity.  This does not by itself
identify the full nonperturbative partition function with a Fredholm
determinant.  The Fredholm formula below is proved for Gaussian
algebras on the uniform-weight scalar lane.

\begin{theorem}[Fredholm scalar partition function on the Gaussian
scalar lane; \ClaimStatusProvedHere]%
% label removed: thm:thqg-X-class-G-fredholm%
\index{partition function!class G}%
\index{scalar saturation!Fredholm determinant}%
Let $\cA$ be a class-G Gaussian chiral algebra on the proved
uniform-weight scalar lane
\textup{(}shadow depth $r_{\max} = 2$\textup{)} with
modular characteristic $\kappa = \kappa(\cA)$ and central
charge $c(\cA)$.  Assume HS-sewing
\textup{(}Definition~\textup{\ref{def:hs-sewing}}\textup{)}.
Then the scalar genus-$g$ partition function is a Fredholm
determinant:
\begin{equation}% label removed: eq:thqg-X-class-G-formula
Z_g^{\mathrm{scal}}(\cA;\,\Omega)
\;=\;
Z_g^{\mathrm{G}}(\kappa;\,\Omega)
\;:=\;
\bigl(\det\operatorname{Im}\Omega\bigr)^{-\kappa/2}
\cdot
\bigl(\zetareg\Delta_{\Sigma_g}\bigr)^{-\kappa/2},
\end{equation}
up to a $\kappa$-dependent prefactor encoding the vacuum energy
$q^{-c/24}$.  This formula depends on $\cA$ only through
$\kappa(\cA)$ on that scalar lane.  For general multi-weight
class-G algebras, the identification of the scalar coefficient with
$\kappa(\cA)\Lambda$ remains open, so no full Fredholm theorem is
claimed here.
\end{theorem}

\begin{proof}
On the proved uniform-weight Gaussian lane,
Theorem~\ref{thm:theta-direct-derivation} identifies the minimal
scalar package as
$\Theta_\cA^{\min} = \kappa(\cA)\eta\otimes\Lambda$.
The scalar genus-$g$ partition function is therefore the
$\kappa(\cA)$-multiple of the Heisenberg scalar partition function.
The HS-sewing condition ensures convergence of the Fredholm
determinant (by Theorem~\ref{thm:general-hs-sewing}, which
requires polynomial OPE growth and subexponential sector growth;
class-G algebras satisfy both).

The Fredholm determinant formula then follows from the
Heisenberg computation
(Theorem~\ref{thm:thqg-X-heisenberg-sewing-full}\,(IV))
with $\kappa$ replaced by $\kappa(\cA)$.
\end{proof}

\begin{example}[Lattice vertex algebras]%
% label removed: ex:thqg-X-lattice-fredholm%
\index{lattice algebra!Fredholm determinant}%
Let $V_\Lambda$ be a lattice vertex algebra associated to an
even integral lattice $\Lambda$ of rank $r$.  Then
$\kappa(V_\Lambda) = r$ (the rank of the lattice), and the
genus-$g$ partition function is
\begin{equation}% label removed: eq:thqg-X-lattice-formula
Z_g(V_\Lambda;\,\Omega)
\;=\;
\bigl(\det\operatorname{Im}\Omega\bigr)^{-r/2}
\cdot
\bigl(\zetareg\Delta_{\Sigma_g}\bigr)^{-r/2}
\cdot
\Theta_\Lambda^{(g)}(\Omega),
\end{equation}
where $\Theta_\Lambda^{(g)}(\Omega)
= \sum_{\lambda \in \Lambda^g}
\exp(\pi i \lambda^T \Omega \lambda)$ is the genus-$g$
lattice theta function.  The Fredholm part (first two factors)
carries the oscillator modes and is universal in $r$;
the theta function carries the zero-mode sum over the
lattice and is lattice-specific.

The HS-sewing condition is verified: $V_\Lambda$ has sector
growth $\dim (V_\Lambda)_n \leq |\Lambda^*/\Lambda| \cdot p(n)^r$
(polynomial times partition function growth), and OPE
coefficients are bounded by $\kappa = r$ times the lattice
norm (Corollary~\ref{cor:hs-sewing-standard-landscape}).
\end{example}

\begin{remark}[Curvature-braiding orthogonality for lattice algebras]%
% label removed: rem:thqg-X-lattice-curvature-braiding%
\index{lattice algebra!curvature-braiding orthogonality}%
By Theorem~\ref{thm:lattice:curvature-braiding-orthogonal},
$\kappa(V_\Lambda^{N,q}) = \mathrm{rank}(\Lambda)$
independent of the cocycle~$q$.  This means the Fredholm
determinant part of the partition function is insensitive to
the choice of cocycle; only the theta function
$\Theta_\Lambda^{(g)}$ depends on the lattice embedding.
This orthogonality between the oscillator (Fredholm) sector
and the zero-mode (theta) sector is a direct consequence of
the class-G property.
\end{remark}

\begin{proposition}[Rank additivity;
\ClaimStatusProvedHere]%
% label removed: prop:thqg-X-rank-additivity%
\index{Fredholm determinant!rank additivity}%
If $\cA = \cA_1 \otimes \cA_2$ is a tensor product of
class-G algebras, then
\begin{equation}% label removed: eq:thqg-X-rank-additivity
Z_g(\cA;\,\Omega)
= Z_g(\cA_1;\,\Omega)
\cdot Z_g(\cA_2;\,\Omega),
\end{equation}
reflecting the additivity of the modular characteristic
$\kappa(\cA) = \kappa(\cA_1) + \kappa(\cA_2)$
and the multiplicativity of the Fredholm determinant
$\det(1 - T_1 \oplus T_2)
= \det(1 - T_1) \cdot \det(1 - T_2)$.
\end{proposition}

\begin{proof}
By Proposition~\ref{prop:independent-sum-factorization},
the modular characteristic is additive:
$\kappa(\cA_1 \otimes \cA_2)
= \kappa(\cA_1) + \kappa(\cA_2)$.
The sewing operator for a tensor product is the direct sum:
$\sewker_g(\cA_1 \otimes \cA_2)
= \sewker_g(\cA_1) \oplus \sewker_g(\cA_2)$
on $A^2(D)^{\oplus g_1} \oplus A^2(D)^{\oplus g_2}$.
The Fredholm determinant of a direct sum is the product:
$\det(1 - T_1 \oplus T_2) = \det(1-T_1)\det(1-T_2)$.
\end{proof}

% =====================================================================
%
%  SECTION 5: NON-FREDHOLM CORRECTIONS FOR CLASSES L, C, M
%
% =====================================================================

\subsection{Non-Fredholm corrections for classes L, C, M}
% label removed: subsec:thqg-X-non-fredholm

For algebras of shadow depth $r_{\max} > 2$, the partition
function is no longer a pure Fredholm determinant.  The higher
shadow data ($\fC$, $\mathfrak{Q}$, and all higher
$o_r$ for $r \geq 3$) produce perturbative corrections
to the Gaussian (Fredholm) base, organized as Feynman
integrals weighted by the graph-sum formula of
Definition~\ref*{V1-def:modular-bar-hamiltonian}.

\begin{definition}[Non-Fredholm correction]%
% label removed: def:thqg-X-non-fredholm-correction%
\index{non-Fredholm correction|textbf}%
\index{partition function!non-Fredholm correction}%
For a chiral algebra $\cA$ with shadow depth $r_{\max}$, define
the \emph{genus-$g$ non-Fredholm correction} as
\begin{equation}% label removed: eq:thqg-X-non-fredholm-correction
\delta Z_g(\cA;\,\Omega)
\;:=\;
Z_g(\cA;\,\Omega)
\;-\;
Z_g^{\mathrm{G}}(\kappa(\cA);\,\Omega),
\end{equation}
where $Z_g^{\mathrm{G}}$ is the Gaussian (Fredholm) base of
\eqref{eq:thqg-X-class-G-formula}.  The correction
$\delta Z_g$ encodes the contribution of all shadow data
beyond the modular characteristic $\kappa$.
\end{definition}

\begin{theorem}[Feynman integral expansion;
\ClaimStatusProvedHere]%
% label removed: thm:thqg-X-feynman-expansion%
\index{Feynman integral!non-Fredholm correction}%
For a modular Koszul chiral algebra $\cA$ satisfying HS-sewing,
the genus-$g$ non-Fredholm correction admits a Feynman integral
expansion:
\begin{equation}% label removed: eq:thqg-X-feynman-expansion
\delta Z_g(\cA;\,\Omega)
\;=\;
\sum_{r=3}^{r_{\max}}
\sum_{\Gamma \in \mathcal{G}^{\mathrm{conn}}_{g,r}}
\frac{1}{|\Aut(\Gamma)|}
\int_{\Sigma_g^{v(\Gamma)}}
\omega_\Gamma(\Omega),
\end{equation}
where:
\begin{itemize}
\item $\mathcal{G}^{\mathrm{conn}}_{g,r}$ is the set of
  connected stable graphs of genus~$g$ with vertices of
  arity at most~$r$;
\item $v(\Gamma)$ is the number of vertices of~$\Gamma$;
\item $\omega_\Gamma(\Omega)$ is the integrand determined by
  the graph~$\Gamma$, with vertex labels from the transferred
  cyclic minimal model $\ell_k^{\mathrm{tr}}$ of $\cA$ and
  edge contractions by the complementarity propagator
  $P_\cA = H_\cA^{-1}$ (the inverse of the Hodge Laplacian).
\end{itemize}
The Feynman expansion converges absolutely for $|q_\alpha| < 1$
(all sewing parameters inside the unit disk).
\end{theorem}

\begin{proof}
The genus-$g$ partition function is computed by the shadow
Postnikov tower: $Z_g = \langle e^{\Theta_\cA^{(g)}} \rangle_g$.
The Gaussian part $Z_g^{\mathrm{G}}$ is the contribution of
$\Theta_\cA^{\leq 2}$, which consists of the modular
characteristic $\kappa$ contracted with the Hodge cycle.
The non-Fredholm correction is the contribution of
$\Theta_\cA - \Theta_\cA^{\leq 2}$, which involves the
higher-arity terms $\ell_k^{\mathrm{tr}}$ for $k \geq 3$.

Expanding $e^{\Theta}$ in the graph-sum formula
(Definition~\ref*{V1-def:modular-bar-hamiltonian}), each
graph $\Gamma$ with a vertex of arity $r > 2$ contributes to
$\delta Z_g$.  The weight of a graph is
$w(\Gamma) = \sum_{v \in V(\Gamma)} \mathrm{arity}(v)$,
and graphs with $w > 2|V(\Gamma)|$ contribute to the
non-Fredholm correction.

Convergence: the HS-sewing condition
(Theorem~\ref{thm:general-hs-sewing}) ensures that each
Feynman integral converges absolutely.  The propagator
$P_\cA = H_\cA^{-1}$ on $\Sigma_g$ has the same exponential
decay as the Bergman composition operator
(Lemma~\ref{lem:thqg-X-composition-decay}), and the vertex
labels $\ell_k^{\mathrm{tr}}$ are bounded by the polynomial
OPE growth.  The arity-cutoff lemma (Lemma~\ref{lem:arity-cutoff})
bounds the number of contributing graphs at each order.
\end{proof}

\subsubsection{Class L: cubic corrections from Lie algebras}

\begin{proposition}[Class-L Feynman integrals;
\ClaimStatusProvedHere]%
% label removed: prop:thqg-X-class-L-feynman%
\index{class L!Feynman integrals}%
For a class-L algebra $\cA$ \textup{(}$r_{\max} = 3$,
e.g.\ $V_k(\fg)$ at generic level\textup{)}, the
non-Fredholm correction at genus~$g$ involves only
trivalent graphs:
\begin{equation}% label removed: eq:thqg-X-class-L-correction
\delta Z_g^{\mathrm{L}}(\cA;\,\Omega)
\;=\;
\sum_{\Gamma \in \mathcal{G}^{3\text{-val}}_{g}}
\frac{1}{|\Aut(\Gamma)|}
\int_{\Sigma_g^{v(\Gamma)}}
\omega_\Gamma^{(3)}(\Omega),
\end{equation}
where $\mathcal{G}^{3\text{-val}}_g$ denotes the set of
connected $3$-valent stable graphs of genus~$g$, and
$\omega_\Gamma^{(3)}$ is the integrand with cubic vertex
labels $\ell_3^{\mathrm{tr}}$ (the transferred $L_\infty$
ternary bracket).

A $3$-valent graph of genus~$g$ with $v$ vertices has
$e = 3v/2$ edges and loop number $g = e - v + 1 = v/2 + 1$,
so $v = 2(g-1)$ and $e = 3(g-1)$.  In particular:
\begin{center}
\begin{tabular}{c|c|c|c}
$g$ & $v$ & $e$ & \# connected $3$-valent graphs \\
\hline
$1$ & $0$ & $0$ & $0$ (no correction) \\
$2$ & $2$ & $3$ & $2$ (theta, sunset) \\
$3$ & $4$ & $6$ & \text{finitely many}
\end{tabular}
\end{center}
The first non-trivial correction appears at genus~$2$
with the theta graph and the sunset graph.
\end{proposition}

\begin{proof}
For class-L algebras, $\ell_k^{\mathrm{tr}} = 0$ for $k \geq 4$
(the shadow obstruction tower terminates at arity~$3$).  Hence only graphs
with vertices of arity exactly~$3$ contribute to $\delta Z_g$.
These are the $3$-regular (trivalent) graphs.

The graph count follows from the Euler relation: a connected
graph with $v$ vertices, $e$ edges, and first Betti number
$g$ satisfies $g = e - v + 1$.  For $3$-regular graphs,
$3v = 2e$ (every edge has two endpoints; every vertex has
three), so $e = 3v/2$ and $g = v/2 + 1$.
\end{proof}

\begin{example}[Affine $\AffKM{g}_k$ at genus~$2$]%
% label removed: ex:thqg-X-affine-genus-2%
\index{affine Kac--Moody!genus-2 correction}%
For $V_k(\fg)$ at generic level $k \neq -h^\vee$, the
genus-$2$ non-Fredholm correction involves two Feynman graphs.
Let $f_{abc}$ denote the structure constants of $\fg$
and $P(z,w;\Omega)$ the genus-$2$ propagator.  Then:
\begin{enumerate}[label=\textup{(\alph*)}]
\item \emph{Theta graph}
  $(\Gamma_\theta$: two vertices, three edges$)$:
  \begin{equation}% label removed: eq:thqg-X-theta-graph
  I_{\theta}
  = \frac{1}{12}
  f_{abc}\,f_{a'b'c'}\,
  \langle a,a'\rangle\langle b,b'\rangle\langle c,c'\rangle
  \int_{\Sigma_2 \times \Sigma_2}
  P(z,w;\Omega)^3\,dz\,dw.
  \end{equation}
\item \emph{Sunset graph}
  $(\Gamma_{\mathrm{sun}}$: two vertices, three edges,
  one edge connects vertex to itself$)$:
  \begin{equation}% label removed: eq:thqg-X-sunset-graph
  I_{\mathrm{sun}}
  = \frac{1}{4}
  f_{abc}\,f_{a'b'c'}\,
  \langle a,a'\rangle
  \int_{\Sigma_2}
  \langle b, c\rangle_z
  \,P(z,w;\Omega)\,
  \langle b', c'\rangle_w\,
  dz\,dw.
  \end{equation}
\end{enumerate}
The total correction is
$\delta Z_2^{\mathrm{L}}(V_k(\fg);\,\Omega)
= I_\theta + I_{\mathrm{sun}}$.
At integrable level $k \in \Z_{\geq 0}$, these Feynman
integrals reduce to combinatorial sums via the Verlinde
formula (Example~\ref{eq:verlinde-general}).

% ----------------------------------------------------------------
% GENUS-2 FEYNMAN GRAPHS FOR CLASS L
% ----------------------------------------------------------------
\begin{center}
\begin{tikzpicture}[x=1cm, y=1cm]
% ---- Theta graph ----
\node[font=\footnotesize\scshape] at (0,1.2) {$\Gamma_\theta$};
\node[gv0, label={[font=\tiny]below left:$f_{abc}$}]
  (t1) at (-0.4,0) {};
\node[gv0, label={[font=\tiny]below right:$f_{a'b'c'}$}]
  (t2) at (0.4,0) {};
\draw[iedge] (t1) to[bend left=40] node[font=\tiny, above]
  {$P$} (t2);
\draw[iedge] (t1) -- (t2);
\draw[iedge] (t1) to[bend right=40] (t2);
\node[font=\tiny, text=black!50] at (0,-0.8) {$I_\theta
  = \frac{1}{12}\!\int P(z,w)^3$};
%
% ---- Plus sign ----
\node[font=\normalsize] at (2.2,0) {$+$};
%
% ---- Sunset graph ----
\node[font=\footnotesize\scshape] at (4.2,1.2) {$\Gamma_{\mathrm{sun}}$};
\node[gv0, label={[font=\tiny]below left:$f_{abc}$}]
  (s1) at (3.8,0) {};
\node[gv0, label={[font=\tiny]below right:$f_{a'b'c'}$}]
  (s2) at (4.6,0) {};
\draw[iedge] (s1) -- (s2);
\draw[iedge] (s1) to[bend left=40] (s2);
\draw[sloop] (s1) to[out=150,in=210,looseness=5] (s1);
\node[font=\tiny, text=black!50] at (4.2,-0.8)
  {$I_{\mathrm{sun}} = \frac{1}{4}\!\int
    \langle b,c\rangle_z\, P\, \langle b',c'\rangle_w$};
%
% ---- Equals ----
\node[font=\normalsize] at (6.6,0) {$=$};
\node[font=\footnotesize] at (8.2,0)
  {$\delta Z_2^{\mathrm{L}}(V_k(\fg);\,\Omega)$};
\end{tikzpicture}
\end{center}
\end{example}

\subsubsection{Class C: quartic contact corrections}

\begin{proposition}[Class-C quartic contact;
\ClaimStatusProvedHere]%
% label removed: prop:thqg-X-class-C-quartic%
\index{class C!quartic contact}%
For a class-C algebra $\cA$ with $r_{\max} = 4$
\textup{(}e.g.\ the $\beta\gamma$ system\textup{)},
the non-Fredholm correction at genus~$g$ involves
graphs with vertices of arity $3$ and $4$.
The quartic vertex $\ell_4^{\mathrm{tr}}$ is governed by the
quartic resonance class
$\mathfrak{Q} \in H^2(\cA^{\mathrm{sh}}_{4,0})$
of Appendix~\textup{\ref*{V1-app:nonlinear-modular-shadows}}.

The first quartic contribution appears at genus~$1$:
\begin{equation}% label removed: eq:thqg-X-quartic-genus-1
\delta Z_1^{(4)}(\cA;\,\tau)
\;=\;
\mathfrak{Q}(\cA)\,
\int_{E_\tau}
P(z;\tau)^2\,dz,
\end{equation}
where $P(z;\tau) = G_\tau(z)$ is the genus-$1$ propagator.
For the $\beta\gamma$ system,
$\mathfrak{Q}_{\beta\gamma} = 0$
\textup{(}Corollary~\textup{\ref*{V1-cor:nms-betagamma-mu-vanishing}}\textup{)},
so the quartic contact invariant vanishes and the first
genuine quartic correction appears at genus~$2$.
\end{proposition}

\begin{proof}
The quartic vertex contributes Feynman graphs with at least one
$4$-valent vertex.  At genus~$1$, the simplest such graph is a
single $4$-valent vertex with two self-loops, contributing
$\delta Z_1^{(4)}$.  The quartic resonance class $\mathfrak{Q}$
is the weight-$4$ component of the shadow algebra $\cA^{\mathrm{sh}}$
(Corollary~\ref*{V1-cor:shadow-extraction}).

For the $\beta\gamma$ system: the quartic contact invariant
$\mu_{\beta\gamma} = \langle \eta, m_3(\eta,\eta,\eta)\rangle = 0$
by Corollary~\ref*{V1-cor:nms-betagamma-mu-vanishing} (rank-one abelian
rigidity on the weight-changing line).
\end{proof}

\subsubsection{Class M: the infinite shadow obstruction tower}

\begin{remark}[Class-M non-Fredholm corrections]%
% label removed: rem:thqg-X-class-M-corrections%
\index{class M!non-Fredholm corrections}%
For class-M algebras ($r_{\max} = \infty$, e.g.\ $\mathrm{Vir}_c$
and $\cW_N$ for $N \geq 3$), the Feynman expansion involves
vertices of arbitrarily high arity.  The genus-$g$ non-Fredholm
correction is an infinite series:
\[
\delta Z_g^{\mathrm{M}}(\cA;\,\Omega)
= \sum_{r=3}^{\infty}
\sum_{\Gamma \in \mathcal{G}^{\mathrm{conn}}_{g,r}}
\frac{1}{|\Aut(\Gamma)|}
\int_{\Sigma_g^{v(\Gamma)}}
\omega_\Gamma(\Omega).
\]
Convergence of this series is guaranteed by the HS-sewing
condition (Theorem~\ref{thm:general-hs-sewing}), which provides
uniform bounds on the vertex labels $\ell_k^{\mathrm{tr}}$, and
the arity-cutoff lemma (Lemma~\ref{lem:arity-cutoff}), which
ensures that the contribution of high-arity graphs is
exponentially suppressed.

For the Virasoro algebra at genus~$1$, the non-Fredholm
correction involves the quartic contact invariant
$\mathfrak{Q}_{\mathrm{Vir}} = 10/[c(5c+22)]$
and all higher obstruction classes $o_r$ for $r \geq 5$
(the quintic obstruction $o_5$ is forced to be nonzero
by Theorem~\ref{thm:quintic-forced}).  The full partition
function $Z_1(\mathrm{Vir}_c;\,\tau)$ at genus~$1$ is:
\[
Z_1(\mathrm{Vir}_c;\,\tau)
= \dedeta(\tau)^{-(c-1)}
= q^{-(c-1)/24}\prod_{n=2}^{\infty}(1-q^n)^{-1},
\]
the character of the Virasoro vacuum module, which is
\emph{not} a pure Fredholm determinant on a single Bergman space
(it is missing the $n=1$ factor from the product).
\end{remark}

\begin{proposition}[Virasoro: Fredholm base + corrections;
\ClaimStatusProvedHere]%
% label removed: prop:thqg-X-virasoro-decomposition%
\index{Virasoro!Fredholm decomposition}%
The genus-$1$ Virasoro partition function admits the decomposition
\begin{equation}% label removed: eq:thqg-X-virasoro-decomposition
Z_1(\mathrm{Vir}_c;\,\tau)
\;=\;
\underbrace{|\dedeta(\tau)|^{-2c}}_{\text{Gaussian base}
(\kappa = c/2)}
\;\cdot\;
\underbrace{|\dedeta(\tau)|^{2}
\cdot (1-q)}_{\text{null-state correction}},
\end{equation}
where the Gaussian base is the Fredholm determinant at
$\kappa = c/2$ and the correction factor removes the
contribution of the $L_{-1}$ null state.
\end{proposition}

\begin{proof}
The Virasoro character at genus~$1$ is
$\chi_c(\tau) = q^{(1-c)/24}/(1-q)\prod_{n=2}^{\infty}(1-q^n)^{-1}
= q^{-c/24}\,\dedeta(\tau)^{-1}\cdot(1-q)^{-1}\cdot\dedeta(\tau)/q^{-1/24}$\,.
More directly:
\begin{align}
\chi_c(\tau)
&= q^{(1-c)/24} \prod_{n=2}^{\infty}(1-q^n)^{-1}
\notag\\
&= q^{(1-c)/24}
\cdot \frac{\prod_{n=1}^{\infty}(1-q^n)^{-1}}{(1-q)^{-1}}
\notag\\
&= q^{(1-c)/24}
\cdot (1-q) \cdot \prod_{n=1}^{\infty}(1-q^n)^{-1}
\notag\\
&= (1-q)\cdot q^{-c/24}\cdot q^{1/24}
\cdot \prod_{n=1}^{\infty}(1-q^n)^{-1}
\notag\\
&= (1-q)\cdot q^{-c/24}\cdot \dedeta(\tau)^{-1}.
% label removed: eq:thqg-X-virasoro-factorization
\end{align}
The factor $(1-q)$ removes the $L_{-1}$ descendant (which is
null for the vacuum module), and $\dedeta(\tau)^{-1}$ is the
oscillator contribution $\prod_{n=1}^{\infty}(1-q^n)^{-1}$
up to the $q^{1/24}$ prefactor.  The Gaussian base
$|\dedeta|^{-2c}$ arises from $c$ free-field copies at
$\kappa = 1$ each, giving $\kappa_{\mathrm{total}} = c/2$.
\end{proof}

\begin{remark}[Resurgent structure for class-M algebras]%
% label removed: rem:thqg-X-resurgent%
\index{resurgence!class M}%
For class-M algebras, the Feynman expansion
$\delta Z_g^{\mathrm{M}}$ is expected to be resurgent:
the asymptotic (divergent) perturbative series in the shadow
tower is Borel summable, with non-perturbative corrections
controlled by the resonance rank $\rho(\cA)$
(Definition~\ref{def:resonance-rank}).  The Stokes phenomenon
as $\tau$ crosses walls of marginal stability corresponds to
the wall-crossing of the MC moduli of
$\Theta_\cA \in \MC(\gAmod)$.  This resurgent structure is
conjectural beyond the Virasoro case, where the Zamolodchikov
recursion provides an independent control.
\end{remark}

% =====================================================================
%
%  SECTION 6: ANALYTIC ASPECTS
%
% =====================================================================

\subsection{Analytic aspects: the analytic bar coalgebra,
coderived categories, and open problems}
% label removed: subsec:thqg-X-analytic-aspects

\subsubsection{The analytic bar coalgebra}

\begin{definition}[Analytic bar coalgebra; review]%
% label removed: def:thqg-X-analytic-bar-review%
\index{bar coalgebra!analytic!review}%
Recall from Definition~\ref{def:analytic-bar-coalgebra} that the
analytic bar coalgebra $\anbar(\cA)$ is defined by closing
the algebraic bar construction inside the sewing envelope:
\[
\anbar_n(\cA)
\;=\;
\overline{
\algcore^{\otimes n}
\otimes \Omega^*_{\log}(C_n(X))
}^{\,\graphnorm(d_{\mathrm{bar}})},
\]
where the graph norm of $d_{\mathrm{bar}}$ is the norm
$\|x\|_{\mathrm{graph}}^2 := \|x\|^2 + \|d_{\mathrm{bar}} x\|^2$.
This is the analytic closure of the algebraic bar construction
(Definition~\ref*{V1-def:bar-differential-complete}).
\end{definition}

\begin{proposition}[Analytic bar differential is bounded;
\ClaimStatusProvedHere]%
% label removed: prop:thqg-X-analytic-bar-bounded%
\index{bar coalgebra!analytic!boundedness}%
Under HS-sewing \textup{(}Definition~\textup{\ref{def:hs-sewing}}\textup{)},
the bar differential $d_{\mathrm{bar}} \colon \anbar_n(\cA) \to
\anbar_{n-1}(\cA)$ is a bounded operator for the graph-norm
topology, with operator norm
\begin{equation}% label removed: eq:thqg-X-bar-bounded
\|d_{\mathrm{bar}}\|_{\mathrm{op}}
\;\leq\;
\binom{n}{2}\,K\,\sup_q \|m_q\|_{\hscl},
\end{equation}
where $K$ is the polynomial OPE growth constant and $m_q$ is the
$q$-weighted multiplication.
\end{proposition}

\begin{proof}
The bar differential $d_{\mathrm{bar}} = \sum_{i<j}
\operatorname{Res}_{z_i=z_j} \circ Y_{ij}$ has
$\binom{n}{2}$ summands, each of which is a composition of
the OPE operator $Y_{ij}$ (bounded by $K$ on the sewing envelope)
and the residue operator (norm $1$ on the graph-norm closure).
The HS-sewing condition ensures $\|m_q\|_{\hscl} < \infty$,
which controls the operator norm of $Y_{ij}$ on the weighted
completion $\sewenv_q$.
\end{proof}

\begin{proposition}[Analytic coproduct;
\ClaimStatusProvedHere]%
% label removed: prop:thqg-X-analytic-coproduct%
\index{bar coalgebra!analytic!coproduct}%
The analytic bar coalgebra $\anbar(\cA)$ carries a continuous
coproduct
\begin{equation}% label removed: eq:thqg-X-analytic-coproduct
\bar{\Delta}^{\mathrm{an}} \colon
\anbar(\cA)
\;\to\;
\anbar(\cA)
\;\widehat{\otimes}\;
\anbar(\cA),
\end{equation}
where $\widehat{\otimes}$ is the completed tensor product
in the graph-norm topology.  The coproduct extends the algebraic
coproduct $\bar{\Delta}$ of the algebraic bar construction
by continuity.
\end{proposition}

\begin{proof}
The algebraic coproduct $\bar{\Delta}$ on $\barB_n(\cA)$
sends $\barB_n \to \bigoplus_{p+q=n} \barB_p \otimes \barB_q$
by deshuffle.  Each deshuffle component is a restriction map
(restricting meromorphic forms from $C_n$ to
$C_p \times C_q$), which is continuous for the graph-norm
topology.  The completed coproduct is the closure of the
algebraic coproduct.
\end{proof}

\begin{remark}[The bar Laplacian and Hodge decomposition]%
\label{rem:bar-laplacian-hodge}%
\index{bar Laplacian|textbf}%
\index{Hodge decomposition!bar complex}
For a unitary VOA~$\cA$ satisfying~(AU1), the formal adjoint
$d_B^*$ is densely defined on~$B^{\mathrm{an}}(\cA)$
(it is the BPZ-conjugate creation operator, reversing the
collision contraction).  The \emph{bar Laplacian}
$\Delta_B := d_B\,d_B^* + d_B^*\,d_B$
is essentially self-adjoint on the algebraic bar complex, because
each weight space $B_n(\cA)_{\mathrm{weight}=N}$ is
finite-dimensional (a basic VOA axiom), and a symmetric operator
on a finite-dimensional space is automatically self-adjoint.
The full operator
$\Delta_B = \bigoplus_{N,n} \Delta_B|_{N,n}$
is an orthogonal direct sum of finite-dimensional self-adjoint
operators, hence essentially self-adjoint on the algebraic core.

Consequences: (i)~discrete spectrum with finite multiplicities,
accumulating only at~$+\infty$;
(ii)~the bar Hodge theorem
$\ker(\Delta_B) = H^*(B(\cA))$
(harmonic forms represent bar cohomology);
(iii)~for chirally Koszul~$\cA$, bar cohomology concentrates on
the diagonal (weight $=$ bar degree), giving a positive spectral
gap off the diagonal.  Compare the screening Hodge theorem
(Theorem~\ref{thm:screening-hodge-theorem} in
\S\ref{sec:ym-spectral-gap}) for the analogous statement in the
Yang--Mills setting.
\end{remark}

\subsubsection{Coderived categories and genus $g \geq 1$}

\begin{remark}[Curvature forces coderived passage; review]%
% label removed: rem:thqg-X-coderived-review%
\index{coderived category!Fredholm context}%
At genus~$g \geq 1$, the fiberwise bar differential satisfies
$\dfib^{\,2} = \kappa(\cA) \cdot \omega_g \neq 0$
(the curvature of the Hodge connection;
Theorem~\ref{thm:genus-graded-bar}).  This means that the
bar complex at genus~$g$ is \emph{curved}: $d^2 \neq 0$,
and ordinary cohomology is not defined.  The correct
receptacle is the \emph{coderived category}
$\codercat(\anbar_g(\cA)\text{-comod})$
(Definition~\ref{def:analytic-shadow} and
Remark~\ref{rem:curvature-coderived-analytic}), not the
ordinary derived category.

For the Fredholm partition functions of this section, the
coderived passage manifests as follows: the genus-$g$
partition function $Z_g$ is the \emph{supertrace} of the
sewing operator on the coderived category of the analytic bar
coalgebra, not an ordinary cohomological invariant.  The
Fredholm determinant formula
$Z_g = \det(1 - \sewker_g)^{-\kappa}$ computes this
supertrace directly via the one-particle reduction, bypassing
the coderived formalism for the Heisenberg and class-G cases.
For classes L, C, M, the coderived category is essential: the
non-Fredholm corrections $\delta Z_g$ are coderived
invariants that vanish in the ordinary derived category.
\end{remark}

\begin{definition}[Coderived shadow partition function]%
% label removed: def:thqg-X-coderived-shadow-pf%
\index{partition function!coderived shadow}%
For a chiral algebra $\cA$ satisfying HS-sewing, define the
\emph{coderived shadow partition function} at genus~$g$ as
\begin{equation}% label removed: eq:thqg-X-coderived-shadow-pf
Z_g^{\mathrm{co}}(\cA;\,\tau)
\;=\;
\operatorname{Str}_{Q_g^{\mathrm{an}}(\cA)}
\!\bigl(\sigma_g\,q^{L_0 - c/24}\bigr),
\end{equation}
where $Q_g^{\mathrm{an}}(\cA)
= \codercat(\anbar_g(\cA)\text{-comod})$ is the
genus-$g$ analytic shadow
\textup{(}Definition~\textup{\ref{def:analytic-shadow}}\textup{)},
$\sigma_g$ is the complementarity involution, and
$\operatorname{Str}$ is the supertrace.
\end{definition}

\begin{proposition}[Coderived = Fredholm for class~G;
\ClaimStatusProvedHere]%
% label removed: prop:thqg-X-coderived-fredholm-G%
\index{coderived category!Fredholm for class G}%
For class-G algebras, the coderived shadow partition function
coincides with the Fredholm partition function:
\begin{equation}% label removed: eq:thqg-X-coderived-fredholm-agreement
Z_g^{\mathrm{co}}(\cA;\,\tau)
\;=\;
Z_g(\cA;\,\Omega)
\;=\;
\det(1 - \sewker_g)^{-\kappa(\cA)}.
\end{equation}
\end{proposition}

\begin{proof}
For class-G algebras, the curvature is scalar
($\dfib^2 = \kappa \cdot \omega_g$), and the coderived category
reduces to the ordinary derived category twisted by the Hodge
line bundle $\hodgebdle^{\kappa/2}$.  The supertrace on
the twisted derived category equals the Fredholm determinant
of the sewing kernel, by the Heisenberg sewing theorem
(Theorem~\ref{thm:thqg-X-heisenberg-sewing-full}).
\end{proof}

\subsubsection{The analytic Koszul pair}

\begin{remark}[Analytic Koszul duality]%
% label removed: rem:thqg-X-analytic-koszul%
\index{Koszul pair!analytic}%
An analytic Koszul pair $(\cA, \cA^!)$ consists of an
algebraic Koszul pair equipped with compatible sewing
envelopes and analytic bar coalgebras
(Definition~\ref{def:analytic-koszul-pair}).  For the
Heisenberg algebra, the analytic Koszul pair is
$(\cH_\kappa, \cH_\kappa^!)$ where
$\cH_\kappa^! = \Sym^{\mathrm{ch}}(V^*)$ with
$\kappa(\cH_\kappa^!) = -\kappa$
(note: $\cH_\kappa^! \neq \cH_{-\kappa}$; the Koszul dual
is a commutative chiral algebra, not a Heisenberg algebra).
The analytic bar coalgebra of $\cH_\kappa^!$
is $\anbar(\cH_\kappa^!) \cong
\overline{\operatorname{Sym} A^2(D)}^{\mathsf{co}}$, the
cofree coalgebra completion of the Bergman Fock space.

The genus-$g$ complementarity decomposition
$Z_g(\cH_\kappa) + Z_g(\cH_\kappa^!)
= (\det\operatorname{Im}\Omega)^{-\kappa/2}
\cdot (\zetareg\Delta)^{-\kappa/2}
+ (\det\operatorname{Im}\Omega)^{\kappa/2}
\cdot (\zetareg\Delta)^{\kappa/2}$
is a direct Fredholm manifestation of Theorem~C
(complementarity): what $\cH_\kappa$ sees as the positive
curvature $+\kappa \cdot \omega_g$, its Koszul dual $\cH_\kappa^!$
sees as negative curvature $-\kappa \cdot \omega_g$.
\end{remark}

\subsubsection{The platonic chain and the sewing hierarchy}

\begin{remark}[The platonic chain for Fredholm partition functions]%
% label removed: rem:thqg-X-platonic-chain%
\index{platonic ideal form!Fredholm partition functions}%
The platonic chain of Remark~\ref{rem:platonic-ideal-form}
has the following concrete incarnation for Fredholm partition
functions:
\[
\underbrace{\algcore}_{\text{algebraic OPE data}}
\;\subset\;
\underbrace{\sewenv}_{\text{sewing envelope}}
\;\rightsquigarrow\;
\underbrace{\operatorname{Sym} A^2(D)}_{\text{Bergman Fock space}}
\;\rightsquigarrow\;
\underbrace{Z_g = \det(1 - \sewker_g)^{-\kappa}}_{\text{Fredholm
partition function}}.
\]
The first arrow (dense embedding) is the mode--Bergman
correspondence; the second arrow (Hilbert completion) uses
the sewing envelope topology; the third arrow (Fredholm
determinant) uses the one-particle reduction and the
Polyakov formula.
Each stage is necessary:
\begin{itemize}
\item $\algcore$ alone gives only formal power series in~$q$
  (the characters, which converge for $|q| < 1$ but have no
  intrinsic analytic continuation).
\item $\sewenv$ promotes the formal power series to operators
  on a Hilbert space, with explicit spectral data.
\item The Fredholm determinant provides the non-perturbative
  completion: it defines $Z_g$ as a function on $\Siegel_g$
  (the Siegel upper half-space), not merely as a formal
  $q$-series.
\end{itemize}
\end{remark}

\subsubsection{Open problems and the analytic frontier}

\begin{openproblem}[Lattice sewing envelope]%
% label removed: prob:thqg-X-lattice-sewing%
\index{lattice algebra!sewing envelope}%
Construct the sewing envelope of a lattice vertex algebra
$V_\Lambda$ explicitly as a completion of the Bergman Fock
space tensored with the lattice group algebra.
Conjecture~\textup{\ref{conj:lattice-sewing}} predicts that
\begin{equation}% label removed: eq:thqg-X-lattice-sewing-conj
V_\Lambda^{\mathrm{sew}}
\;\cong\;
\operatorname{Sym} A^2(D)^{\otimes r}
\;\otimes\;
\ell^2(\Lambda^*/\Lambda),
\end{equation}
where $r = \mathrm{rank}(\Lambda)$ and
$\ell^2(\Lambda^*/\Lambda)$ is the Hilbert space of the
charge lattice.
\end{openproblem}

%% ---- The lattice Steinberg ----

\subsubsection{The lattice Steinberg}
% label removed: sec:lattice-steinberg
\index{Steinberg object!lattice}%
\index{lattice algebra!Steinberg object}%

\providecommand{\Steinberg}{\mathfrak{S}}

\begin{construction}[Lattice Steinberg object]%
% label removed: con:lattice-steinberg%
For a lattice vertex algebra $V_\Lambda$ with $\Lambda$ an
even lattice of rank~$r$, define the \emph{lattice Steinberg
object}
\begin{equation}% label removed: eq:lattice-steinberg-object
\Steinberg_\Lambda
\;=\;
\bigl\{(x_i, \lambda_i)
  \in \FM_k(\bC) \times \Lambda^k
  : \textstyle\sum_i \lambda_i = 0\bigr\}
\;\times\;
\Conf_k^{<}(\bR).
\end{equation}
The constraint $\sum_i \lambda_i = 0$ is charge conservation:
the total lattice momentum flowing through a bar element of
degree~$k$ vanishes.  The Borel--Moore homology decomposes by
coset of the discriminant group:
\begin{equation}% label removed: eq:lattice-steinberg-BM
H_*^{\mathrm{BM}}(\Steinberg_\Lambda)
\;=\;
\bigoplus_{[\mu] \in \Lambda^*/\Lambda}
H_*^{\mathrm{BM}}\!\bigl(\Steinberg_\Lambda^{[\mu]}\bigr),
\end{equation}
where $\Steinberg_\Lambda^{[\mu]}$ restricts the lattice
labels to the coset~$[\mu]$.
\end{construction}

\begin{conjecture}[Conditional on
Conjecture~\textup{\ref{conj:lattice-sewing}};
\ClaimStatusConjectured]%
% label removed: prop:lattice-steinberg-sewing%
Assume the lattice sewing envelope conjecture
\textup{(}Conjecture~\textup{\ref{conj:lattice-sewing})}.
Then the sewing envelope
$\operatorname{Sym} A^2(D)^{\otimes r}
\otimes \ell^2(\Lambda^*/\Lambda)$
is the $L^2$-completion of
$H_*^{\mathrm{BM}}(\Steinberg_\Lambda)$.
The factor $\ell^2(\Lambda^*/\Lambda)$ is the space of
functions on the component group
$\pi_0(\Steinberg_\Lambda) = \Lambda^*/\Lambda$
(the discriminant form).
\end{conjecture}

\begin{proof}
The FM$\times$Conf product decomposition of
$\Steinberg_\Lambda$ gives a direct-sum splitting indexed by
$\Lambda^*/\Lambda$.  On each component the continuous part is
$\FM_k(\bC) \times \Conf_k^{<}(\bR)$ with the charge-zero
sublattice fibre; its BM homology is the $k$-th bar piece of
$\operatorname{Sym} A^2(D)^{\otimes r}$ by the Bergman--Fock
identification of~\S\ref{sec:bergman-fock}.  Summing over
components produces the tensor factor
$\ell^2(\Lambda^*/\Lambda)$, and the $L^2$-completion follows
from the sewing-envelope topology
of~\eqref{eq:thqg-X-lattice-sewing-conj}.
\end{proof}

\begin{remark}[Non-unimodular extension]%
% label removed: rem:lattice-steinberg-non-unimodular%
When $\Lambda$ is not unimodular, $\Lambda^*/\Lambda$ is
nontrivial and contributes additional Stokes data.  The
discriminant form carries a quadratic refinement
$q \colon \Lambda^*/\Lambda \to \bQ/\bZ$ that controls the
phases in the modular transformation; this is the additional
data needed for the non-unimodular sewing.
\end{remark}

\begin{openproblem}[Boundary bar duality]%
% label removed: prob:thqg-X-boundary-bar%
\index{bar coalgebra!boundary duality}%
Establish a Fredholm-level analytic Koszul duality:
for every analytic Koszul pair $(\cA, \cA^!)$ satisfying
HS-sewing, construct a Fredholm equivalence
\begin{equation}% label removed: eq:thqg-X-boundary-bar-conj
\codercat(\anbar(\cA)\text{-comod})
\;\simeq\;
\contrcat(\cA^!\text{-mod})
\end{equation}
between the coderived category of analytic bar comodules and
the contraderived category of $\cA^!$-modules.
This would be the analytic avatar of bar-cobar inversion
(Theorem~B) at genus~$g \geq 1$.
\end{openproblem}

\begin{openproblem}[Non-Fredholm corrections: closed-form expressions]%
% label removed: prob:thqg-X-non-fredholm-closed%
\index{non-Fredholm correction!closed-form}%
For class-L algebras (e.g.\ $V_k(\fg)$), express the
genus-$g$ non-Fredholm correction $\delta Z_g^{\mathrm{L}}$
in terms of Siegel modular forms and conformal blocks.
At integrable level, the Verlinde formula
\eqref{eq:verlinde-general} provides an independent
computation; the problem is to derive the Verlinde formula
directly from the Feynman integral expansion
\eqref{eq:thqg-X-class-L-correction}.
\end{openproblem}

\begin{openproblem}[Analytic completion for W-algebras]%
% label removed: prob:thqg-X-W-algebra-analytic%
\index{W-algebra!analytic completion}%
Construct the sewing envelope and analytic bar coalgebra for
class-M algebras ($\mathrm{Vir}_c$, $\cW_N$) in the
resonance-filtered topology of
Theorem~\ref{thm:resonance-filtered-bar-cobar}.
The resonance rank $\rho(\cA)$
(Definition~\ref{def:resonance-rank}) controls the
difficulty: for $\rho < \infty$ (the resonance completion
theorem, Theorem~\ref{thm:platonic-completion}), the
analytic completion reduces to a finite resonance problem;
for $\rho = \infty$ (if such algebras exist among modular
Koszul chiral algebras), the completion requires a genuinely
infinite-dimensional analytic framework beyond the Fredholm
theory of this section.
\end{openproblem}

\begin{remark}[Summary: the Fredholm partition function as
a projection of $\Theta_\cA$]%
% label removed: rem:thqg-X-fredholm-summary%
\index{Fredholm determinant!as MC projection}%
The genus-$g$ Fredholm partition function
$Z_g(\cA;\,\Omega) = \det(1 - \sewker_g)^{-\kappa(\cA)}$
is the analytic realization of the scalar genus tower
$F_g(\cA) = \kappa(\cA) \cdot \lambda_g^{\mathrm{FP}}$
of Theorem~D.  It is the bottom projection of the single
universal MC element
$\Theta_\cA \in \MC(\gAmod)$
through the analytic completion:
\[
\Theta_\cA
\;\xrightarrow{\;\mathrm{arity}\;2\;}
\kappa(\cA)
\;\xrightarrow{\;\mathrm{sewing}\;}
\sewker_g
\;\xrightarrow{\;\det\;}
Z_g.
\]
For class-G algebras, this projection is lossless: $Z_g$
captures all the information in~$\Theta_\cA$.  For classes
L, C, M, the non-Fredholm corrections $\delta Z_g$ carry the
information from the higher-arity projections of~$\Theta_\cA$
(the cubic shadow $\fC$, the quartic resonance class
$\mathfrak{Q}$, and all higher obstruction classes $o_r$).
The full partition function
$Z_g^{\mathrm{full}} = Z_g^{\mathrm{G}} + \delta Z_g$
is the analytic realization of the complete MC element, and
the Feynman expansion \eqref{eq:thqg-X-feynman-expansion}
is the graph-sum formula of the shadow obstruction tower
evaluated on the analytic sewing data.

This completes the development of Result~\textbf{(G10)}: the
genus-$g$ partition function of a modular Koszul chiral algebra,
computed as a nonperturbative Fredholm determinant on the
one-particle Bergman space, with explicit Feynman corrections
for non-Gaussian shadow obstruction towers.
\end{remark}

%% ----------------------------------------------------------------
%% Analytic Borel--Moore homology and Fredholm Koszul duality
%% ----------------------------------------------------------------

\subsection{Analytic Borel--Moore homology and Fredholm Koszul duality}
% label removed: subsec:analytic-bm-steinberg

\providecommand{\fSb}{\mathfrak{S}_b}

The Steinberg variety $\fSb$ attached to a modular Koszul chiral
algebra $\cA$ parametrises composable pairs of bar elements:
a point of~$\fSb$ is a pair $(b_1, b_2)$ in
$\anbar(\cA) \times_{\cA} \anbar(\cA)$ meeting along a common
$\cA$-output.  The convolution product on~$\fSb$ encodes
the bar differential in the holomorphic direction and the
deconcatenation coproduct in the topological direction,
so that the resulting algebra is the Fredholm-level avatar
of the Swiss-cheese operad action.

\begin{construction}[Analytic Borel--Moore homology of $\fSb$]%
% label removed: constr:analytic-bm-steinberg%
\index{Steinberg variety!analytic BM homology}%
\index{Borel--Moore homology!analytic}%
Complete the Steinberg object $\fSb$ in the analytic topology
appropriate to the shadow class of~$\cA$:
\begin{itemize}
\item \emph{Class-F} (free field): the nuclear Fr\'echet topology
  on the Bergman one-particle space, so that the FM strata
  of~$\fSb$ carry the topology of nuclear Fr\'echet--Stein
  algebras.
\item \emph{Class-M} (W-algebras, resonance): the
  resonance-filtered topology of
  Theorem~\ref{thm:resonance-filtered-bar-cobar}, in which
  $\fSb$ is exhausted by Banach-analytic subvarieties of
  finite resonance rank~$\rho(\cA)$.
\end{itemize}
In either case, define the \emph{analytic Borel--Moore homology}
$H_*^{\mathrm{BM,an}}(\fSb)$ as the homology of the complex of
locally finite analytic chains on~$\fSb$ with coefficients in
the dualising sheaf.  The convolution product
\[
H_i^{\mathrm{BM,an}}(\fSb)
\otimes
H_j^{\mathrm{BM,an}}(\fSb)
\;\xrightarrow{\;\mu\;}
H_{i+j}^{\mathrm{BM,an}}(\fSb)
\]
makes $H_*^{\mathrm{BM,an}}(\fSb)$ into an associative graded
algebra, the \emph{Fredholm-level convolution algebra} of~$\cA$.
\end{construction}

\begin{conjecture}[\ClaimStatusConjectured]%
% label removed: prop:analytic-steinberg-koszul%
\index{Koszul duality!Fredholm-level Steinberg}%
\index{coderived category}%
\index{contraderived category}%
For a modular Koszul chiral algebra $\cA$ satisfying HS-sewing,
there is an equivalence of triangulated categories
\begin{equation}% label removed: eq:analytic-steinberg-koszul
\codercat\bigl(H_*^{\mathrm{BM,an}}(\fSb)\textup{-comod}\bigr)
\;\simeq\;
\contrcat(\cA^!\textup{-mod}).
\end{equation}
This is the analytic enhancement of the algebraic Steinberg
presentation (Theorem~\ref{thm:steinberg-presentation}):
$H_*^{\mathrm{BM,an}}(\fSb)$ presents the Fredholm-level
bar coalgebra, and
\eqref{eq:analytic-steinberg-koszul} is the analytic lift
of bar-cobar inversion (Theorem~B).
The coderived/contraderived asymmetry reflects the choice
of completion topology on~$\fSb$: the coderived side
carries the direct-limit (nuclear) topology of bar comodules,
while the contraderived side carries the inverse-limit
(pro-finite-dimensional) topology of Koszul dual modules.
\end{conjecture}

\begin{remark}[Non-Fredholm corrections from strata at infinity]%
% label removed: rem:steinberg-non-fredholm%
\index{non-Fredholm correction!strata at infinity}%
\index{Siegel modular form}%
For class-L algebras, the non-Fredholm correction
$\delta Z_g$ of
\S\ref{subsec:thqg-non-fredholm} receives a Steinberg
interpretation: it is the contribution from the
\emph{non-compact strata} of~$\fSb$, the strata at
infinity in the FM compactification
$\overline{\mathrm{FM}}_k(\Sigma_g)$.  The genus-$g$
Steinberg variety carries a natural fibration
\[
\fSb^{(g)} \;\longrightarrow\; \mathbb{H}_g
\]
over the Siegel upper half-space $\mathbb{H}_g$, where
the fibre over $\Omega \in \mathbb{H}_g$ is the Steinberg
variety of the sewing kernel
$\sewker_g(\Omega)$.  The compact strata produce the
Fredholm determinant $Z_g^{\mathrm{G}}$; the non-compact
strata produce
\[
\delta Z_g^{\mathrm{L}}
\;=\;
\int_{\mathbb{H}_g \smallsetminus \mathcal{F}_g}
\operatorname{ch}\bigl(
  H_*^{\mathrm{BM,an}}(\fSb^{(g)}_\Omega)
\bigr)
\,d\mu_{\mathrm{Siegel}},
\]
where $\mathcal{F}_g \subset \mathbb{H}_g$ is a Siegel
fundamental domain and $\operatorname{ch}$ denotes the
Chern character of the convolution algebra fibre.  At
integrable level, this integral localises to the
Verlinde formula via the finite-dimensionality of the
conformal-block bundle.
\end{remark}

% ============================================================
\subsection{Prime-locality and quantum gravity}%
\label{subsec:prime-locality-quantum-gravity}%
\index{prime-locality!quantum gravity interface|textbf}

The prime-locality programme
(Vol~I, Conjecture~conj:prime-locality-transfer and
\S\ref{sec:prime-locality-assessment})
asks whether the shadow amplitudes at genus~$1$ decompose
into Hecke eigencomponents with multiplicative Fourier
coefficients.  The question has a precise translation into the
$3$d quantum gravity framework of this volume: prime-locality
is the question of whether the holomorphic-topological partition
function, viewed on the genus-$1$ moduli space, respects the
arithmetic structure of the Hecke correspondences.

\begin{remark}[What the $3$d perspective adds to prime-locality]%
\label{rem:3d-prime-locality}%
\index{prime-locality!3d perspective}
The $3$d holomorphic-topological framework provides three
structural inputs that bear on the prime-locality question:

\textbf{Input~1: the sewing-Hecke collapse is a genus-$1$
phenomenon.}
The sewing operator $\sewop_q$ at genus~$1$ is diagonal in the
Fock basis: $\sewop_q|n\rangle = q^n|n\rangle$.  This
diagonality forces the sewing--Hecke reciprocity
(Vol~I, Theorem~thm:sewing-hecke-reciprocity), which collapses
the two-variable spectral data into a single-variable Dirichlet
series.  At genus~$\ge 2$, the off-diagonal entries $R_{ij}$
in the multi-handle sewing kernel
(\S\ref{subsec:thqg-genus2-fredholm}) break this diagonality.
Consequently, the arithmetic content at genus~$\ge 2$ is
genuinely multi-variable: the Rankin--Selberg unfolding on
$\mathrm{Sp}_{2g}(\mathbb{Z}) \backslash \mathfrak{H}_g$
produces Siegel modular $L$-functions that do not factor
as products of elliptic $L$-functions.

\textbf{Consequence:} the prime-locality question is tied to
the genus-$1$ sewing operator.  At higher genus, the relevant
arithmetic structure is not prime-locality but
\emph{Siegel--Hecke equivariance}, a genuinely different (and
harder) problem.  The genus-$2$ B\"ocherer bridge
(Vol~I, genus-$2$ arithmetic frontier) is the first probe
of this regime.

\textbf{Input~2: the gravitational dichotomy constrains the
Hecke defect.}
At $c = 26$, the matter-ghost system has $\kappa_{\mathrm{eff}}
= \kappa(\mathrm{matter}) + \kappa(\mathrm{ghost}) = 0$,
and the transgression algebra degenerates from Clifford to
exterior (Vol~I, the critical string dichotomy).  At
$c \neq 26$, the transgression algebra is a matrix algebra
(Morita-trivial), and the Koszul duality pairing
$\cA \leftrightarrow \cA^!$ is non-degenerate.

For the Hecke defect class
$[\delta_p] \in H^1(\mathfrak{g}_\cA^{\mathrm{mod}},
\mathrm{ad}_{T_p})$,
the Morita triviality at $c \neq 26$ means that the
convolution algebra has ``fewer'' cohomological directions
for the defect to live in.  At $c = 26$, the degeneration
of the transgression algebra opens new cohomological
directions.  Whether this structural difference affects the
vanishing of $[\delta_p]$ is a concrete question that
connects prime-locality to the critical string dichotomy.

\textbf{Input~3: the planted-forest correction is
tau-dependent at genus~$\ge 2$.}
The planted-forest differential $d_{\mathrm{pf}}$ of
Vol~I involves codimension-$\ge 2$ strata of the log-FM
compactification.  At genus~$1$, the planted-forest
correction $\delta_{\mathrm{pf}}^{(1,r)}$ contributes to the
shadow amplitudes but does not introduce new tau-dependence
beyond $E_2^*(\tau)$.  At genus~$\ge 2$, the planted-forest
corrections involve Siegel modular functions, and the
Hecke equivariance question becomes the question of whether
the planted-forest graph sums on $\overline{\mathcal{M}}_{g,n}$
respect the Siegel--Hecke decomposition.
\end{remark}

\begin{remark}[The honest assessment from the $3$d perspective]%
\label{rem:3d-prime-locality-assessment}%
\index{prime-locality!honest assessment}
The $3$d holomorphic-topological framework does \emph{not}
resolve the prime-locality question.  It clarifies the
\emph{landscape}:
\begin{enumerate}[label=\textup{(\roman*)}]
\item Prime-locality is proved for lattice VOAs
  (all ranks, all levels) and for rational VOAs with diagonal
  modular invariant.  The proof uses the Hecke module structure
  of the theta function or the Franc--Mason eigendecomposition,
  not the $3$d framework.

\item Prime-locality at the shadow \emph{coefficient} level
  (tau-independent OPE data) is proved for all algebraic families
  by tau-independence alone.  This is a trivial consequence of
  the fact that the Hecke operator acts as the identity on
  constants.

\item The open question is prime-locality at the shadow
  \emph{amplitude} level for irrational VOAs: whether the
  genus-$1$ graph sums
  $\operatorname{Sh}_r^{(1)}(\tau) =
  \sum_\Gamma |\mathrm{Aut}(\Gamma)|^{-1}
  \prod_{v} S_{n_v} \prod_{e} E_2^*(\tau)$
  respect the Hecke decomposition of the ambient space of
  quasi-modular forms.  The vertex weights $S_{n_v}$ are
  Hecke-equivariant (tau-independent); the edge weights
  $E_2^*(\tau)$ are individually Hecke-equivariant (the
  Fourier coefficients $\sigma_1(n)$ are multiplicative); but
  the \emph{products} of $E_2^*$ are not automatically
  Hecke-equivariant, because $T_p$ is not a ring homomorphism
  on quasi-modular forms.

\item The $3$d framework contributes one genuine insight:
  the prime-locality question is \emph{intrinsically genus-$1$}.
  At higher genus, the relevant arithmetic object is not a
  Hecke eigendecomposition on $\mathcal{M}_{1,1}$ but a
  Siegel--Hecke structure on $\mathcal{M}_g$, and the
  sewing-kernel non-diagonality opens different
  (and potentially richer) arithmetic channels.
\end{enumerate}

The Riccati algebraicity theorem
(Vol~I, Theorem~thm:riccati-determinacy-assessment)
provides a well-defined ``algebraic prime-locality'': the
shadow obstruction tower is finitely determined by three OPE invariants
$(\kappa, \alpha, S_4)$.  This is a strong constraint on the
algebraic structure but does not address the arithmetic question
of Hecke eigenvalues.  The term ``prime-locality'' should be
reserved for the arithmetic sense; the algebraic determinacy
is better called \emph{Riccati determinacy} or
\emph{shadow-tower finite determination}.
\end{remark}

\begin{remark}[Siegel modular forms and the genus-$2$ HT partition function]%
\label{rem:thqg-siegel-genus2}%
\index{Siegel modular form!3d HT partition function}%
\index{Fredholm determinant!genus-2}%
The genus-$2$ Fredholm determinant
$\freddet(1 - K_2)$ of
Proposition~\textup{thm:genus2-non-diagonal}
\textup{(}Vol~I\textup{)} produces a function on
$\Siegel_2 = \Sp(4,\bZ) \backslash \mathfrak{H}_2$
whose Fourier expansion connects to the Igusa ring
$M_*(\Sp(4,\bZ)) = \bC[E_4, E_6, \chi_{10}, \chi_{12}]$
of Siegel modular forms.
For lattice VOAs of rank~$d$, the genus-$2$ partition
function is the Siegel theta series
$\Theta_\Lambda^{(2)}(\Omega)$, a Siegel modular form of
weight~$d/2$.  The three-dimensional HT perspective adds
a structural interpretation: the off-diagonal coupling
$R_{12}$ in the Schur complement factorization is the
genus-$2$ manifestation of the bar complex coproduct
\textup{(}the $\bR$-direction factorization of the
Swiss-cheese structure\textup{)}, and the cusp-form
component of $\Theta_\Lambda^{(2)}$ encodes
genuinely genus-$2$ bulk--boundary correlations that
are invisible to the genus-$1$ shadow obstruction tower.

For the Leech lattice \textup{(}$d = 24$\textup{)}:
the cusp projection onto the Saito--Kurokawa lift
$\chi_{12} = \mathrm{SK}(f_{22})$ is nonzero
\textup{(}Vol~I, Proposition~\textup{prop:leech-cusp-nonvanishing}\textup{)},
and the B\"ocherer coefficient provides algebraic access
to the central $L$-value $L(1/2, \pi_{f_{22}} \times \chi_D)$
\textup{(}Vol~I, Theorem~\textup{thm:bocherer-bridge}\textup{)}.
In the 3d HT framework, this is the genus-$2$ version of
the holographic modular Koszul datum:
the genus-$2$ MC equation determines the Siegel theta
series from genus-$1$ shadow data and boundary
contributions, and the B\"ocherer factorization converts
this algebraic determination into critical-line access.
\end{remark}

\begin{remark}[MZV periods and the genus-$0$ bar complex
in 3d HT]%
\label{rem:thqg-mzv-bar}%
\index{multiple zeta values!3d HT}%
\index{bar complex!MZV periods in 3d HT}%
The genus-$0$ bar amplitudes of the holomorphic colour
produce multiple zeta value (MZV) periods via iterated
integrals on $\cM_{0,n}$
\textup{(}Vol~I, Theorem~\textup{thm:shadow-mzv-dictionary}\textup{)}.
In the 3d HT interpretation: the MZV periods are the
perturbative contributions to the holomorphic Wilson line,
computed by the bar complex on $\overline{\mathrm{FM}}_n(\bC)$.
The shadow depth classification $\mathbf{G}/\mathbf{L}/\mathbf{C}/\mathbf{M}$
controls which MZV weights appear:
\begin{enumerate}[label=\textup{(\roman*)}]
\item For 3d HT theories with Heisenberg boundary
\textup{(}class~$\mathbf{G}$\textup{)}: only $\zeta(2)$
appears, and the perturbative holomorphic Wilson line
terminates at one loop.
\item For 3d HT theories with affine Kac--Moody boundary
\textup{(}class~$\mathbf{L}$\textup{)}: $\zeta(2)$ and
$\zeta(3)$ appear.  The Drinfeld associator
\textup{(}Vol~I,
Theorem~\textup{thm:drinfeld-associator-bar-transport}\textup{)}
is the parallel transport of the genus-$0$ bar
connection, and its MZV coefficients encode the
perturbative holomorphic data.
\item For 3d gravity \textup{(}Virasoro boundary,
class~$\mathbf{M}$\textup{)}: MZVs at all weights
appear, reflecting the infinite tower of perturbative
corrections to the holomorphic graviton exchange.
The Zagier--Brown dimension sequence
$d_r = d_{r-2} + d_{r-3}$ counts the independent
perturbative MZV contributions at each order.
\end{enumerate}
\end{remark}

